% !TeX program = pdflatex
% Algebra abstrakcyjna: AA-01 do AA-61 stabilne; AA-62 w audycie.
\documentclass[11pt,a4paper,oneside]{book}
\usepackage[T1]{fontenc}
\usepackage[utf8]{inputenc}
\usepackage[polish]{babel}
\usepackage{lmodern}
\usepackage{amsmath,amssymb,amsthm,mathtools}
\usepackage{microtype}
\usepackage{tikz}
\usetikzlibrary{arrows.meta}
\usepackage{xcolor}
\usepackage{geometry}
\usepackage{enumitem}
\usepackage[unicode]{hyperref}
\hypersetup{colorlinks=true,linkcolor=blue!45!black,urlcolor=blue!45!black}
\geometry{margin=2.6cm}
\setlength{\parindent}{0pt}
\setlength{\parskip}{4pt}
\setlength{\emergencystretch}{2em}
\setlist[enumerate]{topsep=4pt,itemsep=2pt}
\numberwithin{equation}{chapter}
\theoremstyle{plain}
\newtheorem{twierdzenie}{Twierdzenie}[chapter]
\newtheorem{lemat}[twierdzenie]{Lemat}
\newtheorem{wniosek}[twierdzenie]{Wniosek}
\newtheorem{stwierdzenie}[twierdzenie]{Stwierdzenie}
\theoremstyle{definition}
\newtheorem{definicja}[twierdzenie]{Definicja}
\newtheorem{przyklad}[twierdzenie]{Przykład}
\theoremstyle{remark}
\newtheorem{uwaga}[twierdzenie]{Uwaga}
\title{Algebra abstrakcyjna\\\large Grupy, pierścienie i moduły}
\author{Biblioteka matematyczna}
\date{Rozdziały 01--61 stabilne; rozdział 62 w audycie}
\begin{document}
\frontmatter
\hypersetup{pageanchor=false}
\maketitle
\clearpage
\hypersetup{pageanchor=true}
\tableofcontents
\mainmatter
% AA-01 | Algebra abstrakcyjna | kanoniczny rozdzial: grupy
% Zrodla i kontynuacja: TOP-02, BT-01. Zachowano oddzielnie wszystkie oryginalne moduly.
\chapter{Grupy, podgrupy i grupy cykliczne}
\label{aa:ch:grupy}

\section{Dlaczego wprowadzamy grupy?}
W arytmetyce dodajemy liczby całkowite, w geometrii składamy obroty,
a w kombinatoryce składamy permutacje. Te działania są z pozoru różne,
ale wszystkie pozwalają wykonywać kolejne operacje, nie troszcząc się
o nawiasy, mają operację neutralną i pozwalają cofnąć każdą operację.
\emph{Teoria grup} bada konsekwencje dokładnie tych własności, bez
zakładania, że elementami grupy są liczby albo macierze.

Przez $\mathbb N$ oznaczamy w tym rozdziale zbiór $\{0,1,2,\ldots\}$.
Dla zbioru $X$ zapis $|X|$ oznacza liczbę jego elementów, jeśli jest skończony,
a w ogólności jego moc. Funkcje składamy od prawej do lewej:
$(f\circ g)(x)=f(g(x))$.

\begin{definicja}[Działanie wewnętrzne, półgrupa i monoid]
\label{aa:def:operacje}
Niech $S$ będzie zbiorem. \emph{Działanie wewnętrzne} na $S$ to
odwzorowanie $\star:S\times S\to S$. Para $(S,\star)$ jest
\emph{półgrupą}, jeżeli $(x\star y)\star z=x\star(y\star z)$
dla wszystkich $x,y,z\in S$. Jest \emph{monoidem}, gdy dodatkowo
istnieje $e\in S$ takie, że $e\star x=x\star e=x$ dla każdego $x\in S$.
\end{definicja}

\begin{przyklad}[Gdzie kończy się monoid?]
\label{aa:prz:monoidy}
Zbiór $(\mathbb N,+)$ jest monoidem z elementem neutralnym $0$,
lecz nie jest grupą: równanie $2+x=0$ nie ma rozwiązania w $\mathbb N$.
Dla niepustego zbioru $X$ wszystkie funkcje $X\to X$ tworzą monoid
ze składaniem, a elementem neutralnym jest $\operatorname{id}_X$.
Funkcja stała na zbiorze co najmniej dwuelementowym nie posiada
odwrotności względem składania: jej obraz ma jeden element,
a obraz złożenia z dowolną funkcją nie może zawierać całego $X$.
\end{przyklad}

\begin{definicja}[Grupa]
\label{aa:def:grupa}
\emph{Grupą} nazywamy parę $(G,\cdot)$, gdzie $G$ jest zbiorem,
$\cdot:G\times G\to G$ jest działaniem i zachodzą następujące warunki:
\begin{enumerate}
\item $(xy)z=x(yz)$ dla wszystkich $x,y,z\in G$ (łączność);
\item istnieje $e\in G$ takie, że $ex=xe=x$ dla każdego $x\in G$
(element neutralny);
\item dla każdego $x\in G$ istnieje $y\in G$ takie, że $xy=yx=e$
(element odwrotny).
\end{enumerate}
Grupa jest \emph{abelowa} (\emph{przemienna}), jeżeli ponadto $xy=yx$
dla wszystkich $x,y\in G$. Jej \emph{rząd} to moc $|G|$;
mówimy o grupie skończonej, gdy $|G|<\infty$.
W grupie abelowej często stosujemy zapis addytywny $x+y$, $0$, $-x$.
\end{definicja}

\begin{lemat}[Jednoznaczność oraz podstawowe rachunki]
\label{aa:lem:jedynka-odwrotnosc}
Niech $G$ będzie grupą. Wówczas:
\begin{enumerate}
\item element neutralny jest jedyny i każdy $x\in G$ ma dokładnie
jeden element odwrotny, oznaczany $x^{-1}$;
\item dla $x,y\in G$ zachodzi $(xy)^{-1}=y^{-1}x^{-1}$ oraz
$(x^{-1})^{-1}=x$;
\item dla dowolnych $a,b,c\in G$ implikacje
$ab=ac\Rightarrow b=c$ oraz $ba=ca\Rightarrow b=c$ są prawdziwe;
\item równania $ax=b$ oraz $ya=b$ mają odpowiednio dokładnie jedno
rozwiązanie $x=a^{-1}b$ oraz $y=ba^{-1}$.
\end{enumerate}
\end{lemat}
\begin{proof}
Jeżeli $e$ i $f$ są neutralne, to z neutralności $e$ otrzymujemy
$ef=f$, a z neutralności $f$ otrzymujemy $ef=e$. Zatem $e=f$.
Jeśli $u$ i $v$ są odwrotnościami $x$, to
$u=u(xv)=(ux)v=v$. Odwrotność jest więc jednoznaczna.

Mamy $(xy)(y^{-1}x^{-1})=x(yy^{-1})x^{-1}=e$ oraz
$(y^{-1}x^{-1})(xy)=y^{-1}(x^{-1}x)y=e$.
Jednoznaczność odwrotności daje pierwszy wzór w (2).
Ponieważ $xx^{-1}=x^{-1}x=e$, element $x$ jest odwrotnością
$x^{-1}$, co daje drugi wzór.

Jeżeli $ab=ac$, to mnożymy obie strony z lewej przez $a^{-1}$:
$a^{-1}(ab)=a^{-1}(ac)$, więc po użyciu łączności $b=c$.
Jeżeli $ba=ca$, to mnożymy obie strony z prawej przez $a^{-1}$:
$(ba)a^{-1}=(ca)a^{-1}$, a więc $b=c$.
Wreszcie $a(a^{-1}b)=b$ i $(ba^{-1})a=b$,
a każde inne rozwiązanie odpowiedniego równania musi się z nim
pokrywać na mocy wykazanej skracalności.
\end{proof}

\begin{uwaga}[Na co naprawdę pozwala łączność?]
Zapis $abc$ nie wymaga nawiasów dzięki łączności, lecz zwykle
$abc\ne acb$. Przestawienie czynników wymaga przemienności,
której nie zakładamy. W szczególności $(ab)^{-1}$ ma odwróconą
kolejność odwrotności. To rozróżnienie będzie kluczowe w grupach
permutacji, w grupach macierzy i w grupach wolnych.
\end{uwaga}

\section{Potęgi, permutacje, macierze i symetrie}
\begin{definicja}[Potęgi całkowite]
Niech $G$ będzie grupą, a $g\in G$. Definiujemy $g^0=e$,
$g^{n+1}=g^ng$ dla $n\ge0$ oraz $g^{-n}=(g^{-1})^n$
dla $n>0$. W grupie addytywnej piszemy odpowiednio $ng$.
\end{definicja}

\begin{lemat}[Prawa działań na potęgach]
\label{aa:lem:potegi}
Niech $G$ będzie grupą. Dla każdego $g\in G$ i $m,n\in\mathbb Z$ zachodzą
$g^{m+n}=g^mg^n$ oraz $(g^m)^n=g^{mn}$.
Jeżeli $g,h\in G$ spełniają $gh=hg$, to $(gh)^n=g^nh^n$
dla każdego $n\in\mathbb Z$.
\end{lemat}
\begin{proof}
Najpierw indukcja po $n\ge0$ daje $g^{m+n}=g^mg^n$
dla $m\ge0$. Dla $n>0$ elementy $g$ i $g^{-1}$ komutują, ponieważ
$gg^{-1}=e=g^{-1}g$. Indukcja pokazuje zatem, że
$g^n(g^{-1})^n=e=(g^{-1})^ng^n$; w szczególności
$g^{-n}=(g^n)^{-1}$.
Rozważmy teraz wszystkie znaki wykładników. Dla $a,b\ge0$
możemy przepisać $g^ag^{-b}$ jako iloczyn $a$ czynników
$g$ i $b$ czynników $g^{-1}$. Sąsiadujące pary $gg^{-1}$
można usuwać. Jeśli $a\ge b$, pozostaje $g^{a-b}$;
jeśli $a<b$, pozostaje $(g^{-1})^{b-a}=g^{a-b}$.
Analogicznie $g^{-a}g^b=g^{b-a}$ (czynniki $g^{-1}g$
też są równe $e$). Wreszcie
$g^{-a}g^{-b}=(g^{-1})^{a+b}=g^{-(a+b)}$,
bo dodatnie potęgi elementu $g^{-1}$ spełniają już
udowodnione prawo. Łącznie obejmuje to wszystkie pary
$m,n\in\mathbb Z$ i dowodzi $g^{m+n}=g^mg^n$. Dla $n\ge0$
druga wynika przez indukcję z pierwszej:
$(g^m)^{n+1}=g^{mn}g^m=g^{m(n+1)}$;
dla $n=-q<0$, gdzie $q>0$, mamy
$(g^m)^n=((g^m)^{-1})^q=(g^{-m})^q=g^{-mq}=g^{mn}$.
Jeżeli $gh=hg$, to indukcja daje
$(gh)^{n+1}=g^nh^ngh=g^{n+1}h^{n+1}$;
ujemne potęgi uzyskujemy, odwracając równość i używając
$(gh)^{-1}=g^{-1}h^{-1}$ (ostatnia równość wymaga przemienności).
\end{proof}

\begin{przyklad}[Grupy liczbowe]
\label{aa:prz:liczbowe}
$(\mathbb Z,+)$, $(\mathbb Q,+)$ i $(\mathbb R,+)$ są grupami
abelowymi. Zbiór dodatnich liczb rzeczywistych $\mathbb R_{>0}$
z mnożeniem także tworzy grupę; odwrotnością $t$ jest $1/t$.
Natomiast $(\mathbb Z,\cdot)$ nie jest grupą, ponieważ $2$
nie ma w $\mathbb Z$ odwrotności względem mnożenia.
\end{przyklad}

\begin{definicja}[Grupa permutacji]
Dla zbioru $X$ oznaczamy przez $\operatorname{Sym}(X)$ grupę
wszystkich bijekcji $X\to X$ ze składaniem. Dla $n\ge1$ i
$X=\{1,\ldots,n\}$ używamy oznaczenia $S_n$.
Dla liczby całkowitej $k\ge2$ i parami różnych $a_1,\ldots,a_k\in X$
cykl $(a_1\,a_2\,\ldots\,a_k)$ wysyła $a_i$ na $a_{i+1}$
(dla $i<k$), $a_k$ na $a_1$, a inne punkty pozostawia bez zmian.
\end{definicja}

\begin{przyklad}[Nieprzemienność w $S_3$]
\label{aa:prz:s3}
Niech $s=(1\,2)$ i $t=(2\,3)$. Wobec przyjętej konwencji
składania najpierw wykonujemy permutację po prawej.
Obliczenie $st$ daje $1\mapsto2$, $2\mapsto3$,
$3\mapsto1$, czyli $st=(1\,2\,3)$. Natomiast
$ts=(1\,3\,2)$, zatem $st\ne ts$.
Wszystkie sześć elementów $S_3$ można wypisać jako
$e,(1\,2),(1\,3),(2\,3),(1\,2\,3),(1\,3\,2)$.
Widzimy wprost, że $|S_3|=6$ i grupa nie jest abelowa.
\end{przyklad}

\begin{definicja}[Znak permutacji i grupa alternująca]
\label{aa:def:alternujaca}
Niech $n\ge1$. Dla $\sigma\in S_n$ \emph{inwersją} nazywamy parę
$(i,j)$ z $1\le i<j\le n$ i $\sigma(i)>\sigma(j)$.
Liczbę inwersji oznaczamy $I(\sigma)$, a \emph{znak} permutacji
definiujemy jako $\operatorname{sgn}(\sigma)=(-1)^{I(\sigma)}$.
Permutacja jest \emph{parzysta}, jeżeli jej znak to $1$,
a \emph{nieparzysta}, gdy jej znak to $-1$.
Grupą \emph{alternującą} nazywamy
$A_n=\{\sigma\in S_n:\operatorname{sgn}(\sigma)=1\}$.
\end{definicja}

\begin{lemat}[Znak iloczynu permutacji]
\label{aa:lem:znak}
Niech $n\ge1$. Dla $\sigma,\tau\in S_n$ mamy
$\operatorname{sgn}(\sigma\tau)
=\operatorname{sgn}(\sigma)\operatorname{sgn}(\tau)$.
Zbiór $A_n$ jest podgrupą $S_n$. Jeśli $n\ge2$,
to $|A_n|=n!/2$.
\end{lemat}
\begin{proof}
Niech $s_i=(i\ i+1)$ oznacza transpozycję sąsiednich
numerów. Pomnożenie permutacji $\sigma$ z prawej przez $s_i$
zamienia w ciągu $\sigma(1),\ldots,\sigma(n)$ dwa sąsiednie
wyrazy. Inwersja utworzona przez te wyrazy powstaje albo znika;
liczba inwersji dotyczących trzeciego wyrazu pozostaje bez zmian,
ponieważ oba zamieniane wyrazy stoją po tej samej stronie tego
trzeciego wyrazu. Zatem liczba inwersji zmienia się dokładnie o $1$
lub $-1$, a jej parzystość ulega odwróceniu. W szczególności
$\operatorname{sgn}(\sigma s_i)=-\operatorname{sgn}(\sigma)$.

Każdą permutację można otrzymać przez składanie transpozycji
sąsiednich. Aby to zobaczyć, zaczynamy od dowolnego ciągu
$\tau(1),\ldots,\tau(n)$ i zamieniamy sąsiednie wyrazy,
które są w odwrotnej kolejności, aż ciąg stanie się rosnący.
Jeżeli ciąg nie jest rosnący, istnieje w nim przynajmniej jedna
para sąsiednich wyrazów w odwrotnej kolejności.
Każdy taki krok zmniejsza liczbę inwersji dokładnie o $1$,
więc procedura się kończy. Odwracając kolejne kroki,
przedstawiamy $\tau$ jako iloczyn transpozycji sąsiednich.
Jeżeli jest ich $m$, to $(-1)^m=\operatorname{sgn}(\tau)$,
bo zaczynamy od permutacji identycznościowej, której znak to $1$.

Przy składaniu $\sigma\tau$ możemy więc kolejno mnożyć
$\sigma$ przez owe $m$ transpozycji. Znak zmienia się w każdym
kroku, toteż
$\operatorname{sgn}(\sigma\tau)
=\operatorname{sgn}(\sigma)(-1)^m
=\operatorname{sgn}(\sigma)\operatorname{sgn}(\tau)$.
Ponadto każda transpozycja $(a\,b)$, gdzie $1\le a<b\le n$,
ma dokładnie $2(b-a)-1$ inwersji: jedną parę $(a,b)$ oraz
po dwie pary dla każdego z $b-a-1$ numerów pomiędzy $a$ i $b$.
Jej znak wynosi więc $-1$, również gdy transponowane numery
nie są sąsiednie.
W szczególności iloczyn permutacji parzystych jest parzysty;
ponadto z $\operatorname{sgn}(\sigma)\operatorname{sgn}(\sigma^{-1})=\operatorname{sgn}(e)=1$ wynika, że $\operatorname{sgn}(\sigma^{-1})=\operatorname{sgn}(\sigma)^{-1}$,
więc również odwrotność permutacji parzystej jest parzysta.
Identyczność ma znak $1$, zatem $A_n$ jest podgrupą.

Dla $n\ge2$ ustalmy transpozycję $s=(1\ 2)$.
Mnożenie z prawej przez $s$ jest bijekcją z $S_n$ na siebie,
która zamienia znak permutacji. Daje więc bijekcję
między permutacjami parzystymi i nieparzystymi.
Ponieważ $|S_n|=n!$ (pierwszy obraz można wybrać na $n$
sposobów, drugi na $n-1$, itd.), połowa elementów należy do $A_n$.
\end{proof}

\begin{przyklad}[Dwanaście elementów $A_4$]
\label{aa:prz:a4-opis}
W grupie $S_4$ mamy $24$ permutacje, więc
$|A_4|=12$. Aby opisać je jawnie, zauważmy, że
$3$-cykl jest iloczynem dwóch transpozycji, a iloczyn
dwóch rozłącznych transpozycji także jest parzysty.
Każdy cykl $(a\,b\,c)$ można zapisać jako
$(a\,c)(a\,b)$, więc rzeczywiście jest parzysty.
Takich $3$-cykli jest $4\cdot3\cdot2/3=8$:
wybieramy trzy różne symbole i dzielimy przez trzy,
ponieważ dokładnie trzy cykliczne przesunięcia zapisu
$(a\,b\,c)$ oznaczają tę samą permutację.
Iloczyny dwóch rozłącznych transpozycji to dokładnie
$(1\,2)(3\,4)$, $(1\,3)(2\,4)$ i $(1\,4)(2\,3)$.
Razem z identycznością daje to $1+8+3=12$ różnych
permutacji parzystych, czyli całą grupę $A_4$.
\end{przyklad}

\begin{przyklad}[Grupy macierzowe]
Niech $K$ będzie ciałem, a $n\ge1$.
Macierze odwracalne stopnia $n$ nad $K$ tworzą grupę
$\operatorname{GL}_n(K)$ z mnożeniem macierzy;
macierz jednostkowa jest neutralna, a odwrotności istnieją
z definicji. Dla każdego ciała $K$ i $n\ge2$ grupa ta nie jest abelowa.
Istotnie, macierze
$A=\begin{psmallmatrix}1&1\\0&1\end{psmallmatrix}$
i $B=\begin{psmallmatrix}1&0\\1&1\end{psmallmatrix}$ mają wyznacznik $1$,
a ich iloczyny wynoszą
$AB=\begin{psmallmatrix}2&1\\1&1\end{psmallmatrix}$
i $BA=\begin{psmallmatrix}1&1\\1&2\end{psmallmatrix}$.
Nie są równe nad żadnym ciałem, ponieważ $2\ne1$ w $K$.
Dla $n>2$ można do obu macierzy dołączyć blok jednostkowy
stopnia $n-2$, otrzymując niekomutujące macierze w $\operatorname{GL}_n(K)$.
Grupa $\operatorname{SL}_n(K)$
składa się z macierzy o wyznaczniku $1$;
ponieważ $\det(AB)=\det A\det B$ i
$\det(A^{-1})=(\det A)^{-1}$, jest ona podgrupą $\operatorname{GL}_n(K)$.
\end{przyklad}

\begin{przyklad}[Symetrie kwadratu]
\label{aa:prz:d4}
Niech $r$ oznacza obrót kwadratu o $90^\circ$,
a $s$ odbicie względem ustalonej osi symetrii przechodzącej
przez środek. Mamy $r^4=e$, $s^2=e$ i $srs=r^{-1}$.
Uzasadnijmy relację $srs=r^{-1}$. Oznaczmy kolejne wierzchołki
przez $v_j=r^j(v_0)$, gdzie indeksy rozumiemy modulo $4$.
Odbicie odwraca ich porządek cykliczny, zatem dla pewnego
$c\in\{0,1,2,3\}$ zachodzi $s(v_j)=v_{c-j}$ dla każdego $j$.
Wobec tego
$(srs)(v_j)=s(r(v_{c-j}))=s(v_{c-j+1})=v_{j-1}=r^{-1}(v_j)$.
Symetrie kwadratu są wyznaczone przez obrazy wierzchołków,
więc $srs=r^{-1}$, a równoważnie $sr=r^{-1}s$.
Stąd każde złożenie obrotów i odbić można przekształcić do
postaci $r^j$ lub $r^js$, gdzie $0\le j<4$.
Te osiem przekształceń jest różnych: cztery pierwsze zachowują
orientację, a cztery pozostałe ją odwracają; w obrębie każdej
czwórki różnią się obrazem wybranego wierzchołka.
Każda symetria kwadratu jest wyznaczona przez obraz ustalonego
wierzchołka (cztery możliwości) i obraz jednego z jego sąsiadów
(dwie możliwości dla każdego wyboru), więc symetrii jest co
najwyżej osiem. Otrzymujemy całą grupę symetrii kwadratu,
czyli grupę dihedralną $D_4$ rzędu $8$.
W szczególności $sr\ne rs$ (bo $r\ne r^{-1}$).
\end{przyklad}

\section{Podgrupy i generowanie}
\begin{definicja}[Podgrupa]
Niech $G$ będzie grupą. Podzbiór $H\subseteq G$ jest
\emph{podgrupą}, co zapisujemy $H\le G$, jeżeli
$H$ jest grupą z działaniem odziedziczonym z $G$.
Podgrupę $H$ nazywamy \emph{właściwą}, gdy $H\ne G$.
\end{definicja}

\begin{twierdzenie}[Kryterium podgrupy]
\label{aa:tw:kryterium-podgrupy}
Niech $G$ będzie grupą, a $H\subseteq G$. Następujące
warunki są równoważne:
\begin{enumerate}
\item $H$ jest podgrupą $G$;
\item $H\ne\varnothing$ i dla każdych $x,y\in H$ zachodzi $xy^{-1}\in H$;
\item $e\in H$ i $H$ jest zamknięte na mnożenie i branie odwrotności.
\end{enumerate}
Jeśli $H$ jest niepustym, skończonym podzbiorem $G$ zamkniętym na mnożenie,
to $H$ także jest podgrupą.
\end{twierdzenie}
\begin{proof}
Jeśli $H$ jest podgrupą, to zawiera jej element neutralny.
Neutralny element $H$ musi być równy $e$ w $G$:
dla $h\in H$ mamy $e_Hh=h=eh$, więc skracalność w $G$
daje $e_H=e$. Podobnie odwrotność elementu w $H$
musi być jego odwrotnością w $G$. Stąd (1) implikuje (2).

Załóżmy (2) i wybierzmy $h\in H$. Ponieważ $hh^{-1}=e$,
mamy $e\in H$. Podstawiając $x=e$ i $y=h$ do (2),
otrzymujemy $h^{-1}\in H$. Dla $x,y\in H$ wiemy już,
że $y^{-1}\in H$, więc z (2) wynika $xy=x(y^{-1})^{-1}\in H$.
To daje (3). Warunek (3) zapewnia wszystkie aksjomaty grupy
na $H$ (łączność jest dziedziczona), czyli (1).

Na koniec niech $H$ będzie skończony, niepusty i zamknięty
na mnożenie. Dla ustalonego $h\in H$ wszystkie dodatnie
potęgi $h,h^2,\ldots$ należą do $H$. Dwie z nich są równe,
powiedzmy $h^i=h^j$ przy $1\le i<j$.
Po pomnożeniu w $G$ przez $h^{-i}$ otrzymujemy $h^{j-i}=e$.
Zatem $e\in H$ i $h^{-1}=h^{j-i-1}\in H$, gdzie
przy $j-i=1$ wykładnik zero oznacza $e$.
Ponieważ $h$ było dowolne, zachodzi (3).
\end{proof}

\begin{lemat}[Przecięcia podgrup]
\label{aa:lem:przeciecia}
Jeżeli $(H_i)_{i\in I}$ jest dowolną rodziną podgrup $G$,
to $\bigcap_{i\in I}H_i$ jest podgrupą $G$.
Dla $I=\varnothing$ przyjmujemy $\bigcap_{i\in I}H_i=G$.
\end{lemat}
\begin{proof}
Każde $H_i$ zawiera $e$, więc przecięcie jest niepuste.
Dla $x,y$ z przecięcia i każdego $i\in I$ mamy
$x,y\in H_i$, a stąd $xy^{-1}\in H_i$.
Zatem $xy^{-1}$ należy do przecięcia; stosujemy kryterium podgrupy.
Dla rodziny pustej wynik jest z definicji równy $G$.
\end{proof}

\begin{definicja}[Podgrupa generowana przez zbiór]
\label{aa:def:generatory}
Niech $G$ będzie grupą, a $S\subseteq G$. Definiujemy $\langle S\rangle$ jako
przecięcie wszystkich podgrup $G$ zawierających $S$.
Jest to najmniejsza podgrupa zawierająca $S$.
Dla $S=\{g_1,\ldots,g_k\}$, gdzie $k\ge1$, piszemy także
$\langle g_1,\ldots,g_k\rangle$.
\end{definicja}

\begin{twierdzenie}[Opis elementów podgrupy generowanej]
\label{aa:tw:slowa}
Dla grupy $G$ i $S\subseteq G$ zbiór $\langle S\rangle$ składa się
dokładnie ze wszystkich skończonych iloczynów
$s_1^{\varepsilon_1}\cdots s_m^{\varepsilon_m}$,
gdzie $m\ge0$, $s_i\in S$, $\varepsilon_i\in\{-1,1\}$,
a iloczyn pusty ($m=0$) jest równy $e$.
\end{twierdzenie}
\begin{proof}
Niech $W$ oznacza opisany zbiór iloczynów. Zawiera $e$
i każdy element $s\in S$ (jako iloczyn długości jeden).
Złączenie dwóch iloczynów daje iloczyn tego samego rodzaju.
Odwrotność $s_1^{\varepsilon_1}\cdots s_m^{\varepsilon_m}$
wynosi $s_m^{-\varepsilon_m}\cdots s_1^{-\varepsilon_1}$,
więc także należy do $W$. Z kryterium podgrupy $W\le G$.
Każda podgrupa zawierająca $S$ zawiera odwrotności jego
elementów oraz wszystkie ich skończone iloczyny; zatem zawiera $W$.
W szczególności $W$ jest zawarty we wszystkich podgrupach
występujących w definicji $\langle S\rangle$, a sam jest jedną
z nich. Stąd $W=\langle S\rangle$.
\end{proof}

\begin{przyklad}[Suma podgrup abelowych i iloczyn podgrup nieabelowych]
W każdej grupie abelowej $(A,+)$ suma dwóch podgrup
$H+K=\{h+k:h\in H,\ k\in K\}$ jest podgrupą:
zawiera $0=0+0$, a różnica $(h+k)-(h' +k')$
ma postać $(h-h')+(k-k')\in H+K$.
W grupie addytywnej $\mathbb Z$ zbiory $2\mathbb Z$ i $3\mathbb Z$
są podgrupami. Ich suma algebraiczna
$2\mathbb Z+3\mathbb Z=\{2a+3b:a,b\in\mathbb Z\}$
jest podgrupą i równa się $\mathbb Z$, ponieważ $1=-2+3$.
W grupie nieabelowej iloczyn zbiorów $HK=\{hk:h\in H,k\in K\}$
nie musi być podgrupą. W $S_3$ weźmy
$H=\{e,(1\,2)\}$ oraz $K=\{e,(2\,3)\}$.
Wtedy $HK=\{e,(1\,2),(2\,3),(1\,2\,3)\}$,
a odwrotność $(1\,2\,3)$ to $(1\,3\,2)$, której nie ma w $HK$.
Nie można więc bez sprawdzenia zamkniętości traktować $HK$ jak grupy.
W przypadku grup nieabelowych pojęcie \emph{sumy podgrup}
wymaga osobnej interpretacji: podstawową konstrukcją jest
podgrupa $\langle H\cup K\rangle$, która zawsze istnieje,
lecz nie musi być równa zbiorowi $HK$.
\end{przyklad}

\section{Grupy cykliczne i rząd elementu}
\begin{definicja}[Podgrupa cykliczna i rząd elementu]
\label{aa:def:rzad-elementu}
Niech $G$ będzie grupą, a $g\in G$. Podgrupę $\langle g\rangle=\{g^n:n\in\mathbb Z\}$
nazywamy \emph{cykliczną}. Jeżeli $\langle g\rangle=G$,
mówimy, że $G$ jest cykliczna, a $g$ jest jej generatorem.
\emph{Rzędem} $g$, oznaczanym $\operatorname{ord}(g)$,
jest liczba elementów $\langle g\rangle$.
\end{definicja}

\begin{lemat}[Najmniejsza dodatnia potęga równa jedynce]
\label{aa:lem:rzad-cykliczny}
Niech $G$ będzie grupą, a $g\in G$.
Jeżeli istnieje liczba całkowita $n\ge1$ z $g^n=e$, to zbiór dodatnich
wykładników o tej własności ma najmniejszy element $d$,
a wtedy dla wszystkich $m,a,b\in\mathbb Z$ zachodzą:
\begin{enumerate}
\item $g^m=e$ dokładnie wtedy, gdy $d\mid m$;
\item $g^a=g^b$ dokładnie wtedy, gdy $a\equiv b\pmod d$;
\item $\langle g\rangle=\{e,g,\ldots,g^{d-1}\}$ i
$\operatorname{ord}(g)=d$.
\end{enumerate}
Jeżeli nie istnieje dodatnie $n$ z $g^n=e$, to wszystkie potęgi
$g^m$, $m\in\mathbb Z$, są różne i $\operatorname{ord}(g)=\infty$.
\end{lemat}
\begin{proof}
Niepusty zbiór dodatnich wykładników ma najmniejszy element $d$.
Dla dowolnego $m\in\mathbb Z$ dzielimy $m$ z resztą:
$m=qd+r$, $0\le r<d$. Ponieważ $g^d=e$, mamy
$g^m=(g^d)^qg^r=g^r$. Gdyby $g^m=e$ i $r>0$,
to $g^r=e$ przeczyłoby minimalności $d$.
Zatem $r=0$, czyli $d\mid m$. Kierunek odwrotny jest jasny,
bo $g^{qd}=e$. Równość $g^a=g^b$ po pomnożeniu przez $g^{-b}$
jest równoważna $g^{a-b}=e$; zastosowanie (1) daje (2).
Każda potęga równa się jednej z $d$ wymienionych w (3),
a ich różność wynika z (2). W przypadku bez dodatniej
potęgi równej $e$ równość dwóch różnych potęg dawałaby
$g^{a-b}=e$; odwracając w razie potrzeby wykładnik,
otrzymalibyśmy potęgę dodatnią równą $e$, sprzeczność.
\end{proof}

Przez \emph{izomorfizm grup} rozumiemy bijekcję między ich
zbiorami elementów, która zachowuje odpowiednie działania.

\begin{twierdzenie}[Klasyfikacja grup cyklicznych]
\label{aa:tw:grupy-cykliczne}
Każda grupa cykliczna $G=\langle g\rangle$ jest abelowa.
Jeżeli $g$ ma rząd nieskończony, odwzorowanie
$\mathbb Z\to G$, $n\mapsto g^n$, jest izomorfizmem
(grupy $\mathbb Z$ rozumianej addytywnie).
Jeżeli $g$ ma skończony rząd $d$, to wzór
$[n]_d\mapsto g^n$ określa izomorfizm
$\mathbb Z/d\mathbb Z\to G$, gdzie
$[n]_d$ oznacza klasę reszt liczby $n$ modulo $d$.
\end{twierdzenie}
\begin{proof}
Dla dowolnych $g^a,g^b\in G$ mamy
$g^ag^b=g^{a+b}=g^{b+a}=g^bg^a$.
Z własności potęg odwzorowanie $n\mapsto g^n$ zachowuje działanie,
a z definicji generowania jest surjektywne. Przy rzędzie
nieskończonym jest różnowartościowe, bo potęgi są różne.
Przy rzędzie skończonym równość $g^a=g^b$ zachodzi dokładnie
wtedy, gdy $d$ dzieli $a-b$, co jest równoważne równości klas
$[a]_d=[b]_d$. To jednocześnie dowodzi, że wzór na klasach
nie zależy od reprezentanta i jest różnowartościowy.
Ponadto $g^{a+b}=g^ag^b$ zapewnia zachowanie dodawania klas
w postaci mnożenia w $G$, a zatem odwzorowanie jest izomorfizmem.
\end{proof}

\begin{lemat}[Podgrupy $\mathbb Z$]
\label{aa:lem:podgrupy-z}
Każda podgrupa $H\le(\mathbb Z,+)$ ma postać $d\mathbb Z$
dla dokładnie jednej liczby $d\in\mathbb N$.
Gdy $H\ne\{0\}$, liczba $d$ jest najmniejszym dodatnim elementem $H$.
\end{lemat}
\begin{proof}
Dla $H=\{0\}$ bierzemy $d=0$; tylko $d=0$ daje $d\mathbb Z=\{0\}$.
Jeśli $H$ zawiera niezerową liczbę, to zawiera dodatnią
(bo z $h\in H$ mamy $-h\in H$). Niech $d$ będzie najmniejszą
dodatnią liczbą w $H$. Zamykanie podgrupy na dodawanie i
przeciwne gwarantuje $d\mathbb Z\subseteq H$.
Dla $h\in H$ dzielimy $h=qd+r$ z $0\le r<d$.
Wtedy $r=h-qd\in H$. Minimalność $d$ daje $r=0$.
Tak więc $h\in d\mathbb Z$ i $H=d\mathbb Z$.
Jednoznaczność $d$ wynika z jego minimalności albo z $d=0$
w przypadku trywialnym.
\end{proof}

\begin{twierdzenie}[Wszystkie podgrupy grupy cyklicznej]
\label{aa:tw:podgrupy-cykliczne}
Każda podgrupa grupy cyklicznej jest cykliczna.
Dokładniej, jeśli $G=\langle g\rangle$ ma rząd $n<\infty$,
to dla każdego dodatniego dzielnika $d$ liczby $n$ istnieje
dokładnie jedna podgrupa rzędu $d$, mianowicie
$\langle g^{n/d}\rangle$; innych podgrup nie ma.
\end{twierdzenie}
\begin{proof}
Niech $\pi:\mathbb Z\to G$ będzie odwzorowaniem
$\pi(m)=g^m$. Dla podgrupy $H\le G$ zbiór
$\pi^{-1}(H)$ jest podgrupą $\mathbb Z$:
zawiera $0$, ponieważ $g^0=e\in H$, a jeżeli
$a,b\in\pi^{-1}(H)$, to $g^{a-b}=g^a(g^b)^{-1}\in H$,
więc $a-b\in\pi^{-1}(H)$.
Z poprzedniego lematu $\pi^{-1}(H)=k\mathbb Z$ dla pewnego $k\ge0$.
Ponieważ $\pi$ jest surjektywne i $H$ jest podgrupą $G$,
otrzymujemy $H=\pi(\pi^{-1}(H))
=\{g^{kj}:j\in\mathbb Z\}=\langle g^k\rangle$.
To dowodzi pierwszego zdania również przy nieskończonym rzędzie.

Przy $|G|=n$ zbiór $\pi^{-1}(H)$ zawiera $n\mathbb Z$,
bo $g^n=e\in H$. Zatem $k\ne0$ i $k$ dzieli $n$:
rzeczywiście $n\in k\mathbb Z$.
Wtedy $\langle g^k\rangle$ zawiera dokładnie
$n/k$ elementów $e,g^k,\ldots,g^{(n/k-1)k}$,
ponieważ dla $0\le i<j<n/k$ równość $g^{ik}=g^{jk}$
prowadziłaby do $n\mid(j-i)k$, czyli $n/k\mid(j-i)$,
co jest niemożliwe przy $0<j-i<n/k$. Ponadto
$(g^k)^{n/k}=g^n=e$, więc dalsze potęgi się powtarzają.
Jeżeli $|H|=d$, to $k=n/d$, co daje podany wzór i jednoznaczność.
Odwrotnie, każdy dzielnik $d\mid n$ prowadzi do takiej podgrupy.
\end{proof}

% Ten lemat musi poprzedzać pierwsze użycie tożsamości Bézouta.
\begin{lemat}[Tożsamość Bézouta]
\label{aa:lem:bezout}
Dla liczb całkowitych $a,b$, z których przynajmniej jedna
jest niezerowa, istnieją $u,v\in\mathbb Z$ takie, że
$ua+vb=\gcd(a,b)$.
\end{lemat}
\begin{proof}
Rozważmy zbiór $I=\{ua+vb:u,v\in\mathbb Z\}$.
Zawiera liczbę dodatnią: jeśli $a\ne0$, to $|a|\in I$,
a w przeciwnym razie $|b|\in I$. Niech $d$ będzie najmniejszą
dodatnią liczbą w $I$. Z definicji istnieją $u,v$ z $d=ua+vb$.
Dzieląc $a$ przez $d$ otrzymujemy $a=qd+r$, gdzie $0\le r<d$.
Ponieważ $a,d\in I$ oraz $I$ jest zamknięty na odejmowanie,
mamy $r=a-qd\in I$. Minimalność $d$ wymusza $r=0$,
więc $d\mid a$. Analogicznie $d\mid b$.
Każdy wspólny dodatni dzielnik $c$ liczb $a,b$ dzieli też
$ua+vb=d$, więc $c\le d$. Ponieważ $d$ samo jest
wspólnym dodatnim dzielnikiem, jest największym z nich,
czyli $d=\gcd(a,b)$.
\end{proof}

\begin{twierdzenie}[Rząd potęgi i generatory]
\label{aa:tw:rzad-potegi}
Niech $G$ będzie grupą i $g\in G$. Jeśli $\operatorname{ord}(g)=n<\infty$ i $k\in\mathbb Z$,
to $\operatorname{ord}(g^k)=n/\gcd(n,k)$,
przy konwencji $\gcd(n,0)=n$. W szczególności $g^k$ generuje
$\langle g\rangle$ wtedy i tylko wtedy, gdy $\gcd(n,k)=1$.
\end{twierdzenie}
\begin{proof}
Niech $d=\gcd(n,k)$; zapiszmy $n=dn'$ i $k=dk'$,
gdzie $\gcd(n',k')=1$. Dla dodatniego $m$ warunek
$(g^k)^m=e$ jest równoważny $n\mid km$, czyli $n'\mid k'm$.
Ponieważ $n'$ i $k'$ są względnie pierwsze, to
$n'\mid k'm$ dokładnie wtedy, gdy $n'\mid m$.
Dla jasności: z tożsamości Bézouta $un'+vk'=1$
po pomnożeniu przez $m$ wynika
$m=un'm+vk'm$, więc podzielność $n'\mid k'm$
pociąga $n'\mid m$. Najmniejszym dodatnim $m$ jest więc $n'$.
Podgrupa $\langle g^k\rangle$ jest zawarta w $\langle g\rangle$.
Ponieważ obie są skończone, zachodzi $\langle g^k\rangle=\langle g\rangle$
dokładnie wtedy, gdy ich rzędy są równe, czyli gdy $n'=n$,
a więc $d=1$.
\end{proof}

\begin{przyklad}[Wszystkie podgrupy i generatory $\mathbb Z/12\mathbb Z$]
\label{aa:prz:z12}
Grupa addytywna $\mathbb Z/12\mathbb Z$ jest cykliczna,
generowana przez $[1]$. Jej podgrupy odpowiadają dodatnim dzielnikom
$1,2,3,4,6,12$ liczby $12$:
$\{[0]\}$, $\langle[6]\rangle$,
$\langle[4]\rangle$, $\langle[3]\rangle$,
$\langle[2]\rangle$ i cała grupa.
Podgrupa $\langle[4]\rangle=\{[0],[4],[8]\}$ ma rząd $3$.
Rząd $[8]$ wynosi $12/\gcd(12,8)=3$.
Generatory całej grupy to $[1],[5],[7],[11]$,
bo dokładnie te klasy mają reprezentantów względnie pierwszych z $12$.
\end{przyklad}

\begin{uwaga}[Granica wnioskowania]
Wiemy już, że rząd grupy cyklicznej wyznacza możliwe rzędy jej
podgrup. W dowolnej grupie skończonej, niekoniecznie cyklicznej,
twierdzenie Lagrange'a daje warunek konieczny:
rząd podgrupy dzieli rząd grupy. Nie każdy dzielnik
rzędu grupy musi być rzędem jakiejś podgrupy; na przykład
grupa $A_4$ rzędu $12$ nie ma podgrupy rzędu $6$
(udowodnimy to po wprowadzeniu grup ilorazowych).
\end{uwaga}

\section{Podsumowanie zależności pojęć}
Poznane przykłady pokazują trzy różne sytuacje: struktura może
nie mieć odwrotności (monoid), mieć odwrotności, lecz nie być
przemienna ($S_3$ i $D_4$), lub być przemienną grupą cykliczną.
W następnym rozdziale zbadamy klasy elementów względem podgrupy.
Z ich rozkładu wynika twierdzenie Lagrange'a, a z niego
klasyczne twierdzenia arytmetyczne Eulera i Fermata.

% AA-02 | Algebra abstrakcyjna | kanoniczny rozdzial: warstwy, izomorfizmy
% Zrodla porownawcze: TOP-02, BT-01, ALG-00. Oryginaly pozostaja niezmienione.
\chapter{Warstwy, twierdzenie Lagrange'a i grupy ilorazowe}
\label{aa:ch:lagrange-ilorazy}

\section{Warstwy: rozbijanie grupy na części tej samej wielkości}
Niech $G$ będzie grupą, a $H\le G$ jej podgrupą.
Będziemy badać, na ile elementy $G$ różnią się od siebie
z punktu widzenia $H$. Kluczowa jest obserwacja, że przesunięte
kopie $H$ są albo identyczne, albo rozłączne.

\begin{definicja}[Warstwy lewo- i prawostronne, indeks]
\label{aa:def:warstwy}
Dla $g\in G$ definiujemy \emph{warstwę lewostronną}
$gH=\{gh:h\in H\}$ oraz \emph{warstwę prawostronną}
$Hg=\{hg:h\in H\}$.
Zbiór warstw lewostronnych oznaczamy $G/H$;
na tym etapie jest to jedynie zbiór, niekoniecznie grupa.
Liczbę warstw lewostronnych nazywamy \emph{indeksem} $H$ w $G$
i zapisujemy $[G:H]$.
\end{definicja}

\begin{lemat}[Kiedy dwie warstwy są równe?]
\label{aa:lem:rowne-warstwy}
Niech $G$ będzie grupą, a $H\le G$. Dla $g_1,g_2\in G$
zachodzą równoważności
\[
 g_1H=g_2H \quad\Longleftrightarrow\quad
 g_2^{-1}g_1\in H\quad\Longleftrightarrow\quad
 g_1H\cap g_2H\ne\varnothing.
\]
Analogicznie $Hg_1=Hg_2$ dokładnie wtedy, gdy
$g_1g_2^{-1}\in H$. Każda warstwa $gH$ ma tyle elementów co $H$.
\end{lemat}
\begin{proof}
Załóżmy $g_1H=g_2H$. Ponieważ $e\in H$, mamy
$g_1=g_1e\in g_2H$, a więc $g_1=g_2h$ dla pewnego $h\in H$.
Po pomnożeniu przez $g_2^{-1}$ dostajemy $g_2^{-1}g_1=h\in H$.

Odwrotnie, jeśli $g_2^{-1}g_1=h_0\in H$, to $g_1=g_2h_0$,
a stąd $g_1H=g_2h_0H=g_2H$. Uzasadnijmy ostatnią równość:
$h_0H=H$, gdyż mnożenie z lewej przez $h_0$ ma odwrotność
mnożenia przez $h_0^{-1}$, a oba należą do $H$.

Jeśli przecięcie warstw jest niepuste, to istnieją
$h_1,h_2\in H$ z $g_1h_1=g_2h_2$.
Przekształcenie daje $g_2^{-1}g_1=h_2h_1^{-1}\in H$,
więc warstwy są równe. Równe warstwy mają oczywiście
niepuste przecięcie, ponieważ $H\ne\varnothing$.
Dla warstw prawostronnych, jeśli $Hg_1=Hg_2$, to
$g_1\in Hg_2$ i istnieje $h\in H$ takie, że $g_1=hg_2$.
Stąd $g_1g_2^{-1}=h\in H$. Odwrotnie, gdy
$g_1g_2^{-1}=h\in H$, to $Hg_1=Hhg_2=Hg_2$,
ponieważ $Hh=H$. Jeśli natomiast $Hg_1\cap Hg_2\ne\varnothing$,
to z $h_1g_1=h_2g_2$ dla pewnych $h_1,h_2\in H$
wynika $g_1g_2^{-1}=h_1^{-1}h_2\in H$;
a zatem także dwie różne warstwy prawostronne są rozłączne.

Odwzorowanie $H\to gH$, $h\mapsto gh$, jest surjekcją
z definicji warstwy, a jego różnowartościowość wynika
ze skracalności: $gh_1=gh_2$ implikuje $h_1=h_2$.
Jest więc bijekcją i $|gH|=|H|$.
\end{proof}

\begin{twierdzenie}[Rozkład na warstwy]
\label{aa:tw:podzial-warstwy}
Niech $G$ będzie grupą, a $H\le G$. Warstwy lewostronne
względem $H$ tworzą rozkład $G$ na niepuste, parami rozłączne
podzbiory. Analogiczne stwierdzenie zachodzi dla warstw prawostronnych.
Ponadto istnieje bijekcja między zbiorami warstw lewostronnych
i prawostronnych zadana przez $gH\mapsto Hg^{-1}$.
\end{twierdzenie}
\begin{proof}
Każda warstwa $gH$ zawiera $g$, bo $e\in H$, zatem jest niepusta.
Każdy $g\in G$ należy do $gH$, więc suma wszystkich warstw daje $G$.
Z lematu o równości warstw każde dwie różne warstwy są rozłączne.
Dla warstw prawostronnych mamy $g=eg\in Hg$, każda jest więc
niepusta, a ich suma wynosi $G$. Z dowodu poprzedniego lematu
dwie warstwy prawostronne są albo równe, albo rozłączne.

Sprawdźmy dobrze określoność podanej bijekcji.
Jeżeli $gH=g'H$, to $g'^{-1}g\in H$,
a po odwróceniu $g^{-1}g'\in H$.
Z kryterium warstw prawostronnych wynika
$Hg^{-1}=Hg'^{-1}$. Definiujemy więc funkcję na warstwach.
Odwzorowanie odwrotne to $Hg\mapsto g^{-1}H$;
identyczny rachunek pokazuje, że jest dobrze określone.
Obie kompozycje są identycznościami, gdyż $(g^{-1})^{-1}=g$.
\end{proof}

\begin{przyklad}[Warstwy w $\mathbb Z$]
W grupie addytywnej $G=\mathbb Z$ i dla $H=4\mathbb Z$
warstwy to $0+H$, $1+H$, $2+H$, $3+H$.
Na przykład $1+H=\{\ldots,-7,-3,1,5,9,\ldots\}$.
Zupełny podział na cztery warstwy to po prostu podział liczb
całkowitych według reszty z dzielenia przez $4$.
Ponieważ dodawanie jest przemienne, warstwy lewo- i prawostronne
nie różnią się od siebie.
\end{przyklad}

\begin{przyklad}[Warstwy w grupie nieabelowej]
\label{aa:prz:warstwy-s3}
W $S_3$ niech $H=\{e,(1\,2)\}$. Mamy trzy warstwy lewostronne:
\[
 H,\quad (1\,3)H=\{(1\,3),(1\,2\,3)\},\quad
 (2\,3)H=\{(2\,3),(1\,3\,2)\}.
\]
Wystarczy wykonać składanie zgodnie z konwencją od prawej do lewej.
Przykładowo $(1\,3)(1\,2)=(1\,2\,3)$.
Natomiast $H(1\,3)=\{(1\,3),(1\,3\,2)\}$,
więc $(1\,3)H\ne H(1\,3)$.
Warstwy istnieją bez normalności podgrupy, ale nie muszą być
takie same po lewej i po prawej stronie.
\end{przyklad}

\begin{figure}[htbp]
\centering
\begin{tikzpicture}[font=\small,
  every node/.style={align=center},
  box/.style={draw,rounded corners=2mm,minimum height=1.40cm,minimum width=3.85cm}]
  \node[box] at (0,0) {$e,\ (1\,2)$};
  \node at (0,1.04) {$H$};
  \node[box] at (4.55,0) {$(1\,3),\ (1\,2\,3)$};
  \node at (4.55,1.04) {$(1\,3)H$};
  \node[box] at (9.10,0) {$(2\,3),\ (1\,3\,2)$};
  \node at (9.10,1.04) {$(2\,3)H$};
\end{tikzpicture}
\caption{Rozkład $S_3$ na trzy rozłączne warstwy lewostronne względem
$H=\{e,(1\,2)\}$. Każda warstwa ma dwa elementy.}
\label{aa:fig:warstwy-s3}
\end{figure}

\section{Twierdzenie Lagrange'a i jego konsekwencje}
\begin{twierdzenie}[Lagrange'a]
\label{aa:tw:lagrange}
Niech $G$ będzie grupą skończoną, a $H\le G$. Wtedy
$|G|=[G:H]\,|H|$. W szczególności $|H|$ dzieli $|G|$.
\end{twierdzenie}
\begin{proof}
Z twierdzenia o rozkładzie $G$ na warstwy każda warstwa
lewostronna jest niepusta, a dwie różne warstwy są rozłączne.
Z lematu o warstwach każda z nich ma dokładnie $|H|$ elementów.
Ponieważ $G$ jest skończona, jest tylko skończenie wiele warstw;
oznaczmy ich liczbę przez $k=[G:H]$.
Możemy wypisać $G=g_1H\,\dot\cup\,\cdots\,\dot\cup\,g_kH$,
gdzie symbol $\dot\cup$ oznacza sumę rozłączną.
Licząc elementy sumy rozłącznej otrzymujemy
$|G|=\sum_{i=1}^{k}|g_iH|=\sum_{i=1}^{k}|H|=k|H|$.
Jest to dokładnie żądany wzór. Ponieważ $k$ jest dodatnią
liczbą całkowitą, $|H|$ dzieli $|G|$.
\end{proof}

\begin{wniosek}[Rząd elementu dzieli rząd grupy]
\label{aa:wn:rzad-dzieli}
Niech $G$ będzie grupą skończoną. Dla każdego $g\in G$
liczba $\operatorname{ord}(g)$ dzieli $|G|$, a ponadto $g^{|G|}=e$.
\end{wniosek}
\begin{proof}
Podgrupa $\langle g\rangle\le G$ ma z definicji
$\operatorname{ord}(g)$ elementów. Twierdzenie Lagrange'a
zastosowane do tej podgrupy daje
$\operatorname{ord}(g)\mid |G|$.
Zapiszmy $|G|=k\operatorname{ord}(g)$.
Z definicji rzędu i twierdzenia o grupach cyklicznych
$g^{\operatorname{ord}(g)}=e$, zatem
$g^{|G|}=(g^{\operatorname{ord}(g)})^k=e$.
\end{proof}

\begin{wniosek}[Grupa rzędu będącego liczbą pierwszą]
\label{aa:wn:p-cykliczna}
Niech $p$ będzie liczbą pierwszą, a $G$ grupą rzędu $p$.
Wówczas $G$ jest cykliczna.
W szczególności jest abelowa.
\end{wniosek}
\begin{proof}
Ponieważ $p>1$, istnieje $g\in G$ z $g\ne e$.
Podgrupa $\langle g\rangle$ ma co najmniej dwa elementy,
więc jej rząd jest dzielnikiem $p$ większym od $1$.
Jedynym takim dzielnikiem liczby pierwszej jest $p$.
Zatem $|\langle g\rangle|=|G|$ i $\langle g\rangle=G$.
Grupa cykliczna jest abelowa na mocy twierdzenia z rozdziału poprzedniego.
\end{proof}

\begin{przyklad}[Co wynika z rzędu, a co nie wynika?]
Twierdzenie Lagrange'a podaje warunek konieczny, ale nie
zapewnia istnienia elementu każdego rzędu będącego dzielnikiem
rzędu grupy. Na przykład w grupie $A_4$ rzędu $12$
nie ma elementu rzędu $6$, choć $6\mid12$.
Rzeczywiście elementy $A_4$ to tożsamość, osiem $3$-cykli
oraz trzy iloczyny dwóch rozłącznych transpozycji.
Ich rzędy to odpowiednio $1$, $3$ i $2$.
Zatem żaden z nich nie ma rzędu $6$.
Nie utożsamiamy też istnienia elementu rzędu $d$ z istnieniem
podgrupy rzędu $d$: kierunek od elementu do podgrupy zachodzi
przez $\langle g\rangle$, lecz odwrotnie podgrupa może nie być cykliczna.
\end{przyklad}

\begin{wniosek}[Grupa rzędu $2p$ dla $p$ pierwszego: obserwacja]
Niech $p$ będzie liczbą pierwszą, a $G$ grupą rzędu $2p$.
Wtedy każdy element $G$ ma rząd należący do
$\{1,2,p,2p\}$, z pominięciem powtórzeń, jeśli $p=2$.
Jest to tylko ograniczenie arytmetyczne: nie oznacza,
że elementy wszystkich wymienionych rzędów muszą istnieć,
i nie klasyfikuje struktury grupy.
\end{wniosek}
\begin{proof}
Dzielnikami $2p$ są wymienione liczby, a z poprzedniego wniosku
rząd każdego elementu musi być jednym z nich.
\end{proof}

\section{Grupa jednostek modulo \texorpdfstring{$n$}{n} i twierdzenie Eulera}
Przy obliczaniu odwrotności modulo $n$ wykorzystamy tożsamość
Bézouta, którą udowodniliśmy w Lemacie~\ref{aa:lem:bezout}.
Pokazuje ona, że dla $\gcd(a,n)=1$ istnieją
$u,v\in\mathbb Z$ z $ua+vn=1$.


\begin{definicja}[Jednostki modulo $n$]
\label{aa:def:jednostki}
Dla liczby całkowitej $n\ge2$ definiujemy
\[
 (\mathbb Z/n\mathbb Z)^\times
 =\{[a]_n\in\mathbb Z/n\mathbb Z:\gcd(a,n)=1\}.
\]
Działaniem jest mnożenie klas modulo $n$. Liczbę tych klas
oznaczamy $\varphi(n)$ i nazywamy \emph{funkcją Eulera}.
\end{definicja}

\begin{lemat}[Dlaczego jednostki są grupą?]
\label{aa:lem:jednostki-grupa}
Dla każdej liczby całkowitej $n\ge2$ zbiór $(\mathbb Z/n\mathbb Z)^\times$
stanowi skończoną grupę abelową o rzędzie $\varphi(n)$.
\end{lemat}
\begin{proof}
Zapis $[a]_n$ zależy tylko od reszty z dzielenia przez $n$.
Jeżeli $a\equiv b\pmod n$, to $\gcd(a,n)=\gcd(b,n)$,
ponieważ wspólne dzielniki $a,n$ są tymi samymi liczbami,
które dzielą $b=a+kn$ i $n$.
Warunek wzajemnej pierwszości jest więc dobrze określony dla klas.

Jeśli $\gcd(a,n)=\gcd(b,n)=1$, to $\gcd(ab,n)=1$.
Wykażmy to bez nieudowodnionego jeszcze lematu Euklidesa.
Z tożsamości Bézouta wybierzmy $u,v,u',v'\in\mathbb Z$ takie,
że $ua+vn=1$ i $u'b+v'n=1$. Mnożąc równości,
otrzymujemy
\[
1=uu'ab+n(uav'+vu'b+vv'n).
\]
Każdy wspólny dzielnik $ab$ i $n$ dzieli prawą stronę
tej równości, a więc dzieli $1$. Stąd $\gcd(ab,n)=1$
i iloczyn dwóch jednostek ponownie jest jednostką.
Klasa $[1]_n$ jest neutralna i należy do tego zbioru.

Z tożsamości Bézouta istnieją liczby całkowite $u,v$ z
$ua+vn=1$. Redukując modulo $n$, otrzymujemy
$[u]_n[a]_n=[1]_n$. Ponadto $\gcd(u,n)=1$:
każdy wspólny dzielnik $u$ i $n$ dzieliłby $ua+vn=1$.
Zatem $[u]_n$ także należy do zbioru jednostek i jest
odwrotnością $[a]_n$. Łączność i przemienność dziedziczą się
z mnożenia liczb całkowitych. Rząd grupy jest z definicji równy
liczbie klas o reprezentancie względnie pierwszym z $n$,
czyli $\varphi(n)$.
\end{proof}

\begin{twierdzenie}[Eulera]
\label{aa:tw:euler}
Niech $n\in\mathbb Z$, $n\ge2$, oraz $a\in\mathbb Z$ spełniają
$\gcd(a,n)=1$. Wtedy $a^{\varphi(n)}\equiv1\pmod n$.
\end{twierdzenie}
\begin{proof}
Z poprzedniego lematu klasa $[a]_n$ jest elementem
skończonej grupy $(\mathbb Z/n\mathbb Z)^\times$,
której rząd wynosi $\varphi(n)$. Z wniosku z twierdzenia
Lagrange'a dla dowolnego elementu grupy $U$ rzędu $|U|$
mamy $u^{|U|}=e_U$. Tutaj $e_U=[1]_n$,
zatem $[a]_n^{\varphi(n)}=[1]_n$.
Z definicji mnożenia klas lewa strona to
$[a^{\varphi(n)}]_n$, a równość klas jest równoważna
kongruencji zapisanej w tezie.
\end{proof}

\begin{twierdzenie}[Małe twierdzenie Fermata]
\label{aa:tw:fermat}
Jeżeli $p$ jest liczbą pierwszą i $a\in\mathbb Z$,
to $a^p\equiv a\pmod p$. Jeśli dodatkowo $p\nmid a$,
to $a^{p-1}\equiv1\pmod p$.
\end{twierdzenie}
\begin{proof}
Gdy $p\nmid a$, każda z klas $[1]_p,\ldots,[p-1]_p$
jest odwracalna: dla $1\le k<p$ mamy $\gcd(k,p)=1$.
Stąd $\varphi(p)=p-1$, więc twierdzenie Eulera daje
$a^{p-1}\equiv1\pmod p$. Po pomnożeniu przez $a$
otrzymujemy $a^p\equiv a\pmod p$.
Jeżeli $p\mid a$, to $a\equiv0\pmod p$, a więc zarówno
$a^p$, jak i $a$ mają resztę zero; pierwszy wzór zachodzi również wtedy.
\end{proof}

\begin{twierdzenie}[Wzór na funkcję Eulera]
\label{aa:tw:phi-wzor}
Niech $n\ge2$ będzie liczbą całkowitą, a
$p_1,\ldots,p_r$ wszystkimi różnymi dzielnikami pierwszymi $n$.
Wtedy
$\varphi(n)=n\prod_{i=1}^{r}(1-1/p_i)$.
W szczególności $\varphi(p^k)=p^k-p^{k-1}$
dla każdej liczby pierwszej $p$ i $k\in\mathbb N$, $k\ge1$.
\end{twierdzenie}
\begin{proof}
Wśród $n$ kolejnych liczb $1,2,\ldots,n$ dokładnie
$n/p_i$ jest podzielnych przez $p_i$; to wielokrotności
$p_i,2p_i,\ldots,(n/p_i)p_i$.
Liczba jest względnie pierwsza z $n$ dokładnie wtedy,
gdy nie jest podzielna przez żaden z $p_i$.
Z zasady włączeń i wyłączeń liczba takich liczb wynosi
\[
 n-\sum_i\frac n{p_i}+\sum_{i<j}\frac n{p_ip_j}
 -\cdots+(-1)^r\frac n{p_1\cdots p_r}.
\]
Każdy iloczyn wybranych spośród różnych liczb pierwszych
$p_1,\ldots,p_r$ dzieli $n$. Stąd wszystkie składniki
w powyższym rachunku są całkowitymi liczbami zliczanych wielokrotności. Rozwinięcie iloczynu
$\prod_i(1-1/p_i)$ daje dokładnie powyższą sumę podzieloną
przez $n$. Dla potęgi $p^k$ występuje tylko jeden
różny czynnik pierwszy, więc
$\varphi(p^k)=p^k(1-1/p)=p^k-p^{k-1}$.
\end{proof}

\begin{przyklad}[Liczenie dużych potęg modulo $15$]
\label{aa:prz:euler15}
Ponieważ $15=3\cdot5$, mamy
$\varphi(15)=15(1-1/3)(1-1/5)=8$.
Liczby względnie pierwsze z $15$ w zakresie od $1$ do $15$
to $1,2,4,7,8,11,13,14$, więc grupę jednostek można też
bezpośrednio wypisać. Twierdzenie Eulera mówi, że
$2^8\equiv1\pmod{15}$. Stąd
$2^{100}=2^{8\cdot12+4}\equiv(2^8)^{12}2^4
\equiv16\equiv1\pmod{15}$.
Zauważmy, że jest to poprawna redukcja potęgi $2$
ze względu na $\gcd(2,15)=1$.
Przy $a=3$ tego twierdzenia użyć nie wolno:
$3^8\equiv6\pmod{15}$, a nie $1$.
\end{przyklad}

\begin{przyklad}[Odwrotność z tożsamości Bézouta]
W grupie $(\mathbb Z/26\mathbb Z)^\times$ klasa $[7]$
jest odwracalna, bo $\gcd(7,26)=1$.
Algorytm Euklidesa daje $26=3\cdot7+5$,
$7=1\cdot5+2$, $5=2\cdot2+1$, a więc
$1=5-2\cdot2=3\cdot5-2\cdot7=3\cdot26-11\cdot7$.
Zatem $[7]^{-1}=[-11]=[15]$ modulo $26$;
sprawdzenie: $7\cdot15=105\equiv1\pmod{26}$.
\end{przyklad}

\begin{uwaga}[Kongruencja Fermata nie jest testem rozstrzygającym pierwszość]
Dla liczby całkowitej $m\ge2$ kongruencja
$a^{m-1}\equiv1\pmod m$ nie oznacza, że $m$ jest liczbą pierwszą.
Na przykład $341=11\cdot31$
jest złożona, ale $2^{10}=1024=3\cdot341+1$,
więc $2^{340}=(2^{10})^{34}\equiv1\pmod{341}$.
Jest to przykład liczby pseudopierwszej przy podstawie $2$.
Sam test jednej podstawy nie dowodzi pierwszości liczby.
\end{uwaga}

\section{Homomorfizmy, jądra i obrazy}
\begin{definicja}[Homomorfizm grup]
\label{aa:def:hom}
Niech $(G,\cdot)$ i $(K,\star)$ będą grupami.
Funkcję $f:G\to K$ nazywamy \emph{homomorfizmem}, jeśli
$f(xy)=f(x)\star f(y)$ dla wszystkich $x,y\in G$.
Jeżeli jest przy tym różnowartościowa, nazywamy ją
\emph{monomorfizmem}; jeśli jest surjektywna ---
\emph{epimorfizmem w sensie surjektywnego homomorfizmu grup};
a jeśli jest bijektywna --- \emph{izomorfizmem}.
Izomorfizm $G\to G$ to \emph{automorfizm} grupy $G$.
Przy istnieniu izomorfizmu $G\to K$ piszemy $G\cong K$.
\emph{Jądro} i \emph{obraz} homomorfizmu to odpowiednio
$\ker f=\{g\in G:f(g)=e_K\}$ oraz
$\operatorname{im}f=\{f(g):g\in G\}$.
\end{definicja}

\begin{lemat}[Działanie homomorfizmu na element neutralny i odwrotności]
\label{aa:lem:hom-podstawy}
Dla homomorfizmu $f:G\to K$ mamy $f(e_G)=e_K$,
$f(g^{-1})=f(g)^{-1}$ i $f(g^m)=f(g)^m$
dla wszystkich $g\in G$ i $m\in\mathbb Z$.
\end{lemat}
\begin{proof}
Z tożsamości $e_Ge_G=e_G$ wynika
$f(e_G)f(e_G)=f(e_G)$. Mnożąc przez odwrotność
$f(e_G)$, otrzymujemy $f(e_G)=e_K$.
Ponieważ $gg^{-1}=g^{-1}g=e_G$,
$ f(g)f(g^{-1})=f(g^{-1})f(g)=e_K$;
jednoznaczność odwrotności daje drugi wzór.
Dla $m=0$ trzeci wzór ma postać $f(e_G)=e_K$.
Jeżeli $f(g^m)=f(g)^m$ dla $m\ge0$, to
$f(g^{m+1})=f(g^m)f(g)=f(g)^{m+1}$.
Wreszcie dla $m=-k<0$, $k\ge1$, mamy
$f(g^m)=f((g^{-1})^k)=f(g^{-1})^k=(f(g)^{-1})^k=f(g)^m$.
\end{proof}

\begin{twierdzenie}[Jądro i obraz]
\label{aa:tw:jadro-obraz}
Dla każdego homomorfizmu grup $f:G\to K$ jądro jest
podgrupą $G$, a obraz podgrupą $K$. Homomorfizm $f$
jest różnowartościowy dokładnie wtedy, gdy $\ker f=\{e_G\}$.
Ponadto dla wszystkich $g\in G$ i $x\in\ker f$
zachodzi $gxg^{-1}\in\ker f$.
\end{twierdzenie}
\begin{proof}
Mamy $f(e_G)=e_K$, więc jądro nie jest puste.
Jeśli $x,y\in\ker f$, to
$f(xy^{-1})=f(x)f(y)^{-1}=e_Ke_K=e_K$,
zatem $xy^{-1}\in\ker f$. Z kryterium podgrupy jądro jest podgrupą.
Podobnie obraz zawiera $e_K$, a dla
$u=f(x)$ i $v=f(y)$ mamy $uv^{-1}=f(xy^{-1})$
należące do obrazu.

Jeżeli $f$ jest różnowartościowy i $f(x)=e_K=f(e_G)$,
to $x=e_G$, więc jądro jest trywialne.
Jeżeli jądro jest trywialne i $f(x)=f(y)$,
to $f(y^{-1}x)=e_K$; zatem $y^{-1}x=e_G$ i $x=y$.
Wreszcie $f(gxg^{-1})=f(g)e_Kf(g)^{-1}=e_K$.
\end{proof}

\begin{przyklad}[Homomorfizm wyznacznika]
Dla każdego ciała $K$ i $n\ge1$ wzór
$\det:\operatorname{GL}_n(K)\to K^\times$
jest homomorfizmem, gdyż $\det(AB)=\det A\det B$.
Jego jądrem jest $\operatorname{SL}_n(K)$.
Jego obraz to całe $K^\times$: liczba $a\ne0$
jest wyznacznikiem macierzy diagonalnej
$\operatorname{diag}(a,1,\ldots,1)$.
Nie jest to homomorfizm addytywny macierzy;
strukturą grupową po obu stronach jest mnożenie.
\end{przyklad}

\section{Podgrupy normalne i konstrukcja grupy ilorazowej}
\begin{definicja}[Podgrupa normalna]
\label{aa:def:normalna}
Niech $G$ będzie grupą, a $N\le G$. Podgrupę $N$
nazywamy \emph{normalną}, zapis $N\trianglelefteq G$,
jeżeli $gNg^{-1}=N$ dla każdego $g\in G$.
\end{definicja}

\begin{twierdzenie}[Równoważne kryteria normalności]
\label{aa:tw:normalna-warstwy}
Niech $G$ będzie grupą, a $N\le G$.
Następujące warunki są równoważne:
\begin{enumerate}
\item $N\trianglelefteq G$;
\item $gN=Ng$ dla każdego $g\in G$;
\item $gng^{-1}\in N$ dla wszystkich $g\in G$ i $n\in N$.
\end{enumerate}
\end{twierdzenie}
\begin{proof}
Równość $gNg^{-1}=N$ po przemnożeniu obu stron z prawej
przez $g$ daje $gN=Ng$, a odwrotnie pomnożenie przez $g^{-1}$
daje (1). Stąd (1) i (2) są równoważne.
(1) implikuje (3) przez definicję. Jeśli zachodzi (3),
to $gNg^{-1}\subseteq N$; stosując (3) do $g^{-1}$,
otrzymujemy $g^{-1}Ng\subseteq N$. Po sprzężeniu przez $g$
wyprowadzamy $N\subseteq gNg^{-1}$, co daje równość (1).
\end{proof}

\begin{twierdzenie}[Iloraz grupy przez podgrupę normalną]
\label{aa:tw:iloraz}
Niech $G$ będzie grupą, a $N\trianglelefteq G$.
Na zbiorze warstw $G/N$ wzór
$(gN)(hN)=(gh)N$ jest dobrze określony i czyni $G/N$
grupą. Element neutralny to $N=eN$,
a odwrotnością $gN$ jest $g^{-1}N$.
Odwzorowanie $q:G\to G/N$, $q(g)=gN$, jest surjektywnym
homomorfizmem, a jego jądro to $N$.
Jeżeli zaś dla dowolnej podgrupy $H\le G$ wzór
$(gH)(hH)=ghH$ jest dobrze określony na warstwach,
to $H$ musi być normalna.
\end{twierdzenie}
\begin{proof}
Załóżmy najpierw normalność $N$.
Niech $gN=g'N$ i $hN=h'N$.
Z lematu o warstwach istnieją $n_1,n_2\in N$ takie,
że $g'=gn_1$ i $h'=hn_2$. Wówczas
$g'h'=gn_1hn_2=gh(h^{-1}n_1h)n_2$.
Ponieważ $N$ jest normalna, mamy $h^{-1}n_1h\in N$;
ponieważ jest podgrupą, iloczyn $(h^{-1}n_1h)n_2$ też należy do $N$.
Otrzymujemy zatem $g'h'N=ghN$.
Wynik nie zależy od wybranych reprezentantów.

Łączność działania warstw wynika z łączności w $G$:
$((gN)(hN))(kN)=((gh)k)N=(g(hk))N=(gN)((hN)(kN))$.
Mamy $(eN)(gN)=gN=(gN)(eN)$ oraz
$(gN)(g^{-1}N)=eN=(g^{-1}N)(gN)$.
Zatem $G/N$ jest grupą.

Równość $q(gh)=ghN=(gN)(hN)=q(g)q(h)$
dowodzi, że $q$ jest homomorfizmem. Każda warstwa
jest postaci $q(g)$, więc jest on surjektywny.
Ponadto $q(g)=N$ dokładnie wtedy, gdy $g\in N$,
zatem $\ker q=N$.

Na koniec niech wzór działa dla pewnej podgrupy $H$.
Dla każdego $g\in G$, $h\in H$ mamy $hH=H$,
więc z poprawności iloczynu
$(gH)(hH)(g^{-1}H)=(gH)(H)(g^{-1}H)=H$.
Ale lewa strona według wzoru to $(ghg^{-1})H$.
Stąd $ghg^{-1}\in H$. Ponieważ dotyczy to każdego $g,h$,
uzyskujemy normalność $H$ z poprzedniego twierdzenia.
\end{proof}

\begin{wniosek}[Podgrupy normalne jako jądra]
\label{aa:wn:normalna-jadro}
Niech $G$ będzie grupą, a $N\le G$. Wówczas $N\trianglelefteq G$
dokładnie wtedy, gdy istnieje grupa $K$ i homomorfizm
$f:G\to K$, dla którego $N=\ker f$.
\end{wniosek}
\begin{proof}
Jeżeli $N$ jest normalna, twierdzenie o grupie ilorazowej
konstruuje homomorfizm $q:G\to G/N$ z jądrem $N$.
Jeżeli natomiast $N=\ker f$, to dla każdego $g\in G$
i $n\in N$ zachodzi
$f(gng^{-1})=f(g)f(n)f(g)^{-1}=e_K$,
zatem $gng^{-1}\in N$. Kryterium normalności dowodzi tezy.
\end{proof}

\begin{przyklad}[Dlaczego naturalny iloczyn warstw może nie istnieć]
Niech $H=\{e,(1\,2)\}\le S_3$.
Mamy $(1\,3)(1\,2)(1\,3)=(2\,3)\notin H$,
więc $H$ nie jest normalna.
Pokażmy zależność od reprezentantów bezpośrednim rachunkiem.
Niech $g=(1\,3)$, $g'=(1\,2\,3)$ i $k=(1\,3)$.
Ponieważ $g'=g(1\,2)$, mamy $gH=g'H$. Jeśli jednak
zastosujemy proponowany wzór do tych samych dwóch warstw,
wybierając jako reprezentanta pierwszej odpowiednio $g$ albo $g'$,
otrzymamy
\[
(gk)H=eH=H,\qquad (g'k)H=(2\,3)H\ne H.
\]
Ostatnia nierówność wynika z $(2\,3)\notin H$.
Zatem ten sam iloczyn warstw miałby dwie różne wartości:
normalność jest warunkiem koniecznym jego dobrej określoności.
\end{przyklad}

\begin{przyklad}[Iloraz $S_3/A_3$]
Niech $A_3=\{e,(1\,2\,3),(1\,3\,2)\}$.
Jest to podgrupa rzędu $3$; składa się z permutacji parzystych.
Znak permutacji $\operatorname{sgn}:S_3\to\{\pm1\}$
jest homomorfizmem (znak iloczynu to iloczyn znaków),
a jego jądrem jest $A_3$. Wynika stąd normalność $A_3$.
Z twierdzenia Lagrange'a indeks $[S_3:A_3]=6/3=2$,
więc iloraz ma dwie klasy. Funkcja
$gA_3\mapsto\operatorname{sgn}(g)$ jest dobrze określona:
przedstawiciele tej samej warstwy różnią się o element jądra
znaku. Jest bijektywna i zachowuje mnożenie, a więc
$S_3/A_3\cong\{\pm1\}\cong\mathbb Z/2\mathbb Z$.
Warto podkreślić, że warstwa nie jest pojedynczą permutacją:
jedna zawiera trzy permutacje parzyste, druga trzy nieparzyste.
\end{przyklad}

\begin{wniosek}[Podgrupy indeksu dwa są normalne]
\label{aa:wn:indeks-dwa}
Niech $G$ będzie grupą, a $H\le G$ podgrupą o indeksie $2$.
Wtedy $H\trianglelefteq G$.
\end{wniosek}
\begin{proof}
Istnieją dokładnie dwie warstwy lewostronne:
$H$ oraz $G\setminus H$, bo warstwy rozkładają $G$
i jedna z nich jest równa $H$.
Z bijekcji warstw lewo- i prawostronnych otrzymanej wcześniej
wynika, że istnieją dokładnie dwie warstwy prawostronne:
$H$ oraz $G\setminus H$.
Dla $g\in H$ mamy $gH=H=Hg$.
Dla $g\notin H$ warstwa $gH$ jest różna od $H$,
więc musi być równa $G\setminus H$; tak samo $Hg$.
Zatem $gH=Hg$ dla wszystkich $g$,
co jest równoważne normalności.
\end{proof}

\begin{przyklad}[Dlaczego $A_4$ nie ma podgrupy rzędu $6$?]
\label{aa:prz:a4-bez-szostki}
Przypomnijmy, że $A_4$ składa się z dwunastu permutacji
parzystych czterech elementów: identyczności, ośmiu $3$-cykli
i trzech produktów dwóch rozłącznych transpozycji.
Załóżmy, że $H\le A_4$ ma rząd $6$.
Wtedy $[A_4:H]=2$, więc na mocy poprzedniego wniosku
$H\trianglelefteq A_4$ i iloraz $A_4/H$ jest grupą rzędu $2$.
Każdy $3$-cykl $c$ spełnia $c^3=e$, a jego klasa
$cH$ spełnia $(cH)^3=H$.
Z twierdzenia Lagrange'a rząd $cH$ dzieli $2$;
jednocześnie z $(cH)^3=H$ wynika, że dzieli także $3$.
Jedynym wspólnym dodatnim dzielnikiem $2$ i $3$ jest $1$,
więc $cH=H$, czyli $c\in H$.
W ten sposób $H$ zawierałaby wszystkie osiem $3$-cykli
oraz element neutralny, czyli co najmniej dziewięć elementów.
To sprzeczne z $|H|=6$. Zatem taka podgrupa nie istnieje.
Ten przykład pokazuje, że odwrotność twierdzenia Lagrange'a
jest fałszywa nawet dla stosunkowo małych grup.
\end{przyklad}

\section{Twierdzenia o izomorfizmie i odpowiedniości}
\begin{twierdzenie}[Pierwsze twierdzenie o izomorfizmie grup]
\label{aa:tw:izo-pierwsze}
Jeżeli $f:G\to K$ jest homomorfizmem grup, to istnieje
dokładnie jeden izomorfizm
$\overline f:G/\ker f\to\operatorname{im}f$
spełniający $\overline f(g\ker f)=f(g)$ dla każdego $g\in G$.
\end{twierdzenie}
\begin{proof}
Niech $N=\ker f$. Z twierdzenia o jądrach wynika normalność $N$,
a więc $G/N$ jest grupą.
Aby sprawdzić, że zadany wzór określa funkcję,
załóżmy $gN=hN$. Z lematu o warstwach $h^{-1}g\in N$,
więc $f(h)^{-1}f(g)=f(h^{-1}g)=e_K$ i $f(g)=f(h)$.
Wartość wzoru nie zależy zatem od reprezentanta klasy.

Jest to homomorfizm, bo
$\overline f((gN)(hN))=f(gh)=f(g)f(h)
=\overline f(gN)\overline f(hN)$.
Jest surjektywny na $\operatorname{im}f$ z definicji obrazu.
Jeśli $\overline f(gN)=e_K$, to $f(g)=e_K$, a więc
$g\in N$ i $gN=N$; jądro $\overline f$ jest trywialne.
Z kryterium iniektywności jest więc izomorfizmem.
Wreszcie każda klasa ma postać $gN$, toteż wymagany
wzór wyznacza $\overline f$ jednoznacznie.
\end{proof}

\begin{przyklad}[Od dzielenia z resztą do izomorfizmu]
Homomorfizm grup addytywnych $f:\mathbb Z\to\mathbb Z/12\mathbb Z$
ma postać $f(k)=[k]_{12}$. Jego jądro to $12\mathbb Z$,
a obraz to cała $\mathbb Z/12\mathbb Z$. Pierwsze twierdzenie
o izomorfizmie daje $\mathbb Z/12\mathbb Z\cong\operatorname{im}f$.
Ten przykład może wydawać się tautologiczny, ale pokazuje schemat:
grupę ilorazową uzyskujemy przez utożsamienie elementów,
których obrazy pod homomorfizmem są takie same.
\end{przyklad}

\begin{twierdzenie}[Drugie twierdzenie o izomorfizmie]
\label{aa:tw:izo-drugie}
Niech $G$ będzie grupą, $H\le G$, $N\trianglelefteq G$.
Wtedy $HN=\{hn:h\in H,n\in N\}$ jest podgrupą $G$,
$N\trianglelefteq HN$, $H\cap N\trianglelefteq H$ i istnieje
naturalny izomorfizm $H/(H\cap N)\cong HN/N$.
\end{twierdzenie}
\begin{proof}
Zbiór $HN$ zawiera $e=ee$. Niech $h_1n_1,h_2n_2\in HN$.
Mamy
$(h_1n_1)(h_2n_2)^{-1}
=h_1n_1n_2^{-1}h_2^{-1}
=(h_1h_2^{-1})(h_2n_1n_2^{-1}h_2^{-1})$.
Pierwszy czynnik należy do $H$, drugi do $N$ dzięki normalności.
Stąd wynik należy do $HN$ i kryterium podgrupy daje $HN\le G$.
Normalność $N$ w $HN$ wynika z normalności w $G$.

Rozważmy homomorfizm $f:H\to HN/N$,
$f(h)=hN$. Mnożenie warstw jest dobrze określone,
ponieważ $N\trianglelefteq HN$. Każdy element $HN/N$
ma postać $hnN=hN$, więc $f$ jest surjektywny.
Jądro to $\{h\in H:hN=N\}=H\cap N$.
Pierwsze twierdzenie o izomorfizmie daje
$H/(H\cap N)\cong HN/N$. Ponieważ $H\cap N$ jest
jądrem homomorfizmu $f$, jest normalna w $H$.
\end{proof}

\begin{twierdzenie}[Trzecie twierdzenie o izomorfizmie]
\label{aa:tw:izo-trzecie}
Niech $G$ będzie grupą, $N\trianglelefteq G$ i $M\trianglelefteq G$,
przy $N\subseteq M$. Wtedy $M/N\trianglelefteq G/N$ oraz
$(G/N)/(M/N)\cong G/M$.
\end{twierdzenie}
\begin{proof}
Zdefiniujmy $F:G/N\to G/M$ wzorem $F(gN)=gM$.
Jeżeli $gN=hN$, to $h^{-1}g\in N\subseteq M$,
więc $gM=hM$; wzór jest dobrze określony.
Ponadto $F((gN)(hN))=ghM=(gM)(hM)$,
zatem $F$ jest homomorfizmem.
Jest surjektywny, ponieważ każda klasa $gM$ ma
przeciwobraz $gN$. Jego jądro składa się z tych $gN$,
dla których $gM=M$, czyli $g\in M$.
Zatem $\ker F=M/N$. Jako jądro jest to normalna podgrupa $G/N$.
Pierwsze twierdzenie o izomorfizmie, zastosowane do $F$,
daje dokładnie żądany izomorfizm.
\end{proof}

\begin{twierdzenie}[Odpowiedniość podgrup]
\label{aa:tw:odpowiedniosc}
Niech $G$ będzie grupą, $N\trianglelefteq G$ i niech $q:G\to G/N$
będzie homomorfizmem ilorazowym. Odwzorowania
$H\mapsto H/N$ oraz $K\mapsto q^{-1}(K)$ są wzajemnie
odwrotnymi bijekcjami między podgrupami $H\le G$
zawierającymi $N$ a podgrupami $K\le G/N$.
Odpowiedniość zachowuje inkluzję, a $H$ jest normalna w $G$
dokładnie wtedy, gdy $H/N$ jest normalna w $G/N$.
\end{twierdzenie}
\begin{proof}
Niech $H$ będzie podgrupą $G$ zawierającą $N$.
Wtedy $N\trianglelefteq H$, więc $H/N$ jest podgrupą $G/N$.
Przeciwobraz podgrupy przez homomorfizm jest podgrupą:
rzeczywiście zawiera element neutralny i jest zamknięty
na $xy^{-1}$, ponieważ $q(xy^{-1})=q(x)q(y)^{-1}$.
Zatem oba przyporządkowania mają wskazane przeciwdziedziny.

Dla $H$ zachodzi $q^{-1}(H/N)=H$.
Dla każdego $h\in H$ klasa $q(h)=hN$ należy do $H/N$,
co daje inkluzję $H\subseteq q^{-1}(H/N)$.
Jeżeli $g\in q^{-1}(H/N)$, to $gN=hN$ dla pewnego $h\in H$,
a więc $h^{-1}g\in N\subseteq H$.
Ponieważ $h\in H$, otrzymujemy $g\in H$.

Dla podgrupy $K\le G/N$ mamy
$q(q^{-1}(K))=K$, bo $q$ jest surjektywny.
Ta równość oznacza dokładnie $(q^{-1}(K))/N=K$.
Oba przyporządkowania są więc wzajemnie odwrotne.
Oczywiście zachowują inkluzję, ponieważ zarówno obraz,
jak i przeciwobraz zachowują inkluzje odpowiednich zbiorów.

Jeśli $H\trianglelefteq G$, to dla $g\in G,h\in H$
mamy $(gN)(hN)(gN)^{-1}=(ghg^{-1})N\in H/N$,
więc $H/N\trianglelefteq G/N$.
W drugą stronę, gdy $H/N$ jest normalna, to
$(ghg^{-1})N\in H/N$, a z równości przeciwobrazu
$q^{-1}(H/N)=H$ wynika $ghg^{-1}\in H$.
Zatem $H\trianglelefteq G$.
\end{proof}

\begin{przyklad}[Podgrupy $\mathbb Z/12\mathbb Z$ z odpowiedniości]
Podgrupy $\mathbb Z$ zawierające $12\mathbb Z$ są dokładnie
$d\mathbb Z$ dla $d\mid12$:
z klasyfikacji podgrup $\mathbb Z$ wiemy, że są to
$\mathbb Z,2\mathbb Z,3\mathbb Z,4\mathbb Z,6\mathbb Z,12\mathbb Z$.
Po przejściu do ilorazu otrzymujemy sześć podgrup
$\mathbb Z/12\mathbb Z$. Jest to druga droga do listy
podgrup otrzymanej przy badaniu grup cyklicznych.
\end{przyklad}

\section{Sprzężenia, komutatory i abelianizacja}
\begin{definicja}[Sprzężenie, centrum, komutator]
Niech $G$ będzie grupą. Dla $g,x\in G$ element $gxg^{-1}$
nazywamy \emph{sprzężonym} do $x$. \emph{Centrum} grupy to
$Z(G)=\{z\in G:zx=xz\text{ dla każdego }x\in G\}$.
\emph{Komutatorem} elementów $x,y$ nazywamy
$[x,y]=xyx^{-1}y^{-1}$.
\emph{Podgrupa pochodna} lub \emph{komutatorowa}
$[G,G]$ to podgrupa generowana przez wszystkie komutatory
$[x,y]$ dla $x,y\in G$.
\end{definicja}

\begin{lemat}[Własności centrum i podgrupy komutatorowej]
\label{aa:lem:komutatory}
Dla każdej grupy $G$ zbiory $Z(G)$ i $[G,G]$ są normalnymi
podgrupami $G$. Ponadto $G/[G,G]$ jest abelowa i jeśli
$N\trianglelefteq G$ oraz $G/N$ jest abelowa, to $[G,G]\subseteq N$.
\end{lemat}
\begin{proof}
Centrum zawiera $e$. Jeśli $z,w\in Z(G)$ i $x\in G$,
to z równości $wx=xw$ po pomnożeniu z lewej i z prawej
przez $w^{-1}$ dostajemy $w^{-1}x=xw^{-1}$. Wobec tego
$(zw^{-1})x=z(w^{-1}x)=z(xw^{-1})=(xz)w^{-1}
=x(zw^{-1})$, więc centrum jest podgrupą.
Dla $z\in Z(G)$ i $g\in G$ mamy $gzg^{-1}=z$,
więc centrum jest normalne.

Dla dowolnego $g\in G$ sprzężenie komutatora jest komutatorem:
$g[x,y]g^{-1}=[gxg^{-1},gyg^{-1}]$.
Zatem sprzężenie przez $g$ zachowuje zbiór generatorów
$[G,G]$ i jego podgrupę; stosując także $g^{-1}$,
otrzymujemy równość sprzężonej podgrupy z $[G,G]$.

W ilorazie przez $[G,G]$ zachodzi dla dowolnych $x,y\in G$
\[
 (x[G,G])(y[G,G])(x[G,G])^{-1}(y[G,G])^{-1}
 =[x,y][G,G]=[G,G].
\]
To znaczy, że dowolne dwie klasy komutują.
Odwrotnie, jeśli $G/N$ jest abelowa, to
$[xN,yN]=N$ i stąd $[x,y]\in N$ dla wszystkich $x,y$.
Ponieważ $N$ jest podgrupą zawierającą wszystkie komutatory,
zawiera też podgrupę przez nie generowaną.
\end{proof}

\begin{definicja}[Abelianizacja]
\label{aa:def:abelianizacja}
Grupę abelową $G^{\mathrm{ab}}=G/[G,G]$ nazywamy
\emph{abelianizacją} grupy $G$.
\end{definicja}

\begin{twierdzenie}[Własność uniwersalna abelianizacji]
\label{aa:tw:abelianizacja}
Niech $G$ będzie grupą, a $A$ grupą abelową.
Każdy homomorfizm $f:G\to A$ ma dokładnie jedno
faktoryzowanie przez rzut $q:G\to G^{\mathrm{ab}}$,
tzn. istnieje dokładnie jeden homomorfizm
$\overline f:G^{\mathrm{ab}}\to A$ taki, że $f=\overline f\circ q$.
\end{twierdzenie}
\begin{proof}
Dla dowolnych $x,y\in G$ w grupie abelowej $A$ mamy
$f([x,y])=f(x)f(y)f(x)^{-1}f(y)^{-1}=e_A$.
Zatem wszystkie komutatory należą do $\ker f$.
Ponieważ jądro jest podgrupą, $[G,G]\subseteq\ker f$.
Możemy więc zdefiniować
$\overline f(g[G,G])=f(g)$.
Jeżeli $g[G,G]=h[G,G]$, to $h^{-1}g\in[G,G]\subseteq\ker f$,
więc $f(g)=f(h)$. Wzór nie zależy od reprezentanta.
Dla dowolnych $g,h\in G$ zachodzi
$\overline f((g[G,G])(h[G,G]))=f(gh)=f(g)f(h)$,
czyli $\overline f$ zachowuje mnożenie.
Równość $f=\overline f\circ q$
zachodzi dla każdego $g$ z konstrukcji.
Jednoznaczność wynika z surjektywności $q$: homomorfizm
$\overline f$ musi mieć wartość $f(g)$ na każdej klasie $q(g)$.
\end{proof}

\begin{przyklad}[Abelianizacja $S_3$]
\label{aa:prz:ab-s3}
Homomorfizm znaku $\operatorname{sgn}:S_3\to\{\pm1\}$
ma jądro $A_3$, więc iloraz $S_3/A_3$ jest abelowy.
Z lematu o komutatorach wynika $[S_3,S_3]\subseteq A_3$.
Z drugiej strony dla $x=(1\,2)$, $y=(2\,3)$ mamy
$[x,y]=xyxy=(xy)^2=(1\,3\,2)$,
ponieważ $x^2=y^2=e$ i $xy=(1\,2\,3)$.
Ten komutator generuje całe $A_3$, więc
$A_3\subseteq[S_3,S_3]$. Otrzymujemy
$[S_3,S_3]=A_3$ i $S_3^{\mathrm{ab}}\cong\mathbb Z/2\mathbb Z$.
\end{przyklad}

\section{Twierdzenie Cayleya i zamknięcie pierwszego etapu}
\begin{twierdzenie}[Cayleya]
\label{aa:tw:cayley}
Każda grupa $G$ jest izomorficzna z podgrupą grupy
$\operatorname{Sym}(G)$ wszystkich permutacji swojego zbioru elementów.
Jeżeli $|G|=n<\infty$, jest więc izomorficzna z podgrupą $S_n$.
\end{twierdzenie}
\begin{proof}
Dla każdego $g\in G$ zdefiniujmy funkcję
$L_g:G\to G$ przez $L_g(x)=gx$.
Jest bijekcją: $L_{g^{-1}}(L_g(x))=g^{-1}(gx)=x$
i $L_g(L_{g^{-1}}(x))=g(g^{-1}x)=x$.
Zatem $L_g\in\operatorname{Sym}(G)$.
Zdefiniujmy $\Lambda:G\to\operatorname{Sym}(G)$,
$\Lambda(g)=L_g$. Dla $g,h,x\in G$ mamy
$(L_g\circ L_h)(x)=g(hx)=(gh)x=L_{gh}(x)$.
Zatem $\Lambda(gh)=\Lambda(g)\Lambda(h)$:
$\Lambda$ jest homomorfizmem.

Aby wykazać różnowartościowość, nie trzeba porównywać
wszystkich wartości funkcji: jeśli $L_g=L_h$, to ich
wartości w $e$ są równe, czyli $g=L_g(e)=L_h(e)=h$.
Obraz $\Lambda(G)$ jest podgrupą $\operatorname{Sym}(G)$,
a $\Lambda:G\to\Lambda(G)$ jest izomorfizmem.
Dla $n$-elementowego $G$ wybierzmy bijekcję
$b:G\to\{1,\ldots,n\}$. Przyporządkowanie
$\sigma\mapsto b\circ\sigma\circ b^{-1}$
jest izomorfizmem $\operatorname{Sym}(G)\to S_n$,
a więc obraz $\Lambda(G)$ odpowiada podgrupie $S_n$.
\end{proof}

\begin{przyklad}[Regularna reprezentacja $\mathbb Z/3\mathbb Z$]
Niech $G=(\mathbb Z/3\mathbb Z,+)$, z elementami $[0],[1],[2]$.
Translacja $L_{[1]}$ daje cykl
$[0]\mapsto[1]\mapsto[2]\mapsto[0]$,
a więc po ponumerowaniu elementów jest $3$-cyklem.
Translacja $L_{[2]}=L_{[1]}^2$ jest cyklem przeciwnym,
a $L_{[0]}$ identycznością. W ten sposób
$G$ reprezentuje się jako trójelementowa podgrupa $S_3$.
Twierdzenie nie mówi, że $G=S_3$; mówi o podgrupie
izomorficznej z $G$.
\end{przyklad}

\begin{uwaga}[Co dalej?]
W kolejnym rozdziale działania grup, orbity i stabilizatory
dostarczą silniejszego języka do liczenia elementów i analizowania
sprzężeń; wyprowadzimy też równanie klas. W następnym rozdziale
przejdziemy do twierdzenia Cauchy'ego i twierdzeń Sylowa.
Zapis iloczynu tensorowego $A\otimes_{\mathbb Z}B$
dla \emph{grup abelowych} $A,B$ wprowadzimy w rozdziale o modułach.
Jest to iloczyn tensorowy $\mathbb Z$-modułów; dla ogólnych
nieabelowych grup standardowy iloczyn tensorowy modułów nie jest
tą samą konstrukcją i wymagałby innej definicji.
\end{uwaga}

% AA-03 | Algebra abstrakcyjna | kanoniczny rozdzial: dzialania grup
% Zachowano istniejace etykiety, twierdzenia, dowody i przyklady.
\chapter{Działania grup, orbity i równanie klas}
\label{aa:ch:dzialania}

\section{Dlaczego grupy działają na zbiorach?}
Grupa koduje działania odwracalne, lecz sama informacja o jej mnożeniu
nie mówi jeszcze, \emph{na czym} wykonujemy te operacje. Grupa
permutacji przestawia elementy zbioru, obroty przesuwają punkty sfery,
a macierze odwracalne przekształcają wektory. Wspólną definicję
otrzymamy, gdy narzucimy zgodność kolejnych operacji z mnożeniem w grupie.

\begin{definicja}[Lewe działanie grupy]
\label{aa:def:dzialanie}
Niech $G$ będzie grupą z elementem neutralnym $e$, a $X$ zbiorem.
\emph{Lewym działaniem} $G$ na $X$ nazywamy odwzorowanie
$G\times X\to X$, $(g,x)\mapsto g\cdot x$, spełniające
$e\cdot x=x$ i $(gh)\cdot x=g\cdot(h\cdot x)$
dla wszystkich $g,h\in G$ i $x\in X$.
Zbiór $X$ z takim działaniem nazywamy \emph{$G$-zbiorem}.
Punkt $x\in X$ jest \emph{stały} dla działania całej grupy,
jeżeli $g\cdot x=x$ dla wszystkich $g\in G$.
\end{definicja}

\begin{lemat}[Działanie to homomorfizm do grupy permutacji]
\label{aa:lem:akcja-hom}
Dane jest lewe działanie grupy $G$ na zbiorze $X$.
Przez $\operatorname{Sym}(X)$ oznaczamy grupę wszystkich bijekcji
$X\to X$ z działaniem składania odwzorowań. Dla każdego
$g\in G$ odwzorowanie $T_g:X\to X$, $T_g(x)=g\cdot x$,
jest bijekcją, a funkcja
$\rho:G\to\operatorname{Sym}(X)$, $\rho(g)=T_g$,
jest homomorfizmem grup.
Odwrotnie, każdy homomorfizm $\rho:G\to\operatorname{Sym}(X)$
wyznacza lewe działanie przez $g\cdot x=\rho(g)(x)$.
\end{lemat}
\begin{proof}
Z aksjomatu działania $T_e(x)=e\cdot x=x$, więc $T_e=\operatorname{id}_X$.
Ponadto $T_g(T_{g^{-1}}(x))=(gg^{-1})\cdot x=x$;
po zamianie $g$ z $g^{-1}$ dostajemy także
$T_{g^{-1}}(T_g(x))=x$. Stąd $T_g^{-1}=T_{g^{-1}}$,
a więc $T_g$ jest bijekcją.
Dla wszystkich $g,h,x$ mamy
$T_{gh}(x)=(gh)\cdot x=g\cdot(h\cdot x)=(T_g\circ T_h)(x)$.
Równość dla każdego $x$ dowodzi $\rho(gh)=\rho(g)\rho(h)$.
W drugą stronę homomorfizm zachowuje element neutralny,
więc $\rho(e)=\operatorname{id}_X$, a równość
$\rho(gh)=\rho(g)\circ\rho(h)$ daje oba aksjomaty działania.
\end{proof}

\begin{definicja}[Działanie wierne i przechodnie]
Niech grupa $G$ działa na zbiorze $X$, a
$\rho:G\to\operatorname{Sym}(X)$ będzie homomorfizmem
odpowiadającym temu działaniu. Działanie jest \emph{wierne}, jeżeli
$\rho$ jest różnowartościowy;
jest \emph{przechodnie}, jeżeli dla wszystkich $x,y\in X$
istnieje $g\in G$ z $g\cdot x=y$.
Jądro działania to $\ker\rho$.
\end{definicja}

\begin{przyklad}[Trzy działania jednej grupy]
\label{aa:prz:trzy-akcje}
Niech $G=S_3$.
\begin{enumerate}
\item $S_3$ działa na $\{1,2,3\}$ przez $\sigma\cdot i=\sigma(i)$.
Jest to działanie wierne i przechodnie: każdą liczbę można
przenieść na dowolną inną, a permutacja ustalająca wszystkie
trzy liczby jest identycznością.
\item $S_3$ działa na zbiorze własnych elementów przez mnożenie
z lewej: $\sigma\cdot\tau=\sigma\tau$. To także działanie wierne
i przechodnie: aby przenieść $\tau$ na $\eta$, bierzemy
$\sigma=\eta\tau^{-1}$. Działanie to ma sześciopunktową,
a nie trzypunktową przestrzeń.
\item $S_3$ działa na sobie przez sprzężenie
$\sigma\cdot\tau=\sigma\tau\sigma^{-1}$.
Nie jest przechodnie: element neutralny przechodzi tylko na siebie.
Każda permutacja w centrum działa trywialnie; pokażemy,
że właśnie centrum stanowi jądro tego działania.
\end{enumerate}
To, że grupy są te same, nie oznacza, że ich działania są takie same.
\end{przyklad}

\begin{przyklad}[Działanie przechodnie nie musi być wierne]
\label{aa:prz:akcja-c4-dwa-punkty}
Niech $G=\mathbb Z/4\mathbb Z$, a $X=\{u,v\}$.
Dla klasy $[k]\in G$ określamy $[k]\cdot u=u$ i
$[k]\cdot v=v$, jeśli $k$ jest parzyste, oraz
$[k]\cdot u=v$, $[k]\cdot v=u$, jeśli $k$ jest nieparzyste.
Parzystość nie zależy od reprezentanta klasy modulo $4$.
Zgodność działania z dodawaniem wynika stąd, że złożenie
zamiany punktów ze sobą daje identyczność, a parzystość sumy
jest sumą parzystości modulo $2$.
Zatem $[0]$ działa identycznościowo, a
$([k]+[\ell])\cdot x=[k]\cdot([\ell]\cdot x)$ dla
$k,\ell\in\mathbb Z$ i $x\in X$.
Działanie jest przechodnie: $[1]\cdot u=v$ i $[1]\cdot v=u$,
a $[0]$ ustala oba punkty.
Nie jest wierne: $[2]\ne[0]$ ustala oba punkty,
a dokładnie klasy parzyste działają trywialnie,
więc $\ker\rho=\{[0],[2]\}$.
Dokładnie dwie klasy $[0]$ i $[2]$ ustalają punkt $u$,
a jego możliwymi obrazami są oba punkty $u,v$.
Zależność $2\cdot2=4$ wyjaśnimy ogólnie w twierdzeniu
o orbicie i stabilizatorze. Przykład pokazuje, że
przechodniość nie pociąga wierności.

\begin{center}
\begin{tikzpicture}[>=Latex,baseline=(current bounding box.center)]
\node[draw,circle,minimum size=9mm] (u) at (0,0) {$u$};
\node[draw,circle,minimum size=9mm] (v) at (3,0) {$v$};
\draw[->,bend left=24] (u) to node[above] {$[1]$} (v);
\draw[->,bend left=24] (v) to node[below] {$[1]$} (u);
\node[align=center] at (1.5,-1.55) {$[2]$ ustala oba punkty};
\end{tikzpicture}
\end{center}
\end{przyklad}

\section{Orbity i stabilizatory}
\begin{definicja}[Orbita i stabilizator]
\label{aa:def:orbita-stabilizator}
Niech $G\curvearrowright X$ będzie działaniem, a $x\in X$.
\emph{Orbitą} punktu $x$ nazywamy zbiór
$G\cdot x=\{g\cdot x:g\in G\}$.
\emph{Stabilizatorem} punktu $x$ nazywamy
$G_x=\{g\in G:g\cdot x=x\}$.
\end{definicja}

\begin{lemat}[Stabilizator jest podgrupą, orbity rozkładają zbiór]
\label{aa:lem:orbity-podzial}
Niech $G$ działa na zbiorze $X$.
Dla każdego $x\in X$ zbiór $G_x$ jest podgrupą $G$.
Relacja $x\sim y\iff y\in G\cdot x$ jest relacją równoważności,
a jej klasy to dokładnie orbity działania.
\end{lemat}
\begin{proof}
Neutralny element $e$ ustala $x$, więc $e\in G_x$.
Jeżeli $g,h\in G_x$, to
$(gh^{-1})\cdot x=g\cdot(h^{-1}\cdot x)=g\cdot x=x$:
istotnie, skoro $h\cdot x=x$, zastosowanie $h^{-1}$
po obu stronach daje $h^{-1}\cdot x=x$.
Kryterium podgrupy z Rozdziału~\ref{aa:ch:grupy} daje $G_x\le G$.

Relacja jest zwrotna, bo $x=e\cdot x$.
Jeśli $y=g\cdot x$, to $x=g^{-1}\cdot y$, zatem jest symetryczna.
Jeśli ponadto $z=h\cdot y$, to $z=(hg)\cdot x$, zatem
jest przechodnia. Jej klasą równoważności punktu $x$
jest zbiór wszystkich $g\cdot x$, czyli $G\cdot x$.
Klasy równoważności są albo równe, albo rozłączne
i pokrywają $X$.
\end{proof}

\begin{twierdzenie}[Orbita i stabilizator]
\label{aa:tw:orbita-stabilizator}
Niech grupa $G$ działa na zbiorze $X$, a $x\in X$.
Odwzorowanie
$\Theta_x:G/G_x\to G\cdot x$, $\Theta_x(gG_x)=g\cdot x$,
jest dobrze określoną bijekcją zbiorów. W szczególności,
gdy $G$ jest skończona, zachodzi
\[
 |G\cdot x|=[G:G_x]=\frac{|G|}{|G_x|}.
\]
\end{twierdzenie}
\begin{proof}
Najpierw porównajmy dwa reprezentanty warstwy:
$gG_x=hG_x$ dokładnie wtedy, gdy $h^{-1}g\in G_x$,
czyli $(h^{-1}g)\cdot x=x$.
Z aksjomatów działania ostatnia równość jest równoważna
$g\cdot x=h\cdot x$: zastosowanie $h$ daje implikację
w prawo, a zastosowanie $h^{-1}$ daje odwrotną.
Wykazaliśmy jednocześnie, że wartość $\Theta_x$
nie zależy od reprezentanta warstwy i że równe obrazy
pochodzą tylko z równych warstw; jest więc różnowartościowa.
Każdy element orbity ma postać $g\cdot x$ i jest obrazem
warstwy $gG_x$, zatem $\Theta_x$ jest surjektywna.
Dla skończonego $G$ stosujemy twierdzenie Lagrange'a
(\ref{aa:tw:lagrange}) do $G_x\le G$.
\end{proof}

\begin{przyklad}[Naturalne działanie $S_4$]
Niech $S_4$ działa na $\{1,2,3,4\}$ w zwykły sposób.
Orbita punktu $1$ to cały zbiór czteroelementowy:
permutacją można wysłać $1$ na dowolną liczbę $i$.
Stabilizator $1$ składa się z permutacji pozostałych trzech
liczb i ma $3!=6$ elementów. Twierdzenie daje
$|S_4\cdot1|=|S_4|/|(S_4)_1|=24/6=4$.
Widać, co dzieje się w dowodzie: sześć permutacji
wysyłających $1$ na określoną liczbę tworzy jedną warstwę
stabilizatora; cztery możliwe obrazy to cztery takie warstwy.
\end{przyklad}

\begin{przyklad}[Nieskończona grupa z orbitą skończoną]
\label{aa:prz:z-dziala-na-zn}
Niech $n\ge2$ będzie liczbą naturalną.
Grupa $(\mathbb Z,+)$ działa na $X=\mathbb Z/n\mathbb Z$
przez $k\cdot[a]=[a+k]$.
Jeśli $[a]=[b]$, to $n$ dzieli $a-b$, a zatem dzieli
$(a+k)-(b+k)$, więc działanie jest dobrze określone.
Ponadto $0\cdot[a]=[a]$ i
$(k+\ell)\cdot[a]=[a+k+\ell]=k\cdot(\ell\cdot[a])$.
Jest przechodnie: aby przenieść $[a]$ na $[b]$, wystarczy
wybrać $k=b-a$.
Stabilizator $[0]$ to $n\mathbb Z$, gdyż
$k\cdot[0]=[0]$ dokładnie wtedy, gdy $n\mid k$.
Warstwy $\mathbb Z/n\mathbb Z$ grupy addytywnej
$\mathbb Z$ względem $n\mathbb Z$ są dokładnie $n$
klasami reszt, więc
$|\mathbb Z\cdot[0]|=[\mathbb Z:n\mathbb Z]=n$.
Nie stosujemy tutaj ilorazu $|\mathbb Z|/|n\mathbb Z|$:
oba zbiory są nieskończone, natomiast ich indeks jest skończony.
\end{przyklad}

\begin{przyklad}[Działanie na klasach warstw]
\label{aa:prz:akcja-kosety}
Niech $H\le G$. Grupa $G$ działa na zbiorze $G/H$
przez $g\cdot(aH)=(ga)H$.
Wzór jest dobrze określony: jeśli $aH=bH$, to
$b^{-1}a\in H$, a więc
$(gb)^{-1}(ga)=b^{-1}a\in H$, czyli $(ga)H=(gb)H$.
Aksjomaty działania wynikają z
$e\cdot(aH)=aH$ oraz
$(gh)\cdot(aH)=((gh)a)H=g\cdot((ha)H)
=g\cdot(h\cdot(aH))$.
Jest to działanie przechodnie, ponieważ
$(ba^{-1})\cdot(aH)=bH$.
Stabilizator warstwy $H$ to dokładnie $H$, gdyż
$gH=H\iff g\in H$.
Stabilizator warstwy $aH$ wynosi $aHa^{-1}$:
$g\cdot aH=aH$ wtedy i tylko wtedy, gdy
$a^{-1}ga\in H$.
To ważny model wszystkich działań przechodnich.
\end{przyklad}

\begin{wniosek}[Każdy przechodni $G$-zbiór jest przestrzenią warstw]
\label{aa:wn:przechodnie-warstwy}
Niech grupa $G$ działa przechodnio na niepustym zbiorze $X$.
Dla dowolnego $x\in X$ istnieje bijekcja
$G/G_x\to X$, $gG_x\mapsto g\cdot x$,
która jest zgodna z działaniem $G$ po obu stronach.
\end{wniosek}
\begin{proof}
Z przechodniości $G\cdot x=X$, więc bijekcja wynika
z twierdzenia o orbicie i stabilizatorze.
Dla $a,g\in G$ zgodność działania oznacza
$\Theta_x(a\cdot(gG_x))=\Theta_x((ag)G_x)
=(ag)\cdot x=a\cdot(g\cdot x)=a\cdot\Theta_x(gG_x)$.
\end{proof}

\section{Jądro działania i normalność}
\begin{twierdzenie}[Jądro i efektywne działanie]
\label{aa:tw:jadro-akcji}
Jeżeli grupa $G$ działa na zbiorze $X$ przez homomorfizm
$\rho:G\to\operatorname{Sym}(X)$, to
\[
 \ker\rho=\{g\in G:g\cdot x=x\text{ dla każdego }x\in X\}
 =\bigcap_{x\in X}G_x.
\]
Zbiór ten jest podgrupą normalną, a działanie jest wierne
dokładnie wtedy, gdy to jądro jest trywialne.
Działanie przechodzi na wierne działanie $G/\ker\rho$
na tym samym zbiorze $X$.
\end{twierdzenie}
\begin{proof}
Z definicji $g\in\ker\rho$ oznacza
$\rho(g)=\operatorname{id}_X$, czyli
$\rho(g)(x)=x$ dla każdego $x\in X$.
To jest równoważne przynależności do wszystkich stabilizatorów.
Jądro homomorfizmu jest normalne, a homomorfizm jest
różnowartościowy dokładnie wtedy, gdy jego jądro jest trywialne
(Rozdział~\ref{aa:ch:lagrange-ilorazy}).
Niech $N=\ker\rho$. Definiujemy
$(gN)\cdot x=g\cdot x$.
Jeśli $gN=hN$, to $h^{-1}g\in N$;
więc $(h^{-1}g)\cdot x=x$, co po zastosowaniu $h$
daje $g\cdot x=h\cdot x$. Działanie jest dobrze określone. Ponadto
$(eN)\cdot x=e\cdot x=x$ oraz
$((gN)(hN))\cdot x=(gh)\cdot x
=g\cdot(h\cdot x)=(gN)\cdot((hN)\cdot x)$.
Jeśli warstwa $gN$ ustala wszystkie $x$, to $g\in N$
i $gN=N$, więc nowe działanie jest wierne.
\end{proof}

\begin{przyklad}[Działanie na warstwach i rdzeń podgrupy]
Niech $H\le G$ i rozważmy działanie na $G/H$.
Element $g$ działa trywialnie wtedy i tylko wtedy, gdy
$g$ ustala wszystkie warstwy $aH$.
Stabilizator $aH$ to $aHa^{-1}$, stąd jądro działania
wynosi $\bigcap_{a\in G}aHa^{-1}$.
Jest to \emph{rdzeń normalny} podgrupy $H$, czyli
największa normalna podgrupa $G$ zawarta w $H$.
Rzeczywiście jest normalna jako jądro i zawiera się
w $H$ (bierzemy $a=e$).
Jeżeli $N\trianglelefteq G$ oraz $N\subseteq H$, to
$aNa^{-1}=N\subseteq aHa^{-1}$ dla każdego $a$,
a więc $N$ zawiera się w tym przecięciu.
\end{przyklad}

\begin{przyklad}[Wierne działanie na warstwach nie wymaga normalności]
\label{aa:prz:s3-core}
W $G=S_3$ niech $H=\langle(2\,3)\rangle$, czyli
stabilizator punktu $1$ w naturalnym działaniu na
$\{1,2,3\}$. Zbiór warstw $G/H$ ma trzy elementy.
Odwzorowanie $gH\mapsto g(1)$ jest bijekcją zgodną
z działaniem $G$; wynika to z Twierdzenia~\ref{aa:tw:orbita-stabilizator}.
Jądro działania na trzech warstwach jest trywialne, bo
permutacja ustalająca jednocześnie punkty $1,2,3$
jest identycznością. Zatem rdzeń normalny $H$ jest
równy $\{e\}$. Sama $H$ nie jest jednak normalna:
$(1\,2)(2\,3)(1\,2)=(1\,3)\notin H$.
W szczególności zbiór $G/H$ ma naturalne działanie $G$,
lecz nie ma dobrze określonego mnożenia warstw
$(gH)(hH)=(gh)H$, które tworzyłoby grupę ilorazową.
Rozróżnienie to wykorzystamy w następnym rozdziale,
w dowodzie drugiego twierdzenia Sylowa.
\end{przyklad}

\section{Sprzężenia, centralizatory i normalizatory}
\begin{definicja}[Centralizator i normalizator]
\label{aa:def:centralizator-normalizator}
Niech $G$ będzie grupą, $x\in G$, $H\le G$.
\emph{Centralizator} $x$ to
$C_G(x)=\{g\in G:gx=xg\}$, a \emph{normalizator} $H$ to
$N_G(H)=\{g\in G:gHg^{-1}=H\}$.
\emph{Klasa sprzężoności} $x$ to
$\operatorname{Cl}_G(x)=\{gxg^{-1}:g\in G\}$.
\end{definicja}

\begin{lemat}[Centralizatory i normalizatory jako stabilizatory]
\label{aa:lem:centralizator-normalizator}
Dla dowolnej grupy $G$, elementu $x\in G$ i podgrupy $H\le G$
zbiory $C_G(x)$ i $N_G(H)$ są podgrupami $G$.
Ponadto $H\trianglelefteq N_G(H)$ oraz
\[
 |\operatorname{Cl}_G(x)|=[G:C_G(x)]
\]
przy skończonym $G$.
\end{lemat}
\begin{proof}
Rozważmy działanie $G$ na sobie przez sprzężenie:
$g\cdot x=gxg^{-1}$.
Aksjomaty działania wynikają z
$e x e^{-1}=x$ oraz
$(gh)x(gh)^{-1}=g(hxh^{-1})g^{-1}$.
Stabilizatorem $x$ jest zbiór tych $g$, dla których
$gxg^{-1}=x$, czyli $gx=xg$.
Zatem $C_G(x)=G_x$ i jest podgrupą, a formuła
na liczebność klasy sprzężoności wynika z twierdzenia
o orbicie i stabilizatorze.

Grupa $G$ działa również na zbiorze wszystkich podgrup $G$
przez $g\cdot K=gKg^{-1}$.
Aksjomaty działania sprawdza się jak wyżej,
a sprzężona podgrupa jest podgrupą dzięki zachowaniu
mnożenia i odwrotności przez sprzężenie.
Stabilizatorem podgrupy $H$ jest $N_G(H)$, stąd jest podgrupą.
Dla każdego $h\in H$ zachodzi $hHh^{-1}=H$,
więc $H\le N_G(H)$. Ponadto każdy element $N_G(H)$
sprzęga $H$ na $H$, czyli $H\trianglelefteq N_G(H)$.
\end{proof}

\begin{przyklad}[Klasy sprzężoności w $S_3$]
W $S_3$ są: identyczność, trzy transpozycje i dwa $3$-cykle.
Przez sprzężenie permutacje zachowują strukturę cykli:
dla bijekcji $g$ zachodzi
$g(a_1\,\ldots\,a_k)g^{-1}=(g(a_1)\,\ldots\,g(a_k))$,
co sprawdzamy, stosując obie strony do $g(a_i)$;
poza punktami $g(a_1),\ldots,g(a_k)$ obie strony
ustalają pozostałe punkty. Dla dowolnych dwóch cykli
tej samej długości możemy wybrać permutację $g$, która
przenosi kolejne elementy pierwszego cyklu na kolejne
elementy drugiego. Stąd trzy transpozycje są sprzężone,
a dwa $3$-cykle również są sprzężone.
Centralizator transpozycji $(1\,2)$ zawiera tylko
$e$ i $(1\,2)$: jeśli $g(1\,2)g^{-1}=(1\,2)$,
to $g$ musi zachować dwuelementowy zbiór $\{1,2\}$,
co pozostawia dwa wybory. Ma więc rząd $2$ i klasa
ma $6/2=3$ elementy.
Centralizator $(1\,2\,3)$ to jego trzy potęgi:
możemy wysłać $1$ na jeden z trzech punktów,
a komutowanie z cyklem wymusza obrazy $2$ i $3$.
Zatem klasa $3$-cykli ma $6/3=2$ elementy.
Ostatecznie klasy sprzężoności mają rozmiary $1,3,2$.
\end{przyklad}

\begin{przyklad}[Jądro działania przez sprzężenie]
Dla dowolnej grupy $G$ jądro działania przez sprzężenie
składa się z elementów $g$ spełniających
$gxg^{-1}=x$ dla wszystkich $x\in G$, czyli z $Z(G)$.
Zatem $G/Z(G)$ działa wiernie na $G$ przez sprzężenie.
Przez $\operatorname{Aut}(G)$ oznaczamy grupę
wszystkich automorfizmów $G$ ze składaniem odwzorowań.
Otrzymujemy przy tym homomorfizm
$G/Z(G)\to\operatorname{Aut}(G)$ dany wzorem
$gZ(G)\mapsto(x\mapsto gxg^{-1})$. Jest on dobrze określony
na mocy opisu jądra działania. Sprzężenie przez $g$
zachowuje mnożenie, a jego odwrotnością jest sprzężenie
przez $g^{-1}$, więc rzeczywiście jest automorfizmem $G$.
Zgodność z mnożeniem warstw wynika ze wzoru
$(gh)x(gh)^{-1}=g(hxh^{-1})g^{-1}$.
\end{przyklad}

\section{Równanie klas i zastosowania do \texorpdfstring{$p$}{p}-grup}
\begin{twierdzenie}[Równanie klas]
\label{aa:tw:rownanie-klas}
Niech $G$ będzie grupą skończoną.
Wybierzmy reprezentantów $x_1,\ldots,x_r$ wszystkich
klas sprzężoności niezawartych w centrum $Z(G)$.
Wówczas
\[
 |G|=|Z(G)|+\sum_{i=1}^{r}[G:C_G(x_i)].
\]
Każdy indeks w występującej sumie jest liczbą całkowitą
nie mniejszą niż $2$.
\end{twierdzenie}
\begin{proof}
Orbity działania przez sprzężenie tworzą rozkład $G$.
Orbita $\operatorname{Cl}_G(x)$ ma jeden element wtedy i tylko wtedy,
gdy $gxg^{-1}=x$ dla wszystkich $g$, co jest równoważne
$x\in Z(G)$. Orbity jednoelementowe składają się zatem
z dokładnie $|Z(G)|$ elementów.
Pozostałe orbity są parami rozłączne i mają rozmiary
$[G:C_G(x_i)]$, ponieważ centralizator jest stabilizatorem.
Sumując liczby elementów wszystkich orbit, otrzymujemy wzór.
Jeśli $x_i\notin Z(G)$, to jego orbita ma więcej niż jeden element.
\end{proof}

\begin{definicja}[Grupa o rzędzie będącym potęgą liczby pierwszej]
Niech $p$ będzie liczbą pierwszą, a $a\ge0$ liczbą całkowitą.
Grupę o rzędzie $p^a$ nazywamy \emph{$p$-grupą}.
Jeżeli $a=0$, jest to grupa trywialna, ponieważ ma tylko
jeden element. W dalszych twierdzeniach o nietrywialnym
centrum zakładamy $a\ge1$.
\end{definicja}

\begin{twierdzenie}[Centrum niebanalnej skończonej $p$-grupy]
\label{aa:tw:centrum-pgrupy}
Jeżeli $G$ jest grupą rzędu $p^a$, gdzie $p$ jest liczbą pierwszą
i $a\ge1$, to $|Z(G)|$ jest podzielne przez $p$.
W szczególności centrum zawiera element różny od $e$.
\end{twierdzenie}
\begin{proof}
Dla każdego $x\in G$ centralizator $C_G(x)$ jest podgrupą,
więc jego rząd jest potęgą $p$ na mocy twierdzenia Lagrange'a.
Jeśli $x\notin Z(G)$, to $C_G(x)$ jest właściwą podgrupą:
rzeczywiście gdyby $C_G(x)=G$, element $x$ komutowałby
ze wszystkimi elementami grupy i należałby do centrum.
A zatem indeks $[G:C_G(x)]$ jest dodatnią potęgą $p$,
i wynosi co najmniej $p$.
Równanie klas daje
$|G|=|Z(G)|+\sum_i [G:C_G(x_i)]$.
Obie liczby $|G|=p^a$ i każdy składnik sumy są podzielne
przez $p$, więc ich różnica $|Z(G)|$ również.
Ponieważ $e\in Z(G)$, centrum nie jest puste;
jeśli jego liczebność jest dodatnią wielokrotnością $p$,
to wynosi co najmniej $p\ge2$.
\end{proof}

\begin{lemat}[Cykliczny iloraz przez centrum wymusza przemienność]
\label{aa:lem:cykliczny-centrum}
Jeśli $G$ jest grupą i $G/Z(G)$ jest cykliczna,
to $G$ jest abelowa.
\end{lemat}
\begin{proof}
Niech $gZ(G)$ generuje $G/Z(G)$.
Dla dowolnych $x,y\in G$ istnieją liczby całkowite $m,n$
i elementy $z,w\in Z(G)$ z $x=g^m z$, $y=g^n w$.
Korzystając z tego, że $z,w$ komutują ze wszystkimi
elementami $G$, otrzymujemy
$xy=g^{m+n}zw=g^{n+m}wz=yx$.
Ponieważ $x,y$ były dowolne, grupa jest abelowa.
\end{proof}

\begin{twierdzenie}[Grupy rzędu $p^2$ są abelowe]
\label{aa:tw:p2-abelowa}
Niech $p$ będzie liczbą pierwszą, a $G$ grupą rzędu $p^2$.
Wtedy $G$ jest abelowa.
\end{twierdzenie}
\begin{proof}
Poprzednie twierdzenie daje $p\mid |Z(G)|$,
a twierdzenie Lagrange'a daje $|Z(G)|\mid p^2$.
Stąd $|Z(G)|$ wynosi $p$ albo $p^2$.
Jeśli $|Z(G)|=p^2$, to $Z(G)=G$ i grupa jest abelowa.
Jeśli $|Z(G)|=p$, to iloraz $G/Z(G)$ ma rząd $p$,
a więc jest cykliczny. Poprzedni lemat ponownie daje
przemienność $G$.
\end{proof}

\begin{przyklad}[Dwie grupy abelowe rzędu cztery]
Grupa $\mathbb Z/4\mathbb Z$ ma element $[1]$ rzędu $4$,
natomiast $(\mathbb Z/2\mathbb Z)^2$ ma trzy niezerowe
elementy, z których każdy ma rząd $2$.
Obie grupy mają rząd $4$ i są abelowe,
lecz nie są izomorficzne: izomorfizm zachowuje rzędy elementów.
Z samej abelowości nie wynika więc cykliczność.
\end{przyklad}

\section{Lemat Burnside'a o orbitach}
\begin{definicja}[Zbiór punktów stałych elementu grupy]
Niech grupa $G$ działa na zbiorze $X$.
Dla $g\in G$ definiujemy
$X^g=\{x\in X:g\cdot x=x\}$.
Nie mylimy $X^g$ ze zbiorem punktów stałych całego działania,
który jest przecięciem wszystkich $X^g$.
\end{definicja}

\begin{twierdzenie}[Lemat Burnside'a]
\label{aa:tw:burnside}
Niech skończona grupa $G$ działa na skończonym zbiorze $X$.
Niech $X/G=\{G\cdot x:x\in X\}$ oznacza zbiór orbit.
Wtedy liczba orbit $|X/G|$ spełnia
\[
 |X/G|=\frac{1}{|G|}\sum_{g\in G}|X^g|.
\]
\end{twierdzenie}
\begin{proof}
Policzymy na dwa sposoby zbiór par
$P=\{(g,x)\in G\times X:g\cdot x=x\}$.
Ustalając najpierw $g$, otrzymujemy
$|P|=\sum_{g\in G}|X^g|$.
Ustalając natomiast $x$, widzimy, że odpowiadające mu
pierwsze współrzędne tworzą stabilizator $G_x$,
a więc $|P|=\sum_{x\in X}|G_x|$.

Podzielmy teraz $X$ na orbity. Jeżeli $O=G\cdot x$
jest jedną z nich, to wszystkie stabilizatory punktów z $O$
mają ten sam rząd: dla $y=h\cdot x$ mamy
$G_y=hG_xh^{-1}$, co wynika z rachunku
$g\cdot(h\cdot x)=h\cdot x\iff
(h^{-1}gh)\cdot x=x$.
Dlatego
$\sum_{y\in O}|G_y|=|O|\,|G_x|$.
Twierdzenie o orbicie i stabilizatorze daje
$|O|\,|G_x|=|G|$.
Każda orbita wnosi zatem dokładnie $|G|$ do sumy
$\sum_{x\in X}|G_x|$.
Jeśli orbit jest $k$, to $|P|=k|G|$.
Porównując oba sposoby liczenia $P$ i dzieląc przez $|G|$,
otrzymujemy żądany wzór.
\end{proof}

\begin{przyklad}[Kolorowanie wierzchołków kwadratu dwoma kolorami]
\label{aa:prz:burnside-kwadrat}
Ustalmy cztery wierzchołki kwadratu i dwa kolory: czarny i biały.
Kolorowanie to funkcja ze zbioru czterech wierzchołków
do zbioru dwuelementowego, więc jest ich $2^4=16$.
Dwa kolorowania uznajemy za to samo, gdy jedno można otrzymać
z drugiego przez symetrię kwadratu. Działa grupa $D_4$
o ośmiu elementach. Jeśli $V$ oznacza zbiór wierzchołków,
a $C=\{\text{czarny},\text{biały}\}$ zbiór kolorów,
to przestrzenią kolorowań jest zbiór funkcji $C^V$.
Dla symetrii $g\in D_4$, kolorowania $c:V\to C$ oraz $v\in V$
definiujemy $(g\cdot c)(v)=c(g^{-1}v)$.
To jest działanie: $e\cdot c=c$, zaś
$((gh)\cdot c)(v)=c(h^{-1}g^{-1}v)
=(g\cdot(h\cdot c))(v)$.
Kolorowanie $c$ jest ustalane przez $g$ wtedy i tylko wtedy,
gdy jest stałe na każdym cyklu permutacji wierzchołków
wyznaczonej przez $g$. Jeżeli cykli jest $k$, można
wybrać jeden z dwóch kolorów niezależnie dla każdego cyklu;
zatem $| (C^V)^g |=2^k$.
Obliczmy dokładnie liczby kolorowań zachowanych przez symetrie.
\begin{enumerate}
\item Identyczność ustala wszystkie $16$ kolorowań.
\item Każdy z dwóch obrotów o $90^\circ$ i $270^\circ$
cyklicznie permutuje wszystkie cztery wierzchołki,
więc kolorowanie jest stałe na całym kwadracie:
każdy z tych obrotów ustala $2$ kolorowania.
\item Obrót o $180^\circ$ wymienia dwie pary wierzchołków,
więc można niezależnie wybrać kolor każdej pary:
ustala $2^2=4$ kolorowania.
\item Każde z dwóch odbić względem przekątnych ustala
dwa wierzchołki i zamienia pozostałe dwa.
Wymaga to jednego wspólnego koloru dla zamienianej pary,
a kolory dwóch ustalanych wierzchołków można wybrać swobodnie:
każde takie odbicie ustala $2^3=8$ kolorowań.
\item Każde z dwóch odbić względem osi przechodzących przez
środki przeciwległych boków zamienia dwie pary wierzchołków.
Ustala $2^2=4$ kolorowania.
\end{enumerate}
W każdym przypadku liczba kolorowań jest potęgą $2$
o wykładniku równym liczbie cykli odpowiedniej permutacji.
Łączna liczba punktów stałych wynosi
$16+2+2+4+8+8+4+4=48$.
Na mocy lematu Burnside'a liczba różnych kolorowań z dokładnością
do symetrii wynosi $48/8=6$.
Można to sprawdzić, klasyfikując kolorowania według liczby czarnych
wierzchołków: $0$ i $4$ dają po jednej klasie, $1$ i $3$ także
po jednej, a dla dokładnie dwóch czarnych wierzchołków
są dwa rodzaje (sąsiednie lub przeciwległe). Razem sześć.
\end{przyklad}

\begin{uwaga}[Konkretni reprezentanci sześciu orbit]
W obliczonym przykładzie sześć orbit można reprezentować przez
kolorowania z odpowiednio: zerem czarnych wierzchołków,
jednym czarnym, dwoma sąsiednimi czarnymi,
dwoma naprzeciwległymi czarnymi, trzema czarnymi i czterema
czarnymi wierzchołkami. Odbicia i obroty nie zmieniają
liczby czarnych wierzchołków, a w przypadku dokładnie dwóch
czarnych wierzchołków zachowują również ich sąsiedztwo.
Dlatego sześć wymienionych przypadków jest parami nierównoważnych.
Z drugiej strony symetrie kwadratu przeprowadzają dowolne
kolorowanie danego rodzaju na wybranego reprezentanta,
co daje niezależną kontrolę wyniku lematu Burnside'a.
\end{uwaga}

\begin{uwaga}[Kolejny krok: twierdzenia Sylowa]
W całym rozdziale wykorzystaliśmy działania grup na punktach,
warstwach, podgrupach i samych grupach.
W następnym rozdziale zastosujemy je do podgrup o rzędach
będących potęgami liczb pierwszych. W ten sposób wyprowadzimy
klasyczne twierdzenia Cauchy'ego i Sylowa.
\end{uwaga}

% AA-04 | Algebra abstrakcyjna | Iteracja 02
\chapter{Twierdzenie Cauchy'ego i trzy twierdzenia Sylowa}
\label{aa:ch:sylow}

\section{Od twierdzenia Lagrange'a do problemu istnienia podgrup}
Twierdzenie Lagrange'a mówi, że jeśli $H\le G$ i $G$ jest
skończona, to $|H|$ dzieli $|G|$. Nie wynika z niego, że
każdy dzielnik $|G|$ jest rzędem podgrupy: grupa $A_4$ rzędu
$12$ nie ma podgrupy rzędu $6$.
Nasuwa się jednak słabsze pytanie: jeśli liczba pierwsza $p$
dzieli $|G|$, czy istnieje w $G$ element rzędu $p$?
Odpowiedź jest twierdząca. Co więcej, jeśli $p^a$ jest
największą potęgą $p$ dzielącą $|G|$, to istnieje nawet
podgrupa rzędu $p^a$. Dowiedziemy obu faktów bez zakładania
jakiejkolwiek klasyfikacji grup skończonych.

\begin{definicja}[Podgrupa Sylowa]
\label{aa:def:sylow}
Niech $G$ będzie skończoną grupą, $p$ liczbą pierwszą,
a $|G|=p^a m$, gdzie $a$ jest nieujemną liczbą całkowitą,
$m$ jest dodatnią liczbą całkowitą i $p\nmid m$.
\emph{Podgrupą $p$-Sylowa} grupy $G$ nazywamy podgrupę
rzędu $p^a$. Zbiór wszystkich takich podgrup oznaczamy
$\operatorname{Syl}_p(G)$, a ich liczbę przez $n_p(G)$.
Jeżeli $p\nmid|G|$, jedyną podgrupą $p$-Sylowa jest $\{e\}$.
\end{definicja}

\begin{uwaga}[Dlaczego rząd $p^a$ jest maksymalny?]
Jeżeli podgrupa $Q\le G$ ma rząd będący potęgą $p$,
to twierdzenie Lagrange'a daje $|Q|\mid p^a m$.
Stąd $|Q|=p^b$ dla pewnego $b\le a$. Podgrupa Sylowa
osiąga największy możliwy rząd wśród takich podgrup.
Samo pojęcie jest sensowne przed dowodem istnienia,
lecz dowód istnienia jest główną treścią pierwszego twierdzenia Sylowa.
\end{uwaga}

\section{Twierdzenie Cauchy'ego}
\begin{lemat}[Orbity działania grupy rzędu $p$]
\label{aa:lem:orbit-cp}
Niech $p$ będzie liczbą pierwszą, a cykliczna grupa $C_p$
rzędu $p$ niech działa na skończonym zbiorze $X$.
Każda orbita ma dokładnie $1$ albo $p$ elementów.
Jeżeli $F$ oznacza zbiór punktów ustalanych przez wszystkie
elementy $C_p$, to $|X|\equiv |F|\pmod p$.
\end{lemat}
\begin{proof}
Stabilizator każdego punktu jest podgrupą $C_p$.
Z twierdzenia Lagrange'a jego rząd dzieli $p$,
więc wynosi $1$ albo $p$.
Z twierdzenia o orbicie i stabilizatorze
rozmiar orbity to odpowiednio $p$ albo $1$.
Orbita jest jednoelementowa wtedy i tylko wtedy,
gdy wszystkie elementy grupy ustalają jej jedyny punkt.
Zbiór $X$ rozpada się zatem na $|F|$ orbit jednoelementowych
i pewną liczbę orbit $p$-elementowych.
Licząc elementy, dostajemy $|X|=|F|+kp$
przy pewnym $k\ge0$, co jest równoważne kongruencji.
\end{proof}

\begin{twierdzenie}[Cauchy'ego]
\label{aa:tw:cauchy}
Niech $G$ będzie grupą skończoną, a $p$ liczbą pierwszą
dzielącą $|G|$. Wtedy w $G$ istnieje element rzędu $p$.
W konsekwencji $G$ zawiera podgrupę rzędu $p$.
\end{twierdzenie}
\begin{proof}
Rozważmy zbiór
\[
 X=\{(g_1,\ldots,g_p)\in G^p:g_1g_2\cdots g_p=e\}.
\]
Pierwsze $p-1$ współrzędnych możemy wybrać zupełnie dowolnie.
Ostatnia współrzędna jest wówczas jednoznacznie wyznaczona:
$g_p=(g_1\cdots g_{p-1})^{-1}$.
Wobec tego $|X|=|G|^{p-1}$, a ponieważ $p\mid |G|$,
liczba $|X|$ jest podzielna przez $p$.

Niech generator cyklicznej grupy $C_p$ działa na $X$
przez cykliczne przesunięcie współrzędnych:
$T(g_1,\ldots,g_p)=(g_2,\ldots,g_p,g_1)$.
Musimy sprawdzić, że $T$ zachowuje $X$, zwłaszcza że $G$
nie musi być przemienna. Jeśli $g_1\cdots g_p=e$, to
\[
 g_2\cdots g_p g_1
 =g_1^{-1}(g_1\cdots g_p)g_1
 =g_1^{-1}eg_1=e.
\]
Zatem $T:X\to X$ jest dobrze określone.
Po $p$ przesunięciach wracamy do początkowej krotki,
więc $T^p=\operatorname{id}_X$ i potęgi $T$ wyznaczają
działanie $C_p$ na $X$ (nawet gdy to działanie nie jest wierne).

Punkty ustalone przez generator $T$ to dokładnie krotki
$(g,g,\ldots,g)$, dla których $g^p=e$.
Rzeczywiście równość $T(g_1,\ldots,g_p)=(g_1,\ldots,g_p)$
wymusza $g_1=\cdots=g_p=g$,
a warunek przynależności do $X$ daje $g^p=e$.
Ponieważ $C_p$ jest generowana przez $T$,
są to również punkty stałe całego działania.
Oznaczmy ich liczbę przez $r$.
Z lematu o orbitach $r\equiv|X|\equiv0\pmod p$.

Jedną z tych krotek jest $(e,\ldots,e)$, zatem $r\ge1$.
Jako dodatnia wielokrotność $p$ liczba $r$ musi być
co najmniej $p$, więc jest też krotka $(g,\ldots,g)$
z $g\ne e$ i $g^p=e$.
Rząd $g$ dzieli $p$ na mocy własności rzędów elementów,
a ponieważ $g\ne e$, nie jest równy $1$.
Zatem $\operatorname{ord}(g)=p$.
Podgrupa $\langle g\rangle$ składa się z dokładnie $p$ elementów.
\end{proof}

\begin{przyklad}[Jak twierdzenie Cauchy'ego działa w $S_3$?]
Grupa $S_3$ ma rząd $6=2\cdot3$.
Dla $p=2$ twierdzenie daje element rzędu $2$,
na przykład transpozycję $(1\,2)$.
Dla $p=3$ daje element rzędu $3$,
na przykład $(1\,2\,3)$.
Twierdzenie nie mówi, że istnieje element rzędu $6$:
w $S_3$ żaden element nie ma tego rzędu, mimo że
istnieją elementy rzędów $2$ i $3$.
Przyczyna jest strukturalna: iloczyn dwóch takich elementów
nie musi mieć rzędu będącego iloczynem ich rzędów,
gdyż elementy na ogół nie komutują.
\end{przyklad}

\begin{przyklad}[Liczba punktów stałych w konstrukcji Cauchy'ego]
\label{aa:prz:cauchy-z6}
Niech $G=(\mathbb Z/6\mathbb Z,+)$.
W dowodzie twierdzenia Cauchy'ego warunek iloczynu równego
$e$ staje się warunkiem sumy współrzędnych równej $[0]$.
Dla $p=2$ mamy $|X|=6$, a punkty stałe przesunięcia
$(x_1,x_2)\mapsto(x_2,x_1)$ są postaci $([a],[a])$
z $2[a]=[0]$. Są dokładnie dwa: $a=0$ oraz $a=3$.
Pozostałe $6-2=4$ krotki tworzą dwie orbity dwuelementowe,
więc istotnie $|X|=2+2\cdot2$.
Dla $p=3$ mamy $|X|=6^2=36$.
Punkty stałe cyklicznego przesunięcia są postaci
$([a],[a],[a])$ z $3[a]=[0]$, czyli $a\in\{0,2,4\}$.
Istnieją zatem trzy punkty stałe i
$(36-3)/3=11$ orbit trójelementowych.
W obu przypadkach dodatkowe nietrywialne rozwiązanie
równania $p[a]=[0]$ daje element rzędu $p$.
Dla ogólnej grupy nieprzemiennej rachunek przebiega
analogicznie, lecz należy użyć iloczynów zamiast sum.
\end{przyklad}

\begin{wniosek}[Centralny element rzędu $p$ w $p$-grupie]
\label{aa:wn:central-p}
Niech $G$ będzie skończoną grupą rzędu $p^a$, gdzie
$p$ jest liczbą pierwszą i $a\ge1$.
Wtedy w centrum $Z(G)$ istnieje element rzędu $p$.
W szczególności $G$ ma normalną podgrupę rzędu $p$.
\end{wniosek}
\begin{proof}
Z Twierdzenia~\ref{aa:tw:centrum-pgrupy} wiemy,
że $p\mid|Z(G)|$. Centrum jest skończoną grupą,
więc twierdzenie Cauchy'ego zastosowane do $Z(G)$
daje element $z$ rzędu $p$.
Podgrupa $\langle z\rangle$ ma rząd $p$.
Każda jej potęga należy do centrum, więc dla wszystkich
$g\in G$ mamy $g\langle z\rangle g^{-1}=\langle z\rangle$:
podgrupa jest normalna.
\end{proof}

\section{Lemat kombinatoryczny przygotowujący dowód Sylowa}
Pierwsze twierdzenie Sylowa wyprowadzimy z działania
na zbiorze podzbiorów grupy określonej liczności.
Kluczowe będzie wykazanie, że pewien współczynnik
dwumianowy nie dzieli się przez $p$. Poniższy lemat
uzasadnia to bez odwołań do nieudowodnionych wyników.

\begin{lemat}[Współczynnik dwumianowy modulo $p$]
\label{aa:lem:dwumian-p}
Niech $p$ będzie liczbą pierwszą, $a\ge0$ i $m\ge1$,
gdzie $p\nmid m$. Dla $n=p^a m$ zachodzi
\[
 \binom{n}{p^a}\equiv m\pmod p.
\]
W szczególności $p$ nie dzieli $\binom{n}{p^a}$.
\end{lemat}
\begin{proof}
Pracujemy w pierścieniu wielomianów
$(\mathbb Z/p\mathbb Z)[T]$; współczynniki traktujemy
jako klasy reszt modulo $p$.
Dla $1\le j\le p-1$ zachodzi $p\mid\binom pj$.
Aby to wykazać, pomnóżmy wzór
$j\binom pj=p\binom{p-1}{j-1}$.
Prawa strona jest podzielna przez $p$.
Ponieważ $p$ jest pierwsza i $p\nmid j$,
to $p\mid\binom pj$ (wynika to również z tożsamości
Bézouta, gdyż istnieją $u,v$ z $uj+vp=1$).
Po redukcji współczynników dwumian Newtona daje zatem
$(1+T)^p=1+T^p$ w naszym pierścieniu wielomianów.

Zastosujemy tę tożsamość wielokrotnie.
Dla $a=0$ mamy $(1+T)^{p^0}=1+T$.
Jeśli $(1+T)^{p^a}=1+T^{p^a}$, to podnosząc obie
strony do potęgi $p$ otrzymujemy
$(1+T)^{p^{a+1}}=(1+T^{p^a})^p=1+T^{p^{a+1}}$.
Indukcja dowodzi wzoru dla wszystkich $a\ge0$.

Wobec tego
$(1+T)^n=((1+T)^{p^a})^m=(1+T^{p^a})^m$
w pierścieniu $(\mathbb Z/p\mathbb Z)[T]$.
Porównajmy współczynniki przy $T^{p^a}$.
Po lewej jest to klasa $\binom{n}{p^a}$.
Po prawej wzór dwumianowy
$\sum_{j=0}^m\binom mj T^{jp^a}$ pokazuje,
że jedynym składnikiem o żądanym stopniu jest ten dla $j=1$,
zatem współczynnik wynosi $\binom m1=m$.
Otrzymujemy kongruencję z tezy. Ponieważ $p\nmid m$,
współczynnik dwumianowy nie dzieli się przez $p$.
\end{proof}

\section{Pierwsze twierdzenie Sylowa: istnienie}
\begin{twierdzenie}[Pierwsze twierdzenie Sylowa]
\label{aa:tw:sylow1}
Niech $G$ będzie grupą skończoną, $p$ liczbą pierwszą,
a $|G|=p^a m$, gdzie $a$ jest nieujemną liczbą całkowitą,
$m$ jest dodatnią liczbą całkowitą i $p\nmid m$.
Wtedy istnieje podgrupa $P\le G$ o rzędzie $p^a$,
czyli podgrupa $p$-Sylowa.
\end{twierdzenie}
\begin{proof}
Dla $a=0$ wystarczy wziąć $P=\{e\}$.
Załóżmy więc $a\ge1$ i oznaczmy $r=p^a$.
Niech $\mathcal X$ będzie zbiorem wszystkich $r$-elementowych
podzbiorów grupy $G$. Ponieważ $r\le|G|$, jest to zbiór niepusty,
a jego liczebność wynosi
$|\mathcal X|=\binom{|G|}{r}=\binom{p^a m}{p^a}$.
Na mocy poprzedniego lematu $p\nmid|\mathcal X|$.

Grupa $G$ działa na $\mathcal X$ przez przesunięcia z lewej:
$g\cdot A=\{ga:a\in A\}$.
To działanie jest dobrze określone, bo lewa translacja
$G\to G$, $a\mapsto ga$, jest bijekcją, więc zachowuje
liczbę elementów podzbioru.
Neutralny element ustala każdy podzbiór,
a $g\cdot(h\cdot A)=(gh)\cdot A$.
Zatem faktycznie mamy działanie grupy.

Orbity tworzą rozkład $\mathcal X$.
Gdyby każda orbita miała liczebność podzielną przez $p$,
cała suma ich liczebności, czyli $|\mathcal X|$,
też byłaby podzielna przez $p$. To przeczy lematowi.
Istnieje więc orbita $\mathcal O=G\cdot A$
o liczebności niepodzielnej przez $p$.

Niech $H=G_A=\{g\in G:gA=A\}$ będzie stabilizatorem
wybranego podzbioru $A$.
Z twierdzenia o orbicie i stabilizatorze
$|\mathcal O|=[G:H]=|G|/|H|$.
Liczba $|\mathcal O|$ nie dzieli się przez $p$, a
$|G|=p^a m$. Stąd $p^a$ musi dzielić $|H|$:
wszystkie $a$ czynniki $p$ występujące w $|G|$
pozostają w $|H|$, skoro w ilorazie $|G|/|H|$
nie występuje żaden czynnik $p$.

Z drugiej strony $H$ działa na samym zbiorze $A$ przez
mnożenie z lewej: dla $h\in H$ mamy $hA=A$.
Działanie to jest \emph{swobodne}: jeżeli $h\cdot x=x$
dla pewnego $x\in A$, to $hx=x$ w grupie $G$,
a po pomnożeniu przez $x^{-1}$ otrzymujemy $h=e$.
Zatem stabilizator każdego punktu $x\in A$ jest trywialny.
Na mocy twierdzenia o orbicie i stabilizatorze
każda orbita tego działania ma dokładnie $|H|$ elementów.
Ponieważ orbity rozkładają $A$ na rozłączne części,
$|H|$ dzieli $|A|=p^a$.

Połączywszy oba wnioski, mamy $p^a\mid|H|$ oraz
$|H|\mid p^a$; dlatego $|H|=p^a$.
Stabilizator $H$ jest więc szukaną podgrupą $p$-Sylowa.
\end{proof}

\begin{przyklad}[Istnienie podgrup Sylowa w $S_3$]
Rząd $S_3$ wynosi $6=2\cdot3$.
Pierwsze twierdzenie Sylowa zapewnia istnienie podgrupy rzędu $2$
i podgrupy rzędu $3$.
Przykładami są
$\langle(1\,2)\rangle=\{e,(1\,2)\}$
oraz $\langle(1\,2\,3)\rangle=\{e,(1\,2\,3),(1\,3\,2)\}$.
W tym przypadku obie podgrupy są cykliczne,
ale ogólnie podgrupy Sylowa nie muszą być cykliczne:
w $A_4$ podgrupa rzędu $4$ może być izomorficzna
z $(\mathbb Z/2\mathbb Z)^2$.
\end{przyklad}

\begin{przyklad}[Wydobycie podgrupy Sylowa ze stabilizatora]
\label{aa:prz:sylow-stabilizator-s3}
Rozpatrzmy dowód istnienia podgrup Sylowa w $G=S_3$.
Dla $p=2$ bierzemy zbiór $\mathcal X$ dwuelementowych
podzbiorów $G$. Ma on $\binom62=15$ elementów,
czyli liczebność nieparzystą.
Wybierzmy $A=\{e,(1\,2)\}$.
Stabilizatorem $A$ w działaniu przez lewe translacje jest
$G_A=A$: jeśli $gA=A$, to $g=ge\in A$; z drugiej strony
każdy element dwuelementowej podgrupy $A$ permutuje $A$
przez mnożenie. Orbita $A$ ma zatem $6/2=3$ elementy.
Są nimi $A$, $(1\,3)A$ i $(2\,3)A$, które są parami
różne, gdyż reprezentują różne lewe warstwy $G/A$.
Znaleźliśmy orbitę rozmiaru niepodzielnego przez $2$,
a jej stabilizator jest poszukiwaną podgrupą $2$-Sylowa.

Dla $p=3$ wybieramy podzbiór
$B=\{e,(1\,2\,3),(1\,3\,2)\}=A_3$.
Zbiór trójelementowych podzbiorów ma $\binom63=20$
elementów, a orbita $B$ składa się z dwóch warstw grupy
$S_3$ względem $A_3$. Jej rozmiar wynosi $2$,
nie dzieli się przez $3$ i stabilizator $G_B=B$
ma rząd $3$. Ważne: nie każda orbita musi mieć liczebność
niepodzielną przez $p$; dowód istnienia zapewnia
co najmniej jedną taką orbitę, a konkretny wybór
$A$ lub $B$ pokazuje mechanizm konstrukcji.
\end{przyklad}

\section{Drugie twierdzenie Sylowa: zawieranie i sprzężoność}
\begin{lemat}[Orbity $p$-grupy na skończonym zbiorze]
\label{aa:lem:orbity-pgrupy}
Niech $Q$ będzie skończoną grupą rzędu $p^b$, gdzie $p$
jest liczbą pierwszą, a $b\ge0$, i niech $Q$ działa
na skończonym zbiorze $X$.
Jeżeli $p\nmid|X|$, to działanie ma punkt stały.
\end{lemat}
\begin{proof}
Orbita punktu $x$ ma rozmiar $[Q:Q_x]$.
Na mocy twierdzenia Lagrange'a jest to potęga $p$,
czyli $1,p,p^2,\ldots$.
Jeśli nie byłoby orbit jednoelementowych, każda orbita
miałaby liczebność podzielną przez $p$,
a więc $|X|$, jako suma ich liczebności,
też dzieliłoby się przez $p$. Otrzymalibyśmy sprzeczność.
Istnieje zatem orbita jednoelementowa,
a jej punkt jest ustalany przez całe $Q$.
\end{proof}

\begin{twierdzenie}[Drugie twierdzenie Sylowa]
\label{aa:tw:sylow2}
Niech $G$ będzie grupą skończoną, $p$ liczbą pierwszą,
a $P\in\operatorname{Syl}_p(G)$.
Dla każdej podgrupy $Q\le G$ o rzędzie będącym potęgą $p$
istnieje $g\in G$ takie, że $g^{-1}Qg\subseteq P$.
W szczególności każda taka podgrupa zawiera się w pewnej
podgrupie $p$-Sylowa, a wszystkie podgrupy $p$-Sylowa
są wzajemnie sprzężone.
\end{twierdzenie}
\begin{proof}
Zapiszmy $|G|=p^a m$ z $p\nmid m$.
Ponieważ $|P|=p^a$, zbiór warstw lewostronnych
$G/P$ ma dokładnie $m$ elementów.
Grupa $Q$ działa na $G/P$ przez mnożenie z lewej:
$q\cdot(gP)=(qg)P$. Poprawność tego działania
wykazaliśmy w Przykładzie~\ref{aa:prz:akcja-kosety}.
Ponieważ $Q$ jest $p$-grupą, a liczba $m$ nie jest
podzielna przez $p$, poprzedni lemat zapewnia warstwę
$gP$ ustalaną przez wszystkie elementy $q\in Q$.
Dla każdego $q\in Q$ mamy więc
$qgP=gP$, czyli $g^{-1}qg\in P$.
Oznacza to dokładnie zawieranie $g^{-1}Qg\subseteq P$.
Po sprzężeniu obustronnie przez $g$ otrzymujemy
$Q\subseteq gPg^{-1}$, a grupa $gPg^{-1}$ jest
podgrupą $p$-Sylowa, gdyż sprzężenie jest izomorfizmem.

Jeśli $Q$ także jest podgrupą $p$-Sylowa, to
$|Q|=|P|=p^a$.
Sprzężona grupa $g^{-1}Qg$ ma tyle samo elementów,
co $Q$. Zawiera się w $P$ i ma taką samą liczebność,
zatem $g^{-1}Qg=P$, czyli $Q=gPg^{-1}$.
Dowodzi to sprzężoności wszystkich podgrup Sylowa.
\end{proof}

\begin{uwaga}[Dlaczego używamy lewych warstw?]
Dla dowolnej podgrupy $P\le G$ przestrzeń $G/P$
jest zbiorem warstw lewostronnych, ale nie musi mieć
struktury grupy. W dowodzie wystarcza działanie $Q$
na \emph{zbiorze} warstw, nie zakładamy normalności $P$.
To celowo słabsza konstrukcja niż grupa ilorazowa.
\end{uwaga}

\section{Trzecie twierdzenie Sylowa: liczba podgrup}
\begin{twierdzenie}[Trzecie twierdzenie Sylowa]
\label{aa:tw:sylow3}
Niech $G$ będzie grupą skończoną i niech
$|G|=p^a m$, gdzie $p$ jest liczbą pierwszą,
$a$ jest nieujemną liczbą całkowitą, a $m$ jest dodatnią
liczbą całkowitą niepodzielną przez $p$.
Wtedy liczba $n_p(G)$ podgrup $p$-Sylowa spełnia
\[
 n_p(G)\mid m,\qquad n_p(G)\equiv1\pmod p.
\]
Ponadto dla dowolnej podgrupy $P\in\operatorname{Syl}_p(G)$
mamy $n_p(G)=[G:N_G(P)]$.
\end{twierdzenie}
\begin{proof}
Pierwsze twierdzenie Sylowa zapewnia, że
$\operatorname{Syl}_p(G)$ jest niepusty.
Grupa $G$ działa na zbiorze podgrup Sylowa przez sprzężenie:
$g\cdot Q=gQg^{-1}$. Sprzężenie nie zmienia rzędu,
zatem działanie jest dobrze określone.
Drugie twierdzenie Sylowa mówi, że jest ono przechodnie.
Stabilizatorem podgrupy $P$ jest z definicji jej normalizator
$N_G(P)$. Z twierdzenia o orbicie i stabilizatorze dostajemy
$n_p(G)=[G:N_G(P)]$.
Ponieważ $P\le N_G(P)\le G$, możemy zastosować
wzór Lagrange'a dwukrotnie:
\[
 [G:P]=[G:N_G(P)]\,[N_G(P):P].
\]
Lewa strona wynosi $m$, zatem $n_p(G)$ dzieli $m$.

Pozostaje wykazać kongruencję modulo $p$.
Niech $P$ działa na $\operatorname{Syl}_p(G)$
przez sprzężenie. Każda orbita tego działania ma
rozmiar będący potęgą $p$, ponieważ jej rozmiar dzieli
$|P|=p^a$. Wobec tego liczba wszystkich podgrup Sylowa
jest przystająca modulo $p$ do liczby orbit jednoelementowych,
czyli do liczby podgrup Sylowa ustalanych przez każde $u\in P$.

Pokażemy, że jedyną taką podgrupą jest $P$.
Oczywiście $P$ jest ustalana przez wszystkie elementy $P$,
gdyż $uPu^{-1}=P$ dla $u\in P$.
Niech teraz $Q\in\operatorname{Syl}_p(G)$ będzie ustalana
przez całe $P$. Oznacza to $uQu^{-1}=Q$ dla każdego $u\in P$,
czyli $P\subseteq N_G(Q)$.
Połóżmy $N=N_G(Q)$. Grupa $Q$ jest normalna w $N$,
a jej rząd to $p^a$.
Ponieważ $Q\le N\le G$, na mocy twierdzenia Lagrange'a
mamy $p^a=|Q|\mid|N|$ oraz $|N|\mid|G|=p^a m$.
Ponieważ $p\nmid m$, największą potęgą $p$ dzielącą
$|N|$ jest dokładnie $p^a$.
A zatem $P$ i $Q$ są obie podgrupami $p$-Sylowa \emph{grupy $N$}.
Stosując drugie twierdzenie Sylowa do $N$, otrzymujemy
$n\in N$ z $nPn^{-1}=Q$.
Normalność $Q$ w $N$ daje $n^{-1}Qn=Q$.
Sprzęgając poprzednią równość przez $n^{-1}$,
uzyskujemy $P=n^{-1}Qn=Q$.
Tak więc jedynym punktem stałym działania $P$
na $\operatorname{Syl}_p(G)$ jest $P$.
Pozostałe orbity mają rozmiary podzielne przez $p$,
więc $n_p(G)\equiv1\pmod p$.
\end{proof}

\begin{wniosek}[Kiedy podgrupa Sylowa jest normalna?]
\label{aa:wn:normalna-sylow}
Niech $G$ będzie grupą skończoną, $p$ liczbą pierwszą
dzielącą $|G|$, a $P\in\operatorname{Syl}_p(G)$.
Wtedy $P\trianglelefteq G$ dokładnie wtedy, gdy $n_p(G)=1$.
\end{wniosek}
\begin{proof}
Jeśli $P$ jest normalna, to każde sprzężenie $gPg^{-1}$
jest równe $P$. Ponieważ wszystkie podgrupy Sylowa są sprzężone,
nie istnieje inna podgrupa Sylowa, więc $n_p(G)=1$.
Odwrotnie, jeżeli istnieje tylko jedna podgrupa Sylowa,
to każde $gPg^{-1}$, będące podgrupą Sylowa, musi być jej równe.
To właśnie definicja normalności $P$.
\end{proof}

\section{Przykłady: jak rzeczywiście korzystać z trzech twierdzeń}
\begin{przyklad}[Podgrupy Sylowa w $S_3$]
\label{aa:prz:sylow-s3}
Mamy $|S_3|=6$.
Dla $p=3$ liczba $n_3$ dzieli $2$ i spełnia
$n_3\equiv1\pmod3$. Jedyna możliwość to $n_3=1$.
Rzeczywiście podgrupa $A_3=\langle(1\,2\,3)\rangle$
jest normalna i ma rząd $3$.
Dla $p=2$ mamy $n_2\mid3$ oraz $n_2\equiv1\pmod2$.
Możliwości są dwie: $1$ albo $3$.
Wypiszmy podgrupy rzędu $2$:
$\langle(1\,2)\rangle$, $\langle(1\,3)\rangle$
i $\langle(2\,3)\rangle$.
Są trzy i są wzajemnie sprzężone, więc $n_2=3$.
To pokazuje, że same warunki liczbowe Sylowa mogą pozostawiać
więcej niż jedną możliwość; trzeba wykorzystać strukturę grupy.
\end{przyklad}

\begin{przyklad}[Podgrupy Sylowa w $S_4$ i podziały na pary]
Niech $G=S_4$, $|G|=24=2^3\cdot3$.
Podgrupa $3$-Sylowa ma rząd $3$, a każda zawiera dwa
nieidentycznościowe $3$-cykle. Ponieważ w $S_4$ takich
cykli jest $\binom43\cdot2=8$, a dwie różne podgrupy
rzędu $3$ nie mają wspólnego nieidentycznościowego elementu,
otrzymujemy $n_3(S_4)=8/2=4$.

Dla $p=2$ trzecie twierdzenie Sylowa daje
$n_2\mid3$ i $n_2\equiv1\pmod2$, więc $n_2=1$ albo $3$.
Rozważmy trzy podziały zbioru $\{1,2,3,4\}$ na dwie
nieuporządkowane pary:
$\{\{1,2\},\{3,4\}\}$,
$\{\{1,3\},\{2,4\}\}$ oraz
$\{\{1,4\},\{2,3\}\}$.
Stabilizator podziału $\{\{1,2\},\{3,4\}\}$ to zbiór
permutacji, które albo zachowują obie pary z osobna,
albo je zamieniają. W pierwszym przypadku mamy $2\cdot2=4$
permutacje: w każdej parze niezależnie wybieramy, czy jej
elementy zamienić. W drugim przypadku obraz elementu $1$
można wybrać spośród $3,4$, a obraz $2$ jest wtedy
drugim z tych dwóch elementów; podobnie obrazy $3,4$
można ustalić na dwa sposoby w parze $\{1,2\}$.
Daje to $2\cdot2=4$ kolejne permutacje.
Obie rodziny są rozłączne, toteż stabilizator ma dokładnie
osiem elementów i jest podgrupą $2$-Sylowa.
Stabilizatory pozostałych dwóch podziałów również mają
osiem elementów, ponieważ podziały są wzajemnie
przenoszone przez permutacje $S_4$, a ich stabilizatory
są sprzężone.
Stabilizatory są parami różne. Istotnie transpozycja
$(1\,2)$ należy do stabilizatora pierwszego podziału,
lecz przestawia dwa pozostałe podziały i nie należy do ich
stabilizatorów. Analogicznie $(1\,3)$ należy do stabilizatora
drugiego podziału, lecz nie należy do stabilizatorów pierwszego
i trzeciego, a $(1\,4)$ należy wyłącznie do stabilizatora
trzeciego podziału. Mamy więc trzy różne podgrupy
$2$-Sylowa, czyli $n_2\ge3$.
Z ograniczenia Sylowa $n_2\mid3$ wynika $n_2=3$.
Ten przykład pokazuje, jak przejść od samych warunków
liczbowych do faktycznego zliczania podgrup.
\end{przyklad}

\begin{przyklad}[Podgrupy Sylowa w $A_4$]
\label{aa:prz:sylow-a4}
Rząd $A_4$ wynosi $12=2^2\cdot3$.
Jej elementy to identyczność, osiem $3$-cykli
i trzy iloczyny dwóch rozłącznych transpozycji.
Każda podgrupa rzędu $3$ zawiera dokładnie dwa
nieidentycznościowe elementy, będące wzajemnie odwrotnymi
$3$-cyklami. Dwie różne takie podgrupy nie mogą współdzielić
nieidentycznościowego elementu: element rzędu $3$
generuje jednoznacznie swoją trójelementową podgrupę.
Osiem $3$-cykli dzieli się więc na cztery pary,
czyli $n_3(A_4)=8/2=4$.
Jest to zgodne z warunkami $n_3\mid4$ i $n_3\equiv1\pmod3$.

Rozważmy zbiór
$V=\{e,a,b,c\}$, gdzie
$a=(1\,2)(3\,4)$, $b=(1\,3)(2\,4)$ oraz
$c=(1\,4)(2\,3)$.
Mamy $a^2=b^2=c^2=e$, ponieważ każdy z tych elementów
jest iloczynem rozłącznych transpozycji.
Przy konwencji składania od prawej do lewej bezpośrednie
obliczenie obrazów $1,2,3,4$ daje
$ab=ba=c$, $bc=cb=a$ oraz $ca=ac=b$.
Na przykład $ab$ wysyła $1\mapsto4$, $4\mapsto1$,
$2\mapsto3$, $3\mapsto2$, a więc $ab=c$.
Wszystkie iloczyny elementów $V$ znów należą do $V$.
Ponieważ $V$ zawiera identyczność i odwrotność każdego
ze swych elementów, jest podgrupą $A_4$.
Funkcja $(\mathbb Z/2\mathbb Z)^2\to V$ określona wzorem
$([i],[j])\mapsto a^i b^j$ jest homomorfizmem,
ponieważ $a$ i $b$ komutują. Jej obrazy czterech par
to kolejno $e,a,b,c$, więc jest bijekcją.
Dlatego $V\cong(\mathbb Z/2\mathbb Z)^2$.
Jest normalna w $A_4$, ponieważ sprzężenie permutacją
zachowuje długości cykli i rozłączność, zatem przestawia
trzy nieidentycznościowe elementy $V$.
Podgrupa $V$ ma rząd $4=2^2$, więc jest podgrupą $2$-Sylowa.
Z normalności i wniosku o normalności Sylowa wynika $n_2(A_4)=1$.
Nie można jednak wyciągnąć stąd wniosku, że $A_4$ jest
przemienna: choć $V$ jest przemienna i normalna,
w $A_4$ znajdują się także $3$-cykle, które nie komutują
ze wszystkimi elementami $V$. Dla przykładu
$(1\,2\,3)(1\,2)(3\,4)(1\,2\,3)^{-1}
=(1\,4)(2\,3)\ne (1\,2)(3\,4)$.
Obliczenia te wyjaśniają, dlaczego podgrupy Sylowa
nie muszą być cykliczne, a ich liczby nie muszą być jednakowe.
\end{przyklad}

\begin{twierdzenie}[Grupy rzędu $pq$ przy braku szczególnej kongruencji]
\label{aa:tw:pq-cykliczna}
Niech $p<q$ będą liczbami pierwszymi i niech $G$ będzie grupą
rzędu $pq$. Jeśli $p\nmid(q-1)$, to $G$ jest cykliczna.
\end{twierdzenie}
\begin{proof}
Dla podgrup $q$-Sylowa mamy $n_q\mid p$ oraz
$n_q\equiv1\pmod q$. Ponieważ $n_q\in\{1,p\}$
i $1<p<q$, liczba $p$ nie jest przystająca do $1$ modulo $q$.
Zatem $n_q=1$, czyli jedyna podgrupa $Q$ rzędu $q$ jest normalna.
Dla podgrup $p$-Sylowa mamy $n_p\mid q$,
więc $n_p\in\{1,q\}$, i $n_p\equiv1\pmod p$.
Założenie $p\nmid(q-1)$ wyklucza $n_p=q$.
Zatem $n_p=1$, czyli podgrupa $P$ rzędu $p$ także jest normalna.

Podgrupy $P$ i $Q$ są cykliczne, ponieważ mają rzędy pierwsze.
Ich przekrój jest trywialny: rząd $P\cap Q$ dzieli jednocześnie
$p$ i $q$, a te są różnymi liczbami pierwszymi.
Niech $x\in P$, $y\in Q$. Rozważmy komutator
$[x,y]=xyx^{-1}y^{-1}$.
Normalność $P$ daje $yx^{-1}y^{-1}\in P$,
więc $[x,y]=x(yx^{-1}y^{-1})\in P$.
Normalność $Q$ daje $xyx^{-1}\in Q$,
więc $[x,y]=(xyx^{-1})y^{-1}\in Q$.
Stąd $[x,y]\in P\cap Q=\{e\}$, a więc $xy=yx$.
Każdy element $P$ komutuje z każdym elementem $Q$.

Wybierzmy generatory $x$ grupy $P$ i $y$ grupy $Q$.
Ponieważ komutują i mają rzędy $p,q$, element $xy$ ma
rząd $pq$. Istotnie, jeśli $(xy)^k=e$, to
$x^k=y^{-k}\in P\cap Q$, więc $x^k=y^k=e$.
Zatem $p\mid k$ i $q\mid k$, a ponieważ $p,q$ są względnie
pierwsze, $pq\mid k$ (z tożsamości Bézouta).
Z drugiej strony $(xy)^{pq}=x^{pq}y^{pq}=e$.
Najmniejszym dodatnim takim $k$ jest więc $pq$.
Podgrupa $\langle xy\rangle$ ma $pq=|G|$ elementów,
toteż jest równa $G$. Grupa jest cykliczna.
\end{proof}

\begin{przyklad}[Dlaczego warunek $p\nmid(q-1)$ jest istotny?]
Niech $q=7$ i $p=3$; wówczas $3\mid(7-1)$.
Zbudujemy konkretną nieabelową grupę rzędu $21$.
Dla $a\in\mathbb Z/7\mathbb Z$ i $i\in\{0,1,2\}$
rozważmy permutację zbioru $\mathbb Z/7\mathbb Z$
zadaną przez
$f_{a,i}(x)=2^i x+a\pmod7$.
Każda jest odwracalna, bo $2^i$ jest niezerowe modulo $7$.
Funkcje $f_{a,i}$ są parami różne: $a=f_{a,i}(0)$,
a $2^i=f_{a,i}(1)-f_{a,i}(0)$, przy czym $1,2,4$
są różnymi klasami modulo $7$.
Jest ich zatem $7\cdot3=21$.
Składanie daje
$f_{a,i}\circ f_{b,j}=f_{a+2^ib,\,i+j\ (\mathrm{mod}\ 3)}$,
ponieważ $2^3\equiv1\pmod7$.
Identycznością jest $f_{0,0}$, a odwrotność też należy
do tej rodziny: można wziąć
$f_{-2^{-i}a,\,-i\ (\mathrm{mod}\ 3)}$,
przy czym $2^{-i}$ oznacza odwrotność $2^i$ modulo $7$.
Zatem te $21$ funkcji tworzy grupę względem składania.
Nie jest abelowa:
$f_{0,1}\circ f_{1,0}=f_{2,1}$,
natomiast $f_{1,0}\circ f_{0,1}=f_{1,1}$.
Z tego powodu nie jest cykliczna.

Podgrupa translacji $T=\{f_{a,0}:a\in\mathbb Z/7\mathbb Z\}$
ma rząd $7$ i jest normalna: dla $f_{b,j}$ sprzężenie
$f_{b,j}\circ f_{a,0}\circ f_{b,j}^{-1}$ jest translacją
o przesunięcie $2^j a$.
Podgrupa $\langle f_{0,1}\rangle$ ma rząd $3$,
lecz nie jest normalna: sprzężenie przez $f_{1,0}$
wysyła $f_{0,1}$ na $f_{-1,1}$, co nie należy
do $\langle f_{0,1}\rangle=\{f_{0,0},f_{0,1},f_{0,2}\}$.
Warunki Sylowa dają $n_7=1$ oraz $n_3\in\{1,7\}$;
z nienormalności podgrupy rzędu $3$ wynika $n_3=7$.
\end{przyklad}

\begin{twierdzenie}[Nie istnieje grupa prosta rzędu $30$]
\label{aa:tw:30-not-simple}
Niech $G$ będzie grupą rzędu $30$. Wtedy $G$ ma
właściwą nietrywialną podgrupę normalną.
Inaczej mówiąc, żadna grupa rzędu $30$ nie jest prosta.
\end{twierdzenie}
\begin{proof}
Przypomnijmy, że grupa \emph{prosta} to grupa $G\ne\{e\}$,
której jedynymi normalnymi podgrupami są $\{e\}$ i $G$.
Załóżmy przeciwnie, że $G$ jest prosta.
Ponieważ $|G|=2\cdot3\cdot5$,
trzecie twierdzenie Sylowa daje
$n_5\mid6$, $n_5\equiv1\pmod5$;
stąd $n_5\in\{1,6\}$.
Gdyby $n_5=1$, podgrupa rzędu $5$ byłaby normalna,
co przeczy prostocie. Zatem $n_5=6$.
Analogicznie $n_3\mid10$, $n_3\equiv1\pmod3$,
więc $n_3\in\{1,10\}$ i prostota wymusza $n_3=10$.

Każda podgrupa rzędu $5$ jest cykliczna i ma cztery
elementy różne od jedynki, wszystkie rzędu $5$.
Dwie różne podgrupy rzędu $5$ nie mają wspólnego
elementu nieidentycznościowego: taki element sam generowałby
obie, co byłoby sprzecznością.
Zatem w $G$ jest co najmniej $6\cdot4=24$
elementów rzędu $5$.
Podobnie dziesięć podgrup rzędu $3$ dostarcza
$10\cdot2=20$ różnych elementów rzędu $3$.
Żaden element nie może mieć równocześnie rzędu $3$ i $5$,
więc obie rodziny są rozłączne.
Dostalibyśmy co najmniej $24+20=44$ nieidentycznościowe
elementy w grupie rzędu $30$, co jest niemożliwe.
Zatem $G$ nie jest prosta.
\end{proof}

\section{Dodatkowe klasyczne zastosowanie: podgrupy małego indeksu}
\begin{twierdzenie}[Normalność podgrupy o indeksie najmniejszego dzielnika pierwszego]
\label{aa:tw:index-minprime}
Niech $G$ będzie grupą skończoną, a $p$ najmniejszą liczbą
pierwszą dzielącą $|G|$. Jeżeli $H\le G$ i $[G:H]=p$,
to $H$ jest normalna w $G$.
\end{twierdzenie}
\begin{proof}
Niech $G$ działa na $p$ warstwach lewostronnych $G/H$.
Działanie wyznacza homomorfizm
$\rho:G\to S_p$ i jego obraz $K=\rho(G)$.
Ponieważ działanie na warstwach jest przechodnie,
$K$ działa przechodnio na zbiorze $p$ elementów.
Z twierdzenia o orbicie i stabilizatorze, zastosowanego do $K$,
wynika, że $p\mid|K|$.
Z pierwszego twierdzenia o izomorfizmie
$|K|=|G|/|\ker\rho|$, więc $|K|$ dzieli $|G|$.
Ponadto $K\le S_p$, stąd $|K|\mid p!$.

Jeśli liczba pierwsza $q$ dzieli $|K|$, to dzieli także
$|G|$, a z minimalności $p$ mamy $q\ge p$.
Z drugiej strony $q\mid p!$ implikuje $q\le p$.
Dlatego jedynym dzielnikiem pierwszym $|K|$ może być $p$.
W liczbie $p!$ czynnik $p$ występuje tylko raz:
jedyną jego wielokrotnością wśród $1,2,\ldots,p$
jest liczba $p$.
Ponieważ $p\mid |K|$ i $|K|\mid p!$,
otrzymujemy $|K|=p$.

Zatem $|\ker\rho|=|G|/p=|H|$.
Jądro działania na warstwach zawiera się w stabilizatorze
warstwy $H$, którym jest $H$.
Mamy więc $\ker\rho\subseteq H$ i równość liczebności,
zatem $\ker\rho=H$.
Jako jądro homomorfizmu $H$ jest normalna w $G$.
\end{proof}

\begin{uwaga}[Jak korzystać z twierdzeń bez nieuprawnionych odwrotności?]
Wnioski Sylowa mówią o istnieniu podgrup określonego rzędu,
ich sprzężoności oraz ograniczeniach liczbowych ich liczby.
Nie twierdzą, że istnieje podgrupa dowolnego rzędu dzielącego $|G|$,
ani nie klasyfikują wszystkich grup danego rzędu.
Przy stosowaniu warunków $n_p\mid m$ i $n_p\equiv1\pmod p$
należy najpierw wyznaczyć wszystkie możliwości arytmetyczne,
a dopiero potem rozstrzygać je za pomocą dalszych informacji
o konkretnych elementach lub działaniach grupy.
\end{uwaga}

% AA-05 | Algebra abstrakcyjna | Iteracja 03
\chapter[Szeregi podnormalne i nilpotentność]{Szeregi podnormalne, grupy rozwiązalne i nilpotentne}
\label{aa:ch:szeregi}

\section{Dlaczego rozkładamy grupę na czynniki?}
Klasyfikowanie grup wyłącznie na podstawie ich rzędu rzadko się udaje.
Możemy jednak badać grupę przez podgrupę normalną $N\trianglelefteq G$
i grupę ilorazową $G/N$. Proces ten przypomina rozkład liczby
na czynniki pierwsze, lecz czynniki grupowe mogą być nieprzemienne,
a ich kolejność i sposób ponownego składania mają znaczenie.
W tym rozdziale udowodnimy, w jakim sensie lista najprostszych czynników
jest niezależna od wybranego rozkładu; następnie rozróżnimy
\emph{rozwiązalność}, która kontroluje kolejne komutatory,
i silniejszą własność \emph{nilpotentności}.

\begin{definicja}[Grupa prosta]
\label{aa:def:prosta}
Grupę $G\ne\{e\}$ nazywamy \emph{prostą}, jeśli jedynymi
podgrupami normalnymi $G$ są $\{e\}$ i $G$.
Podgrupa $N\trianglelefteq G$ jest \emph{maksymalną właściwą podgrupą
normalną}, jeżeli $N\ne G$ oraz nie istnieje
$K\trianglelefteq G$ z $N\subsetneq K\subsetneq G$.
\end{definicja}
\begin{lemat}[Proste ilorazy]
\label{aa:lem:maxnormal}
Niech $N\trianglelefteq G$ będzie właściwą podgrupą normalną.
Wtedy $N$ jest maksymalna wśród właściwych podgrup normalnych $G$
dokładnie wtedy, gdy $G/N$ jest grupą prostą.
\end{lemat}
\begin{proof}
Twierdzenie o odpowiedniości podgrup
(\ref{aa:tw:odpowiedniosc}) daje bijekcję między podgrupami
$K$ takimi, że $N\subseteq K\le G$, a podgrupami $G/N$.
Ponadto $K$ jest normalna w $G$ dokładnie wtedy, gdy
$K/N$ jest normalna w $G/N$: równość
$(gN)(K/N)(gN)^{-1}=K/N$ oznacza $gKg^{-1}N=K$.
Ponieważ $gNg^{-1}=N\subseteq gKg^{-1}$, lewa strona
tej równości jest równa $gKg^{-1}$; otrzymujemy więc
warunek $gKg^{-1}=K$.
Zatem pośrednie właściwe podgrupy normalne są dokładnie
niebanalnymi właściwymi podgrupami normalnymi ilorazu.
Ich brak jest równoważny prostocie $G/N$.
\end{proof}

\begin{definicja}[Szereg podnormalny i kompozycyjny]
\label{aa:def:kompozycyjny}
Dla grupy $G$ \emph{szeregiem podnormalnym} nazywamy skończony ciąg
$G=G_0\trianglerighteq G_1\trianglerighteq\cdots\trianglerighteq G_r=\{e\}$,
w którym $G_{i+1}\trianglelefteq G_i$ dla każdego $i$.
Jego \emph{czynniki} to $G_i/G_{i+1}$.
Jeżeli wszystkie zawierania są właściwe, a każdy czynnik jest prosty,
mówimy o \emph{szeregu kompozycyjnym}.
Nie wymagamy, by $G_i$ były normalne w całej grupie $G$;
normalność dotyczy kolejnych sąsiadujących wyrazów.
\end{definicja}

\begin{twierdzenie}[Istnienie szeregu kompozycyjnego dla grup skończonych]
\label{aa:tw:existcomp}
Każda niebanalna grupa skończona $G$ posiada szereg kompozycyjny.
\end{twierdzenie}
\begin{proof}
Przeprowadzimy indukcję względem $|G|$.
Jeśli $G$ jest prosta, wystarcza $G\triangleright\{e\}$.
Jeśli nie jest prosta, zbiór właściwych podgrup normalnych
$G$ zawierających $\{e\}$ jest niepusty i skończony,
więc możemy wybrać podgrupę $N$ maksymalną względem zawierania.
Z poprzedniego lematu $G/N$ jest prosta.
Ponieważ $|N|<|G|$, założenie indukcyjne daje szereg
kompozycyjny $N=N_0\triangleright\cdots\triangleright N_s=\{e\}$
(gdy $N=\{e\}$, stosujemy szereg długości zero).
Dopisanie $G\triangleright N$ przed tym szeregiem daje
szereg kompozycyjny $G$, gdyż pierwszy czynnik $G/N$ jest prosty,
a wszystkie pozostałe czynniki pochodzą z szeregu $N$.
\end{proof}

\section{Twierdzenie Jordana--Höldera}
\begin{lemat}[Porównanie dwóch maksymalnych podgrup normalnych]
\label{aa:lem:two-max}
Niech $G$ będzie skończoną grupą, a $N,M$ dwiema różnymi
maksymalnymi właściwymi podgrupami normalnymi $G$.
Jeśli $L=N\cap M$, to $NM=G$,
$N/L\cong G/M$ oraz $M/L\cong G/N$.
W szczególności $N/L$ i $M/L$ są grupami prostymi.
\end{lemat}
\begin{proof}
Ponieważ $N$ i $M$ są normalne w $G$, zbiór
$NM=\{nm:n\in N,\ m\in M\}$ jest podgrupą normalną:
$(nm)(n'm')=n(mn'm^{-1})(mm')\in NM$.
Ponadto $(nm)^{-1}=m^{-1}n^{-1}
=(m^{-1}n^{-1}m)m^{-1}\in NM$, a $e\in NM$.
Sprzężenie przez dowolne $g\in G$ zachowuje oba czynniki,
zatem $NM$ jest podgrupą normalną $G$.
Gdyby $M\subseteq N$, maksymalność właściwej podgrupy normalnej
$M$ i normalność $N$ wymuszałyby $N=M$ albo $N=G$,
co przeczy założeniom. Zatem $NM$ zawiera właściwie $N$.
Z maksymalności $N$ dostajemy $NM=G$.

Rozważmy homomorfizm $f:N\to G/M$ dany wzorem $f(n)=nM$.
Jest surjektywny: każdy $g\in G=NM$ ma postać $nm$,
a wtedy $gM=nM$. Jądro to $N\cap M=L$.
Pierwsze twierdzenie o izomorfizmie daje $N/L\cong G/M$.
W ten sam sposób $M/L\cong G/N$.
Oba ilorazy po prawej są proste na mocy
Lematu~\ref{aa:lem:maxnormal}, więc również ilorazy po lewej.
\end{proof}

\begin{twierdzenie}[Jordana--Höldera]
\label{aa:tw:jordan-holder}
Niech $G$ będzie grupą skończoną. Jeśli
$G=G_0\triangleright G_1\triangleright\cdots\triangleright G_r=\{e\}$
oraz $G=H_0\triangleright H_1\triangleright\cdots\triangleright H_s=\{e\}$
są dwoma szeregami kompozycyjnymi, to $r=s$ i istnieje
permutacja $\sigma$ zbioru $\{0,\ldots,r-1\}$ taka, że
$G_i/G_{i+1}\cong H_{\sigma(i)}/H_{\sigma(i)+1}$
dla wszystkich $i$.
Nie twierdzimy, że czynniki w obu szeregach występują w tej samej kolejności.
\end{twierdzenie}
\begin{proof}
Przeprowadzimy indukcję względem $|G|$. Dla grupy trywialnej
oba szeregi mają długość zero. Dla grupy prostej każdy szereg
kompozycyjny ma długość jeden i tezę widać bezpośrednio.
Niech teraz $G$ będzie nieprosta i niech $N=G_1$, $M=H_1$.
Podgrupy $N,M$ są maksymalnymi właściwymi podgrupami normalnymi
$G$, bo $G/N$ i $G/M$ są proste.

Jeśli $N=M$, początkowe czynniki $G/N$ i $G/M$ są jednakowe,
a pozostałe części obu szeregów są szeregami kompozycyjnymi $N$.
Ponieważ $|N|<|G|$, założenie indukcyjne utożsamia ich
czynniki z dokładnością do permutacji. Dołączając wspólny czynnik
$G/N$, otrzymujemy tezę.

Niech $N\ne M$ i $L=N\cap M$.
Poprzedni lemat daje proste ilorazy
$N/L\cong G/M$ oraz $M/L\cong G/N$.
Wybierzmy szereg kompozycyjny $L$: jeśli $L=\{e\}$,
przyjmujemy szereg długości zero, a jeśli $L\ne\{e\}$,
jego istnienie zapewnia Twierdzenie~\ref{aa:tw:existcomp}.
Oznaczmy wielozbiór \emph{typów izomorfizmu} jego czynników
przez $\mathcal F(L)$. Zapis $\{\!\{Q_1,\ldots,Q_k\}\!\}$ oznacza wielozbiór
z zachowaniem krotności, a $\uplus$ oznacza sumę takich wielozbiorów.
Dopisanie $N\triangleright L$ do tego szeregu daje szereg
kompozycyjny $N$, którego lista czynników składa się
z $N/L$ i czynników $\mathcal F(L)$.
Ponieważ $|N|<|G|$, założenie indukcyjne zastosowane do dwóch
szeregów kompozycyjnych $N$ (oryginalnego ogona
$G_1\triangleright\cdots\triangleright G_r$
i właśnie skonstruowanego) pokazuje, że czynniki oryginalnego
ogona są, z dokładnością do permutacji, równe
$\{\!\{N/L\}\!\}\uplus\mathcal F(L)$.
Analogiczne rozumowanie dla $M$ daje
$\{\!\{M/L\}\!\}\uplus\mathcal F(L)$.

Dołączamy pierwsze czynniki $G/N$ i $G/M$.
Dla pierwszego szeregu dostajemy wielozbiór
$\{\!\{G/N,\ N/L\}\!\}\uplus\mathcal F(L)$,
a dla drugiego
$\{\!\{G/M,\ M/L\}\!\}\uplus\mathcal F(L)$.
Izomorfizmy z poprzedniego lematu zamieniają parami
$G/N\cong M/L$ i $G/M\cong N/L$, więc oba wielozbiory
są takie same z dokładnością do izomorfizmu. Ich liczności
są zatem równe, co daje również $r=s$.
\end{proof}

\begin{lemat}[Proste grupy abelowe skończone]
\label{aa:lem:finite-abelian-simple}
Niech $G$ będzie niebanalną skończoną grupą abelową.
Wtedy $G$ jest grupą prostą dokładnie wtedy, gdy
$G\cong C_p$ dla pewnej liczby pierwszej $p$.
\end{lemat}
\begin{proof}
Załóżmy, że $G$ jest prosta. Wybierzmy element $g\ne e$.
Podgrupa $\langle g\rangle$ jest niebanalna. Każda podgrupa
grupy abelowej jest normalna, zatem prostota $G$ wymusza
$\langle g\rangle=G$: grupa $G$ jest cykliczna.
Niech $n=|G|$ będzie rzędem jej generatora.
Gdyby $n$ było liczbą złożoną, istniałby dzielnik $d$
spełniający $1<d<n$. Element $g^{n/d}$ miałby rząd $d$:
$(g^{n/d})^d=e$, a dla $0<k<d$ nierówność
$0<kn/d<n$ wyklucza $(g^{n/d})^k=e$.
Wówczas $\langle g^{n/d}\rangle$ byłaby niebanalną,
właściwą podgrupą normalną $G$, co jest sprzeczne
z prostotą. Zatem $n=p$ jest liczbą pierwszą.

Odwrotnie, w grupie cyklicznej $C_p$ rzędu pierwszego
rząd każdej podgrupy dzieli $p$ na mocy twierdzenia Lagrange'a.
Może więc wynosić tylko $1$ albo $p$, a zatem
$C_p$ nie ma niebanalnych właściwych podgrup normalnych.
\end{proof}

\begin{przyklad}[Ta sama lista czynników, inna kolejność]
\label{aa:prz:jordan-c6}
Rozważmy addytywną grupę $C_6=\mathbb Z/6\mathbb Z$.
Ma ona podgrupy $H_2=\{[0]_6,[3]_6\}$ rzędu $2$
oraz $H_3=\{[0]_6,[2]_6,[4]_6\}$ rzędu $3$.
W grupie abelowej obie są normalne. Istnieją zatem
różne szeregi kompozycyjne
\[
 C_6\triangleright H_2\triangleright\{0\},\qquad
 C_6\triangleright H_3\triangleright\{0\}.
\]
Pierwszy ma czynniki $C_6/H_2\cong C_3$ i
$H_2\cong C_2$, drugi zaś $C_6/H_3\cong C_2$ i
$H_3\cong C_3$. Izomorfizm ilorazu z $C_3$
w pierwszym szeregu jest indukowany przez redukcję
klas modulo $3$; analogicznie redukcja modulo $2$
opisuje drugi. Czynniki występują w innej kolejności,
lecz ich wielozbiór z dokładnością do izomorfizmu
pozostaje niezmieniony, jak przewiduje
Twierdzenie~\ref{aa:tw:jordan-holder}.
\end{przyklad}

\begin{figure}[htbp]
\centering
\begin{tikzpicture}[every node/.style={font=\small},line width=.65pt]
 \node (G) at (0,2.8) {$C_6$};
 \node (H2) at (-3.2,1.35) {$H_2\cong C_2$};
 \node (H3) at (3.2,1.35) {$H_3\cong C_3$};
 \node (E) at (0,-.1) {$\{0\}$};
 \draw (G) -- node[above left,pos=.55] {$C_3$} (H2);
 \draw (G) -- node[above right,pos=.55] {$C_2$} (H3);
 \draw (H2) -- node[below left,pos=.5] {$C_2$} (E);
 \draw (H3) -- node[below right,pos=.5] {$C_3$} (E);
\end{tikzpicture}
\caption{Dwie różne drogi od $C_6$ do podgrupy trywialnej.
 Etykieta każdej krawędzi oznacza izomorficzny typ odpowiedniego
 ilorazu. Obie drogi mają czynniki $C_2$ i $C_3$, ale w odwrotnej
 kolejności.}
\label{aa:fig:c6-composition-chains}
\end{figure}

\begin{przyklad}[Szereg kompozycyjny $S_3$ i kolejność czynników]
\label{aa:prz:jordan-s3}
Oznaczmy przez $A_3=\langle(1\,2\,3)\rangle$ podgrupę
rzędu $3$ grupy $S_3$. Szereg
$S_3\triangleright A_3\triangleright\{e\}$
ma czynniki $S_3/A_3\cong C_2$ oraz $A_3\cong C_3$.
Nie możemy zacząć od żadnej podgrupy rzędu $2$,
ponieważ nie jest ona normalna w $S_3$.
To przykład, w którym kolejność czynników jest wymuszona przez strukturę grupy.
\end{przyklad}

\begin{przyklad}[Szereg w $A_4$ i subtelność normalności]
\label{aa:prz:jordan-a4}
Niech $V_4=\{e,a,b,c\}$, gdzie
$a=(1\,2)(3\,4)$, $b=(1\,3)(2\,4)$,
$c=(1\,4)(2\,3)$.
Jest to normalna podgrupa $A_4$ rzędu $4$
(Przykład~\ref{aa:prz:sylow-a4}).
Wybierzmy $K=\{e,(1\,2)(3\,4)\}$.
Ponieważ $V_4$ jest abelowa, $K\trianglelefteq V_4$.
Zatem
$A_4\triangleright V_4\triangleright K\triangleright\{e\}$
jest szeregiem kompozycyjnym o czynnikach
$C_3,C_2,C_2$ (w tej kolejności).
Jednak $K$ nie jest normalna w $A_4$: sprzężenie
przez $(1\,2\,3)$ przekształca $(1\,2)(3\,4)$
w $(1\,4)(2\,3)$, co nie należy do $K$.
Właśnie dlatego definicja szeregu wymaga normalności
w wyrazie poprzednim, a nie w całym $G$.
\end{przyklad}

\begin{uwaga}[Czego twierdzenie nie mówi?]
Czynniki kompozycyjne nie wyznaczają grupy z dokładnością do
izomorfizmu. Na przykład $C_4$ oraz $C_2\times C_2$
mają po dwa czynniki kompozycyjne izomorficzne z $C_2$,
a jednak nie są izomorficzne: pierwsza grupa ma element
rzędu $4$, druga nie. W kolejnym rozdziale zobaczymy,
że sposób składania grup z tych samych czynników jest
odrębną informacją, opisywaną przez rozszerzenia grup.
\end{uwaga}

\section{Komutatory i grupy rozwiązalne}
\begin{definicja}[Komutator i podgrupa pochodna]
\label{aa:def:pochodna}
Dla $x,y\in G$ kładziemy $[x,y]=xyx^{-1}y^{-1}$.
Dla podgrup $A,B\le G$ oznaczamy przez $[A,B]$ podgrupę
wygenerowaną przez wszystkie $[a,b]$, gdzie $a\in A$, $b\in B$.
\emph{Podgrupą pochodną} $G'=[G,G]$ nazywamy podgrupę
generowaną przez wszystkie komutatory elementów $G$.
\end{definicja}
\begin{lemat}[Najmniejsza normalna podgrupa o abelowym ilorazie]
\label{aa:lem:derived-abelian}
Dla każdej grupy $G$ podgrupa $G'$ jest charakterystyczna
(to znaczy: zachowuje ją każdy automorfizm $G$),
a więc normalna w $G$. Iloraz $G/G'$ jest abelowy.
Jeżeli $N\trianglelefteq G$, to $G/N$ jest abelowa
dokładnie wtedy, gdy $G'\subseteq N$.
\end{lemat}
\begin{proof}
Każdy automorfizm $\alpha:G\to G$ spełnia
$\alpha([x,y])=[\alpha(x),\alpha(y)]$, zatem
przekształca zbiór wszystkich komutatorów na ten sam zbiór.
Dlatego $\alpha(G')=G'$, co jest charakterystycznością.
Każde sprzężenie jest automorfizmem, więc $G'\trianglelefteq G$.
W ilorazie $G/G'$ równość $[x,y]\in G'$ daje
$(xG')(yG')=(yG')(xG')$, czyli przemienność.
Dla ogólnej $N\trianglelefteq G$ analogicznie:
$G/N$ jest abelowa wtedy i tylko wtedy, gdy
$(xyx^{-1}y^{-1})N=N$ dla wszystkich $x,y$,
czyli gdy wszystkie komutatory, a więc $G'$, należą do $N$.
\end{proof}

\begin{definicja}[Szereg pochodny i rozwiązalność]
\label{aa:def:rozwiazalna}
Definiujemy $G^{(0)}=G$ oraz $G^{(i+1)}=[G^{(i)},G^{(i)}]$
dla $i\ge0$. Jest to \emph{szereg pochodny} $G$.
Grupa $G$ jest \emph{rozwiązalna}, jeśli istnieje $r\ge0$
z $G^{(r)}=\{e\}$.
Najmniejsze takie $r$ nazywamy jej długością pochodną.
\end{definicja}

\begin{twierdzenie}[Kryterium rozwiązalności przez ilorazy abelowe]
\label{aa:tw:solvable-series}
Grupa $G$ jest rozwiązalna wtedy i tylko wtedy, gdy istnieje
skończony szereg podnormalny
$G=G_0\trianglerighteq G_1\trianglerighteq\cdots
\trianglerighteq G_r=\{e\}$,
w którym każdy iloraz $G_i/G_{i+1}$ jest abelowy.
\end{twierdzenie}
\begin{proof}
Jeśli $G^{(r)}=\{e\}$, szereg pochodny od $G^{(0)}$
do $G^{(r)}$ jest żądanym szeregiem:
$G^{(i+1)}$ jest normalna w $G^{(i)}$,
a $G^{(i)}/G^{(i+1)}$ jest abelowy
z poprzedniego lematu.

Odwrotnie, niech dany będzie szereg o abelowych czynnikach.
Pokażemy indukcyjnie $G^{(i)}\subseteq G_i$.
Dla $i=0$ jest to równość. Gdy $G^{(i)}\subseteq G_i$,
to
$G^{(i+1)}=[G^{(i)},G^{(i)}]\subseteq[G_i,G_i]$.
Ponieważ $G_i/G_{i+1}$ jest abelowy,
Lemat~\ref{aa:lem:derived-abelian} zastosowany do $G_i$
daje $[G_i,G_i]\subseteq G_{i+1}$.
Otrzymujemy więc $G^{(i+1)}\subseteq G_{i+1}$.
Dla $i=r$ mamy $G^{(r)}\subseteq G_r=\{e\}$,
co kończy dowód.
\end{proof}

\begin{lemat}[Podgrupy, ilorazy i rozszerzenia grup rozwiązalnych]
\label{aa:lem:solvable-extension}
Niech $G$ będzie grupą, $H\le G$ i $N\trianglelefteq G$.
\begin{enumerate}
\item Jeśli $G$ jest rozwiązalna, to $H$ i $G/N$ są rozwiązalne.
\item Jeśli $N$ i $G/N$ są rozwiązalne, to $G$ jest rozwiązalna.
\end{enumerate}
\end{lemat}
\begin{proof}
Dla podgrupy $H$ indukcja daje
$H^{(i)}\subseteq G^{(i)}$, bo komutatory elementów
$H^{(i)}$ są komutatorami elementów $G^{(i)}$.
Jeśli więc $G^{(r)}=\{e\}$, także $H^{(r)}=\{e\}$.
Niech $\pi:G\to G/N$ będzie naturalnym rzutem.
Dla każdego $i$ mamy
$\pi(G^{(i)})=(G/N)^{(i)}$:
w bazie $i=0$ jest to surjektywność $\pi$, a krok indukcyjny
wynika z $\pi([x,y])=[\pi(x),\pi(y)]$ i zachowania
generowania podgrupy przez homomorfizm.
Jeśli $G^{(r)}=\{e\}$, to $(G/N)^{(r)}=\{e\}$.

Załóżmy teraz $N^{(s)}=\{e\}$ i $(G/N)^{(r)}=\{e\}$.
Powyższa równość obrazów oznacza $G^{(r)}\subseteq N$.
Dla dowolnej podgrupy $A\le B$ zachodzi $A^{(j)}\subseteq B^{(j)}$,
zatem $(G^{(r)})^{(s)}\subseteq N^{(s)}=\{e\}$.
Z definicji szeregu pochodnego
$(G^{(r)})^{(s)}=G^{(r+s)}$, co dowodzi rozwiązalności $G$.
\end{proof}

\begin{twierdzenie}[Czynniki kompozycyjne a rozwiązalność]
\label{aa:tw:solvable-composition-factors}
Niech $G$ będzie grupą skończoną. Następujące warunki są
równoważne:
\begin{enumerate}
\item $G$ jest rozwiązalna;
\item w każdym szeregu kompozycyjnym $G$ wszystkie czynniki
są cykliczne o rzędach pierwszych;
\item istnieje szereg kompozycyjny $G$, w którym wszystkie
czynniki są cykliczne o rzędach pierwszych.
\end{enumerate}
Dla $G=\{e\}$ dwa ostatnie warunki rozumiemy jako prawdziwe:
pusty szereg nie ma czynników.
\end{twierdzenie}
\begin{proof}
Dla grupy trywialnej teza wynika bezpośrednio z definicji.
Niech $G$ będzie niebanalna; na mocy
Twierdzenia~\ref{aa:tw:existcomp} posiada szereg kompozycyjny.

Załóżmy (1) i wybierzmy dowolny szereg kompozycyjny
$G=G_0\triangleright\cdots\triangleright G_r=\{e\}$.
Z Lematu~\ref{aa:lem:solvable-extension} wynika, że
każda podgrupa $G_i$ jest rozwiązalna, a ponieważ
$G_{i+1}\trianglelefteq G_i$, również iloraz
$S_i=G_i/G_{i+1}$ jest rozwiązalny. Jest to grupa prosta.
Podgrupa pochodna $S_i'$ jest w $S_i$ normalna, zatem
$S_i'=\{e\}$ albo $S_i'=S_i$. Drugi przypadek
powodowałby $S_i^{(k)}=S_i\ne\{e\}$ dla każdego $k\ge1$,
co przeczy rozwiązalności. Wobec tego $S_i'=\{e\}$.
Z Lematu~\ref{aa:lem:derived-abelian} wynika, że
$S_i$ jest abelowa, a z
Lematu~\ref{aa:lem:finite-abelian-simple}, że $S_i\cong C_{p_i}$
dla pewnej liczby pierwszej $p_i$. Ponieważ wybrany szereg
był dowolny, otrzymaliśmy (2).

Implikacja (2)$\Rightarrow$(3) wynika z istnienia
szeregu kompozycyjnego. Przy (3) mamy szereg podnormalny
z abelowymi ilorazami, gdyż każda grupa $C_{p_i}$
jest abelowa. Twierdzenie~\ref{aa:tw:solvable-series}
daje rozwiązalność $G$, czyli (1).
\end{proof}

\begin{przyklad}[Grupy $S_3$, $S_4$ i $A_5$]
Dla $S_3$ podgrupa komutatorowa to $A_3$.
Zawieranie $S_3'\subseteq A_3$ wynika z abelowości
$S_3/A_3\cong C_2$. Ponadto $S_3$ nie jest abelowa,
więc $S_3'\ne\{e\}$. Ponieważ $A_3$ ma rząd pierwszy,
musi być $S_3'=A_3$; dalej $A_3'=\{e\}$.
Grupa $S_3$ jest zatem rozwiązalna o długości pochodnej $2$.
Grupa $S_4$ ma szereg normalności kolejnych wyrazów
$S_4\triangleright A_4\triangleright V_4\triangleright\{e\}$
o abelowych ilorazach $C_2,C_3,V_4$, więc jest rozwiązalna.
Możemy tu dokładniej wyznaczyć szereg pochodny.
Homomorfizm znaku pokazuje $S_4'\subseteq A_4$.
Komutator dwóch transpozycji o jednym wspólnym punkcie,
na przykład $[(1\,2),(2\,3)]$, jest $3$-cyklem.
Ponieważ $S_4'$ jest normalna w $S_4$, zawiera więc
wszystkie $3$-cykle, które są sprzężone w $S_4$.
Każdą parzystą permutację można zapisać jako iloczyn
parzystej liczby transpozycji, a następnie pogrupować je w pary.
Każdą taką parę zastępujemy iloczynem $3$-cykli:
parę jednakowych transpozycji usuwamy, parę różnych
transpozycji o jednym wspólnym punkcie zastępujemy
$3$-cyklem, a parę rozłączną, $(a\,b)(c\,d)$, zapisujemy
jako $(a\,c\,b)(a\,c\,d)$, przy przyjętym składaniu
od prawej do lewej. Zatem $S_4'=A_4$.

Dalej $A_4'\subseteq V_4$, gdyż $A_4/V_4\cong C_3$.
Jednak $A_4$ nie jest abelowa: $3$-cykle $(1\,2\,3)$
i $(1\,2\,4)$ nie komutują, zatem $A_4'\ne\{e\}$.
Podgrupa $A_4'$ jest normalna w $A_4$, więc wraz
z dowolnym niezerowym elementem $V_4$ zawiera jego
wszystkie sprzężenia w $A_4$. Trzy nieidentycznościowe
elementy $V_4$ są sprzężone przez $3$-cykle,
a więc $A_4'=V_4$. Ponieważ $V_4$ jest abelowa,
$V_4'=\{e\}$. Ostatecznie szereg pochodny ma postać
$S_4\triangleright A_4\triangleright V_4\triangleright\{e\}$,
a długość pochodna $S_4$ wynosi dokładnie $3$.
Dla $A_5$ sytuacja jest inna: pokażemy poniżej,
że to grupa prosta i nieabelowa, co wykluczy rozwiązalność.
\end{przyklad}

\section{Nieabelowa grupa prosta: \texorpdfstring{$A_5$}{A5}}
\begin{twierdzenie}[Prostota grupy alternującej $A_5$]
\label{aa:tw:a5-simple}
Grupa $A_5$, składająca się z parzystych permutacji pięciu punktów,
ma rząd $60$ i jest nieabelową grupą prostą.
\end{twierdzenie}
\begin{proof}
Mamy $|A_5|=5!/2=60$ na mocy homomorfizmu znaku
(\ref{aa:lem:znak}). Wypiszemy wszystkie klasy sprzężoności
w $A_5$, wyprowadzając ich rozmiary zamiast przyjmować je jako dane.
Elementami $A_5$ są: identyczność, $20$ cykli długości $3$,
$15$ iloczynów dwóch rozłącznych transpozycji i $24$ cykle długości $5$.
Rachunki są następujące: $\binom53\cdot2=20$,
$\binom54\cdot3=15$, a $5!/5=24$.
Nie ma innych typów parzystych permutacji pięciu punktów.

Rozważmy $c=(1\,2\,3)$. Permutacja $g\in S_5$ komutująca
z $c$ musi zachowywać nośnik $\{1,2,3\}$ i działać na nim
jak jedna z trzech potęg $c$; na dwóch pozostałych punktach
może wykonać dowolną permutację. Zatem $|C_{S_5}(c)|=3\cdot2=6$,
a dokładnie trzy takie permutacje są parzyste,
więc $|C_{A_5}(c)|=3$. Klasa $c$ w $A_5$ ma $60/3=20$
elementów; obejmuje zatem wszystkie cykle długości $3$.

Dla $d=(1\,2)(3\,4)$ element centralizatora w $S_5$
musi ustalać $5$ i permutować dwie dwuelementowe pary
$\{1,2\}$, $\{3,4\}$: może zamieniać punkty wewnątrz każdej
pary (cztery możliwości) i ewentualnie zamienić całe pary
(dwie możliwości). Mamy więc $|C_{S_5}(d)|=8$.
Wśród tych ośmiu permutacji dokładnie cztery są parzyste:
mnożenie przez nieparzystą transpozycję $(1\,2)$
przełącza parzystość i pozostaje w centralizatorze.
Stąd $|C_{A_5}(d)|=4$ i klasa ma $60/4=15$ elementów.
Zawiera zatem wszystkie iloczyny dwóch rozłącznych transpozycji.

Dla pięciocyklu $u=(1\,2\,3\,4\,5)$ permutacja komutująca
z $u$ jest wyznaczona przez obraz punktu $1$;
jego obraz może być jednym z pięciu punktów, a następne obrazy
są wymuszone przez równość $gu=ug$.
Centralizator w $S_5$ to $\langle u\rangle$ rzędu $5$.
Wszystkie potęgi $u$ są parzyste, więc
$|C_{A_5}(u)|=5$. Klasa w $A_5$ ma $60/5=12$ elementów.
Ponieważ pięciocykli jest $24$, tworzą dwie takie klasy.
Ostatecznie klasy sprzężoności $A_5$ mają rozmiary
$1,15,20,12,12$.

Jeśli $N\trianglelefteq A_5$, to $N$ jest sumą rozłącznych
klas sprzężoności i zawiera identyczność. Dlatego $|N|$
ma postać $1$ powiększonego o sumę pewnych liczb z listy
$15,20,12,12$. Wszystkie możliwe wartości to
$1,13,16,21,25,28,33,36,40,45,48,60$.
Z twierdzenia Lagrange'a $|N|$ musi dzielić $60$,
a z wypisanej listy jedynie $1$ i $60$ są dzielnikami $60$.
Stąd $N=\{e\}$ lub $N=A_5$, czyli $A_5$ jest prosta.
Jest nieabelowa: dla $u=(1\,2\,3)$ i $v=(1\,2\,4)$
mamy $uvu^{-1}=(2\,3\,4)\ne(1\,2\,4)=v$,
zatem $uv\ne vu$.
\end{proof}

\begin{wniosek}[Dlaczego $A_5$ nie jest rozwiązalna?]
\label{aa:wn:a5-not-solvable}
Grupa $A_5$ nie jest rozwiązalna.
\end{wniosek}
\begin{proof}
Gdyby $(A_5)'=\{e\}$, grupa byłaby abelowa.
Skoro jest nieabelowa, $(A_5)'\ne\{e\}$.
Podgrupa pochodna jest normalna, a $A_5$ jest prosta,
więc $(A_5)'=A_5$. Indukcja daje $(A_5)^{(i)}=A_5$
dla wszystkich $i$, zatem szereg pochodny nigdy nie osiąga
podgrupy trywialnej.
\end{proof}

\section{Grupy nilpotentne i szeregi centralne}
\begin{definicja}[Dolny szereg centralny]
\label{aa:def:dolny-centralny}
Dla grupy $G$ definiujemy
$\gamma_1(G)=G$ i $\gamma_{i+1}(G)=[\gamma_i(G),G]$ dla $i\ge1$.
Grupa $G$ jest \emph{nilpotentna}, jeśli $\gamma_{c+1}(G)=\{e\}$
dla pewnej liczby całkowitej $c\ge0$.
Najmniejszą taką liczbę $c$ nazywamy \emph{klasą nilpotentności}.
Dla grupy trywialnej klasa wynosi $0$; dla niebanalnej grupy
abelowej wynosi $1$.
\end{definicja}
\begin{lemat}[Własności dolnego szeregu centralnego]
\label{aa:lem:lower-properties}
Dla dowolnej grupy $G$ wszystkie podgrupy $\gamma_i(G)$
są charakterystyczne w $G$ i tworzą ciąg malejący.
Jeżeli $c\ge1$, $\gamma_{c+1}(G)=\{e\}$ i
$\gamma_c(G)\ne\{e\}$, to $\gamma_c(G)\subseteq Z(G)$.
Ponadto każda grupa nilpotentna jest rozwiązalna.
\end{lemat}
\begin{proof}
Indukcyjnie załóżmy, że $\gamma_i(G)$ jest charakterystyczna.
Każdy automorfizm przenosi zbiór komutatorów
$[\gamma_i(G),G]$ na siebie, zatem
$\gamma_{i+1}(G)$ jest charakterystyczna.
Ponieważ $\gamma_1(G)=G$, dowodzi to pierwszej części.
Dla $i=1$ mamy $\gamma_2(G)\subseteq G=\gamma_1(G)$.
Gdy $\gamma_i(G)\subseteq\gamma_{i-1}(G)$,
każdy generator $[x,g]$, $x\in\gamma_i(G)$, należy również
 do $[\gamma_{i-1}(G),G]$, więc
$\gamma_{i+1}(G)\subseteq\gamma_i(G)$.
Jeśli $\gamma_{c+1}(G)=[\gamma_c(G),G]=\{e\}$,
to każdy element $\gamma_c(G)$ komutuje z każdym elementem $G$;
stąd $\gamma_c(G)\subseteq Z(G)$.

Dla rozwiązalności zauważmy, że $G^{(0)}=\gamma_1(G)$.
Gdy $G^{(i)}\subseteq\gamma_{i+1}(G)$, mamy
$G^{(i+1)}=[G^{(i)},G^{(i)}]
\subseteq[\gamma_{i+1}(G),\gamma_{i+1}(G)]
\subseteq[\gamma_{i+1}(G),G]=\gamma_{i+2}(G)$.
Tak więc $G^{(i)}\subseteq\gamma_{i+1}(G)$ dla wszystkich $i$.
W szczególności $G^{(c)}=\{e\}$, gdy $G$ jest nilpotentna
klasy $c$.
\end{proof}

\begin{definicja}[Górny szereg centralny]
\label{aa:def:gorny-centralny}
Definiujemy $Z_0(G)=\{e\}$, a następnie rekurencyjnie
$Z_{i+1}(G)$ jako pełny przeciwobraz centrum grupy $G/Z_i(G)$
w naturalnym rzucie $G\to G/Z_i(G)$.
Zatem $Z_{i+1}(G)/Z_i(G)=Z(G/Z_i(G))$.
Konstrukcja jest poprawnie określona: $Z_0(G)$ jest normalna,
a jeśli $Z_i(G)\trianglelefteq G$, to przeciwobraz centrum
ilorazu jest również normalny w $G$.
W szczególności $Z_1(G)=Z(G)$.
\end{definicja}
\begin{lemat}[Porównanie dolnego i górnego szeregu centralnego]
\label{aa:lem:upper-lower}
Dla grupy $G$ i liczby całkowitej $c\ge0$ następujące warunki
są równoważne: (i) $\gamma_{c+1}(G)=\{e\}$;
(ii) $Z_c(G)=G$.
\end{lemat}
\begin{proof}
Najpierw wyjaśnijmy, że dla dowolnej podgrupy normalnej
$K\trianglelefteq G$ jej przeciwobraz centrum ilorazu
$G/K$ składa się dokładnie z tych $g\in G$, dla których
$[g,x]\in K$ dla każdego $x\in G$.
Wynika to bezpośrednio z równości
$[gK,xK]=[g,x]K$.

Załóżmy $Z_c(G)=G$. Dowodzimy przez indukcję po $j$, że
$\gamma_{j+1}(G)\subseteq Z_{c-j}(G)$ dla $0\le j\le c$.
Dla $j=0$ jest to $G\subseteq Z_c(G)$.
Jeśli $j<c$ i $\gamma_{j+1}(G)\subseteq Z_{c-j}(G)$,
to każdy komutator z elementem tej podgrupy i z $G$
należy do $Z_{c-j-1}(G)$, ponieważ
$[Z_{k+1}(G),G]\subseteq Z_k(G)$.
Stąd $\gamma_{j+2}(G)\subseteq Z_{c-j-1}(G)$.
W przypadku $j=c$ otrzymujemy $\gamma_{c+1}(G)\subseteq Z_0(G)=\{e\}$.

W drugą stronę załóżmy $\gamma_{c+1}(G)=\{e\}$.
Pokażemy indukcyjnie dla $0\le j\le c$ zawierania
$\gamma_{c-j+1}(G)\subseteq Z_j(G)$.
Dla $j=0$ jest to założenie. Jeśli $j<c$ i
$\gamma_{c-j+1}(G)\subseteq Z_j(G)$,
to $[\gamma_{c-j}(G),G]=\gamma_{c-j+1}(G)\subseteq Z_j(G)$.
Z podanego na początku opisu przeciwobrazu centrum
wynika $\gamma_{c-j}(G)\subseteq Z_{j+1}(G)$.
Dla $j=c$ dostajemy $\gamma_1(G)=G\subseteq Z_c(G)$,
więc $Z_c(G)=G$.
\end{proof}

\begin{twierdzenie}[Każda skończona $p$-grupa jest nilpotentna]
\label{aa:tw:pgroup-nilpotent}
Niech $G$ będzie skończoną grupą rzędu $p^a$, gdzie $p$
jest liczbą pierwszą, a $a\ge0$. Wtedy $G$ jest nilpotentna.
\end{twierdzenie}
\begin{proof}
Dla $a=0$ mamy $G=\{e\}$, a zatem już
$\gamma_1(G)=\{e\}$; grupa ta jest nilpotentna klasy $0$.
Niech $a\ge1$. Rozważmy górny szereg centralny.
Jeśli $Z_i(G)\ne G$, iloraz $G/Z_i(G)$ jest niebanalną
skończoną $p$-grupą. Z Twierdzenia~\ref{aa:tw:centrum-pgrupy}
jego centrum jest nietrywialne. Z definicji $Z_{i+1}(G)$
pełny przeciwobraz tego centrum zawiera właściwie $Z_i(G)$.
Wzrastający ciąg podgrup nie może mieć więcej niż $a$
ścisłych zawierań, ponieważ przy każdym ścisłym zawieraniu
rząd podgrupy rośnie co najmniej $p$ razy, a $|G|=p^a$.
Zatem po skończenie wielu krokach $Z_c(G)=G$.
Lemat~\ref{aa:lem:upper-lower} daje nilpotentność $G$.
\end{proof}

\begin{przyklad}[Grupa diedralna rzędu osiem]
\label{aa:prz:d8-nilpotent}
Niech $D_8=\langle r,s\mid r^4=e,\ s^2=e,\ srs=r^{-1}\rangle$
będzie grupą symetrii kwadratu, opisaną przez obrót o $90^\circ$
i odbicie (por. Przykład~\ref{aa:prz:d4}).
Zbiór jej elementów ma postać
$\{e,r,r^2,r^3,s,rs,r^2s,r^3s\}$.
Widać, że $[r,s]=rsr^{-1}s^{-1}=r^2$;
ponieważ $r^2$ komutuje z $r$ i $s$, jest centralny.
Iloraz przez $\langle r^2\rangle$ jest generowany przez
komutujące elementy rzędu $2$, więc jest abelowy.
Wobec tego $D_8'=\langle r^2\rangle$,
a $\gamma_3(D_8)=[\langle r^2\rangle,D_8]=\{e\}$.
Grupa $D_8$ jest nilpotentna klasy dokładnie $2$, lecz nie jest abelowa.
\end{przyklad}

\begin{przyklad}[Grupa rozwiązalna, która nie jest nilpotentna]
\label{aa:prz:s3-notnil}
Grupa $S_3$ jest rozwiązalna, jak dowiedliśmy powyżej,
lecz $Z(S_3)=\{e\}$: żaden nieidentycznościowy
$3$-cykl nie komutuje z transpozycją, a każda transpozycja
nie komutuje z inną transpozycją.
Gdyby $S_3$ była nilpotentna, ostatni nietrywialny wyraz
jej dolnego szeregu centralnego byłby, na mocy
Lematu~\ref{aa:lem:lower-properties}, zawarty w centrum.
Centrum jest jednak trywialne; sprzeczność.

Można też prześledzić dokładnie dolny szereg centralny.
Niech $r=(1\,2\,3)$ i $s=(1\,2)$.
Wtedy $srs^{-1}=r^{-1}$, a więc
$[r,s]=rsr^{-1}s^{-1}=r^2$.
W szczególności $[A_3,S_3]$ zawiera generator $r^2$
grupy $A_3$. Z drugiej strony
$[A_3,S_3]\subseteq A_3$, ponieważ $A_3$ jest normalna.
Zatem $[A_3,S_3]=A_3$. Wiemy już, że
$\gamma_2(S_3)=S_3'=A_3$, toteż rekurencyjnie
$\gamma_i(S_3)=A_3$ dla wszystkich $i\ge2$.
Szereg dolny nie osiąga zatem grupy trywialnej.
Rozwiązalność nie implikuje nilpotentności.
\end{przyklad}

\enlargethispage{3\baselineskip}
\begin{uwaga}[Co zostanie użyte przy konstrukcjach grup?]
W następnym rozdziale zobaczymy, jak z dwóch grup budować
iloczyny proste i półproste oraz jak interpretować rozszczepiające
się krótkie ciągi dokładne. Wykorzystamy te konstrukcje
do klasyfikacji grup rzędu $pq$ i do charakterystyki
skończonych grup nilpotentnych przez ich podgrupy Sylowa.
\end{uwaga}

% AA-06 | Algebra abstrakcyjna | Iteracja 03
\chapter[Iloczyny i rozszerzenia grup]{Iloczyny proste, półproste i rozszerzenia grup}
\label{aa:ch:konstrukcje-grup}

\section{Iloczyn prosty zewnętrzny}
\begin{definicja}[Iloczyn prosty grup]
\label{aa:def:directproduct}
Niech $G,H$ będą grupami o elementach neutralnych $e_G,e_H$.
Na zbiorze kartezjańskim $G\times H$ definiujemy mnożenie
$(g,h)(g',h')=(gg',hh')$.
Tak otrzymany obiekt nazywamy \emph{zewnętrznym iloczynem prostym}
$G\times H$.
\end{definicja}
\begin{twierdzenie}[Aksjomaty i struktura iloczynu prostego]
\label{aa:tw:direct-product}
Dla dowolnych grup $G,H$ iloczyn $G\times H$ jest grupą,
której jedynką jest $(e_G,e_H)$, a odwrotnością
$(g,h)^{-1}=(g^{-1},h^{-1})$.
Odwzorowania $\pi_G(g,h)=g$, $\pi_H(g,h)=h$ są surjektywnymi
homomorfizmami. Podgrupy $G\times\{e_H\}$ i
$\{e_G\}\times H$ są normalne, mają trywialne przecięcie
i komutują element po elemencie.
Ponadto $Z(G\times H)=Z(G)\times Z(H)$.
\end{twierdzenie}
\begin{proof}
Łączność wynika z równości
\begin{align*}
((g,h)(g',h'))(g'',h'')&=((gg')g'',(hh')h''),\\
(g,h)((g',h')(g'',h''))&=(g(g'g''),h(h'h'')).
\end{align*}
Obie pary są równe z łączności mnożenia w $G$ i $H$.
Mnożenie przez $(e_G,e_H)$ pozostawia każdą parę bez zmiany,
a podana odwrotność jest obustronna, gdyż w każdej współrzędnej
mamy iloczyn elementu i jego odwrotności.
Projekcje zachowują mnożenie z definicji. Są surjektywne,
ponieważ $\pi_G(g,e_H)=g$ dla każdego $g\in G$
oraz $\pi_H(e_G,h)=h$ dla każdego $h\in H$.
Jądro $\pi_H$ to $G\times\{e_H\}$,
a jądro $\pi_G$ to $\{e_G\}\times H$, więc podgrupy
te są normalne na mocy twierdzenia o jądrze homomorfizmu.
Ich przekrój to jedynka. Równości
$(g,e_H)(e_G,h)=(g,h)=(e_G,h)(g,e_H)$
pokazują przemienność obu czynników.
Wreszcie $(g,h)$ komutuje z każdą $(x,y)$ dokładnie wtedy,
gdy $gx=xg$ dla każdego $x\in G$ oraz $hy=yh$ dla każdego
$y\in H$. To oznacza $g\in Z(G)$ i $h\in Z(H)$.
\end{proof}

\begin{stwierdzenie}[Własność odwzorowań do iloczynu]
\label{aa:st:product-universal}
Dla grup $K,G,H$ oraz homomorfizmów $f:K\to G$,
$g:K\to H$ istnieje dokładnie jeden homomorfizm
$u:K\to G\times H$ spełniający $\pi_G\circ u=f$
i $\pi_H\circ u=g$, mianowicie $u(k)=(f(k),g(k))$.
\end{stwierdzenie}
\begin{proof}
Dla $a,b\in K$ mamy
$u(ab)=(f(a)f(b),g(a)g(b))=u(a)u(b)$.
Jeśli $v$ spełnia warunki, jego obie współrzędne są wymuszone:
$v(k)=(\pi_G(v(k)),\pi_H(v(k)))=(f(k),g(k))=u(k)$.
\end{proof}

\begin{przyklad}[Dlaczego $C_6\cong C_2\times C_3$, ale $C_4\not\cong C_2\times C_2$?]
Niech $C_n=\mathbb Z/n\mathbb Z$ z dodawaniem.
Odwzorowanie $f:C_6\to C_2\times C_3$,
$f([k]_6)=([k]_2,[k]_3)$, jest dobrze określone,
bo przystawanie modulo $6$ pociąga przystawanie
modulo $2$ i modulo $3$; ponadto jest homomorfizmem.
Jeśli $f([k]_6)=(0,0)$, to $2\mid k$ i $3\mid k$,
więc $6\mid k$. Jądro jest trywialne, toteż $f$ jest
różnowartościowy. Obie grupy mają sześć elementów,
zatem $f$ jest izomorfizmem.
Element $(1,1)$ ma rząd $6$.
Natomiast każdy element $C_2\times C_2$ spełnia
$2(x,y)=(0,0)$, więc ma rząd co najwyżej $2$.
W $C_4$ istnieje element $[1]_4$ rzędu $4$, stąd
$C_4\not\cong C_2\times C_2$.
To szczególny przypadek chińskiego twierdzenia o resztach,
które udowodnimy przy teorii pierścieni.
\end{przyklad}

\section{Iloczyn prosty wewnętrzny: kiedy grupa rozkłada się na czynniki?}
\begin{lemat}[Iloczyn podgrup i jego liczność]
\label{aa:lem:product-subgroups}
Niech $N,H\le G$, przy czym $N\trianglelefteq G$.
Wtedy $NH=\{nh:n\in N,\ h\in H\}$ jest podgrupą $G$.
Jeśli obie podgrupy są skończone, to
$|NH|=|N|\,|H|/|N\cap H|$.
\end{lemat}
\begin{proof}
Iloczyn dwóch elementów $nh,n'h'$ zapisujemy jako
$n(hn'h^{-1})(hh')$. Ponieważ $N$ jest normalna,
$hn'h^{-1}\in N$, a więc wynik należy do $NH$.
Odwrotność $(nh)^{-1}=h^{-1}n^{-1}$ możemy zapisać jako
$(h^{-1}n^{-1}h)h^{-1}\in NH$. Jedynka też należy do $NH$;
stąd kryterium podgrupy daje $NH\le G$.

Policzmy włókna odwzorowania zbiorów $N\times H\to NH$,
$(n,h)\mapsto nh$. Dwie pary mają ten sam obraz wtedy i tylko wtedy,
gdy $n^{-1}n'=hh'^{-1}\in N\cap H$.
Dla dowolnej pary $(n,h)$ i $k\in N\cap H$ para
$(nk,k^{-1}h)$ ma ten sam iloczyn $nh$.
Odwrotnie, powyższa równość pokazuje, że każda para
z tym samym obrazem jest tej postaci dla dokładnie jednego $k$.
Zatem każde włókno ma dokładnie $|N\cap H|$ elementów.
Policzenie par daje $|N|\,|H|=|NH|\,|N\cap H|$.
\end{proof}

\begin{twierdzenie}[Kryterium wewnętrznego iloczynu prostego]
\label{aa:tw:internal-direct}
Niech $G$ będzie grupą, a $N,H\trianglelefteq G$
normalnymi podgrupami takimi, że $N\cap H=\{e\}$
oraz $NH=G$. Wtedy odwzorowanie
$\Phi:N\times H\to G$, $\Phi(n,h)=nh$, jest izomorfizmem.
Odwrotnie, obrazy obu kanonicznych czynników w dowolnym
izomorfizmie $N\times H\cong G$ spełniają te warunki.
\end{twierdzenie}
\begin{proof}
Najpierw wykażemy, że elementy $N$ komutują z elementami $H$.
Dla $n\in N$, $h\in H$ komutator
$[n,h]=nhn^{-1}h^{-1}$ należy do $H$, ponieważ
$nhn^{-1}\in H$ z normalności $H$.
Należy też do $N$, gdyż $hn^{-1}h^{-1}\in N$ z normalności $N$,
więc $[n,h]=n(hn^{-1}h^{-1})\in N$.
Stąd $[n,h]\in N\cap H=\{e\}$, czyli $nh=hn$.
Dla $(n,h),(n',h')\in N\times H$ mamy zatem
$\Phi((n,h)(n',h'))=nn'hh'=nhn'h'=
\Phi(n,h)\Phi(n',h')$. Mapa jest homomorfizmem.
Jest surjektywna, gdyż $NH=G$.
Jej jądro to pary spełniające $nh=e$, czyli $n=h^{-1}\in N\cap H$.
Z założenia $n=h=e$, więc jądro jest trywialne:
$\Phi$ jest izomorfizmem.
Kierunek odwrotny wynika z własności kanonicznych podgrup
iloczynu prostego udowodnionych w Twierdzeniu~\ref{aa:tw:direct-product};
izomorfizm zachowuje normalność, przecięcia i generowanie grupy.
\end{proof}

\begin{uwaga}[Dlaczego normalność obu czynników jest istotna?]
Jeżeli tylko $N\trianglelefteq G$, a $H$ jest dowolną podgrupą
z $NH=G$ i $N\cap H=\{e\}$, to czynniki nie muszą komutować.
W takiej sytuacji otrzymamy \emph{iloczyn półprosty}.
Przykładem jest $S_3$ z normalną podgrupą $A_3$ rzędu $3$
i podgrupą $\langle(1\,2)\rangle$ rzędu $2$.
\end{uwaga}

\section{Iloczyn półprosty zewnętrzny}
\begin{definicja}[Działanie przez automorfizmy i iloczyn półprosty]
\label{aa:def:semidirect}
Niech $N,H$ będą grupami i niech
$\varphi:H\to\operatorname{Aut}(N)$ będzie homomorfizmem,
gdzie $\operatorname{Aut}(N)$ oznacza grupę automorfizmów $N$
ze składaniem odwzorowań. Dla $h\in H$ piszemy
$\varphi_h=\varphi(h)$.
Na zbiorze $N\times H$ definiujemy
\[
 (n,h)\ast(n',h')=(n\,\varphi_h(n'),\ hh').
\]
Grupę tę oznaczamy $N\rtimes_\varphi H$
i nazywamy \emph{iloczynem półprostym}.
\end{definicja}
\begin{twierdzenie}[Iloczyn półprosty jest grupą]
\label{aa:tw:semidirect-group}
Powyższe działanie jest łączne, ma jedynkę $(e_N,e_H)$,
a odwrotność pary $(n,h)$ wynosi
$(\varphi_{h^{-1}}(n^{-1}),h^{-1})$.
Podzbiór $N\times\{e_H\}$ jest normalną podgrupą izomorficzną z $N$,
a $\{e_N\}\times H$ jest podgrupą izomorficzną z $H$.
Ich przekrój jest trywialny, a razem generują całą grupę.
Jeżeli działanie $\varphi$ jest trywialne, to otrzymujemy iloczyn prosty.
\end{twierdzenie}
\begin{proof}
Porównajmy oba możliwe nawiasowania trzech elementów.
Mamy
\begin{align*}
((n,h)\ast(n',h'))\ast(n'',h'')
 &= (n\varphi_h(n')\varphi_{hh'}(n''),\ hh'h''),\\
(n,h)\ast((n',h')\ast(n'',h''))
 &= (n\varphi_h(n'\varphi_{h'}(n'')),\ hh'h'')\\
 &= (n\varphi_h(n')\varphi_h(\varphi_{h'}(n'')),\ hh'h'').
\end{align*}
Ponieważ $\varphi$ jest homomorfizmem,
$\varphi_{hh'}=\varphi_h\circ\varphi_{h'}$,
więc wyrażenia są równe. To dowodzi łączności.
Mamy $\varphi_{e_H}=\operatorname{id}_N$;
ponadto każdy automorfizm przenosi jedynkę na jedynkę.
Stąd $(e_N,e_H)$ jest jedynką dwustronną.
Dla wskazanego elementu odwrotnego obliczamy
\[
 (n,h)\ast(\varphi_{h^{-1}}(n^{-1}),h^{-1})
 =(n\varphi_h(\varphi_{h^{-1}}(n^{-1})),e_H)=(e_N,e_H).
\]
W drugą stronę pierwszy czynnik wynosi
$\varphi_{h^{-1}}(n^{-1})\varphi_{h^{-1}}(n)=e_N$,
więc otrzymujemy również jedynkę.

Włączenie $n\mapsto(n,e_H)$ jest homomorfizmem,
podobnie $h\mapsto(e_N,h)$.
Normalność pierwszego czynnika wynika z obliczenia
\[
 (n,h)\ast(n',e_H)\ast(n,h)^{-1}
 =\bigl(n\varphi_h(n')n^{-1},e_H\bigr).
\]
Wprost widać trywialność przecięcia czynników i równość
$(n,e_H)\ast(e_N,h)=(n,h)$.
Przy trywialnym $\varphi_h=\operatorname{id}_N$
wzór na mnożenie staje się współrzędnościowy,
czyli jest dokładnie mnożeniem w $N\times H$.
\end{proof}

\begin{uwaga}[Znaczenie działania]
Iloczyn półprosty pamięta nie tylko parę grup $N,H$,
lecz także homomorfizm $\varphi:H\to\operatorname{Aut}(N)$.
Różne działania mogą dawać różne grupy.
Relacja
$(e_N,h)(n,e_H)(e_N,h)^{-1}=(\varphi_h(n),e_H)$
pokazuje, że działanie $H$ na $N$ jest dokładnie
sprzężeniem przez drugi czynnik.
W odróżnieniu od iloczynu prostego dwa czynniki nie muszą komutować.
\end{uwaga}

\section{Wewnętrzny iloczyn półprosty i rozszczepione rozszerzenia}
\begin{twierdzenie}[Kryterium wewnętrznego iloczynu półprostego]
\label{aa:tw:internal-semidirect}
Niech $N\trianglelefteq G$ i $H\le G$ spełniają
$NH=G$ oraz $N\cap H=\{e\}$.
Niech $\varphi:H\to\operatorname{Aut}(N)$ będzie działaniem
przez sprzężenie, $\varphi_h(n)=hnh^{-1}$.
Wtedy $G\cong N\rtimes_\varphi H$ przez odwzorowanie
$(n,h)\mapsto nh$.
\end{twierdzenie}
\begin{proof}
Ponieważ $N$ jest normalna, sprzężenie przez element $h\in H$
ogranicza się do automorfizmu $N$.
Z równości $(hh')n(hh')^{-1}=h(h'nh'^{-1})h^{-1}$
wynika $\varphi_{hh'}=\varphi_h\circ\varphi_{h'}$.
Dla par obliczamy
\[
 (nh)(n'h')=n(hn'h^{-1})hh'=
 n\varphi_h(n')hh',
\]
co jest obrazem ich iloczynu półprostego. Mapa jest więc
homomorfizmem. Jest surjektywna z $NH=G$.
Jeśli $nh=e$, mamy $n=h^{-1}\in N\cap H=\{e\}$,
więc jądro jest trywialne; jest to izomorfizm.
Jednoznaczność przedstawienia $g=nh$ wynika też z równości
$n_1h_1=n_2h_2\implies n_2^{-1}n_1=h_2h_1^{-1}\in N\cap H$.
\end{proof}

\begin{definicja}[Krótki ciąg dokładny i rozszczepienie]
\label{aa:def:split-exact}
\emph{Krótkim ciągiem dokładnym grup} nazywamy układ
$1\longrightarrow N\xrightarrow{i}G\xrightarrow{p}Q\longrightarrow1$,
w którym $i$ jest iniektywnym homomorfizmem, $p$ surjektywnym
homomorfizmem i $\operatorname{im}i=\ker p$.
Ciąg \emph{rozszczepia się}, jeżeli istnieje homomorfizm
$s:Q\to G$ z $p\circ s=\operatorname{id}_Q$.
Takie $s$ nazywamy \emph{przekrojem homomorficznym}.
\end{definicja}

\begin{twierdzenie}[Rozszczepienie a iloczyn półprosty]
\label{aa:tw:split-semidirect}
Dla krótkiego ciągu dokładnego
$1\to N\xrightarrow{i}G\xrightarrow{p}Q\to1$
następujące warunki są równoważne:
\begin{enumerate}
\item istnieje przekrój homomorficzny $s:Q\to G$;
\item istnieje podgrupa $H\le G$ taka, że
$G=i(N)H$ i $i(N)\cap H=\{e\}$;
\item przy pewnym działaniu $\varphi:Q\to\operatorname{Aut}(N)$
istnieje izomorfizm $F:N\rtimes_\varphi Q\to G$
taki, że $F(n,e_Q)=i(n)$ dla każdego $n\in N$ oraz
$p(F(n,q))=q$ dla wszystkich $n\in N$, $q\in Q$.
\end{enumerate}
\end{twierdzenie}
\begin{proof}
Załóżmy (1) i połóżmy $H=s(Q)$. Ponieważ $s$ jest homomorfizmem,
$H$ jest podgrupą. Jeżeli $s(q)\in i(N)=\ker p$,
to $q=p(s(q))=e_Q$, zatem $s(q)=e_G$.
Mamy więc $i(N)\cap H=\{e\}$.
Dla dowolnego $g\in G$ element $s(p(g))$ ma ten sam obraz
w $Q$ co $g$; stąd $g s(p(g))^{-1}\in\ker p=i(N)$,
czyli $g\in i(N)H$. Zachodzi (2).

Zakładając (2), zauważamy, że ograniczenie $p|_H:H\to Q$
jest różnowartościowe, ponieważ jego jądro to
$H\cap\ker p=H\cap i(N)=\{e\}$.
Jest surjektywne: jeśli $q=p(g)$, zapiszmy $g=i(n)h$,
a wtedy $q=p(h)$. Zatem $p|_H$ jest izomorfizmem.
Jego odwrotność $s:Q\to H\subseteq G$ jest przekrojem
homomorficznym i wracamy do (1).

Przy (1) definiujemy $\varphi_q(n)$ jako jednoznaczny element
$N$, dla którego $i(\varphi_q(n))=s(q)i(n)s(q)^{-1}$.
Działanie jest dobrze określone, ponieważ $i(N)=\ker p$
jest normalne w $G$, a $i$ jest izomorfizmem na obraz.
Z Twierdzenia~\ref{aa:tw:internal-semidirect}
i podstawienia $h=s(q)$ otrzymujemy izomorfizm
$F(n,q)=i(n)s(q)$. Wprost mamy
$F(n,e_Q)=i(n)$ i $p(F(n,q))=q$, więc spełnia on (3).
Odwrotnie, izomorfizm z (3) przenosi standardowe włączenie
$q\mapsto(e_N,q)$ na przekrój $Q\to G$, więc (3) daje (1).
\end{proof}

\section{Przykłady: grupy diedralne, symetryczne i afiniczne}
\begin{przyklad}[Grupa diedralna jako iloczyn półprosty]
\label{aa:prz:dihedral-semidirect}
Niech $n\ge3$. Oznaczmy przez $r$ obrót foremnego $n$-kąta
o kąt $2\pi/n$, a przez $s$ odbicie w osi symetrii przechodzącej
przez jego środek. Zachodzą relacje
$r^n=e$, $s^2=e$ oraz $srs^{-1}=r^{-1}$.
Grupa $D_{2n}$ wszystkich symetrii ma $2n$ elementów:
$n$ obrotów $r^i$ oraz $n$ odbić $r^i s$.
Każdą symetrię można zapisać dokładnie w jednej z tych postaci;
obroty zachowują orientację płaszczyzny, a odbicia ją odwracają.
Podgrupa $N=\langle r\rangle\cong C_n$ jest normalna,
a $H=\langle s\rangle\cong C_2$ jest jej dopełnieniem,
gdyż $NH=D_{2n}$ i $N\cap H=\{e\}$.
Sprzężenie przez $s$ działa na $C_n$ przez inwersję:
$\varphi_s(r^i)=r^{-i}$.
Jest to automorfizm rzędu co najwyżej $2$, więc wyznacza
homomorfizm $C_2\to\operatorname{Aut}(C_n)$.
Twierdzenie~\ref{aa:tw:internal-semidirect} daje
$D_{2n}\cong C_n\rtimes_{\mathrm{inv}}C_2$.
Grupa nie jest na ogół iloczynem prostym, bo $srs^{-1}\ne r$
dla $n\ge3$.

Dla $n=4$ dostajemy symetrie kwadratu:
$D_8\cong C_4\rtimes C_2$, zgodnie z opisem
w Przykładzie~\ref{aa:prz:d8-nilpotent}.
Uwaga o notacji: stosujemy $D_{2n}$ na oznaczenie grupy
\emph{rzędu} $2n$; w części literatury tę samą grupę oznacza się $D_n$.
\end{przyklad}

\begin{przyklad}[Struktura $S_3$ i znaczenie normalności]
\label{aa:prz:s3-semidirect}
Niech $N=A_3=\langle r\rangle$,
$r=(1\,2\,3)$, oraz $H=\langle s\rangle$,
$s=(1\,2)$.
Podgrupa $N$ jest normalna, bo ma indeks dwa;
$|N|=3$, $|H|=2$, $N\cap H=\{e\}$.
Z Lematu~\ref{aa:lem:product-subgroups} otrzymujemy
$|NH|=3\cdot2=6=|S_3|$, więc $NH=S_3$.
Bezpośrednie obliczenie na punktach $1,2,3$ daje
$srs^{-1}=r^{-1}$.
Zatem $S_3\cong C_3\rtimes_{\mathrm{inv}}C_2$,
a nie $C_3\times C_2$: pierwsza grupa jest nieabelowa,
druga cykliczna rzędu $6$.
Nietrywialne działanie przez sprzężenie odróżnia tę grupę
od iloczynu prostego, mimo że w obu przypadkach rząd wynosi $6$.
\end{przyklad}

\begin{przyklad}[Grupa przekształceń afinicznych prostej nad ciałem]
\label{aa:prz:affine-group}
Niech $K$ będzie ciałem. Dla $a\in K$ i $u\in K^\times$
rozważmy przekształcenie $f_{a,u}:K\to K$,
$f_{a,u}(x)=ux+a$. Jest ono bijekcją, ponieważ
$u$ ma odwrotność; zachodzi
$f_{a,u}^{-1}(x)=u^{-1}(x-a)$.
Składając przekształcenia otrzymujemy
$f_{a,u}\circ f_{b,v}=f_{a+ub,uv}$.
Zatem grupa tych przekształceń jest iloczynem półprostym
$(K,+)\rtimes K^\times$, w którym
$u\in K^\times$ działa na $(K,+)$ przez mnożenie $x\mapsto ux$.

Dla $K=\mathbb F_7$ podgrupa
$\{f_{a,2^i}:a\in\mathbb F_7,\ i\in\{0,1,2\}\}$
ma $21$ elementów, ponieważ $2^3=1$ w $\mathbb F_7$.
Odzyskujemy więc dokładnie nieabelową grupę rzędu $21$
z Przykładu w Rozdziale~\ref{aa:ch:sylow}.
\end{przyklad}

\section{Przykład rozszerzenia, które się nie rozszczepia}
\begin{przyklad}[Rozszerzenie $C_2$ przez $C_2$ dające $C_4$]
\label{aa:prz:c4-nonsplit}
Rozważmy addytywny ciąg grup
\[
 0\longrightarrow C_2\xrightarrow{i}C_4
 \xrightarrow{p}C_2\longrightarrow0,
 \qquad i([a]_2)=[2a]_4,\quad p([b]_4)=[b]_2.
\]
Wszystkie odwzorowania są homomorfizmami:
$i$ jest iniektywne, $p$ surjektywne,
a $\ker p=\{[0]_4,[2]_4\}=\operatorname{im}i$.
Ciąg jest więc dokładny.

Załóżmy, że istnieje przekrój homomorficzny
$s:C_2\to C_4$ z $p\circ s=\operatorname{id}_{C_2}$.
Wówczas $p(s([1]_2))=[1]_2$, czyli $s([1]_2)$
musi być jedną z klas $[1]_4$ lub $[3]_4$.
Obie mają rząd $4$. Z drugiej strony $[1]_2$ ma rząd $2$,
a homomorfizm musi spełniać
$2s([1]_2)=s(2[1]_2)=s([0]_2)=[0]_4$,
co wyklucza element rzędu $4$. Otrzymujemy sprzeczność.
Ciąg nie rozszczepia się, choć wszystkie trzy grupy są abelowe.

Dla porównania $C_2\times C_2$ daje rozszczepiony ciąg
$0\to C_2\to C_2\times C_2\to C_2\to0$,
gdzie włączenie i rzut dotyczą odpowiednio pierwszego
i drugiego czynnika, a przekrojem jest $[b]\mapsto(0,[b])$.
Mamy więc dwa nieizomorficzne rozszerzenia tych samych grup.
\end{przyklad}

\begin{przyklad}[Nierozszczepiające się rozszerzenie z nieabelową grupą środkową]
\label{aa:prz:q8-nonsplit}
Otrzymamy inny przykład, w którym iloraz i jądro są abelowe,
ale cała grupa nie jest abelowa. Niech $E$ będzie macierzą
jednostkową stopnia $2$, $\mathrm i\in\mathbb C$ jednostką
urojoną, a
\[
 I=\begin{pmatrix}\mathrm i&0\\0&-\mathrm i\end{pmatrix},
 \qquad
 J=\begin{pmatrix}0&1\\-1&0\end{pmatrix}.
\]
Bezpośrednie mnożenie macierzy daje
$I^2=J^2=-E$ i $JI=-IJ$. Z tego wynika, że
$Q_8=\{\pm E,\pm I,\pm J,\pm IJ\}$
jest grupą z mnożeniem macierzy: iloczyn dowolnych dwóch
elementów, przy użyciu podanych relacji, sprowadza się
do jednego z ośmiu wypisanych; odwrotności wynikają
z $I^{-1}=-I$, $J^{-1}=-J$.
Osiem macierzy jest różnych, zatem $|Q_8|=8$.

Podgrupa $Z=\{E,-E\}$ jest centralna i izomorficzna z $C_2$.
Klasy $IZ,JZ$ w $Q_8/Z$ mają rząd $2$, ponieważ
$I^2=J^2=-E\in Z$. Komutują, gdyż $JI=-IJ$, a ich
klasy są różne i nie są neutralne: $I,J\notin Z$ i
$I\notin JZ$. Ponadto $(IZ)(JZ)=IJZ$ jest różne
od $Z$, $IZ$ i $JZ$, ponieważ osiem wypisanych macierzy
tworzy cztery rozłączne dwuelementowe warstwy modulo $Z$.
Te cztery klasy wyczerpują $Q_8/Z$, zatem iloraz
jest generowany przez dwa komutujące elementy rzędu $2$, skąd
$Q_8/Z\cong C_2\times C_2$. Powstaje krótki ciąg dokładny
\[
 1\longrightarrow Z\longrightarrow Q_8
 \longrightarrow Q_8/Z\longrightarrow1.
\]
Ten ciąg nie rozszczepia się. Istotnie, przekrój homomorficzny
$s:Q_8/Z\to Q_8$ byłby iniektywny, ponieważ
$p\circ s=\operatorname{id}_{Q_8/Z}$. Musiałby więc
umieszczać w $Q_8$ podgrupę izomorficzną z
$C_2\times C_2$, która ma trzy różne elementy rzędu $2$.
Tymczasem jedynym elementem rzędu $2$ w $Q_8$ jest $-E$:
dla każdego $X\in\{\pm I,\pm J,\pm IJ\}$ mamy
$X^2=-E$, a więc $X$ ma rząd $4$.
To sprzeczność. Rozszczepienie jest więc własnością
konkretnego ciągu dokładnego, nie tylko rzędów jego grup.
\end{przyklad}

\begin{uwaga}[Słowo \emph{rozszczepienie} jest ważniejsze niż sam rozmiar]
Każda surjekcja zbiorów skończonych ma przekrój
\emph{mnogościowy}, ale nie każdy taki przekrój jest homomorfizmem.
Dla iloczynu półprostego potrzebujemy przekroju zachowującego
mnożenie. Przykład~\ref{aa:prz:c4-nonsplit} pokazuje,
że tego dodatkowego wymagania nie wolno pomijać.
\end{uwaga}

\section{Grupy rzędu iloczynu dwóch różnych liczb pierwszych}
\begin{lemat}[Automorfizmy grupy cyklicznej rzędu pierwszego]
\label{aa:lem:aut-cq}
Jeżeli $q$ jest liczbą pierwszą, to
$\operatorname{Aut}(C_q)\cong(\mathbb Z/q\mathbb Z)^\times$.
W szczególności $|\operatorname{Aut}(C_q)|=q-1$.
Dla każdej liczby pierwszej $p$ dzielącej $q-1$
grupa $\operatorname{Aut}(C_q)$ ma dokładnie jedną
podgrupę rzędu $p$.
\end{lemat}
\begin{proof}
Grupa $C_q$ jest generowana przez $[1]_q$.
Homomorfizm $f:C_q\to C_q$ jest zatem wyznaczony przez
$f([1]_q)=[a]_q$ i ma postać $f([k]_q)=[ak]_q$.
Jest automorfizmem dokładnie wtedy, gdy $[a]_q\ne[0]_q$,
ponieważ dla $q$ pierwszego każda niezerowa klasa jest
odwracalna. Składanie automorfizmów odpowiada mnożeniu klas,
co daje izomorfizm i liczbę $q-1$.

Niech $p\mid(q-1)$. Twierdzenie Cauchy'ego zapewnia istnienie
elementu rzędu $p$, a więc podgrupy rzędu $p$.
Jej wszystkie elementy są rozwiązaniami równania $x^p=1$
w ciele $\mathbb F_q$. Wielomian $X^p-1$ ma nad ciałem
co najwyżej $p$ pierwiastków: dla dowolnego pierwiastka $a$
można wyłączyć czynnik $(X-a)$, a następnie stosować ten
argument indukcyjnie do wielomianu o mniejszym stopniu.
Podgrupa rzędu $p$ zawiera już $p$ różnych pierwiastków
tego wielomianu, więc jest całym zbiorem rozwiązań $x^p=1$.
Każda inna podgrupa rzędu $p$ musiałaby być zawarta w tym
samym zbiorze i miałaby $p$ elementów, więc byłaby równa pierwszej.
\end{proof}

\begin{twierdzenie}[Klasyfikacja grup rzędu $pq$]
\label{aa:tw:classify-pq}
Niech $p<q$ będą różnymi liczbami pierwszymi.
\begin{enumerate}
\item Jeśli $p\nmid(q-1)$, to z dokładnością do izomorfizmu
istnieje tylko jedna grupa rzędu $pq$, mianowicie $C_{pq}$.
\item Jeśli $p\mid(q-1)$, to z dokładnością do izomorfizmu
istnieją dokładnie dwie grupy rzędu $pq$:
cykliczna $C_{pq}$ i nieabelowa $C_q\rtimes C_p$,
w której generator $C_p$ działa na $C_q$
przez automorfizm rzędu $p$.
\end{enumerate}
\end{twierdzenie}
\begin{proof}
Niech $G$ będzie grupą rzędu $pq$.
Z trzeciego twierdzenia Sylowa $n_q\mid p$ i
$n_q\equiv1\pmod q$. Ponieważ $p<q$, jedyna możliwość
to $n_q=1$. Podgrupa $Q$ rzędu $q$ jest więc normalna.
Z pierwszego twierdzenia Sylowa istnieje podgrupa $P$ rzędu $p$.
Obie są cykliczne, a $P\cap Q=\{e\}$, bo rząd przekroju
dzieli jednocześnie $p$ i $q$.
Lemat~\ref{aa:lem:product-subgroups} daje
$|QP|=|Q|\,|P|=pq=|G|$, więc $QP=G$.
Z Twierdzenia~\ref{aa:tw:internal-semidirect}
$G\cong C_q\rtimes_\varphi C_p$ dla pewnego homomorfizmu
$\varphi:C_p\to\operatorname{Aut}(C_q)$.

Ponieważ $C_p$ jest generowana przez element rzędu $p$,
obraz $\varphi$ ma rząd $1$ albo $p$.
Pierwszy przypadek oznacza działanie trywialne,
więc $G\cong C_q\times C_p\cong C_{pq}$;
ostatni izomorfizm wynika z Przykładu o $C_6$,
a jego dowód działa bez zmiany dla różnych liczb pierwszych.
Jeśli obraz ma rząd $p$, to twierdzenie Lagrange'a zastosowane
do $\operatorname{Aut}(C_q)$ daje $p\mid q-1$.
Stąd przy $p\nmid(q-1)$ istnieje tylko grupa cykliczna.

Załóżmy $p\mid(q-1)$. Z Lematu~\ref{aa:lem:aut-cq}
istnieje automorfizm $\alpha$ rzędu $p$.
Wysłanie generatora $C_p$ na $\alpha$ określa nietrywialne
działanie, a więc grupę $C_q\rtimes_\alpha C_p$ rzędu $pq$.
Grupa ta nie jest abelowa: wybierzmy $x\in C_q$ taki,
że $\alpha(x)\ne x$ (istnieje, bo $\alpha$ nie jest identycznością).
Wtedy $(e_{C_q},1)(x,0)(e_{C_q},1)^{-1}=(\alpha(x),0)\ne(x,0)$.

Pokażemy, że wszystkie nietrywialne działania dają grupy
izomorficzne. Z Lematu~\ref{aa:lem:aut-cq} każdy ich obraz
jest tą samą podgrupą $U\le\operatorname{Aut}(C_q)$ rzędu $p$.
Dwa homomorfizmy $\varphi,\psi:C_p\to U$ o obrazach równych $U$
są izomorfizmami na $U$.
Zatem $\theta=\varphi^{-1}\circ\psi:C_p\to C_p$ jest
automorfizmem i $\psi(h)=\varphi(\theta(h))$.
Wzór $(x,h)\mapsto(x,\theta(h))$ jest wtedy izomorfizmem
$C_q\rtimes_\psi C_p\to C_q\rtimes_\varphi C_p$:
po obu stronach pierwszy czynnik iloczynu jest
$x\,\psi_h(y)=x\,\varphi_{\theta(h)}(y)$,
a drugi jest zgodny, bo $\theta$ jest homomorfizmem.
Izomorfizm jest bijektywny, gdyż $\theta$ jest automorfizmem.
To kończy klasyfikację.
\end{proof}

\begin{przyklad}[Wszystkie grupy rzędów $15$ i $21$]
Dla $|G|=15=3\cdot5$ mamy $3\nmid(5-1)=4$,
więc $G\cong C_{15}$: nie ma nieabelowej grupy rzędu $15$.
Dla $|G|=21=3\cdot7$ zachodzi $3\mid(7-1)$,
więc są dokładnie dwie grupy: $C_{21}$ i
$C_7\rtimes C_3$. W drugiej generator $C_3$ działa
na $C_7$ przez mnożenie przez $2$, ponieważ
$2^3\equiv1\pmod7$ i $2\not\equiv1\pmod7$.
To dokładnie grupa przekształceń $f_{a,i}(x)=2^i x+a$
z poprzedniego przykładu~\ref{aa:prz:affine-group}.
Dla pełnego rachunku możemy zapisać jej elementy jako
pary $(a,i)\in C_7\times C_3$ z mnożeniem
\[
 (a,i)(b,j)=(a+2^i b,\ i+j),
\]
gdzie pierwszą współrzędną redukujemy modulo $7$, drugą
modulo $3$. Wzór jest dobrze określony: zastąpienie
$i$ przez $i+3$ nie zmienia mnożnika $2^i$ modulo $7$,
bo $2^3\equiv1\pmod 7$. Elementami $r=(1,0)$
i $s=(0,1)$ generujemy całą grupę: potęgi $r$ dają
wszystkie $(a,0)$, a potęgi $s$ dają $(0,i)$.
Obliczamy $srs^{-1}=(2,0)=r^2$, a zatem
$sr\ne rs$. Przykład pokazuje zarówno sposób kodowania
rozszerzenia przez działanie, jak i źródło jego nieprzemienności.
\end{przyklad}

\section{Skończone grupy nilpotentne i iloczyny podgrup Sylowa}
\begin{lemat}[Warunek normalizatora w grupie nilpotentnej]
\label{aa:lem:nilpotent-normalizer}
Niech $G$ będzie grupą nilpotentną i $H<G$ właściwą
podgrupą. Wtedy $H$ jest właściwą podgrupą swojego
normalizatora: $H<N_G(H)$.
\end{lemat}
\begin{proof}
Ponieważ $G$ jest nilpotentna, istnieje $c$ z $Z_c(G)=G$
na mocy Lematu~\ref{aa:lem:upper-lower}.
Wybierzmy najmniejsze $i\ge1$, dla którego $Z_i(G)$
nie zawiera się w $H$. Taki indeks istnieje, ponieważ
$Z_0(G)=\{e\}\subseteq H$, natomiast $Z_c(G)=G\nsubseteq H$.
Wówczas $Z_{i-1}(G)\subseteq H$.
Wybierzmy $z\in Z_i(G)\setminus H$.
Z definicji $Z_i(G)$ dla każdego $h\in H$ zachodzi
$[z,h]\in Z_{i-1}(G)\subseteq H$.
Stąd $zhz^{-1}=[z,h]h\in H$.
Ten sam argument stosujemy do $z^{-1}\in Z_i(G)$:
$z^{-1}hz\in H$ dla wszystkich $h\in H$.
Zatem $zHz^{-1}\subseteq H$ i $z^{-1}Hz\subseteq H$.
Sprzężenie pierwszej inkluzji przez $z^{-1}$ daje
$H\subseteq z^{-1}Hz$, więc w rzeczywistości
$z^{-1}Hz=H$, a tym samym $zHz^{-1}=H$.
Wynika stąd $z\in N_G(H)\setminus H$.
Każdy element $H$ normalizuje $H$, zatem
$H\subsetneq N_G(H)$.
\end{proof}

\begin{twierdzenie}[Charakterystyka skończonych grup nilpotentnych]
\label{aa:tw:nilpotent-sylow-product}
Dla skończonej grupy $G$ następujące warunki są równoważne:
\begin{enumerate}
\item $G$ jest nilpotentna;
\item każda podgrupa Sylowa grupy $G$ jest normalna;
\item $G$ jest izomorficzna z iloczynem prostym swoich podgrup
Sylowa, po jednej dla każdego dzielnika pierwszego $|G|$.
\end{enumerate}
\end{twierdzenie}
\begin{proof}
Jeżeli $G=\{e\}$, wszystkie warunki są spełnione, a pusty
iloczyn prosty interpretujemy jako grupę trywialną.
W dalszym dowodzie zakładamy $|G|>1$.

\emph{(1) $\Rightarrow$ (2).}
Niech $P$ będzie podgrupą $p$-Sylowa i
$H=N_G(P)$ jej normalizatorem.
Mamy $P\trianglelefteq H$ z definicji normalizatora.
Ponieważ $P\le H\le G$, największa potęga $p$
dzieląca $|H|$ jest taka sama jak dla $|G|$.
Zatem $P$ jest podgrupą $p$-Sylowa także w $H$.
Przypuśćmy $H\ne G$.
Z poprzedniego lematu istnieje element
$x\in N_G(H)\setminus H$.
Ponieważ $x$ normalizuje $H$, podgrupa $xPx^{-1}$
zawiera się w $H$ i ma ten sam rząd co $P$;
jest więc podgrupą $p$-Sylowa $H$.
Ale $P$ jest normalna w $H$, a wszystkie podgrupy $p$-Sylowa
w $H$ są sprzężone (drugie twierdzenie Sylowa), więc
$P$ jest jedyną podgrupą $p$-Sylowa $H$.
Musi więc zachodzić $xPx^{-1}=P$, czyli $x\in N_G(P)=H$,
wbrew wyborowi $x$. Zatem $H=G$, więc $P\trianglelefteq G$.
Ponieważ $p$ i $P$ były dowolne, wszystkie podgrupy Sylowa
są normalne.

\emph{(2) $\Rightarrow$ (3).}
Zapiszmy $|G|=p_1^{a_1}\cdots p_r^{a_r}$ dla parami różnych
liczb pierwszych i wybierzmy podgrupy $P_i$ rzędu $p_i^{a_i}$.
Wszystkie są normalne z założenia.
Dla $i\ne j$ mamy $P_i\cap P_j=\{e\}$,
gdyż rząd przecięcia dzieli obie względnie pierwsze potęgi.
Dla $x\in P_i$ i $y\in P_j$ komutator $[x,y]$ należy
zarówno do $P_i$, jak i do $P_j$, ponieważ obie podgrupy
są normalne. Zatem $[x,y]\in P_i\cap P_j=\{e\}$,
czyli elementy $P_i$ i $P_j$ komutują.

Indukcyjnie zbudujemy podgrupę $K_j=P_1\cdots P_j$
o rzędzie $p_1^{a_1}\cdots p_j^{a_j}$, izomorficzną
z iloczynem prostym $P_1\times\cdots\times P_j$.
Dla $j=1$ jest to oczywiste.
Jeśli własność zachodzi dla $j$, to $K_j$ jest normalna,
gdyż iloczyn normalnych podgrup pozostaje normalny.
Rzędy $K_j$ oraz $P_{j+1}$ są względnie pierwsze,
więc $K_j\cap P_{j+1}=\{e\}$.
Twierdzenie~\ref{aa:tw:internal-direct} daje
$K_jP_{j+1}\cong K_j\times P_{j+1}$
i rząd równy iloczynowi rzędów.
Po $r$ krokach mamy podgrupę $K_r$ o rzędzie
$\prod_i p_i^{a_i}=|G|$, więc $K_r=G$.
To daje (3).

\emph{(3) $\Rightarrow$ (1).}
Każda $P_i$ jest skończoną $p_i$-grupą,
więc jest nilpotentna na mocy
Twierdzenia~\ref{aa:tw:pgroup-nilpotent}.
Centrum iloczynu prostego to iloczyn centrów czynników
(Twierdzenie~\ref{aa:tw:direct-product}).
Indukcja po $j$ i definicja górnego szeregu centralnego
dają
$Z_j(P_1\times\cdots\times P_r)
=Z_j(P_1)\times\cdots\times Z_j(P_r)$:
dla $j=0$ jest to równość grup trywialnych;
przy kroku indukcyjnym iloraz iloczynu przez $Z_j$
identyfikuje się z iloczynem ilorazów $P_i/Z_j(P_i)$,
a centrum tego iloczynu jest iloczynem centrów.
Dla każdego $i$ istnieje $c_i$ takie, że $Z_{c_i}(P_i)=P_i$.
Biorąc $c=\max_i c_i$, otrzymujemy $Z_c(G)=G$,
a stąd nilpotentność na mocy
Lematu~\ref{aa:lem:upper-lower}.
\end{proof}

\begin{przyklad}[Rozwiązalność bez nilpotentności jeszcze raz]
Grupa $S_3$ ma podgrupy Sylowa rzędów $3$ i $2$.
Podgrupa rzędu $3$ jest normalna, lecz trzy podgrupy
rzędu $2$ są wzajemnie sprzężone i żadna nie jest normalna.
Na mocy poprzedniego twierdzenia $S_3$ nie jest nilpotentna,
choć jest rozwiązalna. To daje drugi, niezależny od centrum,
sposób rozpoznania różnicy pomiędzy tymi dwiema własnościami.
\end{przyklad}

\section{Związek z iloczynem tensorowym grup abelowych}
\label{aa:sec:tensor-group-plan}
Dla \emph{grup abelowych} $A,B$ istnieje naturalna konstrukcja
$A\otimes_{\mathbb Z}B$, ale nie jest ona ani iloczynem prostym,
ani opisanym wyżej iloczynem półprostym.
Podstawą jest obserwacja, że każda grupa abelowa jest
kanonicznie modułem nad $\mathbb Z$: definiujemy
$n\cdot a$ jako $n$-krotną sumę $a$ dla $n>0$,
$0\cdot a=0$ oraz $(-n)\cdot a=-(n\cdot a)$.
Wówczas iloczyn tensorowy grup abelowych jest po prostu
iloczynem tensorowym $\mathbb Z$-modułów.
Jego własność uniwersalna opisuje odwzorowania
addytywne w każdej zmiennej osobno.

Dla dodatnich liczb całkowitych $m,n$ mamy
$C_m\otimes_{\mathbb Z}C_n\cong C_{\gcd(m,n)}$.
Przykładowo $C_2\otimes_{\mathbb Z}C_3=0$, mimo że
$C_2\times C_3\cong C_6$. Natomiast
$C_4\otimes_{\mathbb Z}C_6\cong C_2$, choć
$C_4\times C_6$ ma $24$ elementy.

W tym skrypcie konstrukcję przez generatory i relacje oraz
dowód powyższego wzoru przedstawiamy już w
Rozdziale~\ref{aa:ch:grupy-abelowe},
w szczególności w
Twierdzeniu~\ref{aa:tw:abelian-tensor-cyclic}.
Ogólną konstrukcję dla modułów nad pierścieniem przemiennym
oraz jej własności rozwijamy dalej w
Rozdziale~\ref{aa:ch:module-tensor-product}.
Odsyłacze te prowadzą do rzeczywistych rozdziałów tego
samego skryptu; nie są zależnością od obcego modułu ani
od zewnętrznego pliku pomocniczego.

% AA-07 | Algebra abstrakcyjna | Iteracja IV
% Reinterpretacja i integracja dowodu z BT-01; bez cwiczen.
\chapter[Grupy wolne i twierdzenie Nielsena--Schreiera]{Grupy wolne, prezentacje i twierdzenie Nielsena--Schreiera}
\label{aa:ch:grupy-wolne}

\providecommand{\atlasref}[2]{#2}


\section{Dlaczego grupy wolne są obiektem uniwersalnym?}

\noindent\emph{Powiązania tematyczne w Bibliotece matematycznej:} \atlasref{KAT-14::def:adjunction}{KAT-14: sprzężenia funktorów i konstrukcje wolne}.

W grupie cyklicznej jeden generator może spełniać relację $a^m=e$.
W grupie abelowej dwa generatory zawsze komutują. Pytamy, co otrzymamy,
jeżeli wskażemy pewien zbiór generatorów, ale \emph{nie narzucimy im żadnych
relacji poza aksjomatami grupy}. Odpowiedź prowadzi do grup wolnych.
Są one przydatne do opisu grup za pomocą generatorów i relacji, grup
podstawowych przestrzeni topologicznych oraz konstrukcji paradoksalnych.
Ostatnie zastosowanie wykorzystamy przy dowodzie twierdzenia
Banacha--Tarskiego.

\begin{definicja}[Własność uniwersalna grupy wolnej]
\label{aa:def:free-universal}
\label{def:grupa-wolna}
Niech $S$ będzie zbiorem. Para $(F,\iota)$, gdzie $F$ jest grupą,
a $\iota:S\to F$ jest funkcją, nazywa się \emph{grupą wolną na $S$},
gdy dla każdej grupy $G$ i każdej funkcji $f:S\to G$ istnieje
\emph{dokładnie jeden} homomorfizm $\widetilde f:F\to G$ taki, że
$\widetilde f\circ\iota=f$. Grupę taką oznaczamy przez $F(S)$.
\end{definicja}

Strzałka przerywana oznacza jedyny morfizm (odwzorowanie) wyznaczony przez własność uniwersalną.
\begin{center}
\begin{tikzpicture}[>=stealth]
 \node (a) at (0,2.0) {$S$};
 \node (b) at (4.4,2.0) {$F(S)$};
 \node (c) at (2.2,0) {$G$};
 \draw[->] (a) -- node[above] {$\iota$} (b);
 \draw[->] (a) -- node[below left] {$f$} (c);
 \draw[->,dashed] (b) -- node[below right] {$\widetilde f$} (c);
\end{tikzpicture}
\end{center}


Warunek dotyczy dowolnej funkcji $f$: obrazy generatorów mogą mieć
jakiekolwiek rzędy i nie muszą generować całego $G$. Jeśli $S$ ma
dwa elementy $a,b$, to w $F(S)$ wyrazy $ab$ i $ba$ będą różne.
W grupie $G$ ich obrazy mogą się natomiast pokrywać; oznacza to
jedynie, że $\widetilde f$ ma nietrywialne jądro.

\begin{lemat}[Jednoznaczność z dokładnością do izomorfizmu]
\label{aa:lem:free-unique}
Jeśli $(F,\iota)$ i $(F',\iota')$ spełniają powyższą własność
uniwersalną dla tego samego zbioru $S$, to istnieje dokładnie jeden
izomorfizm $u:F\to F'$ spełniający $u\circ\iota=\iota'$.
\end{lemat}
\begin{proof}
Własność uniwersalna grupy $F$ zastosowana do $\iota':S\to F'$ daje
homomorfizm $u:F\to F'$, a własność uniwersalna grupy $F'$ zastosowana
do $\iota:S\to F$ daje $v:F'\to F$.
Złożenie $vu:F\to F$ spełnia $(vu)\iota=\iota$,
tak samo jak identyczność $\operatorname{id}_F$. Jednoznaczność
z definicji daje $vu=\operatorname{id}_F$.
Analogicznie $uv=\operatorname{id}_{F'}$.
Zatem $u$ jest izomorfizmem o odwrotności $v$;
jego jednoznaczność była już częścią konstrukcji.
\end{proof}

\section{Słowa zredukowane: budowa grupy bez ukrytych założeń}
Niech $S^{-1}$ będzie rozłączną kopią $S$; dla $s\in S$ zapisujemy
odpowiadający mu symbol przez $s^{-1}$. Formalny alfabet to
$\mathcal A=S\sqcup S^{-1}$. Przyporządkowanie liter odwrotnych
jest inwolucją: $(s^{-1})^{-1}=s$; na tym etapie nie jest to jeszcze
operacja w grupie.

\begin{definicja}[Słowo, redukcja]
\label{aa:def:free-word}
\emph{Słowo} to skończony ciąg liter $w=s_1\cdots s_n$ z $\mathcal A$;
dopuszczamy słowo puste $e$. \emph{Elementarna redukcja} usuwa
sąsiednią parę $ss^{-1}$ albo $s^{-1}s$.
Słowo jest \emph{zredukowane}, gdy nie zawiera takiej pary.
\end{definicja}

\begin{przyklad}[Co wolno skreślić, a czego nie]
\label{aa:prz:free-reduction}
Dla $S=\{a,b\}$ słowo $a a^{-1} b b^{-1} b a$ redukuje się
do $ba$. Słowo $aba^{-1}b^{-1}$ jest już zredukowane i nie jest
słowem pustym. Nie wolno przestawiać liter, aby usunąć $aa^{-1}$:
możliwość przestawiania liter oznaczałaby dodatkową relację
przemienności, której w grupie wolnej nie wprowadzamy.
\end{przyklad}

\begin{definicja}[Algorytm stosu]
\label{aa:def:free-stack}
Czytamy słowo od lewej do prawej, przechowując jego aktualnie
zredukowany prefiks $r$. Początkowo $r=e$.
Po przeczytaniu litery $s$ usuwamy ostatnią literę $r$, jeżeli jest
ona równa $s^{-1}$; w przeciwnym razie dopisujemy $s$ na końcu.
Wynik oznaczamy $\operatorname{red}(w)$.
\end{definicja}

\begin{lemat}[Wstawienie pary odwrotnej nie zmienia stosu]
\label{aa:lem:free-stack-pair}
\label{lem:stos}
Niech przed przeczytaniem dwóch kolejnych liter $s,s^{-1}$
na stosie leży słowo zredukowane $r$. Po przeczytaniu obu liter
na stosie ponownie leży $r$.
\end{lemat}
\begin{proof}
Jeśli $r$ nie kończy się literą $s^{-1}$, pierwsze czytanie
dopisuje $s$, a drugie natychmiast kasuje dopisaną literę.
Jeśli $r=us^{-1}$, pierwsze czytanie usuwa końcową literę
$s^{-1}$ i pozostawia $u$. Ponieważ $r$ było zredukowane,
$u$ nie kończy się literą $s$; wobec tego drugie czytanie
dopisuje $s^{-1}$, odtwarzając $r=us^{-1}$.
Przypadek pustego $u$ jest taki sam.
\end{proof}

\begin{twierdzenie}[Jednoznaczność redukcji]
\label{aa:tw:free-reduction}
\label{tw:redukcja}
Każde słowo $w$ można doprowadzić przez elementarne redukcje
do dokładnie jednego słowa zredukowanego. Jest nim
$\operatorname{red}(w)$ otrzymane algorytmem stosu.
\end{twierdzenie}
\begin{proof}
Każda elementarna redukcja zmniejsza długość o dwa, więc każdy
ciąg redukcji musi się zakończyć. Algorytm stosu również można
zrealizować za pomocą elementarnych redukcji: po redukcji prefiksu
dopisujemy następną literę i ewentualnie kasujemy jedną parę na końcu.
Indukcja po długości przeczytanego prefiksu pokazuje, że otrzymujemy
osiągalne słowo zredukowane.

Pozostaje wykazać niezależność od wybranej kolejności skreśleń.
Jeśli $w=pss^{-1}q$, to algorytm stosu dla $w$ i dla $pq$
daje ten sam wynik: po prefiksie $p$ stan jest jednakowy,
a dwie kolejne litery $s,s^{-1}$ go nie zmieniają na mocy
Lematu~\ref{aa:lem:free-stack-pair}; sufiks $q$ jest następnie
taki sam. Zatem każda elementarna redukcja zachowuje
$\operatorname{red}(w)$. Jeśli dowolny ciąg redukcji kończy się
słowem zredukowanym $z$, to
$\operatorname{red}(w)=\operatorname{red}(z)=z$.
Wszystkie możliwe wyniki końcowe są więc identyczne.
\end{proof}

\begin{wniosek}[Redukcja nie zależy od wcześniejszych skreśleń]
\label{aa:wn:free-concat}
\label{wn:redukcja-iloczyn}
Dla dowolnych słów $u,v$ zachodzą równości
$\operatorname{red}(\operatorname{red}(u)v)=\operatorname{red}(uv)$
oraz $\operatorname{red}(u\operatorname{red}(v))=\operatorname{red}(uv)$.
\end{wniosek}
\begin{proof}
W słowie $uv$ możemy najpierw wykonać redukcje
w całym fragmencie $u$ (albo $v$), a dopiero później w pozostałej
części. Jednoznaczność końcowego wyniku dowodzi równości.
\end{proof}

\begin{twierdzenie}[Istnienie grupy wolnej]
\label{aa:tw:free-exists}
\label{tw:uniwersalnosc}
\label{st:grupa-slow}
Zbiór wszystkich słów zredukowanych w alfabecie
$S\sqcup S^{-1}$ jest grupą przy mnożeniu
$u\cdot v=\operatorname{red}(uv)$.
Ta grupa wraz z funkcją $\iota(s)=s$ spełnia własność
uniwersalną z Definicji~\ref{aa:def:free-universal}.
\end{twierdzenie}
\begin{proof}
Wynik redukcji jest słowem zredukowanym, więc działanie
jest wewnętrzne. Na mocy
Wniosku~\ref{aa:wn:free-concat} mamy
$(u\cdot v)\cdot w=\operatorname{red}(uvw)
=u\cdot(v\cdot w)$, czyli łączność.
Słowo puste jest jedynką. Dla $u=s_1\cdots s_n$ słowo
$u^{-1}=s_n^{-1}\cdots s_1^{-1}$ jest zredukowane;
redukcja $uu^{-1}$ i $u^{-1}u$ od środka prowadzi do $e$.
Są to więc odwrotności obustronne.

Niech teraz $f:S\to G$ będzie dowolna funkcja do grupy $G$.
Dla zredukowanego słowa $s_1\cdots s_n$ definiujemy
$\widetilde f(s_1\cdots s_n)$ jako iloczyn obrazów liter,
przy czym obrazem formalnego $s^{-1}$ jest $f(s)^{-1}$.
Elementarna redukcja nie zmienia tego iloczynu, bo usuwa
czynniki wzajemnie odwrotne. Wobec tego
$\widetilde f(u\cdot v)=\widetilde f(u)\widetilde f(v)$.
Jest to homomorfizm przedłużający $f$.
Każdy taki homomorfizm musi działać w ten sposób na słowach,
ponieważ zachowuje mnożenie i odwrotności; zatem jest jedyny.
\end{proof}

\begin{przyklad}[Grupa wolna rangi zero, jeden i dwa]
\label{aa:prz:free-ranks}
Dla $S=\varnothing$ jedynym słowem jest słowo puste, więc
$F(\varnothing)=\{e\}$. Dla $S=\{a\}$ zredukowane słowa mają
postać $a^n$, $n\in\mathbb Z$, a mnożenie odpowiada dodawaniu
wykładników, więc $F(\{a\})\cong\mathbb Z$.
Dla $S=\{a,b\}$ słowa $ab$ i $ba$ są różne, ponieważ są
różnymi słowami zredukowanymi. Zatem $F(\{a,b\})$ nie jest
abelowa. Słowo $aba^{-1}b^{-1}\ne e$ jest jej nietrywialnym
komutatorem.
\end{przyklad}

\begin{stwierdzenie}[Liczba słów o danej długości]
\label{aa:st:free-count}
\label{st:przeliczalnosc-f2}
Dla liczb całkowitych $r\ge2$ i $n\ge1$ grupa wolna
$F_r=F(\{s_1,\ldots,s_r\})$ ma dokładnie
$2r(2r-1)^{n-1}$ słów zredukowanych długości $n$.
Jest przeliczalnie nieskończona.
\end{stwierdzenie}
\begin{proof}
Pierwszą literę można wybrać na $2r$ sposobów;
każda następna ma $2r-1$ możliwości, bo zabroniona jest
jedynie odwrotność poprzedniej. Dla każdego $n$ liczba jest
skończona, a zliczalna suma zbiorów skończonych jest przeliczalna.
Nieskończoność wynika z różności słów $s_1^n$ dla $n\ge0$.
\end{proof}


\begin{przyklad}[Dwa rodzaje zależności między generatorami]
\label{aa:prz:free-nonfree-pairs}
Grupa abelowa $\mathbb Z^2$ jest generowana przez
$(1,0)$ i $(0,1)$, ale nie jest grupą wolną na tych dwóch
generatorach: komutator dwóch generatorów jest w niej elementem neutralnym,
podczas gdy odpowiadające mu słowo $aba^{-1}b^{-1}$
w $F(a,b)$ jest zredukowane i niepuste.
Z drugiej strony sama nieprzemienność również nie wystarcza,
aby para elementów była parą wolnych generatorów.
W $S_3$ transpozycje $(1\,2)$ i $(2\,3)$ nie komutują,
ale spełniają relację $(1\,2)^2=e$.
Homomorfizm $F(a,b)\to S_3$ wysyłający $a$ na $(1\,2)$,
a $b$ na $(2\,3)$ ma więc w jądrze nietrywialne słowo $a^2$.
W dowodzie paradoksu Banacha--Tarskiego trzeba będzie
wykluczyć nie tylko przemienność lub relacje torsyjne,
lecz każdą możliwą nietrywialną relację.
\end{przyklad}

\section{Generatory, relacje i ilorazy grup wolnych}
\begin{definicja}[Prezentacja grupy]
\label{aa:def:presentation}
Dla zbioru generatorów $S$ i zbioru słów $R\subseteq F(S)$
niech $\langle\!\langle R\rangle\!\rangle$ oznacza
\emph{domknięcie normalne} $R$, czyli najmniejszą
podgrupę normalną $F(S)$ zawierającą $R$.
Definiujemy
\[
\langle S\mid R\rangle
 =F(S)/\langle\!\langle R\rangle\!\rangle.
\]
Przy $R=\varnothing$ otrzymujemy $F(S)$.
\end{definicja}

\begin{lemat}[Co oznacza domknięcie normalne]
\label{aa:lem:normal-closure}
Dla grupy $G$ i podzbioru $R\subseteq G$ domknięcie normalne
$\langle\!\langle R\rangle\!\rangle$ jest podgrupą generowaną
przez wszystkie elementy $grg^{-1}$, gdzie $g\in G$, $r\in R$.
\end{lemat}
\begin{proof}
Niech $N$ będzie podgrupą generowaną przez wskazane sprzężenia.
Jest normalna, bo sprzężenie $grg^{-1}$ przez $h\in G$
ma postać $(hg)r(hg)^{-1}$, a sprzężenie iloczynu jest iloczynem
sprzężeń. Zawiera $R$, biorąc $g=e$.
Każda inna podgrupa normalna zawierająca $R$ musi zawierać
wszystkie jego sprzężenia, a więc zawiera $N$. Jest więc
$N$ najmniejszą taką podgrupą.
\end{proof}

\begin{twierdzenie}[Własność uniwersalna prezentacji]
\label{aa:tw:presentation-universal}
Niech $S$ będzie zbiorem, $R\subseteq F(S)$ zbiorem słów,
$G$ grupą, a $f:S\to G$ dowolną funkcją.
Istnieje homomorfizm $\langle S\mid R\rangle\to G$
przypisujący klasie generatora $s$ element $f(s)$ wtedy
i tylko wtedy, gdy homomorfizm $\widetilde f:F(S)\to G$
indukowany przez $f$ spełnia $\widetilde f(r)=e_G$
dla każdego $r\in R$. Takie przedłużenie jest jedyne.
\end{twierdzenie}
\begin{proof}
Jeśli istnieje zadany homomorfizm z ilorazu, każda relacja
$r\in R$ jest w nim równa jedynce; w szczególności
$\widetilde f(r)=e_G$.
Odwrotnie, jeśli $\widetilde f$ zeruje wszystkie relacje,
to normalność $\ker\widetilde f$ oznacza, że zawiera ich
domknięcie normalne. Definiujemy więc na ilorazie odwzorowanie
$w\langle\!\langle R\rangle\!\rangle\mapsto\widetilde f(w)$.
Jest ono dobrze określone: równość dwóch warstw oznacza,
że różnica ich reprezentantów (w sensie iloczynu z odwrotnością)
należy do domknięcia normalnego, a więc do jądra $\widetilde f$.
Zachowywanie mnożenia wynika z tego, że $\widetilde f$ jest
homomorfizmem.
Jednoznaczność bierze się z faktu, że klasy generatorów
generują cały iloraz.
\end{proof}

\begin{przyklad}[Najważniejsze prezentacje elementarne]
\label{aa:prz:presentations-basic}
Dla każdej liczby całkowitej $n\ge1$ grupa cykliczna $C_n$
ma prezentację $\langle a\mid a^n\rangle$, ponieważ słowa na jednym generatorze
są potęgami $a$, a relacja $a^n=e$ utożsamia wykładniki
różniące się o wielokrotność $n$.
Grupa $\mathbb Z^2$ ma prezentację
$\langle a,b\mid aba^{-1}b^{-1}\rangle$:
relacja wymusza $ab=ba$, więc każde słowo daje $a^m b^n$;
jednoznaczność pary $(m,n)$ wynika z homomorfizmu do
$\mathbb Z^2$ przypisującego $a\mapsto(1,0)$ i
$b\mapsto(0,1)$.
Grupa diedralna $D_{2n}$ ma prezentację
$\langle r,s\mid r^n,s^2,srs^{-1}r\rangle$ dla $n\ge3$.
Z relacji $sr=r^{-1}s$ każde słowo można sprowadzić do jednej
z postaci $r^i$ albo $r^is$, $0\le i<n$.
Daje to co najwyżej $2n$ elementów; rzeczywista grupa
symetrii foremnego $n$-kąta ma dokładnie $2n$ elementów,
i spełnia relację, więc homomorfizm z prezentacji na tę grupę
jest bijekcją.
\end{przyklad}

\begin{twierdzenie}[Abelianizacja grupy wolnej]
\label{aa:tw:free-abelianization}
Dla dowolnego zbioru $S$ grupa abelowa $F(S)_{\mathrm{ab}}=
F(S)/[F(S),F(S)]$ jest naturalnie izomorficzna z grupą
$\mathbb Z^{(S)}$ skończonych kombinacji formalnych
$\sum_{s\in S}n_s e_s$, $n_s\in\mathbb Z$.
\end{twierdzenie}
\begin{proof}
Niech $u:F(S)\to\mathbb Z^{(S)}$ wysyła generator $s$ na
$e_s$. Istnieje i jest jedyny na mocy uniwersalności grupy
wolnej; jest surjektywny, bo obrazy generatorów generują
$\mathbb Z^{(S)}$. Ponieważ grupa docelowa jest abelowa,
$[F(S),F(S)]\subseteq\ker u$, więc $u$ indukuje surjekcję
$\bar u:F(S)_{\mathrm{ab}}\to\mathbb Z^{(S)}$.
Odwrotnie, przyporządkowanie $e_s\mapsto[s]$ można
rozszerzyć do homomorfizmu z $\mathbb Z^{(S)}$ do
$F(S)_{\mathrm{ab}}$, ponieważ ten ostatni iloraz jest abelowy.
Oba złożenia zachowują generatory $e_s$ i $[s]$, a więc
są identycznościami. Zatem $\bar u$ jest izomorfizmem.
\end{proof}


\begin{stwierdzenie}[Niezmienniczość skończonej rangi grupy wolnej]
\label{aa:st:free-rank-invariant}
Jeśli $F_r\cong F_s$ dla liczb całkowitych $r,s\ge0$,
to $r=s$.
\end{stwierdzenie}
\begin{proof}
Niech $C_2$ będzie dwuelementową grupą cykliczną.
Każdy homomorfizm $F_r\to C_2$ jest jednoznacznie wyznaczony
przez niezależny wybór obrazów $r$ wolnych generatorów,
a każdy taki wybór daje homomorfizm na mocy własności
uniwersalnej. Istnieje więc dokładnie $2^r$ takich homomorfizmów.
Izomorfizm $F_r\cong F_s$ przez złożenie w jedną stronę
ustanawia bijekcję pomiędzy zbiorami homomorfizmów
$F_r\to C_2$ oraz $F_s\to C_2$.
Zatem $2^r=2^s$, skąd $r=s$.
\end{proof}

\section{Graf warstw: geometria ukryta w podgrupie grupy wolnej}
Niech $F=F(S)$, $H\le F$ będzie dowolną podgrupą.
Zbudujemy graf, który pozwala rozpoznać elementy $H$ jako
zamknięte ścieżki oraz znaleźć dla $H$ nowe wolne generatory.
Nie zakładamy, że $H$ ma skończony indeks.

\begin{definicja}[Graf Schreiera]
\label{aa:def:schreier-graph}
\emph{Grafem Schreiera} podgrupy $H\le F(S)$ nazywamy graf
skierowany $\Gamma(H\backslash F,S)$, którego wierzchołki
to prawe warstwy $Hg$, $g\in F$. Dla każdego $v=Hg$ i $s\in S$
umieszczamy krawędź $e(v,s)$ od $v$ do $v s=Hgs$
o etykiecie $s$; dopuszczamy pętle i krawędzie wielokrotne.
Krawędź wolno przejść także w przeciwnym kierunku, z etykietą $s^{-1}$.
Wierzchołek $H$ nazywamy \emph{bazowym}.
Przy rozpatrywaniu drzew rozpinających zapominamy o orientacji,
ale zachowujemy rozróżnienie poszczególnych krawędzi:
przeciwne skierowanie dwóch różnych krawędzi z etykietą $s$
nie utożsamia ich ze sobą. Odwrotne przejście danej krawędzi
nie jest natomiast nową krawędzią grafu.
\end{definicja}

\begin{lemat}[Słowa, drogi i pętle]
\label{aa:lem:schreier-path}
Dla dowolnego słowa $w\in F(S)$ istnieje dokładnie jedna droga
w grafie Schreiera wychodząca z każdego wskazanego wierzchołka
i mająca kolejno etykiety liter zredukowanego słowa $w$.
Droga rozpoczęta w $Hg$ kończy się w $Hgw$.
W szczególności droga z bazowego $H$ jest zamknięta wtedy
i tylko wtedy, gdy $w\in H$.
Ponadto redukcja dwóch kolejnych, przeciwnie skierowanych przejść
tej samej krawędzi odpowiada skreśleniu sąsiednich liter odwrotnych.
\end{lemat}
\begin{proof}
Z każdego wierzchołka wychodzi dokładnie jedna krawędź
z każdą dodatnią etykietą $s$, a dokładnie jedna odwrotność
krawędzi kończącej się w danym wierzchołku ma etykietę $s^{-1}$:
jest to odwrotność $e(Hgs^{-1},s)$.
Indukcja po kolejnych literach pokazuje, że od wierzchołka
$Hg$ przechodzimy do $Hgs_1$, potem do $Hgs_1s_2$ itd.,
ostatecznie do $Hgw$. Jeśli $g=e$, koniec jest równy $H$
dokładnie wtedy, gdy $Hw=H$, czyli $w\in H$.
Przejście krawędzi i natychmiastowy powrót to dokładnie
przejście litery i jej odwrotności.
\end{proof}

Dla dowolnego zbioru generatorów $S$, również nieskończonego,
pojęcie grafu uwzględnia krawędzie wielokrotne i pętle.
Ważne jest to, że warstwy $Hg$ są jedynie elementami zbioru
$H\backslash F$; bez normalności $H$ nie ma on struktury grupy
ilorazowej pochodzącej z $F$.

Graf Schreiera jest spójny: każdy wierzchołek $Hg$ można
osiągnąć z $H$ drogą oznaczoną dowolnym słowem reprezentującym $g$.
Uzasadnimy teraz, dlaczego można wybrać w nim drzewo rozpinające.
W przypadku skończonego grafu wystarczy kolejno dodawać krawędzie,
które docierają do jeszcze nieodwiedzonego wierzchołka.
Dla dowolnego $S$ można dobrze uporządkować alfabet
$S\sqcup S^{-1}$ (korzystając z aksjomatu wyboru).
Słowa porządkujemy najpierw według długości,
a słowa tej samej długości leksykograficznie.
Jest to dobre uporządkowanie słów skończonych. Dla każdej
warstwy wybieramy pierwsze w tym porządku słowo, które
do niej prowadzi. Każdy prefiks wybranego słowa jest wybranym
słowem dla swojego wierzchołka: zastąpienie prefiksu słowem
krótszym albo leksykograficznie wcześniejszym dałoby
wcześniejsze słowo prowadzące do tej samej warstwy.
Wybrane słowa są zredukowane, gdyż skreślenie sąsiednich
liter odwrotnych prowadziłoby do krótszego słowa.
Łączymy każdy niebazowy wierzchołek z jego poprzednikiem
na wybranej drodze. Długość wybranego słowa maleje o jeden
przy przejściu do poprzednika, więc otrzymane krawędzie
nie tworzą cyklu i prowadzą każdy wierzchołek do $H$.
Tworzą zatem drzewo rozpinające $T$.

\begin{lemat}[Pętle związane z krawędziami spoza drzewa]
\label{aa:lem:tree-paths}
Niech $\Gamma$ będzie spojnym grafem z wyróżnionym
wierzchołkiem $v_0$ i drzewem rozpinającym $T$.
Dla każdego wierzchołka $v$ niech $t_v$ będzie jedyną
zredukowaną drogą w $T$ od $v_0$ do $v$.
Dla każdej zorientowanej krawędzi $e:v\to w$ nieleżącej
w $T$ definiujemy pętlę $b_e=t_v e t_w^{-1}$.
Wybieramy dokładnie jedną orientację każdej krawędzi
spoza $T$. Wtedy każda zamknięta droga od $v_0$ jest,
z dokładnością do usuwania natychmiastowych powrotów,
iloczynem pętli $b_e$ i ich odwrotności.
\end{lemat}
\begin{proof}
Przejdźmy po kolei krawędzie zamkniętej drogi
$v_0=v_0',v_1',\dots,v_m'=v_0$.
Dla krawędzi $e_i:v_{i-1}'\to v_i'$ wstawmy
między kolejne kroki drogi ścieżki $t_{v_i'}^{-1}t_{v_i'}$.
Otrzymamy iloczyn pętli
$(t_{v_0'}e_1t_{v_1'}^{-1})\cdots
 (t_{v_{m-1}'}e_mt_{v_m'}^{-1})$.
Jeśli $e_i$ należy do drzewa, pętlą
$t_{v_{i-1}'}e_it_{v_i'}^{-1}$ redukuje się do drogi pustej,
ponieważ w drzewie pomiędzy dwiema końcówkami istnieje
jedyna zredukowana droga.
Jeśli $e_i$ nie należy do drzewa, dostajemy
$b_{e_i}$ albo $b_{e_i}^{-1}$, w zależności od kierunku przejścia.
Po usunięciu pustych pętli pozostaje zadany iloczyn.
\end{proof}

\begin{lemat}[Brak dodatkowych relacji między pętlami]
\label{aa:lem:tree-free}
Niech $\Gamma$ będzie spójnym grafem z wierzchołkiem
bazowym $v_0$, $T$ jego drzewem rozpinającym, a $b_e$
pętlami przypisanymi po jednej do każdej krawędzi spoza $T$
przy wybranych orientacjach, jak w Lemacie~\ref{aa:lem:tree-paths}.
Wtedy żadne niepuste, zredukowane słowo w symbolach
$b_e^{\pm1}$ nie jest równoważne pustej drodze
po usuwaniu natychmiastowych powrotów.
\end{lemat}
\begin{proof}
Rozważmy zredukowane słowo
$b_{e_1}^{\varepsilon_1}\cdots b_{e_k}^{\varepsilon_k}$, $k\ge1$,
$\varepsilon_i\in\{1,-1\}$.
Po zastąpieniu każdego czynnika jego drogą
$t_v e_i^{\varepsilon_i}t_w^{-1}$ oraz skreśleniu
sąsiednich krawędzi drzewa pozostają \emph{kolejno}
$k$ krawędzi spoza drzewa, zorientowanych jak
$e_i^{\varepsilon_i}$; żadna z nich nie może zniknąć
w wyniku redukcji wewnątrz drogi leżącej w drzewie.
Jedyny sposób usunięcia pierwszej krawędzi spoza drzewa
w dalszej redukcji wymagałby, aby kolejna krawędź
spoza drzewa była dokładnie tą samą krawędzią przebieganą
w przeciwnym kierunku oraz aby łącząca je droga w $T$
zredukowała się do drogi pustej.
Ponieważ drzewo ma jednoznaczne zredukowane drogi,
zdarza się to tylko wtedy, gdy sąsiednie symbole
$b_{e_i}^{\varepsilon_i}$ i $b_{e_{i+1}}^{\varepsilon_{i+1}}$
są wzajemnie odwrotne. Przeczy to zredukowaniu
rozpatrywanego słowa. Zatem co najmniej jedna
krawędź spoza drzewa pozostaje; droga nie może
zredukować się do pustej.
\end{proof}

\begin{twierdzenie}[Nielsena--Schreiera]
\label{aa:tw:nielsen-schreier}
Niech $S$ będzie dowolnym zbiorem, $F(S)$ grupą wolną na $S$,
a $H\le F(S)$ dowolną podgrupą. Wtedy $H$ jest grupą wolną.
Dokładniej, dla wybranego drzewa rozpinającego $T$
w grafie Schreiera $\Gamma(H\backslash F(S),S)$ elementy
$h_e\in H$ odpowiadające słowom etykiet pętli $t_v e t_w^{-1}$,
po jednym dla każdej krawędzi poza $T$ i po wybraniu jej orientacji,
tworzą wolną bazę $H$.
\end{twierdzenie}
\begin{proof}
Dla krawędzi $e:v\to w$ o etykiecie $s$ zapiszmy
przez $u_v\in F(S)$ element reprezentowany przez słowo
etykiet drogi $t_v$. Ponieważ $t_v$ kończy się w warstwie $v$,
mamy $v=Hu_v$; analogicznie $w=Hu_w$.
Relacja $w=vs$ daje $Hu_v s=Hu_w$, a więc
$h_e=u_v s u_w^{-1}\in H$. Dla krawędzi drzewa
odpowiednia droga zamknięta redukuje się do pustej,
więc $h_e=e$. Dla krawędzi poza drzewem zachowujemy $h_e$.

Niech $h\in H$ będzie reprezentowany przez zredukowane
słowo $w\in F(S)$. Z Lematu~\ref{aa:lem:schreier-path}
ścieżka o etykiecie $w$ jest pętlą w bazowym wierzchołku.
Lemat~\ref{aa:lem:tree-paths} rozkłada ją na
pętle $b_e^{\pm1}$, zatem $h$ jest iloczynem elementów
$h_e^{\pm1}$. Elementy $h_e$ generują $H$.

Wolna grupa na formalnych symbolach krawędzi spoza $T$
ma więc naturalny surjektywny homomorfizm na $H$.
Gdyby jego jądro było nietrywialne, istniałoby niepuste
zredukowane słowo w symbolach $h_e^{\pm1}$ reprezentujące
jedynkę $F(S)$. Rozwinięcie tego słowa daje drogę zamkniętą
w grafie Schreiera. Słowo etykiet jest równe jedynce w
$F(S)$, więc na mocy jednoznaczności redukcji
(Twierdzenie~\ref{aa:tw:free-reduction}) można usuwać
kolejne sąsiednie pary liter odwrotnych aż do słowa pustego.
W grafie Schreiera każde takie skreślenie odpowiada
przejściu daną krawędzią i natychmiastowemu powrotowi
tą samą krawędzią: dla ustalonego wierzchołka i litery
istnieje dokładnie jedno możliwe przejście.
Zatem również droga zredukowałaby się przez usuwanie
powrotów do drogi pustej.
Ale Lemat~\ref{aa:lem:tree-free} wyklucza taką redukcję.
Sprzeczność dowodzi różnowartościowości homomorfizmu.
Skoro jest on także surjektywny, jest izomorfizmem.
\end{proof}

\begin{uwaga}[Co w dowodzie jest istotne]
Dowód nie zakłada, że $H$ jest normalna w $F(S)$ ani że
$H$ jest skończonego indeksu. Kluczowa jest swoboda
skreślania liter, a nie rachunek rzędów.
Drzewo rozpinające eliminuje wszystkie relacje pochodzące
z kroków „w obie strony” i pozostawia niezależne generatory
odpowiadające pozostałym krawędziom.
\end{uwaga}

\section{Podgrupy skończonego indeksu i wzór Schreiera}
\begin{twierdzenie}[Wzór Schreiera na rangę]
\label{aa:tw:schreier-index}
Niech $F_r$ będzie grupą wolną o skończonej randze $r\ge1$,
a $H\le F_r$ podgrupą o skończonym indeksie
$d=[F_r:H]\ge1$. Wtedy $H$ jest wolna rangi
$1+d(r-1)$.
\end{twierdzenie}
\begin{proof}
Graf Schreiera ma dokładnie $d$ wierzchołków (po jednym
na prawą warstwę). Dla każdego z $r$ generatorów z każdego
wierzchołka wychodzi dokładnie jedna dodatnio zorientowana
krawędź. Mamy więc $dr$ różnych krawędzi (każda ma wybraną
dodatnią orientację), liczonych z krotnościami, również
wtedy, gdy mają te same końce. Przejście wstecz po danej
krawędzi nie jest doliczane jako nowa krawędź.
Dowolne drzewo rozpinające ma dokładnie $d-1$ krawędzi:
można to wykazać indukcją po liczbie wierzchołków, usuwając
liść i jego jedyną krawędź; przypadek $d=1$ jest pusty.
Z Twierdzenia~\ref{aa:tw:nielsen-schreier} pozostałe
$dr-(d-1)=1+d(r-1)$ krawędzie dają wolną bazę $H$.
\end{proof}

\begin{przyklad}[Jądro homomorfizmu $F(a,b)\to C_2$]
\label{aa:prz:schreier-index-two}
Niech $\varepsilon:F(a,b)\to C_2$ wysyła $a$ na $1$,
a $b$ na $0$ (zapisujemy $C_2$ addytywnie).
Dla $H=\ker\varepsilon$ graf Schreiera ma dwa wierzchołki
$H$ oraz $Ha$. Krawędzie o etykiecie $a$ łączą
obie warstwy w obie strony (są to dwie różne krawędzie,
a nie dwa przejścia tej samej krawędzi), zaś przy każdym
wierzchołku istnieje pętla
z etykietą $b$. Jeśli za drzewo wybierzemy pierwszą
krawędź $a:H\to Ha$, to trzy krawędzie poza drzewem
dają odpowiednio elementy $b$, $a^2$ oraz $aba^{-1}$.
Zatem
\[
H=\langle b,a^2,aba^{-1}\rangle
\quad\text{jest grupą wolną rangi }3.
\]
Tego wyniku nie należy mylić z twierdzeniem, że same
$a^2,b$ generują $H$: bez trzeciego generatora ta
konkretna podgrupa jest mniejsza: z udowodnionej wolności
bazy wynika, że trzeci generator nie należy do podgrupy
generowanej przez dwa pozostałe.
\end{przyklad}

\begin{figure}[ht]
\centering
\begin{tikzpicture}[>=Latex, every node/.style={font=\small}]
  \node[draw,circle,minimum size=10mm] (v0) at (0,0) {$H$};
  \node[draw,circle,minimum size=10mm] (v1) at (4.3,0) {$Ha$};
  \draw[->,thick] (v0) to[bend left=21] node[above] {$a$} (v1);
  \draw[->,thick] (v1) to[bend left=21] node[below] {$a$} (v0);
  \draw[->,thick] (v0) to[loop left,looseness=9] node[left] {$b$} (v0);
  \draw[->,thick] (v1) to[loop right,looseness=9] node[right] {$b$} (v1);
\end{tikzpicture}
\caption{Graf Schreiera jądra $F(a,b)\to C_2$: obie krawędzie oznaczone $a$
są różnymi krawędziami, a pętle $b$ odpowiadają dwóm różnym generatorom podgrupy.}
\label{aa:fig:schreier-two}
\end{figure}

\begin{przyklad}[Podgrupa wolna nieskończonej rangi w $F_2$]
\label{aa:prz:infinite-rank-subgroup}
Niech $\psi:F(a,b)\to\mathbb Z$ będzie homomorfizmem
zadanym przez $\psi(a)=1$, $\psi(b)=0$.
Dla $K=\ker\psi$ wierzchołkami grafu Schreiera są liczby
całkowite $n$, odpowiadające warstwom $Ka^n$.
Krawędź $a$ łączy $n$ z $n+1$, a przy każdym $n$
znajduje się pętla o etykiecie $b$.
Wszystkie krawędzie $a$ tworzą drzewo rozpinające (dwustronną
nieskończoną ścieżkę), a pętle $b$ są poza nim.
Droga w drzewie od $0$ do $n$ ma etykietę $a^n$;
odpowiadający jej generator to $a^nba^{-n}$.
Dostajemy więc wolną bazę
$\{a^nba^{-n}:n\in\mathbb Z\}$ grupy $K$.
Pokazuje to, że podgrupa grupy wolnej rangi dwa może mieć
nieskończoną rangę. Nie przeczy to wzorowi Schreiera,
ponieważ $[F(a,b):K]=\infty$.
\end{przyklad}

\section{Paradoks Banacha--Tarskiego: cel i plan dowodu}
\label{aa:sec:banach-tarski}

Korzystamy z konstrukcji $F_2=F(\{a,b\})$ z
Twierdzenia~\ref{aa:tw:free-exists}, z przeliczalności
(Stwierdzenie~\ref{aa:st:free-count}) oraz z teorii działań
grup z wcześniejszej części skryptu.
Przyjmujemy oznaczenia $\bar a=a^{-1}$ i $\bar b=b^{-1}$.
Nie powtarzamy dowodu jednoznaczności redukcji.


Słowo \emph{paradoks} oznacza tutaj twierdzenie sprzeczne z intuicją, a nie sprzeczność logiczną. Pracujemy w zwykłej teorii mnogości ZFC, czyli z aksjomatem wyboru. Kawałkiem nazywamy dowolny podzbiór kuli: nie musi być bryłą, nie musi być spójny ani mierzalny w sensie Lebesgue'a.

Zakładamy znajomość zbiorów, funkcji, liczb rzeczywistych oraz elementarnego rachunku macierzy i wektorów. Korzystamy z przedstawionej wcześniej teorii grup wolnych i działań grup; teorię rozkładów paradoksalnych wprowadzamy poniżej. Nie konstruujemy liczb rzeczywistych ani nie wyprowadzamy podstaw algebry liniowej z aksjomatów. Teoria miary pojawia się dopiero w komentarzu po dowodzie i nie jest narzędziem samej konstrukcji.

\subsection{Co dokładnie zostanie udowodnione?}
Przez $\|x\|$ oznaczamy długość euklidesową wektora $x\in\mathbb R^3$. Niech
\[
 B^3=\{x\in\mathbb R^3:\|x\|\le1\},\qquad S^2=\{x\in\mathbb R^3:\|x\|=1\}.
\]
To odpowiednio kula domknięta i jej brzeg, czyli sfera. Kula zawiera również punkty wnętrza; nie wolno jej utożsamiać ze sferą.

\begin{twierdzenie}[Banach--Tarski: podwojenie kuli]\label{tw:banach-tarski}
Istnieje skończona rodzina parami rozłącznych zbiorów $P_1,\ldots,P_N\subset B^3$ oraz ruchy sztywne $f_i(x)=Q_ix+v_i$, gdzie $Q_i\in\mathrm{SO}(3)$, takie że
\[
 B^3=\bigsqcup_{i=1}^{N}P_i,\qquad
 \bigsqcup_{i=1}^{N}f_i(P_i)
 =(B^3+(-3,0,0))\sqcup(B^3+(3,0,0)).
\]
\end{twierdzenie}
W obu zapisach występują sumy rozłączne. Każdy kawałek jest użyty dokładnie raz; nie rozciągamy go, nie pomniejszamy i nie tworzymy jego dodatkowej kopii. Konstrukcja podana niżej używa co najwyżej $64$ kawałków; nie jest to twierdzenie o najmniejszej możliwej liczbie kawałków.

Grupę $\mathrm{SO}(3)$, czyli grupę obrotów zachowujących orientację i punkt $0$, zdefiniujemy macierzowo w podsekcji~\ref{subsec:so3}. Dodanie $v_i$ pozwala przesuwać kawałki. Odległość między środkami dwóch kul docelowych wynosi $6>2$, więc rzeczywiście są one rozłączne.

\subsection{Siedem etapów}
\begin{enumerate}
\item Wykorzystamy skonstruowaną wcześniej grupę wolną $F_2$ z liter $a,a^{-1},b,b^{-1}$.
\item Podzielimy $F_2$ na cztery zbiory, z których dwa odtworzą pierwszą kopię $F_2$, a pozostałe dwa drugą.
\item Pokażemy dwa konkretne obroty, które nie spełniają żadnej nietrywialnej relacji. Dzięki temu otrzymamy monomorfizm $F_2\to\mathrm{SO}(3)$.
\item Usuniemy ze sfery przeliczalny zbiór punktów stałych $D$. Na pozostałym zbiorze $X$ przeniesiemy rozkład grupy na każdą orbitę.
\item Wykażemy, że dodanie $D$ nie przeszkadza: sfera i $X$ są równoważne przez rozkład na dwa kawałki.
\item Przedłużymy rozkład sfery wzdłuż promieni, otrzymując rozkład kuli bez środka, a następnie osobno wchłoniemy środek.
\item Umieścimy dwie formalne kopie w dwóch różnych miejscach przestrzeni. Sprawdzimy skończoność, rozłączność i pokrycie.
\end{enumerate}
Najtrudniejszy krok algebraiczny, czyli dowód, że wskazane obroty są wolnymi generatorami, przeprowadzimy rachunkiem macierzowym modulo $3$. Nie odwołujemy się do twierdzenia o istnieniu „odpowiednich obrotów”.

\section{Skończone rozkłady i dwie kopie zbioru}
\subsection{Działanie grupy i dozwolone ruchy}
\begin{definicja}[Działanie grupy]
Grupa $G$ działa na zbiorze $Y$, jeśli każdej parze $(g,x)\in G\times Y$ przyporządkowano $gx\in Y$ tak, że $ex=x$ i $(gh)x=g(hx)$. Każde przekształcenie $x\mapsto gx$ jest bijekcją, z odwrotnością $x\mapsto g^{-1}x$. Dla $C\subseteq Y$ piszemy $gC=\{gx:x\in C\}$.
\end{definicja}
Przykłady to działanie grupy na sobie przez mnożenie z lewej, działanie obrotów na sferze oraz działanie ruchów sztywnych na $\mathbb R^3$.
\begin{definicja}[Równoważność przez skończony rozkład]\label{def:rownorozklad}
Niech $E,F\subseteq Y$, gdzie $G$ działa na $Y$. Piszemy $E\mathrel{\equiv_G} F$, jeżeli istnieją skończona podział $E=E_1\sqcup\cdots\sqcup E_k$ oraz $g_1,\ldots,g_k\in G$ takie, że $F=g_1E_1\sqcup\cdots\sqcup g_kE_k$. Słowo \emph{podział} oznacza, że części są parami rozłączne i pokrywają cały zbiór. Puste części można usuwać. Gdy grupa ruchów jest ustalona, opuszczamy indeks $G$.
\end{definicja}
Istnienie takiego rozkładu jest równoważne istnieniu bijekcji $f:E\to F$, która na każdej ze skończenie wielu części $E_i$ ma postać $f(x)=g_ix$. Rozłączność obrazów daje różnowartościowość, a pokrycie $F$ --- surjektywność.

\textbf{Przykład (dwie części).} Niech grupa $(\mathbb Z,+)$ działa na $\mathbb Z$ przez przesunięcia. Zbiór $E=\{0,1\}$ jest równoważny zbiorowi $F=\{0,3\}$: część $\{0\}$ pozostawiamy na miejscu, a część $\{1\}$ przesuwamy o $2$. Obrazy $\{0\}$ i $\{3\}$ są rozłączne i pokrywają $F$. Ten przykład pokazuje, że na różnych częściach mogą działać różne elementy grupy.

\begin{lemat}[Składanie rozkładów]\label{lem:skladanie}
Relacja $\mathrel{\equiv_G}$ jest zwrotna, symetryczna i przechodnia. Złożenie rozkładów na odpowiednio $p$ i $q$ części wymaga co najwyżej $pq$ części.
\end{lemat}
\begin{proof}
Zwrotność daje jedna część i ruch $e$. Dla symetrii bierzemy części $g_iE_i$ zbioru docelowego i ruchy $g_i^{-1}$. Przejdźmy do przechodniości. Niech $f:E\to F$ działa jako $g_i$ na $E_i$, a $h:F\to H$ jako $h_j$ na $F_j$. Dzielimy $E$ na zbiory
\begin{equation}\label{eq:3-1}
 E_{ij}'=E_i\cap g_i^{-1}(F_j).
\end{equation}
Dla każdego $x\in E$ istnieje dokładnie jedno $i$ z $x\in E_i$, a następnie dokładnie jedno $j$ z $g_ix\in F_j$. Zatem zbiory $E_{ij}'$ tworzą podział. Na $E_{ij}'$ złożenie $h\circ f$ jest pojedynczym ruchem $h_jg_i$. Ponieważ złożenie bijekcji jest bijekcją, obrazy tych części są rozłączne i pokrywają $H$. Liczba par $(i,j)$ wynosi $pq$.
\end{proof}

\subsection{Dwie kopie oznaczone etykietami}
Symbol $E\sqcup E$ oznacza sumę rozłączną dwóch kopii $E$. Będziemy ją przedstawiać jako $2E=E\times\{1,2\}$. Pary $(x,1)$ i $(x,2)$ oznaczają dwa różne miejsca docelowe. Etykieta nie jest dodatkową współrzędną euklidesową.

W rozkładzie między kopiami oznaczonymi etykietami każda część ma jedną ustaloną etykietę źródłową i jedną docelową; przekształcenie ma postać $(x,i)\mapsto(gx,j)$. Etykiety służą jedynie do rozróżnienia kopii. Na końcu umieścimy kopie w różnych miejscach $\mathbb R^3$ za pomocą przesunięć.

Dowód lematu~\ref{lem:skladanie} działa bez zmian, gdy w przecięciach \eqref{eq:3-1} uwzględniamy zgodność etykiety pośredniej. Jeżeli $E\mathrel{\equiv_G} F$, możemy wykonać ten sam rozkład osobno w obu kopiach, otrzymując $2E\mathrel{\equiv_G}2F$.
\begin{definicja}[Zbiór paradoksalny]\label{def:paradoksalny}
Zbiór $E$ jest \emph{paradoksalny} dla danego działania grupy, jeśli $E\mathrel{\equiv_G}2E$. Równoważnie, istnieje podział $E=A_1\sqcup\cdots\sqcup A_m\sqcup B_1\sqcup\cdots\sqcup B_n$ i elementy $g_i,h_j\in G$ takie, że
\[
 E=\bigsqcup_{i=1}^{m}g_iA_i,\qquad E=\bigsqcup_{j=1}^{n}h_jB_j.
\]
Pierwsze równanie odtwarza kopię z etykietą $1$, drugie kopię z etykietą $2$. Źródłowe części nie są kopiowane.
\end{definicja}
Niepusty zbiór paradoksalny musi być nieskończony: dla niepustego zbioru skończonego bijekcja $E\to2E$ wymagałaby $|E|=2|E|$. W dalszej części zbudujemy taki rozkład dla $F_2$, a następnie przeniesiemy go na sferę i kulę.
\begin{wniosek}[Przenoszenie paradoksalności]\label{wn:przenoszenie}
Jeśli $E\mathrel{\equiv_G} F$ i $F\mathrel{\equiv_G}2F$, to $E\mathrel{\equiv_G}2E$.
\end{wniosek}
\begin{proof}
Składamy trzy skończone rozkłady: $E\mathrel{\equiv_G} F\mathrel{\equiv_G}2F\mathrel{\equiv_G}2E$. Ostatnia równoważność jest dwukrotnym zastosowaniem odwrotnego rozkładu $F\mathrel{\equiv_G} E$. Lemat~\ref{lem:skladanie} zapewnia, że po złożeniu nadal mamy skończenie wiele części i pojedynczy dozwolony ruch na każdej.
\end{proof}
\subsection{Przesunięcie wzdłuż nieskończonego łańcucha}
\begin{lemat}[Przesunięcie łańcucha]\label{lem:lancuch}
Niech $D\subseteq E\subseteq Y$, a $T:Y\to Y$ będzie przekształceniem wyznaczonym przez pewien element $G$. Załóżmy, że zbiory $D,T(D),T^2(D),\ldots$ są parami rozłączne i wszystkie zawierają się w $E$. Wówczas $E\mathrel{\equiv_G} E\setminus D$ przy użyciu co najwyżej dwóch części.
\end{lemat}
\begin{proof}
Niech $L=\bigcup_{n=0}^{\infty}T^n(D)$. Rozłączność poziomów daje $T(L)=\bigcup_{n=1}^{\infty}T^n(D)=L\setminus D$. Dzielimy $E$ na dwie części $L$ i $E\setminus L$. Pierwszą przenosimy przez $T$, drugą pozostawiamy bez zmian. Obrazy to odpowiednio $L\setminus D$ i $E\setminus L$; są rozłączne i ich suma wynosi $E\setminus D$.
\end{proof}
\begin{uwaga}
Użyliśmy nieskończonej sumy do zdefiniowania jednego kawałka $L$. Nie wykonujemy nieskończenie wielu ruchów na różnych poziomach: na całym $L$ działa jeden i ten sam ruch $T$. Rozkład ma więc dwie części, mimo że jedna z nich jest nieskończoną unią.
\end{uwaga}

\section{Paradoksalny rozkład grupy wolnej}
\subsection{Podział według pierwszej litery}
Dla $s\in\mathcal A$ niech $W(s)$ oznacza zbiór niepustych słów zredukowanych zaczynających się literą $s$. Zatem
\begin{equation}\label{eq:4-1}
 F_2=\{e\}\sqcup W(a)\sqcup W(\bar a)\sqcup W(b)\sqcup W(\bar b).
\end{equation}
Każde słowo niepuste ma dokładnie jedną pierwszą literę, więc zbiory po prawej stronie są parami rozłączne i pokrywają $F_2$.
\begin{lemat}[Kasowanie pierwszej litery]\label{lem:pierwsza-litera}
Mamy $aW(\bar a)=F_2\setminus W(a)$ oraz $bW(\bar b)=F_2\setminus W(b)$.
\end{lemat}
\begin{proof}
Sprawdzimy obie inkluzje pierwszej równości. Słowo z $W(\bar a)$ ma postać $\bar au$, gdzie $u$ jest puste albo nie zaczyna się literą $a$; w przeciwnym razie nie byłoby zredukowane. Mnożenie przez $a$ z lewej kasuje parę $a\bar a$ i daje właśnie $u$. Wynik nie należy do $W(a)$.

Odwrotnie, jeśli $u\notin W(a)$, to $u=e$ albo jego pierwsza litera nie jest $a$. Słowo $\bar au$ jest wtedy zredukowane i należy do $W(\bar a)$, a $a(\bar au)=u$. Dotyczy to także $u=e$: jego poprzednikiem jest słowo $\bar a$. Drugą równość dowodzimy tak samo, zastępując $a$ przez $b$.
\end{proof}
Już teraz mamy dwa równania $F_2=W(a)\sqcup aW(\bar a)$ oraz $F_2=W(b)\sqcup bW(\bar b)$. Nie jest to jeszcze pełny rozkład paradoksalny, ponieważ cztery źródłowe zbiory $W(s)$ nie zawierają $e$. Nie wolno zgubić tego jednego elementu.

\subsection{Włączenie słowa pustego bez dodatkowej kopii}
Połóżmy $P=\{a^n:n=0,1,2,\ldots\}$, gdzie $a^0=e$. Zauważmy dokładnie, że
\begin{equation}\label{eq:4-2}
 aP=P\setminus\{e\}=P\cap W(a).
\end{equation}
Pierwsza równość wynika z $a\cdot a^n=a^{n+1}$ i z tego, że potęgi $a^0,a^1,\ldots$ są parami różne. Druga mówi, że spośród tych potęg tylko $a^0$ nie zaczyna się literą $a$.
\begin{twierdzenie}[Cztery kawałki grupy]\label{tw:cztery-kawalki}
Zbiory
\begin{equation}\label{eq:4-3}
 A_1=W(a)\setminus P,\quad A_2=W(\bar a)\cup P,
 \qquad B_1=W(b),\quad B_2=W(\bar b)
\end{equation}
tworzą podział $F_2$ i spełniają
\begin{equation}\label{eq:4-4}
 F_2=A_1\sqcup aA_2,\qquad F_2=B_1\sqcup bB_2.
\end{equation}
W szczególności $F_2$ jest zbiorem paradoksalnym względem działania na sobie przez mnożenie z lewej.
\end{twierdzenie}
\begin{proof}
\emph{Partycja źródła.} W podziału \eqref{eq:4-1} wydzielamy dodatnie potęgi $a$ ze zbioru $W(a)$ i razem z $e$ dołączamy je do $W(\bar a)$, tworząc $A_2=W(\bar a)\cup P$. Pozostała część $W(a)$ to $A_1$. Nic nie zmieniamy w $W(b)$ i $W(\bar b)$. Zbiory pozostają rozłączne, a teraz pokrywają także $e$.

\emph{Pierwsza kopia.} Korzystając z lematu~\ref{lem:pierwsza-litera} i \eqref{eq:4-2}, otrzymujemy $aA_2=aW(\bar a)\sqcup aP=(F_2\setminus W(a))\sqcup(P\setminus\{e\})$. Pierwszy składnik nie zawiera słów zaczynających się od $a$, a drugi składa się właśnie z dodatnich potęg $a$. Ich suma jest rozłączna. Zbiór $A_1=W(a)\setminus P$ jest rozłączny z obydwoma składnikami i uzupełnia je do całego $F_2$.

\emph{Druga kopia.} Lemat~\ref{lem:pierwsza-litera} od razu daje $B_1\sqcup bB_2=W(b)\sqcup(F_2\setminus W(b))=F_2$. W obu kopiach sprawdziliśmy zarówno pokrycie, jak i rozłączność.
\end{proof}
\begin{center}
\begin{tabular}{c|c|c}
 kawałek źródłowy&ruch&etykieta celu\\\hline
 $A_1$&$e$&$1$\\ $A_2$&$a$&$1$\\ $B_1$&$e$&$2$\\ $B_2$&$b$&$2$
\end{tabular}
\end{center}
\begin{przyklad}[Sprawdzenie kilku punktów]
$e,a,a^2$ należą do $A_2$ i po ruchu $a$ trafiają odpowiednio do $a,a^2,a^3$ w pierwszej kopii. Element $\bar a\in A_2$ trafia do $e$ w pierwszej kopii. Słowo $ab\in A_1$ pozostaje $ab$ w tej samej kopii. Element $\bar b\in B_2$ trafia do $e$ w drugiej kopii. Dwie kopie $e$ w zbiorze docelowym mają różnych poprzedników i różne etykiety.
\end{przyklad}

\section{Dwa konkretne obroty generujące grupę wolną}
\subsection{Grupa obrotów i punkty stałe jej elementów}\label{subsec:so3}
\begin{definicja}
Macierz $Q\in M_3(\mathbb R)$ jest ortogonalna, jeżeli $Q^TQ=I$. Definiujemy $\mathrm{SO}(3)=\{Q\in M_3(\mathbb R):Q^TQ=I,\ \det Q=1\}$. Macierze te zachowują długości i orientację. Działają na $\mathbb R^3$ przez $x\mapsto Qx$ oraz, przez ograniczenie, na $S^2$.
\end{definicja}
Rzeczywiście $\|Qx\|^2=x^TQ^TQx=\|x\|^2$. Iloczyn dwóch takich macierzy i odwrotność $Q^{-1}=Q^T$ również spełniają warunki definicji. Mamy więc grupę.
\begin{lemat}[Oś nietrywialnego obrotu]\label{lem:os-obrotu}
Jeśli $Q\in\mathrm{SO}(3)$ i $Q\ne I$, to zbiór wektorów ustalanych przez $Q$ jest prostą przez $0$. W szczególności $Q$ ma dokładnie dwa punkty stałe na $S^2$.
\end{lemat}
\begin{proof}
Najpierw wykażemy istnienie niezerowego wektora stałego. Mamy
\[
 \det(Q-I)=\det Q\det(I-Q^{-1})=\det(I-Q^T)
 =\det(I-Q)=-\det(Q-I).
\]
Ostatni znak minus wynika z nieparzystego wymiaru $3$. Zatem $\det(Q-I)=0$, więc układ $(Q-I)x=0$ ma rozwiązanie $x\ne0$.

Gdyby $Q$ ustalało dwa liniowo niezależne wektory, ustalałoby całą rozpiętą przez nie płaszczyznę. Wybierzmy w niej ortonormalną bazę $u,v$ i jednostkowy wektor $w$ prostopadły do tej płaszczyzny. Ortogonalność implikuje $Qw\perp u,v$ i $\|Qw\|=1$, zatem $Qw=\pm w$. W bazie $u,v,w$ wyznacznik wynosi odpowiednio $1$ albo $-1$. Warunek $\det Q=1$ wymusza $Qw=w$, a wtedy $Q=I$. To sprzeczność. Przestrzeń punktów stałych ma więc wymiar dokładnie jeden. Prosta przez $0$ przecina sferę w dwóch punktach antypodycznych.
\end{proof}
\subsection{Wybór generatorów}
Rozważmy
\begin{equation}\label{eq:5-1}
 A=\frac13\begin{pmatrix}3&0&0\\0&1&-2\sqrt2\\0&2\sqrt2&1\end{pmatrix},\qquad
 B=\frac13\begin{pmatrix}1&-2\sqrt2&0\\2\sqrt2&1&0\\0&0&3\end{pmatrix}.
\end{equation}
W obu macierzach blok $2\times2$ ma postać $\begin{psmallmatrix}c&-s\\s&c\end{psmallmatrix}$, gdzie $c=1/3$, $s=2\sqrt2/3$ i $c^2+s^2=1$. Jego kolumny są ortonormalne i wyznacznik wynosi $1$. Wobec tego $A,B\in\mathrm{SO}(3)$. Macierz $A$ obraca wokół osi $x$, a $B$ wokół osi $z$, o kąt o cosinusie $1/3$.

Z własności uniwersalnej istnieje homomorfizm $\rho:F_2\to\mathrm{SO}(3)$, $\rho(a)=A$, $\rho(b)=B$. Musimy udowodnić, że jest \emph{różnowartościowy}. To znacznie mocniejsze niż $AB\ne BA$.

\subsection{Rachunek modulo 3 i usunięcie pierwiastków}
Dla liczb całkowitych $m\equiv n\pmod3$ oznacza, że $m-n$ jest podzielne przez $3$. Występują tylko trzy reszty, które wygodnie zapisujemy jako $0,1,-1$. Mamy $2\equiv-1$. Redukcja modulo $3$ zachowuje sumy i iloczyny. Dotyczy to również iloczynów macierzy całkowitych, ponieważ ich wpisy powstają przez sumowanie iloczynów wpisów. Reszty tworzą ciało $\mathbb F_3$; potrzebujemy jedynie tego, że iloczyn dowolnej liczby czynników $1$ i $-1$ nie jest zerem. Nie będziemy dzielić przez $3$ po redukcji modulo $3$.

Dla litery $s\in\{a,\bar a,b,\bar b\}$ oznaczmy przez $R_s$ odpowiednio $A,A^{-1},B,B^{-1}$. Macierze $3R_s$ mają wpisy z pierwiastkiem $\sqrt2$. Aby dostać macierze całkowite, wprowadzamy pomocniczo $D_0=\operatorname{diag}(1,\sqrt2,1)$ oraz $N_s=3D_0^{-1}R_sD_0$. Bezpośrednie mnożenie daje
\begin{equation}\label{eq:5-2}
 N_a=\begin{pmatrix}3&0&0\\0&1&-2\\0&4&1\end{pmatrix},\qquad
 N_{\bar a}=\begin{pmatrix}3&0&0\\0&1&2\\0&-4&1\end{pmatrix},
\end{equation}
\begin{equation}\label{eq:5-3}
 N_b=\begin{pmatrix}1&-4&0\\2&1&0\\0&0&3\end{pmatrix},\qquad
 N_{\bar b}=\begin{pmatrix}1&4&0\\-2&1&0\\0&0&3\end{pmatrix}.
\end{equation}
Na przykład wpis $(2,3)$ w $D_0^{-1}(3A)D_0$ wynosi $(1/\sqrt2)(-2\sqrt2)\cdot1=-2$, a wpis $(3,2)$ wynosi $1\cdot2\sqrt2\cdot\sqrt2=4$. Macierze odwrotnych obrotów otrzymujemy z $A^{-1}=A^T$, $B^{-1}=B^T$, co wyjaśnia zmianę znaków.
\begin{uwaga}[Rola pomocniczej zmiany bazy]
Macierz $D_0$ nie jest obrotem i służy wyłącznie do obliczeń. Sprawdzamy za jej pomocą relacje między obrotami $A,B$. W konstrukcji kawałki sfery i kuli przemieszczamy jedynie obrotami i przesunięciami.
\end{uwaga}

\subsection{Cztery wektory i jedna tabela}
W rachunku modulo $3$ przyjmijmy wektory kolumnowe
\[
 w_a=\begin{pmatrix}0\\1\\1\end{pmatrix},\quad
 w_{\bar a}=\begin{pmatrix}0\\1\\-1\end{pmatrix},\quad
 w_b=\begin{pmatrix}1\\-1\\0\end{pmatrix},\quad
 w_{\bar b}=\begin{pmatrix}1\\1\\0\end{pmatrix}.
\]
Iloczyn $uv^T$ oznacza macierz, której wpis $(i,j)$ to $u_iv_j$. Z \eqref{eq:5-2}--\eqref{eq:5-3} otrzymujemy
\begin{equation}\label{eq:5-4}
 \overline{N_s}=w_sw_s^T,
\end{equation}
gdzie kreska oznacza redukcję wszystkich wpisów modulo $3$. Na przykład $\overline{N_a}=\begin{psmallmatrix}0&0&0\\0&1&1\\0&1&1\end{psmallmatrix}=w_aw_a^T$ i $\overline{N_b}=\begin{psmallmatrix}1&-1&0\\-1&1&0\\0&0&0\end{psmallmatrix}=w_bw_b^T$. Pozostałe dwie równości sprawdzamy identycznie.

Tabela iloczynów skalarnych $w_s^Tw_t$, już modulo $3$, ma postać
\begin{equation}\label{eq:5-5}
\begin{array}{c|rrrr}
 s\backslash t&a&\bar a&b&\bar b\\\hline
 a&-1&0&-1&1\\
 \bar a&0&-1&-1&1\\
 b&-1&-1&-1&0\\
 \bar b&1&1&0&-1
\end{array}
\end{equation}
Na przekątnej jest $1+1=2\equiv-1$. Dla pary $a,\bar a$ mamy $1-1=0$, a dla $b,\bar b$ także $1-1=0$. Przy dwóch różnych osiach tylko jedna współrzędna daje niezerowy wkład, równy $1$ albo $-1$. Otrzymujemy ważną równoważność:
\begin{equation}\label{eq:5-6}
 w_s^Tw_t=0\quad\Longleftrightarrow\quad t=s^{-1} \qquad(\text{w }\mathbb F_3).
\end{equation}
\begin{lemat}[Iloczyn pomocniczych macierzy słowa zredukowanego]\label{lem:liczniki}
Jeżeli $s_1\cdots s_n$ jest niepustym słowem zredukowanym, to $\overline{N_{s_1}\cdots N_{s_n}}$ nie jest macierzą zerową.
\end{lemat}
\begin{proof}
Dla wektorów $u,v,x,y$ zachodzi $(uv^T)(xy^T)=u(v^Tx)y^T$. Wpis $(i,j)$ po lewej stronie jest równy $\sum_k u_iv_kx_ky_j=u_i(\sum_kv_kx_k)y_j$, co dowodzi wzoru bez żadnej dodatkowej teorii. Stosując go kolejno do \eqref{eq:5-4}, dostajemy
\begin{equation}\label{eq:5-7}
 \overline{N_{s_1}\cdots N_{s_n}}=
 \left(\prod_{j=1}^{n-1}w_{s_j}^Tw_{s_{j+1}}\right)w_{s_1}w_{s_n}^T.
\end{equation}
Dla $n=1$ iloczyn skalarnych czynników jest pusty i z definicji równy $1$. Ponieważ słowo jest zredukowane, $s_{j+1}\ne s_j^{-1}$. Według \eqref{eq:5-6} każdy czynnik skalarny jest $1$ albo $-1$, więc ich iloczyn również jest niezerowy. Każdy z wektorów $w_s$ ma pewną współrzędną równą $1$. Macierz $w_{s_1}w_{s_n}^T$ ma zatem co najmniej jeden wpis równy $1$. Pomnożenie przez niezerowy skalar nie może wyzerować tego wpisu.
\end{proof}
\begin{twierdzenie}[Wolna podgrupa obrotów]\label{tw:wolne-obroty}
Homomorfizm $\rho:F_2\to\mathrm{SO}(3)$ określony przez $\rho(a)=A$, $\rho(b)=B$ jest różnowartościowy.
\end{twierdzenie}
\begin{proof}
Niech $w=s_1\cdots s_n$ będzie niepustym słowem zredukowanym. Z definicji $N_s$ i kolejnych skróceń $D_0D_0^{-1}$ wynika $N_{s_1}\cdots N_{s_n}=3^nD_0^{-1}\rho(w)D_0$. Gdyby $\rho(w)=I$, prawa strona byłaby równa $3^nI$. Ponieważ $n\ge1$, jej redukcja modulo $3$ byłaby zerowa, wbrew lematowi~\ref{lem:liczniki}. Zatem żadne niepuste słowo zredukowane nie przechodzi na $I$.

Jeśli $\rho(u)=\rho(v)$, to $\rho(v^{-1}u)=I$. Jedynym elementem mającym taki obraz jest $e$, więc $v^{-1}u=e$ i $u=v$. To dokładnie różnowartościowość.
\end{proof}
\begin{uwaga}[Czego rachunek modulo 3 nie dowodzi]
Nie zanurzyliśmy nieskończonej grupy w skończonej grupie macierzy modulo $3$. Macierze $\overline{N_s}$ są osobliwe. Ponadto $N_sN_{s^{-1}}=9I$, a nie $I$. Badaliśmy podzielność wpisów liczników iloczynu obrotów, a nie homomorfizm grupy obrotów uzyskany przez redukcję ułamka $1/3$.
\end{uwaga}

\section{Od grupy do sfery bez punktów wyjątkowych}
Od tej chwili oznaczamy przez $G=\rho(F_2)\le\mathrm{SO}(3)$ wolną podgrupę z twierdzenia~\ref{tw:wolne-obroty}. Dzięki różnowartościowości możemy utożsamiać elementy $G$ z ich jedynymi słowami zredukowanymi. Litery $a,b$ w działaniu na sferze oznaczają obroty $A,B$.

\subsection{Orbity, stabilizatory i działanie swobodne}
\begin{definicja}
Orbita punktu $x$ to $\operatorname{Orb}(x)=Gx=\{gx:g\in G\}$. Stabilizator punktu to $\operatorname{Stab}(x)=\{g\in G:gx=x\}$. Działanie jest \emph{swobodne}, jeśli $\operatorname{Stab}(x)=\{e\}$ dla każdego punktu rozważanego zbioru.
\end{definicja}
Relacja „należą do tej samej orbity” jest równoważnością: $x=ex$, równanie $y=gx$ daje $x=g^{-1}y$, a z $y=gx$ i $z=hy$ mamy $z=(hg)x$. Orbity są więc parami rozłączne i pokrywają zbiór, na którym działa $G$.
\begin{uwaga}[Grupa wolna a działanie swobodne]
Wolna grupa nie musi działać swobodnie! Nasza grupa $G$ jest wolna, ale obrót $A$ ustala dwa punkty osi $x$ na sferze. Aby punkt sfery miał jednoznaczny opis przez słowo i reprezentanta orbity, trzeba usunąć wszystkie punkty ustalane przez element $G$ różny od identyczności.
\end{uwaga}

\subsection{Przeliczalny zbiór wyjątkowy}
Definiujemy
\begin{equation}\label{eq:6-1}
 D=\bigcup_{g\in G\setminus\{e\}}\operatorname{Fix}_{S^2}(g),\qquad X=S^2\setminus D,
\end{equation}
gdzie $\operatorname{Fix}_{S^2}(g)=\{x\in S^2:gx=x\}$.
\begin{lemat}\label{lem:wyjatkowy}
Zbiór $D$ jest co najwyżej przeliczalny i niezmienniczy względem $G$. Zbiór $X$ jest niezmienniczy, a działanie $G$ na $X$ jest swobodne.
\end{lemat}
\begin{proof}
Grupa $G$ jest przeliczalna na mocy stwierdzenia~\ref{st:przeliczalnosc-f2} i monomorfizmu $\rho$. Każdy $g\ne e$ jest nietrywialnym obrotem i ma dokładnie dwa punkty stałe na sferze, zgodnie z lematem~\ref{lem:os-obrotu}. Przeliczalna suma zbiorów dwuelementowych jest co najwyżej przeliczalna: wyliczamy elementy grupy i dla każdego wypisujemy oba punkty osi. Można najpierw wypisać punkt, którego pierwsza niezerowa współrzędna jest dodatnia, więc ten porządek nie wymaga nowych arbitralnych wyborów.

Jeśli $gx=x$, to dla dowolnego $h\in G$ mamy $(hgh^{-1})(hx)=hx$. Sprzężony element $hgh^{-1}$ nie jest neutralny, gdyż równość $hgh^{-1}=e$ implikowałaby $g=e$. Zatem $hD\subseteq D$. Stosując to samo do $h^{-1}$, dostajemy $D\subseteq hD$, więc $hD=D$. Ponieważ $h$ jest bijekcją sfery, także $hX=X$.

Wreszcie, jeśli $x\in X$ i $gx=x$, to $g\ne e$ oznaczałoby $x\in D$, sprzecznie z definicją $X$. Stabilizator każdego punktu $X$ jest więc trywialny.
\end{proof}

\subsection{Wybór reprezentantów: dokładne miejsce aksjomatu wyboru}
\begin{definicja}[Aksjomat wyboru w potrzebnej postaci]
Dla każdej rodziny niepustych zbiorów istnieje funkcja, która każdemu zbiorowi z tej rodziny przyporządkowuje jeden z jego elementów.
\end{definicja}
Zastosujmy ten aksjomat do rodziny wszystkich orbit działania $G$ na $X$. Wybieramy zbiór $T\subseteq X$ zawierający dokładnie po jednym punkcie z każdej orbity. To jest zasadniczy niekonstruktywny wybór w dowodzie. Nie podajemy algorytmu wyznaczającego $T$.
\begin{lemat}[Jednoznaczne współrzędne orbitowe]\label{lem:wspolrzedne-orbitowe}
Każdy punkt $x\in X$ ma dokładnie jedno przedstawienie $x=gt$, gdzie $g\in G$ i $t\in T$.
\end{lemat}
\begin{proof}
Istnienie wynika z wyboru reprezentanta orbity $x$. Jeśli $gt=hu$, to $t$ i $u$ należą do tej samej orbity, bo $t=g^{-1}hu$. Z wyboru jednego reprezentanta wynika $t=u$. Równość $gt=ht$ daje $(h^{-1}g)t=t$. Swoboda działania wymusza $h^{-1}g=e$, czyli $g=h$.
\end{proof}
\begin{uwaga}
To, że każda orbita jest przeliczalna, nie daje automatycznie przepisu wyboru jednego punktu z każdej orbity. Orbit może być nieprzeliczalnie wiele. Aksjomat wyboru nie jest tu skrótem dla wyliczenia wszystkich punktów sfery.
\end{uwaga}

\subsection{Cztery kawałki na wszystkich orbitach jednocześnie}
Dla $C\subseteq G$ definiujemy $C^T=\{gt:g\in C,\ t\in T\}\subseteq X$. Z jednoznaczności w lemacie~\ref{lem:wspolrzedne-orbitowe} wynika, że rozłączne podzbiory $G$ dają rozłączne zbiory $C^T$. Ponadto
\begin{equation}\label{eq:6-2}
 h(C^T)=(hC)^T\qquad(h\in G),
\end{equation}
ponieważ $h(gt)=(hg)t$. Na wszystkich orbitach działa ten sam obrót $h$, niezależny od wybranego reprezentanta.
\begin{twierdzenie}[Paradoksalność $X$]\label{tw:paradoksalnosc-x}
Zbiór $X=S^2\setminus D$ jest paradoksalny pod działaniem obrotów. Wystarczają cztery części.
\end{twierdzenie}
\begin{proof}
Przenosimy zbiory z \eqref{eq:4-3} na orbity: $C_1=A_1^T$, $C_2=A_2^T$, $C_3=B_1^T$, $C_4=B_2^T$. Tworzą podział $X$, ponieważ odpowiadające im zbiory grupowe tworzą podział $G$, a każdy punkt ma jednoznaczne współrzędne $gt$. Z \eqref{eq:4-4} i \eqref{eq:6-2} otrzymujemy $X=C_1\sqcup aC_2$ oraz $X=C_3\sqcup bC_4$. To rozkład $X\mathrel{\equiv_G}2X$. Ruchy to $h_1=I$, $h_2=A$, $h_3=I$, $h_4=B$, a etykiety docelowe to odpowiednio $1,1,2,2$.
\end{proof}
Nie mnożymy liczby kawałków przez liczbę orbit: $C_1$ jest jednym zbiorem zawierającym stosowne punkty wszystkich orbit, i tak samo dla $C_2,C_3,C_4$.

\section{Przywrócenie wszystkich punktów sfery}
Pozostało uzasadnić $S^2\equiv X$. Nie wystarcza zdanie „zbiór przeliczalny jest mały”. Potrzebujemy skończonego rozkładu i obrotów przypisanych jego częściom.

\subsection{Wybór kąta poza zbiorem przeliczalnym}
\begin{lemat}\label{lem:nieprzeliczalnosc}
Żaden niepusty otwarty przedział liczb rzeczywistych nie jest przeliczalny. W szczególności można w nim wybrać liczbę poza dowolnym podanym zbiorem co najwyżej przeliczalnym.
\end{lemat}
\begin{proof}
Niech $x_1,x_2,\ldots$ będzie dowolną listą punktów przedziału. Wybierzmy w nim niejednopunktowy przedział domknięty $I_0$. Indukcyjnie wybieramy niejednopunktowy przedział domknięty $I_n\subset I_{n-1}$ o długości co najwyżej połowy poprzedniej, taki że $x_n\notin I_n$. Możemy go umieścić po jednej stronie $x_n$; jeśli $x_n$ nie należy do $I_{n-1}$, dowolny mniejszy podprzedział wystarczy.

Przedziały są zstępujące, domknięte i ich długości dążą do zera. Z zupełności liczb rzeczywistych ich przecięcie zawiera dokładnie jeden punkt $x$: jest nim wspólna granica lewych i prawych końców. Dla każdego $n$ mamy $x\in I_n$, więc $x\ne x_n$. Lista nie obejmowała całego przedziału.
\end{proof}
Sfera też jest nieprzeliczalna: zawiera półokrąg parametryzowany różnowartościowo przez $\theta\mapsto(\cos\theta,\sin\theta,0)$ dla $0<\theta<\pi$. Przeliczalność zbioru trójek indeksów naturalnych wynika z wyliczania ich według sumy indeksów; dla każdej sumy jest tylko skończenie wiele trójek. Dlatego suma przeliczalnej rodziny skończonych zbiorów wykluczonych kątów jest przeliczalna.

\subsection{Obrót rozdzielający przeliczalny zbiór}
\begin{lemat}\label{lem:obrot-rozdzielajacy}
Dla dowolnego co najwyżej przeliczalnego $D\subseteq S^2$ istnieje obrót $R\in\mathrm{SO}(3)$, dla którego zbiory $D,RD,R^2D,\ldots$ są parami rozłączne.
\end{lemat}
\begin{proof}
Jeśli $D=\varnothing$, teza jest oczywista. Niech $D\ne\varnothing$. Wybieramy punkt $u\in S^2$ poza $D\cup(-D)$, co jest możliwe, bo ten ostatni zbiór jest przeliczalny. Żaden punkt $D$ nie leży na osi $\mathbb R u$.

Niech $R_u(\theta)$ oznacza obrót wokół tej osi o kąt $\theta$. Można go zdefiniować, wybierając ortonormalną bazę $u,v,w$ o dodatniej orientacji i przyjmując
\[
 R_u(\theta)(\alpha u+\beta v+\gamma w)=\alpha u+
 (\beta\cos\theta-\gamma\sin\theta)v+
 (\beta\sin\theta+\gamma\cos\theta)w.
\]
Ze wzorów dodawania kątów wynika $R_u(\theta)^n=R_u(n\theta)$.

Ustalmy $d,e\in D$ i $n\ge1$. Ponieważ $d$ nie leży na osi, jego składowa w płaszczyźnie $u^\perp$ jest niezerowa. Obracając tę składową wokół osi, otrzymujemy okrąg; każde położenie na nim wyznacza kąt jednoznacznie modulo $2\pi$. Zatem albo nie istnieje obrót wokół tej osi przenoszący $d$ na $e$, albo istnieje dokładnie jedna klasa kątów $\alpha$ modulo $2\pi$ o tej własności.

Równanie $R_u(n\theta)d=e$ sprowadza się w drugim przypadku do $n\theta\equiv\alpha\pmod{2\pi}$. W przedziale $0<\theta<2\pi$ ma najwyżej $n$ rozwiązań. Dla każdej trójki $(d,e,n)$ wykluczamy więc skończenie wiele kątów. Trójek jest co najwyżej przeliczalnie wiele. Na mocy lematu~\ref{lem:nieprzeliczalnosc} możemy wybrać $\theta$ poza unią wszystkich wykluczonych kątów.

Dla $R=R_u(\theta)$ mamy teraz
\begin{equation}\label{eq:7-1}
 R^nD\cap D=\varnothing\qquad(n\ge1).
\end{equation}
Jeśli dla $0\le i<j$ istniałby punkt w $R^iD\cap R^jD$, to jego obraz przez $R^{-i}$ należałby do $D\cap R^{j-i}D$, wbrew \eqref{eq:7-1}. Dostajemy żądaną rozłączność wszystkich poziomów.
\end{proof}
\begin{uwaga}
Nowy obrót $R$ nie musi należeć do $G$. Ponieważ $D$ jest $G$-niezmienniczy, żaden element $G$ nie mógłby przenieść niepustego $D$ całkowicie poza $D$. Od tego miejsca dla równoważności całej sfery dopuszczamy pełną grupę obrotów $\mathrm{SO}(3)$.
\end{uwaga}
\begin{twierdzenie}[Pełna sfera jest paradoksalna]\label{tw:pelna-sfera}
Sfera $S^2$ jest równoważna przez skończony rozkład dwóm swoim etykietowanym kopiom, przy użyciu samych obrotów wokół $0$. Wystarcza co najwyżej $16$ części.
\end{twierdzenie}
\begin{proof}
Stosujemy lemat~\ref{lem:obrot-rozdzielajacy} do zbioru $D$ z \eqref{eq:6-1}. Niech $L=\bigcup_{n=0}^\infty R^nD$. Wszystkie poziomy leżą na sferze, więc lemat~\ref{lem:lancuch} daje dwuczęściowy rozkład $S^2\equiv X$. Twierdzenie~\ref{tw:paradoksalnosc-x} daje $X\equiv2X$. Przeniesienie paradoksalności z wniosku~\ref{wn:przenoszenie} kończy dowód istnienia rozkładu $S^2\equiv2S^2$.

Poniżej podajemy dokładne części, aby samo odwołanie do przechodniości nie ukrywało konstrukcji ani oszacowania ich liczby. Niech $E_0=L$, $E_1=S^2\setminus L$, $g_0=R$, $g_1=I$. Bijekcja $f:S^2\to X$ ma na $E_i$ postać $g_i$. Zbiory $g_0E_0=L\setminus D$ i $g_1E_1=S^2\setminus L$ tworzą podział $X$.

Oznaczmy przez $C_j,h_j,\ell(j)$ cztery części, ruchy i etykiety z dowodu twierdzenia~\ref{tw:paradoksalnosc-x}. Dla $i,k\in\{0,1\}$ i $j\in\{1,2,3,4\}$ połóżmy
\begin{equation}\label{eq:7-2}
 H_{ijk}=\{x\in E_i:g_ix\in C_j,\ h_jg_ix\in g_kE_k\}.
\end{equation}
Na tej części działa obrót
\begin{equation}\label{eq:7-3}
 q_{ijk}=g_k^{-1}h_jg_i,
\end{equation}
z etykietą docelową $\ell(j)$. Dla każdego $x$ indeksy $i,j,k$ są wyznaczone jednoznacznie, w tej kolejności. Zatem części \eqref{eq:7-2} tworzą podział sfery.

Sprawdźmy jeszcze obrazy. Całe odwzorowanie jest złożeniem bijekcji $f$, czteroczęściowej bijekcji $X\to2X$ i odwrotnej bijekcji $f^{-1}$ w każdej kopii. Jest zatem bijekcją $S^2\to2S^2$. Stąd obrazy kawałków są rozłączne wewnątrz każdej etykiety i pokrywają obie sfery. Indeksów jest $2\cdot4\cdot2=16$. Każde $q_{ijk}$ jest złożeniem obrotów, a więc nadal jednym obrotem w $\mathrm{SO}(3)$.
\end{proof}

\section{Od sfery do pełnej kuli}
\subsection{Przedłużenie kawałków wzdłuż promieni}
Niech $B^3_*=B^3\setminus\{0\}$. Dla $H\subseteq S^2$ definiujemy jego stożek promieniowy w kuli bez środka:
\[
 \widehat H=\{rx:x\in H,\ 0<r\le1\}.
\]
Każdy $y\in B^3_*$ ma dokładnie jedno przedstawienie $y=rx$ z $0<r\le1$, $x\in S^2$: mianowicie $r=\|y\|$, $x=y/\|y\|$. Dlatego podział sfery daje podział $B^3_*$ przez odpowiednie stożki.
\begin{lemat}[Obroty komutują z przedłużeniem promieniowym]\label{lem:stozki}
Dla $Q\in\mathrm{SO}(3)$ i $H\subseteq S^2$ mamy $Q\widehat H=\widehat{QH}$.
\end{lemat}
\begin{proof}
Liniowość macierzy daje $Q(rx)=rQx$, a ortogonalność daje $\|Qx\|=1$. To dowodzi jednej inkluzji. Jeśli $ry$ należy do $\widehat{QH}$, to $y=Qx$ dla pewnego $x\in H$, więc $ry=Q(rx)\in Q\widehat H$. Mamy także drugą inkluzję.
\end{proof}
\begin{wniosek}\label{wn:kula-bez-srodka}
$B^3_*\equiv2B^3_*$ za pomocą co najwyżej $16$ kawałków i obrotów.
\end{wniosek}
\begin{proof}
Bierzemy stożki nad częściami $H_{ijk}$ z \eqref{eq:7-2} i stosujemy te same obroty $q_{ijk}$. Jednoznaczność kierunku i promienia zapewnia rozłączność części wyjściowych, a lemat~\ref{lem:stozki} przenosi rozłączność i pokrycie obrazów z każdej sfery na całą kulę bez środka. Nie zmieniamy promienia żadnego punktu podczas tego etapu.
\end{proof}
\begin{uwaga}[Dlaczego nie dopisać po prostu środka?]
Wszystkie obroty wokół $0$ ustalają $0$. Z jednego takiego punktu nie otrzymamy dwóch punktów oznaczonych różnymi etykietami przez dopisanie go do dwóch części --- przestałyby być rozłączne. Środek wymaga osobnego argumentu i ruchu wokół osi nieprzechodzącej przez $0$.
\end{uwaga}

\subsection{Wchłonięcie środka za pomocą obrotu}
W płaszczyźnie $z=0$ weźmy okrąg
\[
 \Gamma=\{(1/4,0,0)+(1/4)(\cos t,\sin t,0):t\in\mathbb R\}.
\]
Ma środek $c=(1/4,0,0)$ i promień $1/4$. Punkt $0$ należy do $\Gamma$ (dla $t=\pi$). Każdy punkt okręgu ma normę nie większą niż $\|c\|+1/4=1/2$, więc cały okrąg leży wewnątrz $B^3$.

Niech $U$ będzie obrotem wokół osi $z$ o kąt $\alpha=2\pi\sqrt2$. Określmy ruch sztywny
\begin{equation}\label{eq:8-1}
 T_0(x)=c+U(x-c)=Ux+(c-Uc).
\end{equation}
Jest to obrót wokół prostej $c+\mathbb R(0,0,1)$; zachowuje okrąg $\Gamma$.
\begin{lemat}\label{lem:wchloniecie-srodka}
Punkty $0,T_0(0),T_0^2(0),\ldots$ są parami różne. W konsekwencji $B^3\equiv B^3_*$ za pomocą dwóch części i ruchów sztywnych zachowujących orientację.
\end{lemat}
\begin{proof}
Gdyby $T_0^m(0)=T_0^n(0)$ dla $m>n$, to $T_0^{m-n}(0)=0$. Punkt $0$ leży na okręgu o dodatnim promieniu, więc wraca na swoje miejsce dokładnie wtedy, gdy kąt obrotu jest całkowitą wielokrotnością $2\pi$. Dostajemy $(m-n)\alpha=2\pi k$, a więc $(m-n)\sqrt2=k$ dla pewnego $k\in\mathbb Z$. To sprzeczność z niewymiernością $\sqrt2$.

Dla kompletności: gdyby $\sqrt2=p/q$ w postaci nieskracalnej, to $p^2=2q^2$, więc $p$ byłoby parzyste. Zapis $p=2r$ dawałby $q^2=2r^2$, więc także $q$ byłoby parzyste, wbrew nieskracalności.

Niech $M=\{T_0^n(0):n\ge0\}$. Wszystkie te punkty leżą na $\Gamma\subset B^3$. Stosujemy lemat~\ref{lem:lancuch} dla $D=\{0\}$, $E=B^3$, $T=T_0$. Częściami są $M$ i $B^3\setminus M$, a ruchami $T_0$ i identyczność. Ich obrazy to $M\setminus\{0\}$ oraz $B^3\setminus M$, czyli podział $B^3_*$.
\end{proof}
\begin{uwaga}
Nie twierdzimy, że $T_0(B^3)=B^3$. Nie jest to potrzebne. Obracamy wyłącznie kawałek $M$, który leży na zachowywanym okręgu. Reszta kuli nie jest poruszana. Dowolny ruch nie musi zachowywać całej kuli, tylko poprawnie przenosić przypisany mu kawałek.
\end{uwaga}

\subsection{Złożenie wszystkich etapów}
\begin{proof}[Dowód twierdzenia~\ref{tw:banach-tarski}]
Lemat~\ref{lem:wchloniecie-srodka} i wniosek~\ref{wn:kula-bez-srodka} dają $B^3\equiv B^3_*\equiv2B^3_*\equiv2B^3$. Przechodniość została udowodniona przez jawne rozdrobnienie części w lemacie~\ref{lem:skladanie}. Nie używamy więc żadnego nieskończonego składania ruchów.

Podajmy również bezpośredni wzór na końcowe części. Zanumerujmy niepuste części rozkładu sfery jako $H_v$, ich obroty jako $q_v$, a etykiety jako $\ell(v)\in\{1,2\}$, gdzie $v$ przebiega zbiór najwyżej $16$ indeksów. Niech $\widehat H_v$ oznacza ich stożki promieniowe. Przyjmijmy $J_0=M$, $J_1=B^3\setminus M$, $t_0=T_0$, $t_1=\mathrm{id}$. Bijekcja $k:B^3\to B^3_*$ działa na $J_s$ jako $t_s$. Zbiory $t_0J_0$ i $t_1J_1$ tworzą podział $B^3_*$.

Dla $s,t\in\{0,1\}$ i indeksu $v$ definiujemy
\begin{equation}\label{eq:8-2}
 P_{svt}=\{x\in J_s:t_sx\in\widehat H_v,\quad q_vt_sx\in t_tJ_t\}.
\end{equation}
Na tym zbiorze złożenie bijekcji $k$, bijekcji wyznaczonej przez paradoksalny rozkład $B^3_*$ i odwrotnej bijekcji $k^{-1}$ ma postać jednego ruchu
\begin{equation}\label{eq:8-3}
 x\longmapsto t_t^{-1}q_vt_sx,
\end{equation}
z etykietą $\ell(v)$. Dla każdego $x\in B^3$ istnieje jedyne $s$ z $x\in J_s$, jedyne $v$ z $t_sx\in\widehat H_v$ i jedyne $t$ określające, w której części $t_tJ_t$ leży $q_vt_sx$. To dowodzi, że \eqref{eq:8-2} jest podziałem całej $B^3$.

Ponieważ każdy etap jest bijekcją między odpowiednimi zbiorami etykietowanymi, obrazy pod ruchami \eqref{eq:8-3} tworzą podział dwóch formalnych kul. Liczba części jest nie większa niż $2\cdot16\cdot2=64$. Nie ma dodatkowego czynnika $2$ za etykiety: każdy indeks $v$ ma już jedną, określoną etykietę docelową.

Teraz umieszczamy te kopie w przestrzeni. Dla etykiety $1$ dodajemy $v_1=(-3,0,0)$, a dla etykiety $2$ dodajemy $v_2=(3,0,0)$. Ostateczny ruch dla części $P_{svt}$ to $f_{svt}(x)=t_t^{-1}q_vt_sx+v_{\ell(v)}$. Każde $t_s$ jest ruchem $x\mapsto Qx+v$ z $Q\in\mathrm{SO}(3)$, a każde $q_v$ obrotem. Odwrotności i złożenia takich ruchów znów mają tę postać: jeśli $f(x)=Qx+v$, $g(x)=Rx+w$, to $g(f(x))=RQx+(Rv+w)$ oraz $f^{-1}(x)=Q^{-1}x-Q^{-1}v$.

Zatem wszystkie $f_{svt}$ są dozwolonymi ruchami sztywnymi. Ich obrazy wypełniają dokładnie dwie wskazane kule i są parami rozłączne: wewnątrz każdej kuli wynika to z bijekcji, a między kulami z ich rozłącznego położenia. To kończy dowód.
\end{proof}

\subsection{Kula o innym środku i promieniu}
\begin{wniosek}\label{wn:dowolna-kula}
Każdą kulę domkniętą o promieniu $r>0$ można podzielić na skończenie wiele części, a następnie za pomocą ruchów sztywnych zachowujących orientację złożyć z nich dwie rozłączne kule tego samego promienia.
\end{wniosek}
\begin{proof}
Dla kuli o środku $c$ użyjmy pomocniczego podobieństwa $S(x)=c+rx$ do zdefiniowania części $S(P_i)$ z dowodu dla kuli jednostkowej. Jeśli ruch jednostkowy miał postać $f_i(x)=Q_ix+v_i$, to $Sf_iS^{-1}(y)=Q_iy+(c-Q_ic+rv_i)$. To nadal izometria: liniowa część wynosi $Q_i$, nie $rQ_i$. Podobieństwo służy tylko przeniesieniu opisu konstrukcji; rzeczywiste kawałki nowej kuli poruszamy powyższymi izometriami, bez skalowania.
\end{proof}

\section{Dlaczego nie przeczy to objętości?}
\subsection{Miara jest określona tylko na pewnych zbiorach}
Miara Lebesgue'a w $\mathbb R^3$ przypisuje zbiorom mierzalnym objętość $\lambda_3(E)\in[0,\infty]$. Do poniższego argumentu potrzebne są tylko jej standardowe własności:
\begin{itemize}
\item kule są mierzalne i $0<\lambda_3(B^3)<\infty$;
\item izometrie zachowują mierzalność i objętość;
\item skończona suma rozłącznych zbiorów mierzalnych ma objętość równą sumie ich objętości.
\end{itemize}
W standardowej normalizacji $\lambda_3(B^3)=4\pi/3$, ale konkretna liczba nie jest istotna. Twierdzenie o istnieniu i własnościach miary Lebesgue'a jest tu wyłącznie tłem interpretacyjnym; główny dowód w poprzednich sekcjach nie korzystał z teorii miary.
\begin{stwierdzenie}[Co najmniej jeden kawałek jest niemierzalny]\label{st:niemierzalnosc}
W rozkładzie z twierdzenia~\ref{tw:banach-tarski} co najmniej jeden kawałek jest niemierzalny w sensie Lebesgue'a.
\end{stwierdzenie}
\begin{proof}
Załóżmy przeciwnie, że wszystkie $P_i$ są mierzalne. Ponieważ tworzą podział $B^3$, mamy $\lambda_3(B^3)=\sum_{i=1}^{N}\lambda_3(P_i)$. Obrazy $f_i(P_i)$ też są mierzalne, mają te same objętości i tworzą podział dwóch rozłącznych kul przystających do $B^3$. Zatem
\[
 2\lambda_3(B^3)=\sum_{i=1}^{N}\lambda_3(f_i(P_i))=
 \sum_{i=1}^{N}\lambda_3(P_i)=\lambda_3(B^3).
\]
Jest to niemożliwe dla dodatniej skończonej liczby $\lambda_3(B^3)$. Sprzeczne było założenie o mierzalności wszystkich części, a nie twierdzenie Banacha--Tarskiego.
\end{proof}
\begin{uwaga}
Sformułowanie „objętość wzrosła dwukrotnie podczas obracania kawałków” jest mylące. Co najmniej jednemu kawałkowi nie można przypisać objętości Lebesgue'a, więc nie można zsumować objętości wszystkich części. Nie mamy prawa zastosować addytywności miary do dowolnych podzbiorów tylko dlatego, że ich suma jest kulą.
\end{uwaga}

\subsection{Co jest, a co nie jest konstruktywne?}
Generatory obrotów są podane jawnymi macierzami. Cztery zbiory słów są zdefiniowane jawnie. Wzory na składanie kawałków też są jawne. Jednak zbiór reprezentantów $T$ został wybrany aksjomatem wyboru. Nie otrzymaliśmy algorytmu, który dla dowolnego punktu kuli rozstrzyga, do którego kawałka należy.

Nie twierdzimy też, że twierdzenie Banacha--Tarskiego jest równoważne pełnemu aksjomatowi wyboru ani że każdy fragment użytego wyboru jest logicznie niezbędny. Wykazaliśmy konkretnie: w ZFC istnieją zbiory i ruchy o opisanych własnościach. Taka jest treść tego dowodu.

Kula jako model ciągłego zbioru punktów nie jest fizycznym przedmiotem zbudowanym z atomów. Rozkład nie opisuje procedury możliwej do wykonania nożem ani technologii powielania materii. Kawałki nie muszą mieć wnętrza ani regularnego brzegu.

\section{Kontrola dowodu i najczęstsze pułapki}
\subsection{Łańcuch zależności}
Każda strzałka w poniższym schemacie została w tekście uzasadniona:
\[
\begin{gathered}
\text{słowa zredukowane}\xrightarrow{\text{tw. \ref{tw:redukcja}, stw. \ref{st:grupa-slow}}}F_2
\xrightarrow{\text{tw. \ref{tw:cztery-kawalki}}}2F_2,\\
F_2\xrightarrow{\text{tw. \ref{tw:wolne-obroty}}}G\le\mathrm{SO}(3)
\xrightarrow{\text{lem. \ref{lem:wyjatkowy}, wybór }T}X,\\
X\xrightarrow{\text{tw. \ref{tw:paradoksalnosc-x}}}2X,
\qquad S^2\xrightarrow{\text{lem. \ref{lem:lancuch}, \ref{lem:obrot-rozdzielajacy}}}X,\\
S^2\xrightarrow{\text{tw. \ref{tw:pelna-sfera}}}2S^2,\\
\text{stożki (lem. \ref{lem:stozki}):}\qquad B^3_*\equiv2B^3_*,\\
B^3\xrightarrow{\text{lem. \ref{lem:wchloniecie-srodka}}}B^3_*
\equiv2B^3_*\equiv2B^3
\xrightarrow{\text{przesunięcia}}\text{dwie rozłączne kule w }\mathbb R^3.
\end{gathered}
\]
\subsection{Dziesięć miejsc wymagających szczególnej uwagi}
\begin{enumerate}
\item \textbf{Redukcja nie polega na przestawianiu liter.} Usuwamy tylko sąsiadujące odwrotności. Jednoznaczność wyniku została udowodniona, a nie przyjęta intuicyjnie.
\item \textbf{Słowo puste musi należeć do podziału.} Włącza je zbiór $P$ w konstrukcji \eqref{eq:4-3}.
\item \textbf{Dwa niekomutujące obroty nie wystarczają.} Rachunek modulo $3$ wyklucza wszystkie niepuste słowa zredukowane.
\item \textbf{Nie redukujemy ułamka $1/3$ modulo $3$.} Redukujemy całkowite macierze liczników $N_s$.
\item \textbf{Grupa wolna to nie działanie swobodne.} Po usunięciu punktów ze zbioru $D$ wszystkie pozostałe stabilizatory są równe $\{e\}$.
\item \textbf{Wybór reprezentantów nie jest automatyczny.} Wskazaliśmy użycie aksjomatu wyboru przed lematem~\ref{lem:wspolrzedne-orbitowe}.
\item \textbf{Przeliczalne wyjątki nie znikają same.} Wchłania je dwuczęściowy rozkład oparty na obrocie spoza $G$.
\item \textbf{Kawałek złożony z nieskończenie wielu poziomów to nadal jeden kawałek.} Na całej sumy poziomów $L$ działa jeden obrót $R$.
\item \textbf{Środek kuli wymaga osobnego ruchu.} Obroty wokół $0$ nigdy go nie przemieszczą; użyliśmy obrotu wokół osi przesuniętej w \eqref{eq:8-1}.
\item \textbf{Formalne kopie nie są jeszcze dwiema kulami w przestrzeni.} Dopiero końcowe przesunięcia o $v_1,v_2$ realizują etykiety jako rzeczywiste rozłączne położenia.
\end{enumerate}
\subsection{Przykłady zastosowania konstrukcji}
\begin{enumerate}
\item \textbf{Do której części należy $a^5$?} Należy do $P$, więc do $A_2$, a nie do $A_1$. Po ruchu $a$ daje $a^6$ w pierwszej kopii. Źródłem $e$ w tej kopii jest $\bar a$.
\item \textbf{Dlaczego z $hgt=hku$ wynika $g=k$ i $t=u$?} Najpierw skracamy działanie $h$, otrzymując $gt=ku$. Punkty $t,u$ reprezentują tę samą orbitę, zatem są równe. Wtedy $(k^{-1}g)t=t$ i swoboda działania daje $k^{-1}g=e$. Bez swobody ostatni wniosek nie byłby uzasadniony.
\item \textbf{Gdzie w rachunku macierzowym wykrywa się redukcję?} Dokładnie w zerowych wpisach tabeli \eqref{eq:5-5}. Sąsiednia para $ss^{-1}$ daje czynnik $w_s^Tw_{s^{-1}}=0$. Brak takich par gwarantuje, że wszystkie czynniki skalarne w \eqref{eq:5-7} są niezerowe.
\item \textbf{Dlaczego nie wystarczy policzyć samej liczby punktów?} Sama bijekcja $B^3\to2B^3$ nie musi być realizowana skończonym rozkładem, w którym każdy kawałek przenosimy ruchem sztywnym. W twierdzeniu istotne są jednocześnie: skończona liczba kawałków, rozłączność i ograniczenie do izometrii.
\item \textbf{Czy możemy powtórzyć podwojenie?} Tak. Po pierwszym podwojeniu stosujemy ten sam rozkład do każdej kuli osobno i otrzymujemy cztery kopie, potem osiem itd. Każda ustalona liczba powtórzeń nadal daje skończony rozkład, bo mnożymy tylko skończone liczby kawałków. Nie twierdzimy, że graniczna procedura z nieskończenie wielu powtórzeń jest skończonym rozkładem w sensie definicji~\ref{def:rownorozklad}.
\end{enumerate}

\subsection{Uwagi bibliograficzne}
Przedstawiony tekst jest samodzielnym opracowaniem klasycznego argumentu: grupa wolna, swobodne działanie na części sfery, wchłanianie przeliczalnych wyjątków i przedłużenie promieniowe. Nie zakłada ogólniejszej postaci twierdzenia o równoważności dowolnych ograniczonych zbiorów o niepustym wnętrzu. Udowodniona tutaj wersja to podwojenie kuli domkniętej. Szerszy kontekst i inne warianty rozkładów omawiają:
\begin{itemize}
\item Stefan Banach, Alfred Tarski, \emph{Sur la décomposition des ensembles de points en parties respectivement congruentes}, Fundamenta Mathematicae \textbf{6} (1924), 244--277, \href{https://doi.org/10.4064/fm-6-1-244-277}{DOI: 10.4064/fm-6-1-244-277}.
\item Stan Wagon, \emph{The Banach--Tarski Paradox}, Cambridge University Press, 1985, \href{https://www.cambridge.org/core/books/the-banach-tarski-paradox/EA35091E69D798AE0FAE5A4E8386DD08}{katalog wydawcy}.
\end{itemize}

% AA-08 | Algebra abstrakcyjna | Iteracja IV, uzupelnienie definicji modulow
\setlength{\parskip}{0pt}
\chapter{Grupy abelowe, torsja i iloczyn tensorowy}
\label{aa:ch:grupy-abelowe}

\section{Od grup abelowych do modułów}
\label{aa:sec:wprowadzenie-moduly}

\subsection{Dlaczego potrzebujemy nowego pojęcia?}
Do tej pory badaliśmy grupy, których podstawową strukturą jest jedno
odwracalne działanie. W grupie abelowej działanie zapisujemy addytywnie,
a zatem możemy nie tylko dodawać elementy, lecz również tworzyć ich
całkowite wielokrotności: $a+a+a=3a$, $-a=(-1)a$.
Widać tu zaczątek \emph{mnożenia elementu grupy przez skalar}.
W algebrze liniowej skalary pochodzą z ciała, na przykład z $\mathbb R$;
jednak nie zawsze potrzebujemy dzielenia przez każdy niezerowy skalar.
Chcemy więc zachować dodawanie wektorów i zgodność ze skalarami,
ale dopuścić, aby skalary pochodziły z ogólnego pierścienia.
Właśnie to opisuje pojęcie modułu.

Przypomnijmy: pierścień $R$ z jedynką jest zbiorem z dodawaniem
i mnożeniem, w którym $(R,+)$ jest grupą abelową, mnożenie
jest łączne oraz rozdzielne względem dodawania, a element $1_R$
jest jedynką mnożenia. Nie zakładamy tutaj przemienności mnożenia.
Pełną teorię pierścieni rozwiniemy w kolejnych rozdziałach.
Taki krótki opis wystarcza do sformułowania definicji modułu.

\subsection{Definicja modułu}
\begin{definicja}[Lewy moduł nad pierścieniem]
\label{aa:def:module-introduction}
Niech $R$ będzie \emph{pierścieniem z jedynką} $1_R$,
niekoniecznie przemiennym. \emph{Lewym modułem nad $R$}
(nazywanym także \emph{lewym $R$-modułem}) nazywamy
\textbf{grupę abelową} $(M,+,0_M)$ wraz z odwzorowaniem
\[
  R\times M\longrightarrow M,\qquad (r,m)\longmapsto rm,
\]
zwanym \emph{mnożeniem przez skalar}, które dla wszystkich
$r,s\in R$ i $m,n\in M$ spełnia cztery warunki:
\begin{enumerate}[label=\textup{(M\arabic*)}]
\item $r(m+n)=rm+rn$;
\item $(r+s)m=rm+sm$;
\item $(rs)m=r(sm)$;
\item $1_Rm=m$.
\end{enumerate}
Elementy $R$ nazywamy \emph{skalarami}, a elementy $M$
\emph{elementami modułu}. Gdy pierścień jest przemienny,
możemy zwykle mówić po prostu o $R$-module.
\end{definicja}
Homomorfizm lewych $R$-modułów $f:M\to N$ oznacza homomorfizm
grup addytywnych spełniający $f(rm)=rf(m)$ dla wszystkich
$r\in R$ i $m\in M$.

\begin{uwaga}[Co jest istotne w definicji?]
\label{aa:uw:module-crucial}
\textbf{Warunek przemienności dodawania w $M$ jest częścią definicji.}
Nie wystarczy dowolna grupa z dodatkową operacją: dodawanie
w module musi być łączne, przemienne, mieć zero i odwrotności.
Pierścień $R$ nie musi być ciałem ani nawet być przemienny.
Jeśli $R$ nie jest przemienny, rozróżniamy moduły lewe i prawe:
w prawym module mnożenie ma postać $M\times R\to M$.
W dalszej części tego rozdziału wystarczy nam przemienny
pierścień $\mathbb Z$.
\end{uwaga}

\begin{lemat}[Pierwsze konsekwencje aksjomatów]
\label{aa:lem:module-zero-minus}
Niech $R$ będzie pierścieniem z jedynką, a $M$ lewym
$R$-modułem. Dla wszystkich $r\in R$ i $m\in M$ zachodzą
\[
 0_Rm=0_M,\qquad r0_M=0_M,\qquad (-r)m=-(rm),
 \qquad r(-m)=-(rm).
\]
\end{lemat}
\begin{proof}
Z warunku (M2) mamy $0_Rm=(0_R+0_R)m=0_Rm+0_Rm$.
Odejmując $0_Rm$ w grupie addytywnej $M$, dostajemy
$0_Rm=0_M$. Podobnie (M1) daje
$r0_M=r(0_M+0_M)=r0_M+r0_M$, a więc $r0_M=0_M$.
Teraz $(-r)m+rm=((-r)+r)m=0_Rm=0_M$, toteż
$(-r)m=-(rm)$. Analogicznie
$r(-m)+rm=r((-m)+m)=r0_M=0_M$,
co dowodzi ostatniej równości.
\end{proof}

\begin{przyklad}[Przestrzeń wektorowa jest modułem]
Jeśli $K$ jest ciałem, a $V$ przestrzenią wektorową nad $K$,
to $V$ jest $K$-modułem: aksjomaty modułu są dokładnie
odpowiednimi aksjomatami przestrzeni wektorowej.
W szczególności $\mathbb C$ jest jednocześnie modułem nad
$\mathbb C$ oraz, po ograniczeniu skalarów, nad $\mathbb R$.
Pierwsza z tych struktur ma jedną współrzędną zespoloną,
druga --- dwie współrzędne rzeczywiste.
Wybór pierścienia skalarów jest więc częścią danych, a nie
jedynie zmianą zapisu.
\end{przyklad}

\begin{przyklad}[Moduły, które nie są przestrzeniami wektorowymi]
Grupa $\mathbb Z/6\mathbb Z$ jest modułem nad $\mathbb Z$.
Jej element $[2]$ nie jest zerowy, lecz
$3[2]=[6]=[0]$. Nie możemy z tego wywnioskować
$[2]=[0]$, bo skalar $3$ nie jest odwracalny w $\mathbb Z$.
Takie zjawisko nazywamy \emph{torsją} i będziemy je badać
w tym rozdziale. Dla porównania, w przestrzeni wektorowej
nad ciałem $K$ z $r\ne0$ z równania $rv=0$ wynika $v=0$,
gdyż można pomnożyć je przez $r^{-1}$.
\end{przyklad}

\subsection[Dlaczego grupa abelowa jest modułem nad liczbami całkowitymi?]{Dlaczego każda grupa abelowa jest modułem nad $\mathbb Z$?}
Przyjmujemy teraz \emph{notację addytywną} w grupie abelowej $A$:
jej element neutralny oznaczamy przez $0$, a odwrotność $a$
przez $-a$. Mnożenie przez liczbę całkowitą nie wymaga
żadnego dodatkowego wyboru; jest już określone przez
strukturę grupy.

\begin{definicja}[Grupa abelowa i działanie liczb całkowitych]
\label{aa:def:abelian-zmodule}
Niech $(A,+,0)$ będzie grupą abelową.
Dla $n\in\mathbb Z$ i $a\in A$ definiujemy $na$ jako sumę
$n$ kopii $a$, jeśli $n>0$, jako $0$ dla $n=0$ oraz jako
$-((-n)a)$ dla $n<0$.
\end{definicja}
\begin{stwierdzenie}[Tożsamości modułowe]
\label{aa:st:abelian-module-laws}
Niech $A$ będzie grupą abelową, a mnożenie przez liczby całkowite
niech będzie określone jak w definicji~\ref{aa:def:abelian-zmodule}.
Dla wszystkich $m,n\in\mathbb Z$ i $a,b\in A$ zachodzą
\[
(m+n)a=ma+na,\qquad m(a+b)=ma+mb,\qquad
(mn)a=m(na),\qquad 1a=a.
\]
Działanie z poprzedniej definicji czyni $A$ modułem nad
$\mathbb Z$. Jest ono \emph{jedynym} możliwym działaniem
$\mathbb Z$ spełniającym aksjomaty modułu. Ponadto
homomorfizmy grup abelowych są dokładnie homomorfizmami
$\mathbb Z$-modułów.
\end{stwierdzenie}
\begin{proof}
Dla nieujemnych $m,n$ pierwsza i trzecia równość wynikają
z łączności dodawania, a druga z łączności i przemienności:
na przykład $2(a+b)=a+b+a+b=a+a+b+b=2a+2b$.
Definicja wielokrotności ujemnych pozwala przeprowadzić
te same rachunki po zastąpieniu składników ich odwrotnościami.
Dla przykładu, jeśli $m\ge n>0$, to
$m a+(-n)a=(m-n)a$, gdyż w sumie $m$ kopii $a$
i $n$ kopii $-a$ możemy usunąć $n$ par przeciwnych.
Pozostałe przypadki znaków otrzymujemy analogicznie.
Równość $1a=a$ wynika wprost z definicji.

Jeżeli $A$ byłoby już $\mathbb Z$-modułem z jakimś działaniem,
to z (M4) i (M2) musielibyśmy mieć
$n a=(1+\cdots+1)a=a+\cdots+a$ dla $n>0$,
a także $0a=0$ i $(-n)a=-(na)$ na mocy lematu
\ref{aa:lem:module-zero-minus}. Działanie jest zatem
jednoznacznie wyznaczone przez dodawanie w $A$.

Jeśli $f:A\to B$ jest homomorfizmem grup abelowych,
to zachowuje sumy, zero i odwrotności. Stąd dla każdego
$n\in\mathbb Z$ mamy $f(na)=nf(a)$; $f$ jest więc
homomorfizmem $\mathbb Z$-modułów. Odwrotnie,
homomorfizm modułów z definicji zachowuje dodawanie,
więc jest homomorfizmem grup abelowych.
\end{proof}

\begin{przyklad}[Obliczenia w module nad $\mathbb Z$]
W $A=\mathbb Z/6\mathbb Z$ mamy
$(-2)[4]=[-8]=[4]$ oraz $3[2]=[0]$.
W $A=\mathbb Q/\mathbb Z$ element $[1/5]$ ma rząd $5$,
bo $5[1/5]=[1]=[0]$, lecz żadna z liczb
$1,2,3,4$ nie zeruje tej klasy. Obie grupy są
$\mathbb Z$-modułami, choć pierwsza jest skończona,
a druga nieskończona.
\end{przyklad}

\begin{uwaga}[Co zrobimy dalej z modułami?]
\label{aa:uw:module-roadmap}
Na tym etapie potrzebujemy tylko definicji modułu i faktu,
że \textbf{grupy abelowe to dokładnie moduły nad $\mathbb Z$}.
Dlatego możemy później zdefiniować iloczyn tensorowy grup
abelowych jako $A\otimes_{\mathbb Z}B$. Pełną teorię
modułów --- podmoduły, moduły wolne i ilorazowe, bazy,
homomorfizmy, ciągi dokładne, moduły nad pierścieniami
wielomianów oraz klasyfikację nad PID --- opracujemy
w osobnej części poświęconej modułom.
Nie należy utożsamiać iloczynu tensorowego grup abelowych
z odrębnymi konstrukcjami zwanymi \emph{nieabelowymi
iloczynami tensorowymi grup}.
\end{uwaga}

\section{Suma prosta, iloczyn prosty i skończone generowanie}
\begin{definicja}[Suma prosta rodziny grup abelowych]
\label{aa:def:abelian-directsum}
Niech $(A_i)_{i\in I}$ będzie rodziną grup abelowych.
Ich \emph{iloczyn prosty} $\prod_{i\in I}A_i$ składa się
ze wszystkich rodzin $(a_i)_{i\in I}$ z działaniem współrzędnym.
\emph{Suma prosta} $\bigoplus_{i\in I}A_i$ to podgrupa
tylko tych rodzin, dla których $a_i=0$ poza skończonym
zbiorem indeksów.
\end{definicja}
\begin{lemat}[Różnica pomiędzy sumą a iloczynem]
\label{aa:lem:abelian-sum-product}
Niech $(A_i)_{i\in I}$ będzie rodziną grup abelowych.
Jeśli zbiór indeksów $I$ jest skończony, to suma prosta
$\bigoplus_{i\in I}A_i$ jest równa iloczynowi
$\prod_{i\in I}A_i$ jako jego podgrupa.
Jeśli $I=\{1,2,3,\ldots\}$ i $A_i=\mathbb Z/2\mathbb Z$
dla każdego $i\in I$, to suma prosta jest przeliczalna,
a iloczyn prosty jest nieprzeliczalny.
\end{lemat}
\begin{proof}
Przy skończonym $I$ każdy element iloczynu ma co najwyżej
$|I|$ niezerowych współrzędnych, więc spełnia definicję sumy.
Przy $I=\mathbb N$ suma jest sumą przeliczalnie wielu
zbiorów skończonych: elementy mające nośnik w
$\{1,\ldots,n\}$ tworzą grupę o $2^n$ elementach.
Cała suma jest przeliczalna i nieskończona.
Dla iloczynu funkcje $\mathbb N\to\{0,1\}$ są jego elementami.
Gdyby można je było wypisać jako $x^{(1)},x^{(2)},\ldots$,
funkcja $y$ o $n$-tej współrzędnej przeciwnej do
$n$-tej współrzędnej $x^{(n)}$ różniłaby się od każdego
wyrazu listy. Jest to sprzeczność, więc iloczyn jest
nieprzeliczalny.
\end{proof}

\begin{definicja}[Grupa skończenie generowana]
\label{aa:def:abelian-fg}
Grupa abelowa $A$ jest \emph{skończenie generowana}, jeśli
istnieją $a_1,\ldots,a_r\in A$ takie, że każdy $a\in A$
można zapisać jako $\sum_{i=1}^r n_i a_i$,
gdzie $n_i\in\mathbb Z$.
\end{definicja}
\begin{lemat}[Opis przez grupę wolną abelową]
\label{aa:lem:abelian-fg-quotient}
Grupa abelowa $A$ jest skończenie generowana wtedy i tylko wtedy,
gdy istnieje $r\ge0$ oraz surjekcja $\mathbb Z^r\to A$.
Każda taka grupa jest izomorficzna z ilorazem
$\mathbb Z^r/K$ przez pewną podgrupę $K\le\mathbb Z^r$.
\end{lemat}
\begin{proof}
Jeśli $a_1,\ldots,a_r$ generują $A$, wzór
$f(n_1,\ldots,n_r)=\sum_i n_i a_i$ określa homomorfizm,
który jest surjektywny z definicji generowania.
Odwrotnie, obrazy standardowych wektorów bazowych
surjekcji $\mathbb Z^r\to A$ generują $A$.
Pierwsze twierdzenie o izomorfizmie grup zastosowane do
$f$ daje $A\cong\mathbb Z^r/\ker f$.
\end{proof}

\section{Rząd, wykładnik i podgrupa torsyjna}
\begin{definicja}[Torsja]
\label{aa:def:abelian-torsion}
Niech $A$ będzie grupą abelową.
Element $a\in A$ jest \emph{torsyjny}, jeśli $na=0$
dla pewnej liczby całkowitej $n>0$.
Zbiór wszystkich takich elementów oznaczamy
$\operatorname{Tors}(A)$ i nazywamy \emph{podgrupą torsyjną}.
Grupa jest \emph{torsyjna}, jeśli $\operatorname{Tors}(A)=A$,
a \emph{beztorsyjna}, jeśli $\operatorname{Tors}(A)=\{0\}$.
\end{definicja}
\begin{lemat}[Torsja jest podgrupą charakterystyczną]
\label{aa:lem:torsion-characteristic}
Dla każdej grupy abelowej $A$ zbiór $\operatorname{Tors}(A)$
jest podgrupą charakterystyczną $A$, a iloraz
$A/\operatorname{Tors}(A)$ jest beztorsyjny.
\end{lemat}
\begin{proof}
Jeśli $ma=0$ i $nb=0$, gdzie $m,n>0$, to
$mn(a-b)=n(ma)-m(nb)=0$; zbiór elementów torsyjnych
jest zatem zamknięty na różnice i zawiera $0$, więc jest
podgrupą. Każdy automorfizm zachowuje równości $na=0$,
stąd podgrupa torsyjna jest charakterystyczna.

Załóżmy, że klasa $a+\operatorname{Tors}(A)$ jest torsyjna
w ilorazie. Wtedy dla pewnego $n>0$ element $na$ należy
do $\operatorname{Tors}(A)$, czyli $m(na)=0$ dla pewnego $m>0$.
A zatem $a$ jest torsyjny i jego klasa jest zerowa.
Iloraz nie ma więc niezerowych elementów torsyjnych.
\end{proof}

\begin{definicja}[Wykładnik grupy]
\label{aa:def:abelian-exponent}
Niech $A$ będzie grupą abelową. Jej \emph{wykładnikiem} nazywamy najmniejszą
liczbę całkowitą $n>0$ taką, że $na=0$ dla wszystkich
$a\in A$, jeśli istnieje. Gdy takiej liczby nie ma,
mówimy, że grupa nie ma skończonego wykładnika.
Dla skończonej grupy abelowej wykładnik istnieje i jest
najmniejszą wspólną wielokrotnością rzędów jej elementów.
\end{definicja}
\begin{proof}
W skończonej grupie wszystkie elementy mają skończone rzędy.
Liczba $n$ zeruje każdy element dokładnie wtedy, gdy jest
podzielna przez rząd każdego z nich. Najmniejszą taką
liczbą jest z definicji ich najmniejsza wspólna wielokrotność.
\end{proof}

\begin{przyklad}[Torsja nie musi być skończona]
\label{aa:prz:torsion-infinite}
W $\mathbb Z\oplus\mathbb Z/12\mathbb Z$ podgrupa torsyjna
to $\{0\}\oplus\mathbb Z/12\mathbb Z$.
W $\mathbb Q/\mathbb Z$ każdy element jest torsyjny:
dla liczby wymiernej $a/b$ mamy
$b(a/b+\mathbb Z)=0+\mathbb Z$.
Jednak $\mathbb Q/\mathbb Z$ jest nieskończona i nie ma
skończonego wykładnika, bo zawiera elementy rzędu $n$
dla każdego $n\ge1$, na przykład $1/n+\mathbb Z$.
\end{przyklad}

\begin{uwaga}[Dlaczego przemienność jest potrzebna?]
W dowolnej grupie nieprzemiennej zbiór elementów skończonego
rzędu nie musi być podgrupą. W nieskończonej grupie
diedralnej generowanej przez dwa odbicia każdy z tych
odbiciowych generatorów ma rząd dwa, ale ich iloczyn
jest obrotem nieskończonego rzędu. Nie można więc
bez zmiany założeń przenosić definicji podgrupy torsyjnej
z grup abelowych na wszystkie grupy.
\end{uwaga}

\section{Części pierwsze i rozkład grupy torsyjnej}
\begin{definicja}[Podgrupa $p$-pierwotna]
\label{aa:def:abelian-primary}
Niech $A$ będzie grupą abelową, a $p$ liczbą pierwszą. Definiujemy
$A[p^\infty]=\{a\in A: p^k a=0\text{ dla pewnego }k\ge0\}$.
Jest to \emph{$p$-pierwotna podgrupa torsyjna} $A$.
Dla $k\ge1$ oznaczamy również
$A[p^k]=\{a\in A:p^k a=0\}$.
\end{definicja}

\begin{twierdzenie}[Rozkład torsji na części pierwsze]
\label{aa:tw:torsion-primary}
Dla każdej grupy abelowej $A$ jej podgrupa torsyjna jest
wewnętrzną sumą prostą podgrup $A[p^\infty]$.
Równoważnie, odwzorowanie sumowania wyznacza kanoniczny izomorfizm
\[
\bigoplus_{p\text{ -- liczba pierwsza}}A[p^\infty]
 \;\cong\;\operatorname{Tors}(A).
\]
W szczególności każda skończona grupa abelowa jest
iloczynem prostym swoich części $p$-pierwotnych.
\end{twierdzenie}
\begin{proof}
Weźmy element $a$ skończonego rzędu $n$.
Jeśli $n=1$, to $a=0$, więc otrzymujemy sumę pustą.
Dalej zakładamy $n>1$ i rozkładamy $n$ na potęgi różnych liczb pierwszych:
$n=n_1\cdots n_r$, gdzie $n_i=p_i^{k_i}$.
Dla każdego $i$ liczba $m_i=n/n_i$ jest względnie pierwsza
z $n_i$, więc z tożsamości Bézouta istnieją liczby całkowite
$c_i,d_i$ z $c_i m_i+d_i n_i=1$.
Połóżmy $e_i=c_i m_i$. Mamy $e_i\equiv1\pmod{n_i}$,
natomiast $e_i\equiv0\pmod{n_j}$ dla $j\ne i$.
Stąd $\sum_i e_i\equiv1\pmod {n_i}$ dla każdego $i$.
Ponieważ liczby $n_i$ są parami względnie pierwsze,
wielokrotna wersja lematu Euklidesa pokazuje, że
$n\mid(1-\sum_i e_i)$.
Wobec $na=0$ otrzymujemy
$a=\sum_i e_i a$.
Dla każdego $i$ zachodzi $n_i(e_i a)=0$:
liczba $n_i e_i$ jest podzielna przez $n$, bo $e_i$
jest podzielne przez wszystkie pozostałe $n_j$.
Zatem $e_i a\in A[p_i^\infty]$; udowodniliśmy, że
każdy element torsyjny jest sumą skończenie wielu elementów
z części pierwszych.

Sprawdźmy jednoznaczność. Jeśli
$\sum_{i=1}^r a_i=0$ dla $a_i\in A[p_i^\infty]$
i różnych liczb pierwszych $p_i$, dobierzmy potęgi
$n_i=p_i^{k_i}$ zerujące $a_i$. Mnożąc równość przez
$e_j$ z poprzedniej konstrukcji otrzymujemy
$0=e_j\sum_i a_i=a_j$, gdyż
$e_j\equiv1\pmod {n_j}$ i
$e_j\equiv0\pmod{n_i}$ dla $i\ne j$.
Wszystkie składniki są zatem zerowe.
Sumowanie jest więc izomorfizmem sumy prostej na
$\operatorname{Tors}(A)$. Jeśli $A$ jest skończona,
wszystkie jej elementy są torsyjne, a tylko skończenie
wiele liczb pierwszych dzieli $|A|$.
\end{proof}

\begin{przyklad}[Rozkład $C_{360}$]
\label{aa:prz:c360-primary}
Ponieważ $360=2^3\cdot3^2\cdot5$, mamy
$C_{360}\cong C_8\oplus C_9\oplus C_5$.
Przy generatorze $[1]_{360}$ podgrupę rzędu $8$
generuje $[45]_{360}$, podgrupę rzędu $9$ generuje
$[40]_{360}$, a podgrupę rzędu $5$ generuje
$[72]_{360}$. Suma podgrup rzędów $8$ i $9$ ma rząd
$8\cdot9=72$, ponieważ ich przekrój jest trywialny.
Przekrój tej sumy z podgrupą rzędu $5$ jest również trywialny,
ponieważ ich rzędy są względnie pierwsze. Wobec tego suma
trzech podgrup ma rząd $72\cdot5=360$ i jest całą grupą.
\end{przyklad}

\begin{wniosek}[Związek z podgrupami Sylowa]
\label{aa:wn:abelian-sylow}
Jeżeli $A$ jest skończoną grupą abelową, to dla każdej
liczby pierwszej $p$ dzielącej $|A|$ podgrupa $A[p^\infty]$
jest jej jedyną podgrupą $p$-Sylowa.
\end{wniosek}
\begin{proof}
Twierdzenie~\ref{aa:tw:torsion-primary} rozkłada $A$
na sumę skończonych $p$-grup. Rzędy tych czynników są
potęgami odpowiednich liczb pierwszych, a ich iloczyn
jest równy $|A|$. Stąd czynnik $p$-pierwotny ma rząd
największej potęgi $p$ dzielącej $|A|$ i jest podgrupą
$p$-Sylowa. Ponieważ $A$ jest abelowa, wszystkie jej podgrupy
są normalne; drugie twierdzenie Sylowa zapewnia, że
wszystkie podgrupy $p$-Sylowa są sprzężone, a zatem równe.
\end{proof}

\section{Skończone grupy abelowe: zarys klasyfikacji i niezmienniki}
\begin{uwaga}[Klasyfikacja skończonych grup abelowych --- zapowiedź twierdzenia]
\label{aa:tw:finite-abelian-classification}
Dla każdej skończonej grupy abelowej $A$ i każdej liczby
pierwszej $p$ istnieje jednoznacznie wyznaczona liczba całkowita
$r_p\ge0$ oraz, gdy $r_p>0$, jednoznacznie wyznaczone
(z dokładnością do kolejności) liczby całkowite
$e_{p,1}\ge\cdots\ge e_{p,r_p}\ge1$ takie, że
\[
A\cong\bigoplus_p\bigoplus_{j=1}^{r_p}C_{p^{e_{p,j}}}.
\]
Przy $r_p=0$ odpowiedni czynnik jest trywialny, a dla wszystkich
poza skończenie wieloma liczbami pierwszymi zachodzi $r_p=0$.
Powyższego twierdzenia nie wykorzystujemy w obecnym rozdziale
bez dowodu. Pełny dowód istnienia i jednoznaczności przeprowadzimy
w rozdziale o klasyfikacji skończenie generowanych
modułów nad dziedziną ideałów głównych, gdzie otrzymamy
zarazem postać Smitha i klasyfikację grup abelowych.
\end{uwaga}

\begin{lemat}[Co już można odczytać bez klasyfikacji]
\label{aa:lem:abelian-p-rank}
Niech $p$ będzie liczbą pierwszą, $r\ge0$ liczbą całkowitą,
$e_1,\ldots,e_r$ dodatnimi liczbami całkowitymi oraz
$A=\bigoplus_{j=1}^r C_{p^{e_j}}$.
Wtedy $A[p]\cong C_p^r$, a dla każdego $k\ge1$
liczba elementów z $A[p^k]$ wynosi
$p^{\sum_{j=1}^r\min(k,e_j)}$.
\end{lemat}
\begin{proof}
W grupie $C_{p^{e}}$ równanie $p^k x=0$ jest spełnione
dokładnie przez klasy wielokrotności $p^{\max(e-k,0)}$.
Ich liczba wynosi $p^{\min(k,e)}$.
W sumie prostej warunek $p^k x=0$ zachodzi dokładnie wtedy,
gdy jest spełniony w każdej współrzędnej. Mnożąc liczności
dostajemy wskazany wzór. Dla $k=1$ każdy czynnik daje
jedną kopię $C_p$, co dowodzi pierwszego izomorfizmu.
\end{proof}

\begin{przyklad}[Dwie grupy rzędu $16$]
\label{aa:prz:abelian-order16}
Grupy $C_8\oplus C_2$ oraz $C_4\oplus C_4$ mają ten sam
rząd $16$ i po cztery elementy zerowane przez $2$
(bo w obu przypadkach $A[2]\cong C_2^2$).
Nie są jednak izomorficzne: pierwsza zawiera element
rzędu $8$, druga nie. Wzór poprzedniego lematu daje też
$|A[4]|=8$ dla pierwszej grupy i $|A[4]|=16$ dla drugiej.
To pokazuje, że znajomość samego rzędu i liczby elementów
rzędu dwa nie wystarcza do rozpoznania grupy abelowej.
\end{przyklad}

\section{Grupy torsyjne i beztorsyjne poza przypadkiem skończonym}
\begin{definicja}[Podzielność i grupa Prüfera]
\label{aa:def:divisible-prufer}
Grupa abelowa $D$ jest \emph{podzielna}, gdy dla każdego
$n>0$ odwzorowanie $D\to D$, $x\mapsto nx$, jest surjektywne.
Dla liczby pierwszej $p$ \emph{grupą Prüfera} $C_{p^\infty}$
nazywamy podgrupę $\mathbb Q/\mathbb Z$ złożoną z klas
$\frac{a}{p^k}+\mathbb Z$, gdzie $a\in\mathbb Z$, $k\ge0$.
\end{definicja}
\begin{stwierdzenie}[Przykłady nieskończonej torsji]
\label{aa:st:prufer-properties}
Niech $p$ będzie liczbą pierwszą i niech
$C_{p^\infty}\le\mathbb Q/\mathbb Z$ będzie grupą Prüfera.
Dla każdego $k\ge1$ określmy
$H_k=\{a/p^k+\mathbb Z:a\in\mathbb Z\}$.
Wówczas $H_k\cong C_{p^k}$, zachodzą ścisłe zawierania
$H_1\subsetneq H_2\subsetneq H_3\subsetneq\cdots$ oraz
$C_{p^\infty}=\bigcup_{k\ge1}H_k$.
Grupa $C_{p^\infty}$ jest torsyjna, nieskończona i podzielna.
Ponadto $\mathbb Q/\mathbb Z\cong
\bigoplus_{q\text{ -- liczba pierwsza}} C_{q^\infty}$.
\end{stwierdzenie}
\begin{proof}
Klasa $a/p^k+\mathbb Z$ jest zerowana przez $p^k$,
więc jest torsyjna. Klasy $1/p^k+\mathbb Z$ mają kolejno
rzędy $p^k$, zatem są różne i grupa jest nieskończona.
Dla $n>0$ i klasy $a/p^k+\mathbb Z$ możemy rozwiązać
$nx=a/p^k+\mathbb Z$ wewnątrz tej grupy.
Zapiszmy $n=p^s m$, gdzie $\gcd(m,p)=1$.
Wybierzmy całkowite $u$ z $um\equiv a\pmod{p^{k+s}}$.
Wtedy $x=u/p^{k+s}+\mathbb Z$ spełnia
$n x=p^s m u/p^{k+s}+\mathbb Z
= m u/p^k+\mathbb Z=a/p^k+\mathbb Z$.
Istnienie $u$ wynika z tożsamości Bézouta.
Dla każdego $k\ge1$ zbiór klas $a/p^k+\mathbb Z$, gdzie
$a\in\mathbb Z$, jest podgrupą cykliczną rzędu $p^k$,
generowaną przez $1/p^k+\mathbb Z$: dwie takie klasy
są równe dokładnie wtedy, gdy ich liczniki są przystające
modulo $p^k$. Każdy element $C_{p^\infty}$ należy do jednej
z tych podgrup, a podgrupa dla $k$ zawiera się w podgrupie
dla $k+1$, gdyż $a/p^k=pa/p^{k+1}$.
Zawieranie jest ścisłe: $1/p^{k+1}+\mathbb Z$ ma rząd $p^{k+1}$,
natomiast każdy element $H_k$ jest zerowany przez $p^k$.
Dla $x\in\mathbb Q/\mathbb Z$ warunek $p^k x=0$ oznacza,
że istnieje reprezentant wymierny $q$ klasy $x$ taki, że
$p^kq\in\mathbb Z$. Wobec tego $q=a/p^k$ dla pewnego
$a\in\mathbb Z$, co dowodzi równości
$(\mathbb Q/\mathbb Z)[p^\infty]=C_{p^\infty}$.
Ostatni izomorfizm wynika teraz z
Twierdzenia~\ref{aa:tw:torsion-primary} zastosowanego
do grupy torsyjnej $\mathbb Q/\mathbb Z$.
\end{proof}

\begin{przyklad}[Beztorsyjność nie oznacza istnienia bazy nad $\mathbb Z$]
\label{aa:prz:torsionfree-notfree}
Grupa addytywna $\mathbb Z[1/2]=\{m/2^k:m\in\mathbb Z,k\ge0\}$
jest beztorsyjna jako podgrupa $\mathbb Q$.
Nie jest jednak wolnym $\mathbb Z$-modułem.
Gdyby była wolna na pewnej bazie, nie mogłaby ona mieć
dwóch różnych elementów: dla niezerowych liczb wymiernych
$x=a/b$ oraz $y=c/d$ zachodzi niebanalna relacja całkowita
$bc\,x-ad\,y=0$. Ponieważ grupa jest niezerowa,
baza musiałaby składać się z dokładnie jednego elementu
$x=m/2^k\ne0$.
Wówczas grupa byłaby cykliczna, generowana przez $x$.
Jednak element $x/2$ należy do $\mathbb Z[1/2]$, a nie jest
całkowitą wielokrotnością $x$ (wymagałoby to współczynnika
$1/2$). Sprzeczność dowodzi, że moduł nie jest wolny.
W przypadku przestrzeni wektorowych nad ciałem taka różnica
nie występuje: każda przestrzeń wektorowa ma bazę.
\end{przyklad}

\section{Iloczyn tensorowy grup abelowych}
\label{aa:sec:abelian-tensor}
Trzeba odróżniać trzy różne konstrukcje: iloczyn prosty
$A\times B$, iloczyn półprosty (zależny od działania grupy
na grupie) oraz iloczyn tensorowy $A\otimes_{\mathbb Z}B$.
Iloczyn prosty $A\times B$ składa się z par i ma działanie
wykonywane na każdej współrzędnej. Iloczyn tensorowy
$A\otimes_{\mathbb Z}B$ jest grupą abelową zbudowaną
z generatorów $a\otimes b$ i relacji wyrażających addytywność
w każdej zmiennej. Jego własność uniwersalna opisuje
odwzorowania dwuliniowe; nie należy traktować tych dwóch
konstrukcji jako zamiennych.

\begin{definicja}[Odwzorowanie dwuliniowe grup abelowych]
\label{aa:def:abelian-bilinear}
Niech $A,B,C$ będą grupami abelowymi. Odwzorowanie
$b:A\times B\to C$ jest \emph{dwuliniowe}, jeśli
$b(a+a',v)=b(a,v)+b(a',v)$ i
$b(a,v+v')=b(a,v)+b(a,v')$
dla wszystkich $a,a'\in A$, $v,v'\in B$.
\end{definicja}

\begin{twierdzenie}[Istnienie i własność uniwersalna iloczynu tensorowego]
\label{aa:tw:abelian-tensor-universal}
Dla każdej pary grup abelowych $A,B$ istnieje grupa abelowa
$A\otimes_{\mathbb Z}B$ i odwzorowanie dwuliniowe
$\otimes:A\times B\to A\otimes_{\mathbb Z}B$ takie, że
każde odwzorowanie dwuliniowe $b:A\times B\to C$ do grupy
abelowej $C$ ma jednoznaczną postać
$b=\widetilde b\circ\otimes$, gdzie
$\widetilde b:A\otimes_{\mathbb Z}B\to C$ jest homomorfizmem.
Para ta jest jednoznaczna z dokładnością do jedynego
izomorfizmu zgodnego z prostymi tensorami $a\otimes v$.
\end{twierdzenie}
\begin{proof}
Niech $F$ będzie wolną grupą abelową o bazie wszystkich
symboli $[a,v]$, gdzie $(a,v)\in A\times B$.
Niech $N\le F$ będzie podgrupą generowaną przez relacje
\[
[a+a',v]-[a,v]-[a',v],\qquad
[a,v+v']-[a,v]-[a,v'].
\]
Połóżmy $A\otimes_{\mathbb Z}B=F/N$ i
$a\otimes v=[a,v]+N$. Obie relacje dokładnie oznaczają
addytywność w każdej zmiennej. Dla dowolnego dwuliniowego
$b$ istnieje dokładnie jeden homomorfizm $B_0:F\to C$
określony na bazie przez $B_0([a,v])=b(a,v)$.
Równania dwuliniowości pokazują, że $N\subseteq\ker B_0$;
$B_0$ przechodzi więc na homomorfizm $\widetilde b:F/N\to C$.
Jego jednoznaczność wynika z generowania ilorazu przez
proste tensory $a\otimes v$.

Jeśli $(T,t)$ jest drugą parą o tej samej własności,
obie własności uniwersalne dają jedyne homomorfizmy
$f:T\to A\otimes_{\mathbb Z}B$ i $g:A\otimes_{\mathbb Z}B\to T$
spełniające $f\circ t=\otimes$ i $g\circ\otimes=t$.
Wówczas $(g\circ f)\circ t=t$ oraz
$(f\circ g)\circ\otimes=\otimes$.
Jednoznaczność faktoryzacji odpowiednio przez $t$ i przez
$\otimes$ daje $g\circ f=\mathrm{id}_T$ i
$f\circ g=\mathrm{id}_{A\otimes_{\mathbb Z}B}$.
\end{proof}

\begin{twierdzenie}[Iloczyny tensorowe grup cyklicznych]
\label{aa:tw:abelian-tensor-cyclic}
Dla grupy abelowej $B$ i liczby całkowitej $m\ge1$ istnieje
naturalny izomorfizm
\[
(\mathbb Z/m\mathbb Z)\otimes_{\mathbb Z} B\cong B/mB.
\]
W szczególności dla $m,n\ge1$ zachodzi
$C_m\otimes_{\mathbb Z}C_n\cong C_{\gcd(m,n)}$,
a ponadto $\mathbb Z\otimes_{\mathbb Z}B\cong B$.
\end{twierdzenie}
\begin{proof}
Niech $\bar1$ oznacza klasę $1$ modulo $m$.
Homomorfizm $\phi:B\to C_m\otimes_{\mathbb Z}B$, $b\mapsto\bar1\otimes b$,
jest surjektywny: dowolny prosty tensor
$\bar a\otimes b=a(\bar1\otimes b)=\bar1\otimes ab$.
Ponieważ $m\bar1=0$, mamy
$\phi(mb)=\bar1\otimes(mb)=m(\bar1\otimes b)=0$.
Zatem $mB\subseteq\ker\phi$ i powstaje surjekcja
$\overline\phi:B/mB\to C_m\otimes_{\mathbb Z}B$.

W drugą stronę zdefiniujmy
$\beta:C_m\times B\to B/mB$ przez
$\beta(\bar a,b)=ab+mB$.
Wzór jest poprawny, bo zastąpienie $a$ przez $a+km$
zmienia wynik o element $km b\in mB$;
jest też addytywny w obu zmiennych.
Własność uniwersalna iloczynu tensorowego daje homomorfizm
$\widetilde\beta:C_m\otimes_{\mathbb Z}B\to B/mB$.
Na generatorach sprawdzamy
$\widetilde\beta(\overline\phi(b+mB))=b+mB$
oraz $\overline\phi(\widetilde\beta(\bar a\otimes b))
=\bar1\otimes(ab)=\bar a\otimes b$.
Mapy są wzajemnie odwrotne. Izomorfizm jest naturalny
względem $B$: dla dowolnego homomorfizmu grup abelowych
$f:B\to B'$ homomorfizmy indukowane na ilorazach oraz
iloczynach tensorowych spełniają
$\overline f(b+mB)=f(b)+mB'$ i
$(\mathrm{id}_{C_m}\otimes f)(\bar1\otimes b)=\bar1\otimes f(b)$.
Przez konstrukcję $\overline\phi$ obie drogi od
$B/mB$ do $C_m\otimes_{\mathbb Z}B'$ mają więc tę samą
wartość na każdej klasie $b+mB$.

Dla $B=C_n$ podgrupa $mC_n$ ma indeks $\gcd(m,n)$:
obraz mnożenia przez $m$ w $C_n$ ma rząd
$n/\gcd(m,n)$, gdyż jądro ma $\gcd(m,n)$ elementów.
Iloraz jest cykliczny odpowiedniego rzędu, stąd wzór.
Dla $\mathbb Z\otimes_{\mathbb Z}B$ analogiczne pary map są dane przez
$n\otimes b\mapsto nb$ oraz $b\mapsto1\otimes b$.
\end{proof}

\begin{przyklad}[Iloczyn tensorowy nie jest iloczynem prostym]
\label{aa:prz:abelian-tensor-contrasts}
Mamy $C_2\otimes_{\mathbb Z}C_3=0$, choć oba czynniki
są niezerowe, zaś $C_2\times C_3\cong C_6$.
Natomiast $C_4\otimes_{\mathbb Z}C_6\cong C_2$,
bo $\gcd(4,6)=2$.
W przeciwieństwie do iloczynu prostego, iloczyn tensorowy
może więc być grupą trywialną nawet wtedy, gdy oba czynniki
są nietrywialne. Jest to zasadnicza różnica między tymi
dwoma rodzajami iloczynów.
Ogólną konstrukcję $M\otimes_R N$ dla modułów nad
pierścieniem $R$, a także jej dokładność i związki
z funktorami $\operatorname{Tor}$, opracujemy w części
poświęconej modułom.
\end{przyklad}

\setlength{\parskip}{4pt}

% AA-09 | Algebra abstrakcyjna | Iteracja V
\setlength{\parskip}{0pt}
\chapter{Pierścienie, ideały i homomorfizmy}
\label{aa:ch:pierscienie}

\section{Dlaczego jedno działanie już nie wystarcza?}
Grupa pozwala mówić o jednym działaniu i jego odwracalności.
Jednak w arytmetyce liczb całkowitych i w rachunku wielomianów
mamy \emph{dwa} działania: dodawanie i mnożenie. Zachowują się one
inaczej: każdą liczbę całkowitą można odjąć, lecz przez niezerową
liczbę całkowitą nie zawsze można podzielić z wynikiem całkowitym.
Z drugiej strony macierze można dodawać i mnożyć, ale mnożenie
macierzy na ogół nie jest przemienne. Szukamy struktury, która
obejmie wszystkie te przykłady i zachowa reguły rozdzielności.

\begin{przyklad}[Jedno równanie w trzech pierścieniach]
\label{aa:prz:pierscienie-rownanie}
Rozważmy równanie $x^2=1$. W pierścieniu $\mathbb Z$
ma ono dokładnie dwa rozwiązania: $1$ i $-1$. Istotnie,
$(x-1)(x+1)=0$; ponieważ z iloczynu dwóch liczb całkowitych
równego zeru wynika zerowość jednego czynnika, mamy $x=1$ albo
$x=-1$. W pierścieniu $\mathbb Z/8\mathbb Z$ rozwiązaniami
są klasy $[1],[3],[5],[7]$. Dla nieparzystej liczby $x=2k+1$
liczba $x^2-1=4k(k+1)$ jest podzielna przez $8$, gdyż jedna
z liczb $k,k+1$ jest parzysta. Klasy parzyste nie rozwiązują
równania, ponieważ ich kwadraty są parzyste. Wreszcie w pierścieniu
$\mathbb R[\varepsilon]/(\varepsilon^2)$ klasa wielomianu
$\varepsilon$ nie jest zerem (bo $\varepsilon\notin(\varepsilon^2)$,
co wynika ze stopni wielomianów), lecz jej kwadrat jest zerem.
Pojęcie pierścienia
musi zatem dopuszczać dzielniki zera i elementy nilpotentne.
\end{przyklad}

\section{Aksjomaty pierścienia i podstawowe konsekwencje}
\begin{definicja}[Pierścień z jedynką]
\label{aa:def:pierscien}
\emph{Pierścieniem z jedynką} nazywamy zbiór $R$ wyposażony
w działania $+:R\times R\to R$ oraz $\cdot:R\times R\to R$
takie, że:
\begin{enumerate}[label=\textup{(P\arabic*)}]
\item $(R,+)$ jest grupą abelową z elementem neutralnym $0_R$;
\item mnożenie jest łączne: $(ab)c=a(bc)$ dla wszystkich $a,b,c\in R$;
\item istnieje $1_R\in R$ takie, że $1_Ra=a1_R=a$ dla każdego $a\in R$;
\item $a(b+c)=ab+ac$ i $(a+b)c=ac+bc$ dla wszystkich $a,b,c\in R$.
\end{enumerate}
Pierścień jest \emph{przemienny}, jeśli $ab=ba$ dla wszystkich
$a,b\in R$. Nie zakładamy tego, jeśli nie zostało to wskazane.
Dopuszczamy pierścień zerowy $R=\{0\}$, w którym $1_R=0_R$;
we wszystkich twierdzeniach o dziedzinach i ciałach wykluczymy go jawnie.
\end{definicja}

\begin{lemat}[Rachunek znaków]
\label{aa:lem:pierscien-znaki}
Niech $R$ będzie pierścieniem z jedynką i niech $a,b\in R$.
Wtedy $a0_R=0_Ra=0_R$, $(-a)b=-(ab)=a(-b)$ oraz
$(-a)(-b)=ab$. Jeżeli $1_R=0_R$, to $R=\{0_R\}$.
\end{lemat}
\begin{proof}
Z rozdzielności otrzymujemy $a0_R=a(0_R+0_R)=a0_R+a0_R$.
Po dodaniu elementu odwrotnego do $a0_R$ dostajemy $a0_R=0_R$.
Analogicznie $0_Ra=0_R$. Następnie
$(-a)b+ab=((-a)+a)b=0_Rb=0_R$, więc $(-a)b=-(ab)$.
Z kolei $a(-b)+ab=a((-b)+b)=0_R$, zatem $a(-b)=-(ab)$.
Stosując pierwszą z tych równości do $-b$, dostajemy
$(-a)(-b)=-(a(-b))=-(-(ab))=ab$.
Jeżeli $1_R=0_R$, to każdy $a\in R$ spełnia
$a=1_Ra=0_Ra=0_R$, co dowodzi ostatniego zdania.
\end{proof}


\section{Konstrukcja pierścienia liczb całkowitych z liczb naturalnych}
\label{aa:sec:konstrukcja-z}
Wcześniej korzystaliśmy z liczb całkowitych jako znanego przykładu
pierścienia. Możemy teraz uzasadnić ich istnienie, nie zakładając
z góry, że odejmowanie liczb naturalnych jest zawsze możliwe.
W tej sekcji $\mathbb N_0=\{0,1,2,\ldots\}$ jest półpierścieniem
zwykłych liczb naturalnych, to znaczy ma łączne i przemienne dodawanie
z zerem, łączne i przemienne mnożenie z jedynką oraz rozdzielność.
Korzystamy z własności skracania dodawania w $\mathbb N_0$:
$a+t=b+t$ pociąga $a=b$; wynika ona z definicji dodawania
liczb naturalnych i można ją wykazać indukcyjnie względem $t$.

\begin{definicja}[Formalne różnice liczb naturalnych]
\label{aa:def:formal-differences}
Na $\mathbb N_0\times\mathbb N_0$ określamy relację
\[
 (a,b)\sim(c,d)\quad\Longleftrightarrow\quad a+d=b+c.
\]
Zbiór klas tej relacji oznaczymy przez $Z$; klasę pary $(a,b)$
zapisujemy $[a,b]$. Para $(a,b)$ ma wyrażać \emph{formalną różnicę}
$a-b$, lecz zapis ten jest na razie tylko intuicją: liczby ujemne
nie zostały jeszcze zdefiniowane.
\end{definicja}

\begin{lemat}[Relacja formalnych różnic]
\label{aa:lem:formal-differences-equivalence}
Niech $\sim$ oznacza warunek $(a,b)\sim(c,d)$ wtedy i tylko wtedy,
  gdy $a+d=b+c$ dla par z $\mathbb N_0\times\mathbb N_0$.
  Jest to relacja równoważności.
\end{lemat}
\begin{proof}
Dla każdej pary $(a,b)$ zachodzi $a+b=b+a$, więc relacja
jest zwrotna. Jeśli $a+d=b+c$, to $c+b=d+a$,
więc jest symetryczna. Wreszcie z $(a,b)\sim(c,d)$
oraz $(c,d)\sim(e,f)$ otrzymujemy
$a+d=b+c$ i $c+f=d+e$. Dodając $f$ do pierwszej
równości, a następnie podstawiając drugą, mamy
$a+d+f=b+c+f=b+d+e$. Dzięki skracaniu składnika $d$
dostajemy $a+f=b+e$, czyli $(a,b)\sim(e,f)$.
\end{proof}

\begin{definicja}[Działania na formalnych różnicach]
\label{aa:def:z-operations}
Niech $Z=(\mathbb N_0\times\mathbb N_0)/\sim$, gdzie
$(a,b)\sim(c,d)$ wtedy i tylko wtedy, gdy $a+d=b+c$.
Dla wszystkich $a,b,c,d\in\mathbb N_0$ określamy na $Z$
działania wzorami
\[
 [a,b]+[c,d]=[a+c,b+d],\qquad
 [a,b]\,[c,d]=[ac+bd,ad+bc].
\]
Pierwszy odpowiada dodawaniu różnic, a drugi tożsamości
$(a-b)(c-d)=ac+bd-(ad+bc)$, użytej jedynie jako
wskazówka przy wyborze wzoru.
\end{definicja}

\begin{lemat}[Działania są dobrze określone]
\label{aa:lem:z-well-defined}
Niech $Z=(\mathbb N_0\times\mathbb N_0)/\sim$, gdzie
$(a,b)\sim(c,d)$ wtedy i tylko wtedy, gdy $a+d=b+c$.
Dodawanie i mnożenie klas zadane wzorami z
Definicji~\ref{aa:def:z-operations} nie zależą od wyboru reprezentantów.
\end{lemat}
\begin{proof}
Niech $(a,b)\sim(a',b')$, a więc $a+b'=b+a'$.
Dla dowolnej pary $(c,d)$ otrzymujemy
$(a+c)+(b'+d)=(b+d)+(a'+c)$, co dowodzi
niezależności dodawania od reprezentanta pierwszego składnika.
Analogicznie postępujemy z drugim składnikiem.

Dla mnożenia należy wykazać
$(ac+bd,ad+bc)\sim(a'c+b'd,a'd+b'c)$.
Warunek ten oznacza równość
\[
 (ac+bd)+(a'd+b'c)=(ad+bc)+(a'c+b'd).
\]
Lewą stronę przestawiamy na $(a+b')c+(b+a')d$,
a prawą na $(b+a')c+(a+b')d$; są one równe, bo
$a+b'=b+a'$. Dowodzi to niezależności od
reprezentanta pierwszego czynnika. Ponieważ podany wzór
na mnożenie jest symetryczny względem par czynników,
ten sam argument stosuje się do drugiego czynnika.
Po dwóch kolejnych zmianach reprezentantów otrzymujemy
niezależność obu działań jednocześnie.
\end{proof}

\begin{twierdzenie}[Pierścień liczb całkowitych]
\label{aa:tw:integers-as-ring}
Niech $Z=(\mathbb N_0\times\mathbb N_0)/\sim$, gdzie
$(a,b)\sim(c,d)$ wtedy i tylko wtedy, gdy $a+d=b+c$.
Niech dla wszystkich $a,b,c,d\in\mathbb N_0$ zachodzą wzory
$[a,b]+[c,d]=[a+c,b+d]$ oraz
$[a,b][c,d]=[ac+bd,ad+bc]$.
Wówczas $Z$ jest niezerowym pierścieniem przemiennym z jedynką.
Jego elementem zerowym
jest $[0,0]$, jedynką $[1,0]$, a przeciwnym do $[a,b]$
jest $[b,a]$. Odwzorowanie
\[
 \iota:\mathbb N_0\longrightarrow Z,\qquad \iota(n)=[n,0]
\]
jest różnowartościowe oraz zachowuje dodawanie,
mnożenie, zero i jedynkę. Pierścień $Z$ utożsamiamy
odtąd z pierścieniem $\mathbb Z$.
\end{twierdzenie}
\begin{proof}
Najpierw sprawdzamy strukturę addytywną. Łączność,
przemienność i neutralność $[0,0]$ wynikają z
odpowiednich reguł dodawania w $\mathbb N_0$.
Przy tym
$[a,b]+[b,a]=[a+b,a+b]=[0,0]$,
bo $(a+b)+0=(a+b)+0$. Każda klasa ma więc element
przeciwny i $(Z,+)$ jest grupą abelową.

Przemienność mnożenia wynika ze wzoru i przemienności
mnożenia w $\mathbb N_0$. Dla łączności porównajmy
współrzędne obu iloczynów trzech par. Pierwsza współrzędna
$(([a,b][c,d])[e,f])$ wynosi
$ace+bde+adf+bcf$, a druga
$acf+bdf+ade+bce$. W iloczynie
$[a,b]([c,d][e,f])$ otrzymujemy dokładnie te same
składniki, jedynie w innej kolejności. Oba wyniki są więc
równe już jako pary. Ponadto
$[a,b][1,0]=[a,b]$. Wreszcie rozdzielność
mnożenia względem dodawania wynika z rachunku
\[
 [a,b]([c,d]+[e,f])
 =[a(c+e)+b(d+f),\,a(d+f)+b(c+e)]
 =[a,b][c,d]+[a,b][e,f].
\]
Rozdzielność z prawej wynika z przemienności mnożenia.
Zatem mamy przemienny pierścień z jedynką.

Elementy $[1,0]$ i $[0,0]$ nie są równe, bo równość
ich klas wymagałaby $1+0=0+0$ w $\mathbb N_0$.
Odwzorowanie $\iota$ zachowuje oba działania i elementy
wyróżnione bezpośrednio ze wzorów. Jeśli
$[m,0]=[n,0]$, to $m+0=0+n$, zatem $m=n$;
$\iota$ jest różnowartościowe.
\end{proof}

\begin{twierdzenie}[Własność uniwersalna dopełnienia grupowego]
\label{aa:tw:z-grothendieck-universal}
Niech $\mathbb Z=(\mathbb N_0\times\mathbb N_0)/\sim$, gdzie
$(a,b)\sim(c,d)$ wtedy i tylko wtedy, gdy $a+d=b+c$,
oraz niech $\iota:\mathbb N_0\to\mathbb Z$ będzie dane wzorem
$\iota(n)=[n,0]$. Niech $A$ będzie dowolną grupą abelową
zapisywaną addytywnie,
a $f:\mathbb N_0\to A$ odwzorowaniem spełniającym
$f(m+n)=f(m)+f(n)$ i $f(0)=0$. Istnieje dokładnie jeden
homomorfizm grup abelowych $\overline f:\mathbb Z\to A$
taki, że $\overline f\circ\iota=f$. Jest nim
\[
 \overline f([a,b])=f(a)-f(b).
\]
Co więcej, dla każdego przemiennego pierścienia $R$
z jedynką jedyny homomorfizm pierścieni
$\Phi:\mathbb Z\to R$ ma postać
$\Phi([a,b])=a1_R-b1_R$.
\end{twierdzenie}
\begin{proof}
Jeśli $(a,b)\sim(c,d)$, to $a+d=b+c$.
Po zastosowaniu addytywności $f$ dostajemy
$f(a)+f(d)=f(b)+f(c)$, stąd
$f(a)-f(b)=f(c)-f(d)$. Wzór nie zależy więc
od reprezentanta. Dodawanie klas przechodzi na dodawanie
w $A$, zatem $\overline f$ jest homomorfizmem,
a $\overline f([n,0])=f(n)$. Każda klasa jest różnicą
$[a,0]-[b,0]$, dlatego każdy homomorfizm spełniający
ten warunek musi przyjmować podane wartości. To dowodzi
jednoznaczności.

Jeśli $R$ jest pierścieniem z jedynką, homomorfizm
$\mathbb Z\to R$ musi wysłać $[1,0]$ na $1_R$, a stąd
$[n,0]$ na $n1_R$. Z poprzedniej części ma zatem
jedyną możliwą postać. Sprawdzenie mnożenia daje
\[
 (a1_R-b1_R)(c1_R-d1_R)
 =(ac+bd)1_R-(ad+bc)1_R,
\]
co jest obrazem produktu klas. Wzór zachowuje więc
również mnożenie i jedynkę. Jest homomorfizmem pierścieni.
\end{proof}

\begin{przyklad}[Rachunek na klasach formalnych różnic]
\label{aa:prz:z-pairs}
Klasy $[7,3]$ i $[4,0]$ są równe, ponieważ
$7+0=3+4$. Klasa $[2,9]$ odpowiada liczbie $-7$;
ściślej, $[2,9]=[0,7]$, bo $2+7=9+0$.
Mnożąc klasy $[2,5]$ i $[7,3]$, otrzymujemy
$[2,5][7,3]=[29,41]=[0,12]$, czyli $(-3)\cdot4=-12$.
Każda klasa ma postać $[n,0]$ lub $[0,n]$ dla pewnego
$n\in\mathbb N_0$: gdy $a\ge b$, piszemy $a=b+n$,
a gdy $b\ge a$, piszemy $b=a+n$ i stosujemy definicję
relacji $\sim$.
\end{przyklad}

\begin{uwaga}[Dlaczego to nie jest zwykły pierścień ilorazowy?]
\label{aa:uw:z-vs-quotient-ring}
Konstrukcja zaczyna się od półpierścienia $\mathbb N_0$
i relacji równoważności na parach liczb naturalnych;
nie dzielimy istniejącego pierścienia przez ideał.
W języku algebry ogólnej jest to \emph{dopełnienie grupowe}
przemiennego monoidu $(\mathbb N_0,+)$, a poprzednie
twierdzenie wyraża jego własność uniwersalną.
W późniejszej części skryptu skonstruujemy $\mathbb Q$
jako ciało ułamków pierścienia $\mathbb Z$.
Konstrukcję ciała $\mathbb R$ --- w szczególności uzasadnienie
jego zupełności i porządku --- pozostawiamy skryptowi analizy.
Tutaj, przyjmując ciało $\mathbb R$ za dane, skonstruujemy
$\mathbb C$ na dwa algebraiczne sposoby:
jako $\mathbb R[x]/(x^2+1)$ oraz na zbiorze $\mathbb R^2$,
a następnie udowodnimy ich izomorfizm.
\end{uwaga}

\begin{definicja}[Jednostki, dzielniki zera i elementy nilpotentne]
\label{aa:def:pierscien-jednostki}
Niech $R$ będzie pierścieniem z jedynką.
Element $u\in R$ jest \emph{jednostką}, jeśli ma odwrotność
obustronną $v$, tj. $uv=vu=1_R$. Zbiór jednostek oznaczamy $R^\times$.
W niezerowym pierścieniu przemiennym element $a\ne0_R$ jest
\emph{dzielnikiem zera}, jeśli istnieje $b\ne0_R$ z $ab=0_R$.
Element $a$ jest \emph{nilpotentny}, jeśli $a^n=0_R$ dla pewnego
$n\ge1$, a \emph{idempotentny}, jeśli $a^2=a$.
\end{definicja}
\begin{stwierdzenie}[Grupa jednostek]
\label{aa:st:unit-group}
Dla dowolnego pierścienia $R$ jego jednostki tworzą grupę
$R^\times$ względem mnożenia. W pierścieniu przemiennym
niezerowym jednostka nie jest dzielnikiem zera.
\end{stwierdzenie}
\begin{proof}
Dla jednostek $u,v$ zachodzą równości
$(uv)(v^{-1}u^{-1})=(v^{-1}u^{-1})(uv)=1_R$.
Wynika z nich, że $uv$ ma odwrotność obustronną,
zatem jest jednostką. Element $1_R$ należy do $R^\times$,
a wraz z każdą jednostką jej odwrotność także należy do $R^\times$.
Łączność mnożenia jest dziedziczona z pierścienia.
Jeśli $ua=0_R$, mnożenie przez $u^{-1}$ daje
$a=u^{-1}(ua)=0_R$, dlatego jednostka nie jest dzielnikiem zera.
\end{proof}

\begin{definicja}[Dziedzina całkowitości i ciało]
\label{aa:def:dziedzina-cialo}
Niezerowy pierścień przemienny z jedynką, w którym
$ab=0_R$ pociąga $a=0_R$ lub $b=0_R$, nazywamy
\emph{dziedziną całkowitości}. \emph{Ciałem} nazywamy
niezerowy pierścień przemienny z jedynką, w którym każdy
niezerowy element jest jednostką.
\end{definicja}
\begin{lemat}[Skracanie i skończone dziedziny]
\label{aa:lem:skracanie-pierscien}
Niech $R$ będzie dziedziną całkowitości, a $a,b,c\in R$.
Jeśli $a\ne0_R$ i $ab=ac$, to $b=c$.
Każda skończona dziedzina całkowitości jest ciałem.
\end{lemat}
\begin{proof}
Z $ab=ac$ mamy $a(b-c)=0_R$. Brak dzielników zera
oraz $a\ne0_R$ dają $b-c=0_R$. Teraz niech $R$ będzie
skończona, a $a\ne0_R$. Odwzorowanie $m_a:R\to R$,
$m_a(x)=ax$, jest różnowartościowe: jeśli $ax=ay$, to $x=y$
na mocy pierwszej części. Każde różnowartościowe odwzorowanie
skończonego zbioru w siebie jest surjektywne. Istnieje więc
$b\in R$ z $ab=1_R$. Ponieważ $R$ jest przemienny,
$b$ jest odwrotnością $a$. Wszystkie niezerowe elementy są jednostkami.
\end{proof}

\begin{przyklad}[Rozpoznawanie struktur]
\label{aa:prz:pierscien-struktury}
Pierścienie $\mathbb Z$, $\mathbb Q[x]$ i $\mathbb R[x,y]$
są dziedzinami całkowitości. W $\mathbb Z$ jednostkami są
$1,-1$. W pierścieniu $K[x]$ nad ciałem $K$ jednostkami
są dokładnie niezerowe wielomiany stałe: dla niezerowych
wielomianów $f,g$ zachodzi $\deg(fg)=\deg f+\deg g$,
zatem iloczyn równy $1$ wymusza stopnie zero.
Pierścień $M_2(\mathbb R)$ jest nieprzemienny:
$E_{12}E_{21}=E_{11}$, podczas gdy $E_{21}E_{12}=E_{22}$.
Ma także niezerowe macierze, których iloczyn jest macierzą zerową,
na przykład $E_{12}^2=0$. Nie jest więc dziedziną całkowitości.
W $\mathbb Z/6\mathbb Z$ zachodzi $[2][3]=[0]$, lecz
żaden czynnik nie jest zerem. W $\mathbb Z/8\mathbb Z$
klasa $[2]$ jest nilpotentna, bo $[2]^3=[8]=[0]$,
chociaż $[2]\ne[0]$ i $[2]^2=[4]\ne[0]$.
W pierścieniu $\mathbb R\times\mathbb R$ element $(1,0)$
jest idempotentny i jest dzielnikiem zera; jego dopełnienie
$(0,1)$ też jest idempotentne, a ich suma jest jedynką.
\end{przyklad}

\begin{definicja}[Charakterystyka pierścienia]
\label{aa:def:char-ring}
Niech $R$ będzie niezerowym pierścieniem z jedynką.
\emph{Charakterystyką} $R$ jest najmniejsza liczba całkowita
$n>0$ spełniająca $n1_R=0_R$, o ile istnieje;
w przeciwnym razie charakterystyka wynosi $0$.
\end{definicja}
\begin{stwierdzenie}[Charakterystyka dziedziny]
\label{aa:st:char-domain}
Niech $R$ będzie dziedziną całkowitości. Jeśli
$\operatorname{char}R>0$, to $\operatorname{char}R$ jest liczbą pierwszą.
Dla każdej liczby całkowitej
$n\ge2$ pierścień $\mathbb Z/n\mathbb Z$ ma charakterystykę $n$.
\end{stwierdzenie}
\begin{proof}
Jeśli charakterystyka dziedziny wynosi $n=ab$, gdzie
$1<a,b<n$, to $0_R=n1_R=(a1_R)(b1_R)$.
Ponieważ $R$ jest niezerowy, $n\ne1$. Minimalność $n$ zapewnia
$a1_R,b1_R\ne0_R$, co przeczy brakowi dzielników zera.
Zatem $n$ jest liczbą pierwszą.
W $\mathbb Z/n\mathbb Z$ równość $m[1]=[0]$ zachodzi
dokładnie wtedy, gdy $n$ dzieli $m$; najmniejsze dodatnie $m$
to $n$.
\end{proof}

\section{Podpierścienie, homomorfizmy i jądra}
\begin{definicja}[Podpierścień i homomorfizm]
\label{aa:def:ring-hom}
Niech $R$ i $T$ będą pierścieniami z jedynką.
\emph{Podpierścieniem z jedynką} pierścienia $R$ jest podzbiór
$S\subseteq R$ zawierający $1_R$ i zamknięty na dodawanie,
przeciwne elementy i mnożenie. \emph{Homomorfizmem pierścieni}
$\varphi:R\to T$ nazywamy odwzorowanie zachowujące dodawanie,
mnożenie i jedynkę: $\varphi(a+b)=\varphi(a)+\varphi(b)$,
$\varphi(ab)=\varphi(a)\varphi(b)$ i $\varphi(1_R)=1_T$.
Izomorfizm jest bijektywnym homomorfizmem pierścieni.
\end{definicja}
\begin{lemat}[Własności homomorfizmu]
\label{aa:lem:ring-hom-properties}
Niech $R,T$ będą pierścieniami z jedynką,
a $\varphi:R\to T$ homomorfizmem pierścieni.
Wtedy $\varphi(0_R)=0_T$ oraz
$\varphi(-a)=-\varphi(a)$ dla każdego $a\in R$.
Jeżeli $u\in R^\times$, to $\varphi(u)\in T^\times$.
Obraz homomorfizmu jest podpierścieniem $T$ z jedynką.
\end{lemat}
\begin{proof}
Z $\varphi(0_R)=\varphi(0_R+0_R)=2\varphi(0_R)$
po odjęciu $\varphi(0_R)$ otrzymujemy $\varphi(0_R)=0_T$.
Z $\varphi(a)+\varphi(-a)=\varphi(0_R)=0_T$
wynika $\varphi(-a)=-\varphi(a)$. Dla jednostki $u$
i odwrotności $v$ obrazy spełniają
$\varphi(u)\varphi(v)=\varphi(v)\varphi(u)=1_T$.
Obraz jest zamknięty na działania, przeciwny element,
a dzięki zachowaniu jedynki zawiera $1_T$.
\end{proof}

\begin{przyklad}[Podstawienie w wielomianie]
\label{aa:prz:poly-eval}
Niech $K$ będzie ciałem, a $a\in K$. Wzór
$\operatorname{ev}_a:K[x]\to K$, $f\mapsto f(a)$,
określa homomorfizm pierścieni, ponieważ podstawienie
zachowuje sumy, iloczyny i jedynkę. Jego jądro składa się
z wielomianów mających pierwiastek $a$. Dla każdego $k\ge1$ wielomian $x^k-a^k$ jest podzielny
przez $x-a$, ponieważ
$x^k-a^k=(x-a)\sum_{j=0}^{k-1}x^{k-1-j}a^j$.
Stąd dla każdego $f\in K[x]$ różnica $f(x)-f(a)$ jest
wielokrotnością $x-a$. Zatem
$\ker\operatorname{ev}_a=(x-a)$, a obraz jest całym $K$,
ponieważ zawiera wszystkie wielomiany stałe.
Ten przykład będzie za chwilę przypadkiem ogólnego twierdzenia
o izomorfizmie pierścieni.
\end{przyklad}

\section{Ideały jako warunek poprawności konstrukcji ilorazu}
\begin{definicja}[Ideały lewostronne, prawostronne i dwustronne]
\label{aa:def:ideal}
Niech $R$ będzie pierścieniem. Podgrupa addytywna
$I\le(R,+)$ jest \emph{ideałem lewostronnym}, jeśli
$ra\in I$ dla wszystkich $r\in R$, $a\in I$, oraz
\emph{ideałem prawostronnym}, jeśli $ar\in I$ dla tych par.
Jest \emph{ideałem dwustronnym}, jeśli ma obie własności.
W pierścieniu przemiennym te pojęcia są identyczne i mówimy
po prostu \emph{ideał}. Ideał jest \emph{właściwy}, jeśli $I\ne R$.
\end{definicja}
\begin{lemat}[Jądro homomorfizmu jest ideałem]
\label{aa:lem:kernel-ideal}
Niech $R,T$ będą pierścieniami z jedynką, a $\varphi:R\to T$
ich homomorfizmem. Wtedy zbiór
$\ker\varphi=\{a\in R:\varphi(a)=0_T\}$
jest ideałem dwustronnym $R$.
\end{lemat}
\begin{proof}
Ponieważ $\varphi$ zachowuje dodawanie i przeciwne elementy,
$\ker\varphi$ jest podgrupą addytywną. Jeśli
$a\in\ker\varphi$ i $r\in R$, to
$\varphi(ra)=\varphi(r)\varphi(a)=\varphi(r)0_T=0_T$.
Analogicznie $\varphi(ar)=0_T$. Obie własności ideału zachodzą.
\end{proof}

\begin{lemat}[Jedynka w ideale]
\label{aa:lem:ideal-unit}
Niech $R$ będzie pierścieniem z jedynką, a $I\subseteq R$
jego ideałem dwustronnym. Wówczas $I=R$
wtedy i tylko wtedy, gdy $1_R\in I$. Wystarczy też,
aby $I$ zawierał dowolną jednostkę $u\in R^\times$.
\end{lemat}
\begin{proof}
Jeśli $1_R\in I$, to dla każdego $r\in R$ mamy
$r=r1_R\in I$, zatem $I=R$. Jeśli $u\in I$ jest jednostką,
to $1_R=u^{-1}u\in I$. Jeżeli $I=R$, to $1_R\in I$, co dowodzi również implikacji odwrotnej.
\end{proof}

\begin{definicja}[Ideał generowany przez zbiór]
\label{aa:def:generated-ideal}
Niech $R$ będzie pierścieniem przemiennym z jedynką,
a $a_1,\ldots,a_s\in R$ dla pewnego $s\ge1$.
Przez $(a_1,\ldots,a_s)$ oznaczamy zbiór wszystkich sum $r_1a_1+\cdots+r_sa_s$,
$r_i\in R$. Jest to najmniejszy ideał zawierający
$a_1,\ldots,a_s$. Ideał $(a)$ nazywamy \emph{głównym}.
Dla dowolnego zbioru $S\subseteq R$ ideał $(S)$ składa się
ze skończonych sum $\sum_i r_i s_i$, gdzie $r_i\in R$, $s_i\in S$;
dla $S=\varnothing$ jest to ideał zerowy.
\end{definicja}
\begin{proof}
Przedstawiony zbiór zawiera zero, jest zamknięty na sumy,
przeciwne elementy i mnożenie przez dowolny element $R$,
zatem jest ideałem. Zawiera generatory, ponieważ
$a_i=1_Ra_i$. Każdy ideał zawierający te generatory musi
zawierać wszystkie ich wielokrotności i skończone sumy,
więc musi zawierać cały opisany zbiór.
\end{proof}

\begin{definicja}[Suma, przecięcie i iloczyn ideałów]
\label{aa:def:ideal-ops}
Dla ideałów $I,J$ przemiennego pierścienia $R$ definiujemy
$I+J=\{a+b:a\in I,b\in J\}$, $I\cap J$ w zwykłym
sensie teoriomnogościowym oraz \emph{iloczyn ideałów}
$I J$ jako ideał składający się z \emph{wszystkich skończonych
sum} elementów $ab$ z $a\in I$, $b\in J$.
W szczególności nie należy zastępować iloczynu ideałów
przez zbiór samych pojedynczych iloczynów.
\end{definicja}
\begin{stwierdzenie}[Relacje między działaniami na ideałach]
\label{aa:st:ideals-ops}
Niech $R$ będzie pierścieniem przemiennym z jedynką,
a $I,J\subseteq R$ jego ideałami. Wtedy
$IJ\subseteq I\cap J\subseteq I+J$. Ponadto $I+J$
jest najmniejszym ideałem zawierającym $I\cup J$.
\end{stwierdzenie}
\begin{proof}
Każdy składnik $ab$ z $a\in I$, $b\in J$ leży jednocześnie
w $I$ (bo $I$ jest ideałem) i w $J$. Suma takich składników
również leży w obu ideałach. Jeśli $x\in I\cap J$,
to $x=x+0_R\in I+J$. Opis minimalności $I+J$ wynika
z tego, że każdy ideał zawierający $I$ i $J$ jest zamknięty
na ich sumy.
\end{proof}

\begin{twierdzenie}[Pierścień ilorazowy]
\label{aa:tw:ring-quotient}
Niech $I$ będzie ideałem dwustronnym pierścienia $R$.
Na grupie abelowej $R/I$ istnieje pierścień z jedynką,
którego mnożenie jest zadane wzorem
$(a+I)(b+I)=ab+I$. Odwzorowanie
$\pi:R\to R/I$, $a\mapsto a+I$, jest surjektywnym
homomorfizmem pierścieni o jądrze $I$.
\end{twierdzenie}
\begin{proof}
Trzeba przede wszystkim sprawdzić niezależność od reprezentantów.
Jeśli $a'=a+i$ i $b'=b+j$, gdzie $i,j\in I$, to
$a'b'-ab=aj+ib+ij\in I$; pierwsze dwa składniki należą
do $I$ dzięki własnościom ideału dwustronnego,
a trzeci dzięki dowolnej z nich. Zatem $a'b'+I=ab+I$.
Łączność obu działań, przemienność dodawania i oba prawa
rozdzielności przechodzą z $R$ do ilorazu.
Elementem neutralnym dodawania jest $I$,
a jedynką mnożenia $1_R+I$. Projekcja zachowuje oba działania
i jedynkę, zaś $a+I=I$ wtedy i tylko wtedy, gdy $a\in I$.
\end{proof}

\begin{twierdzenie}[Pierwsze twierdzenie o izomorfizmie pierścieni]
\label{aa:tw:ring-isomorphism}
Niech $R,T$ będą pierścieniami z jedynką i niech
$\varphi:R\to T$ będzie homomorfizmem. Wtedy indukowane odwzorowanie
$\overline\varphi:R/\ker\varphi\to\operatorname{im}\varphi$,
$a+\ker\varphi\mapsto\varphi(a)$, jest izomorfizmem pierścieni.
\end{twierdzenie}
\begin{proof}
Jeżeli $a+\ker\varphi=b+\ker\varphi$, to
$a-b\in\ker\varphi$, więc $\varphi(a)=\varphi(b)$;
odwzorowanie jest dobrze określone. Zachowuje sumy,
iloczyny i jedynkę, ponieważ $\varphi$ je zachowuje.
Jest surjektywne z definicji obrazu. Jeśli
$\overline\varphi(a+\ker\varphi)=0_T$, to
$a\in\ker\varphi$, czyli $a+\ker\varphi$ jest zerem
ilorazu. Zatem jest też różnowartościowe.
\end{proof}

\begin{twierdzenie}[Odpowiedniość ideałów przy ilorazie]
\label{aa:tw:ring-ideal-correspondence}
Niech $R$ będzie pierścieniem z jedynką, $I\subseteq R$
jego ideałem dwustronnym i niech
$\pi:R\to R/I$, $\pi(a)=a+I$, będzie projekcją kanoniczną.
Wówczas przyporządkowanie $J\mapsto J/I$ jest bijekcją
między ideałami dwustronnymi $J$ pierścienia $R$
zawierającymi $I$ a ideałami dwustronnymi pierścienia $R/I$.
Odwrotną bijekcją jest $L\mapsto\pi^{-1}(L)$.
Konstrukcja zachowuje zawieranie.
\end{twierdzenie}
\begin{proof}
Przeciwobraz ideału dwustronnego przez homomorfizm pierścieni
jest ideałem dwustronnym: przeciwobraz jest podgrupą addytywną,
a mnożenie elementu przeciwobrazu przez dowolny element
$R$ przenosi się pod $\pi$ na mnożenie elementu $L$
przez element $R/I$. Ponadto $I=\ker\pi\subseteq\pi^{-1}(L)$.
Dla ideału $J$ z $I\subseteq J$ jego obraz $\pi(J)=J/I$
jest ideałem ilorazu, gdyż $\pi$ jest surjektywne.
Jeżeli $a\in\pi^{-1}(\pi(J))$, to $\pi(a)=\pi(j)$
dla pewnego $j\in J$, a więc $a-j\in I\subseteq J$;
stąd $a\in J$. Z drugiej strony $\pi(\pi^{-1}(L))=L$
z surjektywności $\pi$. Otrzymaliśmy obie odwrotne bijekcje.
Jeśli $J_1\subseteq J_2$, to każda klasa elementu $J_1$
należy również do $J_2/I$, zatem $J_1/I\subseteq J_2/I$.
Analogicznie z $L_1\subseteq L_2$ wynika
$\pi^{-1}(L_1)\subseteq\pi^{-1}(L_2)$.
\end{proof}

\begin{przyklad}[Wszystkie ideały w $\mathbb Z/12\mathbb Z$]
\label{aa:prz:ideals-z12}
Wszystkie ideały pierścienia $\mathbb Z/12\mathbb Z$ można przedstawić
przez przeciwobrazy w $\mathbb Z$. Uzasadnijmy najpierw, że każdy ideał $J\subseteq\mathbb Z$ ma postać
$d\mathbb Z$. Jeśli $J=\{0\}$, bierzemy $d=0$. W przeciwnym
razie w $J$ istnieje dodatnia liczba całkowita, gdyż dla
$z\in J\setminus\{0\}$ również $-z\in J$,
a jedna z liczb $z,-z$ jest dodatnia. Wybieramy najmniejszą
taką liczbę $d$. Dla dowolnego $a\in J$ dzielenie z resztą daje
$a=qd+r$ z $0\le r<d$. Ponieważ $r=a-qd\in J$,
minimalność $d$ wymusza $r=0$. Zatem $J=d\mathbb Z$.
Teraz zawieranie $12\mathbb Z\subseteq d\mathbb Z$ oznacza
$d\mid12$. Warunek ten wyklucza $d=0$.
Dla dodatniego dzielnika $d$ obraz ideału $d\mathbb Z$ ma
$12/d$ elementów, reprezentowanych przez klasy
$[0],[d],\ldots,[(12/d-1)d]$. Zatem mamy dokładnie sześć ideałów:
$\{0\}$, $([6])$, $([4])$, $([3])$, $([2])$ i cały pierścień.
Na przykład $([4])=\{[0],[4],[8]\}$, a
$([6])=\{[0],[6]\}$. Ich indeksy i liczności wynikają
z dzielenia $12$ przez odpowiednie dzielniki.
\end{przyklad}

\section{Ideały pierwsze i maksymalne}
\begin{definicja}[Ideał pierwszy i maksymalny]
\label{aa:def:prime-max-ideal}
Niech $R$ będzie pierścieniem przemiennym z jedynką.
Ideał właściwy $P\subsetneq R$ jest \emph{pierwszy}, gdy
dla wszystkich $a,b\in R$ warunek $ab\in P$ pociąga
$a\in P$ lub $b\in P$.
Ideał właściwy $M\subsetneq R$ jest \emph{maksymalny}, gdy
nie istnieje ideał $J$ spełniający
$M\subsetneq J\subsetneq R$.
\end{definicja}
\begin{twierdzenie}[Kryteria przez pierścień ilorazowy]
\label{aa:tw:prime-max-quotient}
Niech $R$ będzie pierścieniem przemiennym z jedynką,
a $I\subsetneq R$ jego ideałem właściwym.
Wówczas $I$ jest pierwszy wtedy i tylko wtedy, gdy
$R/I$ jest dziedziną całkowitości. Ponadto $I$ jest
maksymalny wtedy i tylko wtedy, gdy $R/I$ jest ciałem.
W szczególności każdy ideał maksymalny jest pierwszy.
\end{twierdzenie}
\begin{proof}
W ilorazie $R/I$ mamy $(a+I)(b+I)=I$ dokładnie wtedy,
gdy $ab\in I$. Klasa $a+I$ jest zerem dokładnie wtedy,
gdy $a\in I$. Ponieważ $I$ jest właściwy, iloraz nie jest
pierścieniem zerowym. Warunek braku dzielników zera jest
zatem równoważny definicji ideału pierwszego.

Jeżeli $I$ jest maksymalny, to z twierdzenia o odpowiedniości
ideałów wynika, że w $R/I$ jedynymi ideałami są $\{0\}$
i cały pierścień. Dla dowolnego $a+I\ne0$ ideał
$(a+I)$ nie jest zerowy, więc jest całym $R/I$.
Istnieje zatem klasa $b+I$ z $(a+I)(b+I)=1_R+I$,
czyli każda klasa niezerowa jest odwracalna. Iloraz jest ciałem.
Odwrotnie, w ciele każdy ideał niezerowy zawiera element
odwracalny, a więc jest całym ciałem. Twierdzenie
o odpowiedniości daje maksymalność $I$.
Wreszcie każde ciało jest dziedziną całkowitości:
z $ab=0$ i $a\ne0$ wynika $b=a^{-1}(ab)=0$.
\end{proof}

\begin{przyklad}[Ideały w $\mathbb Z$ i wielomiany]
\label{aa:prz:prime-max-examples}
Dla liczby pierwszej $p$ ideał $p\mathbb Z$ jest
maksymalny, ponieważ $\mathbb Z/p\mathbb Z$ jest ciałem.
Ideał $(0)$ w $\mathbb Z$ jest pierwszy, ale nie maksymalny:
iloraz $\mathbb Z/(0)$ to $\mathbb Z$, który nie jest ciałem.
W $K[x]$ ideał $(x-a)$ jest maksymalny dla $a\in K$,
ponieważ przez podstawienie w $a$ mamy izomorfizm
$K[x]/(x-a)\cong K$. Natomiast w $K[x,y]$ ideał $(x)$ jest pierwszy,
ale nie maksymalny. Rzeczywiście, podstawienie $x=0$
określa surjektywny homomorfizm $K[x,y]\to K[y]$ o jądrze
$(x)$: każdy wielomian ma postać $f(0,y)+xg(x,y)$.
Z pierwszego twierdzenia o izomorfizmie wynika
$K[x,y]/(x)\cong K[y]$, a ten ostatni pierścień
jest dziedziną całkowitości, lecz nie jest ciałem.
\end{przyklad}

\section{Chińskie twierdzenie o resztach dla pierścieni}
\begin{definicja}[Ideały względnie pierwsze]
\label{aa:def:comaximal}
Dwa ideały $I,J$ przemiennego pierścienia $R$ nazywamy
\emph{względnie pierwszymi}, jeśli $I+J=R$.
Jest to równoważne istnieniu $u\in I$ i $v\in J$ z $u+v=1_R$.
\end{definicja}
\begin{lemat}[Przecięcie i iloczyn ideałów względnie pierwszych]
\label{aa:lem:comaximal-intersection}
Niech $R$ będzie pierścieniem przemiennym z jedynką,
a $I,J\subseteq R$ ideałami. Jeśli $I+J=R$, to $I\cap J=IJ$.
\end{lemat}
\begin{proof}
Inkluzja $IJ\subseteq I\cap J$ została już udowodniona.
Niech $u\in I$, $v\in J$ oraz $u+v=1_R$.
Jeśli $z\in I\cap J$, to $z=zu+zv$.
Pierwszy składnik leży w $IJ$, ponieważ $z\in J$, $u\in I$;
drugi też leży w $IJ$, ponieważ $z\in I$, $v\in J$.
Zatem $z\in IJ$ i otrzymujemy równość.
\end{proof}
\begin{twierdzenie}[Chińskie twierdzenie o resztach]
\label{aa:tw:crt-ring}
Niech $R$ będzie pierścieniem przemiennym z jedynką,
a $I_1,\ldots,I_s$ parami względnie pierwszymi ideałami
($s\ge2$). Wówczas homomorfizm
\[
\Phi:R\longrightarrow \prod_{i=1}^{s}R/I_i,
\qquad r\longmapsto(r+I_1,\ldots,r+I_s)
\]
jest surjektywny, ma jądro $I_1\cap\cdots\cap I_s$
i indukuje izomorfizm pierścieni
\[
R/(I_1\cdots I_s)\ \cong\ \prod_{i=1}^{s}R/I_i.
\]
\end{twierdzenie}
\begin{proof}
Jądro $\Phi$ to elementy należące do wszystkich ideałów,
czyli ich przecięcie. Udowodnimy surjektywność jawnie.
Dla $i\ne j$ wybierzmy $u_{ij}\in I_i$ i $v_{ij}\in I_j$
takie, że $u_{ij}+v_{ij}=1_R$. Dla ustalonego $i$
połóżmy $e_i=\prod_{j\ne i}v_{ij}$. Mamy wtedy
$e_i\equiv1_R\pmod{I_i}$, ponieważ każdy czynnik
$v_{ij}=1_R-u_{ij}$ ma tę własność, oraz
$e_i\equiv0_R\pmod{I_j}$ dla każdego $j\ne i$, bo w
iloczynie występuje czynnik $v_{ij}\in I_j$.
Dla dowolnie wybranych klas $r_i+I_i$ element
$r=\sum_i r_i e_i$ ma obraz $\Phi(r)=(r_i+I_i)_i$.
To dowodzi surjektywności.

Pozostaje pokazać równość przecięcia z iloczynem dla wielu
ideałów. Jeśli $I+J=R$ i $I+K=R$, wybieramy
$u+v=1_R$ i $u'+w=1_R$ z $u,u'\in I$, $v\in J$, $w\in K$.
Mnożenie daje $1_R=uu'+uw+vu'+vw$, przy czym
pierwsze trzy składniki należą do $I$, a ostatni do $JK$.
Zatem $I+JK=R$. Stosując ten fakt wielokrotnie, otrzymujemy,
że $I_s$ jest względnie pierwszy z $I_1\cdots I_{s-1}$.
Indukcja i poprzedni lemat pokazują
$I_1\cap\cdots\cap I_s=I_1\cdots I_s$.
Pierwsze twierdzenie o izomorfizmie pierścieni kończy dowód.
\end{proof}

\begin{przyklad}[Rachunek modulo $12$ i idempotenty]
\label{aa:prz:crt12}
Ideały $3\mathbb Z$ i $4\mathbb Z$ są względnie pierwsze,
ponieważ $4-3=1$. Chińskie twierdzenie o resztach daje
$\mathbb Z/12\mathbb Z\cong\mathbb Z/3\mathbb Z\times\mathbb Z/4\mathbb Z$.
Wprost: klasa $[4]_{12}$ przechodzi na parę $([1]_3,[0]_4)$,
a klasa $[9]_{12}$ na $([0]_3,[1]_4)$.
Obie są idempotentne, bo $4^2\equiv4\pmod{12}$ i
$9^2\equiv9\pmod{12}$, a ich iloczyn jest zerowy.
Dla układu $x\equiv2\pmod3$, $x\equiv3\pmod4$ dostajemy
$x\equiv 2\cdot4+3\cdot9=35\equiv11\pmod{12}$.
Wynik można sprawdzić: $11-2=9$ jest podzielne przez $3$,
a $11-3=8$ jest podzielne przez $4$.
Natomiast dla ideałów $2\mathbb Z$ i $4\mathbb Z$
warunek względnej pierwszości nie zachodzi; układ
$x\equiv0\pmod2$, $x\equiv1\pmod4$ nie ma rozwiązania.
\end{przyklad}

\begin{uwaga}[Od pierścieni do geometrii i modułów]
\label{aa:uw:ring-next}
Ideały pierścieni wielomianów opisują układy równań
algebraicznych, a ilorazy pierścieni pozwalają zachować
informację nie tylko o zbiorach zer, lecz również o krotnościach
i elementach nilpotentnych. Moduły z poprzedniego rozdziału
są strukturami addytywnymi, na których działa pierścień.
W kolejnym rozdziale sprawdzimy, kiedy podzielność w pierścieniu
jest na tyle dobrze zorganizowana, aby każdy ideał był główny.
\end{uwaga}

\setlength{\parskip}{4pt}

% AA-10 | Algebra abstrakcyjna | Iteracja V
\chapter[Podzielność i dziedziny ideałów głównych]{Podzielność, dziedziny euklidesowe i dziedziny ideałów głównych}
\label{aa:ch:podzielnosc}
\setlength{\parskip}{2.8pt}

\section{Od dzielenia liczb do struktury ideałów}
W pierścieniu $\mathbb Z$ możemy dzielić z resztą, wyznaczać
największy wspólny dzielnik, rozkładać liczby na czynniki
pierwsze i rozwiązywać równania liniowe z całkowitymi
współczynnikami. Okazuje się, że za tymi zjawiskami stoją
\emph{różne poziomy założeń}. Dzielenie z resztą jest
silniejszą własnością niż samo istnienie największych wspólnych
dzielników; z kolei jednoznaczność rozkładu na czynniki
nie wystarcza, aby każdy ideał był generowany przez jeden element.
Rozdział wyjaśnia, które wnioski otrzymujemy na każdym poziomie.

W całym rozdziale, jeśli nie wskazano inaczej, $R$ jest
dziedziną całkowitości z jedynką, czyli niezerowym pierścieniem
przemiennym bez dzielników zera. Założenie to jest ważne:
pozwala skracać niezerowe czynniki i porównywać dzielniki.

\begin{definicja}[Podzielność i elementy stowarzyszone]
\label{aa:def:divides-associates}
Niech $R$ będzie dziedziną całkowitości z jedynką.
Dla $a,b\in R$ zapis $a\mid b$ oznacza, że istnieje
$c\in R$ z $b=ac$. Elementy $a,b$ są \emph{stowarzyszone},
jeśli $a=ub$ dla pewnej jednostki $u\in R^\times$.
Dla $a,b\in R$ \emph{największym wspólnym dzielnikiem}
nazywamy element $d\in R$ taki, że $d\mid a$, $d\mid b$
oraz każdy wspólny dzielnik $a,b$ dzieli $d$.
Jeżeli taki element istnieje, oznaczamy go przez
$\gcd(a,b)$, pamiętając, że jest wyznaczony tylko
z dokładnością do stowarzyszenia.
\end{definicja}
\begin{lemat}[Stowarzyszenie i jednoznaczność największego wspólnego dzielnika]
\label{aa:lem:gcd-unique}
Niech $R$ będzie dziedziną całkowitości. Elementy
$a,b\in R\setminus\{0\}$ są stowarzyszone wtedy i tylko
wtedy, gdy $a\mid b$ oraz $b\mid a$.
Dla dowolnych $a,b\in R$, jeżeli $d,d'\in R$ są największymi
wspólnymi dzielnikami pary $a,b$, to są stowarzyszone.
\end{lemat}
\begin{proof}
Z $a=ub$ i odwracalności $u$ natychmiast wynika
$a\mid b$ i $b\mid a$. W drugą stronę, jeśli $a=bc$
i $b=ad$, to $a=acd$. Ponieważ $a\ne0$, skracanie
daje $cd=1_R$, a więc $c$ i $d$ są jednostkami.
Dla dwóch największych wspólnych dzielników $d,d'$
warunek ich maksymalności w sensie podzielności
daje $d\mid d'$ i $d'\mid d$. Jeśli $a=b=0$, to każdy element dzieli zarówno $a$, jak i $b$.
W szczególności $0$ jest wspólnym dzielnikiem, więc z definicji
największego wspólnego dzielnika każdy taki element $d$ musi
spełniać $0\mid d$, czyli $d=0$. Element $0$ sam spełnia
warunek największego wspólnego dzielnika, ponieważ każdy element
dzieli $0$. Zatem $\gcd(0,0)=0$ także w sensie definicji.
W pozostałych przypadkach $d,d'$ są niezerowe i stosujemy
pierwszą część.
\end{proof}

\begin{przyklad}[Jednostki decydują o jednoznaczności]
\label{aa:prz:units-gcd}
W $\mathbb Z$ elementy stowarzyszone z $6$ to $6$ oraz $-6$;
wybieramy dodatni największy wspólny dzielnik.
W pierścieniu $\mathbb Q[x]$ elementy $x^2-1$ oraz
$7(x^2-1)$ są stowarzyszone, ponieważ $7$ jest jednostką.
Największy wspólny dzielnik wielomianów zwykle wybiera się
\emph{unormowany}, to znaczy o najwyższym współczynniku
równym $1$, jeśli wielomian nie jest zerowy.
\end{przyklad}

\begin{definicja}[Element nierozkładalny i element pierwszy]
\label{aa:def:irreducible-prime}
Niech $R$ będzie dziedziną całkowitości z jedynką i niech
$p\in R$ będzie niezerowy i nie będzie jednostką.
Element $p$ jest \emph{nierozkładalny}, jeśli z $p=ab$
wynika, że $a$ lub $b$ jest jednostką.
Element $p$ jest \emph{pierwszy}, jeśli dla wszystkich
$a,b\in R$ z $p\mid ab$ wynika $p\mid a$ lub $p\mid b$.
\end{definicja}
\begin{lemat}[Pierwszość implikuje nierozkładalność]
\label{aa:lem:prime-irred}
Niech $R$ będzie dziedziną całkowitości z jedynką.
Każdy element pierwszy $p\in R$ jest nierozkładalny.
\end{lemat}
\begin{proof}
Niech $p=ab$. Wtedy $p\mid ab$, a pierwszość
pozwala przyjąć, po ewentualnej zamianie $a,b$, że $p\mid a$.
Istnieje $c$ z $a=pc$, zatem $p=pcb$.
Skracając $p\ne0$, otrzymujemy $1_R=cb$.
Element $b$ jest jednostką. To dokładnie definicja nierozkładalności.
\end{proof}

\section{Dziedziny euklidesowe}
\begin{definicja}[Dziedzina euklidesowa]
\label{aa:def:euclidean-domain}
\emph{Dziedziną euklidesową} nazywamy dziedzinę całkowitości
$R$, dla której istnieje funkcja
$\delta:R\setminus\{0\}\to\mathbb N_0$ taka, że dla
każdych $a\in R$ i $b\in R\setminus\{0\}$ istnieją
$q,r\in R$ spełniające $a=bq+r$ oraz
$r=0$ albo $\delta(r)<\delta(b)$.
Funkcję $\delta$ nazywamy \emph{funkcją euklidesową}.
Nie zakładamy jednoznaczności $q,r$.
\end{definicja}

\begin{przyklad}[Liczby całkowite i dzielenie z resztą]
\label{aa:prz:integer-division}
Dla $R=\mathbb Z$ funkcja $\delta(b)=|b|$, określona dla
$b\in\mathbb Z\setminus\{0\}$, jest euklidesowa.
Jeżeli $b>0$, dzielenie z resztą daje jednoznaczne
$a=bq+r$ przy $0\le r<b$.
Gdy $b<0$, przechodzimy do $|b|$ i zmieniamy znak ilorazu.
Dla $a=-17$, $b=5$ otrzymujemy
$-17=5(-4)+3$, a $0\le3<5$. Wybór $q=-3$ dałby resztę
$-2$, która też spełnia $|{-2}|<5$, ale nie standardowy
warunek $0\le r<5$. Dlatego funkcja euklidesowa nie narzuca
jednoznaczności dzielenia, jeżeli nie ustalimy dodatkowej
konwencji dla reszt.
\end{przyklad}

\begin{przyklad}[Wielomiany jednej zmiennej]
\label{aa:prz:polynomial-division}
Niech $K$ będzie ciałem. W $K[x]$ funkcja $\delta(f)=\deg f$
dla $f\ne0$ jest euklidesowa.
Wyjaśnijmy mechanizm. Jeśli $f=0$ lub $\deg f<\deg g$, wystarczy wziąć
$q=0$ i $r=f$. W pozostałym przypadku $f\ne0$ oraz $\deg f\ge\deg g$.
Niech $a,b\in K^\times$ będą najwyższymi współczynnikami
odpowiednio $f$ i $g$. Odejmujemy od $f$ wielomian
$(a/b)x^{\deg f-\deg g}g$.
Najwyższy wyraz znika; nowy wielomian ma mniejszy stopień.
Powtarzanie kroku kończy się po skończonej liczbie operacji,
bo stopnie nie mogą maleć w nieskończoność.
Dostajemy $f=gq+r$, gdzie $r=0$ lub $\deg r<\deg g$.
Na przykład nad $\mathbb Q$:
\[
x^3+2x+1=(x+1)(x^2-x+3)-2.
\]
Rachunek można sprawdzić przez wymnożenie:
$(x+1)(x^2-x+3)=x^3+2x+3$.
Stała reszta $-2$ ma stopień $0<1$, więc spełnia wymóg.
\end{przyklad}

\begin{przyklad}[Liczby całkowite Gaussa]
\label{aa:prz:gaussian-euclidean}
Pierścień $\mathbb Z[i]=\{a+bi:a,b\in\mathbb Z\}$
jest dziedziną euklidesową względem funkcji euklidesowej
$\delta(z)=N(z)$ dla $z\ne0$, gdzie
$N(a+bi)=a^2+b^2$ jest określona dla wszystkich
$a+bi\in\mathbb Z[i]$, w tym dla zera.
Jeśli $\alpha,\beta\in\mathbb Z[i]$ i $\beta\ne0$,
zapiszmy $\alpha/\beta=x+iy$ z $x,y\in\mathbb R$.
Wybierzmy liczby całkowite $m,n$ najbliższe odpowiednio
$x,y$; wtedy $|x-m|\le\tfrac12$ i $|y-n|\le\tfrac12$.
Dla $q=m+in$ oraz $r=\alpha-\beta q$ zachodzi
$\alpha=\beta q+r$ i
\[
N(r)=N(\beta)\,N(\alpha/\beta-q)
\le N(\beta)(1/4+1/4)=N(\beta)/2<N(\beta).
\]
Ponieważ $r\in\mathbb Z[i]$, uzyskaliśmy wymagane dzielenie
z resztą. Własność multiplikatywną normy można sprawdzić
bez odwołania do analizy: dla $z=a+bi$, $w=c+di$
mamy $zw=(ac-bd)+(ad+bc)i$, a zatem
\[
N(zw)=(ac-bd)^2+(ad+bc)^2
=(a^2+b^2)(c^2+d^2)=N(z)N(w).
\]
Tutaj $i^2=-1$; konstrukcję ciała liczb zespolonych
z danej struktury ciała liczb rzeczywistych podamy
w dalszym rozdziale o konstrukcjach algebraicznych.
\end{przyklad}

\section{Algorytm Euklidesa i tożsamość Bézouta}
\begin{twierdzenie}[Algorytm Euklidesa]
\label{aa:tw:euclidean-algorithm}
Niech $R$ będzie dziedziną euklidesową z funkcją euklidesową
$\delta$ i niech $a,b\in R$, $b\ne0$. Połóżmy
$r_{-1}=a$ oraz $r_0=b$. Istnieje liczba całkowita $k\ge0$ i istnieją elementy
$q_0,\ldots,q_k,r_1,\ldots,r_{k+1}\in R$
takie, że
\[
r_{j-1}=r_jq_j+r_{j+1}\quad(0\le j\le k),
\qquad r_{k+1}=0,
\]
$r_j\ne0$ dla $0\le j\le k$ oraz
$\delta(r_{j+1})<\delta(r_j)$ dla $0\le j<k$.
Ostatni niezerowy element $r_k$ jest największym wspólnym
dzielnikiem $a$ i $b$. W szczególności, gdy już pierwsza
reszta jest zerem, mamy $k=0$ i $r_k=b$.
Co więcej, istnieją $u,v\in R$ takie, że
$ua+vb=\gcd(a,b)$, z dokładnością do wyboru elementu
stowarzyszonego z największym wspólnym dzielnikiem.
\end{twierdzenie}
\begin{proof}
Dzieląc kolejne pary elementów, w każdym kroku otrzymujemy
resztę równą zeru albo o mniejszej wartości $\delta$.
Wartości $\delta$ należą do $\mathbb N_0$, więc nie istnieje
nieskończony ściśle malejący ciąg wartości. Algorytm się kończy.
Z równości $a=bq_0+r_1$ wynika, że element $d$ dzieli
jednocześnie $a,b$ wtedy i tylko wtedy, gdy dzieli $b,r_1$:
jeśli $d\mid a,b$, to $d\mid a-bq_0=r_1$;
jeśli $d\mid b,r_1$, to $d\mid bq_0+r_1=a$.
Taką samą równoważność stosujemy na każdym kolejnym kroku.
Jeżeli pierwsza reszta jest zerem, ostatnim niezerowym
elementem jest $b$, którego dzielniki są wspólnymi dzielnikami
$a$ i $b$. W przeciwnym razie wspólni dzielnicy $a,b$ są
dokładnie dzielnikami ostatniej niezerowej reszty $r_k$: dzielnik $r_k$
dzieli także poprzednią resztę $r_{k-1}=r_k q_k$,
a wszystkie wcześniejsze na podstawie odwróconych równości.
To dowodzi charakterystyki największego wspólnego dzielnika.

Aby uzyskać tożsamość Bézouta, wykonujemy podstawienia wsteczne.
Pierwsza reszta ma postać $r_1=a-bq_0$,
następna $r_2=b-r_1q_1$, a więc też jest kombinacją
$R$-liniową $a,b$. Indukcyjnie każda kolejna reszta
jest kombinacją liniową tych samych dwóch elementów.
W przypadku zerowej pierwszej reszty mamy $b=0a+1b$.
W przeciwnym razie $r_k=ua+vb$ dla pewnych $u,v\in R$.
\end{proof}

\begin{przyklad}[Rozszerzony algorytm Euklidesa dla $252$ i $198$]
\label{aa:prz:extended-euclid}
W $\mathbb Z$ wykonujemy dzielenie:
\[
252=1\cdot198+54,\qquad
198=3\cdot54+36,\qquad
54=1\cdot36+18,\qquad36=2\cdot18.
\]
Ostatnia niezerowa reszta wynosi $18$.
Podstawiamy od końca:
\[
\begin{aligned}
18&=54-36=54-(198-3\cdot54)\\
  &=4\cdot54-198=4(252-198)-198\\
  &=4\cdot252-5\cdot198.
\end{aligned}
\]
Otrzymujemy nie tylko $\gcd(252,198)=18$, lecz także
konkretne współczynniki Bézouta. Ponieważ $18$ jest
największym wspólnym dzielnikiem, równanie
$252x+198y=18$ ma rozwiązanie całkowite $(x,y)=(4,-5)$.
\end{przyklad}

\begin{wniosek}[Lemat Euklidesa w dziedzinie euklidesowej]
\label{aa:wn:euclid-lemma}
Niech $a,b,c$ należą do dziedziny euklidesowej $R$.
Jeżeli $\gcd(a,b)$ jest jednostką oraz $a\mid bc$, to $a\mid c$.
\end{wniosek}
\begin{proof}
Tożsamość Bézouta daje $ua+vb=1_R$ dla pewnych $u,v\in R$
(po pomnożeniu przez odwrotność największego wspólnego dzielnika).
Mnożąc przez $c$, otrzymujemy $c=uac+vbc$.
Pierwszy składnik jest podzielny przez $a$ z definicji,
a drugi na mocy $a\mid bc$. Zatem $a\mid c$.
\end{proof}

\begin{przyklad}[Największy wspólny dzielnik wielomianów]
\label{aa:prz:poly-gcd}
W $\mathbb Q[x]$ rozważmy $f=x^3-1$ i $g=x^2-1$.
Mamy $f=(x-1)(x^2+x+1)$ oraz $g=(x-1)(x+1)$.
Dzielenie $x^2+x+1$ przez $x+1$ daje resztę $1$,
bo podstawienie $x=-1$ daje $1$. Zatem oba czynniki
$x^2+x+1$ i $x+1$ są względnie pierwsze.
Tożsamość Bézouta dla tych dwóch wielomianów daje
$u(x)(x^2+x+1)+v(x)(x+1)=1$ dla pewnych $u,v\in\mathbb Q[x]$.
Po pomnożeniu przez $x-1$ otrzymujemy
$u(x)f(x)+v(x)g(x)=x-1$. Każdy wspólny dzielnik
$f$ i $g$ dzieli więc $x-1$, który sam dzieli oba wielomiany.
Otrzymujemy unormowany największy wspólny dzielnik
$\gcd(f,g)=x-1$.
\end{przyklad}

\section{Dziedziny ideałów głównych}
\begin{definicja}[Dziedzina ideałów głównych]
\label{aa:def:pid}
Dziedzinę całkowitości $R$, w której każdy ideał ma postać
$(d)=\{rd:r\in R\}$ dla pewnego $d\in R$, nazywamy
\emph{dziedziną ideałów głównych}. Oznacza to, że nawet
ideał generowany przez wiele elementów można wygenerować
jednym elementem. Nie znaczy to, że wszystkie elementy
pierścienia są wzajemnie podzielne.
\end{definicja}
\begin{twierdzenie}[Każda dziedzina euklidesowa jest dziedziną ideałów głównych]
\label{aa:tw:euclidean-pid}
Jeżeli $R$ jest dziedziną euklidesową, to każdy ideał
$R$ jest główny.
\end{twierdzenie}
\begin{proof}
Niech $I$ będzie ideałem $R$. Jeśli $I=\{0\}$,
to $I=(0)$. W przeciwnym razie rozważmy niepusty zbiór
$\{\delta(a):a\in I,\ a\ne0\}\subseteq\mathbb N_0$.
Ma on element najmniejszy. Wybierzmy niezerowy $d\in I$
o minimalnej wartości $\delta(d)$. Dla dowolnego $a\in I$
wykonajmy dzielenie $a=dq+r$, gdzie $r=0$
lub $\delta(r)<\delta(d)$. Ponieważ $a,d\in I$
oraz $I$ jest ideałem, $r=a-dq\in I$.
Gdyby $r\ne0$, przeczyłoby to minimalności $\delta(d)$.
Zatem $r=0$, czyli $a=dq\in(d)$. Udowodniliśmy $I\subseteq(d)$.
Odwrotna inkluzja zachodzi, bo $d\in I$ i $I$ jest ideałem.
Stąd $I=(d)$.
\end{proof}

\begin{twierdzenie}[Największy wspólny dzielnik w dziedzinie ideałów głównych]
\label{aa:tw:pid-bezout}
Jeżeli $R$ jest dziedziną ideałów głównych, to dla
każdych $a,b\in R$, z których co najmniej jeden jest niezerowy,
ideał $(a,b)$ ma postać $(d)$, gdzie $d$ jest największym
wspólnym dzielnikiem $a,b$. Istnieją $u,v\in R$ z
$d=ua+vb$.
\end{twierdzenie}
\begin{proof}
Z definicji dziedziny ideałów głównych istnieje $d\in R$
z $(a,b)=(d)$. Ponieważ $a,b\in(d)$, mamy $d\mid a$
i $d\mid b$. Z drugiej strony $d\in(a,b)$, więc
dla pewnych $u,v\in R$ zachodzi $d=ua+vb$.
Jeśli $c$ jest wspólnym dzielnikiem $a,b$, to dzieli każdą
ich kombinację liniową, a zatem dzieli $d$.
W ten sposób wykazaliśmy zarówno własność największego
wspólnego dzielnika, jak i tożsamość Bézouta.
\end{proof}

\begin{wniosek}[Pierwszość elementu nierozkładalnego w dziedzinie ideałów głównych]
\label{aa:wn:irred-is-prime-pid}
Niech $R$ będzie dziedziną ideałów głównych i niech
$p\in R$ będzie elementem nierozkładalnym. Wtedy $p$ jest pierwszy.
\end{wniosek}
\begin{proof}
Niech $p$ będzie nierozkładalny i niech $p\mid ab$.
Jeśli $p\nmid a$, to największy wspólny dzielnik $p,a$
musi być jednostką: dzieli $p$, a nierozkładalność $p$
oznacza, że każdy jego dzielnik jest jednostką albo jest
stowarzyszony z $p$; druga możliwość wymuszałaby $p\mid a$.
Tożsamość Bézouta daje $up+va=1_R$.
Mnożąc przez $b$, otrzymujemy $b=upb+vab$;
oba składniki są podzielne przez $p$. A zatem $p\mid b$.
\end{proof}

\section{Dlaczego rozkład na czynniki musi się kończyć?}
\begin{definicja}[Dziedzina z jednoznacznym rozkładem]
\label{aa:def:ufd}
Dziedzinę całkowitości $R$ nazywamy
\emph{dziedziną z jednoznacznym rozkładem na czynniki},
jeśli każdy niezerowy element niebędący jednostką można
zapisać jako skończony iloczyn elementów nierozkładalnych,
a dwa takie zapisy różnią się co najwyżej kolejnością
czynników i zastąpieniem każdego czynnika elementem stowarzyszonym.
\end{definicja}

\begin{lemat}[Warunek stabilizacji wstępujących ciągów ideałów]
\label{aa:lem:pid-ascending-chain}
Niech $R$ będzie dziedziną ideałów głównych.
Każdy wstępujący ciąg ideałów $R$,
$I_1\subseteq I_2\subseteq\cdots$, staje się stały.
\end{lemat}
\begin{proof}
Suma zbiorów $I=\bigcup_{n\ge1}I_n$ jest ideałem:
jeżeli $a\in I_j$ i $b\in I_k$, to obydwa należą
do $I_{\max(j,k)}$, więc $a-b$ należy do tego ideału;
to samo dotyczy mnożenia przez dowolny element pierścienia.
Ponieważ $R$ jest dziedziną ideałów głównych, $I=(d)$.
Generator $d$ należy do $I_N$ dla pewnego $N$,
zatem $(d)\subseteq I_N$. Z drugiej strony $I_N\subseteq I$.
Wynika stąd $I_N=I$, a dla $n\ge N$ mamy $I_n=I$.
\end{proof}

\begin{twierdzenie}[Dziedziny ideałów głównych mają jednoznaczny rozkład na czynniki]
\label{aa:tw:pid-ufd}
Niech $R$ będzie dziedziną ideałów głównych.
Wówczas $R$ jest dziedziną z jednoznacznym rozkładem
na czynniki.
\end{twierdzenie}
\begin{proof}
Najpierw udowodnimy istnienie rozkładu.
Przypuśćmy przeciwnie, że istnieje niezerowy nieodwracalny
element $a_0$, którego nie można przedstawić jako skończonego
iloczynu elementów nierozkładalnych. Nie jest on nierozkładalny,
bo wtedy iloczyn złożony z jednego czynnika byłby odpowiednim
rozkładem. Istnieją więc nieodwracalne $a_1,b_1$ z
$a_0=a_1b_1$. Co najmniej jeden z czynników nie ma rozkładu
na nierozkładalne: gdyby oba miały, ich rozkłady po pomnożeniu
dałyby rozkład $a_0$. Oznaczmy ten czynnik przez $a_1$
(po ewentualnej zamianie czynników).
Powtarzając argument, otrzymujemy ciąg
$a_0=a_1b_1$, $a_1=a_2b_2$, itd.,
gdzie każde $b_i$ nie jest jednostką. Stąd
$(a_0)\subseteq(a_1)\subseteq(a_2)\subseteq\cdots$.
Każde zawieranie jest ścisłe: gdyby $(a_{j-1})=(a_j)$,
to dla pewnego $d\in R$ mielibyśmy $a_j=a_{j-1}d$.
Z równości $a_{j-1}=a_jb_j$ wynikałoby wtedy
$a_j=a_jb_jd$. Ponieważ $a_j\ne0$, skracanie daje
$b_jd=1_R$, czyli $b_j$ jest jednostką, wbrew konstrukcji.
Otrzymujemy nieskończony ściśle rosnący ciąg ideałów,
sprzeczny z poprzednim lematem. Zatem rozkład istnieje.

Teraz udowodnimy jednoznaczność. Każdy element nierozkładalny
jest pierwszy na mocy Wniosku~\ref{aa:wn:irred-is-prime-pid}.
Jeżeli $p_1\cdots p_r=q_1\cdots q_s$ są dwoma rozkładami,
to $p_1$ dzieli iloczyn po prawej. Wielokrotne zastosowanie
pierwszości daje, że $p_1\mid q_j$ dla pewnego $j$.
Ponieważ $q_j$ jest nierozkładalny, $p_1$ i $q_j$
są stowarzyszone. Po przestawieniu czynników możemy
założyć $j=1$ i zapisać $q_1=up_1$ dla jednostki $u$.
Skracając niezerowe $p_1$, otrzymujemy równość iloczynu
pozostałych $p_i$ z iloczynem pozostałych $q_j$ pomnożonym
przez jednostkę $u$. Jeżeli po obu stronach są jeszcze czynniki,
wchłaniamy $u$ w jeden z nich; nie zmienia to nierozkładalności.
Przypadek, gdy po obu stronach nie pozostał żaden czynnik,
kończy indukcję. Gdy po skróceniu pozostałyby
po jednej stronie czynniki nierozkładalne, a po drugiej
jednostka, powstałaby sprzeczność: iloczyn niejednostek
w dziedzinie całkowitości nie może być jednostką.
Indukcja daje $r=s$ oraz stowarzyszenie odpowiadających sobie
czynników. Gdy jeden z rozkładów jest pusty, obie strony
są jednostkami i nie mogą jednocześnie zawierać czynników
nierozkładalnych. To kończy dowód.
\end{proof}

\begin{uwaga}[Łańcuch implikacji i pytania odwrotne]
\label{aa:uw:euclidean-pid-ufd}
Udowodniliśmy implikacje
\[
\begin{gathered}
\text{dziedzina euklidesowa}\ \Longrightarrow\
\text{dziedzina ideałów głównych}\\
\Longrightarrow\ \text{dziedzina z jednoznacznym rozkładem na czynniki}.
\end{gathered}
\]
Odwrotne implikacje nie wynikają z powyższych dowodów.
Nie każdy pierścień z jednoznacznym rozkładem na czynniki
ma wszystkie ideały główne. Przykład taki uzyskamy przy
pierścieniu $K[x,y]$, gdy udowodnimy jego własność
jednoznacznego rozkładu w następnym rozdziale.
Istnieją również dziedziny ideałów głównych, które nie są
euklidesowe, lecz ich konstrukcja i dowód tego faktu
wykraczają poza obecny rozdział.
\end{uwaga}

\begin{przyklad}[Dlaczego pierścień wielomianów dwóch zmiennych nie jest dziedziną ideałów głównych]
\label{aa:prz:two-polynomials-not-pid}
Niech $K$ będzie dowolnym ciałem. W pierścieniu $K[x,y]$
rozważmy ideał $I=(x,y)$. Załóżmy, że $I=(f)$.
Wtedy $f$ dzieli jednocześnie $x$ i $y$.
Ponieważ $K$ jest ciałem, pierścień $K[x,y]$ jest
dziedziną całkowitości. Rzeczywiście, traktując niezerowe wielomiany $f,g$ jako
wielomiany zmiennej $y$ o współczynnikach w $K[x]$,
ich najwyższe współczynniki względem $y$ są niezerowymi
wielomianami w $x$. Ich iloczyn jest niezerowy, gdyż
iloczyn dwóch niezerowych wielomianów jednej zmiennej
nad ciałem ma niezerowy najwyższy współczynnik.
To dowodzi, że $K[x,y]$ nie ma dzielników zera.
Ponadto stopień całkowity niezerowego iloczynu
jest sumą stopni całkowitych: jeśli $f_d$ i $g_e$
są niezerowymi częściami jednorodnymi najwyższych
stopni całkowitych $d$ i $e$, to częścią jednorodną
stopnia $d+e$ iloczynu jest $f_dg_e$, która jest
niezerowa na mocy właśnie wykazanej własności. Z podzielności $f\mid x$ wynika,
że $f$ ma stopień najwyżej $1$. Jeśli ma stopień $1$,
to musi być stowarzyszony z $x$; analogicznie z $f\mid y$
wynikałoby stowarzyszenie z $y$. Nie jest to możliwe, bo jednostkami
$K[x,y]$ są niezerowe stałe, a równość $x=cy$ z
$c\in K^\times$ obala podstawienie $y=0$.
Zatem $f$ ma stopień $0$, jest niezerową stałą i jest jednostką.
Wtedy $I=K[x,y]$, czyli $1\in(x,y)$, ale to sprzeczne
z podstawieniem $x=y=0$ do równości $1=xf_1+yf_2$.
Wniosek: ideał $(x,y)$ nie jest główny.
\end{przyklad}

\section{Gdy zawodzi jednoznaczny rozkład}
\begin{przyklad}[Dwa różne rozkłady liczby $6$]
\label{aa:prz:zsqrtminus5}
Rozważmy dziedzinę całkowitości
$R=\mathbb Z[\sqrt{-5}]=\{a+b\sqrt{-5}:a,b\in\mathbb Z\}$.
Jest ona podpierścieniem $\mathbb C$, więc nie ma dzielników
zera. Norma $N(a+b\sqrt{-5})=a^2+5b^2$ ma wartości w
$\mathbb N_0$ i jest multiplikatywna:
$N(zw)=N(z)N(w)$. Rzeczywiście, dla sprzężenia
$\overline{a+b\sqrt{-5}}=a-b\sqrt{-5}$ mamy
$N(z)=z\overline z$, a więc
$N(zw)=zw\overline z\,\overline w=N(z)N(w)$. W tej dziedzinie
\[
6=2\cdot3=(1+\sqrt{-5})(1-\sqrt{-5}).
\]
Wykażemy, że wszystkie cztery czynniki są nierozkładalne.
Najpierw wyznaczmy wszystkie jednostki. Jeśli $u$ ma
odwrotność w $R$, to $N(u)N(u^{-1})=1$. Obie normy są
dodatnimi liczbami całkowitymi, więc $N(u)=1$.
Równanie $a^2+5b^2=1$ wymusza $b=0$ i $a=\pm1$.
Odwrotnie, $1$ oraz $-1$ są jednostkami, bo są swoimi
własnymi odwrotnościami. Są to zatem wszystkie jednostki.
Nie istnieją elementy o normie $2$ ani $3$:
przy $b\ne0$ norma jest co najmniej $5$, a przy $b=0$
musiałaby być kwadratem liczby całkowitej.
Jeśli $2=ab$ z nieodwracalnymi $a,b$, to
$4=N(2)=N(a)N(b)$, a oba czynniki normowe byłyby
większe od $1$; jedyna możliwość to $2$ i $2$, niemożliwa.
Tak samo $N(3)=9$ wymagałoby dla rozkładu na niejednostki
norm $3$ i $3$, których nie ma. Dla
$1\pm\sqrt{-5}$ norma wynosi $6$, więc rozkład
na niejednostki wymagałby norm $2$ i $3$, również niemożliwych.
Wreszcie elementy $2,3,1+\sqrt{-5},1-\sqrt{-5}$
nie są parami stowarzyszone: mają różne normy, z wyjątkiem
ostatniej pary, ale $1+\sqrt{-5}\ne\pm(1-\sqrt{-5})$.
Otrzymaliśmy dwa rozkłady na nierozkładalne, które nie
różnią się jedynie kolejnością i jednostkami.
Pierścień $\mathbb Z[\sqrt{-5}]$ nie jest więc dziedziną
z jednoznacznym rozkładem na czynniki.
\end{przyklad}

\begin{przyklad}[Element nierozkładalny, który nie jest pierwszy]
\label{aa:prz:irred-not-prime}
W powyższym pierścieniu liczba $2$ jest nierozkładalna,
ale nie jest elementem pierwszym. Mamy
$2\mid (1+\sqrt{-5})(1-\sqrt{-5})$, ponieważ iloczyn
wynosi $6=2\cdot3$. Tymczasem $2$ nie dzieli żadnego
z czynników: równości
$1\pm\sqrt{-5}=2(a+b\sqrt{-5})$ dla $a,b\in\mathbb Z$
wymagałyby $2a=1$, co nie ma rozwiązania całkowitego.
Widać dokładnie, dlaczego implikacja
„nierozkładalny $\Rightarrow$ pierwszy” wymaga dodatkowych
założeń o pierścieniu.
\end{przyklad}

\begin{uwaga}[Związek z modułami i algebrą liniową]
\label{aa:uw:pid-module-link}
Dziedzina ideałów głównych umożliwia nie tylko jednoznaczny
rozkład elementów. Jej głębszą cechą jest to, że podmoduły
modułów wolnych skończonej rangi również są wolne i mają
kontrolowaną postać. Z tego wyprowadzimy klasyfikację
skończenie generowanych modułów nad dziedziną ideałów
głównych oraz klasyfikację skończenie generowanych grup
abelowych. Następny rozdział rozwinie pierścienie wielomianów
i kryteria nierozkładalności, potrzebne również w teorii Galois.
\end{uwaga}
\setlength{\parskip}{4pt}

% AA-11 | Algebra abstrakcyjna | Iteracja VI
\chapter{Pierścienie wielomianów i kryteria nierozkładalności}
\chaptermark{Wielomiany i nierozkładalność}
\label{aa:ch:wielomiany}

\section{Dlaczego rozróżniamy wielomian i funkcję wielomianową?}
Wielomian jest formalną skończoną sumą współczynników i potęg
symbolu $x$. Jego wartość w danym elemencie pierścienia otrzymujemy
przez podstawienie, lecz wielomian nie jest z definicji funkcją.
To rozróżnienie jest ważne już nad ciałami skończonymi.

\begin{definicja}[Pierścień wielomianów jednej zmiennej]
\label{aa:def:poly-ring}
Niech $R$ będzie przemiennym pierścieniem z jedynką.
Elementem $R[x]$ jest ciąg $(a_0,a_1,\ldots)$ elementów $R$,
w którym tylko skończenie wiele współczynników jest niezerowych.
Zapisujemy go jako $f=\sum_{j=0}^n a_jx^j$.
Dodawanie jest współczynnikowe, natomiast współczynnik przy $x^k$
iloczynu $fg$ jest równy $\sum_{i+j=k}a_i b_j$.
Wielomian zerowy ma wszystkie współczynniki równe zeru.
Jeśli $f\ne0$, jego \emph{stopniem} jest największy indeks
niezerowego współczynnika, a \emph{współczynnikiem wiodącym}
--- współczynnik o tym indeksie. Wielomian jest \emph{unormowany},
gdy jego współczynnik wiodący jest równy $1_R$.
\end{definicja}

\begin{lemat}[Aksjomaty pierścienia wielomianów]
\label{aa:lem:poly-ring-axioms}
Dla każdego przemiennego pierścienia $R$ z jedynką
wzory podane w Definicji~\ref{aa:def:poly-ring}
wyznaczają przemienny pierścień $R[x]$ z jedynką.
Odwzorowanie $a\mapsto(a,0,0,\ldots)$ jest zanurzeniem
pierścieni $R\to R[x]$ (także dla pierścienia zerowego).
\end{lemat}
\begin{proof}
Suma dwóch ciągów o skończonym nośniku ma skończony
nośnik. Jeżeli któryś z czynników jest wielomianem
zerowym, ich iloczyn jest zerowy. W przeciwnym razie
współczynniki iloczynu przy $x^k$ dla $k$ większego
od sumy największych niezerowych indeksów obu czynników
są zerowe. Zatem iloczyn także ma skończony nośnik.
Aksjomaty dodawania wynikają współczynnikowo z aksjomatów
$(R,+)$; elementem zerowym jest ciąg zerowy.
Przy mnożeniu przemienność wynika z wymiany indeksów
$i,j$ w sumie $\sum_{i+j=k}a_ib_j$. Przy łączności
współczynnik przy $x^k$ iloczynu $(fg)h$ ma postać
$\sum_{i+j+\ell=k} a_i b_j c_\ell$, tak samo
jak współczynnik przy $x^k$ iloczynu $f(gh)$.
Wszystkie sumy są skończone, zatem wolno zmienić kolejność
sumowania. Ciąg $(1_R,0,\ldots)$ jest jedynką,
a równość $f(g+h)=fg+fh$ można sprawdzić przy każdym
współczynniku, stosując rozdzielność w $R$.
Stałe wielomiany zachowują oba działania i jedynkę.
Jeśli dwa stałe wielomiany są równe, ich współczynniki
przy $x^0$ są równe, więc odwzorowanie jest różnowartościowe
również dla pierścienia zerowego.
\end{proof}

\begin{lemat}[Stopień iloczynu nad dziedziną]
\label{aa:lem:poly-degree}
Niech $R$ będzie dziedziną całkowitości i niech
$f,g\in R[x]$ będą niezerowe. Wówczas
$\deg(fg)=\deg f+\deg g$. Jednostkami pierścienia $R[x]$
są dokładnie wielomiany stałe, których jedyny współczynnik
jest jednostką w $R$.
\end{lemat}
\begin{proof}
Niech $n=\deg f$, $m=\deg g$ oraz $a_n,b_m$ będą
współczynnikami wiodącymi. Przy potędze $x^{n+m}$ w iloczynie
może wystąpić tylko składnik $a_nb_m$. Jest on niezerowy,
ponieważ $R$ nie ma dzielników zera. Wyższe współczynniki są
zerowe, gdyż nie istnieje para indeksów $i\le n$, $j\le m$
z sumą większą niż $n+m$. To dowodzi wzoru na stopień.
Jeśli $fg=1_R$, obydwa czynniki są niezerowe, a suma
ich stopni jest zerem. Oba stopnie wynoszą zatem zero,
więc $f,g$ są stałe i ich współczynniki są wzajemnymi
odwrotnościami. Odwrotnie, każda stała jednostka $R$
jest jednostką $R[x]$.
\end{proof}

\begin{definicja}[Wartość wielomianu]
\label{aa:def:poly-evaluation}
Niech $R$ będzie przemiennym pierścieniem z jedynką,
$a\in R$ oraz $f=\sum_{j=0}^n c_jx^j\in R[x]$.
Definiujemy
$f(a)=\sum_{j=0}^n c_ja^j$. Odwzorowanie
$\operatorname{ev}_a:R[x]\to R$, $f\mapsto f(a)$,
nazywamy \emph{homomorfizmem podstawienia}.
Jest on rzeczywiście homomorfizmem: wzór na współczynniki
iloczynu i rozdzielność w $R$ dają
$(fg)(a)=f(a)g(a)$, a zgodność z dodawaniem i jedynką
wynika bezpośrednio z definicji.
\end{definicja}
\begin{lemat}[Własność uniwersalna pierścienia wielomianów]
\label{aa:lem:poly-universal}
Niech $R,S$ będą przemiennymi pierścieniami z jedynką,
$\varphi:R\to S$ homomorfizmem pierścieni i $s\in S$.
Istnieje dokładnie jeden homomorfizm pierścieni
$\widetilde\varphi:R[x]\to S$, którego ograniczenie
do $R$ jest równe $\varphi$ oraz dla którego
$\widetilde\varphi(x)=s$. Ma on postać
$\widetilde\varphi(\sum_j a_jx^j)=\sum_j\varphi(a_j)s^j$.
\end{lemat}
\begin{proof}
Homomorfizm musi spełniać podany wzór, ponieważ zachowuje
sumy, iloczyny i jedynkę. Wzór jest dobrze określony:
każdy wielomian ma dokładnie jeden skończony zestaw współczynników.
Zachowanie dodawania wynika ze współczynnikowego dodawania.
Dla mnożenia mamy, z przemienności w $S$,
\[
 \sum_k\varphi\Bigl(\sum_{i+j=k}a_i b_j\Bigr)s^k
 =\sum_{i,j}\varphi(a_i)s^i\varphi(b_j)s^j,
\]
co jest iloczynem obrazów dwóch wielomianów.
Obraz stałej $1_R$ jest $1_S$. Otrzymaliśmy homomorfizm;
jego jednoznaczność została wykazana na początku.
\end{proof}

\begin{przyklad}[Dwa różne wielomiany wyznaczające tę samą funkcję]
\label{aa:prz:poly-finite-field}
W ciele $\mathbb F_p=\mathbb Z/p\mathbb Z$,
gdzie $p$ jest liczbą pierwszą, wielomian $x^p-x$
nie jest wielomianem zerowym: ma dwa niezerowe współczynniki.
Jednak dla każdej klasy $a\in\mathbb F_p$ zachodzi
$a^p=a$, na mocy małego twierdzenia Fermata dla $a\ne0$
i bezpośrednio dla $a=0$. Funkcja zadana przez $x^p-x$
jest więc identycznie zerowa. Tym samym różne wielomiany
nie muszą wyznaczać różnych funkcji na ciele skończonym.
\end{przyklad}



\section{Dzielenie z resztą, pierwiastki i czynniki liniowe}
\begin{twierdzenie}[Dzielenie wielomianów przez wielomian o współczynniku wiodącym odwracalnym]
\label{aa:tw:poly-division}
Niech $R$ będzie przemiennym pierścieniem z jedynką,
$f,g\in R[x]$, $g\ne0$, a współczynnik wiodący $g$
niech będzie jednostką $R$. Istnieją dokładnie jedne
$q,r\in R[x]$ takie, że $f=qg+r$ oraz $r=0$ albo
$\deg r<\deg g$. W szczególności twierdzenie zachodzi
dla dowolnego $g\ne0$ nad ciałem $K$.
\end{twierdzenie}
\begin{proof}
Niech $d=\deg g$ i niech $u$ będzie odwracalnym
współczynnikiem wiodącym $g$.
Zaczynamy od $f_0=f$ i $q_0=0$. Jeśli $f_0=0$
lub $\deg f_0<d$, kończymy. W przeciwnym razie,
przy $\deg f_0=n$ i współczynniku wiodącym $a$,
odejmujemy wielomian $au^{-1}x^{n-d}g$ i dopisujemy
$au^{-1}x^{n-d}$ do ilorazu. Współczynnik przy
$x^n$ znika, więc stopień nowej reszty jest mniejszy
od $n$. Procedura kończy się po skończonej liczbie kroków,
bo stopnie niezerowych reszt są nieujemnymi liczbami całkowitymi
i ściśle maleją. Sumując wykonane odejmowania,
otrzymujemy $f=qg+r$ z wymaganym ograniczeniem.

Dla jednoznaczności załóżmy, że
$f=qg+r=q'g+r'$ z $\deg r,\deg r'<d$ dla niezerowych
reszt. Wtedy $(q-q')g=r'-r$. Jeśli $q\ne q'$,
wyraz wiodący $(q-q')g$ jest niezerowy: współczynnik
wiodący $g$ jest jednostką, więc nie może wyzerować
niezerowego współczynnika wiodącego $q-q'$.
Stąd $\deg((q-q')g)=\deg(q-q')+d\ge d$,
co jest sprzeczne z tym, że $r'-r$ jest wielomianem
zerowym albo ma stopień mniejszy niż $d$. Zatem
$q=q'$ i $r=r'$. 
\end{proof}

\begin{wniosek}[Twierdzenie o czynniku liniowym]
\label{aa:wn:poly-factor}
Jeśli $R$ jest przemiennym pierścieniem z jedynką,
$f\in R[x]$ i $a\in R$, to $f(a)=0$ wtedy i tylko wtedy,
gdy $x-a$ dzieli $f$ w $R[x]$.
\end{wniosek}
\begin{proof}
Jeśli $R$ jest pierścieniem zerowym, to $a=0_R$,
$f=0_{R[x]}$ i $x-a=0_{R[x]}$; zachodzą zarówno
$f(a)=0_R$, jak i $(x-a)\mid f$, gdyż $0=0\cdot0$. Załóżmy zatem,
że $R$ jest niezerowy. Wtedy $x-a$ jest niezerowym
wielomianem o współczynniku wiodącym $1_R$,
więc poprzednie twierdzenie daje $f=(x-a)q+r$,
gdzie $r$ jest stałą. Po podstawieniu $x=a$ dostajemy
$f(a)=r$. Zatem $f(a)=0$ jest równoważne $r=0$,
czyli podzielności przez $x-a$.
\end{proof}

\begin{twierdzenie}[Ograniczenie liczby pierwiastków]
\label{aa:tw:poly-roots}
Jeżeli $R$ jest dziedziną całkowitości, a
$f\in R[x]$ niezerowym wielomianem stopnia $n$, to
$f$ ma co najwyżej $n$ różnych pierwiastków w $R$.
\end{twierdzenie}
\begin{proof}
Indukujemy po $n$. Dla $n=0$ niezerowa stała
nie ma pierwiastków. Niech $n\ge1$. Jeśli $f$ nie ma
pierwiastków, twierdzenie zachodzi. W przeciwnym razie
wybierzmy pierwiastek $a$. Z wniosku poprzedniego wynika
$f=(x-a)g$, a porównanie stopni daje $\deg g=n-1$.
Jeśli $b$ jest innym pierwiastkiem $f$, to
$0=f(b)=(b-a)g(b)$. Ponieważ $b-a\ne0$ i $R$ nie
ma dzielników zera, otrzymujemy $g(b)=0$.
Z założenia indukcyjnego wielomian $g$ ma najwyżej
$n-1$ różnych pierwiastków, a dodanie $a$ daje
co najwyżej $n$ pierwiastków $f$.
\end{proof}

\begin{twierdzenie}[Funkcje wielomianowe na ciele skończonym]
\label{aa:tw:finite-field-polynomial-functions}
Dla liczby pierwszej $p$ każda funkcja
$F:\mathbb F_p\to\mathbb F_p$ jest zadana przez
dokładnie jeden wielomian zerowy lub stopnia mniejszego niż $p$.
Jądro homomorfizmu
\[
 \Theta:\mathbb F_p[x]\longrightarrow
 \{F:\mathbb F_p\to\mathbb F_p\},\qquad
 \Theta(f)(a)=f(a),
\]
gdzie po prawej działania wykonuje się punktowo,
jest ideałem $(x^p-x)$.
\end{twierdzenie}
\begin{proof}
Dla $a\in\mathbb F_p$ rozważmy wielomian Lagrange'a
$L_a(x)=\prod_{b\in\mathbb F_p,\,b\ne a}
 (x-b)/(a-b)$. Wszystkie mianowniki są niezerowe,
więc odwracalne. Mamy $L_a(a)=1$ oraz $L_a(b)=0$
dla $b\ne a$. Wielomian
$P_F(x)=\sum_{a\in\mathbb F_p}F(a)L_a(x)$
ma stopień co najwyżej $p-1$ i spełnia
$P_F(b)=F(b)$ dla każdego $b\in\mathbb F_p$.

Jeżeli dwa wielomiany stopnia mniejszego niż $p$
wyznaczają tę samą funkcję, ich różnica ma co najmniej
$p$ różnych pierwiastków. Na mocy
Twierdzenia~\ref{aa:tw:poly-roots} ta różnica musi być
wielomianem zerowym.

Wielomian $x^p-x$ zeruje się we wszystkich punktach,
więc jego wszystkie wielokrotności należą do jądra.
Dla dowolnego $f$ dzielenie z resztą przez unormowany
$x^p-x$ daje $f=q(x^p-x)+r$ z $r=0$ lub
$\deg r<p$. Jeśli $f$ jest w jądrze, to $r$
zeruje się we wszystkich $p$ punktach, a więc $r=0$.
Odwrotne zawieranie jest oczywiste, co kończy dowód.
\end{proof}

\begin{przyklad}[Dlaczego nie można opuścić założenia o dziedzinie]
\label{aa:prz:poly-roots-z8}
W $\mathbb Z/8\mathbb Z$ wielomian $x^2-1$
ma cztery różne pierwiastki: $[1],[3],[5],[7]$.
Nie przeczy to twierdzeniu, ponieważ pierścień
$\mathbb Z/8\mathbb Z$ nie jest dziedziną całkowitości.
Różnica wynika z zerowania iloczynu $(b-a)g(b)$
przy obu niezerowych czynnikach.
\end{przyklad}

\begin{wniosek}[Wielomian nad nieskończoną dziedziną jest określony przez wartości]
\label{aa:wn:poly-evaluation-infinite-domain}
Niech $R$ będzie nieskończoną dziedziną całkowitości.
Jeśli $f,g\in R[x]$ spełniają $f(a)=g(a)$ dla każdego
$a\in R$, to $f=g$. Równoważnie, odwzorowanie przypisujące
wielomianowi jego funkcję wartości na $R$ jest różnowartościowe.
\end{wniosek}
\begin{proof}
Połóżmy $h=f-g$. Z założenia $h(a)=0$ dla każdego
$a\in R$. Gdyby $h\ne0$, to wielomian $h$ miałby
skończony stopień $n$ i, na mocy
Twierdzenia~\ref{aa:tw:poly-roots}, najwyżej $n$
różnych pierwiastków w $R$. Tymczasem każdy element
nieskończonego zbioru $R$ byłby jego pierwiastkiem.
Otrzymana sprzeczność daje $h=0$, czyli $f=g$.
\end{proof}
\begin{uwaga}
Założenia o nieskończoności i braku dzielników zera
pełnią różne funkcje. Nad ciałem skończonym odwzorowanie
wartościowania może nie być różnowartościowe
(Przykład~\ref{aa:prz:poly-finite-field}); w pierścieniu
z dzielnikami zera zawodzi oszacowanie liczby pierwiastków
(Przykład~\ref{aa:prz:poly-roots-z8}).
\end{uwaga}

\section{Wielomiany pierwotne i lemat Gaussa}
\begin{definicja}[Ciało ułamków]
\label{aa:def:fracfield}
Niech $R$ będzie dziedziną całkowitości. Jej
\emph{ciałem ułamków} jest ciało $K=\operatorname{Frac}(R)$
przedstawione przez klasy par $(a,b)$, gdzie $a\in R$
i $b\in R\setminus\{0\}$, przy relacji
$(a,b)\sim(c,d)$ wtedy i tylko wtedy, gdy $ad=bc$.
Klasę pary $(a,b)$ oznaczamy $a/b$.
Dodawanie i mnożenie ułamków wyznaczają wzory
$a/b+c/d=(ad+bc)/(bd)$ oraz $(a/b)(c/d)=ac/(bd)$.
Odwzorowanie $R\to K$, $a\mapsto a/1$, jest zanurzeniem.
\end{definicja}

\begin{lemat}[Uzasadnienie konstrukcji ciała ułamków]
\label{aa:lem:fracfield-well-defined}
Niech $R$ będzie dziedziną całkowitości. Relacja
$(a,b)\sim(c,d)$, określona przez $ad=bc$ na
$R\times(R\setminus\{0\})$, jest relacją równoważności.
Wzory na dodawanie i mnożenie ułamków nie zależą od
wyboru przedstawicieli, a zbiór klas z tymi działaniami
jest ciałem zawierającym izomorficzną kopię $R$.
\end{lemat}
\begin{proof}
Zwrotność i symetria są bezpośrednie. Jeśli
$ad=bc$ oraz $cf=de$, to $adf=bcf=bde$.
Ponieważ $d\ne0$ i $R$ nie ma dzielników zera,
możemy skrócić $d$, otrzymując $af=be$;
to dowodzi przechodniości.

Sprawdźmy niezależność dodawania od przedstawicieli.
Załóżmy $ab'=a'b$ i $cd'=c'd$ z niezerowymi
$b,b',d,d'$. Wówczas
\[
 (ad+bc)b'd'=ab'dd'+bb'cd'
             =a'bd'd+bb'c'd
             =(a'd'+b'c')bd.
\]
Stąd $(ad+bc)/(bd)=(a'd'+b'c')/(b'd')$.
Dla mnożenia odpowiednia równość ma postać
$acb'd'=(ab')(cd')=(a'b)(c'd)=a'c'bd$.
Mianowniki $bd$ są niezerowe, gdyż $R$ jest dziedziną.

Łączność, przemienność i rozdzielność sprowadzają się
po rozwinięciu liczników do odpowiednich własności $R$.
Dla przykładu
$(a/b+c/d)+e/f=(adf+bcf+ebd)/(bdf)$, niezależnie
od nawiasowania. Zerem jest $0/1$, jedynką $1/1$,
a przeciwnym do $a/b$ element $(-a)/b$.
Z równości $a/b=0/1$ wynika $a=0$.
Jeśli $a\ne0$, to $b/a$ jest poprawnym ułamkiem
i $(a/b)(b/a)=1/1$. Zatem otrzymujemy ciało.
Ponadto $a\mapsto a/1$ zachowuje działania i jest
różnowartościowe, bo $a/1=c/1$ wymusza $a=c$.
Dalsze własności tej konstrukcji omawiamy w
dalszej części skryptu.
\end{proof}

\begin{definicja}[Zawartość wielomianu]
\label{aa:def:poly-content}
Niech $R$ będzie dziedziną z jednoznacznym rozkładem
na czynniki i niech $f=\sum_{j=0}^n a_jx^j\in R[x]$
będzie niezerowy. \emph{Zawartością} $f$ nazywamy
największy wspólny dzielnik jego niezerowych współczynników,
wyznaczony z dokładnością do stowarzyszenia; oznaczamy go
$c(f)$. Wielomian $f$ jest \emph{pierwotny}, gdy $c(f)$
jest jednostką, czyli żaden element pierwszy nie dzieli
wszystkich współczynników $f$. W szczególności
$f=c(f)f_0$ dla pewnego pierwotnego $f_0\in R[x]$.
\end{definicja}

\begin{lemat}[Iloczyn wielomianów pierwotnych]
\label{aa:lem:gauss-primitive}
Niech $R$ będzie dziedziną z jednoznacznym rozkładem
na czynniki. Jeżeli $f,g\in R[x]$ są pierwotne,
to $fg$ jest pierwotny.
\end{lemat}
\begin{proof}
Przypuśćmy, że $fg$ nie jest pierwotny. Wtedy
istnieje element pierwszy $p\in R$ dzielący każdy
współczynnik iloczynu. Ponieważ $p$ jest pierwszy,
ideał $(p)$ jest pierwszy: jeśli $ab\in(p)$, to
$p\mid ab$, zatem $p\mid a$ lub $p\mid b$.
Na mocy kryterium z rozdziału~\ref{aa:ch:pierscienie}
iloraz $R/(p)$ jest dziedziną całkowitości.
Zredukujmy współczynniki $f,g$ modulo $(p)$.
Otrzymamy wielomiany $\bar f,\bar g\in(R/(p))[x]$.
Pierwotność oznacza, że przynajmniej jeden współczynnik
każdego z nich nie jest podzielny przez $p$,
więc $\bar f\ne0$ i $\bar g\ne0$.
Ale wszystkie współczynniki $fg$ są podzielne przez $p$,
więc $\bar f\bar g=0$. To niemożliwe w pierścieniu
wielomianów nad dziedziną całkowitości na mocy
Lematu~\ref{aa:lem:poly-degree}. Otrzymujemy sprzeczność.
\end{proof}

\begin{twierdzenie}[Lemat Gaussa o zawartości]
\label{aa:tw:gauss-content}
Niech $R$ będzie dziedziną z jednoznacznym rozkładem
na czynniki i niech $f,g\in R[x]$ będą niezerowe.
Wówczas $c(fg)\sim c(f)c(g)$, gdzie $\sim$ oznacza
stowarzyszenie w $R$.
\end{twierdzenie}
\begin{proof}
Wybierzmy zawartości $c(f),c(g)$ i zapiszmy
$f=c(f)f_0$, $g=c(g)g_0$ z pierwotnymi $f_0,g_0$.
Wtedy $fg=c(f)c(g)(f_0g_0)$.
Poprzedni lemat mówi, że $f_0g_0$ jest pierwotny.
Zatem największy wspólny dzielnik współczynników iloczynu
jest stowarzyszony z $c(f)c(g)$, co daje tezę.
\end{proof}

\begin{lemat}[Usuwanie mianowników z rozkładu]
\label{aa:lem:gauss-denominators}
Niech $R$ będzie dziedziną z jednoznacznym rozkładem
na czynniki, $K=\operatorname{Frac}(R)$, a
$f\in R[x]$ wielomianem pierwotnym. Jeżeli
$f=gh$ dla wielomianów $g,h\in K[x]$ o dodatnich stopniach,
to istnieją wielomiany $g_0,h_0\in R[x]$ o dodatnich
stopniach takie, że $f=g_0h_0$.
\end{lemat}
\begin{proof}
Każdy niezerowy wielomian $g\in K[x]$ można zapisać
jako $g=\alpha G$, gdzie $\alpha\in K^\times$,
a $G\in R[x]$ jest pierwotny: najpierw mnożymy przez
wspólny niezerowy mianownik współczynników, a potem
dzielimy uzyskany wielomian przez jego zawartość.
Podobnie $h=\beta H$ z pierwotnym $H\in R[x]$.
Zatem $f=\lambda GH$, gdzie $\lambda=\alpha\beta\in K^\times$,
a $GH$ jest pierwotny na mocy Lematu~\ref{aa:lem:gauss-primitive}.

Zapiszmy $\lambda=a/b$ z $a,b\in R$, $b\ne0$,
przy czym $a$ i $b$ nie mają wspólnego czynnika nierozkładalnego.
Z równości $bf=aGH$ wynika, że $b$ dzieli każdy
współczynnik wielomianu $aGH$. Jeśli $b$ nie byłby
jednostką, miałby czynnik pierwszy $p$. Wówczas $p$
dzieliłby każdy współczynnik $aGH$; ponieważ $p\nmid a$,
pierwszość $p$ wymuszałaby, że $p$ dzieli wszystkie
współczynniki $GH$, wbrew pierwotności $GH$.
Stąd $b$ jest jednostką, czyli $\lambda\in R$.
Jeśli $\lambda$ nie byłaby jednostką, dowolny jej
czynnik pierwszy dzieliłby wszystkie współczynniki $f$,
co przeczyłoby pierwotności $f$. Ostatecznie
$\lambda\in R^\times$ i $f=(\lambda G)H$ jest rozkładem
w $R[x]$ na dwa wielomiany dodatniego stopnia.
\end{proof}

\begin{lemat}[Pierwotne wielomiany stowarzyszone nad ciałem ułamków]
\label{aa:lem:primitive-associates}
Niech $R$ będzie dziedziną z jednoznacznym rozkładem
na czynniki, a $K=\operatorname{Frac}(R)$.
Jeżeli dwa niezerowe pierwotne wielomiany
$f,g\in R[x]$ spełniają $f=\lambda g$ dla
pewnego $\lambda\in K^\times$, to $\lambda\in R^\times$.
\end{lemat}
\begin{proof}
Wybierzmy $\lambda=a/b$, gdzie $a,b\in R$ są
niezerowe i nie mają wspólnego czynnika nierozkładalnego.
Równość $bf=ag$ zachodzi współczynnikowo.
Gdyby $b$ nie było jednostką, dowolny czynnik pierwszy
$p\mid b$ dzieliłby każdy współczynnik $ag$.
Ponieważ $p\nmid a$, każdy współczynnik $g$ byłby
podzielny przez $p$, wbrew pierwotności $g$.
Zatem $b\in R^\times$ i $\lambda\in R$.
Gdyby $a$ nie było jednostką, pewien czynnik pierwszy
$p\mid a$ dzieliłby wszystkie współczynniki $bf$,
a ponieważ $p\nmid b$, również wszystkie współczynniki
$f$, znów sprzecznie z pierwotnością.
Ostatecznie $a,b\in R^\times$, a więc $\lambda\in R^\times$.
\end{proof}

\begin{twierdzenie}[Kryterium nierozkładalności Gaussa]
\label{aa:tw:gauss-irreducibility}
Niech $R$ będzie dziedziną z jednoznacznym rozkładem
na czynniki, $K=\operatorname{Frac}(R)$, a
$f\in R[x]$ pierwotnym wielomianem dodatniego stopnia.
Wówczas $f$ jest nierozkładalny w $R[x]$ wtedy i tylko
wtedy, gdy jest nierozkładalny w $K[x]$.
\end{twierdzenie}
\begin{proof}
Jeżeli $f$ rozkłada się na dwa czynniki dodatniego stopnia
w $K[x]$, to poprzedni lemat daje taki sam rodzaj rozkładu
w $R[x]$, więc $f$ nie jest nierozkładalny w $R[x]$.

Odwrotnie, przypuśćmy, że $f=uv$ w $R[x]$, gdzie
$u,v$ są niejednostkami $R[x]$. Gdyby jeden z nich
był stały, byłby nieodwracalnym elementem $R$, który
dzieliłby wszystkie współczynniki $f$. To przeczy
pierwotności. Zatem $u$ i $v$ mają dodatnie stopnie,
więc stanowią właściwy rozkład $f$ także w $K[x]$.
W ten sposób dowiedliśmy obu implikacji.
\end{proof}

\begin{twierdzenie}[Twierdzenie Gaussa o pierścieniu wielomianów]
\label{aa:tw:ufd-polynomial}
Jeżeli $R$ jest dziedziną z jednoznacznym rozkładem
na czynniki, to pierścień $R[x]$ także jest dziedziną
z jednoznacznym rozkładem na czynniki. W konsekwencji,
jeżeli $K$ jest ciałem, to $K[x_1,\ldots,x_n]$
ma jednoznaczny rozkład na czynniki dla każdego $n\ge1$.
\end{twierdzenie}
\begin{proof}
Pierścień $R[x]$ jest dziedziną całkowitości na mocy
Lematu~\ref{aa:lem:poly-degree}. Ustalmy
$K=\operatorname{Frac}(R)$. Niezerowy wielomian
$f\in R[x]$ zapisujemy jako $f=c(f)f_0$,
gdzie $f_0$ jest pierwotny. Element $c(f)$ rozkłada się
w $R$ na czynniki nierozkładalne, jeśli nie jest jednostką.
Z kolei $K[x]$ jest dziedziną euklidesową
(Rozdział~\ref{aa:ch:podzielnosc}), a więc ma
jednoznaczny rozkład na czynniki. Rozłóżmy $f_0$ w $K[x]$
na wielomiany nierozkładalne, przy czym dla stałego
pierwotnego $f_0$ przyjmujemy pusty iloczyn i $r=0$.
Wybierając dla każdego czynnika pierwotnego reprezentanta
w $R[x]$, otrzymujemy $f_0=\lambda g_1\cdots g_r$
z pierwotnymi $g_i\in R[x]$ i $\lambda\in K^\times$.
Dla $r=0$ mamy $f_0=\lambda\in R^\times$.
Lemat~\ref{aa:lem:primitive-associates}, zastosowany do
pierwotnych $f_0$ oraz $g_1\cdots g_r$, pokazuje,
że $\lambda\in R^\times$.
Z kryterium nierozkładalności Gaussa każdy $g_i$ jest
nierozkładalny w $R[x]$. W ten sposób otrzymaliśmy
istnienie rozkładu na czynniki nierozkładalne.

Sprawdźmy jednoznaczność. W każdym rozkładzie na czynniki
nierozkładalne w $R[x]$ czynniki stałe odpowiadają
czynnikom nierozkładalnym zawartości $c(f)$; pozostałe
czynniki mają dodatnie stopnie i muszą być pierwotne,
gdyż niezerowa niejednostka dzieląca wszystkie ich
współczynniki dawałaby właściwy rozkład.
Takie czynniki są nierozkładalne również w $K[x]$.
Jednoznaczność rozkładu zawartości w $R$ i rozkładu
w $K[x]$ zapewnia sparowanie czynników.
Jeśli dwa pierwotne wielomiany różnią się mnożnikiem
$\lambda\in K^\times$, to na mocy
Lematu~\ref{aa:lem:primitive-associates} mamy
$\lambda\in R^\times$. Czynniki odpowiadają sobie
zatem z dokładnością do jednostek $R[x]$.

Wreszcie ciało $K$ samo jest dziedziną z jednoznacznym
rozkładem, a powtarzając wykazany wynik kolejno dla
$K[x_1]$, $K[x_1][x_2]$, $\ldots$ otrzymujemy tezę
o dowolnej skończonej liczbie zmiennych.
\end{proof}

\begin{przyklad}[Jednoznaczny rozkład bez główności wszystkich ideałów]
\label{aa:prz:ufd-not-pid}
Dla dowolnego ciała $K$ pierścień $K[x,y]$ ma
jednoznaczny rozkład na czynniki na mocy poprzedniego
twierdzenia. Jednak ideał $(x,y)$ nie jest główny:
wykazaliśmy to już szczegółowo w
Przykładzie~\ref{aa:prz:two-polynomials-not-pid}.
Oznacza to, że implikacja
„dziedzina z jednoznacznym rozkładem na czynniki
$\Rightarrow$ dziedzina ideałów głównych” jest fałszywa.
\end{przyklad}

\section{Kryterium Eisensteina i badanie redukcji modulo liczbę pierwszą}
\begin{twierdzenie}[Kryterium Eisensteina]
\label{aa:tw:eisenstein}
Niech $R$ będzie dziedziną z jednoznacznym rozkładem
na czynniki, $K=\operatorname{Frac}(R)$, a
$f=a_nx^n+\cdots+a_1x+a_0\in R[x]$ wielomianem
stopnia $n\ge1$. Załóżmy, że istnieje element pierwszy
$p\in R$ taki, że $p\nmid a_n$, $p\mid a_j$ dla
$0\le j<n$ oraz $p^2\nmid a_0$.
Wtedy $f$ jest nierozkładalny w $K[x]$.
\end{twierdzenie}
\begin{proof}
Niech $c=c(f)$ będzie wybraną zawartością $f$
i niech $f_0=f/c\in R[x]$. Z warunku
$p\nmid a_n$ wynika $p\nmid c$, a $f_0$ jest pierwotny.
Ponieważ $a_j=c\,b_j$, gdzie $b_j$ jest współczynnikiem
$f_0$, oraz $p\nmid c$, pierwszość $p$ pokazuje,
że $f_0$ spełnia wszystkie trzy warunki kryterium:
$p$ pozostaje dzielnikiem niższych współczynników,
nie dzieli współczynnika wiodącego, a $p^2$ nie dzieli
wyrazu wolnego, ponieważ $p\nmid c$.

Gdyby $f$ był rozkładalny w $K[x]$, byłby rozkładalny
także $f_0$. Z kryterium nierozkładalności Gaussa
otrzymalibyśmy $f_0=gh$ w $R[x]$, gdzie obydwa czynniki
mają dodatnie stopnie. Zredukujmy współczynniki modulo
ideał $(p)$. Ponieważ $R/(p)$ jest dziedziną,
redukcja daje równość
$\overline g\,\overline h=\overline{a_n/c}\,x^n$.
Współczynnik wiodący po prawej jest niezerowy,
a oba czynniki po lewej mają dodatnie stopnie:
ich współczynniki wiodące nie mogą być podzielne przez
$p$, gdyż ich iloczyn jest współczynnikiem wiodącym $f_0$.

W dziedzinie całkowitości każdy dzielnik monomianu
$c x^n$ z $c\ne0$ jest sam monomianem z niezerowym
współczynnikiem. Można to także zobaczyć bezpośrednio:
jeśli $\bar g$ miałby niezerowy wyraz niższego stopnia,
to jego najniższy niezerowy współczynnik, pomnożony
przez najniższy niezerowy współczynnik $\bar h$,
dałby niezerowy wyraz iloczynu o wykładniku mniejszym
niż $n$. Zatem $\bar g$ i $\bar h$ są monomianami
dodatnich stopni. Ich wyrazy wolne są więc zerowe,
co oznacza $p\mid g(0)$ i $p\mid h(0)$. Wtedy
$p^2\mid g(0)h(0)=f_0(0)$, sprzecznie z założeniem.
Stąd $f$ jest nierozkładalny w $K[x]$.
\end{proof}

\begin{przyklad}[Wielomiany $x^n-p$]
\label{aa:prz:xn-p}
Dla liczby pierwszej $p$ i $n\ge1$ wielomian
$x^n-p\in\mathbb Z[x]$ spełnia kryterium Eisensteina
z elementem pierwszym $p$. Współczynnik wiodący
wynosi $1$, wszystkie pozostałe współczynniki
są podzielne przez $p$, a $p^2$ nie dzieli wyrazu
wolnego $-p$. Zatem $x^n-p$ jest nierozkładalny
nad $\mathbb Q$. W szczególności $x^3-2$ jest
wielomianem minimalnym liczby $\sqrt[3]{2}$ nad $\mathbb Q$.
\end{przyklad}

\begin{przyklad}[Zmiana zmiennej ujawniająca kryterium]
\label{aa:prz:x4plus1}
Wielomian $f(x)=x^4+1$ nie spełnia wprost kryterium
Eisensteina z liczbą pierwszą $2$, ponieważ jego
wyraz wolny nie jest podzielny przez $2$. Rozważmy jednak
$f(x+1)=x^4+4x^3+6x^2+4x+2$.
Wszystkie współczynniki oprócz wiodącego są parzyste,
a $4$ nie dzieli wyrazu wolnego $2$.
Kryterium Eisensteina daje nierozkładalność $f(x+1)$
w $\mathbb Q[x]$. Odwzorowanie
$\mathbb Q[x]\to\mathbb Q[x]$, $g(x)\mapsto g(x+1)$,
jest automorfizmem pierścieni o odwrotności
$h(x)\mapsto h(x-1)$. Automorfizm zachowuje jednostki
i właściwe rozkłady; dlatego także $x^4+1$ jest
nierozkładalny nad $\mathbb Q$.
\end{przyklad}

\begin{twierdzenie}[Kryterium nierozkładalności przez redukcję modulo $p$]
\label{aa:tw:modp-irred}
Niech $f\in\mathbb Z[x]$ będzie pierwotnym wielomianem
dodatniego stopnia, a $p$ liczbą pierwszą, która nie
dzieli współczynnika wiodącego $f$. Jeśli wielomian
$\bar f\in\mathbb F_p[x]$ otrzymany przez redukcję
współczynników modulo $p$ jest nierozkładalny, to
$f$ jest nierozkładalny nad $\mathbb Q$.
\end{twierdzenie}
\begin{proof}
Gdyby $f$ był rozkładalny nad $\mathbb Q$, lemat
Gaussa dostarczyłby właściwego rozkładu
$f=gh$ z $g,h\in\mathbb Z[x]$ dodatnich stopni.
Współczynniki wiodące $g$ i $h$ mnożą się do
współczynnika wiodącego $f$, który nie jest podzielny
przez $p$. Żaden z tych dwóch współczynników nie jest
więc podzielny przez $p$. Po redukcji modulo $p$
stopnie $\bar g,\bar h$ pozostają dodatnie i nie zmieniają
się, a $\bar f=\bar g\bar h$ stanowi właściwy rozkład
w $\mathbb F_p[x]$. Sprzeczność dowodzi tezy.
\end{proof}

\begin{przyklad}[Kryterium redukcyjne jest jednostronne]
\label{aa:prz:modp-not-necessary}
Wielomian $x^2+1$ jest nierozkładalny nad $\mathbb Q$:
gdyby rozkładał się na czynniki liniowe, miałby pierwiastek
w $\mathbb Q$, a równanie $x^2=-1$ nie ma rozwiązania
nawet w $\mathbb R$. Tymczasem modulo $2$ mamy
$x^2+1=(x+1)^2$ w $\mathbb F_2[x]$.
Rozkładalność redukcji nie oznacza zatem
rozkładalności oryginalnego wielomianu.
\end{przyklad}

\begin{przyklad}[Nierozkładalność stopnia trzeciego wykazana modulo $2$]
\label{aa:prz:cubic-mod2-irreducibility}
Weźmy $f(x)=x^3+x+1\in\mathbb Z[x]$.
Współczynnik wiodący $1$ nie jest podzielny przez $2$.
Redukcja $\bar f\in\mathbb F_2[x]$ nie ma pierwiastków:
$\bar f(0)=1$, $\bar f(1)=1+1+1=1$.
Wielomian stopnia $3$ nad ciałem jest rozkładalny
wtedy i tylko wtedy, gdy ma czynnik stopnia $1$:
przy rozkładzie na dwa czynniki dodatniego stopnia
ich stopnie sumują się do $3$, więc jeden wynosi $1$.
Brak pierwiastków oznacza zatem, że $\bar f$
jest nierozkładalny. Kryterium redukcyjne daje
nierozkładalność $x^3+x+1$ nad $\mathbb Q$.
\end{przyklad}

\begin{przyklad}[Dlaczego trzeba zachować stopień przy redukcji]
\label{aa:prz:modp-leading-coefficient-necessary}
Wielomian $f(x)=2x^2+3x+1$ jest pierwotny:
jego współczynniki mają największy wspólny dzielnik
równy $1$. Mimo to $f$ rozkłada się nad
$\mathbb Q$ jako
$f(x)=(2x+1)(x+1)$, gdzie oba czynniki mają
stopień dodatni. Redukcja modulo $2$ daje natomiast
$\bar f(x)=x+1\in\mathbb F_2[x]$, który jest
nierozkładalny jako wielomian stopnia pierwszego.
Nie ma sprzeczności z
Twierdzeniem~\ref{aa:tw:modp-irred}: liczba $2$
dzieli współczynnik wiodący $f$, więc stopień
spada z dwóch do jednego. Po redukcji pierwszy
czynnik rozkładu staje się stałą $1$.
\end{przyklad}

\section{Pierwiastki wymierne i analiza konkretnych wielomianów}
\begin{twierdzenie}[Twierdzenie o pierwiastkach wymiernych]
\label{aa:tw:rational-root}
Niech $f=a_nx^n+\cdots+a_1x+a_0\in\mathbb Z[x]$
będzie wielomianem stopnia $n\ge1$ i niech $a_n\ne0$.
Jeśli $u/v\in\mathbb Q$ jest pierwiastkiem $f$,
gdzie $u\in\mathbb Z$, $v\in\mathbb N$ jest dodatnie
oraz $\gcd(u,v)=1$, to $u\mid a_0$ i $v\mid a_n$.
W przypadku $u=0$ pierwszy warunek oznacza $a_0=0$.
\end{twierdzenie}
\begin{proof}
Równość $f(u/v)=0$ po pomnożeniu przez $v^n$
daje
\[
 a_nu^n+a_{n-1}u^{n-1}v+\cdots+a_1uv^{n-1}+a_0v^n=0.
\]
Wyłączając pierwszy składnik, stwierdzamy, że
$v$ dzieli $a_nu^n$. Ponieważ $\gcd(u,v)=1$,
żaden element pierwszy dzielący $v$ nie dzieli $u$.
Lemat Euklidesa zastosowany wielokrotnie do potęg
czynników pierwszych $v$ daje $v\mid a_n$.
Analogicznie po wyłączeniu ostatniego składnika
otrzymujemy $u\mid a_0v^n$. Jeżeli $u\ne0$,
to $\gcd(u,v)=1$ daje $u\mid a_0$. Jeżeli $u=0$,
równość $f(0)=0$ oznacza $a_0=0$.
\end{proof}

\begin{przyklad}[Od listy kandydatów do pełnego rozkładu]
\label{aa:prz:rational-root-cubic}
Rozważmy $f(x)=2x^3-3x+1$. Każdy pierwiastek wymierny
w postaci nieskracalnej ma licznik dzielący $1$
i dodatni mianownik dzielący $2$. Kandydatami są
$\pm1$, $\pm\tfrac12$. Podstawienie $x=1$ daje
$f(1)=2-3+1=0$, więc z twierdzenia o czynniku liniowym
wiemy, że $x-1$ dzieli $f$.
Dzielenie wielomianów daje
$f(x)=(x-1)(2x^2+2x-1)$.
Równanie kwadratowe $2x^2+2x-1=0$ ma pierwiastki
$(-1\pm\sqrt3)/2$. Zatem otrzymaliśmy wszystkie
pierwiastki nad $\mathbb R$. Czynnika kwadratowego
nie można rozłożyć nad $\mathbb Q$, ponieważ
$\sqrt3\notin\mathbb Q$: gdyby $\sqrt3=m/n$
dla względnie pierwszych $m,n\in\mathbb Z$, $n\ne0$,
to $m^2=3n^2$. Stąd $3\mid m$, zatem po zapisaniu
$m=3k$ także $3\mid n$; byłaby to sprzeczność
z $\gcd(m,n)=1$.
\end{przyklad}

\begin{przyklad}[Brak pierwiastków wymiernych nie wystarcza w stopniu czwartym]
\label{aa:prz:rational-root-insufficient}
Wielomian $f(x)=x^4+4$ nie ma pierwiastków rzeczywistych,
ponieważ $x^4+4>0$ dla każdego $x\in\mathbb R$,
a więc nie ma też pierwiastków wymiernych.
Mimo to w $\mathbb Q[x]$ jest rozkładalny:
\[
 x^4+4=(x^2+2)^2-(2x)^2
       =(x^2-2x+2)(x^2+2x+2).
\]
Oba czynniki są niestałe. Twierdzenie o pierwiastkach
wymiernych pozwala więc wykluczyć rozkład stopnia $2$
lub $3$ (wtedy rozkład musi mieć czynnik liniowy),
lecz samo w sobie nie dowodzi nierozkładalności
wielomianów stopnia co najmniej $4$.
\end{przyklad}

\section{Pierścienie wielomianów wielu zmiennych i ich ilorazy}
\begin{definicja}[Pierścień wielomianów wielu zmiennych]
\label{aa:def:multi-polynomial}
Niech $R$ będzie przemiennym pierścieniem z jedynką
i niech $n\ge1$ będzie liczbą całkowitą.
Przyjmujemy $R[x_1]$ z Definicji~\ref{aa:def:poly-ring},
a dla $n\ge2$ definiujemy indukcyjnie
$R[x_1,\ldots,x_n]=(R[x_1,\ldots,x_{n-1}])[x_n]$.
Jego elementy są skończonymi sumami
$a_{\alpha}x_1^{\alpha_1}\cdots x_n^{\alpha_n}$,
gdzie $a_{\alpha}\in R$, a $\alpha_i\in\mathbb N_0$.
\end{definicja}

\begin{twierdzenie}[Pierścień współrzędnych paraboli]
\label{aa:tw:parabola-quotient}
Niech $K$ będzie ciałem. Homomorfizm pierścieni
$\psi:K[x,y]\to K[x]$, zadany wzorem
$\psi(F(x,y))=F(x,x^2)$, jest surjektywny, a jego
jądrem jest ideał $(y-x^2)$. W szczególności
$K[x,y]/(y-x^2)\cong K[x]$.
\end{twierdzenie}
\begin{proof}
Homomorfizm $\psi$ zachowuje sumy i iloczyny, ponieważ
jest podstawieniem, a jedynkę wysyła na jedynkę.
Jest surjektywny: każdy $h(x)\in K[x]$ jest obrazem
tego samego wielomianu traktowanego jako element
$K[x,y]$ niezależny od $y$.

Ponadto $\psi(y-x^2)=0$, więc ideał $(y-x^2)$ jest
zawarty w jądrze $\psi$. Aby udowodnić odwrotne
zawieranie, potraktujmy $F\in K[x,y]$ jako wielomian
zmiennej $y$ o współczynnikach z pierścienia $K[x]$.
Wielomian $y-x^2$ jest unormowany względem $y$.
Twierdzenie~\ref{aa:tw:poly-division} daje jednoznacznie
$F(x,y)=Q(x,y)(y-x^2)+H(x)$, ponieważ reszta ma
stopień względem $y$ mniejszy niż jeden.
Po podstawieniu $y=x^2$ otrzymujemy
$F(x,x^2)=H(x)$. Jeżeli $F\in\ker\psi$, to $H=0$
i $F\in(y-x^2)$. Mamy więc równość jądra i ideału.
Pierwsze twierdzenie o izomorfizmie pierścieni
(Rozdział~\ref{aa:ch:pierscienie}) daje wówczas
$K[x,y]/(y-x^2)\cong K[x]$.
\end{proof}

\begin{przyklad}[Równania opisujące ten sam zbiór nie muszą dawać tych samych pierścieni]
\label{aa:prz:poly-nilpotent-double}
W pierścieniu $\mathbb R[x]$ rozważmy ideały
$(x)$ oraz $(x^2)$. Zbiory rzeczywistych zer wielomianów
$x$ i $x^2$ są identyczne: w obu przypadkach składają
się z jednego punktu $0$. Jednak
$\mathbb R[x]/(x)\cong\mathbb R$ jest ciałem,
natomiast w $\mathbb R[x]/(x^2)$ klasa $x+(x^2)$ jest
niezerowa i nilpotentna stopnia dwa.
Niezerowość tej klasy wynika z faktu, że $x$ nie jest
wielokrotnością $x^2$. Oba pierścienie nie mogą być
izomorficzne: izomorfizm zachowywałby niezerowe elementy
nilpotentne, których w ciele nie ma. W następnym rozdziale
wyjaśnimy ten przykład za pomocą pierwiastnika ideału.
\end{przyklad}

\begin{uwaga}[Przejście do teorii Galois]
\label{aa:uw:poly-galois}
Kryteria nierozkładalności rozstrzygają, kiedy iloraz
$K[x]/(f)$ jest ciałem: na mocy
Twierdzenia~\ref{aa:tw:prime-max-quotient} zachodzi to
dokładnie wtedy, gdy ideał $(f)$ jest maksymalny.
W $K[x]$ nierozkładalność wielomianu $f$ dodatniego
stopnia jest równoważna maksymalności ideału $(f)$;
wynika to z algorytmu Euklidesa, który daje
odwrotność klasy każdego wielomianu $g$ niepodzielnego
przez $f$. Ta obserwacja pozwoli konstruować
rozszerzenia ciał i przejść do teorii Galois.
\end{uwaga}

% AA-12 | Algebra abstrakcyjna | Iteracja VI
\chapter{Pierwiastniki ideałów i lokalizacja pierścieni}
\chaptermark{Pierwiastniki ideałów i lokalizacja}
\label{aa:ch:lokalizacja}
\setlength{\parskip}{2pt}

\section{Dlaczego zbiór rozwiązań nie wyznacza ideału?}
W Rozdziale~\ref{aa:ch:wielomiany} porównaliśmy ideały
$(x)$ oraz $(x^2)$ w $\mathbb R[x]$. Oba opisują ten sam
zbiór zer, ale tylko drugi iloraz ma niezerowy element
nilpotentny. Pierwiastnik ideału pozwala oddzielić część
równania widoczną przez zwykłe zerowanie od dodatkowych
informacji algebraicznych. Lokalizacja odpowiada na inne
pytanie: jakie ułamki wolno dołączyć do pierścienia, jeśli
chcemy badać zachowanie algebraiczne w wybranym miejscu?

W całym rozdziale, jeśli nie podano inaczej, $R$ jest
\emph{przemiennym pierścieniem z jedynką}, a ideały
są dwustronne. Dopuszczamy pierścień zerowy tylko tam,
gdzie jest to jawnie dozwolone.

\section{Pierwiastnik ideału i nilpotentność}
\begin{definicja}[Pierwiastnik ideału]
\label{aa:def:ideal-radical}
Niech $R$ będzie przemiennym pierścieniem z jedynką,
a $I\subseteq R$ jego ideałem. \emph{Pierwiastnikiem}
ideału $I$ nazywamy zbiór
\[
 \sqrt I=\{a\in R:\ \exists n\in\mathbb N,\ n\ge1,\ a^n\in I\}.
\]
Ideał $I$ nazywamy \emph{radykalnym}, jeśli $\sqrt I=I$.
Pierwiastnik ideału zerowego $\sqrt{(0)}$ nazywamy
\emph{nilradykałem} pierścienia $R$.
\end{definicja}

\begin{twierdzenie}[Pierwiastnik ideału jest ideałem]
\label{aa:tw:radical-ideal}
Niech $R$ będzie przemiennym pierścieniem z jedynką,
a $I\subseteq R$ jego ideałem. Wówczas zbiór $\sqrt I$
jest ideałem zawierającym $I$.
Ponadto $\sqrt{\sqrt I}=\sqrt I$, a dla każdego
ideału $J\subseteq R$ z $I\subseteq J$ zachodzi
$\sqrt I\subseteq\sqrt J$.
\end{twierdzenie}
\begin{proof}
Jeśli $a\in I$, to $a^1\in I$, więc $I\subseteq\sqrt I$.
Dla $a\in\sqrt I$ i $r\in R$ istnieje $m\ge1$ z
$a^m\in I$. Wtedy $(ra)^m=r^ma^m\in I$,
więc $ra\in\sqrt I$. Również $-a\in\sqrt I$, gdyż
$(-a)^m=(-1)^ma^m\in I$.

Pozostaje dowieść zamkniętości na dodawanie.
Niech $a^m\in I$ i $b^n\in I$ dla pewnych $m,n\ge1$.
Rozwińmy $(a+b)^{m+n-1}$ za pomocą dwumianu Newtona.
W każdym składniku $\binom{m+n-1}{k}a^k b^{m+n-1-k}$
albo $k\ge m$, albo $m+n-1-k\ge n$: w przeciwnym
razie mielibyśmy $k\le m-1$ i
$m+n-1-k\le n-1$, co po dodaniu daje sprzeczność.
Pierwszy rodzaj składników należy do $I$, bo jest
podzielny przez $a^m$, a drugi należy do $I$, bo jest
podzielny przez $b^n$. Cała suma należy do $I$.
Wobec tego $a+b\in\sqrt I$ i $\sqrt I$ jest ideałem.

Jeśli $a\in\sqrt{\sqrt I}$, to $a^m\in\sqrt I$
dla pewnego $m\ge1$, a następnie $(a^m)^n\in I$
dla pewnego $n\ge1$. Zatem $a^{mn}\in I$,
czyli $a\in\sqrt I$. Odwrotne zawieranie wynika
z $\sqrt I\subseteq\sqrt{\sqrt I}$. Monotoniczność
jest bezpośrednia: $a^n\in I\subseteq J$
pociąga $a\in\sqrt J$.
\end{proof}

\begin{twierdzenie}[Pierwiastnik iloczynu i przecięcia ideałów]
\label{aa:tw:radical-product}
Niech $R$ będzie przemiennym pierścieniem z jedynką,
a $I,J\subseteq R$ jego ideałami. Wtedy zachodzi
$\sqrt{IJ}=\sqrt{I\cap J}=\sqrt I\cap\sqrt J$.
\end{twierdzenie}
\begin{proof}
Ponieważ $IJ\subseteq I\cap J$, z monotoniczności
pierwiastnika otrzymujemy
$\sqrt{IJ}\subseteq\sqrt{I\cap J}\subseteq\sqrt I\cap\sqrt J$.
Jeśli $a\in\sqrt I\cap\sqrt J$, to istnieją
$m,n\ge1$ takie, że $a^m\in I$ oraz $a^n\in J$.
Wówczas $a^{m+n}=a^m a^n\in IJ$. Oznacza to
$a\in\sqrt{IJ}$, co daje brakujące zawieranie
i wszystkie równości.
\end{proof}

\begin{stwierdzenie}[Nilradykał i pierścienie zredukowane]
\label{aa:st:nilradical-reduced}
Niech $R$ będzie przemiennym pierścieniem z jedynką,
$I\subseteq R$ jego ideałem, a $a\in R$. Wówczas element
$a+I\in R/I$ jest nilpotentny wtedy i tylko wtedy,
gdy $a\in\sqrt I$. W szczególności
$R/I$ nie ma niezerowych elementów nilpotentnych wtedy
i tylko wtedy, gdy $I$ jest radykalny.
Pierścień bez niezerowych elementów nilpotentnych
nazywamy \emph{pierścieniem zredukowanym}.
\end{stwierdzenie}
\begin{proof}
Z definicji działań w ilorazie
$(a+I)^n=a^n+I$. Jest to klasa zerowa dokładnie wtedy,
gdy $a^n\in I$. Istnienie takiego $n\ge1$ jest
równoważne $a\in\sqrt I$. Jeśli wszystkie elementy
nilpotentne ilorazu są zerowe, to każdy $a\in\sqrt I$
należy do $I$, zatem $\sqrt I=I$. Odwrotnie,
gdy $I=\sqrt I$, dowolny element nilpotentny
$a+I$ ma reprezentanta $a\in I$, więc jest zerowy.
\end{proof}

\begin{przyklad}[Trzy ideały opisujące podobne równania]
\label{aa:prz:radical-examples}
Niech $K$ będzie ciałem. W $K[x]$ mamy
$\sqrt{(x^m)}=(x)$ dla każdego $m\ge1$.
Istotnie, $x^m\in(x^m)$, więc $x\in\sqrt{(x^m)}$
i $(x)\subseteq\sqrt{(x^m)}$.
Z drugiej strony $(x^m)\subseteq(x)$, a $(x)$ jest
ideałem pierwszym: $K[x]/(x)\cong K$ jest ciałem.
Ideał pierwszy jest radykalny, ponieważ
z $a^n\in(x)$ jego pierwszość daje kolejno $a\in(x)$.
Stąd $\sqrt{(x^m)}\subseteq(x)$.

W $K[x,y]$ rozważmy $I=(x^2,xy)$. Zachodzi
$I\subseteq(x)$ i $x^2\in I$. Prowadząc dokładnie ten
sam argument, otrzymujemy $\sqrt I=(x)$.
Jednak $I\ne(x)$: gdyby $x=x^2F+xyG$, to
po wyłączeniu $x$ i skróceniu w dziedzinie całkowitości
mielibyśmy $1=xF+yG$, co jest niemożliwe po
podstawieniu $x=y=0$. W szczególności ideał $I$
nie jest radykalny.
\end{przyklad}

\begin{przyklad}[Nilradykał pierścienia $\mathbb Z/12\mathbb Z$]
\label{aa:prz:nilradical-z12}
W pierścieniu $R=\mathbb Z/12\mathbb Z$ klasa $[6]$
jest niezerowym elementem nilpotentnym, gdyż
$[6]^2=[36]=[0]$. Dokładnie elementy ideału $([6])$
są nilpotentne. Istotnie, jeśli $[a]^n=[0]$, to
$12\mid a^n$, a więc $2\mid a$ i $3\mid a$.
Stąd $6\mid a$ i $[a]\in([6])$. Odwrotnie,
$([6])=\{[0],[6]\}$ i obie klasy są nilpotentne.
Ponadto $([2])$ i $([3])$ są ideałami pierwszymi,
bo ich ilorazy są odpowiednio izomorficzne z
$\mathbb F_2$ i $\mathbb F_3$; ich przecięcie wynosi
$([6])$. Są to wszystkie ideały pierwsze $R$:
jeśli $P$ jest ideałem pierwszym, to
$[2]^2[3]=[12]=[0]\in P$, więc $[2]\in P$
lub $[3]\in P$. W pierwszym przypadku
$([2])\subseteq P$, a w drugim $([3])\subseteq P$;
ponieważ oba te ideały są maksymalne i $P\ne R$,
mamy odpowiednio $P=([2])$ lub $P=([3])$.
Otrzymujemy zatem konkretną tożsamość
$\sqrt{(0)}=\bigcap_{P\ \mathrm{pierwszy}}P=([6])$.
\end{przyklad}

\section{Pierwiastnik jako przecięcie ideałów pierwszych}
W następnym twierdzeniu użyjemy lematu Kuratowskiego--Zorna:
jeżeli w niepustym zbiorze częściowo uporządkowanym każdy
łańcuch ma ograniczenie górne, to istnieje w nim element
maksymalny. Potrzebny nam argument jest konstrukcją
maksymalnego ideału unikającego wskazanych potęg elementu.

\begin{lemat}[Ideał pierwszy rozłączny ze zbiorem multiplikatywnym]
\label{aa:lem:prime-avoid-S}
Niech $R$ będzie przemiennym pierścieniem z jedynką,
$I\subseteq R$ ideałem, a $S\subseteq R$ niepustym zbiorem
zawierającym $1_R$ i zamkniętym na mnożenie.
Jeżeli $I\cap S=\varnothing$, to istnieje ideał
pierwszy $P$ taki, że $I\subseteq P$ oraz
$P\cap S=\varnothing$.
\end{lemat}
\begin{proof}
Rozważmy rodzinę $\mathcal F$ wszystkich ideałów
$J$ z $I\subseteq J$ i $J\cap S=\varnothing$,
uporządkowaną przez zawieranie. Rodzina jest niepusta,
bo $I\in\mathcal F$. Pusty łańcuch ma ograniczenie
górne $I$. Jeśli natomiast $\mathcal C\subseteq\mathcal F$
jest niepustym łańcuchem, to suma mnogościowa
$\bigcup_{J\in\mathcal C}J$ jest ideałem: dla dwóch jej elementów wybieramy ideał
z łańcucha zawierający oba, co jest możliwe dzięki
liniowemu uporządkowaniu łańcucha; w nim sprawdzamy
zamkniętość na odejmowanie i mnożenie przez pierścień.
Suma zawiera $I$ i nie przecina $S$, bo żaden z jej
składników nie przecina $S$. Jest zatem ograniczeniem
górnym dla łańcucha. Lemat Kuratowskiego--Zorna daje
maksymalny element $P$ rodziny $\mathcal F$.

Pokażemy, że $P$ jest pierwszy. Jest właściwy, ponieważ
$1_R\in S$, a $P\cap S=\varnothing$.
Załóżmy przeciwnie, że $ab\in P$, lecz $a\notin P$
i $b\notin P$. Ideały $P+(a)$ i $P+(b)$ ściśle
zawierają $P$. Na mocy maksymalności $P$ żaden
z nich nie może należeć do $\mathcal F$.
Istnieją więc $s\in(P+(a))\cap S$ oraz
$t\in(P+(b))\cap S$. Zapiszmy
$s=p+ra$, $t=q+ub$ dla $p,q\in P$, $r,u\in R$.
Wtedy $st=pq+pub+qra+ruab\in P$, ponieważ
pierwsze trzy składniki zawierają czynnik z $P$,
a ostatni należy do $P$ z założenia $ab\in P$.
Jednak $st\in S$ dzięki zamkniętości $S$ na mnożenie.
Otrzymaliśmy sprzeczność z $P\cap S=\varnothing$.
Zatem $P$ jest ideałem pierwszym.
\end{proof}

\begin{twierdzenie}[Pierwiastnik ideału jako przecięcie ideałów pierwszych]
\label{aa:tw:radical-intersection-primes}
Niech $R$ będzie przemiennym pierścieniem z jedynką
i niech $I$ będzie jego ideałem. Wtedy
\[
 \sqrt I=\bigcap_{\substack{P\text{ -- ideał pierwszy}\\ I\subseteq P}}P.
\]
Jeżeli nie ma żadnego ideału pierwszego zawierającego
$I$, przecięcie po pustej rodzinie oznacza cały $R$.
\end{twierdzenie}
\begin{proof}
Jeżeli $I=R$, obie strony są równe $R$ na mocy
przyjętej konwencji. Załóżmy więc $I\ne R$.
Niech $a\in\sqrt I$; dla pewnego $n\ge1$ mamy
$a^n\in I$. Jeśli $P$ jest ideałem pierwszym zawierającym
$I$, to $a^n\in P$. Wielokrotne zastosowanie własności
ideału pierwszego daje $a\in P$. Zatem $a$ należy
do wszystkich takich $P$.

Dla odwrotnego zawierania weźmy $a\notin\sqrt I$.
Wtedy żadna dodatnia potęga $a$ nie należy do $I$;
ponieważ $I$ jest właściwy, również $1_R\notin I$.
Zbiór $S=\{1_R,a,a^2,\ldots\}$ jest multiplikatywny
i nie przecina $I$. Poprzedni lemat daje ideał
pierwszy $P\supseteq I$ rozłączny z $S$.
W szczególności $a\notin P$. Zatem $a$ nie należy
do przecięcia wszystkich ideałów pierwszych zawierających
$I$. Otrzymujemy wymaganą równość.
\end{proof}

\begin{uwaga}[Interpretacja geometryczna]
\label{aa:uw:radical-geometry}
Jeżeli $f\in K[x_1,\ldots,x_n]$ i $f^m$ zanika
w punkcie $z\in K^n$, to także $f$ zanika w tym punkcie,
gdyż w ciele nie ma niezerowych elementów nilpotentnych.
Dlatego ideały $I$ i $\sqrt I$ mają ten sam zbiór
wspólnych zer. Nie twierdzimy tu jednak, że każdy
wielomian znikający na takim zbiorze należy do $\sqrt I$:
taka równość wymaga dodatkowych założeń i prowadzi
do twierdzenia Hilberta o zerach.
\end{uwaga}

\begin{przyklad}[Dlaczego ciało algebraicznie domknięte ma znaczenie]
\label{aa:prz:real-zero-set-not-radical}
Dla $K=\mathbb R$ rozważmy $I=(x^2+1)\subset\mathbb R[x]$.
Zbiór rzeczywistych zer jest pusty: dla każdej
liczby rzeczywistej $a$ zachodzi $a^2+1>0$.
Z tego powodu każdy wielomian znika na wszystkich
punktach tego zbioru (warunek jest spełniony
pusto), więc ideał wielomianów znikających na zbiorze
zer $I$ jest całym $\mathbb R[x]$.
Tymczasem $I$ jest ideałem maksymalnym,
ponieważ $\mathbb R[x]/(x^2+1)$ jest ciałem,
izomorficznym z $\mathbb C$. Ideał maksymalny
jest pierwszy, a każdy ideał pierwszy jest radykalny;
zatem $\sqrt I=I\ne\mathbb R[x]$.
Sama znajomość rzeczywistych zer nie pozwala więc
odzyskać pierwiastnika dowolnego ideału.
\end{przyklad}

\section{Lokalizacja: konstrukcja przez pary i usuwanie mianowników}
\begin{definicja}[Zbiór multiplikatywny]
\label{aa:def:multiplicative-set}
Niech $R$ będzie przemiennym pierścieniem z jedynką.
Podzbiór $S\subseteq R$ jest \emph{multiplikatywny},
jeśli $1_R\in S$ oraz $st\in S$ dla wszystkich $s,t\in S$.
W konstrukcji niezerowej lokalizacji zakładamy dodatkowo
$0_R\notin S$. Gdyby $0_R\in S$, lokalizacja byłaby
pierścieniem zerowym.
\end{definicja}

Nie wystarczy przyjąć zwykłej relacji $as'=a's$
na parach $(a,s)$. W pierścieniach z dzielnikami zera
tak zdefiniowana relacja nie musi być przechodnia.
Należy dopuścić dodatkowy czynnik należący do $S$.

\begin{przyklad}[Zawodność naiwnej relacji na ułamkach]
\label{aa:prz:naive-local-relation}
W $R=\mathbb Z/6\mathbb Z$ weźmy
$S=\{[1],[2],[4]\}$; jest to zbiór multiplikatywny
niezawierający zera. W relacji zadanej wyłącznie
równością iloczynów krzyżowych mamy
$( [0],[1])\sim([0],[2])$, ponieważ $0=0$,
oraz $([0],[2])\sim([3],[1])$, ponieważ
$[0]\cdot[1]=[3]\cdot[2]=[0]$.
Jednak $([0],[1])\not\sim([3],[1])$, gdyż
$[0]\cdot[1]\ne[3]\cdot[1]$.
Relacja nie jest więc przechodnia. W poprawnej
konstrukcji mnożnik $[2]\in S$ zeruje różnicę
$[0]\cdot[1]-[3]\cdot[1]=[3]$.
\end{przyklad}

\begin{definicja}[Lokalizacja pierścienia]
\label{aa:def:localization}
Niech $R$ będzie przemiennym pierścieniem z jedynką,
a $S\subseteq R$ zbiorem multiplikatywnym niezawierającym zera.
Na $R\times S$ definiujemy relację
\[
 (a,s)\sim(b,t)
 \quad\Longleftrightarrow\quad
 \exists u\in S:\ u(at-bs)=0_R.
\]
Klasę pary $(a,s)$ oznaczamy $a/s$. Zbiór klas
nazywamy \emph{lokalizacją} $R$ względem $S$ i oznaczamy
$S^{-1}R$. Na klasach określamy działania
\[
 \frac as+\frac bt=\frac{at+bs}{st},\qquad
 \frac as\cdot\frac bt=\frac{ab}{st}.
\]
\end{definicja}

\begin{twierdzenie}[Poprawność konstrukcji lokalizacji]
\label{aa:tw:localization-well-defined}
Niech $R$ będzie przemiennym pierścieniem z jedynką,
a $S\subseteq R$ zbiorem multiplikatywnym niezawierającym zera.
Dla $(a,s),(b,t)\in R\times S$ określmy
$(a,s)\sim(b,t)$ wtedy i tylko wtedy, gdy istnieje $u\in S$
takie, że $u(at-bs)=0_R$. Wówczas $\sim$ jest relacją
równoważności na $R\times S$, a wzory
\[
\frac as+\frac bt=\frac{at+bs}{st},\qquad
\frac as\cdot\frac bt=\frac{ab}{st}
\]
są dobrze określone na klasach i czynią $S^{-1}R$
przemiennym pierścieniem z jedynką. Odwzorowanie
$\iota:R\to S^{-1}R$, $a\mapsto a/1_R$, jest
homomorfizmem pierścieni; każdy $s/1_R$ dla $s\in S$
jest jednostką lokalizacji o odwrotności $1_R/s$.
\end{twierdzenie}
\begin{proof}
\emph{Relacja równoważności.} Zwrotność wynika z
$1_R(as-sa)=0_R$, a symetria z zamiany znaku
różnicy. Dla przechodniości załóżmy
$(a,s)\sim(b,t)$ oraz $(b,t)\sim(c,v)$.
Istnieją $u,w\in S$ takie, że
$u(at-bs)=0_R$ i $w(bv-ct)=0_R$.
Pomnóżmy pierwszą równość przez $wv$,
a drugą przez $us$. Po dodaniu otrzymujemy
\[
 uwt(av-cs)=0_R.
\]
Ponieważ $uwt\in S$, zachodzi $(a,s)\sim(c,v)$.

\emph{Niezależność działań od reprezentantów.}
Załóżmy $(a,s)\sim(a',s')$ i $(b,t)\sim(b',t')$.
Wtedy pewne $u,w\in S$ zerują odpowiednio
$as'-a's$ oraz $bt'-b't$. Dla dodawania różnica
liczników po sprowadzeniu wyników do wspólnego mianownika
jest równa
\[
 (at+bs)s't'-(a't'+b's')st
 =tt'(as'-a's)+ss'(bt'-b't).
\]
Po pomnożeniu przez $uw\in S$ oba składniki
po prawej znikają, więc sumy wyznaczają tę samą klasę.
Podobnie dla mnożenia różnica wynosi
\[
 ab s't'-a'b' st
 =bt'(as'-a's)+a's(bt'-b't),
\]
która również znika po pomnożeniu przez $uw$.
Mianowniki wynikowych ułamków należą do $S$.

\emph{Aksjomaty pierścienia.} Łączność i przemienność
obu działań oraz rozdzielność wynikają z odpowiednich
tożsamości w $R$ po sprowadzeniu ułamków do wspólnych
mianowników. Dla przykładu
$((a/s)+(b/t))+(c/v)=(atv+bsv+cst)/(stv)$,
co nie zależy od nawiasowania. Elementem zerowym
jest $0_R/1_R$, elementem przeciwnym do $a/s$
jest $(-a)/s$, a jedynką $1_R/1_R$.

Wzory na działania pokazują, że $\iota$ zachowuje
dodawanie, mnożenie i jedynkę. Dla $s\in S$ mamy
$(s/1_R)(1_R/s)=s/s=1_R/1_R$,
ponieważ $s\cdot1_R-1_R\cdot s=0_R$.
Stąd $s/1_R$ jest odwracalny.
\end{proof}

\begin{stwierdzenie}[Kiedy element staje się zerem?]
\label{aa:st:local-zero-kernel}
Niech $R$ będzie przemiennym pierścieniem z jedynką,
a $S\subseteq R$ zbiorem multiplikatywnym niezawierającym
zera. Niech $\iota:R\to S^{-1}R$ będzie homomorfizmem
kanonicznym, $\iota(a)=a/1_R$. Dla $a\in R$ zachodzi
$\iota(a)=0$ w $S^{-1}R$ wtedy i tylko wtedy,
gdy istnieje $s\in S$ z $sa=0_R$.
W szczególności $\iota$ jest różnowartościowe, jeśli
żaden element $S$ nie jest dzielnikiem zera.
Gdy $R$ jest dziedziną całkowitości, warunek ten zachodzi
automatycznie dla $0_R\notin S$.
\end{stwierdzenie}
\begin{proof}
Równość $a/1_R=0_R/1_R$ oznacza z definicji relacji
istnienie $s\in S$ z
$s(a\cdot1_R-0_R\cdot1_R)=sa=0_R$.
Jeśli żaden element $s\in S$ nie jest dzielnikiem zera,
równość ta wymusza $a=0_R$,
więc jądro $\iota$ jest zerowe.
\end{proof}

\begin{przyklad}[Lokalizacja liczb całkowitych]
\label{aa:prz:local-z}
Dla $R=\mathbb Z$ i $S=\{2^n:n\ge0\}$
lokalizacja jest pierścieniem
$\mathbb Z[1/2]=\{a/2^n:a\in\mathbb Z,\ n\ge0\}\subset\mathbb Q$.
Dołączyliśmy odwrotność liczby $2$, ale nie
odwrotność liczby $3$. Gdyby $1/3=a/2^n$,
to $3a=2^n$, co wymagałoby $3\mid 2^n$,
a to jest niemożliwe.

Przy $S=\{2^a3^b:a,b\in\mathbb N_0\}$ otrzymujemy
$\mathbb Z[1/6]$, a jego jednostkami są dokładnie
liczby $\pm2^m3^n$, gdzie $m,n\in\mathbb Z$.
Rzeczywiście, te liczby mają odwrotności tej samej postaci.
Jeżeli ułamek $u/(2^a3^b)$ jest jednostką, to po
pomnożeniu przez odwrotność również należącą do
$\mathbb Z[1/6]$ jego licznik $u$ musi dzielić pewną
liczbę postaci $2^c3^d$ w $\mathbb Z$. Jednoznaczny
rozkład na czynniki w $\mathbb Z$ wymusza
$u=\pm2^r3^s$ dla pewnych $r,s\ge0$.
\end{przyklad}

\section{Własność uniwersalna lokalizacji i zachowanie ideałów}
\begin{twierdzenie}[Własność uniwersalna lokalizacji]
\label{aa:tw:localization-universal}
Niech $R,T$ będą przemiennymi pierścieniami z jedynką,
$S\subseteq R$ zbiorem multiplikatywnym niezawierającym
zera i $f:R\to T$ homomorfizmem takim, że
$f(s)\in T^\times$ dla każdego $s\in S$.
Niech $\iota:R\to S^{-1}R$ oznacza homomorfizm
kanoniczny $\iota(a)=a/1_R$. Istnieje dokładnie jeden
homomorfizm pierścieni $\bar f:S^{-1}R\to T$
spełniający $\bar f\circ\iota=f$.
Jest zadany wzorem $\bar f(a/s)=f(a)f(s)^{-1}$.
\end{twierdzenie}
\begin{proof}
Najpierw sprawdzamy niezależność od reprezentantów.
Jeśli $a/s=b/t$, to pewne $u\in S$ spełnia
$u(at-bs)=0_R$. Po zastosowaniu $f$ dostajemy
$f(u)(f(a)f(t)-f(b)f(s))=0_T$.
Ponieważ $f(u)$ jest jednostką, możemy pomnożyć
przez $f(u)^{-1}$, a następnie przez
$f(s)^{-1}f(t)^{-1}$, otrzymując
$f(a)f(s)^{-1}=f(b)f(t)^{-1}$.
Zatem wzór ma jednoznaczny sens na klasach.
Zachowanie dodawania i mnożenia wynika ze zwykłych
wzorów na działania na ułamkach; jedynka przechodzi
na $f(1_R)f(1_R)^{-1}=1_T$.

Dla jednoznaczności zauważmy, że w lokalizacji
$a/s=(a/1_R)(s/1_R)^{-1}$.
Każdy homomorfizm przedłużający $f$ musi więc
wysyłać $a/s$ na $f(a)f(s)^{-1}$. Otrzymana mapa
jest jedynym możliwym przedłużeniem.
\end{proof}

\begin{definicja}[Rozszerzenie ideału do lokalizacji]
\label{aa:def:ideal-extension-local}
Niech $R$ będzie przemiennym pierścieniem z jedynką,
$S\subseteq R$ zbiorem multiplikatywnym niezawierającym
zera, a $I\subseteq R$ ideałem. Przez $S^{-1}I$ oznaczamy
zbiór klas $a/s\in S^{-1}R$,
gdzie $a\in I$, $s\in S$. Jest to ideał pierścienia
$S^{-1}R$ i nazywamy go \emph{rozszerzeniem ideału} $I$.
\end{definicja}
\begin{lemat}[Rozszerzenie ideału jest ideałem]
\label{aa:lem:local-ideal-is-ideal}
Niech $R$ będzie przemiennym pierścieniem z jedynką,
$S\subseteq R$ zbiorem multiplikatywnym niezawierającym
zera, a $I\subseteq R$ ideałem. Wtedy zbiór $S^{-1}I$
jest ideałem pierścienia $S^{-1}R$.
\end{lemat}
\begin{proof}
Element $0/1$ należy do $S^{-1}I$, bo $0\in I$.
Jeśli $a/s,b/t\in S^{-1}I$, gdzie $a,b\in I$, to
$(a/s)-(b/t)=(at-bs)/(st)$, a $at-bs\in I$,
ponieważ $I$ jest ideałem. Zatem zbiór jest podgrupą
addytywną. Dla dowolnego $c/u\in S^{-1}R$ mamy
$(c/u)(a/s)=ca/(us)$, a $ca\in I$.
To dowodzi zamkniętości na mnożenie przez elementy lokalizacji.
\end{proof}

\begin{lemat}[Przynależność do rozszerzenia ideału]
\label{aa:lem:local-ideal-membership}
Niech $R$ będzie przemiennym pierścieniem z jedynką,
$S\subseteq R$ zbiorem multiplikatywnym niezawierającym
zera, a $I\subseteq R$ ideałem. Dla $a\in R$ i $s\in S$ zachodzi
$a/s\in S^{-1}I$ wtedy i tylko wtedy, gdy
istnieje $u\in S$ takie, że $ua\in I$.
Ponadto $S^{-1}I$ jest ideałem właściwym wtedy
i tylko wtedy, gdy $I\cap S=\varnothing$.
\end{lemat}
\begin{proof}
Jeżeli $a/s=b/t$ dla $b\in I$, $t\in S$,
to istnieje $v\in S$ z $v(at-bs)=0$.
Wobec tego $(vt)a=(vs)b\in I$, a $vt\in S$.
Odwrotnie, jeśli $ua\in I$, to
$a/s=(ua)/(us)$, ponieważ iloczyny krzyżowe są
równe, zatem ułamek należy do $S^{-1}I$.

Ideał $S^{-1}I$ jest całym pierścieniem dokładnie wtedy,
gdy zawiera jedynkę $1_R/1_R$. Z pierwszej części
jest to równoważne istnieniu $u\in S$ z
$u\cdot1_R=u\in I$, czyli $I\cap S\ne\varnothing$.
\end{proof}

\begin{twierdzenie}[Kontrakcja rozszerzonego ideału i nasycenie]
\label{aa:tw:localization-saturation}
Niech $R$ będzie przemiennym pierścieniem z jedynką,
$S\subseteq R$ zbiorem multiplikatywnym niezawierającym
zera, a $I\subseteq R$ ideałem. Niech
$\iota:R\to S^{-1}R$, $\iota(a)=a/1_R$, będzie
homomorfizmem kanonicznym. Przeciwobraz ideału
$S^{-1}I$ ma postać
\[
 \iota^{-1}(S^{-1}I)=I^S
 :=\{a\in R:\ \exists s\in S\ \ sa\in I\}.
\]
W szczególności $I^S$ jest ideałem zawierającym $I$.
Równość $I^S=I$ nie musi zachodzić.
\end{twierdzenie}
\begin{proof}
Z Lematu~\ref{aa:lem:local-ideal-membership} wynika,
że $a/1\in S^{-1}I$ dokładnie wtedy, gdy $sa\in I$
dla pewnego $s\in S$. To jest dokładnie definicja
przeciwobrazu oraz zbioru $I^S$.
Przeciwobraz ideału przez homomorfizm jest ideałem:
jeśli $a,b\in I^S$, to $sa,tb\in I$ dla pewnych
$s,t\in S$, zatem $st(a-b)=t(sa)-s(tb)\in I$;
ponadto $s(ra)=r(sa)\in I$ dla każdego $r\in R$.
Wreszcie $I\subseteq I^S$, gdyż $1\in S$.
\end{proof}

\begin{przyklad}[Nasycenie może powiększyć ideał]
\label{aa:prz:localization-saturation-z}
Niech $R=\mathbb Z$, $S=\{2^n:n\ge0\}$
oraz $I=(6)$. Mamy $I^S=(3)$.
Jeśli $a\in I^S$, to dla pewnego $n$ liczba $6$
dzieli $2^na$, co wymusza $3\mid a$.
Odwrotnie, gdy $3\mid a$, liczba $2a$ jest
podzielna przez $6$, więc $a\in I^S$.
W lokalizacji $\mathbb Z[1/2]$ ideały generowane
przez $6$ i przez $3$ są jednak identyczne,
ponieważ $6=2\cdot3$ i $2$ stało się jednostką.
\end{przyklad}

\begin{twierdzenie}[Pierwiastnik ideału a lokalizacja]
\label{aa:tw:localization-radical}
Niech $R$ będzie przemiennym pierścieniem z jedynką,
$S\subseteq R$ zbiorem multiplikatywnym
niezawierającym zera, a $I\subseteq R$ ideałem.
W pierścieniu $S^{-1}R$ zachodzi równość ideałów
\[
 \sqrt{S^{-1}I}=S^{-1}(\sqrt I).
\]
\end{twierdzenie}
\begin{proof}
Weźmy $a/s\in S^{-1}(\sqrt I)$.
Z definicji rozszerzenia ideału istnieją
$b\in\sqrt I$ i $t\in S$ takie, że $a/s=b/t$.
Ponieważ $b\in\sqrt I$, istnieje liczba naturalna
$n\ge1$ z $b^n\in I$. Stąd
$(a/s)^n=b^n/t^n\in S^{-1}I$, czyli
$a/s\in\sqrt{S^{-1}I}$. To daje pierwsze zawieranie.

Dla zawierania odwrotnego niech
$a/s\in\sqrt{S^{-1}I}$. Dla pewnego $n\ge1$
mamy $a^n/s^n\in S^{-1}I$.
Na mocy Lematu~\ref{aa:lem:local-ideal-membership}
istnieje $u\in S$ takie, że $ua^n\in I$.
Ponieważ $I$ jest ideałem, także
$(ua)^n=u^{n-1}(ua^n)\in I$.
Z definicji pierwiastnika wynika $ua\in\sqrt I$.
Równość ułamków $a/s=(ua)/(us)$ jest poprawna,
gdyż oba iloczyny krzyżowe są równe $aus$.
Ponieważ $ua\in\sqrt I$ i $us\in S$,
otrzymujemy $a/s\in S^{-1}(\sqrt I)$.
Równość ideałów została udowodniona.
\end{proof}

\begin{przyklad}[Pierwiastniki po odwróceniu dwójki]
\label{aa:prz:localization-radical-example}
Niech $R=\mathbb Z$, $I=(12)$ oraz
$S=\{2^n:n\in\mathbb N_0\}$.
W pierścieniu $\mathbb Z$ zachodzi
$\sqrt{(12)}=(6)$: liczba całkowita ma pewną
potęgę podzielną przez $12=2^2\cdot3$ wtedy
i tylko wtedy, gdy sama jest podzielna przez
$2$ i $3$. Po odwróceniu liczby $2$ mamy
$S^{-1}(12)=(3)$ i $S^{-1}(6)=(3)$,
gdzie ideały zapisano już w $\mathbb Z[1/2]$.
Istotnie, $12=3\cdot 2^2$ i $6=3\cdot2$,
a potęgi $2$ są jednostkami lokalizacji.
Twierdzenie daje więc
$\sqrt{(3)}=(3)$ w $\mathbb Z[1/2]$,
co zgadza się także z faktem, że
$\mathbb Z[1/2]/(3)\cong\mathbb F_3$ jest ciałem.
\end{przyklad}

\section{Ideały pierwsze po lokalizacji i pierścienie lokalne}
\begin{twierdzenie}[Odpowiedniość ideałów pierwszych przy lokalizacji]
\label{aa:tw:prime-local-correspondence}
Niech $R$ będzie przemiennym pierścieniem z jedynką,
a $S\subseteq R$ zbiorem multiplikatywnym niezawierającym zera.
Niech $\iota:R\to S^{-1}R$ oznacza homomorfizm
kanoniczny $\iota(a)=a/1_R$. Przyporządkowania
$P\mapsto S^{-1}P$ oraz $Q\mapsto\iota^{-1}(Q)$
są wzajemnie odwrotnymi
bijekcjami między ideałami pierwszymi $P$ pierścienia $R$
spełniającymi $P\cap S=\varnothing$ a ideałami
pierwszymi $Q$ lokalizacji $S^{-1}R$.
\end{twierdzenie}
\begin{proof}
Niech $P$ będzie ideałem pierwszym pierścienia $R$
rozłącznym z $S$. Rozszerzenie $S^{-1}P$ jest właściwe
na mocy Lematu~\ref{aa:lem:local-ideal-membership}.
Jeśli iloczyn $(a/s)(b/t)=ab/(st)$ należy do tego
ideału, istnieje $u\in S$ z $uab\in P$.
Ponieważ $u\notin P$ i $P$ jest pierwszy, otrzymujemy
$ab\in P$, a następnie $a\in P$ lub $b\in P$.
Zatem przynajmniej jeden z ułamków $a/s$, $b/t$
należy do $S^{-1}P$, co dowodzi jego pierwszości.
Ponadto $a/1_R\in S^{-1}P$ dokładnie wtedy,
gdy $ua\in P$ dla pewnego $u\in S$.
Pierwszość $P$ i $u\notin P$ dają wtedy $a\in P$.
W ten sposób przeciwobraz $S^{-1}P$ jest równy $P$.

Odwrotnie, niech $Q$ będzie ideałem pierwszym $S^{-1}R$
i połóżmy $P=\iota^{-1}(Q)$. Przeciwobraz ideału
pierwszego przez homomorfizm jest pierwszy:
$ab\in P$ oznacza $\iota(a)\iota(b)\in Q$, więc
$\iota(a)\in Q$ lub $\iota(b)\in Q$.
Ponieważ elementy $\iota(s)$ dla $s\in S$ są jednostkami,
nie mogą należeć do właściwego ideału $Q$, zatem
$P\cap S=\varnothing$.
Wreszcie dla każdego $a/s\in S^{-1}R$ zachodzi
\[
 \frac as\in Q
 \quad\Longleftrightarrow\quad
 \frac a1=\frac s1\frac as\in Q
 \quad\Longleftrightarrow\quad a\in P.
\]
W pierwszej równoważności wykorzystaliśmy odwracalność
$s/1$: z przynależności $a/1$ do ideału wynika
przynależność $(s/1)^{-1}(a/1)=a/s$.
Stąd $Q=S^{-1}P$, a obie konstrukcje są odwrotne.
\end{proof}

\begin{definicja}[Pierścień lokalny]
\label{aa:def:local-ring}
Niech $R$ będzie przemiennym, niezerowym pierścieniem
z jedynką. Pierścień $R$ nazywamy \emph{lokalnym},
jeśli ma dokładnie jeden ideał maksymalny.
Dla ideału pierwszego $P\subsetneq R$ definiujemy
\emph{lokalizację w ideale pierwszym}
$R_P=(R\setminus P)^{-1}R$.
Zbiór $R\setminus P$ jest multiplikatywny:
$1_R\notin P$, a jeśli $a,b\notin P$, to
$ab\notin P$ na mocy pierwszości $P$.
\end{definicja}

\begin{twierdzenie}[Lokalizacja w ideale pierwszym jest pierścieniem lokalnym]
\label{aa:tw:localization-at-prime}
Niech $R$ będzie przemiennym pierścieniem z jedynką,
a $P\subsetneq R$ ideałem pierwszym. Wtedy
$R_P=(R\setminus P)^{-1}R$ jest pierścieniem lokalnym,
a jego jedynym ideałem maksymalnym jest
$PR_P=(R\setminus P)^{-1}P$.
Element $a/s\in R_P$ jest jednostką wtedy i tylko wtedy,
gdy $a\notin P$.
\end{twierdzenie}
\begin{proof}
Dla $s\notin P$ element $s/1$ jest jednostką.
Jeśli $a\notin P$, także $a/1$ jest jednostką,
a więc $a/s=(a/1)(s/1)^{-1}$ jest jednostką.
Jeżeli $a\in P$, to $a/s\in PR_P$, które jest
ideałem właściwym na mocy rozłączności
$P\cap(R\setminus P)=\varnothing$.
Żadna jednostka nie leży w ideale właściwym,
więc taki ułamek nie jest jednostką.

Udowodniliśmy, że zbiór elementów nieodwracalnych
pierścienia $R_P$ jest dokładnie $PR_P$.
Jeśli $J\subsetneq R_P$ jest dowolnym ideałem właściwym,
nie może zawierać jednostki, więc wszystkie jego
elementy należą do $PR_P$. Stąd $J\subseteq PR_P$.
Wynika z tego, że $PR_P$ jest maksymalny i że każdy
inny ideał maksymalny musi być równy $PR_P$.
\end{proof}

\begin{twierdzenie}[Ciało reszt lokalizacji]
\label{aa:tw:localization-residue-field}
Niech $R$ będzie przemiennym pierścieniem z jedynką,
a $P\subsetneq R$ jego ideałem pierwszym. Niech
$R_P=(R\setminus P)^{-1}R$ i $PR_P=(R\setminus P)^{-1}P$.
Wówczas
\[
 R_P/PR_P\ \cong\ \operatorname{Frac}(R/P).
\]
Po prawej stronie występuje ciało ułamków dziedziny
całkowitości $R/P$.
\end{twierdzenie}
\begin{proof}
Ponieważ $P$ jest pierwszy, $R/P$ jest dziedziną,
a więc ma ciało ułamków opisane w
Lemacie~\ref{aa:lem:fracfield-well-defined}.
Dla $a\in R$ i $s\in R\setminus P$ określmy
\[
 \Phi(a/s)=(a+P)/(s+P)\in\operatorname{Frac}(R/P).
\]
Mianownik jest niezerowy. Jeśli $a/s=b/t$,
to dla pewnego $u\notin P$ mamy $u(at-bs)=0$.
Po przejściu do $R/P$ i skróceniu niezerowej klasy
$u+P$ otrzymujemy $(a+P)(t+P)=(b+P)(s+P)$;
zatem $\Phi$ jest dobrze określone. Ze wzorów na
ułamki wynika, że jest homomorfizmem pierścieni.
Każdy element $\operatorname{Frac}(R/P)$ ma postać
$(a+P)/(s+P)$ z $s\notin P$, więc $\Phi$ jest
surjektywny. Ułamek $a/s$ należy do jądra dokładnie
wtedy, gdy $a+P=0$, czyli $a\in P$; to z kolei
jest równoważne $a/s\in PR_P$ na mocy
Twierdzenia~\ref{aa:tw:localization-at-prime}.
Pierwsze twierdzenie o izomorfizmie kończy dowód.
\end{proof}

Poniższy diagram komutatywny przedstawia jednocześnie
lokalizację, przejście do ilorazu i zanurzenie w ciało
ułamków. Strzałka $\Phi$ ma jądro $PR_P$, a dolna
strzałka jest zanurzeniem dziedziny $R/P$ w jej ciało
ułamków.
\begin{center}
\begin{tikzpicture}[>=Stealth]
\node (R) at (0,1.8) {$R$};
\node (Loc) at (5.0,1.8) {$R_P$};
\node (Quot) at (0,0) {$R/P$};
\node (Frac) at (5.0,0) {$\operatorname{Frac}(R/P)$};
\draw[->] (R) -- node[above] {$\iota$} (Loc);
\draw[->] (R) -- node[left] {$\pi$} (Quot);
\draw[->] (Loc) -- node[right] {$\Phi$} (Frac);
\draw[->] (Quot) -- node[below] {$\jmath$} (Frac);
\end{tikzpicture}
\end{center}
W obu drogach z $R$ do $\operatorname{Frac}(R/P)$
element $a$ przechodzi w $(a+P)/(1+P)$.

\begin{przyklad}[Pierścień liczb wymiernych o mianownikach niepodzielnych przez $5$]
\label{aa:prz:z-local-five}
Dla $R=\mathbb Z$ i $P=(5)$ mamy
$\mathbb Z_{(5)}=\{a/b:\ a,b\in\mathbb Z,\ 5\nmid b\}\subset\mathbb Q$.
Na przykład $7/12$ i $1/3$ należą do tego pierścienia,
a $1/5$ nie należy: gdyby $1/5=a/b$ z $5\nmid b$,
z równości $b=5a$ wynikałaby sprzeczność.
Element $a/b$ jest jednostką dokładnie wtedy,
gdy $5\nmid a$, bo wówczas jego odwrotność $b/a$
również ma mianownik niepodzielny przez $5$.
Jedynym ideałem maksymalnym jest
$5\mathbb Z_{(5)}$, a ciało reszt jest izomorficzne
z $\mathbb F_5$ poprzez $a/b\mapsto[a][b]^{-1}$.
\end{przyklad}

\begin{przyklad}[Lokalizacja wielomianów w ideale $(x)$]
\label{aa:prz:local-kxy}
Weźmy $R=K[x,y]$ nad ciałem $K$ oraz
$P=(x)$. Jest to ideał pierwszy, bo
$K[x,y]/(x)\cong K[y]$ jest dziedziną całkowitości.
Pierścień $R_P$ zawiera odwrotności wszystkich
wielomianów spoza $(x)$, czyli tych, których
wartość po podstawieniu $x=0$ jest niezerowym
wielomianem zmiennej $y$.
W szczególności $y$ staje się jednostką w $R_P$,
ponieważ $y\notin(x)$, ale $x$ pozostaje nieodwracalny.
Jedynym ideałem maksymalnym $R_P$ jest $(x)R_P$.
Nie należy mylić tej konstrukcji z $K[x,x^{-1}]$,
w którym odwracamy potęgi $x$. W tym pierścieniu
$x$ jest jednostką, ale pierścień nie jest lokalny.
Jeśli ciało $K$ ma co najmniej trzy elementy,
wybierzmy $c\in K\setminus\{0,1\}$. Ideały $(x-1)$
i $(x-c)$ są różne i maksymalne, ponieważ podstawienie
odpowiednio $x=1$ i $x=c$ daje surjekcje
$K[x,x^{-1}]\to K$ o tych właśnie jądrach.
Jeśli $K=\mathbb F_2$, można zamiast drugiego użyć
ideału $(x^2+x+1)$. Wielomian $p=x^2+x+1$ nie ma
pierwiastka w $\mathbb F_2$, bo dla $x=0$ i $x=1$
przyjmuje wartość $1$; jako wielomian stopnia dwa jest
zatem nierozkładalny. W ciele $\mathbb F_2[x]/(p)$
zachodzi $[x]([x]+[1])=[1]$, gdyż
$x(x+1)=p-1=p+1$ w charakterystyce $2$.
Stąd klasa $x$ jest odwracalna, a więc
$\mathbb F_2[x,x^{-1}]/(p)\cong\mathbb F_2[x]/(p)$
jest ciałem. Ideał $(p)$ w pierścieniu Laurenta
jest więc maksymalny i różny od $(x-1)$, ponieważ
$p(1)=1\ne0$.
Istotna różnica jest następująca:
\emph{odwrócenie jednego elementu nie oznacza lokalizacji
w jednym ideale pierwszym}.
\end{przyklad}

\begin{uwaga}[Sprostowanie dotyczące pierścieni z dzielnikami zera]
\label{aa:uw:localization-zerodivisors}
W pierścieniu z dzielnikami zera lokalizacja może
zidentyfikować dwa różne elementy. Na przykład
w $R=\mathbb Z/6\mathbb Z$ dla
$S=\{[1],[2],[4]\}$ mamy $[3]\cdot[2]=[0]$,
więc klasa $[3]/[1]$ w $S^{-1}R$ jest zerowa.
Lokalizacja jest tu izomorficzna z $\mathbb F_3$:
redukcja modulo $3$ wysyła każdy element $S$ na
jednostkę $\mathbb F_3$, więc z własności uniwersalnej
otrzymujemy surjektywny homomorfizm
$S^{-1}R\to\mathbb F_3$. Jeżeli $[a]/[2^n]$ ma obraz
zerowy, to $3\mid a$; wtedy $[a]$ w $R$ jest klasą
$[0]$ lub $[3]$, a obie stają się zerem w lokalizacji.
Homomorfizm jest więc także różnowartościowy.
Nie wolno zatem w ogólnym przypadku zastępować relacji
z Definicji~\ref{aa:def:localization} przez zwykłą
równość iloczynów krzyżowych bez dodatkowego czynnika.
\end{uwaga}

\section{Skończone generowanie ideałów i przygotowanie do teorii Noethera}
\begin{definicja}[Pierścień noetherowski]
\label{aa:def:noetherian-ring-intro}
Przemienny pierścień z jedynką $R$ jest
\emph{noetherowski}, jeśli każdy jego ideał jest
skończenie generowany, to znaczy ma postać
$(a_1,\ldots,a_r)$ dla pewnego $r\in\mathbb N_0$
i pewnych $a_1,\ldots,a_r\in R$; dla $r=0$
przyjmujemy ideał zerowy $(0)$.
\end{definicja}

\begin{twierdzenie}[Warunek łańcucha i skończone generowanie ideałów]
\label{aa:tw:noetherian-acc}
Dla przemiennego pierścienia $R$ z jedynką
następujące warunki są równoważne:
\begin{enumerate}[label=\textup{(\roman*)}]
\item każdy ideał $R$ jest skończenie generowany;
\item każdy wstępujący ciąg ideałów
$I_1\subseteq I_2\subseteq\cdots$ od pewnego miejsca jest stały.
\end{enumerate}
\end{twierdzenie}
\begin{proof}
Załóżmy (i) i rozważmy wstępujący ciąg ideałów.
Jego suma mnogościowa $I=\bigcup_{n\ge1}I_n$ jest
ideałem: jeśli $a\in I_j$ i $b\in I_k$, to
$a-b\in I_{\max(j,k)}$, a dla dowolnego $r\in R$
mamy $ra\in I_j$. Istnieją generatory
$a_1,\ldots,a_m$ ideału $I$, gdzie $m\ge0$.
Jeśli $m=0$, to $I=(0)$, a ponieważ każdy
$I_n\subseteq I$, cały ciąg składa się z ideałów
zerowych i jest stały. Załóżmy więc $m\ge1$.
Każdy generator $a_i$ należy do pewnego $I_{n_i}$;
dla $N=\max\{n_1,\ldots,n_m\}$ wszystkie generatory
leżą w $I_N$. Stąd $I=(a_1,\ldots,a_m)\subseteq I_N$.
Odwrotne zawieranie jest oczywiste, zatem $I_N=I$
i dla wszystkich $n\ge N$ mamy $I_n=I_N$.

Załóżmy (ii). Przypuśćmy, że istnieje ideał $I$,
który nie jest skończenie generowany. Wybierzmy
$a_1\in I$ i ideał $I_1=(a_1)$ różny od $I$.
Jest to możliwe, bo $I$ nie jest generowany przez
jeden element. Następnie wybierzmy
$a_2\in I\setminus I_1$ i połóżmy
$I_2=(a_1,a_2)$. Ponieważ $I$ nie jest skończenie
generowany, $I_2\ne I$ i można wybrać
$a_3\in I\setminus I_2$. Powtarzając konstrukcję,
otrzymamy ściśle rosnący ciąg
$I_1\subsetneq I_2\subsetneq I_3\subsetneq\cdots$,
sprzeczny z (ii). Zatem wszystkie ideały są
skończenie generowane.
\end{proof}

\begin{przyklad}[Pierścienie noetherowskie i nie-noetherowskie]
\label{aa:prz:noetherian-example}
Niech $K$ będzie ciałem. Każda dziedzina ideałów
głównych jest noetherowska, ponieważ każdy jej ideał
ma jeden generator. W szczególności $\mathbb Z$
i $K[x]$ są noetherowskie.

Pierścień $K[x_1,x_2,\ldots]$ wielomianów w przeliczalnie
wielu zmiennych nie jest noetherowski. Istotnie, ciąg
ideałów $I_n=(x_1,\ldots,x_n)$ jest ściśle rosnący.
Gdyby $x_{n+1}\in I_n$, istniałyby wielomiany $f_i$
takie, że $x_{n+1}=\sum_{i=1}^n x_if_i$.
Po podstawieniu $x_1=\cdots=x_n=0$ prawa strona
wynosiłaby zero, a lewa $x_{n+1}$, co jest sprzecznością
w pierścieniu wielomianów pozostałych zmiennych.
\end{przyklad}

\begin{uwaga}[Dalsza kolejność pojęć]
\label{aa:uw:ring-further}
W kolejnym etapie rozbudujemy teorię pierścieni
noetherowskich i artinowskich, twierdzenie Hilberta
o bazie, rozkłady pierwotne ideałów i ich zastosowania.
Wcześniejsze wyniki o pierwiastnikach, lokalizacji,
ideałach pierwszych i maksymalnych stanowią do tego
bezpośrednie przygotowanie; nie wymagają pełnej
teorii modułów.
\end{uwaga}

\setlength{\parskip}{4pt}

% AA-13 | Algebra abstrakcyjna | Iteracja VII
\chapter[Pierścienie noetherowskie]{Pierścienie noetherowskie i twierdzenie Hilberta o bazie}
\label{aa:ch:noether-hilbert}

\section{Dlaczego skończona liczba generatorów jest istotna?}
W pierścieniu $\mathbb Z$ każdy ideał jest wyznaczony przez
jeden element, natomiast w $K[x,y]$, gdzie $K$ jest ciałem,
nie każdy ideał jest główny.
Nie oznacza to jednak, że do opisania ideału potrzeba zawsze
nieskończenie wielu generatorów. Twierdzenie Hilberta o bazie
mówi, że przejście od pierścienia noetherowskiego do pierścienia
wielomianów jednej zmiennej zachowuje skończone generowanie
\emph{wszystkich} ideałów.

W rozdziale pracujemy z przemiennymi pierścieniami z jedynką,
jeśli nie wskazano inaczej. W przypadku pierścienia zerowego
stwierdzenia o skończonym generowaniu są natychmiastowe;
w przykładach geometrycznych przyjmujemy pierścienie niezerowe.
Przypomnijmy, że ideał
$I\triangleleft R$ jest skończenie generowany, jeśli
$I=(a_1,\ldots,a_r)$ dla pewnej skończonej rodziny elementów
$a_i\in R$. Dopuszczamy $r=0$; wówczas rodzina generatorów jest pusta,
a generowany przez nią ideał jest zerowy.

\section{Skończone generowanie a łańcuchy ideałów}
\begin{definicja}[Pierścień noetherowski]
\label{aa:def:noether-full}
Przemienny pierścień $R$ z jedynką jest \emph{noetherowski},
jeżeli każdy jego ideał jest skończenie generowany.
Równoważnie, na mocy
Twierdzenia~\ref{aa:tw:noetherian-acc},
każdy wstępujący ciąg ideałów
$I_1\subseteq I_2\subseteq\cdots$ stabilizuje się,
to znaczy istnieje $N$ takie, że $I_n=I_N$ dla każdego
$n\ge N$.
\end{definicja}

\begin{lemat}[Maksymalny element rodziny ideałów]
\label{aa:lem:noether-max-ideal-family}
Niech $R$ będzie pierścieniem noetherowskim, a
$\mathcal F$ niepustą rodziną ideałów $R$.
Wówczas $\mathcal F$ zawiera ideał maksymalny względem
zawierania w rodzinie $\mathcal F$.
\end{lemat}
\begin{proof}
Przypuśćmy, że żaden ideał $I_1\in\mathcal F$ nie jest
maksymalny. Wtedy istnieje $I_2\in\mathcal F$ ściśle
zawierający $I_1$. Skoro również $I_2$ nie jest maksymalny,
wybieramy $I_3\in\mathcal F$ z $I_2\subsetneq I_3$.
Powtarzając tę konstrukcję, otrzymujemy nieskończony
ściśle wstępujący ciąg ideałów. Przeczy to warunkowi
stabilizacji. Zatem pewien element $\mathcal F$ musi
być maksymalny.
\end{proof}

\begin{twierdzenie}[Ilorazy pierścieni noetherowskich]
\label{aa:tw:noether-quotients}
Jeśli $R$ jest pierścieniem noetherowskim i
$I\triangleleft R$ jest ideałem, to $R/I$ jest
pierścieniem noetherowskim.
\end{twierdzenie}
\begin{proof}
Niech $J$ będzie dowolnym ideałem $R/I$, a
$\pi:R\to R/I$ kanonicznym homomorfizmem ilorazowym.
Przeciwobraz $K=\pi^{-1}(J)$ jest ideałem $R$
zawierającym $I$. Z noetherowskości $R$ istnieją
$a_1,\ldots,a_r\in K$ takie, że $K=(a_1,\ldots,a_r)$.
Dla każdego $y\in J$ wybieramy $x\in K$ z $\pi(x)=y$
i zapisujemy $x=\sum_{i=1}^r t_ia_i$ z $t_i\in R$.
Po zastosowaniu $\pi$ otrzymujemy
$y=\sum_i\pi(t_i)\pi(a_i)$.
Zatem $J=(\pi(a_1),\ldots,\pi(a_r))$ w $R/I$.
Każdy ideał ilorazu ma skończenie wielu generatorów.
\end{proof}

\begin{twierdzenie}[Lokalizacje pierścieni noetherowskich]
\label{aa:tw:noether-localizations}
Niech $R$ będzie pierścieniem noetherowskim, a
$S\subseteq R$ zbiorem multiplikatywnym niezawierającym
zera. Wtedy lokalizacja $S^{-1}R$ jest noetherowska.
\end{twierdzenie}
\begin{proof}
Niech $J\triangleleft S^{-1}R$ i oznaczmy przez
$\iota:R\to S^{-1}R$ kanoniczny homomorfizm.
Przeciwobraz $I=\iota^{-1}(J)$ jest ideałem $R$,
więc $I=(a_1,\ldots,a_r)$ dla pewnych $a_i\in R$.
Wykażemy, że $J$ jest generowany przez ułamki $a_i/1$.
Każdy $a_i/1$ należy do $J$, więc zawieranie ideału
przez nie generowanego w $J$ jest oczywiste.
W drugą stronę, jeżeli $b/s\in J$, to
$b/1=(s/1)(b/s)\in J$, zatem $b\in I$.
Istnieją $c_1,\ldots,c_r\in R$ z
$b=\sum_i c_ia_i$. Dzieląc przez $s$ w lokalizacji,
otrzymujemy $b/s=\sum_i (c_i/s)(a_i/1)$.
Oznacza to, że $J=(a_1/1,\ldots,a_r/1)$.
\end{proof}

\begin{przyklad}[Pierścień lokalny uzyskany z wielomianów]
\label{aa:prz:noether-local-ring}
Niech $K$ będzie ciałem. Ponieważ $K[x,y]$ będzie noetherowski na mocy twierdzenia
Hilberta, jego lokalizacja
$K[x,y]_{(x,y)}$ także jest noetherowska.
Jest ona jednocześnie pierścieniem lokalnym: jedyny ideał
maksymalny $\mathfrak m=(x,y)K[x,y]_{(x,y)}$ jest generowany
przez obrazy $x$ i $y$.
Nie jest to ciało, ponieważ $x$ należy do $\mathfrak m$
i nie jest jednostką. Ponadto $\mathfrak m$ nie jest główny.
Istotnie, przestrzeń $\mathfrak m/\mathfrak m^2$ nad
$K\cong K[x,y]_{(x,y)}/\mathfrak m$ ma bazę złożoną z klas
$x$ i $y$: każdy licznik elementu $\mathfrak m$ ma
wyraz stały równy zeru, a mianownik ma niezerowy wyraz
stały. Dokładniej, jeśli $s=s_0+q$ z $s_0\in K^\times$
i $q\in(x,y)$, to
$s^{-1}-s_0^{-1}=-s_0^{-1}q/s\in\mathfrak m$.
Po pomnożeniu przez licznik należący do $\mathfrak m$
widzimy, że modulo $\mathfrak m^2$ ułamek jest kombinacją
liniową klas $x$ i $y$. Klasy są liniowo niezależne:
jeśli $ax+by\in\mathfrak m^2$ dla $a,b\in K$, to z
kryterium przynależności do rozszerzonego ideału
(Lemat~\ref{aa:lem:local-ideal-membership}) wynika,
że pewien wielomian $s$ z $s(0,0)\ne0$ spełnia
$s(ax+by)\in(x^2,xy,y^2)$. Porównanie składników
stopnia jeden daje $s(0,0)a=s(0,0)b=0$, a zatem
$a=b=0$. Zatem
$\dim_K(\mathfrak m/\mathfrak m^2)=2$. Gdyby
$\mathfrak m=(h)$, ten iloraz byłby generowany przez
jedną klasę $h$ nad $K$, więc miałby wymiar najwyżej
jeden, sprzeczność. Noetherowskość nie oznacza zatem
ani główności ideałów, ani odwracalności wszystkich
elementów niezerowych.
\end{przyklad}

\section{Moduły nad pierścieniami noetherowskimi}
Najpierw wyprowadzimy narzędzie dotyczące podmodułów
modułów skończenie generowanych. Następnie przedstawimy
bezpośredni dowód twierdzenia Hilberta o bazie oparty na
współczynnikach wiodących. Wyniki o podmodułach wykorzystamy
ponadto do badania algebr skończonych jako moduły.

\begin{definicja}[Moduł skończenie generowany]
\label{aa:def:noether-fg-module}
Niech $R$ będzie przemiennym pierścieniem z jedynką,
a $M$ modułem nad $R$. Mówimy, że $M$ jest
\emph{skończenie generowany}, jeśli dla pewnej liczby całkowitej
$r\ge0$ istnieją elementy $m_1,\ldots,m_r\in M$
takie, że każdy $m\in M$ jest
postaci $m=\sum_{i=1}^r a_im_i$ dla $a_i\in R$.
Gdy $r=0$, suma jest pusta i $M=0$.
\end{definicja}

\begin{lemat}[Podmoduły modułów wolnych skończonej rangi]
\label{aa:lem:noether-submodules-free}
Jeśli $R$ jest pierścieniem noetherowskim, to każdy
podmoduł $N$ modułu $R^r$ dla $r\ge0$ jest
skończenie generowany.
\end{lemat}
\begin{proof}
Indukujemy względem $r$. Dla $r=0$ moduł $R^0$ jest zerowy,
więc jedynym podmodułem jest moduł zerowy.
Dla $r=1$ podmoduły $R$ to dokładnie ideały, które są
skończenie generowane na mocy założenia.

Niech $r\ge2$ i niech $N\subseteq R^r$.
Rozważmy odwzorowanie $p:R^r\to R$, które wybiera
ostatnią współrzędną. Obraz $p(N)$ jest ideałem $R$,
więc ma generatory $a_1,\ldots,a_s$.
Dla każdego $a_i$ wybieramy $v_i\in N$ z
$p(v_i)=a_i$. Z drugiej strony
$N\cap(R^{r-1}\times\{0\})$ jest podmodułem
$R^{r-1}$ po pominięciu ostatniej współrzędnej,
więc na mocy założenia indukcyjnego ma generatory
$w_1,\ldots,w_t$. Dla dowolnego $v\in N$ zapisujemy
$p(v)=\sum_i c_i a_i$. Wtedy
$v-\sum_i c_i v_i\in N$ ma ostatnią współrzędną zerową,
więc jest kombinacją $w_j$. W rezultacie elementy
$v_i,w_j$ generują cały $N$.
\end{proof}

\begin{twierdzenie}[Podmoduły modułów skończenie generowanych]
\label{aa:tw:noether-submodules-fg}
Niech $R$ będzie przemiennym pierścieniem noetherowskim
z jedynką. Jeśli $M$ jest skończenie generowanym
$R$-modułem, to każdy podmoduł $N\subseteq M$
jest skończenie generowany.
\end{twierdzenie}
\begin{proof}
Wybierzmy generatory $m_1,\ldots,m_r$ modułu $M$.
Homomorfizm $\varphi:R^r\to M$ dany przez
$\varphi(a_1,\ldots,a_r)=\sum_i a_im_i$ jest
surjektywny. Przeciwobraz $\varphi^{-1}(N)$ jest
podmodułem $R^r$, więc poprzedni lemat daje jego
skończony układ generatorów $u_1,\ldots,u_s$.
Ich obrazy $\varphi(u_i)$ generują $N$: dla każdego
$n\in N$ wybieramy $u\in R^r$ z $\varphi(u)=n$,
rozkładamy $u$ na sumę kombinacji $u_i$ i stosujemy
$\varphi$. Tak otrzymujemy wymagane przedstawienie.
\end{proof}

\begin{przyklad}[Dlaczego wymagana jest noetherowskość]
\label{aa:prz:noether-need}
Niech $K$ będzie ciałem i niech $R=K[x_1,x_2,\ldots]$ będzie pierścieniem
wielomianów w przeliczalnie wielu zmiennych.
Jako moduł nad sobą $R$ jest generowany przez $1$.
Podmoduł $I=(x_1,x_2,\ldots)$ nie jest jednak
skończenie generowany. Gdyby generowały go
$f_1,\ldots,f_r$, to wszystkie te wielomiany
występowałyby jedynie w skończonej liczbie zmiennych
$x_1,\ldots,x_N$. Po podstawieniu tych zmiennych
równych zeru obrazy generatorów byłyby zerowe
(każdy $f_i$ ma wyraz wolny równy zeru), podczas gdy
obraz $x_{N+1}$ pozostałby niezerowy. Sprzeczność.
\end{przyklad}


\begin{twierdzenie}[Skończone prezentacje modułów]
\label{aa:tw:noether-finite-presentation}
Niech $R$ będzie przemiennym pierścieniem noetherowskim
z jedynką, a $M$ skończenie generowanym $R$-modułem.
Wtedy istnieją liczby całkowite $r,s\geq0$ oraz
homomorfizmy $R$-modułów
$R^s\xrightarrow{\,A\,}R^r\xrightarrow{\,\varphi\,}M$
takie, że $\varphi$ jest surjekcją oraz
$\operatorname{im} A=\ker\varphi$.
Innymi słowy, moduł $M$ ma skończony układ generatorów
i skończony układ relacji między nimi.
\end{twierdzenie}
\begin{proof}
Wybierzmy generatory $m_1,\ldots,m_r$ modułu $M$.
Zdefiniujmy $\varphi:R^r\to M$ wzorem
$\varphi(c_1,\ldots,c_r)=\sum_{i=1}^r c_i m_i$.
Definicja generowania gwarantuje, że $\varphi$ jest
surjekcją. Jądro $K=\ker\varphi$ jest podmodułem
$R^r$, a z Lematu~\ref{aa:lem:noether-submodules-free}
wynika, że istnieją jego generatory $v_1,\ldots,v_s$.
Niech $A:R^s\to R^r$ odwzorowuje $j$-ty standardowy
wektor bazowy na $v_j$. Wówczas obraz $A$ jest
podmodułem generowanym przez $v_j$, czyli dokładnie
$K$. Otrzymaliśmy $\operatorname{im}A=\ker\varphi$,
co jest żądaną skończoną prezentacją.
Jeśli $K=0$, można przyjąć $s=0$.
\end{proof}

\begin{przyklad}[Prezentacja modułu $\mathbb Z/6\mathbb Z$]
\label{aa:prz:noether-finite-presentation-z6}
Dla pierścienia $R=\mathbb Z$ i $M=\mathbb Z/6\mathbb Z$
bierzemy $\varphi:\mathbb Z\to M$, $n\mapsto\overline n$.
Jego jądrem jest $6\mathbb Z$, więc
$A:\mathbb Z\to\mathbb Z$, $A(n)=6n$, ma obraz
równy $\ker\varphi$. Otrzymujemy ciąg dokładny
$\mathbb Z\xrightarrow{\,\cdot6\,}\mathbb Z
\xrightarrow{\,\varphi\,}\mathbb Z/6\mathbb Z\to0$.
Jeden generator $\overline1$ wystarcza, lecz spełnia
nietrywialną relację $6\overline1=0$; z tego powodu
moduł nie jest wolny rangi jeden.
\end{przyklad}

\section{Twierdzenie Hilberta o bazie}
\begin{twierdzenie}[Twierdzenie Hilberta o bazie]
\label{aa:tw:hilbert-basis}
Jeżeli $R$ jest przemiennym pierścieniem noetherowskim
z jedynką, to pierścień wielomianów $R[x]$ jest
noetherowski.
\end{twierdzenie}
\begin{proof}
Niech $I\triangleleft R[x]$ będzie dowolnym ideałem.
Jeżeli $I=(0)$, to jest generowany przez pustą rodzinę.
Załóżmy więc $I\ne(0)$. Zdefiniujmy $L\subseteq R$
jako zbiór zawierający $0$ oraz wszystkie współczynniki
wiodące niezerowych wielomianów należących do $I$.

Pokażemy, że $L$ jest ideałem $R$. Jeśli $a,b\in L$
i nie są zerami, wybierzmy $f,g\in I$ o współczynnikach
wiodących odpowiednio $a,b$ i stopniach $m,n$.
Mnożąc $f$ przez $x^{\max(m,n)-m}$, a $g$ przez
$x^{\max(m,n)-n}$, otrzymujemy wielomiany
tego samego stopnia, których suma ma
współczynnik wiodący $a+b$, gdy $a+b\ne0$.
W przypadku $a+b=0$ wynik należy do $L$ z definicji.
Podobnie dla $r\in R$ współczynnik wiodący $rf$
równa się $ra$, jeśli $ra\ne0$, a w przeciwnym razie
$ra=0\in L$. Również elementy przeciwne należą do $L$.
Zatem $L$ jest ideałem.

Ponieważ $R$ jest noetherowski, mamy
$L=(a_1,\ldots,a_s)$. Wybieramy dla każdego
niezerowego generatora $a_i$ wielomian $f_i\in I$
o współczynniku wiodącym $a_i$ i stopniu $d_i$.
Generatory zerowe można pominąć. Niech
$D=\max(d_1,\ldots,d_s)$. Ponieważ $I\ne(0)$,
ideał $L$ jest niezerowy i lista wybranych wielomianów
nie jest pusta.

Dla każdego $j$ z $0\le j<D$ definiujemy
$L_j\subseteq R$ jako zbiór zawierający $0$ oraz
współczynniki wiodące wielomianów $h\in I$ o
stopniu dokładnie $j$. Powyższy argument bez
przesuwania stopni pokazuje, że każdy $L_j$ jest
ideałem (gdy dwa współczynniki sumują się do zera,
zero już należy do zbioru). Wybieramy skończony układ
generatorów $b_{j,1},\ldots,b_{j,r_j}$ ideału $L_j$
i wielomiany $g_{j,k}\in I$ stopnia $j$ o odpowiednich
współczynnikach wiodących. Jeśli $L_j=0$, nie wybieramy
żadnego wielomianu.

Udowodnimy, że skończona rodzina wszystkich $f_i$ i
$g_{j,k}$ generuje $I$. Weźmy niezerowy wielomian
$h\in I$ o stopniu $n$ i współczynniku wiodącym $c$.
Jeżeli $n\ge D$, piszemy $c=\sum_i r_i a_i$,
a następnie odejmujemy od $h$ wielomian
$\sum_i r_i x^{n-d_i}f_i\in I$. Wyrazy stopnia $n$
się redukują, więc pozostały wielomian jest zerowy
albo ma stopień mniejszy od $n$.
Jeśli natomiast $n<D$, to $c\in L_n$, zatem
$c=\sum_k r_k b_{n,k}$. Odejmując
$\sum_k r_k g_{n,k}$, ponownie obniżamy stopień.

W obu przypadkach uzyskujemy wielomian w $I$
o ściśle mniejszym stopniu. Powtarzamy krok aż do
zera, co musi nastąpić, bo stopnie niezerowych
wielomianów są nieujemne i nie mogą maleć bez końca.
Odjęte wielomiany należą do ideału generowanego przez
skończoną rodzinę $f_i,g_{j,k}$; zatem także
początkowy $h$ należy do tego ideału. Dowodzi to,
że $I$ jest skończenie generowany. Ponieważ $I$
był dowolny, $R[x]$ jest noetherowski.
\end{proof}

\begin{wniosek}[Pierścienie wielomianów skończenie wielu zmiennych]
\label{aa:wn:hilbert-finite-vars}
Jeśli $R$ jest pierścieniem noetherowskim, to
$R[x_1,\ldots,x_n]$ jest noetherowski dla każdego
$n\ge0$. W szczególności $K[x_1,\ldots,x_n]$
jest noetherowski dla każdego ciała $K$.
\end{wniosek}
\begin{proof}
Dla $n=0$ otrzymujemy sam pierścień $R$.
Gdy wynik jest prawdziwy dla $n-1$, zastosowanie
twierdzenia Hilberta do
$R[x_1,\ldots,x_{n-1}]$ daje noetherowskość
$(R[x_1,\ldots,x_{n-1}])[x_n]$, czyli
$R[x_1,\ldots,x_n]$. Ciało jest noetherowskie,
ponieważ jego jedynymi ideałami są $(0)$ i całe ciało.
\end{proof}

\begin{przyklad}[Skończone generowanie nie oznacza jednego generatora]
\label{aa:prz:hilbert-not-pid}
Niech $K$ będzie ciałem. Ideał $(x,y)\triangleleft K[x,y]$ ma dwa generatory.
Nie ma jednego generatora, co wykazaliśmy w
Przykładzie~\ref{aa:prz:two-polynomials-not-pid}.
Twierdzenie Hilberta gwarantuje jedynie, że dowolny
ideał $K[x,y]$ posiada \emph{jakąś skończoną}
rodzinę generatorów. Nie gwarantuje, że wystarczy jeden.
\end{przyklad}


\begin{twierdzenie}[Noetherowskość algebr skończonych jako moduły]
\label{aa:tw:noether-module-finite-algebra}
Niech $R$ będzie przemiennym pierścieniem noetherowskim
z jedynką, a $S$ przemiennym pierścieniem z jedynką,
wyposażonym w homomorfizm pierścieni
$\eta:R\to S$ zachowujący jedynkę. Jeśli $S$ jest
skończenie generowanym $R$-modułem dla działania
$r\cdot s=\eta(r)s$, to $S$ jest pierścieniem
noetherowskim.
\end{twierdzenie}
\begin{proof}
Niech $J$ będzie dowolnym ideałem pierścienia $S$.
Ponieważ $\eta(r)s\in J$ dla $r\in R$ i $s\in J$,
ideał $J$ jest też $R$-podmodułem $S$.
Twierdzenie~\ref{aa:tw:noether-submodules-fg}
stosuje się do $R$-modułu $S$, ponieważ jest on
skończenie generowany. Otrzymujemy
$J=Ra_1+\cdots+Ra_m$ dla pewnych $a_i\in J$.
Każde $a_i$ należy do $J$, zatem ideał $S$
generowany przez te elementy jest zawarty w $J$.
Z drugiej strony
$J=\sum_i\eta(R)a_i\subseteq\sum_i Sa_i$.
Stąd $J=(a_1,\ldots,a_m)$ jako ideał $S$.
Ponieważ $J$ był dowolny, $S$ jest noetherowski.
Nie wymagamy, by $\eta$ był iniektywny.
\end{proof}

\begin{przyklad}[Podalgebra generowana przez $t^2$ i $t^3$]
\label{aa:prz:noether-cusp-algebra}
Niech $K$ będzie ciałem, $R=K[t^2]$ i
$S=K[t^2,t^3]\subseteq K[t]$. Pierścień $R$
jest izomorficzny z $K[u]$ przez odwzorowanie
$u\mapsto t^2$, zatem jest noetherowski.
Pokażemy, że $S=R+Rt^3$ jako $R$-moduł.
Każdy iloczyn potęg generatorów $t^2,t^3$ ma postać
$t^{2a+3b}$ dla $a,b\ge0$. Gdy $b$ jest parzyste,
$t^{2a+3b}\in R$; gdy $b$ jest nieparzyste,
$t^{2a+3b}=t^3(t^2)^{a+3(b-1)/2}\in Rt^3$.
W drugą stronę zarówno $R$, jak i $Rt^3$ są
podzbiorami $S$. Stąd $S=R+Rt^3$ i twierdzenie
daje noetherowskość $S$.
Jednocześnie $t\notin S$: żaden z jednomianów
generujących $S$ nie ma stopnia jeden, więc
wielomian $t$ nie może być ich kombinacją liniową.
To pokazuje, że skończone generowanie $S$ nad $R$
nie wymaga równości $S=K[t]$.
\end{przyklad}

\begin{przyklad}[Dlaczego same współczynniki wiodące nie wystarczają]
\label{aa:prz:noether-leading-coefficients-example}
Rozważmy ideał $I=(2,x)$ w $\mathbb Z[x]$.
Wszystkie współczynniki wiodące niezerowych
wielomianów z $I$ generują ideał $L=\mathbb Z$:
ponieważ $x\in I$ ma współczynnik wiodący równy $1$,
wystarczy jeden generator $1$.
Wybrany dla niego wielomian $f=x$ ma stopień $D=1$,
lecz sam $x$ nie generuje całego $I$, ponieważ
$2\notin(x)$.
W konstrukcji z dowodu twierdzenia Hilberta potrzeba
jeszcze ideału współczynników stałych $L_0=(2)$,
któremu odpowiada wielomian $g=2$.
Rzeczywiście, jeśli $h=c_0+c_1x+\cdots+c_nx^n\in I$,
to $c_0=h(0)$ jest parzyste, gdyż $h=2a+xb$
dla pewnych $a,b\in\mathbb Z[x]$. Odwrotnie,
gdy $c_0=2k$, wielomian $q=(h-c_0)/x$ należy
do $\mathbb Z[x]$ (także dla $h$ stałego, gdy $q=0$),
a więc $h=2k+xq\in(2,x)$.
Przykładowo $3x^2+5x+6=x(3x+5)+3\cdot2$.
To obliczenie wyjaśnia rolę osobnych generatorów
$L_j$ dla stopni mniejszych niż $D$.
\end{przyklad}

\section{Ideały pierwotne: pierwsze przykłady}
\begin{definicja}[Ideał pierwotny]
\label{aa:def:primary-ideal}
Niech $R$ będzie przemiennym pierścieniem z jedynką.
Właściwy ideał $Q\subsetneq R$ nazywamy
\emph{pierwotnym}, jeśli dla wszystkich $a,b\in R$
z warunków $ab\in Q$ oraz $a\notin Q$ wynika,
że $b^n\in Q$ dla pewnego $n\ge1$.
\end{definicja}

\begin{lemat}[Pierwiastnik ideału pierwotnego]
\label{aa:lem:primary-radical-prime}
Niech $R$ będzie przemiennym pierścieniem z jedynką,
a $Q\subsetneq R$ ideałem pierwotnym. Wówczas $\sqrt Q$
jest ideałem pierwszym pierścienia $R$.
\end{lemat}
\begin{proof}
Pierwiastnik $\sqrt Q$ jest ideałem na mocy
Twierdzenia~\ref{aa:tw:radical-ideal}; jest właściwy,
bo $1\in\sqrt Q$ oznaczałoby $1\in Q$.
Niech $ab\in\sqrt Q$. Wtedy $(ab)^n\in Q$
dla pewnego $n\ge1$. Jeśli $a\in\sqrt Q$, teza
o pierwszości jest spełniona. Załóżmy więc
$a\notin\sqrt Q$, czyli $a^n\notin Q$.
Ponieważ $a^n b^n\in Q$ i $Q$ jest pierwotny,
istnieje $k\ge1$ takie, że $(b^n)^k\in Q$.
Stąd $b\in\sqrt Q$. To dowodzi pierwszości.
\end{proof}

\begin{przyklad}[Rozkład ideału na ideały pierwotne]
\label{aa:prz:primary-decomposition-explicit}
W pierścieniu $K[x,y]$ nad ciałem $K$ zachodzi
\[
 (x^2,xy)=(x)\cap(x^2,y).
\]
Uzasadnijmy obie inkluzje. Generatory $x^2$ i $xy$
należą do $(x)$ oraz $(x^2,y)$, więc lewa strona
jest zawarta w przecięciu. W drugą stronę, jeśli
$f\in(x)\cap(x^2,y)$, to $f=xg$ i
$f=x^2 a+yb$ dla pewnych $g,a,b\in K[x,y]$.
Po podstawieniu $x=0$ otrzymujemy
$0=y\,b(0,y)$ w dziedzinie $K[y]$, więc $b(0,y)=0$.
Oznacza to, że $x$ dzieli $b$, czyli $b=xc$.
Wówczas $f=x^2 a+xyc\in(x^2,xy)$.

Ideał $(x)$ jest pierwszy, ponieważ
$K[x,y]/(x)\cong K[y]$ jest dziedziną całkowitości.
Ideał $(x^2,y)$ jest pierwotny. Istotnie, jego iloraz
jest izomorficzny z $K[x]/(x^2)$. W tym ostatnim
pierścieniu każda klasa ma jednoznaczną postać
$a+bx+(x^2)$ z $a,b\in K$. Jeśli $a\ne0$, klasa
jest jednostką: jej odwrotność wynosi
$a^{-1}-a^{-2}bx+(x^2)$. Jeśli $a=0$, klasa ma
kwadrat równy zeru. Jeśli iloczyn klas $u,v$ jest zerowy,
a $u$ jest niezerowa, to $v$ nie może być jednostką,
bo pomnożenie przez $v^{-1}$ dałoby $u=0$.
Zatem $v^2=0$, co jest warunkiem pierwotności
ideału $(x^2,y)$.

Jest to \emph{konkretny przykład} rozkładu na ideały
pierwotne, a nie jeszcze dowód ogólnego twierdzenia
o rozkładzie pierwotnym ideałów pierścieni noetherowskich.
Pełną teorię tej dekompozycji rozwiniemy dopiero po
ustaleniu dalszych narzędzi modułowych.
\end{przyklad}

% AA-14 | Algebra abstrakcyjna | Iteracja VII
\chapter[Pierścienie artinowskie i lokalne]{Pierścienie artinowskie i lokalne oraz lemat Nakayamy}
\label{aa:ch:artin-local}
\setlength{\parskip}{2pt}

\section{Dwa przeciwne warunki łańcucha}
Pierścienie noetherowskie charakteryzują się stabilizacją
\emph{wstępujących} ciągów ideałów. Można też wymagać,
aby stabilizowały się ciągi zstępujące. Te dwa warunki
nie są równoważne, a ich konsekwencje są bardzo różne.
W całym rozdziale pierścienie są przemienne, mają jedynkę
i są niezerowe, chyba że wskazano inaczej. Warunki łańcucha
rozważamy również dla pierścienia zerowego, w którym zachodzą
trywialnie.

\begin{definicja}[Pierścień artinowski]
\label{aa:def:artinian-ring}
Przemienny pierścień $R$ z jedynką nazywamy
\emph{artinowskim}, jeśli każdy
zstępujący ciąg jego ideałów
$I_1\supseteq I_2\supseteq\cdots$ jest od pewnego miejsca
stały, to znaczy istnieje $N$ takie, że
$I_n=I_N$ dla wszystkich $n\ge N$.
\end{definicja}

\begin{przyklad}[Pierścienie skończone, ciała i liczby całkowite]
\label{aa:prz:artin-basic}
Każdy pierścień skończony jest artinowski: w ściśle
zstępującym ciągu podzbiorów skończonego zbioru
liczności musiałyby ściśle maleć bez końca.
Dowolne ciało $K$ także jest artinowskie, nawet gdy
jest nieskończone, ponieważ ma tylko ideały $(0)$ i $K$.
Natomiast $\mathbb Z$ nie jest artinowski:
\[
 (2)\supsetneq(2^2)\supsetneq(2^3)\supsetneq\cdots.
\]
Każde zawieranie jest ścisłe, bo $2^n$ nie jest
wielokrotnością $2^{n+1}$ w $\mathbb Z$.
Analogicznie $K[x]$ nie jest artinowski dzięki
łańcuchowi $(x)\supsetneq(x^2)\supsetneq\cdots$.
Oba pierścienie są jednak noetherowskie.
\end{przyklad}

\begin{lemat}[Ilorazy pierścieni artinowskich]
\label{aa:lem:artin-quotient}
Niech $R$ będzie przemiennym pierścieniem artinowskim
z jedynką, a $I\triangleleft R$ jego ideałem.
Wówczas $R/I$ jest artinowski, również gdy $I=R$.
\end{lemat}
\begin{proof}
Niech $J_1\supseteq J_2\supseteq\cdots$ będzie
zstępującym ciągiem ideałów $R/I$.
Przeciwobrazy $\pi^{-1}(J_n)$ przez kanoniczny
homomorfizm $\pi:R\to R/I$ tworzą zstępujący ciąg
ideałów $R$, więc od pewnego miejsca są równe.
Ponieważ $\pi$ jest surjekcją, równość przeciwobrazów
implikuje równość odpowiadających im ideałów $J_n$.
\end{proof}

\section{Dziedziny artinowskie są ciałami}
\begin{twierdzenie}[Dziedzina całkowitości artinowska jest ciałem]
\label{aa:tw:artin-domain-field}
Każda dziedzina całkowitości $D$, która jest
pierścieniem artinowskim, jest ciałem.
\end{twierdzenie}
\begin{proof}
Niech $a\in D$ będzie niezerowy. Rozważmy ciąg
ideałów $(a)\supseteq(a^2)\supseteq(a^3)\supseteq\cdots$.
Z artinowskości istnieje $n\ge1$ takie, że
$(a^n)=(a^{n+1})$. W szczególności $a^n$ należy
do ideału $(a^{n+1})$, zatem dla pewnego $b\in D$
mamy $a^n=a^{n+1}b$. Przepisujemy tę równość jako
$a^n(1-ab)=0$. Ponieważ $D$ jest dziedziną i
$a^n\ne0$, możemy skrócić $a^n$, otrzymując
$ab=1$. Każdy niezerowy element $D$ ma więc
odwrotność; przemienność daje ciało.
\end{proof}

\begin{wniosek}[Ideały pierwsze pierścienia artinowskiego]
\label{aa:wn:artin-primes-maximal}
Niech $R$ będzie przemiennym pierścieniem artinowskim
z jedynką, a $P$ jego ideałem pierwszym.
Wówczas $P$ jest ideałem maksymalnym.
\end{wniosek}
\begin{proof}
Dla ideału pierwszego $P$ iloraz $R/P$ jest
dziedziną całkowitości. Na mocy
Lematu~\ref{aa:lem:artin-quotient} jest też artinowski,
więc poprzednie twierdzenie pokazuje, że jest ciałem.
Ideał $P$ jest maksymalny wtedy i tylko wtedy, gdy
iloraz $R/P$ jest ciałem.
\end{proof}

\section{Pierścienie lokalne i radykał Jacobsona}
\begin{definicja}[Radykał Jacobsona]
\label{aa:def:jacobson-radical}
Niech $R$ będzie niezerowym przemiennym pierścieniem
z jedynką. \emph{Radykałem Jacobsona} nazywamy ideał
\[
 J(R)=\bigcap_{\mathfrak m\in\operatorname{Max}(R)}\mathfrak m,
\]
gdzie $\operatorname{Max}(R)$ oznacza rodzinę wszystkich
ideałów maksymalnych pierścienia $R$.
Każdy niezerowy pierścień przemienny z jedynką ma co najmniej
jeden ideał maksymalny: stosujemy lemat Kuratowskiego--Zorna
do rodziny ideałów właściwych uporządkowanych przez zawieranie.
Rodzina jest niepusta, bo zawiera ideał zerowy. Dla
łańcucha pustego jego ograniczeniem górnym jest ideał
zerowy. Suma mnogościowa niepustego łańcucha ideałów
właściwych jest ideałem: dla dwóch jej elementów
wybieramy z łańcucha ideał zawierający oba. Pozostaje
on właściwy, gdyż obecność $1_R$ w tej sumie oznaczałaby,
że $1_R$ należy do jednego z ideałów łańcucha,
wbrew jego właściwości.
\end{definicja}

\begin{twierdzenie}[Kryterium jednostek dla radykału Jacobsona]
\label{aa:tw:jacobson-unit-criterion}
Niech $R$ będzie niezerowym przemiennym pierścieniem
z jedynką i $a\in R$. Następujące warunki są równoważne:
\begin{enumerate}[label=\textup{(\roman*)}]
\item $a\in J(R)$;
\item dla każdego $b\in R$ element $1-ab$ jest jednostką $R$.
\end{enumerate}
\end{twierdzenie}
\begin{proof}
Załóżmy $a\in J(R)$ i wybierzmy $b\in R$.
Gdyby $1-ab$ nie był jednostką, ideał $(1-ab)$
byłby właściwy i zawierałby się w pewnym ideale
maksymalnym $\mathfrak m$. Ale $a\in\mathfrak m$,
a więc także $ab\in\mathfrak m$; stąd
$1=(1-ab)+ab\in\mathfrak m$, sprzeczność.
Oznacza to, że $1-ab$ jest jednostką.

Odwrotnie, przypuśćmy, że $a\notin J(R)$.
Istnieje wtedy ideał maksymalny $\mathfrak m$ taki,
że $a\notin\mathfrak m$. Ponieważ ideał
$\mathfrak m+(a)$ ściśle zawiera $\mathfrak m$,
maksymalność daje $\mathfrak m+(a)=R$.
Istnieją zatem $c\in\mathfrak m$ i $b\in R$
takie, że $1=c+ab$. Wtedy $1-ab=c\in\mathfrak m$
i nie jest jednostką. Warunek (ii) zawodzi.
\end{proof}

\begin{twierdzenie}[Charakterystyka pierścieni lokalnych]
\label{aa:tw:local-nonunits}
Dla niezerowego przemiennego pierścienia $R$ z jedynką
następujące warunki są równoważne:
\begin{enumerate}[label=\textup{(\roman*)}]
\item $R$ ma dokładnie jeden ideał maksymalny $\mathfrak m$;
\item zbiór wszystkich elementów nieodwracalnych $R$
      jest ideałem właściwym.
\end{enumerate}
Jeżeli te warunki zachodzą, to element $a\in R$ jest
jednostką wtedy i tylko wtedy, gdy $a\notin\mathfrak m$.
\end{twierdzenie}
\begin{proof}
Załóżmy (i). Żadna jednostka nie należy do ideału
właściwego, więc każdy element $\mathfrak m$ jest
nieodwracalny. Jeśli $a\notin\mathfrak m$ byłby
nieodwracalny, ideał $(a)$ byłby właściwy i zawierałby
się w pewnym ideale maksymalnym. Jedynym takim ideałem
jest $\mathfrak m$, co przeczyłoby $a\notin\mathfrak m$.
Zbiór elementów nieodwracalnych jest więc dokładnie
$\mathfrak m$, czyli ideałem właściwym.

Załóżmy (ii) i oznaczmy ten ideał przez $N$.
Każdy ideał właściwy nie zawiera jednostek, więc
zawiera się w $N$. W szczególności $N$ jest maksymalny
i każdy inny ideał maksymalny musi być równy $N$.
\end{proof}

\begin{przyklad}[Lokalna algebra z elementami nilpotentnymi]
\label{aa:prz:local-artin-examples}
Dla ciała $K$ i liczby całkowitej $n\ge2$ pierścień
$R=K[t]/(t^n)$ jest lokalny. Element ma postać
$a_0+a_1t+\cdots+a_{n-1}t^{n-1}+(t^n)$.
Jeśli $a_0\ne0$, to jest jednostką:
po wyłączeniu $a_0$ piszemy go jako $a_0(1+u)$,
gdzie $u$ jest wielokrotnością klasy $t$ i $u^n=0$.
Odwrotnością $1+u$ jest skończona suma
$1-u+u^2-\cdots+(-u)^{n-1}$. Jeżeli $a_0=0$,
element jest nilpotentny, a zatem nie może być
jednostką w niezerowym pierścieniu. Zbiór
nieodwracalnych elementów jest więc ideałem $(t)$.
Jest to jedyny ideał maksymalny.

Pierścień $R$ jest także artinowski. Istotnie,
ideały $R$ odpowiadają ideałom $K[t]$ zawierającym
$(t^n)$. Ponieważ $K[t]$ jest dziedziną ideałów
głównych, każdy taki ideał ma postać $(f)$.
Zawieranie $(t^n)\subseteq(f)$ oznacza, że
$f$ dzieli $t^n$. Wielomian $t$ jest nierozkładalny
w $K[t]$, zatem jednoznaczność rozkładu na czynniki
nierozkładalne daje $f=ct^j$ dla pewnego $c\in K^\times$
i $0\le j\le n$. Otrzymujemy dokładnie ideały
$(t^j)$ (włącznie z całym pierścieniem i ideałem
zerowym po przejściu do ilorazu).
Mamy więc tylko skończenie wiele ideałów.
\end{przyklad}

\section{Skończoność zbioru ideałów maksymalnych}
\begin{twierdzenie}[Skończenie wiele ideałów maksymalnych]
\label{aa:tw:artin-finite-maximals}
Niezerowy przemienny pierścień artinowski $R$ ma
skończenie wiele ideałów maksymalnych.
\end{twierdzenie}
\begin{proof}
Załóżmy przeciwnie, że istnieje nieskończenie wiele
różnych ideałów maksymalnych
$\mathfrak m_1,\mathfrak m_2,\ldots$.
Połóżmy $I_n=\bigcap_{j=1}^n\mathfrak m_j$.
Jest to zstępujący ciąg ideałów. Pokażemy, że każde
zawieranie $I_{n+1}\subseteq I_n$ jest ścisłe.
Dla każdego $j\le n$ mamy
$\mathfrak m_j\not\subseteq\mathfrak m_{n+1}$:
gdyby zachodziło zawieranie, maksymalność wymuszałaby
równość tych ideałów. Wybierzmy więc
$a_j\in\mathfrak m_j\setminus\mathfrak m_{n+1}$.
Iloczyn $a_1\cdots a_n$ należy do wszystkich
$\mathfrak m_j$ dla $j\le n$, a więc do $I_n$.
Nie należy jednak do $\mathfrak m_{n+1}$, ponieważ
każdy ideał maksymalny w pierścieniu przemiennym
jest pierwszy. Stąd
$a_1\cdots a_n\in I_n\setminus I_{n+1}$.
Otrzymaliśmy nieskończony ściśle zstępujący ciąg
ideałów, sprzeczny z artinowskością $R$.
\end{proof}

\section{Nilpotentność radykału Jacobsona}
\begin{twierdzenie}[Nilpotentność radykału Jacobsona]
\label{aa:tw:artin-jacobson-nilpotent}
Niech $R$ będzie niezerowym przemiennym pierścieniem
artinowskim z jedynką. Wówczas istnieje liczba
całkowita $N\ge1$ taka, że $J(R)^N=(0)$.
\end{twierdzenie}
\begin{proof}
Piszmy $J=J(R)$. Zstępujący ciąg
$J\supseteq J^2\supseteq\cdots$ stabilizuje się,
więc dla pewnego $N\ge1$ mamy
$J^N=J^{N+1}$. Mnożąc tę równość przez kolejne
potęgi $J$, otrzymujemy $J^N=J^{N+k}$ dla każdego
$k\ge0$. Połóżmy $I=J^N$. W szczególności
$I^2=J^{2N}=J^N=I$.

Załóżmy, że $I\ne(0)$. Rozważmy niepustą rodzinę
$\mathcal F$ wszystkich ideałów $K\subseteq I$
takich, że $IK\ne(0)$. Rodzina jest niepusta, ponieważ
$I\in\mathcal F$ na mocy $I^2=I\ne(0)$.
Artinowskość zapewnia, że $\mathcal F$ ma element
minimalny $K$: gdyby żaden element nie był minimalny,
można byłoby zbudować nieskończony ściśle zstępujący
ciąg ideałów w tej rodzinie.

Skoro $IK\ne0$, istnieje $a\in K$ z $Ia\ne0$.
Ideał $(a)$ należy do $\mathcal F$, ponieważ
$I(a)=Ia\ne0$, i zawiera się w $K$, zatem
minimalność daje $K=(a)$.
Ponieważ $IK\subseteq K$, a
$I(IK)=I^2K=IK\ne0$, także $IK$ należy do
$\mathcal F$. Z minimalności $K$ wynika $IK=K$.
Stąd $a\in K=I(a)$, czyli $a=ba$ dla pewnego
$b\in I$. Ponieważ $I=J^N\subseteq J$, mamy
$b\in J$. Z
Twierdzenia~\ref{aa:tw:jacobson-unit-criterion}
wynika, że $1-b$ jest jednostką.
Równość $(1-b)a=0$ wymusza więc $a=0$, co przeczy
$Ia\ne0$. Zatem $I=(0)$, czyli $J^N=(0)$.
\end{proof}

\begin{wniosek}[Iloraz przez radykał Jacobsona]
\label{aa:wn:artin-semisimple-quotient}
Niech $R$ będzie niezerowym przemiennym pierścieniem
artinowskim, a $\mathfrak m_1,\ldots,\mathfrak m_r$
jego wszystkimi ideałami maksymalnymi. Wtedy
\[
 R/J(R)\cong\prod_{i=1}^r R/\mathfrak m_i,
\]
a każdy czynnik po prawej jest ciałem.
\end{wniosek}
\begin{proof}
Różne ideały maksymalne są parami względnie pierwsze:
$\mathfrak m_i+\mathfrak m_j=R$ dla $i\ne j$, ponieważ
suma jest ideałem ściśle zawierającym każdy z nich.
Chińskie twierdzenie o resztach daje surjekcję
$R\to\prod_i R/\mathfrak m_i$ z jądrem
$\bigcap_i\mathfrak m_i=J(R)$. Pierwsze twierdzenie
o izomorfizmie pierścieni kończy dowód.
\end{proof}

\section{Struktura przemiennych pierścieni artinowskich}
\begin{twierdzenie}[Rozkład na pierścienie artinowskie lokalne]
\label{aa:tw:artin-local-product}
Niech $R$ będzie niezerowym przemiennym pierścieniem
artinowskim z jedynką, a
$\mathfrak m_1,\ldots,\mathfrak m_r$ jego wszystkimi
ideałami maksymalnymi. Wtedy dla pewnego $N\ge1$
zachodzi izomorfizm pierścieni
\[
 R\cong\prod_{i=1}^r R/\mathfrak m_i^N.
\]
Każdy pierścień $R/\mathfrak m_i^N$ jest artinowski
i lokalny, a jego jedyny ideał maksymalny jest obrazem
$\mathfrak m_i$.
\end{twierdzenie}
\begin{proof}
Na mocy poprzednich twierdzeń ideałów maksymalnych
jest skończenie wiele, a $J=\bigcap_i\mathfrak m_i$
jest nilpotentny, powiedzmy $J^N=0$.
Przypomnijmy, dlaczego dla ideałów względnie pierwszych
$I+K=R$ zachodzi $I\cap K=IK$. Wybierzmy $u\in I$
i $v\in K$ z $u+v=1$. Jeśli $z\in I\cap K$, to
$z=zu+zv\in IK$, a odwrotne zawieranie jest oczywiste.
Ponadto, jeśli ideał $L$ jest względnie pierwszy z każdym
z ideałów $I_1,\ldots,I_q$, to jest względnie pierwszy
z ich iloczynem: z równości $u_j+v_j=1$, gdzie
$u_j\in I_j$ i $v_j\in L$, po pomnożeniu wszystkich
tych równości otrzymujemy
$1=\prod_j u_j+w$ dla pewnego $w\in L$.
Indukcja daje więc równość przecięcia i iloczynu
dla skończonej rodziny ideałów parami względnie pierwszych.
W szczególności $J=\prod_i\mathfrak m_i$.
Wobec tego
$\prod_i\mathfrak m_i^N=J^N=(0)$.

Potęgi różnych ideałów maksymalnych pozostają
względnie pierwsze. Jeśli $a\in\mathfrak m_i$,
$b\in\mathfrak m_j$ i $a+b=1$, to w rozwinięciu
$(a+b)^{2N-1}=1$ każdy składnik zawiera albo
$a^N$, albo $b^N$. Stąd
$1\in\mathfrak m_i^N+\mathfrak m_j^N$.
Chińskie twierdzenie o resztach daje więc
\[
 R/\bigcap_i\mathfrak m_i^N
 \cong\prod_{i=1}^r R/\mathfrak m_i^N.
\]
Ponieważ ideały po lewej są parami względnie pierwsze,
przecięcie jest ich iloczynem, czyli $(0)$.
Otrzymujemy żądany izomorfizm.

Każdy czynnik jest artinowski jako iloraz $R$.
Aby pokazać lokalność czynnika dla ustalonego $i$,
rozważmy jego dowolny ideał maksymalny. Jest obrazem
pewnego ideału maksymalnego $\mathfrak m_j$ pierścienia
$R$ zawierającego $\mathfrak m_i^N$.
Ponieważ $\mathfrak m_j$ jest pierwszy i
$\mathfrak m_i^N\subseteq\mathfrak m_j$, dostajemy
$\mathfrak m_i\subseteq\mathfrak m_j$:
dla $a\in\mathfrak m_i$ mamy $a^N\in\mathfrak m_j$,
a więc $a\in\mathfrak m_j$. Maksymalność daje
$\mathfrak m_i=\mathfrak m_j$.
Tylko ten ideał maksymalny występuje w ilorazie,
co dowodzi lokalności.
\end{proof}

\begin{przyklad}[Rozkład pierścienia $\mathbb Z/12\mathbb Z$]
\label{aa:prz:artin-z12}
Pierścień $\mathbb Z/12\mathbb Z$ jest artinowski,
gdyż ma skończenie wiele elementów. Jego ideały
odpowiadają ideałom $\mathbb Z$ zawierającym $(12)$,
czyli ideałom $(d)$ dla dodatnich dzielników $d\mid12$.
Redukcja modulo $d$ wyznacza surjektywny homomorfizm
$\mathbb Z/12\mathbb Z\to\mathbb Z/d\mathbb Z$
o jądrze $(\overline d)$. Stąd iloraz przez
$(\overline d)$ jest izomorficzny z
$\mathbb Z/d\mathbb Z$, a więc jest ciałem dokładnie
wtedy, gdy $d$ jest liczbą pierwszą. Stąd mamy dokładnie
dwa ideały maksymalne: $(\overline2)$ i $(\overline3)$.
Ich przecięcie to ideał $(\overline6)$, składający
się z klas $\overline0,\overline6$.
Jest to radykał Jacobsona, a jego kwadrat wynosi
$(\overline{36})=(\overline0)$. Chińskie twierdzenie o resztach daje izomorfizm
$\mathbb Z/12\mathbb Z\cong
\mathbb Z/4\mathbb Z\times\mathbb Z/3\mathbb Z$,
$\overline a\mapsto(a\bmod4,a\bmod3)$.
Jego odwrotność wysyła parę $(u,v)$ na klasę
$9u+4v$ modulo $12$, ponieważ $9\equiv1\pmod4$,
$9\equiv0\pmod3$, $4\equiv0\pmod4$ i
$4\equiv1\pmod3$.
W opisie z Twierdzenia~\ref{aa:tw:artin-local-product}
możemy przyjąć $N=2$: dla ideałów maksymalnych
$\mathfrak m_2=(\overline2)$ i $\mathfrak m_3=(\overline3)$
mamy $\mathfrak m_2^2=(\overline4)$ oraz
$\mathfrak m_3^2=(\overline9)=(\overline3)$.
Dlatego czynniki $R/\mathfrak m_2^2$ i
$R/\mathfrak m_3^2$ są izomorficzne odpowiednio
z $\mathbb Z/4\mathbb Z$ i $\mathbb Z/3\mathbb Z$.
Pierwszy czynnik, $\mathbb Z/4\mathbb Z$, jest lokalny
z ideałem maksymalnym generowanym przez klasę $2$
modulo $4$, a drugi jest ciałem.
Ten przykład pokazuje, jak nilpotentność radykału
współistnieje z rozkładem pierścienia na lokalne czynniki.
\end{przyklad}

\section{Dlaczego pierścień artinowski jest noetherowski?}
\begin{lemat}[Moduły artinowskie nad skończonym iloczynem ciał]
\label{aa:lem:artin-mod-over-fields}
Niech $r\ge1$ będzie liczbą całkowitą, a
$A=K_1\times\cdots\times K_r$ skończonym iloczynem ciał i niech $M$ będzie $A$-modułem
spełniającym warunek stabilizacji zstępujących
ciągów podmodułów. Wtedy $M$ jest skończenie generowany
jako $A$-moduł.
\end{lemat}
\begin{proof}
Niech $e_i\in A$ będzie elementem równym $1$
na $i$-tej współrzędnej i $0$ na pozostałych.
Wtedy $e_i^2=e_i$, $e_ie_j=0$ dla $i\ne j$
oraz $\sum_i e_i=1$. Każdy $m\in M$ ma więc
jednoznaczny rozkład $m=\sum_i e_im$, a
$M=\bigoplus_i e_iM$.
Moduł $e_iM$ jest przestrzenią wektorową nad $K_i$:
działanie $A$ na nim zależy wyłącznie od $i$-tej
współrzędnej skalara. Każda jego podprzestrzeń $W\subseteq e_iM$ jest też
$A$-podmodułem $M$: dla $a=(a_1,\ldots,a_r)\in A$
i $w\in W$ zachodzi $aw=a_iw\in W$. Dlatego
każdy zstępujący ciąg takich podprzestrzeni stabilizuje się.
Gdyby $e_iM$ było nieskończenie wymiarowe,
można byłoby wybrać liniowo niezależne wektory
$v_1,v_2,\ldots$. Podprzestrzenie
$\operatorname{span}\{v_n,v_{n+1},\ldots\}$
tworzyłyby nieskończony ściśle zstępujący ciąg,
co jest sprzecznością. Zatem każdy $e_iM$ ma
skończony wymiar nad $K_i$, a złączenie skończonych
baz tych przestrzeni generuje $M$ jako $A$-moduł.
\end{proof}

\begin{twierdzenie}[Pierścień artinowski jest noetherowski]
\label{aa:tw:artin-implies-noether}
Każdy przemienny pierścień artinowski z jedynką
jest noetherowski.
\end{twierdzenie}
\begin{proof}
Przypadek pierścienia zerowego jest natychmiastowy;
załóżmy więc, że $R\ne0$. Niech $J=J(R)$.
Na mocy Twierdzenia~\ref{aa:tw:artin-jacobson-nilpotent}
przy pewnym $N$ zachodzi $J^N=0$.
Z Wniosku~\ref{aa:wn:artin-semisimple-quotient}
wiemy, że $A=R/J$ jest skończonym iloczynem ciał.

Ustalmy dowolny ideał $I\subseteq R$.
Rozważmy skończoną filtrację
$ I=I\cap J^0\supseteq I\cap J\supseteq\cdots
\supseteq I\cap J^N=0$.
Dla każdego $k=0,\ldots,N-1$ moduł
\[
 M_k=(I\cap J^k)/(I\cap J^{k+1})
\]
można utożsamić z podmodułem $J^k/J^{k+1}$:
odwzorowanie $a+(I\cap J^{k+1})\mapsto a+J^{k+1}$
jest dobrze określone i ma jądro zerowe, gdyż
$a\in I\cap J^k$ oraz $a\in J^{k+1}$ oznaczają
$a\in I\cap J^{k+1}$.
Mnożenie przez element $J$ zeruje każdą klasę w $M_k$,
ponieważ $J(I\cap J^k)\subseteq I\cap J^{k+1}$.
Działanie $R$ przechodzi więc na $A=R/J$.

Moduł $M_k$ spełnia warunek stabilizacji zstępujących
ciągów $A$-podmodułów. Każdy taki podmoduł jest też
$R$-podmodułem, ponieważ działanie $R$ odbywa się
poprzez homomorfizm $R\to A$. Jego przeciwobraz przy
rzutowaniu $I\cap J^k\to M_k$ jest zatem ideałem
pierścienia $R$ zawierającym $I\cap J^{k+1}$.
Zstępujące ciągi tych przeciwobrazów stabilizują się,
bo $R$ jest artinowski. Surjektywność rzutowania
zapewnia wówczas stabilizację wyjściowego ciągu.
Z Lematu~\ref{aa:lem:artin-mod-over-fields} wynika,
że każdy $M_k$ ma skończenie wielu generatorów.

Wybierzmy dla każdego $k$ skończenie wiele reprezentantów
w $I\cap J^k$ generatorów $M_k$. Ich łączna rodzina
jest skończona i generuje $I$ jako ideał: dla danego
$a\in I$ wyrażamy jego klasę w $M_0$ jako kombinację
obrazów reprezentantów ze współczynnikami w $A$,
a współczynniki podnosimy do $R$. Po odjęciu tej
kombinacji pozostaje element $I\cap J$.
Powtarzając to samo dla jego klasy w $M_1$, otrzymujemy
element $I\cap J^2$, i tak dalej.
Po co najwyżej $N$ krokach reszta należy do
$I\cap J^N=0$. Wszystkie odejmowane składniki należą
do ideału generowanego przez wybrane reprezentanty.
Zatem każdy ideał $I$ jest skończenie generowany.
\end{proof}

\section{Lemat Nakayamy i minimalne układy generatorów}
\begin{twierdzenie}[Lemat Nakayamy dla pierścienia lokalnego]
\label{aa:tw:nakayama-local}
Niech $(R,\mathfrak m)$ będzie przemiennym pierścieniem
lokalnym z jedynką i jedynym ideałem maksymalnym
$\mathfrak m$. Niech $M$ będzie skończenie generowanym
modułem nad $R$. Jeżeli $\mathfrak m M=M$,
to $M=0$.
\end{twierdzenie}
\begin{proof}
Załóżmy, że $M\ne0$. Spośród wszystkich skończonych
układów generatorów $M$ wybierzmy taki, który ma
najmniejszą liczbę $r\ge1$ elementów,
powiedzmy $m_1,\ldots,m_r$.
Ponieważ $M=\mathfrak m M$, element $m_1$
można zapisać jako
$m_1=a_1m_1+\cdots+a_rm_r$ z $a_i\in\mathfrak m$.
Stąd
$(1-a_1)m_1=a_2m_2+\cdots+a_rm_r$.

Element $1-a_1$ nie należy do $\mathfrak m$:
gdyby należał, to $1=(1-a_1)+a_1\in\mathfrak m$.
Na mocy Twierdzenia~\ref{aa:tw:local-nonunits}
element $1-a_1$ jest więc jednostką.
Mnożąc ostatnią równość przez jego odwrotność,
otrzymujemy, że $m_1$ należy do podmodułu
generowanego przez $m_2,\ldots,m_r$.
Wtedy te pozostałe elementy generują całe $M$,
co przeczy minimalności $r$.
Wobec tego $M=0$.
\end{proof}

\begin{wniosek}[Wersja Nakayamy z podmodułem]
\label{aa:wn:nakayama-submodule}
Niech $(R,\mathfrak m)$ będzie niezerowym przemiennym
pierścieniem lokalnym z jedynką i jedynym ideałem
maksymalnym $\mathfrak m$, niech $M$ będzie skończenie
generowanym $R$-modułem, a
$N\subseteq M$ podmodułem takim, że
$M=N+\mathfrak m M$. Wtedy $M=N$.
\end{wniosek}
\begin{proof}
Moduł ilorazowy $M/N$ jest skończenie generowany,
bo obrazy dowolnych skończonych generatorów $M$
generują $M/N$. Założenie
$M=N+\mathfrak m M$ po przejściu do ilorazu daje
$M/N=\mathfrak m(M/N)$. Twierdzenie Nakayamy
zastosowane do $M/N$ prowadzi do $M/N=0$.
Jest to równoważne $M=N$.
\end{proof}

\begin{twierdzenie}[Minimalna liczba generatorów modułu]
\label{aa:tw:nakayama-minimal-generators}
Niech $(R,\mathfrak m)$ będzie przemiennym pierścieniem
lokalnym z jedynką, $K=R/\mathfrak m$ jego ciałem
reszt, a $M$ skończenie generowanym $R$-modułem.
Wtedy $M/\mathfrak m M$ jest skończenie wymiarową
przestrzenią wektorową nad $K$, a najmniejsza możliwa
liczba generatorów $M$ jest równa
$\dim_K(M/\mathfrak m M)$.
\end{twierdzenie}
\begin{proof}
Ponieważ $\mathfrak m$ działa zerowo na
$M/\mathfrak m M$, mnożenie przez $R$ zależy tylko
od klasy skalara modulo $\mathfrak m$.
W ten sposób iloraz jest przestrzenią wektorową nad $K$.
Obrazy skończonej rodziny generatorów $M$ rozpinają
ten iloraz, więc jego wymiar jest skończony, powiedzmy $d$.

Jeżeli $M$ ma $r$ generatorów, ich obrazy rozpinają
przestrzeń wymiaru $d$, zatem $r\ge d$.
Aby wykazać, że $d$ generatorów wystarczy,
wybierzmy bazę $\bar m_1,\ldots,\bar m_d$
przestrzeni $M/\mathfrak m M$ i jej reprezentantów
$m_i\in M$. Niech $N=(m_1,\ldots,m_d)$.
Każda klasa modulo $\mathfrak m M$ jest kombinacją
$\bar m_i$, więc $M=N+\mathfrak m M$.
Z Wniosku~\ref{aa:wn:nakayama-submodule} otrzymujemy
$M=N$. Zatem $d$ jest zarazem osiągalną i najmniejszą
liczbą generatorów.
\end{proof}

\begin{przyklad}[Różnica między bazą a generatorami modułu]
\label{aa:prz:nakayama-example}
Niech $K$ będzie ciałem,
$R=K[x]_{(x)}$ lokalizacją pierścienia wielomianów
w jego ideale maksymalnym $(x)$. Jest to pierścień
lokalny z ideałem maksymalnym $\mathfrak m=(x)R$.
Dla modułu wolnego $M=R^2$ mamy
$M/\mathfrak m M\cong K^2$.
Poprzednie twierdzenie mówi, że $M$ wymaga
co najmniej dwóch generatorów i standardowe dwa
wektory bazowe rzeczywiście wystarczają.

Z kolei dla $M=R/(x^3)R$ iloraz
$M/\mathfrak m M$ jest izomorficzny z $K$,
więc wystarczy jeden generator: klasa $1$.
Nie jest to jednak wolny moduł rangi jeden:
niezerowy element $x^3\in R$ anihiluje cały moduł
$M$, bo $x^3(q+(x^3)R)=0$ dla każdego $q\in R$.
Natomiast w wolnym module rangi jeden, izomorficznym
z $R$, żaden niezerowy element $r\in R$ nie anihiluje
całego modułu, ponieważ $r\cdot1=r\ne0$.
Należy odróżniać \emph{posiadanie
jednego generatora} od \emph{posiadania bazy złożonej
z jednego elementu}.
\end{przyklad}


\begin{twierdzenie}[Lemat Nakayamy dla radykału Jacobsona]
\label{aa:tw:nakayama-jacobson-general}
Niech $R$ będzie niezerowym przemiennym pierścieniem
z jedynką, $J(R)$ jego radykałem Jacobsona, a $M$ skończenie
generowanym $R$-modułem. Jeśli ideał
$I\subseteq J(R)$ spełnia $IM=M$, to $M=0$.
\end{twierdzenie}
\begin{proof}
Przypuśćmy, że $M\neq0$, i wybierzmy układ
generatorów $m_1,\ldots,m_r$ o minimalnej
liczbie elementów $r\ge1$. Ponieważ $IM=M$,
istnieją $a_1,\ldots,a_r\in I$ takie, że
$m_1=\sum_{i=1}^r a_im_i$.
Po przeniesieniu pierwszego składnika otrzymujemy
$(1-a_1)m_1=\sum_{i=2}^r a_im_i$.
Ponieważ $a_1\in I\subseteq J(R)$, z
Twierdzenia~\ref{aa:tw:jacobson-unit-criterion}
(dla $b=1$) wynika, że $1-a_1$ jest jednostką.
Dlatego $m_1$ jest kombinacją $m_2,\ldots,m_r$,
a cały moduł $M$ byłby generowany przez $r-1$
elementów. Jest to sprzeczne z minimalnością $r$.
Zatem $M=0$.
\end{proof}

\begin{przyklad}[Skończone generowanie jest niezbędne]
\label{aa:prz:nakayama-non-fg-counterexample}
Niech $K$ będzie ciałem,
$R=K[x]_{(x)}$ lokalnym pierścieniem funkcji
wymiernych regularnych w zerze,
a $F=K(x)$ jego ciałem ułamków.
Rozważmy $R$-moduł $M=F/R$.
Jest on niezerowy: klasa $1/x+R$ jest niezerowa,
ponieważ $x$ nie ma odwrotności w $R$.
Ideałem maksymalnym $R$ jest $\mathfrak m=xR$.
Dla każdego $q+R\in M$ mamy
$q+R=x(q/x+R)$, więc $\mathfrak m M=M$.
Mimo to $M\ne0$, co nie przeczy lematowi Nakayamy,
gdyż $M$ nie jest skończenie generowany.

Zweryfikujmy ostatnie stwierdzenie bez odwoływania
się do lematu. Każdy element $F$ jest postaci
$r/x^n$ dla pewnego $r\in R$ i $n\ge0$:
po wyłączeniu największej potęgi $x$ z mianownika
pozostały mianownik jest jednostką $R$.
Gdyby klasy $q_1+R,\ldots,q_s+R$ generowały $M$,
można byłoby dobrać wspólne $N\ge0$ tak,
aby $x^Nq_i\in R$ dla wszystkich $i$.
Wówczas $x^NM=0$, ponieważ mnożenie przez $x^N$
zerowałoby każdy generator i wszystkie jego
$R$-kombinacje. Jednak
$x^N(1/x^{N+1}+R)=1/x+R\ne0$.
Otrzymana sprzeczność dowodzi, że $M$ nie ma
skończonego układu generatorów.
\end{przyklad}

\begin{przyklad}[Dwuwymiarowa przestrzeń generatorów ideału]
\label{aa:prz:nakayama-two-generators-square-zero}
Niech $K$ będzie ciałem i
$R=K[x,y]/(x,y)^2$; oznaczmy klasy $x,y$
przez $\bar x,\bar y$. Każda klasa $R$ ma
jednoznaczną postać $a+b\bar x+c\bar y$
z $a,b,c\in K$, ponieważ jednomiany łącznego stopnia
co najmniej dwa znikają w ilorazie.
Element o $a\ne0$ jest jednostką: jeżeli
$u=b\bar x+c\bar y$, to $u^2=0$ i
$(a+u)^{-1}=a^{-1}-a^{-2}u$.
Elementy z $a=0$ tworzą ideał
$\mathfrak m=(\bar x,\bar y)$, a
$\mathfrak m^2=0$. Są więc nilpotentne i nieodwracalne;
razem z wcześniejszą obserwacją oznacza to, że zbiór
elementów nieodwracalnych jest dokładnie ideałem
$\mathfrak m$. Z Twierdzenia~\ref{aa:tw:local-nonunits}
wynika, że $R$ jest lokalny, a $\mathfrak m$ jest
jego jedynym ideałem maksymalnym.

Traktujmy teraz $M=\mathfrak m$ jako $R$-moduł.
Wówczas $\mathfrak m M=\mathfrak m^2=0$, stąd
$M/\mathfrak m M=\mathfrak m$.
Jako przestrzeń wektorowa nad
$R/\mathfrak m\cong K$ ma ona bazę
$\bar x,\bar y$ i wymiar dwa.
Twierdzenie~\ref{aa:tw:nakayama-minimal-generators}
wykazuje, że ideał $\mathfrak m$ wymaga
\emph{dokładnie dwóch} generatorów.
Jeden element $a\bar x+b\bar y$ nie może go
generować, ponieważ jego wielokrotności przez $R$
tworzą najwyżej jednowymiarową podprzestrzeń
$K\bar x+K\bar y$: składniki z $\mathfrak m$
zerują $\mathfrak m$.
\end{przyklad}

\begin{uwaga}[Przejście do szerszej teorii modułów]
\label{aa:uw:artin-to-modules}
Lemat Nakayamy jest pierwszym przykładem wyniku,
w którym struktura pierścienia lokalnego kontroluje
minimalną liczbę generatorów modułu.
W dalszych rozdziałach rozszerzymy tę teorię o
moduły wolne, prezentacje, rozszczepienia, iloczyn
tensorowy modułów, postać Smitha i klasyfikację modułów
nad dziedzinami ideałów głównych. Nie utożsamiamy
iloczynu tensorowego z iloczynem prostym.
\end{uwaga}

\setlength{\parskip}{4pt}

% AA-15 | Algebra abstrakcyjna | Iteracja VIII
\chapter{Pierścienie nieprzemienne i algebry grupowe}
\label{aa:ch:noncommutative-rings}

\section{Dlaczego nie można zakładać przemienności?}
Dotychczas główne twierdzenia dotyczące ideałów, podzielności,
lokalizacji oraz pierścieni noetherowskich sformułowaliśmy dla
pierścieni przemiennych. Istnieją jednak naturalne pierścienie, w których
kolejność mnożenia ma znaczenie. Już dla przestrzeni wektorowej
$V$ nad ciałem $K$ pierścień $\operatorname{End}_K(V)$ z dodawaniem
odwzorowań i mnożeniem przez składanie jest zwykle nieprzemienny.
Nawet jeśli $K$ jest ciałem, macierze mogą nie komutować.

W tym rozdziale wszystkie pierścienie mają jedynkę $1_R\ne0_R$,
a homomorfizmy pierścieni zachowują jedynkę. Nie zakładamy
przemienności mnożenia, chyba że wyraźnie ją zaznaczamy.

\begin{definicja}[Pierścień niekoniecznie przemienny]
\label{aa:def:nc-ring}
Pierścień z jedynką to zbiór $R$ z działaniami $+$ i $\cdot$,
zerem $0_R$ i jedynką $1_R$ taki, że $(R,+,0_R)$ jest grupą
abelową, mnożenie jest łączne, $1_Ra=a1_R=a$ dla każdego
$a\in R$ i zachodzą obie rozdzielności:
$ a(b+c)=ab+ac$ oraz $(a+b)c=ac+bc$ dla wszystkich
$a,b,c\in R$. Pierścień jest przemienny, jeśli ponadto
$ab=ba$ dla wszystkich $a,b\in R$.
\end{definicja}

\begin{przyklad}[Macierze nie komutują]
\label{aa:prz:nc-matrix}
Niech $K$ będzie ciałem, a $E_{ij}$ oznacza macierz $2\times2$
mającą $1$ na pozycji $(i,j)$ i zera poza nią. Wówczas
$E_{12}E_{21}=E_{11}$, lecz $E_{21}E_{12}=E_{22}$.
Dlatego $M_2(K)$ nie jest pierścieniem przemiennym.
Jednocześnie obie macierze $E_{12}$ i $E_{21}$ są niezerowe,
a $E_{12}^2=E_{21}^2=0$. Przykład wyjaśnia, dlaczego
argumenty wykorzystujące skracanie czynników w dziedzinie
całkowitości nie przenoszą się bez dodatkowych założeń
na dowolne pierścienie macierzy.
\end{przyklad}

\begin{definicja}[Centrum i komutator elementów pierścienia]
\label{aa:def:nc-center}
Dla pierścienia $R$ definiujemy jego \emph{centrum}
$Z(R)=\{z\in R:\ za=az\text{ dla każdego }a\in R\}$.
Dla $a,b\in R$ ich \emph{komutatorem addytywnym}
nazywamy $[a,b]=ab-ba$. Nie należy utożsamiać tego
komutatora z komutatorem multiplikatywnym elementów grupy.
\end{definicja}

\begin{stwierdzenie}[Centrum jest podpierścieniem przemiennym]
\label{aa:st:nc-center-subring}
Niech $R$ będzie pierścieniem z jedynką. Wtedy $Z(R)$
jest przemiennym podpierścieniem zawierającym $1_R$.
\end{stwierdzenie}
\begin{proof}
Dla każdego $a\in R$ mamy $0_Ra=a0_R$ i $1_Ra=a1_R$.
Jeżeli $z,w\in Z(R)$ i $a\in R$, to
$(z-w)a=za-wa=az-aw=a(z-w)$ oraz
$(zw)a=z(wa)=z(aw)=(za)w=(az)w=a(zw)$.
Zatem $z-w,zw\in Z(R)$, co dowodzi zamkniętości
na działania i obecności jedynki. Dla $z,w\in Z(R)$
równość $zw=wz$ wynika z definicji centrum, ponieważ
$w$ także należy do $R$. Mnożenie w centrum jest
więc przemienne.
\end{proof}

\section{Ideały lewostronne, prawostronne i dwustronne}
\begin{definicja}[Trzy rodzaje ideałów]
\label{aa:def:nc-ideals}
Niech $R$ będzie pierścieniem. Addytywną podgrupę
$I\subseteq R$ nazywamy \emph{ideałem lewostronnym}, jeśli
$ri\in I$ dla każdego $r\in R$ i $i\in I$;
\emph{ideałem prawostronnym}, jeśli $ir\in I$;
\emph{ideałem dwustronnym}, jeśli zachodzą obie własności.
Dla $a\in R$ oznaczamy przez $Ra$, $aR$ i $RaR$
odpowiednio najmniejszy ideał lewostronny, prawostronny
i dwustronny zawierający $a$. Dokładniej,
$Ra=\{ra:r\in R\}$, $aR=\{ar:r\in R\}$ oraz
$RaR=\{\sum_{j=1}^n r_ja s_j:n\ge0,\ r_j,s_j\in R\}$,
gdzie dla $n=0$ suma ma wartość $0$.
Zbiory $Ra$ i $aR$ są addytywnymi podgrupami,
odpowiednio zamkniętymi na mnożenie z lewej lub prawej,
a dzięki jedynce zawierają $a$. Każdy odpowiedni ideał
zawierający $a$ zawiera także te zbiory.
Analogicznie $RaR$ jest grupą addytywną, zamkniętą
na mnożenie z lewej i prawej dzięki łączności mnożenia.
Zawiera $a=1_Ra1_R$, a każdy ideał dwustronny
zawierający $a$ zawiera wszystkie wyrazy $ras$ i ich
skończone sumy. To dowodzi minimalności $RaR$.
\end{definicja}

\begin{przyklad}[Ideał jednostronny nie musi być dwustronny]
\label{aa:prz:nc-one-sided}
Niech $R=M_2(K)$ i niech $I$ składa się z macierzy
mających drugą kolumnę równą zeru. Jest to ideał
lewostronny: dla $A\in R$ oraz $B\in I$ druga
kolumna $AB$ jest równa $A\cdot0=0$.
Nie jest jednak ideałem prawostronnym, ponieważ
$E_{11}\in I$, a $E_{11}E_{12}=E_{12}\notin I$.
Podobnie macierze o drugim wierszu równym zeru
tworzą ideał prawostronny: mnożenie z prawej
zachowuje zerowy drugi wiersz. Nie jest to ideał
lewostronny, gdyż $E_{11}$ do niego należy,
a $E_{21}E_{11}=E_{21}$ już nie.
Ograniczenie do ideałów dwustronnych nie jest zatem
jedynie zmianą nazwy: tylko one pozwalają utworzyć
pierścień ilorazowy.
\end{przyklad}

\begin{twierdzenie}[Kryterium istnienia pierścienia ilorazowego]
\label{aa:tw:nc-quotient}
Niech $R$ będzie pierścieniem z jedynką i $1_R\ne0_R$,
a $I$ jego addytywną podgrupą. Wzór
$(a+I)(b+I)=ab+I$ określa poprawnie mnożenie
klas w $R/I$ wtedy i tylko wtedy, gdy $I$ jest
ideałem dwustronnym. Wówczas $R/I$ ma łączne
mnożenie z jedynką $1_R+I$, a naturalny rzut
$R\to R/I$ zachowuje jedynkę. Przy przyjętej
w tym rozdziale konwencji $1\ne0$ iloraz jest
pierścieniem w tej samej klasie obiektów dokładnie wtedy,
gdy $I\ne R$. Dla $I=R$ otrzymujemy jednoelementowy
\emph{pierścień zerowy}, w którym $1=0$; jeśli
pierścienie zerowe są dopuszczone, także ten iloraz
jest pierścieniem.
\end{twierdzenie}
\begin{proof}
Załóżmy najpierw dwustronność $I$. Jeśli
$a'=a+i$ oraz $b'=b+j$ dla $i,j\in I$, to
$a'b'-ab=aj+ib+ij\in I$. Wynik mnożenia nie zależy
więc od reprezentantów. Łączność, rozdzielność oraz
jedynka $1_R+I$ przechodzą z $R$ na iloraz.
Przy tym $1_R+I=0_R+I$ jest równoważne
$1_R\in I$. Dla ideału dwustronnego zawierającego
jedynkę mamy $r=r1_R\in I$ dla każdego $r\in R$,
a więc $I=R$. To dowodzi także ostatniego
zastrzeżenia w tezie.

W przeciwną stronę załóżmy, że mnożenie klas
jest dobrze określone. Dla każdego $i\in I$ mamy
$i+I=0+I$. Jeśli $r\in R$, możemy pomnożyć tę
równość z lewej i z prawej strony przez $r+I$.
Dostajemy $ri+I=I$ i $ir+I=I$, a więc $ri,ir\in I$.
Zatem $I$ jest ideałem dwustronnym. Wzory na rzut
zachowują dodawanie, mnożenie i jedynkę z definicji.
\end{proof}

\begin{twierdzenie}[Pierwsze twierdzenie o izomorfizmie pierścieni]
\label{aa:tw:nc-first-isomorphism}
Niech $R,S$ będą pierścieniami z jedynką, a
$f:R\to S$ homomorfizmem zachowującym jedynkę.
Wówczas $\ker f$ jest ideałem dwustronnym i istnieje
kanoniczny izomorfizm $R/\ker f\cong f(R)$.
\end{twierdzenie}
\begin{proof}
Dla $x\in\ker f$ oraz $r\in R$ otrzymujemy
$f(rx)=f(r)f(x)=0$ i $f(xr)=f(x)f(r)=0$.
Jądro jest podgrupą addytywną, więc jest ideałem
dwustronnym. Określamy $\bar f(r+\ker f)=f(r)$.
Jeśli $r-r'\in\ker f$, to $f(r)=f(r')$, zatem
$\bar f$ jest poprawnie określonym homomorfizmem.
Obraz $f(R)$ jest podpierścieniem $S$: jest zamknięty
na różnice i iloczyny, a z warunku $f(1_R)=1_S$
wynika, że zawiera jedynkę $S$.
Odwzorowanie $\bar f$ jest surjektywne na ten podpierścień.
Ponadto $\bar f(r+\ker f)=0$ wtedy i tylko wtedy,
gdy $r\in\ker f$, czyli gdy $r+\ker f$ jest klasą
zerową. Stąd $\bar f$ jest iniektywne, a więc jest
izomorfizmem pierścieni z jedynką.
\end{proof}

\section{Pierścienie macierzy i ich ideały}
\begin{lemat}[Mnożenie jednostek macierzowych]
\label{aa:lem:nc-matrix-units}
Dla pierścienia $R$ z jedynką i jednostek macierzowych
$E_{ij}\in M_n(R)$, $n\ge1$, zachodzi
$E_{ij}E_{k\ell}=\delta_{jk}E_{i\ell}$, gdzie
$\delta_{jk}=1_R$, gdy $j=k$, oraz $\delta_{jk}=0_R$,
gdy $j\ne k$; macierz $E_{ij}$ ma $1_R$ w pozycji
$(i,j)$ i zera na pozostałych pozycjach.
\end{lemat}
\begin{proof}
Współczynnik na pozycji $(p,q)$ iloczynu wynosi
$\sum_{t=1}^n(E_{ij})_{pt}(E_{k\ell})_{tq}$.
Niezerowy składnik może wystąpić wyłącznie przy
$p=i$, $t=j=k$ i $q=\ell$; ma wtedy wartość $1_R$.
W pozostałych przypadkach suma jest zerowa.
\end{proof}

\begin{twierdzenie}[Dwustronne ideały pełnego pierścienia macierzy]
\label{aa:tw:nc-matrix-ideals}
Niech $R$ będzie pierścieniem z jedynką, a $n\ge1$.
Dla każdego ideału dwustronnego $J\triangleleft M_n(R)$
istnieje dokładnie jeden ideał dwustronny $I\triangleleft R$
taki, że $J=M_n(I)$, gdzie $M_n(I)$ oznacza zbiór
macierzy, których wszystkie współczynniki należą do $I$.
\end{twierdzenie}
\begin{proof}
Najpierw, jeśli $I$ jest ideałem dwustronnym $R$,
to $M_n(I)$ jest podgrupą addytywną. Dla
$A\in M_n(R)$ i $B\in M_n(I)$ współczynniki
$(AB)_{ij}=\sum_k a_{ik}b_{kj}$ oraz
$(BA)_{ij}=\sum_k b_{ik}a_{kj}$ należą do $I$.
Stąd $M_n(I)$ jest ideałem dwustronnym.

Niech teraz $J\triangleleft M_n(R)$. Niech $I$ będzie
zbiorem tych $a\in R$, dla których macierz $aE_{11}$
należy do $J$. Jest to podgrupa addytywna. Dla
$r\in R$ i $a\in I$ iloczyny
$(rE_{11})(aE_{11})=(ra)E_{11}$ oraz
$(aE_{11})(rE_{11})=(ar)E_{11}$ należą do $J$.
Zatem $I$ jest ideałem dwustronnym $R$.

Jeśli $B=(b_{ij})\in J$, to
$E_{1i}BE_{j1}=b_{ij}E_{11}\in J$: mnożenie przez
jednostki macierzowe wybiera współczynnik $b_{ij}$,
przy zachowaniu jego pozycji w iloczynie, również
gdy $R$ jest nieprzemienny. Zatem
$b_{ij}\in I$ dla wszystkich $i,j$. Dostaliśmy
$J\subseteq M_n(I)$. Odwrotnie, dla $a\in I$ mamy
$(aE_{11})\in J$, a zatem
$E_{i1}(aE_{11})E_{1j}=aE_{ij}\in J$. Każda
macierz z $M_n(I)$ jest sumą takich macierzy,
więc $M_n(I)\subseteq J$. Jednoznaczność $I$
wynika z warunku $aE_{11}\in J$.
\end{proof}

\begin{wniosek}[Prostota pierścienia macierzy nad ciałem]
\label{aa:wn:nc-matrix-simple}
Dla ciała $K$ i $n\ge1$ pierścień $M_n(K)$
nie posiada ideałów dwustronnych innych niż $(0)$
i cały pierścień, czyli jest \emph{pierścieniem prostym}.
Nie oznacza to, że jego ideały jednostronne są trywialne.
\end{wniosek}
\begin{proof}
Ciało $K$ ma dokładnie dwa ideały: $(0)$ i $K$.
Z Twierdzenia~\ref{aa:tw:nc-matrix-ideals} wynika,
że ideały dwustronne $M_n(K)$ mają postać
$M_n(0)$ lub $M_n(K)$. Dla $n\ge2$ przykład
z macierzami o jednej zerowej kolumnie daje
nietrywialny ideał lewostronny.
\end{proof}

\begin{twierdzenie}[Centrum pierścienia macierzy]
\label{aa:tw:nc-matrix-center}
Dla dowolnego pierścienia $R$ z jedynką i $n\ge1$
zachodzi $Z(M_n(R))=\{zI_n:z\in Z(R)\}$.
\end{twierdzenie}
\begin{proof}
Jeżeli $z\in Z(R)$, to dla każdej macierzy
$A=(a_{ij})$ zachodzi $(zI_n)A=(za_{ij})$ oraz
$A(zI_n)=(a_{ij}z)$, więc $zI_n$ należy do centrum.

Niech odwrotnie $B=(b_{ij})$ komutuje ze wszystkimi
macierzami. Komutując w szczególności z $E_{kk}$,
otrzymujemy $BE_{kk}=E_{kk}B$. Dla indeksów
$i\ne k$ i $j=k$ z tej równości wynika $b_{ik}=0$,
a dla $i=k$, $j\ne k$ wynika $b_{kj}=0$.
Ponieważ $k$ jest dowolne, macierz $B$ jest diagonalna.
Z równości $BE_{ij}=E_{ij}B$ dla $i\ne j$ wynika,
że wszystkie elementy diagonalne są jednakowe:
$b_{ii}=b_{jj}=z$. W przypadku $n=1$ jest to
prawdziwe bez dalszego argumentu. Wreszcie $B$
komutuje z macierzami $aE_{11}$ dla każdego $a\in R$,
co daje $za=az$. Tak więc $z\in Z(R)$ i $B=zI_n$.
\end{proof}


\begin{przyklad}[Ideały dwustronne a ideały jednostronne w macierzach]
\label{aa:prz:nc-matrix-one-sided-contrast}
Niech $K$ będzie ciałem i $R=M_2(K)$. Podzbiór
$I=\{\left(\begin{smallmatrix}a&0\\b&0\end{smallmatrix}\right):a,b\in K\}$
jest niezerowym właściwym ideałem lewostronnym. Istotnie,
$E_{11}\in I$, lecz $E_{22}\notin I$, a dla każdej macierzy
$A\in R$ i $B\in I$ druga kolumna $AB$ jest zerowa.
Nie jest to ideał dwustronny, gdyż $E_{11}E_{12}=E_{12}\notin I$.
Z drugiej strony Twierdzenie~\ref{aa:tw:nc-matrix-ideals} pokazuje,
że $R$ nie ma właściwych niezerowych ideałów dwustronnych.
Nie można więc zastąpić słowa „dwustronny” słowem „lewostronny”
w twierdzeniu o prostocie pierścienia macierzy.
\end{przyklad}

\section{Pierścień przeciwny i moduły prawostronne}
\begin{definicja}[Pierścień przeciwny]
\label{aa:def:nc-opposite}
Dla pierścienia $R$ definiujemy pierścień przeciwny
$R^{\mathrm{op}}$ o tej samej grupie addytywnej i
mnożeniu $a\star b=ba$, gdzie po prawej mnożymy w $R$.
Jedynka pozostaje ta sama. Łączność jest zachowana,
ponieważ $(a\star b)\star c=c(ba)=(cb)a=
a\star(b\star c)$; obie rozdzielności wynikają
z rozdzielności w $R$ po odwróceniu kolejności
czynników.
\end{definicja}
\begin{stwierdzenie}[Zmiana strony działania]
\label{aa:st:nc-opposite-module}
Prawy $R$-moduł $M$ jest równoważnie lewym
$R^{\mathrm{op}}$-modułem poprzez działanie
$a\star m:=ma$. Ta zgodność działa także dla
homomorfizmów modułów.
\end{stwierdzenie}
\begin{proof}
Prawy moduł spełnia $(ma)b=m(ab)$.
W $R^{\mathrm{op}}$ mamy $a\star b=ba$,
a po stronie modułu
$(a\star b)\star m=m(ba)$.
Z drugiej strony
$a\star(b\star m)=(mb)a=m(ba)$, więc
łączność działania jest zachowana. Rozdzielności
oraz działanie jedynki przechodzą bezpośrednio
z aksjomatów prawego modułu. Homomorfizm, który
zachowuje prawostronne mnożenie przez skalary,
zachowuje z definicji odpowiadające mu lewostronne.
\end{proof}

\section{Pierścienie grupowe i odwzorowania uniwersalne}
\begin{definicja}[Pierścień grupowy]
\label{aa:def:nc-group-ring}
Niech $R$ będzie przemiennym pierścieniem z jedynką,
a $G$ dowolną grupą. \emph{Pierścień grupowy} $R[G]$
to zbiór wszystkich formalnych skończonych sum
$\sum_{g\in G}r_g[g]$, gdzie $r_g\in R$, a niemal
wszystkie współczynniki są zerowe. Dodawanie jest
współrzędne, mnożenie określa się przez
\[
\left(\sum_g r_g[g]\right)\left(\sum_h s_h[h]\right)
 =\sum_{k\in G}\left(\sum_{gh=k}r_gs_h\right)[k].
\]
Suma wewnętrzna jest skończona, ponieważ obie rodziny
współczynników mają skończony nośnik. Jedynką jest
$1_R[1_G]$. W przypadku $R=K$, gdzie $K$ jest
ciałem, pierścień $K[G]$ nazywa się również
\emph{algebrą grupową} grupy $G$ nad $K$.
\end{definicja}
\begin{stwierdzenie}[Aksjomaty pierścienia grupowego]
\label{aa:st:nc-group-ring-ring}
Dla przemiennego pierścienia $R$ z jedynką i dowolnej
grupy $G$ powyższe działania czynią $R[G]$ pierścieniem.
Ponadto $[g][h]=[gh]$ oraz $[g]^{-1}=[g^{-1}]$.
Jeżeli $R\ne0$, to $R[G]$ jest przemienny wtedy
i tylko wtedy, gdy $G$ jest abelowa.
\end{stwierdzenie}
\begin{proof}
Skończone nośniki gwarantują, że każdy współczynnik
iloczynu jest sumą skończoną. Łączność wynika z
obliczenia współczynnika przy $[k]$ dla trzech sum:
zarówno $(xy)z$, jak i $x(yz)$ mają współczynnik
$\sum_{abc=k}r_as_bt_c$. Jest to suma skończona,
bo indeksy $a,b,c$ należą do skończonych nośników
trzech czynników. Zmiana nawiasowania jest
więc poprawna dzięki łączności mnożenia w $R$
i w $G$. Rozdzielności sprawdzamy współczynnikowo,
a element $1_R[1_G]$ zachowuje wszystkie współczynniki.
To dowodzi aksjomatów pierścienia; tożsamości
$[g][h]=[gh]$ i $[g][g^{-1}]=[1_G]$ są bezpośrednie.
Jeśli $G$ jest abelowa i $R$ przemienny, wzór na
mnożenie pokazuje, że zamiana kolejności czynników
nie zmienia iloczynu. Odwrotnie, jeśli $R[G]$
jest przemienny, to $[gh]=[hg]$. Ponieważ
$R\ne0$, formalne symbole $[g]$ są różnymi
elementami wolnego modułu dla różnych $g$.
Zatem $gh=hg$, czyli $G$ jest abelowa.
\end{proof}

\begin{twierdzenie}[Własność uniwersalna pierścienia grupowego]
\label{aa:tw:nc-group-ring-universal}
Niech $R$ będzie przemiennym pierścieniem z jedynką,
a $G$ grupą. Niech $S$ będzie pierścieniem z jedynką,
$f:R\to S$ homomorfizmem pierścieni takim, że
$f(R)\subseteq Z(S)$, i niech $u:G\to S^\times$
będzie homomorfizmem grup do grupy jednostek $S$.
Wtedy istnieje dokładnie jeden homomorfizm pierścieni
$F:R[G]\to S$ spełniający $F(r[1_G])=f(r)$
oraz $F([g])=u(g)$ dla wszystkich $r\in R$, $g\in G$.
Jest nim odwzorowanie
$F(\sum_g r_g[g])=\sum_g f(r_g)u(g)$.
\end{twierdzenie}
\begin{proof}
Ponieważ zapisy w $R[G]$ mają jednoznaczne współczynniki,
podany wzór określa odwzorowanie. Zachowuje dodawanie,
bo sumowanie jest współrzędne. Dla dwóch składników
$r[g]$, $s[h]$ mamy
\[
F((r[g])(s[h]))=f(rs)u(gh)
=f(r)f(s)u(g)u(h).
\]
Warunek $f(R)\subseteq Z(S)$ pozwala przestawić
$f(s)$ z $u(g)$, więc prawa strona jest równa
$f(r)u(g)f(s)u(h)=F(r[g])F(s[h])$.
Rozdzielność i skończoność sum rozszerzają to na
dowolne elementy. Zachowanie jedynki daje
$F(1_R[1_G])=f(1_R)u(1_G)=1_S$.
Jeśli $H$ jest innym homomorfizmem o tych wartościach,
to z równości $r[g]=(r[1_G])[g]$ i addytywności
wynika $H(\sum r_g[g])=\sum f(r_g)u(g)$.
Zatem $H=F$.
\end{proof}

\begin{przyklad}[Grupa cykliczna i algebra wielomianów]
\label{aa:prz:nc-cyclic-ring}
Niech $n\ge1$ będzie liczbą całkowitą, $K$ ciałem,
a $C_n=\langle g\rangle$ grupą cykliczną rzędu
$n$ (czyli $g^n=1$ i nie istnieje mniejszy dodatni
wykładnik z tą własnością). Wówczas otrzymujemy izomorfizm
$K[C_n]\cong K[t]/(t^n-1)$. Homomorfizm
$K[t]\to K[C_n]$ dany przez $t\mapsto[g]$
jest surjektywny, gdyż każda potęga $[g]^j$
pochodzi od $t^j$. Ponieważ $[g]^n=[1]$,
wielomian $t^n-1$ należy do jądra. Indukowany
homomorfizm $K[t]/(t^n-1)\to K[C_n]$ jest więc
surjektywny. Po lewej każdy element ma jednoznacznego
reprezentanta stopnia mniejszego niż $n$, zatem
wymiar nad $K$ wynosi $n$. Po prawej bazą są
$[1],[g],\ldots,[g^{n-1}]$, więc wymiar także wynosi
$n$. Surjektywne odwzorowanie liniowe między
przestrzeniami tego samego skończonego wymiaru
jest izomorfizmem.
\end{przyklad}

\begin{przyklad}[Grupa nieabelowa daje pierścień nieprzemienny]
\label{aa:prz:nc-s3-ring}
Niech $K$ będzie dowolnym ciałem i $G=S_3$.
Elementy $[(12)]$ i $[(23)]$ są jednostkami pierścienia
$K[S_3]$, ale nie komutują, ponieważ
$(12)(23)\ne(23)(12)$ w $S_3$.
Jednocześnie $K[S_3]$ ma wymiar sześć nad $K$:
bazą jest sześć formalnych symboli $[g]$.
To prosty przykład sześciowymiarowej algebry
nieprzemiennej powstałej bez udziału macierzy.
\end{przyklad}


\begin{przyklad}[Nilpotenty w pierścieniu grupowym]
\label{aa:prz:nc-group-ring-nilpotent}
Niech $p$ będzie liczbą pierwszą, $K=\mathbb F_p$,
a $G=C_p=\langle g\rangle$. W pierścieniu $K[G]$
element $u=[g]-[1_G]$ jest niezerowy, ponieważ $g\ne1_G$
i formalne symbole $[g]$, $[1_G]$ są liniowo niezależne nad $K$.
Ponieważ $[g]$ i $[1_G]$ komutują, dwumian Newtona daje
$u^p=[g]^p-[1_G]^p=[g^p]-[1_G]=0$:
wszystkie pośrednie współczynniki dwumianowe są podzielne
przez $p$. Zatem nawet pierścień grupowy nad ciałem
może zawierać niezerowe elementy nilpotentne. Założenie
$\operatorname{char}K=p$ jest istotne: dla $K=\mathbb Q$
i $G=C_p$ w rzeczywistości $u^p\ne0$.
Istotnie, we wzorze dwumianowym współczynnik przy
$[g]$ wynosi $(-1)^{p-1}p$, ponieważ jedynymi
potęgami $g$ równymi $g$ wśród $g^0,\ldots,g^p$
są $g^1$. Współczynnik ten jest niezerowy w
$\mathbb Q$, więc $u^p\ne0$.
\end{przyklad}

\section{Wolne pierścienie nieprzemienne}
\begin{definicja}[Pierścień wolnych wielomianów]
\label{aa:def:nc-free-associative}
Niech $K$ będzie ciałem, a $X$ zbiorem. Pierścień
$K\langle X\rangle$ składa się ze skończonych
formalnych kombinacji $\sum_w a_ww$ słów $w$ w
alfabecie $X$, włącznie ze słowem pustym, będącym
jedynką. Iloczyn słów to ich konkatenacja,
a iloczyn sum określamy przez rozdzielność.
Skalary z $K$ komutują ze wszystkimi słowami.
\end{definicja}

\begin{twierdzenie}[Własność uniwersalna wolnego pierścienia]
\label{aa:tw:nc-free-universal}
Niech $K$ będzie ciałem, $X$ zbiorem, a $S$ pierścieniem
z jedynką wyposażonym w homomorfizm $f:K\to Z(S)$.
Dla każdej funkcji $u:X\to S$ istnieje dokładnie jeden
homomorfizm pierścieni $F:K\langle X\rangle\to S$
taki, że $F(a)=f(a)$ dla $a\in K$ oraz
$F(x)=u(x)$ dla $x\in X$.
\end{twierdzenie}
\begin{proof}
Wartość $F$ na słowie $x_1\cdots x_m$ jest wymuszona
przez zachowanie iloczynu i wynosi
$u(x_1)\cdots u(x_m)$; słowo puste musi przejść na
$1_S$. Rozszerzamy ten przepis $K$-liniowo:
$F(\sum_w a_ww)=\sum_w f(a_w)F(w)$.
Skończoność nośnika zapewnia poprawność definicji.
Dla słów $v,w$ zachodzi $F(vw)=F(v)F(w)$ z
łączności w $S$. Dla kombinacji liniowych własność
wynika z rozdzielności; warunek centralności $f(K)$
gwarantuje, że współczynniki można przesuwać
przez obrazy słów. Jednoznaczność wynika z tego,
że elementy $X$ wraz ze skalarami generują cały
pierścień $K\langle X\rangle$.
\end{proof}

\begin{przyklad}[Dlaczego wielomiany nieprzemienne są potrzebne?]
\label{aa:prz:nc-free-example}
Dla $X=\{x,y\}$ słowa $xy$ i $yx$ są różnymi
wektorami bazowymi $K\langle x,y\rangle$; dlatego
$xy-yx\ne0$. W pierścieniu przemiennym $K[x,y]$
ten sam formalny komutator znika. Kanoniczny
homomorfizm $K\langle x,y\rangle\to K[x,y]$
utożsamia dokładnie relacje wymuszające przemienność
$x$ i $y$: jego jądrem jest ideał dwustronny
$J=(xy-yx)$ generowany przez $xy-yx$.
Wykażemy obie inkluzje przez skonstruowanie
izomorfizmu. Odwzorowanie $\pi:K\langle x,y\rangle
\to K[x,y]$, przypisujące słowu przemienny
jednomian o tych samych wykładnikach $x,y$,
jest surjektywnym homomorfizmem, a $\pi(xy-yx)=0$.
Stąd $J\subseteq\ker\pi$ i dostajemy surjekcję
$\bar\pi:Q\to K[x,y]$, gdzie
$Q=K\langle x,y\rangle/J$.
W pierścieniu $Q$ klasy $\bar x,\bar y$ komutują.
Określmy $\psi:K[x,y]\to Q$ wzorem
$\psi(\sum_{a,b}c_{ab}x^ay^b)=
\sum_{a,b}c_{ab}\bar x^a\bar y^b$.
Wzór jest poprawnie określony dzięki jednoznacznemu
zapisowi wielomianu oraz wyznacza homomorfizm,
ponieważ $\bar x\bar y=\bar y\bar x$.
Ponadto $\bar\pi\circ\psi$ jest identycznością
na $K[x,y]$. Z kolei $\psi\circ\bar\pi$ jest
identycznością na $Q$: obie strony zgadzają się
na klasach $x,y$ i skalarach z $K$, które generują
cały $Q$ jako pierścień.
Można to sprawdzić także bezpośrednio na słowach:
w $Q$ każda zamiana sąsiednich liter $yx$ na $xy$
nie zmienia klasy słowa, ponieważ dla dowolnych słów
$v,w$ różnica $v(yx-xy)w$ należy do $J$.
Powtarzanie takich zamian ściśle zmniejsza liczbę par,
w których litera $y$ poprzedza literę $x$.
Po skończenie wielu krokach każde słowo o $a$ literach
$x$ i $b$ literach $y$ przechodzi więc w $x^ay^b$.
Zatem $\bar\pi$ jest izomorfizmem i $\ker\pi=J$,
co dowodzi tezy.
To pokazuje formalną różnicę między pierścieniem
przemiennym a nieprzemiennym.
\end{przyklad}

% AA-16 | Algebra abstrakcyjna | audyt II
\setlength{\parskip}{3pt}
\chapter{Moduły, homomorfizmy i ciągi dokładne}
\label{aa:ch:modules-introduction}

\section{Dlaczego po grupach i pierścieniach pojawiają się moduły?}
Przestrzeń wektorowa nad ciałem $K$ jest grupą abelową,
na której elementy $K$ działają jak skalary.
Zastępując ciało pierścieniem z jedynką, otrzymujemy
pojęcie modułu. Ta zmiana nie jest kosmetyczna: nawet
najprostsze wyniki teorii przestrzeni wektorowych,
takie jak istnienie bazy dla każdej przestrzeni, mogą
przestać obowiązywać. Z drugiej strony jeden język
obejmie grupy abelowe, przestrzenie wektorowe,
ideały pierścieni, przedstawienia grup oraz układy
równań liniowych o współczynnikach z pierścienia.

W tym rozdziale $R$ jest pierścieniem z jedynką,
niekoniecznie przemiennym. W definicjach lewego
modułu zapisujemy działanie skalara z lewej strony.

\begin{definicja}[Lewy i prawy moduł nad pierścieniem]
\label{aa:def:modules-full}
\emph{Lewym $R$-modułem} nazywamy grupę abelową
$(M,+,0_M)$ wraz z działaniem $R\times M\to M$,
$(r,m)\mapsto rm$, takim, że dla wszystkich
$r,s\in R$ i $m,n\in M$:
\begin{enumerate}[label=\textup{(M\arabic*)}]
\item $(r+s)m=rm+sm$;
\item $r(m+n)=rm+rn$;
\item $(rs)m=r(sm)$;
\item $1_Rm=m$.
\end{enumerate}
\emph{Prawym $R$-modułem} nazywamy grupę abelową $M$
z działaniem $M\times R\to M$, $(m,r)\mapsto mr$,
dla którego, dla dowolnych $m,n\in M$ i $r,s\in R$,
zachodzą
$(m+n)r=mr+nr$, $m(r+s)=mr+ms$,
$(mr)s=m(rs)$ oraz $m1_R=m$.
Jeśli $R$ jest przemienny, z każdego lewego
$R$-modułu otrzymujemy prawy, kładąc $mr=rm$:
warunek $(mr)s=m(rs)$ wynika wtedy z
$(sr)m=s(rm)$ i przemienności skalarów.
\end{definicja}

\begin{stwierdzenie}[Pierwsze konsekwencje aksjomatów]
\label{aa:st:modules-zero-laws}
Niech $M$ będzie lewym $R$-modułem. Dla każdego
$r\in R$ i $m\in M$ zachodzą
$0_Rm=0_M$, $r0_M=0_M$,
$(-r)m=-(rm)$ oraz $r(-m)=-(rm)$.
\end{stwierdzenie}
\begin{proof}
Z $(0_R+0_R)m=0_Rm+0_Rm$ otrzymujemy
$0_Rm=0_M$ po odjęciu $0_Rm$ w grupie addytywnej
$M$. Analogicznie
$r(0_M+0_M)=r0_M+r0_M$ daje $r0_M=0_M$.
Ponieważ $0_Rm=(r+(-r))m=rm+(-r)m$, element
$(-r)m$ jest przeciwny do $rm$. Ostatnia równość
wynika tak samo z $r(m+(-m))=r0_M=0_M$.
\end{proof}

\begin{przyklad}[Cztery zasadniczo różne moduły]
\label{aa:prz:modules-four}
\begin{enumerate}[label=\textup{\arabic*)}]
\item Każda grupa abelowa $A$ jest kanonicznie
$\mathbb Z$-modułem, ponieważ $na$ oznacza
$n$-krotną sumę elementu $a$. Homomorfizmy takich
modułów są dokładnie homomorfizmami grup abelowych.
\item Każda przestrzeń wektorowa nad ciałem $K$
jest $K$-modułem. Dla $K$ wszystkie moduły mają bazy,
co udowadnia standardowe twierdzenie o bazie
przestrzeni wektorowej.
\item Każdy ideał lewostronny $I\subseteq R$ jest
lewym $R$-modułem z mnożeniem odziedziczonym z $R$.
Przykładowo $2\mathbb Z$ jest $\mathbb Z$-modułem
izomorficznym z $\mathbb Z$, mimo że jako podzbiór
nie zawiera elementu $1$.
\item Jeśli $R=\mathbb Z/4\mathbb Z$, to
$M=2R=\{\bar0,\bar2\}$ jest $R$-modułem.
Nie jest wolnym $R$-modułem: dla dowolnego
niezerowego $m\in M$ zachodzi $\bar2m=0$,
więc żaden niezerowy element $M$ nie może stanowić
jednoelementowej bazy. Co więcej, każdy niezerowy
moduł wolny nad $R$ zawiera kopię $R$: jeśli $e$
należy do jego bazy, odwzorowanie $r\mapsto re$ jest
różnowartościowe dzięki jednoznaczności współrzędnych.
Taki moduł ma zatem co najmniej cztery elementy,
podczas gdy $M$ ma tylko dwa.
\end{enumerate}
\end{przyklad}

\begin{definicja}[Podmoduł i homomorfizm modułów]
\label{aa:def:modules-sub-hom}
Niech $M,N$ będą lewymi $R$-modułami.
\emph{Podmoduł} $L\subseteq M$ to podgrupa addytywna
zamknięta na działanie każdego $r\in R$.
\emph{Homomorfizm $R$-modułów} to odwzorowanie
$f:M\to N$ spełniające $f(m+m')=f(m)+f(m')$
i $f(rm)=rf(m)$ dla wszystkich $m,m'\in M$, $r\in R$.
Zbiór wszystkich takich homomorfizmów oznaczamy
$\operatorname{Hom}_R(M,N)$. W szczególności każdy homomorfizm
jest addytywny, więc $f(0_M)=0_N$ oraz $f(-m)=-f(m)$:
pierwsza równość wynika z $f(0_M)=f(0_M)+f(0_M)$,
a druga z $f(m)+f(-m)=f(0_M)$.
\end{definicja}

\begin{stwierdzenie}[Grupa homomorfizmów modułów]
\label{aa:st:modules-hom-abelian}
Dla lewych $R$-modułów $M,N$ zbiór
$\operatorname{Hom}_R(M,N)$ jest grupą abelową przy
punktowym dodawaniu $(f+g)(m)=f(m)+g(m)$.
Jeśli $R$ jest przemienny, staje się dodatkowo
$R$-modułem dla $(rf)(m)=r f(m)$.
\end{stwierdzenie}
\begin{proof}
Dla homomorfizmów $f,g$ ich suma zachowuje dodawanie,
ponieważ dla $m,m'\in M$ mamy
$(f+g)(m+m')=f(m)+f(m')+g(m)+g(m')$
$=(f+g)(m)+(f+g)(m')$; środkowa zamiana
składników jest dopuszczalna, gdyż $N$ jest grupą abelową.
Dla $r\in R$ oraz $m\in M$ mamy
$(f+g)(rm)=f(rm)+g(rm)=rf(m)+rg(m)
=r(f(m)+g(m))$, zatem $f+g$ jest $R$-liniowe.
Odwzorowanie zerowe i odwzorowanie $-f$ są także
$R$-liniowe. Łączność i przemienność dodawania
pochodzą z grupy addytywnej $N$.
Jeśli $R$ jest przemienny, to
$(af)(rm)=a f(rm)=a(rf(m))=r(af(m))$;
pozostałe aksjomaty modułu wynikają z aksjomatów $N$.
Przemienność jest tu istotna: dla nieprzemiennego
$R$ wzór $(af)(m)=af(m)$ może nie zachowywać
$R$-liniowości, ponieważ $ar$ nie musi być równe $ra$.
\end{proof}

\begin{lemat}[Kryterium podmodułu]
\label{aa:lem:modules-submodule-test}
Niepusty podzbiór $L\subseteq M$ jest podmodułem
lewego $R$-modułu $M$ wtedy i tylko wtedy, gdy
$\ell+\ell'\in L$ i $r\ell\in L$ dla
wszystkich $\ell,\ell'\in L$ i $r\in R$.
\end{lemat}
\begin{proof}
Konieczność wynika z definicji. Przy założonej
zamkniętości wybierzmy $\ell\in L$.
Ponieważ $0_R\ell=0_M$, do $L$ należy zero.
Dalej $-\ell=(-1_R)\ell$ należy do $L$.
Zamkniętość na dodawanie i elementy przeciwne
pokazuje, że $L$ jest podgrupą addytywną;
zamkniętość na działanie $R$ już założyliśmy.
\end{proof}

\begin{twierdzenie}[Jądro, obraz i pierwszy izomorfizm modułów]
\label{aa:tw:modules-first-isomorphism}
Dla homomorfizmu lewych $R$-modułów $f:M\to N$
zbiór $\ker f$ jest podmodułem $M$, a
$\operatorname{im}f$ podmodułem $N$. Wzór
$\bar f(m+\ker f)=f(m)$ wyznacza izomorfizm
$R$-modułów $M/\ker f\cong\operatorname{im}f$.
\end{twierdzenie}
\begin{proof}
Jeżeli $m,m'\in\ker f$ i $r\in R$, to
$f(m+m')=f(m)+f(m')=0$ oraz $f(rm)=rf(m)=0$.
Jądro zawiera $0_M$, gdyż $f(0_M)=0_N$.
Jeśli $n=f(m)$ i $n'=f(m')$ należą do obrazu,
to $n+n'=f(m+m')$ i $rn=f(rm)$ także do niego należą.
Obraz jest niepusty, ponieważ zawiera $f(0_M)=0_N$.
Lemat~\ref{aa:lem:modules-submodule-test} pozwala zatem
zastosować kryterium podmodułu zarówno do jądra,
jak i do obrazu.

Grupa addytywna $M/\ker f$ jest ilorazem przez
podgrupę; definiujemy działanie $r(m+\ker f)=rm+\ker f$.
Jeśli $m'=m+k$ z $k\in\ker f$, to
$rm'=rm+rk$ oraz $rk\in\ker f$, więc działanie
nie zależy od reprezentanta. Aksjomaty modułu
przechodzą z $M$ na iloraz. Ponieważ
$f(m)=f(m')$ wtedy i tylko wtedy, gdy
$m-m'\in\ker f$, odwzorowanie $\bar f$
jest dobrze określone i różnowartościowe.
Jest na obraz z definicji. Zachowuje dodawanie
i skalary, ponieważ $f$ je zachowuje.
\end{proof}

\begin{twierdzenie}[Drugi i trzeci izomorfizm modułów]
\label{aa:tw:modules-second-third}
Niech $M$ będzie lewym $R$-modułem.
\begin{enumerate}[label=\textup{(\roman*)}]
\item Dla podmodułów $A,B\subseteq M$ zachodzi
$(A+B)/B\cong A/(A\cap B)$.
\item Jeśli $A\subseteq B\subseteq M$ są podmodułami,
to $(M/A)/(B/A)\cong M/B$.
\end{enumerate}
\end{twierdzenie}
\begin{proof}
W (i) suma $A+B=\{a+b:a\in A,\ b\in B\}$
jest podmodułem $M$: suma dwóch takich elementów,
ich elementy przeciwne i wszystkie wielokrotności
przez $R$ znów należą do $A+B$.
Rozważmy homomorfizm $A\to(A+B)/B$,
$a\mapsto a+B$. Jest on na, ponieważ każda klasa
$a+b+B$ jest równa $a+B$. Jego jądro to dokładnie
$A\cap B$. Z pierwszego twierdzenia o izomorfizmie
otrzymujemy $(i)$. W (ii) homomorfizm
$M/A\to M/B$ dany przez $m+A\mapsto m+B$
jest dobrze określony, bo $A\subseteq B$, i jest
surjektywny. Jego jądro składa się dokładnie z
klas $b+A$, $b\in B$, czyli jest równe $B/A$.
Pierwsze twierdzenie o izomorfizmie daje (ii).
\end{proof}

\section{Moduły wolne i prezentacje}
\begin{definicja}[Moduł wolny]
\label{aa:def:modules-free}
Rodzina $(e_s)_{s\in S}$ elementów lewego $R$-modułu
$F$ jest \emph{bazą}, jeśli każdy $x\in F$ ma dokładnie
jedno przedstawienie $x=\sum_{s\in S}r_se_s$,
w którym tylko skończenie wiele współczynników
$r_s\in R$ jest niezerowych. Moduł mający bazę
nazywamy \emph{wolnym}. Symbol $R^{(S)}$ oznacza
zbiór funkcji $S\to R$ o skończonym nośniku,
z dodawaniem i lewym działaniem $R$ określonymi
współrzędnie. Dla $s\in S$ niech $e_s$ przyjmuje
wartość $1_R$ w punkcie $s$ i $0_R$ poza nim.
Wówczas każdy $x\in R^{(S)}$ ma jednoznaczny zapis
$x=\sum_{s\in S}x(s)e_s$; równość współrzędnych
w punkcie $s$ dowodzi jednoznaczności. W szczególności
$(e_s)_{s\in S}$ jest bazą, a dla $S=\varnothing$
otrzymujemy moduł zerowy.
\end{definicja}

\begin{twierdzenie}[Własność uniwersalna modułu wolnego]
\label{aa:tw:modules-free-universal}
Niech $S$ będzie zbiorem, a $M$ lewym $R$-modułem.
Dla każdej funkcji $u:S\to M$ istnieje dokładnie
jeden homomorfizm $R$-modułów
$\bar u:R^{(S)}\to M$, który wysyła standardowy
wektor bazowy $e_s$ na $u(s)$ dla każdego $s\in S$.
\end{twierdzenie}
\begin{proof}
Dla funkcji $x:S\to R$ o skończonym nośniku
określamy $\bar u(x)=\sum_{s\in S}x(s)u(s)$.
Suma zawiera tylko skończenie wiele niezerowych
składników, więc jest dobrze określona.
Jej addytywność wynika z rozdzielności działania,
a $\bar u(rx)=r\bar u(x)$ z aksjomatu
$(rs)m=r(sm)$: istotna jest tu kolejność
mnożenia skalarów w nieprzemiennym $R$.
Współrzędne $e_s$ dają $\bar u(e_s)=u(s)$.
Jeżeli $v$ jest innym takim homomorfizmem,
to dla $x=\sum_s x(s)e_s$ mamy
$v(x)=\sum_s x(s)v(e_s)=\bar u(x)$.
Stąd jednoznaczność.
\end{proof}

\begin{definicja}[Prezentacja modułu przez generatory i relacje]
\label{aa:def:modules-presentation}
\emph{Prezentacją} lewego $R$-modułu $M$ nazywamy
ciąg homomorfizmów
$R^{(T)}\xrightarrow{d}R^{(S)}\xrightarrow{p}M\to0$
taki, że $p$ jest surjekcją i $\operatorname{im}d=\ker p$.
Oznacza to, że elementy $p(e_s)$ generują $M$,
a elementy $d(e_t)$ są relacjami generującymi
wszystkie zależności między generatorami.
\end{definicja}

\begin{twierdzenie}[Każdy moduł ma prezentację wolną]
\label{aa:tw:modules-presentation-exists}
Dla każdego lewego $R$-modułu $M$ istnieją zbiory
$S,T$ i prezentacja $R^{(T)}\to R^{(S)}\to M\to0$.
\end{twierdzenie}
\begin{proof}
Weźmy $S=M$ jako zbiór i funkcję identycznościową
$u:S\to M$. Własność uniwersalna daje homomorfizm
$p:R^{(S)}\to M$ wysyłający $e_m$ na $m$.
Jest on surjektywny, bo każdy $m\in M$ jest obrazem
$e_m$. Niech $K=\ker p$ i weźmy $T=K$ jako zbiór.
Własność uniwersalna modułu wolnego daje
$d:R^{(T)}\to R^{(S)}$ określony na standardowych
wektorach przez $d(e_k)=k$. Wtedy
$\operatorname{im}d=K$; istotnie, obrazy należą do $K$,
a każdy $k\in K$ jest obrazem $e_k$.
Otrzymujemy wymaganą prezentację.
\end{proof}

\begin{przyklad}[Prezentacje modułów nad liczbami całkowitymi]
\label{aa:prz:modules-presentations-z}
Dla $n\ge2$ grupa cykliczna $\mathbb Z/n\mathbb Z$
ma prezentację
$\mathbb Z\xrightarrow{\cdot n}\mathbb Z\to\mathbb Z/n\mathbb Z\to0$.
Generator to klasa liczby $1$, a jedyna relacja
ma postać $n[1]=0$. Dla
$\mathbb Z/6\mathbb Z\oplus\mathbb Z/10\mathbb Z$
możemy użyć dwóch generatorów $u,v$ i relacji
$6u=0$, $10v=0$: homomorfizm
$\mathbb Z^2\to\mathbb Z^2$ o macierzy
$\operatorname{diag}(6,10)$ opisuje wszystkie relacje:
jej obraz wynosi $6\mathbb Z\oplus10\mathbb Z$,
czyli jądro epimorfizmu
$(a,b)\mapsto(a+6\mathbb Z,b+10\mathbb Z)$.
Gdy przejdziemy do postaci Smitha, znajdziemy
kanoniczny sposób upraszczania takich prezentacji.
\end{przyklad}

\begin{przyklad}[Moduł z jednym generatorem nie musi być wolny]
\label{aa:prz:modules-cyclic-not-free}
Niech $K$ będzie ciałem, a $R=K[t]/(t^2)$.
Moduł $M=R/(t)$ jest generowany przez klasę $1$,
ale $t\ne0$ w $R$ i $t\cdot(1+(t))=0$ w $M$.
W pierścieniu $R$ klasa $t$ jest niezerowa,
gdyż $t\notin(t^2)\subseteq K[t]$. Zeruje ona cały
$M$, ponieważ $t\cdot(r+(t))=tr+(t)=0$
dla dowolnego $r\in R$ (ideał $(t)$ jest dwustronny).
Gdyby $M$ był wolnym $R$-modułem z niepustą bazą,
to żaden niezerowy skalar nie mógłby zerować
wszystkich elementów: dla elementu bazowego $e_s$
równość $te_s=0$ przeczyłaby jednoznaczności
współrzędnych. Zatem moduł nie jest wolny. Pokazuje to, że generator
i element bazy to różne pojęcia.
\end{przyklad}

\section{Suma prosta, iloczyn prosty i ciągi dokładne}
\begin{definicja}[Suma prosta i iloczyn prosty modułów]
\label{aa:def:modules-direct-sum}
Dla rodziny lewych $R$-modułów $(M_i)_{i\in I}$
\emph{iloczyn prosty} $\prod_{i\in I}M_i$ składa się
ze wszystkich rodzin elementów, z działaniem
$R$ i dodawaniem określonymi współrzędnie.
\emph{Suma prosta} $\bigoplus_{i\in I}M_i$ jest
podmodułem tych rodzin, które mają skończony nośnik.
Istotnie, nośnik sumy dwóch rodzin jest zawarty
w sumie mnogościowej ich nośników, a nośnik
rodziny pomnożonej przez skalar jest zawarty
w nośniku wyjściowej rodziny. Oba pozostają więc
skończone. Rodzina zerowa należy do sumy prostej,
a element przeciwny nie zwiększa nośnika.
Dla skończonego zbioru indeksów obie konstrukcje
są równe, ale dla nieskończonego mogą się różnić.
Na przykład dla $I=\mathbb N$ i $M_i=\mathbb Z$
rodzina $(1,1,\ldots)$ należy do
$\prod_{i\in\mathbb N}\mathbb Z$, ale nie do
$\bigoplus_{i\in\mathbb N}\mathbb Z$, ponieważ
ma nieskończony nośnik.
\end{definicja}

\begin{definicja}[Dokładność ciągu modułów]
\label{aa:def:modules-exact}
Ciąg $L\xrightarrow{f}M\xrightarrow{g}N$
homomorfizmów lewych $R$-modułów jest \emph{dokładny
w $M$}, jeśli $\operatorname{im}f=\ker g$.
\emph{Krótki ciąg dokładny} ma postać
\[
0\longrightarrow A\xrightarrow{i}B
 \xrightarrow{p}C\longrightarrow0,
\]
gdzie $i$ jest iniekcją, $p$ jest surjekcją oraz
$\operatorname{im}i=\ker p$.
\end{definicja}

\begin{twierdzenie}[Kryteria rozszczepienia ciągu dokładnego]
\label{aa:tw:modules-splitting}
Dla krótkiego ciągu dokładnego lewych $R$-modułów
$0\to A\xrightarrow{i}B\xrightarrow{p}C\to0$
następujące warunki są równoważne:
\begin{enumerate}[label=\textup{(\roman*)}]
\item istnieje $R$-liniowe $s:C\to B$ takie, że
$p\circ s=\operatorname{id}_C$;
\item istnieje $R$-liniowe $r:B\to A$ takie, że
$r\circ i=\operatorname{id}_A$;
\item istnieje izomorfizm $B\cong A\oplus C$
zgodny z danym ciągiem, to znaczy przeprowadzający
$i(a)$ na $(a,0)$ i $p$ na rzut na drugą współrzędną.
\end{enumerate}
\end{twierdzenie}
\begin{proof}
Załóżmy (i). Dla $b\in B$ element
$b-s(p(b))$ należy do $\ker p=\operatorname{im}i$.
Ponieważ $i$ jest iniekcją, istnieje dokładnie
jedno $a\in A$ takie, że
$i(a)=b-s(p(b))$. Wzór
$\Phi(a,c)=i(a)+s(c)$ określa $R$-liniowe
odwzorowanie $A\oplus C\to B$. Jest ono na:
dla $b\in B$ wybieramy wyżej opisane $a$ i
$c=p(b)$. Jeśli $\Phi(a,c)=0$, po zastosowaniu
$p$ dostajemy $c=0$, a następnie $i(a)=0$,
więc $a=0$. Zatem $\Phi$ jest izomorfizmem.
Jego odwrotność ma postać
$\Phi^{-1}(b)=(i^{-1}(b-s(p(b))),p(b))$,
gdzie $i^{-1}$ jest odwrotnością $i:A\to\operatorname{im}i$;
argument pierwszej współrzędnej należy do
$\operatorname{im}i$ na mocy dokładności ciągu.
Ponadto $\Phi(a,0)=i(a)$ i $p\Phi(a,c)=c$.
Dostaliśmy (iii).

Z (iii) otrzymujemy (ii): jeśli
$\Psi:B\to A\oplus C$ jest izomorfizmem zgodnym
z ciągiem, to definiujemy $r=\operatorname{pr}_A\circ\Psi$.
Wtedy $r(i(a))=\operatorname{pr}_A(a,0)=a$.

Załóżmy teraz (ii). Dla $b\in B$ element
$b-i(r(b))$ należy do $\ker r$. Dlatego
każdy $b\in B$ zapisuje się jako
$b=i(r(b))+(b-i(r(b)))$, przy czym drugi składnik
leży w $\ker r$. Rozkład jest jednoznaczny:
$\ker r\cap i(A)=0$, gdyż dla $i(a)\in\ker r$
mamy $0=r(i(a))=a$. Zatem
$B=i(A)\oplus\ker r$. Ograniczenie
$p|_{\ker r}:\ker r\to C$ jest iniektywne,
ponieważ jego jądro to
$\ker r\cap\ker p=\ker r\cap\operatorname{im}i=0$.
Jest też na: dla $c\in C$ wybieramy $b\in B$
z $p(b)=c$ i mamy
$p(b-i(r(b)))=c$. Niech
$s$ będzie odwrotnością tego ograniczenia,
po której stosujemy włączenie $\ker r\hookrightarrow B$.
Odwrotność jest $R$-liniowa: jeśli $c_1=p(u_1)$,
$c_2=p(u_2)$ dla $u_1,u_2\in\ker r$, to
$p(u_1+u_2)=c_1+c_2$ oraz $p(tu_1)=tc_1$
dla każdego $t\in R$. Jednoznaczność przeciwobrazów
w $\ker r$ zapewnia zachowanie dodawania i skalarów.
Wtedy $p\circ s=\operatorname{id}_C$, co daje (i).
\end{proof}

\begin{przyklad}[Krótki ciąg, który się nie rozszczepia]
\label{aa:prz:modules-nonsplit-z}
Rozważmy ciąg $\mathbb Z$-modułów
\[
0\longrightarrow\mathbb Z\xrightarrow{\cdot2}\mathbb Z
 \xrightarrow{\pi}\mathbb Z/2\mathbb Z\longrightarrow0.
\]
Jest dokładny: mnożenie przez $2$ jest iniektywne,
$\pi$ jest surjekcją, a $\ker\pi=2\mathbb Z$.
Gdyby istniał przekrój liniowy
$s:\mathbb Z/2\mathbb Z\to\mathbb Z$, to
$\pi(s([1]))=[1]$, więc $s([1])$ byłby nieparzysty.
Jednak $2s([1])=s(2[1])=s([0])=0$,
co w $\mathbb Z$ wymusza $s([1])=0$.
Sprzeczność dowodzi, że ciąg się nie rozszczepia.
W szczególności nie każdy podmoduł jest
składnikiem sumy prostej.
\end{przyklad}

\section{Moduły rzutowe}
\begin{definicja}[Moduł rzutowy]
\label{aa:def:modules-projective}
Niech $R$ będzie pierścieniem z jedynką,
niekoniecznie przemiennym. Lewy $R$-moduł $P$ jest
\emph{modułem rzutowym}, jeśli dla dowolnych lewych
$R$-modułów $E,N$, każdego surjektywnego homomorfizmu
$R$-modułów $q:E\to N$ i każdego homomorfizmu
$R$-modułów $f:P\to N$
istnieje homomorfizm $\widetilde f:P\to E$
spełniający $q\circ\widetilde f=f$.
Mówimy, że $\widetilde f$ jest \emph{podniesieniem}
$f$ względem $q$.
\end{definicja}

\begin{twierdzenie}[Charakterystyka modułów rzutowych]
\label{aa:tw:modules-projective-summand}
Niech $R$ będzie pierścieniem z jedynką,
niekoniecznie przemiennym, a $P$ lewym $R$-modułem.
Następujące warunki są równoważne:
\begin{enumerate}[label=\textup{(\roman*)}]
\item $P$ jest modułem rzutowym;
\item każdy krótki ciąg dokładny
$0\to A\to B\to P\to0$ rozszczepia się;
\item $P$ jest izomorficzny ze składnikiem sumy prostej
pewnego wolnego lewego $R$-modułu.
\end{enumerate}
\end{twierdzenie}
\begin{proof}
Załóżmy (i) i rozważmy krótki ciąg dokładny
$0\to A\to B\xrightarrow{q}P\to0$.
Stosując definicję modułu rzutowego, wybierzmy $f=\operatorname{id}_P$.
Otrzymujemy $s:P\to B$ z $q\circ s=\operatorname{id}_P$.
Twierdzenie~\ref{aa:tw:modules-splitting} daje rozszczepienie,
a więc (ii).

Załóżmy (ii). Z Twierdzenia~\ref{aa:tw:modules-presentation-exists}
istnieje surjekcja $p:F\to P$ z wolnego modułu $F$.
Krótki ciąg $0\to\ker p\to F\xrightarrow{p}P\to0$
rozszczepia się, zatem $F\cong\ker p\oplus P$.
Daje to (iii).

Załóżmy (iii). Najpierw pokażemy, że wolny moduł
$F=R^{(S)}$ jest modułem rzutowym. Dla surjekcji
$q:E\to N$ i homomorfizmu $f:F\to N$ wybierzmy
dla każdego $s\in S$ element $e_s'\in E$ taki, że
$q(e_s')=f(e_s)$. Przy nieskończonym $S$ korzystamy
tu z aksjomatu wyboru dla niepustych włókien surjekcji.
Własność uniwersalna modułu
wolnego daje $\widetilde f:F\to E$ z
$\widetilde f(e_s)=e_s'$. Wtedy $q\widetilde f$
i $f$ są równe na bazie $F$, więc są równe
na całym $F$.

Niech teraz $F=P\oplus Q$; przypadek modułu
izomorficznego z $P$ rozpatrujemy po przeniesieniu
struktury przez izomorfizm. Dla surjekcji $q:E\to N$
i homomorfizmu $f:P\to N$ rozważmy
$h:F\to N$, $h(p,q')=f(p)$. Wolność $F$
pozwala podnieść $h$ do $\widetilde h:F\to E$.
Ograniczając $\widetilde h$ do $P\oplus\{0\}$
dostajemy podniesienie $f$. Zatem $P$ jest modułem rzutowym i zachodzi (i).
\end{proof}

\begin{przyklad}[Własność bycia modułem rzutowym zależy od pierścienia]
\label{aa:prz:modules-projectivity-example}
Każdy moduł nad ciałem $K$ jest wolny,
a więc jest modułem rzutowym. Moduł $\mathbb Z/2\mathbb Z$
nie jest modułem rzutowym nad $\mathbb Z$, ponieważ
krótki ciąg dokładny z
Przykładu~\ref{aa:prz:modules-nonsplit-z} nie rozszczepia
się. Z drugiej strony dla pierścienia
$R=K\times K$ moduł $P=K\times\{0\}$ jest
rzutowy: $R=P\oplus(\{0\}\times K)$ jako
moduł nad sobą. Nie jest on wolny jako $R$-moduł:
$P\ne0$, a niezerowy element $(0,1)\in R$ zeruje
wszystkie elementy $P$. Gdyby $P$ miał niepustą bazę,
dla jej elementu $e$ mielibyśmy $(0,1)e\ne0$,
gdyż współczynnik $(0,1)$ jest niezerowy i
przedstawienie w bazie jest jednoznaczne.
Otrzymujemy sprzeczność. Jest to konkretny
przykład modułu rzutowego, który nie jest wolny.
\end{przyklad}


\begin{przyklad}[Rozszczepienie nie wynika z dokładności]
\label{aa:prz:modules-split-versus-exact}
Dla każdego pierścienia $R$ z jedynką i lewych $R$-modułów
$A,C$ ciąg
\[
0\longrightarrow A\xrightarrow{i}A\oplus C
\xrightarrow{p}C\longrightarrow0,
\qquad i(a)=(a,0),\quad p(a,c)=c,
\]
jest dokładny i rozszczepiony: odwzorowanie $s(c)=(0,c)$
spełnia $p\circ s=\operatorname{id}_C$. W przeciwieństwie
do tego ciąg z Przykładu~\ref{aa:prz:modules-nonsplit-z}
jest dokładny, ale nie ma takiego przekroju. Dokładność
wymaga jedynie równości obrazu i jądra, nie istnienia
uzupełniającego podmodułu.
\end{przyklad}

\begin{przyklad}[Dlaczego nie każdy homomorfizm można mnożyć przez skalar z lewej]
\label{aa:prz:modules-hom-noncommutative}
Niech $K$ będzie ciałem, $R=M_2(K)$, a $M=N={}_RR$.
Identyczność $f=\operatorname{id}_R$ jest lewym
$R$-liniowym endomorfizmem. Ustalmy $a=E_{12}$ i
rozważmy odwzorowanie $h(x)=af(x)=ax$. Gdyby $h$
było lewym $R$-liniowym homomorfizmem, to dla
$r=E_{21}$ i $x=1_R$ mielibyśmy $h(rx)=rh(x)$.
Tymczasem $h(E_{21})=E_{12}E_{21}=E_{11}$,
natomiast $E_{21}h(1_R)=E_{21}E_{12}=E_{22}$.
Ponieważ $E_{11}\ne E_{22}$, $h$ nie jest
$R$-liniowe. Wyjaśnia to, dlaczego w
Stwierdzeniu~\ref{aa:st:modules-hom-abelian}
przemienność pierścienia jest założeniem, a nie
jedynie wygodną konwencją.
\end{przyklad}

\section{Dlaczego strony modułu mają znaczenie?}
\begin{stwierdzenie}[Endomorfizmy lewego modułu regularnego]
\label{aa:st:modules-end-regular}
Dla pierścienia $R$ z jedynką zachodzi izomorfizm
pierścieni $\operatorname{End}_R({}_RR)\cong R^{\mathrm{op}}$,
gdzie ${}_RR$ oznacza lewy $R$-moduł dany przez
mnożenie w pierścieniu, a w pierścieniu endomorfizmów
mnożeniem jest składanie odwzorowań.
\end{stwierdzenie}
\begin{proof}
Jeśli $f:R\to R$ jest lewym $R$-liniowym endomorfizmem,
to dla każdego $r\in R$ mamy
$f(r)=f(r\cdot1_R)=rf(1_R)$. Zatem $f$
jest jednoznacznie wyznaczony przez element
$a=f(1_R)$ i ma postać $f_a(r)=ra$.
Każdy taki wzór określa lewy $R$-liniowy endomorfizm,
bo $f_a(sr)=(sr)a=s(ra)=sf_a(r)$.
Dla $a,b\in R$ mamy
$(f_a\circ f_b)(r)=f_a(rb)=(rb)a=r(ba)=f_{ba}(r)$.
Zatem odwzorowanie $a\mapsto f_a$ zachowuje
mnożenie w pierścieniu przeciwnym $R^{\mathrm{op}}$,
a także dodawanie i jedynkę. Jest bijekcją,
więc jest izomorfizmem pierścieni.
\end{proof}

\begin{uwaga}[Dalszy program teorii modułów]
\label{aa:uw:modules-future}
Zdefiniowaliśmy moduły, podmoduły, prezentacje,
krótkie ciągi dokładne i moduły rzutowe, nie zakładając
przemienności pierścienia. W dalszym ciągu kursu
opracujemy klasyfikację modułów nad dziedzinami
ideałów głównych oraz postać Smitha. Osobno
rozwiniemy \emph{iloczyn tensorowy modułów}, który
nie jest ani sumą prostą, ani iloczynem prostym.
Dla nieprzemiennego pierścienia iloczyn tensorowy
nad $R$ wymaga co do zasady prawego $R$-modułu
i lewego $R$-modułu; przyjmiemy wtedy odpowiednią
własność zbilansowania.
\end{uwaga}

\setlength{\parskip}{4pt}

% AA-17 | Algebra abstrakcyjna | audyt I
\chapter[Postać Smitha i klasyfikacja modułów]{Moduły nad dziedzinami ideałów głównych i postać Smitha}
\label{aa:ch:pid-modules-smith}

\section{Od prezentacji do klasyfikacji}
W Rozdziale~\ref{aa:ch:modules-introduction} pokazaliśmy, że każdy moduł
jest ilorazem modułu wolnego, a każdy moduł skończenie generowany jest
obrazem pewnego wolnego modułu skończonej rangi. W tym rozdziale pytamy, czy relacje między
nimi można uprościć tak, by otrzymać jednoznaczny opis modułu. Dla
ogólnych pierścieni odpowiedź jest przecząca. Dla \emph{dziedziny ideałów
głównych} sytuacja okazuje się jednak wyjątkowa: odpowiednią macierz
relacji można sprowadzić do postaci diagonalnej przez odwracalne
operacje na wierszach i kolumnach.

\begin{definicja}[Dziedzina ideałów głównych]
\label{aa:def:pid17}
Dziedziną ideałów głównych nazywamy niezerowy pierścień przemienny
$R$ z jedynką, bez dzielników zera, w którym każdy ideał ma postać
$(a)=aR$ dla pewnego $a\in R$. Element $u\in R$ jest \emph{jednostką},
jeśli istnieje $v\in R$ z $uv=1$. Elementy $a,b$ są \emph{stowarzyszone},
gdy $a=ub$ dla pewnej jednostki $u$. Dalej $R$ jest dziedziną ideałów
głównych, o ile nie zaznaczymy inaczej.
\end{definicja}

Przykłady obejmują $\mathbb Z$ oraz $K[t]$ dla dowolnego ciała $K$;
dla $K[t]$ stosujemy Twierdzenie~\ref{aa:tw:euclidean-pid}.
Pierścień $\mathbb Z[t]$ nie jest dziedziną ideałów głównych.
Istotnie, gdyby $(2,t)=(f)$, to wielomian $f$ dzieliłby
zarówno $2$, jak i $t$ w $\mathbb Z[t]$. Dzielnik $2$ jest
stały i równy $\pm1$ lub $\pm2$; ponieważ $2$ nie dzieli
$t$, otrzymalibyśmy $f=\pm1$ i $(2,t)=\mathbb Z[t]$.
Jest to sprzeczne z faktem, że homomorfizm
$\mathbb Z[t]\to\mathbb Z/2\mathbb Z$,
$g(t)\mapsto g(0)\bmod2$, zeruje $2$ i $t$, ale
nie zeruje $1$.

\begin{lemat}[Warunek stabilizacji rosnących ciągów ideałów]
\label{aa:lem:pid-acc17}
W dziedzinie ideałów głównych $R$ każdy rosnący ciąg ideałów
$I_1\subseteq I_2\subseteq\cdots$ od pewnego miejsca jest stały.
\end{lemat}
\begin{proof}
Suma mnogościowa $I=\bigcup_{n\ge1}I_n$ jest ideałem: dla $x\in I_m$
i $y\in I_n$ mamy $x-y\in I_{\max(m,n)}$, a mnożenie przez dowolny
element $R$ zachowuje tę samą składową. Ponieważ $R$ jest dziedziną
ideałów głównych, $I=(a)$. Generator $a$ należy do pewnego $I_N$.
Zatem $I=(a)\subseteq I_N\subseteq I$, czyli $I_N=I$ i wszystkie
$I_n$ dla $n\ge N$ są równe $I$. Nie użyliśmy funkcji euklidesowej;
argument działa w dowolnej dziedzinie ideałów głównych.
\end{proof}

\section{Podmoduły modułów wolnych i prezentacje skończone}
\begin{twierdzenie}[Podmoduł wolnego modułu skończonej rangi]
\label{aa:tw:pid-submodule-free}
Niech $R$ będzie dziedziną ideałów głównych, a $N\subseteq R^n$
podmodułem. Wtedy $N$ jest wolnym $R$-modułem o bazie zawierającej
co najwyżej $n$ elementów.
\end{twierdzenie}
\begin{proof}
Postępujemy indukcyjnie względem $n$. Dla $n=0$ moduł $R^0=0$,
a więc $N=0$ jest wolny z bazą pustą. Załóżmy twierdzenie dla $n-1$
i rozważmy projekcję $\pi:R^n\to R$ na pierwszą współrzędną.
Obraz $\pi(N)$ jest ideałem, ma zatem postać $aR$.

Jeżeli $a=0$, to $N\subseteq\{0\}\oplus R^{n-1}$ i teza wynika
z założenia indukcyjnego. Niech więc $a\ne0$. Wybierzmy $u\in N$
z pierwszą współrzędną $a$. Dla każdego $x\in N$ istnieje $r\in R$
takie, że $\pi(x)=ra$; wówczas $x-ru\in N\cap\ker\pi$.
Zatem $N=Ru+(N\cap\ker\pi)$. Suma jest prosta: jeśli
$ru\in\ker\pi$, to $ra=0$, a brak dzielników zera i $a\ne0$
dają $r=0$. Moduł $N\cap\ker\pi$ jest podmodułem $R^{n-1}$,
więc na mocy indukcji jest wolny z bazą co najwyżej $n-1$ elementów.
Dołączenie $u$ do tej bazy daje bazę $N$ długości co najwyżej $n$.
\end{proof}

\begin{wniosek}[Skończona prezentacja]
\label{aa:wn:pid-finite-presentation}
Niech $R$ będzie dziedziną ideałów głównych, a $M$
skończenie generowanym lewym $R$-modułem.
Wówczas $M$ ma prezentację
$R^m\xrightarrow{A}R^n\longrightarrow M\longrightarrow0$,
gdzie $m,n$ są skończone, a $A$ jest macierzą o współczynnikach w $R$.
\end{wniosek}
\begin{proof}
Wybierzmy generatory $m_1,\dots,m_n$ modułu $M$. Homomorfizm
$p:R^n\to M$ zadany przez $p(e_i)=m_i$ jest surjektywny.
Jądro $N=\ker p$ jest podmodułem $R^n$, a zatem jest wolne
z bazą $u_1,\dots,u_m$, gdzie $m\le n$. Odwzorowanie
$A:R^m\to R^n$ wysyłające $j$-ty wektor bazowy na $u_j$
ma obraz $N$. Otrzymujemy wymagany ciąg dokładny.
\end{proof}

\begin{uwaga}[Dlaczego założenie o dziedzinie jest istotne?]
Dla $R=K[\varepsilon]/(\varepsilon^2)$ ideał $(\varepsilon)$
nie jest wolnym $R$-modułem, choć jest podmodułem wolnego modułu $R$.
Rzeczywiście $\varepsilon(\varepsilon)=0$ i $\varepsilon\ne0$,
więc niezerowy skalar zeruje każdy jego element; w niezerowym
module wolnym nie zeruje każdego elementu bazowego. W pierścieniu
$R$ występują dzielniki zera, więc twierdzenie nie ma zastosowania.
\end{uwaga}

\section{Operacje elementarne i postać Smitha}
\begin{definicja}[Macierz relacji i operacje elementarne]
\label{aa:def:smith-operations}
Niech $A\in M_{m\times n}(R)$ wyznacza homomorfizm
$A:R^n\to R^m$. Dopuszczamy następujące operacje na wierszach
i kolumnach: zamianę miejscami, pomnożenie jednego wiersza lub
jednej kolumny przez jednostkę pierścienia oraz dodanie wielokrotności
jednego wiersza do innego (analogicznie dla kolumn).
Każda z tych operacji jest mnożeniem przez macierz odwracalną
z lewej lub z prawej strony. Ponadto w dowodzie postaci Smitha
będziemy wykonywać odwracalne zmiany baz na dwóch wybranych
wierszach lub kolumnach, opisane w
Lemacie~\ref{aa:lem:smith-bezout}. Są one również
mnożeniem przez macierz odwracalną, ale nad dowolną
dziedziną ideałów głównych nie zakładamy, że każdą taką
zmianę można rozłożyć na trzy powyższe typy prostych
operacji elementarnych. Rozróżniamy zatem
\emph{operacje elementarne} i ogólne \emph{odwracalne
zmiany baz}.
\end{definicja}

\begin{lemat}[Dwa elementy można zastąpić największym wspólnym dzielnikiem]
\label{aa:lem:smith-bezout}
Niech $R$ będzie dziedziną ideałów głównych, $a,b\in R$
i $(a,b)=(d)$, gdzie $d\ne0$.
Istnieje odwracalna macierz $U\in\operatorname{GL}_2(R)$ taka, że
$U(a,b)^T=(d,0)^T$.
\end{lemat}
\begin{proof}
Ponieważ $d$ generuje ideał $(a,b)$, istnieją $u,v\in R$
z $ua+vb=d$. Ponieważ $d$ dzieli $a$ i $b$, mamy $a=da'$,
$b=db'$ z $a',b'\in R$. Zapiszmy
\[
U=\begin{pmatrix}u&v\\-b'&a'\end{pmatrix}.
\]
Wyznacznik $U$ wynosi $ua'+vb'=(ua+vb)/d=1$, gdzie
skracanie przez niezerowy $d$ jest dozwolone, ponieważ $R$ nie
ma dzielników zera. Macierz $U$ jest zatem odwracalna: jej odwrotność to
$\begin{pmatrix}a'&-v\\b'&u\end{pmatrix}$.
Ponadto $U(a,b)^T=(d,-b'a+a'b)^T=(d,0)^T$.
Dla pary wpisanej w wiersz stosujemy macierz $U^T$ po prawej
stronie. Są to odwracalne zmiany baz, których poprawność
nie zależy od istnienia algorytmu dzielenia z resztą w $R$.
Nie utożsamiamy ich bez dowodu ze złożeniem operacji
elementarnych; nad dziedziną euklidesową taką redukcję
można natomiast realizować algorytmem Euklidesa.
\end{proof}

\begin{twierdzenie}[Postać normalna Smitha]
\label{aa:tw:smith-normal-form}
Niech $R$ będzie dziedziną ideałów głównych i
$A\in M_{m\times n}(R)$, gdzie $m,n\ge0$.
Istnieją odwracalne macierze $U\in\operatorname{GL}_m(R)$
i $V\in\operatorname{GL}_n(R)$ takie, że $UAV$ jest macierzą
prostokątną diagonalną, której niezerowe elementy diagonalne
$d_1,\dots,d_s$ spełniają
$d_1\mid d_2\mid\cdots\mid d_s$.
Pozostałe elementy są równe zeru. Jeśli $R$ jest dziedziną
euklidesową i dysponujemy efektywnym dzieleniem z resztą,
przekształcenia można wykonywać algorytmicznie.
Dla ogólnej dziedziny
ideałów głównych twierdzenie zapewnia istnienie
odwracalnych macierzy $U,V$, ale samo założenie o
ideałach głównych nie dostarcza efektywnej procedury
obliczania generatorów wszystkich występujących ideałów.
\end{twierdzenie}
\begin{proof}
Dla pustej macierzy lub macierzy zerowej twierdzenie jest
oczywiste. Załóżmy, że $A$ zawiera niezerowy współczynnik;
zamiana wierszy i kolumn ustawia go na pozycji $(1,1)$.
Nazwijmy ten element $a\ne0$.

Najpierw zbudujemy etap, w którym $a$ dzieli wszystkie
pozostałe współczynniki macierzy, a pierwsza kolumna i pierwszy
wiersz poza $a$ są zerowe. Jeśli pewien wpis $b$ pierwszej
kolumny nie jest podzielny przez $a$, zastosujmy do odpowiednich
dwóch wierszy Lemat~\ref{aa:lem:smith-bezout}. Powstały nowy
wpis w pozycji $(1,1)$ jest generatorem ideału $(a,b)$.
Ponieważ $a$ nie dzieli $b$, mamy ścisłe zawieranie
$(a)\subsetneq(a,b)$. Tak samo dla współczynnika pierwszego
wiersza, stosując operacje na dwóch kolumnach.

Po każdej takiej zmianie wpisu wiodącego ponownie sprawdzamy
pierwszy wiersz i kolumnę. Gdy $a$ dzieli już wszystkie ich
wpisy, przez odejmowanie
odpowiednich wielokrotności pierwszego wiersza od dalszych
wierszy zerujemy pierwszą kolumnę poniżej $a$. Następnie
odejmując wielokrotności pierwszej kolumny od pozostałych
kolumn zerujemy pierwszy wiersz poza $a$. Dostajemy
$\operatorname{diag}(a,B)$ dla pewnej macierzy $B$.

Jeśli $a$ dzieli wszystkie wpisy $B$, pierwszy etap jest zakończony.
Jeżeli natomiast w $B$ istnieje wpis $b$ niepodzielny przez $a$,
dodajmy jego wiersz do pierwszego wiersza całej macierzy.
Wpis na pozycji $(1,1)$ pozostaje równy $a$, ponieważ poniżej
niego są już zera; w pierwszym wierszu pojawia się za to
wpis $b$. Zastosowanie lematu do dwóch odpowiednich kolumn
zastępuje $a$ generatorem $(a,b)$ i ściśle powiększa ideał
wyznaczony przez wpis wiodący. Po tej operacji wracamy do
kontroli podzielności w pierwszym wierszu i kolumnie;
nie zakładamy, że wcześniejsze zera musiały się zachować.

W każdym kroku, w którym wpis wiodący się zmienia,
powstaje ściśle rosnący ideał główny $(a)$: jeśli
$a$ nie dzieli wybranego $b$, to $(a)\subsetneq(a,b)$.
Lemat~\ref{aa:lem:pid-acc17} wyklucza nieskończony
ciąg takich zmian. Przy ustalonym wpisie wiodącym
sprawdzamy skończenie wiele pozycji w pierwszym wierszu
i kolumnie, a następnie skończenie wiele wpisów bloku
$B$. Gdy któryś wpis nie jest podzielny przez $a$,
po najwyżej skończenie wielu działaniach otrzymujemy
nowy wpis wiodący generujący ściśle większy ideał.
Zatem po skończenie wielu operacjach otrzymujemy
macierz $\operatorname{diag}(d_1,B)$, gdzie $d_1\ne0$ dzieli
każdy współczynnik $B$.

Teraz postępujemy indukcyjnie z macierzą $B$ o wymiarach
$(m-1)\times(n-1)$. Operacje w jej wierszach i kolumnach
nie wpływają na pierwszy element diagonalny ani na zera poza
nim. Ponieważ wszystkie współczynniki $B$ są wielokrotnościami
$d_1$, każda późniejsza kombinacja liniowa tych współczynników
też jest wielokrotnością $d_1$. W szczególności uzyskany
niezerowy drugi element diagonalny $d_2$ jest podzielny przez
$d_1$. Powtarzając argument dostajemy żądaną postać.
Wszystkie operacje są odwracalne, stąd istnieją $U,V$.
Dla dziedziny euklidesowej wybór kombinacji i dzielenie
są dodatkowo efektywne.
\end{proof}

\begin{uwaga}[Ważne: dziedzina ideałów głównych nie musi być euklidesowa]
Dowód wykorzystuje stabilizację rosnących ciągów ideałów
oraz tożsamości Bézouta. W ogólnej dziedzinie ideałów
głównych stosujemy także odwracalne bloki macierzowe
$2\times2$; nie twierdzimy, że można je zawsze rozłożyć
na same operacje elementarne. W żadnym miejscu nie zakładamy
istnienia funkcji euklidesowej. Dlatego postać Smitha istnieje dla
każdej dziedziny ideałów głównych, nie tylko dla $\mathbb Z$
i $K[t]$.
\end{uwaga}

\section{Niezmienniki i jednoznaczność postaci Smitha}
\begin{definicja}[Ideały wyznaczone przez minory]
\label{aa:def:determinantal-ideals}
Niech $R$ będzie dziedziną ideałów głównych,
$A\in M_{m\times n}(R)$, gdzie $m,n\ge0$.
Dla $k\ge1$ oznaczamy przez $D_k(A)$ ideał
generowany przez wyznaczniki wszystkich podmacierzy kwadratowych
stopnia $k$. Kładziemy $D_0(A)=R$ oraz $D_k(A)=(0)$,
gdy $k>\min(m,n)$. Ideały $D_k(A)$ nazywamy
\emph{ideałami minorowymi} macierzy.
\end{definicja}
\begin{lemat}[Niezmienniczość ideałów minorowych]
\label{aa:lem:minor-invariant}
Niech $R$ będzie dziedziną ideałów głównych,
$A\in M_{m\times n}(R)$, $m,n\ge0$.
Dla dowolnych $U\in\operatorname{GL}_m(R)$ i
$V\in\operatorname{GL}_n(R)$ zachodzi $D_k(UAV)=D_k(A)$
dla każdego $k\ge0$.
\end{lemat}
\begin{proof}
Dla $k=0$ i $k>\min(m,n)$ teza wynika z definicji.
Niech zatem $1\le k\le\min(m,n)$. Wykorzystamy
rozwinięcie wyznacznika względem wyboru wierszy i kolumn.
Niech $I$ będzie zbiorem $k$ indeksów wierszy,
a $J$ zbiorem $k$ indeksów kolumn. Z wieloliniowości
wyznacznika i faktu, że wyznacznik z dwoma identycznymi
wierszami jest zerowy, otrzymujemy wzór
\[
 \det((UA)_{I,J})=
 \sum_{K:\,|K|=k}\det(U_{I,K})\det(A_{K,J}),
\]
gdzie $K$ przebiega $k$-elementowe zbiory indeksów
pośrednich. Aby zobaczyć ten wzór, rozpisujemy każdy wiersz
podmacierzy $(UA)_{I,J}$ jako kombinację wierszy $A$,
rozwijamy wyznacznik według wieloliniowości, odrzucamy
wyrazy z powtórzonym indeksem wiersza, a pozostałe
grupiujemy według rosnących ciągów indeksów $K$.
Współczynnik przy $\det(A_{K,J})$ jest wtedy dokładnie
$\det(U_{I,K})$, z uwzględnieniem znaków permutacji.

Każdy minor $UA$ jest więc kombinacją liniową minorów $A$,
stąd $D_k(UA)\subseteq D_k(A)$. Ponieważ $U$ jest
odwracalna, ten sam argument dla $U^{-1}$ daje
$D_k(A)\subseteq D_k(UA)$. Otrzymujemy równość.
Przeprowadzając analogiczne rozwinięcie dla kolumn
macierzy $AV$ i używając $V^{-1}$, dostajemy
$D_k(AV)=D_k(A)$. Złożenie obu równości kończy dowód.

\end{proof}

\begin{twierdzenie}[Jednoznaczność postaci Smitha]
\label{aa:tw:smith-uniqueness}
Niech $R$ będzie dziedziną ideałów głównych,
$A\in M_{m\times n}(R)$, $m,n\ge0$.
Niech $\operatorname{diag}(d_1,\dots,d_s,0,\dots)$ będzie
postacią Smitha macierzy $A$, gdzie $s\ge0$, $d_i\ne0$
i $d_i\mid d_{i+1}$ dla $1\le i<s$. Wówczas dla $1\le k\le s$
zachodzi $D_k(A)=(d_1\cdots d_k)$, natomiast
$D_k(A)=(0)$ dla $k>s$. Zatem liczba $s$ oraz
wszystkie $d_i$ są wyznaczone przez $A$ jednoznacznie
z dokładnością do stowarzyszenia.
\end{twierdzenie}
\begin{proof}
Z Lematu~\ref{aa:lem:minor-invariant} możemy liczyć
ideały minorowe dla postaci diagonalnej. Każdy niezerowy
minor stopnia $k$ jest, z dokładnością do znaku,
iloczynem $k$ różnych elementów diagonalnych.
Warunek $d_1\mid\dots\mid d_s$ sprawia, że
$d_1\cdots d_k$ dzieli każdy taki iloczyn, a sam
jest minorem z pierwszych $k$ wierszy i kolumn.
Dlatego ideał wszystkich minorów stopnia $k$ jest
$(d_1\cdots d_k)$. Dla $k>s$ wszystkie minory są zerowe.

Jeżeli dwie postacie Smitha mają współczynniki $d_i$
i $e_i$, z równości ideałów $D_1$ wynika, że $d_1$
i $e_1$ są stowarzyszone. Indukcyjnie porównujemy kolejne ideały
$(d_1\cdots d_k)=(e_1\cdots e_k)$. Generatory
równego niezerowego ideału są stowarzyszone, więc
$d_1\cdots d_k=u(e_1\cdots e_k)$ dla pewnej jednostki $u$.
Z założenia indukcyjnego
$d_1\cdots d_{k-1}=v(e_1\cdots e_{k-1})$
dla pewnej jednostki $v$. Skracając niezerowy iloczyn
$e_1\cdots e_{k-1}$ w dziedzinie całkowitości,
dostajemy $d_k=uv^{-1}e_k$. Zatem również $d_k$
i $e_k$ są stowarzyszone. Liczba
niezerowych współczynników to największe $k$
z $D_k(A)\ne(0)$, więc również jest niezmiennikiem.
\end{proof}

\begin{przyklad}[Rachunek postaci Smitha nad liczbami całkowitymi]
\label{aa:prz:smith-integer-steps}
Niech $A=\begin{pmatrix}4&6\\10&14\end{pmatrix}$.
Odejmując dwukrotność pierwszego wiersza od drugiego,
a następnie zamieniając wiersze, otrzymujemy
$\begin{pmatrix}2&2\\4&6\end{pmatrix}$.
Odejmując dwukrotność pierwszego wiersza od drugiego,
dostajemy $\begin{pmatrix}2&2\\0&2\end{pmatrix}$.
Na koniec odejmujemy pierwszą kolumnę od drugiej:
$\operatorname{diag}(2,2)$. W szczególności
$\mathbb Z^2/A\mathbb Z^2\cong C_2\oplus C_2$.
Bez wykonania operacji wynik kontrolujemy przez
$\gcd(4,6,10,14)=2$ i $|\det A|=4$:
pierwszy współczynnik Smitha ma wartość bezwzględną $2$,
a iloczyn dwóch współczynników ma wartość bezwzględną $4$.
\end{przyklad}

\begin{przyklad}[Dwa równania relacji i iloraz]
\label{aa:prz:smith-cokernel-90}
Dla $A=\begin{pmatrix}6&9\\0&15\end{pmatrix}$
mamy $D_1(A)=(3)$, $D_2(A)=(90)$. Dlatego
postać Smitha ma współczynniki stowarzyszone z $3$ i $30$.
Moduł o dwóch generatorach $u,v$ i relacjach
$6u=0$, $9u+15v=0$ jest izomorficzny z
$\mathbb Z/3\mathbb Z\oplus\mathbb Z/30\mathbb Z$.
Iloraz ma $90$ elementów, ponieważ ma strukturę
$C_3\oplus C_{30}$. Nie należy utożsamiać liczby
elementów z liczbą generatorów: wystarczają dwa.
W rzeczywistości jeden nie wystarczy, gdyż
$M/3M\cong(\mathbb Z/3\mathbb Z)^2$ ma wymiar dwa
nad $\mathbb Z/3\mathbb Z$.
\end{przyklad}

\section{Klasyfikacja skończenie generowanych modułów}
\begin{twierdzenie}[Twierdzenie o strukturze modułów nad dziedziną ideałów głównych]
\label{aa:tw:pid-structure}
Niech $R$ będzie dziedziną ideałów głównych, a $M$
skończenie generowanym $R$-modułem. Istnieją liczby całkowite
$r,t\ge0$ oraz niezerowe, nieodwracalne elementy
$d_1,\dots,d_t\in R$ spełniające $d_1\mid\cdots\mid d_t$
(dla $t=0$ rodzina jest pusta) takie, że
\[
 M\cong R^r\oplus\bigoplus_{i=1}^t R/(d_i).
\]
Liczba $r$ jest wyznaczona jednoznacznie, a ciąg
$d_1,\dots,d_t$ jest jednoznaczny z dokładnością do
pomnożenia każdego elementu przez jednostkę $R$.
\end{twierdzenie}
\begin{proof}
Najpierw dowiedziemy istnienia. Z
Wniosku~\ref{aa:wn:pid-finite-presentation} mamy
$M\cong\operatorname{coker}(A)$ dla pewnego
$A:R^m\to R^n$. Z Twierdzenia~\ref{aa:tw:smith-normal-form}
otrzymujemy $UAV=D$, gdzie $U$ i $V$ są odwracalne.
Ponieważ $V$ jest automorfizmem modułu źródłowego,
$\operatorname{im}(AV)=\operatorname{im}A$.
Ponieważ $U$ jest automorfizmem modułu docelowego,
indukuje izomorfizm $R^n/\operatorname{im}A\cong
R^n/\operatorname{im}D$. Dla macierzy diagonalnej
iloraz rozkłada się współrzędnie na
$R/(d_1)\oplus\cdots\oplus R/(d_s)\oplus R^{n-s}$.
Jeśli $d_i$ jest jednostką, odpowiedni czynnik jest zerowy;
pozostałe niezerowe współczynniki zachowują warunek
podzielności. To dowodzi istnienia wskazanego rozkładu.

Dla jednoznaczności najpierw oznaczmy przez $T(M)$ zbiór
elementów $m\in M$, dla których $am=0$ dla pewnego
$a\in R\setminus\{0\}$. Jeżeli $am=0$ i $bn=0$ dla niezerowych $a,b\in R$,
to $ab(m-n)=0$, przy czym $ab\ne0$. Dla dowolnego
$c\in R$ mamy ponadto $a(cm)=c(am)=0$, dlatego
$T(M)$ jest podmodułem. Każdy czynnik $R/(d_i)$
jest zerowany przez $d_i$, a cała skończona suma tych
czynników przez iloczyn $d_1\cdots d_t$ (dla $t=0$
przez $1$). Natomiast niezerowy wektor w $R^r$ nie
może zostać wyzerowany przez niezerowy skalar, ponieważ
$R$ nie ma dzielników zera. Stąd w podanym rozkładzie
$T(M)\cong\bigoplus_i R/(d_i)$ i $M/T(M)\cong R^r$.
Dwie rangi wolnych modułów nad $R$ są równe także wtedy,
gdy $R$ nie jest ciałem. Niech $K=\operatorname{Frac}(R)$
będzie ciałem ułamków dziedziny $R$. Izomorfizm
$R^r\cong R^s$ jest reprezentowany przez macierz
$A\in M_{s\times r}(R)$, a jego odwrotność przez
$B\in M_{r\times s}(R)$, z $BA=I_r$ i $AB=I_s$.
Traktując współczynniki jako elementy $K$, dostajemy
wzajemnie odwrotne $K$-liniowe odwzorowania
$K^r\to K^s$ i $K^s\to K^r$. Z równości
wymiarów izomorficznych przestrzeni wektorowych
wynika $r=s$. Argument ten działa dla każdej
dziedziny całkowitości i nie wymaga wcześniejszej
teorii iloczynu tensorowego.

Pozostaje jednoznaczność części torsyjnej. Z
Twierdzenia~\ref{aa:tw:pid-ufd} każdy niezerowy
nieodwracalny element $R$ rozkłada się jednoznacznie
na czynniki nierozkładalne z dokładnością do jednostek
i kolejności. Niech $p\in R$ będzie elementem
nierozkładalnym. Ideał $(p)$ jest maksymalny:
jeśli $(p)\subseteq(a)\subsetneq R$, to $a$ dzieli
$p$, a ponieważ $a$ nie jest jednostką, z nierozkładalności
$p$ wynika, że $a$ jest stowarzyszony z $p$.
Stąd $R/(p)$ jest ciałem.
Dla dowolnego modułu torsyjnego $T$ oznaczmy
$T[p^k]=\{x\in T:p^kx=0\}$ dla $k\ge1$ oraz
$T[p^0]=0$. Iloraz $T[p^k]/T[p^{k-1}]$ jest
przestrzenią wektorową nad ciałem $R/(p)$:
mnożenie przez $p$ wysyła $T[p^k]$ do $T[p^{k-1}]$,
więc działa na ilorazie jako zero.
Dla pojedynczego czynnika $L=R/(p^e)$, gdzie $e\ge1$,
zbiór $L[p^k]$ jest, dla $1\le k\le e$, podmodułem
generowanym przez klasę $p^{e-k}$. Istotnie,
$p^k(x+(p^e))=0$ dokładnie wtedy, gdy $p^e\mid p^kx$;
przez skracanie w dziedzinie otrzymujemy
$p^{e-k}\mid x$. Dla $k>e$ mamy $L[p^k]=L$.
Gdy $1\le k\le e$, odwzorowanie
\[
 R/(p)\longrightarrow L[p^k]/L[p^{k-1}],\qquad
 c+(p)\longmapsto cp^{e-k}+L[p^{k-1}]
\]
jest dobrze określone: zastąpienie $c$ przez $c+pa$
zmienia reprezentanta o $ap^{e-k+1}$, który należy do
$L[p^{k-1}]$. Odwzorowanie jest na z opisu $L[p^k]$
i ma jądro zerowe, ponieważ $cp^{e-k}$ należy do
$L[p^{k-1}]$ dokładnie wtedy, gdy $p\mid c$.
Jest więc izomorfizmem $R/(p)$-przestrzeni wektorowych.
Wobec tego rozważany iloraz ma wymiar jeden dla
$1\le k\le e$ oraz zero dla $k>e$.
Każdy $R/(d_i)$ rozkłada się na sumę składników
$R/(p^{e_{i,p}})$ po rozłożeniu $d_i$ na potęgi
parami niestowarzyszonych elementów nierozkładalnych
i zastosowaniu chińskiego twierdzenia o resztach
(Twierdzenie~\ref{aa:tw:crt-ring}): ideały
$(p^e)$ i $(q^f)$ dla niestowarzyszonych
nierozkładalnych $p,q$ są względnie pierwsze.
Istotnie, ich suma ma postać $(g)$, a $g$ dzieli
$p^e$ i $q^f$, więc na mocy jednoznaczności
rozkładu $g$ jest jednostką. Zatem
$(p^e)+(q^f)=R$. Jeżeli $p\nmid d_i$, to z jednoznaczności rozkładu na
czynniki mamy $(p^k,d_i)=R$. Istnieją zatem $u,v\in R$
takie, że $up^k+vd_i=1$. Jeśli $p^k$ zeruje klasę
$x+(d_i)$, to po pomnożeniu tej tożsamości przez $x$
widzimy, że $x\in(d_i)$; stąd $R/(d_i)$ nie ma
niezerowych elementów zerowanych przez $p^k$.
W konsekwencji liczba
\[
 h_{p,k}=\dim_{R/(p)}\bigl(T[p^k]/T[p^{k-1}]\bigr)
\]
jest równa liczbie tych indeksów $i$, dla których
$e_{i,p}\ge k$. Wartości wszystkich $h_{p,k}$ są
niezmiennikami modułu, ponieważ $T[p^k]$ definiujemy
bez użycia wybranego rozkładu, a każdy izomorfizm
modułów zachowuje te podmoduły i wywołuje izomorfizmy
ich kolejnych ilorazów. Pozwalają one odtworzyć liczbę
wystąpień wykładnika dokładnie $k$ jako
$h_{p,k}-h_{p,k+1}$.

Dla każdego $p$ odzyskujemy więc multizbiór wykładników
$e_{i,p}>0$. Warunek $d_1\mid\cdots\mid d_t$ oznacza,
że dla każdego $p$ wykładniki w kolejnych $d_i$
są niemalejące; zerowe wykładniki dopełniamy od lewej.
Jeśli $T(M)=0$, to $t=0$ i nie ma czynników torsyjnych.
W przeciwnym przypadku $t$ jest największą liczbą
dodatnich wykładników wśród różnych $p$: ponieważ
$d_1$ nie jest jednostką, ma przynajmniej jeden
czynnik nierozkładalny $p$, który dzieli wszystkie
$d_i$; dla tego $p$ wszystkie $t$ wykładników są dodatnie.
Tak ustalone ciągi wykładników określają wszystkie $d_i$
z dokładnością do jednostek pierścienia. To kończy dowód.
\end{proof}

\begin{wniosek}[Klasyfikacja skończenie generowanych grup abelowych]
\label{aa:wn:abelian-fg-structure}
Każda skończenie generowana grupa abelowa $A$ jest izomorficzna z
$\mathbb Z^r\oplus C_{n_1}\oplus\cdots\oplus C_{n_t}$,
gdzie $n_i\ge2$ i $n_1\mid\cdots\mid n_t$.
Liczba $r$ i wartości $n_i$ są jednoznaczne.
\end{wniosek}
\begin{proof}
Grupy abelowe są dokładnie modułami nad $\mathbb Z$.
Ponieważ $\mathbb Z$ jest dziedziną ideałów głównych,
stosujemy Twierdzenie~\ref{aa:tw:pid-structure}.
W $\mathbb Z$ jednostkami są $1$ i $-1$,
więc współczynniki $d_i$ można wybrać dodatnie.
Otrzymane czynniki $\mathbb Z/(d_i)$ to grupy cykliczne
$C_{d_i}$, co daje tezę.
\end{proof}

\begin{przyklad}[Grupy abelowe rzędu $72$]
\label{aa:prz:abelian-order72}
Ponieważ $72=2^3\cdot3^2$, najpierw rozkładamy grupę
na jej części $2$-pierwotną i $3$-pierwotną.
Dla liczby $8$ możliwe rozkłady potęgowo-pierwotne to
$C_8$, $C_4\oplus C_2$ i $C_2^3$; odpowiadają one
podziałom wykładnika $3$: $3$, $2+1$ i $1+1+1$.
Dla liczby $9$ są dwie możliwości: $C_9$ oraz $C_3^2$.
Wobec jednoznaczności klasyfikacji istnieje dokładnie
$3\cdot2=6$ klas izomorfizmu grup abelowych rzędu $72$.
W postaci czynników niezmienniczych reprezentują je
$C_{72}$, $C_3\oplus C_{24}$, $C_2\oplus C_{36}$,
$C_6\oplus C_{12}$, $C_2\oplus C_2\oplus C_{18}$ oraz
$C_2\oplus C_6\oplus C_6$. Każdy z tych rozkładów
spełnia wymagany warunek podzielności.
Przykładowo z chińskiego twierdzenia o resztach mamy
$C_4\oplus C_2\oplus C_9\cong C_2\oplus C_{36}$.
Ta grupa ma cztery elementy o rzędzie dzielącym $2$,
podczas gdy $C_{72}$ ma tylko dwa, więc grupy
nie są izomorficzne.
\end{przyklad}

\begin{przyklad}[Dlaczego hipotezy twierdzenia są potrzebne?]
Moduł $\mathbb Z/2\mathbb Z$ nie jest wolny nad $\mathbb Z$,
ale jest sumą czynników cyklicznych zgodnie z twierdzeniem.
Nad pierścieniem $S=K[t_1,t_2]$, który nie jest dziedziną
ideałów głównych, ideał $I=(t_1,t_2)$ jest skończenie generowanym
modułem, ale nie jest wolny. Rzeczywiście, gdyby $I$ miał bazę zawierającą
przynajmniej dwa elementy $u,v$, to niezerowa relacja
$vu-uv=0$ (ze współczynnikami $v,-u\in S$)
przeczyłaby liniowej niezależności tej bazy.
Niezerowy wolny podmoduł $I\subset S$ musiałby zatem
mieć rangę jeden, czyli być generowany przez
jeden element jako ideał. Tymczasem wspólnymi dzielnikami
$t_1$ oraz $t_2$ w $S$ są tylko jednostki, a
$1\notin(t_1,t_2)$, co widać po podstawieniu
$t_1=t_2=0$. Ideał $I$ nie jest zatem główny.
Nie wolno bez dodatkowych założeń przenosić twierdzenia
o strukturze na wszystkie pierścienie przemienne.
\end{przyklad}

\section{Dalszy związek z algebrą liniową}
Dla $K$ będącego ciałem pierścień $K[t]$ jest dziedziną
euklidesową, a więc dziedziną ideałów głównych. Jeśli
$V$ jest skończenie wymiarową przestrzenią wektorową nad $K$
i $T:V\to V$ jest operatorem liniowym, możemy traktować $V$
jako $K[t]$-moduł przez $t\cdot v=T(v)$.
Wówczas postać Smitha i twierdzenie o strukturze modułów
prowadzą do postaci kanonicznej Frobeniusa oraz do teorii
czynników niezmienniczych operatora. Pełne dowody tych
rezultatów umieścimy w późniejszym skrypcie algebry liniowej.

% AA-18 | Algebra abstrakcyjna | Iteracja IX
\chapter{Iloczyn tensorowy modułów i jego własności}
\label{aa:ch:module-tensor-product}

\section{Dlaczego potrzebujemy nowej konstrukcji?}
Rozważmy dwie grupy abelowe $M,N$. Iloczyn prosty $M\times N$
przechowuje po prostu parę ich elementów. Iloczyn tensorowy jest
czymś innym: pozwala zamienić \emph{dwuliniowe} odwzorowania
z $M\times N$ na homomorfizmy z jednej grupy. W szczególności
iloczyn tensorowy dwóch niezerowych modułów może być zerowy.

Ważne jest rozróżnienie trzech konstrukcji. W iloczynie prostym
występują uporządkowane pary $(m,n)$. W sumie prostej występują
rodziny elementów o skończonym nośniku. W iloczynie tensorowym
występują klasy formalnych symboli $m\otimes n$, podlegające
konkretnym relacjom addytywności i zgodności ze skalarami.
Nie utożsamiamy żadnej z tych konstrukcji z pozostałymi.

Jeśli pierścień $R$ nie jest przemienny, strony działania
skalarów mają znaczenie. Do podstawowej konstrukcji nad $R$
potrzebujemy \emph{prawego} $R$-modułu $M$ i \emph{lewego}
$R$-modułu $N$. Otrzymany iloczyn tensorowy jest w ogólności
grupą abelową; bez dodatkowych struktur nie musi mieć
kanonicznego działania pierścienia $R$.

\section{Odwzorowania zbilansowane i konstrukcja ilorazowa}
\begin{definicja}[Odwzorowanie addytywne i zbilansowane]
\label{aa:def:tensor-balanced}
Niech $R$ będzie pierścieniem z jedynką, niekoniecznie
przemiennym, $M$ prawym $R$-modułem, $N$ lewym $R$-modułem,
a $A$ grupą abelową. Odwzorowanie $b:M\times N\to A$
nazywamy \emph{addytywnym w każdej zmiennej i zbilansowanym},
gdy dla wszystkich $m,m'\in M$, $n,n'\in N$ i $r\in R$ zachodzą
\[
\begin{split}
b(m+m',n)&=b(m,n)+b(m',n),\\
b(m,n+n')&=b(m,n)+b(m,n'),\\
b(mr,n)&=b(m,rn).
\end{split}
\]
Trzeci warunek pozwala przenieść skalar działający
z prawej strony na $M$ do działania z lewej strony na $N$.
Jest konieczny do zdefiniowania iloczynu tensorowego nad $R$.
\end{definicja}

\begin{twierdzenie}[Istnienie iloczynu tensorowego i własność uniwersalna]
\label{aa:tw:tensor-universal}
Niech $R$ będzie pierścieniem z jedynką, niekoniecznie
przemiennym, $M$ prawym $R$-modułem, a $N$ lewym
$R$-modułem. Wtedy istnieje grupa abelowa $M\otimes_R N$ oraz odwzorowanie
$\tau:M\times N\to M\otimes_R N$ addytywne w każdej
zmiennej i zbilansowane, takie, że dla każdej grupy abelowej
$A$ i każdego odwzorowania $b:M\times N\to A$ o tych
własnościach istnieje dokładnie jeden homomorfizm grup
$\bar b:M\otimes_R N\to A$ z $\bar b\circ\tau=b$.
Ponadto para $(M\otimes_RN,\tau)$ jest wyznaczona
jednoznacznie z dokładnością do jedynego izomorfizmu
zgodnego z $\tau$.
\end{twierdzenie}
\begin{proof}
Niech $F$ będzie \emph{wolną grupą abelową} na zbiorze
symboli $[m,n]$, gdzie $(m,n)\in M\times N$.
Oznacza to, że elementy $F$ są formalnymi skończonymi
sumami całkowitych wielokrotności tych symboli; różne
symbole są niezależnymi generatorami.

Niech $L$ będzie podgrupą $F$ generowaną przez elementy
\[
\begin{split}
 &[m+m',n]-[m,n]-[m',n],\\
 &[m,n+n']-[m,n]-[m,n'],\\
 &[mr,n]-[m,rn]
\end{split}
\]
dla wszystkich dopuszczalnych $m,m',n,n',r$.
Definiujemy $M\otimes_RN=F/L$ i
$m\otimes n=[m,n]+L$. Wymienione relacje zapewniają,
że $(m,n)\mapsto m\otimes n$ spełnia dokładnie
trzy wymagane warunki. Z addytywności wynika też
$0\otimes n=m\otimes0=0$ oraz
$(-m)\otimes n=-(m\otimes n)$.

Niech teraz $b:M\times N\to A$ będzie jak w tezie.
Ponieważ $F$ jest wolną grupą abelową na zbiorze par,
istnieje dokładnie jeden homomorfizm $B:F\to A$
taki, że $B([m,n])=b(m,n)$. Trzy założenia o $b$
mówią, że $B$ zeruje wszystkie generatory $L$.
Zatem $L\subseteq\ker B$ i wzór
$\bar b(x+L)=B(x)$ określa homomorfizm na ilorazie.
Z definicji $\bar b(m\otimes n)=b(m,n)$.
Każdy element $F/L$ jest skończoną sumą całkowitych
wielokrotności tensorów prostych; stąd homomorfizm
o wskazanych wartościach jest jednoznaczny.

Jeśli $(T,t)$ jest inną parą o tej samej własności,
uniwersalność daje homomorfizmy $T\to M\otimes_RN$
oraz $M\otimes_RN\to T$, zachowujące obrazy wszystkich
par $(m,n)$. Ich złożenia zachowują te same generatory,
więc są identycznościami. Homomorfizmy są odwrotne
i dają wymagany jednoznaczny izomorfizm.
\end{proof}

\begin{uwaga}[Oznaczenie tensoru prostego]
Element $m\otimes n$ nazywamy \emph{tensorem prostym}.
Nie każdy element $M\otimes_RN$ musi być tensorem prostym:
na ogół jest sumą takich elementów. Na przykład dla dowolnego ciała $K$, przestrzeni
$V=W=K^2$ i baz $e_1,e_2$ oraz $f_1,f_2$ tensor
$e_1\otimes f_1+e_2\otimes f_2$ nie jest prosty.
Gdyby był równy $(ae_1+be_2)\otimes(cf_1+df_2)$,
to porównanie współczynników przy czterech tensorach
bazowych dawałoby $ac=bd=1$ oraz $ad=bc=0$.
Pierwsze dwie równości wymuszają $a,b,c,d\ne0$,
co przeczy pozostałym dwóm. Wykorzystaliśmy tu
niezależność tensorów bazowych, udowodnioną
w Stwierdzeniu~\ref{aa:st:tensor-vspace-dimension}.
Równość $m\otimes n=m'\otimes n'$ nie oznacza, że
$m=m'$ i $n=n'$. W szczególności z $m\otimes n=0$
nie wynika, że któryś z czynników jest zerowy;
porównajmy Przykład~\ref{aa:prz:tensor-cyclic}.
\end{uwaga}

\section{Przemienny pierścień współczynników}
\begin{twierdzenie}[Iloczyn tensorowy jako moduł]
\label{aa:tw:tensor-commutative-module}
Niech $R$ będzie pierścieniem przemiennym z jedynką,
a $M,N$ lewymi $R$-modułami. Traktując $M$ jako prawy
moduł przez $mr=rm$, otrzymujemy $R$-moduł
$M\otimes_R N$ z działaniem
$r(m\otimes n)=(rm)\otimes n=m\otimes(rn)$.
Dla każdego $R$-modułu $P$ odwzorowania $R$-dwuliniowe
$M\times N\to P$ odpowiadają bijektywnie homomorfizmom
$R$-modułów $M\otimes_RN\to P$.
\end{twierdzenie}
\begin{proof}
Dla $r\in R$ na wolnej grupie abelowej $F$ z poprzedniego
dowodu określmy $T_r([m,n])=[rm,n]$ i przedłużmy
addytywnie. Sprawdzimy, że $T_r(L)\subseteq L$.
Dla pierwszego rodzaju relacji addytywnych otrzymujemy
$[r(m+m'),n]-[rm,n]-[rm',n]$; z równości
$r(m+m')=rm+rm'$ wynika, że jest to generator
pierwszego rodzaju. Dla drugiego rodzaju
otrzymujemy bezpośrednio relację
$[rm,n+n']-[rm,n]-[rm,n']$. Dla relacji zbilansowania
obraz elementu $[ms,n]-[m,sn]$ wynosi
$[r(ms),n]-[rm,sn]=[(rm)s,n]-[rm,sn]$,
ponieważ $R$ jest przemienny. Jest to generator $L$.
Zatem $T_r$ przechodzi na grupę ilorazową.
Wzory $T_{r+s}=T_r+T_s$, $T_{rs}=T_rT_s$
i $T_1=\operatorname{id}$ sprawdzamy na generatorach
$[m,n]$, używając odpowiednio addytywności działania,
$(rs)m=r(sm)$ oraz $1_Rm=m$. Przechodzą one na
iloraz i dają wszystkie aksjomaty $R$-modułu.
Równość $(rm)\otimes n=m\otimes(rn)$ jest
jedną z relacji zbilansowania.

Jeśli $b:M\times N\to P$ jest $R$-dwuliniowe,
jest addytywne w każdej zmiennej i zbilansowane,
gdyż $b(rm,n)=rb(m,n)=b(m,rn)$.
Własność uniwersalna daje homomorfizm grup
$\bar b:M\otimes_RN\to P$. Na tensorach prostych
$\bar b(r(m\otimes n))=b(rm,n)=rb(m,n)$,
więc $\bar b$ jest $R$-liniowy. Odwrotnie, złożenie
$R$-liniowego homomorfizmu z kanonicznym odwzorowaniem
$(m,n)\mapsto m\otimes n$ jest $R$-dwuliniowe.
Obie konstrukcje są wzajemnie odwrotne.
\end{proof}

\begin{przyklad}[Iloczyn tensorowy grup abelowych]
\label{aa:prz:tensor-abelian-example}
Każda grupa abelowa jest $\mathbb Z$-modułem.
Zatem dla grup abelowych $A,B$ zapis
$A\otimes_{\mathbb Z}B$ oznacza ich iloczyn tensorowy
jako modułów nad $\mathbb Z$. Jest to dokładnie
ta sama konstrukcja, którą wprowadziliśmy w
Rozdziale~\ref{aa:ch:grupy-abelowe}, teraz
umieszczona w ogólnej teorii modułów.
\end{przyklad}

\section{Podstawowe izomorfizmy}
\begin{twierdzenie}[Element neutralny dla iloczynu tensorowego]
\label{aa:tw:tensor-unit}
Niech $R$ będzie pierścieniem z jedynką, a $N$
lewym $R$-modułem. Odwzorowanie
$R\otimes_RN\to N$, $r\otimes n\mapsto rn$,
jest izomorfizmem grup abelowych. Jeśli $R$ jest
przemienny, jest też izomorfizmem $R$-modułów.
Analogicznie dla prawego $R$-modułu $M$ mamy
$M\otimes_RR\cong M$.
\end{twierdzenie}
\begin{proof}
Odwzorowanie $(r,n)\mapsto rn$ jest addytywne
w każdej zmiennej i zbilansowane: $(rs)n=r(sn)$.
Uniwersalność daje homomorfizm $F:R\otimes_RN\to N$.
Zdefiniujmy $G:N\to R\otimes_RN$ przez $G(n)=1\otimes n$.
Jest to homomorfizm grup. Wtedy
$F(G(n))=1n=n$, natomiast
$G(F(r\otimes n))=1\otimes(rn)=r\otimes n$.
Ostatnia równość wynika z relacji zbilansowania.
Ponieważ tensory proste generują całą grupę,
$G\circ F=\operatorname{id}$. Dla przemiennego $R$
oba homomorfizmy zachowują działanie skalarów.
Dla prawego $R$-modułu $M$ druga teza wynika wprost
z odwzorowań $m\otimes r\mapsto mr$ oraz
$m\mapsto m\otimes1_R$. Ich złożenia są identycznościami,
ponieważ $(mr)\otimes1_R=m\otimes(r1_R)=m\otimes r$.
\end{proof}

\begin{twierdzenie}[Suma prosta i iloczyn tensorowy]
\label{aa:tw:tensor-direct-sum}
Niech $R$ będzie przemiennym pierścieniem z jedynką,
a $M,N,P$ jego modułami. Wówczas istnieje kanoniczny izomorfizm
\[
 (M\oplus N)\otimes_R P\cong
 (M\otimes_RP)\oplus(N\otimes_RP).
\]
Taki sam wzór zachodzi dla dowolnej rodziny modułów
z sumą prostą po lewej i sumą prostą ich iloczynów
tensorowych po prawej.
\end{twierdzenie}
\begin{proof}
Dwuliniowe odwzorowanie
$((m,n),p)\mapsto(m\otimes p,n\otimes p)$
indukuje homomorfizm $F$ z lewej strony na prawą.
W odwrotnym kierunku definiujemy homomorfizmy
$M\otimes_RP\to(M\oplus N)\otimes_RP$ przez
$m\otimes p\mapsto(m,0)\otimes p$ oraz
$N\otimes_RP\to(M\oplus N)\otimes_RP$ przez
$n\otimes p\mapsto(0,n)\otimes p$.
Są one poprawne dzięki własności uniwersalnej.
Ich suma daje homomorfizm $G$ z prawej strony na lewą.
Na tensorach prostych mamy
$GF((m,n)\otimes p)=(m,0)\otimes p+(0,n)\otimes p
=(m,n)\otimes p$.
Z drugiej strony $FG$ zachowuje osobno oba składniki,
więc jest identycznością. Dla dowolnej rodziny $(M_i)_{i\in I}$ definiujemy
$F((m_i)_{i\in I}\otimes p)=(m_i\otimes p)_{i\in I}$.
Wyrażenie po prawej ma skończony nośnik, gdyż
$(m_i)$ należy do sumy prostej. Jest dwuliniowe,
więc daje homomorfizm z iloczynu tensorowego.
W przeciwną stronę każda kanoniczna inkluzja
$M_i\hookrightarrow\bigoplus_{j\in I}M_j$ daje
homomorfizm $M_i\otimes_RP\to
(\bigoplus_{j\in I}M_j)\otimes_RP$.
Dla rodziny tensorów o skończonym nośniku sumujemy
ich obrazy po tych inkluzjach; otrzymujemy $G$.
Na tensorach prostych obie kompozycje są
identycznościami, a na prawej sumie prostej
sprawdzamy to oddzielnie na każdym składniku.
Zatem $F$ i $G$ są wzajemnie odwrotne.
\end{proof}

\begin{twierdzenie}[Przemienność i łączność iloczynu tensorowego]
\label{aa:tw:tensor-sym-assoc}
Niech $R$ będzie przemiennym pierścieniem z jedynką,
a $M,N,P$ jego modułami. Istnieją kanoniczne izomorfizmy
\[
 M\otimes_RN\cong N\otimes_RM,
 \qquad (M\otimes_RN)\otimes_RP\cong
 M\otimes_R(N\otimes_RP),
\]
wyznaczone odpowiednio przez $m\otimes n\mapsto n\otimes m$
oraz $(m\otimes n)\otimes p\mapsto m\otimes(n\otimes p)$.
\end{twierdzenie}
\begin{proof}
Odwzorowanie $(m,n)\mapsto n\otimes m$ jest
$R$-dwuliniowe, więc daje homomorfizm dla pierwszego
izomorfizmu. To samo odwzorowanie w odwrotnej kolejności
jest jego odwrotnością na generatorach.

Przy drugim izomorfizmie sprawdzenie poprawności wymaga
dwóch kolejnych zastosowań własności uniwersalnej.
Dla ustalonego $p$ odwzorowanie $(m,n)\mapsto
m\otimes(n\otimes p)$ jest $R$-dwuliniowe, więc
daje homomorfizm $\alpha_p:M\otimes_RN\to
M\otimes_R(N\otimes_RP)$. Dla ustalonego $x$ odwzorowanie
$p\mapsto\alpha_p(x)$ jest liniowe: dla
$x=m\otimes n$ wynika to z liniowości
$p\mapsto n\otimes p$, a ogólny $x$ jest sumą
skończenie wielu tensorów prostych. Podobnie
$\alpha_{rp}(m\otimes n)=r\alpha_p(m\otimes n)$
i $\alpha_p(rx)=r\alpha_p(x)$, zatem
$(x,p)\mapsto\alpha_p(x)$ jest $R$-dwuliniowe.
Przez kolejne zastosowanie własności uniwersalnej
otrzymujemy homomorfizm
$(M\otimes_RN)\otimes_RP\to M\otimes_R(N\otimes_RP)$.
W odwrotną stronę konstruujemy analogicznie
$m\otimes(n\otimes p)\mapsto(m\otimes n)\otimes p$.
Na elementach $(m\otimes n)\otimes p$ i
$m\otimes(n\otimes p)$ złożenia tych homomorfizmów
są identycznościami; ponieważ generują one odpowiednie
moduły, homomorfizmy są wzajemnie odwrotne.
\end{proof}

\section{Iloczyn tensorowy z modułem ilorazowym}
\begin{twierdzenie}[Iloczyn tensorowy z ilorazem pierścienia]
\label{aa:tw:tensor-quotient}
Niech $R$ będzie przemiennym pierścieniem z jedynką,
$I\subseteq R$ ideałem, a $N$ lewym $R$-modułem.
Przez $IN$ oznaczamy podmoduł generowany przez wszystkie
$in$, gdzie $i\in I$ i $n\in N$. Wtedy istnieje naturalny
izomorfizm $R$-modułów
\[
 (R/I)\otimes_R N\cong N/IN,
 \qquad (r+I)\otimes n\longmapsto rn+IN.
\]
\end{twierdzenie}
\begin{proof}
Definiujemy odwzorowanie $b:(R/I)\times N\to N/IN$
przez $b(r+I,n)=rn+IN$. Jeśli $r'=r+i$ dla $i\in I$,
to $(r'-r)n=in\in IN$, więc wzór nie zależy od
reprezentanta $r+I$. Zachowuje dodawanie w obu
zmiennych i spełnia $b((r+I)s,n)=b(r+I,sn)$.
Dlatego daje homomorfizm $F$ z iloczynu tensorowego
na $N/IN$.

W drugą stronę definiujemy $G:N/IN\to (R/I)\otimes_RN$
przez $G(n+IN)=(1+I)\otimes n$.
Jeżeli $n'=n+\sum_j i_j n_j$ z $i_j\in I$, to
$G(n')-G(n)=\sum_j(1+I)\otimes i_jn_j
=\sum_j(i_j+I)\otimes n_j=0$.
Zatem $G$ jest dobrze określony i $R$-liniowy.
Dla $n\in N$ mamy $FG(n+IN)=n+IN$.
Dla generatora $(r+I)\otimes n$ z drugiej strony
$GF((r+I)\otimes n)=(1+I)\otimes rn
=(r+I)\otimes n$.
Obie kompozycje są więc identycznościami.
\end{proof}

\begin{wniosek}[Dwa ilorazy pierścienia]
\label{aa:wn:tensor-two-quotients}
Niech $R$ będzie przemiennym pierścieniem z jedynką,
a $I,J\subseteq R$ jego ideałami. Wówczas istnieje
naturalny izomorfizm $R$-modułów
$(R/I)\otimes_R(R/J)\cong R/(I+J)$.
\end{wniosek}
\begin{proof}
Z Twierdzenia~\ref{aa:tw:tensor-quotient} dostajemy
$(R/I)\otimes_R(R/J)\cong(R/J)/I(R/J)$.
Obrazem podmodułu $I(R/J)$ są dokładnie klasy
elementów ideału $I+J$ modulo $J$.
Trzecie twierdzenie o izomorfizmie modułów daje
$(R/J)/((I+J)/J)\cong R/(I+J)$.
\end{proof}

\begin{przyklad}[Iloczyn tensorowy grup cyklicznych]
\label{aa:prz:tensor-cyclic}
Dla liczb całkowitych $m,n\ge1$ oraz $d=\gcd(m,n)$
zachodzi
\[
 (\mathbb Z/m\mathbb Z)\otimes_{\mathbb Z}
 (\mathbb Z/n\mathbb Z)\cong \mathbb Z/d\mathbb Z.
\]
Rzeczywiście $(m)+(n)=(d)$, więc wystarczy
zastosować Wniosek~\ref{aa:wn:tensor-two-quotients}.
Dla $m=12$, $n=18$ otrzymujemy grupę cykliczną
rzędu sześć. Natomiast dla $m=2$, $n=3$ iloczyn
tensorowy jest zerowy. Jest to wyraźnie różna sytuacja
od iloczynu prostego $C_2\times C_3\cong C_6$.
\end{przyklad}

\begin{przyklad}[Iloczyn tensorowy z ciałem liczb wymiernych]
\label{aa:prz:tensor-rational-torsion}
Dla $n\ge1$ mamy $\mathbb Q\otimes_{\mathbb Z}
(\mathbb Z/n\mathbb Z)=0$. Istotnie elementy
$\mathbb Q$ można dzielić przez $n$, więc
$n\mathbb Q=\mathbb Q$. Z Twierdzenia
\ref{aa:tw:tensor-quotient}, po zamianie czynników,
otrzymujemy $(\mathbb Z/n\mathbb Z)\otimes_{\mathbb Z}
\mathbb Q\cong\mathbb Q/n\mathbb Q=0$.
Nie oznacza to, że $\mathbb Q$ lub $\mathbb Z/n\mathbb Z$
są modułami zerowymi; zerowy jest ich iloczyn tensorowy.
\end{przyklad}

\section{Iloczyn tensorowy i krótkie ciągi dokładne}
\begin{twierdzenie}[Prawostronna dokładność funktora iloczynu tensorowego]
\label{aa:tw:tensor-right-exact}
Niech $R$ będzie przemiennym pierścieniem z jedynką,
a $N$ jego modułem. Jeśli ciąg homomorfizmów
$A\xrightarrow{f}B\xrightarrow{g}C\to0$ jest dokładny,
to dokładny jest również ciąg
\[
 A\otimes_RN\xrightarrow{f\otimes\operatorname{id}_N}
 B\otimes_RN\xrightarrow{g\otimes\operatorname{id}_N}
 C\otimes_RN\longrightarrow0.
\]
Dla dowolnego homomorfizmu $f$ odwzorowanie
$f\otimes\operatorname{id}_N$ jest określone przez
$a\otimes n\mapsto f(a)\otimes n$.
\end{twierdzenie}
\begin{proof}
Najpierw poprawność ostatniego wzoru wynika z własności
uniwersalnej, ponieważ $(a,n)\mapsto f(a)\otimes n$
jest dwuliniowe. Z dokładności pierwotnego ciągu
mamy $g$ surjektywne i $\ker g=\operatorname{im}f$.
Każdy tensor prosty $c\otimes n\in C\otimes_RN$
ma przedobraz $b\otimes n$ z $g(b)=c$; ponieważ
tensory proste generują, odwzorowanie
$g\otimes\operatorname{id}_N$ jest surjektywne.

Oznaczmy $I=\operatorname{im}f\subseteq B$.
Z pierwszego twierdzenia o izomorfizmie modułów
mamy $C\cong B/I$. Rozważmy teraz iloraz
$D=(B\otimes_RN)/H$, gdzie $H$ jest podmodułem
generowanym przez wszystkie $i\otimes n$ z
$i\in I$, $n\in N$. Odwzorowanie
$(b+I,n)\mapsto b\otimes n+H$
jest dobrze określone: zmiana $b$ o $i\in I$
zmienia jego wartość o $i\otimes n\in H$.
Jest dwuliniowe, więc otrzymujemy
$(B/I)\otimes_RN\to D$.
Z drugiej strony odwzorowanie
$b\otimes n\mapsto(b+I)\otimes n$ zeruje $H$
i daje homomorfizm $D\to(B/I)\otimes_RN$.
Oba złożenia zachowują tensory proste, a zatem są
identycznościami. Tak więc
$D\cong(B/I)\otimes_RN\cong C\otimes_RN$.
Oznacza to, że jądro $g\otimes\operatorname{id}_N$
jest dokładnie $H$. Ale $H$ jest obrazem
$f\otimes\operatorname{id}_N$: obraz tego homomorfizmu
jest podmodułem, zawiera każdy generator $i\otimes n$
modułu $H$, gdyż $i=f(a)$ dla pewnego $a\in A$,
a jednocześnie obrazy wszystkich tensorów prostych
$a\otimes n$ należą do $H$. Ponieważ tensory proste
generują $A\otimes_RN$, oba podmoduły są równe.
Dostaliśmy żądaną dokładność.
\end{proof}

\begin{przyklad}[Dlaczego iniekcja może zniknąć]
\label{aa:prz:tensor-not-left-exact}
Rozważmy krótki ciąg dokładny $\mathbb Z$-modułów
\[
0\longrightarrow\mathbb Z\xrightarrow{\cdot2}\mathbb Z
\longrightarrow\mathbb Z/2\mathbb Z\longrightarrow0.
\]
Po wzięciu iloczynu tensorowego z $\mathbb Z/2\mathbb Z$
i użyciu $\mathbb Z\otimes_{\mathbb Z}\mathbb Z/2\mathbb Z
\cong\mathbb Z/2\mathbb Z$ otrzymujemy odwzorowanie
$\mathbb Z/2\mathbb Z\xrightarrow{\cdot2}
\mathbb Z/2\mathbb Z$. Jest ono zerowe, ponieważ
$(2a)\otimes\bar b=a\otimes(2\bar b)=0$ dla
każdych $a\in\mathbb Z$, $\bar b\in\mathbb Z/2\mathbb Z$.
Nie jest ono iniekcją. W ogólności \emph{iloczyn tensorowy
nie zachowuje iniektywności}, choć zachowuje surjektywność.
To odróżnia go również od iloczynu prostego modułów.
\end{przyklad}

\section{Moduły płaskie i moduły rzutowe}
\begin{definicja}[Moduł płaski]
\label{aa:def:flat-module}
Niech $R$ będzie przemiennym pierścieniem z jedynką.
$R$-moduł $N$ nazywamy \emph{płaskim}, jeżeli dla
każdego krótkiego ciągu dokładnego
$0\to A\to B\to C\to0$ ciąg uzyskany przez
wzięcie iloczynu tensorowego z $N$,
\[
0\longrightarrow A\otimes_RN\longrightarrow B\otimes_RN
\longrightarrow C\otimes_RN\longrightarrow0,
\]
jest także dokładny.
\end{definicja}

\begin{twierdzenie}[Moduły wolne i rzutowe są płaskie]
\label{aa:tw:projective-flat}
Niech $R$ będzie przemiennym pierścieniem z jedynką.
Każdy wolny $R$-moduł jest płaski; każdy składnik sumy
prostej modułu płaskiego jest płaski. W szczególności
każdy moduł rzutowy jest płaski.
\end{twierdzenie}
\begin{proof}
Jeśli $F=R^{(S)}$ jest modułem wolnym na zbiorze $S$,
rozdzielność iloczynu tensorowego względem sum prostych
z Twierdzenia~\ref{aa:tw:tensor-direct-sum} i izomorfizm
$M\otimes_RR\cong M$ pokazują, że
$M\otimes_RF\cong M^{(S)}$.
Suma prosta kopii krótkiego ciągu dokładnego jest
dokładna współrzędnie: element sumy prostej ma skończony
nośnik, więc badanie jąder i obrazów odbywa się w skończonej
liczbie współrzędnych. Zatem $F$ jest płaski.

Załóżmy, że $N=P\oplus Q$ jest płaski. Dla każdej
iniekcji $f:A\to B$ homomorfizm
$f\otimes\operatorname{id}_N$ jest iniektywny.
Pod identyfikacjami
$A\otimes_RN\cong(A\otimes_RP)\oplus(A\otimes_RQ)$
i odpowiednio dla $B$ jest on sumą prostą
$f\otimes\operatorname{id}_P$ oraz
$f\otimes\operatorname{id}_Q$. Każda z tych map
musi być iniektywna, ponieważ ich suma prosta jest
iniektywna. W połączeniu z prawostronną dokładnością
dowodzi to płaskości $P$ oraz $Q$.
W Rozdziale~\ref{aa:ch:modules-introduction} udowodniliśmy,
że moduł rzutowy jest składnikiem sumy prostej modułu
wolnego; dlatego każdy moduł rzutowy jest płaski.
\end{proof}

\begin{przyklad}[Moduł niepłaski]
Moduł $\mathbb Z/2\mathbb Z$ nie jest płaski nad
$\mathbb Z$, ponieważ przykład
\ref{aa:prz:tensor-not-left-exact} pokazuje,
że iloczyn tensorowy z nim nie zachowuje iniekcji.
\end{przyklad}

\begin{twierdzenie}[Ciało ułamków dziedziny całkowitości jest modułem płaskim]
\label{aa:tw:frac-flat}
Niech $R$ będzie dziedziną całkowitości i $K$
jej ciałem ułamków. Wtedy $K$ jest płaskim
$R$-modułem.
\end{twierdzenie}
\begin{proof}
Użyjemy jawnego opisu lokalizacji. Dla dowolnego
$R$-modułu $M$ niech $S=R\setminus\{0\}$.
Na parach $(m,s)\in M\times S$ przyjmujemy
$(m,s)\sim(m',s')$ wtedy i tylko wtedy, gdy
istnieje $t\in S$ takie, że $t(s'm-sm')=0$.
Jest to relacja równoważności: zwrotność i symetria
są natychmiastowe. Dla przechodniości załóżmy, że
$t(s'm-sm')=0$ i $u(s''m'-s'm'')=0$.
Po pomnożeniu pierwszej równości przez $us''$,
a drugiej przez $ts$, i dodaniu otrzymujemy
$tu s'(s''m-sm'')=0$; element $tu s'$ należy do $S$.

Na zbiorze klas $S^{-1}M$ określamy działania
$[m,s]+[m',s']=[s'm+sm',ss']$ oraz
$r[m,s]=[rm,s]$. Są one niezależne od reprezentantów.
Istotnie, z $t(s_1m-sm_1)=0$ i
$u(s'_1m'-s'm'_1)=0$ wynika
\[
\begin{aligned}
&s_1s'_1(s'm+sm')-ss'(s'_1m_1+s_1m'_1)\\
&\quad=s'_1s'(s_1m-sm_1)+s_1s(s'_1m'-s'm'_1).
\end{aligned}
\]
Po pomnożeniu prawej strony przez $tu\in S$
otrzymujemy zero. Jest to dokładnie warunek równoważności
sum dwóch par. Analogicznie $t(s_1rm-srm_1)=0$
uzasadnia poprawność mnożenia klasy przez $r$.
Łączność dodawania wynika z tego, że suma trzech
klas, niezależnie od nawiasowania, ma przedstawiciela
$[s_2s_3m_1+s_1s_3m_2+s_1s_2m_3,s_1s_2s_3]$.
Przemienność dodawania wynika z przemienności $M$.
Elementem zerowym jest $[0,1]$, bo
$[m,s]+[0,1]=[m,s]$, a przeciwnym do $[m,s]$
jest $[-m,s]$, gdyż $[0,s^2]=[0,1]$.
Dla $r,q\in R$ i klas $x=[m,s]$, $y=[n,t]$
zachodzą $r(x+y)=rx+ry$ oraz
$(r+q)x=rx+qx$, co otrzymujemy przez sprowadzenie
do wspólnego mianownika. Ponadto
$r(qx)=[rqm,s]=(rq)x$ i $1_Rx=x$.
Wszystkie aksjomaty $R$-modułu są więc spełnione.
Dla homomorfizmu $f:M\to N$ wzór
$S^{-1}f([m,s])=[f(m),s]$ jest poprawny,
ponieważ $f(t(s'm-sm'))=t(s'f(m)-sf(m'))$.

Przypomnijmy, że $K=S^{-1}R$ jest ciałem ułamków.
Własność uniwersalna iloczynu tensorowego daje
homomorfizm $K\otimes_RM\to S^{-1}M$ zadany przez
$(r/s)\otimes m\mapsto[rm,s]$. Wzór jest poprawny
względem reprezentanta $r/s$: z $rs'=r's$
wynika $[rm,s]=[r'm,s']$. Addytywność względem
pierwszego argumentu sprawdzamy na wspólnym mianowniku:
para $(r/s+r'/s',m)$ jest wysyłana na
$[(rs'+r's)m,ss']$, a klasa ta jest równa
$[rm,s]+[r'm,s']$.
Addytywność względem $m$ wynika z
$[r(m+m'),s]=[rm,s]+[rm',s]$.
Zbilansowanie wynika z $[(rq)m,s]=[r(qm),s]$
dla każdego $q\in R$. Własność uniwersalna daje
zatem homomorfizm grup abelowych.
Odwrotność definiujemy przez
$[m,s]\mapsto(1/s)\otimes m$.
Jeśli $[m,s]=[m',s']$, to dla pewnego $t\in S$
mamy $t(s'm-sm')=0$. Wtedy
$(1/s)\otimes m=(1/(tss'))\otimes(t s'm)$ oraz
$(1/s')\otimes m'=(1/(tss'))\otimes(t sm')$,
więc wzór nie zależy od reprezentanta.
Jest addytywny, gdyż
$(1/(ss'))\otimes(s'm+sm')=
(1/s)\otimes m+(1/s')\otimes m'$, i zachowuje
mnożenie przez skalary. Oba złożenia zachowują
elementy postaci $[m,s]$ i $(r/s)\otimes m$,
a więc są identycznościami.
Sprawdźmy ponadto naturalność tego izomorfizmu.
Oznaczmy go przez $\Theta_M:K\otimes_RM\to S^{-1}M$.
Dla każdego homomorfizmu $f:M\to N$ i tensora prostego
$(r/s)\otimes m$ zachodzi
$\Theta_N((\operatorname{id}_K\otimes f)((r/s)\otimes m))
=[rf(m),s]=S^{-1}f(\Theta_M((r/s)\otimes m))$.
Ponieważ tensory proste generują $K\otimes_RM$,
otrzymujemy równość homomorfizmów
$\Theta_N\circ(\operatorname{id}_K\otimes f)
=(S^{-1}f)\circ\Theta_M$.

Zobaczmy, że lokalizacja zachowuje dokładność.
Jeżeli $0\to A\xrightarrow{f}B\xrightarrow{g}C\to0$
jest dokładny, to $S^{-1}g$ jest surjektywny:
każdy $[c,s]$ ma przedobraz $[b,s]$ dla pewnego
$b$ z $g(b)=c$. Jeśli $[b,s]$ przechodzi na zero,
to istnieje $t\in S$ takie, że $t g(b)=0$,
a więc $tb\in\ker g=\operatorname{im}f$.
Wybierzmy $a$ z $f(a)=tb$. Wtedy
$[b,s]=[tb,ts]=[f(a),ts]$, czyli klasa należy
do obrazu $S^{-1}f$. Wreszcie, jeśli
$[a,s]\in S^{-1}A$ przechodzi na zero pod
$S^{-1}f$, to $tf(a)=0$ dla pewnego $t\in S$.
Iniektywność $f$ daje $ta=0$, czyli $[a,s]=0$.
Zatem ciąg lokalizowany jest dokładny.
Ponieważ $K\otimes_R-$ jest naturalnie izomorficzny
z lokalizacją, $K$ jest modułem płaskim.
\end{proof}

\begin{przyklad}[Moduł płaski, który nie jest rzutowy]
\label{aa:prz:flat-not-projective}
Grupa addytywna $\mathbb Q$ jest płaskim
$\mathbb Z$-modułem na mocy poprzedniego twierdzenia.
Nie jest jednak modułem rzutowym. Gdyby nim była,
na mocy Twierdzenia~\ref{aa:tw:modules-projective-summand}
byłaby składnikiem sumy prostej pewnego modułu wolnego
$\mathbb Z^{(S)}$. W szczególności istniałby
niezerowy homomorfizm $\mathbb Q\to\mathbb Z^{(S)}$
będący włączeniem składnika.

Tymczasem każdy homomorfizm $h:\mathbb Q\to\mathbb Z$
jest zerowy: dla dowolnego $n\ge1$ mamy
$h(1)=n h(1/n)$, zatem liczba całkowita $h(1)$
jest podzielna przez każde $n$, co wymusza $h(1)=0$.
Dla $a\in\mathbb Z$, $b\ge1$ mamy dalej
$b\,h(a/b)=a\,h(1)=0$. Ponieważ $\mathbb Z$
nie ma torsji, $h(a/b)=0$ dla każdej liczby wymiernej.
Każda projekcja współrzędna homomorfizmu
$\mathbb Q\to\mathbb Z^{(S)}$ jest więc zerowa,
stąd sam homomorfizm jest zerowy. Otrzymaliśmy sprzeczność.
\end{przyklad}

\section{Powiązanie iloczynu tensorowego z modułami homomorfizmów}
\begin{twierdzenie}[Związek iloczynu tensorowego z homomorfizmami]
\label{aa:tw:tensor-hom-adjunction}
Niech $R$ będzie przemiennym pierścieniem z jedynką,
a $M,N,P$ modułami nad $R$. Istnieje naturalny
izomorfizm $R$-modułów
\[
\operatorname{Hom}_R(M\otimes_RN,P)\cong
\operatorname{Hom}_R\bigl(M,\operatorname{Hom}_R(N,P)\bigr).
\]
Homomorfizm $F:M\otimes_RN\to P$ odpowiada
odwzorowaniu $m\mapsto(n\mapsto F(m\otimes n))$.
\end{twierdzenie}
\begin{proof}
Niech $F:M\otimes_RN\to P$ będzie $R$-liniowy.
Dla ustalonego $m\in M$ odwzorowanie
$n\mapsto F(m\otimes n)$ jest $R$-liniowe
ze względu na dwuliniowość kanonicznego odwzorowania.
Wobec tego $\Phi(F)(m)$ należy do
$\operatorname{Hom}_R(N,P)$. Dla $m,m'\in M$
i $r\in R$ mamy
$\Phi(F)(m+m')=\Phi(F)(m)+\Phi(F)(m')$ oraz
$\Phi(F)(rm)=r\Phi(F)(m)$, co sprawdzamy na
każdym $n$. Zatem $\Phi(F)$ jest $R$-liniowe.

Odwrotnie, jeśli $G:M\to\operatorname{Hom}_R(N,P)$
jest $R$-liniowe, określmy $b(m,n)=G(m)(n)$.
Jest ono $R$-dwuliniowe: addytywność i zgodność
ze skalarami w pierwszej zmiennej wynika z
liniowości $G$, a w drugiej z liniowości $G(m)$.
Własność uniwersalna daje jedyny homomorfizm
$\Psi(G):M\otimes_RN\to P$ z
$\Psi(G)(m\otimes n)=G(m)(n)$.
Dla każdego $F$ i $G$ porównanie wartości na
wszystkich tensorach prostych pokazuje, że
$\Psi\Phi(F)=F$ oraz $\Phi\Psi(G)=G$.
Obie operacje zachowują dodawanie i skalary,
więc stanowią izomorfizm $R$-modułów.
\end{proof}

\begin{przyklad}[Obliczenie grupy homomorfizmów cyklicznych]
Dla dodatnich $m,n$ odwzorowanie
$f:\mathbb Z/m\mathbb Z\to\mathbb Z/n\mathbb Z$
jest wyznaczone przez $y=f([1])$, który musi spełniać
$my=0$. Niech $d=\gcd(m,n)$ i zapiszmy $m=dm'$, $n=dn'$,
gdzie $\gcd(m',n')=1$. Warunek $my=0$ oznacza,
że $n\mid m\widetilde y$, gdzie $\widetilde y\in\mathbb Z$
reprezentuje $y$. Jest on równoważny warunkowi
$n'\mid\widetilde y$. Zatem możliwe wartości $y$
są dokładnie klasami wielokrotności $n'$, tworzącymi
podgrupę cykliczną rzędu $d$. Dlatego
$\operatorname{Hom}_{\mathbb Z}(C_m,C_n)\cong C_d$.
Warto porównać ten rezultat z
Przykładem~\ref{aa:prz:tensor-cyclic}:
grupy otrzymane w tych dwóch obliczeniach
są izomorficzne, ale definicje konstrukcji są różne.
\end{przyklad}

\section{Związek z iloczynem tensorowym przestrzeni wektorowych}
\begin{stwierdzenie}[Baza iloczynu tensorowego skończenie wymiarowych przestrzeni]
\label{aa:st:tensor-vspace-dimension}
Niech $K$ będzie ciałem, a $V,W$ przestrzeniami wektorowymi
nad $K$ z bazami $(v_i)_{i=1}^m$ i $(w_j)_{j=1}^n$.
Wtedy rodzina $(v_i\otimes w_j)_{1\le i\le m,\,1\le j\le n}$
jest bazą $V\otimes_KW$; w szczególności
$\dim_K(V\otimes_KW)=mn$.
\end{stwierdzenie}
\begin{proof}
Każdy $v\in V$ i $w\in W$ ma postać
$v=\sum_i a_iv_i$, $w=\sum_j b_jw_j$.
Dwuliniowość daje
$v\otimes w=\sum_{i,j}a_ib_j(v_i\otimes w_j)$,
więc wskazana rodzina rozpina iloczyn tensorowy.

Dla niezależności niech $v_i^*:V\to K$ i
$w_j^*:W\to K$ oznaczają funkcjonały współrzędne,
zadane przez $v_i^*(v_k)=\delta_{ik}$ oraz
$w_j^*(w_\ell)=\delta_{j\ell}$.
Odwzorowanie $(v,w)\mapsto v_i^*(v)w_j^*(w)$
jest dwuliniowe. Z własności uniwersalnej wyznacza
funkcjonał $F_{ij}:V\otimes_KW\to K$, który
spełnia $F_{ij}(v_k\otimes w_\ell)=\delta_{ik}\delta_{j\ell}$.
Jeśli $\sum_{k,\ell}c_{k\ell}(v_k\otimes w_\ell)=0$,
to zastosowanie $F_{ij}$ daje $c_{ij}=0$.
Wszystkie współczynniki są zerowe, więc rodzina jest
liniowo niezależna i stanowi bazę.
\end{proof}

\begin{uwaga}[Co dalej?]
Iloczyn tensorowy jest prawostronnie dokładny, ale na ogół
nie zachowuje iniekcji. Dalszym algebraicznym narzędziem
badania tego zjawiska jest funktor $\operatorname{Tor}$,
zdefiniowany z wykorzystaniem rozdzielczości modułów.
Analogicznie funktor $\operatorname{Ext}$ bada dokładność
odwzorowań $\operatorname{Hom}$. Oba funktory zostaną
zdefiniowane i zbadane w części poświęconej algebrze
homologicznej; nie używamy ich tutaj bez uprzedniej definicji.
\end{uwaga}

% AA-19 | Algebra abstrakcyjna | Iteracja X
\chapter[Ciała i ciała ułamków]{Ciała, ciała ułamków i konstrukcje algebraiczne liczb}
\label{aa:ch:fields-fractions}
\setlength{\parskip}{2pt}

\section{Dlaczego pojawia się pojęcie ciała?}
W rozdziałach o pierścieniach pokazaliśmy, że nie każdą niezerową
liczbę całkowitą można odwrócić wewnątrz $\mathbb Z$. Dodanie
odwrotności mnożeniowych prowadzi do $\mathbb Q$, ale konstrukcja
musi zachować wcześniejsze działania i nie może utożsamić
różnych liczb całkowitych. Jednocześnie ta sama idea działa
dla dowolnej dziedziny całkowitości, na przykład dla $K[t]$.
Właśnie to wyjaśnia konstrukcja ciała ułamków.

\begin{definicja}[Ciało, podciało i homomorfizm ciał]
\label{aa:def:field}
\emph{Ciałem} nazywamy niezerowy pierścień przemienny $K$ z jedynką,
w którym każdy element $a\in K\setminus\{0\}$ ma odwrotność
mnożeniową $a^{-1}\in K$. Zbiór $K^\times=K\setminus\{0\}$
jest wówczas grupą abelową względem mnożenia:
jeśli $a,b\ne0$ i $ab=0$, pomnożenie przez $a^{-1}$
dałoby $b=0$, sprzeczność. Jedynka i odwrotności
należą do $K^\times$, a łączność i przemienność
mnożenia dziedziczymy z $K$.
\emph{Podciało} $F\subseteq K$ jest podzbiorem zawierającym
$0,1$, zamkniętym na dodawanie, odejmowanie, mnożenie
oraz odwrotności elementów niezerowych. Dla ciał $K,L$ \emph{homomorfizmem ciał}
nazywamy odwzorowanie $f:K\to L$ zachowujące
$0,1,+,\cdot$. Zachowuje ono również odwrotności:
dla $a\ne0$ równość $f(a)f(a^{-1})=f(1)=1$
pokazuje, że $f(a^{-1})=f(a)^{-1}$.
\end{definicja}

\begin{stwierdzenie}[Homomorfizm ciał jest różnowartościowy]
\label{aa:st:field-embedding-injective}
Niech $K,L$ będą ciałami, a $f:K\to L$
homomorfizmem ciał. Wówczas $f$ jest różnowartościowy.
\end{stwierdzenie}
\begin{proof}
Niech $a\in\ker f$ i załóżmy, że $a\ne0$. Wówczas
$1_K=aa^{-1}$, a z zachowywania mnożenia mamy
$1_L=f(1_K)=f(a)f(a^{-1})=0$, co przeczy definicji ciała.
Zatem $\ker f=\{0\}$. Jeśli $f(x)=f(y)$, to
$f(x-y)=0$ i stąd $x-y=0$, czyli $x=y$.
\end{proof}

\begin{uwaga}[Nie mylić ciała z pierścieniem z dzieleniem]
W ciele zakładamy przemienność mnożenia. W nieprzemiennym
pierścieniu z dzieleniem każdy element niezerowy również
jest odwracalny, ale może nie komutować z innymi.
Przykładem jest algebra kwaternionów Hamiltona $\mathbb H$,
która jest pierścieniem z dzieleniem, lecz nie jest ciałem
w przyjętym tu, przemiennym znaczeniu.
\end{uwaga}

\section{Ciało ułamków dowolnej dziedziny całkowitości}
\begin{definicja}[Relacja ułamków formalnych]
\label{aa:def:fractions-relation}
Niech $R$ będzie dziedziną całkowitości, a
$S=R\setminus\{0\}$. Na $R\times S$ definiujemy
$(a,b)\sim(c,d)$ wtedy i tylko wtedy, gdy $ad=bc$.
Klasę $(a,b)$ oznaczamy $a/b$; na razie jest to
symbol formalny, a nie dzielenie wykonane w $R$.
\end{definicja}

\begin{lemat}[Poprawność relacji równoważności]
\label{aa:lem:fractions-equivalence}
Niech $R$ będzie dziedziną całkowitości i
$S=R\setminus\{0\}$. Relacja z
Definicji~\ref{aa:def:fractions-relation} jest relacją
równoważności na $R\times S$.
\end{lemat}
\begin{proof}
Zwrotność wynika z $ab=ba$. Jeśli $ad=bc$, to
$cb=da$, więc zachodzi symetria. Dla przechodniości
załóżmy $(a,b)\sim(c,d)$ oraz $(c,d)\sim(e,f)$.
Oznacza to $ad=bc$ oraz $cf=de$. Pierwszą równość
mnożymy przez $f$, a drugą przez $b$:
$adf=bcf=bde$. Stąd $d(af-be)=0$. Ponieważ
$d\in S$, brak dzielników zera w $R$ daje $af=be$.
Jest to dokładnie warunek $(a,b)\sim(e,f)$.
\end{proof}

\begin{twierdzenie}[Istnienie ciała ułamków]
\label{aa:tw:fractions-field}
Niech $R$ będzie dziedziną całkowitości,
$S=R\setminus\{0\}$ i niech $\sim$ będzie relacją
z Definicji~\ref{aa:def:fractions-relation}. Zbiór klas
$K=(R\times S)/{\sim}$ z działaniami
\[
 \frac ab+\frac cd=\frac{ad+bc}{bd},\qquad
 \frac ab\cdot\frac cd=\frac{ac}{bd}
\]
jest ciałem. Zero i jedynkę reprezentują odpowiednio
$0/1$ i $1/1$, a dla $a/b\ne0/1$ odwrotnością jest $b/a$.
Odwzorowanie $\iota:R\to K$, $a\mapsto a/1$,
jest różnowartościowym homomorfizmem pierścieni.
\end{twierdzenie}
\begin{proof}
\emph{(1) Zamkniętość.} Ponieważ $b,d$ są niezerowe,
a $R$ jest dziedziną całkowitości, również $bd\ne0$.
Oba prawe wyrażenia wyznaczają więc poprawne pary.

\emph{(2) Niezależność od reprezentantów.}
Niech $a/b=a'/b'$ oraz $c/d=c'/d'$,
czyli $ab'=a'b$ i $cd'=c'd$. Wówczas
$ac\,b'd'=(ab')(cd')=(a'b)(c'd)=a'c'\,bd$,
co dowodzi poprawności mnożenia.
W przypadku dodawania sprawdzamy równość
$(ad+bc)b'd'=(a'd'+b'c')bd$:
po rozwinięciu pierwszego składnika mamy
$ad b'd'=ab'\,dd'=a'b\,dd'=a'd'\,bd$;
dla drugiego składnika analogicznie
$bc b'd'=bb'\,cd'=bb'\,c'd=b'c'\,bd$.
Dodając obie równości dostajemy żądany wynik.

\emph{(3) Aksjomaty pierścienia.} Przemienność i łączność
mnożenia wynikają z odpowiednich własności $R$:
$(a/b)(c/d)(e/f)=ace/(bdf)$, niezależnie od nawiasowania.
Dodawanie jest przemienne, a dla trzech ułamków
obydwa nawiasowania sumy $a/b$, $c/d$, $e/f$ dają
$(adf+cbf+ebd)/(bdf)$. Każdy składnik licznika
powstaje przez pomnożenie odpowiedniego licznika
przez dwa pozostałe mianowniki. Element $0/1$ jest zerem, gdyż
$(a/b)+(0/1)=a/b$. Przeciwieństwem $a/b$
jest $(-a)/b$, ponieważ ich suma jest $0/b^2=0/1$.
Ponadto $(a/b)(1/1)=a/b$, więc $1/1$ jest jedynką.
Różnią się one, gdyż $0/1=1/1$ oznaczałoby $0=1$
w dziedzinie $R$.
Rozdzielność sprawdzamy bez skrótów:
$(a/b)((c/d)+(e/f))=a(cf+ed)/(bdf)$
i $(a/b)(c/d)+(a/b)(e/f)=(acf+aed)/(bdf)$.
Licznikami są w obu przypadkach $acf+aed$.

\emph{(4) Odwrotności.} Warunek $a/b=0/1$ oznacza $a=0$.
Jeśli więc $a/b\ne0/1$, to $a\ne0$ i $b/a$ jest
dopuszczalnym ułamkiem. Iloczyn wynosi
$(a/b)(b/a)=ab/(ab)=1/1$, ponieważ $ab\ne0$
i $(ab)\cdot1=1\cdot(ab)$.
Ponadto wzór nie zależy od reprezentanta: jeśli
$a/b=c/d\ne0/1$, to $a,c\ne0$, a równość
$ad=bc$ jest także warunkiem równoważności par
$(b,a)$ i $(d,c)$. Zatem $b/a=d/c$.

\emph{(5) Zanurzenie.} Wzory na działania dają
$\iota(a+c)=\iota(a)+\iota(c)$ i
$\iota(ac)=\iota(a)\iota(c)$; ponadto
$\iota(1)=1/1$. Jeśli $\iota(a)=\iota(c)$,
to $a\cdot1=c\cdot1$, zatem $a=c$.
\end{proof}


\begin{definicja}[Ciało ułamków]
\label{aa:def:field-of-fractions}
Niech $R$ będzie dziedziną całkowitości.
Ciało skonstruowane dla $R$ w
Twierdzeniu~\ref{aa:tw:fractions-field} nazywamy
\emph{ciałem ułamków} dziedziny $R$ i oznaczamy
$\operatorname{Frac}(R)$.
\end{definicja}

\begin{twierdzenie}[Własność uniwersalna ciała ułamków]
\label{aa:tw:fractions-universal}
Niech $R$ będzie dziedziną całkowitości, $L$ ciałem,
a $f:R\to L$ różnowartościowym homomorfizmem pierścieni.
Istnieje dokładnie jeden homomorfizm ciał
$\overline f:\operatorname{Frac}(R)\to L$ taki, że
$\overline f(a/1)=f(a)$ dla każdego $a\in R$.
Jest nim odwzorowanie $\overline f(a/b)=f(a)f(b)^{-1}$.
\end{twierdzenie}
\begin{proof}
Ponieważ $b\ne0$ oraz $f$ jest różnowartościowy,
$f(b)\ne0$ i $f(b)^{-1}$ istnieje. Jeśli $a/b=c/d$,
to $ad=bc$, a więc $f(a)f(d)=f(b)f(c)$.
Mnożąc przez $f(b)^{-1}f(d)^{-1}$ dostajemy
$f(a)f(b)^{-1}=f(c)f(d)^{-1}$. Wzór nie zależy
od reprezentantów. Dla $a/b,c/d$ otrzymujemy
\[
\overline f\left(\frac ab+\frac cd\right)
=\frac{f(a)f(d)+f(b)f(c)}{f(b)f(d)}
=\frac{f(a)}{f(b)}+\frac{f(c)}{f(d)},
\]
a dla mnożenia, korzystając z przemienności w $L$,
mamy
\[
\overline f\left(\frac ab\cdot\frac cd\right)
=f(ac)f(bd)^{-1}
=\bigl(f(a)f(b)^{-1}\bigr)
 \bigl(f(c)f(d)^{-1}\bigr).
\]
Ponadto $\overline f(0/1)=f(0)=0$ i
$\overline f(1/1)=f(1)=1$, więc jest to homomorfizm ciał.
Dla jednoznaczności zauważmy, że $b/1$ jest odwracalne
w $\operatorname{Frac}(R)$ i że
$a/b=(a/1)(b/1)^{-1}$. Każdy homomorfizm ciał
przedłużający $f$ musi zachować ten iloczyn oraz odwrotność,
a zatem musi działać według podanego wzoru.
Różnowartościowość $f$ jest niezbędna dla tego wzoru:
bez niej jakiś niezerowy mianownik $b$ mógłby mieć
$f(b)=0$, którego nie można odwrócić w $L$.
Na przykład redukcja $\mathbb Z\to\mathbb F_p$
nie przedłuża się do homomorfizmu ciał
$\mathbb Q\to\mathbb F_p$: odwrotność liczby
$p/1$ musiałaby być wysłana na odwrotność zera.
\end{proof}

\begin{przyklad}[Liczby wymierne jako ciało ułamków]
\label{aa:prz:construct-rationals}
Wybierzmy $R=\mathbb Z$. Wówczas otrzymujemy
\[
\mathbb Q\cong\operatorname{Frac}(\mathbb Z)
=(\mathbb Z\times(\mathbb Z\setminus\{0\}))/\!\sim,
\qquad (a,b)\sim(c,d)\iff ad=bc.
\]
Na przykład $(2,3)\sim(4,6)$, lecz $(2,3)\not\sim(3,4)$.
Równość $1/2+1/3=5/6$ wynika z formalnego wzoru
$(1\cdot3+1\cdot2)/(2\cdot3)$, a nie z wcześniej
założonej arytmetyki ułamków. Różne pary reprezentują tę
samą liczbę dokładnie wtedy, gdy zachodzi równość
iloczynów krzyżowych. Element $-5/7$ można zapisać
jako $5/(-7)$, ponieważ $(-5)\cdot(-7)=5\cdot7$.
W szczególności wcześniejsza konstrukcja $\mathbb Z$
z $\mathbb N_0\times\mathbb N_0$ oraz obecna konstrukcja
$\mathbb Q$ tworzą łańcuch zanurzeń
$\mathbb N_0\hookrightarrow\mathbb Z\hookrightarrow\mathbb Q$.
\end{przyklad}

\begin{przyklad}[Funkcje wymierne]
\label{aa:prz:rational-functions}
Dla dowolnego ciała $K$ pierścień $K[t]$ jest dziedziną
całkowitości, a jego ciało ułamków oznaczamy $K(t)$.
Elementem $K(t)$ jest klasa pary wielomianów $(p,q)$,
$q\ne0$, którą zapisujemy $p(t)/q(t)$.
Przykładowo $t/(t+1)$ i $t^2/(t^2+t)$ są równe
w $K(t)$, gdyż $t(t^2+t)=t^2(t+1)$;
oba mianowniki są niezerowymi wielomianami.
Wspólna klasa tych ułamków nie należy do podpierścienia
$K[t]\subseteq K(t)$. Gdyby istniał wielomian $h\in K[t]$
taki, że $t/(t+1)=h$, to w $K[t]$ mielibyśmy
$t=h(t)(t+1)$; podstawienie $t=-1$ dałoby $-1=0$,
co przeczy $1\ne0$ w ciele $K$.
Jeżeli $p/q=0$, to $p=0$, gdyż $K[t]$ nie ma
dzielników zera.
\end{przyklad}

\begin{przyklad}[Dlaczego nie wystarczy pierścień z dzielnikami zera]
\label{aa:prz:frac-zero-divisors}
W $R=\mathbb Z/6\mathbb Z$ mamy $[2][3]=0$ mimo
$[2]\ne0$ i $[3]\ne0$. Gdyby można było zanurzyć
$R$ w ciało $K$, to obrazy tych dwóch elementów byłyby
niezerowe, lecz ich iloczyn równy zeru. Jest to niemożliwe
w ciele. Nie istnieje więc ciało zawierające $R$ jako
podpierścień. Można badać \emph{całkowity pierścień ułamków}
po odwróceniu wszystkich elementów niebędących dzielnikami
zera, lecz nie jest on na ogół ciałem.
\end{przyklad}

\section{Charakterystyka i ciało pierwsze}
\begin{definicja}[Charakterystyka pierścienia]
\label{aa:def:field-characteristic}
Dla pierścienia $R$ z jedynką i $1_R\ne0_R$
\emph{charakterystyką} nazywamy najmniejszą dodatnią
liczbę całkowitą $n$, dla której $n1_R=0_R$, jeśli
taka istnieje; w przeciwnym przypadku charakterystyka
wynosi zero. Zapis $n1_R$ oznacza sumę $n$ kopii $1_R$.
\end{definicja}

\begin{twierdzenie}[Charakterystyka ciała]
\label{aa:tw:field-characteristic-prime}
Charakterystyka dowolnego ciała $K$ jest równa zeru
albo liczbie pierwszej $p$.
\end{twierdzenie}
\begin{proof}
Niech $n>0$ będzie charakterystyką $K$.
Gdyby $n=ab$ dla $1<a,b<n$, to z minimalności $n$
wynikałoby $a1_K\ne0$ i $b1_K\ne0$.
Jednak $(a1_K)(b1_K)=(ab)1_K=n1_K=0$.
Jest to niemożliwe, ponieważ ciało nie ma dzielników zera.
Zatem $n$ nie jest złożone, a minimalność i $1_K\ne0$
wykluczają $n=1$. Stąd $n$ jest liczbą pierwszą.
\end{proof}

\begin{twierdzenie}[Najmniejsze podciało]
\label{aa:tw:prime-subfield}
Każde ciało $K$ zawiera dokładnie jedno najmniejsze podciało.
Jest ono izomorficzne z $\mathbb Q$, jeśli
$\operatorname{char}K=0$, oraz z $\mathbb F_p=\mathbb Z/p\mathbb Z$,
jeśli $\operatorname{char}K=p>0$.
\end{twierdzenie}
\begin{proof}
Rozważmy homomorfizm pierścieni
$\phi:\mathbb Z\to K$, $n\mapsto n1_K$.
Jego jądro jest ideałem $\mathbb Z$, a więc ma
postać $m\mathbb Z$ dla pewnego $m\ge0$.
Jeśli charakterystyka jest zerowa, to $m=0$ i $\phi$
jest różnowartościowy. Na mocy
Twierdzenia~\ref{aa:tw:fractions-universal} istnieje
jednoznaczne zanurzenie $\mathbb Q\to K$.
Jeżeli charakterystyka jest dodatnia, to z poprzedniego
twierdzenia wynika $m=p$; pierwsze twierdzenie o
izomorfizmie daje zanurzenie
$\mathbb Z/p\mathbb Z\to K$. Iloraz ten jest ciałem,
ponieważ dla każdego $a$ niepodzielnego przez $p$
tożsamość Bézouta $ua+vp=1$ daje odwrotność klasy $[u]$.

Każde podciało $F\subseteq K$ zawiera $1_K$, a zatem
wszystkie wielokrotności $n1_K$ i odwrotności wszystkich
niezerowych spośród nich. W charakterystyce zero zawiera
też wszystkie ilorazy $(m1_K)(n1_K)^{-1}$ dla $n\ne0$,
czyli cały obraz $\mathbb Q$. Musi więc zawierać opisany obraz
$\mathbb Q$ lub $\mathbb F_p$. Obraz sam jest podciałem,
stąd jest najmniejszy i jedyny.
\end{proof}

\begin{twierdzenie}[Endomorfizm Frobeniusa]
\label{aa:tw:field-frobenius}
Niech $K$ będzie ciałem charakterystyki $p>0$.
Odwzorowanie $F:K\to K$, $F(x)=x^p$, jest
różnowartościowym homomorfizmem ciał. Jeśli $K$
jest skończone, $F$ jest automorfizmem.
\end{twierdzenie}
\begin{proof}
Wzór dwumianowy daje
$(x+y)^p=\sum_{j=0}^p\binom pj x^j y^{p-j}$.
Dla $1\le j\le p-1$ współczynnik $\binom pj$
jest podzielny przez $p$: w równości całkowitoliczbowej
$\binom pj\,j!(p-j)!=p!$ żaden z czynników $j!$ ani $(p-j)!$ nie jest
podzielny przez $p$.
Ponieważ $p$ jest liczbą pierwszą, z lematu Euklidesa wynika
$p\mid\binom pj$.
W charakterystyce $p$ pośrednie składniki znikają,
więc $(x+y)^p=x^p+y^p$. Równość
$(xy)^p=x^py^p$ jest oczywista z przemienności.
Ponieważ również $F(1)=1^p=1$,
$F$ jest homomorfizmem ciał, a
Twierdzenie~\ref{aa:st:field-embedding-injective}
daje jego różnowartościowość.
Na zbiorze skończonym każda funkcja różnowartościowa
jest surjektywna, więc $F$ jest automorfizmem.
\end{proof}

\begin{przyklad}[Frobenius nie zawsze jest na]
\label{aa:prz:frobenius-not-surjective}
Niech $p$ będzie liczbą pierwszą. W ciele
$K=\mathbb F_p(t)$ element $t$ nie jest
$p$-tą potęgą żadnego elementu. Gdyby
$(a(t)/b(t))^p=t$ dla wielomianów $a,b\in\mathbb F_p[t]$,
$b\ne0$, to $a(t)\ne0$ i $a(t)^p=t b(t)^p$.
Po porównaniu stopni niezerowych wielomianów
otrzymalibyśmy $p\deg a=1+p\deg b$, co jest
niemożliwe modulo $p$. Endomorfizm Frobeniusa
nie jest zatem surjektywny w tym nieskończonym ciele.
\end{przyklad}

\section{Dwie algebraiczne konstrukcje ciała liczb zespolonych}
W tej sekcji \emph{przyjmujemy ciało liczb rzeczywistych}
$\mathbb R$ wraz ze znanym porządkiem i jego własnościami.
Nie konstruujemy tu $\mathbb R$: konstrukcja ciała
uporządkowanego zupełnego należy do skryptu analizy.
Pokazujemy natomiast, że operacje na parach liczb
rzeczywistych i odpowiedni pierścień ilorazowy prowadzą
do tego samego ciała.

\begin{definicja}[Liczby zespolone jako pary]
\label{aa:def:complex-as-pairs}
Na zbiorze $\mathbb R^2$ określamy
\[
(a,b)+(c,d)=(a+c,b+d),\qquad
(a,b)(c,d)=(ac-bd,ad+bc).
\]
Element neutralny dodawania to $(0,0)$, a jedynka
mnożenia to $(1,0)$. Oznaczamy $i=(0,1)$.
\end{definicja}

\begin{twierdzenie}[Struktura ciała na $\mathbb R^2$]
\label{aa:tw:complex-pairs-field}
Działania z Definicji~\ref{aa:def:complex-as-pairs}
czynią $\mathbb R^2$ ciałem,
w którym $i^2=(-1,0)$ oraz
$(a,b)=(a,0)+(b,0)i$.
Dla $(a,b)\ne(0,0)$ odwrotność jest równa
\[
(a,b)^{-1}=\left(\frac a{a^2+b^2},
 -\frac b{a^2+b^2}\right).
\]
\end{twierdzenie}
\begin{proof}
Wzory na dodawanie dają grupę abelową z przeciwnym
elementem $(-a,-b)$. Przemienność mnożenia
wynika z $ac=ca$, $bd=db$ i $ad+bc=cb+da$.
Dla łączności mnożenia obliczmy
$((a,b)(c,d))(e,f)$: jego pierwsza współrzędna to
$(ac-bd)e-(ad+bc)f=ace-bde-adf-bcf$,
a druga to $(ac-bd)f+(ad+bc)e=acf-bdf+ade+bce$.
Obliczając $(c,d)(e,f)=(ce-df,cf+de)$, otrzymujemy
z $(a,b)((c,d)(e,f))$ pierwszą współrzędną
$a(ce-df)-b(cf+de)=ace-adf-bcf-bde$
i drugą $a(cf+de)+b(ce-df)=acf+ade+bce-bdf$.
Są to te same wyrażenia, więc mnożenie jest łączne.
Dla rozdzielności po lewej stronie porównajmy
$(a,b)((c,d)+(e,f))$ z sumą dwóch iloczynów:
w obu przypadkach otrzymujemy parę
$(ac+ae-bd-bf,ad+af+bc+be)$.
Rozdzielność po prawej wynika następnie z przemienności. Bezpośredni rachunek daje
$(1,0)(a,b)=(a,b)$, $(0,1)^2=(-1,0)$
oraz $(a,0)+(b,0)(0,1)=(a,b)$.

Jeżeli $(a,b)\ne(0,0)$, to $a^2+b^2>0$, a więc
mianownik w podanym wzorze nie jest zerem. Iloczyn
$(a,b)(a,-b)=(a^2+b^2,0)$; podzielenie współrzędnych
przez $a^2+b^2$ daje wskazaną odwrotność.
Ponieważ istnieją odwrotności wszystkich elementów
niezerowych, otrzymaliśmy ciało.
\end{proof}

\begin{twierdzenie}[Liczby zespolone jako pierścień ilorazowy]
\label{aa:tw:complex-quotient}
Wielomian $X^2+1\in\mathbb R[X]$ jest nierozkładalny.
Pierścień $\mathbb R[X]/(X^2+1)$ jest ciałem i istnieje
izomorfizm ciał
\[
\mathbb R[X]/(X^2+1)\longrightarrow\mathbb R^2,
\qquad [a+bX]\longmapsto(a,b).
\]
Izomorfizm jest wyznaczony przez przesłanie klasy
$[X]$ na element $i=(0,1)$, dla którego $i^2=(-1,0)$.
\end{twierdzenie}
\begin{proof}
Wielomian drugiego stopnia nad ciałem jest rozkładalny
wtedy i tylko wtedy, gdy ma pierwiastek w tym ciele:
rozkład nietrywialny ma czynniki liniowe, a pierwiastek
$x$ daje rozkład $f(X)=(X-x)g(X)$ na mocy
twierdzenia o reszcie.
Gdyby $X^2+1$ miał pierwiastek $x\in\mathbb R$,
to $x^2=-1$, co przeczyłoby $x^2\ge0$.
Zatem $X^2+1$ jest nierozkładalny. Ponieważ $\mathbb R[X]$ jest dziedziną euklidesową,
dla każdego $g\notin(X^2+1)$ największy wspólny dzielnik
$g$ i $X^2+1$ jest jednostką: nierozkładalność $X^2+1$
wyklucza nieodwracalny wspólny dzielnik.
Tożsamość Bézouta daje wówczas wielomiany $u,v$ z
$ug+v(X^2+1)=1$, czyli $[u][g]=[1]$ w ilorazie.
Każda niezerowa klasa jest odwracalna, zatem iloraz
jest ciałem, a ideał $(X^2+1)$ jest maksymalny.

Dzielenie z resztą przez unormowany wielomian
$X^2+1$ pokazuje, że każda klasa ma reprezentanta
postaci $a+bX$. Jest on jednoznaczny: różnica dwóch
takich reprezentantów o stopniu mniejszym niż $2$
nie może być niezerową wielokrotnością $X^2+1$. W ilorazie
$[X]^2=[-1]$, więc
$(a+b[X])(c+d[X])=(ac-bd)+(ad+bc)[X]$.
Wskazany wzór zachowuje działania i jedynkę.
Jest różnowartościowy dzięki jednoznaczności reszty,
a surjektywny, bo każda para $(a,b)$ jest obrazem
klasy $a+b[X]$. Otrzymaliśmy izomorfizm ciał. Powstałe ciało
$\mathbb R^2$ z powyższymi działaniami oznaczamy
$\mathbb C$.
\end{proof}

\begin{przyklad}[Zespolone równanie kwadratowe bez analizy]
\label{aa:prz:complex-square-root}
W skonstruowanym ciele $\mathbb C$ równanie
$z^2+1=0$ ma dwa rozwiązania: $i$ i $-i$.
Pierwszy wniosek to $i^2=-1$. Ponieważ
$z^2+1=(z-i)(z+i)$ i $\mathbb C$ nie ma dzielników
zera, każde rozwiązanie musi być równe $i$ lub $-i$.
Dwa rozwiązania są różne, gdyż ciało uporządkowane
$\mathbb R$ ma charakterystykę zero:
$n\cdot1>0$ dla każdego $n\ge1$, zatem $2i\ne0$.
\end{przyklad}

\section{Granice opisu konstrukcyjnego}
Ta sama ogólna idea pojawi się w kolejnym rozdziale:
iloraz $K[X]/(f)$ pozwala \emph{dołączyć pierwiastek}
nierozkładalnego wielomianu $f$. Iloraz przez ideał generowany przez wielomian rozkładalny
nie musi być ciałem.
Dla przykładu $K[X]/(X^2)$ zawiera niezerową klasę
$[X]$ o kwadracie równym zeru. Nie jest więc ciałem.
Taka konfrontacja przykładów z twierdzeniem wyjaśnia,
dlaczego nierozkładalność wielomianu jest istotnym
założeniem konstrukcji rozszerzeń ciał.

\setlength{\parskip}{4pt}

% AA-20 | Algebra abstrakcyjna | Iteracja X
\chapter{Rozszerzenia ciał i ciała rozkładu wielomianów}
\label{aa:ch:field-extensions}

\section{Rozszerzenie ciała jako przestrzeń wektorowa}
Gdy rozważamy rozwiązania równań wielomianowych, nawet nad
ciałem może zabraknąć pierwiastków. Przykładowo
$X^2-2$ nie ma pierwiastka w $\mathbb Q$, choć ma go
w większym ciele. Chcemy systematycznie konstruować
większe ciała i mierzyć, ile nowych elementów trzeba dodać.

\begin{definicja}[Rozszerzenie ciał i jego stopień]
\label{aa:def:field-extension-degree}
Niech $K,L$ będą ciałami, a $K\subseteq L$ podciałem.
Mówimy, że $L/K$ jest \emph{rozszerzeniem ciał}.
Mnożenie w $L$ ograniczone do $K\times L$ nadaje $L$
strukturę przestrzeni wektorowej nad $K$.
Jej wymiar oznaczamy $[L:K]$ i nazywamy
\emph{stopniem rozszerzenia}. Rozszerzenie jest
\emph{skończone}, jeśli $[L:K]<\infty$.
W przypadku zanurzenia $K\hookrightarrow L$ utożsamiamy
$K$ z jego obrazem, aby móc stosować ten sam zapis.
\end{definicja}

\begin{przyklad}[Stopień rozszerzenia $\mathbb C/\mathbb R$]
\label{aa:prz:complex-degree}
Z konstrukcji $\mathbb C\cong\mathbb R^2$ w
Rozdziale~\ref{aa:ch:fields-fractions} wynika, że
każdy $z\in\mathbb C$ ma jednoznacznie postać
$z=a+bi$ z $a,b\in\mathbb R$. Wektory $1,i$
tworzą więc bazę $\mathbb C$ nad $\mathbb R$, a zatem
$[\mathbb C:\mathbb R]=2$. Ten fakt nie wymaga
pełnej teorii funkcji holomorficznych ani konstruowania
liczb rzeczywistych na nowo.
\end{przyklad}

\begin{twierdzenie}[Twierdzenie o wieży rozszerzeń]
\label{aa:tw:field-tower}
Niech $K\subseteq E\subseteq L$ będą ciałami.
Jeśli $[E:K]=r<\infty$ i $[L:E]=s<\infty$,
to $[L:K]=rs$. Ponadto, jeśli $[L:K]<\infty$,
to oba stopnie po prawej stronie są skończone.
\end{twierdzenie}
\begin{proof}
Wybierzmy bazę $u_1,\dots,u_r$ przestrzeni $E$ nad $K$
oraz bazę $v_1,\dots,v_s$ przestrzeni $L$ nad $E$.
Pokażemy, że $rs$ iloczynów $u_iv_j$ tworzy bazę $L$ nad $K$.

\emph{Rozpinanie.} Dla $x\in L$ istnieją $b_j\in E$ takie,
że $x=\sum_j b_jv_j$. Dla każdego $j$ istnieją $a_{ij}\in K$
z $b_j=\sum_i a_{ij}u_i$. Wobec tego
$x=\sum_{j=1}^s\sum_{i=1}^r a_{ij}(u_iv_j)$,
czyli wskazane iloczyny rozpinają $L$.

\emph{Niezależność.} Jeśli
$\sum_{j=1}^s\sum_{i=1}^r a_{ij}u_iv_j=0$,
to traktując $\sum_i a_{ij}u_i$ jako współczynnik w $E$
i korzystając z niezależności $v_j$ nad $E$ dostajemy
$\sum_i a_{ij}u_i=0$ dla każdego $j$.
Niezależność $u_i$ nad $K$ daje wszystkie $a_{ij}=0$.
Stąd $[L:K]=rs$.

Jeśli natomiast $[L:K]=n<\infty$, to $E$ jest
$K$-podprzestrzenią $L$, a zatem $[E:K]\le n$.
Każda rodzina liniowo niezależna elementów $L$ nad $E$
jest też niezależna nad $K$: każda relacja o współczynnikach
w $K$ byłaby również relacją nad $E$. Stąd $[L:E]\le n$.
Możemy więc zastosować udowodnioną pierwszą część.
\end{proof}

\begin{wniosek}[Podzielność stopni]
\label{aa:wn:tower-divisibility}
Jeśli $L/K$ jest rozszerzeniem skończonym i
$K\subseteq E\subseteq L$, to $[E:K]$ dzieli $[L:K]$.
\end{wniosek}
\begin{proof}
Równość $[L:K]=[L:E][E:K]$ przedstawia $[L:K]$
jako iloczyn dwóch dodatnich liczb całkowitych.
\end{proof}

\section{Elementy algebraiczne i wielomiany minimalne}
\begin{definicja}[Algebraiczność i przestępność]
\label{aa:def:algebraic-element}
Niech $L/K$ będzie rozszerzeniem ciał, a $\alpha\in L$.
Element $\alpha$ jest \emph{algebraiczny nad $K$}, jeśli
istnieje niezerowy wielomian $f\in K[X]$ z $f(\alpha)=0$.
W przeciwnym przypadku jest \emph{przestępny nad $K$}.
Przez $K[\alpha]$ oznaczamy najmniejszy podpierścień $L$
zawierający $K$ i $\alpha$, a przez $K(\alpha)$
najmniejsze podciało $L$ zawierające te elementy.
\end{definicja}

\begin{lemat}[Wielomian minimalny]
\label{aa:lem:minimal-polynomial}
Niech $L/K$ będzie rozszerzeniem ciał i niech
$\alpha\in L$ będzie algebraiczny nad $K$. Wówczas istnieje
jedyny wielomian unormowany $m_{\alpha,K}\in K[X]$
o najmniejszym dodatnim stopniu spełniający
$m_{\alpha,K}(\alpha)=0$. Wielomian ten jest
nierozkładalny i dzieli każdy $f\in K[X]$
taki, że $f(\alpha)=0$.
\end{lemat}
\begin{proof}
Rozważmy ideał
$I_\alpha=\{f\in K[X]:f(\alpha)=0\}$.
Jest to jądro homomorfizmu podstawiania
$\operatorname{ev}_\alpha:K[X]\to L$.
Ponieważ $\alpha$ jest algebraiczny, ideał zawiera
niezerowy wielomian. Nie zawiera jednak niezerowej
stałej: dla $a\in K$ mamy $a(\alpha)=a$, a $K$ jest
podciałem $L$. Wybierzmy w nim wielomian $m$
o najmniejszym możliwym stopniu, a następnie podzielmy
przez jego niezerowy współczynnik wiodący, aby był unormowany.

Dla dowolnego $f\in I_\alpha$ dzielenie z resztą
w $K[X]$ daje $f=qm+r$, gdzie $r=0$ albo
$\deg r<\deg m$. Po podstawieniu $\alpha$
otrzymujemy $r(\alpha)=f(\alpha)-q(\alpha)m(\alpha)=0$.
Minimalność stopnia $m$ wymusza $r=0$, więc $m\mid f$.
Jeżeli $m=uv$ dla wielomianów dodatniego stopnia,
to $0=m(\alpha)=u(\alpha)v(\alpha)$.
Ponieważ $L$ jest ciałem, jeden z tych dwóch czynników
musi być zerowy. Byłby to wielomian mniejszego stopnia
niż $m$, sprzecznie z jego wyborem. Wielomian $m$ jest
zatem nierozkładalny.

Jeżeli $m'$ jest innym unormowanym wielomianem o żądanych
własnościach, to $m\mid m'$ i $m'\mid m$.
Ich stopnie są równe, a unormowanie wymusza $m=m'$.
\end{proof}

\begin{twierdzenie}[Proste rozszerzenie algebraiczne]
\label{aa:tw:simple-algebraic-extension}
Niech $L/K$ będzie rozszerzeniem ciał i
$\alpha\in L$ elementem algebraicznym nad $K$.
Jeśli $d=\deg m_{\alpha,K}$, to
\[
K[\alpha]=K(\alpha)\cong K[X]/(m_{\alpha,K}),
\qquad [K(\alpha):K]=d.
\]
Bazę $K(\alpha)$ nad $K$ stanowi rodzina
$1,\alpha,\ldots,\alpha^{d-1}$.
\end{twierdzenie}
\begin{proof}
Rozważmy homomorfizm
$\operatorname{ev}_\alpha:K[X]\to L$ z
$f\mapsto f(\alpha)$. Jego obrazem jest $K[\alpha]$,
a jądro jest równe $(m_{\alpha,K})$ na mocy
Lematu~\ref{aa:lem:minimal-polynomial}.
Pierwsze twierdzenie o izomorfizmie pierścieni daje
$K[X]/(m_{\alpha,K})\cong K[\alpha]$.

Dlaczego iloraz jest ciałem? Jeśli $[f]\ne0$, to
$m_{\alpha,K}\nmid f$. Nierozkładalność $m_{\alpha,K}$
i algorytm Euklidesa w $K[X]$ dają
$\gcd(f,m_{\alpha,K})=1$. Z tożsamości Bézouta
istnieją $u,v\in K[X]$ z
$uf+vm_{\alpha,K}=1$. W ilorazie mamy
$[u][f]=[1]$, więc każda niezerowa klasa ma odwrotność.
Stąd obraz $K[\alpha]$ jest już ciałem, a więc
$K[\alpha]=K(\alpha)$.

Dzielenie wielomianu przez $m_{\alpha,K}$ daje
jedyny reprezentant stopnia mniejszego niż $d$.
Dlatego każdy element ma postać
$\sum_{j=0}^{d-1}a_j\alpha^j$ z $a_j\in K$.
Jeśli taka suma jest zerowa, to odpowiedni wielomian
leży w $(m_{\alpha,K})$ i ma stopień mniejszy niż $d$;
musi więc być zerowy. Współczynniki są zerowe,
a wskazana rodzina jest bazą.
\end{proof}

\begin{przyklad}[Rozszerzenie przez pierwiastek z dwóch]
\label{aa:prz:sqrt-two-extension}
Wielomian $X^2-2$ jest nierozkładalny nad
$\mathbb Q$: gdyby miał wymierny pierwiastek $a/b$
w postaci nieskracalnej, to z $a^2=2b^2$ wynikałoby,
że $a$ jest parzyste, a potem także $b$, sprzecznie
z nieskracalnością. Dlatego
$[\mathbb Q(\sqrt2):\mathbb Q]=2$, a każdy element
ma jednoznacznie postać $a+b\sqrt2$ z $a,b\in\mathbb Q$.
Dla $a^2-2b^2\ne0$ odwrotność otrzymujemy przez
pomnożenie przez sprzężenie:
\[
(a+b\sqrt2)^{-1}=\frac{a-b\sqrt2}{a^2-2b^2}.
\]
Mianownik nie może być zerem dla $a+b\sqrt2\ne0$,
bo równość $a^2=2b^2$ przy $b\ne0$ dawałaby
wymierny pierwiastek z dwóch.
\end{przyklad}

\begin{przyklad}[Rozszerzenie o dwa pierwiastki kwadratowe]
\label{aa:prz:biquadratic-extension}
Niech $E=\mathbb Q(\sqrt2,\sqrt3)\subseteq\mathbb R$.
Wiemy, że $[\mathbb Q(\sqrt2):\mathbb Q]=2$.
Sprawdźmy, że $\sqrt3\notin\mathbb Q(\sqrt2)$.
Gdyby $\sqrt3=a+b\sqrt2$ z $a,b\in\mathbb Q$,
to po podniesieniu do kwadratu mielibyśmy
$3=a^2+2b^2+2ab\sqrt2$. Niewymierność $\sqrt2$
wymusza $ab=0$. Jeśli $b=0$, mielibyśmy
$3=a^2$ z $a\in\mathbb Q$, co jest niemożliwe:
wykładnik liczby pierwszej $3$ w rozkładzie kwadratu
liczby wymiernej jest parzysty, a w $3$ wynosi $1$.
Jeśli $a=0$, mielibyśmy $3=2b^2$, czyli
$b^2=3/2$, również niemożliwe, ponieważ
wykładnik liczby pierwszej $2$ w rozkładzie
kwadratu liczby wymiernej jest parzysty,
a w $3/2$ jest nieparzysty.
Wobec tego $X^2-3$ jest nierozkładalny nad
$\mathbb Q(\sqrt2)$ i
$[E:\mathbb Q(\sqrt2)]=2$.
Twierdzenie o wieży daje $[E:\mathbb Q]=4$,
a bazą jest $1,\sqrt2,\sqrt3,\sqrt6$.
\end{przyklad}

\begin{stwierdzenie}[Rozszerzenie skończone jest algebraiczne]
\label{aa:st:finite-implies-algebraic}
Niech $L/K$ będzie rozszerzeniem ciał stopnia
$[L:K]=d<\infty$. Wówczas każdy element $\alpha\in L$
jest algebraiczny nad $K$, a
$\deg m_{\alpha,K}\le d$.
\end{stwierdzenie}
\begin{proof}
Rodzina $1,\alpha,\ldots,\alpha^d$ ma $d+1$
elementów w przestrzeni wektorowej wymiaru $d$,
więc jest liniowo zależna. Istnieją
$a_0,\dots,a_d\in K$, nie wszystkie zerowe,
takie, że $\sum_{j=0}^d a_j\alpha^j=0$.
Oznacza to, że $\alpha$ zeruje niezerowy wielomian
stopnia co najwyżej $d$. Wielomian minimalny dzieli
ten wielomian, więc jego stopień nie przekracza $d$.
\end{proof}

\begin{wniosek}[Dołączenie skończenie wielu elementów algebraicznych]
\label{aa:wn:finite-algebraic-generators}
Niech $L/K$ będzie rozszerzeniem ciał, $r\ge1$,
a $\alpha_1,\ldots,\alpha_r\in L$ będą algebraiczne nad $K$.
Wówczas
$K(\alpha_1,\ldots,\alpha_r)/K$ jest rozszerzeniem
skończonym. Nie oznacza to, że każde rozszerzenie
algebraiczne ma skończony stopień.
\end{wniosek}
\begin{proof}
Niech $K_0=K$, $K_j=K(\alpha_1,\ldots,\alpha_j)$.
Wielomian o współczynnikach w $K$ zerujący
$\alpha_j$ jest także wielomianem nad $K_{j-1}$.
Zatem $\alpha_j$ jest algebraiczny nad $K_{j-1}$
i $[K_j:K_{j-1}]$ jest skończony.
Stosując kolejno twierdzenie o wieży dostajemy
$[K_r:K]=\prod_{j=1}^r[K_j:K_{j-1}]<\infty$.
Natomiast rozszerzenie
$E=\mathbb Q(2^{1/n}:n\ge1)\subseteq\mathbb R$
jest algebraiczne, lecz nieskończonego stopnia.
Istotnie, każdy jego element należy do ciała
$\mathbb Q(2^{1/n_1},\ldots,2^{1/n_k})$ dla
pewnej skończonej rodziny indeksów i jest algebraiczny
na mocy właśnie udowodnionego faktu. Z drugiej strony
dla każdego $n\ge1$ wielomian $X^n-2$ jest
nierozkładalny nad $\mathbb Q$ na mocy
Twierdzenia~\ref{aa:tw:eisenstein} (kryterium Eisensteina)
dla $p=2$: wszystkie współczynniki poza wiodącym
są podzielne przez $2$, a wyraz wolny $-2$
nie jest podzielny przez $4$. Stąd
$[\mathbb Q(2^{1/n}):\mathbb Q]=n$.
Gdyby $[E:\mathbb Q]$ było skończone, twierdzenie
o wieży ograniczałoby wszystkie te stopnie przez
$[E:\mathbb Q]$, co jest niemożliwe dla dowolnie
dużych $n$.
\end{proof}

\begin{twierdzenie}[Konstrukcja przez dołączenie formalnego pierwiastka]
\label{aa:tw:adjoin-root-abstract}
Niech $K$ będzie ciałem, a $f\in K[X]$ wielomianem
nierozkładalnym dodatniego stopnia. Wówczas
$E=K[X]/(f)$ jest ciałem zawierającym kanoniczny obraz
$K$, a $\alpha=[X]\in E$ spełnia $f(\alpha)=0$.
Ponadto $[E:K]=\deg f$. Dla każdego rozszerzenia
$L/K$ i pierwiastka $\beta\in L$ wielomianu $f$
istnieje dokładnie jedno $K$-zanurzenie
$E\to L$ wysyłające $\alpha$ na $\beta$.
Przez $K$-zanurzenie rozumiemy różnowartościowy
homomorfizm ciał, który na $K$ jest tożsamością.
\end{twierdzenie}
\begin{proof}
Nierozkładalność $f$ i algorytm Euklidesa w $K[X]$
pokazują, tak jak w dowodzie
Twierdzenia~\ref{aa:tw:simple-algebraic-extension},
że każdy niezerowy element ilorazu $K[X]/(f)$
ma odwrotność. Włączenie wielomianów stałych
$K\hookrightarrow K[X]/(f)$ jest różnowartościowe:
żaden niezerowy wielomian stały nie jest wielokrotnością
wielomianu dodatniego stopnia $f$.
Z definicji ilorazu $[f(X)]=0$, więc $f([X])=0$.
Jednoznaczność reszt stopnia mniejszego niż $\deg f$
daje bazę $1,\alpha,\dots,\alpha^{\deg f-1}$.

Jeżeli $f(\beta)=0$, to homomorfizm podstawiania
$\operatorname{ev}_\beta:K[X]\to L$ zeruje $(f)$.
Zatem przechodzi na homomorfizm
$\overline{\operatorname{ev}}_\beta:K[X]/(f)\to L$,
który wysyła $[X]$ na $\beta$ i pozostawia
współczynniki z $K$ bez zmian. Jest różnowartościowy,
ponieważ jego dziedzina jest ciałem, więc jądro jest
ideałem równym $0$ albo całemu ciału; drugi przypadek
wyklucza warunek $\overline{\operatorname{ev}}_\beta(1)=1$.
Każdy homomorfizm
o wskazanych wartościach musi wysyłać klasę
$\sum_j a_j[X]^j$ na $\sum_j a_j\beta^j$,
co daje jednoznaczność.
\end{proof}

\begin{przyklad}[Ciało czteroelementowe]
\label{aa:prz:four-element-field}
W $\mathbb F_2[X]$ wielomian $f=X^2+X+1$ nie ma
pierwiastków, bo $f(0)=f(1)=1$; dla stopnia dwa
oznacza to nierozkładalność. Definiujemy
$\mathbb F_4=\mathbb F_2[X]/(f)$ i
$\alpha=[X]$. Wtedy $\alpha^2=\alpha+1$,
ponieważ w charakterystyce dwa mamy $-1=1$.
Ciało zawiera dokładnie cztery elementy
$0,1,\alpha,1+\alpha$, a mnożenie wyznacza
relacja $\alpha^2=\alpha+1$.
W szczególności
$\alpha(\alpha+1)=\alpha^2+\alpha=1$,
więc $\alpha^{-1}=\alpha+1$. Co więcej,
$\alpha^3=\alpha(\alpha+1)=1$, zatem $\alpha$
ma rząd trzy w grupie multiplikatywnej, ponieważ
$\alpha\ne1$ i $3$ jest liczbą pierwszą.
\end{przyklad}

\section{Pierwiastki i ciała rozkładu}
\begin{definicja}[Ciało rozkładu]
\label{aa:def:splitting-field}
Niech $K$ będzie ciałem, a $f\in K[X]$ niezerowym wielomianem
dodatniego stopnia. Rozszerzenie $L/K$ nazywamy
\emph{ciałem rozkładu wielomianu $f$}, jeśli
$f$ rozkłada się w $L[X]$ na czynniki liniowe
(z uwzględnieniem jego niezerowego współczynnika wiodącego)
oraz $L$ jest generowane nad $K$ przez wszystkie
pierwiastki $f$ w $L$.
\end{definicja}

\begin{twierdzenie}[Istnienie ciała rozkładu]
\label{aa:tw:splitting-existence}
Dla każdego ciała $K$ i każdego niezerowego
$f\in K[X]$ dodatniego stopnia istnieje skończone
rozszerzenie $L/K$ będące ciałem rozkładu $f$.
\end{twierdzenie}
\begin{proof}
Przeprowadzimy indukcję po stopniu wielomianu.
Jeśli $\deg f=1$, ma on pierwiastek w $K$,
a jako ciało rozkładu możemy wziąć $L=K$.
Załóżmy, że twierdzenie jest prawdziwe dla
wszystkich wielomianów o stopniu mniejszym niż
$n$ i niech $\deg f=n\ge2$.
Wybierzmy nierozkładalny czynnik $g$ wielomianu $f$
w $K[X]$; taki czynnik istnieje, gdyż przy każdym
nietrywialnym rozkładzie stopnie czynników są mniejsze,
więc kolejne rozkładanie musi się zakończyć. Następnie
zbudujmy ciało
$E=K[X]/(g)$ z pierwiastkiem $\alpha=[X]$.
Twierdzenie~\ref{aa:tw:adjoin-root-abstract}
pokazuje, że $[E:K]=\deg g<\infty$.
W $E[X]$ możemy podzielić $f$ przez $X-\alpha$,
ponieważ $f(\alpha)=0$. Otrzymujemy
$f=(X-\alpha)h$ z $h\in E[X]$ i
$\deg h=n-1$.
Z założenia indukcyjnego istnieje skończone
rozszerzenie $L/E$, w którym $h$ rozkłada się
na czynniki liniowe i które jest generowane przez
pierwiastki $h$. Wtedy w $L$ rozkłada się również
$f$, a $L$ jest generowane nad $K$ przez
$\alpha$ oraz pierwiastki $h$, czyli przez
pierwiastki $f$. Na mocy twierdzenia o wieży
rozszerzenie $L/K$ jest skończone.
\end{proof}

\begin{przyklad}[Ciało rozkładu wielomianu sześciennego]
\label{aa:prz:cubic-splitting}
Wielomian $f=X^3-2$ jest nierozkładalny nad
$\mathbb Q$ na mocy Twierdzenia~\ref{aa:tw:eisenstein}
(kryterium Eisensteina) dla $p=2$.
Niech $\alpha=\sqrt[3]2\in\mathbb R$ i
$\omega=(-1+\sqrt{-3})/2\in\mathbb C$.
Bezpośredni rachunek daje $\omega^2+\omega+1=0$,
a $\omega\ne1$; z tożsamości
$(X-1)(X^2+X+1)=X^3-1$ wynika $\omega^3=1$.
Pierwiastkami $f$ są
$\alpha,\alpha\omega,\alpha\omega^2$.
Ciało rozkładu ma postać
$L=\mathbb Q(\alpha,\omega)$.
Mamy $[\mathbb Q(\alpha):\mathbb Q]=3$.
Wielomian $X^2+X+1$ nie ma pierwiastków rzeczywistych:
gdyby $x\in\mathbb R$ był jego pierwiastkiem, to
$(2x+1)^2=-3$, co jest niemożliwe. Ponieważ
$\mathbb Q(\alpha)\subseteq\mathbb R$, wielomian ten
nie ma pierwiastka również w $\mathbb Q(\alpha)$;
ma stopień dwa, więc jest nad nim nierozkładalny.
Stąd $[L:\mathbb Q(\alpha)]=2$, a twierdzenie
o wieży daje $[L:\mathbb Q]=6$.
Nie wystarcza samo dołączenie rzeczywistego
pierwiastka: $\mathbb Q(\alpha)$ nie zawiera
pozostałych dwóch pierwiastków.
\end{przyklad}

\section{Rozdzielczość wielomianów}
\begin{definicja}[Pierwiastek wielokrotny i rozdzielczość]
\label{aa:def:separable-polynomial}
Niech $K$ będzie ciałem i $f\in K[X]$ niezerowym
wielomianem dodatniego stopnia. Niech $r\ge1$ będzie
liczbą całkowitą. Pierwiastek $\alpha$
w pewnym rozszerzeniu $L/K$ ma \emph{krotność co najmniej
$r$}, jeśli $(X-\alpha)^r\mid f$ w $L[X]$.
Wielomian $f$ jest \emph{rozdzielczy}, jeśli ma
$\deg f$ różnych pierwiastków w swoim ciele rozkładu.
\end{definicja}

\begin{twierdzenie}[Kryterium rozdzielczości przez pochodną]
\label{aa:tw:separability-derivative}
Niech $K$ będzie ciałem, a $f=\sum_{j=0}^n a_jX^j\in K[X]$
niezerowym wielomianem stopnia $n\ge1$.
Wtedy $f$ jest rozdzielczy dokładnie wtedy, gdy
$\gcd(f,f')=1$ w $K[X]$, gdzie
$f'=\sum_{j=1}^n ja_jX^{j-1}$ jest formalną
pochodną wielomianu $f$.
\end{twierdzenie}
\begin{proof}
Dla $\alpha$ w rozszerzeniu $L/K$ zapiszmy
$f=(X-\alpha)q$. Reguła różniczkowania iloczynu
w pierścieniu wielomianów daje
$f'=q+(X-\alpha)q'$, więc $f'(\alpha)=q(\alpha)$.
Ponieważ $X-\alpha$ dzieli $q$ dokładnie wtedy,
gdy $q(\alpha)=0$, pierwiastek $\alpha$ ma krotność
co najmniej dwa wtedy i tylko wtedy, gdy
$f(\alpha)=f'(\alpha)=0$.

Jeśli $d=\gcd(f,f')$ ma dodatni stopień,
możemy wybrać jego nierozkładalny czynnik $h$.
W ciele rozkładu $h$ ma pierwiastek $\alpha$;
wtedy $f(\alpha)=f'(\alpha)=0$, więc $f$
ma pierwiastek wielokrotny i nie jest rozdzielczy.
Odwrotnie, jeśli $f$ ma pierwiastek wielokrotny
$\alpha$ w ciele rozkładu $L$, to oba wielomiany
$f,f'$ zerują $\alpha$. Wielomian minimalny
$m_{\alpha,K}$ dzieli więc jednocześnie $f$ i $f'$
na mocy Lematu~\ref{aa:lem:minimal-polynomial}.
Stąd $\deg\gcd(f,f')\ge1$. Dowiedliśmy równoważności.
\end{proof}

\begin{wniosek}[Nierozkładalność w charakterystyce zero]
\label{aa:wn:charzero-separable}
Każdy nierozkładalny wielomian dodatniego stopnia
nad ciałem charakterystyki zero jest rozdzielczy.
\end{wniosek}
\begin{proof}
Dla $f\in K[X]$ nierozkładalnego mamy
$f'\ne0$ w charakterystyce zero oraz
$\deg f'<\deg f$. Jedyny wspólny dzielnik
$f$ i $f'$ może być jednostką, bo każdy
niejednostkowy dzielnik $f$ jest stowarzyszony z $f$.
Zatem $\gcd(f,f')=1$ i poprzednie twierdzenie daje
rozdzielczość.
\end{proof}

\begin{przyklad}[Rozszerzenie nierozdzielcze]
\label{aa:prz:inseparable-example}
Niech $K=\mathbb F_p(t)$ i $f=X^p-t\in K[X]$.
Wielomian ten jest nierozkładalny nad $K$:
w $\mathbb F_p[t][X]$ spełnia założenia
Twierdzenia~\ref{aa:tw:eisenstein} (kryterium Eisensteina)
dla elementu pierwszego $t$, ponieważ $t$
dzieli współczynniki niewiodące, lecz $t^2\nmid -t$.
Wspomniane kryterium stwierdza nierozkładalność
właśnie nad ciałem ułamków $K$.
Jednak formalna pochodna $f'=pX^{p-1}=0$,
więc $\gcd(f,f')=f\ne1$. Wielomian jest
nierozdzielczy. Jeśli $\alpha^p=t$ w pewnym
rozszerzeniu, to w charakterystyce $p$ zachodzi
$X^p-t=(X-\alpha)^p$. Wielomian ma tylko jeden
różny pierwiastek w swoim ciele rozkładu,
i występuje on z krotnością $p$.
\end{przyklad}

\section{Ciała skończone}
\begin{twierdzenie}[Liczność ciała skończonego]
\label{aa:tw:finite-field-cardinality}
Każde ciało skończone $F$ ma $p^n$ elementów
dla pewnej liczby pierwszej $p$ i pewnej liczby
całkowitej $n\ge1$.
\end{twierdzenie}
\begin{proof}
Ciało $F$ ma dodatnią charakterystykę $p$,
ponieważ homomorfizm $\mathbb Z\to F$ nie może być
różnowartościowy: $\mathbb Z$ jest nieskończony,
a $F$ skończone. Z
Twierdzenia~\ref{aa:tw:prime-subfield} wynika,
że $F$ zawiera $\mathbb F_p$. Jako skończona
przestrzeń wektorowa nad tym podciałem ma pewien
wymiar $n\ge1$. Wybierając bazę, utożsamiamy
każdy element z jednoznaczną $n$-ką współczynników
z $\mathbb F_p$. Istnieje dokładnie $p^n$ takich
$n$-ek, zatem $|F|=p^n$.
\end{proof}

\begin{twierdzenie}[Grupa multiplikatywna ciała skończonego]
\label{aa:tw:finite-field-multiplicative-cyclic}
Dla każdego skończonego ciała $F$ grupa
$(F^\times,\cdot)$ jest cykliczna.
\end{twierdzenie}
\begin{proof}
Grupa $G=F^\times$ jest skończona i abelowa.
Jeśli $|G|=1$ (w szczególności dla $F=\mathbb F_2$),
to $G=\{1\}$ jest cykliczna. Dalej załóżmy $|G|>1$.
Na mocy Wniosku~\ref{aa:wn:abelian-fg-structure}
o klasyfikacji skończenie generowanych grup abelowych
jest sumą
prostą skończenie wielu grup cyklicznych.
Dla każdej liczby pierwszej $\ell$ dzielącej $|G|$
wybierzmy składnik cykliczny o największym rzędzie
$\ell^{a_\ell}$ w części $\ell$-pierwotnej
i generator $g_\ell$ tego składnika. W części
$\ell$-pierwotnej rząd każdego elementu dzieli
$\ell^{a_\ell}$, gdyż każdy jej składnik cykliczny
ma rząd dzielący $\ell^{a_\ell}$. Dlatego
wykładnik grupy $G$ (najmniejsza dodatnia liczba $e$
spełniająca $x^e=1$ dla wszystkich $x\in G$) jest
równy $e=\prod_\ell\ell^{a_\ell}$: liczba ta
anihiluje wszystkie składniki, a musi dzielić się
przez rząd każdego $g_\ell$. Element
$g=\prod_\ell g_\ell$ ma rząd równy $e$, gdyż
rzędy jego czynników są parami względnie pierwsze:
z $g^k=1$ wynika, po rzutowaniu na każdą część
pierwotną, $g_\ell^k=1$, więc
$\ell^{a_\ell}\mid k$ dla wszystkich $\ell$.
Odwrotnie, $g^e=1$.

Każdy element $G$ jest pierwiastkiem wielomianu
$X^e-1\in F[X]$. Niezerowy wielomian stopnia $e$
nad ciałem ma co najwyżej $e$ różnych pierwiastków,
zatem $|G|\le e$. Z drugiej strony $e$ jest
rzędem pewnego elementu $G$, więc z twierdzenia
Lagrange'a wynika $e\le|G|$. Stąd $e=|G|$,
a wybrany element rzędu $e$ generuje całe $G$.
\end{proof}

\begin{przyklad}[Dlaczego $\mathbb Z/4\mathbb Z$ nie jest ciałem]
\label{aa:prz:not-field-four}
Pierścień $\mathbb Z/4\mathbb Z$ ma cztery elementy,
ale $[2]^2=[0]$ i $[2]\ne[0]$, więc ma dzielniki
zera. Nie jest ciałem, mimo że istnieje
czteroelementowe ciało $\mathbb F_4$ skonstruowane
w Przykładzie~\ref{aa:prz:four-element-field}.
Równa liczba elementów nie oznacza izomorfizmu
struktur algebraicznych.
\end{przyklad}

\begin{uwaga}[Kolejna część: teoria Galois]
\label{aa:uw:field-extensions-future}
Skonstruowaliśmy rozszerzenia przez pierwiastki,
ciała rozkładu i poznaliśmy rolę rozdzielczości.
W następnych rozdziałach zbadamy przedłużanie
zanurzeń, automorfizmy rozszerzeń, normalność
oraz podstawowe twierdzenie teorii Galois.
Do wyjaśnienia, które równania można rozwiązać
przez pierwiastniki, potrzebne będą ponadto grupy
rozwiązalne z wcześniejszej części skryptu.
\end{uwaga}

% AA-21 | Algebra abstrakcyjna | Iteracja XI
\chapter{Zanurzenia ciał, automorfizmy i rozszerzenia Galois}
\label{aa:ch:embeddings-normal-separable}

\section{Dlaczego pierwiastki można przestawiać?}
W Rozdziale~\ref{aa:ch:field-extensions} skonstruowaliśmy ciała
rozkładu wielomianów. Rozpatrzmy $L=\mathbb Q(\sqrt2)$.
Każdy element ma jednoznacznie postać $a+b\sqrt2$, a przyporządkowanie
$a+b\sqrt2\mapsto a-b\sqrt2$ zachowuje dodawanie, mnożenie oraz
wszystkie liczby wymierne. Nie jest to przypadek: każdy homomorfizm
ciał zachowuje równania wielomianowe, zatem musi wysyłać pierwiastek
wielomianu na inny pierwiastek tego samego wielomianu. Nie wszystkie
takie zamiany są jednak możliwe wewnątrz rozważanego ciała.

W całym rozdziale wszystkie ciała są przemienne, a ich homomorfizmy
zachowują jedynkę. Dla stosowania zanurzeń wygodnie jest pracować
w ciele, w którym każdy wielomian ma pierwiastek. Następująca
konstrukcja uzasadnia istnienie takiego rozszerzenia bez zakładania
go jako nieudowodnionego postulatu.

\section{Istnienie domknięcia algebraicznego}
\begin{definicja}[Ciało algebraicznie domknięte i domknięcie algebraiczne]
\label{aa:def:algebraic-closure}
Ciało $\Omega$ jest \emph{algebraicznie domknięte}, jeśli każdy
wielomian dodatniego stopnia z $\Omega[X]$ ma pierwiastek w $\Omega$.
\emph{Domknięciem algebraicznym} ciała $K$ nazywamy rozszerzenie
algebraiczne $\Omega/K$, w którym $\Omega$ jest algebraicznie domknięte.
\end{definicja}

\begin{twierdzenie}[Istnienie domknięcia algebraicznego]
\label{aa:tw:algebraic-closure-exists}
Każde ciało $K$ ma domknięcie algebraiczne.
\end{twierdzenie}
\begin{proof}
\emph{Krok 1: jednoczesne dołączenie pierwiastków wielomianów nad $K$.}
Niech $\mathcal P$ oznacza zbiór wszystkich unormowanych
wielomianów dodatniego stopnia z $K[X]$. Dla każdego
$f\in\mathcal P$ wprowadźmy osobną zmienną $Y_f$.
W pierścieniu wielomianów
$A=K[Y_f:f\in\mathcal P]$ rozważmy ideał
$I=(f(Y_f):f\in\mathcal P)$.
Wykazujemy, że $I$ jest właściwy.
Gdyby $1\in I$, istniałaby skończona relacja
$1=\sum_{j=1}^r h_jf_j(Y_{f_j})$ z $h_j\in A$.
Na mocy konstrukcji przez kolejne dołączanie
pierwiastków (Twierdzenie~\ref{aa:tw:adjoin-root-abstract})
możemy kolejno dołączać pierwiastki $f_1,\ldots,f_r$.
W każdym kroku otrzymujemy rozszerzenie skończone,
a twierdzenie o wieży zapewnia, że wynikowe $E/K$
również jest skończone. Wybierzmy w $E$ pierwiastek
$\alpha_j$ każdego $f_j$. Możemy zdefiniować homomorfizm
$\operatorname{ev}:A\to E$, podstawiając
$Y_{f_j}\mapsto\alpha_j$ i pozostałe zmienne
na przykład na $0$; jeśli dany wielomian występuje
wiele razy, wybieramy dla niego jedną wartość.
Po zastosowaniu $\operatorname{ev}$ otrzymalibyśmy
$1=\sum_j\operatorname{ev}(h_j)f_j(\alpha_j)=0$.
Jest to sprzeczność, więc $I\ne A$.

Rozważmy rodzinę wszystkich ideałów właściwych pierścienia $A$
zawierających $I$, uporządkowaną przez inkluzję. Jest ona
niepusta, a suma mnogościowa każdego niepustego łańcucha
z tej rodziny jest ideałem właściwym: gdyby zawierała $1$, element $1$
należałby już do jednego z ideałów łańcucha. Zatem
lemat Kuratowskiego--Zorna daje element maksymalny tej
rodziny (pusty łańcuch jest ograniczony przez $I$). Taki ideał $\mathfrak m$ jest maksymalnym ideałem
$A$: każdy ideał większy od $\mathfrak m$ zawiera $I$,
więc nie może być właściwy. W konsekwencji iloraz
$E=A/\mathfrak m$ jest ciałem.
Ponieważ $I$ jest właściwy, homomorfizm
$K\to E$ jest różnowartościowy: jego jądro jest
ideałem ciała $K$ i nie może zawierać $1$.
Dla każdego $f\in\mathcal P$ obraz $y_f$ zmiennej
$Y_f$ w $E$ spełnia $f(y_f)=0$.
Ciało $E=A/\mathfrak m$ jest generowane jako algebra nad $K$
przez elementy $y_f$, ponieważ każdy jego element
jest obrazem wielomianu z $A$. Każdy $y_f$ jest
algebraiczny nad $K$, a obraz dowolnego wielomianu
z $A$ zależy tylko od skończenie wielu zmiennych.
Zatem każdy $x\in E$ należy do
$K(y_{f_1},\ldots,y_{f_s})$, skończonego rozszerzenia
$K$ na mocy wielokrotnego zastosowania twierdzenia
o wieży. Każdy taki $x$ jest więc algebraiczny nad $K$.
Otrzymaliśmy rozszerzenie algebraiczne $E/K$, w którym
każdy niezerowy wielomian dodatniego stopnia z $K[X]$
ma pierwiastek.

\emph{Krok 2: iteracja konstrukcji.}
Zaczynamy od $K_0=K$. Dla każdego $n\ge0$ stosujemy
Krok 1 do $K_n$ i otrzymujemy rozszerzenie algebraiczne
$K_{n+1}/K_n$ takie, że każdy wielomian dodatniego
stopnia z $K_n[X]$ ma pierwiastek w $K_{n+1}$.
Utożsamiając kolejne ciała z podciałami większych,
określamy $\Omega=\bigcup_{n\ge0}K_n$.
Jest ono ciałem: każde skończenie wiele jego elementów
należy do wspólnego $K_n$, w którym można wykonywać
działania i brać odwrotności. Rozszerzenie $\Omega/K$
jest algebraiczne, gdyż każdy $x\in\Omega$ należy
do pewnego $K_n$. Uzasadnijmy używaną tu
przechodniość algebraiczności. Jeśli $B/A$ jest
algebraiczne, a $x$ jest algebraiczny nad $B$,
to współczynniki wielomianu, który zeruje $x$,
należą do pewnego skończonego rozszerzenia
$A(b_0,\ldots,b_d)/A$. Zatem $x$ należy do skończonego
rozszerzenia $A(b_0,\ldots,b_d,x)/A$, na mocy
twierdzenia o wieży, i jest algebraiczny nad $A$.
Stosując to skończenie wiele razy do wieży
$K_0\subseteq K_1\subseteq\cdots\subseteq K_n$,
widzimy, że $x$ jest algebraiczny nad $K$.

Niech teraz $f\in\Omega[X]$ ma dodatni stopień.
Wszystkie jego współczynniki należą do pewnego $K_n$,
więc $f$ ma pierwiastek w $K_{n+1}\subseteq\Omega$.
Zatem $\Omega$ jest algebraicznie domknięte i stanowi
domknięcie algebraiczne $K$.
W powyższej konstrukcji istnienie ideału maksymalnego wynika
z lematu Kuratowskiego--Zorna; bez tego kroku nie można zakładać,
że iloraz przez dowolny ideał będzie ciałem.
\end{proof}

W dalszej części przez $\Omega$ oznaczamy odpowiednie
domknięcie algebraiczne zawierające badane skończone
rozszerzenie. Konstrukcja przez wybór kolejnych pierwiastków
pozwala też umieścić każde ustalone rozszerzenie algebraiczne
w takim ciele. Wszystkie użyte niżej przedłużenia zanurzeń
dotyczą rozszerzeń skończonych.

\begin{definicja}[Zanurzenie nad ciałem bazowym i automorfizm]
\label{aa:def:field-k-embedding}
Niech $L/K$ oraz $M/K$ będą rozszerzeniami ciał.
\emph{$K$-zanurzeniem} $L$ w $M$ nazywamy homomorfizm
ciał $\sigma:L\to M$ taki, że $\sigma(a)=a$ dla każdego $a\in K$.
\emph{$K$-automorfizmem} $L$ nazywamy bijektywne $K$-zanurzenie
$L\to L$. Zbiór tych automorfizmów oznaczamy
$\operatorname{Aut}_K(L)$. Jest on grupą względem składania:
składanie i odwrotność bijektywnych homomorfizmów ciał
ponownie zachowują działania, jedynkę i każdy element $K$.
Jeśli $\varphi:K\to K'$ jest dowolnym zanurzeniem, to
$\varphi$ przenosi wielomian $f=\sum a_iX^i\in K[X]$
na $f^\varphi=\sum\varphi(a_i)X^i\in K'[X]$.
\end{definicja}

\begin{stwierdzenie}[Automorfizm zachowuje wielomian minimalny]
\label{aa:st:automorphism-minimal-polynomial}
Jeśli $L/K$ jest rozszerzeniem algebraicznym,
$\sigma\in\operatorname{Aut}_K(L)$ i $\alpha\in L$, to
$m_{\sigma(\alpha),K}=m_{\alpha,K}$.
W szczególności $m_{\alpha,K}(\sigma(\alpha))=0$,
a więc $\sigma(\alpha)$ jest jednym z pierwiastków
$m_{\alpha,K}$ leżących w $L$.
\end{stwierdzenie}
\begin{proof}
Dla $f(X)=\sum_{j=0}^d a_jX^j\in K[X]$ mamy
$\sigma(f(\alpha))=\sum a_j\sigma(\alpha)^j=f(\sigma(\alpha))$,
ponieważ $\sigma(a_j)=a_j$. Stosujemy tę równość do
$f=m_{\alpha,K}$, dla którego $f(\alpha)=0$.
Ponadto $m_{\alpha,K}$ jest unormowany i nierozkładalny,
więc jest także wielomianem minimalnym $\sigma(\alpha)$:
oba wielomiany są unormowane, a minimalny dzieli
$m_{\alpha,K}$.
\end{proof}

\begin{przyklad}[Dwa zanurzenia rozszerzenia kwadratowego]
\label{aa:prz:quadratic-conjugation}
Niech $L=\mathbb Q(\sqrt2)$. Są dokładnie dwa
$\mathbb Q$-automorfizmy $L$: identyczność i
$\tau(a+b\sqrt2)=a-b\sqrt2$. Istotnie, obraz $\sqrt2$
musi należeć do $\{\sqrt2,-\sqrt2\}$, a ten obraz jednoznacznie
wyznacza homomorfizm. Wzór na $\tau$ zachowuje mnożenie, ponieważ
$(a+b\sqrt2)(c+d\sqrt2)=(ac+2bd)+(ad+bc)\sqrt2$.
Wynika stąd $\operatorname{Aut}_{\mathbb Q}(L)\cong C_2$.
\end{przyklad}

\section{Przedłużanie zanurzeń}
\begin{lemat}[Przedłużenie zanurzenia na proste rozszerzenie]
\label{aa:lem:extend-embedding-simple}
Niech $\varphi:K\to K'$ będzie izomorfizmem ciał,
niech $\alpha$ będzie algebraiczny nad $K$, a
$m=m_{\alpha,K}\in K[X]$ jego wielomianem minimalnym.
Jeżeli $M/K'$ jest rozszerzeniem oraz $\beta\in M$ spełnia
$m^\varphi(\beta)=0$, to istnieje dokładnie jedno zanurzenie
$\widetilde\varphi:K(\alpha)\to M$ przedłużające
$\varphi$ i spełniające $\widetilde\varphi(\alpha)=\beta$.
\end{lemat}
\begin{proof}
Z Twierdzenia~\ref{aa:tw:simple-algebraic-extension} mamy
$K(\alpha)\cong K[X]/(m)$. Na przedstawicielach definiujemy
$\widetilde\varphi([f])=f^\varphi(\beta)$.
Jeśli $[f]=[g]$, to $f-g=mh$ dla pewnego $h\in K[X]$;
stąd $f^\varphi(\beta)-g^\varphi(\beta)=
m^\varphi(\beta)h^\varphi(\beta)=0$.
Zatem wzór jest poprawnie określony. Zachowanie działań
i jedynki wynika z tożsamości
$(f+g)^\varphi=f^\varphi+g^\varphi$ oraz
$(fg)^\varphi=f^\varphi g^\varphi$ i z własności
podstawiania wielomianów. Jądro homomorfizmu z ciała
jest ideałem, a ponieważ $1$ przechodzi na $1$, nie
jest całym ciałem. Zatem jądro jest zerowe i odwzorowanie
jest różnowartościowe; ponadto przyjmuje żądane wartości
na $K$ i $\alpha$.
Każdy element $K(\alpha)$ jest postaci $f(\alpha)$ dla
pewnego $f\in K[X]$, więc wartości na $K$ i $\alpha$
wyznaczają całe odwzorowanie jednoznacznie.
\end{proof}

\begin{wniosek}[Przedłużenie zanurzenia na rozszerzenie skończone]
\label{aa:wn:extend-embedding-finite}
Niech $L/K$ będzie rozszerzeniem skończonym,
$\varphi:K\hookrightarrow\Omega$ zanurzeniem w ciało
algebraicznie domknięte $\Omega$. Wtedy istnieje zanurzenie
$\widetilde\varphi:L\hookrightarrow\Omega$ przedłużające
$\varphi$. To samo zachodzi dla zadanego zanurzenia dowolnego
ciała pośredniego $E$, $K\subseteq E\subseteq L$.
\end{wniosek}
\begin{proof}
Ponieważ $[L:K]<\infty$, możemy wybrać skończony zbiór
generatorów $L=K(\alpha_1,\ldots,\alpha_r)$, na przykład
dowolną bazę przestrzeni wektorowej $L$ nad $K$.
Budujemy kolejne ciała $K_i=K(\alpha_1,\ldots,\alpha_i)$.
Jeżeli zanurzenie $\varphi_{i-1}:K_{i-1}\hookrightarrow\Omega$
już skonstruowano, to wielomian minimalny $\alpha_i$
nad $K_{i-1}$ po przeniesieniu współczynników przez
$\varphi_{i-1}$ ma pierwiastek w $\Omega$. Lemat
\ref{aa:lem:extend-embedding-simple}, stosowany z
$K'=\varphi_{i-1}(K_{i-1})$, daje jego przedłużenie
na $K_i$. Po skończonej liczbie kroków dostajemy
zanurzenie całego $L$. Dla ciała pośredniego $E$
powtarzamy argument, biorąc bazę $L$ nad $E$;
stopień $[L:E]$ jest skończony na mocy
Twierdzenia~\ref{aa:tw:field-tower}.
\end{proof}

\begin{twierdzenie}[Liczba zanurzeń i rozdzielczość]
\label{aa:tw:embeddings-count}
Niech $L/K$ będzie rozszerzeniem skończonym,
a $\Omega$ ciałem algebraicznie domkniętym zawierającym
$L$. Liczba różnych $K$-zanurzeń $L\hookrightarrow\Omega$
nie przekracza $[L:K]$. Jest równa $[L:K]$ wtedy
i tylko wtedy, gdy każdy element $L$ jest rozdzielczy
nad $K$, czyli gdy $L/K$ jest rozszerzeniem rozdzielczym.
\end{twierdzenie}
\begin{proof}
Zapiszmy $L=K(\alpha_1,\ldots,\alpha_r)$ oraz
$K_i=K(\alpha_1,\ldots,\alpha_i)$, $K_0=K$.
Oznaczmy $d_i=[K_i:K_{i-1}]$.
Gdy wybrano zanurzenie $\varphi_{i-1}:K_{i-1}\to\Omega$,
jego przedłużenie na $K_i$ jest jednoznacznie wyznaczone
przez obraz $\alpha_i$. Wskazany obraz musi być
pierwiastkiem wielomianu minimalnego $\alpha_i$
nad $K_{i-1}$ z przeniesionymi współczynnikami,
a takich różnych pierwiastków jest co najwyżej $d_i$.
Każdy wybór pierwiastka daje dokładnie jedno przedłużenie
na mocy Lematu~\ref{aa:lem:extend-embedding-simple}
(z izomorfizmem ciała pośredniego na jego obraz),
a różne pierwiastki dają różne przedłużenia.
Wynika stąd oszacowanie liczby zanurzeń przez
$\prod_i d_i=[L:K]$ na mocy twierdzenia o wieży.

Jeśli $L/K$ jest rozdzielcze, każdy $\alpha_i$ ma
rozdzielczy wielomian minimalny nad $K$. Wielomian
minimalny nad $K_{i-1}$ dzieli go w $K_{i-1}[X]$,
a więc także jest rozdzielczy. Po przeniesieniu
współczynników zachowuje stopień i brak pierwiastków
wielokrotnych. Ma więc dokładnie $d_i$ różnych
pierwiastków w $\Omega$ i na każdym etapie istnieje
dokładnie $d_i$ wyborów. Otrzymujemy $\prod_i d_i$
różnych zanurzeń.

Odwrotnie, załóżmy, że liczba wszystkich
$K$-zanurzeń $L\hookrightarrow\Omega$ wynosi $[L:K]$.
Dla dowolnego $\alpha\in L$ rozpocznijmy wieżę
od $K(\alpha)$, a następnie dołączmy skończenie
wielu pozostałych generatorów. Jeśli $m_{\alpha,K}$
ma mniej różnych pierwiastków niż jego stopień $d$,
na pierwszym etapie mamy mniej niż $d$ zanurzeń;
każde z nich przedłuża się do całego $L$ na co
najwyżej $[L:K(\alpha)]$ sposobów. Liczba
zanurzeń byłaby ściśle mniejsza niż
$d[L:K(\alpha)]=[L:K]$, co przeczy założeniu.
Zatem każdy $m_{\alpha,K}$ jest rozdzielczy.
\end{proof}

\begin{przyklad}[Dlaczego brak rozdzielczości zmniejsza liczbę zanurzeń]
\label{aa:prz:inseparable-embeddings}
Niech $K=\mathbb F_p(t)$ i $\alpha^p=t$. Rozszerzenie
$L=K(\alpha)$ ma stopień $p$, ponieważ $X^p-t$
jest nierozkładalny (Przykład~\ref{aa:prz:inseparable-example}).
Jednak $X^p-t=(X-\alpha)^p$ w $L[X]$, więc
każde $K$-zanurzenie $L\hookrightarrow\Omega$ musi
wysyłać $\alpha$ na $\alpha$. Jest zatem tylko jedno
takie zanurzenie, mimo że $[L:K]=p$.
\end{przyklad}

\section{Jednoznaczność ciała rozkładu}
\begin{twierdzenie}[Jednoznaczność ciała rozkładu z dokładnością do izomorfizmu]
\label{aa:tw:splitting-uniqueness}
Niech $\varphi:K\to K'$ będzie izomorfizmem ciał,
$f\in K[X]$ niezerowym wielomianem dodatniego stopnia,
$L/K$ ciałem rozkładu $f$ oraz $L'/K'$ ciałem rozkładu
$f^\varphi$. Wtedy istnieje izomorfizm $\Phi:L\to L'$
przedłużający $\varphi$. Nie musi on być jednoznaczny.
\end{twierdzenie}
\begin{proof}
Przeprowadzimy indukcję po $n=\deg f$. Dla $n=1$
oba ciała rozkładu są odpowiednio $K$ i $K'$, więc
wystarcza $\varphi$. Niech $n\ge2$ i przyjmijmy wynik
dla mniejszych stopni. Wybierzmy nierozkładalny czynnik
$g\in K[X]$ wielomianu $f$, a w $L$ pierwiastek
$\alpha$ wielomianu $g$. W $L'$ wybierzmy pierwiastek
$\beta$ wielomianu $g^\varphi$.
Lemat~\ref{aa:lem:extend-embedding-simple} daje
izomorfizm $\varphi_1:K(\alpha)\to K'(\beta)$
przedłużający $\varphi$ i wysyłający $\alpha$ na $\beta$.
W $K(\alpha)[X]$ mamy $f=(X-\alpha)h$,
gdzie $\deg h=n-1$. Po przeniesieniu współczynników
przez $\varphi_1$ otrzymujemy
$f^\varphi=(X-\beta)h^{\varphi_1}$.
Ciało $L$ jest ciałem rozkładu $h$ nad $K(\alpha)$:
rozszczepia $h$ i jest generowane nad $K(\alpha)$
przez pierwiastki $h$. Analogicznie $L'$ jest ciałem
rozkładu $h^{\varphi_1}$ nad $K'(\beta)$.
Założenie indukcyjne, stosowane do rozszerzenia bazowego
$K(\alpha)\cong K'(\beta)$, dostarcza izomorfizmu $L\to L'$
przedłużającego $\varphi_1$, a więc także $\varphi$.
\end{proof}

\section{Normalność i rozszerzenia Galois}
\begin{definicja}[Rozszerzenie normalne, rozdzielcze i Galois]
\label{aa:def:normal-separable-galois}
Niech $L/K$ będzie rozszerzeniem algebraicznym.
Nazywamy je \emph{normalnym}, jeśli każdy nierozkładalny
wielomian $f\in K[X]$ mający pierwiastek w $L$
rozkłada się w $L[X]$ na czynniki liniowe.
Rozszerzenie jest \emph{rozdzielcze}, jeśli dla każdego
$\alpha\in L$ wielomian minimalny $m_{\alpha,K}$
ma wyłącznie pierwiastki pojedyncze.
Rozszerzenie \emph{Galois} to rozszerzenie algebraiczne,
normalne i rozdzielcze; gdy dodatkowo $[L:K]<\infty$,
mówimy o \emph{skończonym rozszerzeniu Galois}.
\end{definicja}

\begin{twierdzenie}[Kryterium normalności przez zanurzenia]
\label{aa:tw:normal-embeddings}
Niech $L/K$ będzie rozszerzeniem skończonym i
$\Omega$ ciałem algebraicznie domkniętym zawierającym $L$.
Następujące warunki są równoważne:
\begin{enumerate}[label=\textup{(\roman*)}]
\item $L/K$ jest normalne;
\item dla każdego $K$-zanurzenia $\sigma:L\hookrightarrow\Omega$
      zachodzi $\sigma(L)=L$;
\item $L$ jest ciałem rozkładu pewnej skończonej rodziny
      niezerowych wielomianów z $K[X]$.
\end{enumerate}
\end{twierdzenie}
\begin{proof}
\emph{(i)$\Rightarrow$(ii).} Dla $\alpha\in L$ obraz
$\sigma(\alpha)$ zeruje $m_{\alpha,K}$.
Normalność gwarantuje, że wszystkie pierwiastki
tego wielomianu należą do $L$, stąd $\sigma(L)\subseteq L$.
Wtedy $\sigma:L\to L$ jest różnowartościowym
$K$-liniowym endomorfizmem przestrzeni wektorowej
skończonego wymiaru, więc jest surjektywny.

\emph{(ii)$\Rightarrow$(i).} Niech $f\in K[X]$ będzie
nierozkładalny i $\alpha\in L$ jego pierwiastkiem.
Dla dowolnego innego pierwiastka $\beta\in\Omega$
lemat o przedłużeniu zanurzeń daje $K$-zanurzenie
$K(\alpha)\to\Omega$ wysyłające $\alpha$ na $\beta$.
Przedłużamy je do $L$ na mocy
Wniosku~\ref{aa:wn:extend-embedding-finite}.
Z założenia jego obraz to $L$, zatem $\beta\in L$.
Wszystkie pierwiastki $f$ leżą w $L$, więc $f$
rozkłada się na czynniki liniowe w $L[X]$.

\emph{(i)$\Rightarrow$(iii).} Niech
$L=K(\alpha_1,\ldots,\alpha_r)$. Ponieważ $L/K$
jest normalne, każdy wielomian minimalny
$m_{\alpha_i,K}$ rozkłada się w $L[X]$.
Ciało generowane nad $K$ przez wszystkie pierwiastki
tych wielomianów zawiera wszystkie $\alpha_i$, a więc
jest równe $L$.

\emph{(iii)$\Rightarrow$(ii).} Każde $K$-zanurzenie
$\sigma:L\to\Omega$ wysyła pierwiastki odpowiednich
wielomianów na pierwiastki tych samych wielomianów.
Wszystkie one leżą w $L$, więc $\sigma(L)\subseteq L$;
argument wymiarowy z pierwszego kroku daje równość.
\end{proof}

\begin{twierdzenie}[Grupa Galois i stopień rozszerzenia]
\label{aa:tw:galois-group-degree}
Jeśli $L/K$ jest skończonym rozszerzeniem Galois, to
$\operatorname{Gal}(L/K):=\operatorname{Aut}_K(L)$
jest grupą skończoną oraz
$|\operatorname{Gal}(L/K)|=[L:K]$.
\end{twierdzenie}
\begin{proof}
Rozdzielczość i Twierdzenie~\ref{aa:tw:embeddings-count}
dają dokładnie $[L:K]$ różnych $K$-zanurzeń
$L\hookrightarrow\Omega$. Normalność oraz
Twierdzenie~\ref{aa:tw:normal-embeddings} pokazują,
że obraz każdego z nich jest równy $L$. Są to zatem
dokładnie wszystkie $K$-automorfizmy $L$.
\end{proof}

\begin{przyklad}[Rozszerzenie rozdzielcze, które nie jest normalne]
\label{aa:prz:separable-not-normal}
Dla rzeczywistego pierwiastka $\alpha=\sqrt[3]{2}$
rozszerzenie $\mathbb Q(\alpha)/\mathbb Q$ ma stopień trzy,
a więc jest skończone. W charakterystyce zero jest
rozdzielcze, ponieważ wszystkie nierozkładalne wielomiany
są rozdzielcze (Wniosek~\ref{aa:wn:charzero-separable}).
Nie jest normalne: nierozkładalny wielomian $X^3-2$
ma pierwiastek $\alpha$ w tym rozszerzeniu, lecz
pozostałe dwa pierwiastki są nierzeczywiste i nie należą
do $\mathbb Q(\alpha)\subseteq\mathbb R$.
Jedynym $\mathbb Q$-automorfizmem tego rozszerzenia
jest identyczność: obraz $\alpha$ musi być rzeczywistym
pierwiastkiem $X^3-2$, czyli $\alpha$.
\end{przyklad}

\begin{przyklad}[Rozszerzenie normalne, które nie jest rozdzielcze]
\label{aa:prz:normal-not-separable}
Niech $K=\mathbb F_p(t)$ i $L=K(\alpha)$, gdzie
$\alpha^p=t$. Najpierw uzasadnijmy $[L:K]=p$.
W pierścieniu $\mathbb F_p[t]$ wielomian $X^p-t$
traktujemy jako wielomian zmiennej $X$ o współczynnikach
w $\mathbb F_p[t]$. Jest on nierozkładalny nad
$\mathbb F_p(t)$ na mocy
Twierdzenia~\ref{aa:tw:eisenstein}
(kryterium Eisensteina), zastosowanego do dziedziny
euklidesowej $\mathbb F_p[t]$ i jej elementu pierwszego
$t$ (ideał $(t)$ jest pierwszy, gdyż iloraz to $\mathbb F_p$): współczynnik
wiodący jest równy $1$, wszystkie pozostałe
współczynniki nieprowadzące są podzielne przez $t$,
a wyraz wolny $-t$ nie jest podzielny przez $t^2$.
Dlatego $[L:K]=p$. Każdy element $x\in L$ ma $p$-tą potęgę w $K$.
Istotnie, dla $x=u(\alpha)/v(\alpha)$, gdzie
$u,v\in K[X]$ i $v(\alpha)\ne0$, endomorfizm
Frobeniusa daje
$x^p=u^{[p]}(t)/v^{[p]}(t)\in K$;
tutaj $u^{[p]}$ oznacza wielomian powstały przez
podniesienie współczynników $u$ do potęgi $p$.
Mianownik jest niezerowy, ponieważ jego wartość
wynosi $v(\alpha)^p$. Jeśli $m_{x,K}$
ma stopień większy od jednego, to dzieli
$X^p-x^p=(X-x)^p$ w $L[X]$. Wszystkie pierwiastki
$m_{x,K}$ są więc równe $x$, a zatem leżą w $L$:
rozszerzenie jest normalne. Nie jest rozdzielcze,
ponieważ wielomian minimalny $\alpha$ to $X^p-t$
i ma wielokrotny pierwiastek.
\end{przyklad}

\section{Ciała skończone jako rozszerzenia Galois}
\begin{twierdzenie}[Istnienie i jednoznaczność ciał skończonych]
\label{aa:tw:finite-fields-existence-uniqueness}
Dla każdej liczby pierwszej $p$ i każdej liczby całkowitej
$n\ge1$ istnieje ciało o $p^n$ elementach.
Każde dwa ciała o $p^n$ elementach są izomorficzne.
Ciało to oznaczamy $\mathbb F_{p^n}$. Jest ono ciałem
rozkładu $X^{p^n}-X$ nad $\mathbb F_p$.
\end{twierdzenie}
\begin{proof}
Niech $q=p^n$ i niech $L$ będzie ciałem rozkładu
$f=X^q-X\in\mathbb F_p[X]$, istniejącym na mocy
Twierdzenia~\ref{aa:tw:splitting-existence}. Pochodna
$f'=qX^{q-1}-1=-1$, więc $f$ ma dokładnie $q$
różnych pierwiastków w $L$. Oznaczmy ich zbiór przez $F$.
Dla $x,y\in F$ mamy $x^q=x$, $y^q=y$, a więc
$(x+y)^q=x^q+y^q=x+y$, ponieważ
podnoszenie do potęgi $p$ jest addytywne i iterujemy je
$n$ razy. Analogicznie $(xy)^q=xy$ oraz
$(-x)^q=-x$. Elementy $0,1$ należą do $F$,
a jeśli $x\in F\setminus\{0\}$, to
$(x^{-1})^q=(x^q)^{-1}=x^{-1}$.
Zbiór $F$ jest zatem podciałem $L$ i ma dokładnie
$q$ elementów. Zawiera $\mathbb F_p$, a jego elementy
są dokładnie wszystkimi pierwiastkami $f$, więc
$F$ jest ciałem rozkładu $f$.

Niech teraz $E$ będzie dowolnym ciałem o $q$ elementach.
Jego charakterystyka wynosi $p$: na mocy
Twierdzenia~\ref{aa:tw:finite-field-cardinality}
liczba elementów $E$ jest potęgą pewnej liczby
pierwszej $\ell=\operatorname{char}E$.
Równość $\ell^r=p^n$ i jednoznaczność rozkładu
na czynniki pierwsze wymuszają $\ell=p$.
Dla $x\in E^\times$ twierdzenie Lagrange'a daje
$x^{q-1}=1$, zatem $x^q=x$; to samo zachodzi dla $x=0$.
Wszystkie $q$ elementów $E$ są więc pierwiastkami $f$,
a $E$ jest ciałem rozkładu $f$ nad swoim kanonicznym
podciałem $\mathbb F_p$. Z jednoznaczności ciała
rozkładu (Twierdzenie~\ref{aa:tw:splitting-uniqueness})
otrzymujemy $E\cong F$.
\end{proof}

\begin{twierdzenie}[Automorfizmy ciała skończonego]
\label{aa:tw:finite-field-frobenius-group}
Niech $p$ będzie liczbą pierwszą, $n\ge1$ liczbą
całkowitą i $q=p^n$. Wtedy $\mathbb F_q/\mathbb F_p$
jest rozszerzeniem Galois stopnia $n$, a jego grupa
Galois jest cykliczna i generowana przez automorfizm
Frobeniusa $\operatorname{Fr}_p:x\mapsto x^p$.
\end{twierdzenie}
\begin{proof}
Ciało $\mathbb F_q$ jest ciałem rozkładu
$X^q-X$, a wielomian ten ma pochodną $-1$,
więc jest rozdzielczy. Rozszerzenie jest zatem
normalne i rozdzielcze; jego stopień wynosi $n$,
ponieważ jako przestrzeń wektorowa nad $\mathbb F_p$
ma $p^n$ elementów. Z
Twierdzenia~\ref{aa:tw:galois-group-degree} wynika,
że jego grupa Galois ma $n$ elementów.
Endomorfizm Frobeniusa jest automorfizmem
(Twierdzenie~\ref{aa:tw:field-frobenius}) i zachowuje
$\mathbb F_p$. Dla każdego $x\in\mathbb F_q$
mamy $\operatorname{Fr}_p^n(x)=x^q=x$, więc rząd
tego automorfizmu dzieli $n$. Jeśli $0<r<n$
i $\operatorname{Fr}_p^r$ był identycznością,
to wszystkie $q$ elementów byłyby pierwiastkami
niezerowego wielomianu $X^{p^r}-X$ stopnia
$p^r<q$, co jest niemożliwe. Rząd wynosi więc
$n$, czyli $\operatorname{Fr}_p$ generuje całą grupę.
\end{proof}

\begin{przyklad}[Ciało o ośmiu elementach]
\label{aa:prz:eight-field-frobenius}
Niech $F=\mathbb F_2[\alpha]$, gdzie
$\alpha^3+\alpha+1=0$. Wielomian $X^3+X+1$
nie ma pierwiastka w $\mathbb F_2$, więc jako
wielomian stopnia trzy jest nierozkładalny.
Mamy $[F:\mathbb F_2]=3$ i $|F|=8$.
Grupa Galois jest cykliczna rzędu trzy; jej
generatorami są $x\mapsto x^2$ oraz $x\mapsto x^4$.
Nie należy mylić tego faktu z cyklicznością
$F^\times$, która jest \emph{inną} grupą i ma rząd siedem.
\end{przyklad}

\begin{uwaga}[Droga do twierdzenia o odpowiedniości Galois]
\label{aa:uw:embeddings-next}
Wykazaliśmy, że w skończonym rozszerzeniu Galois
liczba automorfizmów jest równa stopniowi rozszerzenia.
W kolejnym rozdziale zbadamy, jakie podgrupy tych
automorfizmów odpowiadają ciałom pośrednim. Kluczowy
będzie dowód twierdzenia Artina o ciele elementów stałych,
bez zakładania z góry odpowiedniości Galois.
\end{uwaga}

\begin{uwaga}[Rola normalności w liczeniu automorfizmów]
Dla rozszerzenia $\mathbb Q(\sqrt[3]{2})/\mathbb Q$
stopień wynosi $3$, lecz jedynym automorfizmem nad $\mathbb Q$
jest identyczność: pozostałe dwa pierwiastki $T^3-2$ nie należą
do ciała rzeczywistego $\mathbb Q(\sqrt[3]{2})$.
Zatem równość liczby automorfizmów ze stopniem wymaga
założeń normalności i rozdzielczości, a nie samej skończoności.
\end{uwaga}

% AA-22 | Algebra abstrakcyjna | Iteracja XI
\chapter[Twierdzenie Artina i odpowiedniość Galois]{Twierdzenie Artina i podstawowe twierdzenie teorii Galois}
\label{aa:ch:galois-correspondence}

\section{Od automorfizmów do ciał elementów stałych}
W Rozdziale~\ref{aa:ch:embeddings-normal-separable} zobaczyliśmy,
że skończone rozszerzenie Galois $L/K$ ma dokładnie
$[L:K]$ automorfizmów pozostawiających elementy $K$ bez zmian.
Teraz odwracamy pytanie: dla zbioru automorfizmów chcemy
odtworzyć podciało, które jest przez nie ustalone, oraz
ustalić, czy \emph{każde} ciało pośrednie można tak opisać.

\begin{definicja}[Ciało elementów stałych]
\label{aa:def:fixed-field}
Niech $L$ będzie ciałem, a $H\le\operatorname{Aut}(L)$
podgrupą grupy automorfizmów ciała. Definiujemy
\emph{ciało elementów stałych} podgrupy $H$ przez
\[
L^H=\{x\in L:\sigma(x)=x\text{ dla każdego }\sigma\in H\}.
\]
Jeśli $L/K$ jest rozszerzeniem ciał i
$H\le\operatorname{Gal}(L/K)$, to $K\subseteq L^H\subseteq L$.
\end{definicja}

\begin{lemat}[Podstawowe własności ciała elementów stałych]
\label{aa:lem:fixed-field}
Niech $L$ będzie ciałem i niech $H\le\operatorname{Aut}(L)$. Wtedy $L^H$ jest
podciałem $L$. Jeśli $H_1,H_2\le\operatorname{Aut}(L)$ i
$H_1\subseteq H_2$, to $L^{H_2}\subseteq L^{H_1}$.
Ponadto dla
$\tau\in\operatorname{Aut}(L)$ zachodzi
$\tau(L^H)=L^{\tau H\tau^{-1}}$.
\end{lemat}
\begin{proof}
Każdy automorfizm zachowuje $0$ i $1$, więc $0,1\in L^H$.
Jeśli $x,y\in L^H$, to dla każdego $\sigma\in H$ mamy
$\sigma(x-y)=x-y$, $\sigma(xy)=xy$, a dla $x\ne0$
także $\sigma(x^{-1})=\sigma(x)^{-1}=x^{-1}$.
Dowodzi to, że $L^H$ jest podciałem. Włączenie dla
$H_1\subseteq H_2$ jest bezpośrednie: warunek
niezmienniczości względem większej grupy jest silniejszy.
Wreszcie $x\in L^H$ jest równoważne równości
$(\tau\sigma\tau^{-1})(\tau(x))=\tau(x)$ dla
wszystkich $\sigma\in H$. Ponieważ $\tau$ jest bijekcją,
daje to ostatnią tożsamość.
\end{proof}

\begin{przyklad}[Pierwsza para: ciało i podgrupa]
\label{aa:prz:fixed-quadratic}
W $L=\mathbb Q(\sqrt2)$ niech $H$ będzie grupą
złożoną z identyczności i automorfizmu
$\tau(a+b\sqrt2)=a-b\sqrt2$. Wtedy
$\tau(a+b\sqrt2)=a+b\sqrt2$ wtedy i tylko wtedy,
gdy $2b\sqrt2=0$, czyli $b=0$.
Dlatego $L^H=\mathbb Q$. Natomiast dla podgrupy
trywialnej $\{\operatorname{id}_L\}$ mamy
$L^{\{\operatorname{id}_L\}}=L$.
Widać zatem, że \emph{większa podgrupa odpowiada
mniejszemu ciału elementów stałych}.
\end{przyklad}

\section{Liniowa niezależność różnych zanurzeń}
Poniższy lemat jest zasadniczym narzędziem algebraicznym.
Nie korzysta ani z twierdzenia o odpowiedniości Galois,
ani z wcześniej znanego wzoru na stopień ciała elementów stałych.

\begin{lemat}[Niezależność homomorfizmów ciał]
\label{aa:lem:artin-independence}
Niech $L,M$ będą ciałami, niech $r\ge1$ będzie liczbą
całkowitą i niech
$\sigma_1,\ldots,\sigma_r:L\to M$ będą parami różnymi
homomorfizmami ciał. Jeśli $c_1,\ldots,c_r\in M$ spełniają
\[
\sum_{i=1}^r c_i\sigma_i(x)=0\qquad\text{dla każdego }x\in L,
\]
to wszystkie $c_i$ są zerowe. Innymi słowy, homomorfizmy
te są liniowo niezależne nad $M$ jako funkcje $L\to M$.
\end{lemat}
\begin{proof}
Przypuśćmy, że istnieje nietrywialna relacja, i wybierzmy
taką relację z najmniejszą możliwą liczbą niezerowych
współczynników. Po pominięciu wyrazów zerowych mamy
$\sum_{i=1}^s c_i\sigma_i(x)=0$ dla każdego $x\in L$,
gdzie wszystkie $c_i\ne0$. Nie może być $s=1$, bo
podstawienie $x=1$ dałoby $c_1=0$.
Ponieważ $\sigma_1\ne\sigma_s$, istnieje $a\in L$ takie,
że $\sigma_1(a)\ne\sigma_s(a)$.
W relacji podstawiamy $ax$ w miejsce $x$, a następnie
odejmujemy od niej pierwotną relację pomnożoną przez
$\sigma_s(a)$. Dla każdego $x\in L$ dostajemy
\[
\sum_{i=1}^{s-1}c_i(\sigma_i(a)-\sigma_s(a))\sigma_i(x)=0.
\]
Jest to relacja o mniejszej liczbie wyrazów, lecz
współczynnik przy $\sigma_1$ jest niezerowy.
Sprzeczność z minimalnością kończy dowód.
\end{proof}

\section{Twierdzenie Artina o ciele elementów stałych}
\begin{twierdzenie}[Twierdzenie Artina]
\label{aa:tw:artin-fixed-field}
Niech $L$ będzie ciałem, a $G$ skończoną podgrupą
grupy $\operatorname{Aut}(L)$. Niech $F=L^G$.
Wtedy $L/F$ jest skończonym rozszerzeniem Galois,
\[
[L:F]=|G|,\qquad \operatorname{Gal}(L/F)=G.
\]
Nie zakładamy z góry, że $L$ ma skończony stopień
nad $F$.
\end{twierdzenie}
\begin{proof}
Oznaczmy $m=|G|$. Najpierw pokażemy $[L:F]\le m$.
Wybierzmy dowolne $m+1$ elementów $x_1,\ldots,x_{m+1}\in L$.
Rozważmy układ $m$ jednorodnych równań o niewiadomych
$a_1,\ldots,a_{m+1}\in L$:
\[
\sum_{j=1}^{m+1}a_j\sigma(x_j)=0
\qquad(\sigma\in G).
\]
Ponieważ liczba niewiadomych jest większa od liczby
równań, układ ma niezerowe rozwiązanie.
Wybierzmy spośród takich rozwiązań jedno z najmniejszą
możliwą liczbą niezerowych współczynników.
Dzieląc wszystkie współczynniki przez dowolny wybrany
niezerowy współczynnik, możemy przyjąć $a_{j_0}=1$
dla pewnego $j_0$.

Dla $\tau\in G$ zastosujmy $\tau$ do równania
odpowiadającego $\tau^{-1}\sigma$.
Ponieważ $\tau((\tau^{-1}\sigma)(x_j))=\sigma(x_j)$,
otrzymujemy
$\sum_j\tau(a_j)\sigma(x_j)=0$ dla każdego $\sigma\in G$.
Tak więc rodzina $(\tau(a_j))_j$ także jest rozwiązaniem.
Różnica $(\tau(a_j)-a_j)_j$ jest rozwiązaniem,
a jej współczynnik przy $j_0$ jest równy
$\tau(1)-1=0$. Jeżeli różnica byłaby niezerowa,
miałaby mniej niezerowych współczynników niż wybrane
minimalne rozwiązanie, co jest niemożliwe.
Stąd $\tau(a_j)=a_j$ dla każdego $j$ i $\tau\in G$,
a więc wszystkie $a_j$ należą do $F$.
Dla $\sigma=\operatorname{id}_L$ dostajemy
$\sum_j a_jx_j=0$, nietrywialną zależność liniową
nad $F$: co najmniej jeden współczynnik $a_j$ jest
niezerowy, a wszystkie należą do $F$.
Ponieważ dowolne $m+1$ elementów jest
zależne nad $F$, mamy $[L:F]\le m$.

Oznaczmy teraz $n=[L:F]$, który jest skończony,
i wybierzmy bazę $b_1,\ldots,b_n$ przestrzeni $L$ nad $F$.
Każdy $\sigma\in G$ jest $F$-liniowy, a więc wyznaczony
przez wartości $\sigma(b_1),\ldots,\sigma(b_n)$.
Gdyby $m>n$, to $m$ wektorów
$(\sigma(b_1),\ldots,\sigma(b_n))\in L^n$ byłoby
liniowo zależnych nad $L$. Istniałyby zatem
współczynniki $c_\sigma\in L$, nie wszystkie zerowe,
takie, że $\sum_{\sigma\in G}c_\sigma\sigma(b_j)=0$
dla każdego $j$. Dzięki $F$-liniowości tych automorfizmów
równość przeniosłaby się na wszystkie $x\in L$, wbrew
Lematowi~\ref{aa:lem:artin-independence}.
Zatem $m\le n$, co wraz z pierwszą nierównością daje
$[L:F]=m$.

Dla $x\in L$ rozważmy jego orbitę $Gx=\{\sigma(x):\sigma\in G\}$.
Wielomian
\[
p_x(X)=\prod_{y\in Gx}(X-y)\in L[X]
\]
ma parami różne pierwiastki. Grupa $G$ permutuje orbitę,
więc każdy współczynnik $p_x$ jest ustalany przez
każdy element $G$; dlatego $p_x\in F[X]$.
Wielomian minimalny $m_{x,F}$ dzieli $p_x$ w $F[X]$,
ponieważ $p_x(x)=0$. Wielomian $p_x$ ma w $L$ same
pierwiastki proste, więc również $m_{x,F}$ ma tam same
pierwiastki proste i rozkłada się na czynniki liniowe.
Ponieważ dotyczy to każdego $x\in L$, rozszerzenie
$L/F$ jest rozdzielcze. Jest też normalne na mocy
kryterium rozkładu wszystkich wielomianów minimalnych
z Twierdzenia~\ref{aa:tw:normal-embeddings}.
Jest skończone, zatem Galois.
Wreszcie $G\subseteq\operatorname{Gal}(L/F)$, ponieważ
każdy automorfizm z $G$ ustala wszystkie elementy $F=L^G$.
Z Twierdzenia~\ref{aa:tw:galois-group-degree} wynika,
że grupa po prawej ma $[L:F]=m=|G|$ elementów.
Obie grupy są więc równe.
\end{proof}

\begin{wniosek}[Odzyskiwanie grupy z jej ciała elementów stałych]
\label{aa:wn:artin-recover-group}
Jeśli $L/K$ jest skończonym rozszerzeniem Galois,
a $H\le\operatorname{Gal}(L/K)$, to
$\operatorname{Gal}(L/L^H)=H$ oraz $[L:L^H]=|H|$.
\end{wniosek}
\begin{proof}
Zastosujmy Twierdzenie~\ref{aa:tw:artin-fixed-field}
do skończonej grupy $H$ automorfizmów ciała $L$.
\end{proof}

\section{Podstawowe twierdzenie teorii Galois}
\begin{twierdzenie}[Podstawowe twierdzenie teorii Galois]
\label{aa:tw:fundamental-galois}
Niech $L/K$ będzie \emph{skończonym rozszerzeniem Galois}
i niech $G=\operatorname{Gal}(L/K)$. Między zbiorem
wszystkich ciał pośrednich $E$, $K\subseteq E\subseteq L$,
a zbiorem wszystkich podgrup $H\le G$ istnieją
wzajemnie odwrotne odpowiedniości
\[
E\longmapsto\operatorname{Gal}(L/E),\qquad
H\longmapsto L^H.
\]
Odwracają one inkluzje. Jeśli $E=L^H$, to
\[
[L:E]=|H|,\qquad [E:K]=[G:H].
\]
Ponadto $E/K$ jest normalne (a zatem jest rozszerzeniem
Galois) wtedy i tylko wtedy, gdy $H$ jest podgrupą normalną
w $G$. W tym przypadku ograniczenie automorfizmów daje
kanoniczny izomorfizm
$\operatorname{Gal}(E/K)\cong G/H$.
\end{twierdzenie}
\begin{proof}
\emph{Krok 1: odpowiedniości są odwrotne.}
Ustalmy ciało algebraicznie domknięte $\Omega$ zawierające $L$,
którego istnienie gwarantuje Twierdzenie~\ref{aa:tw:algebraic-closure-exists}.
Niech $E$ będzie ciałem pośrednim. Rozszerzenie $L/E$
jest skończone, ponieważ na mocy twierdzenia o wieży
(Twierdzenie~\ref{aa:tw:field-tower}) zachodzi
$[L:K]=[L:E][E:K]<\infty$.
Ponadto jest rozdzielcze: dla każdego $x\in L$
wielomian minimalny $x$ nad $E$ dzieli wielomian
minimalny nad $K$ w $E[X]$; ten drugi nie ma
pierwiastków wielokrotnych, więc pierwszy także nie.
Rozszerzenie $L/E$ jest również normalne. Z
Twierdzenia~\ref{aa:tw:normal-embeddings} wynika,
że istnieją wielomiany $f_1,\ldots,f_r\in K[X]$,
których ciałem rozkładu nad $K$ jest $L$.
Oznaczmy przez $S$ zbiór wszystkich ich pierwiastków
w $L$. Wtedy $L=K(S)$; ponieważ
$K\subseteq E\subseteq L$, mamy również $L=E(S)$.
Zatem te same wielomiany należą do $E[X]$, rozkładają
się w $L$ i generują $L$ nad $E$.
Stąd $L/E$ jest ciałem rozkładu tej rodziny nad $E$,
a zatem rozszerzeniem normalnym i rozdzielczym,
czyli Galois.

Niech $H=\operatorname{Gal}(L/E)$. Wszystkie elementy
$H$ ustalają $E$, zatem $E\subseteq L^H$.
Z Twierdzenia~\ref{aa:tw:galois-group-degree} i Wniosku
\ref{aa:wn:artin-recover-group} dostajemy
\[
[L:E]=|H|=[L:L^H].
\]
Ponieważ $E\subseteq L^H$, twierdzenie o wieży
wymusza $[L^H:E]=1$, a więc $E=L^H$.
Odwrotnie, dla każdej podgrupy $H\le G$ z Wniosku
\ref{aa:wn:artin-recover-group} wynika bezpośrednio
$\operatorname{Gal}(L/L^H)=H$.

\emph{Krok 2: inkluzje i stopnie.} Jeśli
$E_1\subseteq E_2$, to każdy automorfizm ustalający
$E_2$ ustala też $E_1$, czyli
$\operatorname{Gal}(L/E_2)\subseteq\operatorname{Gal}(L/E_1)$.
Dla podgrup $H_1\subseteq H_2$
Lemat~\ref{aa:lem:fixed-field} daje
$L^{H_2}\subseteq L^{H_1}$, więc także druga
odpowiedniość odwraca inkluzje. Dla $E=L^H$
Twierdzenie Artina daje $[L:E]=|H|$. Z twierdzenia
o wieży i równości $[L:K]=|G|$ wynika
$[E:K]=|G|/|H|=[G:H]$.

\emph{Krok 3: normalność.} Dla $\sigma\in G$ mamy
$\sigma(E)=L^{\sigma H\sigma^{-1}}$ na mocy
Lematu~\ref{aa:lem:fixed-field}.
Dzięki bijektywności odpowiedniości
$\sigma(E)=E$ wtedy i tylko wtedy, gdy
$\sigma H\sigma^{-1}=H$.
Zatem $H\triangleleft G$ dokładnie wtedy, gdy
$E$ jest zachowywane przez wszystkie automorfizmy
$\sigma\in G$; w obu kierunkach korzystamy z
jednoznaczności podgrupy odpowiadającej $E$.

Wykażmy, że to ostatnie jest równoważne normalności
$E/K$. Jeśli $E/K$ jest normalne, to dowolne
$K$-zanurzenie $E\hookrightarrow L$, w szczególności
ograniczenie $\sigma|_E$, ma obraz równy $E$:
wystarczy zastosować argument z dowodu
Twierdzenia~\ref{aa:tw:normal-embeddings} do ciała $E$.
Odwrotnie, załóżmy, że $\sigma(E)=E$ dla wszystkich
$\sigma\in G$. Dowolne $K$-zanurzenie
$\iota:E\hookrightarrow\Omega$ przedłuża się na
$K$-zanurzenie $\widetilde\iota:L\hookrightarrow\Omega$
na mocy Wniosku~\ref{aa:wn:extend-embedding-finite}.
Normalność $L/K$ gwarantuje, że
$\widetilde\iota\in G$. Zatem
$\iota(E)=\widetilde\iota(E)=E$.
Kryterium normalności przez zanurzenia daje
normalność $E/K$.

\emph{Krok 4: grupa ilorazowa.} Jeśli $H\triangleleft G$,
to $E/K$ jest normalne. Rozdzielczość przechodzi
na podrozszerzenia, więc $E/K$ jest Galois.
Ograniczenie
$\rho:G\to\operatorname{Gal}(E/K)$,
$\rho(\sigma)=\sigma|_E$, jest dobrze określonym
homomorfizmem, bo $\sigma(E)=E$. Jego jądrem są
dokładnie automorfizmy ustalające $E$, czyli $H$.
Jest on surjektywny: każdy $K$-automorfizm $\tau:E\to E$
przedłuża się na $K$-zanurzenie $L\to\Omega$,
które dzięki normalności $L/K$ jest elementem $G$.
Twierdzenie o izomorfizmie grup daje
$G/H\cong\operatorname{Gal}(E/K)$.
\end{proof}

\section{Przykłady odpowiedniości Galois}
\begin{przyklad}[Rozszerzenie dwukwadratowe]
\label{aa:prz:biquadratic-galois-lattice}
Niech $L=\mathbb Q(\sqrt2,\sqrt3)$. W
Przykładzie~\ref{aa:prz:biquadratic-extension} wykazaliśmy
$[L:\mathbb Q]=4$. Ciało $L$ jest ciałem rozkładu
$(X^2-2)(X^2-3)$, a rozszerzenie jest rozdzielcze,
ponieważ $\operatorname{char}\mathbb Q=0$.
Zatem $L/\mathbb Q$ jest rozszerzeniem Galois.

Niech $s$ zmienia znak $\sqrt2$ i zachowuje $\sqrt3$,
a $t$ zachowuje $\sqrt2$ i zmienia znak $\sqrt3$.
Istnienie tych automorfizmów wymaga uzasadnienia, a nie
samego wskazania możliwych obrazów generatorów.
Najpierw $\mathbb Q(\sqrt2)$ ma dwa zanurzenia nad
$\mathbb Q$, wysyłające $\sqrt2$ odpowiednio na
$\sqrt2$ i $-\sqrt2$ (Przykład~\ref{aa:prz:quadratic-conjugation}).
Wielomian minimalny $\sqrt3$ nad $\mathbb Q(\sqrt2)$
ma stopień dwa: gdyby $\sqrt3$ należał do
$\mathbb Q(\sqrt2)$, stopień rozszerzenia
$L/\mathbb Q$ byłby równy dwa, a wynosi cztery.
Jest to zatem $X^2-3$. Po przeniesieniu współczynników
przez każde z dwóch zanurzeń pozostaje on tym samym
wielomianem i ma dwa pierwiastki $\sqrt3,-\sqrt3$
w $L$. Lemat~\ref{aa:lem:extend-embedding-simple}
przedłuża każde zanurzenie na dwa sposoby. Powstałe
cztery zanurzenia $L\to L$ są automorfizmami:
każde jest $\mathbb Q$-liniowe i różnowartościowe,
a $L$ ma skończony wymiar nad $\mathbb Q$, więc
jest również surjektywne. W szczególności istnieją $s$ i $t$
o podanych wartościach na generatorach.
Mamy $s^2=t^2=1$, $st=ts$, a cztery automorfizmy to
$1,s,t,st$. Stąd
$G\cong C_2\times C_2$. Podgrupy rzędu dwa to
$\langle s\rangle$, $\langle t\rangle$ i
$\langle st\rangle$; ich ciała elementów stałych
mają stopień dwa nad $\mathbb Q$ i wynoszą odpowiednio
\[
L^{\langle s\rangle}=\mathbb Q(\sqrt3),\quad
L^{\langle t\rangle}=\mathbb Q(\sqrt2),\quad
L^{\langle st\rangle}=\mathbb Q(\sqrt6).
\]
Ostatnia równość wynika z tego, że $st$ zmienia
oba znaki, a więc zachowuje $\sqrt6$, oraz ze
zgodności stopni obu ciał nad $\mathbb Q$.
Wszystkie podgrupy $G$ są normalne, dlatego każde
ciało pośrednie jest normalne nad $\mathbb Q$.
\end{przyklad}

\begin{przyklad}[Grupa symetryczna i nienormalne ciało pośrednie]
\label{aa:prz:cubic-galois-s3}
Niech $\alpha=\sqrt[3]2\in\mathbb R$,
$\omega\in\mathbb C$ będzie pierwiastkiem trzeciego
stopnia z jedynki różnym od $1$, a
$L=\mathbb Q(\alpha,\omega)$. Z
Przykładu~\ref{aa:prz:cubic-splitting} wiemy,
że $L$ jest ciałem rozkładu $X^3-2$ i
$[L:\mathbb Q]=6$. Wielomian jest rozdzielczy
w charakterystyce zero, zatem grupa $G=\operatorname{Gal}(L/\mathbb Q)$
ma sześć elementów.

Definiujemy dwa automorfizmy na generatorach:
\[
\sigma(\alpha)=\omega\alpha,\quad\sigma(\omega)=\omega,
\qquad
\tau(\alpha)=\alpha,\quad\tau(\omega)=\omega^2.
\]
Aby uzasadnić istnienie $\sigma$, zauważmy, że
$[\mathbb Q(\omega):\mathbb Q]=2$, a
$[L:\mathbb Q]=6$, zatem $[L:\mathbb Q(\omega)]=3$.
Wielomian $X^3-2$ jest więc wielomianem minimalnym
$\alpha$ nad $\mathbb Q(\omega)$ i ma w $L$ również
pierwiastek $\omega\alpha$.
Lemat~\ref{aa:lem:extend-embedding-simple} daje zatem
$\mathbb Q(\omega)$-zanurzenie $L\to L$ o wskazanych
wartościach. Jego obraz zawiera $\omega$ oraz
$\omega\alpha$, więc także
$\alpha=(\omega\alpha)/\omega$, a zatem całe $L$.
Jest to automorfizm $\sigma$. Sprzężenie zespolone
zachowuje $L$, ponieważ $\alpha\in\mathbb R$ oraz
$\overline{\omega}=\omega^2\in L$. Jego ograniczenie
do $L$ jest drugim $\mathbb Q$-automorfizmem $\tau$. Mamy
$\sigma^3=\tau^2=1$ oraz
$\tau\sigma\tau^{-1}=\sigma^{-1}$.
Elementy $1,\sigma,\sigma^2,\tau,\sigma\tau,\sigma^2\tau$
są parami różne: trzy pierwsze ustalają $\omega$,
a trzy ostatnie wysyłają go na $\omega^2$;
w każdej trójce obrazy $\alpha$ są różne.
Ponieważ $|G|=6$, elementy te tworzą całe $G$.
Działanie $G$ na trzech pierwiastkach
$\alpha,\omega\alpha,\omega^2\alpha$ daje homomorfizm
do $S_3$. Jest on różnowartościowy, bo automorfizm
ustalający wszystkie trzy pierwiastki ustala zarówno
$\alpha$, jak i $\omega=(\omega\alpha)/\alpha$,
a więc całe $L$. Obie grupy mają rząd sześć,
zatem $G\cong S_3$.

Podgrupa $\langle\sigma\rangle$ ma rząd trzy,
więc jej ciało elementów stałych ma stopień dwa
nad $\mathbb Q$. Ponieważ $\sigma(\omega)=\omega$,
mamy $L^{\langle\sigma\rangle}=\mathbb Q(\omega)$.
Podgrupa $\langle\tau\rangle$ ma rząd dwa,
a jej ciało elementów stałych to $\mathbb Q(\alpha)$:
$\tau$ zachowuje $\alpha$, a oba ciała mają stopień
trzy nad $\mathbb Q$. Podgrupa $\langle\tau\rangle$
nie jest normalna w $S_3$, więc
$\mathbb Q(\alpha)/\mathbb Q$ nie jest normalne.
To pokazuje na konkretnym przykładzie, dlaczego
w podstawowym twierdzeniu teorii Galois nie można
żądać normalności każdego ciała pośredniego.
\end{przyklad}

\begin{przyklad}[Podgrupy i ciała pośrednie ciała skończonego]
\label{aa:prz:finite-fields-galois-lattice}
Niech $p$ będzie liczbą pierwszą,
$L=\mathbb F_{p^{12}}$ i $K=\mathbb F_p$
(przy czym $K$ utożsamiamy z pierwszym podciałem $L$).
Grupa Galois rozszerzenia $L/K$ jest cykliczna rzędu $12$, generowana
przez $\operatorname{Fr}_p:x\mapsto x^p$.
Dla każdego dzielnika $d\mid12$ istnieje dokładnie
jedna podgrupa indeksu $d$, generowana przez
$\operatorname{Fr}_p^{d}$. Jej ciało elementów stałych
ma stopień $d$ nad $K$ i dlatego $p^d$ elementów.
Zatem dla każdego $d\in\{1,2,3,4,6,12\}$
odpowiedniość Galois daje dokładnie jedno takie
ciało pośrednie. Z twierdzenia o jednoznaczności ciał
skończonych wynika, że jest ono izomorficzne z
$\mathbb F_{p^d}$.
Wszystkie te rozszerzenia są normalne, ponieważ
każda podgrupa grupy cyklicznej jest normalna.
\end{przyklad}

\begin{uwaga}[Granice obecnego etapu]
\label{aa:uw:galois-future}
Podstawowe twierdzenie teorii Galois klasyfikuje ciała
pośrednie skończonych rozszerzeń Galois poprzez podgrupy.
Nie rozstrzyga jeszcze, kiedy dane równanie wielomianowe
jest rozwiązywalne przez pierwiastniki. W następnej części
trzeba zestawić tę odpowiedniość z szeregiem podnormalnym
i rozwiązalnością grup z Rozdziału~\ref{aa:ch:szeregi}.
Osobnego opracowania wymagają również norma, ślad,
rozszerzenia cykliczne i odpowiednie zastosowania teorii Galois.
\end{uwaga}

\begin{uwaga}[Dlaczego odpowiedniość wymaga rozszerzenia Galois]
Niech $L=\mathbb Q(\sqrt[3]{2})$ i $K=\mathbb Q$.
Istotnie, każdy $K$-automorfizm $L$ musi wysłać
$\sqrt[3]{2}$ na pierwiastek jego wielomianu minimalnego
$X^3-2$. Ponieważ $L\subseteq\mathbb R$, jedynym takim
pierwiastkiem w $L$ jest $\sqrt[3]{2}$, więc automorfizm
jest identycznością. Grupa $\operatorname{Aut}_K(L)$
jest zatem trywialna, a jej ciało elementów stałych
to całe $L$, nie zaś $K$.
Nie wolno zatem przenosić równości
$L^{\operatorname{Gal}(L/K)}=K$ na dowolne rozszerzenia
skończone bez założenia normalności.
\end{uwaga}

% AA-23 | Algebra abstrakcyjna | Iteracja XII
\chapter{Ślad, norma i wyróżnik w rozszerzeniach ciał}
\label{aa:ch:trace-norm-discriminant}
\label{aa:ch:trace-norm}

\section{Od mnożenia w ciele do operatora liniowego}
W Rozdziale~\ref{aa:ch:galois-correspondence} opisaliśmy grupy automorfizmów
rozszerzeń Galois. Innym sposobem wydobywania informacji z rozszerzenia
jest potraktowanie mnożenia przez ustalony element jako przekształcenia
liniowego. Otrzymamy w ten sposób \emph{ślad} i \emph{normę}.
Konstrukcja ma sens również wtedy, gdy rozszerzenie nie jest rozdzielcze.

\begin{definicja}[Ślad i norma elementu rozszerzenia]
\label{aa:def:field-trace-norm}
Niech $L/K$ będzie rozszerzeniem ciał skończonego stopnia $n=[L:K]$.
Dla $a\in L$ niech $m_a:L\to L$ oznacza $K$-liniowy endomorfizm
$m_a(x)=ax$. Definiujemy
\[
\operatorname{Tr}_{L/K}(a)=\operatorname{tr}_K(m_a),\qquad
N_{L/K}(a)=\det_K(m_a).
\]
Wyznacznik i ślad po prawej stronie obliczamy w dowolnej bazie
$K$-przestrzeni $L$. Oba należą do $K$.
Ponadto wielomianem charakterystycznym elementu $a$
względem $L/K$ nazywamy $\chi_{a,L/K}(T)=\det(TI-m_a)\in K[T]$.
\end{definicja}

\begin{uwaga}[Zależność od rozszerzenia]\label{aa:uw:trace-relative-extension}
Indeksy $L/K$ są częścią definicji. Jeżeli $a\in F$ dla ciała pośredniego
$K\subseteq F\subseteq L$, to operatory mnożenia przez $a$
na $F$ i na $L$ działają na przestrzeniach o różnych wymiarach;
wobec tego ich ślady i wyznaczniki nie muszą być równe.
Wielomian $\chi_{a,L/K}$ ma zawsze stopień $[L:K]$,
niezależnie od stopnia minimalnego wielomianu $a$ nad $K$.
Przypadek $L=K$ jest również objęty definicją:
wtedy $m_a$ ma macierz $(a)$, a więc
$\operatorname{Tr}_{K/K}(a)=N_{K/K}(a)=a$.
\end{uwaga}

\begin{lemat}[Niezależność od wyboru bazy]
\label{aa:lem:trace-norm-basis}
Niech $L/K$ będzie skończonym rozszerzeniem ciał, a $a\in L$.
Ślad $\operatorname{Tr}_{L/K}(a)$, norma $N_{L/K}(a)$
oraz wielomian charakterystyczny $\chi_{a,L/K}$
nie zależą od wyboru bazy $K$-przestrzeni $L$.
\end{lemat}
\begin{proof}
Jeżeli $A$ i $B$ są macierzami tego samego endomorfizmu $m_a$
w dwóch bazach, to istnieje odwracalna macierz $S$ taka, że
$B=S^{-1}AS$. Stąd $\operatorname{tr}B=\operatorname{tr}A$,
gdyż $\operatorname{tr}(UV)=\sum_{i,j}u_{ij}v_{ji}=\operatorname{tr}(VU)$,
oraz $\det B=(\det S)^{-1}(\det A)\det S=\det A$.
Ponadto $TI-B=S^{-1}(TI-A)S$, a więc
$\det(TI-B)=\det(TI-A)$ w $K[T]$.
Zatem ślad, norma i wielomian charakterystyczny
nie zależą od wyboru bazy.
\end{proof}

\begin{twierdzenie}[Podstawowe własności śladu i normy]
\label{aa:tw:trace-norm-basic}
\label{aa:thm:trace-norm-properties}
Niech $L/K$ będzie skończonym rozszerzeniem stopnia $n$.
Dla wszystkich $a,b\in L$ i $c\in K$ zachodzą
\begin{enumerate}[label=\textup{(\roman*)}]
\item $\operatorname{Tr}_{L/K}(a+b)=\operatorname{Tr}_{L/K}(a)+\operatorname{Tr}_{L/K}(b)$;
\item $\operatorname{Tr}_{L/K}(ca)=c\operatorname{Tr}_{L/K}(a)$;
\item $N_{L/K}(ab)=N_{L/K}(a)N_{L/K}(b)$;
\item $\operatorname{Tr}_{L/K}(c)=nc$ oraz $N_{L/K}(c)=c^n$;
\item $N_{L/K}(a)\ne0$ wtedy i tylko wtedy, gdy $a\ne0$;
      w szczególności $N_{L/K}(a^{-1})=N_{L/K}(a)^{-1}$ dla $a\ne0$.
\end{enumerate}
\end{twierdzenie}
\begin{proof}
Mamy $m_{a+b}=m_a+m_b$ oraz $m_{ca}=c\,m_a$.
Liniowość śladu daje (i) i (ii). Z łączności mnożenia
w $L$ wynika $m_{ab}=m_a\circ m_b$, więc multiplikatywność
wyznacznika daje (iii). Dla $c\in K$ endomorfizm $m_c$
ma w każdej bazie macierz $cI_n$, skąd (iv).
Jeśli $a\ne0$, to mnożenie przez $a$ ma odwrotność $m_{a^{-1}}$,
zatem jego wyznacznik jest niezerowy. Jeśli $a=0$, to $m_0=0$,
więc jego wyznacznik jest zerem. Ostatnia równość wynika z (iii)
i $N_{L/K}(1)=1$.
\end{proof}

\begin{uwaga}[Ślad nie jest homomorfizmem pierścieni]
\label{aa:uw:trace-not-ring}
Ślad jest $K$-liniowy, lecz zazwyczaj nie jest multiplikatywny.
Norma jest multiplikatywna, lecz zazwyczaj nie jest addytywna.
Na przykład dla $L=\mathbb Q(\sqrt2)$ mamy
$\operatorname{Tr}_{L/\mathbb Q}(1)=2$ i
$\operatorname{Tr}_{L/\mathbb Q}(1)^2=4$,
podczas gdy $\operatorname{Tr}_{L/\mathbb Q}(1\cdot1)=2$.
Ponadto $N_{L/\mathbb Q}(1+1)=4$, ale
$N_{L/\mathbb Q}(1)+N_{L/\mathbb Q}(1)=2$.
\end{uwaga}

\section{Macierz mnożenia i wielomian minimalny}
\begin{twierdzenie}[Ślad i norma pierwiastka wielomianu minimalnego]
\label{aa:tw:trace-norm-companion}
Niech $L=K(\alpha)$, gdzie $\alpha$ jest algebraiczny nad $K$,
a jego unormowany wielomian minimalny ma postać
$m_{\alpha,K}(X)=X^n+a_{n-1}X^{n-1}+\cdots+a_1X+a_0$.
Wtedy
\[
\operatorname{Tr}_{L/K}(\alpha)=-a_{n-1},\qquad
N_{L/K}(\alpha)=(-1)^n a_0.
\]
\end{twierdzenie}
\begin{proof}
Na mocy Twierdzenia~\ref{aa:tw:simple-algebraic-extension}
układ $(1,\alpha,\ldots,\alpha^{n-1})$ jest bazą $L$ nad $K$.
Mnożenie przez $\alpha$ przesuwa każdy z pierwszych $n-1$
elementów bazy na następny, a z równości $m_{\alpha,K}(\alpha)=0$
otrzymujemy
$\alpha^n=-a_0-a_1\alpha-\cdots-a_{n-1}\alpha^{n-1}$.
Dla $n=1$ mamy $L=K$, $\alpha=-a_0$ i macierz
$m_\alpha$ jest równa $(-a_0)$; obie formuły
wynikają zatem bezpośrednio z definicji.
Dla $n\ge2$ macierz $m_\alpha$ ma postać blokową
\[
C=\begin{pmatrix}
0_{1\times(n-1)}&-a_0\\
I_{n-1}&-\begin{pmatrix}a_1\\\vdots\\a_{n-1}\end{pmatrix}
\end{pmatrix}.
\]
To jest macierz towarzysząca, poprawnie określona także
dla $n=2$. Jej ślad wynosi $-a_{n-1}$.
Rozwijając wyznacznik względem pierwszego wiersza,
możemy wybrać tylko wpis $-a_0$; jego dopełnienie
jest macierzą jednostkową stopnia $n-1$.
Zatem $\det C=(-a_0)(-1)^{1+n}=(-1)^n a_0$.
\end{proof}

\begin{twierdzenie}[Wielomian charakterystyczny elementu]
\label{aa:thm:trace-norm-minpoly}
Niech $L/K$ będzie skończonym rozszerzeniem ciał,
$a\in L$, $d=[K(a):K]$, $r=[L:K(a)]$, a
$m_{a,K}(T)=T^d+c_{d-1}T^{d-1}+\cdots+c_0$
niech oznacza unormowany wielomian minimalny $a$ nad $K$.
Wtedy wielomian charakterystyczny elementu $a$ względem $L/K$
ma postać $\chi_{a,L/K}(T)=m_{a,K}(T)^r$, a więc
\[
\operatorname{Tr}_{L/K}(a)=-r c_{d-1},\qquad
N_{L/K}(a)=\bigl((-1)^d c_0\bigr)^r.
\]
Nie zakładamy rozdzielczości rozszerzenia.
\end{twierdzenie}
\begin{proof}
Wybierzmy bazę $u_1,\ldots,u_r$ przestrzeni $L$ nad $K(a)$.
Każdy element $L$ można jednoznacznie zapisać jako
kombinację liniową nad $K$ elementów
$a^j u_i$, $0\le j<d$, $1\le i\le r$:
istnienie wynika z rozwinięcia współczynników
w bazie $(1,a,\ldots,a^{d-1})$ ciała $K(a)$ nad $K$,
a niezależność wynika z niezależności $u_i$ nad $K(a)$
i potęg $a^j$ nad $K$. Jest to więc baza $L$ nad $K$.
Mnożenie przez $a$ pozostawia każdą podprzestrzeń
$K$ rozpiętą przez $u_i,au_i,\ldots,a^{d-1}u_i$.
Na tej podprzestrzeni jego macierz jest macierzą
towarzyszącą $m_{a,K}$, gdyż $m_{a,K}(a)=0$.
Wielomian charakterystyczny takiego bloku to $m_{a,K}$.
Aby uzasadnić ten wzór, zauważmy, że w macierzy $TI-C$
na przekątnej stoją wpisy $T$, tuż pod nią wpisy $-1$,
a w ostatniej kolumnie wpisy $c_0,\ldots,c_{d-1}$
(z dodanym $T$ na ostatniej pozycji).
W rozwinięciu wyznacznika względem ostatniej kolumny
składnik pochodzący od $c_{i-1}$ wynosi
$c_{i-1}T^{i-1}$: dopełnienie składa się z $i-1$
czynników $T$ i $d-i$ czynników $-1$, których znaki
znoszą znak dopełnienia algebraicznego.
Wpis $T$ z ostatniego wiersza daje dodatkowo $T^d$.
Dla $d=1$ formuła jest bezpośrednia.
Otrzymujemy zatem $\det(TI-C)=T^d+c_{d-1}T^{d-1}+\cdots+c_0$.
Macierz całego $m_a$ jest blokowo diagonalna
z $r$ takimi blokami. Zatem $\chi_{a,L/K}=m_{a,K}^r$.
Ponieważ dla każdej macierzy $n\times n$ współczynnik
przy $T^{n-1}$ w $\det(TI-A)$ jest równy
$-\operatorname{tr}A$, a wyraz stały to
$(-1)^n\det A$, porównanie współczynników
po obu stronach daje obie formuły.
\end{proof}

\begin{przyklad}[Rozszerzenie kwadratowe]
\label{aa:prz:trace-norm-quadratic}
\label{aa:ex:quadratic-norm}
Niech $d\in\mathbb Q$ będzie liczbą, która nie jest kwadratem
liczby wymiernej, i $L=\mathbb Q(\sqrt d)$.
Dla $a,b\in\mathbb Q$ mnożenie przez $z=a+b\sqrt d$
ma w bazie $(1,\sqrt d)$ macierz
\[
\begin{pmatrix}a&bd\\ b&a\end{pmatrix}.
\]
Otrzymujemy dokładnie
$\operatorname{Tr}_{L/\mathbb Q}(a+b\sqrt d)=2a$,
$N_{L/\mathbb Q}(a+b\sqrt d)=a^2-db^2$.
Zatem element $a+b\sqrt d$ jest niezerowy wtedy i tylko wtedy,
gdy $a^2-db^2\ne0$, a jego odwrotność ma postać
$(a-b\sqrt d)/(a^2-db^2)$. Zauważmy, że obliczenie to
nie korzysta z istnienia automorfizmów rozszerzenia.
\end{przyklad}

\begin{przyklad}[Zmiana ciała, w którym obliczamy normę]
\label{aa:prz:trace-norm-relative-biquadratic}
Niech $K=\mathbb Q$, $F=\mathbb Q(\sqrt2)$ oraz
$L=\mathbb Q(\sqrt2,\sqrt3)$. Ponieważ $[F:K]=2$, sprawdźmy, że $\sqrt3\notin F$.
Gdyby $\sqrt3=u+v\sqrt2$ dla $u,v\in\mathbb Q$, to po
podniesieniu do kwadratu dostalibyśmy
$3=u^2+2v^2+2uv\sqrt2$. Wobec niewymierności
$\sqrt2$ musi być $uv=0$. Dla $v=0$ mielibyśmy
$u^2=3$, a dla $u=0$ mielibyśmy $v^2=3/2$;
żadna z tych liczb nie jest kwadratem liczby wymiernej
(co wynika z rozkładów liczników i mianowników na
czynniki pierwsze). Zatem $\sqrt3\notin F$,
$[L:F]=2$ i $[L:K]=4$.
Dla $a=\sqrt2$ wielomian minimalny nad $K$ to $T^2-2$.
Zatem w rozszerzeniu $F/K$ wzory z
Twierdzenia~\ref{aa:tw:trace-norm-companion} dają
$\operatorname{Tr}_{F/K}(a)=0$ i $N_{F/K}(a)=-2$.
Na $L$ operator mnożenia przez $a$ ma w bazie
$(1,\sqrt2,\sqrt3,\sqrt6)$ dwa identyczne bloki
$\left(\begin{smallmatrix}0&2\\1&0\end{smallmatrix}\right)$:
pierwszy działa na $\operatorname{span}_K(1,\sqrt2)$,
a drugi na $\operatorname{span}_K(\sqrt3,\sqrt6)$.
W konsekwencji $\operatorname{Tr}_{L/K}(a)=0+0=0$,
lecz $N_{L/K}(a)=(-2)^2=4$.
Obliczenie pokazuje, że nie wolno usuwać indeksu $L/K$
ani utożsamiać normy elementu w różnych rozszerzeniach.
\end{przyklad}

\begin{przyklad}[Rozszerzenie trzeciego stopnia]
\label{aa:prz:trace-norm-cubic}
Dla $L=\mathbb Q(\alpha)$, $\alpha=\sqrt[3]2$,
minimalny wielomian $\alpha$ wynosi $X^3-2$:
wielomian ten jest nierozkładalny nad $\mathbb Q$
na mocy kryterium Eisensteina z liczbą pierwszą $2$
(Twierdzenie~\ref{aa:tw:eisenstein}).
Dlatego $\operatorname{Tr}_{L/\mathbb Q}(\alpha)=0$
i $N_{L/\mathbb Q}(\alpha)=2$. Równość dotycząca normy
można również otrzymać z
$N(\alpha)^3=N(\alpha^3)=N(2)=2^3$:
ponieważ $N(\alpha)\in\mathbb Q$, z równości
$N(\alpha)^3=8$ wynika $N(\alpha)=2$.
Bezpośredni wzór z wielomianu minimalnego nie wymaga
rozwiązywania takiego równania.
\end{przyklad}

\begin{przyklad}[Zanik śladu w rozszerzeniu nierozdzielczym]
\label{aa:ex:inseparable-trace}
Niech $p$ będzie liczbą pierwszą,
$K=\mathbb F_p(t)$ i $L=K(s)$, gdzie $s^p=t$.
Wielomian $X^p-t\in\mathbb F_p[t][X]$ jest
nierozkładalny nad $K$ na mocy kryterium Eisensteina
względem elementu pierwszego $t$ w pierścieniu
$\mathbb F_p[t]$ i lematu Gaussa.
Zatem $[L:K]=p$.
Dla dowolnego $a\in L$ mamy $a^p\in K$,
ponieważ $K=\mathbb F_p(t)\subseteq\mathbb F_p(s)$,
$s^p=t$ i stąd $L=\mathbb F_p(s)$;
każdy ułamek wymierny w $s$, podniesiony do potęgi $p$,
jest ułamkiem wymiernym w $s^p=t$ o współczynnikach
z $\mathbb F_p$.
Jeśli $a\in K$, to z podstawowych własności śladu
$\operatorname{Tr}_{L/K}(a)=pa=0$.
Jeśli $a\notin K$, to $[K(a):K]$ jest dzielnikiem
liczby pierwszej $[L:K]=p$, więc jest równy $p$.
Wielomian $X^p-a^p\in K[X]$ ma stopień $p$
i zeruje $a$, zatem jest wielomianem minimalnym $a$.
Nie ma współczynnika przy $X^{p-1}$;
z Twierdzenia~\ref{aa:thm:trace-norm-minpoly}
wynika $\operatorname{Tr}_{L/K}(a)=0$.
Zatem cały funkcjonał śladu $\operatorname{Tr}_{L/K}$
jest zerowy, choć $L/K$ ma dodatni stopień.
Przykład pokazuje, dlaczego w twierdzeniach o
niezdegenerowaniu formy śladu wymagamy rozdzielczości.
\end{przyklad}

\begin{uwaga}[Norma w tym samym nierozdzielczym rozszerzeniu]
\label{aa:uw:pure-inseparable-norm}
W poprzednim przykładzie norma nie znika. Dokładniej,
dla każdego $a\in L$ zachodzi $N_{L/K}(a)=a^p\in K$.
Jeśli $a\in K$, jest to równość $N(a)=a^{[L:K]}=a^p$.
Jeśli $a\notin K$, wielomian minimalny ma postać $X^p-a^p$.
Wzór z Twierdzenia~\ref{aa:thm:trace-norm-minpoly} daje
$N(a)=(-1)^p(-a^p)=a^p$, ponieważ
$(-1)^{p+1}=1$ w ciele charakterystyki $p$
(dla $p=2$ znaki $1$ i $-1$ są identyczne).
W szczególności $N(s)=s^p=t\ne0$, choć
$\operatorname{Tr}_{L/K}(s)=0$.
\end{uwaga}

\section{Element pierwotny rozszerzenia rozdzielczego}
Przejście od wyznacznika do iloczynu wartości zanurzeń
wymaga rozdzielczości. Wyprowadzimy najpierw twierdzenie,
które pozwala opisać skończone rozszerzenie rozdzielcze
przez jeden generator.

\begin{twierdzenie}[Twierdzenie o elemencie pierwotnym]
\label{aa:tw:primitive-element}
Jeżeli $L/K$ jest skończonym rozszerzeniem rozdzielczym,
to istnieje $\gamma\in L$ takie, że $L=K(\gamma)$.
\end{twierdzenie}
\begin{proof}
Najpierw rozważmy przypadek nieskończonego ciała $K$.
Wystarczy udowodnić, że dla dowolnych $\alpha,\beta\in L$
istnieje $c\in K$ z $K(\alpha,\beta)=K(\alpha+c\beta)$;
powtarzanie tego kroku dla skończonej liczby generatorów
daje tezę.
Niech $\Omega$ będzie domknięciem algebraicznym $K$
zawierającym $L$; możemy je wybrać na mocy
Twierdzenia~\ref{aa:tw:algebraic-closure-exists}
i przedłużania zanurzeń. Niech
$\sigma_1,\ldots,\sigma_n:L\hookrightarrow\Omega$
będą wszystkimi $K$-zanurzeniami,
gdzie $n=[L:K]$ na mocy
Twierdzenia~\ref{aa:tw:embeddings-count}. Jeżeli
$(\sigma_i(\alpha),\sigma_i(\beta))\ne
(\sigma_j(\alpha),\sigma_j(\beta))$, to równość
$\sigma_i(\alpha+c\beta)=\sigma_j(\alpha+c\beta)$
może zachodzić dla co najwyżej jednej wartości $c\in K$:
różnica wartości jest liniowa w $c$, a jej współczynniki
nie są jednocześnie zerowe. Ponieważ par indeksów jest
skończenie wiele, wybieramy $c\in K$ omijające wszystkie
wartości powodujące równość dla par zanurzeń różniących się
na $K(\alpha,\beta)$. Wtedy liczba różnych obrazów
$\gamma=\alpha+c\beta$ przez $K$-zanurzenia $L$ jest równa
liczbie różnych ograniczeń tych zanurzeń do
$K(\alpha,\beta)$. Z twierdzenia o liczbie zanurzeń
i przedłużaniu zanurzeń ta liczba wynosi
$[K(\alpha,\beta):K]$. Tymczasem liczba różnych obrazów
$\gamma$ to $[K(\gamma):K]$, ponieważ $\gamma$ jest
rozdzielczy nad $K$. Mamy zatem
$[K(\gamma):K]=[K(\alpha,\beta):K]$.
Ponieważ $K(\gamma)\subseteq K(\alpha,\beta)$,
twierdzenie o wieży daje równość tych ciał.

Równość liczby obrazów z wymiarem $[K(\gamma):K]$
wynika z rozdzielczości: minimalny wielomian $\gamma$
ma dokładnie tyle różnych pierwiastków, ile wynosi jego stopień,
a każdy obraz $\gamma$ pod $K$-zanurzeniem jest jednym z nich.
Z drugiej strony każde $K$-zanurzenie $K(\gamma)$
przedłuża się do $K$-zanurzenia $L$, zatem żaden z tych
pierwiastków nie został pominięty. Analogiczne uzasadnienie
dotyczy ograniczeń do $K(\alpha,\beta)$.
Ponieważ $L/K$ jest skończone, wybór skończonego zbioru
generatorów $L$ nad $K$ nie wymaga dodatkowego twierdzenia:
wystarczy wziąć dowolną bazę $L$ jako przestrzeni wektorowej
i dołączać jej elementy po kolei.

Jeżeli $K$ jest skończone, to również $L$ jest ciałem
skończonym. Na mocy
Twierdzenia~\ref{aa:tw:finite-field-multiplicative-cyclic}
grupa $L^\times$ jest cykliczna. Niech $\gamma$ będzie
jej generatorem. Każdy niezerowy element $L$ jest potęgą
$\gamma$, więc $L=K(\gamma)$.
\end{proof}

\section{Ślad i norma jako sumy oraz iloczyny sprzężeń}
\begin{twierdzenie}[Wzory przez zanurzenia]
\label{aa:tw:trace-norm-embeddings}
Niech $L/K$ będzie skończonym rozszerzeniem rozdzielczym,
$n=[L:K]$, a $\Omega$ domknięciem algebraicznym $K$
zawierającym $L$. Niech $\sigma_1,\ldots,\sigma_n$
będą wszystkimi $K$-zanurzeniami $L\to\Omega$.
Wtedy dla każdego $a\in L$ zachodzą
\[
\operatorname{Tr}_{L/K}(a)=\sum_{i=1}^n\sigma_i(a),\qquad
N_{L/K}(a)=\prod_{i=1}^n\sigma_i(a).
\]
\end{twierdzenie}
\begin{proof}
Z Twierdzenia~\ref{aa:tw:primitive-element} wybieramy
$\gamma\in L$ z $L=K(\gamma)$. Niech $\gamma_i=\sigma_i(\gamma)$;
są one parami różne, ponieważ minimalny wielomian $\gamma$
jest rozdzielczy. W bazie $1,\gamma,\ldots,\gamma^{n-1}$
endomorfizm $m_a$ ma macierz $A=(a_{ij})$ o współczynnikach
w $K$. Dla każdego $j$ mamy
$a\gamma^{j-1}=\sum_i a_{ij}\gamma^{i-1}$.
Po zastosowaniu $\sigma_k$ dostajemy
$\sigma_k(a)\gamma_k^{j-1}=\sum_i a_{ij}\gamma_k^{i-1}$.
Jeśli $V=(\gamma_k^{j-1})_{k,j}$ jest macierzą Vandermonde'a,
to te równości dają
$VA=\operatorname{diag}(\sigma_1(a),\ldots,\sigma_n(a))V$.
Wyznacznik $V$ wynosi
$\prod_{i<j}(\gamma_j-\gamma_i)\ne0$, więc macierze
$A$ i diagonalna są podobne nad $\Omega$. Ich ślady
i wyznaczniki są równe, co daje obie żądane formuły.
\end{proof}

\begin{lemat}[Wyznacznik przekształcenia po ograniczeniu skalarów]
\label{aa:lem:restriction-determinant}
Niech $F/K$ będzie skończonym rozszerzeniem ciał stopnia $d$
i niech $r\ge1$ będzie liczbą całkowitą.
Niech $T:F^r\to F^r$ będzie dowolnym $F$-liniowym endomorfizmem.
Jeżeli $A\in M_r(F)$ jest jego macierzą w bazie standardowej,
to po potraktowaniu $F^r$ jako $K$-przestrzeni zachodzi
$\det_K T=N_{F/K}(\det_F A)$.
\end{lemat}
\begin{proof}
Jeżeli $A$ jest osobliwa, to $T$ ma niezerowy wektor
w jądrze również jako odwzorowanie $K$-liniowe.
Obie strony są wtedy zerowe.
Załóżmy, że $A$ jest odwracalna. Metoda eliminacji Gaussa
rozkłada $A$ na iloczyn macierzy elementarnych trzech rodzajów.
Wystarczy więc wykazać wzór dla każdej z nich,
ponieważ wyznaczniki oraz norma są multiplikatywne.

Jeśli macierz elementarna dodaje do jednej współrzędnej
wielokrotność innej, to po ograniczeniu skalarów
jej macierz jest blokową macierzą elementarną
z identycznościami na przekątnej i jednym blokiem
pozadiagonalnym. Jej wyznacznik nad $K$ wynosi $1$,
tak samo jak norma wyznacznika $1$ nad $F$.
Jeżeli mnożymy jedną współrzędną przez $u\in F^\times$,
to jako $K$-liniowy endomorfizm mamy na tej współrzędnej
mnożenie $m_u:F\to F$, a na pozostałych identyczność.
Wyznacznik nad $K$ jest zatem $N_{F/K}(u)$,
czyli normą wyznacznika nad $F$.
Wreszcie zamiana dwóch współrzędnych $F$ odpowiada
zamianie dwóch bloków o wymiarze $d$ nad $K$.
Jej wyznacznik wynosi $(-1)^{d^2}=(-1)^d$.
Wyznacznik tej operacji nad $F$ jest równy $-1$,
a $N_{F/K}(-1)=(-1)^d$ na mocy
Twierdzenia~\ref{aa:tw:trace-norm-basic}.
Wszystkie trzy rodzaje operacji spełniają wzór,
co kończy dowód.
\end{proof}

\begin{twierdzenie}[Przechodniość śladu i normy]
\label{aa:tw:trace-norm-tower}
\label{aa:thm:trace-norm-tower}
Niech $K\subseteq F\subseteq L$ będą ciałami takimi, że
$[F:K]<\infty$ i $[L:F]<\infty$
(nie wymagamy rozdzielczości).
Wtedy dla każdego $a\in L$ mamy
\[
\operatorname{Tr}_{L/K}(a)=\operatorname{Tr}_{F/K}\bigl(
\operatorname{Tr}_{L/F}(a)\bigr),\qquad
N_{L/K}(a)=N_{F/K}\bigl(N_{L/F}(a)\bigr).
\]
\end{twierdzenie}
\begin{proof}
Niech $r=[L:F]$, $d=[F:K]$ i wybierzmy bazę
$u_1,\ldots,u_r$ przestrzeni $L$ nad $F$.
Macierz $F$-liniowego mnożenia $m_a:L\to L$
oznaczmy przez $A=(a_{ij})\in M_r(F)$.
Jeśli $v_1,\ldots,v_d$ jest bazą $F$ nad $K$,
to $u_iv_j$ dla $1\le i\le r$, $1\le j\le d$
tworzą bazę $L$ nad $K$. W tej bazie macierz
$K$-liniowego mnożenia $m_a$ ma bloki
odpowiadające $K$-liniowym przekształceniom
$m_{a_{ij}}:F\to F$.
Zatem jej ślad to suma śladów bloków diagonalnych:
\[
\operatorname{Tr}_{L/K}(a)
=\sum_{i=1}^r\operatorname{Tr}_{F/K}(a_{ii})
=\operatorname{Tr}_{F/K}\bigl(\sum_{i=1}^r a_{ii}\bigr)
=\operatorname{Tr}_{F/K}(\operatorname{Tr}_{L/F}(a)).
\]
Dla normy stosujemy
Lemat~\ref{aa:lem:restriction-determinant}
do tego samego $F$-liniowego endomorfizmu $m_a$:
\[
N_{L/K}(a)=\det_K m_a
=N_{F/K}(\det_F A)=N_{F/K}(N_{L/F}(a)).
\]
Obie tożsamości są więc prawdziwe również dla
rozszerzeń nierozdzielczych.
\end{proof}

\begin{uwaga}[Wzory przez zanurzenia a rozdzielczość]
\label{aa:uw:trace-norm-inseparable}
Przechodniość śladu i normy jest twierdzeniem
czysto liniowoalgebraicznym, niewymagającym rozdzielczości.
Rozdzielczość jest natomiast konieczna w sformułowaniu
Twierdzenia~\ref{aa:tw:trace-norm-embeddings}
z dokładnie $[L:K]$ \emph{różnymi} zanurzeniami.
\end{uwaga}

\section{Forma śladu i wyróżnik}
\begin{definicja}[Forma śladu i wyróżnik bazy]
\label{aa:def:trace-discriminant}
Niech $L/K$ będzie skończonym rozszerzeniem stopnia $n$.
\emph{Formą śladu} na $L$ nazywamy $K$-dwuliniowe
odwzorowanie $B_{L/K}:L\times L\to K$,
$B_{L/K}(x,y)=\operatorname{Tr}_{L/K}(xy)$.
Dla bazy $b=(b_1,\ldots,b_n)$ definiujemy jej
\emph{wyróżnik}
$\operatorname{disc}_{L/K}(b)=\det\bigl(
\operatorname{Tr}_{L/K}(b_i b_j)\bigr)_{i,j=1}^n\in K$.
\end{definicja}

\begin{twierdzenie}[Niezdegenerowanie formy śladu]
\label{aa:tw:trace-nondegenerate}
Niech $L/K$ będzie skończonym rozszerzeniem rozdzielczym.
Wówczas forma śladu $B_{L/K}$ jest niezdegenerowana: dla każdego
$x\in L\setminus\{0\}$ istnieje $y\in L$ takie, że
$\operatorname{Tr}_{L/K}(xy)\ne0$. W szczególności
wyróżnik dowolnej bazy jest niezerowy.
\end{twierdzenie}
\begin{proof}
Niech $\Omega$ będzie domknięciem algebraicznym $K$
zawierającym $L$. Wybierzmy element pierwotny $\gamma$
i wszystkie $n=[L:K]$ różne zanurzenia
$\sigma_i:L\hookrightarrow\Omega$ na mocy
Twierdzenia~\ref{aa:tw:trace-norm-embeddings}. Dla bazy
$b_j=\gamma^{j-1}$ macierz formy śladu ma
współczynniki
$\operatorname{Tr}_{L/K}(b_jb_k)=\sum_i\sigma_i(b_j)\sigma_i(b_k)$.
Jeśli $V=(\sigma_i(b_j))$, to macierz formy wynosi
$V^{\mathsf T}V$. Zatem jej wyznacznik jest równy
$(\det V)^2$, a wyznacznik Vandermonde'a $V$ jest niezerowy.

Warto podkreślić, że macierz $V$ ma współczynniki w $\Omega$,
a niekoniecznie w $K$. Nie przeszkadza to w argumentacji:
niezerowość wyznacznika macierzy o współczynnikach w $K$
można sprawdzić po rozszerzeniu skalarów do $\Omega$.
W bazie potęgowej dostajemy ponadto konkretną formułę
\[
\operatorname{disc}_{L/K}(1,\gamma,\ldots,\gamma^{n-1})
=\prod_{1\le i<j\le n}(\sigma_j(\gamma)-\sigma_i(\gamma))^2.
\]
Prawa strona nie zależy od wybranej numeracji zanurzeń,
ponieważ zamiana dwóch wierszy macierzy $V$ zmienia
znak $\det V$, ale nie zmienia jego kwadratu.
Dla dowolnej innej bazy macierz formy zmienia się
na $S^{\mathsf T}(V^{\mathsf T}V)S$ dla odwracalnej
macierzy $S\in\operatorname{GL}_n(K)$, więc również
ma niezerowy wyznacznik. Macierz odwracalna oznacza,
że forma ma zerowe jądro, co daje tezę.
\end{proof}

\begin{przyklad}[Wyróżnik bazy kwadratowej]
\label{aa:prz:quadratic-discriminant}
Niech $d\in\mathbb Z\setminus\{0,1\}$ będzie liczbą
bezkwadratową, $L=\mathbb Q(\sqrt d)$, a
$b=(1,\sqrt d)$ bazą $L$ nad $\mathbb Q$.
Wówczas $\operatorname{Tr}(1)=2$,
$\operatorname{Tr}(\sqrt d)=0$ oraz
$\operatorname{Tr}(d)=2d$. Macierz formy śladu to
$\begin{psmallmatrix}2&0\\0&2d\end{psmallmatrix}$,
a wyróżnik jest równy $4d$. W szczególności
nie zanika, gdy $d\ne0$.
\end{przyklad}

\begin{lemat}[Zmiana bazy i wyróżnik]\label{aa:lem:trace-discriminant-base-change}
Niech $L/K$ będzie dowolnym skończonym rozszerzeniem
stopnia $n\ge1$ i niech $b=(b_1,\ldots,b_n)$,
$c=(c_1,\ldots,c_n)$ będą dwiema bazami $L$ nad $K$.
Jeżeli $S=(s_{ij})\in\operatorname{GL}_n(K)$
spełnia $c_j=\sum_{i=1}^n s_{ij}b_i$, to
\[
\operatorname{disc}_{L/K}(c)=(\det S)^2
\operatorname{disc}_{L/K}(b).
\]
Nie zakładamy rozdzielczości; obie strony mogą być zerowe.
\end{lemat}
\begin{proof}
Niech $G_b=(\operatorname{Tr}_{L/K}(b_ib_j))_{i,j}$.
Dla dowolnych $j,k$ dwuliniowość formy śladu daje
\[
\operatorname{Tr}_{L/K}(c_jc_k)
=\sum_{i=1}^n\sum_{h=1}^n
s_{ij}s_{hk}\operatorname{Tr}_{L/K}(b_ib_h).
\]
To właśnie wzór na wyraz $(j,k)$ macierzy
$G_c=S^{\mathsf T}G_bS$. Biorąc wyznaczniki,
otrzymujemy
\[
\det G_c=\det(S^{\mathsf T})\det G_b\det S
=(\det S)^2\det G_b,
\]
jak żądano.
\end{proof}
\begin{przyklad}[Wyróżnik zależy od bazy]
\label{aa:prz:trace-discriminant-other-basis}
Weźmy $L=\mathbb Q(\sqrt5)$. Dla bazy
$b=(1,\sqrt5)$ poprzedni przykład daje $\operatorname{disc}_{L/\mathbb Q}(b)=20$.
Rozważmy inną bazę $c=(1,(1+\sqrt5)/2)$.
Jej macierz przejścia z bazy $b$ wynosi
$S=\bigl(\begin{smallmatrix}1&1/2\\0&1/2\end{smallmatrix}\bigr)$,
więc $\det S=1/2\ne0$. Z lematu otrzymujemy
$\operatorname{disc}_{L/\mathbb Q}(c)=(1/2)^2\cdot20=5$.
Można to skontrolować bezpośrednio:
$\operatorname{Tr}(1)=2$,
$\operatorname{Tr}((1+\sqrt5)/2)=1$, a
$((1+\sqrt5)/2)^2=(3+\sqrt5)/2$, którego ślad wynosi $3$.
Macierz formy w bazie $c$ to
$\bigl(\begin{smallmatrix}2&1\\1&3\end{smallmatrix}\bigr)$
i jej wyznacznik wynosi $6-1=5$.
Zatem wyróżnik \emph{bazy} nie jest niezmiennikiem liczbowym
samego rozszerzenia; jego klasa modulo niezerowe kwadraty
jest niezależna od bazy. Nie należy tej definicji utożsamiać
bez dodatkowych założeń z wyróżnikiem wielomianu.
\end{przyklad}


\section{Ślad i norma w ciałach skończonych}
\begin{twierdzenie}[Jawne wzory nad ciałem pierwszym]
\label{aa:tw:finite-field-trace-norm}
Niech $q=p^n$, gdzie $p$ jest liczbą pierwszą,
a $n\ge1$. Dla $a\in\mathbb F_q$ zachodzą
\[
\operatorname{Tr}_{\mathbb F_q/\mathbb F_p}(a)
=\sum_{j=0}^{n-1}a^{p^j},\qquad
N_{\mathbb F_q/\mathbb F_p}(a)=a^{1+p+\cdots+p^{n-1}}.
\]

Wykładnik w formule normy wynosi
$1+p+\cdots+p^{n-1}=(p^n-1)/(p-1)$.
Dla $a=0$ powyższy wzór na normę również daje $0$,
ponieważ wykładnik jest dodatni (także dla $n=1$).
Oba odwzorowania są surjektywne na odpowiednio
$(\mathbb F_p,+)$ i $(\mathbb F_p^\times,\cdot)$,
przy czym dla normy ograniczamy dziedzinę do $\mathbb F_q^\times$.
\end{twierdzenie}
\begin{proof}
Grupę Galois $\mathbb F_q/\mathbb F_p$ generuje
$\operatorname{Fr}_p:a\mapsto a^p$ na mocy
Twierdzenia~\ref{aa:tw:finite-field-frobenius-group}.
Wzory wynikają z
Twierdzenia~\ref{aa:tw:trace-norm-embeddings}.
Ponieważ $\mathbb F_q^\times$ jest cykliczna rzędu $q-1$,
obraz homomorfizmu $x\mapsto x^{(q-1)/(p-1)}$
ma rząd
$(q-1)/\gcd(q-1,(q-1)/(p-1))=p-1$.
Jest to cała grupa $\mathbb F_p^\times$.
Ślad jest liniowym odwzorowaniem przestrzeni
$\mathbb F_q$ nad $\mathbb F_p$ w jednowymiarową
przestrzeń $\mathbb F_p$. Nie jest zerowy: w przeciwnym
razie każdy z $q$ elementów byłby pierwiastkiem
wielomianu $T(X)=\sum_{j=0}^{n-1}X^{p^j}$,
którego stopień wynosi $p^{n-1}<q$ i który jest
niezerowy. Wobec tego obraz śladu jest niezerową
podprzestrzenią jednowymiarowej przestrzeni $\mathbb F_p$,
a więc całym $\mathbb F_p$.
\end{proof}

\begin{przyklad}[Ślad w ciele czteroelementowym]
\label{aa:prz:finite-field-trace-four}
Niech $\mathbb F_4=\mathbb F_2(\omega)$, gdzie
$\omega^2+\omega+1=0$. Wówczas $\omega^2=\omega+1$
i $\omega^3=1$, więc $\mathbb F_4=\{0,1,\omega,\omega^2\}$.
Z poprzedniego twierdzenia dla $p=n=2$ mamy
$\operatorname{Tr}(x)=x+x^2$ i $N(x)=x^3$.
Dostajemy
\[
\begin{array}{c|cccc}
x&0&1&\omega&\omega^2\\\hline
\operatorname{Tr}(x)&0&0&1&1\\
N(x)&0&1&1&1
\end{array}
\]
Na przykład $\operatorname{Tr}(\omega)=\omega+\omega^2
=\omega+(\omega+1)=1$, a
$N(\omega)=\omega^{1+2}=\omega^3=1$.
Ten rachunek pokazuje, dlaczego w charakterystyce $p$
nie wystarczy sprawdzić $\operatorname{Tr}(1)$:
tutaj $\operatorname{Tr}(1)=2=0$, mimo że ślad jest surjektywny.
\end{przyklad}

\begin{uwaga}[Co dalej?]
\label{aa:uw:trace-norm-next}
Ślad i norma są użyteczne w arytmetyce ciał liczbowych,
teorii form kwadratowych i opisie rozszerzeń cyklicznych.
W następnym rozdziale połączymy grupy Galois z pytaniem,
czy pierwiastki wielomianu można wyrazić za pomocą
skończonego ciągu działań arytmetycznych i wyciągania pierwiastków.
\end{uwaga}

% AA-24 | Algebra abstrakcyjna | Iteracja XII
\chapter[Rozwiązywalność przez pierwiastniki]{Rozwiązywalność równań przez pierwiastniki i twierdzenie Abela--Ruffiniego}
\label{aa:ch:radical-solvability}
\label{aa:ch:radicals}

\section{Co oznacza rozwiązanie równania przez pierwiastniki?}
Wzór kwadratowy, wzory Cardana oraz wzory dla równań
czwartego stopnia wykorzystują działania arytmetyczne i
wyciąganie pierwiastków. Nie oznacza to, że każdy wielomian
piątego stopnia można rozwiązać w ten sposób.
Sformułujemy precyzyjne kryterium rozwiązywalności przez
pierwiastniki dla wielomianów nad ciałem charakterystyki zero.
Jego treść łączy teorię rozszerzeń ciał z teorią grup rozwiązalnych
z Rozdziału~\ref{aa:ch:szeregi}.

\begin{definicja}[Wieża rozszerzeń przez pierwiastniki]
\label{aa:def:radical-tower}
\label{aa:def:radical-extension}
Niech $K$ będzie ciałem charakterystyki zero.
\emph{Wieżą rozszerzeń przez pierwiastniki} nad $K$ nazywamy
skończony ciąg rozszerzeń ciał
$K=E_0\subseteq E_1\subseteq\cdots\subseteq E_r$
taki, że dla każdego $i=1,\ldots,r$ istnieją
$\alpha_i\in E_i$ i liczba całkowita $n_i\ge2$, dla których
$E_i=E_{i-1}(\alpha_i)$ i $\alpha_i^{n_i}\in E_{i-1}$.
Wielomian $f\in K[X]$ nazywamy \emph{rozwiązywalnym przez
pierwiastniki nad $K$}, jeśli istnieje taka wieża, której
ostatnie ciało zawiera wszystkie pierwiastki $f$
w pewnym ustalonym domknięciu algebraicznym ciała $K$.
\end{definicja}

\begin{uwaga}[Dlaczego dopuszczamy pierwiastki z jedynki?]
\label{aa:uw:radicals-roots-unity}
Jeśli $\zeta_m$ jest pierwiastkiem pierwotnym $m$-tego stopnia
z jedynki, to $\zeta_m^m=1\in K$; dołączenie $\zeta_m$
jest zatem krokiem dopuszczonym w wieży. Możemy więc zawsze
najpierw dołączyć skończenie wiele potrzebnych pierwiastków
z jedynki. Nie oznacza to, że rozpatrujemy jedynie
pierwiastki rzeczywiste: rozszerzenia i ich generatory
mogą leżeć w dowolnym domknięciu algebraicznym.
\end{uwaga}

\begin{przyklad}[Równanie trzeciego stopnia]
\label{aa:prz:radical-cubic-simple}
Dla $f=X^3-2\in\mathbb Q[X]$ jego pierwiastki to
$\alpha=\sqrt[3]2$, $\omega\alpha$, $\omega^2\alpha$,
gdzie $\omega^2+\omega+1=0$ i $\omega^3=1$.
Wieża
$\mathbb Q\subseteq\mathbb Q(\alpha)
\subseteq\mathbb Q(\alpha,\omega)$
spełnia definicję, ponieważ $\alpha^3=2$ i $\omega^3=1$.
Jest to przykład równania rozwiązywalnego przez pierwiastniki,
choć jego grupa Galois jest nieabelową grupą $S_3$
(Przykład~\ref{aa:prz:cubic-galois-s3}).
Zatem \emph{abelowość} grupy Galois nie jest
warunkiem koniecznym rozwiązywalności przez pierwiastniki.
\end{przyklad}

\begin{przyklad}[Dlaczego pojedynczy krok radykalny nie musi być Galois]
\label{aa:prz:radicals-not-galois-step}
Niech $\alpha=\sqrt[3]{2}\in\mathbb R$, a $\omega\in\mathbb C$
niech będzie pierwotnym pierwiastkiem trzeciego stopnia z jedynki.
Wielomian $X^3-2$ jest nierozkładalny nad $\mathbb Q$
na mocy kryterium Eisensteina dla liczby $2$, toteż
$[\mathbb Q(\alpha):\mathbb Q]=3$. Ciało
$\mathbb Q(\alpha)$ zawiera się w $\mathbb R$, a więc
nie zawiera nierzeczywistych pierwiastków $\omega\alpha$
i $\omega^2\alpha$. Rozszerzenie $\mathbb Q(\alpha)/\mathbb Q$
nie jest normalne ani Galois, chociaż $\alpha^3\in\mathbb Q$.
Natomiast $L=\mathbb Q(\alpha,\omega)$ jest ciałem rozkładu;
$[L:\mathbb Q(\alpha)]=2$, bo $\omega$ jest nierzeczywisty
i spełnia $\omega^2+\omega+1=0$. Wobec tego $[L:\mathbb Q]=6$.
Wybierając najpierw dołączenie $\omega$, a następnie $\alpha$,
otrzymujemy wieżę o cyklicznych grupach względnych:
$\mathbb Q(\omega)/\mathbb Q$ jest kwadratowe,
a $L/\mathbb Q(\omega)$ jest cykliczne stopnia $3$.
Ostatnie stwierdzenie wynika z rozdzielczości i normalności
$L/\mathbb Q(\omega)$ oraz z faktu, że jego automorfizmy
wysyłają $\alpha$ na $\omega^j\alpha$.
Grupa Galois $L/\mathbb Q$ ma rząd $6$ i zawiera
cykl trzech pierwiastków oraz sprzężenie zespolone,
które zamienia dwa pierwiastki; jest zatem $S_3$.
Przykład uzasadnia potrzebę dołączania pierwiastków z jedynki
przed stosowaniem lematu o cyklicznych krokach Galois.
\end{przyklad}

\section{Rozszerzenia cykliczne i pierwiastki z jedynki}
\begin{lemat}[Rozszerzenie przez pierwiastek przy obecności pierwiastków z jedynki]
\label{aa:lem:kummer-simple-galois}
\label{aa:lem:kummer-cyclic}
Niech $F$ będzie ciałem charakterystyki zero,
$n\ge2$ i niech $F$ zawiera wszystkie $n$-te pierwiastki
z jedynki. Jeśli $\alpha\ne0$ spełnia $\alpha^n=a\in F^\times$
i $E=F(\alpha)$, to $E/F$ jest skończonym rozszerzeniem
Galois, a jego grupa Galois jest izomorficzna z podgrupą
cyklicznej grupy $\mu_n$ pierwiastków $n$-tego stopnia
z jedynki. W szczególności jest abelowa.
\end{lemat}
\begin{proof}
Element $\alpha$ jest algebraiczny nad $F$, ponieważ
zeruje $X^n-a$. Każdy pierwiastek tego wielomianu
ma postać $\zeta\alpha$ dla pewnego $\zeta\in\mu_n\subseteq F$,
więc wszystkie pierwiastki należą do $E$.
Wielomian jest rozdzielczy: jego pochodna to
$nX^{n-1}$, niezerowa w charakterystyce zero i
niemająca wspólnego pierwiastka z $X^n-a$.
Stąd $E/F$ jest ciałem rozkładu wielomianu rozdzielczego,
czyli rozszerzeniem Galois.
Dla $\sigma\in\operatorname{Gal}(E/F)$ istnieje jednoznaczne
$\zeta_\sigma\in\mu_n$ z $\sigma(\alpha)=\zeta_\sigma\alpha$.
Ponieważ $\sigma$ zachowuje wszystkie elementy $\mu_n$,
$\zeta_{\sigma\tau}=\zeta_\sigma\zeta_\tau$.
Wzór $\sigma\mapsto\zeta_\sigma$ jest zatem
homomorfizmem grup. Jego jądro jest trywialne,
ponieważ $\sigma$ ustalający $F$ i $\alpha$ ustala całe $E$.
Jest więc iniekcją w $\mu_n$, które jest grupą cykliczną.
\end{proof}

\begin{lemat}[Cykliczne rozszerzenie stopnia będącego liczbą pierwszą]
\label{aa:lem:kummer-cyclic-prime}
Niech $p$ będzie liczbą pierwszą, $F$ ciałem
charakterystyki zero zawierającym pierwotny $p$-ty
pierwiastek z jedynki $\zeta$, a $E/F$ rozszerzeniem
Galois o cyklicznej grupie rzędu $p$.
Wtedy istnieje $\alpha\in E$ takie, że
$E=F(\alpha)$ oraz $\alpha^p\in F$.
\end{lemat}
\begin{proof}
Niech $\sigma$ generuje grupę Galois $G$ rzędu $p$.
Rozszerzenie ma stopień $[E:F]=p$, a $\sigma$ jest
$F$-liniowym endomorfizmem $E$ spełniającym
$\sigma^p=\operatorname{id}$. Ponieważ $\operatorname{char}F=0$,
liczba $p$ jest odwracalna w $F$. Dla $0\le j<p$ określmy
jawnie operator $F$-liniowy
\[
P_j=\frac1p\sum_{k=0}^{p-1}\zeta^{-jk}\sigma^k:E\longrightarrow E.
\]
Po zastąpieniu indeksu $k$ przez $k+1$ i wykorzystaniu
$\sigma^p=\operatorname{id}$ oraz $\zeta^p=1$ dostajemy
$\sigma P_j=\zeta^j P_j$. Ponadto
\[
\sum_{j=0}^{p-1}P_j
=\frac1p\sum_{k=0}^{p-1}\left(\sum_{j=0}^{p-1}\zeta^{-jk}\right)\sigma^k
=\operatorname{id},
\]
ponieważ wewnętrzna suma wynosi $p$ dla $k=0$
i $0$ dla $1\le k<p$ (jest sumą geometryczną
z ilorazem różnym od $1$). Jeżeli $P_j=0$ dla
wszystkich $1\le j<p$, to $P_0=\operatorname{id}$,
a więc z $\sigma P_0=P_0$ wynikałoby
$\sigma=\operatorname{id}$, wbrew $|G|=p>1$.
Istnieje zatem $j\in\{1,\ldots,p-1\}$, dla którego
$P_j\ne0$. Wybierzmy $x\in E$ tak, aby
$\alpha=P_j(x)\ne0$. Mamy $\sigma(\alpha)=\zeta^j\alpha$.
Podnosząc tę równość do potęgi $p$, otrzymujemy
$\sigma(\alpha^p)=\alpha^p$, a ponieważ ciałem stałym
$E^G$ jest $F$ (Twierdzenie~\ref{aa:tw:artin-fixed-field}),
wnioskujemy $\alpha^p\in F$.

Liczba $\zeta^j$ jest pierwiastkiem pierwotnym stopnia $p$,
gdyż $p$ jest pierwsza i $j\not\equiv0\pmod p$.
Elementy $\sigma^k(\alpha)=\zeta^{jk}\alpha$ dla
$0\le k<p$ są zatem parami różne.
Każdy jest pierwiastkiem wielomianu minimalnego
$m_{\alpha,F}\in F[X]$, ponieważ automorfizmy $G$
ustalają jego współczynniki. Wynika stąd
$\deg m_{\alpha,F}\ge p$. Z drugiej strony
$F(\alpha)\subseteq E$ daje $[F(\alpha):F]\le p$.
Zatem $[F(\alpha):F]=p=[E:F]$,
a z twierdzenia o wieży wynika $E=F(\alpha)$.
\end{proof}

\section{Skończone grupy rozwiązalne i wieże ciał}
\begin{lemat}[Szereg podnormalny, którego kolejne ilorazy mają rząd będący liczbą pierwszą]
\label{aa:lem:solvable-prime-series}
Skończona grupa $G$ jest rozwiązalna wtedy i tylko wtedy,
gdy istnieje szereg podnormalny
$G=G_0\triangleright G_1\triangleright\cdots
\triangleright G_r=\{1\}$, w którym każda grupa
ilorazowa $G_{i-1}/G_i$ jest cykliczna, a jej rząd jest liczbą pierwszą.
\end{lemat}
\begin{proof}
Jeśli istnieje taki szereg, wszystkie jego czynniki
są abelowe, więc $G$ jest rozwiązalna na mocy
Twierdzenia~\ref{aa:tw:solvable-series}.
W drugą stronę wybierzmy szereg z czynnikami abelowymi,
którego istnienie wynika z tego samego twierdzenia.
Każda skończona grupa abelowa ma podgrupę indeksu
pierwszego: wybieramy czynnik cykliczny rzędu podzielnego
przez pewną liczbę pierwszą $p$ i homomorfizm na
$\mathbb Z/p\mathbb Z$. Jego jądro jest normalne,
a iloraz ma rząd $p$. Powtarzając ten argument,
dostajemy szereg podnormalny o ilorazach cyklicznych
o rzędach będących liczbami pierwszymi. Wstawiamy takie szeregi do szeregu
wyjściowego; przeciwobrazy podgrup przez homomorfizmy
$G_{i-1}\to G_{i-1}/G_i$ zapewniają normalność
kolejnych włączeń. Liczba kroków jest skończona,
ponieważ w każdym właściwym kroku rząd grupy maleje.

Dokładniej, jeśli $G_{i-1}/G_i=A_0\triangleright
A_1\triangleright\cdots\triangleright A_t=\{1\}$
jest szeregiem z ilorazami o rzędach będących liczbami pierwszymi, to dla
naturalnej projekcji $\pi:G_{i-1}\to G_{i-1}/G_i$
wstawiamy grupy $\pi^{-1}(A_j)$ między
$G_{i-1}$ i $G_i$. Izomorfizm
$\pi^{-1}(A_{j-1})/\pi^{-1}(A_j)\cong A_{j-1}/A_j$
uzasadnia dokładnie rzędy nowych ilorazów.
Dla grupy abelowej $A$ istnienie takiego szeregu
wynika np. z rozkładu na sumę prostą cyklicznych
$p$-grup: z czynnika $C_{p^e}$ uzyskujemy
kolejno $C_{p^e}\supset C_{p^{e-1}}\supset\cdots
\supset C_p\supset\{1\}$; łączymy je po jednym czynniku.
Wszystkie podgrupy grupy abelowej są normalne.
\end{proof}

\begin{twierdzenie}[Kryterium rozwiązywalności przez pierwiastniki]
\label{aa:tw:radicals-solvable-galois}
\label{aa:thm:radicals-implies-solvable}
Niech $K$ będzie ciałem charakterystyki zero,
$f\in K[X]$ wielomianem niezerowym dodatniego stopnia,
a $L/K$ jego ciałem rozkładu w ustalonym domknięciu
algebraicznym $K$. Wtedy $f$ jest rozwiązywalny
przez pierwiastniki nad $K$ wtedy i tylko wtedy,
gdy skończona grupa $\operatorname{Gal}(L/K)$ jest rozwiązalna.
\end{twierdzenie}
\begin{proof}
W charakterystyce zero wielomian $f$ po usunięciu
powtórzonych czynników ma pierwiastki pojedyncze,
a jego ciało rozkładu $L/K$ jest skończonym rozszerzeniem
Galois. Oznaczmy $G=\operatorname{Gal}(L/K)$.

\emph{Krok 1: grupa rozwiązalna daje wieżę przez pierwiastniki.}
Załóżmy, że $G$ jest rozwiązalna. Niech
$M$ będzie iloczynem wszystkich różnych liczb pierwszych
dzielących $|G|$; jeśli $G=\{1\}$, przyjmujemy $M=1$.
Dołączmy do $K$ pierwiastek pierwotny $M$-tego stopnia
z jedynki i oznaczmy otrzymane ciało przez $K'$.
Mamy $K'=K(\zeta_M)$ i $\zeta_M^M=1\in K$,
więc przejście $K\subseteq K'$ jest dozwolone w wieży
(gdy $M=1$, nie potrzebujemy tego kroku).
Rozważmy ciało złożone $L'=LK'$ we wspólnym domknięciu.
Ponieważ $L/K$ jest Galois, również $L'/K'$ jest Galois.
Ograniczenie automorfizmów daje iniekcję
$G'=\operatorname{Gal}(L'/K')\hookrightarrow G$:
automorfizm ustalający $K'$ i wszystkie elementy $L$
ustala ciało przez nie generowane. Podgrupa grupy rozwiązalnej
jest rozwiązalna; wynika to na przykład z inkluzji
jej kolejnych podgrup pochodnych w odpowiednich
podgrupach pochodnych $G$.

Z Lematu~\ref{aa:lem:solvable-prime-series} wybierzmy
szereg podnormalny
$G'=H_0\triangleright H_1\triangleright\cdots
\triangleright H_s=\{1\}$ z ilorazami cyklicznymi
o rzędach będących liczbami pierwszymi. Odpowiedniość Galois
(Twierdzenie~\ref{aa:tw:fundamental-galois}) daje wieżę
$K'=(L')^{H_0}\subseteq (L')^{H_1}\subseteq\cdots
\subseteq (L')^{H_s}=L'$.
Ponieważ $H_i\triangleleft H_{i-1}$,
każdy kolejny krok jest cyklicznym rozszerzeniem
Galois stopnia $p_i$, gdzie $p_i$ jest liczbą pierwszą. Wszystkie $p_i$
dzielą $|G|$, a więc w $K'$ znajdują się pierwiastki
pierwotne $p_i$-tego stopnia z jedynki.
Lemat~\ref{aa:lem:kummer-cyclic-prime} wykazuje,
że każdy taki krok powstaje przez dołączenie jednego
elementu, którego $p_i$-ta potęga leży w poprzednim ciele.
Jest to wieża przez pierwiastniki zawierająca $L$,
więc zawiera wszystkie pierwiastki $f$.

\emph{Krok 2: wieża przez pierwiastniki daje grupę rozwiązalną.}
Załóżmy, że mamy wieżę
$K=E_0\subseteq E_1\subseteq\cdots\subseteq E_r$
z $E_i=E_{i-1}(\alpha_i)$, $\alpha_i^{n_i}=a_i\in E_{i-1}$,
i $L\subseteq E_r$. Możemy pominąć kroki, w których
$\alpha_i=0$, bo nie zmieniają one ciała.
Wybierzmy wspólną wielokrotność $N$ liczb $n_1,\ldots,n_r$
i niech $K'_0=K(\zeta_N)$ (dla wieży pustej bierzemy $K'_0=K$).
Ciało $K'_0/K$ jest skończonym rozszerzeniem Galois:
jest ciałem rozkładu $X^N-1$ w charakterystyce zero.
Jego grupa Galois jest abelowa, ponieważ dowolny automorfizm
wysyła $\zeta_N$ na $\zeta_N^u$, gdzie
$u\in(\mathbb Z/N\mathbb Z)^\times$, a przyporządkowanie
$\sigma\mapsto u$ jest iniektywnym homomorfizmem
w grupę abelową.

Zbudujemy indukcyjnie skończone rozszerzenia Galois
$K'_i/K$, zawierające $E_i$ i $K'_0$, takie że
$\operatorname{Gal}(K'_i/K'_{i-1})$ jest abelowa.
Załóżmy, że $K'_{i-1}$ jest już skonstruowane.
Rozpatrzmy wszystkie $K$-zanurzenia
$\sigma:E_{i-1}\hookrightarrow K'_{i-1}$;
istnieją, ponieważ $E_{i-1}\subseteq K'_{i-1}$,
a każde takie zanurzenie przedłuża się do automorfizmu
$K'_{i-1}/K$ na mocy normalności i
Wniosku~\ref{aa:wn:extend-embedding-finite}.
Dla wszystkich tych $\sigma$ dołączmy do $K'_{i-1}$
wszystkie pierwiastki wielomianów
$X^{n_i}-\sigma(a_i)$; otrzymane ciało oznaczmy
przez $K'_i$. Zawiera ono $E_i$, ponieważ wśród
zanurzeń jest identyczność na $E_{i-1}$, a $\alpha_i$
jest jednym z dołączonych pierwiastków.
Ponieważ $K'_{i-1}$ zawiera $\mu_{n_i}$, dla każdego
niezerowego $\sigma(a_i)$ wystarczy dołączyć jeden
pierwiastek, by otrzymać wszystkie pozostałe.
Rozszerzenie $K'_i/K'_{i-1}$ jest zatem ciałem złożonym
skończenie wielu rozszerzeń Galois o cyklicznych
grupach na mocy Lematu~\ref{aa:lem:kummer-simple-galois}.
Jego grupa Galois jest abelowa: ograniczenia
automorfizmów do tych prostych rozszerzeń dają
iniekcję w iloczyn prosty ich abelowych grup Galois.

Wykażmy normalność $K'_i/K$.
Każde $K$-zanurzenie $\rho:K'_i\to\Omega$
zachowuje $K'_{i-1}$ jako zbiór, gdyż
$K'_{i-1}/K$ jest normalne. Ograniczenie $\rho$
permutuje wszystkie $K$-zanurzenia
$E_{i-1}\hookrightarrow K'_{i-1}$ poprzez składanie.
Wysyła więc pierwiastki
$X^{n_i}-\sigma(a_i)$ na pierwiastki wielomianów
tej samej rodziny i zachowuje $K'_i$ jako zbiór.
Ponieważ rozszerzenie jest skończone i rozdzielcze,
wynika stąd, że $K'_i/K$ jest Galois.
Otrzymujemy zatem ciąg dokładny grup
\[
1\to\operatorname{Gal}(K'_i/K'_{i-1})
\to\operatorname{Gal}(K'_i/K)
\to\operatorname{Gal}(K'_{i-1}/K)\to1.
\]
Surjektywność ostatniego homomorfizmu wynika z
przedłużania zanurzeń. Skoro jądro jest abelowe,
a grupa po prawej jest rozwiązalna na mocy indukcji,
to grupa środkowa jest rozwiązalna
(Lemat~\ref{aa:lem:solvable-extension}).

W końcu $K'_r/K$ jest skończonym rozszerzeniem Galois
o rozwiązalnej grupie, a $L\subseteq K'_r$.
Ponieważ $L/K$ jest Galois, ograniczenie daje
surjekcję $\operatorname{Gal}(K'_r/K)\twoheadrightarrow
\operatorname{Gal}(L/K)$; każdemu automorfizmowi $L/K$
można przedłużyć działanie do większego ciała.
Obraz homomorficzny grupy rozwiązalnej jest rozwiązalny.
Stąd $G=\operatorname{Gal}(L/K)$ jest rozwiązalna.
\end{proof}

\section{Wyróżnik i znak permutacji pierwiastków}
\begin{definicja}[Wyróżnik wielomianu rozdzielczego]
\label{aa:def:polynomial-discriminant-galois}
Niech $K$ będzie ciałem charakterystyki różnej od $2$,
a $f\in K[X]$ unormowanym wielomianem rozdzielczym
stopnia $n\ge1$. Niech $\alpha_1,\ldots,\alpha_n$
będą wszystkimi jego pierwiastkami w ciele rozkładu $L$.
Definiujemy wyróżnik $f$ przez
$\Delta(f)=\prod_{1\le i<j\le n}(\alpha_i-\alpha_j)^2$.
Wartość $\Delta(f)$ nie zależy od kolejności pierwiastków.
\end{definicja}

\begin{twierdzenie}[Kryterium znaku grupy Galois]
\label{aa:tw:galois-discriminant-sign}
Przy powyższych założeniach niech
$G=\operatorname{Gal}(L/K)\le S_n$ działa na pierwiastkach.
Wtedy $\Delta(f)\in K^\times$ i
$G\subseteq A_n$ wtedy i tylko wtedy, gdy $\Delta(f)$
jest kwadratem elementu $K^\times$.
\end{twierdzenie}
\begin{proof}
Niech $\delta=\prod_{i<j}(\alpha_i-\alpha_j)\in L^\times$.
Każda permutacja pierwiastków zmienia $\delta$
o swój znak: zamiana dwóch pierwiastków zmienia znak,
a dowolna permutacja jest iloczynem transpozycji.
Stąd dla $\sigma\in G$ mamy
$\sigma(\delta)=\operatorname{sgn}(\sigma)\delta$.
Zatem $\sigma(\delta^2)=\delta^2$ dla wszystkich $\sigma$.
Ponieważ $L/K$ jest Galois, jego ciałem elementów
stałych jest $K$ na mocy
Twierdzenia~\ref{aa:tw:fundamental-galois}; stąd
$\Delta(f)=\delta^2\in K^\times$.
Jeżeli $G\subseteq A_n$, każdy automorfizm ustala
$\delta$, więc $\delta\in K$ i $\Delta(f)$ jest kwadratem.
Jeżeli natomiast $\Delta(f)=d^2$ dla $d\in K^\times$,
to $\delta^2=d^2$. W ciele charakterystyki różnej od $2$
wynika stąd $\delta=d$ lub $\delta=-d$, czyli $\delta\in K$.
Każdy automorfizm musi zatem mieć znak dodatni,
a więc $G\subseteq A_n$.
\end{proof}

\begin{przyklad}[Wyróżnik wielomianu trzeciego stopnia]
\label{aa:prz:discriminant-cubic}
Dla $f=X^3-2$ jego pierwiastki to
$\alpha$, $\omega\alpha$, $\omega^2\alpha$,
gdzie $\alpha^3=2$, $\omega^3=1$ i $\omega\ne1$.
Korzystamy z $(1-\omega)(1-\omega^2)=3$
i $(\omega-\omega^2)^2=-3$. Otrzymujemy
\[
\Delta(f)=\alpha^6(1-\omega)^2(1-\omega^2)^2
(\omega-\omega^2)^2=4\cdot9\cdot(-3)=-108.
\]
Liczba $-108$ nie jest kwadratem w $\mathbb Q$,
więc grupa Galois nie zawiera się w $A_3$.
Z wcześniejszej analizy wiemy, że jest równa $S_3$.
\end{przyklad}

\section{Twierdzenie Abela--Ruffiniego i konkretny wielomian}
\begin{twierdzenie}[Rozwiązywalność wielomianów stopnia co najwyżej cztery]
\label{aa:tw:radicals-degree-four}
Niech $K$ będzie ciałem charakterystyki zero.
Każdy niezerowy wielomian $f\in K[X]$ stopnia co najwyżej
cztery jest rozwiązywalny przez pierwiastniki nad $K$.
\end{twierdzenie}
\begin{proof}
Po usunięciu powtórzonych czynników możemy założyć,
że $f$ jest rozdzielczy, nie zmieniając jego ciała rozkładu.
Grupa Galois ciała rozkładu działa wiernie na co najwyżej
czterech pierwiastkach, więc jest podgrupą $S_n$ dla
$n\le4$. Grupa $S_4$ jest rozwiązalna, gdyż posiada szereg
$S_4\triangleright A_4\triangleright V_4\triangleright\{1\}$,
w którym kolejne ilorazy mają odpowiednio rzędy $2$, $3$ i $4$
i są abelowe ($V_4\cong C_2\times C_2$).

Tutaj $V_4=\{\operatorname{id},(12)(34),(13)(24),(14)(23)\}$;
jest podgrupą normalną $A_4$, ponieważ sprzężenie
zachowuje typ cykliczny permutacji. Mamy
$S_4/A_4\cong C_2$, $A_4/V_4\cong C_3$,
a końcowy iloraz $V_4/\{1\}$ jest abelowy.
Nie twierdzimy, że każdy z tych ilorazów musi być cykliczny:
$V_4$ nie jest cykliczna, ale jest abelowa, co wystarcza
w definicji rozwiązalności.
Każda podgrupa grupy rozwiązalnej jest rozwiązalna,
więc grupa Galois jest rozwiązalna.
Twierdzenie~\ref{aa:tw:radicals-solvable-galois} kończy dowód.
\end{proof}

\begin{lemat}[Kryterium grupy symetrycznej stopnia pięć]
\label{aa:lem:s5-criterion}
Jeżeli nierozkładalny wielomian $f\in\mathbb Q[X]$
stopnia pięć ma dokładnie trzy pierwiastki rzeczywiste,
to grupa Galois jego ciała rozkładu jest izomorficzna z $S_5$.
\end{lemat}
\begin{proof}
W charakterystyce zero $f$ ma pięć różnych pierwiastków.
Jego grupa Galois $G\le S_5$ działa na nich przechodnio:
izomorfizm prostych rozszerzeń odpowiadających dwóm
pierwiastkom przedłuża się do automorfizmu ciała rozkładu.
Ponieważ dokładnie dwa pierwiastki są nierzeczywiste,
sprzężenie zespolone zamienia tę parę i ustala trzy
pozostałe pierwiastki. Zatem $G$ zawiera transpozycję.
Utwórzmy graf na pięciu pierwiastkach, w którym
para jest krawędzią dokładnie wtedy, gdy zamieniająca
ją transpozycja należy do $G$.
Graf nie jest pusty, a $G$ permutuje jego składowe spójności,
ponieważ sprzężenie transpozycji przez $g\in G$
jest transpozycją końców przeniesionych przez $g$.
Przechodniość działania oznacza, że wszystkie składowe
mają tyle samo wierzchołków; ich liczność dzieli więc $5$.
Ponieważ istnieje krawędź, składowe nie mogą być jednopunktowe.
Graf jest zatem spójny. Transpozycje krawędzi grafu spójnego
generują pełną grupę symetryczną: na drzewie rozpinającym
przez kolejne sprzęganie transpozycji wzdłuż ścieżki
otrzymujemy transpozycję dowolnej pary wierzchołków.

Aby uzasadnić ostatni krok rachunkiem, niech
$v_0,v_1,\ldots,v_m$ będzie ścieżką bez powtórzeń
i $t_k=(v_{k-1}\,v_k)$ transpozycją odpowiadającą jej
$k$-tej krawędzi. Dla $m\ge2$ element
$h=t_1t_2\cdots t_{m-1}$ wysyła $v_{m-1}$ na $v_0$
i ustala $v_m$, więc
$h t_m h^{-1}=(v_0\,v_m)$.
Dla $m=1$ żądana transpozycja jest po prostu $t_1$.
Każde dwa wierzchołki są połączone ścieżką,
a więc w $G$ leży transpozycja dowolnej pary.
Ponieważ każda permutacja jest iloczynem transpozycji,
$G$ musi być całą grupą $S_5$.
Wobec tego $G=S_5$.
\end{proof}

\begin{figure}[htbp]
\centering
\begin{tikzpicture}[scale=0.86, every node/.style={font=\small}]
  \foreach \i in {0,1,2,3,4}{
    \node[draw,circle,minimum size=6.5mm,inner sep=0pt] (v\i) at (1.55*\i,0) {$v_{\i}$};
  }
  \foreach \i/\j in {0/1,1/2,2/3,3/4}{
    \draw[thick] (v\i) -- (v\j);
  }
  \draw[dashed,thick,bend left=38] (v0) to (v4);
  \node[above] at (3.1,1.45) {transpozycja otrzymana przez sprzężenie};
\end{tikzpicture}
\caption{Przykładowa ścieżka w grafie transpozycji.
Ciągłe krawędzie oznaczają transpozycje należące do $G$.
Linia przerywana nie oznacza nowego założenia:
transpozycja $(v_0\,v_4)$ należy do $G$ na mocy
rachunku sprzężeń z dowodu lematu.}
\label{aa:fig:transpositions-path}
\end{figure}

\begin{twierdzenie}[Konkretny nierozwiązywalny wielomian piątego stopnia]
\label{aa:tw:quintic-s5-example}
\label{aa:ex:quintic-unsolvable}
Wielomian $f(X)=X^5-4X+2\in\mathbb Q[X]$ ma
grupę Galois izomorficzną z $S_5$ i nie jest
rozwiązywalny przez pierwiastniki nad $\mathbb Q$.
\end{twierdzenie}
\begin{proof}
\emph{Krok 1: nierozkładalność.}
Wielomian $f$ spełnia kryterium Eisensteina
dla liczby pierwszej $2$: jego współczynnik
wiodący nie jest podzielny przez $2$,
pozostałe współczynniki są podzielne przez $2$,
a wyraz wolny $2$ nie jest podzielny przez $4$.
Zatem $f$ jest nierozkładalny nad $\mathbb Q$.
Ciało rozkładu $L$ ma grupę Galois $G$ działającą
wiernie i przechodnio na pięciu różnych pierwiastkach:
przechodniość wynika z możliwości przedłużania
$\mathbb Q$-zanurzeń wysyłających dowolny pierwiastek
$f$ na dowolny inny, a różność pierwiastków
z charakterystyki zero. Zatem $G\le S_5$ jest
podgrupą przechodnią.

\emph{Krok 2: liczba pierwiastków rzeczywistych.}
Pochodna wynosi $f'(x)=5x^4-4$.
Ma dokładnie dwa zera rzeczywiste $-c$ i $c$,
gdzie $c=(4/5)^{1/4}>0$.
Funkcja $f$ rośnie na $(-\infty,-c)$, maleje
na $(-c,c)$ i rośnie na $(c,+\infty)$.
Ponadto z $c^5=(4/5)c$ wynika
$f(-c)=2+(16/5)c>0$ oraz
$f(c)=2-(16/5)c<0$;
ostatnia nierówność zachodzi, gdyż $c>5/8$,
co wynika choćby z $c^4=4/5>(5/8)^4$.
Ponieważ $f(x)\to-\infty$ dla $x\to-\infty$
oraz $f(x)\to+\infty$ dla $x\to+\infty$,
na każdym z trzech wskazanych przedziałów
leży dokładnie jeden pierwiastek rzeczywisty.
Pozostałe dwa pierwiastki są nierzeczywiste
i wzajemnie sprzężone zespolenie.
Sprzężenie zespolone, ograniczone do $L$, jest
$\mathbb Q$-automorfizmem, bo $L$ jest ciałem
rozkładu wielomianu o współczynnikach wymiernych.
Jako permutacja pięciu pierwiastków jest transpozycją:
zachowuje trzy pierwiastki rzeczywiste i zamienia
pozostałe dwa. Stąd $G$ zawiera transpozycję.

\emph{Krok 3: przechodnia grupa permutacji pięciu elementów
zawierająca transpozycję jest całą grupą symetryczną.}
Udowodnimy to bez odwoływania się do klasyfikacji
podgrup $S_5$. Na zbiorze pięciu pierwiastków
rozważmy graf, którego krawędziami są pary
$\{i,j\}$ takie, że transpozycja $(ij)$ należy do $G$.
Graf ma co najmniej jedną krawędź. Jeśli
$g\in G$ i $(ij)\in G$, to
$g(ij)g^{-1}=(g(i)\,g(j))\in G$,
więc $G$ permutuje składowe spójności grafu.
Ponieważ działanie $G$ na pięciu wierzchołkach
jest przechodnie, wszystkie składowe mają tę samą
liczbę wierzchołków. Ta liczba dzieli $5$.
Nie może wynosić $1$, gdyż istnieje krawędź,
a więc wynosi $5$: graf jest spójny.
Transpozycje odpowiadające krawędziom dowolnego
grafu spójnego generują pełną grupę symetryczną
jego wierzchołków. Istotnie, dla krawędzi drzewa
rozpinającego można, przez kolejne sprzęganie
transpozycji wzdłuż ścieżki, otrzymać transpozycję
dowolnych dwóch wierzchołków. Każda transpozycja
należy zatem do $G$. Wynika stąd $G=S_5$.

\emph{Krok 4: brak rozwiązywalności.}
Grupa $A_5$ jest nieabelową grupą prostą
(Twierdzenie~\ref{aa:tw:a5-simple}), zatem nie jest
rozwiązalna (Wniosek~\ref{aa:wn:a5-not-solvable}).
Ponieważ $A_5\le S_5$, grupa $S_5$ również nie może
być rozwiązalna: podgrupa grupy rozwiązalnej
zawsze jest rozwiązalna. Kryterium
\ref{aa:tw:radicals-solvable-galois} pokazuje,
że $f$ nie jest rozwiązywalny przez pierwiastniki.
\end{proof}

\begin{wniosek}[Twierdzenie Abela--Ruffiniego]
\label{aa:wn:abel-ruffini}
Nie istnieje ogólna metoda wyrażenia wszystkich pierwiastków
każdego wielomianu piątego stopnia nad $\mathbb Q$
za pomocą skończonego ciągu działań arytmetycznych
i wyciągania pierwiastków, zaczynając od jego współczynników.
Nie oznacza to, że \emph{żaden} wielomian piątego stopnia
nie jest rozwiązywalny przez pierwiastniki.
\end{wniosek}
\begin{proof}
Gdyby taka metoda istniała, stosowałaby się również
do wielomianu $X^5-4X+2$. Byłby on wówczas
rozwiązywalny przez pierwiastniki nad $\mathbb Q$,
w sprzeczności z
Twierdzeniem~\ref{aa:tw:quintic-s5-example}.
Z drugiej strony wielomian $X^5-2$ jest rozwiązywalny
przez pierwiastniki: po dołączeniu $\sqrt[5]2$ oraz
pierwotnego pierwiastka piątego stopnia z jedynki
otrzymujemy wszystkie jego pierwiastki.
\end{proof}

\begin{uwaga}[Granice i zastosowania kryterium]
\label{aa:uw:radicals-next}
W charakterystyce dodatniej wymagane są osobne
ostrożności dotyczące rozdzielczości. Nasz dowód
kryterium rozwiązywalności przez pierwiastniki zakłada
charakterystykę zero. W przyszłych rozdziałach można
rozwinąć teorię rozszerzeń cyklicznych, 90. twierdzenie Hilberta,
rozwiązywalność geometrycznych konstrukcji oraz teorię Galois
wielomianów związanych z działaniem grup na zbiorach pierwiastków.
\end{uwaga}
% AA-25 | Iteracja XIII; szczegolowy audyt historyczny 2026-11-20, partia 02.
\chapter{90. twierdzenie Hilberta i równania normowe}
\label{aa:ch:hilbert90}

\section{Motywacja: kiedy element jest ilorazem swoich sprzężeń?}
Niech $L/K$ będzie skończonym rozszerzeniem Galois, a $G=\operatorname{Gal}(L/K)$
jego grupą Galois. Jeżeli $b\in L^\times$ i $\sigma\in G$, to element
$a=b/\sigma(b)$ ma szczególną własność: gdy $G=\langle\sigma\rangle$, jego
norma z $L$ do $K$ wynosi $1$. W istocie kolejne czynniki
$a,\sigma(a),\ldots,\sigma^{n-1}(a)$ znoszą się po wymnożeniu.
Pojawia się pytanie odwrotne: czy warunek $N_{L/K}(a)=1$ wystarcza,
aby znaleźć takie $b$? Odpowiedź daje \emph{90. twierdzenie Hilberta}.
Wykażemy też wersję addytywną, w której normę zastępuje ślad,
a iloraz $b/\sigma(b)$ różnica $b-\sigma(b)$.

Normę $N_{L/K}$ i ślad $\operatorname{Tr}_{L/K}$ definiowaliśmy
w rozdziale~\ref{aa:ch:trace-norm}. Przypomnijmy, że są one określone
za pomocą wyznacznika i śladu $K$-liniowego odwzorowania
$L\ni x\mapsto ax\in L$; nie zakładamy przy tym, że
$\operatorname{char}K=0$. Dla automorfizmu $\sigma$ symbol $\sigma^j$
oznacza $j$-krotne złożenie, a $\sigma^0=\operatorname{id}_L$.

\section{Niezależność automorfizmów i przygotowanie}
Twierdzenie multiplikatywne będzie wymagało znalezienia elementu,
na którym pewna suma automorfizmów nie znika. Nie można tego
uzasadnić samym faktem, że każdy jej składnik jest niezerowy:
suma niezerowych odwzorowań bywa zerowa. Potrzebujemy silniejszego lematu.

\begin{lemat}[Niezależność różnych homomorfizmów ciał]\label{aa:lem:dedekind-independent}
Niech $F,E$ będą ciałami, $r\geq1$ liczbą całkowitą,
a $\sigma_1,\ldots,\sigma_r:F\to E$ parami różnymi homomorfizmami ciał.
Dla dowolnych $c_1,\ldots,c_r\in E$, jeżeli
$\sum_{i=1}^r c_i\sigma_i(x)=0$ dla każdego $x\in F$,
to $c_1=\cdots=c_r=0$.
\end{lemat}
\begin{proof}
Załóżmy wbrew tezie, że istnieje niezerowy układ współczynników
spełniający wskazaną zależność. Spośród wszystkich takich zależności
wybierzmy tę o najmniejszej liczbie niezerowych współczynników.
Po usunięciu składników zerowych i zmianie numeracji możemy zapisać ją jako
\[
\sum_{i=1}^{s}c_i\sigma_i(x)=0\qquad(x\in F),
\quad c_i\neq0\ (1\leq i\leq s).
\]
Musi być $s\geq2$: dla $s=1$ podstawienie $x=1$ daje
$c_1\sigma_1(1)=c_1=0$, sprzecznie z wyborem.
Homomorfizmy $\sigma_1$ i $\sigma_s$ są różne, zatem
istnieje $u\in F$ spełniające $\sigma_1(u)\neq\sigma_s(u)$.

Dla każdego $x\in F$ zastosujmy początkową zależność do $ux$.
Ponieważ każdy $\sigma_i$ zachowuje mnożenie, otrzymujemy
$\sum_{i=1}^{s}c_i\sigma_i(u)\sigma_i(x)=0$.
Pomnóżmy osobno początkową zależność dla $x$ przez
$\sigma_s(u)\in E$ i odejmijmy ją od otrzymanej równości.
Dostajemy dla wszystkich $x\in F$ zależność
\[
\sum_{i=1}^{s}c_i\bigl(\sigma_i(u)-\sigma_s(u)\bigr)\sigma_i(x)=0.
\]
Współczynnik przy $\sigma_s$ jest zerowy, natomiast przy
$\sigma_1$ wynosi $c_1(\sigma_1(u)-\sigma_s(u))\neq0$,
ponieważ $E$ jest ciałem. Powstała zatem niezerowa zależność
z co najwyżej $s-1$ niezerowymi współczynnikami. Przeczy to minimalności
$s$ i kończy dowód.
\end{proof}

\begin{uwaga}[Zakres lematu]
Homomorfizmy $\sigma_i$ nie muszą być surjektywne ani automorfizmami
tego samego ciała. W dalszych zastosowaniach przyjmiemy $F=E=L$
i wybierzemy $\sigma_i$ spośród elementów grupy Galois.
Wtedy suma $\sum_i c_i\sigma_i$ jest odwzorowaniem $L\to L$,
a lemat stwierdza, że nie może być identycznie zerowa,
jeżeli co najmniej jeden współczynnik $c_i$ jest niezerowy.
\end{uwaga}

\begin{lemat}[Norma i ślad w rozszerzeniu Galois]\label{aa:lem:galois-norm-product}
Niech $L/K$ będzie skończonym rozszerzeniem Galois oraz
$G=\operatorname{Gal}(L/K)$. Wówczas dla każdego $a\in L$ zachodzą
\[
N_{L/K}(a)=\prod_{\tau\in G}\tau(a),\qquad
\operatorname{Tr}_{L/K}(a)=\sum_{\tau\in G}\tau(a).
\]
\end{lemat}
\begin{proof}
Niech $n=[L:K]$. Ponieważ rozszerzenie $L/K$ jest rozdzielcze,
ma dokładnie $n$ różnych $K$-zanurzeń w ustalone domknięcie
algebraiczne $\overline K$. Normalność oznacza, że obraz każdego
z tych zanurzeń jest zawarty w $L$. Każde takie zanurzenie jest
$K$-liniowe i iniektywne; jako iniektywne odwzorowanie skończenie
wymiarowej przestrzeni $L$ w siebie jest również surjektywne,
a zatem jest elementem $G$.
Stąd $|G|=n$. Wzory na normę i ślad jako iloczyn i sumę
obrazów elementu przy wszystkich $K$-zanurzeniach,
udowodnione w rozdziale~\ref{aa:ch:trace-norm},
dają obie żądane równości. Wzór iloczynowy obejmuje również $a=0$.
\end{proof}

\section{Multiplikatywna postać 90. twierdzenia Hilberta}
\begin{twierdzenie}[90. twierdzenie Hilberta — postać multiplikatywna]\label{aa:tw:hilbert90-multiplicative}
Niech $L/K$ będzie skończonym cyklicznym rozszerzeniem Galois,
$G=\langle\sigma\rangle$ grupą rzędu $n\geq1$, a $a\in L^\times$.
Następujące warunki są równoważne:
\begin{enumerate}[label=\textup{(\roman*)}]
\item $N_{L/K}(a)=1$;
\item istnieje $b\in L^\times$ takie, że $a=b/\sigma(b)$.
\end{enumerate}
\end{twierdzenie}
\begin{proof}
\emph{Kierunek (ii)$\Rightarrow$(i).}
Niech $a=b/\sigma(b)$, gdzie $b\neq0$. Każdy automorfizm
ciała zachowuje niezerowość, więc wszystkie mianowniki
$\sigma^j(b)$ są niezerowe. Korzystając z lematu
\ref{aa:lem:galois-norm-product} oraz $\sigma^n=\operatorname{id}_L$,
otrzymujemy
\[
N_{L/K}(a)=\prod_{j=0}^{n-1}\frac{\sigma^j(b)}{\sigma^{j+1}(b)}
=\frac{b}{\sigma^n(b)}=1.
\]

\emph{Kierunek (i)$\Rightarrow$(ii).}
Załóżmy, że $N_{L/K}(a)=1$. Zdefiniujmy
$ p_0=1$ oraz, dla $1\leq j\leq n$,
$ p_j=\prod_{i=0}^{j-1}\sigma^i(a)$.
Ponieważ $a\neq0$, wszystkie $p_j$ są niezerowe;
ponadto z założenia i lematu o normie wynika $p_n=1$.
Dla $0\leq j<n$ obliczamy, osobno wypisując przesunięcie indeksu,
\[
\sigma(p_j)=\prod_{i=0}^{j-1}\sigma^{i+1}(a)
=\frac{p_{j+1}}{a};
\]
przy $j=0$ obie strony są równe $1$. Definiujemy odwzorowanie
$B:L\to L$ wzorem
\[
B(x)=\sum_{j=0}^{n-1}p_j\sigma^j(x),\qquad x\in L.
\]
Automorfizmy $\operatorname{id},\sigma,\ldots,\sigma^{n-1}$
są parami różne. Z lematu~\ref{aa:lem:dedekind-independent}
i faktu, że $p_0=1\neq0$, wynika, że $B$ nie jest odwzorowaniem
identycznie zerowym. Istnieje więc $x\in L$ takie, że
$b:=B(x)\neq0$; tutaj wskazaliśmy zarówno sposób zdefiniowania
$b$, jak i przyczynę jego niezerowości.

Zastosujmy $\sigma$ do każdego składnika $b$ i skorzystajmy
z tożsamości dotyczącej $\sigma(p_j)$:
\[
\begin{aligned}
a\sigma(b)
&=a\sum_{j=0}^{n-1}\sigma(p_j)\sigma^{j+1}(x)
 =\sum_{j=0}^{n-1}p_{j+1}\sigma^{j+1}(x)\\
&=\sum_{j=1}^{n-1}p_j\sigma^j(x)
   +p_n\sigma^n(x)
 =\sum_{j=1}^{n-1}p_j\sigma^j(x)+x=B(x)=b.
\end{aligned}
\]
Ostatnie przejście używa $p_n=1$, $\sigma^n(x)=x$ i $p_0=1$.
Ponieważ $b\neq0$, także $\sigma(b)\neq0$, więc wolno podzielić
przez $\sigma(b)$. Dostajemy $a=b/\sigma(b)$.
Argument obejmuje również $n=1$: wtedy $G$ jest trywialna,
$N_{L/K}(a)=a=1$, a można wybrać $b=1$.
\end{proof}

\begin{przyklad}[Norma zespolona]\label{aa:prz:hilbert90-complex}
Rozpatrujemy rozszerzenie $\mathbb C/\mathbb R$ stopnia $2$;
jego grupę Galois generuje sprzężenie $\sigma(z)=\overline z$.
Dla $z=x+iy$, $x,y\in\mathbb R$, mamy
$N_{\mathbb C/\mathbb R}(z)=(x+iy)(x-iy)=x^2+y^2=|z|^2$.
Warunek normy $1$ jest zatem równoważny $|a|=1$.
Jeśli $a\neq-1$, wybierzmy $b=1+a$. Jest on niezerowy, a ponieważ
$a\overline a=1$, zachodzi $\overline a=a^{-1}$ oraz
\[
\frac{b}{\overline b}
=\frac{1+a}{1+\overline a}
=\frac{1+a}{1+a^{-1}}=a.
\]
Dla $a=-1$ wyrażenie $1+a$ znika, więc trzeba zmienić wybór:
$b=i$ daje $b/\overline b=i/(-i)=-1$.
Przykład pokazuje również różnicę terminologiczną: norma
rozszerzenia $\mathbb C/\mathbb R$ nie jest modułem $z$, lecz
jego kwadratem. Dla $z=2+3i$ norma wynosi $13$, a moduł $\sqrt{13}$.
\end{przyklad}

\begin{przyklad}[Norma w rozszerzeniu kwadratowym]\label{aa:prz:hilbert90-quadratic}
Niech $K$ będzie ciałem charakterystyki różnej od $2$, a
$d\in K^\times$ nie będzie kwadratem w $K$. Wtedy
$L=K(\sqrt d)$ ma bazę $(1,\sqrt d)$ nad $K$, a
$\sigma(u+v\sqrt d)=u-v\sqrt d$. Dla $u,v\in K$ obliczamy
\[
N_{L/K}(u+v\sqrt d)=(u+v\sqrt d)(u-v\sqrt d)=u^2-dv^2.
\]
Jeśli $N(a)=1$ i $a\neq-1$, można \emph{jawnie} wziąć $b=1+a$,
ponieważ $\sigma(a)=a^{-1}$ i ten sam rachunek co w przykładzie
zespolonym daje $b/\sigma(b)=a$. Jeśli $a=-1$, wybieramy
$b=\sqrt d$, dla którego $\sigma(b)=-b$.

Na przykład dla $K=\mathbb Q$, $d=2$ i $a=3+2\sqrt2$ mamy
$N(a)=9-8=1$. Bierzemy $b=1+a=4+2\sqrt2$ i sprawdzamy:
\[
\frac{4+2\sqrt2}{4-2\sqrt2}
=\frac{(4+2\sqrt2)^2}{16-8}
=\frac{24+16\sqrt2}{8}=3+2\sqrt2.
\]
Ważne jest rozróżnienie dwóch pytań: 90. twierdzenie Hilberta
opisuje \emph{jądro} normy, natomiast równanie $N(b)=c$
pyta, czy dane $c\in K^\times$ należy do \emph{obrazu} normy.
Nie wynika stąd, że dowolne $c$ jest normą.
Dla $\mathbb C/\mathbb R$ mamy
$N_{\mathbb C/\mathbb R}(\mathbb C^\times)=\mathbb R_{>0}$,
ponieważ $x^2+y^2>0$ dla $(x,y)\neq(0,0)$,
a każdą dodatnią liczbę $c$ otrzymujemy jako normę $\sqrt c$.
W szczególności $-1$ nie jest taką normą.
\end{przyklad}

\begin{przyklad}[Multiplikatywne równanie w ciele czteroelementowym]\label{aa:prz:hilbert90-f4}
Niech $K=\mathbb F_2$, a $L=\mathbb F_2[\theta]$, gdzie
$\theta^2+\theta+1=0$. Wielomian $X^2+X+1$ nie ma
pierwiastka w $K$, zatem $L/K$ jest rozszerzeniem Galois
stopnia $2$. Jego grupa jest generowana przez automorfizm
Frobeniusa $\sigma(x)=x^2$. Ponieważ $\theta^2=\theta+1$,
mamy $\theta^3=\theta(\theta+1)=\theta^2+\theta=1$.
Grupa $L^\times$ składa się dokładnie z trzech elementów
$1,\theta,\theta^2$. Dla każdego $a\in L^\times$ zachodzi
$N_{L/K}(a)=a\sigma(a)=a^3=1$.

Sprawdźmy konstrukcję z 90. twierdzenia Hilberta dla
$a=\theta^2$. Wybieramy $b=\theta$. Wówczas
$\sigma(b)=\theta^2$ i
\[
\frac{b}{\sigma(b)}=\frac{\theta}{\theta^2}
=\theta^{-1}=\theta^2=a.
\]
W tym przykładzie jądro normy jest całą grupą $L^\times$.
Nie oznacza to, że norma każdego rozszerzenia ciał jest
identycznie równa $1$: zależy to od konkretnego rozszerzenia.
\end{przyklad}

\section{Addytywna postać 90. twierdzenia Hilberta}
\begin{twierdzenie}[90. twierdzenie Hilberta — postać addytywna]\label{aa:tw:hilbert90-additive}
Niech $L/K$ będzie skończonym cyklicznym rozszerzeniem Galois
stopnia $n\geq1$, a $\sigma$ generatorem $\operatorname{Gal}(L/K)$.
Dla każdego $a\in L$ następujące warunki są równoważne:
\begin{enumerate}[label=\textup{(\roman*)}]
\item $\operatorname{Tr}_{L/K}(a)=0$;
\item istnieje $b\in L$ takie, że $a=b-\sigma(b)$.
\end{enumerate}
\end{twierdzenie}
\begin{proof}
\emph{Kierunek (ii)$\Rightarrow$(i).}
Jeśli $a=b-\sigma(b)$, to z lematu
\ref{aa:lem:galois-norm-product} wynika
\[
\operatorname{Tr}_{L/K}(a)
=\sum_{j=0}^{n-1}(\sigma^j(b)-\sigma^{j+1}(b))
=b-\sigma^n(b)=0.
\]

\emph{Kierunek (i)$\Rightarrow$(ii).}
Rozpatrzmy $K$-liniowe odwzorowanie
$T:L\to L$, $T(b)=b-\sigma(b)$.
Jego jądro jest zbiorem elementów ustalonych przez $\sigma$,
czyli $L^{\langle\sigma\rangle}=K$ na mocy odpowiedniości Galois
(zob. twierdzenie~\ref{aa:tw:fundamental-galois}).
Ponieważ $[L:K]=n$, twierdzenie o rzędzie i jądrze daje
$\dim_K\operatorname{im}T=n-1$.

Ślad $\operatorname{Tr}_{L/K}:L\to K$ jest funkcjonałem
$K$-liniowym i jest niezerowy. Uzasadnijmy ten drugi fakt bez
założenia, że $\operatorname{char}K$ nie dzieli $n$:
niezdegenerowanie formy $L\times L\to K$,
$(x,y)\mapsto\operatorname{Tr}_{L/K}(xy)$,
wynikające z rozdziału~\ref{aa:ch:trace-norm}
(twierdzenie~\ref{aa:tw:trace-nondegenerate}), oznacza, że
istnieje $y\in L$ z $\operatorname{Tr}_{L/K}(1\cdot y)\neq0$.
Dlatego $\ker\operatorname{Tr}_{L/K}$ ma wymiar $n-1$.
Pierwszy kierunek dowodu wykazał już
$\operatorname{im}T\subseteq\ker\operatorname{Tr}_{L/K}$.
Obie są podprzestrzeniami tej samej przestrzeni $L$
o jednakowym skończonym wymiarze, a więc są równe.
Dla każdego $a$ o śladzie zerowym istnieje zatem $b\in L$
spełniające $T(b)=a$, zgodnie z tezą.
Dla $n=1$ obie podprzestrzenie są zerowe i dowód nadal działa.
\end{proof}

\begin{przyklad}[Addytywne równanie w charakterystyce dwa]
Niech $K=\mathbb F_2$ i
$L=\mathbb F_2[\theta]$, gdzie $\theta^2+\theta+1=0$.
Wielomian $X^2+X+1$ nie ma pierwiastka w $\mathbb F_2$,
więc $[L:K]=2$ i $L$ jest ciałem czteroelementowym.
Automorfizm $\sigma(x)=x^2$ ma rząd $2$ i ustala $K$;
w szczególności $\sigma(\theta)=\theta^2=\theta+1$.
Mamy $\operatorname{Tr}_{L/K}(1)=1+1=0$ oraz
$\operatorname{Tr}_{L/K}(\theta)=\theta+(\theta+1)=1$.
Stąd $\ker\operatorname{Tr}_{L/K}=\{0,1\}$, a
$\theta-\sigma(\theta)=\theta-(\theta+1)=1$,
ponieważ w charakterystyce $2$ zachodzi $-1=1$.
Zatem równanie $1=b-\sigma(b)$ ma rozwiązanie $b=\theta$,
mimo że $\operatorname{Tr}_{L/K}(1)=0$.
Przykład pokazuje, dlaczego w dowodzie nie wolno uzasadniać
niezerowości śladu wyłącznie równością
$\operatorname{Tr}_{L/K}(1)=[L:K]$.
\end{przyklad}

\begin{uwaga}[Dlaczego potrzebna jest cykliczność?]\label{aa:uw:hilbert90-cyclic}
Jeżeli skończone rozszerzenie Galois $L/K$ ma grupę
$G$, która nie jest cykliczna, żaden pojedynczy element
$\sigma\in G$ nie generuje całej grupy. Wtedy ciało stałe
$L^{\langle\sigma\rangle}$ jest większe od $K$.
Ze wzoru na wymiar obrazu operatora $1-\sigma$ wynika
$\dim_K\operatorname{im}(1-\sigma)<[L:K]-1$,
podczas gdy jądro niezerowego funkcjonału śladu ma wymiar
$[L:K]-1$. Nie można więc zamienić grupy cyklicznej
w twierdzeniu addytywnym na dowolną grupę Galois,
zachowując jedną ustaloną różnicę $b-\sigma(b)$.
\end{uwaga}

\begin{przyklad}[Kontrprzykład dla jednej różnicy w grupie niecyklicznej]\label{aa:prz:hilbert90-noncyclic}
Weźmy $L=\mathbb Q(\sqrt2,\sqrt3)$ oraz $K=\mathbb Q$.
Elementy $1,\sqrt2,\sqrt3,\sqrt6$ tworzą bazę $L$ nad $K$,
a grupa Galois jest izomorficzna z $C_2\times C_2$.
Niech $\sigma$ zmienia znak $\sqrt2$ i ustala $\sqrt3$.
Dla dowolnego
$x=x_0+x_1\sqrt2+x_2\sqrt3+x_3\sqrt6$,
$x_j\in\mathbb Q$, mamy
\[
(1-\sigma)(x)=2x_1\sqrt2+2x_3\sqrt6.
\]
Stąd $\operatorname{im}(1-\sigma)$ jest dwuwymiarową
podprzestrzenią rozpiętą przez $\sqrt2,\sqrt6$.
Z kolei suma czterech sprzężeń elementu $x$ wynosi $4x_0$,
więc $\ker\operatorname{Tr}_{L/K}$ jest trójwymiarową
podprzestrzenią rozpiętą przez $\sqrt2,\sqrt3,\sqrt6$.
W szczególności $a=\sqrt3$ ma ślad równy $0$,
ale nie istnieje $b\in L$ z $a=b-\sigma(b)$:
we wspomnianej bazie każda taka różnica ma współczynnik
przy $\sqrt3$ równy $0$.
Ten kontrprzykład nie przeczy twierdzeniu
\ref{aa:tw:hilbert90-additive}, ponieważ jego założenie
cykliczności grupy Galois nie jest spełnione.
\end{przyklad}
% AA-26 | Iteracja XIII; szczegolowy audyt historyczny 2026-11-20, partia 02.
\chapter[Rozszerzenia cykliczne i teoria Kummera]{Rozszerzenia cykliczne, pierwiastki z jedynki i teoria Kummera}
\label{aa:ch:kummer}

\section{Motywacja: od pierwiastnika do grupy Galois}
Niech $F$ będzie ciałem, $n\geq2$ liczbą całkowitą, a $a\in F^\times$.
Po dołączeniu pierwiastka $\alpha$ równania $\alpha^n=a$
otrzymujemy rozszerzenie $E=F(\alpha)$. Chcielibyśmy odczytać
jego grupę Galois z prostego równania $X^n-a=0$.
Dwa zjawiska mogą jednak temu przeszkadzać: inne pierwiastki
tego wielomianu mogą nie należeć do $E$, a w dodatniej
charakterystyce wielomian może mieć pierwiastki wielokrotne.
Teoria Kummera określa warunki, w których oba problemy znikają,
i wyjaśnia, dlaczego powstają rozszerzenia cykliczne.

Dla ciała $F$ symbol $F^\times$ oznacza jego grupę
multiplikatywną elementów niezerowych. Przez
$\operatorname{char}F\nmid n$ rozumiemy, że obraz liczby całkowitej
$n$ w $F$ jest niezerowy, a więc odwracalny.
Wszystkie pierwiastki wielomianów rozpatrujemy w ustalonym
domknięciu algebraicznym $\overline F$.

\section{Pierwiastki z jedynki i rozszerzenia radykalne}
\begin{definicja}[Grupa pierwiastków z jedynki]\label{aa:def:mu-n}
Niech $F$ będzie ciałem i $n\geq1$ liczbą całkowitą. Określamy
\[
\mu_n(F)=\{u\in F^\times:u^n=1\}.
\]
Jest to podgrupa $F^\times$: $1\in\mu_n(F)$, iloczyn
dwóch elementów spełniających równanie $u^n=1$
spełnia je ponownie, a $u^{-1}$ także jest jego rozwiązaniem.
Element $\zeta\in\mu_n(F)$ nazywamy
\emph{pierwotnym $n$-tym pierwiastkiem z jedynki},
jeżeli jego rząd w $F^\times$ jest dokładnie równy $n$.
\end{definicja}

Najpierw uzasadnimy własność, którą często używa się bez dowodu:
skończona podgrupa grupy multiplikatywnej ciała jest cykliczna.

\begin{lemat}[Skończone podgrupy $F^\times$]
Dla dowolnego ciała $F$ każda skończona podgrupa $H\leq F^\times$
jest cykliczna.
\end{lemat}
\begin{proof}
Grupa $H$ jest abelowa, ponieważ mnożenie w ciele jest przemienne.
Jeśli $H=\{1\}$, teza jest spełniona. Załóżmy, że $|H|>1$
i oznaczmy przez $m$ najmniejszą wspólną wielokrotność
rzędów wszystkich elementów $H$. Jest ona określona,
ponieważ $H$ jest skończona.

Niech $p^e$ będzie dowolną maksymalną potęgą liczby pierwszej
$p$ dzielącą $m$. W definicji najmniejszej wspólnej wielokrotności
musi występować element $h_p\in H$, którego rząd dzieli się przez
$p^e$. Jeśli $\operatorname{ord}(h_p)=p^e q$ z $p\nmid q$,
to element $g_p=h_p^q$ ma rząd dokładnie $p^e$:
$(h_p^q)^{p^e}=1$, a żadna mniejsza dodatnia potęga $p$
nie wystarcza do uzyskania $1$.

Dla różnych dzielników pierwszych $p$ liczby $m$ elementy
$g_p$ komutują, a ich rzędy są parami względnie pierwsze.
W grupie abelowej iloczyn dwóch elementów o względnie
pierwszych rzędach $r,s$ ma rząd $rs$: jeśli
$(uv)^k=1$, to $u^k=v^{-k}$ należy jednocześnie
do podgrup $\langle u\rangle$ i $\langle v\rangle$; ich
przecięcie jest trywialne, bo rząd elementu przecięcia
dzieli jednocześnie $r$ i $s$. Stąd $r\mid k$ i $s\mid k$,
więc $rs\mid k$. Powtarzając ten argument, otrzymujemy,
że iloczyn $g=\prod_{p\mid m}g_p$ ma rząd
$\prod_{p\mid m}p^{v_p(m)}=m$.
W szczególności $H$ zawiera $m$ różnych potęg $g$.
Z drugiej strony rząd każdego elementu $h\in H$ dzieli $m$,
więc $h^m=1$. Wszystkie elementy $H$ są zatem pierwiastkami
wielomianu $X^m-1\in F[X]$. Ponieważ niezerowy wielomian
stopnia $m$ nad ciałem ma najwyżej $m$ różnych pierwiastków,
$|H|\leq m$. Istnienie elementu $g$ rzędu $m$ daje
$|H|\geq m$. Obie nierówności oznaczają
$H=\langle g\rangle$.
\end{proof}

\begin{wniosek}[Kiedy istnieje pierwotny pierwiastek?]
Jeśli $\operatorname{char}F\nmid n$ oraz wielomian $X^n-1$
ma w $F$ wszystkie swoje $n$ pierwiastków, to
$\mu_n(F)$ jest grupą cykliczną rzędu $n$.
\end{wniosek}
\begin{proof}
Pochodna $X^n-1$ jest równa $nX^{n-1}$.
Przy założeniu $\operatorname{char}F\nmid n$ wspólny pierwiastek
$X^n-1$ i $nX^{n-1}$ musiałby jednocześnie być zerem
$X$, co jest niemożliwe dla pierwiastka $X^n-1$.
Zatem wielomian ma $n$ różnych pierwiastków w $\overline F$.
Jeśli wszystkie należą do $F$, to $|\mu_n(F)|=n$.
Z poprzedniego lematu wynika cykliczność tej grupy,
a jej generator jest pierwotnym $n$-tym pierwiastkiem z jedynki.
\end{proof}

\begin{twierdzenie}[Kryterium cykliczności prostego rozszerzenia radykalnego]\label{aa:tw:kummer-radical-cyclic}
Niech $F$ będzie ciałem, $n\geq2$ liczbą całkowitą
odwracalną w $F$ i niech $\mu_n(F)$ ma rząd $n$.
Dla każdego $a\in F^\times$ i każdego $\alpha\in\overline F$
takiego, że $\alpha^n=a$, rozszerzenie $E=F(\alpha)/F$
jest skończonym rozszerzeniem Galois o grupie cyklicznej,
której rząd dzieli $n$.
\end{twierdzenie}
\begin{proof}
\emph{Rozkład wielomianu.}
Ponieważ $a\neq0$, także $\alpha\neq0$. Jeśli $\beta^n=a$,
to $(\beta/\alpha)^n=1$, zatem $\beta=\zeta\alpha$
dla pewnego $\zeta\in\mu_n(\overline F)$.
Na mocy założeń wszystkie $n$-te pierwiastki z jedynki
leżą w $F$, więc każdy pierwiastek $X^n-a$ należy do $E$.
Oznacza to, że $E$ jest ciałem rozkładu tego wielomianu nad $F$.

\emph{Rozdzielczość.}
Pochodna $X^n-a$ jest równa $nX^{n-1}$.
Ponieważ $n\neq0$ w $F$, wspólny pierwiastek obu wielomianów
musiałby być zerem $X$, ale $0^n-a=-a\neq0$.
Wszystkie pierwiastki są więc jednokrotne. Ponieważ $E$
jest ciałem rozkładu wielomianu rozdzielczego,
$E/F$ jest skończonym rozszerzeniem Galois.

\emph{Opis grupy.}
Niech $G=\operatorname{Gal}(E/F)$. Dla każdego $\tau\in G$
mamy $\tau(\alpha)^n=\tau(a)=a$, więc
$\chi(\tau):=\tau(\alpha)/\alpha\in\mu_n(F)$.
Sprawdźmy, że odwzorowanie $\chi:G\to\mu_n(F)$ jest
homomorfizmem, a nie tylko dowolną funkcją.
Jeśli $\tau,\rho\in G$, to $\rho(\alpha)=\chi(\rho)\alpha$,
a $\tau$ ustala element $\chi(\rho)\in F$. Stąd
\[
\chi(\tau\rho)=\frac{\tau(\rho(\alpha))}{\alpha}
=\frac{\chi(\rho)\tau(\alpha)}{\alpha}
=\chi(\rho)\chi(\tau)=\chi(\tau)\chi(\rho),
\]
gdzie ostatnia równość wykorzystuje przemienność grupy
$\mu_n(F)$. Jeśli $\chi(\tau)=1$, to $\tau(\alpha)=\alpha$;
automorfizm $\tau$ ustala też $F$, a zatem każdy element
$F(\alpha)=E$. Wobec tego $\tau=\operatorname{id}_E$,
więc $\chi$ jest iniektywne.
Grupa $G$ jest zatem izomorficzna z podgrupą cyklicznej grupy
$\mu_n(F)$ rzędu $n$. Każda podgrupa grupy cyklicznej
jest cykliczna, a na mocy twierdzenia Lagrange'a jej rząd
dzieli $n$.
\end{proof}

\begin{uwaga}[Dlaczego sam warunek „nie jest $n$-tą potęgą” nie wystarcza?]
Dla złożonego $n$ wielomian $X^n-a$ może być rozkładalny,
mimo że $a$ nie jest $n$-tą potęgą w $F$.
Poprzednie twierdzenie gwarantuje wtedy jedynie, że stopień
rozszerzenia \emph{dzieli} $n$, nie zaś, że jest równy $n$.
Przykład o stopniu $2$ przy $n=4$ podamy niżej.
\end{uwaga}

\begin{twierdzenie}[Odwrotne twierdzenie Kummera dla rozszerzeń cyklicznych]\label{aa:tw:kummer-converse}
Niech $F$ będzie ciałem, $n\geq2$ liczbą całkowitą
odwracalną w $F$, a $\zeta\in F$ pierwotnym
$n$-tym pierwiastkiem z jedynki. Jeśli $E/F$ jest
skończonym rozszerzeniem Galois, a $\operatorname{Gal}(E/F)$
jest grupą cykliczną rzędu $n$, to istnieją
$\alpha\in E^\times$ oraz generator $\sigma$ tej grupy takie, że
\[
E=F(\alpha),\qquad \alpha^n\in F,\qquad
\sigma(\alpha)=\zeta\alpha.
\]
\end{twierdzenie}
\begin{proof}
Wybierzmy dowolny generator $\sigma$ cyklicznej grupy
$G=\operatorname{Gal}(E/F)$. Rozszerzenie jest Galois,
więc $[E:F]=|G|=n$ i $E^G=F$.
Dla każdego $x\in E$ zdefiniujmy
\[
P(x)=\sum_{j=0}^{n-1}\zeta^{-j}\sigma^j(x).
\]
Jest to dobrze określone odwzorowanie $E\to E$, ponieważ
każdy składnik należy do $E$, a $\zeta$ jest niezerowe.
Automorfizmy $\sigma^0,\ldots,\sigma^{n-1}$ są parami różne.
Z lematu o niezależności homomorfizmów
\ref{aa:lem:dedekind-independent} oraz faktu, że pierwszy
współczynnik $\zeta^0=1$, wynika, że $P$ nie może być
identycznie zerowe. Wybierzmy $x\in E$ takie, że
$\alpha=P(x)\neq0$.

Aby sprawdzić, jak działa $\sigma$ na $\alpha$, przesuwamy
indeks w skończonej sumie. Ponieważ $\zeta\in F$ i $\sigma$
ustala $F$, mamy
\[
\begin{aligned}
\sigma(\alpha)
 &=\sum_{j=0}^{n-1}\zeta^{-j}\sigma^{j+1}(x)\\
 &=\zeta\sum_{k=1}^{n}\zeta^{-k}\sigma^k(x)
 =\zeta\left(\sum_{k=1}^{n-1}\zeta^{-k}\sigma^k(x)
       +\zeta^{-n}\sigma^n(x)\right)\\
 &=\zeta\left(\sum_{k=1}^{n-1}\zeta^{-k}\sigma^k(x)+x\right)
 =\zeta\alpha.
\end{aligned}
\]
W przedostatniej równości użyto $\zeta^n=1$
i $\sigma^n=\operatorname{id}_E$.
Z tej równości wynika
$\sigma(\alpha^n)=\sigma(\alpha)^n
=(\zeta\alpha)^n=\alpha^n$.
Ponieważ $\sigma$ generuje $G$, element $\alpha^n$
jest ustalony przez całe $G$, a więc $\alpha^n\in E^G=F$.

Pozostaje uzasadnić, że $\alpha$ generuje całe rozszerzenie.
Dla $0\leq j<n$ zachodzi
$\sigma^j(\alpha)=\zeta^j\alpha$.
Wartości te są parami różne, gdyż $\zeta$ ma rząd $n$
i $\alpha\neq0$.
Każdy automorfizm $F$-liniowy przenosi pierwiastek wielomianu
minimalnego $m_{\alpha,F}$ na inny jego pierwiastek.
Zatem $m_{\alpha,F}$ ma co najmniej $n$ różnych pierwiastków
w $E$, a więc $\deg m_{\alpha,F}\geq n$.
Z drugiej strony $F(\alpha)\subseteq E$, więc
$[F(\alpha):F]\leq[E:F]=n$.
Ponieważ $[F(\alpha):F]=\deg m_{\alpha,F}$,
obie nierówności dają $[F(\alpha):F]=n$ i $F(\alpha)=E$.
\end{proof}

\begin{przyklad}[Rozszerzenie kwadratowe]\label{aa:prz:kummer-square}
Dla $F=\mathbb Q$, $n=2$ oraz $a=2$ mamy
$\mu_2(F)=\{1,-1\}$, a liczba $2$ nie jest kwadratem
w $\mathbb Q$. Wielomian $X^2-2$ nie ma zatem pierwiastka
w $\mathbb Q$, jest nierozkładalny i
$E=\mathbb Q(\sqrt2)$ ma stopień $2$ nad $\mathbb Q$.
Jego oba pierwiastki $\sqrt2$ i $-\sqrt2$ należą do $E$;
automorfizmy Galois to identyczność oraz
$\sigma(u+v\sqrt2)=u-v\sqrt2$ ($u,v\in\mathbb Q$).
Grupa jest cykliczna rzędu $2$, a $\sigma(\sqrt2)=-\sqrt2$.
Jest to przypadek teorii Kummera z pierwotnym pierwiastkiem
$\zeta=-1$.
\end{przyklad}

\begin{przyklad}[Dlaczego potrzebujemy pierwiastków z jedynki?]\label{aa:prz:kummer-missing-roots}
Niech $\alpha=\sqrt[3]{2}$ będzie dodatnim rzeczywistym
pierwiastkiem wielomianu $X^3-2$, a
$\omega\in\mathbb C$ pierwotnym pierwiastkiem trzeciego stopnia
z jedynki, zatem $\omega^2+\omega+1=0$.
Wielomian $X^3-2$ jest nierozkładalny nad $\mathbb Q$
na mocy kryterium Eisensteina dla liczby pierwszej $2$,
więc $[\mathbb Q(\alpha):\mathbb Q]=3$.
Ciało $\mathbb Q(\alpha)$ zawiera się w $\mathbb R$,
a pozostałe pierwiastki $\omega\alpha,\omega^2\alpha$
są nierzeczywiste, zatem do niego nie należą.
Nie jest ono normalne nad $\mathbb Q$, a więc nie jest Galois.

Zastąpmy teraz ciało podstawowe przez $F'=\mathbb Q(\omega)$.
Mamy $[F':\mathbb Q]=2$ i $\mu_3(F')=\langle\omega\rangle$.
Gdyby $X^3-2$ był rozkładalny nad $F'$, to jako wielomian
stopnia $3$ miałby pierwiastek $\omega^j\alpha\in F'$
dla pewnego $j\in\{0,1,2\}$. Wówczas
$\alpha=\omega^{-j}(\omega^j\alpha)\in F'$, co przeczy
$[\mathbb Q(\alpha):\mathbb Q]=3$, gdyż stopień podciała
$\mathbb Q(\alpha)\subseteq F'$ musiałby dzielić
$[F':\mathbb Q]=2$.
Zatem $X^3-2$ jest nierozkładalny nad $F'$, a
$E=F'(\alpha)/F'$ jest cyklicznym rozszerzeniem Galois
stopnia $3$ z automorfizmem $\sigma(\alpha)=\omega\alpha$.
Jego działanie na pierwiastkach przedstawia
rysunek~\ref{aa:fig:kummer-cubic-cycle}.
\end{przyklad}
\begin{figure}[htbp]
\centering
\begin{tikzpicture}[>=Stealth, every node/.style={font=\small}]
  \node (a) at (0,1.6) {$\alpha$};
  \node (b) at (2.6,0) {$\omega\alpha$};
  \node (c) at (-2.6,0) {$\omega^2\alpha$};
  \draw[->,thick] (a) -- node[above right] {$\sigma$} (b);
  \draw[->,thick] (b) -- node[below] {$\sigma$} (c);
  \draw[->,thick] (c) -- node[above left] {$\sigma$} (a);
\end{tikzpicture}
\caption{Cykliczne działanie generatora grupy Galois na trzech pierwiastkach $X^3-2$ nad $\mathbb Q(\omega)$.}
\label{aa:fig:kummer-cubic-cycle}
\end{figure}

\begin{przyklad}[Dlaczego charakterystyka ma znaczenie?]\label{aa:prz:kummer-inseparable}
Niech $p$ będzie dowolną liczbą pierwszą, $F=\mathbb F_p(t)$
(ciało funkcji wymiernych w zmiennej $t$), a $\alpha$
niech spełnia $\alpha^p=t$. Wielomian $X^p-t$ ma
pochodną zerową. Wykażemy, że $\alpha\notin F$.
Gdyby $t=(f(t)/g(t))^p$ dla niezerowych wielomianów
$f,g\in\mathbb F_p[t]$, to porównanie wykładników czynnika
nierozkładalnego $t$ po obu stronach równości
$f(t)^p=tg(t)^p$ dałoby
$p\,v_t(f)=1+p\,v_t(g)$, co jest niemożliwe modulo $p$.
Tutaj $v_t(h)$ oznacza krotność czynnika $t$ w niezerowym
wielomianie $h$.

Ponieważ $t=\alpha^p$ jest przestępne nad $\mathbb F_p$,
także $\alpha$ jest przestępne nad $\mathbb F_p$. Mamy
$F(\alpha)=\mathbb F_p(\alpha)$ i
$F=\mathbb F_p(\alpha^p)$.
Elementy $1,\alpha,\ldots,\alpha^{p-1}$ są liniowo
niezależne nad $F$: po sprowadzeniu ewentualnej relacji
do wspólnego mianownika otrzymalibyśmy niezerowy
wielomian w $\alpha$ o wyrazach, których wykładniki
są różne modulo $p$, sprzecznie z przestępnością $\alpha$.
Z drugiej strony $\alpha^p=t\in F$, więc
wielomiany w $\alpha$ redukują się do kombinacji
$1,\ldots,\alpha^{p-1}$. Zbiór tych kombinacji jest
podpierścieniem $F[\alpha]$ o wymiarze $p$ nad $F$;
ponieważ jest dziedziną skończenie wymiarową nad ciałem,
każdy niezerowy element ma odwrotność: mnożenie przez
niego jest iniektywnym, a więc surjektywnym
$F$-liniowym endomorfizmem $F[\alpha]$.
Zatem $F[\alpha]=F(\alpha)$, a wskazane elementy
tworzą bazę tego ciała nad $F$. W konsekwencji
$[F(\alpha):F]=p$.

Jeśli $\tau$ jest dowolnym $F$-automorfizmem $F(\alpha)$,
to $\tau(\alpha)^p=t=\alpha^p$, zatem
$(\tau(\alpha)-\alpha)^p=0$. W ciele element
o zerowej $p$-tej potędze musi być zerowy,
a więc $\tau(\alpha)=\alpha$ i $\tau=\operatorname{id}$.
Rozszerzenie nie jest Galois mimo stopnia $p$.
Właśnie dlatego nie wolno usuwać założenia
$\operatorname{char}F\nmid n$.
\end{przyklad}

\begin{przyklad}[Złożony wykładnik — stopień rozszerzenia może być mniejszy]
Rozpatrzmy $F=\mathbb Q(i)$, $n=4$ oraz $a=4$.
Ciało $F$ zawiera wszystkie cztery pierwiastki z jedynki
$\{1,-1,i,-i\}$. Wybierając $\alpha=\sqrt2$, mamy
$\alpha^4=4$, ale $[F(\alpha):F]=2$, ponieważ
$\sqrt2\notin\mathbb Q(i)$, gdyż rzeczywiste elementy $\mathbb Q(i)$
są wymierne, a $\sqrt2\notin\mathbb Q$; ponadto $\sqrt2$
spełnia $X^2-2=0.

Co więcej, $4$ nie jest czwartą potęgą w $F$.
Jeśli bowiem $b=x+yi\in\mathbb Q(i)$ ($x,y\in\mathbb Q$)
i $b^4=4$, to $b^2=2$ lub $b^2=-2$.
Ponieważ $b^2=(x^2-y^2)+2xyi$ jest rzeczywiste,
$xy=0$. Jeśli $b^2=2$, to $x=0$ jest niemożliwe,
więc $y=0$ i $x^2=2$. Jeśli $b^2=-2$, to $y=0$
jest niemożliwe, więc $x=0$ i $y^2=2$;
żaden z tych warunków nie zachodzi dla liczb wymiernych.
Tak więc $a$ nie jest czwartą potęgą, a mimo to
$X^4-4=(X^2-2)(X^2+2)$ ma nad $F$ czynniki stopnia $2$.
Grupa $\operatorname{Gal}(F(\sqrt2)/F)$ jest cykliczna
rzędu $2$, a nie $4$. Wyjaśnia to potrzebę dodatkowych
założeń przy klasyfikacji rozszerzeń o zadanym stopniu.
\end{przyklad}

\section{Kiedy dwa generatory dają to samo rozszerzenie?}
Dla liczby pierwszej $p$ sytuacja jest szczególnie przejrzysta:
nie istnieją nietrywialne właściwe podgrupy grupy cyklicznej
rzędu $p$. Dlatego przy założeniach $\operatorname{char}F\neq p$
oraz $|\mu_p(F)|=p$ element niebędący $p$-tą potęgą
daje rozszerzenie dokładnie stopnia $p$.

\begin{twierdzenie}[Klasyfikacja cyklicznych rozszerzeń Kummera stopnia $p$]\label{aa:tw:kummer-prime-classification}
Niech $p$ będzie liczbą pierwszą różną od
$\operatorname{char}F$. Załóżmy, że $F$ zawiera pierwotny
$p$-ty pierwiastek z jedynki. Niech $a,b\in F^\times$
nie będą $p$-tymi potęgami w $F$, a
$\alpha,\beta\in\overline F$ spełniają
$\alpha^p=a$ oraz $\beta^p=b$.
Wówczas równoważne są warunki:
\begin{enumerate}[label=\textup{(\roman*)}]
\item $F(\alpha)=F(\beta)$ jako podciała $\overline F$;
\item istnieją $j\in\{1,\ldots,p-1\}$ oraz
$c\in F^\times$ takie, że $b=c^p a^j$.
\end{enumerate}
\end{twierdzenie}
\begin{proof}
\emph{Przygotowanie: oba rozszerzenia mają stopień $p$.}
Na mocy twierdzenia~\ref{aa:tw:kummer-radical-cyclic}
$F(\alpha)/F$ jest Galois o grupie, której rząd dzieli $p$.
Jeśli rząd byłby $1$, to $[F(\alpha):F]=1$,
a więc $\alpha\in F$ i $a=\alpha^p$ byłoby $p$-tą
potęgą w $F$, sprzecznie z założeniem. Ponieważ $p$ jest
pierwsze, musi zachodzić $[F(\alpha):F]=p$.
Ten sam argument daje $[F(\beta):F]=p$.

\emph{Kierunek (ii)$\Rightarrow$(i).}
Załóżmy $b=c^p a^j$ dla wskazanych $c,j$.
Element $d=\beta/(c\alpha^j)$ jest dobrze określony,
ponieważ $c\alpha^j\neq0$, a jego $p$-ta potęga wynosi
\[
d^p=\frac{\beta^p}{c^p(\alpha^p)^j}
 =\frac{b}{c^pa^j}=1.
\]
Stąd $d\in\mu_p(\overline F)=\mu_p(F)\subseteq F$.
Otrzymujemy $\beta=dc\alpha^j\in F(\alpha)$,
czyli $F(\beta)\subseteq F(\alpha)$.
Obie strony mają stopień $p$ nad $F$, dlatego z twierdzenia
o stopniach w wieży wynika równość tych ciał.

\emph{Kierunek (i)$\Rightarrow$(ii).}
Niech $E=F(\alpha)=F(\beta)$, a $\zeta\in F$ będzie
pierwotnym pierwiastkiem $p$-tego stopnia z jedynki.
Grupa $\operatorname{Gal}(E/F)$ jest cykliczna rzędu $p$.
Każdy jej automorfizm wysyła $\alpha$ na $u\alpha$
dla pewnego $u\in\mu_p(F)$. Ponieważ $\alpha\notin F$,
nie wszystkie automorfizmy ustalają $\alpha$;
możemy więc wybrać generator $\sigma$ taki, że
$\sigma(\alpha)=\zeta\alpha$:
wyjściowy iloraz $u=\zeta^k$ ma $k\not\equiv0\pmod p$,
a odpowiednia potęga automorfizmu, wyznaczona odwrotnością
$k$ modulo $p$, daje żądany generator.

Ponieważ $\beta^p=b\in F$, także
$\sigma(\beta)/\beta$ jest w $\mu_p(F)$.
Istnieje zatem $j\in\{0,\ldots,p-1\}$ spełniające
$\sigma(\beta)=\zeta^j\beta$.
Wartość $j=0$ jest niemożliwa: wtedy $\sigma(\beta)=\beta$,
a ponieważ $\sigma$ generuje grupę Galois,
$\beta\in E^{\langle\sigma\rangle}=F$,
co zaprzecza temu, że $b$ nie jest $p$-tą potęgą w $F$.
Zatem $1\leq j\leq p-1$.
Dla elementu $c=\beta/\alpha^j\in E^\times$ obliczamy
\[
\sigma(c)=\frac{\zeta^j\beta}{(\zeta\alpha)^j}
=c.
\]
Ponieważ ciało elementów stałych całej grupy Galois jest równe $F$,
mamy $c\in F^\times$. Podnosząc równość
$\beta=c\alpha^j$ do potęgi $p$, dostajemy
$b=c^pa^j$, jak należało wykazać.
\end{proof}

\begin{uwaga}[Iloczyn klas a iloczyn tensorowy]\label{aa:uw:kummer-products}
Niech $p$ będzie liczbą pierwszą, a $F$ dowolnym ciałem.
Grupa $F^\times/(F^\times)^p$ składa się z warstw
$[a]=a(F^\times)^p$ dla $a\in F^\times$;
mnożymy je według reguły $[a][b]=[ab]$.
Każda klasa spełnia $[a]^p=[1]$, więc ta grupa abelowa
jest w naturalny sposób przestrzenią liniową nad $\mathbb F_p$:
dodawanie wektorów jest mnożeniem klas, a mnożenie
przez skalar $j\in\mathbb F_p$ jest operacją
$[a]\mapsto[a]^j$. Działania są dobrze określone,
ponieważ zmiana przedstawiciela $a$ o $p$-tą potęgę
nie zmienia klasy $[a]^j$.
Nie jest to ani iloczyn prosty grup, ani iloczyn tensorowy modułów.

Przy założeniu $\operatorname{char}F\neq p$ i obecności
pierwotnego $p$-tego pierwiastka z jedynki
niezerowe jednowymiarowe podprzestrzenie
$F^\times/(F^\times)^p$ odpowiadają dokładnie cyklicznym
rozszerzeniom Galois stopnia $p$ wewnątrz $\overline F$.
Istotnie, każda taka prosta jest generowana przez klasę
$[a]\neq[1]$, a dołączając $\sqrt[p]{a}$ otrzymujemy
rozszerzenie stopnia $p$ na mocy poprzedniego twierdzenia.
Dwie klasy generują tę samą prostą dokładnie wtedy,
gdy $[b]=[a]^j$ dla pewnego $1\leq j\leq p-1$;
poprzednie twierdzenie mówi, że jest to równoważne
równości odpowiednich ciał. Każde cykliczne rozszerzenie
Galois stopnia $p$ pojawia się w ten sposób na mocy
twierdzenia~\ref{aa:tw:kummer-converse}.
\end{uwaga}

\begin{przyklad}[Trzy różne rozszerzenia kwadratowe]\label{aa:prz:kummer-three-quadratic}
Niech $F=\mathbb Q$ i $p=2$. Ponieważ $-1\in\mathbb Q$
jest pierwotnym pierwiastkiem drugiego stopnia z jedynki,
możemy stosować twierdzenie~\ref{aa:tw:kummer-prime-classification}.
Liczby $2$, $3$ i $6$ nie są kwadratami liczb wymiernych.
Istotnie, jeśli $r^2=m$ dla dodatniej liczby całkowitej
$m\in\{2,3,6\}$ i liczby wymiernej $r$, to po zapisaniu
$r=u/v$ w postaci nieskracalnej, gdzie $u\in\mathbb Z$,
$v\in\mathbb Z_{>0}$, z równości $u^2=mv^2$ wynika,
że wykładnik każdej liczby pierwszej w rozkładzie $m$
musiałby być parzysty. Dla każdej z trzech liczb jest
to fałszywe.

Rozszerzenia $\mathbb Q(\sqrt2)$, $\mathbb Q(\sqrt3)$
i $\mathbb Q(\sqrt6)$ są zatem kwadratowe i Galois.
Są parami różne: dla $p=2$ zbiór dopuszczalnych wykładników
$j\in\{1,\ldots,p-1\}$ składa się wyłącznie z $j=1$,
więc równość dwóch z tych ciał wymagałaby, aby iloraz
odpowiednich liczb $2,3,6$ był kwadratem w $\mathbb Q$.
Ilorazy te to $2/3$, $1/3$ i $1/2$ (lub ich odwrotności),
a każdy ma w rozkładzie na czynniki pierwsze wykładnik
nieparzysty. Warunek klasyfikacji nie jest spełniony.

Jednocześnie $\sqrt6=\sqrt2\sqrt3$ należy do
$\mathbb Q(\sqrt2,\sqrt3)$. Iloczyn klas $[2][3]=[6]$
w $\mathbb Q^\times/(\mathbb Q^\times)^2$ odpowiada
trzeciemu podciału kwadratowemu rozszerzenia
$\mathbb Q(\sqrt2,\sqrt3)/\mathbb Q$; nie jest to
iloczyn tensorowy tych ciał ani ich iloczyn prosty.
\end{przyklad}

\begin{przyklad}[Dwie różne klasy dające to samo rozszerzenie]\label{aa:prz:kummer-class-example}
Niech $F=\mathbb Q(\omega)$, gdzie $\omega$ jest pierwotnym
pierwiastkiem trzeciego stopnia z jedynki,
a $a=2$, $b=4=a^2$. W przykładzie
\ref{aa:prz:kummer-missing-roots} wykazaliśmy,
że $2$ nie jest sześcianem w $F$.
Gdyby $[2]=[4]$ w $F^\times/(F^\times)^3$,
to $4/2=2$ byłoby sześcianem w $F$,
co jest niemożliwe. Klasy są więc różne.

Niech $\alpha^3=2$, $\beta=\alpha^2$.
Wtedy $\beta^3=4$ i oczywiście $\beta\in F(\alpha)$.
Z drugiej strony $\beta^2=\alpha^4=2\alpha$,
więc $\alpha=\beta^2/2\in F(\beta)$.
Stąd $F(\alpha)=F(\beta)$, mimo że $[2]\neq[4]$.
Obie klasy należą do tej samej jednowymiarowej podprzestrzeni,
ponieważ $[4]=[2]^2$, a w $\mathbb F_3$ skalar $2$
jest niezerowy.
Jeżeli $\sigma(\alpha)=\omega\alpha$, to
$\sigma(\beta)=\omega^2\beta$:
różne generatory tego samego rozszerzenia mogą więc
transformować się przez różne potęgi pierwotnego
pierwiastka z jedynki.
\end{przyklad}

% AA-27 | Iteracja XIV; szczegółowy audyt historyczny III, 20 XI 2026.
\chapter[Rozszerzenia Artina--Schreiera]{Rozszerzenia Artina--Schreiera w charakterystyce dodatniej}
\label{aa:ch:artin-schreier}

\section{Wielomian addytywny i jego pierwiastki}

W rozdziale~\ref{aa:ch:kummer} do opisu pewnych rozszerzeń cyklicznych
wykorzystano równania potęgowe. Jeżeli jednak $F$ ma charakterystykę $p$,
to pochodna wielomianu $X^p-a$ jest zerowa. Wielomian ten nie może zatem
pełnić tej samej roli w teorii rozszerzeń rozdzielczych stopnia $p$.
Zamiast niego rozważymy $X^p-X-a$: dodatkowy składnik $-X$ zapewnia
niezerową pochodną, a tożsamość $(x+y)^p=x^p+y^p$ pozwala opisać
wszystkie pierwiastki i automorfizmy.

\begin{definicja}\label{aa:def:as-map}
Niech $p$ będzie liczbą pierwszą, a $F$ ciałem charakterystyki $p$.
Przez $\mathbb F_p\subseteq F$ oznaczamy ciało pierwsze. Odwzorowanie
\(\wp_F:(F,+)\longrightarrow(F,+),\quad \wp_F(x)=x^p-x\),
nazywamy \emph{odwzorowaniem Artina--Schreiera}. Jego obraz oznaczamy
przez $\wp(F)=\{x^p-x:x\in F\}$. Wszystkie rozszerzenia algebraiczne
$F$ rozpatrujemy we wspólnym, ustalonym domknięciu algebraicznym
$\overline F$.
\end{definicja}
\begin{proof}[Uzasadnienie własności odwzorowania]
Dla dowolnych $x,y\in F$ dwumian Newtona daje
\((x+y)^p=x^p+y^p\): dla $1\le j\le p-1$ współczynnik
$\binom pj$ jest podzielny przez $p$, a więc jest zerem w $F$.
W konsekwencji $\wp_F(x+y)=\wp_F(x)+\wp_F(y)$, co dowodzi
addytywności. Jeżeli $\wp_F(x)=0$, to $x$ jest pierwiastkiem
$X^p-X$. Każdy $c\in\mathbb F_p$ spełnia $c^p=c$:
dla $c=0$ jest to natychmiastowe, a dla $c\ne0$ wynika z faktu,
że grupa $\mathbb F_p^\times$ ma rząd $p-1$. Wielomian stopnia
$p$ nie ma więcej niż $p$ pierwiastków w ciele. Zatem
$\ker\wp_F=\mathbb F_p$. W szczególności $\wp(F)$ jest podgrupą
addytywnej grupy $F$, a iloraz $F/\wp(F)$ ma sens.
\end{proof}

\begin{twierdzenie}[Konstrukcja rozszerzenia Artina--Schreiera]\label{aa:tw:as-forward}
Niech $p$ będzie liczbą pierwszą, $F$ ciałem charakterystyki $p$,
$a\in F$, zaś $\alpha\in\overline F$ pierwiastkiem
$f_a(X)=X^p-X-a$. Jeżeli $a\notin\wp(F)$, to $f_a$ jest
nierozkładalny nad $F$, rozszerzenie $E=F(\alpha)/F$ jest Galois
stopnia $p$, a przyporządkowanie
\(c\longmapsto\sigma_c,\quad \sigma_c(\alpha)=\alpha+c\)
wyznacza izomorfizm grupy addytywnej $(\mathbb F_p,+)$ na
$\operatorname{Gal}(E/F)$.
\end{twierdzenie}
\begin{proof}
\emph{Krok 1: pierwiastki i rozdzielczość.} Ponieważ
$f_a'(X)=pX^{p-1}-1=-1$, wielomian $f_a$ nie ma pierwiastków
wielokrotnych. Dla każdego $c\in\mathbb F_p$ zachodzi
\((\alpha+c)^p-(\alpha+c)=\alpha^p-\alpha+c^p-c=a\).
Elementy $\alpha+c$ są parami różne, gdyż dodawanie $\alpha$
jest iniekcją. Jest ich dokładnie $p$, czyli tyle, ile wynosi
stopień $f_a$. Wszystkie należą do $E$, więc $E$ jest ciałem
rozkładu rozdzielczego wielomianu $f_a$ nad $F$. Stąd $E/F$
jest skończonym rozszerzeniem Galois.

\emph{Krok 2: zanurzenie grupy Galois.} Niech
$G=\operatorname{Gal}(E/F)$ i $\sigma\in G$.
Ponieważ $\sigma$ ustala współczynniki $f_a$,
element $\sigma(\alpha)$ jest pewnym pierwiastkiem $f_a$.
Istnieje zatem dokładnie jeden $c_\sigma\in\mathbb F_p$,
dla którego $\sigma(\alpha)=\alpha+c_\sigma$.
Odwzorowanie $G\to(\mathbb F_p,+)$ dane przez
$\sigma\mapsto c_\sigma$ jest iniektywne: jeśli
$c_\sigma=c_\tau$, to $\sigma$ i $\tau$ zgadzają się na
$\alpha$ i na $F$, a zatem na całym $F(\alpha)$. Jest także
homomorfizmem, ponieważ dla $\sigma,\tau\in G$
\[
(\sigma\circ\tau)(\alpha)
=\sigma(\alpha+c_\tau)
=\alpha+c_\sigma+c_\tau,
\]
gdzie użyliśmy faktu $c_\tau\in\mathbb F_p\subseteq F$.
W szczególności $|G|$ dzieli $p$, gdyż $G$ jest izomorficzna
z podgrupą grupy addytywnej rzędu $p$.

\emph{Krok 3: eliminacja przypadku trywialnego.}
Jeżeli $|G|=1$, to $[E:F]=|G|=1$, ponieważ $E/F$ jest Galois.
Wtedy $\alpha\in F$, skąd $a=\alpha^p-\alpha\in\wp(F)$,
wbrew założeniu. Zatem $|G|=p$ i $[E:F]=p$.
Wielomian minimalny $\alpha$ nad $F$ dzieli $f_a$ i ma
stopień $[F(\alpha):F]=p$, wobec czego jest równy $f_a$.
Zanurzenie $G\to(\mathbb F_p,+)$ ma teraz obraz rzędu $p$,
a więc jest izomorfizmem. Odwrotne przyporządkowanie daje
wskazane automorfizmy $\sigma_c$.
\end{proof}

\begin{twierdzenie}[Odwrotne twierdzenie Artina--Schreiera]\label{aa:tw:as-converse}
Niech $p$ będzie liczbą pierwszą, $F$ ciałem charakterystyki $p$
i $E/F$ rozszerzeniem Galois, którego grupa jest cykliczna rzędu
$p$. Wówczas istnieją $\alpha\in E$ i $a\in F\setminus\wp(F)$
takie, że $E=F(\alpha)$ oraz $\alpha^p-\alpha=a$.
\end{twierdzenie}
\begin{proof}
Niech $\sigma$ będzie generatorem $G=\operatorname{Gal}(E/F)$.
Ponieważ $E/F$ jest rozdzielcze, jego forma śladu
$(x,y)\mapsto\operatorname{Tr}_{E/F}(xy)$ jest niezdegenerowana
(zob. rozdział~\ref{aa:ch:trace-norm}). W szczególności sam
funkcjonał $x\mapsto\operatorname{Tr}_{E/F}(x)$ nie jest
zerowy: w przeciwnym razie forma śladu byłaby identycznie
zerowa. Jest to niezerowy funkcjonał do jednowymiarowej
$F$-przestrzeni $F$, więc istnieje $u\in E$ o śladzie $1$
(po pomnożeniu dowolnego elementu o niezerowym śladzie przez
odwrotność tego śladu).

Definiujemy
\[
\alpha=-\sum_{j=0}^{p-1}j\,\sigma^j(u),
\]
przy czym $j$ oznacza obraz liczby całkowitej $j$ w
$\mathbb F_p\subseteq F$. Sprawdzimy zmianę $\alpha$ pod
$\sigma$, zamiast przyjmować ją bez uzasadnienia. Po zmianie
indeksu $k=j+1$ oraz użyciu $\sigma^p=\mathrm{id}$ mamy
\[
\sigma(\alpha)=-\sum_{k=1}^{p-1}(k-1)\sigma^k(u)
-(p-1)u,
\qquad
\alpha=-\sum_{k=1}^{p-1}k\sigma^k(u).
\]
Odejmowanie daje
\[
\sigma(\alpha)-\alpha
=\sum_{k=1}^{p-1}\sigma^k(u)-(p-1)u
=\sum_{k=0}^{p-1}\sigma^k(u)
=\operatorname{Tr}_{E/F}(u)=1,
\]
ponieważ $-(p-1)=1$ w charakterystyce $p$.
Następnie
\[
\sigma(\alpha^p-\alpha)
=(\alpha+1)^p-(\alpha+1)=\alpha^p-\alpha.
\]
Zatem $a:=\alpha^p-\alpha$ należy do ciała elementów
stałych $E^G=F$. Ponieważ $\sigma(\alpha)\ne\alpha$,
element $\alpha$ nie należy do $F$. Z drugiej strony
$[E:F]=p$ i $F\subsetneq F(\alpha)\subseteq E$, więc z
multiplikatywności stopni wynika $[F(\alpha):F]=p$ i
$E=F(\alpha)$. Gdyby $a=b^p-b$ dla pewnego $b\in F$,
to $(\alpha-b)^p-(\alpha-b)=0$. Własność jądra $\wp_F$
dowiedziona wyżej (w domknięciu algebraicznym również)
pokazuje, że $\alpha-b\in\mathbb F_p\subseteq F$; otrzymalibyśmy
$\alpha\in F$, sprzeczność. Stąd $a\notin\wp(F)$.
\end{proof}

\section{Przykłady, kontrprzykłady i klasyfikacja}

\begin{przyklad}[Rozszerzenie ciała funkcji wymiernych]\label{aa:prz:as-rational}
Niech $F=\mathbb F_p(t)$, gdzie $t$ jest zmienną
transcendentną nad $\mathbb F_p$. Uzasadnimy, że
$t\notin\wp(F)$. Każdy niezerowy $b\in F$ zapisuje się
jako $P(t)/Q(t)$, gdzie $P,Q\in\mathbb F_p[t]$ są niezerowe.
Niech $m=\deg P-\deg Q$, czyli rząd bieguna
$b$ w punkcie nieskończoności, jeśli $m>0$.
Gdy $m>0$, funkcja $b^p$ ma stopień w nieskończoności $pm$,
ściśle większy od $m$; dlatego $b^p-b$ również ma tam
biegun rzędu $pm$. Gdy $m\le0$, zarówno $b^p$, jak i $b$
nie mają bieguna w nieskończoności, więc również ich różnica
nie ma takiego bieguna. Tymczasem $t$ ma biegun rzędu $1$.
Równość $b^p-b=t$ jest niemożliwa, ponieważ $pm\ne1$
dla żadnej liczby całkowitej $m>0$. Dla $b=0$ jest to
tym bardziej jasne. Twierdzenie~\ref{aa:tw:as-forward}
daje więc cykliczne rozszerzenie Galois
$\mathbb F_p(t,\alpha)/\mathbb F_p(t)$ stopnia $p$, gdzie
$\alpha^p-\alpha=t$. Jego automorfizmy mają postać
$\alpha\mapsto\alpha+c$ dla $c\in\mathbb F_p$.
\end{przyklad}

\begin{przyklad}[Dlaczego warunek na klasę $a$ jest konieczny?]\label{aa:prz:as-trivial}
Niech $b\in F$ i $a=b^p-b$. Wówczas dla każdego
$c\in\mathbb F_p$ mamy $(b+c)^p-(b+c)=a$, skąd
\[
X^p-X-a=\prod_{c\in\mathbb F_p}(X-b-c).
\]
Wszystkie czynniki są różne, bo $b+c=b+d$ oznacza $c=d$.
Wielomian pozostaje więc rozdzielczy, lecz jego ciało rozkładu
to samo $F$. Na przykład nad $F=\mathbb F_2(t)$ dla
$b=t$ otrzymujemy $a=t^2+t$ i rozkład
$X^2+X+t^2+t=(X+t)(X+t+1)$.
\end{przyklad}

\begin{stwierdzenie}[Zależność od klasy addytywnej]\label{aa:prop:as-classes}
Niech $F$ będzie ciałem charakterystyki $p$ i $a,b\in F$.
Jeżeli istnieją $c\in\mathbb F_p^\times$ oraz $d\in F$
takie, że $b=ca+\wp_F(d)$, to ciała rozkładu wielomianów
$X^p-X-a$ i $X^p-X-b$ w ustalonym $\overline F$ są równe.
\end{stwierdzenie}
\begin{proof}
Wybierzmy $\alpha\in\overline F$ o
$\alpha^p-\alpha=a$ i połóżmy $\beta=c\alpha+d$.
Ponieważ $c^p=c$ dla $c\in\mathbb F_p$, otrzymujemy
$\beta^p-\beta=c(\alpha^p-\alpha)+d^p-d=b$.
Wszystkimi pierwiastkami pierwszego wielomianu są
$\alpha+\mathbb F_p$, a drugiego $\beta+\mathbb F_p$.
Ich ciała rozkładu to odpowiednio $F(\alpha)$ i
$F(\beta)$. Z równości $\beta=c\alpha+d$ wnioskujemy
$F(\beta)\subseteq F(\alpha)$. Ponieważ $c\ne0$,
mamy też $\alpha=c^{-1}(\beta-d)\in F(\beta)$;
zatem zachodzi równość.
\end{proof}

\begin{twierdzenie}[Klasyfikacja rozszerzeń cyklicznych stopnia $p$]
\label{aa:tw:as-classification}
Niech $F$ będzie ciałem charakterystyki $p$, a $a,b$ dwoma
elementami $F\setminus\wp(F)$. Niech $E_a,E_b\subset\overline F$
będą ciałami rozkładu odpowiednio $X^p-X-a$ i $X^p-X-b$.
Wówczas $E_a=E_b$ wtedy i tylko wtedy, gdy istnieją
$c\in\mathbb F_p^\times$ i $d\in F$ takie, że
$b=ca+\wp_F(d)$. Innymi słowy, jednowymiarowe
podprzestrzenie przestrzeni wektorowej $F/\wp(F)$ nad
$\mathbb F_p$ odpowiadają bijektywnie cyklicznym
rozszerzeniom Galois $F$ stopnia $p$ wewnątrz $\overline F$.
\end{twierdzenie}
\begin{proof}
Wystarczalność dowiedziono w
stwierdzeniu~\ref{aa:prop:as-classes}. Dla konieczności
załóżmy $E_a=E_b=:E$. Wybierzmy pierwiastki
$\alpha^p-\alpha=a$ i $\beta^p-\beta=b$ w $E$.
Na mocy twierdzenia~\ref{aa:tw:as-forward} istnieje generator
$\sigma$ grupy $\operatorname{Gal}(E/F)$ spełniający
$\sigma(\alpha)=\alpha+1$. Ponieważ $\sigma$ jest
$F$-automorfizmem, element $\sigma(\beta)$ także jest
pierwiastkiem $X^p-X-b$, a więc
$\sigma(\beta)=\beta+c$ dla pewnego $c\in\mathbb F_p$.
Gdyby $c=0$, to $\beta\in E^{\langle\sigma\rangle}=F$,
co dawałoby $b\in\wp(F)$. Zatem $c\ne0$.
Teraz $\sigma(\beta-c\alpha)=\beta-c\alpha$, więc
$d:=\beta-c\alpha\in F$. Z tożsamości
$\beta=c\alpha+d$ wynika
$b=\beta^p-\beta=c(\alpha^p-\alpha)+d^p-d
=ca+\wp_F(d)$.
Na koniec $\wp(F)$ jest podprzestrzenią $F$ nad
$\mathbb F_p$ (addytywność oraz
$\wp_F(cx)=c\wp_F(x)$ dla $c\in\mathbb F_p$).
Zatem warunek $b=ca+\wp_F(d)$ jest dokładnie równością
prostych $\mathbb F_p[a]=\mathbb F_p[b]$ w ilorazie,
gdzie $[a]$ oznacza klasę $a$. Konstrukcja z pierwszego
twierdzenia i jego odwrotność zapewniają istnienie oraz
pełność tej klasyfikacji.
\end{proof}

\begin{przyklad}[Charakterystyka dwa i ciało czteroelementowe]
\label{aa:prz:as-four-elements}
Dla $F=\mathbb F_2$ mamy $\wp_F(0)=\wp_F(1)=0$,
więc $1\notin\wp(F)$. Wielomian $X^2-X-1=X^2+X+1$
nie ma pierwiastka w $\mathbb F_2$ i wyznacza
$E=\mathbb F_2[\alpha]$ o czterech elementach.
W $E$ zachodzi $\alpha^2=\alpha+1$, dlatego
$\sigma(\alpha)=\alpha^2=\alpha+1$ jest nietrywialnym
$F$-automorfizmem. Ponieważ
$\sigma^2(\alpha)=(\alpha+1)^2=\alpha^2+1=\alpha$,
jest to generator grupy cyklicznej rzędu $2$.
W tym przypadku $\alpha^2-\alpha=1$ i wszystkie formuły
Artina--Schreiera można bezpośrednio sprawdzić.
\end{przyklad}

\begin{uwaga}\label{aa:uw:as-vs-kummer}
Warunek $a\notin\wp(F)$ oznacza, że klasa $[a]$ jest
niezerowa w $F/\wp(F)$. W charakterystyce $p$ wielomian
$X^p-a$ ma pochodną zero i, gdy ma pierwiastek w domknięciu
algebraicznym, jest potęgą $(X-\gamma)^p$.
Natomiast $X^p-X-a$ ma pochodną $-1$, niezależnie od $a$.
To wyjaśnia, dlaczego w opisie \emph{rozdzielczych}
rozszerzeń cyklicznych stopnia $p$ stosujemy konstrukcję
Artina--Schreiera, a nie równanie potęgowe z teorii Kummera.
\end{uwaga}

% AA-28 | Iteracja XIV; szczegółowy audyt historyczny III, 20 XI 2026.
\chapter{Ciała cyklotomiczne i wielomiany cyklotomiczne}
\label{aa:ch:cyclotomic}
\section{Pierwiastki pierwotne i wielomiany cyklotomiczne}

W teorii Galois ważnym źródłem rozszerzeń abelowych są ciała
otrzymywane przez dołączenie pierwiastków z jedynki.
Zamiast rozpoczynać od wzorów z funkcją wykładniczą zespoloną,
rozważymy czysto algebraicznie pierwiastki równania $X^n=1$.
To podejście nie wymaga analitycznej konstrukcji $\mathbb R$
ani definicji funkcji wykładniczej. Jeżeli zechcemy narysować
pierwiastki w płaszczyźnie zespolonej, skorzystamy jedynie
z geometrycznej interpretacji $\mathbb C$, traktowanego tutaj
jako uprzednio dane ciało, zob. rozdział~\ref{aa:ch:fields-fractions}.

\begin{definicja}\label{aa:def:cyclotomic}
Niech $n\ge1$ będzie liczbą całkowitą. W ustalonym domknięciu
algebraicznym $\overline{\mathbb Q}$ wybieramy element
$\zeta_n$ o \emph{dokładnym rzędzie multiplikatywnym $n$}:
$\zeta_n^n=1$ oraz $\zeta_n^d\ne1$ dla każdej liczby
całkowitej $1\le d<n$. Taki element nazywamy
\emph{pierwotnym $n$-tym pierwiastkiem z jedynki}.
Przyjmujemy $\zeta_1=1$. Definiujemy
\[
\Phi_n(X)=\prod_{\substack{1\le k\le n\\\gcd(k,n)=1}}
(X-\zeta_n^k)\in\overline{\mathbb Q}[X].
\]
Jest to \emph{$n$-ty wielomian cyklotomiczny}. Symbol
$\varphi(n)$ oznacza liczbę liczb całkowitych
$1\le k\le n$ względnie pierwszych z $n$.
\end{definicja}

\begin{uwaga}[Istnienie i niezależność od wyboru pierwiastka]
W charakterystyce $0$ pochodna $X^n-1$ ma postać
$nX^{n-1}$ i jest względnie pierwsza z $X^n-1$.
W $\overline{\mathbb Q}$ wielomian $X^n-1$ ma więc
$n$ różnych pierwiastków. Stanowią one skończoną podgrupę
multiplikatywną ciała, a każda skończona podgrupa grupy
multiplikatywnej ciała jest cykliczna (zob. wyniki
z rozdziału~\ref{aa:ch:kummer}).
Istnieje zatem jej generator rzędu $n$.
Jeśli $\eta_n$ jest innym generatorem, to
$\eta_n=\zeta_n^u$ z $\gcd(u,n)=1$.
Mnożenie przez $u$ permutuje klasy liczb względnie
pierwszych z $n$, dlatego zbiór czynników w definicji
$\Phi_n$ nie zależy od wyboru $\zeta_n$.
Dla $n=1$ definicja daje $\Phi_1(X)=X-1$.
\end{uwaga}

\begin{twierdzenie}[Rozkład wielomianu $X^n-1$]
\label{aa:tw:cyclotomic-product}
Dla każdej liczby całkowitej $n\ge1$ zachodzi równość
\[
X^n-1=\prod_{d\mid n}\Phi_d(X)
\qquad\text{w }\mathbb Z[X].
\]
Ponadto każdy $\Phi_n$ jest moniczny, należy do $\mathbb Z[X]$
i ma stopień $\varphi(n)$.
\end{twierdzenie}
\begin{proof}
\emph{Krok 1: rozkład według rzędów.} Pierwiastki
$X^n-1$ są różne. Rząd $d$ dowolnego takiego pierwiastka
dzieli $n$. Z kolei każdy pierwiastek dokładnego rzędu
$d\mid n$ jest pierwiastkiem $X^n-1$. Zbiory pierwiastków
poszczególnych dokładnych rzędów są rozłączne, a zbiór
pierwiastków rzędu $d$ ma $\varphi(d)$ elementów:
przy ustalonym generatorze są to dokładnie
$\zeta_d^k$ z $\gcd(k,d)=1$. Ponieważ obie strony są
wielomianami monicznymi i mają te same pierwiastki
z krotnościami $1$, otrzymujemy tożsamość w
$\overline{\mathbb Q}[X]$. Z definicji wynika też
$\deg\Phi_n=\varphi(n)$.

\emph{Krok 2: całkowitość współczynników.}
Przeprowadzimy indukcję po $n$. Dla $n=1$ mamy
$\Phi_1=X-1\in\mathbb Z[X]$. Niech $n>1$ i załóżmy,
że $\Phi_d\in\mathbb Z[X]$ dla wszystkich $d<n$.
Wówczas
\(Q(X)=\prod_{d\mid n,\,d<n}\Phi_d(X)\)
jest monicznym wielomianem w $\mathbb Z[X]$.
Algorytm dzielenia przez moniczny wielomian działa nad
\emph{dowolnym} pierścieniem współczynników: odejmujemy
kolejne wielokrotności $Q$, aby usuwać najwyższe wyrazy.
Otrzymujemy więc $H,R\in\mathbb Z[X]$ takie, że
$X^n-1=QH+R$ i $R=0$ albo $\deg R<\deg Q$.
Jednocześnie krok pierwszy daje $X^n-1=Q\Phi_n$
w $\overline{\mathbb Q}[X]$. Jednoznaczność dzielenia
w tym ciele implikuje $R=0$ i $H=\Phi_n$. Stąd
$\Phi_n\in\mathbb Z[X]$ i jest moniczny.
\end{proof}

Aby dowieść nierozkładalności $\Phi_n$ nad $\mathbb Q$,
nie wystarczy znać jego pierwiastki. Trzeba wykazać, że
minimalny wielomian jednego pierwiastka zawiera również
wszystkie pozostałe pierwiastki pierwotne. Kluczowy jest
poniższy argument z redukcją modulo liczbę pierwszą.

\begin{lemat}[Zamiana pierwiastka przez potęgę pierwszą]
\label{aa:lem:cyclotomic-prime}
Niech $n\ge1$, $\alpha\in\overline{\mathbb Q}$ będzie
pierwotnym $n$-tym pierwiastkiem z jedynki, zaś
$f\in\mathbb Z[X]$ jego monicznym wielomianem minimalnym
nad $\mathbb Q$. Jeżeli $\ell$ jest liczbą pierwszą
względnie pierwszą z $n$, to $f(\alpha^\ell)=0$.
\end{lemat}
\begin{proof}
Załóżmy przeciwnie, że $f(\alpha^\ell)\ne0$.
Ponieważ $\gcd(\ell,n)=1$, element $\alpha^\ell$
ma dokładny rząd $n$. Niech $g\in\mathbb Z[X]$ będzie
jego monicznym wielomianem minimalnym nad $\mathbb Q$.
Jego współczynniki są całkowite: $\alpha^\ell$ jest
pierwiastkiem monicznego $\Phi_n\in\mathbb Z[X]$, a
moniczny czynnik nierozkładalny takiego wielomianu ma
współczynniki w $\mathbb Z$ na mocy lematu Gaussa.
Wielomiany $f$ i $g$ są różne, ponieważ
$f(\alpha^\ell)\ne0$, i oba dzielą $X^n-1$ w
$\mathbb Z[X]$. Są nierozkładalne nad $\mathbb Q$, a więc
względnie pierwsze w $\mathbb Q[X]$. Ich iloczyn
$fg$ dzieli $X^n-1$ w $\mathbb Z[X]$:
najpierw uzyskujemy tę podzielność w $\mathbb Q[X]$,
a następnie stosujemy lemat Gaussa do wielomianów monicznych.

Ponieważ $g(\alpha^\ell)=0$, mamy
$(g(X^\ell))(\alpha)=0$. Minimalność $f$ daje
$f\mid g(X^\ell)$ w $\mathbb Q[X]$ i, ponownie z
moniczności i lematu Gaussa, w $\mathbb Z[X]$.
Oznaczmy redukcję współczynników modulo $\ell$ kreską.
W charakterystyce $\ell$ zachodzi
$\overline{g(X^\ell)}=\overline g(X)^\ell$, toteż
\(\overline f\mid\overline g(X)^\ell\)
w $\mathbb F_\ell[X]$. Wielomian $\overline f$ jest
niezerowy i moniczny. Wybierzmy dowolny jego moniczny
nierozkładalny dzielnik $h$. Ponieważ $h$ jest
pierwszy w pierścieniu $\mathbb F_\ell[X]$, z
$h\mid\overline g^\ell$ wynika $h\mid\overline g$.
Oznacza to, że $h^2$ dzieli
$\overline f\,\overline g=\overline{fg}$, a zatem
$\overline{X^n-1}$. Jest to niemożliwe, ponieważ
$\ell\nmid n$ i
\[
\gcd(X^n-1,\,nX^{n-1})=1
\quad\text{w }\mathbb F_\ell[X].
\]
Rzeczywiście, wielomian $X^n-1$ nie ma wspólnego
pierwiastka z $X^{n-1}$, a $n$ jest niezerowym skalarem
w $\mathbb F_\ell$. Zatem $X^n-1$ jest bez czynników
powtórzonych, w sprzeczności z podzielnością przez $h^2$.
\end{proof}

\begin{twierdzenie}[Nierozkładalność wielomianów cyklotomicznych]
\label{aa:tw:cyclotomic-irreducible}
Dla każdej liczby całkowitej $n\ge1$ wielomian
$\Phi_n\in\mathbb Z[X]$ jest nierozkładalny nad
$\mathbb Q$.
\end{twierdzenie}
\begin{proof}
Dla $n=1$ otrzymujemy wielomian liniowy $X-1$.
Niech $n>1$ i niech $\alpha$ będzie pierwotnym
$n$-tym pierwiastkiem z jedynki. Oznaczmy przez $f$
jego moniczny wielomian minimalny nad $\mathbb Q$.
Wielomian $f$ dzieli $\Phi_n$, a jego współczynniki
należą do $\mathbb Z$ na mocy lematu Gaussa.

Weźmy dowolną liczbę całkowitą $k\ge1$ z
$\gcd(k,n)=1$. Jeżeli $k=1$, to $f(\alpha^k)=0$
z definicji. Gdy $k>1$, rozłóżmy $k$ na iloczyn liczb
pierwszych $k=\ell_1\cdots\ell_r$. Każda $\ell_j$
jest względnie pierwsza z $n$. Stosujemy kolejno
Lemat~\ref{aa:lem:cyclotomic-prime}: najpierw
$f(\alpha^{\ell_1})=0$, następnie
$f(\alpha^{\ell_1\ell_2})=0$ i tak dalej.
Na każdym etapie aktualny pierwiastek ma nadal
rząd $n$, a ponieważ jest pierwiastkiem nierozkładalnego
$f$, jego wielomianem minimalnym także jest $f$;
wszystkie warunki lematu pozostają zatem spełnione.
Otrzymujemy $f(\alpha^k)=0$.

W ten sposób $f$ ma wszystkie $\varphi(n)$ parami
różne pierwiastki $\Phi_n$. Ponieważ $f$ i $\Phi_n$
są moniczne, a $f\mid\Phi_n$, wynika $f=\Phi_n$.
\end{proof}

\section{Grupa Galois i przykłady}

\begin{twierdzenie}[Grupa Galois ciała cyklotomicznego]
\label{aa:tw:cyclotomic-galois}
Dla każdej liczby całkowitej $n\ge1$ rozszerzenie
$\mathbb Q(\zeta_n)/\mathbb Q$ jest Galois stopnia
$\varphi(n)$, a przyporządkowanie
\[
\sigma\longmapsto [k]_n\quad\text{gdzie}\quad
\sigma(\zeta_n)=\zeta_n^k
\]
wyznacza izomorfizm
\(\operatorname{Gal}(\mathbb Q(\zeta_n)/\mathbb Q)
\cong(\mathbb Z/n\mathbb Z)^\times\).
Dla $n=1$ obie grupy są trywialne.
\end{twierdzenie}
\begin{proof}
Z Twierdzenia~\ref{aa:tw:cyclotomic-irreducible} mamy
$[\mathbb Q(\zeta_n):\mathbb Q]=\deg\Phi_n=\varphi(n)$.
Wielomian $\Phi_n$ ma różne pierwiastki, gdyż dzieli
$X^n-1$ w charakterystyce zero. Wszystkie jego
pierwiastki są postaci $\zeta_n^k$ i należą do
$\mathbb Q(\zeta_n)$, więc jest to ciało rozkładu
wielomianu rozdzielczego; rozszerzenie jest Galois.

Niech $\sigma$ będzie automorfizmem tego rozszerzenia.
Ponieważ $\Phi_n(\sigma(\zeta_n))=0$,
istnieje $k$, względnie pierwszy z $n$, taki że
$\sigma(\zeta_n)=\zeta_n^k$. Klasa $[k]_n$ jest
jednoznacznie wyznaczona, bo $\zeta_n$ ma rząd $n$.
Jeżeli dwa automorfizmy mają ten sam obraz $\zeta_n$,
to zgadzają się na generatorze i na $\mathbb Q$, a
więc są równe. Odwzorowanie z tezy jest więc iniektywne.
Dla każdego $[k]_n\in(\mathbb Z/n\mathbb Z)^\times$
element $\zeta_n^k$ jest pierwiastkiem tego samego
nierozkładalnego wielomianu minimalnego $\Phi_n$.
Istnieje zatem $\mathbb Q$-zanurzenie
$\mathbb Q(\zeta_n)\to\mathbb Q(\zeta_n)$
wysyłające $\zeta_n$ na $\zeta_n^k$. Jego obraz jest
całym ciałem, gdyż $k$ jest odwracalny modulo $n$:
$\zeta_n=(\zeta_n^k)^u$ dla odpowiedniego $u$.
Jest to więc automorfizm, co dowodzi surjektywności.
Jeśli $\sigma(\zeta_n)=\zeta_n^k$ oraz
$\tau(\zeta_n)=\zeta_n^\ell$, to
$(\sigma\circ\tau)(\zeta_n)=\zeta_n^{k\ell}$.
Mnożenie klas jest zatem zgodne ze składaniem
automorfizmów, co kończy dowód.
\end{proof}

\begin{figure}[htbp]
\centering
\begin{tikzpicture}[scale=0.95]
  \draw[gray!70] (0,0) circle (1.65);
  \draw[->,gray!70] (-2.05,0)--(2.2,0);
  \draw[->,gray!70] (0,-1.98)--(0,2.02);
  \draw[blue!65!black,thick] (0:1.65)--(72:1.65)--(144:1.65)--(216:1.65)--(288:1.65)--cycle;
  \foreach \k in {0,...,4} {\filldraw[blue!65!black] (72*\k:1.65) circle (2pt);}
  \node[right] at (0:1.65) {$1=\zeta_5^5$};
  \node[above right] at (72:1.65) {$\zeta_5$};
  \node[above left] at (144:1.65) {$\zeta_5^2$};
  \node[below left] at (216:1.65) {$\zeta_5^3$};
  \node[below right] at (288:1.65) {$\zeta_5^4$};
  \node[below right] at (0,0) {$0$};
\end{tikzpicture}
\caption{Pięć pierwiastków równania $X^5=1$ na okręgu jednostkowym
w płaszczyźnie zespolonej. Cztery pierwiastki różne od $1$ są
sprzężone nad $\mathbb Q$ i są pierwiastkami $\Phi_5$.}
\label{aa:fig:cyclotomic-five}
\end{figure}

\begin{przyklad}[Ciało piątych pierwiastków z jedynki]
\label{aa:prz:cyclotomic-five}
Dla $n=5$, ponieważ jedynymi dodatnimi dzielnikami $5$
są $1$ i $5$, z Twierdzenia~\ref{aa:tw:cyclotomic-product}
wynika
\[
\Phi_5(X)=\frac{X^5-1}{X-1}
=X^4+X^3+X^2+X+1.
\]
Element $2\bmod5$ generuje grupę
$(\mathbb Z/5\mathbb Z)^\times$, gdyż jego kolejne
potęgi modulo $5$ to $2,4,3,1$.
Stąd $\operatorname{Gal}(\mathbb Q(\zeta_5)/\mathbb Q)$
jest cykliczna rzędu $4$.
Jej jedyna podgrupa rzędu $2$ to
$H=\{\sigma_1,\sigma_{-1}\}$, gdzie
$\sigma_{-1}(\zeta_5)=\zeta_5^{-1}$; przy zespolonej
realizacji jest to sprzężenie zespolone.

Połóżmy $u=\zeta_5+\zeta_5^{-1}$. Jest on stały pod
$H$. Z równania
$1+\zeta_5+\zeta_5^2+\zeta_5^3+\zeta_5^4=0$
po podzieleniu przez $\zeta_5^2$ otrzymujemy
$1+(\zeta_5+\zeta_5^{-1})
 +(\zeta_5^2+\zeta_5^{-2})=0$.
Ponieważ
$\zeta_5^2+\zeta_5^{-2}=u^2-2$,
odpowiednia równość ma postać
$u^2+u-1=0$. Wyróżnik tego wielomianu jest równy $5$,
które nie jest kwadratem liczby wymiernej; zatem
$[\mathbb Q(u):\mathbb Q]=2$.
Z odpowiedniości Galois wynika, że ciało elementów stałych $E^H$
również ma stopień $4/2=2$ nad $\mathbb Q$.
Ponieważ $\mathbb Q(u)\subseteq E^H$ i obie strony mają
stopień $2$, otrzymujemy
$E^H=\mathbb Q(\zeta_5+\zeta_5^{-1})$.
Rysunek~\ref{aa:fig:cyclotomic-five} pokazuje położenie
pierwiastków, ale obliczenie stopni i ciała stałego jest
czysto algebraiczne.
\end{przyklad}

\begin{przyklad}[Grupa Galois nie musi być cykliczna]
\label{aa:prz:cyclotomic-eight}
Dla $n=8$ pierwiastki rzędu dokładnie $8$ spełniają
$X^4=-1$, a zatem $\Phi_8(X)=X^4+1$. Rzeczywiście,
wszystkie cztery pierwiastki $X^4+1$ są pierwiastkami
$X^8-1$, lecz nie $X^4-1$, więc mają dokładny rząd $8$.
Grupa
$(\mathbb Z/8\mathbb Z)^\times=\{[1],[3],[5],[7]\}$
ma cztery elementy, przy czym
$3^2\equiv5^2\equiv7^2\equiv1\pmod8$.
Nie jest więc cykliczna rzędu $4$ (nie ma elementu
rzędu $4$), lecz jest izomorficzna z $C_2\times C_2$.
To kontrprzykład dla stwierdzenia, że każda grupa
Galois ciała cyklotomicznego jest cykliczna.
\end{przyklad}

\begin{uwaga}[Charakterystyka dodatnia]
\label{aa:uw:cyclotomic-characteristic}
Dowód nierozkładalności $\Phi_n$ dotyczył ciała
$\mathbb Q$. Nie wolno bez dodatkowych założeń twierdzić,
że redukcja $\Phi_n$ modulo liczbę pierwszą $p$ pozostaje
nierozkładalna albo rozdzielcza. Na przykład
\[
\Phi_3(X)=X^2+X+1=(X-1)^2
\qquad\text{w }\mathbb F_3[X].
\]
Równość wynika z rozwinięcia $(X-1)^2=X^2-2X+1$
i kongruencji $-2\equiv1\pmod3$.
Pochodna $X^n-1$ w charakterystyce $p$ jest
$nX^{n-1}$; rozdzielczość jest zapewniona dokładnie
wtedy, gdy $p\nmid n$. Nawet w tym przypadku
nierozkładalność $\Phi_n$ nad $\mathbb Q$ sama w sobie
nie implikuje nierozkładalności po redukcji.
\end{uwaga}

% AA-29 | Iteracja XV; szczegółowy audyt historyczny III, 20 XI 2026.
\chapter[Wielomiany nad ciałami skończonymi]{Struktura ciał skończonych i wielomiany nierozkładalne}
\label{aa:ch:finite-fields-structure}

\section{Frobenius i rozkład wielomianów}

Znajomość liczby elementów ciała skończonego daje narzędzie do
znajdowania jego elementów algebraicznych. Jeżeli ciało ma $q$
elementów, każdy jego element jest pierwiastkiem $X^q-X$.
Uogólnienie do $X^{q^n}-X$ pozwoli scharakteryzować wszystkie
moniczne wielomiany nierozkładalne, których stopnie dzielą $n$.
Należy przy tym odróżnić dwa stwierdzenia: wielomian
$X^{q^n}-X$ jest rozdzielczy, lecz dowolny wielomian
stopnia $n$ nad ciałem $q$-elementowym nie musi być rozdzielczy
ani nierozkładalny.

\begin{twierdzenie}[Pierwiastki wielomianu Frobeniusa]
\label{aa:tw:finite-frobenius-roots}
Niech $p$ będzie liczbą pierwszą, $m,n\ge1$ liczbami
całkowitymi, a $q=p^m$. W ustalonym domknięciu
algebraicznym $\overline{\mathbb F}_p$ połóżmy
\[
E_q=\{x\in\overline{\mathbb F}_p:x^q=x\}.
\]
Wówczas $E_q$ jest ciałem o dokładnie $q$ elementach.
Ponadto w pierścieniu $E_q[X]$ zachodzi równość
\[
X^{q^n}-X=\prod_{\substack{f\in E_q[X]\ \text{$f$ moniczny, nierozkładalny}\\ \deg f\mid n}}f(X),
\]
gdzie każdy czynnik występuje dokładnie jeden raz.
\end{twierdzenie}
\begin{proof}
\emph{Krok 1: ciało $E_q$.} W charakterystyce $p$
iteracja tożsamości $(x+y)^p=x^p+y^p$ daje
$(x+y)^q=x^q+y^q$ dla $q=p^m$.
Wobec tego, jeśli $x^q=x$ oraz $y^q=y$, to
$(x+y)^q=x+y$ i $(xy)^q=xy$.
Ponadto $(-x)^q=-x$ (w charakterystyce $2$ mamy
$-x=x$), a dla $x\ne0$ również
$(x^{-1})^q=(x^q)^{-1}=x^{-1}$.
Elementy $0$ i $1$ należą do $E_q$, więc wszystkie
aksjomaty podciała są spełnione.
Pochodna $X^q-X$ jest równa $-1$, ponieważ
$q=0$ w $\mathbb F_p$; wielomian jest więc rozdzielczy.
W domknięciu algebraicznym ma dokładnie $q$ różnych
pierwiastków, czyli $|E_q|=q$.

\emph{Krok 2: ciało o $q^d$ elementach.}
Niech $f\in E_q[X]$ będzie moniczny i nierozkładalny
stopnia $d\ge1$, a $\alpha\in\overline{\mathbb F}_p$
jego pierwiastkiem. Iloraz $E_q[X]/(f)$ jest ciałem
oraz ma bazę $1,\bar X,\ldots,\bar X^{d-1}$ nad
$E_q$. Wobec tego ma dokładnie $q^d$ elementów.
Ewaluacja $\bar X\mapsto\alpha$ identyfikuje ten
iloraz z $E_q(\alpha)$, zatem
$|E_q(\alpha)|=q^d$.
Jego grupa multiplikatywna ma $q^d-1$ elementów,
więc dla $\alpha\ne0$ twierdzenie Lagrange'a daje
$\alpha^{q^d-1}=1$, a stąd
$\alpha^{q^d}=\alpha$. Dla $\alpha=0$ ta sama
równość jest natychmiastowa.

\emph{Krok 3: minimalny okres Frobeniusa.}
Dla każdej liczby całkowitej $r\ge1$ oznaczmy
$\operatorname{Fr}_q^r(x)=x^{q^r}$, gdzie
$\operatorname{Fr}_q:x\mapsto x^q$ jest iteracją
potęgowania; na każdym ciele $E_{q^s}$ jest automorfizmem,
ponieważ zachowuje $E_{q^s}$ (dla $x^{q^s}=x$ mamy
$(x^q)^{q^s}=x^q$) i jest iniektywny:
z $x^q=y^q$ wynika $(x-y)^q=0$, a więc $x=y$.
Ponieważ $E_{q^s}$ jest skończone, potęgowanie przez
$q$ jest również surjektywne.
Niech $r$ będzie najmniejszą dodatnią liczbą taką, że
$\alpha^{q^r}=\alpha$. Istnieje ona i $r\le d$, na
mocy poprzedniego kroku.
Ponieważ $\alpha\in E_{q^r}$ oraz $E_q\subset E_{q^r}$
(jeśli $x^q=x$, to kolejne potęgowania przez $q$
także pozostawiają $x$ nieruchomy),
$E_q(\alpha)$ jest podciałem $E_{q^r}$.
Porównanie liczności podciała i ciała daje
$q^d\le q^r$, czyli $d\le r$. Łącznie z
$r\le d$ wynika $r=d$.

Teraz wykażemy, że
$\alpha^{q^n}=\alpha$ wtedy i tylko wtedy, gdy $d\mid n$.
Kierunek $d\mid n$ wynika z wielokrotnego zastosowania
$\operatorname{Fr}_q^d(\alpha)=\alpha$.
W kierunku odwrotnym podzielmy $n$ z resztą przez $d$:
$n=ad+s$ dla $a\ge0$ i $0\le s<d$.
Z uwagi na $\operatorname{Fr}_q^d(\alpha)=\alpha$
mamy $\operatorname{Fr}_q^n(\alpha)=
\operatorname{Fr}_q^s(\alpha)$.
Jeżeli lewa strona jest równa $\alpha$, to minimalność
$d$ wymusza $s=0$, czyli $d\mid n$.

\emph{Krok 4: identyfikacja wszystkich czynników.}
Wielomian moniczny nierozkładalny $f$ dzieli
$X^{q^n}-X$ dokładnie wtedy, gdy jego pierwiastek
$\alpha$ spełnia $\alpha^{q^n}=\alpha$, a więc
wtedy i tylko wtedy, gdy $\deg f=d$ dzieli $n$.
Pochodna $X^{q^n}-X$ jest równa $-1$, toteż
w jego rozkładzie nie ma powtórzonych czynników.
Ponieważ każdy niestały wielomian nad ciałem rozkłada się
na czynniki nierozkładalne, otrzymujemy wskazaną
faktoryzację.
\end{proof}

\begin{wniosek}[Tożsamość stopni]\label{aa:cor:finite-degree-count}
Niech $q=p^m$, gdzie $p$ jest liczbą pierwszą i $m\ge1$.
Dla każdej liczby całkowitej $d\ge1$ oznaczmy przez
$N_q(d)$ liczbę monicznych wielomianów nierozkładalnych
stopnia $d$ w $E_q[X]$. Dla każdego $n\ge1$ zachodzi
\[
q^n=\sum_{d\mid n}d\,N_q(d).
\]
\end{wniosek}
\begin{proof}
Po lewej stronie faktoryzacji z
Twierdzenia~\ref{aa:tw:finite-frobenius-roots} znajduje się
wielomian stopnia $q^n$. Dla każdego dzielnika $d\mid n$
po prawej stronie występuje dokładnie $N_q(d)$ czynników
stopnia $d$. Stopień iloczynu niezerowych wielomianów
nad ciałem jest sumą stopni czynników, skąd wzór.
\end{proof}

\section{Odwrócenie Möbiusa i liczba wielomianów}

Potrzebujemy teraz odwrócić sumę po dzielnikach $n$.
Umożliwia to funkcja Möbiusa: jest ona dobrana tak, aby
jej suma po wszystkich dzielnikach była zerowa poza
przypadkiem $n=1$.

\begin{definicja}\label{aa:def:mobius}
Dla każdej liczby całkowitej $n\ge1$ definiujemy
funkcję Möbiusa $\mu(n)$ następująco:
$\mu(1)=1$; jeśli dla pewnej liczby pierwszej $p$
zachodzi $p^2\mid n$, to $\mu(n)=0$;
w pozostałym przypadku $n>1$ jest iloczynem $r\ge1$
\emph{różnych} liczb pierwszych i $\mu(n)=(-1)^r$.
\end{definicja}
\begin{lemat}\label{aa:lem:mobius-sum}
Dla każdej liczby całkowitej $n\ge1$ mamy
\[
\sum_{d\mid n}\mu(d)=
\begin{cases}1,&n=1,\\0,&n>1.\end{cases}
\]
\end{lemat}
\begin{proof}
Dla $n=1$ jedynym dzielnikiem dodatnim jest $1$.
Jeżeli $n>1$, niech $p_1,\ldots,p_r$ będą jego
różnymi dzielnikami pierwszymi, gdzie $r\ge1$.
Składnik $\mu(d)$ jest niezerowy dokładnie dla tych
dzielników $d$, które są iloczynami pewnego podzbioru
zbioru $\{p_1,\ldots,p_r\}$, w tym podzbioru pustego.
Podzbiór $j$-elementowy wnosi $(-1)^j$, a takich
podzbiorów jest $\binom rj$. Otrzymujemy
$\sum_{d\mid n}\mu(d)
=\sum_{j=0}^r\binom rj(-1)^j=(1-1)^r=0$.
\end{proof}

\begin{twierdzenie}[Wzór na liczbę wielomianów nierozkładalnych]
\label{aa:tw:finite-irreducible-count}
Niech $p$ będzie liczbą pierwszą, $m,n\ge1$ liczbami
całkowitymi, $q=p^m$, a $E_q$ ciałem o $q$ elementach.
Niech $N_q(n)$ oznacza liczbę monicznych nierozkładalnych
wielomianów stopnia $n$ w $E_q[X]$. Wówczas
\[
N_q(n)=\frac1n\sum_{d\mid n}\mu(d)q^{n/d}.
\]
Ponadto $N_q(n)>0$: nad $E_q$ istnieje moniczny
wielomian nierozkładalny każdego dodatniego stopnia $n$.
\end{twierdzenie}
\begin{proof}
\emph{Krok 1: odwrócenie sumy po dzielnikach.}
Z Wniosku~\ref{aa:cor:finite-degree-count} dla dowolnego
$k\ge1$ mamy $q^k=\sum_{e\mid k}eN_q(e)$.
Wstawiając tę tożsamość do sumy z tezy, uzyskujemy
\[
\begin{aligned}
\sum_{d\mid n}\mu(d)q^{n/d}
&=\sum_{d\mid n}\mu(d)\sum_{e\mid n/d}eN_q(e)\\
&=\sum_{e\mid n}eN_q(e)\sum_{d\mid n/e}\mu(d)\\
&=nN_q(n).
\end{aligned}
\]
Druga równość jest dopuszczalna, ponieważ sumy są
skończone, a para $(d,e)$ spełnia $de\mid n$ dokładnie
wtedy, gdy $e\mid n$ i $d\mid n/e$. Ostatnia równość
wynika z Lematu~\ref{aa:lem:mobius-sum}: suma wewnętrzna
jest równa $1$ tylko dla $e=n$.
Dzieląc obie strony przez dodatnią liczbę $n$,
otrzymujemy żądany wzór.

\emph{Krok 2: istnienie nierozkładalnego wielomianu.}
Dla $n=1$ wystarczy moniczny wielomian liniowy $X$.
Jest on nierozkładalny nad każdym ciałem.
Niech $n\ge2$. Ciało $E_{q^n}$ ma $q^n$ elementów
i zawiera $E_q$. Rozważmy zbiór tych elementów
$\alpha\in E_{q^n}$, których stopień
$[E_q(\alpha):E_q]$ jest właściwym dzielnikiem $n$.
Jest to rzeczywiście dzielnik $n$ na mocy twierdzenia
o wieży rozszerzeń, bo $E_q\subseteq E_q(\alpha)
\subseteq E_{q^n}$ i $[E_{q^n}:E_q]=n$: ostatnia
równość wynika z $|E_{q^n}|=q^n$ i $|E_q|=q$.
Jeśli ten stopień jest równy $d$, to
$\alpha\in E_{q^d}$, ponieważ każde ciało o $q^d$
elementach spełnia równanie $x^{q^d}=x$.
Zatem liczba wszystkich takich elementów jest
co najwyżej
\[
\sum_{\substack{d\mid n\\d<n}}q^d
\le \sum_{j=1}^{\lfloor n/2\rfloor}q^j
<q^n.
\]
Pierwsza nierówność wynika z faktu, że każdy właściwy
dodatni dzielnik $n$ jest co najwyżej $n/2$.
Druga jest ścisła: dla $k\ge1$ i $q\ge2$
\(\sum_{j=1}^{k}q^j
=q(q^k-1)/(q-1)<q^{k+1}\), a
$\lfloor n/2\rfloor+1\le n$.
Istnieje wobec tego $\alpha\in E_{q^n}$, którego
stopień nad $E_q$ nie jest właściwym dzielnikiem $n$.
Jest to z konieczności stopień $n$. Wielomian minimalny
$\alpha$ nad $E_q$ jest moniczny, nierozkładalny i ma
stopień $n$. Stąd $N_q(n)\ge1$.
\end{proof}

\begin{przyklad}[Stopień dwa nad $\mathbb F_2$]
\label{aa:prz:finite-quadratic}
Dla $q=2$ i $n=2$ wzór daje
$N_2(2)=(2^2-2)/2=1$. Wielomian
$f=X^2+X+1$ nie ma pierwiastka w $\mathbb F_2$:
$f(0)=1$, $f(1)=1+1+1=1$.
Wielomian stopnia $2$ bez pierwiastka jest
nierozkładalny, ponieważ ewentualny rozkład na
czynniki niestałe musiałby być iloczynem dwóch
wielomianów liniowych. Istnieje więc dokładnie jeden
moniczny nierozkładalny wielomian kwadratowy.
Iloraz $\mathbb F_2[X]/(f)$ jest ciałem i jako
$\mathbb F_2$-przestrzeń liniowa ma bazę $1,x$,
gdzie $x$ oznacza klasę $X$. Jego elementami są
$0,1,x,x+1$. Równość $f(x)=0$ daje
$x^2=x+1$, a więc $x^3=x(x+1)=x^2+x=1$.
Dla kontroli rozkładu Frobeniusa mamy
\[
X^4-X=X(X+1)(X^2+X+1)
\quad\text{w }\mathbb F_2[X].
\]
Iloczyn po prawej stronie ma stopień $4$ i można go
sprawdzić przez wymnożenie; odpowiada liniowym
czynnikom stopnia $1$ i jednemu czynnikowi stopnia $2$.
\end{przyklad}

\begin{przyklad}[Kontrprzykład: nie każdy wielomian stopnia $n$ jest nierozkładalny]
\label{aa:prz:finite-reducible}
W $\mathbb F_2[X]$ zachodzi
$X^4+X^2+1=(X^2+X+1)^2$. Jest to wielomian
rozkładalny, którego pochodna wynosi zero.
W pierścieniu
$R=\mathbb F_2[X]/(X^4+X^2+1)$
klasa $v$ wielomianu $X^2+X+1$ spełnia $v^2=0$.
Jest niezerowa: gdyby należała do ideału generowanego
przez wielomian stopnia $4$, byłaby równa zeru albo
miałaby stopień co najmniej $4$, podczas gdy ma stopień
$2$. Zatem $R$ zawiera niezerowy element nilpotentny.
Żadne ciało nie ma takiego elementu, bo z $v^2=0$
i $v\ne0$ można by pomnożyć przez $v^{-1}$ i uzyskać
$v=0$. Chociaż $R$ ma $2^4=16$ elementów jako
czterowymiarowa przestrzeń nad $\mathbb F_2$, nie jest
ciałem szesnastoelementowym.
\end{przyklad}

\begin{przyklad}[Wielomiany trzeciego stopnia nad $\mathbb F_2$]
\label{aa:prz:finite-cubics}
Dla $n=3$ wzór Möbiusa daje
$N_2(3)=(2^3-2)/3=2$.
Sprawdzimy to niezależnie. Wielomian stopnia $3$
nad ciałem jest nierozkładalny dokładnie wtedy,
gdy nie ma pierwiastka w tym ciele: jeśli jest
rozkładalny, stopnie jego niestałych czynników sumują
się do $3$, więc jeden z nich ma stopień $1$.
Każdy moniczny wielomian sześcienny w $\mathbb F_2[X]$
ma postać $X^3+aX^2+bX+c$ dla $a,b,c\in\{0,1\}$.
Aby nie miał pierwiastka $0$, musi być $c=1$.
Wtedy jego wartość dla $X=1$ to
$1+a+b+1=a+b$, więc brak pierwiastka $1$ wymaga
$a+b=1$. Są dokładnie dwie pary współczynników
$(a,b)$ spełniające tę równość: $(0,1)$ i $(1,0)$.
Otrzymujemy wielomiany
$X^3+X+1$ i $X^3+X^2+1$, oba nierozkładalne.
Wszystkie pozostałe wielomiany stopnia $3$ mają
pierwiastek w $\mathbb F_2$, a zatem są rozkładalne.
Istotnie, zgodnie z Twierdzeniem~\ref{aa:tw:finite-frobenius-roots},
\[
X^8-X=X(X+1)(X^3+X+1)(X^3+X^2+1)
\quad\text{w }\mathbb F_2[X].
\]
\end{przyklad}

\begin{przyklad}[Trzy wielomiany kwadratowe nad $\mathbb F_3$]
\label{aa:prz:finite-quadratics-three}
Wzór dla $q=3$ i $n=2$ daje
$N_3(2)=(9-3)/2=3$.
Moniczny wielomian kwadratowy ma postać
$X^2+aX+b$ z $a,b\in\mathbb F_3$.
Jest nierozkładalny dokładnie wtedy, gdy jego
wyróżnik $\Delta=a^2-4b=a^2-b$ nie jest kwadratem
w $\mathbb F_3$: istotnie, ponieważ $2$ jest
odwracalne, równanie $X^2+aX+b=0$ jest równoważne
$(2X+a)^2=\Delta$. Ponieważ kwadratami są tylko $0$ i $1$,
warunek ten jest równoważny $\Delta=2$.
Jeśli $a=0$, dostajemy $b=1$ i $X^2+1$.
Jeśli $a=1$ lub $a=2$, to $a^2=1$, więc $b=2$,
co daje $X^2+X+2$ oraz $X^2+2X+2$.
W ten sposób niezależnie znaleźliśmy wszystkie trzy
wielomiany nierozkładalne stopnia $2$ nad $\mathbb F_3$.
\end{przyklad}

\begin{uwaga}\label{aa:uw:finite-field-previous}
Istnienie i jednoznaczność ciał skończonych były już
omawiane w Rozdziale~\ref{aa:ch:embeddings-normal-separable}.
Tutaj rozwinęliśmy rozkład wielomianów Frobeniusa,
ścisłe wyliczanie liczby wielomianów nierozkładalnych
oraz konstrukcję elementów stopnia dokładnie $n$.
\end{uwaga}

% AA-30 | Iteracja XV
\chapter[Konstrukcje i wielokąty foremne]{Konstrukcje cyrklem i linijką oraz wielokąty foremne}
\label{aa:ch:constructibility}
\section{Algebraiczny opis konstrukcji}

W klasycznej geometrii wolno nam wykonywać wyłącznie skończenie
wiele rysunków prostych i okręgów. Trudność polega na wykazaniu,
że nawet bardzo sprytny wybór kolejnych przecięć nie pozwala
otrzymać wszystkich liczb rzeczywistych. Opiszemy algorytm
geometryczny za pomocą kolejnych rozszerzeń ciał stopnia
co najwyżej dwa. W ten sposób zagadnienia dotyczące
podwajania sześcianu, trysekcji kąta oraz wielokątów
foremnych sprowadzimy do precyzyjnych pytań algebraicznych.

Przyjmujemy ciało liczb rzeczywistych $\mathbb R$ jako już
dane; nie konstruujemy go w tym skrypcie. Punkty płaszczyzny
identyfikujemy z liczbami zespolonymi $x+iy$, gdzie
$x,y\in\mathbb R$ i $i^2=-1$.
\begin{definicja}\label{aa:def:constructible}
Niech $S\subset\mathbb C$ będzie skończonym zbiorem punktów
zawierającym $0$ i $1$. Punkt $z\in\mathbb C$ nazywamy
\emph{konstruowalnym z $S$ cyrklem i linijką}, jeśli istnieje
skończony ciąg zbiorów punktów
$S=S_0\subseteq S_1\subseteq\cdots\subseteq S_m$
taki, że $z\in S_m$, a każdy zbiór $S_j$ powstaje
z $S_{j-1}$ przez dołączenie jednego z \emph{izolowanych}
punktów przecięcia dwóch \emph{różnych} obiektów
następujących rodzajów:
\begin{enumerate}
\item prostych wyznaczonych przez dwa różne punkty
      z $S_{j-1}$;
\item okręgów o środku w punkcie z $S_{j-1}$
      i promieniu równym odległości dwóch punktów
      z $S_{j-1}$ (promień dodatni).
\end{enumerate}
Zaznaczenie izolowanego przecięcia jest dozwolone tylko,
jeżeli przecięcie istnieje. Nie wolno \emph{dowolnie} wybierać
punktu z dwóch pokrywających się prostych lub okręgów:
taki wybór nie stanowi skończonej konstrukcji.
Punkty już narysowane mogą być ponownie wykorzystywane;
wybór jednego z dwóch izolowanych przecięć jest dozwolony.
\end{definicja}
\begin{lemat}[Operacje arytmetyczne w konstrukcjach]
\label{aa:lem:constructibility-field-operations}
Jeżeli liczby rzeczywiste $a,b$ są konstruowalne z $\{0,1\}$,
to konstruowalne są także $a+b$, $a-b$ i $ab$.
Jeżeli ponadto $b\ne0$, konstruowalna jest liczba $a/b$.
Dla $a\ge0$ konstruowalna jest również liczba $\sqrt a\ge0$.
\end{lemat}
\begin{proof}
\textbf{Oś i współrzędne.}
Prosta przez $0,1$ jest osią rzeczywistą. Okrąg
o środku $0$ i promieniu $1$ przecina tę oś
w punktach $1$ i $-1$. Dwa okręgi o promieniu $2$,
równym długości już skonstruowanego odcinka $[-1,1]$,
których środkami są
$1$ i $-1$, przecinają się w dwóch punktach
leżących po przeciwnych stronach osi. Prosta przez
te przecięcia jest symetralną odcinka $[-1,1]$,
ponieważ oba punkty są jednakowo odległe od $1$
i $-1$. Jest to oś prostopadła do osi rzeczywistej
w punkcie $0$. Jej przecięcie z okręgiem jednostkowym
pozwala wybrać punkt $i=(0,1)$.

Ten sam algorytm zastosowany do dwóch punktów
na dowolnej skonstruowanej prostej daje symetralną
ich odcinka. Aby skonstruować prostopadłą przez
punkt $P$ należący do prostej, zaznaczamy dwa
punkty na tej prostej jednakowo odległe od $P$
i konstruujemy ich symetralną. Dla punktu poza
prostą najpierw konstruujemy jego rzut prostopadły
na prostą: okrąg o środku w $P$ przecina prostą
w dwóch punktach (wybieramy dostatecznie duży
konstruowalny promień, np. kolejne podwojenia długości
ustalonego odcinka aż do przekroczenia odległości
$P$ od prostej), a ich symetralna
przechodzi przez $P$. Konstruując następną
prostopadłą do uzyskanej, otrzymujemy równoległą
do pierwszej prostej przechodzącą przez zadany punkt.
Przenoszenie odcinków na osie odbywa się za pomocą
okręgów o odpowiednich promieniach; wybór jednego
z dwóch przecięć ustala znak współrzędnej.
Przykładanie kolejnych odcinków zorientowanych
na tej samej osi daje dodawanie i odejmowanie.

\textbf{Mnożenie.}
Współrzędne $1$, $a$ i $b$ możemy zaznaczyć
na osi rzeczywistej; korzystając z prostopadłych,
umieszczamy także punkt $(0,b)$. Rysujemy prostą
przez punkty $(1,0)$ oraz $(0,b)$, o równaniu
$y=b(1-x)$. Prosta przez $(a,0)$ równoległa
do niej ma równanie $y=b(a-x)$ i przecina
oś pionową w $(0,ab)$. Aby przenieść tę podpisaną
współrzędną na oś poziomą,
konstruujemy prostą $y=x$ przez $0$ i $(1,1)$; ten drugi
punkt jest przecięciem pionowej przez $1$ i poziomej przez $i$.
Odbicie $(0,ab)$ względem prostej $y=x$ jest konstrukcją
cyrklem i linijką (przez rzut prostopadły i odłożenie
tej samej odległości po przeciwnej stronie prostej)
i daje punkt $(ab,0)$.

\textbf{Dzielenie.}
Tym samym twierdzeniem o podobieństwie trójkątów,
stosując prostą o nachyleniu zadanym przez
punkty $(b,0)$, $(0,a)$ i prostą równoległą
przechodzącą przez $(1,0)$, otrzymujemy
na osi pionowej współrzędną $a/b$.
Rzeczywiście nachylenie pierwszej prostej wynosi
$-a/b$, więc równoległa przechodząca przez $(1,0)$
ma punkt przecięcia z osią pionową $(0,a/b)$.
Założenie $b\ne0$ gwarantuje, że użyty odcinek
na osi poziomej jest niezerowy.

\textbf{Pierwiastek kwadratowy.}
Dla $a>0$ kreślimy okrąg o średnicy
$[0,1+a]$ na osi rzeczywistej, a następnie
prostą prostopadłą do osi w punkcie $1$.
Jeśli $(1,y)$ należy do okręgu, to jego równanie daje
\[
 \left(1-\frac{1+a}{2}\right)^2+y^2
 =\left(\frac{1+a}{2}\right)^2,
 \qquad y^2=a.
\]
Wybieramy przecięcie z $y>0$, otrzymując $\sqrt a$.
Dla $a=0$ pierwiastek wynosi $0$ i jest już dany.
W każdej konstrukcji użyto tylko dozwolonych
izolowanych przecięć prostych i okręgów.
\end{proof}

\begin{twierdzenie}[Kryterium wieży kwadratowej]\label{aa:tw:constructibility-tower}
Dla każdego $z=x+iy\in\mathbb C$, gdzie $x,y\in\mathbb R$,
następujące warunki są równoważne:
\begin{enumerate}
\item punkt $z$ jest konstruowalny z $\{0,1\}$;
\item istnieje liczba całkowita $r\ge0$ i wieża podciał
$\mathbb Q=K_0\subseteq K_1\subseteq\cdots\subseteq K_r
\subseteq\mathbb R$ taka, że
$[K_j:K_{j-1}]\le2$ dla każdego $1\le j\le r$
oraz $x,y\in K_r$.
\end{enumerate}
\end{twierdzenie}
\begin{proof}
\textbf{Od konstrukcji do wieży.}
Współrzędne punktów początkowych $0,1$ należą do
$K_0=\mathbb Q$. Załóżmy, że wszystkie otrzymane
punkty mają współrzędne w ustalonym podciele
$K\subseteq\mathbb R$. Prosta przez dwa różne takie
punkty ma równanie $Ax+By=C$, gdzie $A,B,C\in K$
oraz $(A,B)\ne(0,0)$. Okrąg o środku $(a,b)$
i promieniu równym odległości znanych punktów
ma równanie $(x-a)^2+(y-b)^2=R$, gdzie
$a,b,R\in K$, $R>0$.

Dwie różne proste mają co najwyżej jeden punkt
przecięcia. Jeżeli go mają, układ dwóch równań
liniowych ma niezerowy wyznacznik, więc obie
współrzędne punktu należą do $K$.
W przypadku przecięcia prostej z okręgiem
parametryzujemy prostą jako
$(x,y)=(x_0,y_0)+t(u,v)$, gdzie
$x_0,y_0,u,v\in K$ i $(u,v)\ne(0,0)$.
Po podstawieniu do równania okręgu dostajemy
$At^2+Bt+C=0$ z $A=u^2+v^2>0$ i
$A,B,C\in K$. Każdy parametr $t$ wybranego
punktu przecięcia ma stopień co najwyżej $2$
nad $K$; obie współrzędne należą do tego samego
ciała $K(t)$.

Jeżeli przecinamy dwa okręgi o różnych środkach,
to po odjęciu ich równań otrzymujemy równanie
prostej o niezerowym wektorze współczynników
w $K^2$. Punkty przecięcia obu okręgów
należą również do tej prostej, więc wracamy
do poprzedniego przypadku. Okręgi współśrodkowe
i różne nie mają punktów wspólnych; okręgi
pokrywające się nie dają dozwolonego izolowanego
przecięcia. W każdym kroku współrzędne nowego
punktu należą zatem do jednego rzeczywistego
rozszerzenia stopnia $1$ lub $2$. Dołączając te
rozszerzenia kolejno, otrzymujemy wymaganą wieżę.

\textbf{Od wieży do konstrukcji.}
Indukujemy po liczbie poziomów wieży. Każda liczba
w $K_0=\mathbb Q$ jest konstruowalna dzięki
Lematowi~\ref{aa:lem:constructibility-field-operations}.
Przypuśćmy, że skonstruowaliśmy wszystkie liczby
$K_{j-1}$. Jeśli $K_j=K_{j-1}$, niczego nie
trzeba dodawać. W przeciwnym razie rozszerzenie
$K_j/K_{j-1}$ jest kwadratowe, więc istnieje
$u\in K_j$ takie, że $K_j=K_{j-1}(u)$, a jego
wielomian minimalny ma postać
$T^2+bT+c$ dla $b,c\in K_{j-1}$.
Ponieważ $u\in\mathbb R$, wyróżnik
$d=b^2-4c=(2u+b)^2$ jest nieujemny.
Z lematu potrafimy skonstruować
$\sqrt d$ i obie wartości $(-b\pm\sqrt d)/2$,
w tym wybrane $u$. Każdy element $K_j$ ma postać
$a+bu$ dla pewnych $a,b\in K_{j-1}$,
więc jest konstruowalny z indukcji i lematu.
Po $r$ krokach konstruowalne są $x$ i $y$,
a przecięcie prostej pionowej przez $(x,0)$
z poziomą przez $(0,y)$ daje punkt $(x,y)$.
\end{proof}

\begin{wniosek}[Przeszkoda stopnia]\label{aa:cor:constructibility-degree}
Jeżeli liczba rzeczywista $a$ jest konstruowalna z $\{0,1\}$, to stopień $[\mathbb Q(a):\mathbb Q]$ jest potęgą liczby $2$.
\end{wniosek}
\begin{proof}
Z Twierdzenia~\ref{aa:tw:constructibility-tower} istnieje
wieża kwadratowa zawierająca $a$. Prawo wieży daje
$[K_r:\mathbb Q]=\prod_{j=1}^r[K_j:K_{j-1}]=2^m$
dla pewnego $m\ge0$: każdy czynnik jest równy $1$ lub $2$.
Ponieważ $\mathbb Q\subseteq\mathbb Q(a)\subseteq K_r$,
ponowne zastosowanie prawa wieży daje
$[K_r:\mathbb Q]=[K_r:\mathbb Q(a)]\,[\mathbb Q(a):\mathbb Q]$.
Dodatni dzielnik potęgi $2$ również jest potęgą $2$.
\end{proof}
\begin{lemat}[Domknięcie normalne wieży kwadratowej]
\label{aa:lem:quadratic-tower-normal-closure}
Niech $K$ będzie ciałem charakterystyki różnej od $2$ i
$K=K_0\subseteq K_1\subseteq\cdots\subseteq K_r$
skończoną wieżą rozszerzeń, w której
$[K_j:K_{j-1}]\le2$ dla każdego $j$. Wówczas
normalne domknięcie $K_r/K$ jest skończonym rozszerzeniem
Galois, którego grupa Galois ma rząd będący potęgą $2$.
\end{lemat}
\begin{proof}
Postępujemy indukcyjnie po liczbie poziomów wieży.
Dla $r=0$ domknięcie normalne jest równe $K$.
Przypuśćmy, że istnieje skończone rozszerzenie Galois
$E/K$ zawierające $K_{r-1}$ i takie, że
$\operatorname{Gal}(E/K)$ jest $2$-grupą.
Ponieważ charakterystyka jest różna od $2$, każde
właściwe rozszerzenie kwadratowe ma postać
$K_r=K_{r-1}(\sqrt d)$ dla pewnego
$d\in K_{r-1}^{\times}$. Gdy stopień jest równy $1$,
wystarczy przyjąć $E$ bez zmian.
Niech $G=\operatorname{Gal}(E/K)$ i niech
\[
 F=E\bigl(\sqrt{\sigma(d)}:\sigma\in G\bigr)
\]
we wspólnym domknięciu algebraicznym. Każde kolejne
dołączenie pierwiastka kwadratowego zwiększa stopień
co najwyżej dwukrotnie, a $F/E$ jest rozdzielcze i normalne,
ponieważ jest ciałem rozkładu rodziny wielomianów
$T^2-\sigma(d)$. Zatem $\operatorname{Gal}(F/E)$
jest $2$-grupą (jego rząd równa się $[F:E]$).
Ponadto każdy $K$-automorfizm domknięcia algebraicznego
zachowuje $E$ i permutuje wielomiany $T^2-\sigma(d)$,
więc zachowuje ich ciało rozkładu $F$. Stąd $F/K$ jest
normalne. Jest też rozdzielcze, a więc Galois.
Restrykcja automorfizmów daje dokładny ciąg
\[
1\longrightarrow\operatorname{Gal}(F/E)
\longrightarrow\operatorname{Gal}(F/K)
\longrightarrow\operatorname{Gal}(E/K)
\longrightarrow1,
\]
co wynika z normalności obu rozszerzeń. Rząd grupy
środkowej jest iloczynem rzędów skrajnych, czyli potęgą $2$.
Ponieważ $K_r\subseteq F$, jego domknięcie normalne nad $K$
zawiera się w $F$ i ma grupę Galois będącą ilorazem
podgrupy odpowiedniej grupy Galois $F/K$, a więc także
$2$-grupą. To kończy indukcję.
\end{proof}

\begin{uwaga}[Odwrotność jest fałszywa]\label{aa:uw:constructibility-degree-not-sufficient}
Sam fakt, że liczba algebraiczna ma stopień równy potędze $2$,
\emph{nie} zapewnia jej konstruowalności. W celu podania
sprawdzalnego przykładu rozważmy
$f(X)=X^4-X-1\in\mathbb Q[X]$.

\textbf{Nierozkładalność.}
Modulo $2$ otrzymujemy $\overline f=X^4+X+1$.
Nie ma on pierwiastka w $\mathbb F_2$, bo wartości
dla $0$ i $1$ są równe $1$. Jedynym monicznym
nierozkładalnym wielomianem kwadratowym nad
$\mathbb F_2$ jest $q=X^2+X+1$.
W ilorazie $\mathbb F_2[X]/(q)$ mamy $X^2=X+1$
i $X^4=X$, skąd $\overline f\equiv1\pmod q$.
Nie istnieje więc podział $\overline f$ na czynniki
stopni $1+3$ ani $2+2$, co dowodzi nierozkładalności
$\overline f$, a następnie $f$ nad $\mathbb Q$.
Każdy pierwiastek $f$ ma zatem stopień $4$ nad $\mathbb Q$.

\textbf{Działanie Galois na trzech podziałach.}
Niech $a,b,c,d$ oznaczają cztery różne pierwiastki
$f$ w ciele rozkładu $L$. Są one różne,
ponieważ $f$ jest nierozkładalny nad $\mathbb Q$,
a jego pochodna $f'=4X^3-1$ jest niezerowa i ma stopień
mniejszy niż $f$, wielomiany $f$ i $f'$ są względnie
pierwsze. W charakterystyce $0$ oznacza to, że $f$
jest rozdzielczy.
Definiujemy trzy liczby
\[
 u_1=ab+cd,\qquad u_2=ac+bd,\qquad u_3=ad+bc.
\]
Każda permutacja pierwiastków permutuje te liczby:
reprezentują one trzy podziały zbioru czterech
indeksów na dwie nieuporządkowane pary.
Oznaczmy przez $e_j$ elementarną sumę iloczynów
$j$ różnych elementów zbioru $\{a,b,c,d\}$.
Porównując współczynniki w tożsamości
\[
\prod_{t\in\{a,b,c,d\}}(X-t)
=X^4-e_1X^3+e_2X^2-e_3X+e_4=X^4-X-1,
\]
otrzymujemy $e_1=e_2=0$, $e_3=1$, $e_4=-1$.
Bezpośrednie rozwinięcie trzech iloczynów daje
\[
 u_1+u_2+u_3=e_2=0,
\]
\[
 u_1u_2+u_1u_3+u_2u_3=e_1e_3-4e_4=4,
\]
\[
 u_1u_2u_3=e_3^2+e_1^2e_4-4e_2e_4=1.
\]
Stąd $u_1,u_2,u_3$ są pierwiastkami wielomianu
$R(Y)=Y^3+4Y-1$. Ponieważ jedynymi możliwymi
pierwiastkami wymiernymi monicznego wielomianu
$R$ są $1$ i $-1$, a $R(1)=4$ i $R(-1)=-6$,
$R$ jest nierozkładalny nad $\mathbb Q$.
Jest też rozdzielczy: jego pochodna formalna
$R'=3Y^2+4$ jest niezerowym wielomianem stopnia $2$,
a $R$ jest nierozkładalny stopnia $3$ nad ciałem
charakterystyki $0$. Zatem $R$ i $R'$ nie mają
wspólnego nietrywialnego dzielnika, a w konsekwencji
$R$ nie ma pierwiastka wielokrotnego.
Można również sprawdzić, że wyróżnik $R$ wynosi $-283$.
Jego grupa Galois działa zatem przechodnio na
$u_1,u_2,u_3$. W rezultacie grupa
$G=\operatorname{Gal}(L/\mathbb Q)$ ma rząd
podzielny przez $3$ (Twierdzenie~\ref{aa:tw:orbita-stabilizator}).
Działanie $G$ na czterech pierwiastkach $f$
również jest przechodnie, więc $4$ dzieli $|G|$.
Ponieważ $G\subseteq S_4$, jej rząd dzieli $24$;
jedynymi możliwościami są $|G|=12$ lub $|G|=24$.

\textbf{Wykluczenie rzędu $12$.}
Wielomian $f$ ma dokładnie dwa pierwiastki rzeczywiste:
pochodna $f'=4X^3-1$ ma jedno rzeczywiste miejsce
zerowe $x_0>0$. Funkcja $f$ jest ściśle malejąca
na $(-\infty,x_0)$, rosnąca na $(x_0,\infty)$,
a $f(0)=-1$ i $f(x)\to+\infty$ dla $x\to\pm\infty$.
Pozostałe dwa pierwiastki tworzą parę sprzężoną
zespolenie. Sprzężenie zespolone jest automorfizmem
$L/\mathbb Q$, który ustala dwa rzeczywiste
pierwiastki i zamienia dwa nierzeczywiste, czyli
jest transpozycją w $G$.
Gdyby $|G|=12$, $G$ byłaby podgrupą indeksu $2$
w $S_4$, a więc normalną. Wszystkie transpozycje
w $S_4$ są sprzężone i generują $S_4$.
Niezerowy homomorfizm $S_4\to C_2$ musi zatem
wysyłać każdą transpozycję na niezerowy element,
więc jego jądro jest dokładnie $A_4$.
Taka podgrupa nie może zawierać transpozycji.
Dlatego $|G|=24$ i $G=S_4$.

\textbf{Dodatkowa kontrola arytmetyczna.}
Przy redukcji $f$ modulo $7$ otrzymujemy rozkład
\[
 X^4-X-1\equiv (X-3)(X^3+3X^2+2X-2)\pmod 7.
\]
Sześcienny czynnik nie ma pierwiastka w $\mathbb F_7$:
jego wartości dla $0,1,2,3,4,5,6$ wynoszą
odpowiednio $5,4,1,2,6,5,5$. Jest zatem nierozkładalny.
Z kolei wyróżnik wielomianu $f$ jest równy $-283$,
co zgadza się z wyróżnikiem obliczonego wyżej
wielomianu $R$. Te dwa obliczenia dają niezależną
kontrolę danych przykładu, choć nie były potrzebne
do przedstawionego dowodu, że $G=S_4$.

\textbf{Konsekwencja dla konstrukcji.}
Gdyby choć jeden rzeczywisty pierwiastek $a$ wielomianu
$f$ był konstruowalny, należałby do wieży rozszerzeń
kwadratowych nad $\mathbb Q$.
Z Lematu~\ref{aa:lem:quadratic-tower-normal-closure} wynika,
że jej domknięcie normalne ma grupę Galois będącą
$2$-grupą. Ponieważ wielomian minimalny $a$ to $f$,
jego ciało rozkładu $L$ mieści się w tym domknięciu,
a grupa $\operatorname{Gal}(L/\mathbb Q)$ jest ilorazem
odpowiedniej $2$-grupy. Jest to niemożliwe,
ponieważ $|S_4|=24$ nie jest potęgą $2$.
Otrzymujemy więc konkretną liczbę algebraiczną stopnia
$4$, która nie jest konstruowalna cyrklem i linijką.
\end{uwaga}
\section{Zastosowania i granice metody}
\begin{przyklad}[Podwojenie sześcianu]\label{aa:prz:double-cube}
Podwojenie objętości sześcianu o krawędzi $1$ wymaga skonstruowania $\sqrt[3]{2}$. Wielomian $X^3-2$ jest nierozkładalny nad $\mathbb Q$ (kryterium Eisensteina dla liczby pierwszej $2$), zatem stopień tej liczby wynosi $3$, co przeczy Wnioskowi~\ref{aa:cor:constructibility-degree}.
\end{przyklad}
\begin{przyklad}[Trysekcja kąta]\label{aa:prz:trisect-angle}
Nie istnieje ogólna konstrukcja trysekcji kąta. Dla kąta $60^\circ$ trysekcja wymagałaby skonstruowania $c=2\cos20^\circ$. Tożsamość $2\cos3\theta=(2\cos\theta)^3-3(2\cos\theta)$ daje $c^3-3c-1=0$. Wielomian $X^3-3X-1$ nie ma pierwiastka wymiernego (jedynymi kandydatami są $\pm1$), jest więc nierozkładalny stopnia $3$. Nie oznacza to, że żaden kąt nie jest podzielny na trzy równe części: na przykład $90^\circ$ można podzielić na trzy kąty $30^\circ$.
\end{przyklad}
W poniższym twierdzeniu przez konstrukcję foremnego $n$-kąta
rozumiemy skonstruowanie jego wszystkich wierzchołków
na okręgu jednostkowym o środku $0$, z jednym wierzchołkiem
w punkcie $1$. Jest to równoważne zwykłemu sformułowaniu
z zadanym bokiem: środki okręgów, długości promieni
oraz podobieństwa płaszczyzny o konstruowalnych danych
można uzyskać dozwolonymi konstrukcjami.
\begin{twierdzenie}[Kryterium wielokąta foremnego]\label{aa:tw:regular-polygon}
Dla $n\ge3$ foremny $n$-kąt można skonstruować cyrklem i linijką wtedy i tylko wtedy, gdy $\varphi(n)$ jest potęgą liczby $2$. Równoważnie $n=2^a p_1\cdots p_s$, gdzie $a\ge0$, a $p_i$ są różnymi liczbami pierwszymi Fermata, tj. postaci $2^{2^k}+1$.
\end{twierdzenie}
\begin{proof}
\textbf{Konieczność.}
Niech $\zeta_n=\cos(2\pi/n)+i\sin(2\pi/n)$
będzie pierwotnym pierwiastkiem $n$-tego stopnia z jedynki.
Jeśli foremny $n$-kąt jest konstruowalny, można obrócić
jego obraz tak, aby dwa punkty odniesienia odpowiadały
jednostkowemu okręgowi i wierzchołkowi $1$, więc
współrzędne $\zeta_n$ należą do pewnej rzeczywistej
wieży kwadratowej $K_r$ z
Twierdzenia~\ref{aa:tw:constructibility-tower}.
Złożenie ciał $K_r(i)$ zawiera $\zeta_n$.
Ponieważ $K_r\subseteq\mathbb R$, $i\notin K_r$
i $[K_r(i):K_r]=2$. Dlatego
$[K_r(i):\mathbb Q]=2[K_r:\mathbb Q]$ jest potęgą $2$.
Z prawa wieży wynika, że
$[\mathbb Q(\zeta_n):\mathbb Q]$ też jest potęgą $2$.
Na mocy Twierdzenia~\ref{aa:tw:cyclotomic-galois}
stopień ten wynosi $\varphi(n)$.

\textbf{Wystarczalność.}
Przyjmijmy teraz $\varphi(n)=2^m$.
Grupa $G=\operatorname{Gal}(\mathbb Q(\zeta_n)/\mathbb Q)$
jest abelową $2$-grupą. Sprzężenie zespolone
$c:\zeta_n\mapsto\zeta_n^{-1}$ jest w niej elementem
rzędu $2$ dla $n\ge3$. Ponieważ $G$ jest abelowa,
podgrupa $H=\langle c\rangle$ jest normalna,
a iloraz $G/H$ jest abelową $2$-grupą.
Jeśli jest nietrywialny, posiada podgrupę indeksu $2$:
wynika to z rozkładu na czynniki cykliczne,
z których co najmniej jeden ma rząd parzysty.
Powtarzając ten argument, otrzymujemy łańcuch podgrup
$G=H_0\supset H_1\supset\cdots\supset H_s=H$,
w którym wszystkie kolejne indeksy wynoszą $2$.
Na mocy odpowiedniości Galois odpowiada mu wieża
rozszerzeń stopnia $2$ od $\mathbb Q$ do ciała
$\mathbb Q(\zeta_n)^H$.
Każdy element tego ostatniego ciała jest rzeczywisty,
ponieważ jest ustalany przez sprzężenie zespolone.
Ponadto element $u=\zeta_n+\zeta_n^{-1}$ jest
ustalany przez $c$, więc do niego należy.
Odwrotnie $\zeta_n$ spełnia nad $\mathbb Q(u)$
równanie $T^2-uT+1=0$; ponieważ
$\zeta_n\notin\mathbb R$ dla $n\ge3$,
$[\mathbb Q(\zeta_n):\mathbb Q(u)]=2$.
Również $[\mathbb Q(\zeta_n):\mathbb Q(\zeta_n)^H]=2$,
więc $\mathbb Q(u)=\mathbb Q(\zeta_n)^H$.
Zatem $\cos(2\pi/n)=u/2$ należy do rzeczywistej
wieży kwadratowej. Ponieważ
$0<2\pi/n<\pi$ dla $n\ge3$, zachodzi
$\sin(2\pi/n)>0$ i
\[
 \sin(2\pi/n)=\sqrt{1-\cos^2(2\pi/n)}.
\]
Dołączając co najwyżej jeden pierwiastek kwadratowy,
możemy zatem skonstruować obie współrzędne $\zeta_n$.
Następnie długość boku $|1-\zeta_n|$ i przecięcia
kolejnych okręgów tej długości z okręgiem jednostkowym
pozwalają uzyskać wszystkie wierzchołki wielokąta.

\textbf{Postać liczb $n$.}
Rozkładając $n$ na potęgi liczb pierwszych, dostajemy
\[
 \varphi(n)=\prod_{p^e\parallel n}p^{e-1}(p-1),
\]
gdzie $p^e\parallel n$ znaczy $p^e\mid n$ oraz
$p^{e+1}\nmid n$. Jeżeli $\varphi(n)$ jest potęgą $2$,
to dla każdej nieparzystej liczby pierwszej $p\mid n$
musi zachodzić $e=1$ i $p-1=2^m$ dla pewnego $m\ge0$.
Jeśli $2^m+1$ jest liczbą pierwszą, wykładnik $m$
musi być potęgą $2$: gdyby $m=uv$, gdzie $u>1$
jest nieparzyste, to $2^m+1=(2^v)^u+1$
byłoby podzielne przez $2^v+1$, liczbę większą od $1$
i mniejszą od $2^m+1$. Zatem każda nieparzysta
liczba pierwsza dzieląca $n$ jest pierwszą Fermata,
a występuje w $n$ dokładnie raz. W przeciwnym kierunku,
dla $n=2^a p_1\cdots p_s$ ze wskazanymi różnymi
pierwszymi Fermata każdy czynnik wzoru na $\varphi(n)$
jest potęgą $2$. To kończy dowód równoważności.
\end{proof}
\begin{przyklad}[Wielokąty o $5$, $7$ i $17$ bokach]\label{aa:prz:polygon-examples}
Dla $n=5$ i $n=17$ liczby $\varphi(n)=4$ i $16$ są potęgami $2$, więc oba wielokąty są konstruowalne. Dla $n=7$ mamy $\varphi(7)=6$, więc foremny siedmiokąt nie jest konstruowalny. To kontrprzykład dla przekonania, że każdy wielokąt foremny można skonstruować cyrklem i linijką.
\end{przyklad}

\begin{przyklad}[Jawna konstrukcja foremnego pięciokąta]
\label{aa:prz:constructible-pentagon}
Nie poprzestaniemy na sprawdzeniu, że $\varphi(5)=4$.
Niech $\zeta=e^{2\pi i/5}$ oraz
$u=\zeta+\zeta^{-1}=2\cos(2\pi/5)$.
Wielomian cyklotomiczny $\Phi_5(T)=T^4+T^3+T^2+T+1$
zeruje się dla $T=\zeta$. Dzieląc tę równość przez
$\zeta^2\ne0$ i korzystając z
$\zeta^2+\zeta^{-2}=u^2-2$, dostajemy
$u^2+u-1=0$. Ponieważ kąt $2\pi/5=72^\circ$ jest ostry,
$u>0$, a więc $u=(\sqrt5-1)/2$.
Stąd współrzędne pierwszego wierzchołka poza $(1,0)$ są
\[
 \left(\cos\frac{2\pi}{5},\sin\frac{2\pi}{5}\right)
 =\left(\frac{\sqrt5-1}{4},\frac{\sqrt{10+2\sqrt5}}4\right).
\]
Każdą liczbę występującą w tych współrzędnych otrzymujemy
za pomocą działań arytmetycznych i dodatnich pierwiastków
kwadratowych, więc kryterium wieży kwadratowej daje
rzeczywistą konstrukcję punktu cyrklem i linijką.
Obracanie o ten sam kąt można wykonać przez kolejny wybór
punktu przecięcia okręgu jednostkowego z okręgiem o
odpowiedniej długości cięciwy; powtarzając krok pięciokrotnie,
otrzymujemy wszystkie wierzchołki.
W przeciwieństwie do tego $\varphi(7)=6$, więc
analogiczna konstrukcja regularnego siedmiokąta jest niemożliwa.

\begin{center}
\begin{tikzpicture}[scale=1.25]
  \draw[gray!70] (0,0) circle (1);
  \foreach \j in {0,...,4}{
    \coordinate (P\j) at ({cos(72*\j)},{sin(72*\j)});
  }
  \draw[thick,blue!65!black]
     (P0)--(P1)--(P2)--(P3)--(P4)--cycle;
  \draw[gray!70,->] (-1.26,0)--(1.28,0) node[right] {$x$};
  \draw[gray!70,->] (0,-1.25)--(0,1.3) node[above] {$y$};
  \fill (0,0) circle(1.3pt) node[below right] {$0$};
  \foreach \j in {0,...,4}{\fill (P\j) circle(1.5pt);}
  \node[anchor=west] at (1.04,0.12) {$1$};
  \node[anchor=south west] at (P1) {$\zeta_5$};
\end{tikzpicture}
\end{center}
Rysunek pokazuje jednostkowy okrąg oraz wszystkie pięć
wierzchołków otrzymanych z jednego skonstruowanego punktu
$\zeta_5$. Jest to ilustracja położeń, a nie zastępstwo
dowodu konstruowalności współrzędnych.
\end{przyklad}

% AA-31 | Iteracja XVI; szczegolowy audyt historyczny IV
\chapter{Wielomiany symetryczne i wzory Viète'a}\label{aa:ch:symmetric-polynomials}
\section{Działanie grupy symetrycznej}

W obliczeniach z pierwiastkami wielomianu często nie znamy samych
pierwiastków, lecz znamy ich sumę i iloczyn. Zamiana dwóch
pierwiastków nie zmienia tych danych. Powstaje pytanie, czy każdą
wielomianową zależność niezmienniczą przy zamianie pierwiastków
można zapisać wyłącznie za pomocą pewnych podstawowych sum
symetrycznych. Odpowiedź jest twierdząca nawet wtedy, gdy
współczynniki należą do pierścienia z dzielnikami zera.

\begin{definicja}\label{aa:def:symmetric-poly}
Niech $A$ będzie przemiennym pierścieniem z jedynką
i niech $n\ge1$ będzie liczbą całkowitą. Grupa $S_n$ składa się
ze wszystkich bijekcji zbioru $\{1,\ldots,n\}$ na siebie.
Wielomian $f\in A[X_1,\ldots,X_n]$ jest
\emph{symetryczny}, jeśli dla każdej permutacji $\sigma\in S_n$
zachodzi
\[
 f(X_{\sigma(1)},\ldots,X_{\sigma(n)})=f(X_1,\ldots,X_n).
\]
Zbiór wszystkich takich wielomianów, z działaniami odziedziczonymi
z $A[X_1,\ldots,X_n]$, oznaczamy przez
$A[X_1,\ldots,X_n]^{S_n}$.
Dla $0\le k\le n$ oznaczamy przez $e_k$ $k$-ty
\emph{elementarny wielomian symetryczny}:
\[
 e_0=1,\qquad
 e_k=\sum_{1\le i_1<\cdots<i_k\le n}X_{i_1}\cdots X_{i_k}
 \quad(1\le k\le n).
\]
Każdy składnik sumy odpowiada wyborowi $k$ \emph{różnych}
indeksów; suma nie zawiera powtórzeń tego samego podzbioru.
Każda permutacja zmiennych permutuje rodzinę takich podzbiorów,
więc $e_k\in A[X_1,\ldots,X_n]^{S_n}$ dla $0\le k\le n$.
\end{definicja}

\begin{przyklad}[Trzy zmienne]
Dla $n=3$ otrzymujemy
\[
 e_1=X_1+X_2+X_3,\qquad
 e_2=X_1X_2+X_1X_3+X_2X_3,\qquad
 e_3=X_1X_2X_3.
\]
Po zamianie $X_1$ z $X_3$ składniki pierwszej sumy jedynie
zmieniają kolejność, podobnie jak trzy składniki drugiej sumy.
Iloczyn $e_3$ również pozostaje bez zmian. Tak samo działa
każda permutacja, więc $e_k$ należy do pierścienia wielomianów
symetrycznych.
\end{przyklad}

\begin{twierdzenie}[Podstawowe twierdzenie o wielomianach symetrycznych]\label{aa:tw:fundamental-symmetric}
Dla każdego przemiennego pierścienia $A$ z jedynką
i każdego $n\ge1$ homomorfizm pierścieni
\[
 \Phi:A[T_1,\ldots,T_n]\longrightarrow
 A[X_1,\ldots,X_n]^{S_n},\qquad
 \Phi(T_k)=e_k\quad(1\le k\le n),\quad\Phi(a)=a\quad(a\in A),
\]
jest izomorfizmem. W szczególności każdy wielomian symetryczny
$f$ ma \emph{jednoznaczne} przedstawienie
$f=P(e_1,\ldots,e_n)$ dla pewnego
$P\in A[T_1,\ldots,T_n]$.
\end{twierdzenie}

\begin{proof}
Jeżeli $1_A=0$, to $A$ jest pierścieniem zerowym.
Wówczas oba pierścienie występujące w twierdzeniu
są jednopunktowe i $\Phi$ jest izomorfizmem.
W dalszej części przyjmujemy $1_A\ne0$.

\textbf{Krok 1: porządek i pierwszy jednomian.}
Jednomiany $X^a=X_1^{a_1}\cdots X_n^{a_n}$, gdzie
$a=(a_1,\ldots,a_n)\in\mathbb N_0^n$, porządkujemy
leksykograficznie przy $X_1>\cdots>X_n$.
Oznacza to, że $X^a>X^b$, jeśli w pierwszej współrzędnej,
w której $a$ i $b$ się różnią, mamy $a_j>b_j$.
Każdy niezerowy wielomian ma skończenie wiele wyrazów
o niezerowych współczynnikach, a zatem ma największy wyraz,
zwany wiodącym. Jeśli $f$ jest symetryczny i jego wyraz
wiodący jest równy $cX_1^{a_1}\cdots X_n^{a_n}$, gdzie
$c\in A\setminus\{0\}$, to $a_1\ge\cdots\ge a_n$.
Rzeczywiście, gdyby dla pewnego $i<j$ zachodziło $a_i<a_j$,
permutacja zamieniająca $X_i$ z $X_j$ wymuszałaby wystąpienie
w $f$ wyrazu o tej samej niezerowej wartości współczynnika
$c$ i o większym jednomianie. To przeczy wyborowi wyrazu
wiodącego. Wniosek nie wymaga dzielenia przez żaden element $A$.

\textbf{Krok 2: eliminacja wyrazu wiodącego.}
Wyrazem wiodącym $e_k$ jest $X_1\cdots X_k$ z
współczynnikiem $1_A$. Dla $a_1\ge\cdots\ge a_n$ kładziemy
$a_{n+1}=0$ i rozpatrujemy
\[
 M_a=e_1^{a_1-a_2}e_2^{a_2-a_3}\cdots e_n^{a_n-a_{n+1}}.
\]
Wykładniki są nieujemne. Mnożenie wyrazów wiodących
daje dokładnie $X_1^{a_1}\cdots X_n^{a_n}$, ponieważ
wykładnik $X_j$ wynosi
$(a_j-a_{j+1})+\cdots+(a_n-a_{n+1})=a_j$.
Ponieważ współczynniki wiodące wszystkich $e_k$ są
równe $1_A$, współczynnik tego wyrazu w $M_a$ również
jest $1_A$, niezależnie od istnienia dzielników zera w $A$.
Wielomian $f-cM_a$ jest symetryczny, a jego wyraz wiodący
jest ściśle mniejszy od dotychczasowego (chyba że otrzymaliśmy
wielomian zerowy).

\textbf{Krok 3: skończoność algorytmu.}
Wszystkie jednomiany $M_a$ są jednorodne stopnia
$a_1+\cdots+a_n$. W szczególności żaden nowy wyraz
powstały przy odejmowaniu $cM_a$ nie ma stopnia
większego niż początkowy $d=\deg f$.
W $n$ zmiennych istnieje tylko skończenie wiele jednomianów
stopnia całkowitego nie większego niż $d$: są one
indeksowane przez wektory $b\in\mathbb N_0^n$ o
$b_1+\cdots+b_n\le d$. Przy każdym niezakończonym
kroku jednomian wiodący ściśle maleje w tym skończonym
zbiorze. Algorytm zatem się kończy i wyraża $f$ jako
wielomian w $e_1,\ldots,e_n$. Jest to surjektywność $\Phi$.

\textbf{Krok 4: jednoznaczność przedstawienia.}
Niech $P\in A[T_1,\ldots,T_n]$ będzie niezerowy.
Dla każdego jednomianu $T^b$ jego obraz ma współczynnik
wiodący $1_A$, ponieważ wielomiany $e_k$ są moniczne
w ustalonym porządku leksykograficznym.
Jednomianowi $T^b=T_1^{b_1}\cdots T_n^{b_n}$ odpowiada
jednomian wiodący
\[
 X_1^{b_1+\cdots+b_n}X_2^{b_2+\cdots+b_n}
 \cdots X_n^{b_n}
\]
po podstawieniu $T_k=e_k$. Mapa wykładników
\[
 (b_1,\ldots,b_n)\longmapsto
 (b_1+\cdots+b_n,b_2+\cdots+b_n,\ldots,b_n)
\]
jest różnowartościowa: z prawej strony odzyskujemy
$b_j$ jako różnicę $j$-tej i $(j+1)$-szej współrzędnej
(przyjmując ostatnią dodatkową współrzędną równą $0$).
Spośród niezerowych wyrazów $P$ wybierzmy ten, którego
obraz ma największy jednomian wiodący. Inne wyrazy
mają mniejsze jednomiany wiodące, a każdy niewiodący
składnik ich rozwinięcia jest tym bardziej mniejszy.
Zatem wybrany największy jednomian nie może się skasować:
pochodzi tylko z wybranego wyrazu $P$, a jego współczynnik
pozostaje ten sam, niezerowy w $A$.
Otrzymujemy $\Phi(P)\ne0$, co dowodzi iniektywności.
\end{proof}

\begin{przyklad}[Algorytm dla dwóch zmiennych]\label{aa:prz:symmetric-algorithm}
Weźmy $f=X_1^2+X_2^2+X_1X_2$ nad $\mathbb Z$.
Wyrazem wiodącym jest $X_1^2$, więc odejmujemy $e_1^2$:
\[
 f-e_1^2
 =X_1^2+X_2^2+X_1X_2-(X_1^2+2X_1X_2+X_2^2)
 =-X_1X_2=-e_2.
\]
Oznacza to $f=e_1^2-e_2$. Nie zgadywaliśmy współczynników:
wyznaczył je algorytm eliminowania wyrazu wiodącego.
\end{przyklad}

\begin{przyklad}[Pierścień z dzielnikami zera]
Dla $A=\mathbb Z/4\mathbb Z$ i dwóch zmiennych
wielomian $X_1^2+X_2^2$ nadal spełnia
$X_1^2+X_2^2=e_1^2-2e_2$.
Nie można tutaj skrócić dowodu przez przejście do
ciała ułamków $A$: pierścień $A$ ma niezerowe dzielniki zera,
więc nie posiada ciała ułamków w zwykłym znaczeniu.
Dowód powyższego twierdzenia działa bez takiego przejścia,
ponieważ korzysta jedynie ze współczynników wiodących $1_A$.
\end{przyklad}

\begin{przyklad}[Eliminacja wyrazu wiodącego w trzech zmiennych]\label{aa:prz:symmetric-three-elimination}
Niech $A=\mathbb Z$ i
$f=X_1^2X_2+X_1^2X_3+X_2^2X_1+X_2^2X_3+X_3^2X_1+X_3^2X_2$.
Jest to suma wszystkich jednomianów $X_i^2X_j$ z $i\ne j$;
permutacja zmiennych jedynie zmienia kolejność sześciu składników.
Wyraz wiodący w porządku leksykograficznym jest $X_1^2X_2$.
Odpowiada mu wektor wykładników $(2,1,0)$, a więc
$M_a=e_1^{2-1}e_2^{1-0}=e_1e_2$.
Rozwijając iloczyn, otrzymujemy
\[
 e_1e_2=f+3X_1X_2X_3=f+3e_3.
\]
Współczynnik $3$ pochodzi stąd, że jednomian $X_1X_2X_3$
można uzyskać na trzy sposoby: wybierając w $e_1$ kolejno
$X_1$, $X_2$ lub $X_3$, a w $e_2$ iloczyn pozostałych dwóch
zmiennych. Stąd $f=e_1e_2-3e_3$.
Tożsamość pozostaje poprawna po przejściu do dowolnego
przemiennego pierścienia z jedynką, także gdy jego charakterystyka
dzieli $3$.
\end{przyklad}

\section{Pierwiastki i współczynniki}

Każdy wielomian moniczny dodatniego stopnia nad ciałem
rozkłada się w swoim ciele rozkładu na czynniki liniowe.
Zestawienie jego współczynników
z wielomianami symetrycznymi wyliczanymi na pierwiastkach wyjaśnia,
dlaczego nie musimy znać poszczególnych pierwiastków,
aby obliczać niektóre wyrażenia z ich udziałem.

\begin{twierdzenie}[Wzory Viète'a]\label{aa:tw:vieta-general}
Niech $K$ będzie ciałem, $n\ge1$ i
$f(X)=X^n+c_1X^{n-1}+\cdots+c_n\in K[X]$.
Niech $L/K$ będzie ciałem rozkładu $f$, a
$\alpha_1,\ldots,\alpha_n\in L$ jego pierwiastkami,
wypisanymi z krotnościami algebraicznymi.
Wówczas dla każdej liczby $k\in\{1,\ldots,n\}$ zachodzi
$e_k(\alpha_1,\ldots,\alpha_n)=(-1)^kc_k$.
\end{twierdzenie}
\begin{proof}
W $L[X]$ zachodzi
$f(X)=\prod_{j=1}^n(X-\alpha_j)$, gdzie każdy pierwiastek
jest uwzględniony tyle razy, ile wynosi jego krotność.
W rozwinięciu tego iloczynu wyraz stopnia $n-k$ w $X$
otrzymujemy, wybierając $k$ spośród $n$ czynników,
z których bierzemy $-\alpha_j$, a z pozostałych
bierzemy $X$. Zbiór wyborów jest dokładnie zbiorem
$k$-elementowych podzbiorów $\{1,\ldots,n\}$.
Współczynnik przy $X^{n-k}$ jest zatem równy
$(-1)^ke_k(\alpha_1,\ldots,\alpha_n)$.
Porównanie go ze współczynnikiem $c_k$ po lewej stronie
daje $c_k=(-1)^ke_k$. Mnożąc przez $(-1)^k$
(dla dowolnej charakterystyki) otrzymujemy tezę.
\end{proof}

\begin{przyklad}[Suma kwadratów pierwiastków]\label{aa:prz:symmetric-squares}
Dla monicznego wielomianu
$f(X)=X^3-2X^2+3X-4\in\mathbb Q[X]$ mamy ze wzorów
Viète'a $e_1(\alpha_1,\alpha_2,\alpha_3)=2$ oraz
$e_2(\alpha_1,\alpha_2,\alpha_3)=3$.
Rozwijając kwadrat sumy, otrzymujemy
\[
 (\alpha_1+\alpha_2+\alpha_3)^2
 =\alpha_1^2+\alpha_2^2+\alpha_3^2
  +2(\alpha_1\alpha_2+\alpha_1\alpha_3+\alpha_2\alpha_3).
\]
Stąd suma kwadratów wynosi $e_1^2-2e_2=4-6=-2$.
Obliczenie nie wymaga wyznaczania żadnego z trzech pierwiastków.
\end{przyklad}

\begin{przyklad}[Suma sześcianów]
Niech $f(X)=X^3+c_1X^2+c_2X+c_3$ będzie moniczny
nad dowolnym ciałem $K$. Tożsamość w pierścieniu
$\mathbb Z[U,V,W]$ daje
\[
 U^3+V^3+W^3=(U+V+W)^3
 -3(U+V+W)(UV+UW+VW)+3UVW.
\]
Dla pierwiastków $f$, po podstawieniu
$e_1=-c_1$, $e_2=c_2$, $e_3=-c_3$, otrzymujemy
$\sum_j\alpha_j^3=-c_1^3+3c_1c_2-3c_3$.
Ponieważ wykorzystana tożsamość ma współczynniki
całkowite, ten sam wzór jest poprawny również
w charakterystykach $2$ i $3$ (odczytując odpowiednie
współczynniki w $K$).
\end{przyklad}

\begin{przyklad}[Kontrprzykład: symetria nie oznacza liniowości w $e_k$]\label{aa:prz:symmetric-nonlinear}
Wielomian $X_1^2+X_2^2$ jest symetryczny, lecz nie ma
postaci $a e_1+b e_2$ dla $a,b\in\mathbb Z$:
$e_1=X_1+X_2$ ma tylko wyrazy stopnia $1$, a
$e_2=X_1X_2$ jest jednomianem stopnia $2$ bez
wyrazu $X_1^2$. Porównanie współczynnika przy
$X_1^2$ prowadzi do sprzeczności $1=0$.
Poprawne przedstawienie jest nieliniowe:
$X_1^2+X_2^2=e_1^2-2e_2$.
\end{przyklad}

\begin{uwaga}[Ostrzeżenie w dodatniej charakterystyce]
W charakterystyce $2$ tożsamość
$X_1^2+X_2^2=e_1^2-2e_2$ upraszcza się do
$X_1^2+X_2^2=e_1^2$.
Nie wolno zatem próbować odzyskać $e_2$ z sumy
kwadratów przez dzielenie przez $2$: liczba $2$
jest zerem w ciele charakterystyki $2$.
Nie przeczy to podstawowemu twierdzeniu o wielomianach
symetrycznych: jego generatorami są nadal $e_1,e_2$,
a nie jedynie suma kwadratów.
\end{uwaga}

\begin{uwaga}\label{aa:uw:symmetric-galois}
Niech $f\in K[X]$ będzie rozdzielczy i niech $L/K$
będzie jego ciałem rozkładu. Grupa
$G=\operatorname{Gal}(L/K)$ permutuje pierwiastki $f$.
Każdy wielomian symetryczny o współczynnikach w $K$,
obliczony na tych pierwiastkach, należy do $K$:
wynika to bezpośrednio z Twierdzenia~\ref{aa:tw:fundamental-symmetric}
i wzorów Viète'a (Twierdzenie~\ref{aa:tw:vieta-general}),
bez korzystania z odpowiedniości Galois.
Nie trzeba też zakładać rozdzielczości dla tej ostatniej
konkluzji, jeżeli pierwiastki liczymy z krotnościami.
\end{uwaga}

% AA-32 | Iteracja XVI; szczegolowy audyt historyczny IV
\chapter{Wyróżnik wielomianu i grupa Galois}\label{aa:ch:discriminant-galois}
\section{Wyróżnik i znak permutacji}

Jeżeli wielomian ma dwa równe pierwiastki, różnica między
nimi jest zerem. Chcemy otrzymać z tych różnic jedno
wyrażenie, które nie zależy od numeracji pierwiastków
i daje się obliczać ze współczynników wielomianu.
Zastosowanie do grup Galois wymaga starannego rozróżnienia
charakterystyk: w charakterystyce różnej od $2$
wyróżnik wykrywa, czy automorfizmy Galois działają
jako permutacje parzyste; w charakterystyce $2$
wniosek tego rodzaju na ogół jest fałszywy.

\begin{definicja}\label{aa:def:polynomial-discriminant}
Niech $K$ będzie ciałem, $n\ge1$ liczbą całkowitą,
$f\in K[X]$ monicznym wielomianem stopnia $n$,
a $\overline K$ ustalonym domknięciem algebraicznym $K$.
Niech $\alpha_1,\ldots,\alpha_n\in\overline K$
będą wszystkimi pierwiastkami $f$, wypisanymi
z krotnościami, czyli
$f(X)=\prod_{j=1}^n(X-\alpha_j)$ w $\overline K[X]$.
\emph{Wyróżnikiem} $f$ nazywamy element
\[
 \operatorname{disc}(f)=
 \prod_{1\le i<j\le n}(\alpha_i-\alpha_j)^2\in\overline K.
\]
Przy $n=1$ iloczyn pusty ma wartość $1$.
\end{definicja}

\begin{przyklad}[Stopień drugi]
Dla $f(X)=X^2+bX+c$ i pierwiastków $\alpha,\beta$
zachodzą $\alpha+\beta=-b$ oraz $\alpha\beta=c$.
Stąd
\[
 \operatorname{disc}(f)=(\alpha-\beta)^2
 =(\alpha+\beta)^2-4\alpha\beta=b^2-4c.
\]
Nie potrzebujemy wzoru na pierwiastki; rachunek jest
poprawny nad dowolnym ciałem, również w charakterystyce $2$.
\end{przyklad}

\begin{twierdzenie}[Wyróżnik a rozdzielczość]\label{aa:tw:disc-separable}
Dla dowolnego ciała $K$ i każdego monicznego
$f\in K[X]$ stopnia dodatniego zachodzą dwa stwierdzenia:
$\operatorname{disc}(f)\in K$, a ponadto
$\operatorname{disc}(f)\ne0$ wtedy i tylko wtedy,
gdy $f$ jest rozdzielczy nad $K$, to znaczy ma
$\deg f$ różnych pierwiastków w $\overline K$.
\end{twierdzenie}
\begin{proof}
\textbf{Przynależność do $K$.}
Wyrażenie $D(X_1,\ldots,X_n)=
\prod_{i<j}(X_i-X_j)^2$ jest symetryczne:
permutacja zmiennych permutuje czynniki z możliwą
zmianą znaku różnic, lecz znak znika po podniesieniu
do kwadratu. Ponadto $D$ ma współczynniki całkowite,
więc po przejściu do $K$ należy do
$K[X_1,\ldots,X_n]^{S_n}$.
Twierdzenie~\ref{aa:tw:fundamental-symmetric} daje
$D=P(e_1,\ldots,e_n)$ dla pewnego $P\in K[T_1,\ldots,T_n]$.
Podstawiając pierwiastki i korzystając z
Twierdzenia~\ref{aa:tw:vieta-general}, dostajemy
\[
 \operatorname{disc}(f)
 =P(-c_1,c_2,\ldots,(-1)^nc_n)\in K,
\]
gdzie $f(X)=X^n+c_1X^{n-1}+\cdots+c_n$.

\textbf{Kryterium niezerowości.}
W $\overline K$, które jest ciałem, iloczyn
niezerowych elementów jest niezerowy. Wobec tego
$\operatorname{disc}(f)\ne0$ dokładnie wtedy,
gdy dla wszystkich $1\le i<j\le n$ mamy
$\alpha_i-\alpha_j\ne0$. Oznacza to, że wszystkie
wypisane pierwiastki są parami różne, czyli że
$f$ jest rozdzielczy. Argument działa dla każdej
charakterystyki i obejmuje wielomiany stopnia $1$.
\end{proof}

\begin{przyklad}[Wyróżnik zerowy w charakterystyce dwa]\label{aa:prz:disc-inseparable-char2}
Niech $K=\mathbb F_2(t)$, gdzie $t$ jest elementem
przestępnym nad $\mathbb F_2$, i niech $f(X)=X^2+t$.
W domknięciu algebraicznym $K$ istnieje element $\alpha$
taki, że $\alpha^2=t$. Ponieważ charakterystyka wynosi $2$,
mamy $(X-\alpha)^2=X^2+\alpha^2=X^2+t$.
Wielomian ma zatem jeden pierwiastek z krotnością $2$,
a jego wyróżnik wynosi $0$.
Z drugiej strony $f$ jest nierozkładalny nad $K$:
gdyby $t=(A(t)/B(t))^2$ dla względnie pierwszych
niezerowych wielomianów $A,B\in\mathbb F_2[t]$, to
$2\deg A=1+2\deg B$, co jest niemożliwe.
Przykład rozróżnia nierozkładalność od rozdzielczości:
w dodatniej charakterystyce nierozkładalny wielomian
może mieć wielokrotny pierwiastek.
\end{przyklad}

\begin{twierdzenie}[Kryterium grupy alternującej]\label{aa:tw:disc-alternating}
Niech $K$ będzie ciałem charakterystyki różnej od $2$,
$f\in K[X]$ monicznym i rozdzielczym wielomianem
stopnia $n\ge2$, $L$ jego ciałem rozkładu oraz
$G=\operatorname{Gal}(L/K)$.
Ustalmy numerację różnych pierwiastków; daje ona wierne
zanurzenie $G\hookrightarrow S_n$, ponieważ element $G$
jest określony przez obrazy pierwiastków generujących $L$.
Wówczas $G\subseteq A_n$ wtedy i tylko wtedy,
gdy $\operatorname{disc}(f)$ jest kwadratem w $K$.
\end{twierdzenie}
\begin{proof}
Wybieramy kolejność $\alpha_1,\ldots,\alpha_n$ i kładziemy
\[
 \Delta=\prod_{1\le i<j\le n}(\alpha_i-\alpha_j)\in L.
\]
Z rozdzielczości wynika $\Delta\ne0$, a definicja
daje $\Delta^2=\operatorname{disc}(f)$.
Dla dowolnej transpozycji dwóch indeksów
w iloczynie $\Delta$ znak zmienia się dokładnie
raz: czynnik różnicy zamienianych pierwiastków
zmienia znak, natomiast pozostałe czynniki
zostają odpowiednio poprzestawiane parami.
Każdą permutację można przedstawić jako iloczyn
transpozycji, więc dla każdego $\sigma\in G$ zachodzi
\[
 \sigma(\Delta)=\operatorname{sgn}(\sigma)\Delta,
\]
gdzie $\operatorname{sgn}:S_n\to\{1,-1\}$
oznacza znak permutacji.

Jeśli $G\subseteq A_n$, wszystkie elementy $G$
mają znak $1$, zatem każdy z nich ustala $\Delta$.
Ponieważ $L/K$ jest skończonym rozszerzeniem Galois,
twierdzenie o ciele elementów stałych daje
$\Delta\in L^G=K$. W takim razie wyróżnik jest
kwadratem elementu $\Delta\in K$.

Odwrotnie, jeśli istnieje $a\in K$ takie, że
$\operatorname{disc}(f)=a^2$, to
$(\Delta-a)(\Delta+a)=0$ w ciele $L$.
Stąd $\Delta=a$ lub $\Delta=-a$, więc $\Delta\in K$.
Dla każdego $\sigma\in G$ mamy jednocześnie
$\sigma(\Delta)=\Delta$ oraz
$\sigma(\Delta)=\operatorname{sgn}(\sigma)\Delta$.
Ponieważ $\Delta\ne0$ i $1\ne-1$ w $K$, koniecznie
$\operatorname{sgn}(\sigma)=1$. Oznacza to
$\sigma\in A_n$ dla każdego $\sigma\in G$.
\end{proof}

\begin{uwaga}[Dlaczego charakterystyka $2$ wymaga osobnego traktowania?]
W charakterystyce $2$ liczby $1$ i $-1$ są równe.
Wobec tego $\Delta$ jest ustalane przez wszystkie
automorfizmy, nawet te wyznaczające nieparzyste
permutacje. Dla rozdzielczego $f$ oznacza to
$\Delta\in K$, a więc $\operatorname{disc}(f)=\Delta^2$
zawsze jest kwadratem w $K$. Nie stanowi to kryterium
zawierania się grupy Galois w $A_n$; poniżej podajemy
konkretny wielomian z grupą $S_3$.
\end{uwaga}

\section{Wielomiany stopnia trzeciego}

\begin{twierdzenie}[Wyróżnik wielomianu sześciennego]\label{aa:tw:cubic-disc}
Niech $K$ będzie dowolnym ciałem, a $p,q\in K$.
Dla $f(X)=X^3+pX+q\in K[X]$ zachodzi
$\operatorname{disc}(f)=-4p^3-27q^2$,
gdzie liczby całkowite $-4$ i $-27$ rozumiemy
jako ich obrazy w $K$ przez homomorfizm $\mathbb Z\to K$.
\end{twierdzenie}
\begin{proof}
Niech $a,b,c\in\overline K$ będą pierwiastkami
$f$, wypisanymi z krotnościami. Ze wzorów Viète'a
mamy $a+b+c=0$, $ab+ac+bc=p$ oraz $abc=-q$.
Kładziemy $u=a+b$ oraz $v=ab$; z pierwszej równości
wynika $c=-u$. Obliczamy kolejno
\[
 p=ab+(a+b)c=v-u^2,\qquad
 q=-abc=uv,
\]
\[
 (a-b)^2=(a+b)^2-4ab=u^2-4v
\]
oraz
\[
 (a-c)(b-c)=(a+u)(b+u)
 =ab+u(a+b)+u^2=v+2u^2.
\]
Zatem, używając definicji wyróżnika,
\[
 \begin{aligned}
 \operatorname{disc}(f)
 &=(a-b)^2(a-c)^2(b-c)^2\\
 &=(u^2-4v)(2u^2+v)^2\\
 &=(u^2-4v)(4u^4+4u^2v+v^2)\\
 &=4u^6-12u^4v-15u^2v^2-4v^3.
 \end{aligned}
\]
Z drugiej strony, podstawiając $p=v-u^2$ i $q=uv$,
dostajemy
\[
 \begin{aligned}
 -4p^3-27q^2
 &=-4(v-u^2)^3-27u^2v^2\\
 &=-4(v^3-3v^2u^2+3vu^4-u^6)-27u^2v^2\\
 &=4u^6-12u^4v-15u^2v^2-4v^3.
 \end{aligned}
\]
Wyrażenia są równe. Każdy krok był tożsamością
wielomianową o współczynnikach całkowitych:
nie dzieliliśmy przez $2$ ani przez $3$.
Dlatego ten sam dowód działa także w charakterystyce
$2$ lub $3$.
\end{proof}

\begin{twierdzenie}[Grupa Galois nierozkładalnego wielomianu sześciennego]\label{aa:tw:cubic-galois}
Niech $K$ będzie ciałem charakterystyki różnej od $2$,
a $f\in K[X]$ monicznym, nierozkładalnym i rozdzielczym
wielomianem stopnia $3$. Wtedy jego grupa Galois jest
izomorficzna z $A_3\cong C_3$, jeśli wyróżnik jest
kwadratem w $K$, a z $S_3$ w przeciwnym razie.
\end{twierdzenie}
\begin{proof}
Niech $L/K$ będzie ciałem rozkładu $f$ i
$G=\operatorname{Gal}(L/K)\hookrightarrow S_3$.
Nierozkładalność $f$ oznacza, że działanie $G$
na jego trzech pierwiastkach jest przechodnie:
każde dwa pierwiastki wielomianu nierozkładalnego
są sprzężone nad $K$, a odpowiednie zanurzenie
$K(\alpha)\to L$ przedłuża się do automorfizmu
normalnego rozszerzenia $L/K$.
Z Twierdzenia~\ref{aa:tw:orbita-stabilizator} wynika, że $3$ dzieli $|G|$.
Z drugiej strony $G$ jest podgrupą $S_3$,
więc na mocy twierdzenia Lagrange'a $|G|$ dzieli
$|S_3|=6$. Jedynymi możliwościami są $|G|=3$ i $6$.
Podgrupa $S_3$ rzędu $3$ jest równa
$A_3=\{\mathrm{id},(123),(132)\}$, a podgrupa rzędu $6$
jest całym $S_3$. Twierdzenie~\ref{aa:tw:disc-alternating}
rozstrzyga między tymi możliwościami na podstawie
wyróżnika, gdy $\operatorname{char}K\ne2$.
\end{proof}

\begin{przyklad}[Dwa możliwe typy]\label{aa:prz:cubic-groups}
Dla $f(X)=X^3-2\in\mathbb Q[X]$ kryterium Eisensteina
przy liczbie pierwszej $2$ daje nierozkładalność;
$f'(X)=3X^2$ nie ma wspólnego pierwiastka z $f$,
więc $f$ jest rozdzielczy. Mamy $p=0$, $q=-2$,
zatem $\operatorname{disc}(f)=-27\cdot4=-108$.
Kwadrat liczby wymiernej nie może być ujemny,
więc grupa Galois jest izomorficzna z $S_3$.

Dla $g(X)=X^3-3X+1$ każdy ewentualny wymierny
pierwiastek musi być całkowitym dzielnikiem $1$,
czyli $1$ albo $-1$. Tymczasem $g(1)=-1$,
$g(-1)=3$, zatem wielomian nie ma pierwiastka
wymiernego i jako sześcienny jest nierozkładalny.
Jego wyróżnik wynosi
$-4(-3)^3-27=108-27=81=9^2\ne0$.
W konsekwencji $g$ jest rozdzielczy i jego grupa
Galois jest cykliczna rzędu $3$.
\end{przyklad}

\begin{przyklad}[Kontrprzykład w charakterystyce dwa]\label{aa:prz:disc-char2}
Zacznijmy od przykładu, w którym grupa jest parzysta.
Nad $\mathbb F_2$ wielomian $X^3+X+1$ nie ma
pierwiastka w $\mathbb F_2$ (wartości w $0$ i $1$
są równe $1$), jest więc nierozkładalny.
Jego ciało rozkładu to $\mathbb F_8$, a grupa Galois
jest cykliczna rzędu $3$, generowana przez Frobeniusa.
Wyróżnik jest równy $-4-27=-31=1$ w $\mathbb F_2$.
Ten przykład sam w sobie \emph{nie} pokazuje błędności
kryterium: grupa rzeczywiście zawiera się w $A_3$.

Teraz niech $K=\mathbb F_2(t)$ będzie ciałem funkcji
wymiernych w zmiennej przestępnej $t$ nad $\mathbb F_2$
i rozważmy $h(X)=X^3+X+t\in K[X]$.
Pokażemy, że $h$ jest rozdzielczy i ma grupę Galois $S_3$,
choć jego wyróżnik jest kwadratem w $K$.

\textbf{Nierozkładalność.}
Gdyby istniał element $r\in K$ spełniający $r^3+r=t$,
to $r\ne0$, ponieważ $t\ne0$. Zapiszmy więc $r=A(t)/B(t)$,
gdzie $A,B\in\mathbb F_2[t]$ są niezerowe
i względnie pierwsze.
Rząd bieguna $r$ w nieskończoności jest równy
$m=\max\{\deg A-\deg B,0\}$.
Jeśli $m=0$, funkcja $r$ nie ma bieguna w nieskończoności,
więc nie ma go też $r^3+r$; tymczasem $t$ ma biegun
rzędu $1$. Jeśli $m>0$, składnik $r^3$ ma biegun rzędu
$3m$, a składnik $r$ rzędu $m$, więc ich suma ma biegun
rzędu $3m\ne1$. Sprzeczność dowodzi, że $h$ nie ma
pierwiastka w $K$, a zatem, jako wielomian stopnia $3$,
jest nierozkładalny.

\textbf{Rozdzielczość.}
Pochodna formalna to $h'(X)=X^2+1=(X+1)^2$.
Wspólny pierwiastek $h$ i $h'$ musiałby być równy $1$,
lecz $h(1)=t\ne0$. Zatem pierwiastki $h$ są różne.

\textbf{Stopień ciała rozkładu.}
Niech $\alpha$ będzie jednym z pierwiastków $h$.
Z równości $t=\alpha^3+\alpha$ wynika
$K(\alpha)=\mathbb F_2(\alpha)$; przy tym $\alpha$
jest przestępny nad $\mathbb F_2$, ponieważ $t$ jest
przestępny. W pierścieniu $K(\alpha)[X]$ mamy
\[
 h(X)=(X+\alpha)
 \bigl(X^2+\alpha X+\alpha^2+1\bigr).
\]
Drugi czynnik ma pierwiastek w $K(\alpha)$
wtedy i tylko wtedy, gdy dla pewnego
$z\in\mathbb F_2(\alpha)$ zachodzi
$z^2+z=1+\alpha^{-2}$: wystarczy w jego równaniu
podstawić $X=\alpha z$ i podzielić przez
$\alpha^2\ne0$ (ponieważ $h(0)=t\ne0$).
Załóżmy, że takie $z$ istnieje. Prawa strona ma
w punkcie $\alpha=0$ biegun rzędu $2$.
Jeśli $z$ nie ma tam bieguna, to $z^2+z$ również
go nie ma. Jeśli $z$ ma biegun rzędu $m>0$,
to $z^2+z$ ma biegun rzędu $2m$. Musi więc być
$m=1$. Zapisując $z$ w postaci ułamka nieskracalnego,
widzimy, że przy biegunie prostym jego część główna
ma postać $c\alpha^{-1}$ dla pewnego
$c\in\mathbb F_2^\times=\{1\}$.
Zatem $z=\alpha^{-1}+w$, gdzie $w$ nie ma bieguna
przy $\alpha=0$. Otrzymujemy
\[
 z^2+z=\alpha^{-2}+\alpha^{-1}+w^2+w,
\]
co nie może być równe $1+\alpha^{-2}$, gdyż
$w^2+w$ nie ma bieguna, a $\alpha^{-1}$ ma.
Sprzeczność dowodzi nierozkładalności czynnika
kwadratowego nad $K(\alpha)$.
Stopień ciała rozkładu $L$ jest zatem
$[L:K]=[L:K(\alpha)][K(\alpha):K]=2\cdot3=6$.
Skoro $L/K$ jest Galois, jego grupa ma rząd $6$
i działa wiernie na trzech pierwiastkach, więc
$\operatorname{Gal}(L/K)\cong S_3$.

Wreszcie Twierdzenie~\ref{aa:tw:cubic-disc},
poprawne także w charakterystyce $2$, daje
$\operatorname{disc}(h)=-4-27t^2=t^2$,
czyli jest on kwadratem niezerowego elementu $t\in K$.
Otrzymaliśmy więc kontrprzykład w charakterystyce $2$:
wyróżnik jest kwadratem w $K$, lecz grupa $S_3$
\emph{nie} zawiera się w $A_3$.
\end{przyklad}

% AA-33 | Audyt I, 22 XI 2026; kontrola matematyczna i redakcyjna.
\chapter{Wypadkowa wielomianów i macierz Sylvestera}\label{aa:ch:resultant}

\section{Motywacja i definicja}
Dla dwóch równań wielomianowych $f(x)=0$ i $g(x)=0$ chcemy
rozstrzygać istnienie wspólnego rozwiązania bez obliczania pierwiastków.
Nad ciałem odpowiedzi udziela największy wspólny dzielnik, lecz
jego obliczenie wymaga dzielenia wielomianów. Wypadkowa jest natomiast
\emph{wyznacznikiem} macierzy zbudowanej z samych współczynników;
można ją więc obliczać również nad pierścieniami innymi niż ciała.

\begin{definicja}\label{aa:def:sylvester}
Niech $A$ będzie pierścieniem przemiennym z jedynką oraz
$m,n\ge1$ liczbami całkowitymi. Ustalmy dwa ciągi
$a_0,\ldots,a_m,b_0,\ldots,b_n\in A$ i wielomiany
$f=\sum_{i=0}^m a_iX^i$, $g=\sum_{j=0}^n b_jX^j$.
Na tym etapie $m,n$ są \emph{stopniami formalnymi}: wolno mieć
$a_m=0$ albo $b_n=0$. Przez $A[X]_{<k}$ rozumiemy wolny
$A$-moduł wielomianów o stopniu mniejszym niż $k$ (łącznie z zerowym),
z bazą $(1,X,\ldots,X^{k-1})$.
\emph{Macierz Sylvestera} $S_{m,n}(f,g)$ jest macierzą homomorfizmu
\[
 \Phi_{f,g}:A[X]_{<n}\oplus A[X]_{<m}\longrightarrow A[X]_{<m+n},
 \qquad (u,v)\longmapsto uf+vg,
\]
w bazach: najpierw współczynniki $u$, następnie $v$, a w przeciwdziedzinie
współczynniki rosnących potęg $X$.
Definiujemy
\[
 \operatorname{Res}(f,g)=\det S_{m,n}(f,g),\qquad
 R_{m,n}(f,g)=(-1)^{mn}\operatorname{Res}(f,g).
\]
W tym rozdziale symbol $R(f,g)$ bez indeksów oznacza $R_{m,n}(f,g)$,
gdy $m=\deg f$, $n=\deg g$ są rzeczywistymi stopniami dodatnimi.
\end{definicja}

\begin{uwaga}\label{aa:uw:res-sign}
Dla $f=a_1X+a_0$ i $g=b_1X+b_0$ otrzymujemy
\[
 S_{1,1}(f,g)=\begin{pmatrix}a_0&b_0\\a_1&b_1\end{pmatrix},\quad
 \operatorname{Res}(f,g)=a_0b_1-a_1b_0,\quad
 R(f,g)=a_1b_0-a_0b_1.
\]
Zatem $R$ odpowiada konwencji $a_1g(-a_0/a_1)$, gdy $a_1$
jest odwracalne. Nazwa \emph{wypadkowa} będzie dalej odnosiła się do $R$.
Formalne dopisywanie zerowych współczynników wiodących może zmienić
wartość wyznacznika i czynnik znakowy. Na przykład dla
$f=2X+1$ i $g=X+3$ nad $\mathbb Q$ mamy $R_{1,1}(f,g)=5$,
natomiast przy formalnym stopniu $2$ wielomianu $g$ otrzymujemy
$R_{1,2}(f,g)=10$, zaś przy formalnym stopniu $2$ wielomianu
$f$ mamy $R_{2,1}(f,g)=-5$. W kryterium wspólnego czynnika stosujemy więc
rzeczywiste stopnie $m,n$; inny wybór stosujemy jawnie przy pochodnej.
Ponadto dla homomorfizmu pierścieni $\rho:A\to B$ zachodzi
$\rho(R_{m,n}(f,g))=R_{m,n}(\rho(f),\rho(g))$ przy
zachowaniu \emph{stopni formalnych}, ponieważ wyznacznik jest
wielomianem o współczynnikach całkowitych w wyrazach macierzy.
\end{uwaga}

\section{Wspólny czynnik i wspólny pierwiastek}
\begin{twierdzenie}[Kryterium wspólnego czynnika]\label{aa:tw:res-common}
Niech $K$ będzie ciałem, a $f,g\in K[X]$ niezerowymi wielomianami
stopni odpowiednio $m,n\ge1$. Przez $\gcd(f,g)$ rozumiemy
moniczny największy wspólny dzielnik. Następujące warunki są równoważne:
\begin{enumerate}
\item $R(f,g)=0$;
\item $\gcd(f,g)$ ma dodatni stopień;
\item $f$ i $g$ mają wspólny pierwiastek w pewnym rozszerzeniu algebraicznym $K$.
\end{enumerate}
\end{twierdzenie}
\begin{proof}
Macierz Sylvestera jest kwadratowa rozmiaru $m+n$. Ponieważ $K$
jest ciałem, jej wyznacznik znika dokładnie wtedy, gdy
$\Phi_{f,g}$ ma niezerowy element w jądrze. To znaczy, że
istnieją $u\in K[X]_{<n}$ i $v\in K[X]_{<m}$, nieoba zerowe,
takie że $uf+vg=0$.

Jeżeli $d=\gcd(f,g)$ ma stopień $r\ge1$, napiszmy $f=df_1$,
$g=dg_1$. Podstawiając $u=g_1$ i $v=-f_1$, otrzymujemy
$uf+vg=0$ oraz $\deg u=n-r<n$ i $\deg v=m-r<m$.
Ponieważ $f,g$ są niezerowe, para $(u,v)$ jest niezerowa,
więc wyznacznik znika.

W przeciwnym kierunku załóżmy, że $\gcd(f,g)=1$ i
$uf=-vg$. Z lematu Euklidesa w pierścieniu $K[X]$ wynika,
że $f\mid v$. Ponieważ $\deg v<m=\deg f$, mamy $v=0$.
Zatem $uf=0$, co w pierścieniu bez dzielników zera oznacza
$u=0$. Jądro jest zerowe, wyznacznik nie znika.

Na koniec w domknięciu algebraicznym $\overline K$ każdy
nierozkładalny czynnik wielomianu ma pierwiastek. Wspólny
czynnik dodatniego stopnia dostarcza wspólnego pierwiastka.
Odwrotnie, jeśli $f(\alpha)=g(\alpha)=0$ dla $\alpha\in\overline K$,
minimalny wielomian $\alpha$ nad $K$ dzieli zarówno $f$, jak i $g$.
\end{proof}

\section{Wzór przez pierwiastki bez założenia rozdzielczości}
\begin{twierdzenie}[Wzór przez pierwiastki]\label{aa:tw:res-roots}
Niech $K$ będzie ciałem, a $f,g\in K[X]$ wielomianami o stopniach
$m,n\ge1$ i niezerowych współczynnikach wiodących $a_m,b_n$.
Jeżeli w ciele rozkładu $L$ wielomianów $f$ i $g$ mamy
\[
 f=a_m\prod_{i=1}^m(X-\alpha_i),\qquad
 g=b_n\prod_{j=1}^n(X-\beta_j),
\]
gdzie pierwiastki powtarzamy zgodnie z ich krotnościami, to
\[
 R(f,g)=a_m^n\prod_{i=1}^m g(\alpha_i)
 =a_m^nb_n^m\prod_{i=1}^m\prod_{j=1}^n(\alpha_i-\beta_j).
\]
Nie zakładamy, że $f$ lub $g$ są rozdzielcze.
\end{twierdzenie}
\begin{proof}
Rozszerzenie skalarów z $K$ do $L$ nie zmienia wyznacznika,
więc możemy pracować nad $L$. Przyjmijmy następującą kolejność
kroków, w których kontrolujemy również znak i współczynnik wiodący.

\emph{Krok 1: odwrócenie bloków.}
W dziedzinie $\Phi_{f,g}$ przestawiamy bloki tak, aby pierwszym
argumentem było $v\in L[X]_{<m}$, a drugim $u\in L[X]_{<n}$.
Przestawienie bloków długości $m$ i $n$ ma wyznacznik
$(-1)^{mn}$, dlatego macierz $\Psi(v,u)=vg+uf$ ma
wyznacznik $R(f,g)$.

\emph{Krok 2: dzielenie przez $f$.}
Każdy $h\in L[X]_{<m+n}$ daje się jednoznacznie zapisać
$h=r+qf$, gdzie $\deg r<m$ i $\deg q<n$.
Ta zmiana współrzędnych jest liniowym izomorfizmem
\[
 T:L[X]_{<m+n}\longrightarrow L[X]_{<m}\oplus L[X]_{<n},
 \qquad h\longmapsto(r,q).
\]
W bazach rosnących potęg macierz odwrotna do $T$ jest blokowo
trójkątna: kolumny $1,\ldots,X^{m-1}$ oraz
$f,Xf,\ldots,X^{n-1}f$ mają wyznacznik $a_m^n$.
Zatem $\det T=a_m^{-n}$.

\emph{Krok 3: działanie na ilorazie.}
Niech $M_g:L[X]/(f)\to L[X]/(f)$ oznacza mnożenie przez klasę
$g$. Dla $\Psi(v,u)=vg+uf$ reszta z dzielenia przez $f$
zależy tylko od $v$ i jest reprezentantem $M_g([v])$.
Współczynnik ilorazu przy $f$ ma postać $u+Q(v)$ dla pewnego
liniowego $Q$. Macierz $T\circ\Psi$ ma więc postać
$\begin{psmallmatrix}M_g&0\\Q&I_n\end{psmallmatrix}$.
Jej wyznacznik jest $\det M_g$. Z drugiej strony jest równy
$a_m^{-n}R(f,g)$, co daje
$R(f,g)=a_m^n\det M_g$.

\emph{Krok 4: uwzględnienie pierwiastków wielokrotnych.}
Niech $P_0=1$ oraz $P_k=\prod_{i=1}^k(X-\alpha_i)$
dla $1\le k<m$. Ponieważ $P_k$ jest moniczny stopnia $k$,
klasy $[P_0],\ldots,[P_{m-1}]$ tworzą bazę $L[X]/(f)$.
W tej bazie mnożenie przez $X$ spełnia
\[
 X[P_{k-1}]=[P_k]+\alpha_k[P_{k-1}]\quad(1\le k<m),
 \qquad X[P_{m-1}]=\alpha_m[P_{m-1}],
\]
bo $P_m=f/a_m$ jest zerem w ilorazie.
Macierz $M_X$ jest zatem trójkątna o przekątnej
$\alpha_1,\ldots,\alpha_m$, nawet jeśli pierwiastki się powtarzają.
Macierz $M_g=g(M_X)$ jest trójkątna o przekątnej
$g(\alpha_1),\ldots,g(\alpha_m)$, więc
$\det M_g=\prod_i g(\alpha_i)$. Daje to pierwszą równość;
druga wynika z podstawienia rozkładu $g$.
\end{proof}

\begin{wniosek}[Symetria wypadkowej]\label{aa:wn:res-sym}
Niech $K$ będzie ciałem, a $f,g\in K[X]$ niezerowymi wielomianami
stopni odpowiednio $m,n\ge1$. Wówczas
$R(g,f)=(-1)^{mn}R(f,g)$.
\end{wniosek}
\begin{proof}
We wzorze z Twierdzenia~\ref{aa:tw:res-roots} oba wyrażenia
zawierają czynnik $a_m^nb_n^m$. Przy zamianie kolejności
pierwiastków każdy z $mn$ czynników $\alpha_i-\beta_j$
zmienia znak. Stąd dokładnie czynnik $(-1)^{mn}$.
\end{proof}

\section{Pochodna formalna i wyróżnik}
\begin{twierdzenie}[Wyróżnik jako wypadkowa]\label{aa:tw:disc-resultant}
Niech $K$ będzie dowolnym ciałem, a $f\in K[X]$ wielomianem
monicznym stopnia $n\ge2$. Wówczas
\[
 \operatorname{disc}(f)=(-1)^{n(n-1)/2}R_{n,n-1}(f,f').
\]
Stopień $n-1$ przy $f'$ jest tutaj formalny: dopuszczamy,
że rzeczywisty stopień pochodnej jest mniejszy lub że $f'=0$.
\end{twierdzenie}
\begin{proof}
Uzasadnimy użycie stopnia formalnego. Kroki 1--3 dowodu
Twierdzenia~\ref{aa:tw:res-roots} wymagają wyłącznie, aby
$f$ miał rzeczywisty stopień $n$ i niezerowy współczynnik wiodący,
a drugi wielomian miał stopień \emph{co najwyżej równy}
zadanemu stopniowi formalnemu $r\ge1$. Istotnie, przestrzeń $u$ ma wtedy
wymiar $r$, dzielenie przez $f$ nadal daje blok jednostkowy
rozmiaru $r$, a ostatni krok dowodu używa jedynie mnożenia przez
klasę drugiego wielomianu. Wobec tego przy $r=n-1$ i monicznym
$f$ otrzymujemy, także gdy $f'=0$, równość
$R_{n,n-1}(f,f')=\prod_{i=1}^n f'(\alpha_i)$.
Jeśli $f=\prod_{i=1}^n(X-\alpha_i)$, to z reguły Leibniza
\[
 f'(\alpha_i)=\prod_{j\ne i}(\alpha_i-\alpha_j),
\]
równość ta pozostaje poprawna również przy powtarzających się
pierwiastkach. W iloczynie po wszystkich $i$ każda para
$i<j$ daje czynnik
$(\alpha_i-\alpha_j)(\alpha_j-\alpha_i)
=-(\alpha_i-\alpha_j)^2$. Par jest $n(n-1)/2$.
To dokładnie definicja wyróżnika, także w dodatniej charakterystyce.
\end{proof}

\begin{przyklad}[Eliminacja zmiennej bez wyznaczania pierwiastków]\label{aa:prz:res-elimination}
Niech $A=\mathbb Z[s,t]$, $f(X)=X^2-s$ i $g(X)=X-t$.
W ciele ułamków $\mathbb Q(s,t)$ mamy $R(g,f)=f(t)=t^2-s$,
a symetria wypadkowej, ponieważ $2\cdot1$ jest parzyste,
daje $R(f,g)=t^2-s$. Możemy też obliczyć wyznacznik
bez przechodzenia do ciała:
\[
 S_{2,1}(f,g)=
 \begin{pmatrix}-s&-t&0\\0&1&-t\\1&0&1\end{pmatrix},
 \qquad\det S_{2,1}(f,g)=t^2-s.
\]
Aby poprawnie odczytać warunek istnienia wspólnego pierwiastka,
wybierzmy dowolne ciało $F$ i dowolne wartości $s_0,t_0\in F$.
Po specjalizacji $s\mapsto s_0$, $t\mapsto t_0$ para równań
$x^2-s_0=0$, $x-t_0=0$ ma wspólne rozwiązanie w rozszerzeniu $F$
dokładnie wtedy, gdy $s_0=t_0^2$: jedynym kandydatem jest $x=t_0$.
Nie trzeba wybierać pierwiastka z $s_0$. Zgodność tej specjalizacji
z obliczonym wyznacznikiem wynika również z uwagi o homomorfizmach
pierścieni po Definicji~\ref{aa:def:sylvester}.
\end{przyklad}

\begin{przyklad}[Pierwiastek wielokrotny]\label{aa:prz:res-repeated}
Niech $K=\mathbb Q$, $f=(X-1)^2$ i $g=X+2$.
Twierdzenie~\ref{aa:tw:res-roots} daje
$R(f,g)=g(1)^2=9$; jest to prawidłowy wynik mimo
nierozdzielczości $f$. Dla pochodnej $f'=2(X-1)$ mamy
$R(f,f')=f'(1)^2=0$ i $\operatorname{disc}(f)=0$.
Pierwsza wypadkowa nie wykrywa wielokrotności pierwiastka
$f$, lecz drugie obliczenie już tak.
\end{przyklad}

\begin{przyklad}[Kontrprzykład nad pierścieniem z dzielnikami zera]\label{aa:prz:res-zero-divisors}
Niech $A=\mathbb Z/4\mathbb Z$ i $f=g=2X+1$.
Wyznacznik macierzy Sylvestera jest równy zero, ponieważ
jej obie kolumny są identyczne. Jednak żaden $x\in A$
nie spełnia $f(x)=0$: wartość $2x+1$ jest zawsze klasą
nieparzystą. Pierścień $A$ nie jest ciałem, a jego niezerowy
współczynnik wiodący $2$ nie jest odwracalny.
Co więcej, $f$ jest elementem odwracalnym $A[X]$, ponieważ
$(1+2X)(1-2X)=1-4X^2=1$. Nie ma więc pierwiastka nawet
w żadnej niezerowej przemiennej $A$-algebrze: obraz elementu
odwracalnego pozostaje odwracalny. Nie wolno bez dodatkowych
założeń przenosić równoważności z
Twierdzenia~\ref{aa:tw:res-common} na dowolne pierścienie.
\end{przyklad}

% AA-34 | Audyt I, 22 XI 2026. Zrodlo po kontroli niezaleznej. Zachowano etykiety historyczne.
\chapter[Rezolwenta sześcienna i grupy Galois]{Rezolwenta sześcienna i grupy Galois wielomianów czwartego stopnia}\label{aa:ch:quartic-resolvent}

\section{Motywacja: skąd bierze się wielomian stopnia trzeciego?}
Niech $K$ będzie ciałem, $f\in K[X]$ wielomianem rozdzielczym
stopnia czwartego, a $G$ grupą Galois jego ciała rozkładu nad $K$.
Działa ona na czterech różnych pierwiastkach $f$.
Grupa $S_4$ działa nie tylko na czterech punktach, lecz również
na \emph{trzech} podziałach tych punktów na dwie pary.
Z tego działania otrzymamy wielomian stopnia trzeciego,
który pozwoli rozpoznać część informacji o $G$. Pierwiastkami
nowego wielomianu będą odpowiednie wyrażenia z pierwiastków $f$.

\section{Działanie na podziałach czteroelementowego zbioru}
\begin{definicja}\label{aa:def:quartic-pairings}
Niech $\Omega=\{1,2,3,4\}$. Przez $\mathcal P$ oznaczamy zbiór
trzech podziałów $\Omega$ na dwie nieuporządkowane pary:
\[
 P_1=\{\{1,2\},\{3,4\}\},\quad
 P_2=\{\{1,3\},\{2,4\}\},\quad
 P_3=\{\{1,4\},\{2,3\}\}.
\]
Dla $\sigma\in S_4$ i $P=\{\{i,j\},\{k,l\}\}$ definiujemy
$\sigma(P)=\{\{\sigma(i),\sigma(j)\},\{\sigma(k),\sigma(l)\}\}$.
Jest to działanie grupy $S_4$ na $\mathcal P$.
\end{definicja}

\begin{center}
\begin{tikzpicture}[scale=.73, every node/.style={font=\small}]
\foreach \sh/\lab in {0/$P_1$,4.7/$P_2$,9.4/$P_3$}{
  \begin{scope}[xshift=\sh cm]
    \node at (1,2.7){\lab};
    \coordinate (a) at (0,2); \coordinate (b) at (2,2);
    \coordinate (c) at (0,0); \coordinate (d) at (2,0);
    \foreach \nm/\pt in {1/a,2/b,3/c,4/d}{\node[circle,draw,fill=white,inner sep=2pt] at (\pt) {\nm};}
  \end{scope}}
\begin{scope}[line width=1.15pt]
 \draw[blue!65!black] (0.35,2)--(1.65,2);
 \draw[blue!65!black] (0.35,0)--(1.65,0);
 \draw[blue!65!black] (5.05,1.7)--(5.05,.3);
 \draw[blue!65!black] (6.35,1.7)--(6.35,.3);
 \draw[blue!65!black] (9.75,1.7)--(11.05,.3);
 \draw[blue!65!black] (9.75,.3)--(11.05,1.7);
\end{scope}
\end{tikzpicture}
\end{center}
\begin{uwaga}
Na rysunku każdy odcinek łączy dwa elementy jednej pary.
Przecięcie odcinków w $P_3$ nie oznacza dodatkowego wierzchołka.
\end{uwaga}

\begin{twierdzenie}[Jądro działania na podziałach]\label{aa:tw:s4-v4}
Homomorfizm $\rho:S_4\to S(\mathcal P)\cong S_3$ jest surjektywny,
a jego jądro jest grupą Kleina
\[
 V_4=\{1,(12)(34),(13)(24),(14)(23)\}.
\]
W konsekwencji $S_4/V_4\cong S_3$.
\end{twierdzenie}
\begin{proof}
Permutacja $(12)$ ustala $P_1$ i zamienia $P_2$ z $P_3$.
Permutacja $(23)$ ustala $P_3$ i zamienia $P_1$ z $P_2$.
Zatem obraz $\rho$ zawiera transpozycje $(P_2\ P_3)$
oraz $(P_1\ P_2)$. Dwie sąsiednie transpozycje generują
całe $S(\mathcal P)$: ich iloczyny dają oba cykle długości trzy,
a ich sprzężenie daje trzecią transpozycję. To dowodzi surjektywności.
Każda z trzech podwójnych transpozycji zachowuje wszystkie trzy
podziały, co można sprawdzić, stosując podaną definicję działania.
Otrzymujemy więc $V_4\subseteq\ker\rho$. Z twierdzenia o homomorfizmie
$|\ker\rho|=|S_4|/|S_3|=24/6=4$.
Ponieważ $V_4$ ma dokładnie cztery elementy, jest całym jądrem.
\end{proof}

\section{Konstrukcja i jawne obliczenie rezolwenty}
\begin{definicja}\label{aa:def:quartic-resolvent}
Niech $K$ będzie ciałem, a
$f(X)=X^4+aX^3+bX^2+cX+d\in K[X]$ wielomianem monicznym.
Wybierzmy pierwiastki $\alpha_1,\ldots,\alpha_4$ wielomianu $f$
w pewnym ciele rozkładu, z uwzględnieniem krotności. Kładziemy
\[
 \beta_1=\alpha_1\alpha_2+\alpha_3\alpha_4,\quad
 \beta_2=\alpha_1\alpha_3+\alpha_2\alpha_4,\quad
 \beta_3=\alpha_1\alpha_4+\alpha_2\alpha_3.
\]
Wielomian $r_f(Y)=\prod_{i=1}^{3}(Y-\beta_i)$ nazywamy
\emph{rezolwentą sześcienną} $f$.
Zmiana numeracji $\alpha_i$ permutuje trzy wartości $\beta_i$,
więc nie zmienia samego wielomianu $r_f$.
\end{definicja}

\begin{twierdzenie}[Wzór na rezolwentę]\label{aa:tw:quartic-resolvent-formula}
Dla każdego ciała $K$ i każdego monicznego
$f=X^4+aX^3+bX^2+cX+d\in K[X]$ zachodzi
\[
 r_f(Y)=Y^3-bY^2+(ac-4d)Y-(a^2d+c^2-4bd).
\]
W szczególności $r_f\in K[Y]$ bez względu na charakterystykę $K$.
\end{twierdzenie}
\begin{proof}
Wprowadzamy elementarne funkcje symetryczne $e_1,e_2,e_3,e_4$
zmiennych $\alpha_1,\ldots,\alpha_4$. Ze wzorów Viète'a
(Twierdzenie~\ref{aa:tw:vieta-general}) otrzymujemy
$e_1=-a$, $e_2=b$, $e_3=-c$ oraz $e_4=d$.
Wyznaczymy kolejno trzy współczynniki $r_f$.

\emph{Suma pierwiastków.} W $\beta_1+\beta_2+\beta_3$
każdy iloczyn $\alpha_i\alpha_j$ ($i<j$) występuje dokładnie raz.
Zatem $\beta_1+\beta_2+\beta_3=e_2$.

\emph{Suma iloczynów parami.} Na przykład
\[
\begin{aligned}
 \beta_1\beta_2={}&\alpha_1^2\alpha_2\alpha_3
 +\alpha_1\alpha_2^2\alpha_4
 +\alpha_1\alpha_3^2\alpha_4
 +\alpha_2\alpha_3\alpha_4^2.
\end{aligned}
\]
Dwa pozostałe iloczyny powstają przez permutacje indeksów.
Łącznie otrzymujemy wszystkie $12$ jednomianów postaci
$\alpha_i^2\alpha_j\alpha_k$, gdzie $i,j,k$ są parami różne,
każdy dokładnie raz. Z kolei w $e_1e_3$ powstają te same
$12$ jednomianów oraz $4e_4$: w iloczynie $\alpha_i$ z sumą
czterech jednomianów trzeciego stopnia jeden z nich nie zawiera
$\alpha_i$, a pozostałe trzy go zawierają.
Stąd $\sum_{i<j}\beta_i\beta_j=e_1e_3-4e_4$.

\emph{Iloczyn trzech pierwiastków.} Oznaczmy
$T_i=\prod_{j\ne i}\alpha_j$ dla $1\le i\le4$.
Po wymnożeniu trzech nawiasów $\beta_1\beta_2\beta_3$
dostajemy osiem jednomianów: cztery postaci $T_i^2$ i cztery
postaci $e_4\alpha_i^2$. Dlatego
\[
 \beta_1\beta_2\beta_3=\sum_{i=1}^4 T_i^2
 +e_4\sum_{i=1}^4\alpha_i^2.
\]
Z tożsamości
$e_1^2=\sum_i\alpha_i^2+2e_2$ wynika
$\sum_i\alpha_i^2=e_1^2-2e_2$.
Ponieważ $T_iT_j=e_4\alpha_k\alpha_l$ dla
$\{i,j,k,l\}=\{1,2,3,4\}$,
mamy też
$e_3^2=\sum_iT_i^2+2e_4e_2$.
Po podstawieniu i zebraniu wyrazów otrzymujemy
\[
 \beta_1\beta_2\beta_3=e_1^2e_4+e_3^2-4e_2e_4.
\]
Wszystkie te równości są tożsamościami wielomianowymi
o współczynnikach całkowitych, więc obowiązują również w
charakterystyce $2$ i $3$. Zastępując $e_i$ przez współczynniki
$f$ oraz stosując wzory Viète'a do $r_f$, dostajemy tezę.
\end{proof}

\begin{twierdzenie}[Rezolwenta i grupa Galois]\label{aa:tw:quartic-resolvent-action}
Niech $K$ będzie ciałem charakterystyki różnej od $2$,
a $f\in K[X]$ monicznym, rozdzielczym wielomianem stopnia czwartego.
Niech $L$ będzie ciałem rozkładu $f$, a $G=\operatorname{Gal}(L/K)$
traktujmy jako podgrupę $S_4$ przez działanie na
$\alpha_1,\ldots,\alpha_4$. Wówczas $\beta_1,\beta_2,\beta_3$
są parami różne, a grupa Galois rezolwenty jest izomorficzna
z $\rho(G)\le S_3$. Jądrem działania $G$ na pierwiastkach
rezolwenty jest $G\cap V_4$.
\end{twierdzenie}
\begin{proof}
Wprost z definicji mamy
\[
 \beta_1-\beta_2=(\alpha_1-\alpha_4)(\alpha_2-\alpha_3),\quad
 \beta_1-\beta_3=(\alpha_1-\alpha_3)(\alpha_2-\alpha_4),
\]
a różnicę $\beta_2-\beta_3$ otrzymujemy analogicznie.
Ponieważ wielomian $f$ jest rozdzielczy, każda różnica
dwóch różnych $\alpha_i$ jest niezerowa. Ciało nie ma dzielników
zera, więc również $\beta_i\ne\beta_j$ dla $i\ne j$.

Dla $\sigma\in G$ zachodzi $\sigma(\beta_i)=\beta_j$ dokładnie
wtedy, gdy $\sigma$ wysyła podział $P_i$ na $P_j$.
Ponieważ $\beta_i$ są różne, oba działania mają to samo jądro:
$G\cap\ker\rho=G\cap V_4$.
Ciało $M=K(\beta_1,\beta_2,\beta_3)$ jest ciałem rozkładu
$r_f$ zawartym w $L$. Zatem $M/K$ jest rozszerzeniem Galois.
Ograniczenie automorfizmów $L/K$ do $M/K$ jest surjektywne
na mocy odpowiedniości Galois, a jego jądrem jest $G\cap V_4$.
Twierdzenie o izomorfizmie daje
$\operatorname{Gal}(M/K)\cong G/(G\cap V_4)\cong\rho(G)$.
\end{proof}

\begin{wniosek}[Równość wyróżników]\label{aa:wn:quartic-resolvent-disc}
Niech $K$ będzie dowolnym ciałem, a $f\in K[X]$ monicznym
wielomianem stopnia czwartego. Dla rezolwenty $r_f$ z
Definicji~\ref{aa:def:quartic-resolvent} zachodzi
\[
\operatorname{disc}(r_f)=\operatorname{disc}(f).
\]
W szczególności $f$ jest rozdzielczy wtedy i tylko wtedy,
gdy $r_f$ jest rozdzielczy.
\end{wniosek}
\begin{proof}
Dla pierwiastków $\alpha_1,\ldots,\alpha_4$ wielomianu $f$,
uwzględnionych z krotnościami, mamy
\[
\begin{aligned}
\beta_1-\beta_2&=(\alpha_1-\alpha_4)(\alpha_2-\alpha_3),\\
\beta_1-\beta_3&=(\alpha_1-\alpha_3)(\alpha_2-\alpha_4),\\
\beta_2-\beta_3&=(\alpha_1-\alpha_2)(\alpha_3-\alpha_4).
\end{aligned}
\]
Po pomnożeniu kwadratów tych trzech równości otrzymujemy
dokładnie sześć czynników tworzących definicję
$\operatorname{disc}(f)$. Jest to tożsamość wielomianowa,
ważna także w charakterystyce $2$ i dla pierwiastków
wielokrotnych. Ostatnie stwierdzenie wynika z
Twierdzenia~\ref{aa:tw:disc-separable}.
\end{proof}

\begin{uwaga}\label{aa:uw:quartic-char2}
Równości dla $\beta_i-\beta_j$ są poprawne także w charakterystyce
$2$. Zatem rozdzielczy $f$ ma trzy różne pierwiastki rezolwenty
również w tej charakterystyce, a opis obrazu $\rho(G)$ pozostaje
prawdziwy. Ograniczenie $\operatorname{char}K\ne2$ jest natomiast
konieczne przy późniejszym wykorzystaniu kryterium wyróżnika
z Twierdzenia~\ref{aa:tw:disc-alternating}.
\end{uwaga}

\section{Możliwe grupy dla nierozkładalnego wielomianu}
\begin{twierdzenie}[Przypadek nierozkładalny]\label{aa:tw:quartic-classification}
Niech $K$ będzie ciałem charakterystyki różnej od $2$,
a $f\in K[X]$ monicznym, nierozkładalnym i rozdzielczym
wielomianem stopnia czwartego. Wtedy jego grupa Galois
jest izomorficzna z jedną z grup $C_4$, $V_4$, $D_4$, $A_4$, $S_4$.
Jeżeli ponadto $r_f$ jest nierozkładalny nad $K$, to
$G\cong A_4$ albo $G\cong S_4$; w takim przypadku
$G\cong A_4$ dokładnie wtedy, gdy $\operatorname{disc}(f)$
jest kwadratem niezerowego elementu $K$.
\end{twierdzenie}
\begin{proof}
Ponieważ $f$ jest nierozkładalny i rozdzielczy, grupa Galois
działa przechodnio na czterech jego pierwiastkach.
Dla jednego pierwiastka $\alpha_1$ twierdzenie o orbicie
i stabilizatorze daje $|G|=4|G_{\alpha_1}|$.
Jako podgrupa $S_4$ grupa $G$ ma rząd dzielący $24$,
więc $|G|$ należy do $\{4,8,12,24\}$.

Jeżeli $|G|=4$, stabilizator pierwiastka jest trywialny,
a działanie jest regularne. Grupa rzędu cztery jest albo
cykliczna $C_4$ (gdy zawiera element rzędu cztery), albo
izomorficzna z $V_4$ (gdy każdy nietrywialny element ma rząd dwa).
W drugim przypadku każdy element $z\in G$ spełnia $z^2=1$.
Dla dowolnych $x,y\in G$ mamy więc
$(xy)^{-1}=xy$, a jednocześnie $(xy)^{-1}=y^{-1}x^{-1}=yx$.
Stąd $xy=yx$, co potwierdza przemienność. Argument obejmuje
także przypadki $x=y$ oraz $x=1$ lub $y=1$.

Jeżeli $|G|=8$, grupa $G$ jest podgrupą Sylowa rzędu $8$
w $S_4$. Grupa symetrii kwadratu wpisanego w cztery wierzchołki
ma osiem elementów: generują ją obrót $(1234)$ i odbicie $(24)$
z relacjami $r^4=s^2=1$, $srs=r^{-1}$.
Na mocy twierdzenia Sylowa każda podgrupa rzędu $8$
jest z nią sprzężona, zatem izomorficzna z $D_4$.

Jeżeli $|G|=12$, podgrupa $G$ ma indeks dwa i jest normalna.
Działanie $S_4$ na dwóch warstwach $S_4/G$ daje
nietrywialny homomorfizm $S_4\to C_2$ z jądrem $G$.
Każda transpozycja ma taki sam obraz, bo transpozycje są sprzężone.
Obraz ten nie może być tożsamością, gdyż transpozycje generują
$S_4$. Homomorfizm jest więc znakiem permutacji, a $G=A_4$.
Jeżeli $|G|=24$, to $G=S_4$.

Załóżmy teraz, że rezolwenta jest nierozkładalna. Z poprzedniego
twierdzenia wiemy, że jej grupa Galois działa przechodnio
na trzech pierwiastkach; ma więc rząd podzielny przez $3$.
Jest ona obrazem $G$, toteż $3\mid |G|$.
Wśród rozpatrywanych rzędów $4,8,12,24$ warunek ten
spełniają tylko $12$ i $24$. Pozostają grupy $A_4$ i $S_4$.
Ponieważ $\operatorname{char}K\ne2$, kryterium wyróżnika
z Twierdzenia~\ref{aa:tw:disc-alternating} mówi, że
$G\le A_4$ dokładnie wtedy, gdy wyróżnik jest kwadratem w $K$.
Rozdzielczość $f$ gwarantuje, że wyróżnik jest niezerowy.
\end{proof}

\begin{wniosek}[Rozkład rezolwenty a grupa Galois]
\label{aa:wn:quartic-resolvent-factorization}
Niech $K$ będzie ciałem charakterystyki różnej od $2$,
a $f\in K[X]$ monicznym, nierozkładalnym i rozdzielczym
wielomianem stopnia czwartego. Wówczas rezolwenta $r_f$
jest rozdzielcza i zachodzi dokładnie jeden z trzech przypadków:
\begin{enumerate}
\item jeśli $r_f$ rozkłada się na trzy czynniki liniowe w $K[Y]$,
      to $\operatorname{Gal}(f/K)\cong V_4$;
\item jeśli $r_f$ ma jeden czynnik liniowy i jeden nierozkładalny
      czynnik kwadratowy, to $\operatorname{Gal}(f/K)\cong C_4$
      albo $D_4$;
\item jeśli $r_f$ jest nierozkładalna, to
      $\operatorname{Gal}(f/K)\cong A_4$ albo $S_4$.
\end{enumerate}
W ostatnim przypadku wyróżnik $f$ rozstrzyga między $A_4$
i $S_4$ zgodnie z Twierdzeniem~\ref{aa:tw:quartic-classification}.
\end{wniosek}
\begin{proof}
Z Twierdzenia~\ref{aa:tw:quartic-resolvent-action} pierwiastki
$r_f$ są różne, a działanie jej grupy Galois to obraz
$\rho(G)\subseteq S_3$, gdzie $G=\operatorname{Gal}(f/K)$.
Trzy typy rozkładu odpowiadają odpowiednio orbitom
$1+1+1$, $1+2$ i $3$ na zbiorze trzech pierwiastków.
W pierwszym przypadku $\rho(G)=1$, zatem $G\subseteq V_4$;
z przechodniości $G$ na czterech pierwiastkach $f$ wynika
$4\mid |G|$, więc $G=V_4$. W drugim przypadku obraz
ma rząd $2$. Na mocy klasyfikacji w
Twierdzeniu~\ref{aa:tw:quartic-classification} może być wtedy
tylko $C_4$ albo $D_4$: obie te grupy mają obraz rzędu $2$,
podczas gdy $V_4$ ma obraz trywialny, $A_4$ obraz rzędu $3$,
a $S_4$ obraz rzędu $6$.
Trzeci przypadek jest już drugą częścią wspomnianego twierdzenia.
Ponieważ rozdzielcza rezolwenta stopnia trzeciego ma tylko
wymienione typy rozkładu, lista jest wyczerpująca.
\end{proof}

\begin{przyklad}[Cykliczna grupa $C_4$]\label{aa:prz:quartic-c4}
Niech $f(X)=\Phi_5(X)=X^4+X^3+X^2+X+1\in\mathbb Q[X]$.
Po podstawieniu $X=T+1$ otrzymujemy
\[
 f(T+1)=T^4+5T^3+10T^2+10T+5.
\]
Kryterium Eisensteina dla liczby pierwszej $5$ dowodzi
nierozkładalności $f$ nad $\mathbb Q$.
Ciałem rozkładu $f$ jest $\mathbb Q(\zeta_5)$, a grupa Galois
jest izomorficzna z
$(\mathbb Z/5\mathbb Z)^\times\cong C_4$;
izomorfizm ten wynika z
Twierdzenia~\ref{aa:tw:cyclotomic-galois}.
Z Twierdzenia~\ref{aa:tw:quartic-resolvent-formula},
przy $a=b=c=d=1$, mamy
\[
 r_f(Y)=Y^3-Y^2-3Y+2=(Y-2)(Y^2+Y-1).
\]
Czynnik kwadratowy jest nierozkładalny nad $\mathbb Q$,
gdyż jego wyróżnik wynosi $5$, niebędące kwadratem
liczby wymiernej. Jest to przypadek $1+2$ poprzedniego
wniosku. Pokazuje on również, że rozkład $1+2$
rezolwenty nie wystarcza do rozróżnienia grup $C_4$ i $D_4$.
\end{przyklad}

\begin{przyklad}[Grupa $S_4$]\label{aa:prz:quartic-s4}
Weźmy $f=X^4-X-1\in\mathbb Q[X]$. Współczynniki
$a=b=0$, $c=d=-1$ dają rezolwentę
$r_f(Y)=Y^3+4Y-1$.
Ponieważ jest ona moniczna o współczynnikach całkowitych,
każdy jej pierwiastek wymierny musiałby należeć do $\{1,-1\}$.
Tymczasem $r_f(1)=4$ i $r_f(-1)=-6$,
więc $r_f$ jest nierozkładalny jako wielomian sześcienny.

Redukcja $f$ modulo $2$ jest równa $X^4+X+1$.
Nie ma pierwiastka w $\mathbb F_2$ (wartości w $0,1$ wynoszą $1$),
a jedyny moniczny nierozkładalny wielomian kwadratowy
nad $\mathbb F_2$ to $X^2+X+1$.
Jego kwadrat wynosi $X^4+X^2+1$, a nie $X^4+X+1$.
Gdyby nasz wielomian stopnia czwartego był rozkładalny,
miałby czynnik liniowy lub dwa czynniki nierozkładalne stopnia dwa;
oba przypadki są wykluczone. Zatem $f$ jest nierozkładalny nad
$\mathbb F_2$, a w konsekwencji nad $\mathbb Q$.

Wyróżnik monicznego wielomianu $X^4+cX+d$ można obliczyć
jako wypadkową z pochodną $4X^3+c$.
Jeśli $t^3=-c/4$, to
$t^4+ct+d=d+(3c/4)t$, więc iloczyn wartości
$f(t)$ po trzech pierwiastkach pochodnej jest
$d^3-27c^4/256$: dla pierwiastków $t_1,t_2,t_3$
wielomianu $X^3+c/4$ mamy $\sum_i t_i=0$,
$\sum_{i<j}t_it_j=0$ i $t_1t_2t_3=-c/4$.
Wzór na wypadkową daje
$R(f',f)=4^4\prod_i f(t_i)$, zaś symetria wypadkowej
(Twierdzenie~\ref{aa:wn:res-sym}) daje
$R(f,f')=R(f',f)$, bo $4\cdot3$ jest parzyste.
Otrzymujemy
$\operatorname{disc}(X^4+cX+d)=256d^3-27c^4$.
Dla $c=d=-1$ jest to $-283$, liczba ujemna, a zatem
niekwadrat w $\mathbb Q$. Na mocy
Twierdzenia~\ref{aa:tw:quartic-classification} mamy
$\operatorname{Gal}(f/\mathbb Q)\cong S_4$.
\end{przyklad}

\begin{przyklad}[Kontrprzykład: rezolwenta rozkładalna, a wielomian nierozkładalny]\label{aa:prz:quartic-v4}
Wielomian $f=X^4-2$ jest nierozkładalny nad $\mathbb Q$
na mocy kryterium Eisensteina dla liczby pierwszej $2$.
Wzór na rezolwentę daje
$r_f(Y)=Y^3+8Y=Y(Y^2+8)$, więc rezolwenta jest rozkładalna.
Niech $\alpha=\sqrt[4]{2}>0$. Ciało rozkładu jest
$L=\mathbb Q(\alpha,i)$, ponieważ pierwiastki $f$
to $\alpha,-\alpha,i\alpha,-i\alpha$.
Mamy $[\mathbb Q(\alpha):\mathbb Q]=4$, a
$\mathbb Q(\alpha)\subset\mathbb R$, zatem $i\notin\mathbb Q(\alpha)$
i $[L:\mathbb Q]=8$.
Ponieważ $[L:\mathbb Q(i)]=8/2=4$, minimalnym
wielomianem $\alpha$ nad $\mathbb Q(i)$ jest $X^4-2$.
Liczba $i\alpha$ jest również jego pierwiastkiem, dlatego
$\mathbb Q(i)$-zanurzenie wyznaczone przez
$\alpha\mapsto i\alpha$ jest automorfizmem $L$.
Oznaczmy je przez $r$. Automorfizm ten ustala $i$,
a sprzężenie zespolone $s$ ustala $\alpha$ i wysyła
$i$ na $-i$. Spełniają one relacje
$r^4=s^2=1$, $srs=r^{-1}$, a osiem automorfizmów
$r^k,sr^k$ ($0\le k<4$) jest różnych.
Grupa Galois ma więc postać $D_4$ i nie jest ani $A_4$, ani $S_4$.
Przykład pokazuje, dlaczego w drugiej części
Twierdzenia~\ref{aa:tw:quartic-classification} trzeba
zakładać nierozkładalność rezolwenty.
\end{przyklad}

\begin{przyklad}[Grupa Kleina i całkowicie rozkładalna rezolwenta]\label{aa:prz:quartic-klein}
Niech $\alpha=\sqrt2+\sqrt3$ oraz
$f=X^4-10X^2+1$. Ponieważ $\alpha^2=5+2\sqrt6$,
liczba $\alpha$ spełnia $(\alpha^2-5)^2=24$,
czyli $f(\alpha)=0$.
Ponadto $\alpha^{-1}=\sqrt3-\sqrt2$, skąd
\[
 \sqrt2=\frac{\alpha-\alpha^{-1}}2,\qquad
 \sqrt3=\frac{\alpha+\alpha^{-1}}2.
\]
Zatem $\mathbb Q(\alpha)=\mathbb Q(\sqrt2,\sqrt3)$.
Jest to rozszerzenie stopnia czwartego: $\sqrt3$ nie należy
do $\mathbb Q(\sqrt2)$, bo równość
$(u+v\sqrt2)^2=3$ dla $u,v\in\mathbb Q$
wymagałaby $uv=0$ i następnie $u^2=3$ lub $v^2=3/2$,
co jest niemożliwe dla liczb wymiernych.
Wobec tego $f$ jest nierozkładalny.
Jego pierwiastki to $\pm\sqrt2\pm\sqrt3$, wszystkie należą
do $\mathbb Q(\sqrt2,\sqrt3)$. Grupa Galois jest więc $V_4$:
niezależnie można zmieniać znaki obu pierwiastków kwadratowych.
Wzór z Twierdzenia~\ref{aa:tw:quartic-resolvent-formula} daje
\[
 r_f(Y)=Y^3+10Y^2-4Y-40=(Y+10)(Y-2)(Y+2).
\]
Wszystkie trzy pierwiastki rezolwenty są wymierne,
co odpowiada trywialnemu obrazowi $V_4$ w $S_3$.
\end{przyklad}

% AA-35 | Audyt historyczny V: pelne uzasadnienie podnoszenia Frobeniusa.
\chapter[Redukcja modulo p i cykle Frobeniusa]{Redukcja modulo liczb pierwszych i cykle Frobeniusa}\label{aa:ch:frobenius-reduction}

\section{Motywacja i dobra redukcja}
Niech $f\in\mathbb Z[X]$ będzie monicznym wielomianem rozdzielczym.
W ciele rozkładu nad $\mathbb Q$ trudno zwykle wypisać automorfizmy
jego pierwiastków. Znacznie łatwiej rozłożyć $f$ modulo niewielką
liczbę pierwszą $p$. Pokażemy, że stopnie czynników takiego rozkładu
wyznaczają typ cykli \emph{pewnego} elementu grupy Galois nad
$\mathbb Q$. Ważne jest zastrzeżenie o dobrej redukcji:
pierwiastki nie mogą się ze sobą zlać po redukcji modulo $p$.

\begin{definicja}\label{aa:def:good-reduction}
Niech $f\in\mathbb Z[X]$ będzie monicznym wielomianem stopnia
$n\ge1$, o niezerowym wyróżniku $\operatorname{disc}(f)$.
Liczbę pierwszą $p$ nazywamy \emph{dobrą dla $f$}, jeśli
$p\nmid\operatorname{disc}(f)$.
Przez $\bar f\in\mathbb F_p[X]$ oznaczamy wielomian powstały
przez redukcję każdego współczynnika $f$ modulo $p$.
Ponieważ $f$ jest moniczny, $\deg\bar f=\deg f$.
\end{definicja}

\begin{lemat}\label{aa:lem:good-squarefree}
Niech $f\in\mathbb Z[X]$ będzie moniczny stopnia $n\ge1$
i niech jego wyróżnik będzie niezerowy. Dla dowolnej liczby
pierwszej $p$ następujące warunki są równoważne:
\begin{enumerate}
\item $p$ jest liczbą pierwszą dobrą dla $f$;
\item $\gcd(\bar f,\bar f')=1$ w $\mathbb F_p[X]$;
\item $\bar f$ nie ma pierwiastka wielokrotnego w domknięciu
algebraicznym $\mathbb F_p$.
\end{enumerate}
\end{lemat}
\begin{proof}
Najpierw uzasadnimy zgodność wyróżnika z redukcją współczynników.
Wielomian
$D(X_1,\ldots,X_n)=\prod_{i<j}(X_i-X_j)^2$
jest symetryczny i ma współczynniki całkowite. Podstawowe twierdzenie
o wielomianach symetrycznych
(Twierdzenie~\ref{aa:tw:fundamental-symmetric}), zastosowane nad
$\mathbb Z$, daje $D=P(e_1,\ldots,e_n)$ z pewnym
$P\in\mathbb Z[T_1,\ldots,T_n]$.
Dla każdego monicznego wielomianu
$g=X^n+c_1X^{n-1}+\cdots+c_n$ nad dowolnym ciałem wzory Viète'a
wyrażają jego wyróżnik jako
$P(-c_1,c_2,\ldots,(-1)^nc_n)$.
Ponieważ $P$ ma współczynniki całkowite, otrzymujemy
\[
 \overline{\operatorname{disc}(f)}
 =\operatorname{disc}(\bar f)\quad\text{w }\mathbb F_p.
\]
Z definicji dobrej redukcji i z kryterium wyróżnika
(Twierdzenie~\ref{aa:tw:disc-separable}) wynika zatem
równoważność warunków (1) i (3).

Pozostaje porównać (2) i (3). Niech $\alpha$ będzie pierwiastkiem
$\bar f$ w domknięciu algebraicznym $\overline{\mathbb F}_p$.
Piszemy $\bar f=(X-\alpha)^m h$, gdzie $m\ge1$ i
$h(\alpha)\ne0$. Z reguły różniczkowania iloczynu wynika,
że gdy $m=1$, to $(\bar f)'(\alpha)=h(\alpha)\ne0$,
a gdy $m\ge2$, to $(\bar f)'(\alpha)=0$.
Zatem $\bar f$ i $(\bar f)'$ mają wspólny pierwiastek
w domknięciu algebraicznym dokładnie wtedy, gdy $\bar f$
ma pierwiastek wielokrotny. Nad ciałem jest to równoważne
temu, że ich największy wspólny dzielnik ma dodatni stopień.
Rozumowanie obejmuje także przypadek $(\bar f)'=0$:
ponieważ $\deg\bar f\ge1$, wówczas
$\gcd(\bar f,0)=\bar f$ nie jest jednostką, a każdy
pierwiastek $\bar f$ jest wielokrotny.

Nie korzystaliśmy z nieuprawnionego utożsamiania
macierzy Sylvestera o różnych rozmiarach: po redukcji
modulo $p$ stopień pochodnej może bowiem się zmniejszyć.
\end{proof}

\section{Cykle Frobeniusa nad ciałami skończonymi}
\begin{twierdzenie}[Cykle Frobeniusa]\label{aa:tw:frob-cycles}
Niech $q=p^r$, gdzie $p$ jest liczbą pierwszą, a $r\ge1$.
Niech $h\in\mathbb F_q[X]$ będzie monicznym, rozdzielczym
wielomianem stopnia $n\ge1$, o rozkładzie
$h=h_1\cdots h_s$ na parami różne moniczne czynniki
nierozkładalne, $\deg h_i=d_i$.
Niech $E$ będzie ciałem rozkładu $h$ nad $\mathbb F_q$.
Wtedy automorfizm $\varphi:E\to E$, $x\mapsto x^q$,
działa na $n$ pierwiastkach $h$ jako iloczyn rozłącznych
cykli długości $d_1,\ldots,d_s$.
\end{twierdzenie}
\begin{proof}
Wybierzmy $i\in\{1,\ldots,s\}$ i pierwiastek $\alpha_i$
wielomianu $h_i$. Ponieważ $h_i$ jest nierozkładalny,
$[\mathbb F_q(\alpha_i):\mathbb F_q]=d_i$.
Z twierdzenia o ciałach skończonych
(Twierdzenie~\ref{aa:tw:finite-frobenius-roots})
$\mathbb F_q(\alpha_i)$ ma $q^{d_i}$ elementów,
a wszystkie jego elementy spełniają $x^{q^{d_i}}=x$.
Stąd $\varphi^{d_i}(\alpha_i)=\alpha_i$.

Pokażemy, że żaden dodatni $k<d_i$ nie spełnia
$\varphi^k(\alpha_i)=\alpha_i$. Gdyby tak było,
$\alpha_i$ należałby do zbioru pierwiastków
$X^{q^k}-X$, czyli do ciała $\mathbb F_{q^k}$.
W ustalonym domknięciu algebraicznym
$\mathbb F_q(\alpha_i)$ byłoby wtedy podciałem
jednoznacznie określonego ciała $\mathbb F_{q^k}$.
Ze wzoru na stopnie wieży otrzymalibyśmy
$d_i\mid k$, co jest niemożliwe dla $0<k<d_i$.
Zatem orbita $\alpha_i$ pod iteracjami $\varphi$
ma dokładnie $d_i$ elementów.

Wszystkie te elementy są pierwiastkami $h_i$:
automorfizm Frobeniusa ustala współczynniki z $\mathbb F_q$.
Ponieważ $h_i$ jest rozdzielczy i ma dokładnie $d_i$
pierwiastków, orbita jest całym zbiorem jego pierwiastków.
Różne czynniki $h_i$ mają rozłączne zbiory pierwiastków,
więc odpowiadające im cykle również są rozłączne.
\end{proof}

\section{Podnoszenie Frobeniusa do grupy Galois nad $\mathbb Q$}
W następnym dowodzie wykorzystamy pierścień generowany
przez \emph{pierwiastki rozpatrywanego wielomianu}.
Dzięki temu nie musimy zakładać bez dowodu struktury całego
pierścienia liczb całkowitych ciała rozkładu.

\begin{lemat}[Skończony pierścień generowany przez pierwiastki]
\label{aa:lem:frob-finite-order}
Niech $f\in\mathbb Z[X]$ będzie monicznym wielomianem
stopnia $n\ge1$, a $\alpha_1,\ldots,\alpha_n$ jego
pierwiastkami w pewnym ciele rozkładu $L/\mathbb Q$,
liczonymi z krotnościami. Wtedy pierścień
$B=\mathbb Z[\alpha_1,\ldots,\alpha_n]$ jest skończenie
generowany jako $\mathbb Z$-moduł, każdy jego element jest
całkowity nad $\mathbb Z$, a $B\cap\mathbb Q=\mathbb Z$.
Jeżeli $f$ jest rozdzielczy i $G=\operatorname{Gal}(L/\mathbb Q)$,
to $B$ jest stabilny względem $G$ i $B^G=\mathbb Z$.
\end{lemat}
\begin{proof}
Każdy $\alpha_i$ spełnia moniczne równanie
$\alpha_i^n+c_{n-1}\alpha_i^{n-1}+\cdots+c_0=0$
z całkowitymi współczynnikami. Pozwala ono zastępować
potęgi $\alpha_i^k$ dla $k\ge n$ kombinacjami całkowitymi
niższych potęg. Stosując to kolejno do wszystkich zmiennych,
otrzymujemy, że $B$ jest generowany nad $\mathbb Z$ przez
najwyżej $n^n$ monomów
$\alpha_1^{e_1}\cdots\alpha_n^{e_n}$, $0\le e_i<n$.

Niech $b\in B$, a $z_1,\ldots,z_N$ będzie skończoną rodziną
generatorów $B$ jako $\mathbb Z$-modułu.
Ponieważ $b z_i\in B$, istnieje macierz
$C=(c_{ij})\in M_N(\mathbb Z)$ taka, że
$b z_i=\sum_j c_{ij}z_j$. Dla wektora kolumnowego
$z=(z_1,\ldots,z_N)^T$ zapisujemy to jako $bz=Cz$.
Niech $P(T)=\det(TI_N-C)\in\mathbb Z[T]$.
Z twierdzenia Cayleya–Hamiltona mamy $P(C)=0$,
a z równości $b^kz=C^kz$ wynika $P(b)z=0$.
Ponieważ $z_i$ generują $B$ i $1\in B$, dostajemy $P(b)=0$.
Wielomian $P$ jest moniczny, więc $b$ jest całkowity
nad $\mathbb Z$.

Jeśli $b=u/v\in B\cap\mathbb Q$, gdzie $u,v\in\mathbb Z$,
$\gcd(u,v)=1$ i $v>0$, to z monicznego równania
$b^d+c_{d-1}b^{d-1}+\cdots+c_0=0$ po pomnożeniu
przez $v^d$ wynika $v\mid u^d$. Względna pierwszość
$u,v$ wymusza $v=1$. Stąd $B\cap\mathbb Q=\mathbb Z$.

Jeśli $f$ jest rozdzielczy, każdy $\sigma\in G$ permutuje
$\alpha_1,\ldots,\alpha_n$, a zatem zachowuje $B$.
Elementy stałe działania $G$ w $L$ tworzą $\mathbb Q$,
więc $B^G=B\cap\mathbb Q=\mathbb Z$.
\end{proof}

\begin{lemat}[Podnoszenie Frobeniusa]
\label{aa:lem:frob-lift}
Niech $L/\mathbb Q$ będzie ciałem rozkładu monicznego,
rozdzielczego wielomianu $f\in\mathbb Z[X]$, niech
$G=\operatorname{Gal}(L/\mathbb Q)$ i
$B=\mathbb Z[\alpha_1,\ldots,\alpha_n]$ jak wyżej.
Dla dowolnej liczby pierwszej $p$ i dowolnego ideału
maksymalnego $\mathfrak p$ pierścienia $B$, zawierającego $p$,
istnieje $\sigma\in G$ spełniające
$\sigma(\mathfrak p)=\mathfrak p$ i
\[
 \overline{\sigma(b)}=\bar b^{\,p}
 \quad\text{dla każdego }b\in B,
\]
gdzie kreska oznacza klasę modulo $\mathfrak p$.
\end{lemat}
\begin{proof}
\emph{Krok 1: skończoność ciał reszt.}
Z Lematu~\ref{aa:lem:frob-finite-order} pierścień $B$
jest skończonym $\mathbb Z$-modułem. Ponieważ $1/p$ nie jest
całkowity nad $\mathbb Z$, nie może należeć do $B$;
zatem $pB\ne B$. Iloraz $B/pB$ jest niezerowym skończonym
pierścieniem, gdyż jest generowany jako grupa addytywna
przez skończenie wiele elementów rzędu dzielącego $p$.
Każdy jego ideał maksymalny ma iloraz będący skończonym
ciałem. W szczególności
$\kappa=B/\mathfrak p$ jest skończonym rozszerzeniem $\mathbb F_p$.
W pierścieniu $B$ nad $p$ leży tylko skończenie wiele
ideałów maksymalnych,
bo odpowiadają one ideałom maksymalnym skończonego pierścienia
$B/pB$.

\emph{Krok 2: zastosowanie chińskiego twierdzenia o resztach.}
Dla różnych ideałów maksymalnych $\mathfrak p_i$, $\mathfrak p_j$
mamy $\mathfrak p_i+\mathfrak p_j=B$.
Wybierzmy dla każdego $j\ne i$ element $t_{ij}\in\mathfrak p_j$
taki, że $t_{ij}\equiv1\pmod{\mathfrak p_i}$.
Iloczyn $e_i=\prod_{j\ne i}t_{ij}$, równy $1$, gdy nie ma
innych ideałów, jest równy $1$ modulo $\mathfrak p_i$
i $0$ modulo każdego innego ideału.
Kombinacje liniowe $e_i$ z wybranymi reprezentantami klas
pokazują, że można zadawać niezależnie reszty modulo wszystkie
ideały maksymalne nad $p$.

\emph{Krok 3: przechodniość działania $G$ na ideałach nad $p$.}
Gdyby pewien ideał maksymalny $\mathfrak q\supset pB$
nie należał do orbity $\mathfrak p$, to
$\mathfrak p\ne\sigma^{-1}(\mathfrak q)$ dla każdego $\sigma\in G$.
Z Kroku 2 można wybrać $b\in B$ takie, że
$b\equiv0\pmod{\mathfrak p}$ oraz
$b\equiv1\pmod{\sigma^{-1}(\mathfrak q)}$
dla wszystkich $\sigma\in G$.
Norma $N=\prod_{\sigma\in G}\sigma(b)$ należy do
$B^G=\mathbb Z$. Ponieważ $b\in\mathfrak p$, także
$N\in\mathfrak p\cap\mathbb Z=p\mathbb Z\subset\mathfrak q$.
Jednocześnie dla każdego $\sigma\in G$ zachodzi
$\sigma(b)\notin\mathfrak q$, gdyż
$b\notin\sigma^{-1}(\mathfrak q)$.
Ideał maksymalny jest pierwszy, więc ich iloczyn
$N\notin\mathfrak q$, sprzeczność. Zatem działanie jest przechodnie.

\emph{Krok 4: podniesienie automorfizmu Frobeniusa.}
Niech $D=\{\sigma\in G:\sigma(\mathfrak p)=\mathfrak p\}$.
Wybierzmy element $c\in\kappa^\times$ generujący
multiplikatywną grupę $\kappa^\times$; wówczas
$\kappa=\mathbb F_p(c)$. Dzięki Krokowi 2 istnieje $b\in B$
o klasie $c$ modulo $\mathfrak p$ i klasie $0$ modulo
każdego innego ideału nad $p$.
Rozpatrzmy
$P(T)=\prod_{\sigma\in G}(T-\sigma(b))\in\mathbb Z[T]$.
Współczynniki tego wielomianu należą do $B$ i są stałe
względem $G$, a więc należą do $B^G=\mathbb Z$.
W $\kappa$ element $c$ jest pierwiastkiem redukcji $P$,
bo dla czynnika z $\sigma=1$ mamy $c-\bar b=0$.
Zatem $c^p$ również jest pierwiastkiem redukcji $P$,
ponieważ jej współczynniki należą do $\mathbb F_p$.
Wobec tego $c^p=\overline{\sigma(b)}$ dla pewnego $\sigma\in G$.
Ponieważ $c^p\ne0$, mamy $b\notin\sigma^{-1}(\mathfrak p)$.
Wszystkie ideały $\sigma^{-1}(\mathfrak p)$ są maksymalne i leżą nad $p$;
z naszego wyboru $b$ wynika więc
$\sigma^{-1}(\mathfrak p)=\mathfrak p$, czyli $\sigma\in D$.
Ponieważ $\sigma(\mathfrak p)=\mathfrak p$, działanie $\sigma$
na $B$ indukuje dobrze określony automorfizm ciała
$\kappa=B/\mathfrak p$; jego odwrotność indukuje
$\sigma^{-1}$, a więc jest on bijektywny. Automorfizm ten
ustala $\mathbb F_p$ i wysyła generator $c$ na $c^p$.
Każdy element $\kappa=\mathbb F_p(c)$ jest funkcją wymierną
$c$ o współczynnikach w $\mathbb F_p$, więc automorfizm ten
zgadza się na całym $\kappa$ z Frobeniusem
$x\mapsto x^p$. To daje żądane $\sigma$ dla wszystkich $b\in B$.
\end{proof}

\begin{twierdzenie}[Redukcja i grupa Galois]\label{aa:tw:dedekind-cycle}
Niech $f\in\mathbb Z[X]$ będzie moniczny, rozdzielczy,
stopnia $n\ge1$, a $L$ jego ciałem rozkładu nad $\mathbb Q$.
Niech $p$ będzie dowolną dobrą liczbą pierwszą dla $f$.
Jeżeli w $\mathbb F_p[X]$ wielomian $\bar f$ rozkłada się
na moniczne czynniki nierozkładalne stopni
$d_1,\ldots,d_s$, to $\operatorname{Gal}(L/\mathbb Q)$
zawiera element, którego działanie na $n$ różnych
pierwiastkach $f$ ma rozkład na rozłączne cykle
długości $d_1,\ldots,d_s$.
\end{twierdzenie}
\begin{proof}
Weźmy pierścień $B$ z Lematu~\ref{aa:lem:frob-finite-order}
i ideał maksymalny $\mathfrak p$ zawierający $pB$;
istnieje on, gdyż $B/pB$ jest niezerowym skończonym pierścieniem.
Z Lematu~\ref{aa:lem:frob-lift} istnieje $\sigma\in G$
stabilizujące $\mathfrak p$ i działające na
$\kappa=B/\mathfrak p$ jako Frobenius $x\mapsto x^p$.

Wszystkie $\alpha_i$ należą do $B$.
Gdyby dwie ich redukcje $\bar\alpha_i,\bar\alpha_j$
przy $i\ne j$ były równe, to $\alpha_i-\alpha_j\in\mathfrak p$.
W iloczynie
$\operatorname{disc}(f)=\prod_{i<j}(\alpha_i-\alpha_j)^2$
znalazłby się czynnik z $\mathfrak p$, więc
$\operatorname{disc}(f)\in\mathfrak p\cap\mathbb Z=p\mathbb Z$,
wbrew dobrej redukcji. Dlatego $\bar\alpha_i$ są
$n$ różnymi pierwiastkami $\bar f$ w $\kappa$.
Ponieważ $\bar f$ ma stopień $n$, są to wszystkie pierwiastki
w dowolnym domknięciu algebraicznym $\mathbb F_p$.

Dla każdego $i$ mamy
$\overline{\sigma(\alpha_i)}=\bar\alpha_i^{\,p}$.
Tak więc bijekcja $\alpha_i\mapsto\bar\alpha_i$
przekształca permutację pierwiastków wywołaną przez $\sigma$
w permutację wywołaną przez Frobeniusa.
Na mocy Twierdzenia~\ref{aa:tw:frob-cycles} jej długości cykli
są dokładnie równe $d_1,\ldots,d_s$.
\end{proof}

\begin{uwaga}\label{aa:uw:dedekind-prereq}
W dowodzie użyto: twierdzenia Cayleya–Hamiltona dla macierzy
nad $\mathbb Z$, podstawowego twierdzenia teorii Galois,
cykliczności grupy multiplikatywnej ciała skończonego,
chińskiego twierdzenia o resztach oraz twierdzeń o pierścieniach
i ciałach ilorazowych. Kluczowe fakty o pierścieniu
$B=\mathbb Z[\alpha_1,\ldots,\alpha_n]$ i o podnoszeniu
Frobeniusa zostały powyżej udowodnione; nie jest wymagane
przyjmowanie skończoności całego pierścienia liczb całkowitych
$\mathcal O_L$ jako nieudowodnionego twierdzenia.
\end{uwaga}

\begin{wniosek}[Parzystość permutacji Frobeniusa]
\label{aa:wn:frob-parity}
Niech $f\in\mathbb Z[X]$ będzie monicznym, rozdzielczym
wielomianem stopnia $n\ge2$ i niech $p$ będzie liczbą pierwszą
niedzielącą $\operatorname{disc}(f)$. Załóżmy, że redukcja
$\bar f\in\mathbb F_p[X]$ ma dokładnie $s$ różnych monicznych
czynników nierozkładalnych, o stopniach $d_1,\ldots,d_s$.
Wtedy grupa Galois wielomianu $f$ nad $\mathbb Q$ zawiera
permutację pierwiastków o znaku $(-1)^{n-s}$. W szczególności,
jeśli $n-s$ jest nieparzyste, wyróżnik $\operatorname{disc}(f)$
nie jest kwadratem liczby wymiernej. Jeżeli wyróżnik jest kwadratem
liczby wymiernej, to dla każdej dobrej liczby pierwszej $p$
liczba $n-s$ musi być parzysta.
\end{wniosek}
\begin{proof}
Z Twierdzenia~\ref{aa:tw:dedekind-cycle} wynika, że istnieje
automorfizm $\sigma$ ciała rozkładu, który działa na zbiorze
pierwiastków jako iloczyn $s$ rozłącznych cykli długości
$d_1,\ldots,d_s$. Każdy cykl długości $d_i$ jest iloczynem
$d_i-1$ transpozycji, a więc ma znak $(-1)^{d_i-1}$.
Znak całej permutacji wynosi zatem
\[
 \operatorname{sgn}(\sigma)=\prod_{i=1}^s(-1)^{d_i-1}
 =(-1)^{d_1+\cdots+d_s-s}=(-1)^{n-s}.
\]
Jeśli $n-s$ jest nieparzyste, grupa Galois nie zawiera się
w $A_n$. Ponieważ $\operatorname{char}\mathbb Q=0$, z kryterium wyróżnika
(Twierdzenie~\ref{aa:tw:disc-alternating}) wynika, że grupa
zawiera się w $A_n$ dokładnie wtedy,
gdy wyróżnik jest kwadratem w $\mathbb Q$.
Stąd wyróżnik nie jest kwadratem. Ostatnie zdanie wynika
przez przeciwstawienie.
\end{proof}

\section{Przykłady i granice stosowalności}
\begin{przyklad}[Dwa rodzaje cykli]\label{aa:prz:modp-cycles}
Niech $f=X^3-X-1\in\mathbb Z[X]$.
Ze wzoru na wyróżnik wielomianu sześciennego
(Twierdzenie~\ref{aa:tw:cubic-disc}) mamy
$\operatorname{disc}(f)=-4(-1)^3-27(-1)^2=4-27=-23$.
Zatem zarówno $2$, jak i $5$ są liczbami dobrymi.

Modulo $2$ dostajemy $X^3+X+1$. Jego wartości dla
$0$ i $1$ wynoszą $1$, więc nie ma pierwiastków
w $\mathbb F_2$. Wielomian stopnia trzy bez pierwiastka
nad ciałem jest nierozkładalny, a więc otrzymujemy cykl
$(123)$ długości trzy.

Modulo $5$ mamy
\[
 \bar f=(X-2)(X^2+2X+3).
\]
Pierwszy czynnik jest liniowy, a drugi ma wyróżnik
$2^2-4\cdot3= -8\equiv2\pmod5$.
Kwadraty w $\mathbb F_5$ to $0,1,4$, więc czynnik
kwadratowy jest nierozkładalny. Otrzymujemy typ cykli
$(1,2)$, czyli transpozycję.
Dla $p=2$ mamy $n=3$, $s=1$, zatem znak otrzymanego
$3$-cyklu jest $(-1)^{3-1}=1$. Dla $p=5$ mamy $s=2$ i
znak permutacji o typie $(1,2)$ wynosi $(-1)^{3-2}=-1$.
Wniosek~\ref{aa:wn:frob-parity} potwierdza więc, że
wyróżnik $-23$ nie jest kwadratem liczby wymiernej.

Grupa Galois jest podgrupą $S_3$ zawierającą element
rzędu $3$ i element rzędu $2$, więc jej rząd
jest podzielny przez $6$. Jako dzielnik $6$ musi być
równy $6$. Zatem $\operatorname{Gal}(f/\mathbb Q)\cong S_3$.
\end{przyklad}

\begin{przyklad}[Dlaczego warunek dobrej redukcji jest konieczny]\label{aa:prz:bad-prime}
Dla $f=X^2-5$ wyróżnik wynosi $20$, więc liczba pierwsza
$2$ jest zła. Modulo $2$ mamy
$\bar f=X^2+1=(X+1)^2$.
Dwa pierwiastki nad $\mathbb Q$ redukują się więc do
jednego pierwiastka podwójnego. Nie wolno odczytywać
z zapisu $(X+1)^2$ cyklu długości dwa: wielomian
nie jest rozdzielczy po redukcji. Twierdzenie
\ref{aa:tw:dedekind-cycle} nie daje dla $p=2$ typu cyklu.

Tak samo dla wielomianu $g=X^3-X-1$ z
Przykładu~\ref{aa:prz:modp-cycles} liczba $23$ jest zła,
gdyż $23\mid\operatorname{disc}(g)=-23$.
W $\mathbb F_{23}[X]$ otrzymujemy
$\bar g=(X-3)(X-10)^2$. Rzeczywiście,
$(X-10)^2=X^2+3X+8$ modulo $23$, a pomnożenie
przez $X-3$ daje $X^3-X-1$.
Podwójny pierwiastek $10$ nie wyznacza cyklu długości dwa:
redukcja skleja tu dwa różne pierwiastki z charakterystyki zero.
\end{przyklad}

\begin{przyklad}[Rozkład na czynniki liniowe nie oznacza grupy trywialnej]
Niech $f=X^3-X-1$ jak wyżej. Weźmy $p=59$; liczba ta jest dobra, bo $59\nmid23$. Rachunek
modulo $59$ daje $\bar f(4)=\bar f(13)=\bar f(42)=0$.
Wszystkie trzy liczby są różne modulo $59$, więc
$\bar f=(X-4)(X-13)(X-42)$ w $\mathbb F_{59}[X]$.
Twierdzenie~\ref{aa:tw:dedekind-cycle}
dostarcza zatem w grupie $S_3$ elementu o typie $(1,1,1)$,
czyli elementu identycznościowego. Nie wynika z tego,
że cała grupa Galois jest trywialna: wiemy już,
że jest izomorficzna z $S_3$. Informacja z jednej liczby
pierwszej dotyczy istnienia \emph{pewnego} elementu,
a nie wszystkich elementów grupy.
\end{przyklad}

\begin{przyklad}[Wielomian rozkładalny nad $\mathbb Q$]
\label{aa:prz:frob-reducible}
Nie trzeba zakładać nierozkładalności $f$ nad $\mathbb Q$.
Rozważmy
$f=(X^2-2)(X-3)=X^3-3X^2-2X+6$.
Wyróżnik iloczynu dwóch monicznych wielomianów, których
zbiory pierwiastków są rozłączne, jest iloczynem
ich wyróżników i kwadratu wypadkowej. W naszym przypadku
$\operatorname{disc}(X^2-2)=8$, a
$R(X^2-2,X-3)=7$, więc
$\operatorname{disc}(f)=8\cdot 7^2=392$.
Dla $p=5$ wyróżnik nie znika modulo $5$, natomiast
$\bar f=(X^2-2)(X-3)$, gdzie $X^2-2$ jest
nierozkładalny nad $\mathbb F_5$, ponieważ $2$ nie
jest kwadratem w tym ciele. Otrzymujemy więc element
grupy Galois o typie cykli $(2,1)$.
Można to też zobaczyć bezpośrednio: ciało rozkładu
jest równe $\mathbb Q(\sqrt 2)$, a jego nietrywialny
automorfizm zamienia $\sqrt2$ z $-\sqrt2$ i ustala $3$.
\end{przyklad}

% AA-36 | Audyt historyczny VI: pełne dowody i przykłady
\chapter{Wielomian ogólny i jego grupa Galois}\label{aa:ch:generic-polynomial}
\section{Algebraiczna niezależność współczynników}
Dla ustalonego wielomianu nad $\mathbb Q$ grupa Galois zależy od
jego współczynników. Na przykład dwa nierozkładalne wielomiany stopnia
trzeciego mogą mieć grupy $S_3$ i $C_3$. Chcemy oddzielić własności
wspólne \emph{wszystkim formalnie niezależnym współczynnikom} od
szczególnych równań, które pojawiają się po podstawieniu konkretnych
liczb. W tym celu zamiast wybierać współczynniki w $k$, wprowadzamy
zmienne transcendentalne nad $k$.

\begin{definicja}\label{aa:def:generic-poly}
Niech $k$ będzie ciałem, $n\ge2$ liczbą całkowitą, a
$t_1,\ldots,t_n$ zmiennymi algebraicznie niezależnymi nad $k$;
znaczy to, że dla każdego niezerowego
$P\in k[T_1,\ldots,T_n]$ zachodzi $P(t_1,\ldots,t_n)\ne0$.
Przez $k(t_1,\ldots,t_n)$ rozumiemy ciało ułamków
$k[t_1,\ldots,t_n]$. Wielomian
\[
 F(X)=X^n-t_1X^{n-1}+t_2X^{n-2}-\cdots+(-1)^nt_n
 \in k(t_1,\ldots,t_n)[X]
\]
nazywamy \emph{ogólnym wielomianem monicznym stopnia $n$}.
Alternatywna konwencja znaków, z niezależnymi współczynnikami
przy kolejnych potęgach $X$, daje ten sam typ wielomianu: zmiana
$t_i\mapsto(-1)^it_i$ jest automorfizmem pierścienia wielomianów.
\end{definicja}

\begin{lemat}[Niezależność elementarnych funkcji symetrycznych]
\label{aa:lem:generic-independence}
Niech $x_1,\ldots,x_n$ będą algebraicznie niezależne nad ciałem $k$.
Dla $1\le j\le n$ oznaczmy
$e_j=\sum_{1\le i_1<\cdots<i_j\le n}x_{i_1}\cdots x_{i_j}$.
Wówczas $e_1,\ldots,e_n$ są algebraicznie niezależne nad $k$.
\end{lemat}
\begin{proof}
Rozważmy homomorfizm
$\varphi:k[T_1,\ldots,T_n]\to k[X_1,\ldots,X_n]$
określony wzorem $\varphi(T_j)=e_j(X_1,\ldots,X_n)$.
Podstawowe twierdzenie o wielomianach symetrycznych
(Twierdzenie~\ref{aa:tw:fundamental-symmetric}) stwierdza nie tylko,
że każdy wielomian symetryczny jest wielomianem w $e_j$, lecz także,
że taki zapis jest \emph{jednoznaczny}. Stąd $\ker\varphi=(0)$.
Ponadto odwzorowanie
$\psi:k[X_1,\ldots,X_n]\to k(x_1,\ldots,x_n)$
określone przez $X_i\mapsto x_i$ jest injektywne: gdyby
$P(x_1,\ldots,x_n)=0$ dla niezerowego $P$, otrzymalibyśmy relację
algebraiczną między $x_i$. Zatem także $\psi\circ\varphi$ jest
injektywne. Jest to dokładnie warunek niezależności algebraicznej
$e_1,\ldots,e_n$.
\end{proof}

\section{Obliczenie grupy Galois}
\begin{twierdzenie}[Grupa Galois wielomianu ogólnego]
\label{aa:tw:generic-sn}
Dla każdego ciała $k$ i każdej liczby całkowitej $n\ge2$ ogólny
wielomian moniczny $F\in k(t_1,\ldots,t_n)[X]$ jest nierozkładalny
i rozdzielczy, a jego grupa Galois nad $k(t_1,\ldots,t_n)$
jest izomorficzna z grupą symetryczną $S_n$.
\end{twierdzenie}
\begin{proof}
\emph{Krok 1: ciało zawierające formalne pierwiastki.}
Weźmy algebraicznie niezależne zmienne $x_1,\ldots,x_n$ nad $k$,
połóżmy $L=k(x_1,\ldots,x_n)$ i
$K=k(e_1(x_1,\ldots,x_n),\ldots,e_n(x_1,\ldots,x_n))$.
Lemat~\ref{aa:lem:generic-independence} daje izomorfizm
$k(t_1,\ldots,t_n)\cong K$, określony przez $t_j\mapsto e_j$.
Będziemy więc obliczać grupę Galois nad $K$.
Ze wzorów Viète'a wynika tożsamość
\[
 f(X)=\prod_{i=1}^n(X-x_i)
 =X^n-e_1X^{n-1}+\cdots+(-1)^ne_n\in K[X].
\]
Każdy $x_i$ jest pierwiastkiem $f$; ponieważ są to różne zmienne,
$x_i-x_j\ne0$ dla $i\ne j$, więc $f$ ma $n$ różnych pierwiastków.
Jest zatem rozdzielczy również w dodatniej charakterystyce.
Ponadto $L=K(x_1,\ldots,x_n)$: z $k\subseteq K$ mamy
$L=k(x_1,\ldots,x_n)\subseteq K(x_1,\ldots,x_n)$,
a z $K\subseteq L$ wynika odwrotne zawieranie.
Tym samym $L$ jest ciałem rozkładu $f$ nad $K$.

\emph{Krok 2: permutacje zmiennych rzeczywiście dają automorfizmy.}
Dla $\sigma\in S_n$ podstawienie $x_i\mapsto x_{\sigma(i)}$
indukuje automorfizm pierścienia $k[x_1,\ldots,x_n]$:
jego odwrotnością jest podstawienie wyznaczone przez $\sigma^{-1}$.
Ponieważ automorfizm przenosi niezerowe wielomiany na niezerowe,
rozszerza się jednoznacznie na ciało ułamków $L$ przez
$P/Q\mapsto\sigma(P)/\sigma(Q)$.
Wielomiany $e_j$ są symetryczne, więc $\sigma(e_j)=e_j$;
automorfizm ustala punktowo $K$.
Różne permutacje dają różne automorfizmy, bo dla pewnego $i$
mają różne wartości na $x_i$. Otrzymaliśmy więc co najmniej
$n!$ elementów $\operatorname{Gal}(L/K)$.

\emph{Krok 3: nie ma innych automorfizmów.}
Jeżeli $\tau$ jest automorfizmem $L$ ustalającym $K$, to
$f(\tau(x_i))=\tau(f(x_i))=0$.
Wobec rozdzielczości $f$ obraz $\tau(x_i)$ jest jednym z $x_j$.
Automorfizm jest injektywny, więc obrazy $n$ różnych pierwiastków
są parami różne. Stąd $\tau$ wyznacza permutację zbioru
$\{x_1,\ldots,x_n\}$; ponieważ ten zbiór generuje $L$ nad $K$,
permutacja całkowicie wyznacza $\tau$. Mamy więc co najwyżej $n!$
automorfizmów, a z Kroku 2 otrzymujemy
$\operatorname{Gal}(L/K)\cong S_n$.

\emph{Krok 4: nierozkładalność.}
Gdyby $f=gh$ z $g,h\in K[X]$ o dodatnich stopniach,
zbiory ich pierwiastków w $L$ byłyby niepuste i właściwe
(ponieważ $f$ jest rozdzielczy).
Dla automorfizmu $\tau\in\operatorname{Gal}(L/K)$ oraz pierwiastka
$a$ wielomianu $g$ mielibyśmy $g(\tau(a))=\tau(g(a))=0$.
Zbiór pierwiastków $g$ byłby więc niezmienniczy względem
całej grupy $S_n$. Jest to niemożliwe, gdyż $S_n$ działa
przechodnio na $\{x_1,\ldots,x_n\}$, a jedyne niezmiennicze
podzbiory zbioru będącego jedną orbitą to zbiór pusty i cały zbiór.
Zatem $f$ jest nierozkładalny.
\end{proof}

\begin{wniosek}[Ogólna niemożliwość rozwiązania przez pierwiastniki]
\label{aa:wn:generic-abel}
Jeżeli $\operatorname{char}k=0$ i $n\ge5$, to ogólny wielomian
stopnia $n$ nie jest rozwiązywalny przez pierwiastniki nad
$k(t_1,\ldots,t_n)$.
\end{wniosek}
\begin{proof}
Z Twierdzenia~\ref{aa:tw:generic-sn} jego grupa Galois to $S_n$.
Dla $n\ge5$ grupa $S_n$ zawiera podgrupę izomorficzną
z $A_5$: parzyste permutacje indeksów $1,\ldots,5$
przedłużamy, pozostawiając indeksy $6,\ldots,n$ nieruchome.
Wniosek~\ref{aa:wn:a5-not-solvable}, udowodniony na podstawie
prostoty i nieprzemienności $A_5$, stwierdza,
że $A_5$ nie jest rozwiązalna.
Każda podgrupa grupy rozwiązalnej jest rozwiązalna:
jej szereg pochodny zawiera się wyraz po wyrazie
w szeregu pochodnym całej grupy.
Zatem $S_n$ nie może być rozwiązalna.
Kryterium z Rozdziału~\ref{aa:ch:radicals} (w charakterystyce zero)
mówi, że rozwiązywalność wielomianu rozdzielczego przez pierwiastniki
pociąga rozwiązalność jego grupy Galois.
Stosując przeciwstawienie, otrzymujemy tezę.
\end{proof}

\begin{wniosek}[Ciało niezmienników i stopień rozszerzenia]
\label{aa:wn:generic-invariants}
Niech $k$ będzie dowolnym ciałem, $n\ge2$ i niech
$x_1,\ldots,x_n$ będą algebraicznie niezależnymi elementami
nad $k$. Oznaczmy $L=k(x_1,\ldots,x_n)$ i
$K=k(e_1,\ldots,e_n)$, gdzie $e_j$ jest $j$-tą elementarną
funkcją symetryczną zmiennych $x_1,\ldots,x_n$.
Wówczas
\[
 k(x_1,\ldots,x_n)^{S_n}=k(e_1,\ldots,e_n),
 \qquad [L:K]=n!,
\]
gdzie $S_n$ działa przez permutacje zmiennych. Innymi słowy,
każda funkcja wymierna $x_1,\ldots,x_n$ niezmiennicza względem
wszystkich ich permutacji jest funkcją wymierną elementarnych
funkcji symetrycznych.
\end{wniosek}
\begin{proof}
W dowodzie Twierdzenia~\ref{aa:tw:generic-sn} wykazano, że
$L/K$ jest skończonym rozszerzeniem Galois i że jego grupa
Galois jest dokładnie grupą wszystkich permutacji zmiennych,
czyli $S_n$. Z podstawowego twierdzenia teorii Galois
(Twierdzenie~\ref{aa:tw:fundamental-galois}) wynika
$L^{\operatorname{Gal}(L/K)}=K$. Ponieważ działanie
grupy $S_n$ pokrywa się z działaniem całej grupy Galois,
pierwsza równość jest udowodniona.
Rozszerzenie $L/K$ jest skończone i galoisowskie,
więc stopień rozszerzenia równa się rzędowi grupy Galois:
$[L:K]=|S_n|=n!$.
\end{proof}

\begin{przyklad}[Dwie funkcje wymierne niezmiennicze względem permutacji]
\label{aa:prz:generic-rational-invariants}
Niech $x_1,x_2,x_3$ będą algebraicznie niezależne nad ciałem
$k$, a $e_1,e_2,e_3$ odpowiadającymi im elementarnymi
funkcjami symetrycznymi. Rozwijając kwadrat sumy, otrzymujemy
\[
 x_1^2+x_2^2+x_3^2=(x_1+x_2+x_3)^2
 -2(x_1x_2+x_1x_3+x_2x_3)=e_1^2-2e_2.
\]
Tożsamość obowiązuje również w charakterystyce $2$,
gdzie jej ostatni składnik znika. Ponieważ algebraiczna
niezależność gwarantuje $x_1x_2x_3=e_3\ne0$, można też
przeprowadzić rachunek
\[
 \frac1{x_1}+\frac1{x_2}+\frac1{x_3}
 =\frac{x_2x_3+x_1x_3+x_1x_2}{x_1x_2x_3}
 =\frac{e_2}{e_3}.
\]
Obie funkcje pozostają niezmienione po dowolnej permutacji
trzech zmiennych; uzyskane wzory jawnie przedstawiają je
jako funkcje wymierne $e_1,e_2,e_3$.
\end{przyklad}

\begin{przyklad}[Kwadratowe ciało pośrednie dla wielomianu ogólnego stopnia trzeciego]
\label{aa:prz:generic-discriminant}
Niech $\operatorname{char}k\ne2$ i $n=3$.
W ciele $L=k(x_1,x_2,x_3)$ połóżmy
$\Delta=(x_1-x_2)(x_1-x_3)(x_2-x_3)$.
Dla każdej permutacji $\sigma\in S_3$ zachodzi
$\sigma(\Delta)=\operatorname{sgn}(\sigma)\Delta$.
W szczególności $\Delta$ jest ustalany przez wszystkie
elementy $A_3$, więc należy do $L^{A_3}$.
Jego kwadrat $D=\Delta^2$ jest niezmienniczy względem
całej grupy $S_3$ i dlatego należy do
$K=k(e_1,e_2,e_3)$. Rozwijając iloczyn różnic,
otrzymujemy jawnie
\[
 D=e_1^2e_2^2-4e_2^3-4e_1^3e_3-27e_3^2+18e_1e_2e_3.
\]
Jednak $\Delta\notin K$:
transpozycja zamienia go na $-\Delta\ne\Delta$.
Ponieważ $\operatorname{char}k\ne2$, jedynymi pierwiastkami
$T^2-D$ w $L$ są $\Delta$ i $-\Delta$.
Żaden z nich nie należy do $K$, bo transpozycja zmienia
znak niezerowego $\Delta$. Wielomian $T^2-D\in K[T]$
nie ma więc pierwiastka w $K$ i, jako wielomian stopnia dwa,
jest nierozkładalny.
Otrzymujemy $[K(\Delta):K]=2$. Podstawowe twierdzenie
teorii Galois daje $[L^{A_3}:K]=[S_3:A_3]=2$, więc
$L^{A_3}=K(\Delta)$. Jest to przykład nietrywialnego
ciała pośredniego, wyznaczonego przez parzyste permutacje.
\end{przyklad}

\section{Przykłady i granice twierdzenia}
\begin{przyklad}[Dwa nierozkładalne wielomiany sześcienne]
\label{aa:prz:generic-cubic}
Dla algebraicznie niezależnych $a,b,c$ nad $\mathbb Q$
wielomian $X^3+aX^2+bX+c$ ma grupę $S_3$ nad $\mathbb Q(a,b,c)$:
w Definicji~\ref{aa:def:generic-poly} wystarczy zastąpić
$t_1=-a$, $t_2=b$, $t_3=-c$.
Po podstawieniu konkretnych współczynników sytuacja się zmienia.
Wielomian $g(X)=X^3-2$ jest nierozkładalny nad $\mathbb Q$
(kryterium Eisensteina dla $2$), a jego wyróżnik wynosi
$-4\cdot0^3-27\cdot(-2)^2=-108$.
Ponieważ $-108$ nie jest kwadratem w $\mathbb Q$,
grupa Galois działająca przechodnio na trzech pierwiastkach
nie zawiera się w $A_3$, więc jest równa $S_3$.
Natomiast $h(X)=X^3-3X+1$ nie ma pierwiastków wymiernych:
jedynymi kandydatami według twierdzenia o pierwiastkach wymiernych
są $\pm1$, a $h(1)=-1$ i $h(-1)=3$.
Stąd $h$ jest nierozkładalny, a jego wyróżnik jest równy
$-4(-3)^3-27\cdot1^2=81=9^2$.
W charakterystyce zero kryterium wyróżnika daje więc podgrupę
$A_3$; z przechodniości wynika, że jest ona całym
$A_3\cong C_3$. Oba wielomiany mają stopień trzy, lecz różne
grupy Galois.
\end{przyklad}

\begin{przyklad}[Specjalizacja może zmniejszyć grupę]
\label{aa:prz:generic-specialization}
Wielomian ogólny stopnia czwartego ma grupę $S_4$.
Dla specjalizacji do $X^4-2$ otrzymujemy natomiast grupę
izomorficzną z grupą dihedralną $D_4$ rzędu $8$:
po wprowadzeniu $\alpha=\sqrt[4]{2}$ ciałem rozkładu jest
$\mathbb Q(\alpha,i)$, a jej osiem automorfizmów opisano
w Przykładzie~\ref{aa:prz:quartic-v4}.
Nie jest to sprzeczność: specjalizacja narzuca zależności między
współczynnikami, więc przestają one być algebraicznie niezależne.
\end{przyklad}

\begin{przyklad}[Dlaczego nie pomijamy rozdzielczości]
\label{aa:prz:generic-inseparable-specialization}
Niech $k=\mathbb F_p(t)$, gdzie $p$ jest liczbą pierwszą,
a $t$ elementem transcendentalnym nad $\mathbb F_p$.
Wielomian $q(X)=X^p-t\in k[X]$ jest nierozkładalny.
Istotnie, rozpatrując go najpierw w $\mathbb F_p[t][X]$,
widzimy, że wszystkie współczynniki oprócz wiodącego są
podzielne przez pierwszy element $t$, a wyraz stały $-t$
nie jest podzielny przez $t^2$. Kryterium Eisensteina
w pierścieniu $\mathbb F_p[t]$ i lemat Gaussa dają
nierozkładalność w jego ciele ułamków $k=\mathbb F_p(t)$.
Jeżeli $\alpha$ jest jego pierwiastkiem w domknięciu
algebraicznym $k$, to z $\alpha^p=t$ wynika
$q(X)=(X-\alpha)^p$, ponieważ w charakterystyce $p$
zachodzi $(X-\alpha)^p=X^p-\alpha^p$.
Stąd $q'(X)=0$ i w ciele rozkładu występuje tylko jeden
różny pierwiastek $\alpha$, mimo że $q$ jest nierozkładalny.
Jego grupa automorfizmów nad $k$
jest trywialna: każdy $k$-automorfizm ciała rozkładu
$k(\alpha)$ musi wysłać $\alpha$ na pierwiastek $q$,
a ten pierwiastek jest jedyny, zatem automorfizm ustala
$\alpha$ i całe $k(\alpha)$. Nie otrzymujemy więc $S_p$.
To przykład \emph{specjalizacji}, a nie kontrprzykład do
Twierdzenia~\ref{aa:tw:generic-sn}: w modelu użytym w jego dowodzie
pierwiastki są różnymi zmiennymi algebraicznie niezależnymi nad $k$,
natomiast przy $n=p$ podstawienie $t_1=\cdots=t_{p-1}=0$,
$t_p=t$ daje właśnie $X^p-t$, który nie jest rozdzielczy.
\end{przyklad}

\begin{uwaga}\label{aa:uw:generic-characteristic}
Twierdzenie~\ref{aa:tw:generic-sn} działa nad ciałem dowolnej
charakterystyki, ponieważ konstrukcja permutacji zmiennych nie
wymaga dzielenia przez $n!$. Natomiast Wniosek
\ref{aa:wn:generic-abel} opiera się na kryterium rozwiązywalności
przez pierwiastniki przyjętym w Rozdziale~\ref{aa:ch:radicals},
stąd jawne założenie charakterystyki zero.
\end{uwaga}

% AA-37 | Audyt I (2026-11-22); etykiety historyczne zachowane
\chapter{Algebraiczna niezależność i bazy transcendencji}
\label{aa:ch:transcendence}
\section{Motywacja: ile niezależnych parametrów ma rozszerzenie?}
Jeżeli $u,v$ są algebraicznie niezależnymi zmiennymi nad ciałem $K$,
to ciało $K(u,v)$ jest algebraiczne nad $K(u+v,uv)$:
oba elementy $u,v$ są pierwiastkami wielomianu
$T^2-(u+v)T+uv$. Ciała te nie muszą być równe.
Chcemy zatem liczyć nie liczbę \emph{wybranych} generatorów ciała,
lecz liczbę parametrów, których nie wiąże żadne niezerowe
równanie wielomianowe o współczynnikach z $K$.

\begin{definicja}\label{aa:def:alg-independent}
Niech $K\subseteq L$ będzie rozszerzeniem ciał, a $S\subseteq L$.
Zbiór $S$ nazywamy \emph{algebraicznie niezależnym nad $K$},
jeżeli dla każdego $m\ge1$, każdych parami różnych
$s_1,\ldots,s_m\in S$ i każdego niezerowego wielomianu
$P\in K[X_1,\ldots,X_m]$ zachodzi
$P(s_1,\ldots,s_m)\ne0$.
Zbiór pusty jest z definicji algebraicznie niezależny.
Zbiór $B\subseteq L$ jest \emph{bazą transcendencji rozszerzenia $L/K$},
jeżeli jest algebraicznie niezależny nad $K$, a każdy element $L$
jest algebraiczny nad ciałem $K(B)$ generowanym przez $K$ i $B$.
\end{definicja}

\begin{uwaga}
Niezależność zbioru jest własnością wszystkich jego
\emph{skończonych} podzbiorów. Z tego powodu definicja ma sens
także przy nieskończenie wielu parametrach; nie potrzebujemy
pierścienia wielomianów z nieskończonymi jednomianami.
Jeśli $B=\varnothing$ jest bazą transcendencji, oznacza to
po prostu, że $L/K$ jest algebraiczne.
\end{uwaga}

\begin{lemat}[Baza jako maksymalny zbiór niezależny]
\label{aa:lem:transcendence-maximal}
Niech $K\subseteq L$ będzie rozszerzeniem ciał i $B\subseteq L$.
Następujące warunki są równoważne:
\begin{enumerate}[label=(\roman*)]
 \item $B$ jest bazą transcendencji $L/K$;
 \item $B$ jest algebraicznie niezależny nad $K$ i nie istnieje
       algebraicznie niezależny $C$ taki, że $B\subsetneq C\subseteq L$.
\end{enumerate}
\end{lemat}
\begin{proof}
Załóżmy najpierw (i). Niech $x\in L\setminus B$.
Ponieważ $L/K(B)$ jest algebraiczne, istnieje moniczny
wielomian $m(T)=T^d+c_1T^{d-1}+\cdots+c_d\in K(B)[T]$
zerujący się w $x$. Każdy współczynnik $c_i$ jest ilorazem
dwóch wielomianów w skończenie wielu elementach $B$.
Po wybraniu wspólnego mianownika $q\ne0$ otrzymujemy
równanie $q x^d+q c_1x^{d-1}+\cdots+q c_d=0$,
którego współczynniki leżą w $K[b_1,\ldots,b_r]$
dla pewnych różnych $b_1,\ldots,b_r\in B$.
Wyrażenie $q$ jest niezerowe w $K(B)$, gdyż mianowniki
reprezentują niezerowe elementy ciała; niezależność $B$
gwarantuje ponadto, że jego zapis jako wielomian w $b_i$
nie jest zerowy. Otrzymujemy więc \emph{niezerowy} wielomian
w $r+1$ zmiennych zerujący się na $(b_1,\ldots,b_r,x)$.
To pokazuje zależność $B\cup\{x\}$. Dowolny właściwy
nadzbiór $B$ zawiera taki element $x$, a więc też jest zależny.

Załóżmy następnie (ii). Dla $x\in L\setminus B$ zbiór
$B\cup\{x\}$ jest zależny. Istnieją więc różne
$b_1,\ldots,b_r\in B$ i niezerowy
$P\in K[X_1,\ldots,X_r,T]$ z
$P(b_1,\ldots,b_r,x)=0$.
Wielomian $P$ musi mieć dodatni stopień względem $T$:
w przeciwnym razie naruszałby niezależność $B$.
Po podstawieniu $b_i$ dostajemy niezerowy wielomian
$P(b_1,\ldots,b_r,T)\in K(B)[T]$;
jego niezerowość wynika z faktu, że co najmniej jeden
niezerowy współczynnik przy potędze $T$ pozostaje niezerowy
po podstawieniu niezależnych $b_i$.
W konsekwencji $x$ jest algebraiczny nad $K(B)$.
Elementy $x\in B$ oczywiście należą do $K(B)$.
Wobec tego każde $x\in L$ jest algebraiczne nad $K(B)$,
co daje (i).
\end{proof}

\begin{twierdzenie}[Istnienie bazy transcendencji]
\label{aa:tw:transcendence-existence}
Każde rozszerzenie ciał $K\subseteq L$ posiada bazę transcendencji.
\end{twierdzenie}
\begin{proof}
Niech $\mathcal I$ będzie rodziną algebraicznie niezależnych
podzbiorów $L$, uporządkowaną przez inkluzję.
Rodzina nie jest pusta, gdyż $\varnothing\in\mathcal I$.
Jeżeli $\mathcal C\subseteq\mathcal I$ jest łańcuchem,
to $U=\bigcup_{C\in\mathcal C}C$ jest jego ograniczeniem górnym.
Aby sprawdzić niezależność $U$, weźmy dowolne
$s_1,\ldots,s_m\in U$. Dla każdego $s_i$ istnieje
$C_i\in\mathcal C$ zawierający $s_i$. Ponieważ $\mathcal C$
jest łańcuchem i lista $C_1,\ldots,C_m$ jest skończona,
jedno z $C_i$ zawiera wszystkie pozostałe. Zatem
$s_1,\ldots,s_m$ leżą razem w pewnym zbiorze niezależnym,
więc nie spełniają niezerowej relacji wielomianowej.
Lemat Zorna, wymagający w ogólności aksjomatu wyboru,
daje maksymalny $B\in\mathcal I$.
Na mocy Lematu~\ref{aa:lem:transcendence-maximal}
jest to baza transcendencji.
\end{proof}

\section{Wymiana elementów}
\begin{lemat}[Wymiana algebraiczna]
\label{aa:lem:transcendence-exchange}
Niech $K\subseteq L$ będzie rozszerzeniem ciał,
$B\subseteq L$ zbiorem algebraicznie niezależnym nad $K$,
a $b\in B$. Jeśli $x\in L$ jest algebraiczny nad $K(B)$,
lecz transcendentalny nad $F=K(B\setminus\{b\})$,
to $b$ jest algebraiczny nad $F(x)$.
\end{lemat}
\begin{proof}
Element $b$ jest transcendentalny nad $F$, ponieważ $B$ jest
niezależny. Z algebraiczności $x$ nad $F(b)$ wynika istnienie
niezerowego wielomianu $P(T)\in F(b)[T]$ zerującego się w $x$.
Usuwając mianowniki współczynników, otrzymujemy
niezerowy wielomian $Q(U,V)\in F[U,V]$
taki, że $Q(b,x)=0$; wielomian pozostaje niezerowy,
ponieważ podstawienie formalnego $U=b$ jest iniektywne
na $F[U]$.
Zapiszmy $Q(U,V)=\sum_{j=0}^r q_j(V)U^j$,
gdzie $q_j\in F[V]$ i co najmniej jeden $q_j$ jest niezerowy.
Ponieważ $x$ jest transcendentalny nad $F$,
niezerowy $q_j$ nie może spełniać $q_j(x)=0$.
Zatem $Q(U,x)\in F(x)[U]$ jest niezerowy i zeruje się
dla $U=b$. To jest równanie algebraiczne elementu $b$
nad $F(x)$.
\end{proof}

\begin{lemat}[Nierozkładalność po dołączeniu parametru]
\label{aa:lem:transcendence-irreducible-pure}
Niech $F$ będzie ciałem, $z$ elementem transcendentalnym nad $F$,
a $h(T)\in F[T]$ wielomianem nierozkładalnym dodatniego stopnia.
Wówczas $h$ pozostaje nierozkładalny w $F(z)[T]$.
\end{lemat}
\begin{proof}
Przypuśćmy, że $h=uv$ w $F(z)[T]$,
gdzie $u,v$ mają dodatnie stopnie względem $T$.
Pierścień $F[z]$ jest pierścieniem wielomianów nad ciałem,
a więc dziedziną z jednoznacznością rozkładu.
Wielomian $h$ jest prymitywny jako element $F[z][T]$:
jego niezerowe współczynniki należą do $F$ i ich
największy wspólny dzielnik jest jednostką.
Z lematu Gaussa wynika istnienie rozkładu
$h=\widetilde u\widetilde v$ w $F[z][T]$,
gdzie oba czynniki mają dodatni stopień względem $T$.
Traktując tę równość jako tożsamość w $F[T][z]$, mamy
$\deg_z h=0$. Ponieważ $F[T]$ jest dziedziną,
$\deg_z(\widetilde u\widetilde v)=
\deg_z\widetilde u+\deg_z\widetilde v$.
Oba stopnie muszą zatem wynosić $0$, czyli
$\widetilde u,\widetilde v\in F[T]$.
Otrzymujemy rozkład $h$ nad $F$, wbrew założeniu.
\end{proof}

\begin{lemat}[Zmienna istotna w funkcji wymiernej]
\label{aa:lem:transcendence-essential-variable}
Niech $D$ będzie zbiorem algebraicznie niezależnym nad ciałem $K$,
$A\subseteq D$, a $r\in K(D)\setminus K(A)$.
Istnieje $d\in D\setminus A$ takie, że
$r\notin K(D\setminus\{d\})$.
\end{lemat}
\begin{proof}
Z definicji ciała generowanego przez $D$ element $r$ jest
ilorazem dwóch wielomianów zależnych od skończenie wielu
elementów $D$. Niech $E\subseteq D\setminus A$ będzie
skończonym zbiorem obejmującym wszystkie elementy spoza $A$
użyte w takim ilorazie. Wtedy $r\in K(A)(E)$.
Wykażemy najpierw własność przecięcia: dla każdego $d\in E$
\[
 K(A)(E)\cap K(D\setminus\{d\})
   =K(A)(E\setminus\{d\}).
\]
Inkluzja od prawej do lewej wynika z definicji.
Dla inkluzji odwrotnej weźmy $u$ należące do przecięcia.
Jako funkcję wymierną zmiennej $d$ możemy je zapisać
$u=P(d)/Q(d)$, gdzie
$P,Q\in K(A)(E\setminus\{d\})[T]$, $Q\ne0$.
Jednocześnie $u\in K(D\setminus\{d\})$.
Ponieważ $d$ jest transcendentalny nad tym ostatnim ciałem,
z równości $P(d)=uQ(d)$ wynika równość każdego
współczynnika tych dwóch wielomianów.
Wybierzmy niezerowy współczynnik $q_j$ wielomianu $Q$.
Wówczas odpowiadający mu współczynnik $p_j$ wielomianu $P$
spełnia $p_j=uq_j$, skąd
$u=p_j/q_j\in K(A)(E\setminus\{d\})$.

Gdyby dla każdego $d\in E$ zachodziło
$r\in K(D\setminus\{d\})$, to powyższa własność pozwalałaby
usuwać kolejno wszystkie elementy skończonego zbioru $E$
z przedstawienia $r$.
Ostatecznie otrzymalibyśmy $r\in K(A)$,
sprzecznie z założeniem. Istnieje zatem wymagane $d$.
\end{proof}

\begin{twierdzenie}[Jednoznaczność liczby elementów bazy]
\label{aa:tw:transcendence-degree}
Niech $K\subseteq L$ będzie dowolnym rozszerzeniem ciał.
Jeżeli $B,C\subseteq L$ są bazami transcendencji $L/K$,
to mają tę samą moc. Moc tę nazywamy
\emph{stopniem transcendencji} i oznaczamy $\operatorname{trdeg}_K L$.
\end{twierdzenie}
\begin{proof}
\emph{Krok 1: zamiana jednego elementu w bazie.}
Niech $D$ będzie bazą transcendencji $L/K$, a $A\subseteq D$
podrodziną, której elementy chcemy zachować. Weźmy
$x\in L$ transcendentalny nad $K(A)$.
Ponieważ $L/K(D)$ jest algebraiczne, element $x$ jest
algebraiczny nad $K(D)$. Jego moniczny wielomian minimalny
$m_x(T)\in K(D)[T]$ nie może mieć wszystkich współczynników
w $K(A)$, gdyż wtedy $x$ byłby algebraiczny nad $K(A)$.
Ponieważ co najmniej jeden współczynnik $m_x$ nie należy
do $K(A)$, Lemat~\ref{aa:lem:transcendence-essential-variable}
pozwala wybrać $d\in D\setminus A$ tak, aby ten
współczynnik nie należał do $F=K(D\setminus\{d\})$.

Gdyby $x$ był algebraiczny nad $F$, jego moniczny
wielomian minimalny $h(T)\in F[T]$ pozostałby nierozkładalny
po dołączeniu transcendentalnego nad $F$ elementu $d$,
na mocy Lematu~\ref{aa:lem:transcendence-irreducible-pure}.
Ponieważ $F(d)=K(D)$, mielibyśmy $m_x=h$, co przeczyłoby
wyborowi $d$. Stąd $x$ jest transcendentalny nad $F$.
Z Lematu~\ref{aa:lem:transcendence-exchange} wynika,
że $d$ jest algebraiczny nad $F(x)$.
Połóżmy $D'=(D\setminus\{d\})\cup\{x\}$.
Niezależność $D'$ wynika z transcendentalności $x$ nad $F$,
a fakt, że $L/K(D')$ jest algebraiczne, z dwóch zdań:
$L$ jest algebraiczne nad $K(D)=F(d)$ oraz $d$ jest algebraiczny
nad $F(x)=K(D')$. Przechodniość algebraiczności kończy dowód,
że $D'$ jest bazą zawierającą $A\cup\{x\}$.

\emph{Krok 2: bazy skończone.}
Załóżmy, że $C$ ma $m<\infty$ elementów.
Dla dowolnych $m+1$ różnych elementów $b_1,\ldots,b_{m+1}\in B$
próbujemy kolejno wprowadzić je do bazy $C$.
Po wprowadzeniu $b_1,\ldots,b_j$ aktualna baza zawiera
$A_j=\{b_1,\ldots,b_j\}$; kolejny element $b_{j+1}$ jest
transcendentalny nad $K(A_j)$, ponieważ jest częścią
algebraicznie niezależnego zbioru $B$.
Z Kroku 1 można więc wymienić element spoza $A_j$ na
$b_{j+1}$, nie zmieniając liczby elementów bazy.
Po $m$ krokach baza ma nadal $m$ elementów, wszystkie należące
do $B$. Element $b_{m+1}$ jest nad nimi transcendentalny,
a więc nie może należeć do algebraicznego rozszerzenia
$K(C')\subseteq L$ wyznaczonego przez tę bazę $C'$.
Sprzeczność dowodzi $|B|\le m$.
W szczególności $B$ jest skończona; zamieniając role $B,C$,
uzyskujemy $|C|\le|B|$, a więc równość.

\emph{Krok 3: bazy nieskończone.}
Jeżeli żadna baza nie jest skończona, to dla każdego $c\in C$
istnieje skończony zbiór $B_c\subseteq B$ taki, że
$c$ jest algebraiczny nad $K(B_c)$: algebraiczność $c$
nad $K(B)$ jest wyrażona równaniem o skończonej liczbie
współczynników, z których każdy zawiera tylko skończenie
wiele elementów $B$.
Niech $B_*=\bigcup_{c\in C}B_c$.
Wtedy każdy $c\in C$ jest algebraiczny nad $K(B_*)$.
Każdy element $L$ jest z kolei algebraiczny nad $K(C)$,
więc, na mocy przechodniości algebraiczności,
$L$ jest algebraiczne nad $K(B_*)$.
Zbiór $B_*$ jest niezależny, bo $B_*\subseteq B$,
a zatem jest bazą transcendencji.
Jeśli $B_*\subsetneq B$, wybór $b\in B\setminus B_*$ dawałby
element transcendentalny nad $K(B_*)$, sprzecznie
z algebraicznością $L/K(B_*)$. Stąd $B=B_*$ i
\[
 |B|\le\sum_{c\in C}|B_c|\le |C|\cdot\aleph_0=|C|,
\]
gdzie ostatnia równość wynika z arytmetyki nieskończonych
liczb kardynalnych przy aksjomacie wyboru.
Ściślej, wybieramy dla każdego $c\in C$ skończony zbiór
$B_c$; ponieważ $C$ jest nieskończony, zbiór par
$(c,b)$ z $c\in C$ i $b\in B_c$ wstrzykuje się
w $C\times\mathbb N$ po wybraniu numeracji każdego
$B_c$. Stąd $|B|\le |C\times\mathbb N|=|C|$.
Po zamianie ról $B,C$ otrzymujemy odwrotną nierówność,
więc $|B|=|C|$.
\end{proof}

\begin{twierdzenie}[Dodawanie stopni transcendencji w wieży ciał]
\label{aa:tw:transcendence-tower}
Niech $K\subseteq E\subseteq L$ będzie wieżą rozszerzeń ciał.
Jeżeli $U\subseteq E$ jest bazą transcendencji $E/K$,
a $V\subseteq L$ jest bazą transcendencji $L/E$,
to $U\cup V$ jest bazą transcendencji $L/K$. W szczególności
\[
 \operatorname{trdeg}_K L
 =\operatorname{trdeg}_K E+\operatorname{trdeg}_E L,
\]
gdzie dodawanie mocy zbiorów rozumiemy także dla baz
nieskończonych.
\end{twierdzenie}
\begin{proof}
\emph{Krok 1: rozłączność i niezależność.}
Element $v\in V$ nie może należeć do $E$: wielomian
$T-v\in E[T]$ przeczyłby niezależności $V$ nad $E$.
Zatem $U\cap V=\varnothing$.
Przypuśćmy, że różne $u_1,\ldots,u_r\in U$ i
$v_1,\ldots,v_s\in V$ spełniają niezerową relację
$P(u_1,\ldots,u_r,v_1,\ldots,v_s)=0$, gdzie
$P\in K[X_1,\ldots,X_r,Y_1,\ldots,Y_s]$.
Zapiszmy $P=\sum_\alpha P_\alpha(X_1,\ldots,X_r)Y^\alpha$.
Ponieważ $P\ne0$, co najmniej jeden $P_\alpha$ jest
niezerowy. Niezależność $U$ nad $K$ gwarantuje
$P_\alpha(u_1,\ldots,u_r)\ne0$.
Po podstawieniu $u_i$ otrzymujemy więc niezerowy
wielomian o współczynnikach w $K(U)\subseteq E$,
który zeruje się na $(v_1,\ldots,v_s)$; sprzeczność
z niezależnością $V$ nad $E$.
Jeżeli $r=0$ lub $s=0$, argument sprowadza się
bezpośrednio do odpowiedniej niezależności $U$ lub $V$.

\emph{Krok 2: algebraiczność nad ciałem parametrów.}
Z definicji bazy transcendencji każdy $e\in E$ jest
algebraiczny nad $K(U)$, a więc również nad większym
ciałem $K(U\cup V)$. Oznacza to, że rozszerzenie
$E(V)/K(U\cup V)$ jest algebraiczne: jego elementy
powstają z elementów $E$ i $V$ przez skończenie wiele
działań ciała, zaś elementy algebraiczne nad danym
ciałem tworzą podciało.
Ponieważ $L/E(V)$ jest algebraiczne, z przechodniości
algebraiczności wynika, że $L/K(U\cup V)$ również
jest algebraiczne. Krok 1 i definicja bazy kończą dowód.
Równość stopni otrzymujemy z jednoznaczności mocy baz
(Twierdzenie~\ref{aa:tw:transcendence-degree}) i rozłączności
$U,V$.
\end{proof}

\begin{przyklad}[Wieża ciał i rozszerzenie algebraiczne]
\label{aa:prz:transcendence-tower-rational}
Niech $u,v$ będą algebraicznie niezależne nad $K$,
a $L=K(u,v)$ i $E=K(u+v)$. Podstawienie
$S\mapsto u+v$ jest iniekcją $K[S]\to K[u,v]$,
gdyż niezerowy wielomian w $u+v$ ma niezerowy
współczynnik przy najwyższej potędze $u$.
Stąd $\operatorname{trdeg}_K E=1$.
Ponieważ $u=(u+v)-v$, otrzymujemy $L=E(v)$.
Element $v$ jest transcendentalny nad $E$. Istotnie,
po oczyszczeniu mianowników ewentualnej relacji
nad $E=K(u+v)$ dostalibyśmy niezerowy wielomian
$P(S,T)\in K[S,T]$ z $P(u+v,v)=0$.
Jednak podstawienie $u\mapsto u+v$, $v\mapsto v$
jest automorfizmem $K[u,v]$ z odwrotnością
$u\mapsto u-v$, $v\mapsto v$; niezerowy $P$
nie może zatem zerować się na $(u+v,v)$.
Wobec tego $\operatorname{trdeg}_E L=1$, co daje $2=1+1$.

Dla porównania rozszerzenie $K(u,v)/K(u+v,uv)$
jest algebraiczne: obie zmienne spełniają równanie
$T^2-(u+v)T+uv=0$. Dodanie rozszerzenia algebraicznego
nie zmienia stopnia transcendencji, co jest przypadkiem
twierdzenia dla $\operatorname{trdeg}_E L=0$.
\end{przyklad}

\begin{przyklad}[Nieskończona baza transcendencji]
\label{aa:prz:transcendence-infinite}
Niech $(t_j)_{j\ge1}$ będzie ciągiem elementów algebraicznie
niezależnych nad ciałem $K$ i niech
$L=K(t_1,t_2,\ldots)$. Zbiór $B=\{t_j:j\ge1\}$
jest bazą transcendencji $L/K$: jest niezależny z założenia,
a każdy element $L$ należy już do $K(B)$.
Ponadto baza ta jest przeliczalnie nieskończona,
więc $\operatorname{trdeg}_K L=\aleph_0$.

Rozważmy również
$C=\{t_1+t_2,t_2,t_3,t_4,\ldots\}$.
Zachodzi $K(C)=L$, bo $t_1=(t_1+t_2)-t_2$.
Aby wykazać niezależność $C$, załóżmy, że pewien niezerowy
wielomian $P\in K[X_1,\ldots,X_m]$ zeruje się po
podstawieniu skończenie wielu różnych elementów $C$.
Wówczas otrzymujemy niezerową relację wielomianową między
skończenie wieloma $t_j$, gdyż podstawienie
$t_1\mapsto t_1+t_2$ i $t_j\mapsto t_j$ dla $j\ge2$
jest automorfizmem pierścienia
$K[t_1,t_2,\ldots]$, o odwrotności
$t_1\mapsto t_1-t_2$ oraz $t_j\mapsto t_j$ dla $j\ge2$.
To przeczy niezależności $B$. Zatem $C$ również jest
bazą transcendencji $L/K$.
Przykład pokazuje też, że nawet dla bazy nieskończonej
sprawdzanie niezależności sprowadza się do relacji
obejmujących tylko skończenie wiele elementów.
\end{przyklad}

\section{Przykłady i kontrprzykłady}
\begin{przyklad}[Dwie różne bazy w tym samym ciele]
\label{aa:prz:transcendence-rational}
Niech $u,v$ będą niezależnymi zmiennymi nad ciałem $K$.
Z definicji $\{u,v\}$ jest bazą transcendencji $K(u,v)/K$.
Połóżmy $s=u+v$ i $p=uv$.
Na mocy jednoznaczności zapisu wielomianu symetrycznego
(Twierdzenie~\ref{aa:tw:fundamental-symmetric}) $s,p$
są algebraicznie niezależne nad $K$.
Ponadto zarówno $u$, jak i $v$ spełniają
$T^2-sT+p=0$ nad $K(s,p)$.
Ciało $K(u,v)$ jest więc algebraiczne nad $K(s,p)$;
zatem $\{s,p\}$ jest drugą bazą transcendencji i
$\operatorname{trdeg}_K K(u,v)=2$.
W charakterystyce $2$ rachunek działa tak samo: współczynnik
$-s$ jest równy $s$, lecz tożsamość iloczynu
$(T-u)(T-v)=T^2-sT+p$ pozostaje prawdziwa.
\end{przyklad}

\begin{przyklad}[Dwie bazy jednoelementowe]
Niech $t$ będzie transcendentalny nad $K$ i $m\ge2$.
Zbiór $\{t^m\}$ jest niezależny, ponieważ
$P(t^m)=0$ dla niezerowego $P\in K[T]$ dawałoby
niezerową relację wielomianową dla $t$.
Element $t$ jest algebraiczny nad $K(t^m)$:
spełnia moniczne równanie $T^m-t^m=0$.
Dlatego $\{t\}$ i $\{t^m\}$ są różnymi bazami
transcendencji tego samego rozszerzenia $K(t)/K$.
\end{przyklad}


\begin{przyklad}[Baza transcendencji a nierozdzielczość]
\label{aa:prz:transcendence-inseparable}
Niech $K$ będzie ciałem charakterystyki $p>0$,
a $t$ elementem transcendentalnym nad $K$.
Ciało $K(t)$ jest algebraiczne nad $K(t^p)$, gdyż
$t$ spełnia $T^p-t^p=0$. Zatem zbiory $\{t\}$
i $\{t^p\}$ są bazami transcendencji $K(t)/K$.
Ponadto $[K(t):K(t^p)]=p$. Istotnie, gdyby
$1,t,\ldots,t^{p-1}$ były liniowo zależne nad $K(t^p)$,
to po oczyszczeniu mianowników mielibyśmy
$\sum_{i=0}^{p-1}A_i(t^p)t^i=0$ z pewnymi
$A_i\in K[T]$, nie wszystkimi zerowymi.
W rozwinięciu tej sumy potęgi $t$ pochodzące
od różnych indeksów $i$ mają różne reszty modulo $p$;
nie mogą się więc skracać. Otrzymujemy sprzeczność
z transcendentalnością $t$. Natomiast równanie
$T^p-t^p=0$ dowodzi nierówności odwrotnej dla stopnia.
Jest to rozszerzenie czysto nierozdzielcze: wielomian
minimalny $T^p-t^p$ ma pochodną formalną równą zeru.
Zależność algebraiczna potrzebna
w definicji bazy nie wymaga rozdzielczości;
w szczególności twierdzenie o jednoznaczności stopnia
transcendencji działa również w dodatniej charakterystyce.
\end{przyklad}

\begin{przyklad}[Kontrprzykład: dwa generatory nie muszą być niezależne]
\label{aa:prz:transcendence-dependent}
W ciele $K(t)$ zbiór $\{t,t^2\}$ generuje to samo
rozszerzenie co $\{t\}$, ale nie jest niezależny:
niezerowy wielomian $X_2-X_1^2\in K[X_1,X_2]$
zeruje się na parze $(t,t^2)$.
Liczba generatorów ciała nie musi być zatem równa
stopniowi transcendencji.
\end{przyklad}


\begin{przyklad}[Jawna wymiana parametru w bazie]
\label{aa:prz:transcendence-exchange-explicit}
Niech $u,v$ będą algebraicznie niezależne nad ciałem $K$,
$L=K(u,v)$ oraz $x=u+v^2$. Zbiory $\{u,v\}$ i
$\{x,v\}$ są bazami transcendencji $L/K$.
Rzeczywiście, dla niezerowego $P(S,T)\in K[S,T]$
założenie $P(x,v)=0$ oznaczałoby
$P(u+v^2,v)=0$. Podstawienie $S\mapsto U+V^2$,
$T\mapsto V$ jest automorfizmem $K[S,T]$ z odwrotnością
$U\mapsto S-T^2$, $V\mapsto T$; dlatego
$P(U+V^2,V)$ pozostaje niezerowym wielomianem.
Algebraiczna niezależność $u,v$ wyklucza jego zerowanie.
Ponadto $u=x-v^2\in K(x,v)$, więc $K(x,v)=L$.
To pokazuje bezpośrednio, że wymiana elementu bazy
może zachować nawet równość generowanych ciał,
a nie tylko algebraiczność ich rozszerzenia.
\end{przyklad}

% AA-38 | Ponowny audyt nocny LVI (2026-11-21); etykiety historyczne zachowane
\chapter{Lemat normalizacyjny Noether}\label{aa:ch:noether-normalization}
\section{Motywacja: algebra skończenie generowana a niezależne parametry}
Niech $A$ będzie skończenie generowaną przemienną algebrą nad ciałem
$K$. Generatory algebry nie muszą być algebraicznie niezależne:
w algebrze ilorazowej $K[X,Y]/(Y^2-X^3)$ klasy $x,y$
zmiennych $X,Y$ są powiązane równaniem $y^2=x^3$. Celem normalizacji Noether jest wykazanie,
że można zawsze wybrać pewną liczbę \emph{niezależnych} parametrów
$y_1,\ldots,y_r\in A$, nad którymi cała algebra jest skończonym
modułem. W dowodzie kluczowa jest zmiana generatorów zmieniająca
wybraną relację na równanie \emph{moniczne}.

Przypomnijmy, że element $a$ algebry $A$ jest \emph{całkowity}
nad podpierścieniem $B\subseteq A$, gdy spełnia równanie
$a^d+c_1a^{d-1}+\cdots+c_d=0$ dla pewnego $d\ge1$
i $c_i\in B$. Jeżeli $A=B[a]$, takie równanie daje
$A=B\cdot1+\cdots+B\cdot a^{d-1}$: wyższe potęgi $a$
redukujemy kolejno, zastępując $a^d$ wyrażeniem
$-c_1a^{d-1}-\cdots-c_d$.

\section{Wybór zmiennych i najwyższej potęgi}
\begin{lemat}[Moniczność po trójkątnej zmianie zmiennych]
\label{aa:lem:noether-monic}
Niech $K$ będzie dowolnym ciałem, $n\ge2$,
a $0\ne f\in K[X_1,\ldots,X_n]$ wielomianem niestałym.
Istnieją dodatnie liczby całkowite $m_1,\ldots,m_{n-1}$
i stała $c\in K\setminus\{0\}$ takie, że po podstawieniu
$X_i=Y_i+Y_n^{m_i}$ ($1\le i<n$), $X_n=Y_n$,
wielomian $c^{-1}f$ staje się moniczny względem $Y_n$,
ze współczynnikami w $K[Y_1,\ldots,Y_{n-1}]$.
Podstawienie jest automorfizmem pierścienia wielomianów.
\end{lemat}
\begin{proof}
\emph{Krok 1: wybór różnych wag.}
Niech $d=\deg f\ge1$ i wybierzmy całkowitą liczbę $M>d+1$.
Dla $1\le i<n$ połóżmy $m_i=M^i$. Każdemu
wielowskaźnikowi $\alpha=(\alpha_1,\ldots,\alpha_n)\in\mathbb N_0^n$
o $|\alpha|\le d$ przypiszmy wagę
$w(\alpha)=\alpha_n+\sum_{i=1}^{n-1}\alpha_i M^i$.
Ponieważ $0\le\alpha_i\le d<M$, liczba $w(\alpha)$
ma w zapisie przy podstawie $M$ cyfry
$\alpha_n,\alpha_1,\ldots,\alpha_{n-1}$.
Jednoznaczność takiego zapisu oznacza, że różne
wielowskaźniki stopnia co najwyżej $d$ mają różne wagi.

\emph{Krok 2: najwyższa potęga dla jednomianu.}
Po wykonaniu podstawienia jednomian
$X^\alpha=\prod_{i=1}^n X_i^{\alpha_i}$
przechodzi w
\[
 Y_n^{\alpha_n}\prod_{i=1}^{n-1}
       (Y_i+Y_n^{m_i})^{\alpha_i}.
\]
Wybranie ze wszystkich nawiasów potęg $Y_n^{m_i\alpha_i}$
daje składnik $Y_n^{w(\alpha)}$ ze współczynnikiem $1$.
Każdy inny składnik zawiera dla pewnego $i<n$
przynajmniej jedną potęgę $Y_i$, a więc ma stopień względem
$Y_n$ ściśle mniejszy od $w(\alpha)$.
Nie ma tu potrzeby dzielenia przez liczby całkowite:
argument działa również wtedy, gdy niektóre współczynniki
dwumianowe znikają w dodatniej charakterystyce.

\emph{Krok 3: brak skasowania najwyższego składnika.}
Zapiszmy $f=\sum_{\alpha}c_\alpha X^\alpha$,
gdzie zbiór indeksów z $c_\alpha\ne0$ jest skończony.
Wybierzmy $\beta$ o największej wadze $w(\beta)$.
Dzięki Krokowi 1 indeks $\beta$ jest jednoznaczny.
Krok 2 pokazuje, że $c_\beta Y_n^{w(\beta)}$ jest
jedynym składnikiem z tak wysoką potęgą $Y_n$ w całym
przekształconym $f$. Jego współczynnik $c_\beta\in K^\times$
nie zależy od $Y_1,\ldots,Y_{n-1}$.
Podzielenie przez $c_\beta$ czyni wielomian monicznym.
Ponieważ $f$ jest niestały, zawiera jednomian o niezerowym
wielowskaźniku, którego waga jest dodatnia. Zatem
$w(\beta)>0$ i uzyskane równanie ma dodatni stopień względem $Y_n$;
bez tego warunku nie wynikałaby całkowitość w zastosowaniu lematu.

\emph{Krok 4: odwracalność podstawienia.}
Przekształcenie $K$-algebry określone przez
$X_i\mapsto Y_i+Y_n^{m_i}$ dla $i<n$, $X_n\mapsto Y_n$
ma odwrotność $Y_i\mapsto X_i-X_n^{m_i}$ i $Y_n\mapsto X_n$.
Zatem jest izomorfizmem $K$-algebr $K[X_1,\ldots,X_n]$
i $K[Y_1,\ldots,Y_n]$, z odwracalną zmianą układu zmiennych.
Nie wprowadza ona nowych relacji ani nie traci istniejących.
\end{proof}

\begin{przyklad}[Dlaczego wybieramy różne wagi]
Niech $f(X_1,X_2)=X_1^2-X_2^3$ i niech $M>4$.
Podstawienie $X_1=Y_1+Y_2^M$, $X_2=Y_2$ daje
\[
 f=Y_1^2+2Y_1Y_2^M+Y_2^{2M}-Y_2^3.
\]
Najwyższa potęga $Y_2$ to $2M$ i jej współczynnik wynosi $1$.
Nawet w charakterystyce $2$, gdy wyraz $2Y_1Y_2^M$
znika, wielomian pozostaje moniczny względem $Y_2$.
Nie wystarcza natomiast dowolna zmiana współrzędnych:
na przykład podstawienie $X_1=Y_1+Y_2$, $X_2=Y_2$
w wielomianie $X_1-X_2$ prowadzi do $Y_1$,
który ma stopień zero względem $Y_2$ i nie daje
monicznego równania dodatniego stopnia dla $Y_2$.
\end{przyklad}

\begin{przyklad}[Jawna normalizacja nad ciałem dwuelementowym]
\label{aa:prz:noether-f2}
Niech $K=\mathbb F_2$ oraz
$A=K[X,Y]/(X^2+XY+Y)$; oznaczmy klasy zmiennych przez
$x,y$. Weźmy $u=x+y^4$. Podstawienie $X=U+V^4$,
$Y=V$ przekształca relację, w charakterystyce $2$, w
\[
 (U+V^4)^2+(U+V^4)V+V
   =V^8+V^5+UV+U^2+V.
\]
Otrzymany wielomian jest moniczny stopnia osiem względem
$V$. W szczególności $y$ spełnia nad $K[u]$ równanie
\[
 y^8+y^5+uy+u^2+y=0,
\]
a więc $A=K[u,y]$ jest generowany jako $K[u]$-moduł
przez $1,y,\ldots,y^7$. Co więcej, $u$ jest
transcendentalny nad $K$: gdyby niezerowy $P\in K[T]$
spełniał $P(u)=0$, to wielomian $P(U)$ należałby
do ideału $(V^8+V^5+UV+U^2+V)$ w $K[U,V]$.
Jest to niemożliwe, ponieważ wielomian tworzący ten ideał
jest moniczny dodatniego stopnia względem $V$,
a każdy jego niezerowy wielokrotny również ma dodatni
stopień względem $V$. Zatem $K[u]\subseteq A$ jest
normalizacją Noether, mimo że ciało bazowe jest skończone.
\end{przyklad}

\section{Twierdzenie o normalizacji}
\begin{twierdzenie}[Lemat normalizacyjny Noether]
\label{aa:tw:noether-normalization}
Niech $K$ będzie ciałem, a $A$ niezerową, przemienną,
skończenie generowaną $K$-algebrą z jedynką,
przy czym odwzorowanie strukturalne $K\to A$ zachowuje jedynkę.
Wówczas istnieją liczba całkowita $r\ge0$ i elementy
$y_1,\ldots,y_r\in A$ takie, że:
\begin{enumerate}[label=(\roman*)]
\item $y_1,\ldots,y_r$ są algebraicznie niezależne nad $K$;
\item $A$ jest skończenie generowanym modułem nad podalgebrą
      $K[y_1,\ldots,y_r]\subseteq A$.
\end{enumerate}
Nie wymagamy, aby $K$ było nieskończone ani aby $A$
była dziedziną całkowitości.
\end{twierdzenie}
\begin{proof}
Niech $n\ge0$ będzie liczbą wybranych generatorów
$A=K[a_1,\ldots,a_n]$. Dowodzimy tezy indukcyjnie względem $n$.

\emph{Krok 1: początek indukcji i przypadek niezależny.}
Jeżeli $n=0$, to $A$ jest obrazem homomorfizmu $K\to A$.
Jądro tego homomorfizmu jest ideałem ciała, a więc jest
albo $(0)$, albo całym $K$. Druga możliwość dawałaby $1_A=0_A$,
co jest sprzeczne z $A\ne0$.
Dlatego $A\cong K$; można przyjąć $r=0$.
Jeśli $n\ge1$ i generatory $a_1,\ldots,a_n$ są algebraicznie
niezależne, to ich podalgebra jest pierścieniem wielomianów,
a $A=K[a_1,\ldots,a_n]$ jest wolnym modułem rzędu jeden
nad tą podalgebrą. Przyjmujemy $r=n$.

\emph{Krok 2: jedna relacja algebraiczna.}
Jeżeli generatory są zależne, to istnieje niezerowy
$f\in K[X_1,\ldots,X_n]$ z $f(a_1,\ldots,a_n)=0$.
Wielomian $f$ nie może być niezerową stałą, ponieważ
$K\hookrightarrow A$ jest injekcją.
Dla $n=1$ dzielimy $f$ przez niezerowy współczynnik
najwyższej potęgi; otrzymujemy moniczne równanie
$a_1^d+c_1a_1^{d-1}+\cdots+c_d=0$ nad $K$.
Zatem $A=K[a_1]$ jest generowany jako $K$-moduł przez
$1,a_1,\ldots,a_1^{d-1}$, co daje tezę z $r=0$.

\emph{Krok 3: redukcja liczby generatorów przy $n\ge2$.}
Zastosujmy Lemat~\ref{aa:lem:noether-monic} do niezerowej
relacji $f$. Dla otrzymanych wykładników $m_i$ definiujemy
$b_i=a_i-a_n^{m_i}$ ($1\le i<n$) i
$B=K[b_1,\ldots,b_{n-1}]\subseteq A$.
Z równości $a_i=b_i+a_n^{m_i}$ wynika
\[
 A=K[a_1,\ldots,a_n]=B[a_n].
\]
Po podstawieniu zmiennych przekształcone równanie $f=0$
przyjmuje, po podzieleniu przez niezerowy skalar,
postać
\[
 a_n^d+c_1(b_1,\ldots,b_{n-1})a_n^{d-1}
 +\cdots+c_d(b_1,\ldots,b_{n-1})=0
\]
z pewnym $d\ge1$, gdzie $c_j\in K[Y_1,\ldots,Y_{n-1}]$.
Zatem $a_n$ jest całkowity nad $B$.
Ten krok nie wymaga, aby $A$ była dziedziną całkowitości:
moniczne równanie pozostaje ważne także w obecności
dzielników zera i elementów nilpotentnych.
Każdy wielomian w $a_n$ można redukować modulo tę relację
do stopnia mniejszego niż $d$.
Dlatego $1,a_n,\ldots,a_n^{d-1}$ generują $A=B[a_n]$
jako $B$-moduł.

\emph{Krok 4: zastosowanie indukcji.}
Algebra $B$ jest niezerową $K$-algebrą generowaną przez
$n-1$ elementów. Z założenia indukcyjnego istnieją
niezależne $y_1,\ldots,y_r\in B$ takie, że $B$ jest
skończonym modułem nad $C=K[y_1,\ldots,y_r]$.
Niech $v_1,\ldots,v_s$ generują $B$ jako $C$-moduł.
Z Kroku 3 wynika, że skończenie wiele elementów
$v_i a_n^j$, gdzie $1\le i\le s$ i $0\le j<d$,
generuje $A$ jako $C$-moduł:
każdy element $A$ jest sumą $B$-wielokrotności $a_n^j$,
a każdy współczynnik z $B$ jest sumą $C$-wielokrotności $v_i$.
Elementy $y_i$ pozostają niezależne w $A$, ponieważ
$B\subseteq A$ jest podpierścieniem i równania
wielomianowe między nimi nie zmieniają się po włączeniu do $A$.
To kończy indukcję i dowód twierdzenia.
\end{proof}

\section{Wymiar transcendencji i przykłady}
\begin{wniosek}\label{aa:wn:noether-dimension}
Jeżeli $A$ jest niezerową, skończenie generowaną dziedziną
całkowitości nad ciałem $K$, to liczba $r$ w normalizacji
Noether jest jednoznacznie wyznaczona i spełnia
\[
 r=\operatorname{trdeg}_K\operatorname{Frac}(A).
\]
\end{wniosek}
\begin{proof}
Niech $B=K[y_1,\ldots,y_r]$ będzie podalgebrą z twierdzenia.
Ponieważ $A$ jest skończonym modułem nad $B$, każdy element
$a\in A$ jest całkowity nad $B$.
Uzasadnijmy to bez dodatkowego założenia: jeżeli
$v_1,\ldots,v_s$ generują $A$ nad $B$, mnożenie przez $a$
wyraża się relacjami
$a v_i=\sum_{j=1}^s c_{ij}v_j$ z $c_{ij}\in B$.
Po zapisaniu $(aI-C)v=0$ i pomnożeniu przez macierz
dołączoną otrzymujemy
$\det(aI-C)v_i=0$ dla każdego $i$.
W szczególności $1=\sum b_i v_i$ dla pewnych $b_i\in B$,
więc $\det(aI-C)\cdot1=0$ w $A$.
Wielomian $\det(TI-C)\in B[T]$ jest moniczny i zeruje się
w $a$. Zatem $a$ jest algebraiczny nad $\operatorname{Frac}(B)$.
Również każdy element $\operatorname{Frac}(A)$ jest ilorazem
dwóch elementów $A$, z których mianownik jest niezerowy;
wszystkie takie ilorazy są algebraiczne nad
$\operatorname{Frac}(B)$, ponieważ elementy algebraiczne
w rozszerzeniu ciał tworzą podciało.
Stąd $\operatorname{Frac}(A)/K(y_1,\ldots,y_r)$ jest
algebraiczne. Niezależność $y_i$ oznacza, że
$\{y_1,\ldots,y_r\}$ jest bazą transcendencji
$\operatorname{Frac}(A)/K$.
Twierdzenie~\ref{aa:tw:transcendence-degree} pokazuje,
że jej liczba elementów nie zależy od wyboru normalizacji.
\end{proof}

\begin{przyklad}[Parabola półsześcienna]
\label{aa:prz:noether-parabola}
Niech $A=K[X,Y]/(Y^2-X^3)$ i niech $x,y$ oznaczają klasy
$X,Y$ w $A$. Relacja $y^2=x^3$ jest moniczna względem $y$.
Dzielenie z resztą w $K[X][Y]$ przez moniczny wielomian
$Y^2-X^3$ daje dla każdego elementu $A$ jednoznaczny zapis
$a(x)+b(x)y$, gdzie $a,b\in K[X]$.
Jednoznaczność wynika stąd, że niezerowy wielomian
$H(X,Y)(Y^2-X^3)$ ma stopień co najmniej $2$ względem $Y$.
Zatem $A=K[x]\oplus K[x]y$ jako $K[x]$-moduł.
Ponadto $x$ jest transcendentalny nad $K$:
gdyby $P(x)=0$ dla niezerowego $P\in K[X]$, to
$P(X)$ należałby do ideału $(Y^2-X^3)$, co jest niemożliwe
z powyższego porównania stopni względem $Y$.
Otrzymujemy normalizację z jednym niezależnym parametrem $x$.
\end{przyklad}

\begin{przyklad}[Hiperbola i nieoczywisty parametr]
Niech $A=K[U,V]/(UV-1)$, a $u,v$ będą klasami zmiennych.
Z relacji $uv=1$ wynika, że $v=u^{-1}$, zatem
$A\cong K[u,u^{-1}]$. Element $u$ jest transcendentalny
nad $K$, bo odwzorowanie $K[U]\to K[U,U^{-1}]$
jest injektywne.
Sam pierścień $K[u]$ nie daje normalizacji: $u^{-1}$
nie jest całkowity nad $K[u]$. Gdyby spełniał moniczne
równanie $(u^{-1})^d+c_1(u^{-1})^{d-1}+\cdots+c_d=0$
ze współczynnikami $c_j\in K[u]$, pomnożenie przez $u^d$
dałoby $1+c_1u+\cdots+c_du^d=0$ w $K[u,u^{-1}]$.
Ponieważ zanurzenie $K[u]\hookrightarrow K[u,u^{-1}]$
jest injektywne, byłaby to tożsamość w $K[u]$.
Podstawienie $u=0$ prowadziłoby do $1=0$ w $K$,
co jest niemożliwe.
Wybierzmy natomiast $s=u+u^{-1}\in A$.
Mamy $u^2-su+1=0$, więc $u$ jest całkowity nad $K[s]$.
Ponadto $u^{-1}=s-u$, stąd $A=K[s,u]$ i
$A=K[s]\cdot1+K[s]\cdot u$.
Element $s$ musi być transcendentalny nad $K$:
gdyby był algebraiczny, to z równania $u^2-su+1=0$
wynikałaby algebraiczność $u$ nad $K$, wbrew powyższemu.
Tak więc $K[s]\subseteq A$ jest normalizacją Noether,
choć $K[u]\subseteq A$ nią nie jest.
\end{przyklad}

\begin{przyklad}[Algebraiczność ciał ułamków nie daje skończoności modułowej]
\label{aa:prz:noether-nonfinite}
Niech $t$ będzie transcendentalny nad $K$.
Rozszerzenie pierścieni $K[t]\subset K(t)$ ma identyczne
ciała ułamków: oba dają $K(t)$, a więc rozszerzenie
ciał ułamków jest algebraiczne stopnia $1$.
Jednak $K(t)$ nie jest skończonym modułem nad $K[t]$.
W przeciwnym razie $1/t$ byłby całkowity nad $K[t]$
(z dowodu Wniosku~\ref{aa:wn:noether-dimension}).
Istniałoby więc równanie
$(1/t)^d+c_1(1/t)^{d-1}+\cdots+c_d=0$ z $c_i\in K[t]$.
Po pomnożeniu przez $t^d$ otrzymujemy
$1+c_1t+\cdots+c_dt^d=0$, co wymusza $1\in tK[t]$;
to niemożliwe, bo każdy wielomian z $tK[t]$ ma wyraz
stały równy zeru.
Kontrprzykład pokazuje, że algebraiczność rozszerzenia
\emph{ciał ułamków} nie wystarcza do skończoności
modułowej rozszerzenia \emph{pierścieni}.
Nie narusza on Twierdzenia~\ref{aa:tw:noether-normalization},
ponieważ $K(t)$ nie jest skończenie generowaną $K$-algebrą.
Aby to udowodnić, przypuśćmy, że
$K(t)=K[f_1(t)/g_1(t),\ldots,f_m(t)/g_m(t)]$,
gdzie $f_i,g_i\in K[t]$, $g_i\ne0$.
Połóżmy $D=g_1\cdots g_m\ne0$.
Wszystkie generatory należą wtedy do $K[t,D^{-1}]$,
więc $K(t)=K[t,D^{-1}]$. Jednak wielomian
$h=1+tD\in K[t]$ jest niestały i względnie pierwszy
z $D$: wspólny dzielnik $h$ i $D$ dzieliłby też
$h-tD=1$. Element $1/h\in K(t)$ nie należy do
$K[t,D^{-1}]$. W przeciwnym razie
$D^N=hP$ dla pewnych $N\ge0$ i $P\in K[t]$;
wówczas $h$ dzieliłby $D^N$, co przeczy względnej
pierwszości $h,D$ i temu, że $h$ nie jest jednostką.
Otrzymana sprzeczność dowodzi zapowiedzianego faktu.
\end{przyklad}



\begin{przyklad}[Normalizacja algebry z elementami nilpotentnymi]
\label{aa:prz:noether-nilpotents}
Niech $A=K[X,E]/(E^2)$, gdzie $K$ jest dowolnym ciałem,
a $x,\varepsilon$ oznaczają klasy zmiennych $X,E$.
Ponieważ $E^2$ jest wielomianem monicznym stopnia dwa
względem $E$, dzielenie w $K[X][E]$ daje jednoznaczny
zapis elementu $A$ jako
$a(x)+b(x)\varepsilon$ dla $a,b\in K[X]$.
W szczególności $A=K[x]\oplus K[x]\varepsilon$ jako
$K[x]$-moduł: ma on bazę $\{1,\varepsilon\}$.
Element $x$ jest transcendentalny nad $K$,
ponieważ niezerowy $P(X)\in K[X]$ nie może być
wielokrotnością $E^2$ w $K[X,E]$.
Zatem $K[x]\subseteq A$ jest normalizacją Noether.
Jednocześnie $\varepsilon^2=0$ i $\varepsilon\ne0$,
więc $A$ nie jest pierścieniem zredukowanym.
Przykład pokazuje, że nie tylko dzielniki zera,
lecz także niezerowe elementy nilpotentne są zgodne
z twierdzeniem o normalizacji.
\end{przyklad}

\begin{przyklad}[Algebra z dzielnikami zera]
\label{aa:prz:noether-zero-divisors}
Niech $A=K[X,Y]/(XY)$, a $x,y$ oznaczają klasy zmiennych.
Element $s=x+y$ jest transcendentalny nad $K$.
Istotnie, homomorfizm $A\to K[X]$ określony przez
$y\mapsto0$, $x\mapsto X$ przesyła $s$ na $X$;
gdyby niezerowy wielomian $P\in K[T]$ zerował się w $s$,
to $P(X)=0$ w $K[X]$, co jest sprzecznością.
Relacja $xy=0$ daje $x(s-x)=0$, czyli
$x^2-sx=0$. Zatem $x$ jest całkowity nad $K[s]$.
Ponieważ $y=s-x$, otrzymujemy $A=K[s,x]$.
Redukcja potęg przez $x^2=sx$ dowodzi
$A=K[s]\cdot1+K[s]\cdot x$, więc $A$ jest skończonym
modułem nad $K[s]$. Jednocześnie $xy=0$, a $x\ne0$
i $y\ne0$: obrazy odpowiednio w ilorazach przez $(y)$
i przez $(x)$ są niezerowymi zmiennymi wielomianowymi.
Przykład potwierdza, że lemat normalizacyjny Noether
nie wymaga braku dzielników zera.
\end{przyklad}

% AA-39 | Audyt I: samodzielny dowód wymiaru K[U,V], doprecyzowanie łańcuchów i odsyłaczy.
\chapter{Wymiar Krulla i łańcuchy ideałów pierwszych}
\label{aa:ch:krull-dimension}

\section{Dlaczego mierzymy długości łańcuchów?}
W przestrzeni afinicznej nad ciałem $K$ prosta ma jeden swobodny
parametr, a płaszczyzna dwa. Dla pierścieni zamiast współrzędnych
można badać ciągi coraz większych ideałów pierwszych. Pozwala to
zdefiniować wymiar także dla pierścieni, które nie są pierścieniami
wielomianów, a nawet mają dzielniki zera. W całym rozdziale pierścienie
są przemienne, mają jedynkę i są niezerowe, o ile wyraźnie nie
postanowimy inaczej. Wprowadzenie ideałów pierwszych i maksymalnych
znajduje się w Rozdziale~\ref{aa:ch:pierscienie}.

\begin{definicja}[Łańcuch ideałów i wymiar Krulla]\label{aa:def:krull-dimension}
Niech $A$ będzie niezerowym pierścieniem przemiennym z jedynką.
\emph{Łańcuchem ideałów pierwszych długości $r$} nazywamy ciąg
\[
 \mathfrak p_0\subsetneq\mathfrak p_1\subsetneq\cdots
 \subsetneq\mathfrak p_r
\]
ideałów pierwszych pierścienia $A$, gdzie $r\in\mathbb N_0$.
\emph{Wymiar Krulla} $\dim A$ jest kresem górnym zbioru długości
wszystkich takich łańcuchów; może mieć wartość $+\infty$.
Dla ideału pierwszego $\mathfrak p$ definiujemy jego \emph{wysokość}
$\operatorname{ht}(\mathfrak p)$ jako kres górny długości łańcuchów
ideałów pierwszych kończących się w $\mathfrak p$.
\end{definicja}

Długość oznacza liczbę ścisłych zawierań, a nie liczbę występujących
ideałów. W szczególności pojedynczy ideał pierwszy tworzy łańcuch
długości $0$. Istnienie ideału maksymalnego w każdym niezerowym
pierścieniu przemiennym z jedynką wynika z lematu Zorna (w ogólności
wykorzystujemy tu aksjomat wyboru), więc $\dim A\ge0$.
Nie wprowadzamy tu wymiaru pierścienia zerowego; wyłączenie go
zapobiega niejasnej konwencji dla pustego spektrum.

Jeżeli $\mathfrak p$ jest ideałem pierwszym, to nie wolno
utożsamiać $\operatorname{ht}(\mathfrak p)$ z liczbą
ideałów pierwszych zawartych w $\mathfrak p$:
wysokość liczy długości \emph{ściśle rosnących łańcuchów}.
Na przykład w $\mathbb Z$ mamy
$\operatorname{ht}((0))=0$ oraz
$\operatorname{ht}((p))=1$ dla każdej liczby pierwszej $p$.
Dla dowolnego pierścienia możliwa jest również wysokość nieskończona.



\begin{wniosek}[Charakterystyka wymiaru zero]
\label{aa:cor:krull-dim-zero}
Niech $A$ będzie niezerowym pierścieniem przemiennym z jedynką.
Wówczas $\dim A=0$ wtedy i tylko wtedy, gdy każdy ideał pierwszy
pierścienia $A$ jest maksymalny.
\end{wniosek}
\begin{proof}
Jeśli $\dim A=0$, nie istnieją dwa różne ideały pierwsze
$\mathfrak p\subsetneq\mathfrak q$. Każdy ideał pierwszy
$\mathfrak p$ zawiera się w pewnym ideale maksymalnym
$\mathfrak m$, który również jest pierwszy; stąd
$\mathfrak p=\mathfrak m$. Odwrotnie, jeśli wszystkie ideały
pierwsze są maksymalne, żadne dwa różne takie ideały nie mogą
pozostawać w relacji ścisłego zawierania. Ponieważ $A$ posiada
co najmniej jeden ideał pierwszy, największa możliwa długość
łańcucha wynosi $0$.
\end{proof}
\begin{przyklad}[Ciało i dziedzina liczb całkowitych]\label{aa:prz:krull-field-integers}
W ciele $K$ jedynym ideałem właściwym jest $(0)$, który jest
pierwszy. Nie można utworzyć łańcucha długości $1$, dlatego
$\dim K=0$.
W $\mathbb Z$ ideał $(0)$ jest pierwszy, ponieważ $\mathbb Z$
nie ma dzielników zera. Każdy niezerowy ideał pierwszy ma postać
$(p)$ dla liczby pierwszej $p$ i jest maksymalny, gdyż
$\mathbb Z/(p)$ jest ciałem. Stąd wszystkie nietrywialne łańcuchy
mają postać $(0)\subsetneq(p)$ i $\dim\mathbb Z=1$.
\end{przyklad}

\section{Ilorazy i lokalizacje}
\begin{twierdzenie}[Ideały pierwsze w ilorazie]\label{aa:tw:krull-quotient-primes}
Niech $A$ będzie pierścieniem przemiennym z jedynką, a $I\subsetneq A$
ideałem. Odwzorowanie $\mathfrak p\mapsto\mathfrak p/I$ jest
zachowującą inkluzję bijekcją między ideałami pierwszymi $A$
zawierającymi $I$ a ideałami pierwszymi $A/I$.
W szczególności $\dim(A/I)$ jest kresem długości łańcuchów
ideałów pierwszych $A$ zawierających $I$.
\end{twierdzenie}
\begin{proof}
Niech $\pi:A\to A/I$ będzie homomorfizmem ilorazowym.
Jeśli $\mathfrak p$ zawiera $I$ i jest pierwszy, to
$(A/I)/(\mathfrak p/I)\cong A/\mathfrak p$ jest dziedziną
całkowitości, a więc $\mathfrak p/I$ jest pierwszy.
Jeśli $\mathfrak q$ jest pierwszy w $A/I$, jego przeciwobraz
$\pi^{-1}(\mathfrak q)$ jest pierwszy w $A$: z
$ab\in\pi^{-1}(\mathfrak q)$ wynika
$\pi(a)\pi(b)\in\mathfrak q$, zatem $\pi(a)\in\mathfrak q$
lub $\pi(b)\in\mathfrak q$. Przeciwobraz zawiera $I$.
Sprawdźmy odwrotność obu konstrukcji. Dla ideału
$\mathfrak p\supseteq I$ równość
$\pi^{-1}(\mathfrak p/I)=\mathfrak p$ wynika z tego, że
$\pi(a)\in\mathfrak p/I$ dokładnie wtedy, gdy
$a\in\mathfrak p$. Dla ideału $\mathfrak q\subseteq A/I$
mamy $\pi(\pi^{-1}(\mathfrak q))=\mathfrak q$, ponieważ
$\pi$ jest surjekcją. Jeśli $\mathfrak p_1\subsetneq
\mathfrak p_2$ zawierają $I$, to ich obrazy pozostają różne:
wybierzmy $a\in\mathfrak p_2\setminus\mathfrak p_1$;
wtedy $a+I$ należy do $\mathfrak p_2/I$, lecz nie należy
do $\mathfrak p_1/I$. Analogicznie działamy w kierunku
przeciwnym przez przeciwobrazy. Tym samym bijekcja zachowuje
i odzwierciedla ścisłe inkluzje, więc przenosi łańcuchy
z zachowaniem długości.
\end{proof}




\begin{wniosek}[Przejście do ilorazu nie zwiększa wymiaru]
\label{aa:cor:krull-quotient-bound}
Niech $A$ będzie niezerowym pierścieniem przemiennym z jedynką,
a $I\subsetneq A$ ideałem. Wówczas $\dim(A/I)\leq\dim A$.
Ogólniej, jeśli $I\subseteq J\subsetneq A$, to
$\dim(A/J)\leq\dim(A/I)$.
\end{wniosek}
\begin{proof}
Każdy łańcuch ideałów pierwszych $A/I$ odpowiada, zgodnie z
Twierdzeniem~\ref{aa:tw:krull-quotient-primes}, łańcuchowi
ideałów pierwszych $A$ zawierających $I$ i mającemu dokładnie
tę samą długość. W szczególności jego długość nie przekracza
$\dim A$. Dla drugiej nierówności rozpatrujemy
$A/J\cong(A/I)/(J/I)$ i stosujemy pierwszą nierówność
do pierścienia $A/I$, który jest niezerowy, gdyż $I\subsetneq A$.
\end{proof}

\begin{wniosek}[Redukcja nie zmienia wymiaru]\label{aa:cor:krull-nilradical}
Niech $A$ będzie niezerowym pierścieniem przemiennym z jedynką,
a $N=\sqrt{(0)}=\{a\in A:\text{ istnieje }m\ge1\text{ takie, że }a^m=0\}$
jego nilradykałem. Wówczas $A/N$ jest niezerowym pierścieniem
i $\dim(A/N)=\dim A$. Ponadto ideały pierwsze tych dwóch
pierścieni odpowiadają sobie bijektywnie, z zachowaniem wszystkich
ścisłych zawierań.
\end{wniosek}
\begin{proof}
Najpierw $N\ne A$: gdyby $1\in N$, dla pewnego $m\ge1$
mielibyśmy $1^m=0$, a wtedy $A$ byłby pierścieniem zerowym.
Niech $\mathfrak p$ będzie dowolnym ideałem pierwszym $A$.
Dla $a\in N$ istnieje $m\ge1$ takie, że $a^m=0\in\mathfrak p$.
Pierwszość $\mathfrak p$, zastosowana kolejno do iloczynów
$a\cdot a^{m-1}$, $a\cdot a^{m-2}$ itd., daje $a\in\mathfrak p$.
Zatem $N\subseteq\mathfrak p$ dla \emph{każdego} ideału
pierwszego $A$. Twierdzenie~\ref{aa:tw:krull-quotient-primes}
ustanawia więc bijekcję między wszystkimi ideałami pierwszymi
$A$ i $A/N$, zachowującą i odzwierciedlającą ścisłe zawieranie.
Przenosi ona każdy łańcuch na łańcuch tej samej długości
w obie strony, skąd wynika równość wymiarów.
\end{proof}
\begin{twierdzenie}[Ideały pierwsze w lokalizacji]\label{aa:tw:krull-localization-primes}
Niech $A$ będzie pierścieniem przemiennym z jedynką,
a $S\subseteq A$ zbiorem multiplikatywnym zawierającym $1$
i niezawierającym $0$. Ideały pierwsze $S^{-1}A$ odpowiadają
bijektywnie ideałom pierwszym $\mathfrak p\subset A$ takim,
że $\mathfrak p\cap S=\varnothing$, przez
$\mathfrak p\mapsto S^{-1}\mathfrak p$ i przeciwobraz względem
$A\to S^{-1}A$. Obie operacje zachowują i odzwierciedlają inkluzję.
\end{twierdzenie}
\begin{proof}
Ułamki formalne i warunek ich równości są opisane
w Rozdziale~\ref{aa:ch:lokalizacja}.
Dla $\mathfrak p\cap S=\varnothing$ oznaczmy przez
$\overline S$ obraz $S$ w $A/\mathfrak p$.
Homomorfizm
\[
\Phi:S^{-1}A\longrightarrow
\overline S^{-1}(A/\mathfrak p),\qquad
\Phi(a/s)=(a+\mathfrak p)/(s+\mathfrak p)
\]
jest dobrze określony: jeśli $a/s=a'/s'$, to istnieje
$t\in S$ z $t(as'-a's)=0$ w $A$ i ten sam warunek
przechodzi do ilorazu. Jest on surjektywny, ponieważ
każdy ułamek po prawej ma postać obrazu $a/s$.
Ponadto $\Phi(a/s)=0$ dokładnie wtedy, gdy dla pewnego
$t\in S$ zachodzi $ta\in\mathfrak p$.
Pierwszość $\mathfrak p$ oraz $t\notin\mathfrak p$ dają
$a\in\mathfrak p$; w drugą stronę jest to bezpośrednie.
Jądro $\Phi$ jest więc równe $S^{-1}\mathfrak p$.
Ponieważ $A/\mathfrak p$ jest dziedziną całkowitości,
a każdy element $\overline S$ jest niezerowy,
lokalizacja $\overline S^{-1}(A/\mathfrak p)$ również
jest niezerową dziedziną. Z izomorfizmu pierwszego wynika,
że $S^{-1}\mathfrak p$ jest ideałem pierwszym.

Jeśli $\mathfrak q$ jest pierwszy w $S^{-1}A$, jego przeciwobraz
$\mathfrak p$ jest pierwszy w $A$ (własność ta wynika z
zachowania iloczynów przez homomorfizm). Nie może zawierać
elementu $s\in S$, ponieważ $s/1$ jest jednostką lokalizacji,
a ideał właściwy nie zawiera jednostek. Dla $a/s\in S^{-1}A$
mamy $a/s\in\mathfrak q$ wtedy i tylko wtedy, gdy
$a/1=(s/1)(a/s)\in\mathfrak q$, czyli $a\in\mathfrak p$.
Stąd $\mathfrak q=S^{-1}\mathfrak p$; ten sam rachunek pokazuje,
że przeciwobraz $S^{-1}\mathfrak p$ równa się $\mathfrak p$:
jeśli $a/1\in S^{-1}\mathfrak p$, istnieje $s\in S$ z
$sa\in\mathfrak p$, a pierwszość i $s\notin\mathfrak p$
dają $a\in\mathfrak p$. Obie operacje są więc odwrotne.
\end{proof}



\begin{wniosek}[Lokalizacja nie zwiększa wymiaru]
\label{aa:cor:krull-localization-bound}
Niech $A$ będzie niezerowym pierścieniem przemiennym z jedynką,
a $S\subseteq A$ zbiorem multiplikatywnym takim, że
$1\in S$ i $0\notin S$. Wtedy
$\dim(S^{-1}A)\le\dim A$.
\end{wniosek}
\begin{proof}
Każdy łańcuch ideałów pierwszych pierścienia $S^{-1}A$
przechodzi przez przeciwobrazy względem homomorfizmu
$A\to S^{-1}A$ na łańcuch ideałów pierwszych $A$.
Na mocy Twierdzenia~\ref{aa:tw:krull-localization-primes}
przeciwobrazy zachowują ścisłość wszystkich zawierań.
Długość łańcucha w lokalizacji nie przekracza zatem
kresu długości łańcuchów w $A$. Przechodząc do kresu
górnego po wszystkich łańcuchach w lokalizacji, otrzymujemy tezę.
\end{proof}
\begin{wniosek}[Wysokość a pierścień lokalny]\label{aa:cor:krull-height-local}
Niech $\mathfrak p$ będzie ideałem pierwszym pierścienia
przemiennego $A$ z jedynką. Wtedy
$\operatorname{ht}(\mathfrak p)=\dim A_{\mathfrak p}$,
gdzie $A_{\mathfrak p}=(A\setminus\mathfrak p)^{-1}A$.
\end{wniosek}
\begin{proof}
Ideały pierwsze $A_{\mathfrak p}$ odpowiadają, na mocy
Twierdzenia~\ref{aa:tw:krull-localization-primes}, dokładnie tym
ideałom pierwszym $\mathfrak q$ pierścienia $A$, które są rozłączne
z $A\setminus\mathfrak p$, czyli spełniają
$\mathfrak q\subseteq\mathfrak p$.
Każdy łańcuch takich ideałów można przedłużyć do
$\mathfrak p$, jeżeli jeszcze na nim nie kończy się.
Stąd kres długości wszystkich takich łańcuchów jest
równy kresowi długości łańcuchów kończących się w $\mathfrak p$.
\end{proof}

\section{Pierwsze obliczenia wymiaru}
\begin{twierdzenie}[Wymiar dziedziny ideałów głównych]\label{aa:tw:krull-pid}
Niech $A$ będzie dziedziną ideałów głównych. Jeżeli $A$ jest
ciałem, to $\dim A=0$; w przeciwnym razie $\dim A=1$.
\end{twierdzenie}
\begin{proof}
W dziedzinie $A$ ideał $(0)$ jest pierwszy. Jeśli $A$ jest
ciałem, innych ideałów pierwszych nie ma.
Jeśli $A$ nie jest ciałem, wybierzmy niezerowy element
nieodwracalny $a\in A$. Właściwy ideał $(a)$ jest zawarty
w pewnym ideale maksymalnym $\mathfrak m$ (lemat Zorna).
Jest on niezerowy, a jako maksymalny również pierwszy.

Pokażemy, że każdy niezerowy ideał pierwszy jest maksymalny.
Niech $\mathfrak p=(p)$ będzie niezerowy i pierwszy,
a $(p)\subseteq (b)\subsetneq A$. Ponieważ $p\in(b)$,
możemy napisać $p=bc$. Pierwszość ideału $(p)$ oznacza,
że $b\in(p)$ lub $c\in(p)$. Jeżeli $b=pu$, to
$p=puc$, a skracanie przez $p\ne0$ daje $1=uc$.
Wtedy $c$ jest jednostką i $(b)=(p)$. Jeżeli zaś
$c=pv$, to podobnie $1=bv$, zatem $(b)=A$,
co jest sprzeczne z założeniem. Nie istnieje więc
właściwy ideał ściśle zawierający $\mathfrak p$.
Każdy niezerowy ideał pierwszy jest maksymalny, więc łańcuchy
mają długość co najwyżej $1$, a $(0)\subsetneq(p)$
pokazuje, że długość $1$ występuje.
\end{proof}

\begin{lemat}[Wymiar pierścienia wielomianów w dwóch zmiennych]
\label{aa:lem:krull-plane-dimension}
Niech $K$ będzie dowolnym ciałem. Wówczas
$\dim K[U,V]=2$. Ponadto każdy niezerowy ideał pierwszy
$\mathfrak p$ pierścienia $K[U,V]$ jest maksymalny albo ma
postać $(f)$ dla pewnego nierozkładalnego $f\in K[U,V]$.
\end{lemat}
\begin{proof}
Oznaczmy $R=K[U,V]$. Ponieważ $K[U]$ jest pierścieniem
z jednoznacznością rozkładu, z lematu Gaussa wynika, że
$R=K[U][V]$ także jest takim pierścieniem. W szczególności
każdy niezerowy element nieodwracalny rozkłada się na
iloczyn elementów nierozkładalnych.

Wykażemy najpierw, że dla niezerowych, nierozkładalnych
i niestowarzyszonych $f,g\in R$ iloraz $R/(f,g)$ jest
skończenie wymiarową przestrzenią liniową nad $K$.
W pierścieniu $K(U)[V]$ wielomiany $f,g$ są względnie
pierwsze. Istotnie, każdy wspólny dzielnik dodatniego stopnia
względem $V$ dawałby na mocy lematu Gaussa wspólny czynnik
w $K[U,V]$; czynnik stopnia zero jest zaś jednostką
w $K(U)[V]$. Z tożsamości Bézouta w pierścieniu
$K(U)[V]$ otrzymujemy wielomiany $a,b\in K(U)[V]$ takie,
że $af+bg=1$. Po usunięciu mianowników dostajemy
niezerowy $h(U)\in K[U]$ należący do $(f,g)$.
Zamieniając role $U$ i $V$, uzyskujemy analogicznie
niezerowy $k(V)\in K[V]$ należący do $(f,g)$.
W ilorazie $R/(f,g)$ można obniżać stopnie osobno względem
$U$ i $V$ za pomocą równań $h(U)=0$ oraz $k(V)=0$:
współczynniki wiodące obu wielomianów są niezerowymi
skalarami z $K$. Zatem klasy jednomianów
$U^iV^j$ dla $0\le i<\deg h$ i $0\le j<\deg k$
rozpinają $R/(f,g)$ nad $K$. Jeżeli któryś z wielomianów
$h,k$ jest stałą niezerową, iloraz jest zerowy, więc
wniosek o skończonym wymiarze pozostaje prawdziwy.

Niech teraz $\mathfrak p\ne(0)$ będzie ideałem pierwszym $R$.
Zawiera on niezerowy element $a$. Rozkładając $a$ na
czynniki nierozkładalne i korzystając z pierwszości
$\mathfrak p$, znajdujemy nierozkładalne $f\in\mathfrak p$.
Jeżeli $\mathfrak p\ne(f)$, istnieje
$g\in\mathfrak p\setminus(f)$. Po rozłożeniu $g$ na
czynniki nierozkładalne przynajmniej jeden z nich,
powiedzmy $g_0$, należy do $\mathfrak p$ i nie jest
stowarzyszony z $f$. Wykazane wyżej rozumowanie daje
$\dim_K R/(f,g_0)<\infty$. Tym bardziej
$R/\mathfrak p$, będący jego ilorazem, jest skończenie
wymiarową $K$-algebrą. Ponieważ $\mathfrak p$ jest
pierwszy, $R/\mathfrak p$ jest dziedziną. Każda niezerowa
skończenie wymiarowa dziedzina nad ciałem jest ciałem:
mnożenie przez dowolny element niezerowy jest injektywnym,
a zatem surjektywnym, endomorfizmem przestrzeni liniowej.
Wobec tego $\mathfrak p$ jest maksymalny.

Każdy niemaksymalny ideał pierwszy $R$ jest więc równy
$(0)$ albo ma postać $(f)$ dla pewnego nierozkładalnego $f$.
Ideał $(f)$ jest pierwszy, ponieważ w pierścieniu z
jednoznacznością rozkładu elementy nierozkładalne są pierwsze.
Jeżeli $(f)\subsetneq\mathfrak q$ dla ideału pierwszego
$\mathfrak q$, to $\mathfrak q$ nie może być innym głównym
ideałem $(g)$ generowanym przez element nierozkładalny:
z $f\in(g)$ wynikałoby, że $f$ i $g$ są stowarzyszone,
a wtedy $(f)=(g)$. Zatem $\mathfrak q$ jest maksymalny.
Wynika stąd, że długość łańcucha ideałów pierwszych
$R$ nie przekracza $2$. Z drugiej strony łańcuch
$(0)\subsetneq(U)\subsetneq(U,V)$ ma długość $2$;
wszystkie trzy ideały są pierwsze na mocy odpowiednich
ilorazów $R$, $K[V]$ i $K$. Zatem $\dim K[U,V]=2$.
\end{proof}

\begin{przyklad}[Krzyżujące się proste]\label{aa:prz:krull-crossed-lines}
Niech $K$ będzie ciałem, $A=K[X,Y]/(XY)$, a $x,y$ oznaczają
klasy zmiennych. W pierścieniu $A$ mamy $xy=0$, choć
$x\ne0$ i $y\ne0$. Jeżeli $\mathfrak p$ jest ideałem
pierwszym $A$, z $xy\in\mathfrak p$ wynika
$x\in\mathfrak p$ lub $y\in\mathfrak p$.
W pierwszym przypadku $\mathfrak p$ pochodzi od ideału
pierwszego w $A/(x)\cong K[Y]$; w drugim od ideału
pierwszego w $A/(y)\cong K[X]$. Rozważmy teraz
\emph{dowolny} łańcuch $\mathfrak p_0\subsetneq\cdots
\subsetneq\mathfrak p_r$. Ponieważ $xy=0\in\mathfrak p_0$,
pierwszość $\mathfrak p_0$ daje $x\in\mathfrak p_0$
lub $y\in\mathfrak p_0$. W pierwszym przypadku każdy
$\mathfrak p_j$ zawiera $x$, w drugim każdy zawiera $y$.
Przez odpowiedniość ideałów pierwszych w ilorazie cały
łańcuch przechodzi więc do łańcucha w $K[Y]$ albo $K[X]$,
którego długość nie przekracza $1$.
Ponieważ $A/(x)\cong K[Y]$ i mamy ścisły łańcuch
$(x)\subsetneq(x,y)$, otrzymujemy $\dim A=1$.
Ideały $(x)$ i $(y)$ są pierwsze, ponieważ odpowiednie
ilorazy są dziedzinami $K[Y]$ i $K[X]$.
Są one minimalne: gdyby ideał pierwszy $\mathfrak q$
był ściśle zawarty w $(x)$, to $y\notin\mathfrak q$
(gdyż $y\notin(x)$), a z $xy=0$ wynikałoby
$x\in\mathfrak q$, zatem $\mathfrak q=(x)$,
sprzeczność. Analogicznie postępujemy z $(y)$.
Dwa składniki geometryczne, opisane równaniami $x=0$
i $y=0$, spotykają się w jednym punkcie, lecz nie zwiększają
wymiaru całości do dwóch. Rysunek~\ref{aa:fig:crossed-lines}
jest schematem dla $K=\mathbb R$; rachunek ideałowy
pozostaje prawdziwy nad dowolnym ciałem.
\end{przyklad}
\begin{figure}[htbp]
\centering
\setlength{\unitlength}{1mm}
\begin{picture}(82,43)
\put(11,19){\vector(1,0){60}}
\put(40,3){\vector(0,1){37}}
\put(73,16){$x$}
\put(43,37){$y$}
\put(54,13){$y=0$}
\put(44,31){$x=0$}
\put(37,15){$0$}
\end{picture}
\caption{Krzywa afiniczna $xy=0$ składa się z dwóch prostych
jednowymiarowych przecinających się w punkcie początkowym.
Nie jest całą dwuwymiarową płaszczyzną afiniczną; rysunek
przedstawia punkty rzeczywiste, a obliczenie wymiaru działa
nad dowolnym ciałem.}
\label{aa:fig:crossed-lines}
\end{figure}


\begin{twierdzenie}[Wymiar iloczynu dwóch pierścieni]
\label{aa:tw:krull-product-dimension}
Niech $A$ i $B$ będą niezerowymi pierścieniami przemiennymi
z jedynką. Wówczas
\[
 \dim(A\times B)=\max\{\dim A,\dim B\}.
\]
Każdy ideał pierwszy $A\times B$ ma dokładnie jedną
z postaci $\mathfrak p\times B$ lub $A\times\mathfrak q$,
gdzie $\mathfrak p$ jest ideałem pierwszym $A$, a
$\mathfrak q$ ideałem pierwszym $B$.
\end{twierdzenie}
\begin{proof}
W pierścieniu $A\times B$ rozważmy elementy
$e=(1_A,0_B)$ i $f=(0_A,1_B)$. Zachodzą równości
$ef=0$ i $e+f=1$. Jeżeli $\mathfrak r$ jest ideałem
pierwszym $A\times B$, to z $ef\in\mathfrak r$ wynika
$e\in\mathfrak r$ lub $f\in\mathfrak r$. Oba nie mogą
należeć do $\mathfrak r$, ponieważ wtedy należałaby doń
jedynka, a ideał pierwszy jest właściwy.

Jeżeli $f\in\mathfrak r$, to $\{0\}\times B\subseteq
\mathfrak r$. Istnieje zatem ideał $\mathfrak p\subseteq A$
taki, że $\mathfrak r=\mathfrak p\times B$; istotnie,
wystarczy wziąć $\mathfrak p=\{a\in A:(a,0)\in\mathfrak r\}$.
Iloraz $(A\times B)/(\mathfrak p\times B)$ jest izomorficzny
z $A/\mathfrak p$, a więc $\mathfrak r$ jest pierwszy
dokładnie wtedy, gdy $\mathfrak p$ jest pierwszy.
Analogicznie przypadek $e\in\mathfrak r$ daje
$\mathfrak r=A\times\mathfrak q$ dla ideału pierwszego
$\mathfrak q$ pierścienia $B$.

Ideały pierwsze należące do różnych opisanych rodzin
nie mogą się zawierać jeden w drugim. Istotnie,
$\mathfrak p\times B$ zawiera $f$, natomiast
$A\times\mathfrak q$ go nie zawiera (gdyż
$1_B\notin\mathfrak q$); w odwrotną stronę stosujemy $e$.
W obrębie każdej rodziny inkluzje ideałów odpowiadają
dokładnie inkluzjom w odpowiednim czynniku, np.
$\mathfrak p_1\times B\subsetneq\mathfrak p_2\times B$
wtedy i tylko wtedy, gdy
$\mathfrak p_1\subsetneq\mathfrak p_2$.
Stąd każdy łańcuch ideałów pierwszych iloczynu leży w jednej
rodzinie, a łańcuchy obu czynników przenoszą się do iloczynu
bez zmiany długości. Biorąc kresy długości, otrzymujemy
wskazaną równość, również gdy któryś wymiar jest nieskończony.
\end{proof}

\begin{przyklad}[Maksymalne ideały o różnych wysokościach]
\label{aa:prz:krull-maximal-varying-heights}
Niech $K$ będzie ciałem,
$A=K[T]$, $B=K[U,V]$, a $R=A\times B$.
Z poprzedniego twierdzenia, obliczenia wymiaru $K[T]$
(Twierdzenie~\ref{aa:tw:krull-pid}) oraz
Lematu~\ref{aa:lem:krull-plane-dimension} wynika,
że $\dim R=\max\{1,2\}=2$.
Rozważmy ideały
$\mathfrak m_1=(T)\times B$ oraz
$\mathfrak m_2=A\times(U,V)$. Są maksymalne, ponieważ
$R/\mathfrak m_i\cong K$ dla $i=1,2$.
Każdy łańcuch ideałów pierwszych kończący się w
$\mathfrak m_1$ ma postać
$\mathfrak p_0\times B\subsetneq\cdots\subsetneq(T)\times B$;
odpowiada więc łańcuchowi w $K[T]$.
Łańcuch $(0)\times B\subsetneq(T)\times B$ istnieje,
a dłuższy nie może istnieć, stąd
$\operatorname{ht}(\mathfrak m_1)=1$.
Podobnie łańcuch
$A\times(0)\subsetneq A\times(U)\subsetneq A\times(U,V)$
pokazuje $\operatorname{ht}(\mathfrak m_2)\ge2$,
a obliczenie wymiaru $K[U,V]$ daje nierówność przeciwną.
Zatem $\operatorname{ht}(\mathfrak m_2)=2$.
W szczególności ideał maksymalny nie musi mieć wysokości
równej wymiarowi całego pierścienia. Zauważmy, że oba
występujące czynniki są dziedzinami, ale ich iloczyn już nie.
\end{przyklad}

\begin{przyklad}[Iloraz może zmniejszyć wymiar]
Niech $A=\mathbb Z$, $I=(6)$. Ideały pierwsze $A/I$
odpowiadają ideałom pierwszym $\mathbb Z$ zawierającym $(6)$.
Wyjaśnijmy, skąd pochodzi pełna lista. Ideał pierwszy
$\mathfrak p$ pierścienia $\mathbb Z$ zawiera $(6)$
wtedy i tylko wtedy, gdy $6\in\mathfrak p$.
Ideał $(0)$ nie spełnia tego warunku, zaś dla
$\mathfrak p=(p)$, gdzie $p$ jest liczbą pierwszą,
warunek $p\mid6$ pozostawia jedynie $p=2$ lub $p=3$.
Stąd $\mathbb Z/(6)$ ma dokładnie dwa ideały pierwsze:
$(2)/(6)$ i $(3)/(6)$. Oba są maksymalne, ponieważ
ich ilorazy są odpowiednio $\mathbb F_2$ i $\mathbb F_3$.
Są też nieporównywalne, bo $2\notin(3)$ i $3\notin(2)$.
Żaden łańcuch nie ma długości dodatniej, więc
$\dim(\mathbb Z/(6))=0$. Równocześnie $\dim\mathbb Z=1$.
Konkretnie klasy $2$ i $3$ są niezerowymi dzielnikami zera,
gdyż $2\cdot3=0$ modulo $6$.
Przykład pokazuje, że przejście do ilorazu może zmniejszyć
wymiar i że pierścień wymiaru zero może mieć dzielniki zera.
\end{przyklad}



\begin{przyklad}[Wymiar zero nie oznacza, że pierścień jest ciałem]
\label{aa:prz:krull-dual-numbers}
Niech $K$ będzie dowolnym ciałem i
$A=K[T]/(T^2)$, gdzie przez
$\varepsilon=T+(T^2)$ oznaczamy klasę zmiennej $T$.
Każdy element $A$ ma jednoznaczną postać
$a+b\varepsilon$, gdzie $a,b\in K$, oraz
$\varepsilon^2=0$, lecz $\varepsilon\ne0$.
Jeżeli $\mathfrak p$ jest ideałem pierwszym $A$, to
$\varepsilon^2=0\in\mathfrak p$ wymusza
$\varepsilon\in\mathfrak p$. Otrzymujemy
$(\varepsilon)\subseteq\mathfrak p$; z drugiej strony
$A/(\varepsilon)\cong K$ jest ciałem, więc $(\varepsilon)$
jest ideałem maksymalnym i nie może być zawarty w większym
ideale właściwym. Zatem $\mathfrak p=(\varepsilon)$.
Jest to jedyny ideał pierwszy pierścienia $A$, wobec czego
$\dim A=0$. Pierścień $A$ nie jest jednak ciałem:
niezerowy element $\varepsilon$ nie jest odwracalny,
gdyż z $\varepsilon u=1$ wynikałoby
$0=\varepsilon^2u=\varepsilon$, co jest sprzecznością.
Wniosek~\ref{aa:cor:krull-nilradical} wyjaśnia ten wynik
również przez równość $\dim A=\dim(A/(\varepsilon))$.
\end{przyklad}

\begin{przyklad}[Pierścień o nieskończonym wymiarze Krulla]
\label{aa:prz:krull-infinite-dimension}
Niech $K$ będzie ciałem, a
$A=K[X_1,X_2,X_3,\ldots]$ pierścieniem wielomianów
w przeliczalnie wielu zmiennych, przy czym każdy wielomian
zależy od skończonej liczby zmiennych. Dla $r\ge1$
połóżmy $\mathfrak p_r=(X_1,\ldots,X_r)$
oraz $\mathfrak p_0=(0)$. Iloraz
$A/\mathfrak p_r\cong K[X_{r+1},X_{r+2},\ldots]$
nie ma dzielników zera, dlatego każdy $\mathfrak p_r$
jest ideałem pierwszym. Zawieranie
$\mathfrak p_r\subsetneq\mathfrak p_{r+1}$ jest ścisłe:
$X_{r+1}$ nie należy do $\mathfrak p_r$, co widać
z jego niezerowej klasy w tym ilorazie.
Dla każdej liczby naturalnej $m\ge1$ otrzymujemy więc
łańcuch $\mathfrak p_0\subsetneq\cdots\subsetneq
\mathfrak p_m$ długości $m$. Długości łańcuchów
nie są ograniczone, czyli $\dim A=+\infty$.
Przykład nie przeczy późniejszemu obliczeniu wymiaru
$K[X_1,\ldots,X_n]$: tam liczba generatorów jest skończona.
\end{przyklad}
\begin{przyklad}[Lokalizacja może obniżyć wymiar]\label{aa:prz:krull-local-example}
W $\mathbb Z$ wybierzmy $S=\mathbb Z\setminus\{0\}$.
Wtedy $S^{-1}\mathbb Z=\mathbb Q$ ma wymiar $0$,
a $\mathbb Z$ ma wymiar $1$. Z drugiej strony
lokalizacja $\mathbb Z_{(p)}$ względem ideału $(p)$
ma dokładnie dwa ideały pierwsze:
$(0)$ oraz $(p)\mathbb Z_{(p)}$. Wynika to z bijekcji
z ideałami pierwszymi $\mathbb Z$ zawartymi w $(p)$.
Dlatego nadal ma wymiar $1$; odpowiadający im łańcuch
jest ścisły. Nie należy więc oczekiwać, że każda
lokalizacja obniża wymiar: może go zachować.
\end{przyklad}

\begin{uwaga}[Związek z bazami transcendencji]
\label{aa:uw:krull-transcendence-preview}
W pierścieniu $K[X_1,\ldots,X_n]$ występuje łańcuch
$(0)\subsetneq(X_1)\subsetneq\cdots\subsetneq
(X_1,\ldots,X_n)$ długości $n$. To dowodzi nierówności
$\dim K[X_1,\ldots,X_n]\ge n$.
Dla $n=1$ równość $\dim K[X_1]=1$ wynika
z Twierdzenia~\ref{aa:tw:krull-pid}, a dla $n=2$
z Lematu~\ref{aa:lem:krull-plane-dimension}.
Ogólny dowód nierówności odwrotnej dla dowolnego
$n\ge3$ wymaga dodatkowego argumentu, który przeprowadzimy
po rozwinięciu teorii rozszerzeń całkowitych.
Tam też zwiążemy wymiar z
$\operatorname{trdeg}_K\operatorname{Frac}(A)$ dla skończenie
generowanych dziedzin będących $K$-algebrami.
\end{uwaga}

% AA-40 | Audyt I: kontrola dowodów i uzupełnienia o włókna skończone oraz rozgałęzienie.
\chapter{Całkowitość i twierdzenia o ideałach pierwszych}
\label{aa:ch:integral-extensions}

\section{Elementy całkowite i skończoność modułowa}
W całym rozdziale zakładamy, że $A\subseteq B$ są
przemiennymi pierścieniami z jedynką i wspólnym elementem $1$.
Pierścienie te nie muszą być dziedzinami całkowitości, chyba że
zostało to osobno zaznaczone. Rozróżniamy dwa pojęcia:
\emph{element algebraiczny} nad ciałem oraz \emph{element całkowity}
nad pierścieniem. Drugie wymaga równania monicznego, a nie
jakiegokolwiek równania wielomianowego. Różnica jest istotna
nawet wewnątrz tego samego ciała ułamków: każdy element
$\mathbb Q$ jest algebraiczny nad $\mathbb Q$, ale
np. $1/2$ nie jest całkowity nad $\mathbb Z$.
Dalej pokażemy, dlaczego moniczność pozwala sprowadzić
wszystkie wysokie potęgi elementu do skończenie wielu
niższych potęg.

\begin{definicja}[Całkowitość]\label{aa:def:integral-element}
Element $b\in B$ jest \emph{całkowity nad $A$}, jeśli istnieją
$n\ge1$ i $a_0,\ldots,a_{n-1}\in A$ takie, że
\[b^n+a_{n-1}b^{n-1}+\cdots+a_1b+a_0=0.\]
Rozszerzenie $A\subseteq B$ jest \emph{całkowite}, jeżeli każdy
element $B$ jest całkowity nad $A$. Mówimy, że $B$ jest
\emph{skończony modułowo nad $A$}, gdy istnieją
$b_1,\ldots,b_m\in B$ takie, że
$B=Ab_1+\cdots+Ab_m$.
\end{definicja}

\begin{przyklad}[Pierwiastki kwadratowe i ułamki]\label{aa:prz:integral-sqrt}
Element $\sqrt2$ jest całkowity nad $\mathbb Z$, gdyż spełnia
moniczne równanie $T^2-2=0$. Stąd także
$3+2\sqrt2$ jest całkowity, co wynika z ogólnego twierdzenia
poniżej; można sprawdzić to bezpośrednio z równania
$(T-3)^2-8=0$, czyli $T^2-6T+1=0$.
Natomiast $1/2\in\mathbb Q$ jest algebraiczny nad
$\mathbb Q$, ale nie jest całkowity nad $\mathbb Z$.
Gdyby spełniał moniczne równanie nad $\mathbb Z$ stopnia $n$,
po pomnożeniu przez $2^n$ otrzymalibyśmy
$1+2a_{n-1}+\cdots+2^na_0=0$, co jest niemożliwe,
ponieważ lewa strona jest liczbą nieparzystą.
\end{przyklad}

\begin{lemat}[Skończony moduł a równanie moniczne]
\label{aa:lem:integrality-determinant}
Niech $b\in B$. Następujące warunki są równoważne:
\begin{enumerate}[label=\textup{(\roman*)}]
\item element $b$ jest całkowity nad $A$;
\item podpierścień $A[b]$ jest skończenie generowanym $A$-modułem;
\item istnieje skończenie generowany $A$-podmoduł $M\subseteq B$
 zawierający $1$ i spełniający $bM\subseteq M$.
\end{enumerate}
\end{lemat}
\begin{proof}
Załóżmy (i). Jeżeli $b^n$ jest $A$-kombinacją niższych
potęg, to po pomnożeniu przez $b$ to samo dotyczy $b^{n+1}$;
indukcja względem wykładnika pokazuje, że każda potęga $b$
należy do $A$-modułu rozpiętego przez
$1,b,\ldots,b^{n-1}$. Stąd (ii). Warunek (ii) implikuje (iii),
ponieważ wystarczy wybrać $M=A[b]$.

Niech zachodzi (iii), $M=Am_1+\cdots+Am_r$, gdzie $1\in M$.
Dla każdego $i$ zapiszmy
$bm_i=\sum_{j=1}^r a_{ij}m_j$ z $a_{ij}\in A$.
Niech $C=(a_{ij})\in M_r(A)$ oraz niech $I_r$ oznacza
macierz jednostkową stopnia $r$. Układ równań możemy zapisać
$(bI_r-C)\mathbf m=0$, gdzie
$\mathbf m=(m_1,\ldots,m_r)^{\mathsf T}$.
To jest równość wektorów w $M^r$, nie zaś równanie
macierzowe nad ciałem: każdy współczynnik macierzy
z pierścienia $B$ działa na odpowiedniej składowej modułu.
Mnożąc przez macierz dołączoną $\operatorname{adj}(bI-C)$
dostajemy $\det(bI-C)m_i=0$ dla każdego $i$.
Ponieważ $1=\sum_i c_im_i$ dla pewnych $c_i\in A$,
wynika stąd $\det(bI-C)\cdot1=0$ w $B$.
Wyznacznik $\det(TI-C)\in A[T]$ jest moniczny stopnia $r$:
w rozwinięciu wyznacznika składnik $T^r$ pochodzi tylko
z iloczynu elementów diagonalnych, natomiast pozostałe
składniki mają stopień mniejszy. Zatem $b$ spełnia
moniczne równanie i zachodzi (i). Argument nie wymaga,
aby $A$ lub $B$ były dziedzinami: tożsamość macierzy
dołączonej i rozwinięcie wyznacznika obowiązują nad każdym
pierścieniem przemiennym. Istotne jest natomiast,
że $M$ zawiera $1$; inaczej równanie mogłoby
anihilować cały $M$ bez anihilowania jedynki pierścienia $B$.
\end{proof}

\begin{twierdzenie}[Domknięcie na działania i przechodniość]
\label{aa:tw:integrality-transitive}
Niech $A\subseteq B\subseteq C$ będą przemiennymi pierścieniami
z jedynką. Wtedy:
\begin{enumerate}[label=\textup{(\roman*)}]
\item jeżeli $b,c\in B$ są całkowite nad $A$, to
$-b$, $b+c$ i $bc$ są całkowite nad $A$;
\item jeżeli $B/A$ i $C/B$ są rozszerzeniami całkowitymi,
 to $C/A$ także jest całkowite.
\end{enumerate}
\end{twierdzenie}
\begin{proof}
Dla (i) z Lematu~\ref{aa:lem:integrality-determinant} wynika,
że $A[b]$ jest skończenie generowanym $A$-modułem.
Ten sam moniczny wielomian dla $c$ ma współczynniki w
$A\subseteq A[b]$, a więc $c$ jest również całkowity
nad $A[b]$. Stąd $A[b,c]=A[b][c]$ jest skończonym
$A[b]$-modułem. Jeżeli
$u_1,\ldots,u_s$ generują $A[b]$ nad $A$, a
$v_1,\ldots,v_t$ generują $A[b,c]$ nad $A[b]$, to
$u_iv_j$ (dla wszystkich $i,j$) generują $A[b,c]$ nad $A$:
współczynniki względem $v_j$ rozwijamy względem $u_i$.
Stosując kryterium (iii) lematu do
$M=A[b,c]$ i do elementów $-b,b+c,bc$ otrzymujemy
ich całkowitość nad $A$.

Dla (ii) weźmy dowolne $z\in C$. Istnieje równanie moniczne
$z^n+b_{n-1}z^{n-1}+\cdots+b_0=0$ z $b_i\in B$.
Każdy $b_i$ jest całkowity nad $A$, więc
$A'=A[b_0,\ldots,b_{n-1}]$ jest skończenie generowanym
$A$-modułem (stosujemy poprzedni argument indukcyjnie).
Z kolei $A'[z]$ jest skończenie generowanym $A'$-modułem.
Przez przechodniość skończoności modułowej $A'[z]$ jest
skończonym $A$-modułem, a lemat pokazuje, że $z$ jest
całkowity nad $A$. Ponieważ $z$ był dowolny, dostajemy (ii).
Zwróćmy uwagę na rolę skończonej liczby współczynników
$b_0,\ldots,b_{n-1}$: nie zakładamy, że cały pierścień
$B$ jest skończonym $A$-modułem. Każdy badany element
$z$ pozwala jednak znaleźć skończoną podalgebrę
$A'\subseteq B$ wystarczającą dla tego dowodu.
\end{proof}


\begin{twierdzenie}[Całkowitość algebry o skończonej liczbie generatorów]
\label{aa:tw:integral-finite-type-finite-module}
Niech $A\subseteq B$ będą przemiennymi pierścieniami z jedynką.
Jeżeli $B=A[b_1,\ldots,b_m]$ dla pewnej skończonej rodziny
elementów $b_i$ całkowitych nad $A$, to $B$ jest skończenie
generowanym $A$-modułem. Odwrotnie, jeśli $B$ jest
skończenie generowanym $A$-modułem, to rozszerzenie $B/A$
jest całkowite.
\end{twierdzenie}
\begin{proof}
Dla $m=0$ mamy $B=A$, czyli moduł jest generowany
przez $1$. Niech $m\ge1$. Dla każdego $i$ wybierzmy
moniczne równanie stopnia $d_i\ge1$ spełnione przez $b_i$,
ze współczynnikami w $A$. Po dołączeniu kolejnych
elementów równanie to pozostaje prawdziwe, ponieważ
$A\subseteq A[b_1,\ldots,b_{i-1}]$. W rezultacie
$A[b_1,\ldots,b_i]$ jest modułem nad
$A[b_1,\ldots,b_{i-1}]$ generowanym przez
$1,b_i,\ldots,b_i^{d_i-1}$. Wykonując indukcyjnie
rozwinięcie współczynników, dostajemy skończoną rodzinę
$\{b_1^{j_1}\cdots b_m^{j_m}:0\le j_i<d_i\}$
generującą $B$ jako $A$-moduł. Rodzina ta nie musi
być liniowo niezależna.

W drugą stronę dla każdego $b\in B$ moduł $M=B$
jest skończenie generowany nad $A$, zawiera jedynkę
i spełnia $bM\subseteq M$. Zatem warunek (iii) z
Lematu~\ref{aa:lem:integrality-determinant} implikuje,
że $b$ jest całkowity nad $A$. Ponieważ $b$ był dowolny,
całe rozszerzenie $B/A$ jest całkowite.
\end{proof}

\begin{przyklad}[Rozszerzenie całkowite nie musi być skończone modułowo]
\label{aa:prz:integral-infinite-extension}
Niech $K=\mathbb F_2$ i niech $L=\overline{\mathbb F}_2$
będzie algebraicznym domknięciem tego ciała. Każdy element
$\alpha\in L$ jest algebraiczny nad $K$, zatem spełnia
moniczne równanie nad $K$: wystarczy podzielić
wielomian anihilujący przez jego niezerowy współczynnik
wiodący. Rozszerzenie $K\subseteq L$ jest więc całkowite.

Nie jest ono jednak skończone modułowo. Dla każdej liczby
$m\ge1$ wewnątrz $L$ istnieje ciało o $2^m$ elementach,
np. zbiór wszystkich pierwiastków wielomianu
$T^{2^m}-T$; pochodna tego wielomianu wynosi $-1$,
więc w ciele $L$ ma on dokładnie $2^m$ różnych pierwiastków.
Dla $q=2^m$ zachodzi $(u+v)^q=u^q+v^q$ oraz
$(uv)^q=u^qv^q$ (charakterystyka $2$). Jeśli $u,v$ są
pierwiastkami, to są nimi także $u+v$, $-u$ i $uv$.
Dla $u\ne0$ równość $u^q=u$ daje
$(u^{-1})^q=u^{-1}$. Zatem zbiór pierwiastków
jest rzeczywiście podciałem $L$ o $2^m$ elementach.
Ponieważ $m$ jest dowolne, $L$ jest nieskończone.
Gdyby $L$ było generowane nad $\mathbb F_2$ przez
$m_0$ elementów jako moduł, miałoby najwyżej $2^{m_0}$
elementów, gdyż każda kombinacja liniowa ma tylko
$2^{m_0}$ możliwych układów współczynników.
Otrzymujemy sprzeczność. Przykład oddziela hipotezę
\emph{całkowitości} od dodatkowego założenia o skończonej
liczbie generatorów algebry w poprzednim twierdzeniu.
\end{przyklad}

\section{Wznoszenie i nieporównywalność ideałów pierwszych}
\begin{lemat}[Ideał maksymalny i jego przeciwobraz]\label{aa:lem:integral-maximal}
Niech $A\subseteq B$ będzie rozszerzeniem całkowitym,
a $\mathfrak q$ ideałem pierwszym $B$.
Wtedy $\mathfrak q$ jest maksymalny w $B$ wtedy i tylko wtedy,
gdy $\mathfrak p=\mathfrak q\cap A$ jest maksymalny w $A$.
\end{lemat}
\begin{proof}
Przede wszystkim odwzorowanie
$A/\mathfrak p\to B/\mathfrak q$ jest
różnowartościowe: jego jądrem jest
$(\mathfrak q\cap A)/\mathfrak p=(0)$.
Rozszerzenie $A/\mathfrak p\subseteq B/\mathfrak q$ jest
całkowite, ponieważ redukcja równania monicznego modulo
$\mathfrak q$ daje moniczne równanie dla dowolnej klasy
$b+\mathfrak q$ ze współczynnikami w $A/\mathfrak p$.
Oba pierścienie ilorazowe są dziedzinami całkowitości.
Jeżeli $\mathfrak q$ jest maksymalny, $B/\mathfrak q$ jest ciałem.
Dla niezerowego $a\in A/\mathfrak p$ jego odwrotność
$a^{-1}\in B/\mathfrak q$ spełnia moniczne równanie
$a^{-n}+c_{n-1}a^{-(n-1)}+\cdots+c_0=0$ nad $A/\mathfrak p$.
Po pomnożeniu przez $a^{n-1}$ dostajemy
$a^{-1}=-(c_{n-1}+c_{n-2}a+\cdots+c_0a^{n-1})\in A/\mathfrak p$.
A więc $A/\mathfrak p$ jest ciałem i $\mathfrak p$ jest maksymalny.

Odwrotnie, jeśli $A/\mathfrak p$ jest ciałem, to
$B/\mathfrak q$ jest dziedziną całkowitości algebraiczną
nad tym ciałem: każdy jej element spełnia niezerowy
wielomian nad $A/\mathfrak p$.
Dla każdego niezerowego $b\in B/\mathfrak q$ istnieje
niezerowy wielomian anihilujący nad ciałem $A/\mathfrak p$.
Spośród wszystkich monicznych wielomianów anihilujących
$b$ wybierzmy $m_b(T)$ o najmniejszym stopniu; możliwość
unormowania współczynnika wiodącego zapewnia fakt,
że pierścień współczynników jest ciałem.
Ma on niezerowy wyraz wolny: gdyby był zerowy,
to $m_b(T)=T h(T)$ i z $0=b h(b)$ w dziedzinie
$B/\mathfrak q$ wynikałoby $h(b)=0$. Ponieważ $b\ne0$,
przeczyłoby to minimalnemu stopniowi $m_b$.
Z równania $m_b(b)=0$ wynika
$b\cdot(-b^{n-1}-c_{n-1}b^{n-2}-\cdots-c_1)=c_0$;
ponieważ $c_0\ne0$ należy do ciała $A/\mathfrak p$,
$b$ ma odwrotność. Zatem $B/\mathfrak q$ jest ciałem. Pokazaliśmy obie
implikacje osobno, wykorzystując założenie pierwszości
$\mathfrak q$ właśnie po to, by iloraz był dziedziną.
\end{proof}

\begin{twierdzenie}[Twierdzenie o istnieniu ideału leżącego nad danym]\label{aa:tw:integral-lying-over}
Jeśli $A\subseteq B$ jest rozszerzeniem całkowitym, to dla
każdego ideału pierwszego $\mathfrak p\subset A$ istnieje
ideał pierwszy $\mathfrak q\subset B$ spełniający
$\mathfrak q\cap A=\mathfrak p$.
\end{twierdzenie}
\begin{proof}
Ustalmy $\mathfrak p$ i połóżmy $S=A\setminus\mathfrak p$.
Lokalizacja $S^{-1}A=A_{\mathfrak p}$ jest pierścieniem
lokalnym o jedynym ideale maksymalnym
$\mathfrak pA_{\mathfrak p}$, co wynika z opisu jednostek:
$a/s$ jest jednostką dokładnie wtedy, gdy $a\notin\mathfrak p$.
Rozszerzenie $A_{\mathfrak p}\subseteq S^{-1}B$ jest
całkowite. Aby to sprawdzić, weźmy dowolny ułamek
$b/s\in S^{-1}B$ i moniczne równanie
$b^n+a_{n-1}b^{n-1}+\cdots+a_0=0$ w $B$.
Po podzieleniu przez $s^n$ dostajemy
\[
(b/s)^n+(a_{n-1}/s)(b/s)^{n-1}
+\cdots+(a_1/s^{n-1})(b/s)+a_0/s^n=0.
\]
Wszystkie współczynniki tego monicznego równania należą
do $A_{\mathfrak p}$, bo $s\in S$.
Jest to równanie w $S^{-1}B$, niezależnie od tego,
czy $B$ ma dzielniki zera.
Pierścień $S^{-1}B$ jest niezerowy: gdyby $1/1=0/1$
w $S^{-1}B$, istniałby $s\in S$ spełniający $s\cdot1_B=0$,
a więc $s=0$ w $B$ i w $A$ z różnowartościowości
włączenia $A\subseteq B$. Przeczyłoby to $0\notin S$.
Wybierzmy jego ideał maksymalny $\mathfrak n$.
Lemat~\ref{aa:lem:integral-maximal} mówi, że
$\mathfrak n\cap A_{\mathfrak p}$ jest maksymalny, więc
równa się $\mathfrak pA_{\mathfrak p}$.
Ideał pierwszy $\mathfrak q$ będący przeciwobrazem
$\mathfrak n$ względem $B\to S^{-1}B$ ma zatem
$\mathfrak q\cap A=\mathfrak p$: jego przecięcie z $A$
jest przeciwobrazem $\mathfrak n\cap A_{\mathfrak p}
=\mathfrak p A_{\mathfrak p}$ względem
$A\to A_{\mathfrak p}$, a ten przeciwobraz to dokładnie
$\mathfrak p$ na mocy
Twierdzenia~\ref{aa:tw:krull-localization-primes}.
\end{proof}

\begin{twierdzenie}[Wznoszenie łańcuchów ideałów pierwszych]\label{aa:tw:integral-going-up}
Niech $A\subseteq B$ będzie rozszerzeniem całkowitym.
Jeśli $\mathfrak p_0\subsetneq\cdots\subsetneq\mathfrak p_r$
jest łańcuchem ideałów pierwszych $A$, a
$\mathfrak q_0$ ideałem pierwszym $B$ z
$\mathfrak q_0\cap A=\mathfrak p_0$, to istnieją ideały pierwsze
$\mathfrak q_0\subsetneq\cdots\subsetneq\mathfrak q_r$
takie, że $\mathfrak q_i\cap A=\mathfrak p_i$ dla każdego $i$.
\end{twierdzenie}
\begin{proof}
Wystarcza podnieść jeden krok; następnie powtórzymy
konstrukcję $r$ razy. Jeśli $\mathfrak q_i$ leży
nad $\mathfrak p_i$, to odwzorowanie
$A/\mathfrak p_i\to B/\mathfrak q_i$ jest
różnowartościowe, ponieważ
$\mathfrak q_i\cap A=\mathfrak p_i$.
Rozszerzenie to jest całkowite: każda klasa elementu $b\in B$
spełnia równanie moniczne otrzymane przez redukcję modulo
$\mathfrak q_i$ równania dla $b$ nad $A$.
Ideał $\mathfrak p_{i+1}/\mathfrak p_i$ jest pierwszy
w $A/\mathfrak p_i$. Twierdzenie~\ref{aa:tw:integral-lying-over}
daje nad nim ideał pierwszy $\overline{\mathfrak q}_{i+1}$
w $B/\mathfrak q_i$. Jego przeciwobraz $\mathfrak q_{i+1}$
w $B$ zawiera $\mathfrak q_i$ i leży nad
$\mathfrak p_{i+1}$. Inkluzja jest ścisła, gdyż kontrakcje
obu ideałów są różne. Powtarzamy to skończenie wiele razy.
\end{proof}

\begin{twierdzenie}[Nieporównywalność ideałów nad tym samym]\label{aa:tw:integral-incomparability}
Niech $A\subseteq B$ będzie rozszerzeniem całkowitym,
a $\mathfrak q_1\subseteq\mathfrak q_2$ ideałami pierwszymi $B$.
Jeśli $\mathfrak q_1\cap A=\mathfrak q_2\cap A$, to
$\mathfrak q_1=\mathfrak q_2$.
\end{twierdzenie}
\begin{proof}
Załóżmy przeciwnie, że inkluzja jest ścisła.
Przechodzimy do całkowitego rozszerzenia dziedzin
$A/\mathfrak p\subseteq B/\mathfrak q_1$: włączenie
wynika z równości kontrakcji $\mathfrak q_1\cap A
=\mathfrak p$, a całkowitość przechodzi na ilorazy, gdzie
$\mathfrak p=\mathfrak q_i\cap A$.
Niezerowy ideał pierwszy $\mathfrak q_2/\mathfrak q_1$
ma przecięcie $0$ z $A/\mathfrak p$.
Wybierzmy w nim niezerowy element $b$. Ponieważ $b$ jest
całkowity, spełnia moniczne równanie minimalnego możliwego
stopnia nad $A/\mathfrak p$:
$b^n+a_{n-1}b^{n-1}+\cdots+a_0=0$.
Ponieważ $b$ należy do ideału, również $a_0$ do niego należy;
a zatem $a_0=0$, gdyż przecięcie z $A/\mathfrak p$ jest zerowe.
Dzieląc równanie przez $b$ w dziedzinie $B/\mathfrak q_1$
(znaczy: skracając niezerowy czynnik $b$), dostajemy moniczne
równanie stopnia $n-1$ dla $b$, wbrew minimalności $n$.
\end{proof}

\begin{twierdzenie}[Skończoność włókien przy rozszerzeniu skończonym modułowo]
\label{aa:tw:integral-finite-fibers}
Niech $A\subseteq B$ będzie rozszerzeniem pierścieni przemiennych
z jedynką, w którym $B$ jest generowany jako $A$-moduł przez
$m\ge1$ elementów. Dla każdego ideału pierwszego
$\mathfrak p\subset A$ istnieje co najmniej jeden, a co najwyżej
$m$ ideałów pierwszych $\mathfrak q\subset B$ spełniających
$\mathfrak q\cap A=\mathfrak p$.
\end{twierdzenie}
\begin{proof}
Z Twierdzenia~\ref{aa:tw:integral-finite-type-finite-module}
rozszerzenie jest całkowite, więc istnienie przynajmniej
jednego takiego ideału wynika z
Twierdzenia~\ref{aa:tw:integral-lying-over}.
Ustalmy $\mathfrak p$ i połóżmy
$S=A\setminus\mathfrak p$ oraz
$\kappa(\mathfrak p)=\operatorname{Frac}(A/\mathfrak p)$.
Rozważmy pierścień
\[
 C=B\otimes_A\kappa(\mathfrak p)
  \cong S^{-1}(B/\mathfrak p B).
\]
Pierwsza postać ilustruje zmianę pierścienia współczynników;
druga wynika z faktu, że lokalizacja i przejście do ilorazu
komutują. Ponieważ $B$ ma $m$ generatorów jako $A$-moduł,
ich obrazy generują $C$ jako przestrzeń liniową nad
$\kappa(\mathfrak p)$. Zatem
$\dim_{\kappa(\mathfrak p)}C\le m$.

Ideały pierwsze $C$ odpowiadają dokładnie ideałom pierwszym
$\mathfrak q$ pierścienia $B$, które zawierają
$\mathfrak p B$ i są rozłączne z $S$. Warunki te są równoważne
równości $\mathfrak q\cap A=\mathfrak p$.
Każdy ideał pierwszy $\mathfrak n$ pierścienia $C$ jest
maksymalny: iloraz $C/\mathfrak n$ jest skończenie wymiarową
dziedziną nad $\kappa(\mathfrak p)$, a więc ciałem.
Istotnie, mnożenie przez dowolny niezerowy element tej
dziedziny jest różnowartościowym endomorfizmem jej
skończenie wymiarowej przestrzeni liniowej, zatem jest
surjektywne i element ma odwrotność.

Gdy $\mathfrak n_1,\ldots,\mathfrak n_r$ są parami różnymi
ideałami maksymalnymi $C$, są parami współmaksymalne.
Chińskie twierdzenie o resztach daje surjekcję
$C\twoheadrightarrow\prod_{j=1}^r C/\mathfrak n_j$.
Każdy czynnik po prawej ma wymiar nad
$\kappa(\mathfrak p)$ co najmniej $1$, więc
$r\le\dim_{\kappa(\mathfrak p)}C\le m$.
Wobec dowolności tej rodziny istnieje co najwyżej
$m$ ideałów maksymalnych $C$, a zatem tyle samo
ideałów pierwszych $B$ leżących nad $\mathfrak p$.
\end{proof}

\begin{uwaga}\label{aa:uw:integral-fibers-finite-module}
Ograniczenie liczby ideałów w powyższym twierdzeniu
korzysta ze \emph{skończoności modułowej}, nie tylko
z całkowitości. Natomiast twierdzenia o istnieniu
ideału leżącego nad danym, wznoszeniu i nieporównywalności
wymagają jedynie całkowitości rozszerzenia.
\end{uwaga}

\begin{wniosek}[Wymiar nie zmienia się w rozszerzeniu całkowitym]\label{aa:cor:integral-dimension}
Jeśli $A\subseteq B$ jest rozszerzeniem całkowitym niezerowych
pierścieni przemiennych z jedynką, to $\dim A=\dim B$.
\end{wniosek}
\begin{proof}
Każdy łańcuch ideałów pierwszych $A$ podnosimy do łańcucha
tej samej długości w $B$ za pomocą
Twierdzeń~\ref{aa:tw:integral-lying-over} i
\ref{aa:tw:integral-going-up}, stąd $\dim A\le\dim B$.
Z kolei weźmy dowolny ścisły łańcuch
$\mathfrak q_0\subsetneq\cdots\subsetneq\mathfrak q_r$
w $B$. Jego kontrakcje $\mathfrak p_i=\mathfrak q_i\cap A$
są ideałami pierwszymi i spełniają
$\mathfrak p_i\subseteq\mathfrak p_{i+1}$.
Żadne dwa sąsiednie ideały nie mogą mieć tej samej kontrakcji,
gdyż przeczyłoby to
Twierdzeniu~\ref{aa:tw:integral-incomparability}.
Stąd $\dim B\le\dim A$.
\end{proof}


\section{Przykłady wznoszenia oraz rola całkowitości}
\begin{przyklad}[Dwa ideały nad $(5)$ w pierścieniu Gaussa]
\label{aa:prz:integral-gaussian-five}
Niech $A=\mathbb Z$ i $B=\mathbb Z[i]$, gdzie $i^2=-1$.
Każdy element $a+bi\in B$ jest całkowity nad $\mathbb Z$:
ma równanie
\[
T^2-2aT+(a^2+b^2)=0,
\]
co sprawdzamy po podstawieniu $T=a+bi$. Zatem $B/A$
jest rozszerzeniem całkowitym.

Rozpatrzmy dwa homomorfizmy
$\varphi_2,\varphi_3:B\to\mathbb F_5$ określone przez
$\varphi_2(a+bi)=\overline a+2\overline b$ oraz
$\varphi_3(a+bi)=\overline a+3\overline b$.
Oba są homomorfizmami pierścieni, ponieważ
$2^2=3^2=-1$ w $\mathbb F_5$.
Są surjektywne: każdą klasę z $\mathbb F_5$ można otrzymać
z odpowiedniej liczby całkowitej. Ich jądra
\[
\mathfrak q_2=(5,i-2),\qquad
\mathfrak q_3=(5,i-3)
\]
są więc ideałami maksymalnymi $B$. Ich kontrakcje do
$\mathbb Z$ wynoszą $(5)$, bo ograniczenie obu
homomorfizmów do $\mathbb Z$ jest redukcją modulo $5$.
Są to dwa różne ideały: element $i-2$ należy do
$\mathfrak q_2$, lecz
$\varphi_3(i-2)=\overline 1\ne0$.
Nie mogą się zawierać jeden w drugim, gdyż oba są
maksymalne. Otrzymujemy konkretny obraz tego, że nad jednym
ideałem pierwszym może leżeć więcej niż jeden ideał pierwszy,
a twierdzenie o nieporównywalności nie orzeka ich jedyności.
\end{przyklad}



\begin{przyklad}[Nierozkładalność nad liczbą pierwszą a liczba ideałów leżących nad nią]
\label{aa:prz:integral-gaussian-three}
Dla $A=\mathbb Z$ i $B=\mathbb Z[i]$ porównajmy
Przykład~\ref{aa:prz:integral-gaussian-five} z ideałem $(3)$.
Rozważmy homomorfizm $\mathbb Z[T]\to B$,
$T\mapsto i$. Jest on surjektywny, a jego jądro jest
ideałem $(T^2+1)$: dzieląc dowolny wielomian przez
moniczny wielomian $T^2+1$, otrzymujemy resztę $a+bT$;
jeżeli $a+bi=0$ w $B\subseteq\mathbb C$, to $a=b=0$.
Zatem
\[
 B/(3)\cong\mathbb F_3[T]/(T^2+1).
\]
W ciele $\mathbb F_3$ wartości $T^2+1$ dla
$T=0,1,2$ są odpowiednio $1,2,2$; wielomian stopnia
$2$ nie ma więc pierwiastka i jest nierozkładalny.
W konsekwencji $B/(3)$ jest ciałem dziewięcioelementowym,
a $(3)\subseteq B$ jest ideałem maksymalnym.
Jeżeli ideał pierwszy $\mathfrak q$ pierścienia $B$ ma
kontrakcję $(3)$ do $\mathbb Z$, to zawiera $(3)$.
Odpowiada więc ideałowi pierwszemu ciała $B/(3)$;
tym ideałem może być wyłącznie $(0)$. Stąd
$\mathfrak q=(3)$ i nad $(3)$ leży dokładnie jeden
ideał pierwszy $B$.
Natomiast nad $(5)$ leżą dwa różne ideały maksymalne,
opisane w poprzednim przykładzie. Ta różnica pokazuje,
że twierdzenie o istnieniu ideału leżącego nad danym
nie orzeka liczby takich ideałów.
\end{przyklad}
\begin{przyklad}[Rozgałęzienie ideału $(2)$ w pierścieniu Gaussa]
\label{aa:prz:integral-gaussian-two}
Dla $B=\mathbb Z[i]$ mamy
\[
 B/(2)\cong\mathbb F_2[T]/(T^2+1)
 =\mathbb F_2[T]/((T+1)^2).
\]
Niech $\varepsilon$ oznacza klasę $T+1$ w ostatnim ilorazie.
Wtedy $\varepsilon^2=0$, ale $\varepsilon\ne0$.
Każdy element tego pierścienia ma jednoznaczną postać
$a+b\varepsilon$ dla $a,b\in\mathbb F_2$;
nieodwracalne są dokładnie elementy ideału
$(\varepsilon)$, gdyż dla $a=1$ zachodzi
$(1+b\varepsilon)^{-1}=1+b\varepsilon$.
Pierścień $B/(2)$ ma zatem dokładnie jeden ideał pierwszy,
równy $(\varepsilon)$, będący jego jedynym ideałem maksymalnym.
Jego przeciwobrazem w $B$ jest
$\mathfrak q=(2,i+1)=(1+i)$; równość wynika z
$2=-i(1+i)^2$ oraz $i+1=1+i$.
Ponadto $(2)=\mathfrak q^2$ jako ideały $B$.

Porównując z Przykładami~\ref{aa:prz:integral-gaussian-three}
i~\ref{aa:prz:integral-gaussian-five}, widzimy trzy
możliwe zachowania liczby pierwszej w tym samym
rozszerzeniu skończonym modułowo:
$(2)$ daje pojedynczy ideał z rozgałęzieniem,
$(3)$ pozostaje ideałem pierwszym, natomiast
$(5)$ ma dwa różne ideały nad sobą.
Nie utożsamiamy przy tym ideału $(2)\subseteq B$
z pierwszym ideałem $\mathfrak q$ leżącym nad $(2)$.
\end{przyklad}

\begin{przyklad}[Bez całkowitości nie ma wznoszenia]
\label{aa:prz:integral-localization-failure}
Włączenie $\mathbb Z\subseteq\mathbb Z[1/2]$ nie jest
całkowite. Rzeczywiście, gdyby $1/2$ spełniał
moniczne równanie o współczynnikach całkowitych,
po przemnożeniu przez odpowiednią potęgę $2$
uzyskalibyśmy sprzeczność między liczbą nieparzystą
i zerem, jak w Przykładzie~\ref{aa:prz:integral-sqrt}.
W pierścieniu $\mathbb Z[1/2]$ liczba $2$ jest jednostką,
gdyż ma odwrotność $1/2$. Żaden ideał właściwy tego
pierścienia nie zawiera zatem $2$. Gdyby istniał
ideał pierwszy $\mathfrak q\subseteq\mathbb Z[1/2]$
z kontrakcją $(2)\subset\mathbb Z$, to $2\in\mathfrak q$,
sprzeczność. Bez założenia o całkowitości nawet twierdzenie
o istnieniu ideału leżącego nad ustalonym może zawieść.
\end{przyklad}

\section{Wymiar algebr skończenie generowanych}
\begin{twierdzenie}[Wymiar pierścienia wielomianów]\label{aa:tw:polynomial-krull-dimension}
Dla każdego ciała $K$ i każdej liczby całkowitej $n\ge0$
zachodzi $\dim K[X_1,\ldots,X_n]=n$.
\end{twierdzenie}
\begin{proof}
Dla $n=0$ otrzymujemy ciało $K$, więc wymiar wynosi $0$.
Dla $n=1$ wynik wynika z
Twierdzenia~\ref{aa:tw:krull-pid}, gdyż $K[X_1]$
jest dziedziną ideałów głównych. Załóżmy, że $n\ge2$
i znamy wynik dla $n-1$ zmiennych.
Łańcuch
$(0)\subsetneq(X_1)\subsetneq(X_1,X_2)\subsetneq\cdots
\subsetneq(X_1,\ldots,X_n)$
składa się z ideałów pierwszych: odpowiednie ilorazy
są pierścieniami wielomianów, czyli dziedzinami.
Wynika stąd $\dim K[X_1,\ldots,X_n]\ge n$.

Dla nierówności odwrotnej weźmy dowolny łańcuch
$\mathfrak p_0\subsetneq\mathfrak p_1\subsetneq\cdots
\subsetneq\mathfrak p_r$. Jeśli $r=0$, nie ma czego dowodzić.
Ideał $\mathfrak p_1$ jest niezerowy, gdyż
$\mathfrak p_0\supseteq(0)$ i inkluzja jest ścisła.
Wybierzmy $0\ne f\in\mathfrak p_1$. Wielomian $f$
nie może być niezerową stałą: stałe niezerowe są jednostkami
w $K[X_1,\ldots,X_n]$, a właściwy ideał pierwszy
$\mathfrak p_1$ nie zawiera jednostki. Dlatego $f$
jest niestały. Stosujemy trójkątną zmianę zmiennych z
Lematu~\ref{aa:lem:noether-monic}. Jest ona automorfizmem:
jeśli $Y_i=X_i-X_n^{m_i}$ dla $i<n$ i $Y_n=X_n$, to
$X_i=Y_i+Y_n^{m_i}$ i $X_n=Y_n$ wyznaczają odwrotną
substytucję. Po tej zmianie niezerowy wielomian $f$
przyjmuje postać moniczną w $Y_n$ (po pomnożeniu przez
niezerową stałą z $K$, co nie zmienia ideału $(f)$). Wobec tego pierścień
$K[X_1,\ldots,X_n]/(f)$ jest skończenie generowanym
modułem nad pierścieniem $K[Y_1,\ldots,Y_{n-1}]$:
klasy $1,Y_n,\ldots,Y_n^{d-1}$ generują go, gdzie $d$
to stopień otrzymanego wielomianu względem $Y_n$.
Włączenie tego pierścienia wielomianów jest różnowartościowe.
Istotnie, jeśli $0\ne h\in K[Y_1,\ldots,Y_{n-1}]$
należałby do ideału $(f)$, to $h=fq$ dla pewnego
$q\in K[Y_1,\ldots,Y_n]$. Jeżeli $q\ne0$, moniczność
$f$ w $Y_n$ daje $\deg_{Y_n}(fq)=d+\deg_{Y_n}q
\ge d>0$, wbrew $\deg_{Y_n}h=0$; jeśli $q=0$,
to $h=0$, także sprzeczność.
Z Lematu~\ref{aa:lem:integrality-determinant} wynika
całkowitość rozszerzenia; Wniosek~\ref{aa:cor:integral-dimension}
i założenie indukcyjne dają
$\dim(K[X_1,\ldots,X_n]/(f))=n-1$.
Łańcuch $\mathfrak p_1\subsetneq\cdots\subsetneq\mathfrak p_r$
po przejściu do ilorazu przez $(f)$ ma długość $r-1$.
Stąd $r-1\le n-1$, czyli $r\le n$.
Ponieważ łańcuch był dowolny, dowód jest zakończony.
\end{proof}

\begin{wniosek}[Wymiar i stopień transcendencji]\label{aa:cor:finite-type-dimension}
Jeżeli $A$ jest niezerową skończenie generowaną przemienną
$K$-algebrą z jedynką, to istnieje $r\ge0$ takie, że
$\dim A=r$ i $A$ jest skończonym modułem nad
podalgebrą $K[y_1,\ldots,y_r]$ izomorficzną z
pierścieniem wielomianów. Jeżeli dodatkowo $A$ jest dziedziną
całkowitości, to
$\dim A=\operatorname{trdeg}_K\operatorname{Frac}(A)$.
\end{wniosek}
\begin{proof}
Lemat normalizacyjny Noether
(Twierdzenie~\ref{aa:tw:noether-normalization}) dostarcza
algebraicznie niezależnych $y_1,\ldots,y_r\in A$,
nad którymi $A$ jest skończenie generowanym modułem.
Włączenie $K[y_1,\ldots,y_r]\subseteq A$ jest więc
całkowite na mocy Lematu~\ref{aa:lem:integrality-determinant}:
każdy element $a\in A$ mnoży skończony moduł $A$ w siebie,
a $1\in A$. Wniosek~\ref{aa:cor:integral-dimension} daje
$\dim A=\dim K[y_1,\ldots,y_r]=r$.
Jeśli $A$ jest dziedziną, Wniosek~\ref{aa:wn:noether-dimension}
identyfikuje $r$ ze stopniem transcendencji jej ciała ułamków.
\end{proof}

\begin{przyklad}[Krzywa o osobliwości ostrzowej]\label{aa:prz:integral-cusp}
Niech $A=K[X,Y]/(Y^2-X^3)$, gdzie $K$ jest dowolnym ciałem.
Oznaczmy klasy zmiennych przez $x,y$. Z równania
$y^2=x^3$ otrzymujemy $A=K[x]+K[x]y$:
każdą wyższą potęgę $y$ redukujemy, zastępując $y^2$ przez $x^3$.
Element $x$ jest transcendentalny nad $K$, gdyż żaden
niezerowy wielomian w $K[X]$ nie należy do ideału
$(Y^2-X^3)$ (jego generator ma dodatni stopień względem $Y$
i jest moniczny). Wobec tego $K[x]\cong K[T]$.
Całkowitość $A/K[x]$ daje $\dim A=1$, chociaż $A$ ma dwa
generatory jako $K$-algebra. Warto porównać ten przykład z
Rozdziałem~\ref{aa:ch:noether-normalization}.
\end{przyklad}



\begin{przyklad}[Normalizacja krzywej ostrzowej — pełne obliczenie]
\label{aa:prz:integral-cusp-normalization}
Niech $K$ będzie dowolnym ciałem, $t$ zmienną,
$A=K[t^2,t^3]\subseteq B=K[t]$. Element $t\in B$
spełnia moniczne równanie $U^2-t^2=0$ o współczynnikach
w $A$. Co więcej, każdy wielomian $h(t)\in K[t]$
rozkłada się na sumę swoich składników o parzystych
i nieparzystych wykładnikach:
$h(t)=a(t^2)+t b(t^2)$ dla pewnych $a,b\in K[U]$.
Wobec $K[t^2]\subseteq A$ oznacza to
$B=A\cdot1+A\cdot t$; zatem $B$ jest skończenie
generowanym $A$-modułem, a rozszerzenie jest całkowite.

Sprawdźmy precyzyjnie związek z równaniem ostrza.
Homomorfizm $K[X,Y]\to K[t]$ określony wzorami
$X\mapsto t^2$, $Y\mapsto t^3$ ma obraz $A$.
Wielomian $Y^2-X^3$ leży w jego jądrze. Dzięki temu,
że $Y^2-X^3$ jest moniczny względem $Y$, każdy
wielomian modulo $(Y^2-X^3)$ można przedstawić jako
$a(X)+Yb(X)$ dla $a,b\in K[X]$. Jeśli jego obraz
$a(t^2)+t^3b(t^2)$ jest zerowy, to oba składniki
muszą znikać oddzielnie: pierwszy zawiera wyłącznie
parzyste potęgi $t$, a drugi wyłącznie nieparzyste
potęgi $t$ (od stopnia $3$). Stąd $a=b=0$.
Jądro homomorfizmu jest zatem \emph{dokładnie}
$(Y^2-X^3)$, a więc
\[
 K[X,Y]/(Y^2-X^3)\cong K[t^2,t^3]=A.
\]
W szczególności $A$ jest dziedziną całkowitości.
Warto zauważyć, że $A$ nie jest izomorficzny z
pierścieniem wielomianów $K[U]$. Istotnie, ideał
$\mathfrak m=(t^2,t^3)$ pierścienia $A$ nie jest główny.
Dla niezerowego $h(t)\in K[t]$ niech
$\nu(h)$ oznacza najmniejszy wykładnik potęgi $t$
występującej w $h$ z niezerowym współczynnikiem.
Wtedy $\nu(hk)=\nu(h)+\nu(k)$, a niezerowe elementy
$A$ mają wartość $\nu$ równą $0$ lub co najmniej $2$.
Gdyby $\mathfrak m=(h)$ dla pewnego $0\ne h\in A$,
z $h\in\mathfrak m$ wynikałoby $\nu(h)\ge2$,
a z $t^2=ha$ dla pewnego $a\in A$ mielibyśmy
$\nu(h)=2$ i $\nu(a)=0$.
Jednocześnie $t^3=hb$ dla pewnego $b\in A$ wymuszałoby
$\nu(b)=1$, co jest niemożliwe.
Ponieważ każdy ideał $K[U]$ jest główny,
$A$ nie może być izomorficzny z $K[U]$.

W pierścieniu $B$ mamy ścisły łańcuch ideałów pierwszych
$(0)\subsetneq(t)$. Kontrakcja $(t)\cap A$ jest
ideałem $(t^2,t^3)$ pierścienia $A$: wielomian w $t^2,t^3$
ma wyraz wolny zerowy dokładnie wtedy, gdy należy do
$(t^2,t^3)$: każda dodatnia potęga $t$ występująca w $A$
ma wykładnik parzysty co najmniej $2$ lub nieparzysty
co najmniej $3$, a zatem jest podzielna w $A$ przez
$t^2$ albo $t^3$. Ten ideał jest maksymalny, gdyż
$A/(t^2,t^3)\cong K$. Widzimy więc konkretny łańcuch
$(0)\subsetneq(t^2,t^3)$ w $A$. Z twierdzenia o
niezmienniczości wymiaru w rozszerzeniu całkowitym wynika
$\dim A=\dim K[t]=1$; pierścień $A$ nie jest wprawdzie
pierścieniem wielomianów w jednej zmiennej, ale ma ten sam
wymiar Krulla co $B$.
\end{przyklad}
\begin{przyklad}[Nie należy mylić algebraiczności z całkowitością]
\label{aa:prz:integral-nonexample}
Rozszerzenie ciał $\mathbb Q/\mathbb Q$ jest algebraiczne,
lecz $1/2\in\mathbb Q$ nie jest całkowity nad $\mathbb Z$.
Dlatego z faktu, że $\operatorname{Frac}(B)$ jest rozszerzeniem
algebraicznym $\operatorname{Frac}(A)$, nie wynika całkowitość
$B/A$. Na przykład $\mathbb Z\subset\mathbb Z[1/2]$
nie jest rozszerzeniem całkowitym, mimo że oba pierścienie mają
to samo ciało ułamków $\mathbb Q$.
\end{przyklad}

% AA-41 | Audyt I: sprawdzenie matematyczne, językowe i przykład w dodatniej charakterystyce.
\chapter{Reprezentacje grup skończonych i twierdzenie Maschkego}
\markboth{Reprezentacje i twierdzenie Maschkego}{Reprezentacje i twierdzenie Maschkego}
\label{aa:ch:representations}
\section{Od działań grup do operatorów liniowych}
Działanie grupy na zbiorze zachowuje jedynie strukturę zbiorową.
Gdy elementy zbioru tworzą przestrzeń wektorową, możemy wymagać,
aby każdy element grupy działał liniowo. Wówczas pytania o niezmiennicze
podprzestrzenie stają się pytaniami o rozkład reprezentacji.
\begin{definicja}[Reprezentacja i podreprezentacja]\label{aa:def:representation}
Niech $G$ będzie grupą, $K$ ciałem, a $V$ przestrzenią wektorową nad $K$.
Reprezentacją $G$ na $V$ nazywamy homomorfizm grup
$\rho:G\to\operatorname{GL}_K(V)$, gdzie $\operatorname{GL}_K(V)$
oznacza grupę wszystkich odwracalnych $K$-liniowych operatorów na $V$.
Podprzestrzeń $W\subseteq V$ jest $G$-niezmiennicza, jeżeli
$\rho(g)w\in W$ dla wszystkich $g\in G$ i $w\in W$.
Reprezentacja jest nieprzywiedlna, gdy $V\ne0$ i jedynymi
podprzestrzeniami niezmienniczymi są $0$ oraz $V$.
\end{definicja}
Warunek homomorfizmu oznacza jawnie $\rho(gh)=\rho(g)\rho(h)$
oraz $\rho(e)=\operatorname{id}_V$. Ponieważ $g^{-1}\in G$,
$\rho(g)^{-1}=\rho(g^{-1})$; zatem $\rho(g)W\subseteq W$
implikuje również $\rho(g)W=W$.
\begin{przyklad}[Reprezentacja znaku]\label{aa:prz:sign-rep}
Niech $G=S_3$, $K=\mathbb Q$, $V=\mathbb Q$.
Definiujemy $\rho(\sigma)(v)=\operatorname{sgn}(\sigma)v$.
Ponieważ znak iloczynu permutacji jest iloczynem znaków,
$\rho$ jest homomorfizmem. Każda podprzestrzeń jednowymiarowego
$V$ jest zerowa albo równa $V$, zatem reprezentacja jest
nieprzywiedlna. Nie jest trywialna: transpozycja działa jako $-1$.
\end{przyklad}
\begin{przyklad}[Reprezentacja permutacyjna i rozkład]\label{aa:prz:permutation-rep}
Niech $S_3$ permutuje bazę $e_1,e_2,e_3$ przestrzeni $\mathbb Q^3$.
Wówczas $L=\mathbb Q(e_1+e_2+e_3)$ jest niezmiennicza.
Także $W=\{(x_1,x_2,x_3):x_1+x_2+x_3=0\}$ jest niezmiennicza,
gdyż permutacja nie zmienia sumy współrzędnych.
Jest ona dwuwymiarowa: wektory $a=e_1-e_2$ i
$b=e_2-e_3$ leżą w $W$ i tworzą jej bazę, ponieważ
$(u,v,w)=(u)a+(-w)b$ dla każdego $(u,v,w)\in W$.
Każdy $v=(x_1,x_2,x_3)$ zapisujemy jako
$v=\frac{x_1+x_2+x_3}{3}(1,1,1)+w$, gdzie $w\in W$.
Ponadto $L\cap W=0$: jeżeli $t(1,1,1)\in W$, to
$3t=0$ w $\mathbb Q$, a więc $t=0$.
Stąd $\mathbb Q^3=L\oplus W$.
Pokażemy też, że podreprezentacja na $W$ jest
nieprzywiedlna nad $\mathbb Q$. Trójcykl
$g=(1\;2\;3)$ przekształca $a$ w $b$, a $b$ w
$-a-b$. Macierz $\rho(g)|_W$ w bazie $(a,b)$ ma więc
postać
\[
\begin{pmatrix}0&-1\\1&-1\end{pmatrix},
\qquad \det(TI-\rho(g)|_W)=T^2+T+1.
\]
Ten wielomian nie ma pierwiastka wymiernego: jego
wyróżnik jest równy $-3$, a kwadrat liczby wymiernej
nie może być ujemny. Gdyby $W$ zawierała niezerową
właściwą podprzestrzeń $S_3$-niezmienniczą, byłaby to
prosta wymierna niezmiennicza również względem $g$.
Operator na tej prostej działałby przez skalar z
$\mathbb Q$, będący wartością własną powyższej macierzy;
takiego skalara nie ma. Zatem $W$ jest nieprzywiedlna.
Widzimy konkretnie sumę reprezentacji wymiarów $1$ i $2$,
a nie tylko dwa abstrakcyjne składniki.
\end{przyklad}
\section{Uśrednianie projekcji}
Jeżeli $W$ jest niezmiennicza, dowolne liniowe dopełnienie
$U$ nie musi być niezmiennicze. Dla grupy skończonej
można poprawić rzut na $W$, uśredniając jego sprzężenia.
\begin{twierdzenie}[Maschke]\label{aa:tw:maschke}
Niech $G$ będzie grupą skończoną, $K$ ciałem takim, że
$\operatorname{char}K$ nie dzieli $|G|$, a $\rho:G\to
\operatorname{GL}_K(V)$ reprezentacją na skończenie wymiarowej
przestrzeni $V$. Dla każdej podprzestrzeni $G$-niezmienniczej
$W\subseteq V$ istnieje podprzestrzeń $G$-niezmiennicza $U$
taka, że $V=W\oplus U$. W konsekwencji każda taka reprezentacja
jest sumą prostą reprezentacji nieprzywiedlnych.
\end{twierdzenie}
\begin{proof}
Wybierzmy bazę $W$ i uzupełnijmy ją do bazy $V$.
Rzut liniowy $p:V\to W$ określony przez $p(w)=w$ dla $w\in W$
oraz $p(u)=0$ na pozostałych wektorach bazy istnieje i spełnia
$p|_W=\operatorname{id}_W$.
Ponieważ $|G|\cdot1_K\ne0$, element $|G|\cdot1_K$ ma
odwrotność w $K$. Określmy
\[
P=\frac1{|G|}\sum_{g\in G}\rho(g)p\rho(g)^{-1}.
\]
Każdy składnik ma obraz w $W$, ponieważ $p(V)\subseteq W$
i $\rho(g)W=W$. Stąd $P(V)\subseteq W$. Ponadto suma jest
$K$-liniowa, ponieważ każdy operator
$\rho(g)p\rho(g)^{-1}$ jest $K$-liniowy.
Jeśli $w\in W$, to $\rho(g)^{-1}w\in W$, więc
$\rho(g)p\rho(g)^{-1}w=w$. Sumując $|G|$ równych składników,
dostajemy $P(w)=w$. W szczególności $P^2=P$.
Dla $h\in G$ przekształcamy sumę przez podstawienie $k=hg$:
\[
\rho(h)P\rho(h)^{-1}=
\frac1{|G|}\sum_{g\in G}\rho(hg)p\rho(hg)^{-1}=P.
\]
Zatem $P\rho(h)=\rho(h)P$. Jądro $U=\ker P$ jest niezmiennicze:
$P(\rho(h)u)=\rho(h)P(u)=0$ dla $u\in U$.
Dla każdego $v\in V$ mamy $v=P(v)+(v-P(v))$, przy czym
$P(v)\in W$ oraz $P(v-P(v))=P(v)-P^2(v)=0$.
Jeżeli $w\in W\cap U$, to $w=P(w)=0$.
To dowodzi $V=W\oplus U$.

Aby uzyskać ostatnią tezę, prowadzimy indukcję względem
$\dim_KV$. Wymiar $0$ daje pustą sumę. Jeżeli $V\ne0$
i jest nieprzywiedlna, stanowi żądaną sumę. W przeciwnym
razie wybieramy niezerową właściwą podprzestrzeń niezmienniczą
$W$. Jej dopełnienie $U$ jest niezmiennicze na mocy pierwszej
części dowodu. Obie przestrzenie mają wymiar mniejszy od
$\dim V$, więc stosujemy do nich założenie indukcyjne.
\end{proof}
\begin{przyklad}[Dlaczego założenie o charakterystyce jest konieczne]
\label{aa:prz:maschke-counter}
Niech $K=\mathbb F_2$, $G=C_2=\{e,s\}$ i $V=K^2$.
Definiujemy $\rho(s)=\begin{pmatrix}1&1\\0&1\end{pmatrix}$.
W charakterystyce $2$ kwadrat tej macierzy jest jednostkowy,
więc otrzymujemy reprezentację $C_2$. Prosta $W=Ke_1$
jest niezmiennicza. Gdyby istniało niezmiennicze dopełnienie,
musiałoby być prostą $K(e_2+ae_1)$ dla pewnego $a\in K$.
Jednak $\rho(s)(e_2+ae_1)=e_2+(a+1)e_1$ nie należy do tej
samej prostej: współczynnik przy $e_2$ wymusza skalar $1$,
a wtedy $a+1\ne a$. Dopełnienie nie istnieje.
\end{przyklad}
\begin{przyklad}[Twierdzenie Maschkego w dodatniej charakterystyce]
\label{aa:prz:maschke-f2-c3}
Niech $K=\mathbb F_2$, $G=C_3=\langle g\rangle$ i niech
$g$ działa na $V=K^3$ przez $g e_1=e_2$, $g e_2=e_3$,
$g e_3=e_1$. Ponieważ $|G|=3$, liczba $3\cdot1_K=1_K$
jest odwracalna, mimo że charakterystyka $K$ jest dodatnia.
Podprzestrzenie
\[
L=K(1,1,1),\qquad
W=\{(a,b,c)\in K^3:a+b+c=0\}
\]
są $G$-niezmiennicze. Mają wymiary odpowiednio $1$ i $2$.
Jeżeli $t(1,1,1)\in W$, to $3t=t=0$ w $K$,
więc $L\cap W=0$ i $V=L\oplus W$.
Zdefiniujmy średnią operatorów działania
$R_G=\frac1{|G|}\sum_{h\in G}\rho(h)$. Ponieważ
$1/3=1$ w $\mathbb F_2$, mamy
\[
R_G(a,b,c)=(a,b,c)+(c,a,b)+(b,c,a)
=(a+b+c)(1,1,1).
\]
Jest to rzut na $L$ wzdłuż $W$.

Pokażmy ponadto, że $W$ jest nieprzywiedlna. Wektory
$a=e_1+e_2$ i $b=e_2+e_3$ tworzą bazę $W$.
Ponieważ $g a=b$ i $g b=e_3+e_1=a+b$, macierz
$g|_W$ w bazie $(a,b)$ wynosi
\[
\begin{pmatrix}0&1\\1&1\end{pmatrix},
\qquad \det(TI-g|_W)=T^2+T+1.
\]
Wielomian $T^2+T+1$ nie ma pierwiastka w $\mathbb F_2$,
bo przy $T=0$ i $T=1$ przyjmuje wartość $1$.
Gdyby $W$ zawierała niezerową właściwą podprzestrzeń
$G$-niezmienniczą, byłaby nią prosta stabilna względem $g$,
a więc istniałaby wartość własna w $K$; sprzeczność.
Otrzymaliśmy zatem rozkład na reprezentacje nieprzywiedlne
w charakterystyce $2$, która nie dzieli rzędu grupy $C_3$.
\end{przyklad}

\begin{stwierdzenie}[Dwie projekcje związane z inwolucją]
\label{aa:prop:involution-projectors}
Niech $K$ będzie ciałem, w którym $2\cdot1_K\ne0$, niech
$V$ będzie skończenie wymiarową przestrzenią wektorową nad $K$,
a $T\in\operatorname{GL}_K(V)$ operatorem spełniającym $T^2=I$.
Wtedy operatory $P_+=(I+T)/2$ i $P_-=(I-T)/2$ są rzutami,
a przestrzeń $V$ jest sumą prostą przestrzeni
$V_+=\ker(T-I)$ oraz $V_-=\ker(T+I)$.
Dokładniej, $\operatorname{im}P_+=V_+$,
$\operatorname{im}P_-=V_-$ i dla każdego $v\in V$ zachodzi
$v=P_+v+P_-v$.
\end{stwierdzenie}
\begin{proof}
Rachunki przeprowadzamy w algebrze operatorów liniowych
na $V$. Z równości $T^2=I$ otrzymujemy
$(I\pm T)^2=I\pm2T+T^2=2(I\pm T)$.
Ponieważ $2$ ma odwrotność w $K$, wynika stąd
$P_+^2=P_+$ oraz $P_-^2=P_-$.
Ponadto $P_++P_-=I$ i
$P_+P_-=(I-T^2)/4=0$.
Dla każdego $v$ mamy $T(P_+v)=P_+v$ i
$T(P_-v)=-P_-v$, gdyż $T$ komutuje z wielomianami od $T$.
Stąd $\operatorname{im}P_+\subseteq V_+$ i
$\operatorname{im}P_-\subseteq V_-$.
Jeżeli $w\in V_+$, to $P_+w=(w+w)/2=w$;
analogicznie $P_-w=w$ dla $w\in V_-$.
Otrzymujemy równości obrazów.
Z $P_++P_-=I$ wynika $V=V_++V_-$.
Gdyby $w\in V_+\cap V_-$, to zarazem $Tw=w$ i
$Tw=-w$, a zatem $2w=0$ i $w=0$.
Suma jest więc prosta.
\end{proof}
\begin{przyklad}[Jawna dekompozycja działania $C_2$]
\label{aa:prz:involution-projectors}
Dla macierzy $T=\begin{pmatrix}0&1\\1&0\end{pmatrix}$
na $\mathbb Q^2$ otrzymujemy
$P_+=\tfrac12\begin{pmatrix}1&1\\1&1\end{pmatrix}$ i
$P_-=\tfrac12\begin{pmatrix}1&-1\\-1&1\end{pmatrix}$.
Pierwszy operator wysyła $(x,y)$ na
$((x+y)/2,(x+y)/2)$, a drugi na
$((x-y)/2,(y-x)/2)$.
Ich obrazy to odpowiednio $\mathbb Q(1,1)$ i
$\mathbb Q(1,-1)$; otrzymujemy w sposób wykonawczy
rozkład z Przykładu~\ref{aa:prz:maschke-explicit}.
W charakterystyce dwa operatory $P_\pm$ nie są
zdefiniowane: nie istnieje element $1/2$ w ciele.
\end{przyklad}

\section{Algebra grupowa}
\begin{definicja}[Algebra grupowa]\label{aa:def:group-algebra-rep}
Dla grupy skończonej $G$ i ciała $K$ algebra $K[G]$ jest
$K$-przestrzenią wektorową z bazą formalnych symboli $[g]$,
$g\in G$, z mnożeniem określonym przez $[g][h]=[gh]$
i przedłużonym dwuliniowo. Elementem jednostkowym jest $[e]$.
\end{definicja}
\begin{stwierdzenie}[Reprezentacje jako moduły]\label{aa:prop:rep-modules}
Dla ustalonych $G,K$ struktury reprezentacji $G$ na
$K$-przestrzeni $V$ są we wzajemnie jednoznacznej odpowiedniości
ze strukturami lewego $K[G]$-modułu na $V$, których działanie
$K$ jest zwykłym mnożeniem przez skalary.
\end{stwierdzenie}
\begin{proof}
Dla reprezentacji $\rho$ określamy
$(\sum_g a_g[g])v=\sum_g a_g\rho(g)v$.
Jest to dobrze określone, ponieważ każdy element $K[G]$
ma jednoznaczne współczynniki. Rozdzielność wynika z liniowości,
a łączność z równości $\rho(g)\rho(h)=\rho(gh)$;
$[e]$ działa identycznościowo. Odwrotnie, z działania
$K[G]$ bierzemy $\rho(g)v=[g]v$. Ponieważ $[g][g^{-1}]=[e]$,
operator $\rho(g)$ jest odwracalny, a reguła mnożenia daje
homomorfizm grup. Obie konstrukcje są wzajemnie odwrotne.
\end{proof}


\begin{przyklad}[Co dzieje się z algebrą grupową w charakterystyce grupy?]
\label{aa:prz:group-algebra-modular}
Niech $G=C_2=\{e,s\}$ i $K=\mathbb F_2$.
W algebrze grupowej $K[G]$ rozpatrzmy
$z=[s]-[e]$. Mamy
\[
z^2=[s]^2-2[s][e]+[e]^2=[e]+[e]=0,
\]
gdyż $s^2=e$ i $2=0$ w $\mathbb F_2$.
Jednocześnie $z\ne0$, ponieważ $[s]$ i $[e]$
są odrębnymi wektorami bazy $K[G]$.
Niezerowy element nilpotentny nie świadczy sam w sobie
o braku półprostoty dowolnego pierścienia nieprzemiennego.
W naszym konkretnym przypadku możemy jednak wykazać
brak niezmienniczego dopełnienia bez takiego wnioskowania.
Lewy moduł regularny $K[G]$ jest dwuwymiarowy, a mnożenie
przez $[s]$ zamienia wektory bazy $[e]$ i $[s]$.
Prosta $W=K([e]+[s])$ jest niezmiennicza, ponieważ
$[s]([e]+[s])=[s]+[e]$.
Każda inna prosta w dwuwymiarowej przestrzeni nad
$\mathbb F_2$ jest równa $K[e]$ albo $K[s]$:
w istocie przestrzeń ma dokładnie trzy wektory niezerowe
$[e]$, $[s]$ i $[e]+[s]$. Żadna z tych dwóch pozostałych
prostych nie jest niezmiennicza względem zamiany bazy.
Stąd $W$ nie ma dopełnienia $G$-niezmienniczego.
Ten rachunek pokazuje mechanizm powstania trudności
w charakterystyce dzielącej rząd grupy i daje
wersję kontrprzykładu~\ref{aa:prz:maschke-counter}
bez użycia macierzy Jordana.
\end{przyklad}

\begin{twierdzenie}[Ostrość warunku twierdzenia Maschkego]
\label{aa:tw:maschke-necessity}
Niech $G$ będzie grupą skończoną i $K$ dowolnym ciałem.
Następujące warunki są równoważne:
\begin{enumerate}
\item $|G|\cdot1_K\ne0$;
\item dla każdej skończenie wymiarowej reprezentacji $G$
nad $K$ i każdej jej podprzestrzeni niezmienniczej
istnieje niezmiennicze dopełnienie;
\item każda skończenie wymiarowa reprezentacja $G$
nad $K$ jest sumą prostą reprezentacji nieprzywiedlnych.
\end{enumerate}
\end{twierdzenie}
\begin{proof}
Implikację $1\Rightarrow2$ daje
Twierdzenie~\ref{aa:tw:maschke}, a $2\Rightarrow3$
jego dowód indukcyjny.
Dla uzasadnienia $3\Rightarrow2$ zauważmy, że każdy
iloraz skończenie wymiarowej reprezentacji, będącej
sumą prostą reprezentacji nieprzywiedlnych, jest ponownie
sumą takich reprezentacji. Istotnie, jeśli
$V=S_1\oplus\cdots\oplus S_r$ z $S_i$ nieprzywiedlnymi,
a $q:V\to V/W$ jest homomorfizmem ilorazowym, to
obrazy $q(S_i)$ są zerowe albo nieprzywiedlne.
Z nich wybieramy maksymalną rodzinę, której suma jest
prosta. Jeżeli obraz $q(S_i)$ nie został wybrany i nie
można go do tej rodziny dołączyć, przecina on niezerowo
sumę wybranych obrazów. Nieprzywiedlność $q(S_i)$
sprawia, że cały ten obraz należy już do owej sumy.
Maksymalna suma zawiera więc wszystkie obrazy $q(S_i)$,
a zatem jest równa $V/W$, gdyż te obrazy generują iloraz.
Każdy niezerowy $q(S_i)$ jest obrazem
$S_i$ przez izomorfizm (jądro ograniczenia jest
podreprezentacją nieprzywiedlnej reprezentacji $S_i$).
Wybór odpowiednich $S_i$ daje więc podprzestrzeń
$U\subseteq V$ taką, że $q|_U:U\to V/W$ jest
izomorfizmem reprezentacji. Wtedy $V=W\oplus U$.


Dla implikacji $3\Rightarrow1$ wygodniejsza jest
kontrapozycja. Załóżmy $|G|\cdot1_K=0$.
Rozważmy reprezentację regularną na $V=K[G]$
z działaniem $\rho(h)([g])=[hg]$.
Wektor $u=\sum_{g\in G}[g]$ jest niezerowy,
gdyż w bazie $([g])_{g\in G}$ ma wszystkie współczynniki
równe $1_K$. Dla każdego $h\in G$ zachodzi
$\rho(h)u=u$, bo mnożenie z lewej permutuje elementy $G$.
Zatem $W=Ku$ jest jednowymiarową podreprezentacją.

Gdyby $W$ miała niezmiennicze dopełnienie $U$,
rzut $P:V=W\oplus U\to W$ wzdłuż $U$ byłby
$G$-równoważny: dla $w\in W$ i $v\in U$
mamy $P(\rho(h)(w+v))=\rho(h)w=\rho(h)P(w+v)$.
Ponieważ $[g]=\rho(g)[e]$ i działanie $G$ na $W$
jest trywialne, istnieje $a\in K$ takie, że
$P([g])=\rho(g)P([e])=au$ dla wszystkich $g\in G$.
Wobec tego
$P(u)=\sum_gP([g])=(|G|\cdot1_K)au=0$.
Ale $P|_W=\operatorname{id}_W$, więc $P(u)=u\ne0$:
sprzeczność.

Gdyby zachodził warunek 3, to reprezentacja regularna
byłaby sumą nieprzywiedlnych reprezentacji. W takim
rozkładzie każda podreprezentacja ma niezmiennicze
dopełnienie, co wynika ze struktury modułu półprostego.
Ponieważ dla konkretnej podreprezentacji $W$ wykazaliśmy
brak dopełnienia, warunek 3 również nie może zachodzić,
na mocy wykazanej już implikacji $3\Rightarrow2$.
Stąd $3\Rightarrow1$, co kończy dowód.
\end{proof}

\section{Dodatkowa kontrola założeń i przykład rozkładu}
\begin{uwaga}[Znaczenie charakterystyki]\label{aa:rem:maschke-char}
Zapis $\operatorname{char}K\nmid |G|$ w Twierdzeniu~\ref{aa:tw:maschke}
oznacza, że $|G|\cdot1_K\ne0$. Dla charakterystyki zero
warunek jest spełniony automatycznie; dla charakterystyki
pierwszej $p$ oznacza $p\nmid|G|$. Bez tego warunku
nie wolno dzielić przez $|G|$ w definicji operatora $P$.
Przykład~\ref{aa:prz:maschke-counter} pokazuje, że problem
nie jest wyłącznie techniczny: teza o dopełnieniu może zawieść.
\end{uwaga}
\begin{przyklad}[Rozkład reprezentacji permutacyjnej $C_2$]
\label{aa:prz:maschke-explicit}
Niech $G=C_2=\{e,s\}$, $K=\mathbb Q$ i niech
$\rho(s)(x,y)=(y,x)$ na $V=\mathbb Q^2$.
Wektory $u=(1,1)$ i $v=(1,-1)$ są liniowo niezależne,
a $\rho(s)u=u$, $\rho(s)v=-v$. Dlatego
$V=\mathbb Qu\oplus\mathbb Qv$ jest sumą dwóch
nieprzywiedlnych reprezentacji jednowymiarowych.
Dla dowolnego $(x,y)$ współczynniki tego rozkładu wynoszą
$(x+y)/2$ przy $u$ i $(x-y)/2$ przy $v$.
W charakterystyce $2$ te dwa wektory byłyby równe,
co unaocznia rolę odwracalności liczby $2$.
Możemy w tym przykładzie także \emph{obliczyć samo
uśrednianie} z dowodu Twierdzenia~\ref{aa:tw:maschke}.
Dla $W=\mathbb Q(1,1)$ wybieramy rzut
$p(x,y)=(x,x)$. Nie jest on przemienny z zamianą
współrzędnych: $\rho(s)p\rho(s)^{-1}(x,y)=(y,y)$.
Uśredniony rzut ma więc postać
\[
P(x,y)=\tfrac12\big((x,x)+(y,y)\big)
=\left(\tfrac{x+y}{2},\tfrac{x+y}{2}\right).
\]
Jego obrazem jest $W$, a jądrem jest
$\mathbb Q(1,-1)$. Otrzymaliśmy dokładnie dopełnienie
niezmiennicze, o którego istnieniu orzeka twierdzenie.
\end{przyklad}

\section{Średnia niezmiennicza i jej dokładny obraz}
\begin{stwierdzenie}[Operator uśredniania wektorów]\label{aa:prop:reynolds-vectors}
Niech $G$ będzie grupą skończoną, $K$ ciałem takim, że
$|G|\cdot1_K\ne0$, oraz $\rho:G\to\operatorname{GL}_K(V)$
reprezentacją na przestrzeni wektorowej $V$ nad $K$ (niekoniecznie
skończenie wymiarowej). Niech
$V^G=\{v\in V:\rho(g)v=v\ \text{dla każdego }g\in G\}$.
Wówczas operator
\[
R_G(v)=\frac1{|G|}\sum_{g\in G}\rho(g)v
\]
jest $K$-liniowym rzutem na $V^G$. W szczególności
$V=V^G\oplus\ker R_G$, a oba składniki są $G$-niezmiennicze.
\end{stwierdzenie}
\begin{proof}
Zbiór $V^G$ jest podprzestrzenią: suma dwóch wektorów
stałych oraz iloczyn wektora stałego przez skalar są stałe,
gdyż każdy $\rho(g)$ jest liniowy. Suma definiująca $R_G$
jest skończona, więc określa operator liniowy. Dla $h\in G$
i $v\in V$ mamy
\[
\rho(h)R_G(v)=\frac1{|G|}\sum_{g\in G}\rho(hg)v
=\frac1{|G|}\sum_{k\in G}\rho(k)v=R_G(v),
\]
gdyż mnożenie z lewej przez $h$ jest bijekcją $G$.
Wynika stąd $\operatorname{im}R_G\subseteq V^G$.
Jeżeli $v\in V^G$, wszystkie składniki sumy są równe $v$,
a więc $R_G(v)=v$. Zatem $\operatorname{im}R_G=V^G$
i $R_G^2=R_G$. Każdy $v\in V$ rozkłada się na
$R_G(v)+(v-R_G(v))$, przy czym drugi składnik należy
do jądra, ponieważ $R_G(v-R_G(v))=0$.
Przecięcie obrazu i jądra jest zerowe, ponieważ $R_G$
działa na swoim obrazie identycznościowo.
Wreszcie $R_G\rho(h)=R_G$ (podstawienie $k=gh$),
więc jądro jest niezmiennicze; obraz składa się z wektorów
stałych, a zatem również jest niezmienniczy.
\end{proof}

\begin{wniosek}[Dokładność brania wektorów niezmienniczych]
\label{aa:cor:invariants-exact}
Niech $G$ będzie grupą skończoną, $K$ ciałem takim, że
$|G|\cdot1_K\ne0$, a
\[
0\longrightarrow U\xrightarrow{\,i\,}V
 \xrightarrow{\,q\,}W\longrightarrow0
\]
krótkim ciągiem dokładnym reprezentacji $G$ nad $K$:
$i$ jest różnowartościowym operatorem przeplatającym,
$q$ jest surjektywnym operatorem przeplatającym oraz
$\operatorname{im}i=\ker q$. Wtedy ciąg
\[
0\longrightarrow U^G\xrightarrow{\,i|_{U^G}\,}V^G
 \xrightarrow{\,q|_{V^G}\,}W^G\longrightarrow0
\]
jest dokładny. Nie zakładamy skończonego wymiaru reprezentacji.
\end{wniosek}
\begin{proof}
Ponieważ $i$ i $q$ przeplatają działania grupy, przekształcają
wektory niezmiennicze w wektory niezmiennicze, toteż ograniczenia
w tezie są dobrze określone. Pierwsze jest różnowartościowe,
a dla $v\in V^G$ warunek $q(v)=0$ jest równoważny
$v\in\ker q=\operatorname{im}i$. Wówczas $v=i(u)$ dla
pewnego $u\in U$. Dla dowolnego $g\in G$ mamy
$i(\rho_U(g)u)=\rho_V(g)i(u)=v=i(u)$;
różnowartościowość $i$ daje $\rho_U(g)u=u$.
Stąd $\ker(q|_{V^G})=i(U^G)$.

Pozostaje surjektywność ostatniej strzałki. Niech $w\in W^G$.
Surjektywność $q$ pozwala wybrać $v\in V$ z $q(v)=w$.
Uśredniając według Stwierdzenia~\ref{aa:prop:reynolds-vectors},
otrzymujemy $R_G^V(v)\in V^G$. Ponieważ $q$ przeplata
działania, suma skończona spełnia
\[
q(R_G^V(v))=\frac1{|G|}\sum_{g\in G}q(\rho_V(g)v)
=\frac1{|G|}\sum_{g\in G}\rho_W(g)q(v)
=\frac1{|G|}\sum_{g\in G}w=w.
\]
Zatem każdy $w\in W^G$ ma niezmienniczy przeciwobraz,
co kończy dowód.
\end{proof}
\begin{przyklad}[Dlaczego branie niezmienników zawodzi w złej charakterystyce]
\label{aa:prz:invariants-not-exact-modular}
Niech $K=\mathbb F_2$ i $G=C_2=\{e,s\}$.
Rozważmy reprezentację regularną $K[G]$ z
Przykładu~\ref{aa:prz:group-algebra-modular} oraz homomorfizm
$G$-modułów $\varepsilon:K[G]\to K$ określony wzorem
$\varepsilon(a[e]+b[s])=a+b$, gdzie $G$ działa na $K$
trywialnie. Mapa $\varepsilon$ jest surjektywna, gdyż
$\varepsilon([e])=1$. Wektor $a[e]+b[s]$ jest niezmienniczy
wtedy i tylko wtedy, gdy $a=b$, zatem
$K[G]^G=K([e]+[s])$. Jednak
$\varepsilon([e]+[s])=1+1=0$ w $\mathbb F_2$, więc
$\varepsilon:K[G]^G\to K^G=K$ jest mapą zerową,
a nie surjekcją. Mimo dokładności ciągu
$0\to\ker\varepsilon\to K[G]\xrightarrow{\varepsilon}K\to0$
ciąg jego niezmienników nie jest dokładny na końcu.
\end{przyklad}

\begin{przyklad}[Średnia współrzędnych w reprezentacji permutacyjnej]
\label{aa:prz:reynolds-permutation}
Niech $G=S_3$ działa na $V=\mathbb Q^3$ przez permutowanie
współrzędnych. Dla $v=(a,b,c)$ każda z liczb $a,b,c$
występuje dokładnie dwa razy na każdej pozycji w sześciu
wektorach $\rho(g)v$. Dlatego
\[
R_G(a,b,c)=\left(\frac{a+b+c}{3},\frac{a+b+c}{3},
\frac{a+b+c}{3}\right).
\]
Obrazem jest prosta $\mathbb Q(1,1,1)$, a jądrem płaszczyzna
$a+b+c=0$. Są to dokładnie składniki rozkładu z
Przykładu~\ref{aa:prz:permutation-rep}. Warunek odwracalności
$|G|$ w ciele nie może być pominięty w ogólnym wzorze;
przykład nad $\mathbb F_2$ z
\ref{aa:prz:maschke-counter} pokazuje, że nawet istnienie
niezmienniczego dopełnienia może wtedy zawieść.
\end{przyklad}

% AA-42 | Reprezentacje zespolone grup skonczonych: Schur, charaktery, ortogonalnosc.
\chapter{Lemat Schura, charaktery i relacje ortogonalności}
\markboth{Lemat Schura i charaktery}{Lemat Schura i charaktery}
\label{aa:ch:schur-characters}

\section{Dlaczego same rozkłady reprezentacji nie wystarczą?}
Twierdzenie Maschkego z Rozdziału~\ref{aa:ch:representations}
zapewnia rozkład każdej skończenie wymiarowej reprezentacji
grupy skończonej na składniki nieprzywiedlne, ale samo nie daje
sposobu obliczenia ich krotności ani rozpoznawania izomorfizmów.
Narzędziami służącymi do tego celu są operatory przeplatające,
lemat Schura i charaktery. Przy badaniu relacji ortogonalności
będziemy pracować nad ciałem \(\mathbb C\); warunek ten
jest istotny i nie może zostać bez komentarza przeniesiony
na dowolne ciało.

\begin{definicja}[Operator przeplatający i równoważność reprezentacji]
\label{aa:def:intertwiner}
Niech \(G\) będzie grupą, \(K\) ciałem, a
\(\rho:G\to\operatorname{GL}_K(V)\) oraz
\(\sigma:G\to\operatorname{GL}_K(W)\) reprezentacjami
na przestrzeniach wektorowych \(V,W\) nad \(K\).
Operatorem przeplatającym te reprezentacje nazywamy
odwzorowanie \(K\)-liniowe \(T:V\to W\), dla którego
\(T\rho(g)=\sigma(g)T\) dla każdego \(g\in G\).
Przestrzeń wszystkich takich operatorów oznaczamy przez
\(\operatorname{Hom}_G(V,W)\).
Reprezentacje są równoważne, gdy istnieje
odwracalny operator z \(\operatorname{Hom}_G(V,W)\).
\end{definicja}

\begin{lemat}[Pierwszy lemat Schura]\label{aa:lem:schur-first}
Niech \(G\) będzie grupą i \(K\) ciałem.
Jeżeli \(V\) i \(W\) są nieprzywiedlnymi, skończenie
wymiarowymi reprezentacjami \(G\) nad \(K\), to każdy
niezerowy \(T\in\operatorname{Hom}_G(V,W)\) jest izomorfizmem.
W szczególności, jeżeli \(V\) i \(W\) nie są równoważne,
to \(\operatorname{Hom}_G(V,W)=\{0\}\).
\end{lemat}
\begin{proof}
Jeżeli \(v\in\ker T\), to
\(T\rho(g)v=\sigma(g)Tv=0\), więc jądro jest
podprzestrzenią \(G\)-niezmienniczą przestrzeni \(V\).
Analogicznie \(\sigma(g)T(V)=T\rho(g)(V)\subseteq T(V)\),
zatem obraz jest \(G\)-niezmienniczy w \(W\).
Nieprzywiedlność daje \(\ker T=0\) albo \(\ker T=V\),
oraz \(T(V)=0\) albo \(T(V)=W\).
Dla \(T\ne0\) zachodzą odpowiednio \(\ker T=0\)
i \(T(V)=W\), co oznacza, że \(T\) jest izomorfizmem.
\end{proof}

\begin{twierdzenie}[Lemat Schura nad \(\mathbb C\)]
\label{aa:tw:schur-complex}
Niech \(G\) będzie grupą, a \(V\ne0\) skończenie
wymiarową, nieprzywiedlną reprezentacją \(G\)
nad \(\mathbb C\). Wówczas każdy operator
\(T\in\operatorname{Hom}_G(V,V)\) jest postaci
\(T=\lambda I_V\) dla pewnej liczby
\(\lambda\in\mathbb C\).
\end{twierdzenie}
\begin{proof}
Wielomian charakterystyczny operatora \(T\)
ma dodatni stopień, więc na mocy zasadniczego
twierdzenia algebry posiada pierwiastek
\(\lambda\in\mathbb C\).
Operator \(T-\lambda I_V\) przeplata działanie \(G\)
i nie jest odwracalny, gdyż ma niezerowe jądro.
Z pierwszego lematu Schura wynika, że gdyby był
niezerowy, musiałby być izomorfizmem.
Dlatego \(T-\lambda I_V=0\).
\end{proof}

\begin{przyklad}[Dlaczego ciało \(\mathbb R\) nie wystarcza]
\label{aa:prz:schur-real}
Niech \(G=C_3=\langle g\rangle\) i niech generator działa
na \(\mathbb R^2\) przez obrót o \(2\pi/3\):
\[
\rho(g)=
\begin{pmatrix}
-\frac12&-\frac{\sqrt3}{2}\\
\frac{\sqrt3}{2}&-\frac12
\end{pmatrix}.
\]
Macierz ma wielomian charakterystyczny \(T^2+T+1\),
bez pierwiastków rzeczywistych. Nie istnieje więc
prosta niezmiennicza, a reprezentacja jest nieprzywiedlna
nad \(\mathbb R\). Operator
\(J=\left(\begin{smallmatrix}0&-1\\1&0\end{smallmatrix}\right)\)
komutuje z \(\rho(g)\), lecz nie jest wielokrotnością
\(I\). Zatem końcowa teza lematu Schura nad \(\mathbb C\)
nie zachodzi w tej postaci nad \(\mathbb R\).
\end{przyklad}

\section{Uśrednianie iloczynu hermitowskiego}
\begin{stwierdzenie}[Unitarność reprezentacji grupy skończonej]
\label{aa:prop:unitarization}
Niech \(G\) będzie grupą skończoną, a
\(\rho:G\to\operatorname{GL}_{\mathbb C}(V)\)
reprezentacją na skończenie wymiarowej przestrzeni
zespolonej \(V\). Istnieje dodatnio określony
iloczyn hermitowski na \(V\), względem którego
wszystkie \(\rho(g)\) są unitarne.
\end{stwierdzenie}
\begin{proof}
Wybierzmy dowolny dodatnio określony iloczyn hermitowski
\(\langle\cdot,\cdot\rangle_0\) na \(V\) i połóżmy
\[
\langle v,w\rangle_G
=\frac1{|G|}\sum_{g\in G}
\langle\rho(g)v,\rho(g)w\rangle_0.
\]
Skończona suma zachowuje własności
półtoraliniowości oraz symetrii hermitowskiej.
Dla \(v\ne0\) każdy składnik
\(\langle\rho(g)v,\rho(g)v\rangle_0\) jest dodatni,
bo \(\rho(g)\) jest odwracalne. Iloczyn jest zatem
dodatnio określony. Dla \(h\in G\) podstawienie
\(k=gh\) daje
\(\langle\rho(h)v,\rho(h)w\rangle_G
=\langle v,w\rangle_G\).
Z definicji oznacza to unitarność każdego \(\rho(h)\).
\end{proof}

\begin{definicja}[Charakter reprezentacji]
\label{aa:def:character}
Niech \(G\) będzie grupą skończoną, a
\(\rho:G\to\operatorname{GL}_{\mathbb C}(V)\)
reprezentacją na skończenie wymiarowej przestrzeni
zespolonej \(V\). Jej charakterem nazywamy funkcję
\(\chi_\rho:G\to\mathbb C\) określoną wzorem
\(\chi_\rho(g)=\operatorname{tr}(\rho(g))\).
Charakter reprezentacji nieprzywiedlnej nazywamy
charakterem nieprzywiedlnym.
Funkcję \(f:G\to\mathbb C\) nazywamy funkcją klasową,
gdy \(f(hgh^{-1})=f(g)\) dla wszystkich \(g,h\in G\).
\end{definicja}

\begin{stwierdzenie}[Podstawowe własności charakterów]
\label{aa:prop:character-basics}
Niech \(G\) będzie grupą skończoną, a
\(\rho:G\to\operatorname{GL}_{\mathbb C}(V)\)
skończenie wymiarową reprezentacją zespoloną
o charakterze \(\chi_\rho\). Wtedy
\(\chi_\rho(e)=\dim_{\mathbb C}V\),
\(\chi_\rho(hgh^{-1})=\chi_\rho(g)\) oraz
\(\chi_\rho(g^{-1})=\overline{\chi_\rho(g)}\).
Ponadto dla reprezentacji \(\rho_1,\rho_2\)
zachodzi \(\chi_{\rho_1\oplus\rho_2}
=\chi_{\rho_1}+\chi_{\rho_2}\).
\end{stwierdzenie}
\begin{proof}
Pierwsza równość wynika z \(\rho(e)=I_V\).
Ponieważ
\(\rho(hgh^{-1})=\rho(h)\rho(g)\rho(h)^{-1}\),
niezmienniczość śladu przy podobieństwie daje drugą.
Możemy przyjąć iloczyn hermitowski ze
Stwierdzenia~\ref{aa:prop:unitarization}.
Wtedy \(\rho(g^{-1})=\rho(g)^*\), a ślad operatora
sprzężonego hermitowsko jest sprzężeniem śladu
wyjściowego. Ostatnia własność wynika z postaci
blokowo diagonalnej macierzy sumy prostej.
\end{proof}

\section{Relacje ortogonalności}
\begin{definicja}[Iloczyn skalarny funkcji na grupie]
\label{aa:def:character-inner-product}
Niech \(G\) będzie grupą skończoną.
Dla funkcji \(f,h:G\to\mathbb C\) określamy
\[
\langle f,h\rangle_G
=\frac1{|G|}\sum_{g\in G}f(g)\overline{h(g)}.
\]
Jest to dodatnio określony iloczyn hermitowski
na zespolonej przestrzeni wszystkich funkcji \(G\to\mathbb C\).
Ogranicza się on do przestrzeni funkcji klasowych.
\end{definicja}

\begin{twierdzenie}[Ortogonalność współczynników macierzowych]
\label{aa:tw:matrix-orthogonality}
Niech \(G\) będzie grupą skończoną oraz
\(\rho_i:G\to U(V_i)\), \(1\le i\le r\),
skończenie wymiarowymi nieprzywiedlnymi reprezentacjami
zespolonymi, parami nierównoważnymi.
Niech \(d_i=\dim_{\mathbb C}V_i\) i wybierzmy
bazę ortonormalną każdej \(V_i\). Dla współczynników
macierzy \(\rho_i(g)\) zachodzi
\[
\frac1{|G|}\sum_{g\in G}
(\rho_i(g))_{ab}\overline{(\rho_j(g))_{cd}}
=\begin{cases}
\dfrac{\delta_{ac}\delta_{bd}}{d_i},&i=j,\\[3pt]
0,&i\ne j,
\end{cases}
\]
dla \(1\le a,b\le d_i\) oraz \(1\le c,d\le d_j\).
\end{twierdzenie}
\begin{proof}
Dla liniowego \(A:V_j\to V_i\) rozważmy operator
\[
\mathcal P_{ij}(A)=\frac1{|G|}\sum_{g\in G}
\rho_i(g)A\rho_j(g)^{-1}.
\]
Podstawienie \(k=hg\) pokazuje, że
\(\rho_i(h)\mathcal P_{ij}(A)
=\mathcal P_{ij}(A)\rho_j(h)\) dla każdego \(h\in G\).
Gdy \(i\ne j\), pierwszy lemat Schura daje
\(\mathcal P_{ij}(A)=0\). Gdy \(i=j\), zespolony
lemat Schura zapewnia \(\mathcal P_{ii}(A)=\lambda I\).
Z niezmienniczości śladu przy sprzężeniu otrzymujemy
\(d_i\lambda=\operatorname{tr}A\), czyli
\(\mathcal P_{ii}(A)=(\operatorname{tr}A/d_i)I\).

Niech teraz \(A\) wysyła \(d\)-ty wektor bazy \(V_j\)
na \(b\)-ty wektor bazy \(V_i\), a pozostałe
wektory bazowe na zero. Wówczas współczynnik
\((a,c)\) operatora \(\mathcal P_{ij}(A)\) wynosi
\[
\frac1{|G|}\sum_{g\in G}
(\rho_i(g))_{ab}(\rho_j(g)^{-1})_{dc}
=\frac1{|G|}\sum_{g\in G}
(\rho_i(g))_{ab}\overline{(\rho_j(g))_{cd}},
\]
gdzie ostatnia równość korzysta z unitarności.
Gdy \(i=j\), mamy
\(\operatorname{tr}A=\delta_{bd}\), a współczynnik
\((a,c)\) operatora \(I\) jest równy \(\delta_{ac}\).
Dostajemy obie postacie żądanej relacji.
\end{proof}

\begin{wniosek}[Ortogonalność charakterów nieprzywiedlnych]
\label{aa:cor:character-orthogonality}
Niech \(G\) będzie grupą skończoną, a
\(\chi_1,\ldots,\chi_r\) charakterami parami
nierównoważnych, skończenie wymiarowych
nieprzywiedlnych reprezentacji zespolonych \(G\).
Wtedy \(\langle\chi_i,\chi_j\rangle_G=\delta_{ij}\).
\end{wniosek}
\begin{proof}
Dzięki unitarności możemy zastosować
Twierdzenie~\ref{aa:tw:matrix-orthogonality}.
Wyrażamy oba ślady jako sumy elementów diagonalnych:
\[
\langle\chi_i,\chi_j\rangle_G
=\sum_{a=1}^{d_i}\sum_{c=1}^{d_j}
\frac1{|G|}\sum_{g\in G}
(\rho_i(g))_{aa}\overline{(\rho_j(g))_{cc}}.
\]
Dla \(i\ne j\) każdy składnik jest zerowy.
Dla \(i=j\) wewnętrzna suma wynosi
\(\delta_{ac}/d_i\); suma po \(a,c\) jest równa \(1\).
\end{proof}

\section{Krotności składników i reprezentacja regularna}
\begin{stwierdzenie}[Występowanie nieprzywiedlnych składników]
\label{aa:prop:regular-contains-all}
Niech \(G\) będzie grupą skończoną.
Każda skończenie wymiarowa nieprzywiedlna reprezentacja
zespolona \(V\) grupy \(G\) jest równoważna składnikowi
reprezentacji regularnej na \(\mathbb C[G]\).
W szczególności istnieje tylko skończenie wiele klas
równoważności takich reprezentacji.
\end{stwierdzenie}
\begin{proof}
Wybierzmy \(0\ne v\in V\) i określmy
\(G\)-równoważne odwzorowanie
\(\Phi:\mathbb C[G]\to V\) przez
\(\Phi([g])=\rho(g)v\), przedłużając liniowo.
Jego obraz jest niezerową podprzestrzenią niezmienniczą,
a więc dzięki nieprzywiedlności jest równy \(V\).
Twierdzenie Maschkego zapewnia niezmiennicze dopełnienie
\(\ker\Phi\) w \(\mathbb C[G]\); ograniczenie \(\Phi\)
do tego dopełnienia daje izomorfizm z \(V\).
Ponadto rozkład skończenie wymiarowej reprezentacji
regularnej na nieprzywiedlne składniki ma skończenie
wiele wyrazów. Każda klasa nieprzywiedlna pojawia
się w tym rozkładzie: w istocie izomorficzny z \(V\)
podmoduł, który przed chwilą skonstruowaliśmy,
ma charakter zawarty w rozkładzie regularnym.
Ściślej, jeśli \(\mathbb C[G]=\bigoplus_j S_j\),
to rzut na \(V\) złożony z włączeniami \(S_j\)
musi być niezerowy dla pewnego \(j\). Pierwszy
lemat Schura daje wtedy \(S_j\simeq V\).
Liczba różnych klas jest więc skończona.
\end{proof}

\begin{twierdzenie}[Krotności z iloczynów skalarnych]
\label{aa:tw:multiplicities}
Niech \(G\) będzie grupą skończoną.
Niech \(V_1,\ldots,V_r\) będą przedstawicielami
wszystkich klas równoważności nieprzywiedlnych
reprezentacji zespolonych \(G\), z charakterami
\(\chi_1,\ldots,\chi_r\). Jeżeli \(V\) jest dowolną
skończenie wymiarową reprezentacją zespoloną
o charakterze \(\chi_V\), to istnieją jednoznaczne
liczby całkowite \(m_i\ge0\) takie, że
\[
V\simeq\bigoplus_{i=1}^r V_i^{\oplus m_i},
\qquad m_i=\langle\chi_V,\chi_i\rangle_G.
\]
W konsekwencji dwie skończenie wymiarowe reprezentacje
zespolone grupy \(G\) są równoważne wtedy i tylko
wtedy, gdy mają te same charaktery.
\end{twierdzenie}
\begin{proof}
Istnienie rozkładu wynika z twierdzenia Maschkego,
a kompletność listy \(V_i\) ze
Stwierdzenia~\ref{aa:prop:regular-contains-all}.
Addytywność śladu daje
\(\chi_V=\sum_i m_i\chi_i\).
Iloczyn skalarny z \(\chi_j\) i ortogonalność
charakterów prowadzą do
\(\langle\chi_V,\chi_j\rangle_G=m_j\).
Stąd krotności są wyznaczone jednoznacznie.
Jeśli charaktery dwóch reprezentacji są równe,
ich krotności wszystkich składników są jednakowe,
więc obie reprezentacje są równoważne.
Odwrotność wynika z niezmienniczości śladu
przy podobieństwie.
\end{proof}

\begin{twierdzenie}[Charakter regularny i suma kwadratów wymiarów]
\label{aa:tw:regular-character}
Niech \(G\) będzie grupą skończoną i niech
\(\rho_{\mathrm{reg}}\) będzie lewą reprezentacją
regularną na \(\mathbb C[G]\). Jej charakter spełnia
\[
\chi_{\mathrm{reg}}(g)=
\begin{cases}|G|,&g=e,\\0,&g\ne e.\end{cases}
\]
Jeżeli \(d_i=\dim_{\mathbb C}V_i=\chi_i(e)\)
dla pełnej rodziny \(V_1,\ldots,V_r\) jak wyżej,
to \(\mathbb C[G]\simeq\bigoplus_{i=1}^r
V_i^{\oplus d_i}\) oraz
\(\sum_{i=1}^r d_i^2=|G|\).
\end{twierdzenie}
\begin{proof}
W bazie \(([h])_{h\in G}\) działanie elementu \(g\)
permutuje bazę według \([h]\mapsto[gh]\).
Jeżeli \(g\ne e\), to żaden wektor bazowy nie
jest stały, bo \(gh=h\) wymusza \(g=e\).
Ślad macierzy permutacyjnej jest liczbą
nieruchomych wektorów bazowych, zatem wynosi \(0\);
dla \(g=e\) wynosi \(|G|\).
Z Twierdzenia~\ref{aa:tw:multiplicities} krotność
\(V_i\) w reprezentacji regularnej to
\[
\langle\chi_{\mathrm{reg}},\chi_i\rangle_G
=\frac1{|G|}|G|\overline{\chi_i(e)}=d_i.
\]
Porównując wymiary obu stron otrzymanego rozkładu,
dostajemy \(|G|=\sum_i d_i^2\).
\end{proof}

\begin{stwierdzenie}[Jawny rzut na składnik izotypiczny]
\label{aa:prop:isotypic-projector}
Niech \(G\) będzie grupą skończoną, a
\(V\) skończenie wymiarową reprezentacją zespoloną
o działaniu \(\rho\). Dla każdego charakteru
nieprzywiedlnego \(\chi_i\), którego wymiar
wynosi \(d_i=\chi_i(e)\), określmy
\[
P_i=\frac{d_i}{|G|}\sum_{g\in G}
\overline{\chi_i(g)}\,\rho(g).
\]
Wtedy \(P_i\) jest \(G\)-równoważnym rzutem
na sumę wszystkich składników \(V\) równoważnych
\(V_i\) w dowolnym rozkładzie na składniki
nieprzywiedlne. Ponadto \(P_iP_j=0\) dla \(i\ne j\)
oraz \(\sum_{i=1}^rP_i=I_V\), gdy \(\chi_1,\ldots,\chi_r\)
stanowią pełną rodzinę.
\end{stwierdzenie}
\begin{proof}
Charakter \(\chi_i\) jest funkcją klasową, zatem
sprzężenie sumy definiującej \(P_i\) przez
\(\rho(h)\) jedynie permutuje jej składniki.
Wobec tego \(P_i\) przeplata działanie \(G\).
Z twierdzenia Maschkego weźmy rozkład
\(V=\bigoplus_\alpha W_\alpha\) na nieprzywiedlne
podreprezentacje. Każda z nich jest zachowana przez
\(P_i\), ponieważ sumujemy operatory \(\rho(g)\),
które osobno zachowują \(W_\alpha\).
Gdy \(W_\alpha\simeq V_j\), zespolony lemat Schura
daje \(P_i|_{W_\alpha}=\lambda I\). Obliczając ślad,
otrzymujemy
\[
d_j\lambda
=\frac{d_i}{|G|}\sum_{g\in G}
\overline{\chi_i(g)}\chi_j(g)
=d_i\delta_{ij}.
\]
Jeżeli \(j=i\), to \(\lambda=1\); w przeciwnym
przypadku \(\lambda=0\). To ustala działanie
operatorów \(P_i\) na każdym składniku rozkładu
i dowodzi wszystkich wskazanych własności.
\end{proof}

\section{Funkcje klasowe i ortogonalność kolumn}
\begin{twierdzenie}[Charaktery jako baza funkcji klasowych]
\label{aa:tw:characters-basis}
Niech \(G\) będzie grupą skończoną.
Wszystkie charaktery nieprzywiedlne zespolone
\(\chi_1,\ldots,\chi_r\) tworzą bazę ortonormalną
przestrzeni funkcji klasowych \(G\to\mathbb C\)
względem iloczynu skalarnego
z Definicji~\ref{aa:def:character-inner-product}.
W szczególności liczba \(r\) klas równoważności
nieprzywiedlnych reprezentacji zespolonych \(G\)
jest równa liczbie klas sprzężoności w \(G\).
\end{twierdzenie}
\begin{proof}
Charaktery są funkcjami klasowymi i są ortonormalne
na mocy Wniosku~\ref{aa:cor:character-orthogonality}.
Aby wykazać zupełność, załóżmy, że funkcja klasowa
\(f:G\to\mathbb C\) jest prostopadła do każdego
\(\chi_i\). Dla dowolnej reprezentacji \(\rho\)
grupy \(G\) rozważmy operator
\[
T_f^\rho=\sum_{g\in G}f(g)\rho(g^{-1}).
\]
Ponieważ \(f\) jest stała na klasach sprzężoności,
operator ten przeplata działanie \(G\).
Na nieprzywiedlnej reprezentacji \(V_i\)
lemat Schura daje \(T_f^{\rho_i}=\lambda_iI\).
Biorąc ślad i korzystając z równości
\(\chi_i(g^{-1})=\overline{\chi_i(g)}\), mamy
\[
d_i\lambda_i=\sum_{g\in G}
f(g)\overline{\chi_i(g)}
=|G|\langle f,\chi_i\rangle_G=0.
\]
Zatem \(T_f^{\rho_i}=0\) dla wszystkich \(i\).
Reprezentacja regularna jest sumą takich składników,
więc \(T_f^{\mathrm{reg}}=0\). Jednak
\(T_f^{\mathrm{reg}}([e])
=\sum_{g\in G}f(g)[g^{-1}]\).
Liniowa niezależność formalnej bazy \(([g])\)
wymusza \(f(g)=0\) dla każdego \(g\in G\).
W przestrzeni skończenie wymiarowej brak
niezerowego wektora prostopadłego do rozpiętości
ortonormalnych charakterów oznacza, że rozpinają
całą przestrzeń funkcji klasowych.
Jej wymiar jest liczbą klas sprzężoności, ponieważ
funkcja klasowa jest jednoznacznie wyznaczona przez
jedną wartość na każdej klasie. Stąd \(r\) jest równy
tej liczbie.
\end{proof}

\begin{wniosek}[Ortogonalność kolumn tablicy charakterów]
\label{aa:cor:column-orthogonality}
Niech \(G\) będzie grupą skończoną,
\(C_1,\ldots,C_r\) jej wszystkimi klasami sprzężoności
i \(g_k\in C_k\) ustalonymi reprezentantami.
Jeżeli \(\chi_1,\ldots,\chi_r\) są wszystkimi
charakterami nieprzywiedlnymi zespolonymi, to
\[
\sum_{i=1}^r\chi_i(g_k)\overline{\chi_i(g_\ell)}
=\begin{cases}
|G|/|C_k|,&k=\ell,\\
0,&k\ne\ell.
\end{cases}
\]
\end{wniosek}
\begin{proof}
W przestrzeni funkcji klasowych iloczyn skalarny ma
postać
\(\langle f,h\rangle_G
=\sum_{k=1}^r (|C_k|/|G|)f(g_k)\overline{h(g_k)}\).
Z Twierdzenia~\ref{aa:tw:characters-basis}
macierz kwadratowa
\(A_{ik}=\sqrt{|C_k|/|G|}\,\chi_i(g_k)\)
ma ortonormalne wiersze. Jest więc unitarna,
a jej kolumny także są ortonormalne.
Równości \((A^*A)_{k\ell}=\delta_{k\ell}\)
dają tezę po usunięciu dodatnich czynników
\(\sqrt{|C_k||C_\ell|}/|G|\).
\end{proof}

\begin{wniosek}[Rozwinięcie funkcji klasowej]
\label{aa:cor:class-fourier}
Niech \(G\) będzie grupą skończoną, a
\(f:G\to\mathbb C\) funkcją klasową.
Dla pełnej rodziny charakterów
nieprzywiedlnych \(\chi_1,\ldots,\chi_r\) zachodzi
\[
f=\sum_{i=1}^r\langle f,\chi_i\rangle_G\,\chi_i.
\]
Współczynniki są wyznaczone jednoznacznie.
\end{wniosek}
\begin{proof}
Twierdzenie~\ref{aa:tw:characters-basis} zapewnia
rozwinięcie w bazie ortonormalnej. Biorąc iloczyn
skalarny obu stron z \(\chi_j\), dostajemy
jednoznacznie współczynnik przy \(\chi_j\).
\end{proof}

\section{Dwie tablice charakterów}
\begin{przyklad}[Wszystkie charaktery grupy \(C_3\)]
\label{aa:prz:character-c3}
Niech \(C_3=\langle g:g^3=e\rangle\) oraz
\(\zeta=\exp(2\pi i/3)\). Trzy reprezentacje
jednowymiarowe określone wzorami
\(\rho_j(g)=\zeta^j\), \(j=0,1,2\), mają charaktery
\(\chi_j(g^k)=\zeta^{jk}\).
Są nierównoważne, gdyż mają różne wartości na \(g\).
Równość
\(\sum_{k=0}^2\zeta^{(j-\ell)k}
=3\delta_{j\ell}\) dowodzi bezpośrednio ich
ortogonalności. Wzór na sumę kwadratów wymiarów
daje \(1^2+1^2+1^2=3\), a zatem nie występują
żadne inne klasy reprezentacji nieprzywiedlnych.
Tablica charakterów ma postać
\[
\begin{array}{c|ccc}
 & e & g & g^2\\ \hline
\chi_0&1&1&1\\
\chi_1&1&\zeta&\zeta^2\\
\chi_2&1&\zeta^2&\zeta
\end{array}.
\]
Ponieważ \(C_3\) jest abelowa, każda klasa
sprzężoności jest jednoelementowa. Ortogonalność
kolumn wynika stąd również przez bezpośrednie
obliczenie iloczynów kolumn tej tablicy.
\end{przyklad}

\begin{przyklad}[Pełna tablica charakterów \(S_3\)]
\label{aa:prz:character-s3}
Grupa \(S_3\) ma trzy klasy sprzężoności:
jednostkę \(e\), trzy transpozycje \(t\) i dwa
trójcykle \(c\). Reprezentacja trywialna ma charakter
\((1,1,1)\), a reprezentacja znaku \((1,-1,1)\).
Dla permutacyjnej reprezentacji \(\mathbb C^3\)
ślad permutacji jest liczbą nieruchomych wektorów
bazy. Jej charakter przyjmuje więc wartości
\((3,1,0)\). Odejmując charakter składnika
trywialnego, otrzymujemy charakter
\(\chi_{\mathrm{std}}=(2,0,-1)\) reprezentacji
na płaszczyźnie \(x_1+x_2+x_3=0\).

W tej sytuacji dla funkcji klasowych
\(\langle f,h\rangle_{S_3}
=(f(e)\overline{h(e)}
+3f(t)\overline{h(t)}
+2f(c)\overline{h(c)})/6\).
Stąd
\(\langle\chi_{\mathrm{std}},
\chi_{\mathrm{std}}\rangle_{S_3}
=(4+2)/6=1\),
więc reprezentacja standardowa jest nieprzywiedlna:
gdyby była sumą co najmniej dwóch nieprzywiedlnych,
jej iloczyn skalarny z samą sobą byłby sumą
kwadratów dodatnich całkowitych krotności,
a więc wynosiłby co najmniej \(2\).
Iloczyny mieszane trzech charakterów są zerowe,
co można sprawdzić na przykładzie
\((2+2(-1))/6=0\) dla charakteru trywialnego
i standardowego.

Pełna tablica jest następująca:
\[
\begin{array}{c|ccc}
& [e]\ (1)&[t]\ (3)&[c]\ (2)\\ \hline
\chi_{\mathrm{triv}}&1&1&1\\
\chi_{\mathrm{sgn}}&1&-1&1\\
\chi_{\mathrm{std}}&2&0&-1
\end{array}.
\]
Liczby w nawiasach to rozmiary klas. Ponieważ
\(1^2+1^2+2^2=6=|S_3|\), tablica obejmuje
wszystkie reprezentacje nieprzywiedlne.
W pierwszej kolumnie suma kwadratów to \(6\),
a w kolumnach transpozycji i trójcykli
wynosi odpowiednio \(2=6/3\) i \(3=6/2\),
zgodnie z ortogonalnością kolumn.
\end{przyklad}

\begin{uwaga}[Granice stosowalności]
\label{aa:rem:characters-limits}
Wnioski o ortonormalności charakterów,
rozpoznawaniu reprezentacji po ich charakterach
i sumie kwadratów wymiarów zostały udowodnione
dla reprezentacji zespolonych grup skończonych.
Nie wolno ich bezpośrednio stosować w
charakterystyce dzielącej rząd grupy:
wtedy twierdzenie Maschkego może zawieść,
a różne nierównoważne reprezentacje mogą mieć
ten sam charakter śladowy. Przykładowo nad
\(\mathbb F_2\) trywialna reprezentacja \(C_2\)
w wymiarze \(2\) i nietrywialne działanie
\(s\mapsto\left(\begin{smallmatrix}1&1\\0&1\end{smallmatrix}\right)\)
mają ślad \(0\) dla obu elementów \(C_2\),
chociaż reprezentacje nie są równoważne:
pierwsza ma działanie trywialne, druga nie.
\end{uwaga}

% AA-43 | Rozdział przygotowany ze starszego skryptu sources/pl/galois(2).tex.
% Dalsze ogniwa: AA-44 — twierdzenie Liouville'a, AA-45 — całki nieelementarne,
% AA-46 — podstawy różniczkowej teorii Galois.
\chapter{Ciała różniczkowe, derywacje i rozszerzenia różniczkowe}
\chaptermark{Ciała różniczkowe i derywacje}
\label{aa:ch:differential-fields}

Równanie algebraiczne można badać przez rozszerzenia ciał i ich
automorfizmy. Badanie równania $DF=f$, gdzie $D$ jest pochodną,
wymaga dodatkowo informacji o tym, jak różniczkowanie zachowuje się
przy dołączaniu nowych elementów. Zaczniemy od derywacji, ciał
stałych oraz rozszerzeń algebraicznych i przestępnych. Zbudowane
narzędzia posłużą w kolejnym rozdziale do sformułowania i dowodu
twierdzenia Liouville'a o całkowaniu w funkcjach elementarnych.

\section{Derywacje i ciała stałych}

\begin{definicja}[Ciało różniczkowe]
\label{aa:def:differential-field}
Niech $K$ będzie ciałem charakterystyki $0$.
\emph{Derywacją} na $K$ nazywamy odwzorowanie $D:K\to K$,
dla którego dla wszystkich $a,b\in K$ zachodzą
$D(a+b)=Da+Db$ oraz $D(ab)=aDb+bDa$.
Parę $(K,D)$ nazywamy \emph{ciałem różniczkowym}, a
\[
C_K=\{a\in K:Da=0\}
\]
jego \emph{ciałem stałych}.
\end{definicja}

\begin{stwierdzenie}[Reguły rachunku różniczkowego]
\label{aa:st:differential-field-rules}
Niech $(K,D)$ będzie ciałem różniczkowym. Wtedy
$D(1)=0$, $D(-a)=-Da$, $D(a^{-1})=-a^{-2}Da$
dla $a\ne0$ oraz
\[
D\!\left(\frac ab\right)=\frac{bDa-aDb}{b^2}
\quad(a,b\in K,\ b\ne0).
\]
Zbiór $C_K$ jest podciałem $K$, zawierającym $\mathbb Q$.
Dla każdej liczby całkowitej $n$ i $a\in K^\times$
zachodzi $D(a^n)=na^{n-1}Da$.
\end{stwierdzenie}
\begin{proof}
Z reguły iloczynu dla $1\cdot1=1$ wynika
$D1=2D1$, stąd $D1=0$. Addytywność daje
$D(-a)=-Da$. Jeśli $a\ne0$, to
$0=D(aa^{-1})=(Da)a^{-1}+aD(a^{-1})$,
co daje regułę odwrotności, a następnie ilorazu.
Jądro addytywnego odwzorowania $D$ jest zamknięte
na dodawanie, odejmowanie i mnożenie; reguła
odwrotności zapewnia też zamkniętość na odwrotności
elementów niezerowych. Każda liczba całkowita ma
pochodną zero, podobnie jak jej odwrotność w $K$,
więc $\mathbb Q\subseteq C_K$.
Regułę dla dodatnich $n$ otrzymujemy przez indukcję
z reguły iloczynu, a dla ujemnych z reguły
odwrotności. Przypadek $n=0$ został już wykazany.
\end{proof}

\begin{przyklad}[Funkcje wymierne]
\label{aa:prz:differential-rational-field}
Niech $C$ będzie ciałem charakterystyki $0$, a
$K=C(x)$ ciałem funkcji wymiernych jednej zmiennej.
Istnieje jedyna derywacja $D$ o własnościach
$D(c)=0$ dla $c\in C$ i $Dx=1$.
Dla $P,Q\in C[x]$, $Q\ne0$, reguła ilorazu przyjmuje postać
\[
D(P/Q)=\frac{P'Q-PQ'}{Q^2}.
\]
Pokażemy, że $C_K=C$. Jeśli niezerowa funkcja $P/Q$,
zapisana w postaci nieskracalnej, spełnia $D(P/Q)=0$,
to $P'Q=PQ'$. Względna pierwszość $P$ i $Q$
w $C[x]$ oznacza, że $P\mid P'$.
Gdyby $\deg P\ge1$, w charakterystyce $0$ mielibyśmy
$\deg P'=\deg P-1$, sprzeczność. Zatem $P$
jest stały; równość $P'Q=PQ'$ implikuje wtedy
$Q'=0$, czyli także $Q$ jest stały.
Funkcja zerowa oczywiście należy do $C$.
\end{przyklad}

\begin{przyklad}[Derywacja zerowa i rola charakterystyki]
\label{aa:prz:zero-derivation-and-positive-characteristic}
Derywacja $D=0$ na dowolnym ciele charakterystyki $0$
ma całe to ciało jako ciało stałych.
Nie można natomiast bez zmiany twierdzeń zastąpić
założenia o charakterystyce $0$ charakterystyką $p>0$:
w $\mathbb F_p(x)$, przy $Dx=1$, mamy
$D(x^p)=p x^{p-1}=0$, choć $x^p$ nie należy
do $\mathbb F_p$. Co więcej, dołączenie pierwiastka
$\alpha$ równania $\alpha^p=x$ nie pozwala przedłużyć
derywacji z $Dx=1$, bo
$1=Dx=D(\alpha^p)=p\alpha^{p-1}D\alpha=0$.
\end{przyklad}

\begin{definicja}[Rozszerzenie różniczkowe]
\label{aa:def:differential-extension}
Niech $(K,D)$ będzie ciałem różniczkowym.
\emph{Rozszerzeniem różniczkowym} $L/K$ nazywamy
rozszerzenie ciał wraz z derywacją, również
oznaczaną przez $D:L\to L$, której ograniczenie
do $K$ jest równe pierwotnej derywacji.
Mówimy, że jest ono \emph{bez nowych stałych},
jeśli $C_L=C_K$.
Homomorfizm $\sigma:L\to M$ rozszerzeń różniczkowych
nad $K$ nazywamy \emph{różniczkowym}, jeśli
$\sigma|_K=\mathrm{id}_K$ i $D\sigma(a)=\sigma(Da)$
dla każdego $a\in L$.
\end{definicja}

Warunek $C_L=C_K$ jest istotny, ponieważ z samej
własności $\sigma|_K=\mathrm{id}_K$ nie wynika,
że funkcja pierwotna ma jednoznaczną stałą całkowania
w $C_K$. Pokażemy niżej, jak zapewnić brak nowych
stałych w konstrukcjach logarytmicznych.

\section{Algebraiczne przedłużanie derywacji}

\begin{twierdzenie}[Jednoznaczne przedłużenie]
\label{aa:tw:algebraic-derivation-extension}
Niech $(K,D)$ będzie ciałem różniczkowym
charakterystyki $0$, a $L/K$ dowolnym rozszerzeniem
algebraicznym. Wówczas istnieje dokładnie jedna
derywacja $L\to L$ przedłużająca $D$.

Jeśli $\alpha\in L$ ma nad $K$ unormowany wielomian
minimalny $m_\alpha(T)\in K[T]$, a
$m_\alpha^D(T)$ powstaje przez zastosowanie $D$
do wszystkich jego współczynników, to
\[
D\alpha=-\frac{m_\alpha^D(\alpha)}{m_\alpha'(\alpha)}.
\]
Mianownik jest niezerowy.
\end{twierdzenie}
\begin{proof}
Każdy nierozkładalny wielomian nad ciałem charakterystyki
$0$ jest rozdzielczy (zob.
Wniosek~\ref{aa:wn:charzero-separable}). W szczególności
$m_\alpha'(\alpha)\ne0$.
Jeśli przedłużenie istnieje, zastosowanie derywacji
do równości $m_\alpha(\alpha)=0$ oraz reguły iloczynu
wymusza
$0=m_\alpha^D(\alpha)+m_\alpha'(\alpha)D\alpha$.
Otrzymujemy wskazany, jedyny możliwy wzór.

Dla ustalonej $\alpha$ zdefiniujmy $D\alpha$ tym wzorem.
Dla wielomianu $f\in K[T]$ połóżmy
$D(f(\alpha))=f^D(\alpha)+f'(\alpha)D\alpha$.
Musimy wykazać niezależność od przedstawienia.
Jeśli $f(\alpha)=g(\alpha)$, to
$f-g=m_\alpha h$ dla pewnego $h\in K[T]$.
Reguła Leibniza dla wielomianów daje
\[
(f-g)^D(\alpha)+(f-g)'(\alpha)D\alpha
=h(\alpha)\bigl(m_\alpha^D(\alpha)
                +m_\alpha'(\alpha)D\alpha\bigr)=0,
\]
ponieważ wszystkie wyrazy zawierające
$m_\alpha(\alpha)$ znikają. Operacja jest więc
dobrze określona na $K[\alpha]$.
Bezpośrednio z definicji i reguły iloczynu
dla wielomianów wynika addytywność i reguła
Leibniza. Ponieważ $\alpha$ jest algebraiczny,
$K[\alpha]=K(\alpha)$ jest ciałem; wzór na
odwrotność zapewnia derywację na całym $K(\alpha)$.

Jeżeli $L/K$ jest skończone, możemy dołączać
skończenie wiele generatorów algebraicznych
jeden po drugim. Na każdym piętrze przedłużenie
jest jednoznaczne, więc wynik nie zależy
od wybranej kolejności generatorów.
W rozszerzeniu dowolnym każdy element $L$
należy do skończonego podrozszerzenia.
Na sumie dwóch takich podrozszerzeń
przedłużenia z poprzedniego kroku są jednoznaczne;
zgadzają się więc na ich części wspólnej.
Definiując pochodną każdego elementu przez
jego dowolne skończone podrozszerzenie,
otrzymujemy spójną derywację na $L$.
Jednoznaczność dla $L$ wynika z tego samego wzoru.
\end{proof}

\begin{wniosek}[Brak nowych stałych przy rozszerzeniu algebraicznym]
\label{aa:cor:algebraic-constants}
Niech $(K,D)$ będzie ciałem różniczkowym
charakterystyki $0$, którego ciało stałych $C=C_K$
jest algebraicznie domknięte. Wtedy dowolne
rozszerzenie algebraiczne $L/K$, wyposażone
w jedyne przedłużenie derywacji, spełnia $C_L=C$.
\end{wniosek}
\begin{proof}
Oczywiście $C\subseteq C_L$. Niech $\alpha\in C_L$,
a $m_\alpha\in K[T]$ będzie unormowanym
wielomianem minimalnym. Z poprzedniego twierdzenia
i $D\alpha=0$ wynika $m_\alpha^D(\alpha)=0$.
Ponieważ współczynnik wiodący $m_\alpha$ wynosi $1$,
mamy $\deg m_\alpha^D<\deg m_\alpha$,
chyba że $m_\alpha^D=0$. Minimalność wielomianu
$m_\alpha$ wymusza $m_\alpha^D=0$.
Wszystkie współczynniki $m_\alpha$ należą zatem do $C$,
więc $\alpha$ jest algebraiczny nad $C$.
Algebraiczna domkniętość $C$ daje $\alpha\in C$.
\end{proof}

\begin{przyklad}[Pochodna funkcji algebraicznej]
\label{aa:prz:derivative-algebraic-sqrt}
Niech $K=\mathbb C(x)$ z derywacją $Dx=1$, a
$\alpha^2=x^2+1$ w rozszerzeniu algebraicznym.
Wielomian $T^2-(x^2+1)$ ma nad $K$ stopień $2$
i jest nierozkładalny: $x^2+1$ ma w $\mathbb C(x)$
proste zera $i$ i $-i$, więc nie jest kwadratem
funkcji wymiernej. Różniczkując relację, otrzymujemy
$2\alpha D\alpha=2x$, czyli
$D\alpha=x/\alpha$. Jest to także wynik wzoru
z poprzedniego twierdzenia.
\end{przyklad}

\begin{przyklad}[Dlaczego domkniętość algebraiczna stałych jest potrzebna?]
\label{aa:prz:algebraic-new-constants}
Niech $K=\mathbb Q(x)$, $Dx=1$, $Dq=0$
dla $q\in\mathbb Q$. Z przykładu
\ref{aa:prz:differential-rational-field} mamy
$C_K=\mathbb Q$, które nie jest algebraicznie
domknięte. W rozszerzeniu
$L=\mathbb Q(i)(x)$ jedyne przedłużenie
derywacji spełnia $D(i)=0$, bo
$i^2+1=0$ i $2iD(i)=0$. Dlatego
$C_L=\mathbb Q(i)\ne\mathbb Q$.
\end{przyklad}

\section{Ślad i norma a derywacja}

\begin{twierdzenie}[Różniczkowanie śladu i normy]
\label{aa:tw:differential-trace-norm}
Niech $(K,D)$ będzie ciałem różniczkowym
charakterystyki $0$, a $L/K$ skończonym rozszerzeniem
algebraicznym z jedynym przedłużeniem $D$.
Dla każdego $v\in L$ i $u\in L^\times$ zachodzi
\[
D\!\left(\operatorname{Tr}_{L/K}(v)\right)
=\operatorname{Tr}_{L/K}(Dv),
\qquad
\frac{D\!\left(N_{L/K}(u)\right)}{N_{L/K}(u)}
=\operatorname{Tr}_{L/K}\!\left(\frac{Du}{u}\right).
\]
Nie zakładamy normalności $L/K$.
\end{twierdzenie}
\begin{proof}
Wybierzmy domknięcie algebraiczne $\overline K$
z jedynym przedłużeniem $D$ z poprzedniego
twierdzenia. Ponieważ $\operatorname{char}K=0$,
rozszerzenie $L/K$ jest rozdzielcze.
Niech $\sigma_1,\ldots,\sigma_d:L\to\overline K$
będą wszystkimi $K$-zanurzeniami, gdzie $d=[L:K]$.
Dla każdego $w\in L$ mamy
$D(\sigma_j(w))=\sigma_j(Dw)$: istotnie,
obie strony są wymuszone przez wielomian minimalny
elementu $w$ nad $K$ i wzór na pochodną jego
pierwiastka. Korzystając ze wzorów na ślad
i normę przez zanurzenia
(Twierdzenie~\ref{aa:tw:trace-norm-embeddings}),
otrzymujemy
\[
\begin{aligned}
D\operatorname{Tr}_{L/K}(v)
 &=D\left(\sum_{j=1}^d\sigma_j(v)\right)
  =\sum_{j=1}^d\sigma_j(Dv)
  =\operatorname{Tr}_{L/K}(Dv),\\
\frac{DN_{L/K}(u)}{N_{L/K}(u)}
 &=\frac{D\prod_{j=1}^d\sigma_j(u)}
         {\prod_{j=1}^d\sigma_j(u)}
  =\sum_{j=1}^d\frac{\sigma_j(Du)}{\sigma_j(u)}
  =\operatorname{Tr}_{L/K}\left(\frac{Du}{u}\right).
\end{aligned}
\]
Wszystkie czynniki normy są niezerowe,
gdyż zanurzenia ciał są różnowartościowe.
\end{proof}

\begin{przyklad}[Ślad i norma w rozszerzeniu kwadratowym]
\label{aa:prz:differential-quadratic-trace-norm}
W przykładzie \ref{aa:prz:derivative-algebraic-sqrt}
mamy $\alpha^2=x^2+1$ i
$\operatorname{Tr}_{L/K}(\alpha)=0$,
$N_{L/K}(\alpha)=-(x^2+1)$.
Pierwsza tożsamość daje
$\operatorname{Tr}_{L/K}(D\alpha)=D0=0$.
Druga ma po lewej stronie
$D(-(x^2+1))/(-(x^2+1))=2x/(x^2+1)$.
Po prawej, ponieważ $D\alpha/\alpha=x/(x^2+1)\in K$,
otrzymujemy
$\operatorname{Tr}_{L/K}(D\alpha/\alpha)=2x/(x^2+1)$.
W ten sposób obie tożsamości można sprawdzić bezpośrednio.
\end{przyklad}

\section{Przestępne dołączanie funkcji pierwotnych}

Niech odtąd $(K,D)$ ma charakterystykę $0$ i
$t$ będzie \emph{przestępny} nad $K$.
Podając $Dt=a\in K$, rozszerzamy $D$ jednoznacznie
na $K(t)$ przez formalną regułę iloczynu i ilorazu.
W następnym lemacie ustalimy dokładnie, kiedy
takie dołączenie zachowuje ciało stałych.

\begin{lemat}[Kryterium braku nowych stałych]
\label{aa:lem:primitive-extension-constants}
Niech $(K,D)$ będzie ciałem różniczkowym
charakterystyki $0$, $C=C_K$, a $t$ będzie
przestępny nad $K$ i spełniał $Dt=a\in K$.
Wtedy $C_{K(t)}=C$ wtedy i tylko wtedy,
gdy nie istnieje $b\in K$ spełniający $Db=a$.
\end{lemat}
\begin{proof}
Jeżeli $Db=a$, to $D(t-b)=0$, a
$t-b\notin K$ z powodu przestępności $t$.
Powstaje zatem nowa stała.

W drugim kierunku przypuśćmy, że
$R\in K(t)\setminus C$ i $DR=0$.
Funkcję wymierną $R$ można rozwinąć
jednoznacznie jako formalny szereg Laurenta
w nieskończoności
$R=\sum_{j\le m}r_jt^j$, gdzie $r_j\in K$
i $r_m\ne0$, a $m\in\mathbb Z$.
Derywacja jest na tym szeregu zgodna z
\[
D(r_jt^j)=(Dr_j)t^j+j\,a\,r_jt^{j-1}.
\]
Tożsamość zachodzi również dla ujemnych $j$,
a współczynnik każdej potęgi pochodzi ze
skończenie wielu wyrazów.

Z porównania współczynników przy $t^m$
w równości $DR=0$ otrzymujemy $Dr_m=0$,
czyli $r_m\in C^\times$.
Jeżeli $m=0$, odejmujemy stałą $r_0\in C$.
Otrzymany niezerowy szereg ma najwyższy
wykładnik ujemny (gdyby był zerowy,
$R$ od początku należałoby do $C$).
Możemy więc w każdym przypadku założyć $m\ne0$.
Współczynnik przy $t^{m-1}$ daje
$Dr_{m-1}+ma r_m=0$; jeżeli w rozwinięciu
nie ma wyrazu $t^{m-1}$, przyjmujemy
$r_{m-1}=0$. Wobec tego
$b=-r_{m-1}/(mr_m)\in K$ spełnia $Db=a$.
To kończy dowód przez kontrapozycję.
\end{proof}

\begin{przyklad}[Dołączenie logarytmu]
\label{aa:prz:logarithmic-extension}
Weźmy $K=\mathbb C(x)$, $Dx=1$ oraz
$a=1/x$. W $K$ nie ma elementu $b$
spełniającego $Db=1/x$: po rozkładzie
funkcji wymiernej na ułamki proste
pochodna każdego składnika ma przy
każdym skończonym biegunie resztę $0$,
natomiast $1/x$ ma w $0$ resztę $1$.
Dołączając przestępny element $\ell$ z
$D\ell=1/x$, otrzymujemy rozszerzenie
różniczkowe $K(\ell)/K$ bez nowych stałych.

Lokalnie analitycznie $\ell$ może być
zarodkiem wybranej gałęzi funkcji $\log x$.
Uzasadnijmy również jej przestępność.
Gdyby taki zarodek był algebraiczny nad $K$,
to dla $L=K(\ell)$ i $d=[L:K]$ poprzednie
twierdzenie dałoby
\[
D\!\left(\frac{\operatorname{Tr}_{L/K}(\ell)}d\right)
=\frac{\operatorname{Tr}_{L/K}(D\ell)}d
=\frac{\operatorname{Tr}_{L/K}(1/x)}d
=\frac1x.
\]
Po lewej byłaby pochodna funkcji wymiernej
z $\mathbb C(x)$, co przeczy sprawdzonej
właściwości jej reszt. Zatem wybrana
lokalna gałąź logarytmu rzeczywiście jest
przestępna nad $\mathbb C(x)$.
\end{przyklad}

\begin{przyklad}[Sztuczne dołączenie nowej stałej]
\label{aa:prz:new-constant-obstruction}
Niech $K=\mathbb C(x)$, $Dx=1$ i dołączmy
przestępny element $t$ z $Dt=1$.
Ponieważ $Dx=1$, element $t-x$
jest nową stałą w $K(t)$.
Nie wolno więc automatycznie nazywać
dowolnej przestępnej zmiennej spełniającej
równanie pierwszego rzędu rozszerzeniem
bez nowych stałych.
\end{przyklad}

\section{Rozszerzenia wykładnicze i elementarne}

\begin{lemat}[Stałe w rozszerzeniu wykładniczym]
\label{aa:lem:exponential-extension-constants}
Niech $(K,D)$ będzie ciałem różniczkowym
charakterystyki $0$, $C=C_K$, a $t$ będzie
przestępny nad $K$ i spełniał $Dt=ht$
dla pewnego $h\in K$. Wtedy $C_{K(t)}=C$
wtedy i tylko wtedy, gdy dla żadnej liczby
całkowitej $m\ne0$ i żadnego $u\in K^\times$
nie zachodzi $Du/u=-mh$.
\end{lemat}
\begin{proof}
Jeśli taka para $(m,u)$ istnieje,
to $D(ut^m)=0$, a $ut^m\notin K$
z powodu przestępności $t$. Jest więc
nową stałą.

Odwrotnie, przypuśćmy, że
$R\in K(t)\setminus C$ i $DR=0$.
Rozwijamy $R$ formalnie w nieskończoności
jako $\sum_{j\le m}r_jt^j$ z $r_m\ne0$.
Tym razem
$D(r_jt^j)=(Dr_j+jhr_j)t^j$.
Jeśli $m=0$, z równania współczynnika
$t^0$ otrzymujemy $Dr_0=0$.
Odejmujemy $r_0\in C$; pozostały niezerowy
szereg ma najwyższą potęgę $m<0$.
Zatem możemy przyjąć $m\ne0$.
Współczynnik przy $t^m$ daje
$Dr_m+mhr_m=0$,
czyli $Dr_m/r_m=-mh$.
Wybór $u=r_m\in K^\times$ spełnia
żądaną tożsamość.
\end{proof}

\begin{przyklad}[Dołączenie wykładniczej]
\label{aa:prz:exponential-extension}
W $K=\mathbb C(x)$, $Dx=1$ nie istnieje
niezerowa funkcja wymierna $u$ spełniająca
$Du/u=c$ dla niezerowej stałej $c\in\mathbb C$.
Rzeczywiście, rozkład licznika i mianownika
na czynniki liniowe daje
\[
\frac{Du}{u}
=\sum_{a\in\mathbb C}\frac{m_a}{x-a},
\qquad m_a\in\mathbb Z,
\]
gdzie suma ma skończoną liczbę niezerowych wyrazów.
Po prawej jest funkcja właściwa, a więc
nie może być niezerową stałą.

Jeśli przestępny element $t$ spełnia
$Dt=t$, poprzedni lemat daje
$C_{K(t)}=\mathbb C$.
Analitycznie możemy wybrać $t=e^x$.
Przestępność tego zarodka można też wykazać
algebraicznie: gdyby $t$ był algebraiczny
nad $K$, skończony stopień $d=[K(t):K]$
i twierdzenie o śladzie i normie dałyby
\[
\frac{DN_{K(t)/K}(t)}{N_{K(t)/K}(t)}
=\operatorname{Tr}_{K(t)/K}(Dt/t)
=\operatorname{Tr}_{K(t)/K}(1)=d,
\]
co jest niemożliwe dla niezerowej stałej $d$
i funkcji wymiernej $N(t)\in\mathbb C(x)^\times$.
\end{przyklad}

\begin{definicja}[Rozszerzenie elementarne bez nowych stałych]
\label{aa:def:elementary-differential-extension}
Niech $(K,D)$ będzie ciałem różniczkowym
charakterystyki $0$.
\emph{Rozszerzeniem elementarnym} $E/K$
nazywamy rozszerzenie różniczkowe, dla
którego istnieje skończona wieża
$K=K_0\subseteq K_1\subseteq\cdots\subseteq K_r=E$,
gdzie $K_j=K_{j-1}(\theta_j)$, a każdy
generator $\theta_j$ spełnia co najmniej
jeden z trzech warunków:
\begin{enumerate}
\item $\theta_j$ jest algebraiczny nad $K_{j-1}$;
\item $D\theta_j=Du/u$ dla pewnego
      $u\in K_{j-1}^\times$ (krok logarytmiczny);
\item $\theta_j\ne0$ i
      $D\theta_j/\theta_j=Dv$ dla pewnego
      $v\in K_{j-1}$ (krok wykładniczy).
\end{enumerate}
Jeżeli ponadto $C_E=C_K$, mówimy, że
rozszerzenie jest \emph{bez nowych stałych}.
Dla $f\in K$ mówimy, że $f$ ma
\emph{elementarną funkcję pierwotną bez nowych stałych},
jeśli istnieje takie $E/K$ i $F\in E$
spełniające $DF=f$.
\end{definicja}

\begin{uwaga}[Granice definicji]
\label{aa:uw:elementary-versus-liouvillian}
Krok logarytmiczny ma postać
$D\theta=Du/u$, a \emph{nie} dowolną postać
$D\theta=a\in K_{j-1}$. Dopuszczenie dowolnej
funkcji pierwotnej już w definicji
prowadzi do większej klasy rozszerzeń
\emph{liouvillowskich}, rozważanych w
różniczkowej teorii Galois.
Wybór lokalnej gałęzi logarytmu lub
pierwiastka nie oznacza, że te funkcje
muszą istnieć jako jednowartościowe
funkcje globalne na całej płaszczyźnie.
Elementarną całkowalność będziemy
badać algebraicznie, a interpretować
lokalnie na wspólnym obszarze
określoności odpowiednich funkcji.
\end{uwaga}

\begin{przyklad}[Automorfizmy dwóch rozszerzeń]
\label{aa:prz:elementary-automorphisms}
Niech $K=\mathbb C(x)$ z $Dx=1$.

Dla $L=K(t)$, gdzie $t$ jest przestępny i $Dt=t$,
każdy automorfizm różniczkowy $\sigma:L\to L$
nad $K$ spełnia
$D(\sigma(t)/t)=0$, więc
$\sigma(t)=ct$ dla pewnego $c\in\mathbb C^\times$,
ponieważ $C_L=\mathbb C$.
Każdy taki wybór odwrotnie wyznacza automorfizm
o odwrotności $t\mapsto c^{-1}t$.
Grupa tych automorfizmów jest więc
izomorficzna z $(\mathbb C^\times,\cdot)$.

Dla $M=K(\ell)$, gdzie $\ell$ jest przestępny
i $D\ell=1/x$, mamy $C_M=\mathbb C$.
Równanie $D(\sigma(\ell)-\ell)=0$ daje
$\sigma(\ell)=\ell+c$ dla $c\in\mathbb C$,
a wszystkie te przesunięcia są automorfizmami.
Ich grupa jest izomorficzna z $(\mathbb C,+)$.
W obu przypadkach warunek komutowania z
derywacją ogranicza automorfizmy zwykłego
rozszerzenia ciał do bardzo konkretnej rodziny.
\end{przyklad}

Równości dla śladu i normy oraz kryteria
braku nowych stałych są dokładnie tymi
narzędziami, które pozwolą w następnym
rozdziale sprowadzić problem całkowania
elementarnego z wieży rozszerzeń do
tożsamości w początkowym ciele różniczkowym.

% AA-44 | Twierdzenie Liouville'a, pełna konieczność i dostateczność.
% Podstawa: AA-43, sources/pl/galois(2).tex; AA-45 dopiero po audytach.
\chapter{Twierdzenie Liouville'a o całkowaniu elementarnym}
\chaptermark{Twierdzenie Liouville'a}
\label{aa:ch:liouville-elementary-integration}

Funkcja mająca funkcję pierwotną nie musi mieć jej w klasie funkcji
elementarnych. Aby uczynić to pytanie ścisłym, funkcje zastępujemy
elementami ciała różniczkowego, a wyrażenia elementarne opisujemy
skończoną wieżą rozszerzeń. Narzędzia konstrukcyjne są już dostępne
w rozdziale~\ref{aa:ch:differential-fields}: jednoznaczne
przedłużenie derywacji do rozszerzenia algebraicznego, pochodne
śladu i normy oraz kryteria braku nowych stałych. Teraz
udowodnimy, że cała wieża może zostać wyeliminowana z
\emph{kryterium istnienia} elementarnej funkcji pierwotnej.

\section{Sformułowanie problemu i rola stałych}

\begin{definicja}[Postać Liouville'a]
\label{aa:def:liouville-form}
Niech $(K,D)$ będzie ciałem różniczkowym charakterystyki $0$,
$C=C_K$, a $f\in K$. Mówimy, że $f$ ma
\emph{postać Liouville'a w $K$}, jeżeli istnieją
$n\ge0$, $v\in K$, $u_1,\ldots,u_n\in K^\times$ oraz
$c_1,\ldots,c_n\in C$ takie, że
\begin{equation}\label{aa:eq:liouville-form}
 f=Dv+\sum_{j=1}^{n}c_j\frac{Du_j}{u_j}.
\end{equation}
Dla $n=0$ suma jest pusta. Współczynniki $c_j$
są \emph{stałymi derywacji}, a nie dowolnymi elementami $K$.
\end{definicja}

Równość~\eqref{aa:eq:liouville-form} nie zawiera
symboli logarytmów, dlatego ma sens w dowolnym
ciele różniczkowym charakterystyki $0$.
Jeżeli wybierzemy lokalne gałęzie logarytmów,
formalnym kandydatem na funkcję pierwotną jest
$v+\sum_j c_j\log u_j$. Trzeba jednak sprawdzić,
czy dołączanie logarytmów nie wprowadza nowych stałych.

\begin{przyklad}[Dwa elementarne przedstawienia]
\label{aa:prz:liouville-rational}
W $K=\mathbb C(x)$ z $Dx=1$ mamy
\[
 \frac{x}{x^2+1}=\frac12\frac{D(x^2+1)}{x^2+1},
 \qquad
 \frac1{x^2+1}
 =\frac1{2i}\left(\frac{D(x-i)}{x-i}
                   -\frac{D(x+i)}{x+i}\right).
\]
Odpowiednimi funkcjami pierwotnymi są lokalne gałęzie
$\frac12\log(x^2+1)$ oraz
$\frac1{2i}\log((x-i)/(x+i))$.
W drugiej tożsamości sprawdzamy jawnie
$\frac1{2i}((x-i)^{-1}-(x+i)^{-1})=(x^2+1)^{-1}$.
W obu przypadkach stałe współczynniki należą do $\mathbb C$.
\end{przyklad}

\section{Formalna część wielomianowa i pochodne logarytmiczne}

\begin{definicja}[Część wielomianowa w nieskończoności]
\label{aa:def:liouville-polynomial-part}
Niech $F$ będzie ciałem i $t$ elementem przestępnym nad $F$.
Każdy $R\in F(t)$ ma jednoznaczne przedstawienie
$R=P+Q$, gdzie $P\in F[t]$, a $Q=0$ albo
$Q=A/B$ dla $A,B\in F[t]$, $B\ne0$ i $\deg A<\deg B$.
Definiujemy $\operatorname{pp}_t(R)=P$.
\end{definicja}

Istnienie wynika z dzielenia z resztą licznika przez mianownik.
Dla jednoznaczności zauważmy, że niezerowy wielomian nie może
być właściwą funkcją wymierną: z $P=A/B$ wynikałoby
$\deg(PB)\ge\deg B>\deg A$.

Równoważnie, rozwijając mianownik w potęgach $t^{-1}$,
otrzymujemy jednoznaczne formalne rozwinięcie
$R=\sum_{j\le m}r_jt^j$ z $r_j\in F$,
gdzie tylko skończenie wiele wykładników jest nieujemnych.
Wówczas $\operatorname{pp}_t(R)=\sum_{j\ge0}r_jt^j$.
Jest to tożsamość w ciele formalnych szeregów Laurenta
$F((t^{-1}))$, a nie twierdzenie o granicy funkcji
analitycznej w nieskończoności.

\begin{lemat}[Części wielomianowe pochodnych]
\label{aa:lem:liouville-polynomial-parts}
Niech $(F,D)$ będzie ciałem różniczkowym
charakterystyki $0$, a $t$ elementem przestępnym nad $F$.
Załóżmy jeden z przypadków: $(\mathrm P)$ $Dt=a\in F$
albo $(\mathrm E)$ $Dt=ht$ dla $h\in F$.
W obu przypadkach dla $R\in F(t)$ zachodzi
\[
 \operatorname{pp}_t(DR)=D(\operatorname{pp}_t(R)).
\]
Jeśli $u\in F(t)^\times$ i
$u=b\,p/q$, gdzie $b\in F^\times$, a $p,q\in F[t]$
są unormowane, to
\begin{equation}\label{aa:eq:liouville-log-pp}
 \operatorname{pp}_t\left(\frac{Du}{u}\right)
 =\begin{cases}
       Db/b,&(\mathrm P),\\
       Db/b+(\deg p-\deg q)h,&(\mathrm E).
   \end{cases}
\end{equation}
W szczególności pochodna logarytmiczna nie wnosi
dodatnich potęg $t$.
\end{lemat}
\begin{proof}
Rozkład $R=P+Q$ z poprzedniej definicji pozwala
rozważyć osobno wielomian $P$ i właściwą funkcję
wymierną $Q=A/B$. W obu przypadkach $DP\in F[t]$,
bo $D(t^j)=ja t^{j-1}$ w $(\mathrm P)$ oraz
$D(t^j)=jht^j$ w $(\mathrm E)$.
Jeżeli $\deg A<\deg B=d$, to w obu przypadkach
$\deg DA\le\deg A$ (dla wielomianu stałego możliwe
jest $DA=0$), a $\deg DB\le d$.
Stąd licznik $(DA)B-A(DB)$ ma stopień co najwyżej
$\deg A+d<2d$, więc $DQ$ jest właściwe.
Dla $Q=0$ wniosek jest natychmiastowy.
To dowodzi pierwszej równości bez jakichkolwiek
założeń o rozkładzie $B$ na czynniki liniowe.

W drugiej części wykorzystujemy równość
$Du/u=Db/b+Dp/p-Dq/q$. Jeśli $p$ jest
unormowany stopnia $d$, to w przypadku $(\mathrm P)$
współczynnik przy $t^d$ w $Dp$ jest zerowy,
więc $\deg Dp<d$. W przypadku $(\mathrm E)$
współczynnik przy $t^d$ w $Dp$ wynosi $dh$,
więc $(Dp)/p-dh$ jest właściwą funkcją wymierną.
Analogiczne zdania zachodzą dla $q$; gdy
$p=1$ lub $q=1$, otrzymujemy wkład zerowy.
Część wielomianowa jest addytywna, skąd
wynika~\eqref{aa:eq:liouville-log-pp}.
\end{proof}

\begin{uwaga}[Dlaczego nie porównujemy dowolnych potęg Laurenta?]
\label{aa:uw:liouville-negative-powers}
Gdy $Dt=ht$, pochodna logarytmiczna
$D(t-1)/(t-1)=ht/(t-1)$ ma formalne rozwinięcie
$h+ht^{-1}+ht^{-2}+\cdots$. Nie wolno zatem
twierdzić, że pochodne logarytmiczne nie mają
ujemnych potęg. Udowodniony lemat dotyczy wyłącznie
\emph{części wielomianowej}, co dokładnie wystarczy
w dowodzie twierdzenia Liouville'a.
\end{uwaga}

\section{Trzy kroki schodzenia po wieży}

Zajmiemy się następującą sytuacją. Element $f$
należy do ciała niższego piętra $F$, natomiast
jego przedstawienie Liouville'a jest na razie znane
tylko w jednym rozszerzeniu $L/F$. Pokażemy, że
można je przenieść do $F$ bez zmiany ciała stałych.

\begin{lemat}[Schodzenie przez rozszerzenie algebraiczne]
\label{aa:lem:liouville-algebraic-descent}
Niech $(F,D)$ będzie ciałem różniczkowym
charakterystyki $0$ o ciele stałych $C$, a
$L/F$ rozszerzeniem skończonym algebraicznym
z przedłużoną derywacją. Jeśli $f\in F$ oraz
\[
 f=DV+\sum_{j=1}^{n}c_j\frac{DU_j}{U_j},
 \quad V\in L,\quad U_j\in L^\times,\quad c_j\in C,
\]
to $f$ ma postać Liouville'a w $F$.
\end{lemat}
\begin{proof}
Niech $d=[L:F]$. Ponieważ $\operatorname{char}F=0$,
liczba $d$ jest odwracalna w $C$.
Stosujemy ślad do obu stron tożsamości.
Ponieważ $f\in F$, mamy
$\operatorname{Tr}_{L/F}(f)=df$.
Na mocy twierdzenia~\ref{aa:tw:differential-trace-norm}
otrzymujemy
\[
 f=D\left(\frac{\operatorname{Tr}_{L/F}(V)}d\right)
   +\sum_{j=1}^{n}\frac{c_j}{d}
       \frac{D N_{L/F}(U_j)}{N_{L/F}(U_j)}.
\]
Ślad należy do $F$, a normy należą do $F^\times$.
Nowe współczynniki $c_j/d$ należą do $C$.
Nie jest potrzebne założenie normalności $L/F$.
\end{proof}

\begin{lemat}[Schodzenie przez krok logarytmiczny]
\label{aa:lem:liouville-log-descent}
Niech $(F,D)$ ma charakterystykę $0$ i ciało
stałych $C$. Niech $L=F(t)$, gdzie $t$ jest
przestępny nad $F$, $C_L=C$ oraz
$Dt=Db/b$ dla pewnego $b\in F^\times$.
Jeśli $f\in F$ ma postać Liouville'a w $L$
ze współczynnikami w $C$, to ma ją w $F$.
\end{lemat}
\begin{proof}
Zapiszmy
$f=DV+\sum_j c_j DU_j/U_j$ z $V\in F(t)$
i $U_j\in F(t)^\times$. W lemacie
\ref{aa:lem:liouville-polynomial-parts}
wybieramy przypadek $(\mathrm P)$ z $a=Db/b$.
Każde $U_j$ zapisujemy w postaci
$U_j=b_j p_j/q_j$, gdzie
$b_j\in F^\times$, a $p_j,q_j$ są unormowane.
Ponieważ $f\in F$, zastosowanie
$\operatorname{pp}_t$ daje
\begin{equation}\label{aa:eq:liouville-log-descent}
 f=DP+\sum_j c_j\frac{Db_j}{b_j},
 \quad P=\operatorname{pp}_t(V)\in F[t].
\end{equation}
W szczególności $DP\in F$.

Jeśli $P$ jest stały jako wielomian w $t$,
to~\eqref{aa:eq:liouville-log-descent} jest
już wymaganym przedstawieniem. W przeciwnym
razie niech $m=\deg_tP\ge1$ i
$P=p_mt^m+\cdots+p_0$, $p_m\ne0$.
Współczynnik przy $t^m$ w $DP$ wynosi
$Dp_m$; musi zatem być zerowy, a więc
$p_m\in C^\times$.
Gdyby $m\ge2$, współczynnik przy $t^{m-1}$
byłby równy $Dp_{m-1}+mp_m Db/b=0$.
Element $w=-p_{m-1}/(mp_m)\in F$
spełniałby $Dw=Db/b=Dt$.
Stąd $D(t-w)=0$, więc $t-w\in C_L=C\subseteq F$,
sprzecznie z przestępnością $t$ nad $F$.
Pozostaje $m=1$, zatem $P=p_0+ct$, $c\in C$.
Podstawiając do~\eqref{aa:eq:liouville-log-descent},
otrzymujemy
\[
 f=Dp_0+c\,\frac{Db}{b}
              +\sum_j c_j\frac{Db_j}{b_j}.
\]
Każdy składnik jest już zdefiniowany nad $F$.
\end{proof}

\begin{lemat}[Schodzenie przez krok wykładniczy]
\label{aa:lem:liouville-exp-descent}
Niech $(F,D)$ ma charakterystykę $0$ i ciało
stałych $C$. Niech $L=F(t)$, gdzie $t$ jest
przestępny nad $F$, $C_L=C$, $t\ne0$ oraz
$Dt/t=Da$ dla pewnego $a\in F$.
Jeśli $f\in F$ ma postać Liouville'a w $L$
ze współczynnikami w $C$, to ma ją w $F$.
\end{lemat}
\begin{proof}
Piszemy $f=DV+\sum_jc_j DU_j/U_j$ i
$U_j=b_jp_j/q_j$ z $b_j\in F^\times$ oraz
unormowanymi wielomianami $p_j,q_j$.
Lemat~\ref{aa:lem:liouville-polynomial-parts},
tym razem w przypadku $(\mathrm E)$ z $h=Da$,
daje
\begin{equation}\label{aa:eq:liouville-exp-descent}
 f=DP+\sum_jc_j\frac{Db_j}{b_j}+A\,Da,
 \qquad
 A=\sum_jc_j(\deg p_j-\deg q_j)\in C,
\end{equation}
gdzie $P=\operatorname{pp}_t(V)=\sum_{k=0}^{m}p_kt^k$.
Po prawej poza $DP$ wszystkie składniki należą do $F$,
więc $DP\in F$. Dla każdego $k\ge1$
współczynnik przy $t^k$ w $DP$ wynosi
$Dp_k+k(Da)p_k$ i musi być zerowy.
Jeśli $p_k\ne0$, to
$D(p_kt^k)=(Dp_k+k(Da)p_k)t^k=0$.
Wobec $C_L=C$ mamy $p_kt^k\in C^\times$,
czyli $t^k\in F$, co przeczy przestępności $t$.
Zatem $p_k=0$ dla wszystkich $k\ge1$,
a $P=p_0\in F$. Ponieważ $A\in C$,
mamy $A\,Da=D(Aa)$. Z~\eqref{aa:eq:liouville-exp-descent}
wynika więc
\[
 f=D(p_0+Aa)+\sum_jc_j\frac{Db_j}{b_j},
\]
czyli przedstawienie nad $F$.
\end{proof}

\section{Twierdzenie Liouville'a}

\begin{twierdzenie}[Liouville'a o całkowaniu elementarnym]
\label{aa:tw:liouville-elementary-integration}
Niech $(K,D)$ będzie ciałem różniczkowym
charakterystyki $0$ o ciele stałych $C=C_K$
i niech $f\in K$.
Następujące warunki są równoważne:
\begin{enumerate}
\item istnieje rozszerzenie elementarne
$E/K$ bez nowych stałych w sensie
Definicji~\ref{aa:def:elementary-differential-extension}
oraz element $F\in E$ taki, że $DF=f$;
\item $f$ ma postać Liouville'a w $K$, to znaczy
istnieją $n\ge0$, $v\in K$, $u_1,\ldots,u_n\in K^\times$
i $c_1,\ldots,c_n\in C$ spełniające
\eqref{aa:eq:liouville-form}.
\end{enumerate}
W szczególności nie trzeba przechodzić do rozszerzenia
algebraicznego, aby zapisać elementy $v,u_j$.
Nie zakładamy algebraicznej domkniętości $C$:
istotnym warunkiem jest \emph{brak nowych stałych}.
\end{twierdzenie}
\begin{proof}
\emph{Konieczność.} Załóżmy (1) i wybierzmy wieżę
elementarną $K=K_0\subseteq K_1\subseteq\cdots
\subseteq K_r=E$. Ponieważ
$C_K\subseteq C_{K_j}\subseteq C_E=C_K$,
wszystkie piętra mają to samo ciało stałych $C$.
Na najwyższym piętrze mamy przedstawienie
$f=DF$ odpowiadające~\eqref{aa:eq:liouville-form}
z $n=0$.

Przypuśćmy indukcyjnie, że $f$ ma postać
Liouville'a w $K_j$ ze współczynnikami w $C$.
Jeżeli generator $\theta_j$ jest algebraiczny
nad $K_{j-1}$, stosujemy lemat
\ref{aa:lem:liouville-algebraic-descent}.
Jeżeli generator jest przestępny i
$D\theta_j=Du/u$, stosujemy lemat
\ref{aa:lem:liouville-log-descent}.
Jeżeli generator jest przestępny,
$\theta_j\ne0$ i $D\theta_j/\theta_j=Dv$,
stosujemy lemat~\ref{aa:lem:liouville-exp-descent}.
Jeśli generator algebraiczny spełnia również
któryś z pozostałych warunków, używamy kroku
algebraicznego: nie ma potrzeby traktować
tych przypadków oddzielnie.
W każdym wariancie otrzymujemy przedstawienie
z elementami $K_{j-1}$ i współczynnikami w $C$.
Indukcja malejąca po $j=r,r-1,\ldots,1$
prowadzi do przedstawienia w $K$, co daje (2).

\emph{Dostateczność.} Załóżmy (2), zatem
$f=Dv+\sum_{j=1}^n c_j Du_j/u_j$.
Zbudujemy skończoną wieżę
$K=E_0\subseteq E_1\subseteq\cdots\subseteq E_n$,
zachowując ciało stałych $C$.
Na etapie $j$ kładziemy $a_j=Du_j/u_j\in K
\subseteq E_{j-1}$. Jeśli istnieje
$w_j\in E_{j-1}$ takie, że $Dw_j=a_j$,
nie powiększamy ciała, kładąc $E_j=E_{j-1}$.
Jeżeli taka funkcja pierwotna nie istnieje,
dołączamy \emph{przestępny} element $w_j$
z derywacją $Dw_j=a_j$ i kładziemy
$E_j=E_{j-1}(w_j)$. Jest to krok
logarytmiczny, bo $u_j\in E_{j-1}^\times$.
Lemat~\ref{aa:lem:primitive-extension-constants}
stwierdza, że w tej drugiej sytuacji
$C_{E_j}=C_{E_{j-1}}=C$. W pierwszej jest
to oczywiste.

Po wszystkich krokach $E=E_n$ jest
rozszerzeniem elementarnym bez nowych stałych.
Element $F=v+\sum_{j=1}^n c_jw_j\in E$
spełnia
$DF=Dv+\sum_jc_jDw_j
=Dv+\sum_jc_j Du_j/u_j=f$.
Dla $n=0$ wystarcza $E=K$ i $F=v$.
Dowodzi to (1).
\end{proof}

\begin{uwaga}[Co dokładnie zostało udowodnione?]
\label{aa:uw:liouville-scope}
Twierdzenie dotyczy istnienia funkcji pierwotnej
w \emph{elementarnym rozszerzeniu bez nowych stałych}
ustalonego ciała różniczkowego. Przy interpretacji
jako zarodków funkcji meromorficznych nad
$\mathbb C$ warunek braku nowych stałych jest
naturalny; gałęzie funkcji algebraicznych i
logarytmów wybieramy lokalnie. Jeżeli zezwalamy
na dowolne powiększanie ciała stałych, przed
zastosowaniem twierdzenia należy osobno uzasadnić
odpowiednią redukcję. Nie jest ona ukrytym
elementem niniejszego dowodu.
\end{uwaga}

\section{Kontrola twierdzenia na konkretnych przykładach}

\begin{przyklad}[Norma usuwa algebraiczny pierwiastek]
\label{aa:prz:liouville-norm-sqrt}
Weźmy $K=\mathbb C(x)$, $\alpha^2=x^2+1$,
$L=K(\alpha)$. Z rozdziału
\ref{aa:ch:differential-fields} wiemy, że
$D\alpha=x/\alpha$, zatem
\[
 \frac{D\alpha}{\alpha}
 =\frac{x}{x^2+1}\in K.
\]
W $L$ prawa strona jest pojedynczą
pochodną logarytmiczną.
Ponieważ $N_{L/K}(\alpha)=-(x^2+1)$,
uśrednienie śladu daje w $K$
\[
 \frac{x}{x^2+1}
 =\frac12\frac{D(-(x^2+1))}{-(x^2+1)}
 =\frac12\frac{D(x^2+1)}{x^2+1}.
\]
Znak normy nie wpływa na pochodną logarytmiczną.
Przykład ilustruje rzeczywistą zmianę elementu
$u_j$ przy schodzeniu po kroku algebraicznym.
\end{przyklad}

\begin{przyklad}[Krok logarytmiczny]
\label{aa:prz:liouville-log-step}
Niech $K=\mathbb C(x)$, $Dx=1$, a $t$ będzie
przestępny z $Dt=1/x$.
Fakt, że $K(t)$ nie ma nowych stałych, wynika
z rozkładu funkcji wymiernych na ułamki proste
i lematu~\ref{aa:lem:primitive-extension-constants}.
Dla $V=t^2+x$ zachodzi $DV=2t/x+1$,
co nie należy do $K$; człon $t^2$ nie może
zostać ukryty w pochodnych logarytmicznych,
gdyż ich części wielomianowe są stałe w $t$.
Natomiast dla $V=3t+x^2$
mamy $DV=3/x+2x\in K$ i twierdzenie
sprowadza $DV$ do postaci
$D(x^2)+3\,Dx/x$.
W ten sposób obliczenie współczynników
przy potęgach $t$ wyjaśnia oba przypadki.
\end{przyklad}

\begin{przyklad}[Krok wykładniczy]
\label{aa:prz:liouville-exp-step}
Niech $K=\mathbb C(x)$, $Dx=1$ oraz
$t=e^x$ rozumiane jako przestępny generator
z $Dt=t$. Brak nowych stałych wykazano
w przykładzie~\ref{aa:prz:exponential-extension}.
Dla $u=xt^2/(t+1)$ otrzymujemy
\[
 \frac{Du}{u}=\frac1x+2-\frac{t}{t+1}.
\]
Ponieważ $t/(t+1)=1-1/(t+1)$,
część wielomianowa tego wyrażenia wynosi
$1/x+1$, zgodnie ze wzorem
\eqref{aa:eq:liouville-log-pp}:
$\deg(xt^2)-\deg(t+1)=1$ po wyłączeniu $x$.
W lemacie wykładniczym otrzymana
stała $1$ jest dokładną pochodną $Dx$.
\end{przyklad}

\begin{przyklad}[Dlaczego zabronione są nowe stałe?]
\label{aa:prz:liouville-new-constants}
Jeżeli do $K=\mathbb C(x)$ dołączymy
przestępny element $t$ z $Dt=1=Dx$,
to $D(t-x)=0$ mimo $t-x\notin K$.
Przykład ten nie stanowi kroku
logarytmicznego bez nowych stałych:
kryterium z AA-43 jest naruszone.
W szczególności argument, że
$D(p_mt^m)=0$ zmusza $p_mt^m$ do należenia
do pierwotnego ciała stałych, byłby wtedy błędny.
\end{przyklad}

\section{Granice wyniku i następny krok}

Twierdzenie Liouville'a sprowadza problem
całkowalności elementarnej do znalezienia
elementów $v,u_j\in K$ i stałych $c_j$.
Samo w sobie nie rozstrzyga jeszcze,
czy dane $f$ posiada takie przedstawienie.
W szczególności nie wolno przyjąć dowolnej
funkcji pierwotnej w kroku elementarnym,
ponieważ wtedy problem stałby się tautologiczny.
W kolejnym rozdziale zastosujemy uzyskane
kryterium do funkcji hiperwykładniczych:
porównując współczynniki przy dodatnich
potęgach przestępnego generatora,
sprowadzimy pytanie o elementarność konkretnych
całek do równań różniczkowych w $\mathbb C(x)$.

% AA-45 | Kryteria hiperwykładnicze i całki nieelementarne.
% Źródła: AA-43, AA-44, sources/pl/galois(2).tex.
% Stan: nowe opracowanie, przed niezależnymi audytami i finalną kontrolą PDF.
\chapter{Kryteria hiperwykładnicze i całki nieelementarne}
\chaptermark{Całki nieelementarne}
\label{aa:ch:nonelementary-integrals}

W rozdziale~\ref{aa:ch:liouville-elementary-integration}
udowodniliśmy twierdzenie Liouville'a: element $f$ ciała
różniczkowego ma elementarną funkcję pierwotną bez nowych stałych
wtedy i tylko wtedy, gdy jest sumą pochodnej elementu tego ciała
i skończonej kombinacji pochodnych logarytmicznych
o stałych współczynnikach. Daje to kryterium strukturalne,
które trzeba przekształcić w równania poddające się badaniu.
Najpierw zrobimy to dla rozszerzeń wykładniczych,
a potem przeanalizujemy cztery różne całki nieelementarne.

\section{Funkcje elementarne i ich lokalna interpretacja}

Przy funkcjach zespolonych używamy zarodków funkcji
meromorficznych w punkcie, w którym skończona liczba
rozważanych wyrażeń jest określona; dla pierwiastków
i logarytmów wybieramy lokalne gałęzie. Ciało takich
zarodków zawierające wszystkie stałe zespolone ma
dokładnie $\mathbb C$ jako ciało stałych:
zarodek o zerowej pochodnej jest stały.
Dlatego dowód nieelementarności w rozszerzeniu
bez nowych stałych dotyczy również zwykłych
lokalnych wzorów elementarnych.

Z rzeczywistego wzoru elementarnego na otwartym
przedziale można w zwykłym punkcie otrzymać
lokalny wzór zespolony, wybierając holomorficzne
gałęzie jego składników. Taką interpretację
stosujemy poniżej: wykluczenie lokalnej
zespolonej funkcji pierwotnej elementarnej
wyklucza wzór na rzeczywistym przedziale.
Nie twierdzimy, że logarytmy i pierwiastki
są globalnie jednowartościowe.

\section{Elementy hiperwykładnicze}

\begin{definicja}[Element hiperwykładniczy]
\label{aa:def:hyperexponential}
Niech $(F,D)$ będzie ciałem różniczkowym
charakterystyki $0$. Niezerowy element $t$
rozszerzenia różniczkowego $L/F$ nazywamy
\emph{hiperwykładniczym nad $F$}, jeśli
$Dt/t=h$ dla pewnego $h\in F$.
W kryteriach poniżej zakładamy dodatkowo,
że $t$ jest przestępny nad $F$ i
$C_{F(t)}=C_F$. Te dwa założenia są
istotne, a nie są częścią samego równania $Dt=ht$.
\end{definicja}

\begin{lemat}[Wykładnicze wielomianów]
\label{aa:lem:exp-polynomial-constant-field}
Niech $p\in\mathbb C[x]$ będzie niestałym
wielomianem i niech $K_0=\mathbb C(x)$
z $Dx=1$. Wówczas element $t=e^{p(x)}$,
rozumiany jako zarodek analityczny,
jest przestępny nad $K_0$, spełnia
$Dt=p'(x)t$, a $C_{K_0(t)}=\mathbb C$.
\end{lemat}
\begin{proof}
Każdy $u\in\mathbb C(x)^\times$ ma,
po rozkładzie licznika i mianownika,
postać $u=c\prod_{\nu=1}^{N}(x-a_\nu)^{m_\nu}$,
gdzie $c\in\mathbb C^\times$, $a_\nu\in\mathbb C$
są parami różne, a $m_\nu\in\mathbb Z$.
Stąd
\[
 \frac{Du}{u}=\sum_{\nu=1}^N\frac{m_\nu}{x-a_\nu}.
\]
Ta pochodna logarytmiczna jest właściwą
funkcją wymierną: jej część wielomianowa
w $x$ jest zerowa. Nie może więc być
niezerowym wielomianem.

Przypuśćmy, że $t$ jest algebraiczny nad
$K_0$ i niech $d=[K_0(t):K_0]$.
Stosując tożsamość dla normy z
Twierdzenia~\ref{aa:tw:differential-trace-norm}
do elementu $t$, otrzymujemy
\[
 \frac{D N_{K_0(t)/K_0}(t)}
        {N_{K_0(t)/K_0}(t)}
 =\operatorname{Tr}_{K_0(t)/K_0}(Dt/t)
 =d\,p'(x).
\]
Po lewej jest pochodna logarytmiczna
niezerowej funkcji wymiernej,
a po prawej niezerowy wielomian:
$p$ jest niestały i charakterystyka wynosi $0$.
Otrzymaliśmy sprzeczność.

Pozostaje zbadać stałe. Na mocy
Lematu~\ref{aa:lem:exponential-extension-constants}
nowa stała w $\mathbb C(x,t)$ istnieje
dokładnie wtedy, gdy
$Du/u=-m p'(x)$ dla pewnego
$u\in\mathbb C(x)^\times$ i
$m\in\mathbb Z\setminus\{0\}$.
Prawa strona jest niezerowym wielomianem,
a lewa właściwą funkcją wymierną.
Takiej pary nie ma, czyli
$C_{\mathbb C(x,t)}=\mathbb C$.
\end{proof}

\begin{uwaga}[Rola założeń]
\label{aa:uw:hyperexponential-obstructions}
Równanie $Dt=ht$ nie gwarantuje samo
przez się zachowania stałych.
Jeśli $K_0=\mathbb C(x)$, $h=1/x$
i $t$ jest przestępny z $Dt=t/x$,
to $D(t/x)=0$, a $t/x\notin K_0$.
Również przy $h=0$ przestępny generator
mający $Dt=0$ jest nową stałą.
Z kolei $e^{p(x)}$ z $p$ stałym
już należy do $\mathbb C$.
\end{uwaga}

\section{Pełne kryterium dla skończonych sum wykładniczych}

\begin{twierdzenie}[Kryterium hiperwykładnicze]
\label{aa:tw:hyperexponential-criterion}
Niech $h\in\mathbb C(x)$, a $t$ będzie
przestępny nad $\mathbb C(x)$ i spełnia
$Dt=ht$ oraz $C_{\mathbb C(x,t)}=\mathbb C$.
Dla $s\in\mathbb C(x)$ następujące warunki
są równoważne:
\begin{enumerate}
\item element $st$ ma elementarną funkcję
pierwotną nad $\mathbb C(x,t)$, bez nowych stałych;
\item istnieje $r\in\mathbb C(x)$ takie, że
\begin{equation}\label{aa:eq:hyperexp-ode}
 r'+hr=s.
\end{equation}
\end{enumerate}
Jeżeli te warunki zachodzą, element $rt$
jest funkcją pierwotną $st$.
\end{twierdzenie}
\begin{proof}
Jeżeli $r'+hr=s$, to
$D(rt)=(r'+hr)t=st$,
co daje elementarną funkcję pierwotną
w samym $\mathbb C(x,t)$.

Załóżmy odwrotnie istnienie elementarnej
funkcji pierwotnej. Z twierdzenia
\ref{aa:tw:liouville-elementary-integration}
dla ciała $\mathbb C(x,t)$ otrzymujemy
\[
 st=Dv+\sum_{j=1}^N c_j\frac{Du_j}{u_j},
 \qquad v\in\mathbb C(x,t),\quad
 u_j\in\mathbb C(x,t)^\times,\quad
 c_j\in\mathbb C.
\]
Niech $P=\operatorname{pp}_t(v)\in\mathbb C(x)[t]$.
Lemat~\ref{aa:lem:liouville-polynomial-parts}
daje po zastosowaniu $\operatorname{pp}_t$:
$st=DP+b(x)$, gdzie $b\in\mathbb C(x)$;
wszystkie pochodne logarytmiczne wnoszą
jedynie wyrazy niezależne od $t$.
Oznaczmy przez $r$ współczynnik przy
$t$ w $P$, przyjmując $r=0$, jeśli go nie ma.
Współczynnik przy $t$ w $DP$ wynosi
$r'+hr$. Przestępność $t$ zapewnia
jednoznaczność porównania współczynników,
więc otrzymujemy~\eqref{aa:eq:hyperexp-ode}.
\end{proof}

\begin{twierdzenie}[Kryterium dla wielomianów Laurenta]
\label{aa:tw:laurent-hyperexp-criterion}
Przy założeniach
Twierdzenia~\ref{aa:tw:hyperexponential-criterion},
niech $S\subset\mathbb Z$ będzie zbiorem
skończonym, a dla $k\in S$ niech
$s_k\in\mathbb C(x)$. Element
$f=\sum_{k\in S}s_k(x)t^k\in\mathbb C(x,t)$
ma elementarną funkcję pierwotną bez nowych
stałych wtedy i tylko wtedy, gdy dla
każdego $k\in S\setminus\{0\}$ istnieje
$r_k\in\mathbb C(x)$ spełniające
\begin{equation}\label{aa:eq:laurent-ode}
 r_k'+khr_k=s_k.
\end{equation}
Jeśli $0\notin S$, przyjmujemy $s_0=0$.
Współczynnik $s_0$ nie nakłada dodatkowego
warunku: każda funkcja wymierna nad
$\mathbb C$ ma elementarną funkcję pierwotną.
\end{twierdzenie}
\begin{proof}
Jeśli istnieją wszystkie $r_k$,
to $D(r_kt^k)=(r_k'+khr_k)t^k$.
Pozostałą funkcję wymierną $s_0$ całkujemy
elementarnie: po rozkładzie na ułamki proste
nad $\mathbb C$ składniki wielomianowe
i bieguny rzędu co najmniej $2$ mają
wymierne funkcje pierwotne, a bieguny proste
dają stałe razy logarytmy. Te ostatnie
można dołączać bez nowych stałych zgodnie
z dowodem dostateczności twierdzenia
\ref{aa:tw:liouville-elementary-integration}.
To dowodzi wystarczalności.

Dla konieczności stosujemy twierdzenie
Liouville'a do $f$. Dla ustalonego $k>0$
porównujemy współczynniki przy $t^k$
w częściach wielomianowych po obu stronach
przedstawienia Liouville'a. Na mocy
Lematu~\ref{aa:lem:liouville-polynomial-parts}
pochodne logarytmiczne nie mają
dodatnich potęg $t$, dlatego
współczynnik $r_k\in\mathbb C(x)$ przy
$t^k$ w $\operatorname{pp}_t(v)$
spełnia~\eqref{aa:eq:laurent-ode}.
Współczynnik może być zerowy, jeśli
brak odpowiedniego wyrazu.

Dla ustalonego $k<0$ wprowadzamy
$w=t^{-1}$. Wtedy
$\mathbb C(x,w)=\mathbb C(x,t)$,
$Dw=-hw$ i $w$ jest przestępny.
Zastosowanie poprzedniego argumentu
do współczynnika przy $w^{-k}$ daje
$r_k'+(-k)(-h)r_k=s_k$, czyli
\eqref{aa:eq:laurent-ode}.
Nie twierdzimy, że pochodne logarytmiczne
nie mają ujemnych potęg w rozwinięciu
względem $t^{-1}$; dlatego osobna zmiana
generatora jest konieczna.
\end{proof}

\begin{przyklad}[Wzory elementarne i ich test]
\label{aa:prz:hyperexp-positive-controls}
Dla $t=e^{-x^2}$ mamy $h=-2x$.
Równanie $r'-2xr=x$ ma rozwiązanie
$r=-1/2$, a więc
$\int x e^{-x^2}\,\mathrm dx=-\frac12e^{-x^2}+C$.
Dla $t=e^{x^2}$ równanie
$r'+2xr=x$ ma rozwiązanie $r=1/2$.
Również
\[
 D\left(\frac{e^x}{x}\right)
 =e^x\left(\frac1x-\frac1{x^2}\right)
\]
jest przykładem dodatnim mimo biegunów
współczynników. Samo występowanie
wykładniczej nie przesądza zatem
o nieelementarności funkcji pierwotnej.
\end{przyklad}

\section{Przeszkody wymierne: bieguny i stopnie}

\begin{lemat}[Przeszkoda bieguna prostego]
\label{aa:lem:simple-pole-obstruction}
Niech $h,s\in\mathbb C(x)$ i niech
$a\in\mathbb C$. Załóżmy, że $h$
nie ma bieguna w $a$, natomiast $s$
ma w $a$ biegun dokładnie rzędu $1$.
Wtedy równanie $r'+hr=s$
nie ma rozwiązania $r\in\mathbb C(x)$.
\end{lemat}
\begin{proof}
Jeśli $r$ jest regularna w $a$,
to $r'+hr$ też jest regularna,
wbrew założeniu o $s$.
Jeśli $r$ ma biegun rzędu $m\ge1$,
to w lokalnej współrzędnej $z=x-a$
jej rozwinięcie zaczyna się od
$c z^{-m}$, $c\ne0$. Pochodna
$r'$ ma biegun rzędu $m+1$ z niezerowym
współczynnikiem $-mc$, gdyż
charakterystyka wynosi $0$.
Ponieważ $h$ jest regularna, $hr$
ma biegun rzędu co najwyżej $m$.
Nie może więc skasować wyrazu
$-mc z^{-m-1}$. Lewa strona ma
biegun rzędu $m+1\ge2$, sprzeczność.
\end{proof}

\begin{lemat}[Przeszkoda stopnia dla wielomianów]
\label{aa:lem:polynomial-degree-obstruction}
Niech $h,s\in\mathbb C[x]$,
$h\ne0$ i $\deg h\ge1$. Jeśli
$r\in\mathbb C(x)$ spełnia
$r'+hr=s$, to $r\in\mathbb C[x]$.
Jeżeli $s\ne0$ i $\deg s<\deg h$,
równanie nie ma rozwiązania wymiernego.
\end{lemat}
\begin{proof}
Gdyby $r$ miała skończony biegun
rzędu $m\ge1$, pochodna $r'$
miałaby biegun rzędu $m+1$,
a $hr$ co najwyżej $m$,
bo $h$ jest wielomianem.
Nie mogłyby się skasować,
tymczasem $s$ nie ma biegunów.
Zatem $r$ nie ma biegunów skończonych
i jako funkcja wymierna jest wielomianem.

Jeżeli $r$ jest niezerowym wielomianem
stopnia $d\ge0$, to $hr$ ma stopień
$d+\deg h$, podczas gdy $r'$
ma stopień co najwyżej $d-1$
(lub jest zerem). Stąd
$\deg(r'+hr)=d+\deg h\ge\deg h$.
Dla $\deg s<\deg h$ otrzymujemy
sprzeczność. Element $r=0$ nie
rozwiązuje równania z $s\ne0$.
\end{proof}

\begin{uwaga}[Założenie regularności $h$]
\label{aa:uw:pole-obstruction-edge}
W Lemacie~\ref{aa:lem:simple-pole-obstruction}
nie można opuścić założenia o regularności $h$.
Dla $r=1/x$ i $h=1/x$ zachodzi
$r'+hr=-1/x^2+1/x^2=0$:
dwa bieguny rzędu $2$ mogą się skasować
wtedy, gdy $h$ sama ma biegun w $0$.
\end{uwaga}

\section{Całka Gaussa}

\begin{twierdzenie}[Nieelementarność całki Gaussa]
\label{aa:tw:gaussian-nonelementary}
Funkcja $e^{-x^2}$ nie ma elementarnej
funkcji pierwotnej na żadnym niepustym
otwartym przedziale rzeczywistym.
Algebraicznie element $t$ z $Dt=-2xt$
nie ma elementarnej funkcji pierwotnej
nad $\mathbb C(x,t)$ bez nowych stałych.
\end{twierdzenie}
\begin{proof}
Lemat~\ref{aa:lem:exp-polynomial-constant-field}
dla $p=-x^2$ uzasadnia przestępność
$t=e^{-x^2}$ oraz brak nowych stałych.
Twierdzenie~\ref{aa:tw:hyperexponential-criterion}
z $s=1$ i $h=-2x$ wymagałoby
rozwiązania $r'-2xr=1$
w $\mathbb C(x)$. Lemat
\ref{aa:lem:polynomial-degree-obstruction}
wyklucza takie rozwiązanie, gdyż
$\deg 1=0<1=\deg(-2x)$.
Nieelementarność lokalnego wzoru
zespolonego wyklucza też wzór
na dowolnym przedziale rzeczywistym.
\end{proof}

\begin{przyklad}[Funkcja błędu i całka oznaczona]
\label{aa:prz:gaussian-erf-value}
Brak elementarnej funkcji pierwotnej nie oznacza,
że funkcja pierwotna nie istnieje.
Można zdefiniować funkcję
$\operatorname{erf}(x)=
(2/\sqrt{\pi})\int_0^x e^{-u^2}\,\mathrm du$.
Jest ona różniczkowalna, a jej pochodna
wynosi $(2/\sqrt{\pi})e^{-x^2}$.

Ponadto całka niewłaściwa
$I=\int_0^\infty e^{-x^2}\,\mathrm dx$
jest zbieżna: dla $x\ge1$ zachodzi
$e^{-x^2}\le e^{-x}$. Tonelli pozwala
pomnożyć dodatnie całki, a zmiana
zmiennych na biegunowe w pierwszej
ćwiartce płaszczyzny daje
\[
 I^2=\int_0^\infty\int_0^\infty
 e^{-(x^2+y^2)}\,\mathrm dx\,\mathrm dy
 =\int_0^{\pi/2}\int_0^\infty
 e^{-\rho^2}\rho\,\mathrm d\rho\,\mathrm d\theta
 =\frac{\pi}{4}.
\]
Ponieważ $I>0$, otrzymujemy
$I=\sqrt{\pi}/2$. Wartość oznaczona
może więc być elementarna mimo
nieelementarności funkcji pierwotnej.
\end{przyklad}

\section{Całka sinusowa}

\begin{twierdzenie}[Nieelementarność całki sinusowej]
\label{aa:tw:sine-integral-nonelementary}
Funkcja $\sin x/x$, z wartością $1$
w punkcie $x=0$, nie ma elementarnej
funkcji pierwotnej na żadnym
niepustym otwartym przedziale rzeczywistym.
\end{twierdzenie}
\begin{proof}
Przyjmijmy $t=e^{ix}$, więc
$Dt=it$. Lemat
\ref{aa:lem:exp-polynomial-constant-field}
z $p(x)=ix$ zapewnia przestępność
i brak nowych stałych.
Korzystamy z pojedynczego generatora:
\[
 \frac{\sin x}{x}
 =\frac{t-t^{-1}}{2ix}.
\]
Nie zakładamy niezależności
$t$ i $t^{-1}$; ich iloczyn wynosi $1$.
Twierdzenie~\ref{aa:tw:laurent-hyperexp-criterion}
dla współczynnika przy dodatniej
potędze $t$ wymuszałoby
$r'+ir=1/(2ix)$ dla
$r\in\mathbb C(x)$.
Po położeniu $R=2ir$
dostalibyśmy $R'+iR=1/x$.
Jest to niemożliwe na mocy
Lematu~\ref{aa:lem:simple-pole-obstruction}
w punkcie $0$ z $h=i$.
Nawet jeśli sam iloraz $\sin x/x$
jest regularny w zerze, wymuszone
równanie wymierne ma tam biegun prosty.
Sprzeczność dowodzi tezy.
\end{proof}

\begin{przyklad}[Całka Dirichleta i funkcja $\operatorname{Si}$]
\label{aa:prz:sine-integral-dirichlet}
Funkcję sinusową definiuje się przez
$\operatorname{Si}(x)=
\int_0^x (\sin u/u)\,\mathrm du$,
przyjmując wartość $1$ funkcji podcałkowej
w zerze. Wyznaczymy ponadto
$\int_0^\infty(\sin x/x)\,\mathrm dx$.

Dla $\varepsilon>0$ oznaczmy
$I(\varepsilon)=\int_0^\infty
e^{-\varepsilon x}(\sin x/x)\,\mathrm dx$.
Całka jest bezwzględnie zbieżna.
Na dowolnym przedziale
$\varepsilon\in[\delta,M]$,
$\delta>0$, pochodną podcałkową
względem $\varepsilon$
dominuje $e^{-\delta x}$.
Możemy zatem różniczkować pod całką:
\[
 I'(\varepsilon)
 =-\int_0^\infty e^{-\varepsilon x}\sin x\,\mathrm dx
 =-\frac1{1+\varepsilon^2}.
\]
Druga równość wynika z części urojonej
całki $\int_0^\infty e^{-(\varepsilon-i)x}
\,\mathrm dx=(\varepsilon-i)^{-1}$.
Ponieważ $|\sin x|/x\le1$ dla $x>0$,
mamy $|I(\varepsilon)|\le1/\varepsilon$.
Całkując równanie różniczkowe
od $\varepsilon$ do nieskończoności, otrzymujemy
$I(\varepsilon)=\pi/2-\arctan\varepsilon$.

Uzasadnijmy przejście do granicy przy
$\varepsilon\downarrow0$. Dla $0<A<B$
funkcja $g_\varepsilon(x)=e^{-\varepsilon x}/x$
jest dodatnia i malejąca do zera dla
każdego $\varepsilon\ge0$.
Całkowanie przez części daje
\[
 \left|\int_A^B g_\varepsilon(x)\sin x\,\mathrm dx\right|
 \le g_\varepsilon(A)+g_\varepsilon(B)
       +\int_A^B|g_\varepsilon'(x)|\,\mathrm dx
 =2g_\varepsilon(A)\le\frac2A.
\]
To oszacowanie jest jednostajne względem
$\varepsilon$ i pozwala kontrolować ogony.
Na każdym $[0,A]$ funkcje podcałkowe
zbiegają jednostajnie, gdy
$\varepsilon\downarrow0$, do
ciągłego przedłużenia $\sin x/x$.
Dzieląc całki na $[0,A]$ i $[A,\infty)$,
przechodzimy więc do granicy i otrzymujemy
\[
 \int_0^\infty\frac{\sin x}{x}\,\mathrm dx
 =\lim_{\varepsilon\downarrow0}I(\varepsilon)
 =\frac{\pi}{2}.
\]
Zbieżność niewłaściwej całki nie oznacza
elementarności jej funkcji pierwotnej.
\end{przyklad}

\section{Logarytm całkowy i całka wykładnicza}

\begin{twierdzenie}[Logarytm całkowy]
\label{aa:tw:logarithmic-integral-nonelementary}
Na żadnym niepustym otwartym przedziale
$I\subset(0,1)$ lub $I\subset(1,\infty)$
funkcja $1/\log x$ nie ma elementarnej
funkcji pierwotnej. Również $e^x/x$
nie ma elementarnej funkcji pierwotnej
na żadnym przedziale nieprzechodzącym
przez $0$.
\end{twierdzenie}
\begin{proof}
Najpierw rozważmy $e^x/x$.
Lemat~\ref{aa:lem:exp-polynomial-constant-field}
zapewnia, że $t=e^x$ jest
przestępny nad $\mathbb C(x)$,
$Dt=t$ i nie dodaje stałych.
Z Twierdzenia~\ref{aa:tw:hyperexponential-criterion}
wynika, że elementarna funkcja pierwotna
wymagałaby $r'+r=1/x$ dla
$r\in\mathbb C(x)$. Lemat
\ref{aa:lem:simple-pole-obstruction}
w punkcie $x=0$, przy regularnym $h=1$,
wyklucza takie rozwiązanie.
Biegun równania wymiernego jest przeszkodą
globalnej tożsamości w $\mathbb C(x)$,
nawet gdy badany przedział nie zawiera zera.

Przypuśćmy teraz, że na przedziale
$I\subset(0,1)$ lub $I\subset(1,\infty)$
istnieje elementarna funkcja pierwotna
$F$ funkcji $1/\log x$.
Podstawmy $x=e^u$ i różniczkujmy:
\[
 \frac{\mathrm d}{\mathrm du}F(e^u)
 =\frac{e^u}{u}.
\]
Skończone wyrażenie elementarne
po takim podstawieniu nadal pozostaje
elementarne; wybrane lokalne gałęzie
funkcji mają sens również w zespolonym
otoczeniu zwykłego punktu.
Dawałoby to elementarną funkcję
pierwotną $e^u/u$, którą właśnie
wykluczyliśmy. Otrzymujemy sprzeczność.
\end{proof}

\begin{przyklad}[Dwie funkcje specjalne]
\label{aa:prz:li-ei-definitions}
Na $x>1$ możemy określić
$\operatorname{Li}(x)=
\int_2^x(\log u)^{-1}\,\mathrm du$.
Odcinek całkowania nie przecina osobliwości
$x=1$; na $(0,1)$ wybiera się punkt
początkowy z tego przedziału.
W podobny sposób na $x>0$ określamy
$\operatorname{Ei}(x)$, z dokładnością
do stałej, jako funkcję pierwotną $e^x/x$.
Z twierdzenia wynika, że obie funkcje
nie mają elementarnych lokalnych wzorów.
Dla kontrastu
\[
 \int\frac1{x\log x}\,\mathrm dx
 =\log|\log x|+C
\]
na każdym z przedziałów $(0,1)$
oraz $(1,\infty)$; pochodną prawej
strony sprawdzamy przez regułę łańcuchową.
\end{przyklad}

\section{Całka eliptyczna jako odmienny rodzaj przeszkody}

Rozważmy funkcję algebraiczną $y$ daną
równaniem $y^2=1-x^4$. Nie jest ona
hiperwykładnicza nad $\mathbb C(x)$
w sensie poprzednich kryteriów:
$Dy/y=-2x^3/(1-x^4)$ wprawdzie należy
do $\mathbb C(x)$, ale $y$ jest
\emph{algebraiczny}, a nie przestępny.
Wymagana będzie odrębna metoda oparta
na formach różniczkowych i rzędach
zer oraz biegunów.

\begin{lemat}[Ciało krzywej i lokalne współrzędne]
\label{aa:lem:elliptic-charts}
Niech $K=\mathbb C(x,y)$, gdzie
$y^2=1-x^4$, $Dx=1$ i
$Dy=-2x^3/y$. Wówczas $[K:\mathbb C(x)]=2$
i $C_K=\mathbb C$. Każdy $v\in K$
ma jednoznaczną postać $A(x)+B(x)y$
z $A,B\in\mathbb C(x)$.
Na gładkiej krzywej tej relacji,
z dołączonymi dwoma punktami
nad nieskończonością, lokalne współrzędne
można wybrać następująco:
\begin{enumerate}
\item gdy $x=a$ i $y\ne0$, współrzędną
jest $z=x-a$;
\item gdy $x=a$, $y=0$ i $a^4=1$,
współrzędną jest $z=y$, a $x-a$
ma rząd $2$;
\item nad nieskończonością, w zmiennych
$s=1/x$, $w=y/x^2$,
mamy $w^2=s^4-1$ i dwa punkty
$(s,w)=(0,i)$, $(0,-i)$.
W obu współrzędną jest $s$, a
$y=w/s^2$ ma biegun rzędu $2$.
\end{enumerate}
\end{lemat}
\begin{proof}
Wielomian $1-x^4$ ma cztery różne
pierwiastki i dlatego nie jest
kwadratem funkcji wymiernej
(kwadraty mają parzyste rzędy
zer i biegunów). Stąd
$T^2-(1-x^4)$ jest nierozkładalny
w $\mathbb C(x)[T]$ i otrzymujemy
stopień rozszerzenia oraz
postać $A+By$. Ponieważ
$\mathbb C$ jest algebraicznie
domkniętym ciałem stałych
$\mathbb C(x)$, wniosek
\ref{aa:cor:algebraic-constants}
zapewnia $C_K=\mathbb C$.

Relacja algebraiczna daje
$2y\,\mathrm dy=-4x^3\,\mathrm dx$.
Gdy $y\ne0$, z twierdzenia
o funkcji uwikłanej można
lokalnie wyrazić $y$ przez $x$.
Gdy $y=0$, zachodzi $x^4=1$;
ponieważ $-4x^3\ne0$, możemy
wyrazić $x$ jako funkcję $y$.
Dokładniej, dla gałęzi mającej
$x(0)=a$ mamy
$x=a(1-y^2)^{1/4}$.
Współczynnik przy $y^2$
wynosi $-a/4\ne0$, więc
$x-a$ ma rząd dokładnie $2$.

Dla $s=1/x$ i $w=y/x^2$
po przekształceniu równania
otrzymujemy $w^2=s^4-1$.
Nad $s=0$ są dokładnie dwa
punkty $w=\pm i$. Ponieważ
$2w\ne0$, $w$ jest tam
holomorficzną, niezerową
funkcją $s$.
Wprowadzone współrzędne
obejmują wszystkie skończone
punkty krzywej oraz wszystkie
punkty nad nieskończonością.
Każdy niezerowy element $K$
ma w każdej z nich
niezerowe rozwinięcie Laurenta
o skończonym rzędzie.
\end{proof}

\begin{lemat}[Funkcje bez biegunów na krzywej]
\label{aa:lem:elliptic-no-poles}
Przy oznaczeniach
Lematu~\ref{aa:lem:elliptic-charts},
jeśli element $v\in K$
nie ma bieguna w żadnym
z opisanych tam punktów,
to $v\in\mathbb C$.
\end{lemat}
\begin{proof}
Niech $v=A(x)+B(x)y$.
Inwolucja $\tau(x,y)=(x,-y)$
przestawia punkty krzywej, także
oba punkty nad nieskończonością.
Funkcja $\tau(v)=A-By$ nie ma
zatem biegunów. To samo dotyczy
$A=(v+\tau(v))/2$ oraz
$By=(v-\tau(v))/2$.

Gdyby funkcja wymierna $A(x)$
miała biegun w pewnym skończonym
$x=a$, miałaby biegun nad $a$
na krzywej, gdyż współrzędna
$x-a$ ma dodatni rząd.
Zatem $A$ nie ma skończonych
biegunów i jest wielomianem.
Przy $s=0$ wielomian
dodatniego stopnia miałby
biegun, więc $A\in\mathbb C$.

Rozważmy $By$. W punkcie
$x=a$, gdzie $1-a^4\ne0$,
funkcja $y$ jest niezerowa
i regularna; biegun $B$
stałby się więc biegunem $By$.
W punktach $a^4=1$ mamy
$\operatorname{ord}(y)=1$
oraz $\operatorname{ord}(x-a)=2$.
Jeśli $B$ miałaby biegun
rzędu $m\ge1$ w $x=a$,
rząd $By$ wynosiłby
$1-2m<0$, znów sprzeczność.
Funkcja $B$ nie ma więc
skończonych biegunów
i jest wielomianem.
Jeżeli $B\ne0$ ma stopień $d$,
to przy $s=0$ element
$By$ ma biegun rzędu $d+2$,
ponieważ $y=w/s^2$, a
$w(0)=\pm i\ne0$.
Zatem $B=0$, a $v=A$
jest stała.
\end{proof}

\begin{lemat}[Forma $\mathrm dx/y$]
\label{aa:lem:elliptic-regular-form}
Na krzywej z
Lematu~\ref{aa:lem:elliptic-charts}
forma różniczkowa
$\omega=\mathrm dx/y$ jest
wszędzie holomorficzna
i nie jest formą zerową.
\end{lemat}
\begin{proof}
Gdy $y\ne0$ i $x$ jest
lokalną współrzędną,
$\omega=\mathrm dx/y$
jest regularna i niezerowa.
Przy $y=0$ z relacji
$2y\,\mathrm dy=-4x^3\,\mathrm dx$
otrzymujemy
$\omega=-\mathrm dy/(2x^3)$;
ponieważ $x^4=1$,
współczynnik jest regularny
i niezerowy.
Nad nieskończonością
$\mathrm dx=-s^{-2}\,\mathrm ds$
i $y=w/s^2$, zatem
$\omega=-\mathrm ds/w$.
Również ten współczynnik
jest regularny i niezerowy
przy $w=\pm i$.
\end{proof}

\begin{twierdzenie}[Nieelementarność całki eliptycznej]
\label{aa:tw:elliptic-integral-nonelementary}
Niech $K=\mathbb C(x,y)$ będzie
ciałem różniczkowym, dla którego
$y^2=1-x^4$, $Dx=1$,
$Dy=-2x^3/y$. Wtedy element
$1/y$ nie ma elementarnej
funkcji pierwotnej nad $K$
bez nowych stałych.
W szczególności funkcja
$1/\sqrt{1-x^4}$ nie ma
elementarnej funkcji pierwotnej
na żadnym niepustym otwartym
przedziale zawartym w $(-1,1)$.
\end{twierdzenie}
\begin{proof}
Załóżmy przeciwnie istnienie
elementarnej funkcji pierwotnej
w rozszerzeniu bez nowych stałych.
Na mocy Twierdzenia
\ref{aa:tw:liouville-elementary-integration}
istnieją $v\in K$,
$u_1,\ldots,u_N\in K^\times$
oraz $c_1,\ldots,c_N\in\mathbb C$,
dla których
\[
 \frac1y=Dv+\sum_{j=1}^Nc_j\frac{Du_j}{u_j}.
\]
Pomnożenie przez $\mathrm dx$
daje tożsamość form na krzywej
\begin{equation}\label{aa:eq:elliptic-liouville-form}
 \omega=\mathrm dv+
       \sum_{j=1}^Nc_j\frac{\mathrm du_j}{u_j}.
\end{equation}
Jest to tożsamość elementów
ciała funkcji krzywej
i ich różniczek, więc
możemy badać ją w każdym
lokalnym układzie współrzędnych.

\emph{Eliminacja części dokładnej.}
W lokalnej współrzędnej $z$
element $u\in K^\times$ ma
postać $u=z^m g(z)$,
gdzie $m\in\mathbb Z$,
$g$ jest holomorficzna
i $g(0)\ne0$.
Wtedy
\[
 \frac{\mathrm du}{u}
 =m\frac{\mathrm dz}{z}
   +\frac{\mathrm dg}{g}.
\]
Pochodne logarytmiczne
mają zatem co najwyżej
bieguny proste.
Jeśli $v$ miałby
biegun rzędu $q\ge1$,
to $\mathrm dv$
miałaby biegun rzędu
$q+1\ge2$:
wyraz wiodący jest
$-q a_{-q}z^{-q-1}\,\mathrm dz\ne0$.
Nie mogłyby usunąć go
pochodne logarytmiczne,
a $\omega$ nie ma biegunów
na mocy
Lematu~\ref{aa:lem:elliptic-regular-form}.
Zatem $v$ nie ma biegunów
w żadnym punkcie. Lemat
\ref{aa:lem:elliptic-no-poles}
daje $v\in\mathbb C$,
czyli $\mathrm dv=0$.

\emph{Eliminacja części logarytmicznej.}
Współczynnik przy
$\mathrm dz/z$, czyli residuum
formy $\mathrm du_j/u_j$
w punkcie $P$, wynosi
$\operatorname{ord}_P(u_j)\in\mathbb Z$.
Ponieważ $\omega$ nie ma
biegunów, z
\eqref{aa:eq:elliptic-liouville-form}
wynika dla każdego punktu $P$
\[
 \sum_{j=1}^Nc_j\operatorname{ord}_P(u_j)=0.
\]
Jeśli $N=0$ albo wszystkie
$c_j=0$, prawa strona
\eqref{aa:eq:elliptic-liouville-form}
już jest zerowa.
W przeciwnym razie
wybierzmy bazę
$\gamma_1,\ldots,\gamma_s$
nad $\mathbb Q$
przestrzeni liniowej
rozpiętej przez
$c_1,\ldots,c_N$.
Zapiszmy
$c_j=\sum_{\nu=1}^s
q_{\nu j}\gamma_\nu$,
$q_{\nu j}\in\mathbb Q$.
Niezależność $\gamma_\nu$
nad $\mathbb Q$ daje
dla wszystkich $P$ i $\nu$
równości
$\sum_jq_{\nu j}\operatorname{ord}_P(u_j)=0$.
Niech $M\ge1$ będzie wspólnym
mianownikiem wszystkich
$q_{\nu j}$.
Dla każdego $\nu$ element
$U_\nu=\prod_{j=1}^N
u_j^{Mq_{\nu j}}\in K^\times$
ma całkowite wykładniki i
$\operatorname{ord}_P(U_\nu)=0$
dla każdego punktu $P$.
Nie ma zatem biegunów;
z Lematu~\ref{aa:lem:elliptic-no-poles}
mamy $U_\nu\in\mathbb C^\times$,
więc $\mathrm dU_\nu=0$.
Wobec tego
\[
 \sum_{j=1}^Nc_j\frac{\mathrm du_j}{u_j}
 =\frac1M\sum_{\nu=1}^s
       \gamma_\nu\frac{\mathrm dU_\nu}{U_\nu}
 =0.
\]
Równość~\eqref{aa:eq:elliptic-liouville-form}
dawałaby $\omega=0$,
wbrew Lematowi
\ref{aa:lem:elliptic-regular-form}.
Otrzymujemy sprzeczność.

Na przedziale $I\subset(-1,1)$
wybór dodatniej gałęzi
$y=\sqrt{1-x^4}$ daje
zanurzenie $K$ w ciele
zarodków funkcji meromorficznych.
Gdyby istniał tam lokalny
elementarny wzór na funkcję
pierwotną, dawałby
elementarne rozszerzenie
różniczkowe bez nowych
stałych nad $K$.
Właśnie wykluczyliśmy
takie rozszerzenie.
\end{proof}

\begin{uwaga}[Rola residuów]
\label{aa:uw:elliptic-residue-warning}
Nie wystarczyłoby powiedzieć,
że holomorficzna forma
ma zerowe residua:
niezerowa forma
$\omega=\mathrm dx/y$
ma właśnie tę własność.
Kluczowy krok dowodu
wykorzystuje \emph{całkowitość}
rzędów zer i biegunów funkcji
$u_j$ oraz niezależność
współczynników $c_j$
nad $\mathbb Q$.
Pozwala to skonstruować
rzeczywiste elementy $U_\nu\in K$,
które są stałe, i w konsekwencji
usunąć całą sumę
pochodnych logarytmicznych.
\end{uwaga}

\section{Przykłady i zakres zastosowań}

\begin{przyklad}[Kontrast z funkcjami wymiernymi]
\label{aa:prz:rational-primitive}
Każda $f\in\mathbb C(x)$
ma lokalną elementarną
funkcję pierwotną.
Istotnie, rozkład
na ułamki proste ma postać
\[
 f=P(x)+\sum_a\sum_{k=1}^{m_a}
     \frac{c_{a,k}}{(x-a)^k},
 \qquad P\in\mathbb C[x].
\]
Wielomian i wyrazy z $k\ge2$
całkujemy wymiernie, korzystając
z $D((x-a)^{1-k})=(1-k)(x-a)^{-k}$.
Dla $k=1$ funkcją pierwotną
jest $c_{a,1}\log(x-a)$
na ustalonej lokalnej gałęzi.
Ta obserwacja uzasadniała
brak warunku dla $s_0$
w Twierdzeniu
\ref{aa:tw:laurent-hyperexp-criterion}.
\end{przyklad}

Zebrane przykłady pokazują
dwie odrębne przeszkody
w elementarnym całkowaniu:
równanie wymierne
$r'+hr=s$ może nie mieć
rozwiązania, nawet gdy
całkowana funkcja jest
całkowita lub ma prosty
opis wykładniczy,
natomiast całka eliptyczna
wymaga analizy form
różniczkowych na
krzywej algebraicznej.
W następnym rozdziale,
dopiero po stabilizacji
niniejszego, wykorzystamy
ciała różniczkowe do
badania automorfizmów
rozszerzeń Picarda--Vessiota
równań liniowych.

% AA-46 | Picard--Vessiot i różniczkowa teoria Galois.
% Opracowanie ze źródła sources/pl/galois(2).tex i literatury:
% M. van der Put, M. Singer, Galois Theory of Linear Differential Equations.
% Twierdzenia głębokie (istnienie i jednoznaczność PV, odpowiedniość
% Galois, algebraiczność grupy, kryterium liouvillowskie) są JAWNIE cytowane.
\chapter{Rozszerzenia Picarda--Vessiota i różniczkowa teoria Galois}
\chaptermark{Teoria Picarda--Vessiota}
\label{aa:ch:picard-vessiot}

Klasyczna teoria Galois przyporządkowuje równaniom algebraicznym
grupy automorfizmów ciał rozkładu. Dla liniowych równań
różniczkowych analogiczną rolę pełnią rozszerzenia
Picarda--Vessiota i ich grupy automorfizmów różniczkowych.
Derywacja ogranicza automorfizmy: obraz rozwiązania musi
pozostać rozwiązaniem \emph{tego samego} równania. Zobaczymy,
jak wyprowadza to reprezentację grupy w $\mathrm{GL}_n(C)$.

Odwołujemy się do konstrukcji i kryteriów z rozdziałów
\ref{aa:ch:differential-fields}--\ref{aa:ch:nonelementary-integrals}.
W przeciwieństwie do pełnego dowodu twierdzenia Liouville'a,
kilka fundamentalnych twierdzeń strukturalnych teorii
Picarda--Vessiota przytoczymy jawnie jako wyniki klasyczne.
Ich pełne dowody wymagają rozwiniętej algebry różniczkowej
i teorii grup algebraicznych. Dowody elementarnych
lematów i wszystkich bezpośrednich rachunków wykonamy tutaj.

\section{Układy liniowe i macierze fundamentalne}

\begin{definicja}[Układ liniowy i macierz fundamentalna]
\label{aa:def:pv-fundamental}
Niech $(K,D)$ będzie ciałem różniczkowym
charakterystyki $0$, $C=C_K$, $n\ge1$
i $A\in M_n(K)$. Przez układ
\begin{equation}\label{aa:eq:pv-system}
 DY=AY
\end{equation}
rozumiemy równanie dla wektorów kolumnowych
$y\in L^n$, gdzie $L/K$ jest rozszerzeniem
różniczkowym, a $Dy$ oznacza różniczkowanie
wszystkich współrzędnych. Macierz
$U\in\mathrm{GL}_n(L)$ spełniającą
$DU=AU$ nazywamy \emph{macierzą fundamentalną}.
Jej kolumny są liniowo niezależnymi
rozwiązaniami nad ciałem stałych $C_L$.
\end{definicja}

\begin{lemat}[Przestrzeń rozwiązań]
\label{aa:lem:pv-solutions-constants}
Niech $L/K$ będzie rozszerzeniem różniczkowym,
$A\in M_n(K)$ i $U\in\mathrm{GL}_n(L)$
spełnia $DU=AU$. Wówczas
\[
 \{y\in L^n:Dy=Ay\}=U(C_L)^n.
\]
W szczególności przestrzeń rozwiązań
jest $n$-wymiarowa nad $C_L$.
Jeśli $V\in\mathrm{GL}_n(L)$ jest inną
macierzą fundamentalną, to $V=UM$
dla dokładnie jednej macierzy
$M\in\mathrm{GL}_n(C_L)$.
\end{lemat}
\begin{proof}
Z reguły różniczkowania odwrotności macierzy
$D(U^{-1})=-U^{-1}(DU)U^{-1}=-U^{-1}A$.
Dla rozwiązania $y$ otrzymujemy
\[
 D(U^{-1}y)=-U^{-1}Ay+U^{-1}Dy=0,
\]
więc współrzędne $U^{-1}y$ należą do $C_L$.
Odwrotnie, jeśli $c\in(C_L)^n$,
to $D(Uc)=(DU)c+U(Dc)=AUc$.
Stąd pierwsza równość.
Stosując ją do kolumn $V$,
otrzymujemy $V=UM$ z $M\in M_n(C_L)$,
a odwracalność $U,V$ wymusza
$M\in\mathrm{GL}_n(C_L)$. Jednoznaczność
wynika z odwracalności $U$.
\end{proof}

\begin{stwierdzenie}[Równanie Abela dla wyznacznika]
\label{aa:st:pv-abel-determinant}
Niech $U\in\mathrm{GL}_n(L)$ spełnia
$DU=AU$ dla $A\in M_n(K)$. Wówczas
\[
 D(\det U)=\operatorname{tr}(A)\det U.
\]
W szczególności, jeśli $\operatorname{tr}(A)=0$,
to $\det U\in C_L^\times$.
\end{stwierdzenie}
\begin{proof}
Wzór na pochodną wyznacznika
można otrzymać przez rozwinięcie
multiliniowe po kolumnach,
a po pomnożeniu przez $U^{-1}$
przyjmuje on postać
$D(\det U)=(\det U)\operatorname{tr}(U^{-1}DU)$.
Ponieważ $DU=AU$ i
$\operatorname{tr}(U^{-1}AU)=\operatorname{tr}(A)$,
otrzymujemy tezę. Gdy ślad jest zerowy,
pochodna wyznacznika znika; wyznacznik
jest niezerowy, bo $U$ jest odwracalna.
\end{proof}

\begin{przyklad}[Równanie skalarne jako układ]
\label{aa:prz:pv-companion}
Niech $a_0,\ldots,a_{n-1}\in K$.
Równanie
$D^ny+a_{n-1}D^{n-1}y+\cdots+a_0y=0$
jest równoważne układowi
$Dz=Az$ dla
$z=(y,Dy,\ldots,D^{n-1}y)^{\mathsf T}$,
gdzie $A$ ma jedynki na nadprzekątnej,
zera w pozostałych wierszach poza ostatnim,
a ostatni wiersz jest równy
$(-a_0,-a_1,\ldots,-a_{n-1})$.
Ostatnie równanie układu jest dokładnie
podanym równaniem skalarnym, a pozostałe
wiążą kolejne pochodne. Dzięki temu
wystarcza teoria macierzowa.
\end{przyklad}

\section{Różniczkowe pierścienie proste i rozszerzenia Picarda--Vessiota}

\begin{definicja}[Pierścień różniczkowy i ideał różniczkowy]
\label{aa:def:pv-differential-ideal}
Niech $R$ będzie przemiennym pierścieniem
z jedynką i derywacją $D:R\to R$.
Ideał $I\subseteq R$ nazywamy
\emph{różniczkowym}, jeśli $D(I)\subseteq I$.
Pierścień $R\ne0$ jest
\emph{różniczkowo prosty}, jeżeli
jego jedynymi ideałami różniczkowymi
są $(0)$ i $R$.
Prostota różniczkowa nie oznacza,
że $R$ musi być ciałem.
\end{definicja}

\begin{definicja}[Pierścień i rozszerzenie Picarda--Vessiota]
\label{aa:def:pv-extension}
Niech $(K,D)$ będzie ciałem różniczkowym
charakterystyki $0$ o algebraicznie
domkniętym ciele stałych $C=C_K$.
Dla $A\in M_n(K)$
\emph{pierścieniem Picarda--Vessiota}
dla układu~\eqref{aa:eq:pv-system}
nazywamy różniczkowo prostą
$K$-algebrę różniczkową $R$,
dla której istnieje
$U\in\mathrm{GL}_n(R)$ spełniająca
$DU=AU$ oraz
\[
 R=K[U_{ij},(\det U)^{-1}:
          1\le i,j\le n].
\]
Jest to zwykła algebra generowana
przez elementy macierzy i odwrotność
wyznacznika, a nie dowolna algebra
różniczkowa z dodatkowymi generatorami.

\emph{Ciałem Picarda--Vessiota} nazywamy
$E=\operatorname{Frac}(R)$, gdy $R$
jest dziedziną całkowitości.
Rozszerzenie $E/K$ nazywamy
\emph{rozszerzeniem Picarda--Vessiota}
dla rozważanego układu.
\end{definicja}

W definicji algebraiczna domkniętość $C$
jest założeniem, a nie automatyczną
własnością dowolnego ciała różniczkowego.
Pozwala sformułować standardową
teorię istnienia, jednoznaczności
i grup algebraicznych bez subtelności
związanych z rozszerzaniem stałych.

\begin{twierdzenie}[Istnienie i jednoznaczność -- wynik klasyczny]
\label{aa:tw:pv-existence-uniqueness}
Niech $(K,D)$ będzie ciałem różniczkowym
charakterystyki $0$ o algebraicznie
domkniętym ciele stałych $C_K=C$.
Dla każdego $n\ge1$ i każdego
$A\in M_n(K)$ istnieje pierścień
Picarda--Vessiota $R$.
Jest on dziedziną całkowitości,
a $E=\operatorname{Frac}(R)$ spełnia
$C_E=C$. Każde dwa takie rozszerzenia
są izomorficzne jako rozszerzenia
różniczkowe nad $K$
(nie kanonicznie).
\end{twierdzenie}

\begin{uwaga}[Co zawiera cytowany wynik?]
\label{aa:uw:pv-existence-cited}
To twierdzenie przytaczamy w klasycznej
postaci (van der Put--Singer,
\emph{Galois Theory of Linear Differential
Equations}, rozdział o teorii
Picarda--Vessiota).
Zarys konstrukcji jest następujący.
Na pierścieniu
$S=K[X_{ij},(\det X)^{-1}]$,
gdzie $X=(X_{ij})$ jest macierzą
formalnych zmiennych, rozszerzamy $D$
przepisem $DX=AX$.
Wybieramy maksymalny właściwy
ideał różniczkowy $I$.
Iloraz $S/I$ jest różniczkowo prosty
i zawiera fundamentalną macierz
obrazów $X_{ij}$.
Głębsze lemata algebry różniczkowej
gwarantują, że w charakterystyce
$0$ taki iloraz jest dziedziną
całkowitości i nie ma nowych stałych,
gdy $C$ jest algebraicznie domknięte.
Jednoznaczność wymaga ponadto
porównania dwóch pierścieni
fundamentalnych i ich
różniczkowych ideałów.
\emph{Właśnie te kroki są tu cytowane,
a nie przedstawiane jako dowiedzione.}
\end{uwaga}

\begin{uwaga}[Dlaczego nie wystarcza dowolne rozwiązanie?]
\label{aa:uw:pv-new-constants}
Dla równania $Dy=0$ i ciała
$K=\mathbb C(x)$ można dołączyć
przestępny element $t$ z $Dt=0$.
Nie jest to rozszerzenie
Picarda--Vessiota tego równania:
powstała nowa stała $t$, a
minimalne ciało rozwiązania to
już $K$. Sam fakt występowania
rozwiązań nie gwarantuje minimalności
ani zachowania ciała stałych.
\end{uwaga}

\section{Różniczkowa grupa Galois i jej reprezentacja}

\begin{definicja}[Różniczkowa grupa Galois]
\label{aa:def:pv-group}
Niech $E/K$ będzie rozszerzeniem
Picarda--Vessiota układu $DY=AY$.
\emph{Różniczkową grupą Galois}
nazywamy grupę
$G=\operatorname{Aut}_D(E/K)$
wszystkich automorfizmów ciała $E$,
które pozostawiają punktowo $K$
i komutują z derywacją:
$D(\sigma z)=\sigma(Dz)$
dla każdego $z\in E$.
\end{definicja}

\begin{twierdzenie}[Macierze automorfizmów]
\label{aa:tw:pv-automorphism-matrices}
Niech $(K,D)$ ma charakterystykę $0$
oraz algebraicznie domknięte ciało
stałych $C$, niech $E/K$ będzie
rozszerzeniem Picarda--Vessiota
dla $A\in M_n(K)$ i niech
$U\in\mathrm{GL}_n(E)$ będzie
macierzą fundamentalną
generującą $E$ nad $K$.
Dla każdego $\sigma\in G$
istnieje dokładnie jedna
$M_\sigma\in\mathrm{GL}_n(C)$ taka, że
\begin{equation}\label{aa:eq:pv-action}
 \sigma(U)=UM_\sigma.
\end{equation}
Odwzorowanie
$\rho_U:G\to\mathrm{GL}_n(C)$,
$\rho_U(\sigma)=M_\sigma$,
jest różnowartościowym
homomorfizmem grup.
Jeśli zmienimy bazę rozwiązań
$U$ na $UQ$ z
$Q\in\mathrm{GL}_n(C)$,
to macierze grupy zamieniają się
na $Q^{-1}M_\sigma Q$.
\end{twierdzenie}
\begin{proof}
Z $\sigma|_K=\mathrm{id}_K$ oraz
$D\sigma=\sigma D$ otrzymujemy
$D(\sigma U)=\sigma(DU)=\sigma(AU)=A\sigma U$.
Macierz $\sigma U$ jest odwracalna,
więc Lemat~\ref{aa:lem:pv-solutions-constants}
z $C_E=C$ daje
$\sigma U=UM_\sigma$ dla dokładnie
jednej macierzy $M_\sigma\in\mathrm{GL}_n(C)$.

Jeśli $\sigma,\tau\in G$, to
$(\sigma\circ\tau)U=\sigma(UM_\tau)
=(\sigma U)M_\tau=UM_\sigma M_\tau$,
ponieważ automorfizmy pozostawiają
$C\subseteq K$ punktowo.
Stąd $M_{\sigma\tau}=M_\sigma M_\tau$.
Jeśli $M_\sigma=I_n$, to $\sigma$
pozostawia wszystkie elementy $U$
nieruchome; ponieważ
$E=K(U_{ij}:1\le i,j\le n)$,
otrzymujemy $\sigma=\mathrm{id}_E$.
Dla zmiany $U'=UQ$ mamy
$\sigma(U')=UM_\sigma Q
=U'Q^{-1}M_\sigma Q$.
\end{proof}

\begin{twierdzenie}[Algebraiczność grupy -- wynik klasyczny]
\label{aa:tw:pv-algebraic-group}
Przy założeniach
Twierdzenia~\ref{aa:tw:pv-automorphism-matrices}
obraz $\rho_U(G)$ jest
\emph{domkniętą w topologii Zariskiego}
podgrupą $\mathrm{GL}_n(C)$,
czyli liniową grupą algebraiczną nad $C$.
Ponadto
$\operatorname{trdeg}_K E=\dim G$.
\end{twierdzenie}

\begin{uwaga}[Zakres twierdzenia]
\label{aa:uw:pv-algebraicity-cited}
Algebraiczność obrazu i formuła
dla stopnia przestępności są
wynikami teorii Picarda--Vessiota,
których tutaj nie dowodzimy.
Poprzednie twierdzenie daje
\emph{samodzielnie tylko} różnowartościowy
homomorfizm do $\mathrm{GL}_n(C)$.
Nie każda podgrupa $\mathrm{GL}_n(C)$
jest algebraiczna; domkniętości
Zariskiego nie wolno wyprowadzać
wyłącznie z działania na macierzy.
\end{uwaga}

\begin{twierdzenie}[Odpowiedniość Galois -- wynik klasyczny]
\label{aa:tw:pv-correspondence}
Niech $E/K$ będzie rozszerzeniem
Picarda--Vessiota w charakterystyce
$0$ z algebraicznie domkniętym
ciałem stałych $C$, a
$G=\operatorname{Aut}_D(E/K)$.
Wówczas $E^G=K$.
Odwzorowania
\[
 H\longmapsto E^H,\qquad
 F\longmapsto\operatorname{Aut}_D(E/F)
\]
ustalają odwracającą zawieranie
bijekcję między domkniętymi
podgrupami algebraicznymi $H\le G$
a pośrednimi ciałami różniczkowymi
$K\subseteq F\subseteq E$.
Pośrednie rozszerzenie $F/K$
jest rozszerzeniem Picarda--Vessiota
wtedy i tylko wtedy, gdy
odpowiadająca mu podgrupa
$H$ jest normalna w $G$;
w tym przypadku
$\operatorname{Aut}_D(F/K)\simeq G/H$.
\end{twierdzenie}

\begin{uwaga}[Dlaczego wymagamy domkniętości?]
\label{aa:uw:pv-correspondence-cited}
To pełne twierdzenie przytaczamy
za teorią klasyczną, a nie dowodzimy
z samego homomorfizmu~\eqref{aa:eq:pv-action}.
Istotne jest ograniczenie do
\emph{podgrup algebraicznych domkniętych
w topologii Zariskiego}. Dla
dowolnej, niedomkniętej podgrupy
odwzorowanie $H\mapsto E^H$
nie musi być różnowartościowe:
podgrupa i jej domknięcie Zariskiego
mogą mieć to samo ciało elementów stałych.
\end{uwaga}

\section{Przykłady: grupy multiplikatywna, addytywna i cykliczna}

\begin{przyklad}[Równanie $Dy=y$: grupa multiplikatywna]
\label{aa:prz:pv-multiplicative}
Niech $K=\mathbb C(x)$, $Dx=1$, a
$Dy=y$. Rozwiązaniem jest $t=e^x$.
Na mocy Lematu
\ref{aa:lem:exp-polynomial-constant-field}
$t$ jest przestępny nad $K$
i $C_{K(t)}=\mathbb C$.
Pierścień $R=K[t,t^{-1}]$
jest różniczkowo prosty.
Udowodnijmy ten ostatni fakt.
Niech $I\ne(0)$ będzie
ideałem różniczkowym $R$.
Wybierzmy niezerowy element
$p=\sum_{j=0}^m a_jt^j\in I$
o najmniejszej liczbie
niezerowych wyrazów; po
pomnożeniu przez odwrotność
monomianu $a_kt^k$ możemy
przyjąć, że $p$ ma
wyraz stały równy $1$.
Wtedy $Dp$ także należy do $I$,
ale wyraz stały jego pochodnej
jest zerowy; minimalność
liczby wyrazów daje $Dp=0$.
Każdy niezerowy współczynnik
$a_j$, $j\ne0$, musiałby
spełniać $Da_j/a_j=-j$.
W $\mathbb C(x)$ jest to
niemożliwe: pochodna logarytmiczna
funkcji wymiernej nie może być
niezerową stałą. Zatem $p=1$,
czyli $I=R$.

Jest to więc pierścień
Picarda--Vessiota, $E=K(t)$.
Każdy automorfizm różniczkowy
spełnia $\sigma(t)=ct$ dla
$c\in\mathbb C^\times$, a
każdy taki wybór daje
automorfizm $K(t)$ komutujący
z derywacją. Stąd
$G\simeq\mathbb G_m(\mathbb C)
=(\mathbb C^\times,\cdot)$.
\end{przyklad}

\begin{przyklad}[Równanie logarytmiczne: grupa addytywna]
\label{aa:prz:pv-additive}
Niech $K=\mathbb C(x)$ i
$t=\log x$ będzie przestępnym
generatorem o pochodnej $Dt=1/x$.
Brak nowych stałych wynika z
Lematu~\ref{aa:lem:primitive-extension-constants},
bo funkcja wymierna nie może
mieć pochodnej $1/x$.
Elementy $1,t$ są rozwiązaniami
równania $D^2y+x^{-1}Dy=0$.
Jeśli $z=(y,Dy)^{\mathsf T}$,
to jest to układ
\[
 Dz=\begin{pmatrix}0&1\\0&-1/x\end{pmatrix}z,
 \qquad
 U=\begin{pmatrix}1&t\\0&1/x\end{pmatrix}.
\]
Sprawdzamy wprost $DU=AU$ i
$\det U=1/x\in K^\times$.
Pierścień $R=K[t]$ jest
różniczkowo prosty:
dla niezerowego ideału
różniczkowego $I$ wybieramy
unormowany wielomian
$p=t^m+\sum_{j<m}a_jt^j\in I$
o najmniejszym stopniu.
Ponieważ $Dt=1/x$,
pochodna $Dp$ ma stopień
co najwyżej $m-1$, więc
z minimalności musi być zerowa.
Dla $m\ge1$ współczynnik
przy $t^{m-1}$ daje
$Da_{m-1}+m/x=0$,
co jest niemożliwe dla
funkcji wymiernej $a_{m-1}$,
gdyż pochodne wymierne
mają zerowe residua.
Zatem $m=0$ i $I=R$.

Automorfizmy mają postać
$t\mapsto t+c$ dla
$c\in\mathbb C$ i wszystkie
takie przesunięcia istnieją.
W bazie macierzy $U$ grupa
jest dana przez
\[
 M_c=\begin{pmatrix}1&c\\0&1\end{pmatrix},
 \qquad c\in\mathbb C,
\]
więc $G\simeq\mathbb G_a(\mathbb C)$.
\end{przyklad}

\begin{przyklad}[Rozszerzenie algebraiczne: grupa $\mu_m$]
\label{aa:prz:pv-cyclic}
Niech $m\ge2$, $K=\mathbb C(x)$
oraz $t^m=x$. Jedyna
przedłużona derywacja
spełnia $Dt=t/(mx)$.
Wielomian $T^m-x$ jest
nierozkładalny nad $\mathbb C(x)$
na mocy kryterium Eisensteina
w pierścieniu $\mathbb C[x]$
względem elementu pierwszego $x$:
współczynniki poza wiodącym
są podzielne przez $x$,
współczynnik stały $-x$
nie jest podzielny przez $x^2$,
a współczynnik wiodący $1$
nie jest podzielny przez $x$.
Samo wykluczenie $m$-tej potęgi
funkcji wymiernej nie wystarczałoby
do udowodnienia nierozkładalności
dla złożonego $m$.
Stąd $[K(t):K]=m$.
Ponieważ $C=\mathbb C$
jest algebraicznie domknięte,
$C_{K(t)}=\mathbb C$.
Pierścień $R=K[t,t^{-1}]=K(t)$
jest ciałem, zatem prostym
pierścieniem różniczkowym.
Rozwiązania równania
$Dy=(mx)^{-1}y$ generują $R$.
Każdy automorfizm wysyła
$t$ na $\zeta t$ z
$\zeta^m=1$, a każdy taki
wybór istotnie jest
automorfizmem różniczkowym.
Grupa wynosi
\[
 G=\mu_m=\{\zeta\in\mathbb C^\times:
                       \zeta^m=1\}.
\]
Ma ona $m$ elementów i wymiar
algebraiczny $0$, zgodnie z
$\operatorname{trdeg}_K K(t)=0$.
\end{przyklad}

\begin{przyklad}[Równanie harmoniczne]
\label{aa:prz:pv-harmonic}
Dla $K=\mathbb C(x)$ rozważmy
$D^2y+y=0$. Po podstawieniu
$t=e^{ix}$ dostajemy rozwiązania
$t$ i $t^{-1}$. Macierz
\[
 U=\begin{pmatrix}t&t^{-1}\\
                   it&-it^{-1}\end{pmatrix}
\]
spełnia $DU=AU$ dla
$A=\left(\begin{smallmatrix}0&1\\-1&0\end{smallmatrix}\right)$,
a $\det U=-2i\ne0$.
Mamy $E=K(t)$,
a grupa działa przez
$t\mapsto ct$, $c\in\mathbb C^\times$,
czyli
$\rho_U(G)=\{
\operatorname{diag}(c,c^{-1}):
c\in\mathbb C^\times\}$.
W szczególności wyznacznik
macierzy grupy wynosi $1$.
Nie wolno skalować rozwiązań
$t$ i $t^{-1}$ niezależnie,
ponieważ ich iloczyn wynosi $1$.
\end{przyklad}

\section{Rozszerzenia liouvillowskie i rozwiązania równań liniowych}

\begin{definicja}[Rozszerzenie liouvillowskie]
\label{aa:def:pv-liouvillian}
Niech $(K,D)$ ma charakterystykę $0$
i ciało stałych $C_K=C$.
Rozszerzenie różniczkowe $L/K$
bez nowych stałych nazywamy
\emph{liouvillowskim}, jeżeli istnieje
skończona wieża
$K=K_0\subseteq\cdots\subseteq K_r=L$,
w której $K_j=K_{j-1}(\theta_j)$,
a dla każdego $j$ zachodzi
co najmniej jeden warunek:
$\theta_j$ jest algebraiczny
nad $K_{j-1}$,
$D\theta_j\in K_{j-1}$ albo
$\theta_j\ne0$ i
$D\theta_j/\theta_j\in K_{j-1}$.
W odróżnieniu od rozszerzeń
elementarnych z
Definicji~\ref{aa:def:elementary-differential-extension}
krok $D\theta_j\in K_{j-1}$
pozwala dołączyć \emph{dowolną}
funkcję pierwotną, niekoniecznie
logarytm elementu niższego piętra.
\end{definicja}

\begin{twierdzenie}[Kryterium liouvillowskie -- wynik klasyczny]
\label{aa:tw:pv-liouvillian-criterion}
Niech $E/K$ będzie rozszerzeniem
Picarda--Vessiota układu $DY=AY$,
gdzie $\operatorname{char}K=0$,
a $C_K$ jest algebraicznie domknięte.
Niech $G=\operatorname{Aut}_D(E/K)$,
a $G^\circ$ oznacza składową spójną
jedności grupy algebraicznej $G$.
Wówczas $E/K$ jest rozszerzeniem
liouvillowskim wtedy i tylko wtedy,
gdy $G^\circ$ jest grupą rozwiązalną.
Dla równania drugiego rzędu
$D^2y=qy$, $q\in K$, istnienie
niezerowego rozwiązania w
pewnym rozszerzeniu liouvillowskim
jest równoważne istnieniu
elementu algebraicznego $r$
nad $K$ spełniającego
\begin{equation}\label{aa:eq:pv-riccati}
 Dr+r^2=q.
\end{equation}
\end{twierdzenie}

\begin{uwaga}[Zakres cytowanego kryterium]
\label{aa:uw:pv-lie-kolchin}
Kryterium to wykorzystuje
twierdzenie Liego--Kołczyna
o trójkątyzacji spójnych
rozwiązalnych grup liniowych
oraz odpowiedniość Galois
z Twierdzenia~\ref{aa:tw:pv-correspondence}.
Nie dowodzimy tu ani tego
twierdzenia, ani ogólnej
równoważności z liouvillowskością.
Prosty kierunek ostatniego
kryterium można jednak
sprawdzić bezpośrednio:
gdy $r$ jest algebraiczne
i spełnia~\eqref{aa:eq:pv-riccati},
dołączamy $r$, a następnie
niezerowy element $y$ z
$Dy/y=r$; otrzymujemy
$D^2y/y=Dr+r^2=q$.
Dla odwrotnego kierunku
potrzebna jest teoria
grup algebraicznych.
\end{uwaga}

\begin{przyklad}[Nieelementarna, ale liouvillowska]
\label{aa:prz:pv-erf-liouvillian}
Z rozdziału~\ref{aa:ch:nonelementary-integrals}
wiemy, że funkcja
$F(x)=\int_0^x e^{-u^2}\,\mathrm du$
nie ma elementarnej funkcji pierwotnej.
Jest jednak liouvillowska
w przyjętym znaczeniu:
najpierw dołączamy
$t=e^{-x^2}$ z
$Dt/t=-2x$, potem
element $F$ z $DF=t$.
Przestępność $t$ i brak nowych
stałych dla $K(t)$ wykazano
w Lemacie~\ref{aa:lem:exp-polynomial-constant-field}.
Jeśli $DF=t$ ma już rozwiązanie
w $K(t)$, nie trzeba nic
dołączać; w przeciwnym przypadku
przestępny krok z $DF=t$
również nie wprowadza nowych
stałych na mocy
Lematu~\ref{aa:lem:primitive-extension-constants}.
Ponadto $D^2F+2xDF=0$.
Jest to konkretny przykład
rozwiązania liniowego równania,
które nie jest elementarne,
choć należy do wieży liouvillowskiej.
\end{przyklad}

\section{Równanie Airy'ego i grupa specjalna liniowa}

\begin{lemat}[Przeszkoda Riccatiego dla równania Airy'ego]
\label{aa:lem:pv-airy-riccati}
Równanie $Dr+r^2=x$,
gdzie $D=\mathrm d/\mathrm dx$,
nie ma rozwiązań algebraicznych
nad $\mathbb C(x)$.
\end{lemat}
\begin{proof}
Przypuśćmy, że istnieje
element algebraiczny $r$
spełniający równanie.
Lokalnie przy każdym skończonym
$x=a$ gałęzie funkcji
algebraicznej mają rozwinięcia
Puiseux w potęgach
$(x-a)^{1/e}$.

Najpierw sprawdźmy, że żadna
taka gałąź nie może być
rozgałęziona w skończonym punkcie.
Jeżeli $r$ jest tam ograniczona,
równanie $r'=x-r^2$
ma holomorficzną prawą stronę
i lokalne rozwiązanie
o zadanej wartości $r(a)$
jest holomorficzne
i jednoznaczne.
Również porównanie najniższej
niecałkowitej potęgi Puiseux
pokazuje, że taka potęga
nie może wystąpić:
różniczkowanie obniża ją
o $1$, podczas gdy
prawa strona jej nie wytwarza
na niższym poziomie.

Jeśli gałąź ma biegun,
kładziemy $u=1/r$.
Z równania wynika
$u'=1-xu^2$,
a $u(a)=0$.
Prawa strona jest holomorficzna,
więc poprzedni argument
zastosowany do $u$
wyklucza rozgałęzienie
również w tym przypadku.
Zatem wszystkie skończone
punkty są nierozgałęzione.

Otrzymujemy spójne skończone
nierozgałęzione przykrycie
płaszczyzny zespolonej,
odpowiadające wielomianowi
minimalnemu $r$.
Płaszczyzna zespolona jest
po prostu spójna, więc takie
przykrycie ma stopień $1$.
W tym kroku używamy klasycznego
twierdzenia o przykryciach
powierzchni Riemanna.
Wynika stąd $r\in\mathbb C(x)$.

Pokażmy, że nie istnieje
nawet rozwiązanie wymierne.
Jeżeli wielomianowa część
$r$ w nieskończoności ma
stopień $d\ge1$, to $r^2$
ma biegun rzędu $2d$,
podczas gdy $r'$ ma biegun
rzędu co najwyżej $d-1$,
więc $r'+r^2$ ma biegun
rzędu parzystego $2d\ge2$.
Funkcja $x$ ma tam
biegun rzędu $1$.
Jeśli $r$ jest regularna
w nieskończoności, to
$r'+r^2$ jest również
regularna, a $x$ nie.
W obu przypadkach sprzeczność.
\end{proof}

\begin{twierdzenie}[Równanie Airy'ego]
\label{aa:tw:pv-airy-group}
Niech $K=\mathbb C(x)$,
$D=\mathrm d/\mathrm dx$.
Równanie $D^2y-xy=0$
nie ma niezerowego rozwiązania
liouvillowskiego nad $K$,
a jego różniczkowa grupa
Galois jest izomorficzna
z $\mathrm{SL}_2(\mathbb C)$.
\end{twierdzenie}
\begin{proof}
Z Lematu~\ref{aa:lem:pv-airy-riccati}
oraz cytowanego kryterium
\ref{aa:tw:pv-liouvillian-criterion}
wynika pierwsze stwierdzenie.

Niech $E/K$ będzie rozszerzeniem
Picarda--Vessiota rozważanego
równania, a $U$ jego macierzą
fundamentalną. Dla towarzyszącego
układu $A=\left(\begin{smallmatrix}
0&1\\x&0\end{smallmatrix}\right)$
mamy $\operatorname{tr}(A)=0$.
Na mocy
Stwierdzenia~\ref{aa:st:pv-abel-determinant}
$\det U\in\mathbb C^\times$.
Wybierając bazę rozwiązań
o wyznaczniku $1$, dostajemy
$G\subseteq\mathrm{SL}_2(\mathbb C)$.

Gdyby składowa spójna
$G^\circ$ była rozwiązalna,
cytowane kryterium
\ref{aa:tw:pv-liouvillian-criterion}
zapewniałoby niezerowe
rozwiązanie liouvillowskie,
wbrew pierwszej części dowodu.
Zatem $G^\circ$ nie jest
rozwiązalna.
Każda właściwa spójna
algebraiczna podgrupa
$\mathrm{SL}_2(\mathbb C)$
jest jednak rozwiązalna:
jej algebra Liego jest
właściwą podalgebrą
trójwymiarowej algebry
$\mathfrak{sl}_2(\mathbb C)$,
ma więc wymiar co najwyżej
$2$, a każda algebra Liego
wymiaru co najwyżej $2$
jest rozwiązalna.
Stąd $G^\circ=\mathrm{SL}_2(\mathbb C)$,
a ponieważ $G$ leży
w $\mathrm{SL}_2(\mathbb C)$,
otrzymujemy
$G=\mathrm{SL}_2(\mathbb C)$.
\end{proof}

\begin{uwaga}[Granice samodzielności dowodu]
\label{aa:uw:pv-airy-deep-inputs}
Dowód lematu o Riccatim
wymaga standardowego faktu
o lokalnych rozwinięciach
Puiseux i prostego twierdzenia
o nierozgałęzionych przykryciach
$\mathbb C$. Przejście od
braku algebraicznych rozwiązań
Riccatiego do grupy Galois
korzysta jawnie z głębokiego
kryterium~\ref{aa:tw:pv-liouvillian-criterion}.
Nie należy przedstawiać
końcowego wyniku
$\mathrm{SL}_2(\mathbb C)$
jako konsekwencji samego
rachunku stopni funkcji
wymiernych.
\end{uwaga}

\section{Co pozostaje poza zakresem rozdziału?}

Odpowiedniość Galois wyjaśnia,
dlaczego struktura rozwiązań
układu liniowego jest związana
z grupą algebraiczną automorfizmów.
Nie daje ona jeszcze
algorytmu rozstrzygającego
dla każdego układu,
jaka jest jego grupa,
ani automatycznej procedury
znajdowania funkcji pierwotnych.

Algorytm Rischa dotyczy
elementarnego całkowania,
a algorytm Kovacica dotyczy
liouvillowskich rozwiązań
jednorodnych równań
drugiego rzędu o wymiernych
współczynnikach. Ich ogólne
schematy i dowody poprawności
wymagają osobnego rozdziału
i nie są objęte niniejszym
opracowaniem. Rozdział ten
kończy przygotowany wcześniej
blok AA-43--AA-46;
ewentualny AA-47 wymaga
osobnego zatwierdzenia zakresu.

% AA-47 | Algorytmy Rischa i Kovacica — zakres zatwierdzony przez kontynuację.
% Dowody pełne: integracja C(x), redukcja Hermite'a, test Rothsteina--Tragera,
% r'+hr=s przy h wielomianowym, wymierna galaz Riccatiego dla q wielomianowego.
% Wyniki cytowane: pelny algorytm Rischa i pelny algorytm Kovacica (przypadki 2, 3).
\chapter{Algorytmy Rischa i Kovacica}
\chaptermark{Algorytmy Rischa i Kovacica}
\label{aa:ch:risch-kovacic}

Twierdzenie Liouville'a z rozdziału
\ref{aa:ch:liouville-elementary-integration}
daje kryterium istnienia elementarnej funkcji pierwotnej,
a teoria Picarda--Vessiota z rozdziału
\ref{aa:ch:picard-vessiot}
wiąże rozwiązania liouvillowskie równań liniowych
z grupami algebraicznymi. Są to twierdzenia
strukturalne. W tym rozdziale przejdziemy
od istnienia do \emph{rozstrzygalnych obliczeń}.

Dla przejrzystości rozdzielamy dwa poziomy:
\emph{pełne dowody} procedur nad
$\mathbb C(x)$, dla równań $r'+hr=s$
z $h\in\mathbb C[x]$, $s\in\mathbb C(x)$,
oraz dla wymiernego równania Riccatiego
z wielomianowym współczynnikiem;
\emph{klasyczne wyniki cytowane}
dotyczą ogólnego algorytmu Rischa
i pozostałych gałęzi algorytmu Kovacica.
Nie utożsamiamy pierwszego zakresu
z pełną implementacją drugiego.

\section{Model obliczeń i znaczenie odpowiedzi algorytmu}

\begin{definicja}[Efektywnie zadane współczynniki]
\label{aa:def:effective-constants}
Przy obliczeniach nad ciałem $C$ charakterystyki
$0$ zakładamy dostęp do dokładnych działań,
testu równości i efektywnego rozkładu
wielomianów używanych w procedurze
na czynniki w odpowiednim skończonym
algebraicznym rozszerzeniu ciała współczynników.
Jako podstawowy przykład można przyjąć
ciało liczb algebraicznych
$\overline{\mathbb Q}\subset\mathbb C$.
Twierdzenia matematyczne poniżej formułujemy
nad $\mathbb C$, a słowo \emph{algorytm}
oznacza skończoną liczbę wskazanych działań
na efektywnie reprezentowanych danych.
\end{definicja}

Przy zadaniu współczynników przez dowolne
nienazwane liczby zespolone bez testu równości
nie otrzymujemy modelu efektywnego.
W szczególności \emph{pełne}
twierdzenia o rozstrzygalności wymagają
jawnego określenia sposobu reprezentacji
ciała współczynników.

\section{Całkowanie wymierne: pełny algorytm}

\begin{twierdzenie}[Rozkład na część wymierną i logarytmy]
\label{aa:tw:rational-integration-algorithm}
Niech $f\in\mathbb C(x)$. Istnieją
$R\in\mathbb C(x)$, parami różne
$a_1,\ldots,a_m\in\mathbb C$
i $c_1,\ldots,c_m\in\mathbb C$ takie, że
\begin{equation}\label{aa:eq:rational-integration}
 f=R'+\sum_{j=1}^{m}\frac{c_j}{x-a_j}.
\end{equation}
Można je wyznaczyć w skończonej liczbie
działań arytmetycznych i rozkładów wielomianów.
Ponadto $f$ ma funkcję pierwotną
w $\mathbb C(x)$ wtedy i tylko wtedy,
gdy wszystkie residua $c_j$ są zerowe.
Niezależnie od tego $f$ ma lokalną
elementarną funkcję pierwotną.
\end{twierdzenie}
\begin{proof}
Dzielimy licznik przez mianownik,
a następnie rozkładamy część właściwą
na ułamki proste:
\[
 f=P(x)+
 \sum_{j=1}^{m}\sum_{k=1}^{n_j}
         \frac{b_{j,k}}{(x-a_j)^k}.
\]
Jest to skończona procedura nad
ciałem, w którym rozkłada się mianownik.
Wielomian $P$ całkujemy wyraz po wyrazie.
Dla $k\ge2$ element
$-b_{j,k}(x-a_j)^{1-k}/(k-1)$
ma pochodną
$b_{j,k}(x-a_j)^{-k}$.
Po odjęciu otrzymanych pochodnych
pozostają tylko składniki
$b_{j,1}/(x-a_j)$.
Daje to~\eqref{aa:eq:rational-integration}
z $c_j=b_{j,1}$.
W każdym skończonym punkcie
residuum pochodnej funkcji wymiernej
jest zerowe: pochodna szeregu
Laurenta $\sum_k v_k(x-a)^k$
nie ma współczynnika przy
$(x-a)^{-1}$, gdyż dla $k=0$
otrzymujemy $k v_k=0$.
Jeśli więc $f$ ma wymierną funkcję
pierwotną, wszystkie $c_j$ muszą być
zerowe. Odwrotnie, przy znikających
residuach $f=R'$, więc taka
funkcja pierwotna istnieje.
W ogólnym przypadku na ustalonych
lokalnych gałęziach można przyjąć
$R+\sum_jc_j\log(x-a_j)$.
\end{proof}

\begin{definicja}[Redukcja Hermite'a]
\label{aa:def:hermite-rational}
\emph{Redukcją Hermite'a funkcji wymiernej}
nazywamy konstrukcję, która odejmuje
pochodne wymiernych części głównych
rzędu co najmniej dwa i sprowadza
$f$ do sumy pochodnej wymiernej
oraz funkcji wymiernej mającej
wyłącznie bieguny proste.
W dowodzie
Twierdzenia~\ref{aa:tw:rational-integration-algorithm}
wykonaliśmy taką redukcję
przez rozkład na ułamki proste.
Warianty bez rozdzielania mianownika
na czynniki liniowe używają
dzielenia wielomianów i ich NWD.
\end{definicja}

\begin{przyklad}[Obliczenie obu części]
\label{aa:prz:hermite-example}
Dla
$f=2x+3/(x-1)^3+5/(x-1)$
otrzymujemy
\[
 R=x^2-\frac{3}{2(x-1)^2},
 \qquad
 f=R'+\frac5{x-1}.
\]
Zatem na dowolnym obszarze,
na którym można wybrać gałąź
logarytmu $x-1$, funkcją
pierwotną jest
$R+5\log(x-1)$.
Nie istnieje wymierna funkcja
pierwotna, ponieważ residuum
w $1$ wynosi $5$.
\end{przyklad}

\begin{twierdzenie}[Wykorzystanie wypadkowej do residuów]
\label{aa:tw:rothstein-trager}
Niech $A,B\in\mathbb C[x]$,
$\deg A<\deg B$, a $B$ będzie
unormowanym wielomianem bez
pierwiastków wielokrotnych.
Dla każdego pierwiastka $a$ wielomianu
$B$ residuum $A/B$ jest równe
$c_a=A(a)/B'(a)$.
Wielomian
\begin{equation}\label{aa:eq:rt-resultant}
 T(z)=\operatorname{Res}_x(B(x),A(x)-zB'(x))
\end{equation}
ma (z uwzględnieniem krotności)
dokładnie te wartości $c_a$ jako pierwiastki.
Jeśli $c$ przebiega różne pierwiastki $T$,
a $P_c=\gcd(B,A-cB')$ jest unormowany,
to zachodzi tożsamość
\begin{equation}\label{aa:eq:rt-reconstruction}
 \frac{A}{B}=\sum_{c:T(c)=0}
              c\frac{P_c'}{P_c}.
\end{equation}
Zatem przy efektywnym rozkładzie
$T$ i obliczaniu NWD można znaleźć
część logarytmiczną bez wcześniejszego
rozkładu $B$ na czynniki liniowe.
\end{twierdzenie}
\begin{proof}
Ponieważ pierwiastki $a$ wielomianu $B$
są proste, lokalnie
$B(x)=(x-a)B'(a)+O((x-a)^2)$,
co daje podane residuum.
Własność wypadkowej dla
unormowanego $B$ daje
\[
 T(z)=\prod_{B(a)=0}
           (A(a)-zB'(a))
     =\left(\prod_{B(a)=0}B'(a)\right)
       \prod_{B(a)=0}(c_a-z).
\]
Zatem pierwiastki wypadkowej są
dokładnie wartościami residuów.
Pierwiastkami $P_c$ są właśnie
te pierwiastki $a$ wielomianu $B$,
dla których $c_a=c$. W szczególności
\[
 \frac{P_c'}{P_c}
 =\sum_{\substack{B(a)=0\\c_a=c}}
             \frac1{x-a}.
\]
Obie strony~\eqref{aa:eq:rt-reconstruction}
mają te same części główne
przy każdym skończonym biegunie;
obie są właściwymi funkcjami wymiernymi
(przy nieskończoności mają granicę
zerową). Ich różnica jest wielomianem
bez części wielomianowej, a więc zerem.
\end{proof}

\begin{przyklad}[Wypadkowa bez rozkładu mianownika]
\label{aa:prz:rt-example}
Dla $A=1$, $B=x^2+1$ mamy
$B'=2x$ i
\[
 T(z)=\operatorname{Res}_x(x^2+1,1-2zx)
     =1+4z^2.
\]
Pierwiastkami są $c=\pm i/2$.
Odpowiednie NWD wybierają czynniki
$x+i$ i $x-i$. Otrzymujemy
\[
 \frac1{x^2+1}
 =\frac1{2i}\left(\frac1{x-i}
                         -\frac1{x+i}\right).
\]
Dla $A=2x$ otrzymalibyśmy
jedno residuum $c=1$, a
$P_1=x^2+1$, skąd
$2x/(x^2+1)=B'/B$.
\end{przyklad}

\section{Efektywna część algorytmu Rischa: krok wykładniczy}

Niech $h\in\mathbb C[x]$ i $s\in\mathbb C(x)$.
Dla niezerowego $h$ wielomian
$p(x)=\int h(x)\,\mathrm dx$
jest niestały; na mocy
Lematu~\ref{aa:lem:exp-polynomial-constant-field}
element $t=e^{p(x)}$ jest
przestępny nad $\mathbb C(x)$
i nie dodaje stałych.
Z Twierdzenia
\ref{aa:tw:hyperexponential-criterion}
wynika, że element
$st$ ma elementarną funkcję pierwotną
wtedy i tylko wtedy, gdy
\begin{equation}\label{aa:eq:risch-first-order}
 r'+hr=s
\end{equation}
ma rozwiązanie $r\in\mathbb C(x)$.
Pokażemy, że w tym zakresie
test jest faktycznie \emph{skończonym
algorytmem decyzyjnym}.

\begin{lemat}[Bieguny wymiernego rozwiązania]
\label{aa:lem:risch-poles}
Niech $h\in\mathbb C[x]$ i
$s,r\in\mathbb C(x)$ spełniają
$r'+hr=s$. Jeżeli $r$ ma
w skończonym punkcie $a$
biegun rzędu $j\ge1$, to
$s$ ma w $a$ biegun dokładnie
rzędu $j+1$. Współczynnik
wiodący części głównej $s$
jest równy $-j$ razy
współczynnik wiodący
części głównej $r$.
W szczególności rozwiązanie
$r$ nie może mieć biegunów
w punktach regularności $s$,
a prosty biegun $s$ uniemożliwia
istnienie rozwiązania.
\end{lemat}
\begin{proof}
W lokalnej zmiennej $z=x-a$
zapiszmy $r=b_{-j}z^{-j}+
O(z^{1-j})$, gdzie $b_{-j}\ne0$.
Wtedy $r'=-j b_{-j}z^{-j-1}
+O(z^{-j})$, natomiast
$hr$ ma biegun rzędu
co najwyżej $j$, ponieważ
$h$ jest regularny w $a$.
Składnik o rzędzie $j+1$
nie może się skasować.
Zatem $s$ ma dokładnie
taki biegun i podany
współczynnik.
\end{proof}

\begin{twierdzenie}[Algorytm wymiernego równania Rischa]
\label{aa:tw:risch-polynomial-h-algorithm}
Niech $h\in\mathbb C[x]\setminus\{0\}$,
$s\in\mathbb C(x)$.
Istnieje skończona procedura,
oparta na rozkładzie mianownika
$s$ oraz rozwiązywaniu
skończonych układów równań liniowych,
która rozstrzyga istnienie
$r\in\mathbb C(x)$ spełniającego
$r'+hr=s$ i, jeśli rozwiązanie
istnieje, wyznacza je.
Dla $h=0$ kryterium
istnienia wymiernej funkcji pierwotnej
jest podane w
Twierdzeniu~\ref{aa:tw:rational-integration-algorithm}.
\end{twierdzenie}
\begin{proof}
Niech $a$ będzie skończonym biegunem
$s$ rzędu $m$. Jeżeli $m=1$,
Lemat~\ref{aa:lem:risch-poles}
wyklucza rozwiązanie.
Gdy $m\ge2$, ewentualne rozwiązanie
ma w $a$ część główną
$R_a(z)=\sum_{j=1}^{m-1}b_jz^{-j}$.
Podstawiamy ją do $r'+hr=s$
i porównujemy kolejno współczynniki
przy $z^{-m},z^{-m+1},\ldots,z^{-2}$.
Układ jest trójkątny:
współczynnik przy $z^{-j-1}$
zawiera $-j b_j$ oraz wyrazy
zależące wyłącznie od
$b_{j+1},\ldots,b_{m-1}$,
bo $h$ ma nieujemne potęgi $z$.
Ponieważ charakterystyka wynosi
$0$, wszystkie $b_j$ są
jednoznacznie wyznaczone.
Pozostaje sprawdzić współczynnik
przy $z^{-1}$. Jeżeli
różni się od odpowiedniego
współczynnika $s$, nie ma
rozwiązania. Jeśli się zgadza,
wyznaczona część główna usuwa
wszystkie bieguny w $a$.
Powtarzamy to dla każdego
skończonego bieguna $s$.
Otrzymujemy wymierną funkcję
$R=\sum_aR_a(x-a)$ i wielomian
$S=s-R'-hR\in\mathbb C[x]$.

Gdyby pierwotne równanie miało
rozwiązanie wymierne $r$,
to Lemat~\ref{aa:lem:risch-poles}
i jednoznaczność porównań
współczynników wymuszają,
że $r-R$ nie ma żadnych
skończonych biegunów.
Zatem $r-R=P\in\mathbb C[x]$.
Pozostaje zdecydować, czy
$P'+hP=S$ ma rozwiązanie wielomianowe.

Jeżeli $h$ ma dodatni stopień
$d$ i $S\ne0$, to
dla niezerowego $P$ mamy
$\deg(P'+hP)=\deg P+d$.
Wobec tego stopień $P$
musi być równy
$\deg S-d$. Jeśli liczba
ta jest ujemna, rozwiązania
nie ma; w przeciwnym razie
podstawiamy wielomian
$P=\sum_{k=0}^{\deg S-d}p_kx^k$
i rozstrzygamy skończony
\emph{liniowy} układ równań
na współczynniki $p_k$.
Jeśli $S=0$, rozwiązaniem
jest $P=0$, a innego
wielomianowego rozwiązania
jednorodnego nie ma.

Jeśli $h=c\in\mathbb C^\times$
jest stały, operator
$P\mapsto P'+cP$ jest
odwracalny na każdym
skończenie wymiarowym
zbiorze wielomianów
stopnia co najwyżej $\deg S$.
Można to zobaczyć bezpośrednio:
\[
 P=\sum_{k=0}^{\deg S}
                 \frac{(-1)^k S^{(k)}}{c^{k+1}}
\]
spełnia $P'+cP=S$, ponieważ
suma teleskopuje, a wyższe
pochodne $S$ znikają.
Dla $S=0$ przyjmujemy $P=0$.
We wszystkich przypadkach
wykonaliśmy skończoną liczbę
działań. Uzasadniliśmy też
konieczność i dostateczność
każdego testu.
\end{proof}

\begin{przyklad}[Obliczenie zakończone sukcesem]
\label{aa:prz:risch-e-success}
Dla $h=1$ i $s=1/x^2-1/x$
jedyny możliwy biegun $r$
jest prosty i znajduje się
w $0$. Z wyrazu $x^{-2}$
wynika $R=-1/x$.
Sprawdzamy
$R'+R=1/x^2-1/x=s$,
zatem $S=0$.
Wynik algorytmu jest
\[
 \int e^x\left(\frac1{x^2}-\frac1x\right)
 \,\mathrm dx=-\frac{e^x}{x}+C.
\]
Dla $h=1$, $s=1/x$
procedura kończy się negatywnie
od razu po wykryciu bieguna prostego.
\end{przyklad}

\begin{przyklad}[Całka Gaussa jako odpowiedź negatywna]
\label{aa:prz:risch-gauss-decision}
Dla $h=-2x$ i $s=1$
nie ma biegunów skończonych.
Pozostało $P'-2xP=1$.
Ponieważ $\deg S=0<1=\deg h$,
algorytm dowodzi nieistnienia
wymiernego $P$.
W połączeniu z kryterium
\ref{aa:tw:hyperexponential-criterion}
jest to dowód braku
elementarnej funkcji pierwotnej
$e^{-x^2}$, nie tylko niepowodzenie
heurystycznego programu.
\end{przyklad}

\begin{uwaga}[Pełny algorytm Rischa -- wynik klasyczny]
\label{aa:uw:full-risch-scope}
Właściwy algorytm Rischa
wychodzi daleko poza pole
$\mathbb C(x,e^{p(x)})$:
przetwarza odpowiednio
reprezentowane wieże
rozszerzeń elementarnych,
redukuje części algebraiczne
i przestępne i rozwiązuje
równania różniczkowe o bardziej
złożonych współczynnikach.
Jego zupełność w pełnym
zakresie przyjmujemy jako
wynik klasyczny (Risch;
Bronstein, \emph{Symbolic Integration I}).
W tym rozdziale udowodniono
kompletną procedurę \emph{tylko}
dla opisanej klasy wymiernej
i dla $h\in\mathbb C[x]$.
Nie twierdzimy, że
implementacje systemów CAS
realizują wszystkie przypadki
ogólnego twierdzenia.
\end{uwaga}

\section{Algorytm Kovacica: redukcja i klasyfikacja}

\begin{lemat}[Usunięcie pierwszej pochodnej]
\label{aa:lem:kovacic-normal-form}
Niech $a,b\in\mathbb C(x)$.
Lokalnie, w punkcie bez osobliwości
odpowiednich funkcji, podstawienie
$y=z\exp(-\tfrac12\int a(x)\,\mathrm dx)$
przekształca równanie
\[
 y''+a(x)y'+b(x)y=0
\]
w równanie $z''=q(x)z$, gdzie
\begin{equation}\label{aa:eq:kovacic-normal-q}
 q=\frac{a'}2+\frac{a^2}4-b
 \in\mathbb C(x).
\end{equation}
Podstawienie jest odwracalne
w lokalnym rozszerzeniu
liouvillowskim.
\end{lemat}
\begin{proof}
Niech $g=\exp(-\tfrac12\int a\,\mathrm dx)$;
wtedy $g'/g=-a/2$.
Dla $y=gz$ mamy
$y'=g(z'-az/2)$ oraz
$y''=g(z''-az'+(a^2/4-a'/2)z)$.
Po podstawieniu w równanie
znikają wyrazy przy $z'$,
a pozostaje
$z''+(b-a'/2-a^2/4)z=0$,
co jest równoważne podanej tezie.
Współczynnik $g$ jest niezerowy,
zatem podstawienie jest odwracalne.
\end{proof}

\begin{lemat}[Równanie Riccatiego]
\label{aa:lem:kovacic-riccati}
Niech $q\in\mathbb C(x)$ i
$z$ będzie niezerowym
lokalnym rozwiązaniem
$z''=qz$. Wówczas
$r=z'/z$ spełnia
$r'+r^2=q$.
Odwrotnie, jeśli element
$r$ dowolnego rozszerzenia
różniczkowego spełnia
$r'+r^2=q$, a
$z\ne0$ spełnia $z'/z=r$,
to $z''=qz$.
\end{lemat}
\begin{proof}
Z reguły ilorazu
$r'=z''/z-(z'/z)^2
=q-r^2$.
Odwrotnie, różniczkując
$z'=rz$ otrzymujemy
$z''=(r'+r^2)z=qz$.
\end{proof}

\begin{twierdzenie}[Trzy przypadki Kovacica -- wynik klasyczny]
\label{aa:tw:kovacic-classification}
Niech $q\in\mathbb C(x)$,
a współczynniki będą zadane
efektywnie. Dla równania
$z''=qz$ klasyczny algorytm
Kovacica rozstrzyga, czy istnieje
niezerowe rozwiązanie
liouvillowskie, i jeśli istnieje,
konstruuje takie rozwiązanie.
Wynik klasyfikacyjny dla
logarytmicznej pochodnej
$r=z'/z$ dopuszcza następujące
rodzaje przypadków:
\begin{enumerate}
\item $r\in\mathbb C(x)$
(przypadek redukowalny);
\item $r$ jest algebraiczne
nad $\mathbb C(x)$ stopnia
$2$ (przypadek grupy imprymitywnej
grupy Galois);
\item $r$ jest algebraiczne
stopnia $4$, $6$ lub $12$
(przypadek skończonej
prymitywnej grupy Galois);
\item nie ma niezerowego
rozwiązania liouvillowskiego.
\end{enumerate}
W przypadkach 1--3
algorytm wyszukuje skończoną
liczbę możliwości wynikających
z rozwinięć lokalnych
i sprawdza tożsamości
wymierne lub algebraiczne.
\end{twierdzenie}

\begin{uwaga}[Źródło i granice klasyfikacji]
\label{aa:uw:kovacic-citation}
Powyższy wynik cytujemy z:
J.~Kovacic, \emph{An algorithm
for solving second order linear
homogeneous differential equations},
\emph{Journal of Symbolic Computation}
2 (1986), 3--43.
Klasyfikacja korzysta z
podgrup algebraicznych
$\mathrm{SL}_2(\mathbb C)$
oraz z teorii Picarda--Vessiota
w rozdziale~\ref{aa:ch:picard-vessiot}.
Nie przedstawiamy rozgałęzień
2 i 3 jako samodzielnie
udowodnionych w tym rozdziale,
ani nie traktujemy testu
gałęzi 1 jako testu
liouvillowskości w całej klasie.
\end{uwaga}

\section{Gałąź wymierna Riccatiego dla wielomianowego współczynnika}

\begin{lemat}[Bieguny wymiernego rozwiązania Riccatiego]
\label{aa:lem:kovacic-rational-poles}
Niech $q\in\mathbb C[x]$
i niech $r\in\mathbb C(x)$
spełnia $r'+r^2=q$.
Każdy skończony biegun $r$
jest prosty i ma residuum $1$.
Istnieją więc
$\omega\in\mathbb C[x]$
oraz unormowany, niezerowy
wielomian $P\in\mathbb C[x]$
bez pierwiastków wielokrotnych
takie, że
\begin{equation}\label{aa:eq:kovacic-rational-r}
 r=\omega+\frac{P'}P.
\end{equation}
Co więcej, $P$ ma dokładnie
wszystkie skończone bieguny $r$
jako pierwiastki.
\end{lemat}
\begin{proof}
Jeśli $r$ ma biegun rzędu
$m\ge2$ w $a$, to $r^2$
ma biegun rzędu $2m$,
a $r'$ rzędu $m+1<2m$,
co przeczy regularności $q$.
Dla bieguna prostego piszemy
$r=c(x-a)^{-1}+O(1)$,
$c\ne0$. Współczynnik
przy $(x-a)^{-2}$ w
$r'+r^2$ wynosi
$c^2-c$; musi być zerowy,
więc $c=1$.
Odejmując od $r$
sumę $1/(x-a)$ po jego
skończonych biegunach
otrzymujemy funkcję wymierną
bez biegunów skończonych,
a więc wielomian $\omega$.
Jeśli $P$ jest iloczynem
$(x-a)$ po tych biegunach
(jeśli ich nie ma, $P=1$),
to $P'/P=\sum_a1/(x-a)$,
co daje tezę.
\end{proof}

\begin{twierdzenie}[Pełny test wymiernego Riccatiego dla $q$ wielomianowego]
\label{aa:tw:kovacic-polynomial-rational-decision}
Niech $q\in\mathbb C[x]$
będzie zadany efektywnie.
Istnieje skończony algorytm,
który rozstrzyga, czy
równanie
\begin{equation}\label{aa:eq:kovacic-poly-riccati}
 r'+r^2=q
\end{equation}
ma rozwiązanie $r\in\mathbb C(x)$.
Jeżeli $q$ jest stałe,
rozwiązaniem jest dowolny
stały pierwiastek kwadratowy $q$.
Jeżeli $\deg q$ jest nieparzysty,
nie ma rozwiązania wymiernego.
Gdy $\deg q=2d\ge2$,
procedura obejmuje dokładnie
dwa możliwe wielomiany
$\omega$ stopnia $d$,
obliczenie jednego wymuszonego
stopnia $n$ dla każdego
z nich, oraz rozwiązanie
skończonego układu
\emph{liniowego} na
współczynniki wielomianu $P$
stopnia $n$.
\end{twierdzenie}
\begin{proof}
W Lemacie~\ref{aa:lem:kovacic-rational-poles}
otrzymaliśmy postać
$r=\omega+P'/P$, gdzie
$\omega\in\mathbb C[x]$
i $P$ jest unormowany.
Podstawienie do Riccatiego daje
\begin{equation}\label{aa:eq:kovacic-polynomial-P}
 P''+2\omega P'
 +(\omega'+\omega^2-q)P=0.
\end{equation}
Jeżeli $\deg P=n$, to
$P'/P=n/x+O(x^{-2})$
w nieskończoności.
Jeśli $\omega$ ma dodatni
stopień $d$, to człony
$\omega'$ oraz
$2\omega(P'/P)$
mają stopień co najwyżej
$d-1$ w części wielomianowej,
a pochodna $P'/P$ i jego
kwadrat mają wyłącznie
ujemne potęgi.
Stąd wielomianowa część
$r'+r^2$ ma stopień $2d$
i współczynnik wiodący
równy kwadratowi
współczynnika wiodącego
$\omega$. W szczególności
$q$ musi mieć dodatni
\emph{parzysty} stopień.
Jeśli $\omega$ jest stały,
wielomianowa część
$r'+r^2$ jest stała,
więc nie otrzymamy
$q$ dodatniego stopnia.
To dowodzi przeszkody
dla nieparzystego stopnia.

Załóżmy teraz $\deg q=2d\ge2$.
Niech $a_d$ oznacza
współczynnik wiodący
$\omega$. Musi zachodzić
$a_d^2=[x^{2d}]q$, są
zatem dwa możliwe wybory
$a_d\in\mathbb C^\times$.
Dla każdego wybieramy
kolejno współczynniki
$a_{d-1},\ldots,a_0$
tak, aby
\[
 \deg(q-\omega^2)\le d-1.
\]
Współczynniki przy
$x^{2d-1},\ldots,x^d$
wyznaczają rekurencyjnie
wszystkie $a_j$,
ponieważ współczynnik
przy nieznanym $a_j$
wynosi $2a_d\ne0$.
Są to dokładnie
dwa możliwe wielomiany
$\omega$.

Porównajmy teraz
współczynnik przy
$x^{d-1}$ w tożsamości
\[
 q-\omega^2-\omega'
 =2\omega\frac{P'}P+
 \left(\frac{P'}P\right)'
 +\left(\frac{P'}P\right)^2.
\]
Ponieważ $P'/P=n/x+O(x^{-2})$,
jego prawa strona ma
współczynnik przy $x^{d-1}$
równy $2a_dn$. Dlatego
\begin{equation}\label{aa:eq:kovacic-degree-n}
 n=\frac{[x^{d-1}](q-\omega^2-\omega')}
          {2a_d}.
\end{equation}
Jeśli otrzymana liczba
nie należy do $\mathbb Z_{\ge0}$,
to wybór $\omega$ jest
niemożliwy.

W pozostałych przypadkach
podstawiamy
$P=x^n+\sum_{j=0}^{n-1}p_jx^j$
do~\eqref{aa:eq:kovacic-polynomial-P}.
Porównanie współczynników
daje skończony
\emph{liniowy} układ równań
dla $p_0,\ldots,p_{n-1}$.
Każde jego rozwiązanie
wyznacza wielomian $P$
spełniający tożsamość.
Nie może on mieć
wielokrotnych pierwiastków:
gdyby $a$ był pierwiastkiem
krotności $m\ge2$, wyraz
$P''$ w~\eqref{aa:eq:kovacic-polynomial-P}
miałby w $a$ rząd
$m-2$ z niezerowym współczynnikiem,
podczas gdy pozostałe
wyrazy miałyby rząd
co najmniej $m-1$.
Stąd konstrukcja $r$
z~\eqref{aa:eq:kovacic-rational-r}
jest poprawna i spełnia
\eqref{aa:eq:kovacic-poly-riccati}.

Na odwrót każde wymierne
rozwiązanie miało postać
z Lematu~\ref{aa:lem:kovacic-rational-poles},
a porównanie współczynników
wymuszało jeden z dwóch
wyborów $\omega$, ten sam
stopień $n$ i rozwiązanie
jednego z badanych układów.
Algorytm nie odrzuci więc
żadnego wymiernego rozwiązania.
Dla $q$ stałego zawsze
wystarcza stałe $r$
z $r^2=q$, co kończy dowód.
\end{proof}

\begin{przyklad}[Dwa wyniki pozytywne]
\label{aa:prz:kovacic-polynomial-success}
Dla $q=x^2+1$ możemy
wybrać $\omega=x$ i
\eqref{aa:eq:kovacic-degree-n}
daje $n=0$; wielomian
$P=1$ spełnia
\eqref{aa:eq:kovacic-polynomial-P}.
Otrzymujemy $r=x$ oraz
rozwiązanie
$z=e^{x^2/2}$.
Dla $q=x^2+3$ ten sam
wybór $\omega=x$
daje $n=1$, a
$P=x$ spełnia
$P''+2xP'+(1+x^2-(x^2+3))P=0$.
Stąd
\[
 r=x+\frac1x,\qquad
 z=xe^{x^2/2},\qquad
 z''=(x^2+3)z.
\]
Obydwa wyniki można
zweryfikować niezależnie
przez różniczkowanie.
\end{przyklad}

\begin{przyklad}[Dlaczego gałąź wymierna nie jest całą procedurą]
\label{aa:prz:kovacic-quadratic-example}
Rozważmy
$q=x+5/(16x^2)$.
W nieskończoności $q$
ma część wielomianową
stopnia nieparzystego,
więc równanie Riccatiego
nie ma \emph{wymiernego}
rozwiązania: gdyby miało,
wielomianowa część
$r'+r^2$ byłaby stopnia
parzystego albo stała.
Jednakże
\[
 r=\sqrt{x}-\frac1{4x}
\]
spełnia $r'+r^2=q$,
co sprawdzamy bezpośrednio:
składniki z $x^{-1/2}$
znoszą się, a przy
$x^{-2}$ zostaje $5/16$.
Jest to element
stopnia $2$ nad
$\mathbb C(x)$. Odpowiada
mu rozwiązanie
\[
 z=x^{-1/4}
       \exp\left(\frac23 x^{3/2}\right)
\]
na lokalnie wybranych
gałęziach. Dlatego
odrzucenie pierwszej gałęzi
nie jest dowodem braku
rozwiązań liouvillowskich.
\end{przyklad}

\begin{przyklad}[Równanie Airy'ego]
\label{aa:prz:kovacic-airy}
Dla $q=x$ mamy
$\deg q=1$, zatem test
wymierny daje odpowiedź negatywną.
Silniejsze stwierdzenie,
że równanie Airy'ego
$z''=xz$ \emph{nie ma żadnego}
niezerowego rozwiązania
liouvillowskiego,
udowodniliśmy już w
Twierdzeniu~\ref{aa:tw:pv-airy-group},
korzystając z kryterium
Picarda--Vessiota i
wykluczenia wszystkich
\emph{algebraicznych}
rozwiązań Riccatiego.
Nie przypisujemy tej
silniejszej konkluzji
samemu testowi wymiernemu.
\end{przyklad}

\section{Zestawienie i granice rozstrzygalności}

\begin{itemize}
\item Dla $f\in\mathbb C(x)$
procedura wymierna
z~\ref{aa:tw:rational-integration-algorithm}
wyznacza pochodną wymierną,
residua i logarytmy, a
test znikania residuów
rozstrzyga istnienie
\emph{wymiernej} funkcji pierwotnej.
\item Dla $s\in\mathbb C(x)$,
$h\in\mathbb C[x]\setminus\{0\}$
procedura~\ref{aa:tw:risch-polynomial-h-algorithm}
i kryterium AA-45
rozstrzygają elementarność
funkcji pierwotnej
$s\exp(\int h\,\mathrm dx)$.
\item Dla $z''=qz$ z
$q\in\mathbb C[x]$
procedura~\ref{aa:tw:kovacic-polynomial-rational-decision}
rozstrzyga istnienie
\emph{wymiernego} rozwiązania
Riccatiego i daje wtedy
rozwiązanie liouvillowskie.
Ogólna procedura Kovacica
dla $q\in\mathbb C(x)$
obejmuje też rozwiązania
algebraiczne stopni
$2,4,6,12$; jej pełna
poprawność jest cytowana,
nie udowodniona.
\end{itemize}

W żadnym z tych problemów
numeryczne przybliżenie
funkcji pierwotnej albo
rozwiązania równania
różniczkowego nie zastępuje
dokładnej odpowiedzi o
elementarności lub
liouvillowskości.

% AA-48: Wymierna gałąź Kovacica dla q w C(x).
% Rozdział samodzielnie dowodzi dokładnej i skończonej procedury
% dla WSZYSTKICH racjonalnych rozwiązań Riccatiego r'+r^2=q.
% Przypadki Kovacica 2 i 3 pozostają wynikami klasycznymi.
\chapter{Wymierna gałąź algorytmu Kovacica dla potencjałów wymiernych}
\chaptermark{Wymierna gałąź Kovacica}
\label{aa:ch:kovacic-rational-potential}

W rozdziale~\ref{aa:ch:risch-kovacic} skonstruowaliśmy
zupełny test wymiernych rozwiązań równania
$r'+r^2=q$ przy założeniu $q\in\mathbb C[x]$.
Dla $q\in\mathbb C(x)$ występują także bieguny
w skończonych punktach, a ich części główne
ograniczają możliwe rozwiązania.
Pokażemy, że również w tym przypadku
\emph{wszystkie wymierne rozwiązania}
można znaleźć przez skończoną liczbę
rozgałęzień algebraicznych i układów
\emph{liniowych}.
Zachowujemy rozróżnienie:
jest to kompletna \emph{pierwsza gałąź}
algorytmu Kovacica, nie pełny algorytm
obejmujący algebraiczne rozwiązania
Riccatiego wyższych stopni.

\section{Konwencja i lokalne dane}

\begin{definicja}[Współczynnik rozwinięcia w nieskończoności]
\label{aa:def:kov-rational-laurent}
Dla $f\in\mathbb C(x)$ przez
$[x^j]_\infty f$ oznaczamy
współczynnik przy $x^j$
w jednoznacznym rozwinięciu
Laurenta funkcji $f$ w nieskończoności.
Dla $j<0$ jest to współczynnik
części właściwej przy
$x^j$, a dla $j\ge0$
współczynnik części wielomianowej.
Przez $\operatorname{pp}_a(f)$
oznaczamy skończoną sumę wyrazów
o ujemnych potęgach $x-a$
w rozwinięciu Laurenta
w skończonym punkcie $a$.
\end{definicja}

\begin{uwaga}[Efektywność obliczeń]
\label{aa:uw:kov-rational-effective}
W twierdzeniu o algorytmie zakładamy,
że współczynniki wielomianów
dopuszczają dokładną arytmetykę,
test równości, rozkład wielomianów
na czynniki nad efektywnym
algebraicznym rozszerzeniem
oraz rozwiązywanie równań
kwadratowych. Wystarcza
na przykład podciało
$\overline{\mathbb Q}$.
Nie utożsamiamy dowolnie
nazwanych liczb zespolonych
bez reprezentacji z danymi
obliczalnymi.
\end{uwaga}

\section{Analiza biegunów skończonych}

\begin{lemat}[Możliwe rzędy biegunów]
\label{aa:lem:kov-rational-local-orders}
Niech $q,r\in\mathbb C(x)$
spełniają $r'+r^2=q$.
W każdym skończonym punkcie
$a\in\mathbb C$ zachodzi:
\begin{enumerate}
\item jeśli $r$ ma biegun rzędu
$m\ge2$, to $q$ ma w $a$
biegun dokładnie rzędu $2m$,
a współczynnik wiodący
$\beta_{-2m}$ jest kwadratem
współczynnika wiodącego
$c_{-m}$ funkcji $r$;
\item jeśli $r$ ma biegun prosty
$r=c_{-1}(x-a)^{-1}+O(1)$,
to $q$ ma biegun rzędu co
najwyżej $2$, a
$[ (x-a)^{-2}]_a q
=c_{-1}(c_{-1}-1)$;
\item jeżeli $q$ jest regularna w $a$,
a $r$ ma tam biegun,
to jest on prosty i ma residuum $1$.
\end{enumerate}
W szczególności, jeśli $q$ ma
w jakimkolwiek skończonym
punkcie biegun nieparzystego
rzędu co najmniej $3$,
równanie nie ma rozwiązania wymiernego.
\end{lemat}
\begin{proof}
Niech $z=x-a$.
Dla $r=c_{-m}z^{-m}+O(z^{-m+1})$,
$m\ge2$, kwadrat $r^2$
ma niezerowy wyraz
$c_{-m}^2z^{-2m}$,
a $r'$ ma biegun rzędu
$m+1<2m$. Wyraz ten
nie może się skasować,
co dowodzi punktu pierwszego.
Dla $m=1$ wyraz przy
$z^{-2}$ w $r'+r^2$
wynosi $-c_{-1}+c_{-1}^2$;
nie ma wyższych ujemnych potęg,
co daje punkt drugi.
Jeżeli $q$ jest regularna,
pierwszy punkt wyklucza
$m\ge2$, a przy $m=1$
otrzymujemy
$c_{-1}(c_{-1}-1)=0$.
Ponieważ $c_{-1}\ne0$,
jest równy $1$.
Jeśli $q$ ma biegun rzędu
nieparzystego co najmniej $3$,
nie może on pochodzić
ani z pierwszego, ani
z drugiego przypadku.
\end{proof}

\begin{lemat}[Skończone listy dopuszczalnych części głównych]
\label{aa:lem:kov-rational-local-candidates}
Niech $a$ będzie biegunem funkcji
$q\in\mathbb C(x)$ stopnia $k$.
Dla każdego wymiernego rozwiązania
$r'+r^2=q$ jego część główna
$\operatorname{pp}_a(r)$ należy
do skończonej, efektywnie
obliczalnej listy:
\begin{enumerate}
\item jeśli $k=1$, lista
ma jeden element $(x-a)^{-1}$;
\item jeśli $k=2$ i
$q=\beta_{-2}(x-a)^{-2}+
O((x-a)^{-1})$, lista składa
się z wyrazów $c(x-a)^{-1}$,
gdzie $c(c-1)=\beta_{-2}$;
\item jeśli $k=2m\ge4$,
lista liczy co najwyżej
dwa wielomiany Laurenta
$\sum_{j=1}^m c_{-j}(x-a)^{-j}$.
Dla każdego z dwóch znaków
$c_{-m}=\pm\sqrt{\beta_{-2m}}$
pozostałe $c_{-m+1},\ldots,c_{-1}$
wyznacza jednoznacznie
porównanie współczynników
przy potęgach
$(x-a)^{-2m+1},\ldots,(x-a)^{-m-1}$
w równaniu
$r'+r^2=q$.
\end{enumerate}
\end{lemat}
\begin{proof}
Jeśli $q$ ma biegun prosty,
to $r$ nie może być regularna,
bo $r'+r^2$ byłaby regularna.
Z Lematu~\ref{aa:lem:kov-rational-local-orders}
wynika prosty biegun $r$,
a brak wyrazu $(x-a)^{-2}$
w $q$ wymusza residuum $1$.
Jeśli $q$ ma biegun drugiego rzędu,
$r$ musi mieć biegun prosty,
którego współczynnik spełnia
podane równanie kwadratowe.

Dla $k=2m\ge4$ biegun $r$
musi mieć rząd $m$.
Niech $u=\operatorname{pp}_a(r)$
i zapiszmy $r=u+v$,
gdzie $v$ jest regularna w $a$.
Wtedy
\[
 r'+r^2=u'+u^2+v'+2uv+v^2.
\]
Ponieważ $u$ ma biegun rzędu $m$,
wyrazy $v'$, $2uv$, $v^2$
mają bieguny rzędu
co najwyżej $m$.
Dlatego współczynniki
przy $z^{-2m},\ldots,z^{-m-1}$
w $q$ zależą wyłącznie
od $u'+u^2$.
Współczynnik przy $z^{-2m}$
daje $c_{-m}^2=\beta_{-2m}$.
W każdym kolejnym równaniu
niewyznaczone dotąd
$c_{-m+j}$ występuje liniowo
ze współczynnikiem
$2c_{-m}\ne0$,
a reszta zależy tylko
od wcześniej wyznaczonych
współczynników.
Tym samym rekurencja
kończy się po $m-1$
krokach i nie wprowadza
dodatkowych wyborów.
\end{proof}

\begin{przyklad}[Biegun czwartego rzędu]
\label{aa:prz:kov-rational-high-pole}
Niech $q=x^{-4}-2x^{-3}$.
W punkcie $0$ mamy $m=2$,
a więc $c_{-2}=\pm1$.
Przy wyborze $c_{-2}=1$
współczynnik przy $x^{-3}$
w $u'+u^2$ wynosi
$-2+2c_{-1}$.
Porównanie z $-2$
daje $c_{-1}=0$,
czyli $u=1/x^2$.
Istotnie,
\[
 (1/x^2)' +(1/x^2)^2
 =-2/x^3+1/x^4=q.
\]
Nie należy traktować
takiego bieguna jako
prostego: jego rząd
determinuje już wiodący
wyraz równania.
\end{przyklad}

\section{Normalna postać wymiernego rozwiązania}

\begin{lemat}[Wydzielenie biegunów ruchomych]
\label{aa:lem:kov-rational-normal-form}
Niech $q,r\in\mathbb C(x)$
i $r'+r^2=q$.
Niech $S\subset\mathbb C$
będzie skończonym zbiorem
wszystkich biegunów $q$.
Dla $a\in S$ połóżmy
$R_a=\operatorname{pp}_a(r)$
oraz $R=\sum_{a\in S}R_a$.
Wówczas istnieją
$\omega\in\mathbb C[x]$
oraz unormowany, niezerowy,
bezkwadratowy $P\in\mathbb C[x]$
takie, że
\begin{equation}\label{aa:eq:kov-rational-normal-r}
 r=\omega+R+\frac{P'}P,
 \qquad
 P(a)\ne0\quad(a\in S).
\end{equation}
Pierwiastki $P$ to dokładnie
bieguny $r$ spoza $S$.
Jeśli $S=\varnothing$, przyjmujemy $R=0$.
\end{lemat}
\begin{proof}
Po odjęciu $R$ funkcja
$r-R$ jest regularna
we wszystkich punktach $S$.
W pozostałych punktach,
w których $q$ jest regularna,
na mocy
Lematu~\ref{aa:lem:kov-rational-local-orders}
każdy biegun $r$ jest prosty
i ma residuum $1$.
Niech $P$ będzie iloczynem
$(x-b)$ po wszystkich takich
pozostałych biegunach $b$;
jeśli ich nie ma, połóżmy $P=1$.
Wielomian $P$ jest unormowany,
bez pierwiastków wielokrotnych,
a żaden jego pierwiastek
nie należy do $S$.
Funkcja $P'/P=\sum_b1/(x-b)$
ma dokładnie części główne
$r-R$ poza $S$.
Zatem $r-R-P'/P$
nie ma skończonych biegunów,
jest więc wielomianem $\omega$.
\end{proof}

\begin{stwierdzenie}[Równanie liniowe dla wielomianu]
\label{aa:st:kov-rational-P-equation}
Niech $q\in\mathbb C(x)$
oraz $u\in\mathbb C(x)$.
Dla każdego niezerowego
$P\in\mathbb C[x]$
równoważne są warunki
\[
 \left(u+\frac{P'}P\right)'
 +\left(u+\frac{P'}P\right)^2=q
\]
oraz
\begin{equation}\label{aa:eq:kov-rational-P-ode}
 P''+2uP'+(u'+u^2-q)P=0.
\end{equation}
Jeśli zarówno $u$, jak i $q$
są regularne w $b$
i $P$ rozwiązuje drugie
równanie, to $P$ nie może
mieć w $b$ pierwiastka
wielokrotnego.
\end{stwierdzenie}
\begin{proof}
Rozwijamy lewą stronę:
\[
 u'+u^2+2u\frac{P'}P+\frac{P''}P;
\]
wyrazy $(P'/P)^2$
z pochodnej ilorazu
i kwadratu znoszą się.
Po pomnożeniu przez $P$
otrzymujemy drugie równanie.
Jeżeli $P$ ma w $b$ zero
rzędu $m\ge2$ i $u$ jest
regularna, to $P''$
ma rząd dokładnie $m-2$,
a pozostałe wyrazy
$2uP'$ i $(u'+u^2-q)P$
mają rzędy co najmniej
$m-1$ i $m$, pod warunkiem,
że $q$ jest regularna
w $b$. W zastosowaniach
poniżej $b$ będzie punktem
spoza biegunów $q$,
więc warunek jest spełniony.
Takie składniki nie mogą się
skasować.
\end{proof}

\section{Kontrola w nieskończoności}

Niech $R$ będzie wybraną sumą
części głównych w punktach
skończonych, a
$\rho=[x^{-1}]_\infty R$.
Wszystkie składniki $R$
są właściwymi funkcjami
wymiernymi, zatem
$R=\rho/x+O(x^{-2})$.
Dla unormowanego $P$
stopnia $n\ge0$ mamy
\begin{equation}\label{aa:eq:kov-rational-logarithmic-infinity}
 \frac{P'}P=\frac n x+O(x^{-2}).
\end{equation}

\begin{lemat}[Skończone możliwości przy nieskończoności]
\label{aa:lem:kov-rational-infinity}
Niech $q\in\mathbb C(x)$,
a $R$ będzie ustaloną sumą
dopuszczalnych części głównych
z Lematu~\ref{aa:lem:kov-rational-local-candidates}.
Załóżmy, że
$r=\omega+R+P'/P$
spełnia $r'+r^2=q$,
gdzie $\omega\in\mathbb C[x]$,
$P$ jest unormowany stopnia $n$
i nie zeruje się w biegunach $q$.
Wtedy zachodzi dokładnie jeden
z poniższych przypadków.
\begin{enumerate}
\item Jeśli część wielomianowa
$q$ ma dodatni stopień,
to jest on parzysty $2d\ge2$,
a $\deg\omega=d$.
Istnieją tylko dwa możliwe
wielomiany $\omega$, wyznaczone
przez warunek
$\deg\operatorname{pp}_\infty(q-\omega^2)
\le d-1$, gdzie
$\operatorname{pp}_\infty$
oznacza część wielomianową,
a ich współczynniki wiodące
są pierwiastkami kwadratowymi
wiodącego współczynnika $q$.
Dla każdego wyboru $\omega$,
gdy $a_d=[x^d]\omega\ne0$,
stopień $P$ jest wymuszony:
\begin{equation}\label{aa:eq:kov-rational-degree-positive}
 n=\frac{[x^{d-1}]_\infty
 \bigl(q-(\omega+R)'-(\omega+R)^2\bigr)}
 {2a_d}.
\end{equation}
\item Jeśli $q=c+O(x^{-1})$
przy $x\to\infty$, $c\ne0$,
to $\omega$ jest stała
i $\omega^2=c$.
Dla każdej z dwóch możliwości
$\omega$ koniecznie
\begin{equation}\label{aa:eq:kov-rational-degree-constant}
 n=\frac{[x^{-1}]_\infty
 \bigl(q-(\omega+R)'-(\omega+R)^2\bigr)}
 {2\omega}.
\end{equation}
\item Jeśli $q=O(x^{-1})$,
to $\omega=0$ i koniecznie
$q=O(x^{-2})$.
Przy $\beta=[x^{-2}]_\infty q$
liczba $n$ musi być
nieujemną liczbą całkowitą
spełniającą
\begin{equation}\label{aa:eq:kov-rational-degree-vanishing}
 (n+\rho)(n+\rho-1)=\beta.
\end{equation}
Istnieją co najwyżej
dwa takie wybory $n$.
\end{enumerate}
\end{lemat}
\begin{proof}
Z~\eqref{aa:eq:kov-rational-logarithmic-infinity}
wynika
$r=\omega+(\rho+n)/x+O(x^{-2})$.
Jeśli $\deg\omega=d\ge1$,
najwyższa potęga w
$r'+r^2$ ma stopień $2d$
i współczynnik będący
kwadratem wiodącego
współczynnika $\omega$.
Wszystkie potęgi
$x^{2d},\ldots,x^d$
pochodzą wyłącznie z
$\omega^2$; człony
$\omega'$, $2\omega R$,
$2\omega P'/P$ mają stopień
co najwyżej $d-1$.
Dopasowujemy więc kolejno
współczynniki $\omega$
od najwyższego w dół,
dla każdego z dwóch
znaków wiodącego pierwiastka.
Przy potędze $x^{d-1}$
pozostały wkład od $P'/P$
wynosi dokładnie $2a_dn$,
a $(P''/P)=O(x^{-2})$.
Stąd~\eqref{aa:eq:kov-rational-degree-positive}.

Jeśli $\omega$ jest niezerową
stałą, to wartość graniczna
$r'+r^2$ wynosi $\omega^2$
i musi być równa $c\ne0$.
Wyraz przy $x^{-1}$ w
równaniu po odjęciu
$(\omega+R)'+(\omega+R)^2$
pochodzi wyłącznie z
$2\omega P'/P$ i jest
równy $2\omega n$.
Daje to drugi wzór.

Jeśli $\omega=0$,
to $r=(\rho+n)/x+O(x^{-2})$.
Zatem $r'+r^2=O(x^{-2})$,
a współczynnik przy $x^{-2}$
wynosi
$-(\rho+n)+(\rho+n)^2$.
To jest dokładnie
\eqref{aa:eq:kov-rational-degree-vanishing}.
Równanie kwadratowe
ma najwyżej dwa rozwiązania.
Wymienione możliwości
wyczerpują wszystkie
stopnie części wielomianowej $r$.
\end{proof}

\section{Główne twierdzenie algorytmiczne}

\begin{twierdzenie}[Zupełny test wymiernej gałęzi Kovacica]
\label{aa:tw:kov-rational-complete-algorithm}
Niech $q\in\mathbb C(x)$,
a współczynniki będą podane
efektywnie w sensie
Uwagi~\ref{aa:uw:kov-rational-effective}.
Istnieje skończony algorytm
rozstrzygający, czy
\begin{equation}\label{aa:eq:kov-rational-main-riccati}
 r'+r^2=q
\end{equation}
ma rozwiązanie $r\in\mathbb C(x)$.
Jeżeli tak, algorytm
wyznacza co najmniej jedno
takie rozwiązanie, a następnie
lokalne niezerowe rozwiązanie
$z''=qz$ postaci
$z=\exp(\int r\,\mathrm dx)$.
\end{twierdzenie}
\begin{proof}
\emph{Krok 1: bieguny skończone.}
Rozkładamy mianownik $q$
i znajdujemy jego skończony
zbiór biegunów $S$
oraz ich rzędy.
Jeśli którykolwiek ma
nieparzysty rząd co najmniej
$3$, zgłaszamy
brak rozwiązania wymiernego,
zgodnie z
Lematem~\ref{aa:lem:kov-rational-local-orders}.
W przeciwnym razie dla
każdego $a\in S$ wyznaczamy
jedną lub dwie możliwe
części główne z
Lematu~\ref{aa:lem:kov-rational-local-candidates}.
Wykonujemy wszystkie skończenie
wiele kombinacji tych wyborów,
a dla każdej z nich
definiujemy $R=\sum_aR_a$.

\emph{Krok 2: nieskończoność.}
Dla wybranego $R$
stosujemy
Lemat~\ref{aa:lem:kov-rational-infinity}.
Uzyskujemy skończoną listę
par $(\omega,n)$.
Odrzucamy wszystkie,
dla których wymagany
stopień $n$ nie jest
nieujemną liczbą całkowitą.
Warunki z lematu są konieczne,
więc niczego nie tracimy.

\emph{Krok 3: liniowy układ dla $P$.}
Dla każdej pozostałej pary
stawiamy $u=\omega+R$ i
$P=x^n+\sum_{j=0}^{n-1}p_jx^j$.
Równanie~\eqref{aa:eq:kov-rational-P-ode},
po pomnożeniu przez
wspólny mianownik,
daje skończenie wiele równań
\emph{liniowych} na
$p_0,\ldots,p_{n-1}$.
Rozstrzygamy ich zgodność
dokładną eliminacją liniową.
Nie wystarcza wybrać
dowolnego rozwiązania
układu: dodatkowo żądamy
$P(a)\ne0$ dla każdego
$a\in S$.

Warunki te także
można sprawdzić efektywnie:
rozwiązania zgodnego
układu liniowego tworzą
afiniczną przestrzeń $V$.
Dla każdego $a\in S$
równanie $P(a)=0$
wyznacza podzbiór afiniczny
będący hiperpłaszczyzną
lub całym $V$.
Jeżeli dla pewnego $a$
całe $V$ jest w niej zawarte,
odrzucamy daną parę
$(\omega,n)$.
Jeśli dla żadnego $a$
tak nie jest, skończona
suma właściwych
hiperpłaszczyzn nie pokrywa
przestrzeni afinicznej
nad nieskończonym ciałem
$\mathbb C$. Możemy więc
wybrać $P\in V$ bez zer
w $S$, na przykład
sprawdzając kolejno
skończenie wiele
specjalizacji parametrów
w efektywnym nieskończonym
podciele współczynników.
Wszystkie pierwiastki $P$
poza $S$ są automatycznie
proste: gdyby $b\notin S$
był pierwiastkiem krotności
co najmniej $2$,
w równaniu~\eqref{aa:eq:kov-rational-P-ode}
człon $P''$ miałby
niekasowalny wyraz wiodący
rzędu o dwa mniejszego
niż rząd zera $P$,
bo $u$ i $q$ są w $b$
regularne.

\emph{Poprawność wyniku pozytywnego.}
Jeśli znaleziono $P$,
to $r=u+P'/P$ jest
funkcją wymierną i
ze Stwierdzenia
\ref{aa:st:kov-rational-P-equation}
spełnia dokładnie
\eqref{aa:eq:kov-rational-main-riccati}.
Z Lematu~\ref{aa:lem:kovacic-riccati}
dostajemy lokalne niezerowe
rozwiązanie $z''=qz$
po dołączeniu elementu
$z$ o $z'/z=r$.

\emph{Zupełność wyniku negatywnego.}
Niech $r\in\mathbb C(x)$
będzie dowolnym rozwiązaniem.
Na mocy
Lematów~\ref{aa:lem:kov-rational-local-orders}
i~\ref{aa:lem:kov-rational-local-candidates}
jego części główne pojawiają
się w jednej z badanych
kombinacji $R$.
Z Lematu~\ref{aa:lem:kov-rational-normal-form}
otrzymujemy $r=\omega+R+P'/P$
z unormowanym $P$
bez zer na $S$.
Lemat~\ref{aa:lem:kov-rational-infinity}
zapewnia, że odpowiednie
$(\omega,\deg P)$ jest
na naszej skończonej liście.
Wreszcie $P$ spełnia
liniowe równanie
\eqref{aa:eq:kov-rational-P-ode}
oraz wszystkie warunki
$P(a)\ne0$.
Algorytm zatem znajdzie
co najmniej jedną
poprawną parę i nie
może zakończyć się odpowiedzią
negatywną.
Dowiedliśmy obu kierunków
oraz skończoności całej procedury.
\end{proof}

\begin{uwaga}[Zakres wyniku]
\label{aa:uw:kov-rational-not-all}
Twierdzenie
\ref{aa:tw:kov-rational-complete-algorithm}
rozstrzyga istnienie
\emph{wymiernej} logarytmicznej
pochodnej rozwiązania
$z''=qz$ dla każdego
$q\in\mathbb C(x)$.
Nie rozstrzyga samo
istnienia rozwiązania
liouvillowskiego w ogólności:
gałęzie algorytmu Kovacica,
w których pochodna
logarytmiczna jest algebraiczna
stopnia $2$, $4$, $6$ lub $12$,
wymagają dalszych procedur.
Ogólną klasyfikację przytoczono
w~\ref{aa:tw:kovacic-classification}
jako twierdzenie klasyczne.
\end{uwaga}

\section{Przykłady z pełnymi rachunkami}

\begin{przyklad}[Biegun prosty potencjału]
\label{aa:prz:kov-rational-simple}
Weźmy
\[
 q=x^2+3+\frac2{x-1}.
\]
Dla $a=1$ potencjał ma
biegun prosty, więc
$R=1/(x-1)$.
Jego część wielomianowa
ma stopień $2$, zatem
$\omega=\pm x$.
Dla $\omega=x$ mamy
$u=x+1/(x-1)$ oraz
\[
 u'+u^2
 =x^2+3+\frac2{x-1}=q.
\]
Lemat o nieskończoności
daje $n=0$ i $P=1$.
Zatem $r=x+1/(x-1)$,
a lokalnym rozwiązaniem
jest $z=(x-1)e^{x^2/2}$.
Sprawdzenie $z''/z=q$
można wykonać bezpośrednio
lub przez równanie Riccatiego.
\end{przyklad}

\begin{przyklad}[Dwa możliwe residua]
\label{aa:prz:kov-rational-double}
Niech $q=2/x^2$.
Przy $a=0$ jest biegun
rzędu $2$ z
$\beta_{-2}=2$.
Równanie $c(c-1)=2$
daje $c=2$ albo $c=-1$.
Dla $R=2/x$ przy
nieskończoności mamy
$\rho=2$, $\beta=2$;
równanie
$(n+2)(n+1)=2$
ma dopuszczalne $n=0$.
Z $P=1$ otrzymujemy
$r=2/x$.
Analogicznie drugi wybór
$R=-1/x$, $n=0$
daje $r=-1/x$.
Oba spełniają
\[
 r'+r^2=\frac2{x^2}.
\]
Odpowiadają rozwiązaniom
$z_1=x^2$ oraz $z_2=x^{-1}$
równania $z''=(2/x^2)z$.
\end{przyklad}

\begin{przyklad}[Brak rozwiązania i przypadek algebraiczny]
\label{aa:prz:kov-rational-negative}
Dla $q=1/x^3$
biegun w $0$ ma rząd
nieparzysty większy od $2$,
więc żadne rozwiązanie
Riccatiego nie jest wymierne.
Nie wynika z tego
nieistnienie wszystkich
rozwiązań liouvillowskich.

Dla $q=x+5/(16x^2)$
pierwsza gałąź również
jest wykluczona,
ponieważ część wielomianowa
$q$ ma nieparzysty
dodatni stopień.
Jednocześnie
$r=\sqrt{x}-1/(4x)$
jest algebraicznym,
niewymiernym rozwiązaniem
Riccatiego i prowadzi
do liouvillowskiego
rozwiązania $z$ opisanego
w Przykładzie
\ref{aa:prz:kovacic-quadratic-example}.
Pokazuje to konieczność
rozróżniania odpowiedzi
pierwszej gałęzi
od odpowiedzi całego
algorytmu Kovacica.
\end{przyklad}

\begin{przyklad}[Biegun czwartego rzędu i zachowanie w nieskończoności]
\label{aa:prz:kov-rational-fourth}
Dla $q=x^{-4}-2x^{-3}$
wybór $R=1/x^2$
pochodzi z analizy
bieguna czwartego rzędu.
Mamy $\rho=0$, a
$q=O(x^{-3})$, więc
$\beta=[x^{-2}]_\infty q=0$.
Równanie
$n(n-1)=0$ daje
kandydatów $n=0,1$.
Przy $n=0$ można
wybrać $P=1$,
a wówczas $r=1/x^2$
jest rozwiązaniem.
Wystarcza jeden
pozytywny kandydat;
algorytm nie musi
wyliczać wszystkich
możliwych rozwiązań,
aby rozstrzygnąć istnienie.
\end{przyklad}

\section{Kontrola połączenia z teorią Galois}

Jeżeli algorytm zwraca
wymierne $r\in\mathbb C(x)$,
to w rozszerzeniu
Picarda--Vessiota równania
$z''=qz$ istnieje
niezerowe rozwiązanie
o logarytmicznej pochodnej
$r$. Jego jednowymiarowa
przestrzeń rozwiązań
jest zachowywana przez
grupę Galois, więc
reprezentacja grupy
w odpowiedniej bazie
jest trójkątna.
Jest to kierunek
\emph{od wymiernego Riccatiego
do redukowalności reprezentacji}.
Przejście odwrotne
w ogólnym języku teorii
Picarda--Vessiota korzysta
z klasycznego twierdzenia
o niezmienniczej prostej
i nie jest tu nowym
dowodem strukturalnym
pełnej teorii Galois.

Dalsze przypadki
imprymitywne i skończone
należą do pełnego algorytmu
Kovacica, a nie do
niniejszego rozdziału.

% AA-49 | Kwadratowy przypadek algorytmu Kovacica
% Klasyczne listy lokalne Ec i Einfty są cytowane z Kovacica (1986).
% Samodzielne dowody: symetryczny kwadrat, nieliniowy warunek śladu,
% odzyskanie obu pierwiastków Riccatiego i liniowość równania pomocniczego.
\chapter{Kwadratowy przypadek algorytmu Kovacica}
\chaptermark{Kwadratowy przypadek Kovacica}
\label{aa:ch:kovacic-quadratic}

W rozdziale~\ref{aa:ch:kovacic-rational-potential}
otrzymaliśmy kompletny algorytm znajdujący
wymierne rozwiązania równania Riccatiego
$r'+r^2=q$ dla każdego $q\in\mathbb C(x)$.
Nie rozstrzyga to istnienia rozwiązań
algebraicznych stopnia drugiego nad
$\mathbb C(x)$, co pokazuje już
Przykład~\ref{aa:prz:kovacic-quadratic-example}.
W niniejszym rozdziale badamy
\emph{drugą gałąź algorytmu Kovacica}:
rozwiązania Riccatiego w kwadratowych
rozszerzeniach $\mathbb C(x)$.
Przygotowujemy równanie pomocnicze
trzeciego rzędu i wykonujemy wszystkie
rachunki odzyskujące rozwiązania.
Skończoność katalogu możliwych
lokalnych danych wejściowych
opiera się natomiast na
\emph{klasycznym twierdzeniu Kovacica},
które jawnie cytujemy; nie przedstawiamy
jego ogólnego dowodu w tym rozdziale.

\section{Ślad i wyróżnik dwóch rozwiązań Riccatiego}

\begin{definicja}[Przypadek kwadratowy Riccatiego]
\label{aa:def:kov2-quadratic}
Niech $K=\mathbb C(x)$, $D=\mathrm d/\mathrm dx$,
a $q\in K$. Rozwiązanie $r$
równania $Dr+r^2=q$ w rozszerzeniu
algebraicznym $L/K$ nazywamy
\emph{ściśle kwadratowym}, jeśli
$[K(r):K]=2$. W takim przypadku
oznaczamy przez $\bar r$ drugi
pierwiastek jego wielomianu minimalnego.
Wówczas $\bar r$ również spełnia
równanie Riccatiego: jedyne
przedłużenie $D$ do
separowalnego rozszerzenia
algebraicznego komutuje z
automorfizmem sprzężenia.
\end{definicja}

\begin{lemat}[Ślad, iloczyn i wyróżnik]
\label{aa:lem:kov2-trace-discriminant}
Niech $r,\bar r$ będą dwoma
różnymi sprzężonymi nad $K$
rozwiązaniami $Dr+r^2=q$,
gdzie $q\in K$.
Oznaczmy
$\phi=r+\bar r\in K$ oraz
$\Delta=(r-\bar r)^2\in K$.
Wówczas
\begin{equation}\label{aa:eq:kov2-discriminant}
 \Delta=4q-2D\phi-\phi^2,\qquad
 D\Delta=-2\phi\Delta,
\end{equation}
a ich wielomian minimalny ma postać
\begin{equation}\label{aa:eq:kov2-polynomial}
 T^2-\phi T+
 \frac{D\phi+\phi^2-2q}{2}.
\end{equation}
Ponadto $\Delta\ne0$ i
$\Delta$ nie jest kwadratem w $K$.
\end{lemat}
\begin{proof}
Ponieważ oba pierwiastki spełniają
Riccatiego,
\[
 D\phi=2q-r^2-\bar r^2
      =2q-\phi^2+2r\bar r,
\]
zatem
$r\bar r=(D\phi+\phi^2-2q)/2$.
Stąd
$\Delta=\phi^2-4r\bar r
=4q-2D\phi-\phi^2$
i otrzymujemy wielomian~\eqref{aa:eq:kov2-polynomial}.
Odejmując równania Riccatiego
dla obu pierwiastków,
otrzymujemy
$D(r-\bar r)=-(r+\bar r)(r-\bar r)$.
Różniczkowanie kwadratu daje
$D\Delta=-2\phi\Delta$.
Gdyby $\Delta$ była
kwadratem w $K$, równanie
\eqref{aa:eq:kov2-polynomial}
rozkładałoby się na czynniki
liniowe nad $K$, wbrew
$[K(r):K]=2$.
\end{proof}

\begin{lemat}[Wielomianowe równanie śladu]
\label{aa:lem:kov2-trace-equation}
Niech $q,\phi\in K$.
Jeżeli $\phi$ jest śladem dwóch
różnych rozwiązań Riccatiego,
to
\begin{equation}\label{aa:eq:kov2-trace}
 D^2\phi+3\phi D\phi+\phi^3
     -4q\phi-2Dq=0.
\end{equation}
Odwrotnie, jeśli $\phi\in K$
spełnia~\eqref{aa:eq:kov2-trace}, to
$\Delta=4q-2D\phi-\phi^2$
spełnia $D\Delta=-2\phi\Delta$.
Jeśli $\Delta\ne0$, dwa elementy
\begin{equation}\label{aa:eq:kov2-root}
 r_\pm=\frac{\phi\pm\sqrt{\Delta}}2
\end{equation}
w rozszerzeniu
$K(\sqrt{\Delta})$ spełniają
równanie Riccatiego.
Jeżeli $\Delta$ nie jest kwadratem
w $K$, oba rozwiązania są
ściśle kwadratowe i sprzężone.
\end{lemat}
\begin{proof}
Z pierwszej równości
\eqref{aa:eq:kov2-discriminant}
otrzymujemy po różniczkowaniu
\[
 D\Delta=4Dq-2D^2\phi-2\phi D\phi.
\]
Po dodaniu $2\phi\Delta$ prawa
strona wynosi
$-2(D^2\phi+3\phi D\phi+\phi^3
-4q\phi-2Dq)$.
Zatem warunek~\eqref{aa:eq:kov2-trace}
jest równoważny
$D\Delta=-2\phi\Delta$.
Załóżmy teraz $\Delta\ne0$
i wybierzmy $\delta^2=\Delta$.
W charakterystyce $0$
jedyna przedłużona derywacja
spełnia
$D\delta=(D\Delta)/(2\delta)
=-\phi\delta$.
Dla $r_\pm=(\phi\pm\delta)/2$ mamy
\[
 Dr_\pm+r_\pm^2
 =\frac{D\phi}{2}
   +\frac{\phi^2+\Delta}{4}=q;
\]
wyrazy zawierające $\delta$
znoszą się dzięki $D\delta=-\phi\delta$.
Dodatkowy warunek niekwadratowości
$\Delta$ jest dokładnie warunkiem
$[K(r_\pm):K]=2$.
\end{proof}

\begin{uwaga}[Weryfikacja negatywna]
\label{aa:uw:kov2-quadratic-versus-rational}
Równanie~\eqref{aa:eq:kov2-trace}
może mieć wymierne rozwiązanie
$\phi$ również wtedy, gdy
$\Delta$ jest kwadratem w $K$
albo gdy $\Delta=0$.
Wtedy
\eqref{aa:eq:kov2-root}
daje rozwiązania wymierne,
które należą już do
gałęzi~\ref{aa:tw:kov-rational-complete-algorithm}.
Dlatego sam test trzeciego
rzędu \emph{bez sprawdzenia
wyróżnika} nie dowodzi,
że otrzymaliśmy nowy
przypadek kwadratowy.
\end{uwaga}

\section{Symetryczny kwadrat równania liniowego}

\begin{lemat}[Równanie trzeciego rzędu]
\label{aa:lem:kov2-symmetric-square}
Niech $q$ należy do
ciała różniczkowego charakterystyki
$0$, a $z_1,z_2$ spełniają
$D^2z_j=qz_j$ w jego
rozszerzeniu różniczkowym.
Wówczas ich iloczyn
$v=z_1z_2$ spełnia
\begin{equation}\label{aa:eq:kov2-third-order}
 D^3v-4qDv-2(Dq)v=0.
\end{equation}
Jeżeli $v\ne0$, to
$\phi=Dv/v$ spełnia
\eqref{aa:eq:kov2-trace}.
\end{lemat}
\begin{proof}
Różniczkując iloczyn,
dostajemy
$D^2v=2qv+2(Dz_1)(Dz_2)$.
Po kolejnym różniczkowaniu:
\[
 D^3v=2(Dq)v+2qDv
 +2q\bigl(z_1Dz_2+(Dz_1)z_2\bigr)
 =2(Dq)v+4qDv.
\]
Ponadto z $Dv=\phi v$ wynika
$D^3v=(D^2\phi+3\phi D\phi
+\phi^3)v$.
Po podstawieniu do poprzedniej
równości i podzieleniu
przez $v$ otrzymujemy
\eqref{aa:eq:kov2-trace}.
\end{proof}

\begin{lemat}[Liniowe równanie dla wielomianu]
\label{aa:lem:kov2-polynomial-equation}
Niech $q,\theta\in\mathbb C(x)$,
a $P\in\mathbb C[x]$ będzie
niezerowy. Połóżmy
$\phi=\theta+P'/P$.
Wówczas $\phi$ spełnia
\eqref{aa:eq:kov2-trace}
wtedy i tylko wtedy, gdy
\begin{equation}\label{aa:eq:kov2-P-equation}
\begin{split}
 P''' +3\theta P''
  +(3\theta^2+3\theta'-4q)P'
  +\bigl(\theta''+3\theta\theta'
         +\theta^3-4q\theta-2q'\bigr)P=0.
\end{split}
\end{equation}
Po pomnożeniu przez
wspólny mianownik jest to
\emph{liniowy} układ równań
na współczynniki $P$
ustalonego stopnia.
\end{lemat}
\begin{proof}
Niech $v$ spełnia $Dv=\theta v$,
$v\ne0$, i rozważmy $w=Pv$.
Wtedy $Dw/v=P'+\theta P$,
a iterując otrzymujemy
\[
 D^3w/v=P'''+3\theta P''
 +(3\theta^2+3\theta')P'
 +(\theta''+3\theta\theta'+\theta^3)P.
\]
Po podstawieniu w równanie
\eqref{aa:eq:kov2-third-order}
widzimy, że $w$ jest jego
rozwiązaniem wtedy i tylko wtedy,
gdy zachodzi
\eqref{aa:eq:kov2-P-equation}.
Ponieważ
$Dw/w=\theta+P'/P=\phi$,
równoważność z
\eqref{aa:eq:kov2-trace}
wynika z poprzedniego lematu
i tożsamości na pochodne
logarytmiczne. Wszystkie
współczynniki równania
\eqref{aa:eq:kov2-P-equation}
są ustalonymi funkcjami
wymiernymi, więc porównanie
współczynników po usunięciu
mianowników daje liniowy
układ algebraiczny.
\end{proof}

\section{Klasyczne lokalne listy Kovacica}

\begin{definicja}[Rząd w nieskończoności]
\label{aa:def:kov2-infinity-order}
Niech $q=p/s\in\mathbb C(x)$,
gdzie $p,s\in\mathbb C[x]$
są względnie pierwsze,
$s\ne0$, $p\ne0$.
\emph{Rzędem $q$ w nieskończoności}
nazywamy liczbę
$\operatorname{ord}_{\infty}q
=\deg s-\deg p$.
Jeżeli $\operatorname{ord}_{\infty}q=2$,
przez $b_\infty$ oznaczamy
współczynnik przy $x^{-2}$
w rozwinięciu Laurenta $q$
w nieskończoności.
Dla $q=0$ stosujemy
najpierw gałąź wymierną
z rozdziału~\ref{aa:ch:kovacic-rational-potential}.
\end{definicja}

\begin{twierdzenie}[Listy lokalne przypadku drugiego -- wynik klasyczny]
\label{aa:tw:kov2-classical-candidates}
Niech $q\in\mathbb C(x)$ będzie
niezerowe, a $\Gamma$ skończonym
zbiorem wszystkich jego biegunów
w $\mathbb C$.
Dla każdego $a\in\Gamma$
definiujemy skończony zbiór
$E_a\subseteq\mathbb Z$:
\begin{enumerate}
\item jeśli biegun ma rząd $1$,
to $E_a=\{4\}$;
\item jeśli biegun ma rząd $2$,
a $b_a$ jest współczynnikiem
przy $(x-a)^{-2}$, to
\[
 E_a=\{2,\;2+2\sqrt{1+4b_a},
                 \;2-2\sqrt{1+4b_a}\}
                 \cap\mathbb Z;
\]
\item jeśli biegun ma rząd $m>2$,
to $E_a=\{m\}$.
\end{enumerate}
Również w nieskończoności
definiujemy zbiór:
\[
 E_\infty=
 \begin{cases}
 \{0,2,4\},
   &\operatorname{ord}_\infty q>2,\\
 \{2,\;2+2\sqrt{1+4b_\infty},
              \;2-2\sqrt{1+4b_\infty}\}
              \cap\mathbb Z,
   &\operatorname{ord}_\infty q=2,\\
 \{\operatorname{ord}_\infty q\},
   &\operatorname{ord}_\infty q<2.
 \end{cases}
\]
Jeśli równanie $r'+r^2=q$
ma ściśle kwadratowe
rozwiązanie nad $\mathbb C(x)$,
to istnieją wybory
$e_a\in E_a$ oraz
$e_\infty\in E_\infty$
takie, że
\begin{equation}\label{aa:eq:kov2-degree}
 n=\frac12\left(e_\infty-
                      \sum_{a\in\Gamma}e_a\right)
 \in\mathbb Z_{\ge0}
\end{equation}
oraz dla
\begin{equation}\label{aa:eq:kov2-theta}
 \theta=\frac12\sum_{a\in\Gamma}
                     \frac{e_a}{x-a}
\end{equation}
istnieje unormowany
$P\in\mathbb C[x]$
stopnia $n$ spełniający
\eqref{aa:eq:kov2-P-equation}.
\end{twierdzenie}

\begin{uwaga}[Status dowodu list lokalnych]
\label{aa:uw:kov2-local-cited}
To twierdzenie cytujemy
w standardowej postaci
drugiego przypadku algorytmu
Kovacica: J.~Kovacic,
\emph{An algorithm for solving
second order linear homogeneous
differential equations},
\emph{Journal of Symbolic Computation}
2 (1986), 3--43.
Dla biegunów drugiego rzędu
liczby $2\pm2\sqrt{1+4b_a}$
wynikają z lokalnych
wykładników równania
$z''=qz$, ale pełny
dowód wyczerpującej listy
dla wszystkich rzędów,
w tym rozgałęzienia
w punktach osobliwych,
wymaga analizy lokalnej
użytej przez Kovacica.
Właśnie ta konieczność
i wyczerpujący charakter
list $E_a,E_\infty$
\emph{nie są tutaj
samodzielnie dowodzone}.
Resztę algorytmu
wyprowadzamy i sprawdzamy
wprost z
Lematów~\ref{aa:lem:kov2-trace-equation}
i~\ref{aa:lem:kov2-polynomial-equation}.
\end{uwaga}

\section{Skończona procedura drugiej gałęzi}

\begin{twierdzenie}[Druga gałąź Kovacica, z cytowanym twierdzeniem lokalnym]
\label{aa:tw:kov2-algorithm}
Niech $q\in\mathbb C(x)$,
a jego współczynniki
będą efektywnie zadane
w sensie
Uwagi~\ref{aa:uw:kov-rational-effective}.
Przyjmując klasyczne
Twierdzenie~\ref{aa:tw:kov2-classical-candidates},
istnieje skończony algorytm
rozstrzygający, czy
$r'+r^2=q$ ma
rozwiązanie ściśle kwadratowe
nad $\mathbb C(x)$,
i konstruujący takie
rozwiązanie, jeśli istnieje.
\end{twierdzenie}
\begin{proof}
Dla $q=0$ dowolne rozwiązanie
wymierne zapewnia gałąź
\ref{aa:tw:kov-rational-complete-algorithm};
ewentualne przypadki
ściśle kwadratowe
wykluczamy bezpośrednio:
jeżeli $r'+r^2=0$ i $r\ne0$,
to $D(1/r)=1$,
stąd $1/r=x+c$, gdzie
$c$ jest stałą w algebraicznym
rozszerzeniu $\mathbb C(x)$.
Ponieważ $\mathbb C$ jest
algebraicznie domknięte,
$c\in\mathbb C$, a więc
$r=1/(x+c)\in\mathbb C(x)$.
Dla $r=0$ również
nie ma stopnia drugiego.

Niech dalej $q\ne0$.
Wyznaczamy bieguny
$q$, ich rzędy oraz
$\operatorname{ord}_\infty q$.
Tworzymy skończone
listy $E_a,E_\infty$
z cytowanego twierdzenia.
Rozważamy wszystkie
skończenie liczne kombinacje
$(e_\infty,(e_a)_{a\in\Gamma})$;
odrzucamy te, dla których
$n$ określone w~\eqref{aa:eq:kov2-degree}
nie należy do
$\mathbb Z_{\ge0}$.

Dla każdej pozostałej
kombinacji obliczamy $\theta$
z~\eqref{aa:eq:kov2-theta}
i podstawiamy
$P=x^n+\sum_{j<n}p_jx^j$
do równania
\eqref{aa:eq:kov2-P-equation}.
Po pomnożeniu przez wspólny
mianownik otrzymujemy
skończony liniowy układ
nad efektywnym ciałem
współczynników. Jeśli nie
ma rozwiązań, odrzucamy
daną kombinację.
Jeśli ma, jego zbiór
rozwiązań stanowi
afiniczną przestrzeń $V$.

Dla każdego $P\in V$
obliczamy
$\phi=\theta+P'/P$
oraz
$\Delta=4q-2\phi'-\phi^2$.
Jeżeli $\Delta\ne0$
i nie jest kwadratem
w $\mathbb C(x)$,
Lemat~\ref{aa:lem:kov2-trace-equation}
daje ściśle kwadratowe
rozwiązania
$r_\pm=(\phi\pm\sqrt{\Delta})/2$.
Test kwadratowości funkcji
wymiernej jest efektywny:
w jej zredukowanym
rozkładzie na
czynniki wykładniki
wszystkich czynników
muszą być parzyste;
ponieważ $\mathbb C$
jest algebraicznie domknięte,
stały współczynnik
nie stanowi przeszkody.

Pozostaje uzasadnić,
że istnienie odpowiedniego
$P\in V$ można
rozstrzygnąć bez
nieskończonego przeszukiwania.
Po podstawieniu
$P=\sum p_jx^j$
i usunięciu mianowników
warunek $\Delta=0$
albo
$\Delta$ jest kwadratem
w $\mathbb C(x)$
opisuje algebraiczne
warunki na parametry
$P$ (można użyć
rozkładu wielomianów,
NWD i kwantyfikatorów
nad algebraicznie
domkniętym ciałem).
Dla ustalonego stopnia $P$
jest to skończony problem
eliminacji algebraicznej:
rozstrzyga się, czy
afininiczny zbiór $V$
zawiera punkt, dla
którego $\Delta$ jest
niezerowa i niewymiernie
pierwiastkowalna.
Równoważnie można
skorzystać z pełnego
rozstrzygania zdań
w teorii ciał algebraicznie
domkniętych charakterystyki
$0$; w warunkach efektywnej
prezentacji współczynników
jest to procedura
kończąca się.

\emph{Dostateczność:}
każde znalezione $P$,
po obliczeniu wyróżnika
i wybraniu pierwiastka,
daje dokładnie rozwiązanie
Riccatiego i lokalne
rozwiązanie $z''=qz$
postaci
$z=\exp(\int r\,\mathrm dx)$,
na mocy
Lematu~\ref{aa:lem:kov2-trace-equation}.
\emph{Konieczność:}
jeśli istnieje rozwiązanie
ściśle kwadratowe,
to klasyczne
Twierdzenie~\ref{aa:tw:kov2-classical-candidates}
zapewnia jedną z
badanych kombinacji
i pewien unormowany
wielomian $P$
spełniający układ.
Jego ślad $\phi$
ma wyróżnik niekwadratowy,
a zatem jest dopuszczony
przez końcowy test.
To dowodzi skończoności
oraz obu kierunków
\emph{przy wyraźnie
cytowanym twierdzeniu lokalnym}.
\end{proof}

\begin{uwaga}[O kosztach i implementacji]
\label{aa:uw:kov2-computation}
Dla danych z prostymi
parametrami w większości
przykładów liniowy układ
wyznacza jedyne $P$.
Dopuszczenie ogólnych
parametrycznych rozwiązań
i rozstrzyganie
kwadratowości wyróżnika
wymaga dodatkowej kontroli,
nie zaś arbitralnego
wyboru pierwszego
rozwiązania układu.
Użycie eliminacji
algebraicznej zapewnia
teoretyczną skończoność,
ale nie jest twierdzeniem
o praktycznej wydajności.
W szczególności
\emph{testy symboliczne
nie zastępują tego dowodu}.
\end{uwaga}

\section{Przykłady z pełnym rachunkiem}

\begin{przyklad}[Kwadratowy pierwiastek i wielomian stopnia zero]
\label{aa:prz:kov2-zero-degree}
Rozważmy
\[
 q=x+\frac5{16x^2}.
\]
W punkcie $0$ występuje
biegun drugiego rzędu
o $b_0=5/16$.
Stąd
$\sqrt{1+4b_0}=3/2$
i $E_0=\{-1,2,5\}$.
W nieskończoności
$\operatorname{ord}_\infty q=-1$,
więc $E_\infty=\{-1\}$.
Wybierając $e_0=-1$
i $e_\infty=-1$,
otrzymujemy $n=0$ oraz
\[
 \theta=-\frac1{2x},\qquad P=1.
\]
Bezpośrednie podstawienie
do równania
\eqref{aa:eq:kov2-P-equation}
daje zero. Zatem
$\phi=-1/(2x)$ oraz
\[
 \Delta=4q-2\phi'-\phi^2=4x.
\]
Nie jest to kwadrat w
$\mathbb C(x)$.
Otrzymujemy dokładnie
\[
 r_\pm=-\frac1{4x}\pm\sqrt{x},
 \qquad r_\pm'+r_\pm^2=q.
\]
Odpowiadające
lokalne rozwiązania
to
$x^{-1/4}\exp(\pm\tfrac23 x^{3/2})$.
\end{przyklad}

\begin{przyklad}[Niezerowy stopień wielomianu pomocniczego]
\label{aa:prz:kov2-positive-degree}
Niech
\[
 q=\frac{4x+5}{16x^2}
   =\frac1{4x}+\frac5{16x^2}.
\]
Przy $x=0$ znów
$E_0=\{-1,2,5\}$,
lecz teraz
$\operatorname{ord}_\infty q=1$
i $E_\infty=\{1\}$.
Dla $e_0=-1$,
$e_\infty=1$
wzór~\eqref{aa:eq:kov2-degree}
daje $n=1$, a
$\theta=-1/(2x)$.
Podstawiamy
$P=x+p_0$ do równania
\eqref{aa:eq:kov2-P-equation}
i porównujemy współczynniki.
Otrzymujemy $p_0=-1$,
czyli $P=x-1$.
W szczególności
\[
 \phi=-\frac1{2x}+\frac1{x-1}
      =\frac{x+1}{2x(x-1)},
 \qquad
 \Delta=\frac{x}{(x-1)^2}.
\]
Wyróżnik nie jest
kwadratem w $\mathbb C(x)$,
ponieważ $x$ ma
w $0$ nieparzystą
krotność zera.
Pierwiastki Riccatiego wynoszą
\[
 r_\pm=-\frac1{4x}
       +\frac1{2(x-1)}
       \pm\frac{\sqrt{x}}{2(x-1)}.
\]
Podstawienie
$r_\pm'+r_\pm^2=q$
jest bezpośrednim
sprawdzeniem wyniku.
Przykład pokazuje,
że pominięcie $P$
pozwalałoby znaleźć
tylko przypadek $n=0$.
\end{przyklad}

\begin{przyklad}[Inna konfiguracja lokalna]
\label{aa:prz:kov2-infinity-one}
Dla
$q=1/x-3/(16x^2)$
mamy
$b_0=-3/16$,
$\sqrt{1+4b_0}=1/2$
i $E_0=\{1,2,3\}$.
W nieskończoności
$\operatorname{ord}_\infty q=1$,
więc $e_\infty=1$.
Wybór $e_0=1$
daje $n=0$,
$\theta=1/(2x)$
i $P=1$.
Następnie
\[
 \Delta=4q-2\theta'-\theta^2=\frac4x,
 \qquad
 r_\pm=\frac1{4x}\pm\frac1{\sqrt{x}}.
\]
Każdy pierwiastek
spełnia Riccatiego,
co można sprawdzić
przez różniczkowanie.
Zatem istnieją lokalne
rozwiązania
$z=x^{1/4}\exp(\pm2\sqrt{x})$.
\end{przyklad}

\begin{przyklad}[Wynik negatywny dla równania Airy'ego]
\label{aa:prz:kov2-airy}
Dla $q=x$ nie ma
biegunów skończonych.
Ponieważ
$\operatorname{ord}_\infty q=-1$,
mamy $E_\infty=\{-1\}$,
a więc
$n=-1/2\notin\mathbb Z_{\ge0}$.
Lista drugiej gałęzi
jest pusta.
W rozdziale~\ref{aa:ch:picard-vessiot}
udowodniono silniejszy
wynik: równanie Airy'ego
nie ma \emph{żadnych}
niezerowych rozwiązań
liouvillowskich.
Sam obecny test
dotyczy tylko drugiej
gałęzi algorytmu.
\end{przyklad}

\section{Związek z grupą Galois}

Niech $r_\pm$
będą ściśle kwadratowymi,
sprzężonymi rozwiązaniami
Riccatiego.
Odpowiadające im
rozwiązania $z_\pm$
równania $z''=qz$
wyznaczają dwie różne
proste w dwuwymiarowej
przestrzeni rozwiązań.
Automorfizm różniczkowy
nad $\mathbb C(x)$
przekształca zbiór
$\{r_+,r_-\}$
w siebie, zatem
przekształca także
parę tych prostych
w siebie. Grupa Galois
zachowuje więc
\emph{nieuporządkowaną}
parę prostych,
co odpowiada
podgrupie normalizatora
torusa diagonalnego
w $\mathrm{SL}_2(\mathbb C)$.
Jeśli odrzuciliśmy
wcześniej przypadek
wymierny z AA-48,
nie istnieje
niezmiennicza pojedyncza
prosta zdefiniowana
nad $\mathbb C(x)$.
Powyższe rozumowanie
jest ilustracją
imprymitywnego
przypadku teorii Galois;
pełna klasyfikacja
grup i odpowiedniość
Galois zostały
sformułowane jako
wyniki klasyczne
w rozdziale~\ref{aa:ch:picard-vessiot}.

Trzeci przypadek algorytmu
Kovacica, związany
z grupami skończonymi
i stopniami algebraicznymi
$4$, $6$ i $12$,
nie należy do zakresu
AA-49. Jego
osobne opracowanie wymaga
nowego rozdziału
i analogicznych audytów.

% AA-50 — Trzeci przypadek Kovacica i skończone grupy Galois.
% Wyniki klasyczne oznaczone jawnie: lokalne listy, kompletność
% i klasyfikacja prymitywnych grup tetra/octa/icosahedralnych.
% Udowodnione wprost: rekurencja i tożsamość przepływu Riccatiego;
% odzyskanie pierwiastka, końcowy test i przykład czworościenny.
\chapter{Trzeci przypadek algorytmu Kovacica i skończone grupy Galois}
\chaptermark{Trzeci przypadek Kovacica}
\label{aa:ch:kovacic-finite-groups}

W rozdziałach~\ref{aa:ch:kovacic-rational-potential}
i~\ref{aa:ch:kovacic-quadratic} poznaliśmy dwa
pierwsze przypadki algorytmu Kovacica.
Pierwszy wykrywa rozwiązania wymierne
równania Riccatiego $w'+w^2=q$,
drugi --- rozwiązania algebraiczne stopnia $2$
nad $\mathbb C(x)$. Trzeci przypadek,
związany ze skończonymi prymitywnymi
grupami Galois, wymaga osobnej konstrukcji.
Jej zasadniczym narzędziem jest
rekurencja prowadząca do wielomianu
$F(x,T)$, którego pierwiastki są
rozwiązaniami Riccatiego.

Rozdzielamy \emph{dowody bezpośrednie}
rekurencji, jej równoważności z
równaniem transportowym i
odtwarzania pierwiastków od
\emph{cytowanych twierdzeń klasycznych}
o wyczerpujących listach danych lokalnych
oraz skończonych grupach Galois.
Przykłady i testy dokładne weryfikują
konkretne konstrukcje, a nie zastępują
ogólnej klasyfikacji.

\section{Postać normalna i warunki wstępne}

Niech $K=\mathbb C(x)$ z derywacją
$D=\mathrm d/\mathrm dx$.
Rozpatrujemy równanie $z''=qz$
dla $q\in K$. Zakładamy
efektywną reprezentację współczynników
w sensie
Uwagi~\ref{aa:uw:kov-rational-effective};
w szczególności dysponujemy dokładnymi
działaniami i rozkładem wielomianów
na czynniki w odpowiednim
algebraicznym rozszerzeniu.

\begin{definicja}[Rząd potencjału w nieskończoności]
\label{aa:def:kov3-infinity-order}
Niech $q=a/b\in\mathbb C(x)\setminus\{0\}$,
gdzie $a,b\in\mathbb C[x]$ są względnie
pierwsze i $b\ne0$.
Kładziemy
$\operatorname{ord}_{\infty}q=\deg b-\deg a$.
Jeśli $\operatorname{ord}_{\infty}q=2$,
przez $b_\infty$ oznaczamy współczynnik
przy $x^{-2}$ w rozwinięciu
Laurenta $q$ przy nieskończoności.
Przy skończonym biegunie $c$ rzędu $2$
oznaczamy
$b_c=[(x-c)^{-2}]_c q$.
\end{definicja}

\begin{definicja}[Prymitywna reprezentacja wymiaru drugiego]
\label{aa:def:kov3-primitive}
Niech $G$ będzie grupą,
a $\rho:G\to\operatorname{GL}(V)$
reprezentacją zespoloną na
dwuwymiarowej przestrzeni $V$.
Mówimy, że jest ona
\emph{prymitywna}, jeśli
jest nieprzywiedlna oraz
nie istnieją różne
jednowymiarowe podprzestrzenie
$L_1,L_2\subset V$ takie, że
$V=L_1\oplus L_2$ i
zbiór $\{L_1,L_2\}$
jest zachowywany przez
$\rho(G)$.
Jeżeli taki rozkład
istnieje, reprezentację
nieprzywiedlną nazywamy
\emph{imprymitywną}.
\end{definicja}

\begin{twierdzenie}[Przypadek trzeci: klasyfikacja -- wynik klasyczny]
\label{aa:tw:kov3-classification}
Niech $q\in\mathbb C(x)$.
Jeśli różniczkowa grupa Galois
równania $z''=qz$ nad $\mathbb C(x)$
jest skończona i jej reprezentacja
wymiaru $2$ jest prymitywna,
to jej obraz w
$\operatorname{PSL}_2(\mathbb C)$
jest izomorficzny z jedną z grup
$A_4$, $S_4$ lub $A_5$.
Odpowiadające im parametry
trzeciego przypadku algorytmu Kovacica
należą do zbioru
\[
 n\in\{4,6,12\}.
\]
W tym przypadku wszystkie
rozwiązania równania liniowego
są algebraiczne nad $\mathbb C(x)$.
Wyczerpujący charakter tej klasyfikacji
jest klasycznym wynikiem teorii
Picarda--Vessiota i algorytmu Kovacica.
\end{twierdzenie}

\begin{uwaga}[Źródło]
\label{aa:uw:kov3-classical}
Korzystamy z pracy J.~Kovacica,
\emph{An algorithm for solving second
order linear homogeneous differential
equations}, \emph{Journal of Symbolic
Computation} 2 (1986), 3--43.
W szczególności nie twierdzimy, że
poniższe rachunki z jednym przykładem
dowodzą klasyfikacji skończonych
podgrup $\operatorname{PSL}_2(\mathbb C)$
lub pełnej odpowiedniości Galois.
Pierwszy i drugi przypadek algorytmu
zostały omówione wcześniej;
sam wynik negatywny trzeciego przypadku
jest konkluzją dla całego algorytmu
dopiero po wykluczeniu obu pierwszych.
\end{uwaga}

\section{Dane lokalne trzeciego przypadku}

\begin{twierdzenie}[Wyczerpujące listy lokalne -- wynik klasyczny]
\label{aa:tw:kov3-local-sets}
Niech $q\in K\setminus\{0\}$.
W trzecim przypadku Kovacica każdy
skończony biegun $q$ ma rząd $1$ albo $2$
oraz $\operatorname{ord}_{\infty}q\ge2$.
Oznaczmy przez $\Gamma$ zbiór
wszystkich skończonych biegunów $q$.
Dla każdego $n\in\{4,6,12\}$
definiujemy następujące
skończone zbiory liczb całkowitych:
\begin{equation}\label{aa:eq:kov3-local-sets}
 E_c(n)=
 \begin{cases}
 \{12\},&\text{jeśli biegun w $c$ jest prosty},\\[2pt]
 \left\{6+\dfrac{12k}{n}\sqrt{1+4b_c}:
       k=-\dfrac n2,\ldots,\dfrac n2\right\}
       \cap\mathbb Z,
       &\text{jeśli ma rząd $2$};
 \end{cases}
\end{equation}
przy nieskończoności przyjmujemy
\begin{equation}\label{aa:eq:kov3-infinity-sets}
 E_\infty(n)=
 \begin{cases}
 \left\{6+\dfrac{12k}{n}:
       k=-\dfrac n2,\ldots,\dfrac n2\right\}
       \cap\mathbb Z,
       &\operatorname{ord}_{\infty}q>2,\\[2pt]
 \left\{6+\dfrac{12k}{n}\sqrt{1+4b_\infty}:
       k=-\dfrac n2,\ldots,\dfrac n2\right\}
       \cap\mathbb Z,
       &\operatorname{ord}_{\infty}q=2 .
 \end{cases}
\end{equation}
Jeśli występuje skończona
prymitywna grupa Galois, to
dla pewnego $n\in\{4,6,12\}$
i pewnego wyboru
$e_c\in E_c(n)$,
$e_\infty\in E_\infty(n)$,
liczba
\begin{equation}\label{aa:eq:kov3-degree}
 d=\frac n{12}
       \left(e_\infty-\sum_{c\in\Gamma}e_c\right)
\end{equation}
jest nieujemną liczbą całkowitą.
Dla takiego wyboru definiujemy
\begin{equation}\label{aa:eq:kov3-theta-S}
 \theta=\frac n{12}\sum_{c\in\Gamma}
                   \frac{e_c}{x-c},
 \qquad S=\prod_{c\in\Gamma}(x-c).
\end{equation}
W klasycznej procedurze
istnieje unormowany
$P\in\mathbb C[x]$ stopnia $d$
spełniający podaną poniżej
rekurencję z warunkiem
$P_{-1}=0$.
\end{twierdzenie}

\begin{uwaga}[Dokładna granica cytatu]
\label{aa:uw:kov3-sets-cited}
Twierdzenie~\ref{aa:tw:kov3-local-sets}
przyjmujemy jako klasyczną,
wyczerpującą część algorytmu
Kovacica. Samodzielnie
nie dowodzimy ani kompletności
list $E_c(n)$ i $E_\infty(n)$,
ani wynikającej z teorii
niezmienników konieczności
istnienia wielomianu $P$.
Natomiast \emph{poprawność
każdego znalezionego wyniku}
wyniknie z dowodu równania
transportowego. Jeśli $q=0$,
rozwiązania $z=1,x$ są już
znane i nie uruchamiamy
przypadku trzeciego.
\end{uwaga}

\section{Rekurencja wielomianowa}

\begin{definicja}[Wielomiany pomocnicze]
\label{aa:def:kov3-recurrence}
Ustalmy $n\in\{4,6,12\}$,
dane $S,\theta$ z
\eqref{aa:eq:kov3-theta-S}
oraz unormowany wielomian
$P$ stopnia $d$.
Definiujemy
$P_{n+1}=0$, $P_n=-P$
i rekurencyjnie dla
$i=n,n-1,\ldots,0$:
\begin{equation}\label{aa:eq:kov3-recurrence}
 P_{i-1}
 =-SP_i'
 +\bigl((n-i)S'-S\theta\bigr)P_i
 -(n-i)(i+1)S^2qP_{i+1}.
\end{equation}
Każdy $P_i$ jest wielomianem,
ponieważ $S\theta$ i
$S^2q$ są wielomianami
przy założeniach
Twierdzenia~\ref{aa:tw:kov3-local-sets}.
Warunkiem końcowym jest
$P_{-1}=0$ jako tożsamość
w $\mathbb C[x]$.
\end{definicja}

Wbrew pozorom indeks $i=-1$
nie oznacza pierwiastka
równania różniczkowego.
Jest to warunek zgodności
ostatniego kroku rekurencji,
a jego współczynniki
dają skończony układ
liniowy na nieznane
współczynniki $P$.

\begin{lemat}[Rekurencja jako równanie transportowe]
\label{aa:lem:kov3-transport}
Niech $q,S,\theta$ będą
funkcjami wymiernymi,
$S\ne0$, a $n\ge1$.
Niech funkcje $P_i\in K$
spełniają
\eqref{aa:eq:kov3-recurrence}
dla $i=n,\ldots,0$,
z $P_{n+1}=0$, $P_n=-P\ne0$
oraz $P_{-1}=0$.
Wprowadźmy
\begin{equation}\label{aa:eq:kov3-polynomial-F}
 F(x,T)=\sum_{i=0}^{n}
          \frac{S^iP_i}{(n-i)!}T^i\in K[T]
\end{equation}
i derywację
$\mathcal D=\partial_x+(q-T^2)\partial_T$
na $K[T]$.
Wtedy
\begin{equation}\label{aa:eq:kov3-transport}
 \mathcal DF
 =\left(-nT+n\frac{S'}S-\theta\right)F.
\end{equation}
Odwrotnie, przy
$P_{n+1}=0$, $P_n=-P\ne0$
tożsamość~\eqref{aa:eq:kov3-transport}
jest równoważna rekurencji
i warunkowi $P_{-1}=0$.
\end{lemat}
\begin{proof}
Oznaczmy
$A_i=S^iP_i/(n-i)!$ dla
$0\le i\le n$, a poza tym
zakresem $A_{-1}=A_{n+1}=0$.
Współczynnik przy $T^i$
w $\mathcal DF$ dla
$0\le i\le n$ wynosi
\[
 A_i'+(i+1)qA_{i+1}-(i-1)A_{i-1}.
\]
Współczynnik przy tej samej
potędze w prawej stronie
\eqref{aa:eq:kov3-transport}
wynosi
$-nA_{i-1}+
 (nS'/S-\theta)A_i$.
Ich równość można
przekształcić do
\[
 (n-i+1)A_{i-1}
 =-A_i'+\left(n\frac{S'}S-\theta\right)A_i
 -(i+1)qA_{i+1}.
\]
Po podstawieniu $A_i$
i pomnożeniu przez
$(n-i)!/S^{i-1}$
otrzymujemy dokładnie
\eqref{aa:eq:kov3-recurrence}.
Dla $i=0$ powstaje
warunek $P_{-1}=0$.
Przy $T^{n+1}$ po obu
stronach otrzymujemy
$-nA_n$, więc nie ma
dodatkowego warunku.
To dowodzi także
równoważności w drugą stronę.
\end{proof}

\begin{lemat}[Algebraiczne pierwiastki rozwiązują Riccatiego]
\label{aa:lem:kov3-roots}
Niech $F\in K[T]\setminus K$
oraz $H\in K[T]$
spełniają $\mathcal DF=HF$,
gdzie
$\mathcal D=\partial_x+(q-T^2)\partial_T$
i $q\in K$.
Wtedy każdy pierwiastek
$w$ dowolnego nierozkładalnego
czynnika wielomianu $F$
nad $K$ spełnia
$w'+w^2=q$ w jedynym
przedłużeniu derywacji
do $K(w)$.
W szczególności otrzymujemy
algebraiczne rozwiązanie
równania Riccatiego.
\end{lemat}
\begin{proof}
Rozłóżmy w $K[T]$
wielomian $F$ na
niezerowe nierozkładalne
czynniki $f_j$ o
krotnościach $m_j$.
Z tożsamości
$\mathcal DF=HF$ wynika
\[
 \sum_jm_j\frac{\mathcal Df_j}{f_j}=H.
\]
W pierścieniu ilorazowym
$K[T]_{(f_j)}$ pozostałe
czynniki są odwracalne.
Ponieważ prawa strona
jest wielomianem,
wyraz $\mathcal Df_j/f_j$
nie może mieć bieguna,
czyli $f_j$ dzieli
$\mathcal Df_j$.
Niech $w$ będzie
pierwiastkiem $f_j$.
W charakterystyce $0$
nierozkładalny $f_j$
jest separowalny, a więc
$\partial_T f_j(x,w)\ne0$.
Zróżniczkowanie
$f_j(x,w)=0$ daje
\[
 0=\partial_x f_j(x,w)
   +\partial_T f_j(x,w)w'.
\]
Jednocześnie
$\mathcal D f_j(x,w)=0$,
a zatem
$\partial_x f_j(x,w)
 +(q-w^2)\partial_T f_j(x,w)=0$.
Odjęcie tych równości
i podzielenie przez
$\partial_Tf_j(x,w)$
daje $w'=q-w^2$.
\end{proof}

\begin{twierdzenie}[Skończona procedura i jej poprawność]
\label{aa:tw:kov3-algorithm}
Niech $q\in\mathbb C(x)$
ma efektywnie zadane
współczynniki.
Przyjmując klasyczne
Twierdzenie~\ref{aa:tw:kov3-local-sets}
i wcześniejszą klasyfikację
Kovacica, trzecia gałąź
jest rozstrzygalna:
\begin{enumerate}
\item odrzucamy $q\ne0$
o skończonym biegunie
rzędu większego niż $2$
lub o
$\operatorname{ord}_{\infty}q<2$;
\item dla kolejnych
$n=4,6,12$ obliczamy
skończone zbiory
$E_c(n)$ i $E_\infty(n)$;
\item rozpatrujemy każdą
kombinację, dla której
$d$ z~\eqref{aa:eq:kov3-degree}
jest nieujemną liczbą
całkowitą, i tworzymy
$S,\theta$;
\item wstawiamy
$P=x^d+\sum_{j=0}^{d-1}p_jx^j$
do~\eqref{aa:eq:kov3-recurrence}
i rozwiązujemy
skończony liniowy układ
powstały z $P_{-1}=0$;
\item dla znalezionego
$P$ tworzymy $F$
z~\eqref{aa:eq:kov3-polynomial-F},
znajdujemy jego
pierwiastek algebraiczny
$w$ i otrzymujemy
lokalne niezerowe rozwiązanie
$z=\exp(\int w\,\mathrm dx)$.
\end{enumerate}
Odpowiedź dodatnia jest
weryfikowalna bez odwołania
do klasyfikacji grup Galois.
Jeśli pierwsze dwie gałęzie
Kovacica zostały wykluczone,
negatywny wynik wszystkich
trzech gałęzi oznacza brak
niezerowych rozwiązań
liouvillowskich.
\end{twierdzenie}
\begin{proof}
Zbiorów danych lokalnych
jest skończenie wiele.
Dla każdej kombinacji
znany jest stopień $d$,
a rekurencja jest liniowa
w $P$ i jego pochodnych.
Ponieważ $S\theta$ i
$S^2q$ są wielomianami,
każdy $P_i$ jest
wielomianem liniowo
zależnym od nieznanych
$p_j$, więc test
$P_{-1}=0$ sprowadza się
do skończonego układu
równań liniowych.
Dla dowolnego dodatniego
wyniku konstrukcja $F$
spełnia równanie
\eqref{aa:eq:kov3-transport}
na mocy
Lematu~\ref{aa:lem:kov3-transport}.
Ponieważ współczynnik
wiodący $F$ przy $T^n$
wynosi $-S^nP\ne0$,
wielomian nie jest stały
w $T$. Lemma
\ref{aa:lem:kov3-roots}
daje pierwiastek
$w$ spełniający Riccatiego.
Stąd $z' = wz$ prowadzi do
$z''=qz$.
Znalezienie pierwiastka
w skończonym rozszerzeniu
algebraicznym jest
dopuszczalne w zadanym
efektywnym modelu.
Natomiast \emph{konieczność}
istnienia jednej z
kombinacji i wielomianu $P$
dla skończonej prymitywnej
grupy Galois wynika
właśnie z cytowanego
Twierdzenia~\ref{aa:tw:kov3-local-sets}.
Ostatni wniosek o braku
wszystkich rozwiązań
liouvillowskich wykorzystuje
pełną klasyfikację czterech
przypadków Kovacica
z~\ref{aa:tw:kovacic-classification},
również cytowaną
jako wynik klasyczny.
\end{proof}

\section{Przykład czworościenny}

\begin{przyklad}[Potencjał typu czworościennego]
\label{aa:prz:kov3-tetrahedral}
Rozważmy
\begin{equation}\label{aa:eq:kov3-tetra-q}
 q=-\frac3{16x^2}
    -\frac2{9(x-1)^2}
    +\frac3{16x(x-1)}.
\end{equation}
Skończone bieguny
w $0$ i $1$ mają rząd $2$,
a w nieskończoności
$\operatorname{ord}_\infty q=2$.
Współczynniki wiodące są
\[
 b_0=-\frac3{16},\qquad
 b_1=-\frac29,\qquad
 b_\infty=-\frac29.
\]
Przy $n=4$ pierwiastki
$\sqrt{1+4b_c}$ wynoszą
odpowiednio $1/2$, $1/3$, $1/3$.
Zatem
\[
 E_0(4)=\{3,6,9\},\qquad
 E_1(4)=E_\infty(4)=\{4,5,6,7,8\}.
\]
Wybór
$e_0=3$, $e_1=4$,
$e_\infty=7$ daje
\[
 d=\frac13(7-3-4)=0,\qquad
 P=1,\quad
 S=x(x-1),\quad
 \theta=\frac1x+\frac4{3(x-1)}.
\]
Rekurencja~\eqref{aa:eq:kov3-recurrence}
przy $n=4$ wyznacza
\[
\begin{aligned}
 P_4&=-1,\\
 P_3&=\frac{7x-3}{3},\\
 P_2&=-\frac{48x^2-41x+9}{12},\\
 P_1&=\frac{320x^3-409x^2+180x-27}{72},\\
 P_0&=-\frac{2048x^4-3484x^3+2313x^2
                -702x+81}{864},\\
 P_{-1}&=0.
\end{aligned}
\]
Tak więc wielomian
\[
 F(x,T)=\sum_{i=0}^{4}
       \frac{S^iP_i}{(4-i)!}T^i
\]
ma stopień $4$ względem $T$,
a każdy pierwiastek
nierozkładalnego czynnika
$F$ nad $\mathbb C(x)$
spełnia $w'+w^2=q$.
Otrzymaliśmy w pełni
weryfikowalny dodatni wynik
procedury bez zapisywania
rozwiniętego wielomianu
o dużych współczynnikach.
Przykład odpowiada
lokalnym różnicom wykładników
$1/2$, $1/3$, $1/3$,
typowym dla grupy czworościanu.
Przypisanie pełnej grupy
Galois korzysta z klasycznej
teorii monodromii
równania hipergeometrycznego.
\end{przyklad}

\begin{przyklad}[Przypadki ośmiościenny i dwudziestościenny]
\label{aa:prz:kov3-octa-icosa}
Dwa kolejne przykłady
mają wspólną postać
\[
 q_m=-\frac3{16x^2}
     -\frac2{9(x-1)^2}
     +\frac{C_m}{x(x-1)},\qquad
 C_4=\frac{101}{576},\quad
 C_5=\frac{611}{3600}.
\]
Dla $m=4$ odpowiadające
różnice wykładników lokalnych
przy $0,1,\infty$ wynoszą
$1/2,1/3,1/4$, czyli
dane typu ośmiościennego.
Przy $n=6$ zbiorami
kandydatów są
\[
 E_0=\{3,4,5,6,7,8,9\},
 \quad E_1=\{4,6,8\},
 \quad E_\infty=\{5,6,7\}.
\]
Wybór $e_0=3$, $e_1=4$,
$e_\infty=7$ daje $d=0$,
$P=1$ i
\[
 S=x(x-1),\qquad
 \theta=\frac3{2x}+\frac2{x-1},
 \qquad S\theta=\frac72x-\frac32,
\]
a także
\[
 S^2q_4
 =-\frac{15}{64}x^2
  +\frac{115}{576}x-\frac3{16}.
\]
Wykonanie skończonej
rekurencji
\eqref{aa:eq:kov3-recurrence}
dla $i=6,\ldots,0$
prowadzi dokładnie
do $P_{-1}=0$.
Współczynniki są wymierne,
zatem wszystkie etapy można
sprawdzić arytmetycznie,
bez przybliżeń.

Dla $m=5$ otrzymujemy
różnice wykładników
$1/2,1/3,1/5$ i dane
typu dwudziestościennego.
Przy $n=12$ mamy
\[
 E_0=\{3,4,5,6,7,8,9\},
 \quad E_1=\{4,5,6,7,8\},
 \quad E_\infty=\{5,6,7\}.
\]
Ponownie $e_0=3$, $e_1=4$,
$e_\infty=7$ daje $d=0$,
$P=1$, lecz tym razem
\[
 \theta=\frac3x+\frac4{x-1},
 \qquad S\theta=7x-3,
\]
oraz
\[
 S^2q_5
 =-\frac6{25}x^2
  +\frac{739}{3600}x-\frac3{16}.
\]
Rekurencja dla
$i=12,\ldots,0$ również
daje $P_{-1}=0$.
Równości dla obu wartości
$n$ są dodatkowo sprawdzane
przez niezależne testy
symboliczne i wymierne.
Przyporządkowanie
pełnej grupy monodromii
typom ośmio- i
dwudziestościennemu
wykorzystuje klasyczną
klasyfikację Schwarzowską,
a nie sam rachunek
$P_{-1}=0$.
\end{przyklad}

\begin{przyklad}[Rola dokładnej liczby całkowitej]
\label{aa:prz:kov3-integer-sets}
Dla $q$ z~\eqref{aa:eq:kov3-tetra-q}
i $n=4$ współczynnik
$b_0=-3/16$ daje
$\sqrt{1+4b_0}=1/2$.
W zbiorze
$\{6+3k/2:k=-2,-1,0,1,2\}$
nie wszystkie liczby są
całkowite:
$\{3,9/2,6,15/2,9\}$.
Lista $E_0(4)$ wynosi
zatem dokładnie $\{3,6,9\}$.
Niedopuszczalne jest
zaokrąglanie $9/2$ czy
$15/2$ do liczb całkowitych.
Podobnie w przypadku
$b_1=-2/9$ pierwiastek
wynosi $1/3$, a
$\{6+k:k=-2,\ldots,2\}$
daje pięć różnych
wartości całkowitych.
\end{przyklad}

\begin{przyklad}[Wcześniejsza gałąź nie znika]
\label{aa:prz:kov3-rational-degenerate}
Dla $q=2/x^2$ już
$w=2/x$ i $w=-1/x$
są rozwiązaniami
\emph{wymiernymi}
Riccatiego, ponieważ
$w'+w^2=q$.
Przypadek trzeci
nie jest tutaj potrzebny.
Ewentualne rozwiązanie
rekurencji~\eqref{aa:eq:kov3-recurrence}
nie upoważniałoby do
przypisania grupie Galois
typu czworościennego:
najpierw należy odróżnić
przypadki wymierny
i kwadratowy od skończonego
prymitywnego.
\end{przyklad}

\begin{przyklad}[Natychmiastowa przeszkoda]
\label{aa:prz:kov3-airy}
Dla równania Airy'ego
$z''=xz$ potencjał
$q=x$ nie ma skończonych
biegunów, ale
$\operatorname{ord}_{\infty}q=-1<2$.
Trzeci przypadek jest więc
wykluczony już przez
warunek wstępny z
Twierdzenia~\ref{aa:tw:kov3-local-sets}.
Łącznie z wynikami
rozdziałów~\ref{aa:ch:kovacic-rational-potential}
i~\ref{aa:ch:kovacic-quadratic}
oraz klasyczną klasyfikacją
Kovacica wyklucza to
liouvillowskość,
zgodnie z analizą
grupy Galois Airy'ego
w~\ref{aa:ch:picard-vessiot}.
\end{przyklad}

\section{Zamknięcie trzech przypadków}

Mamy teraz trzy osobne
procedury wraz z ich
granicami:
pierwszy przypadek bada
$w\in\mathbb C(x)$;
drugi --- $w$ algebraiczne
stopnia $2$;
trzeci --- przypadek
skończonej prymitywnej
grupy Galois, którego
właściwe stopnie algebraiczne
to $4$, $6$ i $12$.
Twierdzenie o wyczerpaniu
tych możliwości i
wniosek o braku rozwiązań
liouvillowskich w czwartym
przypadku są wynikiem
klasycznym, a nie
następstwem skończonej
liczby testów regresyjnych.

W odróżnieniu od
deklaracji teoretycznej
o pełnym algorytmie Rischa,
wszystkie trzy gałęzie
Kovacica są tu
\emph{opisane algorytmicznie},
ale ich matematyczna
kompletność poza
bezpośrednio udowodnionymi
redukcjami wykorzystuje
cytowaną klasyczną
klasyfikację lokalną
i teorię Picarda--Vessiota.

% AA-51 — klasyfikacja typów izomorfizmu, nie klasyfikacja sprzężoności.
\chapter{Skończone podgrupy grupy przekształceń Möbiusa}
\chaptermark{Skończone podgrupy PGL(2)}
\label{aa:ch:finite-pgl2}

W rozdziale~\ref{aa:ch:kovacic-finite-groups} wykorzystaliśmy
klasyfikację skończonych prymitywnych grup rzutowych. Teraz
udowodnimy klasyfikację typów izomorfizmu wszystkich skończonych
podgrup $\operatorname{PGL}_2(\mathbb C)$ oraz wyjaśnimy
pochodzenie liczb $4,6,12$. Nie dowodzimy tutaj jednoznaczności
zanurzeń z dokładnością do sprzężenia ani pełnego twierdzenia
o listach lokalnych algorytmu Kovacica.

\section{Działanie na prostej rzutowej}

\begin{definicja}[Grupa rzutowa]
\label{aa:def:pgl2-projective}
Niech $\mathbb P^1(\mathbb C)$ będzie zbiorem prostych
wektorowych w $\mathbb C^2$. Definiujemy
$\operatorname{PGL}_2(\mathbb C)=
\operatorname{GL}_2(\mathbb C)/\{\lambda I:\lambda\in\mathbb C^\times\}$.
Klasa macierzy $A$ działa przez $[v]\mapsto[Av]$.
Współrzędna $z=v_1/v_2$ identyfikuje prostą rzutową
z $\mathbb C\cup\{\infty\}$, a działanie macierzy
$\left(\begin{smallmatrix}a&b\\c&d\end{smallmatrix}\right)$
z przekształceniem $z\mapsto(az+b)/(cz+d)$,
z naturalnym przedłużeniem w biegunie i w nieskończoności.
\end{definicja}

Działanie jest wierne: macierz zachowująca wszystkie proste
zachowuje proste rozpięte przez $e_1,e_2,e_1+e_2$,
więc jest skalarna. Ponadto każda macierz odwracalna
może być pomnożona przez skalar tak, aby jej wyznacznik
wynosił $1$. Stąd $\operatorname{PGL}_2(\mathbb C)
\simeq\operatorname{SL}_2(\mathbb C)/\{\pm I\}
=\operatorname{PSL}_2(\mathbb C)$.

\begin{lemat}[Punkty stałe i stabilizatory]
\label{aa:lem:pgl2-fixed}
Niech $G$ będzie skończoną podgrupą
$\operatorname{PGL}_2(\mathbb C)$. Każdy $g\in G\setminus\{1\}$
ma dokładnie dwa punkty stałe w $\mathbb P^1(\mathbb C)$,
a stabilizator $G_p$ każdego punktu $p$ jest cykliczny.
\end{lemat}
\begin{proof}
Jeżeli $g$ ma rząd $m>1$ i reprezentant $A$,
to $A^m=\lambda I$ dla pewnego $\lambda\ne0$.
Wielomian $T^m-\lambda$ ma różne pierwiastki,
więc $A$ jest diagonalizowalna. Gdyby miała tylko
jedną wartość własną, byłaby skalarna. Zatem ma dokładnie
dwie proste własne, będące punktami stałymi $g$.

Po zmianie współrzędnych możemy przyjąć $p=\infty$.
Wówczas każdy $h\in G_p$ działa przez $z\mapsto a_hz+b_h$.
Liczba $c=|G_p|^{-1}\sum_{h\in G_p}h(0)$ jest wspólnym
punktem stałym: dla $k\in G_p$ afiniczność daje
$k(c)=|G_p|^{-1}\sum_h(kh)(0)=c$.
Po przesunięciu $c$ do zera grupa działa przez mnożenia
$z\mapsto a_hz$ i zanurza się w $\mathbb C^\times$.
Jeśli jej rząd wynosi $s$, każdy $a_h$ jest pierwiastkiem
równania $T^s=1$. Ponieważ istnieje dokładnie $s$ takich
pierwiastków i tworzą grupę cykliczną, stabilizator
jest cykliczny.
\end{proof}

W szczególności element ustalający trzy różne punkty
jest identycznością. To proste spostrzeżenie będzie
również użyte przy konstrukcji wiernych reprezentacji
permutacyjnych.

\section{Równanie orbit wyjątkowych}

\begin{twierdzenie}[Możliwe dane stabilizatorów]
\label{aa:tw:pgl2-signatures}
Niech $G\le\operatorname{PGL}_2(\mathbb C)$ będzie
skończona, $N=|G|>1$. Zbiór punktów o nietrywialnym
stabilizatorze rozkłada się na $r$ orbit $O_1,\ldots,O_r$.
Wybierzmy $p_i\in O_i$ i oznaczmy $m_i=|G_{p_i}|$.
Po zmianie numeracji przyjmujemy $2\le m_1\le\cdots\le m_r$.
Zachodzi
\begin{equation}\label{aa:eq:pgl2-orbit-count}
 \sum_{i=1}^r\left(1-\frac1{m_i}\right)=2-\frac2N .
\end{equation}
Możliwe są dokładnie następujące dane liczbowe:
\[
\begin{array}{c|c}
(m_1,\ldots,m_r)&N\\ \hline
(N,N)&N\\
(2,2,n),\ n\ge2&2n\\
(2,3,3)&12\\
(2,3,4)&24\\
(2,3,5)&60
\end{array}
\]
\end{twierdzenie}
\begin{proof}
Zbiór punktów o nietrywialnym stabilizatorze jest
sumą zbiorów punktów stałych elementów $g\ne1$.
Na mocy Lematu~\ref{aa:lem:pgl2-fixed} jest skończony;
jest niezmienniczy względem $G$, ponieważ
$G_{hp}=hG_ph^{-1}$ dla $h\in G$.
Zliczamy pary $(g,p)$, gdzie $g\ne1$ i $g(p)=p$.
Z Lematu~\ref{aa:lem:pgl2-fixed} ich liczba to $2(N-1)$.
Orbita o stabilizatorze rzędu $m_i$ zawiera $N/m_i$
punktów i każdy z nich jest ustalany przez $m_i-1$
nietrywialnych elementów. Podzielenie równości liczby
par przez $N$ daje~\eqref{aa:eq:pgl2-orbit-count}.
Każdy składnik lewej strony jest co najmniej $1/2$,
więc $r\le3$. Prawa strona jest co najmniej $1$,
a pojedynczy składnik jest mniejszy od $1$, zatem $r\ge2$.
Dla $r=2$ otrzymujemy $1/m_1+1/m_2=2/N$.
Ponieważ $m_i\le N$, obie liczby muszą równać się $N$.

Dla $r=3$ mamy $1/m_1+1/m_2+1/m_3=1+2/N>1$.
Stąd $m_1=2$; inaczej suma byłaby nie większa od $1$.
Jeśli $m_2=2$, równanie daje $N=2m_3$.
Jeżeli $m_2\ge4$, suma znowu jest nie większa od $1$.
Pozostaje $m_2=3$, a wtedy $1/m_3>1/6$,
czyli $m_3=3,4,5$. Podstawienie daje odpowiednio
$N=12,24,60$. Istnienie grup realizujących te dane
wykażemy poniżej.
\end{proof}

\section{Identyfikacja grup}

\begin{lemat}[Podgrupy indeksu dwa grup symetrycznych]
\label{aa:lem:pgl2-index2}
Dla $k\ge2$ jedyną podgrupą indeksu $2$ grupy $S_k$
jest grupa alternująca $A_k$.
\end{lemat}
\begin{proof}
Podgrupa indeksu $2$ jest normalna i jest jądrem
surjektywnego homomorfizmu do grupy $C_2$.
Wszystkie transpozycje są sprzężone, więc mają ten sam
obraz. Generują $S_k$, a homomorfizm jest surjektywny;
obraz transpozycji jest więc nietrywialny. Homomorfizm
pokrywa się z homomorfizmem znaku, którego jądro to $A_k$.
\end{proof}

\begin{twierdzenie}[Klasyfikacja typów izomorfizmu]
\label{aa:tw:pgl2-classification}
Każda skończona podgrupa $G$ grupy
$\operatorname{PGL}_2(\mathbb C)$ jest izomorficzna
z jedną z grup: cykliczną $C_n$ ($n\ge1$),
dwuścienną $D_n$ rzędu $2n$ ($n\ge2$),
$A_4$, $S_4$ lub $A_5$.
W przypadku $n=2$ grupa $D_2$ jest grupą Kleina.
\end{twierdzenie}
\begin{proof}
Grupa trywialna jest $C_1$. Dla $N>1$ korzystamy
z Twierdzenia~\ref{aa:tw:pgl2-signatures}.

\emph{Dane $(N,N)$.} Grupa ustala punkt, więc sama
jest jego cyklicznym stabilizatorem.

\emph{Dane $(2,2,n)$.} Orbita odpowiadająca stabilizatorowi
rzędu $n$ ma dwa punkty. Po zmianie współrzędnych
są to $0,\infty$. Jądro działania na tej parze jest
stabilizatorem jednego punktu, cyklicznym rzędu $n$,
generowanym przez $a:z\mapsto\zeta z$, gdzie $\zeta$
jest pierwotnym pierwiastkiem $n$-tego stopnia z jedności.
Element $b$ zamieniający punkty ma postać $z\mapsto c/z$,
$c\ne0$. Stąd $b^2=1$ i $bab^{-1}=a^{-1}$.
Wszystkie elementy są postaci $a^j$ albo $a^jb$,
$0\le j<n$, co daje grupę dwuścienną rzędu $2n$.

\emph{Dane $(2,3,3)$.} Jedna z orbit o stabilizatorze
rzędu $3$ ma cztery punkty. Działanie na niej jest wierne,
bo element ustalający cztery punkty jest identycznością.
Obraz w $S_4$ ma rząd $12$, a więc jest $A_4$
na mocy Lematu~\ref{aa:lem:pgl2-index2}.

\emph{Dane $(2,3,4)$.} Istnieje osiem punktów
o stabilizatorze rzędu $3$. Każda podgrupa rzędu $3$
ustala dokładnie dwa punkty (są to punkty stałe jej generatora),
a ich stabilizatory są właśnie tą podgrupą: w pozostałych
orbitach rzędy stabilizatorów nie są podzielne przez $3$.
Zbiory tych dwóch punktów dla różnych podgrup są rozłączne,
bo każdy punkt z tej orbity wyznacza swój stabilizator.
Otrzymujemy cztery podgrupy Sylowa rzędu $3$.
Działanie przez sprzężenie na tych czterech podgrupach
jest przechodnie. Jego jądro $K$ zawiera się w normalizatorze
jednej z nich, rzędu $24/4=6$.
Gdyby $3$ dzieliło $|K|$, twierdzenie Cauchy'ego
i normalność $K$ wymusiłyby zawieranie wszystkich czterech
podgrup rzędu $3$, a więc co najmniej dziewięciu elementów.
To niemożliwe. Zatem $|K|=1$ lub $2$.
W drugim przypadku jedyny nietrywialny element $h\in K$
byłby centralny. Dla dowolnego elementu $a$ rzędu $3$,
element $h$ zachowuje dwa punkty stałe $a$.
Nie może ich zamieniać: po przeniesieniu do $0,\infty$
zamiana sprzęga mnożenie $a$ z jego odwrotnością,
a $a\ne a^{-1}$. Stąd $h$ ustala oba punkty stałe
każdej podgrupy rzędu $3$, czyli osiem różnych punktów,
wbrew Lematowi~\ref{aa:lem:pgl2-fixed}.
Zatem $K=1$ i obraz w $S_4$ ma rząd $24$, czyli jest $S_4$.

\emph{Dane $(2,3,5)$.} Trzy wyjątkowe orbity mają
odpowiednio $30,20,12$ punktów.
Każdy nietrywialny element należy do stabilizatora
jednego ze swoich punktów stałych, więc jego rząd
to $2,3$ albo $5$.
Każda inwolucja ustala dwa punkty z pierwszej orbity,
a każdy z jej punktów wyznacza jedną inwolucję.
Jest zatem $15$ inwolucji; są sprzężone, ponieważ
punkty tworzą jedną orbitę. Centralizator inwolucji
ma rząd $60/15=4$.
Podgrupa Sylowa rzędu $4$ nie ma elementu rzędu $4$,
więc każdy nietrywialny element ma rząd $2$.
Dla dowolnych jej elementów $a,b$ mamy
$ab=(ab)^{-1}=b^{-1}a^{-1}=ba$.
Jeśli $a,b$ są różnymi nietrywialnymi elementami,
to $1,a,b,ab$ są parami różne i podgrupa jest
izomorficzna z $C_2\times C_2$.
Każda inwolucja należy do takiej podgrupy, a każda
taka podgrupa zawierająca $h$ zawiera się w $C_G(h)$,
czyli jest równa $C_G(h)$. Inwolucje są więc rozdzielone
na rozłączne trójki i istnieje pięć podgrup Sylowa rzędu $4$.
Działanie przez sprzężenie na nich jest przechodnie;
jego jądro $K$ ma rząd dzielący $60/5=12$.
Gdyby $2$ dzieliło $|K|$, normalność i twierdzenie
Cauchy'ego wymusiłyby obecność wszystkich $15$ inwolucji,
co jest niemożliwe. Jeśli $3$ dzieli $|K|$, to $K$
zawiera wszystkie podgrupy Sylowa rzędu $3$.
Jest ich dziesięć, gdyż dwadzieścia punktów drugiej
orbity odpowiada parami ich punktom stałym.
Daje to co najmniej $21$ elementów $K$, również niemożliwe.
Zatem $K=1$. Wierny obraz grupy rzędu $60$ w $S_5$
ma indeks $2$, więc jest $A_5$.
\end{proof}

\section{Realizacja wszystkich typów}

\begin{przyklad}[Grupy cykliczne i dwuścienne]
\label{aa:prz:pgl2-cyclic-dihedral}
Mnożenia $z\mapsto\zeta^jz$ dla $\zeta=\exp(2\pi i/n)$
realizują $C_n$. Dołączenie $z\mapsto1/z$ dla $n\ge2$
daje dokładnie $2n$ różnych przekształceń, o relacjach
z dowodu Twierdzenia~\ref{aa:tw:pgl2-classification}.
Dla $n=2$ są to $z,-z,1/z,-1/z$.
\end{przyklad}

\begin{lemat}[Obroty jako przekształcenia rzutowe]
\label{aa:lem:pgl2-rotations}
Grupa $\operatorname{SO}(3)$ jest izomorficzna z podgrupą
$\operatorname{PGL}_2(\mathbb C)$ będącą obrazem
$\operatorname{SU}(2)$.
\end{lemat}
\begin{proof}
Niech $V$ będzie rzeczywistą przestrzenią hermitowskich
macierzy $2\times2$ o śladzie zero.
Macierz
$\left(\begin{smallmatrix}z&x-iy\\x+iy&-z\end{smallmatrix}\right)$
identyfikuje $V$ z $\mathbb R^3$;
jej ujemny wyznacznik to $x^2+y^2+z^2$.
Sprzężenie $H\mapsto UHU^{-1}$ przez $U\in\operatorname{SU}(2)$
zachowuje tę normę i daje homomorfizm do $\operatorname O(3)$.
Macierze $\operatorname{SU}(2)$ mają postać
$\left(\begin{smallmatrix}a&b\\-\overline b&\overline a\end{smallmatrix}\right)$,
$|a|^2+|b|^2=1$. Tworzą więc przestrzeń homeomorficzną
ze sferą $S^3$, która jest spójna.
Wyznacznik indukowanego odwzorowania jest ciągły,
przyjmuje tylko wartości $\pm1$ i w identyczności wynosi $1$,
więc obraz leży w $\operatorname{SO}(3)$.
Jądro składa się z $\pm I$: macierz komutująca z
$\operatorname{diag}(1,-1)$ jest diagonalna, a komutowanie
z $\left(\begin{smallmatrix}0&1\\1&0\end{smallmatrix}\right)$
wymusza równość wpisów diagonalnych; wyznacznik to $1$.

Sprzężenie przez $\operatorname{diag}(e^{-it/2},e^{it/2})$
pozostawia $z$ bez zmiany, a współrzędną $x+iy$
mnoży przez $e^{it}$, daje więc obrót o kąt $t$
wokół osi $z$.
Każda jednostkowa macierz $H\in V$ ma wartości własne
$1,-1$; ortonormalną bazę własną można wybrać tak,
aby macierz zmiany bazy miała wyznacznik $1$ (wystarczy
pomnożyć jeden wektor przez skalar modułu $1$).
Stąd każda oś jednostkowa jest obrazem osi $z$
przez element obrazu $\operatorname{SU}(2)$.
Każdy obrót w $\operatorname{SO}(3)$ jest obrotem wokół osi:
rzeczywista macierz ortogonalna wymiaru trzy o wyznaczniku $1$
ma wartość własną $1$, a na prostopadłej płaszczyźnie
działa przez dwuwymiarowy obrót.
Sprzęgając obrót wokół osi $z$ otrzymujemy więc każdy
element $\operatorname{SO}(3)$. Ostatecznie
$\operatorname{SO}(3)\simeq\operatorname{SU}(2)/\{\pm I\}$,
a ten iloraz zanurza się w $\operatorname{PGL}_2(\mathbb C)$.
\end{proof}

\begin{przyklad}[Czworościan i sześcian]
\label{aa:prz:pgl2-tetra-cube}
Na zbiorze czterech wierzchołków
$(1,1,1),(1,-1,-1),(-1,1,-1),(-1,-1,1)$
działają obroty otrzymane z cyklicznych permutacji
trzech współrzędnych i zmian znaków o iloczynie $1$.
Te $3\cdot4=12$ macierzy tworzy grupę.
Działanie jest wierne, ponieważ wierzchołki rozpinają
$\mathbb R^3$, więc obraz ma indeks dwa w $S_4$
i jest $A_4$.

Macierze permutujące osie współrzędnych i zmieniające znaki,
o wyznaczniku $1$, tworzą grupę $6\cdot4=24$ obrotów
sześcianu. Działają na czterech prostych zawierających
przekątne sześcianu. Element jądra zachowuje proste
o wektorach kierunkowych z poprzedniego akapitu.
Dowolne trzy z nich są liniowo niezależne, a ich suma
jest przeciwna czwartemu. Wektory są zatem wektorami
własnymi o tej samej wartości własnej $\pm1$.
Wyznacznik $1$ wymusza identyczność.
Działanie jest wierne i grupa jest $S_4$.
\end{przyklad}

\begin{przyklad}[Dwudziestościan]
\label{aa:prz:pgl2-icosahedron}
Niech $\phi=(1+\sqrt5)/2$ i niech $X$ składa się
z dwunastu wektorów
$(0,\pm1,\pm\phi),(\pm1,\pm\phi,0),(\pm\phi,0,\pm1)$.
Grupa $G_X=\{R\in\operatorname{SO}(3):R(X)=X\}$
jest skończona, bo $X$ rozpina $\mathbb R^3$ i działanie
na $X$ jest wierne. Jest przechodnia na $X$:
cykliczne permutacje współrzędnych oraz zmiany znaków
o iloczynie $1$ pozwalają przejść między dowolnymi
wierzchołkami (znak zerowej współrzędnej można wybrać
tak, aby zachować wyznacznik).

Dla $v=(0,1,\phi)$ pięć wektorów o iloczynie skalarnym
z $v$ równym $\phi$ to, w kolejności cyklicznej,
\[
 (0,-1,\phi),\quad(\phi,0,1),\quad(1,\phi,0),
 \quad(-1,\phi,0),\quad(-\phi,0,1).
\]
Pozostałe wierzchołki to $v,-v$ i przeciwne do tych
pięciu wektorów. Oznaczmy te pięć wektorów przez $w_0,\ldots,w_4$,
z indeksami liczonymi modulo $5$. Wszystkie mają
kwadrat normy $1+\phi^2$, a z $\phi^2=\phi+1$ wynika
\[
 \langle w_i,w_j\rangle=
 \begin{cases}
  1+\phi^2,& i=j,\\
  \phi,& i-j\equiv\pm1\pmod 5,\\
  -\phi,& i-j\equiv\pm2\pmod 5.
 \end{cases}
\]
Możemy więc skonstruować obrót bez odwoływania się
do wizualnej regularności pięciokąta.
Niech $B=(v,w_0,w_1)$ oraz $C=(v,w_1,w_2)$ będą
macierzami o wskazanych kolumnach. Mamy
$\det B=2\phi^2\ne0$ oraz $B^{\mathsf T}B=C^{\mathsf T}C$.
W szczególności $C$ jest odwracalna. Odwzorowanie
$R=CB^{-1}$ jest ortogonalne, bo
\[
 R^{\mathsf T}R
 =B^{-\mathsf T}C^{\mathsf T}CB^{-1}=I .
\]
Z definicji $Rv=v$, $Rw_0=w_1$, $Rw_1=w_2$.
Dla każdego $i$ wektory $Rw_i$ i $w_{i+1}$ mają
te same iloczyny skalarne z $v,w_1,w_2$:
dla pierwszego wektora wynika to z ortogonalności $R$,
a dla drugiego z podanej tabeli iloczynów skalarnych.
Ponieważ $v,w_1,w_2$ tworzą bazę, wynika stąd
$Rw_i=w_{i+1}$. Zatem $R^5=I$, gdyż $R^5$ ustala
wektory bazy $B$, oraz $R\ne I$.
Rząd $R$ wynosi $5$, więc $(\det R)^5=1$ i $\det R=1$.
W szczególności $R\in G_X$; zachowuje także przeciwne
wektory oraz $\pm v$. Stabilizator $v$ ma zatem
co najmniej pięć elementów. Każdy $Q\in (G_X)_v$
zachowuje iloczyn skalarny z $v$, zatem $Qw_0$
należy do zbioru $\{w_0,\ldots,w_4\}$.
Jeśli $Q_1w_0=Q_2w_0$, to obrót $Q_2^{-1}Q_1$
ustala dwa liniowo niezależne wektory $v,w_0$.
Ustala więc punktowo rozpiętą przez nie płaszczyznę,
a ponieważ jest ortogonalny i ma wyznacznik $1$,
ustala również jej prostopadłą prostą.
Stąd $Q_1=Q_2$. Stabilizator ma więc najwyżej
pięć elementów.
Stąd $|G_X|=12\cdot5=60$.

Przechodniość daje sześć różnych osi przez pary
$\{w,-w\}$, $w\in X$, każda z podgrupą obrotów rzędu $5$.
Są to różne podgrupy, bo nietrywialny obrót ma jedną oś.
Grupa cykliczna rzędu $60$ i grupa dwuścienna rzędu $60$
mają tylko jedną podgrupę rzędu $5$ (w dwuściennej
wszystkie elementy nieparzystego rzędu leżą w grupie
obrotów). Z Twierdzenia~\ref{aa:tw:pgl2-classification}
i Lematu~\ref{aa:lem:pgl2-rotations} wynika $G_X\simeq A_5$.
\end{przyklad}

\begin{uwaga}[Realizowalność danych stabilizatorów]
\label{aa:uw:pgl2-signatures-realized}
Powyższe konstrukcje realizują również każdą sygnaturę
z Twierdzenia~\ref{aa:tw:pgl2-signatures}.
Nietrywialna grupa cykliczna realizuje $(N,N)$,
a grupa dwuścienna rzędu $2n$, $n\ge2$, realizuje
$(2,2,n)$: zachowuje parę $\{0,\infty\}$ i nie ma
wspólnego punktu stałego. Dla $n=2$ wszystkie trzy
wyjątkowe orbity mają dwa punkty.
W grupach $A_4,S_4,A_5$ możliwe rzędy elementów
to odpowiednio $1,2,3$; $1,2,3,4$ oraz $1,2,3,5$
(wynika to z rozkładów permutacji na rozłączne cykle).
W szczególności grupy te nie są cykliczne ani dwuścienne:
grupy cykliczne tych samych rzędów zawierałyby element
rzędu $12,24,60$, a grupy dwuścienne element rzędu
$6,12,30$. Zatem jedynymi możliwymi danymi stabilizatorów
są odpowiednio $(2,3,3)$, $(2,3,4)$ i $(2,3,5)$.
\end{uwaga}

\section{Powrót do różniczkowej teorii Galois}

\begin{wniosek}[Prymitywny przypadek i długości orbit]
\label{aa:wn:pgl2-primitive-orbits}
Niech $H\le\operatorname{SL}_2(\mathbb C)$ będzie skończona,
a jej naturalna reprezentacja na $\mathbb C^2$ prymitywna
w sensie Definicji~\ref{aa:def:kov3-primitive}.
Obraz rzutowy $G$ grupy $H$ jest izomorficzny z $A_4,S_4$
lub $A_5$, a najmniejsze długości orbit na
$\mathbb P^1(\mathbb C)$ wynoszą odpowiednio $4,6,12$.
\end{wniosek}
\begin{proof}
Gdy obraz rzutowy jest trywialny lub cykliczny,
grupa $H$ zachowuje prostą (dla nietrywialnej grupy
cyklicznej wybieramy punkt stały generatora).
Gdy obraz jest dwuścienny, zachowuje dwuelementową
orbitę, czyli nieuporządkowaną parę różnych prostych,
których suma jest całą przestrzenią. Oba przypadki
przeczą prymitywności.
W pozostałych trzech przypadkach minimalne długości
orbit wynoszą $N/\max_i m_i$, bo orbita zwykła ma
$N$ punktów. Z danych $(2,3,3)$, $(2,3,4)$, $(2,3,5)$
otrzymujemy $12/3=4$, $24/4=6$, $60/5=12$.
\end{proof}

Liczby te wyjaśniają geometryczne pochodzenie parametrów
z rozdziału~\ref{aa:ch:kovacic-finite-groups}.
Sama klasyfikacja grup i obliczenie długości orbit
nie dowodzą jednak lokalnych warunków dla potencjału
ani rekurencji koniecznej w algorytmie Kovacica.

% AA-52 — skończone podgrupy SL2(C).
\chapter{Binarne grupy wielościenne i skończone podgrupy SL(2)}
\chaptermark{Binarne grupy wielościenne}
\label{aa:ch:binary-polyhedral}

W rozdziale~\ref{aa:ch:finite-pgl2} poznaliśmy typy
skończonych podgrup grupy rzutowej.
Teraz odtworzymy grupy macierzowe leżące nad nimi.
Istotna różnica polega na tym, że nie każdy obrót rzędu
dwa ma podniesienie rzędu dwa: w grupie specjalnej
liniowej jego podniesienia mają rząd cztery.

\section{Sprowadzenie do grupy unitarnej}

\begin{twierdzenie}[Unitarność po zmianie bazy]
\label{aa:tw:binary-unitary}
Dla każdej skończonej podgrupy
$H\le\operatorname{SL}_2(\mathbb C)$ istnieje
$S\in\operatorname{SL}_2(\mathbb C)$ takie, że
$SHS^{-1}\subset\operatorname{SU}(2)$.
\end{twierdzenie}
\begin{proof}
Niech $A=\sum_{h\in H}h^*h$, gdzie $h^*$ oznacza
macierz sprzężoną hermitowsko. Macierz $A$ jest dodatnio
określona, gdyż dla $v\ne0$ zachodzi
$v^*Av=\sum_{h\in H}\|hv\|^2>0$.
Dla $g\in H$ mamy
$g^*Ag=\sum_{h\in H}(hg)^*(hg)=A$.
Istnieje dodatnio określony hermitowski pierwiastek
$P=A^{1/2}$, otrzymany przez diagonalizację unitarną
$A$ i zastąpienie jej dodatnich wartości własnych
ich dodatnimi pierwiastkami.
Wówczas $(PgP^{-1})^*(PgP^{-1})=I$.
Wyznacznik $PgP^{-1}$ wynosi $1$.
Aby uzyskać macierz zmiany bazy o wyznaczniku $1$,
wystarczy przyjąć $S=(\det P)^{-1/2}P$,
co nie zmienia sprzężenia.
\end{proof}

Z Lematu~\ref{aa:lem:pgl2-rotations} mamy surjekcję
$\pi:\operatorname{SU}(2)\to\operatorname{SO}(3)$,
daną sprzężeniem na rzeczywistej przestrzeni
hermitowskich macierzy o śladzie zero.
Jej jądro to $\{I,-I\}$.
Przez $\rho:\operatorname{SL}_2(\mathbb C)
\to\operatorname{PGL}_2(\mathbb C)$ oznaczamy
homomorfizm rzutowy. Jego ograniczenie do
$\operatorname{SU}(2)$ utożsamia się z $\pi$
za pomocą izomorfizmu z poprzedniego rozdziału.

\begin{lemat}[Elementy rzędu dwa i pełne przeciwobrazy]
\label{aa:lem:binary-full-lift}
Niech $H\le\operatorname{SL}_2(\mathbb C)$ będzie skończona.
Jedynym elementem rzędu $2$ w $\operatorname{SL}_2(\mathbb C)$
jest $-I$. Jeżeli $|H|$ jest parzyste, to
$H=\rho^{-1}(\rho(H))$. Jeśli $|H|$ jest nieparzyste,
to $H$ jest cykliczna.
\end{lemat}
\begin{proof}
Jeśli $M^2=I$, to $M$ jest diagonalizowalna, ponieważ
$T^2-1$ ma różne pierwiastki. Jej wartości własne
należą do $\{1,-1\}$. Warunek $\det M=1$ wymusza
ich równość, więc $M=\pm I$.
Dla parzystego $|H|$ twierdzenie Cauchy'ego daje $-I\in H$.
Jeśli $\rho(M)\in\rho(H)$, istnieje $h\in H$ takie, że
$Mh^{-1}\in\ker\rho=\{\pm I\}\subset H$,
więc $M\in H$.
Dla nieparzystego $|H|$ ograniczenie $\rho|_H$ jest
iniektywne. Z klasyfikacji w
Twierdzeniu~\ref{aa:tw:pgl2-classification} wynika,
że jedyne grupy rzutowe o nieparzystym rzędzie są cykliczne.
\end{proof}

\section{Grupy cykliczne i dwucykliczne}

\begin{definicja}[Grupa dwucykliczna]
\label{aa:def:binary-dicyclic}
Dla całkowitego $n\ge2$ połóżmy
$\zeta=\exp(\pi i/n)$ oraz
\[
 A=\begin{pmatrix}\zeta&0\\0&\zeta^{-1}\end{pmatrix},
 \qquad J=\begin{pmatrix}0&1\\-1&0\end{pmatrix}.
\]
Grupę $\operatorname{Dic}_n=\langle A,J\rangle
\le\operatorname{SU}(2)$ nazywamy grupą dwucykliczną
(binarna grupa dwuścienna).
\end{definicja}

\begin{twierdzenie}[Postać i relacje grupy dwucyklicznej]
\label{aa:tw:binary-dicyclic}
Dla każdego $n\ge2$ grupa $\operatorname{Dic}_n$
ma rząd $4n$, obraz rzutowy izomorficzny z $D_n$
rzędu $2n$ oraz prezentację
\[
 \langle a,j\mid
 a^{2n}=1,\ j^2=a^n,\ jaj^{-1}=a^{-1}\rangle.
\]
\end{twierdzenie}
\begin{proof}
Bezpośrednie mnożenie macierzy daje
$A^{2n}=I$, $J^2=-I=A^n$ oraz $JAJ^{-1}=A^{-1}$.
Relacje pozwalają zapisać każdy wyraz w generatorach
i ich odwrotnościach w postaci $a^k j^\varepsilon$,
gdzie $0\le k<2n$ i $\varepsilon\in\{0,1\}$:
przesuwamy $j$ na prawo przez $ja^k=a^{-k}j$,
a każdą parę $j^2$ zastępujemy przez $a^n$.
Stąd grupa zadana prezentacją ma najwyżej $4n$ elementów.
Macierze $A^k$ są parami różne i diagonalne,
zaś $A^kJ$ są parami różne i mają zera na przekątnej.
Otrzymujemy $4n$ różnych macierzy, zatem naturalny
surjektywny homomorfizm z grupy zadanej prezentacją
do grupy macierzy jest izomorfizmem.
Iloraz przez $\{\pm I\}$ jest generowany przez klasy
$A,J$ o rzędach $n,2$ z relacją odwracania; ma $2n$
elementów, więc jest grupą dwuścienną $D_n$.
\end{proof}

\begin{przyklad}[Grupa kwaternionowa]
\label{aa:prz:binary-q8}
Dla $n=2$ mamy $A=\operatorname{diag}(i,-i)$.
Macierze $A,J,AJ$ mają kwadrat $-I$, a ponadto
$AJ=-JA$. Osiem elementów $\pm I,\pm A,\pm J,\pm AJ$
tworzy grupę kwaternionową $Q_8$.
Iloraz przez $\{\pm I\}$ jest grupą Kleina.
To pokazuje, że grupa nad $D_2$ nie jest grupą
$C_2\times D_2$: ta ostatnia ma siedem inwolucji,
podczas gdy $Q_8$ ma jedną.
\end{przyklad}

\begin{lemat}[Przeciwobrazy grup cyklicznych i dwuściennych]
\label{aa:lem:binary-cyclic-dihedral}
Niech $G\le\operatorname{PGL}_2(\mathbb C)$ będzie
skończona. Jeżeli $G\simeq C_n$, to jej pełny
przeciwobraz względem $\rho$ jest sprzężony
z grupą diagonalną rzędu $2n$.
Jeżeli $G\simeq D_n$ dla $n\ge2$, to jej pełny
przeciwobraz jest sprzężony z $\operatorname{Dic}_n$.
Sprzężenia można wybrać w $\operatorname{SL}_2(\mathbb C)$.
\end{lemat}
\begin{proof}
Dla $n=1$ pierwsza grupa to $\{\pm I\}$.
Dla $n>1$ wybieramy dwa punkty stałe generatora
i przenosimy je do $0,\infty$. Grupa rzutowa ma
wówczas postać $z\mapsto\eta^k z$, $\eta$ pierwotny
pierwiastek $n$-tego stopnia z jedności. Jej pełny
przeciwobraz to macierze
$\operatorname{diag}(\lambda,\lambda^{-1})$ z
$\lambda^{2n}=1$.
W przypadku dwuściennym istnieje orbita dwupunktowa
na mocy danych stabilizatorów z rozdziału 51.
Po przeniesieniu jej do $0,\infty$ grupa rzutowa jest
generowana przez $z\mapsto\eta z$ i $z\mapsto c/z$,
$c\ne0$. Sprzężenie przez $z\mapsto dz$ zastępuje
$c$ przez $d^2c$; wybieramy $d^2=c^{-1}$.
Pełny przeciwobraz jest teraz generowany przez
$A,J$ z Definicji~\ref{aa:def:binary-dicyclic}, gdyż
ich obraz rzutowy generuje $G$, a $-I=A^n$.
Każde użyte sprzężenie rzutowe ma reprezentanta
o wyznaczniku $1$ po pomnożeniu przez odpowiedni skalar.
\end{proof}

\section{Binarne grupy wielościenne}

\begin{definicja}[Binarne grupy obrotów]
\label{aa:def:binary-exceptional}
Niech $T,O,I_{\mathrm{rot}}\le\operatorname{SO}(3)$
oznaczają grupy obrotów czworościanu, sześcianu
i dwudziestościanu z przykładów
w rozdziale~\ref{aa:ch:finite-pgl2}.
Definiujemy
$2T=\pi^{-1}(T)$, $2O=\pi^{-1}(O)$ oraz
$2I=\pi^{-1}(I_{\mathrm{rot}})$.
Nazywamy je odpowiednio binarną grupą czworościenną,
ośmiościenną i dwudziestościenną.
\end{definicja}

Surjektywność $\pi$ i jej jądro dają dokładne ciągi
\[
 \begin{aligned}
 1&\to C_2\to 2T\to A_4\to1,\\
 1&\to C_2\to 2O\to S_4\to1,\\
 1&\to C_2\to 2I\to A_5\to1.
 \end{aligned}
\]
Grupy mają odpowiednio $24,48,120$ elementów.
Żaden z tych ciągów nie rozszczepia się:
sekcja homomorficzna wysyłałaby inwolucję ilorazu
na inwolucję różną od $-I$, ponieważ $-I$ należy
do jądra. Sprzeczność z Lematem~\ref{aa:lem:binary-full-lift}.
Ten sam argument dotyczy ciągu dla $\operatorname{Dic}_n$.

\begin{lemat}[Rzędy podniesień]
\label{aa:lem:binary-orders}
Niech $g\in\operatorname{SO}(3)$ ma rząd $m\ge1$.
Jego dwa przeciwobrazy $u,-u\in\operatorname{SU}(2)$
mają dla parzystego $m$ oba rząd $2m$,
a dla nieparzystego $m$ jeden rząd $m$, drugi $2m$.
\end{lemat}
\begin{proof}
Dla $m=1$ podniesienia to $\pm I$.
Dla $m>1$ po zmianie osi obrót ma kąt $2\pi k/m$,
gdzie $\gcd(k,m)=1$. Na mocy konstrukcji surjekcji
$\pi$ podniesienia są sprzężone z
$\pm\operatorname{diag}(e^{-\pi i k/m},e^{\pi i k/m})$.
Dla każdego podniesienia jego rząd jest wielokrotnością
$m$, a jego $m$-ta potęga należy do $\{\pm I\}$,
więc rząd wynosi $m$ lub $2m$.
Jeśli $m$ jest parzyste, $k$ jest nieparzyste i obie
$m$-te potęgi wynoszą $-I$. Jeśli $m$ jest nieparzyste,
zmiana znaku podniesienia zmienia znak $m$-tej potęgi,
co daje dokładnie jeden element każdego z dwóch rzędów.
\end{proof}

\begin{przyklad}[Rozkłady rzędów elementów]
\label{aa:prz:binary-order-distributions}
Z rozkładów permutacji na rozłączne cykle otrzymujemy
liczby elementów rzędów $1,2,3,4,5$ w grupach
$A_4,S_4,A_5$:
\[
\begin{array}{c|rrrrr}
 &1&2&3&4&5\\ \hline
 A_4&1&3&8&0&0\\
 S_4&1&9&8&6&0\\
 A_5&1&15&20&0&24
\end{array}
\]
Na przykład inwolucji w $A_5$ jest $5\cdot3=15$
(wybór punktu stałego i parowania czterech pozostałych).
Cykli długości $3$ jest $\binom53\cdot2=20$,
a cykli długości $5$ jest $4!=24$.
Z Lematu~\ref{aa:lem:binary-orders} wynika tabela
dla grup binarnych:
\[
\begin{array}{c|rrrrrrr}
 &1&2&3&4&5&6&8\\ \hline
 2T&1&1&8&6&0&8&0\\
 2O&1&1&8&18&0&8&12\\
 2I&1&1&20&30&24&20&0
\end{array}
\]
W ostatnim wierszu należy dodatkowo uwzględnić
$24$ elementy rzędu $10$. Sumy wynoszą $24,48,120$.
\end{przyklad}

\section{Jednoznaczność z dokładnością do sprzężenia}

\begin{lemat}[Charakter reprezentacji obrotów]
\label{aa:lem:binary-rotation-character}
Niech $K$ będzie jedną z grup $A_4,S_4,A_5$.
Jeżeli $r_1,r_2:K\to\operatorname{SO}(3)$ są iniektywnymi
homomorfizmami, to po ewentualnym złożeniu $r_2$
z automorfizmem $K$ ich zespolone reprezentacje
trójwymiarowe mają ten sam charakter.
\end{lemat}
\begin{proof}
Macierz obrotu ma wartości własne $1,e^{it},e^{-it}$.
Dla elementu rzędu $m$ kąt ma postać $2\pi k/m$,
gdzie $\gcd(k,m)=1$. Ślad wynosi $1+2\cos(2\pi k/m)$.
Dla $m=1,2,3,4$ ślad zależy wyłącznie od $m$:
odpowiednio $3,-1,0,1$. To wystarcza dla $A_4,S_4$.

W $A_5$ cykle długości pięć tworzą dwie klasy
sprzężoności po $12$ elementów. Istotnie centralizator
cyklu $c$ długości pięć w $S_5$ jest $\langle c\rangle$:
permutacja komutująca z $c$ jest określona przez obraz
jednego punktu. Wszystkie jego elementy są parzyste,
więc centralizator w $A_5$ ma rząd $5$.
Każda klasa ma zatem $60/5=12$ elementów.
Elementy $c,c^2$ nie są sprzężone w $A_5$:
po oznaczeniu punktów przez $\mathbb Z/5\mathbb Z$
i przyjęciu $c:x\mapsto x+1$ każde sprzężenie $c$ z $c^2$
ma postać $x\mapsto 2x+b$. Mnożenie przez $2$
jest czterocyklem na elementach niezerowych,
a przesunięcie jest parzyste, więc taki sprzęgający
element jest nieparzysty.
Ślady $r_\ell(c)$ i $r_\ell(c^2)$ są różne:
wynoszą w pewnej kolejności
$1+2\cos(2\pi/5)$ oraz $1+2\cos(4\pi/5)$.
Sprzężenie $K=A_5$ przez nieparzystą permutację
daje automorfizm zamieniający te dwie klasy
(gdyby zachowywał klasę $c$, istniałby nieparzysty
element centralizatora $c$).
Możemy więc uzgodnić wartości obu charakterów
również na elementach rzędu pięć.
\end{proof}

\begin{lemat}[Od równoważności zespolonej do sprzężenia obrotów]
\label{aa:lem:binary-real-intertwiner}
Niech $K$ będzie skończoną grupą,
a $r_1,r_2:K\to\operatorname{SO}(3)$ homomorfizmami.
Jeśli ich zespolone reprezentacje mają ten sam charakter,
to istnieje $Q\in\operatorname{SO}(3)$ spełniające
$Qr_1(k)Q^{-1}=r_2(k)$ dla każdego $k\in K$.
\end{lemat}
\begin{proof}
Twierdzenie o charakterach reprezentacji skończonych
nad $\mathbb C$ (rozdział 42) daje odwracalną
zespoloną macierz $Z$ taką, że
$Zr_1(k)=r_2(k)Z$.
Zapiszmy $Z=X+iY$ dla macierzy rzeczywistych $X,Y$.
Obie spełniają te same równania liniowe, ponieważ
$r_1(k),r_2(k)$ są rzeczywiste.
Wielomian $p(t)=\det(X+tY)$ nie jest zerowy,
gdyż $p(i)=\det Z\ne0$. Istnieje zatem $t\in\mathbb R$
takie, że $L=X+tY$ jest odwracalna.
Niech $D=L^{\mathsf T}L$.
Z równania $Lr_1(k)=r_2(k)L$ i ortogonalności
otrzymujemy $r_1(k)^{\mathsf T}Dr_1(k)=D$,
a zatem $D$ komutuje z każdym $r_1(k)$.
Każda przestrzeń własna dodatnio określonej
macierzy symetrycznej $D$ jest przez nie zachowywana;
więc również $D^{-1/2}$ komutuje z $r_1(k)$.
Macierz $Q_0=LD^{-1/2}$ jest ortogonalna
i nadal spełnia równania przeplatania.
Jeśli jej wyznacznik jest ujemny, zastępujemy ją
przez $-Q_0$, co w wymiarze trzy zmienia znak wyznacznika,
ale nie zmienia sprzężenia.
\end{proof}

\begin{twierdzenie}[Klasyfikacja skończonych podgrup specjalnej grupy liniowej]
\label{aa:tw:binary-classification}
Każda skończona podgrupa $\operatorname{SL}_2(\mathbb C)$
jest sprzężona w tej grupie z jedną z następujących:
\[
\begin{gathered}
 C_m^{\mathrm{diag}}=
 \{\operatorname{diag}(\lambda,\lambda^{-1}):\lambda^m=1\},
 \quad m\ge1,\\
 \operatorname{Dic}_n\quad(n\ge2),\qquad 2T,\quad 2O,\quad 2I.
\end{gathered}
\]
Typy w tej liście są parami różne z dokładnością
do izomorfizmu (z wyjątkiem powtórzenia identycznego parametru).
\end{twierdzenie}
\begin{proof}
Na mocy Twierdzenia~\ref{aa:tw:binary-unitary}
możemy przyjąć $H\subset\operatorname{SU}(2)$.
Jeżeli $|H|$ jest nieparzyste, grupa jest cykliczna
na mocy Lematu~\ref{aa:lem:binary-full-lift}; jej generator
jest diagonalizowalny unitarnie, a macierz sprzęgającą
można pomnożyć przez skalar, by miała wyznacznik $1$.
Otrzymujemy $C_{|H|}^{\mathrm{diag}}$.
Jeśli $|H|$ jest parzyste, jest pełnym przeciwobrazem
swego obrazu rzutowego. Dla obrazu cyklicznego lub
dwuściennego stosujemy Lemat~\ref{aa:lem:binary-cyclic-dihedral}.
W pozostałych przypadkach obraz $\pi(H)$ jest
izomorficzny z jedną z grup $K=A_4,S_4,A_5$.
Wybieramy izomorfizmy $K\to\pi(H)$ i $K\to G_0$,
gdzie $G_0$ to odpowiednia grupa obrotów
z Definicji~\ref{aa:def:binary-exceptional}.
Lematy~\ref{aa:lem:binary-rotation-character}
i~\ref{aa:lem:binary-real-intertwiner} dają
$Q\in\operatorname{SO}(3)$ sprzęgające te obrazy
(jako podgrupy; automorfizm $K$ nie zmienia obrazu).
Wybierzmy $U\in\pi^{-1}(Q)$. Wówczas
$UHU^{-1}=\pi^{-1}(G_0)$, ponieważ obie strony
są pełnymi przeciwobrazami $G_0$.
To dowodzi klasyfikacji z dokładnością do sprzężenia.

Grupy cykliczne odróżniamy od pozostałych przez
przemienność i cykliczność. Grupy dwucykliczne
są nieprzemienne dla $n\ge2$, a ich rzędy wyznaczają $n$.
Grupy $2T,2O,2I$ mają różne rzędy.
Gdy grupa dwucykliczna ma jeden z tych rzędów,
ma element rzędu odpowiednio $12,24,60$ (generator
podgrupy diagonalnej). Tabela z
Przykładu~\ref{aa:prz:binary-order-distributions}
wyklucza te rzędy dla odpowiednich grup wielościennych.
\end{proof}

\section{Naturalna reprezentacja i równanie Riccatiego}

\begin{wniosek}[Nieprzywiedlność i prymitywność]
\label{aa:wn:binary-primitive}
Niech $H\le\operatorname{SL}_2(\mathbb C)$ będzie skończona.
Jej naturalna reprezentacja na $\mathbb C^2$ jest
przywiedlna dokładnie dla grup cyklicznych,
nieprzywiedlna imprymitywna dokładnie dla grup
$\operatorname{Dic}_n$ ($n\ge2$) i prymitywna dokładnie
dla $2T,2O,2I$ (z dokładnością do sprzężenia).
\end{wniosek}
\begin{proof}
Grupa cykliczna ma wspólną prostą własną.
Jeżeli $H$ ma prostą niezmienniczą, po sprowadzeniu
do grupy unitarnej jej ortogonalne dopełnienie jest
również niezmiennicze. W odpowiedniej bazie $H$ jest
podgrupą macierzy diagonalnych wyznacznika $1$,
więc zanurza się w $\mathbb C^\times$ i jest cykliczna.
Dla grup dwucyklicznych macierze $A,J$ zachowują
nieuporządkowaną parę osi współrzędnych, lecz grupa
nie jest cykliczna. Pozostałe trzy grupy są
nieprzywiedlne; zachowanie pary prostych dawałoby
orbitę rzutową o co najwyżej dwóch elementach.
Najmniejsze długości orbit, obliczone w
rozdziale~\ref{aa:ch:finite-pgl2}, wynoszą odpowiednio
$4,6,12$. Zatem ich reprezentacje są prymitywne.
\end{proof}

\begin{stwierdzenie}[Wielomian związany z orbitą prostych]
\label{aa:st:binary-orbit-form}
Niech $H\le\operatorname{SL}_2(\mathbb C)$ będzie skończona,
a $\mathcal O$ orbitą działania $H$ na
$\mathbb P^1(\mathbb C)$. Dla każdej prostej
$L\in\mathcal O$ wybierzmy niezerowy liniowy funkcjonał
$\ell_L$ na $\mathbb C^2$ o jądrze $L$.
Wielomian jednorodny
$F_{\mathcal O}(v)=\prod_{L\in\mathcal O}\ell_L(v)$
ma stopień $|\mathcal O|$. Istnieje charakter
jednowymiarowy $\chi:H\to\mathbb C^\times$ taki, że
$F_{\mathcal O}(h^{-1}v)=\chi(h)F_{\mathcal O}(v)$
dla $h\in H$ i $v\in\mathbb C^2$.
Jeśli $-I\in H$, to $\chi(-I)=(-1)^{|\mathcal O|}$.
\end{stwierdzenie}
\begin{proof}
Funkcjonał $\ell_L\circ h^{-1}$ ma jądro $hL$,
więc jest niezerową wielokrotnością $\ell_{hL}$.
Ponieważ $h$ permutuje orbitę, iloczyn tych
funkcjonałów jest skalarną wielokrotnością
$F_{\mathcal O}$. Skalar jest jednoznaczny,
bo wielomian nie jest zerowy.
Działanie $(h\cdot F)(v)=F(h^{-1}v)$ jest reprezentacją
na przestrzeni wielomianów jednorodnych;
porównanie działań $h_1h_2$ oraz $h_1,h_2$ pokazuje
$\chi(h_1h_2)=\chi(h_1)\chi(h_2)$.
Ostatnia równość wynika z jednorodności wielomianu.
\end{proof}

\begin{przyklad}[Para osi dla grupy dwucyklicznej]
Dla $\operatorname{Dic}_n$ orbita osi współrzędnych
daje wielomian $F(x,y)=xy$.
Działanie $A$ zachowuje $F$, zaś
$J^{-1}(x,y)=(-y,x)$ daje $(J\cdot F)(x,y)=-xy$.
Wielomian nie jest więc niezmiennikiem, lecz rozpina
jednowymiarową podprzestrzeń niezmienniczą.
Jego kwadrat jest już niezmiennikiem.
\end{przyklad}

Dla grup $2T,2O,2I$ konstrukcja orbit daje
wielomiany półniezmiennicze stopni $4,6,12$.
Zauważmy, że istnienie takiej podprzestrzeni
w potędze symetrycznej reprezentacji sprzężonej
nie jest równoznaczne z niezmienniczością każdego
wektora tej podprzestrzeni.
W przypadku skończonej prymitywnej różniczkowej
grupy Galois równania $z''=qz$ z $q\in\mathbb C(x)$
Twierdzenie~\ref{aa:tw:binary-classification} pozwala
zidentyfikować pełną grupę macierzową, nie tylko jej
obraz rzutowy. Stopnie orbit są zgodne z parametrami
trzeciego przypadku algorytmu Kovacica.
Nie wyprowadzamy z samej klasyfikacji pełnych
warunków lokalnych algorytmu ani tablic charakterów
grup binarnych.

% AA-53 — pierścienie niezmienników.
\chapter{Pierścienie niezmienników grup skończonych}
\label{aa:ch:finite-invariants}
W rozdziale~\ref{aa:ch:binary-polyhedral} rozróżniliśmy
niezmienniki i półniezmienniki. Teraz wyznaczymy
całe pierścienie niezmienników dla dwóch nieskończonych
rodzin grup oraz metodę obliczania wymiarów ich części
jednorodnych. Pracujemy nad ciałem $\mathbb C$.

\section{Działanie na wielomianach}

\begin{definicja}[Pierścień niezmienników i szereg Hilberta]
\label{aa:def:invariants-hilbert}
Niech $V$ będzie skończenie wymiarową zespoloną
przestrzenią liniową, a $G\le\operatorname{GL}(V)$
grupą skończoną. Niech $S=\mathbb C[V]$ oznacza
pierścień funkcji wielomianowych na $V$,
a $S_d$ jego część jednorodną stopnia $d\ge0$.
Dla $g\in G$, $f\in S$ i $v\in V$ definiujemy
$(g\cdot f)(v)=f(g^{-1}v)$. Jest to działanie grupy,
gdyż $1\cdot f=f$ oraz
$(g\cdot(h\cdot f))(v)=f(h^{-1}g^{-1}v)
=((gh)\cdot f)(v)$.
Pierścieniem niezmienników nazywamy
$S^G=\{f\in S:g\cdot f=f\text{ dla każdego }g\in G\}$.
Jego szeregiem Hilberta nazywamy formalny szereg
$H_G(z)=\sum_{d\ge0}\dim_{\mathbb C}(S_d^G)z^d$.
\end{definicja}

Działanie zachowuje stopień oraz iloczyn, więc
$S^G=\bigoplus_{d\ge0}S_d^G$ jest podpierścieniem
z gradacją. Szereg Hilberta opisuje wymiary, ale sam
nie określa mnożenia ani prezentacji pierścienia.

\begin{lemat}[Operator Reynoldsa]
\label{aa:lem:invariants-reynolds}
Dla $G,V,S$ z Definicji~\ref{aa:def:invariants-hilbert}
odwzorowanie $\mathcal R:S\to S$ określone wzorem
$\mathcal R(f)=|G|^{-1}\sum_{g\in G}g\cdot f$
jest rzutem liniowym na $S^G$, zachowującym stopnie.
Dla $a\in S^G$, $f\in S$ zachodzi
$\mathcal R(af)=a\mathcal R(f)$.
\end{lemat}
\begin{proof}
Dla $h\in G$ lewostronne mnożenie przez $h$ permutuje
składniki sumy, więc $h\cdot\mathcal R(f)=\mathcal R(f)$.
Jeśli $f\in S^G$, każdy składnik sumy jest równy $f$.
Stąd obraz jest równy $S^G$ i $\mathcal R^2=\mathcal R$.
Z liniowości działania na $V$ wynika zachowanie stopni.
Ponadto $g\cdot(af)=a(g\cdot f)$, co po uśrednieniu
daje ostatnią równość.
\end{proof}

\begin{twierdzenie}[Skończona generowalność niezmienników]
\label{aa:tw:invariants-finite}
Niech $V$ będzie skończenie wymiarową zespoloną
przestrzenią liniową, a $G\le\operatorname{GL}(V)$
grupą skończoną. Wówczas pierścień niezmienników
$\mathbb C[V]^G$ jest skończenie generowaną algebrą
nad $\mathbb C$.
\end{twierdzenie}
\begin{proof}
Oznaczmy $S=\mathbb C[V]$. Niech $I$ będzie ideałem
w $S$ generowanym przez
wszystkie jednorodne niezmienniki dodatniego stopnia.
Z twierdzenia Hilberta o bazie pierścień $S$ jest
noetherowski. Istnieje więc skończona rodzina
jednorodnych niezmienników $f_1,\ldots,f_r$ dodatniego
stopnia generująca $I$ jako ideał $S$:
każdy z dowolnych skończonych generatorów $I$
jest skończoną kombinacją pierwotnych generatorów,
wystarczy zebrać wszystkie występujące w niej niezmienniki.
Jeśli $I=0$, pierścień niezmienników składa się ze stałych.

Dla jednorodnego niezmiennika $f$ stopnia $d>0$
mamy $f=\sum_i a_i f_i$. Biorąc składniki jednorodne
stopnia $d$, możemy przyjąć $a_i$ jednorodne stopnia
$d-\deg f_i$ (lub zerowe przy ujemnym stopniu).
Operator Reynoldsa daje $f=\sum_i\mathcal R(a_i)f_i$.
Każdy niezerowy $\mathcal R(a_i)$ ma stopień mniejszy
od $d$. Indukcja względem $d$, rozpoczynająca się od
stałych, daje $f\in\mathbb C[f_1,\ldots,f_r]$.
Rozkład na części jednorodne kończy dowód.
\end{proof}

\section{Wzór Moliena}

\begin{twierdzenie}[Wzór Moliena]
\label{aa:tw:molien}
Niech $V$ będzie skończenie wymiarową zespoloną
przestrzenią liniową, a $G\le\operatorname{GL}(V)$
grupą skończoną. Niech $H_G(z)$ będzie szeregiem
Hilberta pierścienia $\mathbb C[V]^G$ dla gradacji
według stopnia wielomianów. W pierścieniu
$\mathbb C[[z]]$ zachodzi
\[
 H_G(z)=\frac{1}{|G|}\sum_{g\in G}
 \frac{1}{\det(I-zg)}.
\]
\end{twierdzenie}
\begin{proof}
Oznaczmy przez $S_d$ przestrzeń wielomianów
jednorodnych stopnia $d$ na $V$ i połóżmy
$m=\dim_{\mathbb C}V$. Ograniczenie operatora Reynoldsa
do $S_d$ jest rzutem, którego ślad wynosi $\dim S_d^G$:
w bazie dostosowanej do obrazu i jądra rzut ma na
przekątnej jedynki na obrazie i zera na jądrze.
Jeśli wartości własne $g^{-1}$ na $V$, z uwzględnieniem
krotności, to $\lambda_1,\ldots,\lambda_m$,
w bazie funkcjonałów liniowych sprzężonej do bazy
wektorów własnych $g^{-1}$ monom stopnia $d$,
o wykładnikach $\alpha_1,\ldots,\alpha_m$,
ma wartość własną $\lambda_1^{\alpha_1}
\cdots\lambda_m^{\alpha_m}$ dla działania $g$.
Taka baza istnieje, bo $g$ ma skończony rząd
i jego wielomian minimalny dzieli wielomian bez
pierwiastków wielokrotnych $T^{|G|}-1$.
Sumując ślady po stopniach, otrzymujemy
\[
 \sum_{d\ge0}\operatorname{tr}(g|_{S_d})z^d
 =\prod_{j=1}^m(1-\lambda_jz)^{-1}
 =\det(I-zg^{-1})^{-1}.
\]
Iloczyny są formalnymi szeregami geometrycznymi;
w każdym stopniu występuje skończenie wiele składników.
Uśrednienie po $G$ i zastąpienie $g^{-1}$ przez $g$
w skończonej sumie daje wzór.
\end{proof}

Dla $G\le\operatorname{SL}_2(\mathbb C)$ mianownik
ma postać $1-\operatorname{tr}(g)z+z^2$.
W szczególności wystarczy znać liczności elementów
o poszczególnych śladach, aby wyznaczyć szereg Hilberta.

\section{Grupy cykliczne}

\begin{twierdzenie}[Niezmienniki grupy cyklicznej]
\label{aa:tw:cyclic-invariant-ring}
Niech $m\ge1$ będzie całkowita, $\zeta$ pierwotnym
pierwiastkiem $m$-tego stopnia z jedności, a
$G=\langle\operatorname{diag}(\zeta,\zeta^{-1})\rangle$.
Dla działania na $\mathbb C[x,y]$ jak w
Definicji~\ref{aa:def:invariants-hilbert} mamy
\[
 \mathbb C[x,y]^G\simeq
 \mathbb C[U,V,W]/(UV-W^m),
 \quad (U,V,W)\longmapsto(x^m,y^m,xy).
\]
Stopnie generatorów to $m,m,2$ i
\[
 H_G(z)=\frac{1+z^m}{(1-z^m)(1-z^2)}.
\]
\end{twierdzenie}
\begin{proof}
Monom $x^a y^b$ jest niezmiennikiem dokładnie wtedy,
gdy $m$ dzieli $a-b$. Ponieważ monomy tworzą bazę
wektorów własnych generatora, każdy niezmiennik
jest kombinacją takich monomów. Jeśli $a\ge b$,
monom ma postać $(xy)^b(x^m)^{(a-b)/m}$;
dla $a<b$ postępujemy analogicznie z $y^m$.
To dowodzi surjektywności wskazanego homomorfizmu.

Równość $UV=W^m$ pozwala zredukować każdy monom
w ilorazie do jednego z monomów $W^k$,
$W^kU^j$, $W^kV^j$ ($k\ge0$, $j\ge1$).
Ich obrazy w $\mathbb C[x,y]$ są parami różnymi
monomami (rozróżnia je znak i wartość $a-b$
oraz $\min(a,b)$). Są zatem liniowo niezależne.
Dowodzi to iniektywności, a zliczanie ich stopni daje
\[
 H_G(z)=\frac{1}{1-z^2}
 \left(1+2\sum_{j\ge1}z^{mj}\right)
 =\frac{1+z^m}{(1-z^2)(1-z^m)}.
\]
\end{proof}

\begin{przyklad}[Grupa rzędu dwa]
Dla $m=2$ niezmiennikami są dokładnie wielomiany
o monomach parzystego stopnia całkowitego.
Generatory $x^2,y^2,xy$ spełniają jedyną relację
$UV=W^2$. Część stopnia $2k$ ma wymiar $2k+1$,
a część nieparzystego stopnia jest zerowa.
Dla $m=1$ relacja $W=UV$ eliminuje trzeci generator
i odzyskujemy cały pierścień $\mathbb C[x,y]$.
\end{przyklad}

\section{Grupy dwucykliczne}

\begin{twierdzenie}[Niezmienniki grupy dwucyklicznej]
\label{aa:tw:dicyclic-invariant-ring}
Niech $n\ge2$ i niech $\operatorname{Dic}_n$
działa na $\mathbb C^2$ przez macierze $A,J$
z Definicji~\ref{aa:def:binary-dicyclic}.
Połóżmy
\[
 u=(xy)^2,\quad v=x^{2n}+y^{2n},\quad
 w=xy(x^{2n}-y^{2n}).
\]
Wówczas
\[
 \mathbb C[x,y]^{\operatorname{Dic}_n}\simeq
 \mathbb C[U,V,W]/(W^2-UV^2+4U^{n+1}),
\]
gdzie $U,V,W$ przechodzą odpowiednio na $u,v,w$.
Szereg Hilberta wynosi
\[
 H_{\operatorname{Dic}_n}(z)
 =\frac{1+z^{2n+2}}{(1-z^4)(1-z^{2n})}.
\]
\end{twierdzenie}
\begin{proof}
Oznaczmy $t=xy$, $a=x^{2n}$, $b=y^{2n}$.
Z Twierdzenia~\ref{aa:tw:cyclic-invariant-ring}
niezmienniki podgrupy $\langle A\rangle$ mają bazę
$t^k$, $t^ka^j$, $t^kb^j$ ($k\ge0,j\ge1$),
a $ab=t^{2n}$. Działanie $J$ wysyła
$(t,a,b)$ na $(-t,b,a)$, ponieważ $J^{-1}(x,y)=(-y,x)$.
Ponieważ $A,J$ generują $\operatorname{Dic}_n$,
wystarczy wyznaczyć część pierścienia
$\mathbb C[x,y]^{\langle A\rangle}$ ustaloną przez $J$.
Jej bazę tworzą $t^{2k}$ oraz
\[
 t^{2k}(a^j+b^j),\qquad
 t^{2k+1}(a^j-b^j),\quad k\ge0,\ j\ge1.
\]
Istotnie dla każdego $k,j$ działanie na parze
$t^ka^j,t^kb^j$ zamienia wektory z czynnikiem $(-1)^k$;
przestrzeń ustalona jest jednowymiarowa.

Sumy $P_j=a^j+b^j$ spełniają rekurencję
$P_j=vP_{j-1}-u^nP_{j-2}$ dla $j\ge2$,
z $P_0=2,P_1=v$, bo $a+b=v$, $ab=u^n$.
Wielomiany $Q_j=\sum_{i=0}^{j-1}a^{j-1-i}b^i$
spełniają tę samą rekurencję dla $j\ge3$,
z $Q_1=1,Q_2=v$, i $a^j-b^j=(a-b)Q_j$.
Istotnie, dla $j\ge3$ iloczyn $(a+b)Q_{j-1}$
zawiera każdy skrajny monom $a^{j-1},b^{j-1}$
raz, a każdy monom pośredni dwa razy.
Odjęcie $abQ_{j-2}$ usuwa dokładnie jedno
wystąpienie każdego monomu pośredniego.
Każdy element wskazanej bazy jest więc kombinacją
wielomianów w $u,v,w$, a te trzy są niezmiennikami.
Ponadto
$w^2=t^2((a+b)^2-4ab)=uv^2-4u^{n+1}$.

Aby wykluczyć dalsze relacje, najpierw pokażemy
algebraiczną niezależność $u,v$.
Podstawienie $x=X,y=T/X$ w pierścieniu
$\mathbb C[T,X,X^{-1}]$ daje
$u=T^2$, $v=X^{2n}+T^{2n}X^{-2n}$.
Gdyby $\sum_{j=0}^r p_j(u)v^j=0$, $p_r\ne0$,
współczynnik przy najwyższej potędze $X^{2nr}$
wynosiłby $p_r(T^2)\ne0$, sprzeczność.
Każdy element ilorazu przez moniczną relację
kwadratową w $W$ zapisuje się jednoznacznie jako
$P(U,V)+WQ(U,V)$ (dzielenie przez wielomian moniczny).
Jeśli $P(u,v)+wQ(u,v)=0$, zamiana $x$ na $-x$
zachowuje $u,v$, a zmienia znak $w$.
Dodanie i odjęcie obu równości daje
$P(u,v)=0$ i $wQ(u,v)=0$.
Ponieważ $\mathbb C[x,y]$ nie ma dzielników zera,
$w\ne0$, a $u,v$ są algebraicznie niezależne,
otrzymujemy $P=Q=0$. To dowodzi iniektywności.

Monomy $u^iv^j$ i $wu^iv^j$ ($i,j\ge0$)
tworzą więc bazę jednorodną. Ich stopnie to
$4i+2nj$ oraz $2n+2+4i+2nj$, co daje szereg Hilberta.
\end{proof}

\begin{przyklad}[Niezmienniki grupy kwaternionowej]
Dla $n=2$ mamy $\operatorname{Dic}_2\simeq Q_8$.
Generatory $u,v$ mają stopień $4$, a $w$ stopień $6$;
jedyna relacja ma stopień $12$ i postać
$w^2=uv^2-4u^3$. W stopniu $4$ otrzymujemy bazę
$u,v$, w stopniu $6$ bazę $w$, a w stopniu $8$
bazę $u^2,uv,v^2$. Nie istnieją niezmienniki
dodatniego stopnia mniejszego niż $4$.
Natomiast $xy$ jest półniezmiennikiem stopnia $2$.
\end{przyklad}

\begin{wniosek}[Najmniejsze stopnie niezmienników]
\label{aa:wn:dicyclic-invariant-minimum}
Dla każdego $n\ge2$ minimalny dodatni stopień
niezmiennika grupy $\operatorname{Dic}_n$ wynosi $4$,
natomiast minimalny dodatni stopień półniezmiennika
wynosi $2$.
\end{wniosek}
\begin{proof}
Z szeregu Hilberta wynika brak niezmienników
w stopniach $1,2,3$ i istnienie $u$ stopnia $4$.
Wielomian $xy$ rozpina jednowymiarową podprzestrzeń
niezmienniczą, ale półniezmiennik stopnia $1$
dawałby niezmienniczą prostą w przestrzeni sprzężonej,
a więc jej jednowymiarowy anihilator byłby niezmienniczy
w $\mathbb C^2$. Jest to niemożliwe z nieprzywiedlności
w Wniosku~\ref{aa:wn:binary-primitive}.
\end{proof}

Prezentacje dla grup wyjątkowych $2T,2O,2I$ wymagają
dodatkowych argumentów o generatorach i relacjach.
Sam wzór Moliena nie zastępuje tych argumentów.

% AA-54 — wyjątkowe binarne grupy wielościenne
\chapter{Niezmienniki binarnych grup wielościennych}
\label{aa:ch:exceptional-invariants}

W rozdziale~\ref{aa:ch:finite-invariants} wyznaczyliśmy
pierścienie niezmienników grup cyklicznych i dwucyklicznych.
Pozostały trzy grupy wyjątkowe z
Twierdzenia~\ref{aa:tw:binary-classification}.
Pracujemy w modelach kwaternionowych z
Przykładu~\ref{aa:prz:binary-quaternion-generators}.
Jeżeli $G'=LGL^{-1}$ dla $L\in\operatorname{GL}_2(\mathbb C)$,
to $P\mapsto P\circ L^{-1}$ jest izomorfizmem
$\mathbb C[x,y]^G\to\mathbb C[x,y]^{G'}$:
dla $g'=LgL^{-1}$ mamy
$(P\circ L^{-1})(g'v)=P(gL^{-1}v)=P(L^{-1}v)$.
Odwzorowanie odwrotne jest dane przez złożenie z $L$.
Oba odwzorowania zachowują stopnie.

\section{Szeregi Hilberta}

\begin{twierdzenie}[Szeregi Hilberta grup wyjątkowych]
\label{aa:tw:exceptional-hilbert}
Niech $G$ będzie jedną z grup $2T,2O,2I$ działających
naturalnie na $\mathbb C^2$. Dla pierścienia
$\mathbb C[x,y]^G$ z gradacją według stopnia mamy
\[
\begin{aligned}
 H_{2T}(z)&=\frac{1+z^{12}}{(1-z^6)(1-z^8)},\\
 H_{2O}(z)&=\frac{1+z^{18}}{(1-z^8)(1-z^{12})},\\
 H_{2I}(z)&=\frac{1+z^{30}}{(1-z^{12})(1-z^{20})}.
\end{aligned}
\]
\end{twierdzenie}
\begin{proof}
Niech $\phi=(1+\sqrt5)/2$. Ślady elementów
$\operatorname{SU}(2)$ są rzeczywiste.
Z konstrukcji nakrycia w rozdziale 52 wynika,
że dwa podniesienia obrotu o kąt $\theta$
mają ślady $\pm2\cos(\theta/2)$.
Korzystając z rozkładów rzędów w
Przykładzie~\ref{aa:prz:binary-order-distributions}
i z dwóch klas cykli długości pięć w $A_5$,
otrzymujemy tabelę liczności. W kolumnach oznaczonych
$\pm a$ podajemy liczność dla \emph{każdego} znaku.
\[
\begin{array}{c|rrrrrr}
 &\pm2&0&\pm1&\pm\sqrt2&\pm\phi&\pm(\phi-1)\\ \hline
2T&1&6&8&0&0&0\\
2O&1&18&8&6&0&0\\
2I&1&30&20&0&12&12
\end{array}
\]
Istotnie ślady pochodzące od obrotów rzędów
$2,3,4$ to odpowiednio $0,\pm1,\pm\sqrt2$.
Dwie klasy obrotów rzędu pięć mają po dwanaście
elementów i kąty, z dokładnością do znaku,
$2\pi/5$ oraz $4\pi/5$; ich podniesienia mają ślady
$\pm\phi$ oraz $\pm(\phi-1)$.

We wzorze Moliena
(Twierdzenie~\ref{aa:tw:molien}) łączymy składniki
odpowiadające przeciwnym śladom. Po podstawieniu
$s=z^2$ oznaczmy
\[
\begin{aligned}
 A&=\frac{2(1+s)}{(1-s)^2},&
 B&=\frac1{1+s},\\
 C&=\frac{2(1+s)}{1+s+s^2},&
 D&=\frac{2(1+s)}{1+s^2},\\
 E&=\frac{2(1+s)(2+s+2s^2)}
 {1+s+s^2+s^3+s^4}.
\end{aligned}
\]
Ostatni składnik jest sumą dwóch par odpowiadających
śladom $\pm\phi,\pm(\phi-1)$:
dla pary $\pm a$ mianownik po połączeniu to
$(1+s)^2-a^2s$. Iloczyn dwóch takich mianowników
dla $a=\phi,\phi-1$ wynosi
$(1+s)^4-3s(1+s)^2+s^2=1+s+s^2+s^3+s^4$,
gdyż $\phi^2+(\phi-1)^2=3$ oraz
$\phi^2(\phi-1)^2=1$.
Po sprowadzeniu obu par do tego mianownika licznik
wynosi $2(1+s)(2(1+s)^2-3s)$, co daje $E$.
Zatem
\[
\begin{aligned}
24H_{2T}&=A+6B+8C,\\
48H_{2O}&=A+18B+8C+6D,\\
120H_{2I}&=A+30B+20C+12E.
\end{aligned}
\]
Po sprowadzeniu do wspólnego mianownika dostajemy
odpowiednio
\[
 \frac{1+s^6}{(1-s^3)(1-s^4)},\quad
 \frac{1+s^9}{(1-s^4)(1-s^6)},\quad
 \frac{1+s^{15}}{(1-s^6)(1-s^{10})}
\]
dla $H_{2T},H_{2O},H_{2I}$. Równości można sprawdzić
przez wymnożenie obu stron przez iloczyn mianowników;
używa się $1-s^k=(1-s)(1+s+\cdots+s^{k-1})$.
Powrót do $z$ daje tezę.
\end{proof}

\section{Kryterium kompletności generatorów}

\begin{lemat}[Kryterium prezentacji kwadratowej]
\label{aa:lem:exceptional-presentation}
Niech $G\le\operatorname{SL}_2(\mathbb C)$ będzie
skończona, a $a,b,c\in\mathbb C[x,y]^G$ niezerowymi
jednorodnymi niezmiennikami dodatnich stopni
$p,q,r$. Załóżmy, że $a,b$ są algebraicznie niezależne,
$c^2=P(a,b)$ dla $P\in\mathbb C[U,V]$,
a $P(U,V)$ nie jest kwadratem w $\mathbb C(U,V)$.
Jeżeli
\[
 H_G(z)=\frac{1+z^r}{(1-z^p)(1-z^q)},
\]
to homomorfizm
\[
 \mathbb C[U,V,W]/(W^2-P(U,V))
 \longrightarrow\mathbb C[x,y]^G,
 \quad (U,V,W)\longmapsto(a,b,c)
\]
jest izomorfizmem.
\end{lemat}
\begin{proof}
Dzielenie przez wielomian moniczny w $W$ daje
jednoznaczną postać $L(U,V)+WM(U,V)$ elementu
ilorazu. Gdy jej obraz jest zerowy i $M\ne0$,
niezależność algebraiczna daje $M(a,b)\ne0$.
Wówczas $c=-L(a,b)/M(a,b)$, zatem po podniesieniu
do kwadratu $P$ jest kwadratem w $\mathbb C(U,V)$,
sprzecznie z założeniem. Jeśli $M=0$, niezależność
$a,b$ daje $L=0$. Homomorfizm jest więc iniektywny.
Nadajmy zmiennym $U,V,W$ wagi $p,q,r$.
Ponieważ $P(a,b)=c^2$ jest jednorodny stopnia $2r$
i $a,b$ są algebraicznie niezależne, każda część
$P$ o wadze różnej od $2r$ jest zerowa.
Relacja jest więc jednorodna dla tych wag,
a homomorfizm zachowuje gradację.
Jego obraz ma bazę jednorodną $a^ib^j$ i $ca^ib^j$,
$i,j\ge0$, oraz szereg Hilberta równy podanemu
$H_G$. Ponieważ obraz jest podprzestrzenią
przestrzeni niezmienników w każdym stopniu,
równość skończonych wymiarów daje równość
tych przestrzeni w każdym stopniu.
\end{proof}

\begin{lemat}[Dwa kryteria algebraiczne]
\label{aa:lem:exceptional-jacobian}
Niech $a,b\in\mathbb C[x,y]$.
Jeżeli $J(a,b)=a_xb_y-a_yb_x\ne0$, to $a,b$ są
algebraicznie niezależne. Ponadto niezerowy
$P\in\mathbb C[U,V]$ nieparzystego stopnia względem
$V$ nie jest kwadratem w $\mathbb C(U,V)$.
\end{lemat}
\begin{proof}
Gdyby istniała niezerowa relacja $Q(a,b)=0$,
wybralibyśmy $Q$ o najmniejszym stopniu całkowitym.
Nie jest stała. Różniczkowanie względem $x,y$ daje
układ jednorodny z niewiadomymi
$Q_U(a,b),Q_V(a,b)$ i wyznacznikiem $J(a,b)$.
Nad ciałem $\mathbb C(x,y)$ układ jest odwracalny,
więc oba wyrażenia znikają. Co najmniej jedna
pochodna wielomianu $Q$ jest niezerowa
(charakterystyka zero) i ma mniejszy stopień,
co przeczy minimalności.

Dla drugiej części traktujemy $P$ jako wielomian
w $V$ nad ciałem $\mathbb C(U)$. Gdyby
$P=(R/S)^2$ dla niezerowych $R,S\in\mathbb C(U)[V]$,
porównanie stopni w równości $PS^2=R^2$
dałoby $\deg_V P+2\deg_V S=2\deg_V R$,
co jest niemożliwe.
\end{proof}

\section{Grupy czworościenna i ośmiościenna}

Połóżmy
\[
\begin{aligned}
 f&=xy(x^4-y^4),\\
 g&=x^8+14x^4y^4+y^8,\\
 t&=x^{12}-33x^8y^4-33x^4y^8+y^{12}.
\end{aligned}
\]

\begin{lemat}[Sprawdzenie niezmienniczości]
\label{aa:lem:exceptional-small-forms}
W modelach z Przykładu~\ref{aa:prz:binary-quaternion-generators}
wielomiany $f,g,t$ są niezmiennikami $2T$.
Dla generatora $r=(1+\mathbf i)/\sqrt2$ grupy $2O$
podstawienie przez $M(r)$ wysyła je odpowiednio
na $-f,g,-t$. Zachodzi $t^2=g^3-108f^4$.
\end{lemat}
\begin{proof}
Macierze $M(\mathbf i),M(\mathbf j)$ dają
podstawienia $(ix,-iy)$ i $(y,-x)$; każdy z trzech
wielomianów jest przez nie zachowywany.
Dla $h=(1+\mathbf i+\mathbf j+\mathbf k)/2$ mamy
\[
 X=\tfrac{1+i}{2}(x+y),\qquad
 Y=\tfrac{-1+i}{2}(x-y).
\]
Stąd $XY=-(x^2-y^2)/2$,
$X^4=-(x+y)^4/4$, $Y^4=-(x-y)^4/4$.
Różnica czwartych potęg daje
$f(X,Y)=xy(x^4-y^4)$. Dla $g$ korzystamy z
\[
 16g(X,Y)=(x+y)^8+14(x^2-y^2)^4+(x-y)^8
          =16g(x,y).
\]
Dla $t$ mamy
\[
 64t(X,Y)=-(x+y)^{12}-(x-y)^{12}
 +33(x^2-y^2)^4\big((x+y)^4+(x-y)^4\big)
 =64t(x,y).
\]
Ostatnie równości wynikają z rozwinięcia dwumianowego.
Podstawienie przez $M(r)$ jest diagonalne z wpisami
$e^{\pi i/4},e^{-\pi i/4}$, więc monom $x^ay^b$
jest mnożony przez $e^{(a-b)\pi i/4}$. Daje to
wskazane znaki. Sprawdzanie przez podstawienie $M(q)$
zamiast $M(q)^{-1}$ jest równoważne niezmienniczości
względem $q$.

Na koniec połóżmy $A=x^4,B=y^4$.
Wówczas $f^4=AB(A-B)^4$,
$g=A^2+14AB+B^2$, $t=(A+B)(A^2-34AB+B^2)$.
Rozwinięcie obu stron daje
$g^3-t^2=108AB(A-B)^4=108f^4$.
\end{proof}

\begin{twierdzenie}[Prezentacje dla dwóch pierwszych grup]
\label{aa:tw:exceptional-e6-e7}
W oznaczeniach Lematu~\ref{aa:lem:exceptional-small-forms}
zachodzą izomorfizmy
\[
\begin{aligned}
 \mathbb C[x,y]^{2T}
 &\simeq\mathbb C[U,V,W]/(W^2-V^3+108U^4),\\
 \mathbb C[x,y]^{2O}
 &\simeq\mathbb C[U,V,W]/(W^2-UV^3+108U^3),
\end{aligned}
\]
gdzie generatory przechodzą odpowiednio na $(f,g,t)$
oraz $(f^2,g,ft)$.
\end{twierdzenie}
\begin{proof}
Mamy $f_x(1,0)=0$, $f_y(1,0)=1$,
$g_x(1,0)=8$, $g_y(1,0)=0$, stąd
$J(f,g)(1,0)=-8$. Zatem $f,g$ są algebraicznie
niezależne na mocy Lematu~\ref{aa:lem:exceptional-jacobian}.
Wielomian $V^3-108U^4$ ma nieparzysty stopień względem
$V$ i nie jest kwadratem. Stopnie $f,g,t$ to $6,8,12$.
Z Twierdzenia~\ref{aa:tw:exceptional-hilbert} i Lematu
\ref{aa:lem:exceptional-presentation} wynika pierwsza
prezentacja.

Grupa $2O$ jest generowana przez $h,r$, więc
$f^2,g,ft$ są jej niezmiennikami. Ich stopnie
wynoszą $12,8,18$. Niezależność $f^2,g$ wynika
z niezależności $f,g$ (podstawienie $U\mapsto U^2$
jest iniektywne na pierścieniu wielomianów).
Ponadto $(ft)^2=f^2(g^3-108f^4)$, a
$UV^3-108U^3$ ma stopień trzy względem $V$.
Stosujemy ponownie to samo kryterium prezentacji,
teraz z szeregiem Hilberta grupy $2O$.
\end{proof}

\section{Grupa dwudziestościenna}

Dla modelu $2I$ z rozdziału 52 potrzebny jest
inny niezmiennik stopnia dwanaście.
Oznaczmy $\phi=(1+\sqrt5)/2$ i połóżmy
\[
 F=\sqrt5\,t-22f^2,\quad
 H=F_{xx}F_{yy}-F_{xy}^2,\quad K=J(F,H).
\]
Ostatnie dwa wielomiany mają stopnie odpowiednio
$20$ i $30$, o ile nie są zerowe.

\begin{lemat}[Niezmienniczość hesjanu i jakobianu]
\label{aa:lem:exceptional-covariants}
Niech $G\le\operatorname{SL}_2(\mathbb C)$,
a $P,Q\in\mathbb C[x,y]^G$. Wówczas
$\det D^2P$ oraz $J(P,Q)$ są niezmiennikami $G$.
\end{lemat}
\begin{proof}
Dla $L\in G$ równość $P(Lv)=P(v)$ i reguła
łańcuchowa dają
$D^2P(v)=L^{\mathsf T}D^2P(Lv)L$.
Po wzięciu wyznaczników i użyciu $\det L=1$
otrzymujemy niezmienniczość hesjanu. Analogicznie
$D(P,Q)(v)=D(P,Q)(Lv)L$, więc wyznacznik
macierzy pochodnych pierwszego rzędu też jest
niezmiennikiem.
\end{proof}

\begin{lemat}[Wielomian stopnia dwanaście]
\label{aa:lem:exceptional-icosahedral-form}
W modelu $2I$ z Przykładu~\ref{aa:prz:binary-quaternion-generators}
wielomian $F$ jest niezmiennikiem.
\end{lemat}
\begin{proof}
Niezmienniczość względem $h,\mathbf i$ wynika
z Lematu~\ref{aa:lem:exceptional-small-forms}.
Dla trzeciego generatora
$u=\phi/2+(\phi-1)\mathbf j/2+\mathbf k/2$
podstawienie przez $M(u)$ daje
\[
 X=\frac{\phi x+(\phi-1+i)y}{2},\quad
 Y=\frac{(1-\phi+i)x+\phi y}{2}.
\]
Podajemy dokładną procedurę sprawdzenia współczynników,
aby nie zakładać niezmienniczości z samego stopnia.
Niech $A,B,C,D$ będą kolejno współczynnikami
$x,y$ w $2X$ i $2Y$. Dla
$F=\sum_{k=0}^{12}a_kx^{12-k}y^k$ współczynnik
$x^{12-l}y^l$ w $2^{12}F(X,Y)$ wynosi
\[
 \sum_{k=0}^{12}a_k
 \sum_{\substack{j+v=l\\0\le j\le12-k\\0\le v\le k}}
 \binom{12-k}{j}\binom{k}{v}
 A^{12-k-j}B^jC^{k-v}D^v.
\]
Używając $\phi^2=\phi+1$, $i^2=-1$,
$\sqrt5=2\phi-1$, po podzieleniu przez $2^{12}$
otrzymujemy dla parzystych $l$ tabelę
\[
\begin{array}{c|rrrrrrr}
l&0&2&4&6&8&10&12\\ \hline
a_l&\sqrt5&-22&-33\sqrt5&44&-33\sqrt5&-22&\sqrt5
\end{array}
\]
oraz zero dla nieparzystych $l$.
Jest to dokładnie ciąg współczynników $F$.
Zatem $F(X,Y)=F(x,y)$, a wskazane trzy
kwaterniony generują $2I$.
\end{proof}

\begin{twierdzenie}[Prezentacja dla grupy dwudziestościennej]
\label{aa:tw:exceptional-e8}
Niech $F,H,K$ będą wielomianami określonymi
powyżej dla modelu $2I$ z
Przykładu~\ref{aa:prz:binary-quaternion-generators}.
Wówczas
\[
 \mathbb C[x,y]^{2I}\simeq
 \frac{\mathbb C[U,V,W]}
 { (W^2+129534197760U^5+\frac{400}{121}V^3) },
 \qquad (U,V,W)\longmapsto(F,H,K).
\]
\end{twierdzenie}
\begin{proof}
Lematy~\ref{aa:lem:exceptional-covariants}
i~\ref{aa:lem:exceptional-icosahedral-form}
zapewniają niezmienniczość trzech wielomianów.
Relację wyprowadzimy z wymiaru części stopnia $60$
oraz dwóch współczynników. Z definicji pochodnych mamy
\[
\begin{array}{c|rr}
 & [x^d] & [x^{d-2}y^2] \\ \hline
 F,\ d=12 & \sqrt5 & -22\\
 H,\ d=20 & -5808\sqrt5 & -367840
\end{array}
\]
oraz $[x^{30}]K=0$,
$[x^{29}y]K=-13939200\sqrt5$.
Dokładniej, dla pierwszych trzech parzystych
współczynników $F$, oznaczonych $a_0,a_2,a_4$,
mamy
\[
 [x^{20}]H=264a_0a_2,\qquad
 [x^{18}y^2]H=1584a_0a_4-220a_2^2.
\]
Wynika to z wymnożenia pochodnych drugiego rzędu.
Dla $a_0=\sqrt5$, $a_2=-22$, $a_4=-33\sqrt5$
otrzymujemy liczby z tabeli.
Współczynnik przy $x^{29}y$ w $K$ wynosi
$12\sqrt5\cdot2(-367840)-2(-22)\cdot20(-5808\sqrt5)$.
Zatem $K=J(F,H)\ne0$ i $F,H$ są algebraicznie niezależne
na mocy Lematu~\ref{aa:lem:exceptional-jacobian}.

Z szeregu Hilberta $2I$ wynika, że część stopnia $60$
ma wymiar dwa: jedyne rozwiązania
$12i+20j+30e=60$, $i,j\ge0$, $e=0,1$,
to $(5,0,0)$ i $(0,3,0)$.
Wielomiany $F^5,H^3$ są liniowo niezależne
i stanowią więc jej bazę. Istnieją $\alpha,\beta$
takie, że $K^2=\alpha F^5+\beta H^3$.
Porównanie współczynników $x^{60}$ oraz $x^{58}y^2$
daje układ
\[
\begin{aligned}
0&=\alpha(\sqrt5)^5+\beta(-5808\sqrt5)^3,\\
5(13939200)^2
 &=5\alpha(\sqrt5)^4(-22)
   +3\beta(-5808\sqrt5)^2(-367840).
\end{aligned}
\]
Rozwiązanie to $\alpha=-129534197760$, $\beta=-400/121$
(oba niezerowe). Wielomian $\alpha U^5+\beta V^3$
ma nieparzysty stopień względem $V$, więc nie jest kwadratem.
Lemat~\ref{aa:lem:exceptional-presentation} wraz z szeregiem Hilberta
$2I$ kończy dowód.
\end{proof}

\section{Postacie równań typu E}

\begin{wniosek}[Trzy równania powierzchni]
\label{aa:wn:exceptional-ade}
Przez skalowanie generatorów niezerowymi liczbami
zespolonymi prezentacje z
Twierdzeń~\ref{aa:tw:exceptional-e6-e7}
i~\ref{aa:tw:exceptional-e8} można sprowadzić do
\[
\begin{array}{c|l}
 2T & X^2+Y^3+Z^4=0\\
 2O & X^2+Y^3+YZ^3=0\\
 2I & X^2+Y^3+Z^5=0 .
\end{array}
\]
Są to standardowe postacie równań osobliwości
Kleina typu $E_6,E_7,E_8$.
\end{wniosek}
\begin{proof}
Dla $2T,2I$ współczynniki przy trzech potęgach
są niezerowe i mogą być unormowane do jedynki
przez niezależny wybór odpowiednich pierwiastków.
Dla $2O$ w równaniu $W^2-UV^3+108U^3=0$
przyjmujemy $X=W$, $Y=aU$, $Z=bV$,
gdzie $a^3=108$, $ab^3=-1$.
Daje to drugie równanie. W każdym przypadku
gradient wielomianu definiującego znika tylko
w początku. Dla drugiego równania warunki to
$X=0$, $3Y^2+Z^3=0$ i $3YZ^2=0$.
Ostatnia równość daje $Y=0$ lub $Z=0$,
a poprzednia wymusza oba zera. Dla pierwszego
i trzeciego równania wynika to bezpośrednio
z pochodnych potęg.
\end{proof}

Klasyfikacja i rezolucje osobliwości oraz odpowiedniość McKaya
wymagają dalszych konstrukcji geometrycznych i reprezentacyjnych.

% AA-55
\chapter{Grafy McKaya i afiniczne diagramy Dynkina}
\label{aa:ch:mckay}

W rozdziałach~\ref{aa:ch:binary-polyhedral}--\ref{aa:ch:exceptional-invariants}
sklasyfikowaliśmy skończone podgrupy $\operatorname{SL}_2(\mathbb C)$
i wyznaczyliśmy ich pierścienie niezmienników. Teraz
odczytamy tę samą listę typów z teorii reprezentacji.
Nie będziemy zakładać klasyfikacji diagramów Dynkina.

\section{Macierz McKaya}

\begin{definicja}[Graf McKaya]
\label{aa:def:mckay}
Niech $G\le\operatorname{SL}_2(\mathbb C)$ będzie skończona,
$V=\mathbb C^2$ jej naturalną reprezentacją, a
$V_0,\ldots,V_s$ pełną listą parami nierównoważnych
zespolonych reprezentacji nieprzywiedlnych, przy czym
$V_0=\mathbb C$ jest reprezentacją trywialną.
Macierzą McKaya nazywamy macierz $A=(a_{ij})$ określoną przez
\[
 V\otimes V_i\simeq\bigoplus_{j=0}^s V_j^{\oplus a_{ij}}.
\]
Graf McKaya ma wierzchołki $0,\ldots,s$, a liczba
skierowanych krawędzi od $i$ do $j$ wynosi $a_{ij}$.
Gdy $A$ jest symetryczna, dla $i\ne j$ będziemy
rysować $a_{ij}$ krawędzi nieskierowanych; wpis
$a_{ii}$ zachowamy jako krotność pętli.
\end{definicja}

\begin{twierdzenie}[Własności macierzy McKaya]
\label{aa:tw:mckay-properties}
Niech $G,V,V_i,A$ będą jak w
Definicji~\ref{aa:def:mckay}, a $d_i=\dim V_i$.
Wówczas $A$ jest symetryczna, graf jest spójny,
$Ad=2d$ dla $d=(d_i)_{i=0}^s$ oraz
$\sum_i d_i^2=|G|$. Jeżeli $-I\in G$, graf jest dwudzielny
i nie ma pętli.
\end{twierdzenie}
\begin{proof}
Po sprzężeniu możemy przyjąć $G\le\operatorname{SU}(2)$
na mocy Twierdzenia~\ref{aa:tw:binary-unitary}.
Wartości własne $g\in G$ mają postać $\lambda,\lambda^{-1}$
dla $|\lambda|=1$. Zatem charakter $\chi_V(g)$ jest rzeczywisty,
a wartość $2$ przyjmuje tylko dla $g=I$: wówczas
$\operatorname{Re}\lambda=1$ wymusza $\lambda=1$,
a unitarna macierz o obu wartościach własnych równych $1$
jest identycznością.
Z formuły na krotności
(Twierdzenie~\ref{aa:tw:multiplicities}) mamy
\[
 a_{ij}=\frac1{|G|}\sum_{g\in G}
   \chi_V(g)\chi_i(g)\overline{\chi_j(g)}.
\]
Ponieważ $\chi_V$ jest rzeczywisty, $a_{ji}=\overline{a_{ij}}=a_{ij}$.
Porównanie wymiarów w definicji daje $2d_i=\sum_j a_{ij}d_j$.
Rozkład reprezentacji regularnej z
Twierdzenia~\ref{aa:tw:regular-character} daje
$|G|=\sum_i d_i^2$.

Dla spójności załóżmy, że $j$ nie należy do składowej
wierzchołka $0$. Indukcja według $k$ pokazuje, że
każdy składnik nieprzywiedlny $V^{\otimes k}$
jest w tej składowej: dla $k=0$ jest to reprezentacja
trywialna, a iloczyn tensorowy z $V$ przechodzi
wyłącznie do sąsiadów. W konsekwencji dla każdego $k\ge0$
\[
 0=\frac1{|G|}\sum_{g\in G}\chi_V(g)^k\overline{\chi_j(g)}.
\]
Niech $T=\{\chi_V(g):g\in G\}$ i
$p(t)=\prod_{u\in T\setminus\{2\}}(t-u)/(2-u)$.
Wielomian ten przyjmuje wartość $1$ w $2$ i $0$ w pozostałych
punktach $T$. Biorąc kombinację powyższych równości
ze współczynnikami $p$, otrzymujemy
$0=\overline{\chi_j(I)}/|G|=d_j/|G|$, sprzeczność.

Gdy $-I\in G$, jego działanie na $V_i$ jest skalarem
$\varepsilon_i\in\{1,-1\}$ na mocy lematu Schura
i równości $(-I)^2=I$. Na $V\otimes V_i$ działa jako
$-\varepsilon_i$. Jeżeli $a_{ij}>0$, to
$\varepsilon_j=-\varepsilon_i$. Podział wierzchołków
według znaku jest więc podziałem dwudzielnym.
\end{proof}

\begin{lemat}[Podgraf o dodatnim wektorze wymiarów]
\label{aa:lem:mckay-subgraph}
Niech $A$ będzie symetryczną macierzą stopnia $N\ge1$
o nieujemnych współczynnikach całkowitych, której
graf jest spójny.
Załóżmy $Ad=2d$ dla wektora $d>0$ (wszystkie współrzędne
dodatnie). Niech $B$ będzie symetryczną macierzą
podgrafu na niepustym podzbiorze $S$ wierzchołków,
czyli $0\le B_{ij}\le A_{ij}$ dla $i,j\in S$.
Jeśli $Bu=2u$ dla $u>0$, to $S$ jest całym zbiorem
wierzchołków i $B=A$.
\end{lemat}
\begin{proof}
Dla ograniczenia $d_S$ mamy dla każdego $i\in S$
\[
 (2d_S-Bd_S)_i=
 \sum_{j\in S}(A_{ij}-B_{ij})d_j+
 \sum_{j\notin S}A_{ij}d_j\ge0.
\]
Nierówność wynika z założeń o współczynnikach i o $d$. Ponadto
$u^{\mathsf T}(2d_S-Bd_S)=(2u-Bu)^{\mathsf T}d_S=0$.
Każda współrzędna różnicy jest nieujemna, a każda
współrzędna $u$ dodatnia, więc $Bd_S=2d_S$.
Ponieważ wszystkie $d_i>0$, żadna pominięta krawędź
ani zmniejszenie krotności nie jest możliwe.
Zatem z $S$ nie wychodzi żadna krawędź, a spójność daje
$S$ równy całemu zbiorowi.
\end{proof}

\section{Klasyfikacja grafów bez pętli}

\begin{twierdzenie}[Klasyfikacja kombinatoryczna]
\label{aa:tw:mckay-graphs}
Niech $A$ będzie symetryczną macierzą stopnia $N\ge1$
o nieujemnych współczynnikach całkowitych, z zerami na przekątnej,
o spójnym grafie i niech $Ad=2d$ dla $d>0$.
Wówczas graf jest jednym z następujących:
\begin{enumerate}
\item dwa wierzchołki połączone dwiema krawędziami
      ($\widetilde A_1$) lub cykl o $m\ge3$ wierzchołkach
      ($\widetilde A_{m-1}$);
\item gwiazda z czterema liśćmi ($\widetilde D_4$)
      lub ścieżka z $k\ge1$ krawędziami, do której
      każdego końca dołączono po dwa liście
      ($\widetilde D_{k+4}$);
\item drzewo o jednym wierzchołku stopnia trzy,
      z ramionami długości odpowiednio $(2,2,2)$,
      $(1,3,3)$ albo $(1,2,5)$
      ($\widetilde E_6,\widetilde E_7,\widetilde E_8$).
\end{enumerate}
Długość ramienia oznacza liczbę krawędzi od wierzchołka
stopnia trzy do liścia. Każdy wymieniony graf ma
dodatni wektor $d$ spełniający $Ad=2d$, jednoznaczny
z dokładnością do dodatniego mnożnika.
\end{twierdzenie}
\begin{proof}
Jeśli istnieje krawędź krotności co najmniej dwa,
podgraf dwóch wierzchołków z dwiema krawędziami
ma dodatni wektor $(1,1)$. Lemat~\ref{aa:lem:mckay-subgraph}
wymusza równość z tym podgrafem.
Odtąd graf jest prosty. Nie może mieć jednego
wierzchołka: wtedy $A=(0)$ i $Ad=0\ne2d$.
Jeśli zawiera cykl, wybieramy
jego krawędzie i wierzchołki. Wektor samych jedynek
spełnia na cyklu wymagane równanie, więc ten sam
lemat daje równość grafu z cyklem.

Pozostał przypadek drzewa. Jeśli jakiś wierzchołek ma
stopień co najmniej cztery, wybieramy go i czterech
sąsiadów z krawędziami do niego. Wektor o wartości
$2$ w centrum i $1$ w liściach spełnia równanie.
Lemat wyklucza wszystkie dodatkowe krawędzie i wierzchołki.
Jeżeli są przynajmniej dwa wierzchołki stopnia trzy,
wybieramy parę o najmniejszej odległości. Wewnętrzne
wierzchołki ich jedynej ścieżki mają stopień dwa,
gdyż w przeciwnym razie istniałaby bliższa para.
Do każdego końca ścieżki dołączamy jego dwóch sąsiadów
spoza ścieżki; są oni parami różni, ponieważ graf
nie ma cykli. Na tym podgrafie nadajemy wartości
$2$ wszystkim wierzchołkom ścieżki oraz $1$ czterem
nowym liściom. Wewnętrzny wierzchołek ścieżki ma
sumę sąsiadów $2+2=4$, jej koniec $2+1+1=4$,
a liść ma sąsiada o wartości $2$.
Z lematu wynika równość z całym grafem.

Drzewo bez rozgałęzienia byłoby ścieżką.
Po oznaczeniu jej wartości przez $d_1,\ldots,d_l$
i dodaniu $d_0=d_{l+1}=0$, równania
$2d_j=d_{j-1}+d_{j+1}$ wymuszają stałe różnice
$d_j-d_{j-1}$. Z obu zer na końcach wynika, że różnica
jest zero, co przeczy dodatniości.
Pozostaje dokładnie jeden wierzchołek stopnia trzy.
Po jego usunięciu otrzymujemy trzy składowe: sąsiedzi
nie mogą należeć do tej samej składowej, bo powstałby
cykl. Każda składowa jest ścieżką (być może
jednowierzchołkową), a dawny sąsiad wierzchołka
centralnego jest jej końcem, bo utracił jedną krawędź.
Niech wierzchołek centralny ma wartość $h>0$,
a ramiona mają długości $l_1,l_2,l_3\ge1$. Na ramieniu długości $l$ oznaczamy
wartości przez $b_0=h,b_1,\ldots,b_l,b_{l+1}=0$.
Równanie $2b_j=b_{j-1}+b_{j+1}$ wymusza
$b_j=h(l+1-j)/(l+1)$. W centrum otrzymujemy
\[
 2h=\sum_{i=1}^3 h\frac{l_i}{l_i+1},
 \qquad
 \sum_{i=1}^3\frac1{l_i+1}=1.
\]
Uporządkujmy $a=l_1+1\le b=l_2+1\le c=l_3+1$.
Gdy $a\ge4$, suma nie przekracza $3/4$.
Dla $a=3$ wszystkie trzy składniki nie przekraczają
$1/3$, więc $a=b=c=3$. Dla $a=2$ mamy
$1/b+1/c=1/2$. Wartość $b=2$ jest niemożliwa,
a $b\ge5$ daje sumę co najwyżej $2/5$.
Jeśli $b=3$, to $c=6$; jeśli $b=4$, to $c=4$.
Otrzymujemy dokładnie podane trzy zestawy ramion.
Wzór na $b_j$ i równanie centralne dowodzą też istnienia
dodatniego wektora dla każdego z nich. Dla pozostałych
grafów istnienie pokazują skonstruowane wcześniej
wektory: same jedynki dla typu $A$, dwójki na ścieżce
i jedynki w liściach dla typu $D$ (w gwieździe
dwójka tylko w centrum).

Na koniec, jeśli $u,v>0$ spełniają $Au=2u$, $Av=2v$,
to dla $t=\min_i u_i/v_i$ wektor $w=u-tv$ jest
nieujemny, ma współrzędną zerową i spełnia $Aw=2w$.
Jeśli $w_i=0$, to $\sum_j a_{ij}w_j=0$, więc
$w_j=0$ dla wszystkich sąsiadów $i$. Spójność
przenosi zero na wszystkie wierzchołki, zatem $u=tv$.
\end{proof}

\section{Grupy cykliczne i dwucykliczne}

\begin{przyklad}[Grupa cykliczna]
\label{aa:prz:mckay-cyclic}
Niech $G=\langle\operatorname{diag}(\zeta,\zeta^{-1})\rangle$
dla pierwotnego pierwiastka $\zeta$ stopnia $m\ge1$.
Jej reprezentacje nieprzywiedlne to charaktery $\chi_j$,
$0\le j<m$, określone przez $\chi_j(g)=\zeta^j$ dla
podanego generatora $g$. Istotnie macierz generatora
jest diagonalizowalna (jej wielomian minimalny dzieli
$X^m-1$), a każda prosta własna jest podreprezentacją.
Mamy $V\simeq\chi_1\oplus\chi_{-1}$, więc
$V\otimes\chi_j\simeq\chi_{j+1}\oplus\chi_{j-1}$,
gdzie indeksy liczymy modulo $m$.
Dla $m\ge3$ otrzymujemy cykl $\widetilde A_{m-1}$,
dla $m=2$ dwie krawędzie, a dla $m=1$ jeden wierzchołek
z wpisem diagonalnym $2$. Ten ostatni przypadek
rozpatrujemy oddzielnie od klasyfikacji bez pętli.
\end{przyklad}

\begin{lemat}[Liczba jednowymiarowych reprezentacji]
\label{aa:lem:mckay-abelianizations}
Grupa $\operatorname{Dic}_n$ dla $n\ge2$ ma dokładnie
cztery nierównoważne reprezentacje zespolone wymiaru jeden.
Grupy $2T,2O,2I$ mają odpowiednio trzy, dwie i jedną
taką reprezentację.
\end{lemat}
\begin{proof}
Z prezentacji $\operatorname{Dic}_n$
(Twierdzenie~\ref{aa:tw:binary-dicyclic}) charakter
jednowymiarowy określają obrazy $\alpha,\beta$ generatorów
$a,j$. Relacje przyjmują postać
$\alpha^{2n}=1$, $\beta^2=\alpha^n$, $\alpha=\alpha^{-1}$.
Zatem $\alpha=\pm1$, a dla każdego wyboru są dokładnie
dwa różne niezerowe pierwiastki równania dla $\beta$.
Każda z czterech par spełnia wszystkie relacje.

W modelach kwaternionowych grup wyjątkowych
(Przykład~\ref{aa:prz:binary-quaternion-generators})
występują $\mathbf i,\mathbf j$. Dla $2O$ mamy
$r^2=\mathbf i$, a sprzężenie przez $h$ cyklicznie
permutuje jednostki urojone; w $2I$ korzystamy
z $\mathbf i,h$ w ten sam sposób.
Komutator $\mathbf i\mathbf j\mathbf i^{-1}\mathbf j^{-1}=-1$,
więc każdy homomorfizm do $\mathbb C^\times$ zabija $-1$.
Charaktery jednowymiarowe grup wyjątkowych są zatem
w bijekcji z charakterami ich ilorazów $A_4,S_4,A_5$.

W $A_4$ zbiór $K$ złożony z identyczności i trzech
podwójnych transpozycji jest podgrupą: każdy element
niejednostkowy ma kwadrat $1$, a iloczyn dwóch różnych
podwójnych transpozycji jest trzecią.
Sprzężenie przez dowolną permutację zachowuje typ
cykliczny, zatem $K$ jest normalna w $A_4$.
Iloraz ma rząd trzy, więc jest cykliczny. Jeśli $s$ jest trójcyklem, sprzężenie
przez $s$ cyklicznie permutuje trzy niejednostkowe
elementy $K$. Dla $v\ne1$ w $K$ komutator
$svs^{-1}v^{-1}$ jest niejednostkowym elementem $K$,
a jego sprzężenia przez $s$ dają pozostałe dwa.
Zatem podgrupa komutatorów zawiera $K$ i jest w niej
zawarta, ponieważ iloraz jest przemienny.
Stąd abelianizacja jest cykliczna rzędu trzy.

W $S_4$ podgrupa komutatorów jest zawarta w $A_4$
przez homomorfizm znaku. Komutator dwóch transpozycji
o jednym wspólnym punkcie jest trójcyklem; sprzężenia
dają wszystkie trójcykle. Generują one $A_4$:
każda permutacja parzysta jest iloczynem par transpozycji,
a para rozłączna rozkłada się na dwa trójcykle,
na przykład $(ab)(cd)=(acb)(acd)$ (składamy od prawej).
Para o wspólnym punkcie jest trójcyklem lub identycznością.
Zatem podgrupa komutatorów jest $A_4$.

W $A_5$ weźmy $u=(12)(34)$ i $v=(12)(35)$.
Oba są inwolucjami. Ponieważ wspólna transpozycja
$(12)$ komutuje z $(34)$ i $(35)$, mamy
$uv=(34)(35)=(354)$ i $[u,v]=(uv)^2=(345)$.
Ten trójcykl należy do podgrupy komutatorów.
Każde dwa trójcykle są sprzężone w $A_5$:
sprzęgający element $S_5$, jeżeli nieparzysty,
można pomnożyć z prawej przez transpozycję dwóch
punktów spoza nośnika pierwszego trójcyklu.
Podgrupa komutatorów zawiera więc wszystkie trójcykle,
które generują $A_5$ tym samym rozkładem par transpozycji.
Abelianizacje trzech ilorazów to $C_3,C_2,\{1\}$,
co daje żądane liczby charakterów.
\end{proof}

\section{Odpowiedniość grafów}

\begin{twierdzenie}[Grafy McKaya skończonych podgrup specjalnej grupy liniowej]
\label{aa:tw:mckay-ade}
Dla skończonej podgrupy $G\le\operatorname{SL}_2(\mathbb C)$
jej graf McKaya ma następujący typ:
\[
\begin{array}{c|c}
G&\text{graf}\\ \hline
C_m,\ m\ge2&\widetilde A_{m-1}\\
\operatorname{Dic}_n,\ n\ge2&\widetilde D_{n+2}\\
2T&\widetilde E_6\\
2O&\widetilde E_7\\
2I&\widetilde E_8
\end{array}
\]
Dla grupy trywialnej obowiązuje opis z
Przykładu~\ref{aa:prz:mckay-cyclic}.
Liczby przy wierzchołkach na poniższych rysunkach
są wymiarami reprezentacji nieprzywiedlnych.
Wierzchołek trywialny jest jednym z liści o wymiarze
jeden w typach $D,E$.
\end{twierdzenie}
\begin{proof}
Klasyfikacja grup z Twierdzenia~\ref{aa:tw:binary-classification}
i Przykład~\ref{aa:prz:mckay-cyclic} redukują dowód
do czterech rodzin niecyklicznych. Zawierają one $-I$,
więc stosujemy Twierdzenia~\ref{aa:tw:mckay-properties}
i~\ref{aa:tw:mckay-graphs}.
W typie $A$ dodatni wektor jest stały, a ponieważ
$d_0=1$, wszystkie wymiary wynoszą jeden.
Wówczas wszystkie reprezentacje nieprzywiedlne są
jednowymiarowe i reprezentacja regularna jest sumą
charakterów jednowymiarowych. Jej macierze są więc
jednocześnie diagonalne w odpowiedniej bazie i komutują.
Reprezentacja regularna jest wierna, ponieważ
działanie $g$ na wektorze bazowym $[1]$ daje $[g]$.
Stąd z równości macierzy odpowiadających $gh$ i $hg$
wynika $gh=hg$ dla wszystkich $g,h\in G$. Wyklucza to typ $A$ dla tych grup.

W typie $D$ wektor wymiarów ma cztery jedynki
w liściach oraz dwójki na pozostałych wierzchołkach,
gdyż taki wektor spełnia równanie, a dodatni mnożnik
musi być równy jeden (najmniejszy wymiar to jeden).
W typie $\widetilde D_{k+4}$, z $k\ge0$ (dla $k=0$
gwiazda), suma kwadratów wynosi $4+4(k+1)=4(k+2)$.
W typach $E$ wzór z dowodu klasyfikacji daje
wektory o wartościach centralnych $3,4,6$; ich
najmniejsze wartości są jedynkami. Dla $E_6$
każde ramię ma wartości $2,1$; dla $E_7$
ramiona mają wartości $(2)$, $(3,2,1)$, $(3,2,1)$;
dla $E_8$ wartości to $(3)$, $(4,2)$, $(5,4,3,2,1)$.
Są to dokładnie wymiary, na mocy jednoznaczności
i istnienia wymiaru jeden. Liczba jedynek wynosi
odpowiednio trzy, dwa, jeden.
Lemat~\ref{aa:lem:mckay-abelianizations}
identyfikuje typ każdej grupy: $\operatorname{Dic}_n$
ma cztery jedynki, a z $4(k+2)=4n$ dostajemy
$k=n-2$. Trzy pozostałe grupy mają odpowiednio trzy,
dwie i jedną jedynkę, co identyfikuje diagramy $E$.
Wierzchołki o wymiarze jeden w każdym z tych diagramów
są wzajemnie równoważne przez symetrie grafu:
w $D$ zamieniamy liście przy końcach ścieżki
i ewentualnie odwracamy ścieżkę (w $D_4$ dowolnie
permutujemy cztery liście), w $E_6$ permutujemy
ramiona, a w $E_7$ zamieniamy dwa dłuższe ramiona.
W $E_8$ jest tylko jedna jedynka.
Możemy zatem wybrać oznaczenia grafu tak,
aby określony liść o wymiarze jeden był trywialny.
\end{proof}

\begin{center}
\begin{tikzpicture}[v/.style={circle,draw,minimum size=5.5mm,inner sep=0pt,font=\small}]
 \node[v] (a) at (0,0) {1}; \node[v] (b) at (1,0) {1};
 \node[v] (c) at (1,1) {1}; \node[v] (d) at (0,1) {1};
 \draw (a)--(b)--(c)--(d)--(a);
 \node at (.5,-.5) {$\widetilde A_3$};
 \node[v] (e) at (3,0.5) {2};
 \foreach \n/\x/\y in {f/2/1.1,g/2/-.1,h/4/1.1,i/4/-.1}
  {\node[v] (\n) at (\x,\y) {1}; \draw (e)--(\n);}
 \node at (3,-.5) {$\widetilde D_4$};
 \node[v] (j) at (6,0.5) {2}; \node[v] (k) at (7,0.5) {2};
 \draw (j)--(k);
 \foreach \n/\x/\y/\p in {l/5/1.1/j,m/5/-.1/j,n/8/1.1/k,o/8/-.1/k}
  {\node[v] (\n) at (\x,\y) {1}; \draw (\p)--(\n);}
 \node at (6.5,-.5) {$\widetilde D_5$};
\end{tikzpicture}

\small Diagramy cykliczny i dwucykliczne; w większym typie $D$
wydłuża się ścieżka wierzchołków wymiaru dwa.
\end{center}

\begin{center}
\begin{tikzpicture}[v/.style={circle,draw,minimum size=5.5mm,inner sep=0pt,font=\small}]
 \node[v] (c) at (2,0) {3};
 \node[v] (l1) at (1,0) {2}; \node[v] (l2) at (0,0) {1};
 \node[v] (r1) at (3,0) {2}; \node[v] (r2) at (4,0) {1};
 \node[v] (u1) at (2,.8) {2}; \node[v] (u2) at (2,1.6) {1};
 \draw (l2)--(l1)--(c)--(r1)--(r2); \draw (c)--(u1)--(u2);
 \node at (2,-.5) {$\widetilde E_6$};
 \node[v] (d) at (8.2,0) {4};
 \foreach \n/\x/\w in {a/5.8/1,b/6.6/2,e/7.4/3,f/9/3,g/9.8/2,h/10.6/1}
  {\node[v] (\n) at (\x,0) {\w};}
 \node[v] (u) at (8.2,.8) {2};
 \draw (a)--(b)--(e)--(d)--(f)--(g)--(h); \draw (d)--(u);
 \node at (8.2,-.5) {$\widetilde E_7$};
\end{tikzpicture}

\begin{tikzpicture}[v/.style={circle,draw,minimum size=5.5mm,inner sep=0pt,font=\small}]
 \foreach \n/\x/\w in {a/0/2,b/.9/4,c/1.8/6,d/2.7/5,e/3.6/4,f/4.5/3,g/5.4/2,h/6.3/1}
  {\node[v] (\n) at (\x,0) {\w};}
 \node[v] (u) at (1.8,.8) {3};
 \draw (a)--(b)--(c)--(d)--(e)--(f)--(g)--(h); \draw (c)--(u);
 \node at (3.15,-.5) {$\widetilde E_8$};
\end{tikzpicture}

\small Diagramy wyjątkowe z wymiarami reprezentacji.
\end{center}

\begin{przyklad}[Odczytywanie iloczynów tensorowych]
Niech $G=2I$. Oznaczmy reprezentacje wzdłuż najdłuższego
ramienia od liścia przez $L_1,L_2,L_3,L_4,L_5,L_6$,
gdzie $\dim L_i=i$. Pozostałe reprezentacje oznaczmy
przez $M_3,M_4,M_2$, zgodnie z wymiarami i położeniem:
$M_3$ jest krótkim ramieniem, $M_4,M_2$ tworzą
ramię długości dwa. Mamy $L_1=\mathbb C$, $L_2=V$.
Z diagramu otrzymujemy
\[
\begin{aligned}
 V\otimes L_1&\simeq L_2,\\
 V\otimes L_i&\simeq L_{i-1}\oplus L_{i+1}\quad(2\le i\le5),\\
 V\otimes L_6&\simeq L_5\oplus M_3\oplus M_4,\\
 V\otimes M_3&\simeq L_6,\qquad
 V\otimes M_4\simeq L_6\oplus M_2,\qquad
 V\otimes M_2\simeq M_4.
\end{aligned}
\]
Przykładowo w trzecim wierszu $2\cdot6=5+3+4$.
Reprezentacje o tych samych wymiarach są rozróżnione
przez położenie w grafie, nie przez sam wymiar.
\end{przyklad}

\section{Potęgi symetryczne i pierścienie niezmienników}

\begin{twierdzenie}[Rekurencja McKaya]
\label{aa:tw:mckay-recurrence}
W oznaczeniach Definicji~\ref{aa:def:mckay} niech
$m_k$ będzie kolumną krotności reprezentacji $V_i$
w $\operatorname{Sym}^k V$, $k\ge0$.
Wówczas $m_0=e_0$, $m_1=Ae_0$ oraz
$m_{k+1}=Am_k-m_{k-1}$ dla $k\ge1$, gdzie $e_0$
jest wektorem jednostkowym przy reprezentacji trywialnej.
Szereg Hilberta pierścienia $\mathbb C[x,y]^G$ spełnia
\[
 H_G(z)=e_0^{\mathsf T}(I-zA+z^2 I)^{-1}e_0.
\]
Odwrotność rozumiemy w pierścieniu macierzy nad
$\mathbb C[[z]]$.
\end{twierdzenie}
\begin{proof}
Każdy $g\in G$ jest diagonalizowalny: jego wielomian
minimalny dzieli $X^{|G|}-1$, który nad $\mathbb C$
nie ma pierwiastków wielokrotnych. Wyznacznik
$g$ wynosi jeden, więc wartości własne to
$\lambda,\lambda^{-1}$. Charakter $\operatorname{Sym}^k V$ wynosi
$c_k(g)=\sum_{j=0}^k\lambda^{k-2j}$:
monomy stopnia $k$ w bazie wektorów własnych tworzą
bazę potęgi symetrycznej. Mnożenie przez
$\lambda+\lambda^{-1}$ daje wyrazy
$\lambda^{k+1},\lambda^{-k-1}$ po jednym razie,
a wszystkie pośrednie potęgi po dwa razy. Stąd
$\chi_V c_k=c_{k+1}+c_{k-1}$ dla $k\ge1$.
Z jednoznaczności rozkładu charakterów na bazę
charakterów nieprzywiedlnych
(Twierdzenie~\ref{aa:tw:characters-basis}) i symetrii $A$
wynika rekurencja: krotność $V_j$ w iloczynie
$V\otimes\operatorname{Sym}^k V$ jest równa
$\sum_i a_{ij}(m_k)_i=(Am_k)_j$, a po drugiej
stronie to $(m_{k+1})_j+(m_{k-1})_j$.
Dwa warunki początkowe wynikają
z $\operatorname{Sym}^0 V=\mathbb C$ i
$\operatorname{Sym}^1 V=V$.

Dla formalnego szeregu kolumn $M(z)=\sum_{k\ge0}m_kz^k$
mnożenie przez $I-zA+z^2I$ daje $e_0$, gdyż
współczynnik stały to $m_0$, liniowy to $m_1-Am_0=0$,
a wyższe współczynniki znikają z rekurencji.
Macierz ma wyraz stały $I$, więc ma odwrotność
formalną (rozwijamy szereg geometryczny macierzy
$zA-z^2I$). Niezdegenerowana forma
$(u,v)\mapsto\det(u,v)$ jest zachowywana przez $G$,
zatem $V\to V^*$, $u\mapsto\det(u,\cdot)$ jest
izomorfizmem reprezentacji. Jednorodna część
$\mathbb C[x,y]$ stopnia $k$ jest $\operatorname{Sym}^k V^*$,
a więc jest izomorficzna z $\operatorname{Sym}^k V$.
Wymiar jej przestrzeni wektorów ustalonych przez $G$
jest krotnością reprezentacji trywialnej (pozostałe
nieprzywiedlne reprezentacje nie mają wektorów ustalonych).
Stąd $H_G=e_0^{\mathsf T}M(z)$.
\end{proof}

\begin{przyklad}[Pierwsze niezmienniki dwudziestościenne]
W grafie $\widetilde E_8$ oznaczmy wektor jednostkowy
odpowiadający $L_i$ przez $e_{L_i}$, analogicznie
dla $M_j$. Z poprzedniej rekurencji
$m_0=e_{L_1}$, $m_1=e_{L_2}$, a dla $2\le k\le5$
otrzymujemy $m_k=e_{L_{k+1}}$ (odejmujemy jednego
sąsiada z poprzedniego kroku).
Następnie $m_6=e_{M_3}+e_{M_4}$, ponieważ
$Ae_{L_6}=e_{L_5}+e_{M_3}+e_{M_4}$.
Dla przykładu następne dwa kroki to
$m_7=A(e_{M_3}+e_{M_4})-e_{L_6}
=e_{L_6}+e_{M_2}$ oraz
$m_8=A(e_{L_6}+e_{M_2})-(e_{M_3}+e_{M_4})
=e_{L_5}+e_{M_4}$.
Powtarzając te działania otrzymujemy
\[
\begin{aligned}
 m_7&=e_{L_6}+e_{M_2},&
 m_8&=e_{L_5}+e_{M_4},\\
 m_9&=e_{L_4}+e_{L_6},&
 m_{10}&=e_{L_3}+e_{L_5}+e_{M_3},\\
 m_{11}&=e_{L_2}+e_{L_4}+e_{L_6},&
 m_{12}&=e_{L_1}+e_{L_3}+e_{L_5}+e_{M_4}.
\end{aligned}
\]
W szczególności pierwsze wystąpienie reprezentacji
trywialnej po stopniu zero następuje w stopniu $12$,
zgodnie z Twierdzeniem~\ref{aa:tw:exceptional-hilbert}.
\end{przyklad}

Otrzymana odpowiedniość dotyczy grafów reprezentacji
i typów równań pierścieni niezmienników z rozdziałów 53--54.
Geometryczna odpowiedniość między reprezentacjami
a krzywymi wyjątkowymi rezolucji minimalnej wymaga
osobnego dowodu; nie wyprowadzamy jej z samej
zgodności nazw diagramów.

% AA-56
\chapter{Macierze Cartana, układy korzeni i grupy Weyla}
\label{aa:ch:root-systems}

Grafy McKaya z rozdziału~\ref{aa:ch:mckay} prowadzą do
macierzy $2I-A$. W tym rozdziale zbadamy również macierze
dodatnio określone oraz wyjaśnimy ich związek z geometrią
skończonych układów korzeni. Będziemy pracować w rzeczywistych
przestrzeniach euklidesowych skończonego wymiaru.

\section{Macierze Cartana i grafy McKaya}

\begin{definicja}[Uogólniona macierz Cartana]
\label{aa:def:generalized-cartan}
Uogólnioną macierzą Cartana stopnia $n\ge1$ nazywamy
macierz całkowitą $C=(c_{ij})_{i,j=1}^n$ taką, że
$c_{ii}=2$, dla $i\ne j$ mamy $c_{ij}\le0$, a
$c_{ij}=0$ wtedy i tylko wtedy, gdy $c_{ji}=0$.
Macierz jest symetryzowalna, jeżeli istnieje dodatnia
macierz diagonalna $D$ taka, że $DC$ jest symetryczna.
Mówimy, że jest typu skończonego, jeśli jest
symetryzowalna i $DC$ jest dodatnio określona.
W grafie macierzy łączymy różne wierzchołki $i,j$,
jeżeli $c_{ij}\ne0$. Nierozkładalność oznacza
spójność tego grafu.
\end{definicja}

\begin{stwierdzenie}[Postać kwadratowa macierzy McKaya]
\label{aa:prop:mckay-cartan}
Niech $G$ będzie niecykliczną skończoną podgrupą
$\operatorname{SL}_2(\mathbb C)$ albo nietrywialną grupą
cykliczną w tej grupie. Niech $A$ będzie jej macierzą
McKaya, a $d$ kolumną wymiarów reprezentacji
nieprzywiedlnych. Wówczas $C=2I-A$ jest symetryczną
uogólnioną macierzą Cartana, dodatnio półokreśloną,
z jądrem $\mathbb R d$. Usunięcie dowolnego wiersza
i odpowiadającej mu kolumny daje macierz dodatnio określoną.
\end{stwierdzenie}
\begin{proof}
Z Twierdzeń~\ref{aa:tw:mckay-properties}
i~\ref{aa:tw:mckay-ade} oraz
Przykładu~\ref{aa:prz:mckay-cyclic} wiemy, że $A$
jest symetryczna, ma zera na przekątnej, jej graf jest
spójny oraz $Ad=2d$, gdzie $d_i>0$.
Dla $x\in\mathbb R^n$ mamy
\[
\begin{aligned}
 \sum_{i<j}a_{ij}d_id_j
  \left(\frac{x_i}{d_i}-\frac{x_j}{d_j}\right)^2
 &=\sum_i \frac{x_i^2}{d_i}\sum_j a_{ij}d_j
       -2\sum_{i<j}a_{ij}x_ix_j\\
 &=2\sum_i x_i^2-\sum_{i,j}a_{ij}x_ix_j
 =x^{\mathsf T}Cx.
\end{aligned}
\]
Suma jest nieujemna i znika dokładnie wtedy, gdy
$x_i/d_i=x_j/d_j$ na każdej krawędzi. Spójność
oznacza, że wszystkie te ilorazy są równe. Ponieważ
$Cd=0$, otrzymujemy dokładnie wskazane jądro.
Dla macierzy głównej po usunięciu indeksu $i$ rozważamy
wektory $x$ o współrzędnej $x_i=0$. Jeżeli taki
wektor ma zerową formę kwadratową, to $x=td$,
a $0=td_i$ daje $t=0$. Forma ograniczona jest więc
dodatnio określona.
\end{proof}

Grupa trywialna ma macierz McKaya $(2)$, zatem
$2I-A=(0)$ nie spełnia definicji uogólnionej macierzy
Cartana. Jest to osobny przypadek brzegowy.

\section{Skończone układy korzeni}

\begin{definicja}[Zredukowany krystalograficzny układ korzeni]
\label{aa:def:root-system}
Niech $E$ będzie rzeczywistą przestrzenią euklidesową
dodatniego wymiaru. Skończony zbiór $\Phi\subset E\setminus\{0\}$
nazywamy zredukowanym krystalograficznym układem korzeni,
jeżeli:
\begin{enumerate}
\item $\operatorname{span}_{\mathbb R}\Phi=E$;
\item $\Phi\cap\mathbb R\alpha=\{\alpha,-\alpha\}$
      dla każdego $\alpha\in\Phi$;
\item odbicie $s_\alpha(x)=x-
      \frac{2(x,\alpha)}{(\alpha,\alpha)}\alpha$
      spełnia $s_\alpha(\Phi)=\Phi$ dla każdego $\alpha\in\Phi$;
\item $\frac{2(\beta,\alpha)}{(\alpha,\alpha)}\in\mathbb Z$
      dla wszystkich $\alpha,\beta\in\Phi$.
\end{enumerate}
Grupa Weyla $W(\Phi)$ jest podgrupą grupy ortogonalnej
generowaną przez wszystkie $s_\alpha$ dla $\alpha\in\Phi$.
Układ jest nierozkładalny, jeśli nie można rozdzielić
$\Phi$ na dwa niepuste wzajemnie ortogonalne podzbiory.
\end{definicja}

\begin{lemat}[Różnica korzeni]
\label{aa:lem:root-difference}
Niech $\alpha,\beta$ będą różnymi, nieproporcjonalnymi
korzeniami układu $\Phi$ z Definicji~\ref{aa:def:root-system}.
Jeśli $(\alpha,\beta)>0$, to $\alpha-\beta\in\Phi$.
\end{lemat}
\begin{proof}
Liczby $p=2(\alpha,\beta)/(\alpha,\alpha)$ oraz
$q=2(\alpha,\beta)/(\beta,\beta)$ są dodatnimi
liczbami całkowitymi. Ze ścisłej nierówności
Cauchy'ego--Schwarza dla nieproporcjonalnych wektorów
wynika $pq=4(\alpha,\beta)^2/
((\alpha,\alpha)(\beta,\beta))<4$.
Co najmniej jedna z liczb $p,q$ jest równa $1$.
Jeśli $q=1$, to $s_\beta(\alpha)=\alpha-\beta\in\Phi$.
Jeśli $p=1$, to $s_\alpha(\beta)=\beta-\alpha\in\Phi$,
więc także jej przeciwieństwo należy do $\Phi$.
\end{proof}

\begin{definicja}[Baza i korzenie dodatnie]
\label{aa:def:root-base}
Niech $\Phi$ będzie zredukowanym krystalograficznym
układem korzeni w $E$. Bazą układu nazywamy bazę
$\Delta=(\alpha_1,\ldots,\alpha_n)$ przestrzeni $E$
złożoną z korzeni, względem której współrzędne
każdego korzenia są całkowite i albo wszystkie
nieujemne, albo wszystkie niedodatnie.
Korzenie z pierwszej klasy nazywamy dodatnimi,
z drugiej ujemnymi; zbiory te oznaczamy
$\Phi^+$ oraz $\Phi^-$.
Elementy $\Delta$ nazywamy korzeniami prostymi.
\end{definicja}

\begin{twierdzenie}[Istnienie bazy układu]
\label{aa:tw:root-base}
Niech $\Phi$ będzie układem jak w
Definicji~\ref{aa:def:root-system}. Dla każdego
$h\in E$ takiego, że $(h,\alpha)\ne0$ dla
wszystkich $\alpha\in\Phi$, istnieje baza układu,
dla której $\Phi^+=\{\alpha\in\Phi:(h,\alpha)>0\}$.
Różne elementy tej bazy mają niedodatni iloczyn skalarny.
\end{twierdzenie}
\begin{proof}
Taki wektor $h$ istnieje, ponieważ skończona suma
właściwych hiperpłaszczyzn nie może pokryć $E$
(na przykład iloczyn ich niezerowych równań liniowych
jest niezerowym wielomianem).
W zbiorze korzeni o $(h,\alpha)>0$ wybierzmy wszystkie
korzenie nierozkładalne na sumę dwóch korzeni z tego
samego zbioru; niech $\Delta$ będzie ich zbiorem.
Każdy dodatni korzeń jest sumą elementów $\Delta$:
uporządkujmy skończony zbiór dodatnich wartości
$(h,\alpha)$ rosnąco. Jeżeli korzeń nie należy do
$\Delta$, to $\alpha=\beta+\gamma$, gdzie
$(h,\beta),(h,\gamma)>0$ są obie mniejsze od
$(h,\alpha)$. Indukcja po tym uporządkowaniu
daje rozkład obu składników, a więc i $\alpha$.
Każdy ujemny korzeń jest przeciwieństwem dodatniego.

Dla różnych $\alpha,\beta\in\Delta$ nie mogą one być
proporcjonalne, ponieważ układ jest zredukowany,
a oba korzenie dodatnie. Gdyby $(\alpha,\beta)>0$,
Lemat~\ref{aa:lem:root-difference} dałby korzeń
$\alpha-\beta$. Jeśli jest dodatni, rozkład
$\alpha=(\alpha-\beta)+\beta$ przeczy wyborowi $\alpha$.
Jeśli jest ujemny, rozkład
$\beta=(\beta-\alpha)+\alpha$ przeczy wyborowi $\beta$.
Zatem $(\alpha,\beta)\le0$.

Aby dowieść niezależności liniowej $\Delta$, rozważmy
równość $\sum_{\alpha\in\Delta}t_\alpha\alpha=0$.
Rozdzielmy dodatnie i ujemne współczynniki:
$u=\sum_{t_\alpha>0}t_\alpha\alpha
=\sum_{t_\beta<0}(-t_\beta)\beta=v$.
Z niedodatnich iloczynów skalarnych między rozłącznymi
rodzinami wynika $(u,v)\le0$. Ponieważ $u=v$,
mamy $\|u\|^2\le0$, więc $u=v=0$. Każdy niepusty
z tych dwóch dodatnich składników miałby dodatni
iloczyn skalarny z $h$, sprzeczność.
Wszystkie współczynniki są więc zerowe.
Z wcześniejszego rozkładu korzeni oraz faktu,
że $\Phi$ rozpina $E$, wynika, że $\Delta$ rozpina $E$.
Rozkłady są zatem całkowitymi współrzędnymi
w bazie, o wymaganych znakach.
\end{proof}

\begin{twierdzenie}[Odbicia proste i komory Weyla]
\label{aa:tw:weyl-chambers}
Niech $\Delta=(\alpha_1,\ldots,\alpha_n)$ będzie
bazą skończonego układu korzeni $\Phi$ w $E$.
Wówczas każdy korzeń jest obrazem pewnego korzenia
prostego przez iloczyn odbić prostych $s_{\alpha_i}$.
Grupa $W(\Phi)$ jest skończona i generowana przez te
odbicia. Dowolne dwie bazy układu jako zbiory są
przeprowadzane jedna na drugą przez element $W(\Phi)$.
Grupa działa przechodnio na składowych spójności
zbioru $E\setminus\bigcup_{\alpha\in\Phi}\alpha^\perp$,
zwanych komorami Weyla. Jedna z komór ma postać
$\mathcal C=\{x\in E:(x,\alpha_i)>0,\ i=1,\ldots,n\}$.
\end{twierdzenie}
\begin{proof}
Grupa Weyla permutuje skończony zbiór $\Phi$.
Jeśli element grupy ustala każdy korzeń, ustala ich
rozpiętość $E$, czyli jest identycznością. Działanie
jest wierne, zatem grupa jest skończona.

Dla dodatniego korzenia $\beta=\sum_i k_i\alpha_i$
definiujemy wysokość $\sum_i k_i$. Jeśli $\beta$
nie jest prosty, to z
$(\beta,\beta)=\sum_i k_i(\alpha_i,\beta)>0$
wynika istnienie $i$ z $k_i>0$ i
$(\alpha_i,\beta)>0$. Wtedy
$s_{\alpha_i}(\beta)=\beta-m\alpha_i$ dla dodatniej
liczby całkowitej $m$ i jest korzeniem.
Gdyby $\beta$ miał tylko jedną niezerową współrzędną,
byłby dodatnią wielokrotnością $\alpha_i$, a z redukcji
wynikałoby $\beta=\alpha_i$, wbrew założeniu.
Zatem co najmniej jedna współrzędna poza $i$
jest dodatnia i nie zmienia się w odbiciu.
Wszystkie współrzędne odbitego korzenia są wobec tego
nieujemne (nie mogą być wszystkie niedodatnie).
Jego wysokość zmalała o $m>0$.
Indukcja kończy się na korzeniu prostym.
Dla korzenia ujemnego stosujemy argument do jego
przeciwieństwa i korzystamy z $-\alpha_i=s_{\alpha_i}(\alpha_i)$.
Dla izometrii $w$ zachodzi $s_{w\alpha}=ws_\alpha w^{-1}$,
co widać przez podstawienie do wzoru na odbicie.
Każde odbicie korzeniowe należy więc do podgrupy
generowanej przez odbicia proste.

Zbiór $\mathcal C$ jest niepusty, gdyż istnieje
wektor $h$ określony równaniami $(h,\alpha_i)=1$
(baza gwarantuje jednoznaczną rozwiązywalność).
Jest wypukły, a dla każdego korzenia znak iloczynu
z $x\in\mathcal C$ jest ustalony przez znak jego
współrzędnych. Z drugiej strony znak $(x,\alpha_i)$
nie może się zmienić w żadnej składowej spójności
dopełnienia hiperpłaszczyzn. Zatem $\mathcal C$
jest dokładnie jedną składową.
Weźmy dowolny $x$ w tym dopełnieniu. Skończoność grupy
daje element $w$, dla którego $(wx,h)$ jest maksymalne
w orbicie. Jeśli $(wx,\alpha_i)<0$, to
\[
 (s_{\alpha_i}wx,h)-(wx,h)
 =-\frac{2(wx,\alpha_i)(\alpha_i,h)}{(\alpha_i,\alpha_i)}>0,
\]
co przeczy maksymalności. Równość zera jest niemożliwa,
bo grupa permutuje hiperpłaszczyzny korzeniowe.
Stąd $wx\in\mathcal C$. Ponieważ izometrie grupy
przeprowadzają składowe dopełnienia na składowe,
działanie na komorach jest przechodnie.

Niech $\Delta'$ będzie inną bazą, a $\mathcal C'$
jej komorą zdefiniowaną analogicznie. Istnieje
$w\in W(\Phi)$ takie, że $w\mathcal C=\mathcal C'$.
Jeżeli $x\in\mathcal C$, to
$(wx,w\beta)=(x,\beta)$ dla każdego $\beta\in\Phi$.
Zatem $w$ przeprowadza dodatnie korzenie względem
$\Delta$ na dodatnie korzenie względem $\Delta'$.
W dowolnej bazie korzeń prosty nie jest sumą dwóch
dodatnich korzeni, ponieważ jego współrzędne
tworzą wektor jednostkowy. Odwrotnie, dowód
Twierdzenia~\ref{aa:tw:root-base} pokazuje, że wszystkie
nierozkładalne dodatnie korzenie tworzą bazę.
Ponieważ zawierają już wszystkie elementy danej
bazy, nie mogą zawierać żadnych dodatkowych elementów.
Liniowość i odwracalność $w$ zachowują nierozkładalność
na sumę dodatnich korzeni, więc $w\Delta=\Delta'$.
\end{proof}

\section{Macierz układu i diagram Dynkina}

\begin{twierdzenie}[Macierz Cartana układu korzeni]
\label{aa:tw:root-cartan}
Niech $\Phi$ będzie skończonym układem korzeni jak wyżej,
a $\Delta=(\alpha_i)_{i=1}^n$ jego bazą.
Macierz
\[
 c_{ij}=\frac{2(\alpha_i,\alpha_j)}{(\alpha_i,\alpha_i)}
\]
jest uogólnioną macierzą Cartana typu skończonego.
Jeżeli $i\ne j$, to $c_{ij}c_{ji}\in\{0,1,2,3\}$.
Ponadto układ jest nierozkładalny wtedy i tylko wtedy,
gdy graf macierzy jest spójny.
\end{twierdzenie}
\begin{proof}
Dla dowolnej bazy układu w sensie
Definicji~\ref{aa:def:root-base} wektor określony przez
$(h,\alpha_i)=1$ wyznacza dokładnie te same korzenie
dodatnie. Korzenie proste są nierozkładalne na dwa
dodatnie: ich współrzędne są wektorami jednostkowymi.
Każdy inny dodatni korzeń nie może być nierozkładalny,
gdyż dowód Twierdzenia~\ref{aa:tw:root-base} tworzy
bazę z nierozkładalnych, a ta zawiera już wszystkie
$\alpha_i$. Zatem również dla tej bazy różne
korzenie proste mają niedodatnie iloczyny skalarne.
Całkowitość wynika z aksjomatu krystalograficzności.
Dla $D=\operatorname{diag}((\alpha_i,\alpha_i)/2)$
macierz $DC$ jest macierzą Grama bazy, a więc jest
symetryczna i dodatnio określona:
$x^{\mathsf T}DCx=\|\sum_i x_i\alpha_i\|^2>0$
dla $x\ne0$.
Dla różnych $i,j$ ścisła nierówność Cauchy'ego--Schwarza
daje $0\le c_{ij}c_{ji}<4$, co daje listę.

Jeżeli graf jest niespójny, przestrzenie rozpięte
przez korzenie proste z poszczególnych składowych
są wzajemnie ortogonalne. Odbicia proste zachowują
każdą z tych przestrzeni i działają identycznościowo
na pozostałych. Z Twierdzenia~\ref{aa:tw:weyl-chambers}
każdy korzeń jest w orbicie korzenia prostego,
więc należy do jednej z tych przestrzeni.
Otrzymujemy rozkład ortogonalny układu.
Jeśli zaś $\Phi=\Phi_1\sqcup\Phi_2$ jest rozkładem
na dwa niepuste ortogonalne podzbiory, oba zawierają
korzenie proste: w przeciwnym razie jeden z nich
byłby ortogonalny do całej bazy, niemożliwe dla
niezerowego korzenia. Pomiędzy tymi częściami
nie ma krawędzi.
\end{proof}

\begin{definicja}[Diagram Dynkina]
\label{aa:def:dynkin}
Dla macierzy Cartana typu skończonego rysujemy między
$i,j$ krawędź o krotności $c_{ij}c_{ji}$, jeśli
iloczyn jest niezerowy. Przy krawędzi wielokrotnej
strzałka wskazuje koniec odpowiadający krótszemu
korzeniowi. W konwencji z Twierdzenia~\ref{aa:tw:root-cartan}
oznacza to wierzchołek $i$ wtedy, gdy
$|c_{ij}|>|c_{ji}|$. W przypadku abstrakcyjnej
macierzy symetryzowalnej porównujemy dodatnie
wpisy symetryzatora $D$, odpowiadające połowom
kwadratów długości korzeni.
\end{definicja}

\section{Klasyfikacja macierzy typu skończonego}

\begin{twierdzenie}[Klasyfikacja diagramów skończonych]
\label{aa:tw:finite-dynkin}
Każda nierozkładalna uogólniona macierz Cartana typu
skończonego ma, z dokładnością do jednoczesnej
permutacji wierszy i kolumn, jeden z diagramów
$A_n$ ($n\ge1$), $B_n,C_n$ ($n\ge2$),
$D_n$ ($n\ge4$), $E_6,E_7,E_8,F_4,G_2$.
Tutaj $B_2$ i $C_2$ są równoważne po zamianie wierzchołków.
Wszystkie wymienione macierze są typu skończonego.
\end{twierdzenie}
\begin{proof}
Niech $D$ będzie dodatnim symetryzatorem.
Macierz $S=D^{1/2}CD^{-1/2}$ jest symetryczna
i dodatnio określona, bo jest równa
$D^{-1/2}(DC)D^{-1/2}$. Ma przekątną $2$,
a jej pozadiagonalne wpisy wynoszą
$-\sqrt{c_{ij}c_{ji}}$. Dodatnia określoność
dwuwymiarowych macierzy głównych daje
$4-c_{ij}c_{ji}>0$. Każda krawędź ma więc
wagę $1,\sqrt2$ lub $\sqrt3$.
Oznaczmy macierz ważonego sąsiedztwa przez $P$,
tak że $S=2I-P$.

Będziemy korzystać z następującej obserwacji.
Jeśli pewien podgraf (z ewentualnie zmniejszonymi
wagami) ma macierz sąsiedztwa $Q$ i dodatni wektor
$u$ spełniający $Qu=2u$, to po przedłużeniu $u$
przez zera na całym grafie otrzymujemy
$u^{\mathsf T}(2I-P)u\le u^{\mathsf T}(2I-Q)u=0$,
gdzie w ostatnim wyrażeniu $Q$ przedłużamy zerami.
Jest to sprzeczne z dodatnią określonością.

Cykl o wagach jeden, czteroramienna gwiazda o wagach
jeden oraz podgraf dwóch wierzchołków stopnia trzy
połączonych ścieżką i zakończonych po dwa liście
mają dodatni wektor o wartości własnej $2$
(z dowodu Twierdzenia~\ref{aa:tw:mckay-graphs}).
Obserwacja wyklucza te konfiguracje także w grafie
ważonym. Graf jest zatem drzewem o stopniach co
najwyżej trzy i co najwyżej jednym rozgałęzieniu.

Jeśli występuje krawędź wagi $\sqrt3$ i jeszcze
jeden wierzchołek, istnieje ścieżka trzech
wierzchołków z wagami $\sqrt3$ i co najmniej $1$.
Po zmniejszeniu drugiej wagi do $1$ wektor
$(\sqrt3,2,1)$, w kolejności ścieżki, spełnia
$Qu=2u$. Jest to wykluczone, więc pozostaje
jedna krawędź potrójna, czyli $G_2$.

Jeśli występują dwie krawędzie wagi $\sqrt2$,
wybieramy dwie takie krawędzie i jedyną ścieżkę
zawierającą obie jako krawędzie skrajne. Istnieje
ona, ponieważ graf jest drzewem: po usunięciu
każdej z tych krawędzi wybieramy jej koniec leżący
w składowej niezawierającej drugiej krawędzi
i łączymy dwa tak wybrane końce.
Zmniejszamy wszystkie pośrednie wagi do $1$.
Nadajemy jej skrajnym wierzchołkom wartości $1$,
a wewnętrznym $\sqrt2$. Wewnętrzny wierzchołek
przy skrajnej krawędzi ma sumę ważoną $2\sqrt2$,
podobnie wewnętrzne wierzchołki przy dwóch wagach
$1$; gdy krawędzie podwójne są sąsiednie, ich wspólny
wierzchołek ma sumę $2\sqrt2$. Skrajny wierzchołek
ma sumę $2$. Obserwacja wyklucza tę ścieżkę.
Zatem co najwyżej jedna krawędź jest podwójna.

Jeśli w drzewie z krawędzią podwójną jest wierzchołek
stopnia trzy, wybierzmy ścieżkę od niego do
dalszego końca krawędzi podwójnej oraz krawędzie
do dwóch jego sąsiadów spoza ścieżki. Sąsiedzi
ci nie należą do ścieżki, gdyż drzewo nie ma cykli.
Zmniejszamy wagi do $1$ poza wskazaną podwójną.
Na całej ścieżce poza ostatnim wierzchołkiem
nadajemy wartość $2$, na ostatnim $\sqrt2$,
a na dwóch dodatkowych liściach $1$.
W centrum suma ważona wynosi $4$ (również gdy
podwójna krawędź wychodzi wprost z centrum).
Wewnątrz ścieżki suma wynosi $4$, na jej ostatnim
wierzchołku $2\sqrt2$, na liściach $2$.
Znów otrzymujemy niedopuszczalny podgraf.
Drzewo z jedną krawędzią podwójną jest więc ścieżką.

Przyda się rachunek dla ścieżki $m$ wierzchołków
z samymi wagami $1$. Niech $T_m$ ma przekątną
$2$ i wpisy $-1$ przy krawędziach. Z rozwinięcia
wyznacznika ostatnim wierszem mamy
$\det T_m=2\det T_{m-1}-\det T_{m-2}$ dla $m\ge2$,
z warunkami $\det T_0=1$, $\det T_1=2$.
Indukcyjnie $\det T_m=m+1$, więc wszystkie wiodące
minory są dodatnie, a $T_m$ jest dodatnio określona.
Z formuły na odwrotność przez dopełnienia algebraiczne
$(T_m^{-1})_{11}=(T_m^{-1})_{mm}=m/(m+1)$.

Niech podwójna krawędź dzieli ścieżkę na dwie
części o $a,b\ge1$ wierzchołkach. Po uporządkowaniu
macierzy według tych części bloki diagonalne to
$T_a,T_b$, a jedyne połączenie ma wpis $-\sqrt2$.
Blok pozadiagonalny to $-\sqrt2 e_a e_1^{\mathsf T}$,
gdzie $e_a$ jest ostatnim wektorem standardowym
w $\mathbb R^a$. Dopełnienie Schura względem $T_a$ to
$T_b-2a/(a+1)e_1e_1^{\mathsf T}$, gdzie $e_1$
jest pierwszym wektorem standardowym w $\mathbb R^b$.
Dla dodatnio określonej $T$ macierz
$T-cee^{\mathsf T}$ przy $c>0$ jest dodatnio
określona dokładnie wtedy, gdy
$c e^{\mathsf T}T^{-1}e<1$: po sprzężeniu przez
$T^{-1/2}$ pozostaje $I-cvv^{\mathsf T}$,
gdzie $v=T^{-1/2}e$, o wartościach własnych
$1$ na $v^\perp$ oraz $1-c\|v\|^2$ na $\mathbb R v$.
Stąd warunek brzmi
$2ab<(a+1)(b+1)$, czyli $(a-1)(b-1)<2$.
Jeśli jedna z liczb wynosi $1$, mamy typy $B_n,C_n$
zależnie od orientacji. Jeżeli obie są co najmniej
$2$, ich zmniejszone o jeden wartości są dodatnimi
liczbami całkowitymi o iloczynie mniejszym od $2$,
więc $a=b=2$, czyli $F_4$.

W przypadku samych krawędzi pojedynczych brak
rozgałęzienia daje $A_n$. Przy jednym rozgałęzieniu
niech ramiona mają długości $l_1,l_2,l_3\ge1$.
Bloki ramion to $T_{l_i}$, a środek jest połączony
z pierwszym wierzchołkiem każdego ramienia wpisem $-1$.
Dopełnienie Schura względem sumy bloków ramion
jest liczbą
\[
 2-\sum_{i=1}^3\frac{l_i}{l_i+1}
 =\sum_{i=1}^3\frac1{l_i+1}-1.
\]
Dodatnia określoność zachodzi wtedy i tylko wtedy,
gdy ta liczba jest dodatnia. Istotnie, dla symetrycznej macierzy blokowej
$\begin{pmatrix}P&R\\R^{\mathsf T}&Q\end{pmatrix}$
z dodatnio określoną $P$ jej forma w wektorze $(x,y)$
jest równa
\[
 (x+P^{-1}Ry)^{\mathsf T}P(x+P^{-1}Ry)
 +y^{\mathsf T}(Q-R^{\mathsf T}P^{-1}R)y.
\]
Stąd dodatnia określoność całej macierzy jest
równoważna dodatniej określoności dopełnienia:
dla konieczności wybieramy $x=-P^{-1}Ry$,
a dla dostateczności rozpatrujemy osobno $y=0$
i $y\ne0$.
Uporządkujmy $a=l_1+1\le b=l_2+1\le c=l_3+1$.
Warunek $1/a+1/b+1/c>1$ wyklucza $a\ge3$.
Dla $a=2,b=2$ dowolne $c\ge2$ daje ramiona
$(1,1,c-1)$, czyli $D_{c+2}$.
Dla $a=2,b=3$ warunek $1/c>1/6$ daje $c=3,4,5$,
czyli ramiona $(1,2,2),(1,2,3),(1,2,4)$
dla $E_6,E_7,E_8$.
Dla $a=2,b\ge4$ suma jest najwyżej $1$.
Lista jest wyczerpująca.

Wszystkie grafy z listy są drzewami. Dla dowolnej
orientacji krawędzi wielokrotnych definiujemy
$c_{ij},c_{ji}$ jako $-1,-1$ przy pojedynczej,
$-1,-2$ (w odpowiedniej kolejności) przy podwójnej
i $-1,-3$ przy potrójnej. Symetryzator istnieje:
ustalamy jedną dodatnią wartość $D_{ii}$ i wzdłuż
każdej krawędzi narzucamy
$D_{jj}=D_{ii}c_{ij}/c_{ji}$. Jedyna ścieżka
w drzewie zapewnia zgodność definicji.
Dodatnią określoność dla $A,D,E,B,C,F$ wykazują
powyższe dopełnienia Schura, a dla $G_2$ macierz
symetryczna ma przekątną $2$ i wyznacznik $1>0$.
Odwrócenie ścieżki identyfikuje orientacje w $F_4$
i $G_2$, a także w $B_2,C_2$.
\end{proof}

\begin{center}
\begin{tikzpicture}[v/.style={circle,draw,inner sep=0pt,minimum size=4mm}]
 \foreach \i in {0,...,3} {\node[v] (a\i) at (\i*.75,0) {};}
 \draw (a0)--(a1)--(a2)--(a3);
 \node at (1.125,-.45) {$A_4$};
 \foreach \i in {0,...,3} {\node[v] (b\i) at (3.3+\i*.75,0) {};}
 \draw (b0)--(b1)--(b2);
 \draw[double distance=2pt,-{Stealth[length=2mm]}] (b2)--(b3);
 \node at (4.425,-.45) {$B_4$};
 \foreach \i in {0,...,3} {\node[v] (c\i) at (6.6+\i*.75,0) {};}
 \draw (c0)--(c1)--(c2);
 \draw[double distance=2pt,{Stealth[length=2mm]}-] (c2)--(c3);
 \node at (7.725,-.45) {$C_4$};
 \foreach \i in {0,...,3} {\node[v] (f\i) at (3.3+\i*.75,-1.5) {};}
 \draw (f0)--(f1);
 \draw[double distance=2pt,-{Stealth[length=2mm]}] (f1)--(f2);
 \draw (f2)--(f3);
 \node at (4.425,-1.95) {$F_4$};
 \node[v] (d0) at (0,-1.5) {}; \node[v] (d1) at (.75,-1.5) {};
 \node[v] (d2) at (1.5,-1.15) {}; \node[v] (d3) at (1.5,-1.85) {};
 \draw (d0)--(d1)--(d2); \draw (d1)--(d3);
 \node at (.75,-2.2) {$D_4$};
 \node[v] (g0) at (7,-1.5) {}; \node[v] (g1) at (8,-1.5) {};
 \draw[-{Stealth[length=2mm]}] (g0)--(g1);
 \draw (7.2,-1.42)--(7.8,-1.42);
 \draw (7.2,-1.58)--(7.8,-1.58);
 \node at (7.5,-1.95) {$G_2$};
\end{tikzpicture}

\small Przykłady diagramów zwykłych. Strzałka wskazuje krótszy korzeń.
\end{center}

\begin{center}
\begin{tikzpicture}[v/.style={circle,draw,inner sep=0pt,minimum size=4mm}]
 \foreach \n/\r/\y in {6/4/0,7/5/-1.5,8/6/-3}
 {
  \foreach \i in {0,...,\r}
    {\node[v] (e\n-\i) at (.8*\i,\y) {};}
  \foreach \i in {1,...,\r}
    {\pgfmathtruncatemacro{\j}{\i-1}\draw (e\n-\j)--(e\n-\i);}
  \node[v] (u\n) at (1.6,\y+.65) {};
  \draw (e\n-2)--(u\n);
  \node[anchor=west] at (5.8,\y) {$E_{\n}$};
 }
\end{tikzpicture}

\small Diagramy $E_6,E_7,E_8$: ramiona mają długości $(1,2,2)$,
$(1,2,3)$ i $(1,2,4)$.
\end{center}

\section{Realizacja i jednoznaczność układów}

\begin{twierdzenie}[Jednoznaczność na podstawie macierzy]
\label{aa:tw:root-uniqueness}
Dwa nierozkładalne skończone układy korzeni z bazami
o tej samej macierzy Cartana są równoważne przez
podobieństwo liniowe przestrzeni euklidesowych.
\end{twierdzenie}
\begin{proof}
Niech $D,D'$ oznaczają diagonalne macierze kwadratów
długości podzielonych przez $2$ w obu bazach.
Na każdej krawędzi iloraz $D_{ii}/D_{jj}$ jest
wyznaczony przez $c_{ji}/c_{ij}$; to samo dla $D'$.
Spójność daje $D'=tD$ dla pewnego $t>0$.
Zatem macierze Grama baz różnią się mnożnikiem $t$.
Liniowe odwzorowanie $L$ przeprowadzające bazę
na bazę spełnia $(Lx,Ly)'=t(x,y)$ dla dowolnych
$x,y$, ponieważ obie strony są dwuliniowe i
zgodne na parach wektorów bazowych.
Z formuły na odbicie wynika
$Ls_{\alpha_i}=s_{L\alpha_i}L$.
Twierdzenie~\ref{aa:tw:weyl-chambers} mówi,
że każdy korzeń jest obrazem korzenia prostego
przez iloczyn odbić prostych; stosując tę równość
dla $L$ i $L^{-1}$ otrzymujemy $L(\Phi)=\Phi'$.
\end{proof}

\begin{twierdzenie}[Realizacja diagramów jednokrawędziowych]
\label{aa:tw:simply-laced-realization}
Niech $C$ będzie jedną z macierzy $A_n,D_n,E_6,E_7,E_8$
z Twierdzenia~\ref{aa:tw:finite-dynkin}. Na $\mathbb R^n$
z bazą standardową $(e_i)$ określmy iloczyn skalarny
$(x,y)=x^{\mathsf T}Cy$ i zbiór
\[
 \Phi_C=\{x\in\mathbb Z^n:x^{\mathsf T}Cx=2\}.
\]
Jest to zredukowany krystalograficzny układ korzeni
z bazą $(e_i)$ i macierzą Cartana $C$.
\end{twierdzenie}
\begin{proof}
Macierz $C$ jest symetryczna i dodatnio określona
na mocy Twierdzenia~\ref{aa:tw:finite-dynkin}.
Istnieje $\delta>0$ takie, że
$x^{\mathsf T}Cx\ge\delta\|x\|_{\mathrm{eukl}}^2$
(dodatnie minimum na sferze jednostkowej).
Zbiór wektorów całkowitych o ograniczonej normie
jest skończony. Każdy $e_i$ należy do $\Phi_C$,
więc korzenie rozpinają przestrzeń.
Wszystkie korzenie mają długość $\sqrt2$,
więc proporcjonalne korzenie różnią się co najwyżej
znakiem. Dla $\alpha,\beta\in\Phi_C$ liczba
$(\alpha,\beta)=\alpha^{\mathsf T}C\beta$ jest całkowita,
a odbicie $s_\alpha(\beta)=\beta-(\alpha,\beta)\alpha$
jest całkowite i zachowuje długość, więc należy
do $\Phi_C$. Ponieważ odbicie jest inwolucją,
przeprowadza zbiór na siebie.

Forma $x^{\mathsf T}Cx$ przyjmuje na $\mathbb Z^n$
wartości parzyste: przekątna jest równa $2$,
a wpisy pozadiagonalne pojawiają się parami.
Dla niezerowego wektora całkowitego jej wartość
jest więc co najmniej $2$.
Jeśli $\beta$ miałby współrzędne obu znaków,
zapiszmy $\beta=u-v$ dla niezerowych wektorów
całkowitych $u,v$ o nieujemnych współrzędnych
i rozłącznych nośnikach. Wtedy
$(u,v)\le0$, bo ich wspólne składniki diagonalne
znikają, a pozostałe wpisy $C$ są niedodatnie.
Stąd $(\beta,\beta)=(u,u)+(v,v)-2(u,v)\ge4$,
sprzecznie z definicją $\Phi_C$.
Współrzędne każdego korzenia mają jeden znak,
więc $(e_i)$ jest bazą układu.
Macierz Cartana tej bazy to $C$, gdyż
$(e_i,e_i)=2$ i $(e_i,e_j)=c_{ij}$.
\end{proof}

\begin{przyklad}[Modele klasyczne]
\label{aa:prz:classical-roots}
W poniższych modelach $e_i$ to standardowe
ortonormalne wektory, a wszystkie wskazane
znaki można wybierać niezależnie.
\[
\begin{array}{c|l|l}
 &\Phi&\text{korzenie proste}\\ \hline
 A_n & e_i-e_j,\ i\ne j,\ 1\le i,j\le n+1
       & e_i-e_{i+1}\ (1\le i\le n)\\
 B_n & \pm e_i,\ \pm e_i\pm e_j,\ i<j
       & e_i-e_{i+1}\ (i<n),\ e_n\\
 C_n & \pm2e_i,\ \pm e_i\pm e_j,\ i<j
       & e_i-e_{i+1}\ (i<n),\ 2e_n\\
 D_n & \pm e_i\pm e_j,\ i<j
       & e_i-e_{i+1}\ (i<n),\ e_{n-1}+e_n
\end{array}
\]
Dla $A_n$ pracujemy w hiperpłaszczyźnie
$\sum x_i=0$ w $\mathbb R^{n+1}$, dla pozostałych
w $\mathbb R^n$; zakresy $n$ są jak w klasyfikacji.

Odbicie względem $e_i-e_j$ zamienia współrzędne
$i,j$, a względem $e_i+e_j$ zamienia je i zmienia
oba znaki. Odbicie względem $e_i$ lub $2e_i$ zmienia
znak jednej współrzędnej. Zatem każdy zbiór jest
zachowany przez wszystkie własne odbicia.
Układy są zredukowane i całkowitość parowań wynika
z podanych długości: korzenie typu $\pm e_i\pm e_j$
mają kwadrat długości $2$ i całkowite iloczyny
skalarne; $\pm e_i$ mają kwadrat długości $1$;
$\pm2e_i$ mają kwadrat długości $4$, lecz ich
iloczyny z każdym korzeniem są parzyste.
Wszystkie podane rodziny korzeni prostych są bazami:
różnice rozpinają hiperpłaszczyznę sumy zerowej,
a w typach $B,C,D$ ostatni wektor ma niezerową
sumę współrzędnych. Dla $i<j$ mamy
$e_i-e_j=\sum_{k=i}^{j-1}(e_k-e_{k+1})$.
W typie $B$ każdy $e_i$ jest nieujemną sumą
korzeni prostych, zatem również $e_i+e_j$.
W typie $C$ piszemy $2e_i=2(e_i-e_n)+2e_n$,
$e_i+e_j=(e_i-e_n)+(e_j-e_n)+2e_n$.
W typie $D$ dla $i<j<n$ piszemy
$e_i+e_j=(e_i-e_{n-1})+(e_j-e_n)+(e_{n-1}+e_n)$,
a dla $j=n$ mamy
$e_i+e_n=(e_i-e_{n-1})+(e_{n-1}+e_n)$.
Wszystkie składniki są nieujemnymi sumami,
pozostałe korzenie są ich przeciwieństwami.
Macierze tych baz mają podane diagramy,
co wynika z iloczynów skalarnych sąsiednich różnic
i ostatniego korzenia.
\end{przyklad}

\begin{przyklad}[Model $F_4$]
\label{aa:prz:f4-roots}
W $\mathbb R^4$ weźmy
\[
 \Phi=\{\pm e_i\}\cup\{\pm e_i\pm e_j:i<j\}
       \cup\left\{\tfrac12(\pm e_1\pm e_2\pm e_3\pm e_4)\right\}.
\]
Zbiór liczy $8+24+16=48$ korzeni i jest zredukowany;
rozpina przestrzeń, ponieważ zawiera $e_i$.
Dla długiego korzenia $\alpha=\pm e_i\pm e_j$
parowanie $2(\beta,\alpha)/(\alpha,\alpha)=(\beta,\alpha)$
jest całkowite nawet dla korzeni połówkowych,
bo dodajemy lub odejmujemy dwie połówki.
Dla krótkiego korzenia $\alpha$ mamy
$(\alpha,\alpha)=1$; iloczyn z korzeniem
połówkowym, jednostkowym lub długim jest
wielokrotnością $1/2$ (dla dwóch połówkowych
suma czterech liczb $\pm1/4$ jest wielokrotnością
$1/2$). Stąd całkowitość.

Odbicia względem korzeni jednostkowych i długich
są zmianami znaków i permutacjami współrzędnych,
więc zachowują $\Phi$. Dla połówkowych wystarczy
sprawdzić $u=(e_1+e_2+e_3+e_4)/2$, ponieważ
zmiany znaków przeprowadzają $u$ na każdy taki
korzeń. Odbicie ma postać $s_u(x)=x-2(x,u)u$.
Dla $x=e_i$ wynikiem jest $e_i-u$, korzeń
połówkowy; przeciwieństwa dają przeciwieństwa.
Dla korzenia długiego iloczyn z $u$ jest $0$ lub
$\pm1$: w pierwszym przypadku korzeń jest ustalony,
w drugim odbicie zastępuje dwie niezerowe współrzędne
dwiema przeciwnymi na dopełniających indeksach.
Dla połówkowego $x$ iloczyn $(x,u)$ jest
$0,\pm1/2,\pm1$. Jeśli wynosi $0$, korzeń jest
ustalony; jeśli $\pm1$, mamy $x=\pm u$, a odbicie
zmienia znak. Jeśli $1/2$, trzy znaki są dodatnie
i jeden ujemny, więc $x-u=-e_i$ dla ujemnego
indeksu. Przypadek $-1/2$ wynika przez zmianę znaku.
To dowodzi zachowania całego zbioru.

Bazę tworzą
$\alpha_1=e_2-e_3$, $\alpha_2=e_3-e_4$,
$\alpha_3=e_4$, $\alpha_4=(e_1-e_2-e_3-e_4)/2$.
Mamy $e_4=\alpha_3$, $e_3=\alpha_2+\alpha_3$,
$e_2=\alpha_1+\alpha_2+\alpha_3$ i
$e_1=2\alpha_4+e_2+e_3+e_4$.
Te równości dowodzą, że wektory tworzą bazę.
Wszystkie $e_i$ i $e_i\pm e_j$ dla $i<j$
mają nieujemne współrzędne w tej bazie:
różnice kolejnych $e_i$ są nieujemnymi sumami,
a sumy również. Korzeń połówkowy o dodatniej
współrzędnej przy $e_1$ ma postać
$\alpha_4+\sum_{i=2}^4\varepsilon_i e_i$,
gdzie $\varepsilon_i$ to $0$ albo $1$.
Pozostałe są przeciwieństwami. Otrzymujemy bazę
układu, z długimi $\alpha_1,\alpha_2$ i krótkimi
$\alpha_3,\alpha_4$, a jedyne niezerowe iloczyny
między różnymi prostymi to
$(\alpha_1,\alpha_2)=-1$, $(\alpha_2,\alpha_3)=-1$,
$(\alpha_3,\alpha_4)=-1/2$. Diagram jest $F_4$.
\end{przyklad}

\begin{przyklad}[Model $G_2$]
\label{aa:prz:g2-roots}
W płaszczyźnie $x_1+x_2+x_3=0$ w $\mathbb R^3$
weźmy sześć krótkich korzeni $e_i-e_j$ ($i\ne j$)
i sześć długich $\pm(2e_i-e_j-e_k)$, gdzie
$\{i,j,k\}=\{1,2,3\}$. Kwadraty długości wynoszą
odpowiednio $2$ i $6$, a zbiór jest zredukowany:
krótkie korzenie mają współrzędną zerową,
długie jej nie mają. Parowania względem krótkich
są całkowite. Dla długiego
$\alpha=2e_i-e_j-e_k$ i dowolnego $\beta$
w tej płaszczyźnie mamy $(\beta,\alpha)=3\beta_i$,
więc parowanie wynosi $\beta_i\in\mathbb Z$.
Odbicia krótkie zamieniają współrzędne,
zaś odbicie długie jest na płaszczyźnie
ujemnością transpozycji współrzędnych $j,k$:
obie izometrie zmieniają znak $\alpha$ i ustalają
prostopadły do niego $e_j-e_k$.
Zatem wszystkie odbicia zachowują zbiór.
Jako bazę bierzemy $\alpha=e_1-e_2$ oraz
$\beta=-2e_1+e_2+e_3$. Są niezależne, a lista
\[
 \alpha,\ \beta,\ \alpha+\beta,\ 2\alpha+\beta,\
 3\alpha+\beta,\ 3\alpha+2\beta
\]
ma w standardowych współrzędnych postać
\[
 (1,-1,0),\ (-2,1,1),\ (-1,0,1),\
 (0,-1,1),\ (1,-2,1),\ (-1,-1,2).
\]
Są to sześć różnych korzeni, żadne dwa nie są
przeciwieństwami, a ich współrzędne względem
$(\alpha,\beta)$ są nieujemne. Razem z przeciwieństwami
wyczerpują wszystkie dwanaście korzeni.
Iloczyn $(\alpha,\beta)=-3$, więc macierz bazy
to $\begin{pmatrix}2&-3\\-1&2\end{pmatrix}$.
\end{przyklad}

\begin{wniosek}[Klasyfikacja układów korzeni]
\label{aa:cor:root-classification}
Każdy nierozkładalny zredukowany krystalograficzny
skończony układ korzeni jest, z dokładnością do
podobieństwa przestrzeni euklidesowych, jednym
z typów $A_n,B_n,C_n,D_n,E_6,E_7,E_8,F_4,G_2$,
z zakresami jak w Twierdzeniu~\ref{aa:tw:finite-dynkin}.
\end{wniosek}
\begin{proof}
Istnienie bazy i macierzy typu skończonego wynika
z Twierdzeń~\ref{aa:tw:root-base}
i~\ref{aa:tw:root-cartan}. Klasyfikujemy tę macierz
przez Twierdzenie~\ref{aa:tw:finite-dynkin}.
Realizacje jednokrawędziowe daje
Twierdzenie~\ref{aa:tw:simply-laced-realization},
a pozostałe podane modele $B,C,F_4,G_2$.
Spójność każdej otrzymanej macierzy daje
nierozkładalność realizacji. Jednoznaczność dla
ustalonej macierzy daje Twierdzenie~\ref{aa:tw:root-uniqueness}.
Jeśli podobieństwo przeprowadza jeden układ na drugi,
przeprowadza także jego bazę na pewną bazę drugiego
układu i zachowuje wszystkie parowania Cartana,
ponieważ iloczyny skalarne i kwadraty długości
zmieniają się o ten sam dodatni czynnik.
Na mocy Twierdzenia~\ref{aa:tw:weyl-chambers}
wszystkie bazy drugiego układu są równoważne przez
izometrię z jego grupy Weyla. Macierze różnią się
więc tylko jednoczesną permutacją indeksów.
Różne diagramy z listy (po utożsamieniu $B_2$ i $C_2$)
nie mogą zatem opisywać podobnych układów.
\end{proof}

\section{Przykład grupy Weyla i powrót do grafów McKaya}

\begin{przyklad}[Grupy Weyla typów klasycznych]
Dla modelu $A_n$ z Przykładu~\ref{aa:prz:classical-roots}
odbicia są transpozycjami współrzędnych, które
generują $S_{n+1}$. Działanie na hiperpłaszczyźnie
jest wierne, ponieważ permutacja ustalająca wszystkie
różnice $e_i-e_j$ musi ustalać każdą uporządkowaną
parę różnych indeksów. Stąd $W(A_n)\simeq S_{n+1}$.
Dla $B_n,C_n$ odbicia generują wszystkie permutacje
i wszystkie niezależne zmiany znaków współrzędnych.
Każde takie przekształcenie ma jednoznaczny rozkład
na zmianę znaków i permutację, więc grupa jest
iloczynem półprostym $(C_2)^n\rtimes S_n$
o rzędzie $2^n n!$.
Dla $D_n$ odbicia zachowują parzystość liczby zmienionych
znaków. Wszystkie permutacje są generowane przez
odbicia względem różnic, a zmiana znaków dokładnie
dwóch współrzędnych jest iloczynem odbić względem
$e_i-e_j$ i $e_i+e_j$. Każdą parzystą rodzinę
zmian znaków rozkładamy na takie pary.
Zatem $W(D_n)$ jest iloczynem półprostym podgrupy
parzystych zmian znaków $(C_2)^{n-1}$ i $S_n$,
o rzędzie $2^{n-1}n!$.
\end{przyklad}

\begin{wniosek}[Od diagramu afinicznego do zwykłego]
\label{aa:cor:mckay-finite-dynkin}
Niech $G$ będzie nietrywialną skończoną podgrupą
$\operatorname{SL}_2(\mathbb C)$, a $A$ jej macierzą
McKaya. Macierz uzyskana z $2I-A$ przez usunięcie
wiersza i kolumny reprezentacji trywialnej ma typ
$A_{m-1}$ dla $C_m$, $D_{n+2}$ dla $\operatorname{Dic}_n$
oraz odpowiednio $E_6,E_7,E_8$ dla $2T,2O,2I$.
\end{wniosek}
\begin{proof}
Dodatnia określoność wynika ze
Stwierdzenia~\ref{aa:prop:mckay-cartan}.
Dla cyklu $\widetilde A_{m-1}$ usunięcie wierzchołka
daje ścieżkę $m-1$ wierzchołków; dla $m=2$
pozostaje macierz $(2)$ typu $A_1$.
Dla $\widetilde D_4$ usunięcie liścia daje trzy
ramiona długości $1$. Dla $\widetilde D_{n+2}$
z $n\ge3$ trywialny wierzchołek jest liściem
przy końcu ścieżki: po jego usunięciu drugi koniec
jest jedynym rozgałęzieniem, z ramionami
$(1,1,n-1)$, co daje $D_{n+2}$.
Dla $\widetilde E_6,\widetilde E_7,\widetilde E_8$
liść o wymiarze jeden znajduje się na ramieniu
długości odpowiednio $2,3,5$. Skrócenie jednego
takiego ramienia daje $(1,2,2),(1,2,3),(1,2,4)$.
To są dokładnie diagramy zwykłe z
Twierdzenia~\ref{aa:tw:finite-dynkin}.
\end{proof}

Związek jest algebraiczno-kombinatoryczny. Konstrukcja
algebr Liego z układów korzeni, pełna teoria afinicznych
macierzy Cartana i geometryczna odpowiedniość McKaya
wymagają dalszych rozdziałów; nie zakładaliśmy ich
w powyższych dowodach.

% AA-57
\chapter{Algebry Liego. Twierdzenia Engela i Liego}
\label{aa:ch:lie-algebras}

W rozdziale~\ref{aa:ch:root-systems} rozważaliśmy układy
korzeni i grupy Weyla. Do algebraicznej interpretacji tych
obiektów potrzebujemy teorii algebr Liego. Zaczniemy od
rachunku komutatorów, a następnie wykażemy dwa podstawowe
twierdzenia o reprezentacjach algebr nilpotentnych
i rozwiązalnych.

\section{Nawias Liego i przykłady macierzowe}

\begin{definicja}[Algebra Liego]
\label{aa:def:lie-algebra}
Niech $K$ będzie ciałem. Algebrą Liego nad $K$ nazywamy
przestrzeń wektorową $\mathfrak g$ nad $K$ z dwuliniowym
odwzorowaniem $[\cdot,\cdot]:\mathfrak g\times\mathfrak g
\to\mathfrak g$, zwanym nawiasem Liego, spełniającym
dla wszystkich $x,y,z\in\mathfrak g$:
\[
 [x,x]=0,\qquad
 [x,[y,z]]+[y,[z,x]]+[z,[x,y]]=0.
\]
Drugą równość nazywamy tożsamością Jacobiego.
Podalgebrą Liego jest podprzestrzeń zamknięta
na nawias. Dla podprzestrzeni $U,W\subset\mathfrak g$
przez $[U,W]$ rozumiemy rozpiętość liniową wszystkich
$[u,w]$, gdzie $u\in U,w\in W$.
\end{definicja}

Z rozwinięcia $[x+y,x+y]=0$ wynika
$[x,y]=-[y,x]$, także w charakterystyce dwa.
Tożsamość Jacobiego można zapisać w postaci
\[
 [x,[y,z]]=[[x,y],z]+[y,[x,z]].
\]
Istotnie przenosimy dwa pozostałe składniki
na prawą stronę i używamy antysymetrii.

\begin{przyklad}[Algebra łączna z komutatorem]
\label{aa:prz:associative-lie}
Niech $A$ będzie algebrą łączną nad $K$.
Na jej przestrzeni wektorowej połóżmy $[x,y]=xy-yx$.
Nawias jest dwuliniowy i $[x,x]=0$.
Ponadto
\[
\begin{aligned}
 [x,[y,z]]&=xyz-xzy-yzx+zyx,\\
 [y,[z,x]]&=yzx-yxz-zxy+xzy,\\
 [z,[x,y]]&=zxy-zyx-xyz+yxz.
\end{aligned}
\]
Suma trzech prawych stron znika wyraz po wyrazie.
Otrzymujemy algebrę Liego. Dla $A=\operatorname{End}_K(V)$
oznaczamy ją przez $\mathfrak{gl}(V)$, a dla
$V=K^n$ przez $\mathfrak{gl}_n(K)$.
Dla jednostek macierzowych $E_{ij}$ zachodzi
$E_{ij}E_{kl}=\delta_{jk}E_{il}$, ponieważ obie
strony mają to samo działanie na wektorach standardowych.
Zatem
\[
 [E_{ij},E_{kl}]=\delta_{jk}E_{il}-\delta_{li}E_{kj}.
\]
\end{przyklad}

\begin{przyklad}[Algebry klasyczne]
\label{aa:prz:classical-lie}
Niech $K$ będzie ciałem charakterystyki zero
i $n\ge1$. Określamy
\[
\begin{aligned}
 \mathfrak{sl}_n(K)&=\{X\in M_n(K):\operatorname{tr}X=0\},\\
 \mathfrak{so}_n(K)&=\{X\in M_n(K):X^{\mathsf T}+X=0\},\\
 \mathfrak{sp}_{2n}(K)&=\{X\in M_{2n}(K):X^{\mathsf T}J+JX=0\},
 \qquad J=\begin{pmatrix}0&I_n\\-I_n&0\end{pmatrix}.
\end{aligned}
\]
Wszystkie są podprzestrzeniami. Ponieważ
$\operatorname{tr}(XY)=\sum_{i,j}x_{ij}y_{ji}
=\operatorname{tr}(YX)$, każdy komutator ma ślad zero,
więc $\mathfrak{sl}_n$ jest podalgebrą, a nawet
ideałem w $\mathfrak{gl}_n$ w sensie poniższej definicji.

Jeżeli dla ustalonej macierzy $B$ mamy
$X^{\mathsf T}B=-BX$, $Y^{\mathsf T}B=-BY$, to
\[
 [X,Y]^{\mathsf T}B
 =Y^{\mathsf T}X^{\mathsf T}B-X^{\mathsf T}Y^{\mathsf T}B
 =BYX-BXY=-B[X,Y].
\]
Podstawienie $B=I_n$ lub $B=J$ dowodzi zamkniętości
pozostałych dwóch przestrzeni na komutator.

Wymiary wynoszą odpowiednio $n^2-1$, $n(n-1)/2$
oraz $n(2n+1)$. Pierwsza przestrzeń jest jądrem
niezerowego funkcjonału śladu (wartość w $E_{11}$
wynosi $1$). W drugiej wpisy nad przekątną są dowolne,
pod przekątną są ich przeciwieństwami, a przekątna
jest zerowa, ponieważ $2\ne0$ w $K$.
W trzeciej piszemy $X=\begin{pmatrix}P&Q\\R&S\end{pmatrix}$.
Mnożenie bloków daje
\[
 X^{\mathsf T}J+JX
 =\begin{pmatrix}R-R^{\mathsf T}&P^{\mathsf T}+S\\
 -S^{\mathsf T}-P&Q^{\mathsf T}-Q\end{pmatrix}.
\]
Zatem $S=-P^{\mathsf T}$, a $Q,R$ są symetryczne.
Macierz $P$ ma $n^2$ dowolnych współczynników,
$Q,R$ po $n(n+1)/2$, skąd wskazany wymiar.
Przy $n=1$ mamy $\mathfrak{sl}_1=\mathfrak{so}_1=0$
i $\mathfrak{sp}_2=\mathfrak{sl}_2$.
\end{przyklad}

\section{Ideały, ilorazy i reprezentacje}

\begin{definicja}[Ideał i homomorfizm]
\label{aa:def:lie-ideal}
Ideałem algebry Liego $\mathfrak g$ nazywamy
podprzestrzeń $\mathfrak a\subset\mathfrak g$ taką,
że $[\mathfrak g,\mathfrak a]\subset\mathfrak a$.
Homomorfizm algebr Liego nad $K$ jest liniowym
odwzorowaniem $f:\mathfrak g\to\mathfrak h$ spełniającym
$f([x,y])=[f(x),f(y)]$ dla wszystkich $x,y$.
\end{definicja}

\begin{twierdzenie}[Algebra ilorazowa i twierdzenie o izomorfizmie]
\label{aa:tw:lie-quotient}
Dla ideału $\mathfrak a$ algebry Liego $\mathfrak g$
nawias $[x+\mathfrak a,y+\mathfrak a]=[x,y]+\mathfrak a$
określa algebrę Liego na przestrzeni ilorazowej.
Jeśli $f:\mathfrak g\to\mathfrak h$ jest homomorfizmem,
to $\ker f$ jest ideałem, $\operatorname{im}f$
podalgebrą, a odwzorowanie indukowane
$\mathfrak g/\ker f\to\operatorname{im}f$ jest
izomorfizmem algebr Liego.
\end{twierdzenie}
\begin{proof}
Zmiana reprezentantów o $u,v\in\mathfrak a$ zmienia
nawias o $[x,v]+[u,y]+[u,v]\in\mathfrak a$.
Nawias ilorazowy jest więc dobrze określony.
Dwuliniowość, alternowanie i tożsamość Jacobiego
wynikają z odpowiednich własności w $\mathfrak g$
po zastosowaniu rzutu ilorazowego.
Dla $u\in\ker f$ mamy $f([x,u])=[f(x),0]=0$,
a nawias dwóch elementów obrazu jest obrazem nawiasu.
Odwzorowanie $x+\ker f\mapsto f(x)$ jest niezależne
od reprezentanta, surjektywne na obraz i ma jądro
zerowe. Zachowuje działania liniowe i nawias
na mocy definicji, więc jest żądanym izomorfizmem.
\end{proof}

\begin{definicja}[Reprezentacja i moduł]
\label{aa:def:lie-representation}
Reprezentacją algebry Liego $\mathfrak g$ nad $K$
w przestrzeni wektorowej $V$ nad $K$ nazywamy
homomorfizm $\rho:\mathfrak g\to\mathfrak{gl}(V)$.
Równoważnie określamy dwuliniowe działanie
$x\cdot v=\rho(x)v$, spełniające
$[x,y]\cdot v=x\cdot(y\cdot v)-y\cdot(x\cdot v)$.
Przestrzeń $V$ z takim działaniem nazywamy
$\mathfrak g$-modułem. Podmoduł to podprzestrzeń
$U\subset V$ zachowana przez wszystkie $\rho(x)$.
Reprezentacja jest wierna, gdy $\ker\rho=0$.
\end{definicja}

Dla podmodułu $U$ działanie na $V/U$ określamy przez
$x\cdot(v+U)=x\cdot v+U$. Zmiana $v$ o element $U$
zmienia wynik o element $U$, więc działanie jest dobrze
określone. Równość definiującą moduł otrzymujemy
przez przejście do klas w odpowiedniej równości w $V$.
Reprezentacja naturalna dowolnej podalgebry
$\mathfrak{gl}(V)$ jest jej włożeniem w tę algebrę,
a zatem jest wierna.

\begin{definicja}[Centrum i reprezentacja dołączona]
\label{aa:def:adjoint-lie}
Dla algebry Liego $\mathfrak g$ definiujemy
$Z(\mathfrak g)=\{x\in\mathfrak g:[x,y]=0
\text{ dla każdego }y\in\mathfrak g\}$ oraz
$\operatorname{ad}_x(y)=[x,y]$ dla $x,y\in\mathfrak g$.
\end{definicja}

\begin{stwierdzenie}
\label{aa:prop:adjoint-lie}
Dla dowolnej algebry Liego $\mathfrak g$ odwzorowanie
$\operatorname{ad}:\mathfrak g\to\mathfrak{gl}(\mathfrak g)$
jest reprezentacją o jądrze $Z(\mathfrak g)$.
\end{stwierdzenie}
\begin{proof}
Liniowość wynika z dwuliniowości nawiasu.
Dla $x,y,z\in\mathfrak g$ tożsamość Jacobiego daje
\[
 (\operatorname{ad}_x\operatorname{ad}_y-
 \operatorname{ad}_y\operatorname{ad}_x)(z)
 =[x,[y,z]]-[y,[x,z]]=[[x,y],z].
\]
Zatem $\operatorname{ad}_{[x,y]}=
[\operatorname{ad}_x,\operatorname{ad}_y]$.
Warunek $\operatorname{ad}_x=0$ jest dokładnie
warunkiem przynależności do centrum.
\end{proof}

\section{Derywacje i automorfizmy}

\begin{definicja}[Derywacja i automorfizm]
\label{aa:def:lie-derivation}
Derywacją algebry Liego $\mathfrak g$ nazywamy
liniowe odwzorowanie $D:\mathfrak g\to\mathfrak g$
spełniające $D[x,y]=[Dx,y]+[x,Dy]$ dla wszystkich
$x,y\in\mathfrak g$. Zbiór derywacji oznaczamy
$\operatorname{Der}(\mathfrak g)$.
Automorfizmem nazywamy bijektywny homomorfizm
$\mathfrak g\to\mathfrak g$.
\end{definicja}

\begin{stwierdzenie}
\label{aa:prop:lie-derivations}
Dla dowolnej algebry Liego $\mathfrak g$ zbiór
$\operatorname{Der}(\mathfrak g)$ jest podalgebrą
$\mathfrak{gl}(\mathfrak g)$, a
$\operatorname{ad}(\mathfrak g)$ jest w nim ideałem.
\end{stwierdzenie}
\begin{proof}
Warunek derywacji jest liniowy względem operatora.
Dla derywacji $D,E$ rozwijamy
\[
\begin{aligned}
 DE[x,y]&=[DEx,y]+[Ex,Dy]+[Dx,Ey]+[x,DEy],\\
 ED[x,y]&=[EDx,y]+[Dx,Ey]+[Ex,Dy]+[x,EDy].
\end{aligned}
\]
Odejmując, otrzymujemy warunek derywacji dla $[D,E]$.
Tożsamość Jacobiego mówi, że $\operatorname{ad}_x$
jest derywacją. Dla dowolnej derywacji $D$ zachodzi
\[
 [D,\operatorname{ad}_x](y)=D[x,y]-[x,Dy]=[Dx,y].
\]
Zatem $[D,\operatorname{ad}_x]=\operatorname{ad}_{Dx}$,
co dowodzi własności ideału.
\end{proof}

\begin{stwierdzenie}[Eksponenta derywacji]
\label{aa:prop:exp-derivation}
Niech $\mathfrak g$ będzie skończenie wymiarową
algebrą Liego nad $\mathbb C$ lub $\mathbb R$.
Jeśli $D$ jest jej derywacją, to $\exp(tD)$ jest
automorfizmem dla każdego $t\in\mathbb R$.
\end{stwierdzenie}
\begin{proof}
W skończonym wymiarze szereg
$\sum_{k\ge0}t^kD^k/k!$ zbiega w normie operatorowej,
bo normy wyrazów są ograniczone przez
$|t|^k\|D\|^k/k!$. Mnożenie dwóch bezwzględnie
zbieżnych szeregów dla $t,-t$ daje $I$:
w stopniu $k\ge1$ współczynnik wynosi
$\sum_{j=0}^k(-1)^j/(j!(k-j)!)=0$.
Zatem $\exp(tD)$ jest odwracalne.
Indukcja z użyciem własności derywacji i tożsamości
Pascala daje dla każdego $k\ge0$
\[
 D^k[x,y]=\sum_{j=0}^k\binom{k}{j}[D^jx,D^{k-j}y].
\]
Dla $k=0$ jest to równość identyczna. Po zastosowaniu
$D$ każdy składnik daje dwa składniki, których
współczynniki w stopniu $k+1$ sumują się według
tożsamości Pascala, także na końcach.
Nawias jako odwzorowanie dwuliniowe w skończonym
wymiarze jest ciągły i spełnia
$\|[u,v]\|\le c\|u\|\|v\|$ dla pewnego $c$.
Szereg podwójny powstały po rozwinięciu
$[\exp(tD)x,\exp(tD)y]$ jest więc bezwzględnie
zbieżny (majorantą jest $c\|x\|\|y\|e^{2|t|\|D\|}$).
Grupując wyrazy według sumy indeksów i korzystając
z wyprowadzonego wzoru, otrzymujemy
$[\exp(tD)x,\exp(tD)y]=\exp(tD)[x,y]$.
\end{proof}

\section{Nilpotentność i rozwiązalność}

\begin{definicja}[Szeregi algebry Liego]
\label{aa:def:lie-series}
Dla algebry Liego $\mathfrak g$ określamy dolny
szereg centralny przez $\gamma_1(\mathfrak g)=\mathfrak g$,
$\gamma_{k+1}(\mathfrak g)=[\mathfrak g,\gamma_k(\mathfrak g)]$
dla $k\ge1$, a szereg pochodny przez
$\mathfrak g^{(0)}=\mathfrak g$,
$\mathfrak g^{(k+1)}=[\mathfrak g^{(k)},\mathfrak g^{(k)}]$
dla $k\ge0$.
Algebra jest nilpotentna, jeśli pewne $\gamma_k=0$,
i rozwiązalna, jeśli pewne $\mathfrak g^{(k)}=0$.
Niezerowa algebra nilpotentna ma klasę $c$, jeśli
$\gamma_c\ne0$ i $\gamma_{c+1}=0$; algebra zerowa
ma klasę zero.
\end{definicja}

\begin{lemat}
\label{aa:lem:lie-series}
Oba szeregi z Definicji~\ref{aa:def:lie-series}
są malejącymi ciągami ideałów. Dla $i,j\ge1$ zachodzi
$[\gamma_i,\gamma_j]\subset\gamma_{i+j}$.
W szczególności $\mathfrak g^{(k)}\subset\gamma_{2^k}$,
więc każda algebra nilpotentna jest rozwiązalna.
\end{lemat}
\begin{proof}
Jeśli $U,W$ są ideałami, to dla $x\in\mathfrak g$,
$u\in U,w\in W$ mamy
$[x,[u,w]]=[[x,u],w]+[u,[x,w]]\in[U,W]$.
Z liniowości wynika, że $[U,W]$ jest ideałem,
zawartym w $U\cap W$. Indukcja daje pierwszą tezę.
Dla ustalonego $i$ dowodzimy drugiej tezy jednocześnie
dla wszystkich $j$ przez indukcję względem $i$.
Dla $i=1$ jest to definicja. Jeśli teza zachodzi
dla $i$ i wszystkich $j$, to z Jacobiego
\[
 [[\mathfrak g,\gamma_i],\gamma_j]
 \subset[\mathfrak g,[\gamma_i,\gamma_j]]
       +[\gamma_i,[\mathfrak g,\gamma_j]]
 \subset\gamma_{i+j+1}.
\]
Zatem zachodzi dla $i+1$. Dla szeregu pochodnego
indukcja od $\mathfrak g^{(0)}=\gamma_1$ daje
$\mathfrak g^{(k+1)}\subset
[\gamma_{2^k},\gamma_{2^k}]\subset\gamma_{2^{k+1}}$.
Jeśli dolny szereg centralny osiąga zero, to przez
malejący charakter ciągu odpowiedni wyraz szeregu
pochodnego jest również zerowy.
\end{proof}

\begin{stwierdzenie}[Podalgebry i obrazy]
\label{aa:prop:lie-series-images}
Podalgebra algebry nilpotentnej (odpowiednio rozwiązalnej)
jest nilpotentna (odpowiednio rozwiązalna).
Jeśli $f:\mathfrak g\to\mathfrak h$ jest surjektywnym
homomorfizmem, to $f(\gamma_k(\mathfrak g))=
\gamma_k(\mathfrak h)$ oraz $f(\mathfrak g^{(k)})=
\mathfrak h^{(k)}$. W szczególności ilorazy przez
ideały dziedziczą obie własności.
\end{stwierdzenie}
\begin{proof}
Dla podalgebry $\mathfrak b$ indukcja daje
$\gamma_k(\mathfrak b)\subset\gamma_k(\mathfrak g)$
oraz $\mathfrak b^{(k)}\subset\mathfrak g^{(k)}$:
nawias zachowuje inkluzje w obu argumentach.
Dla homomorfizmu i podprzestrzeni $U,W$ zachodzi
$f([U,W])=[f(U),f(W)]$, gdyż liniowość $f$
przeprowadza rozpiętość na rozpiętość obrazów,
a nawias na nawias. Stosując to indukcyjnie
i korzystając z surjektywności w pierwszym wyrazie,
otrzymujemy obie równości.
\end{proof}

\begin{przyklad}[Macierze trójkątne]
\label{aa:prz:triangular-lie}
Niech $\mathfrak b_n$ będzie algebrą macierzy
górnotrójkątnych nad $K$, a $\mathfrak n_n$
jej podprzestrzenią macierzy ściśle górnotrójkątnych.
Obie są zamknięte na iloczyny i komutatory.
Dla $r\ge1$ niech $F_r$ będzie rozpiętością
$E_{ij}$ z $j-i\ge r$; dla $r\ge n$ mamy $F_r=0$.
Wzór na iloczyny jednostek macierzowych daje
$F_rF_s\subset F_{r+s}$ (niezerowy iloczyn wymaga
zgodności środkowych indeksów, a różnice indeksów
wówczas się dodają). Zatem $[F_r,F_s]\subset F_{r+s}$.
Mamy $\gamma_1(\mathfrak n_n)=F_1$. Dla $r\ge1$
każda jednostka $E_{ij}$ z $j-i\ge r+1$ spełnia
$E_{ij}=[E_{i,i+1},E_{i+1,j}]$, gdzie drugi
czynnik należy do $F_r$. To daje
$[F_1,F_r]=F_{r+1}$ i indukcyjnie
$\gamma_r(\mathfrak n_n)=F_r$. Stąd dla $n\ge2$
klasa nilpotentności wynosi $n-1$, a dla $n=1$
algebra jest zerowa.

Komutator dwóch macierzy górnotrójkątnych ma zerową
przekątną, bo na $i$-tym wpisie przekątnej iloczynu
pozostaje tylko iloczyn dwóch $i$-tych wpisów.
Ponadto dla $i<j$ mamy $[E_{ii},E_{ij}]=E_{ij}$,
więc $[\mathfrak b_n,\mathfrak b_n]=\mathfrak n_n$.
Z lematu~\ref{aa:lem:lie-series} algebra
$\mathfrak b_n$ jest rozwiązalna.
Dla $n\ge2$ nie jest nilpotentna: skoro
$E_{12}=[E_{11},E_{12}]$, indukcja daje
$E_{12}\in\gamma_k(\mathfrak b_n)$ dla każdego $k$.
\end{przyklad}

\begin{przyklad}[Algebra Heisenberga i algebra $\mathfrak{sl}_2$]
\label{aa:prz:heisenberg-sl2}
Nad ciałem charakterystyki zero algebra
$\mathfrak n_3$ z bazą $x=E_{12},y=E_{23},z=E_{13}$
spełnia $[x,y]=z$, $[x,z]=[y,z]=0$.
Ma centrum $Kz$: dla $u=ax+by+cz$ warunki
$[u,x]=-bz=0$, $[u,y]=az=0$ dają $a=b=0$.
Jej dolny szereg centralny to
$\mathfrak n_3\supset Kz\supset0$.
Dla porównania $\mathfrak{sl}_2$ ma bazę
$h=E_{11}-E_{22}$, $e=E_{12}$, $f=E_{21}$,
z nawiasami $[h,e]=2e$, $[h,f]=-2f$, $[e,f]=h$.
Każdy element tej bazy należy do
$[\mathfrak{sl}_2,\mathfrak{sl}_2]$, więc szereg
pochodny jest stały i algebra nie jest rozwiązalna.
W szczególności nie jest nilpotentna.
\end{przyklad}

\begin{center}
\begin{tikzpicture}[node distance=6mm, every node/.style={font=\small}]
 \node (a) at (0,0) {$\mathfrak n_3$};
 \node (b) at (0,-.9) {$Kz$};
 \node (c) at (0,-1.8) {$0$};
 \draw[-{Stealth}] (a)--node[right] {$[\mathfrak n_3,\cdot]$} (b);
 \draw[-{Stealth}] (b)-- (c);
 \node (d) at (5,0) {$\mathfrak b_2$};
 \node (e) at (5,-.9) {$KE_{12}$};
 \node (f) at (5,-1.8) {$0$};
 \draw[-{Stealth}] (d)--node[right] {$[\mathfrak b_2,\mathfrak b_2]$} (e);
 \draw[-{Stealth}] (e)-- (f);
 \draw[-{Stealth}] (e.west) to[out=170,in=190,looseness=5]
 node[left] {$[\mathfrak b_2,\cdot]$} (e.west);
\end{tikzpicture}

\small Dolny szereg centralny algebry Heisenberga (po lewej).
Szereg pochodny $\mathfrak b_2$ osiąga zero,
ale jej dolny szereg centralny zatrzymuje się na $KE_{12}$.
\end{center}

\section{Twierdzenie Engela}

\begin{lemat}[Komutator z operatorem nilpotentnym]
\label{aa:lem:engel-ad}
Niech $V$ będzie przestrzenią wektorową nad $K$,
a $X\in\operatorname{End}_K(V)$ niech spełnia $X^m=0$
dla pewnego $m\ge1$. Wówczas operator
$\operatorname{ad}_X$ na $\mathfrak{gl}(V)$ jest
nilpotentny, a dokładniej $(\operatorname{ad}_X)^{2m-1}=0$.
\end{lemat}
\begin{proof}
Operatory $L_X(Y)=XY$ i $R_X(Y)=YX$ komutują
z powodu łączności składania. Mamy
$\operatorname{ad}_X=L_X-R_X$.
Ze wzoru dwumianowego dla komutujących operatorów
\[
 (\operatorname{ad}_X)^k(Y)
 =\sum_{j=0}^k(-1)^j\binom{k}{j}X^{k-j}YX^j.
\]
Jeśli $k=2m-1$, to dla każdego $j$ co najmniej
jedna z liczb $j,k-j$ jest nie mniejsza niż $m$.
Każdy składnik znika.
\end{proof}

\begin{twierdzenie}[Twierdzenie Engela dla operatorów]
\label{aa:tw:engel-linear}
Niech $V\ne0$ będzie skończenie wymiarową przestrzenią
wektorową nad dowolnym ciałem $K$, a $\mathfrak l$
podalgebrą $\mathfrak{gl}(V)$. Załóżmy, że każdy
$X\in\mathfrak l$ jest operatorem nilpotentnym na $V$.
Wtedy istnieje $v\in V\setminus\{0\}$ takie, że
$Xv=0$ dla wszystkich $X\in\mathfrak l$.
\end{twierdzenie}
\begin{proof}
Indukcja przebiega względem wymiaru $\mathfrak l$,
jednocześnie dla wszystkich niezerowych skończenie
wymiarowych przestrzeni reprezentacji.
Dla $\mathfrak l=0$ wybieramy dowolny $v\ne0$.
Niech $\dim\mathfrak l>0$ i załóżmy tezę dla
wszystkich mniejszych wymiarów algebr operatorów.

Wybierzmy właściwą podalgebrę $\mathfrak h$ o największym
możliwym wymiarze; istnieje, bo podalgebra zerowa
jest właściwa, a wymiary są liczbami całkowitymi
ograniczonymi przez $\dim\mathfrak l-1$.
Na niezerowej przestrzeni $\mathfrak l/\mathfrak h$
(jedynie ilorazie przestrzeni wektorowych na tym etapie)
określamy działanie
\[
 Y\cdot(X+\mathfrak h)=[Y,X]+\mathfrak h,
 \qquad Y\in\mathfrak h,\ X\in\mathfrak l.
\]
Jest dobrze określone, ponieważ $\mathfrak h$
jest podalgebrą. Tożsamość Jacobiego daje
warunek reprezentacji. Dla $Y\in\mathfrak h$
operator $\operatorname{ad}_Y$ jest nilpotentny
na $\mathfrak{gl}(V)$ na mocy
Lematu~\ref{aa:lem:engel-ad}; jego ograniczenie
do $\mathfrak l$ zachowuje $\mathfrak h$,
więc operator indukowany na ilorazie jest nilpotentny.
Obraz reprezentacji ma wymiar co najwyżej
$\dim\mathfrak h<\dim\mathfrak l$.
Można więc zastosować założenie indukcyjne
do tego obrazu na $\mathfrak l/\mathfrak h$.
Istnieje $X\in\mathfrak l\setminus\mathfrak h$
takie, że $[Y,X]\in\mathfrak h$ dla każdego
$Y\in\mathfrak h$.

Normalizator $N=\{Z\in\mathfrak l:[Z,\mathfrak h]
\subset\mathfrak h\}$ jest podalgebrą:
dla $Z_1,Z_2\in N,Y\in\mathfrak h$ mamy
$[[Z_1,Z_2],Y]=[Z_1,[Z_2,Y]]-[Z_2,[Z_1,Y]]
\in\mathfrak h$, a warunek jest liniowy.
Zawiera $\mathfrak h$ i $X\notin\mathfrak h$,
więc maksymalność wymiaru $\mathfrak h$ wymusza
$N=\mathfrak l$. Oznacza to, że $\mathfrak h$
jest ideałem. Iloraz $\mathfrak l/\mathfrak h$
ma wymiar jeden: gdyby miał wymiar większy,
dowolna prosta w nim byłaby właściwą niezerową
podalgebrą (nawias na prostej jest zerowy),
a jej przeciwobraz byłby właściwą podalgebrą
ściśle zawierającą $\mathfrak h$.

Z założenia indukcyjnego dla $\mathfrak h$ na $V$
przestrzeń $U=\{v\in V:Yv=0\ \forall Y\in\mathfrak h\}$
jest niezerowa. Wybieramy $X\notin\mathfrak h$;
mamy $\mathfrak l=\mathfrak h+KX$.
Dla $u\in U,Y\in\mathfrak h$ zachodzi
$Y(Xu)=X(Yu)+[Y,X]u=0$,
gdyż $Yu=0$ i $[Y,X]\in\mathfrak h$.
Zatem $X$ zachowuje $U$, a jego ograniczenie
do $U$ jest nilpotentne. Ma niezerowe jądro:
gdyby było iniektywne, każda jego potęga byłaby
iniektywna, co przeczy zerowej potędze na $U\ne0$.
Wybieramy $v\ne0$ w tym jądrze. Wówczas
$\mathfrak h v=0$ i $Xv=0$, więc $\mathfrak l v=0$.
\end{proof}

\begin{wniosek}[Wspólna postać ściśle trójkątna]
\label{aa:cor:engel-flag}
Przy założeniach Twierdzenia~\ref{aa:tw:engel-linear}
istnieje pełna flaga podprzestrzeni
$0=V_0\subset V_1\subset\cdots\subset V_d=V$,
$\dim V_i=i$, gdzie $d=\dim V$, taka, że
$X V_i\subset V_{i-1}$ dla każdego $X\in\mathfrak l$
i $1\le i\le d$. W bazie zgodnej z flagą wszystkie
macierze $X$ są ściśle górnotrójkątne.
\end{wniosek}
\begin{proof}
Indukcja po $d$; dla $d=0$ flaga jest trywialna.
Dla $d\ge1$ twierdzenie daje prostą $V_1=Kv$
anihilowaną przez $\mathfrak l$.
Na $V/V_1$ każdy indukowany operator jest nilpotentny,
ponieważ zerowa potęga operatora przechodzi na iloraz.
Obraz $\mathfrak l$ jest podalgebrą operatorów,
więc indukcja stosuje się do ilorazu wymiaru $d-1$.
Przeciwobrazy kolejnych podprzestrzeni flagi ilorazowej
dają $V_1,\ldots,V_d$, a inkluzje obrazów implikują
$XV_i\subset V_{i-1}$ dla $i\ge2$.
Dla $i=1$ wynika to z wyboru $v$.
Wektor bazowy $v_i\in V_i\setminus V_{i-1}$
ma obraz w rozpiętości wcześniejszych wektorów,
co daje żądaną postać macierzy.
\end{proof}

\begin{twierdzenie}[Twierdzenie Engela]
\label{aa:tw:engel}
Niech $\mathfrak g$ będzie skończenie wymiarową
algebrą Liego nad dowolnym ciałem $K$.
Wówczas $\mathfrak g$ jest nilpotentna wtedy
i tylko wtedy, gdy $\operatorname{ad}_x$ jest
operatorem nilpotentnym dla każdego $x\in\mathfrak g$.
\end{twierdzenie}
\begin{proof}
Jeśli $\gamma_{c+1}=0$, to
$\operatorname{ad}_x(\gamma_i)\subset\gamma_{i+1}$,
więc $(\operatorname{ad}_x)^c(\mathfrak g)=0$.
Przypadek $\mathfrak g=0$ jest trywialny.
Odwrotnie, załóżmy nilpotentność wszystkich
$\operatorname{ad}_x$. Na mocy
Stwierdzenia~\ref{aa:prop:adjoint-lie} ich zbiór jest
podalgebrą $\mathfrak{gl}(\mathfrak g)$.
Z Wniosku~\ref{aa:cor:engel-flag} dla $\mathfrak g$
wymiaru $d>0$ otrzymujemy flagę $V_i$ z warunkiem
$[\mathfrak g,V_i]\subset V_{i-1}$.
Ponieważ $\gamma_1=\mathfrak g=V_d$, indukcja daje
$\gamma_j\subset V_{d+1-j}$ dla $1\le j\le d+1$.
W szczególności $\gamma_{d+1}=0$.
\end{proof}

\section{Twierdzenie Liego}

\begin{twierdzenie}[Twierdzenie Liego o wspólnym wektorze własnym]
\label{aa:tw:lie-eigenvector}
Niech $K$ będzie algebraicznie domkniętym ciałem
charakterystyki zero, $\mathfrak g$ skończenie
wymiarową rozwiązalną algebrą Liego nad $K$,
a $\rho:\mathfrak g\to\mathfrak{gl}(V)$ reprezentacją
w niezerowej skończenie wymiarowej przestrzeni $V$.
Wtedy istnieją $v\ne0$ i funkcjonał liniowy
$\lambda:\mathfrak g\to K$ takie, że
$\rho(x)v=\lambda(x)v$ dla każdego $x\in\mathfrak g$.
\end{twierdzenie}
\begin{proof}
Indukcja względem $\dim\mathfrak g$ jednocześnie
dla wszystkich takich reprezentacji. Dla algebry
zerowej bierzemy dowolny wektor niezerowy i funkcjonał
zerowy. Załóżmy $\dim\mathfrak g>0$.
Rozwiązalność wymusza $[\mathfrak g,\mathfrak g]
\ne\mathfrak g$, gdyż inaczej szereg pochodny
byłby stały. Wybieramy hiperpłaszczyznę $\mathfrak h$
zawierającą $[\mathfrak g,\mathfrak g]$ (uzupełniamy
bazę tej właściwej podprzestrzeni do bazy $\mathfrak g$
i pomijamy jeden z nowych wektorów).
Jest ideałem, bo $[\mathfrak g,\mathfrak h]$
zawiera się w $[\mathfrak g,\mathfrak g]$.
Na mocy Stwierdzenia~\ref{aa:prop:lie-series-images}
jest rozwiązalna, więc do ograniczenia reprezentacji
stosujemy założenie indukcyjne.
Istnieją $v\ne0$ i $\mu:\mathfrak h\to K$ liniowe
takie, że $\rho(y)v=\mu(y)v$ dla $y\in\mathfrak h$.

Wybieramy $x\in\mathfrak g\setminus\mathfrak h$
i oznaczamy $T=\rho(x)$. Pokażemy najpierw
$\mu([x,y])=0$ dla wszystkich $y\in\mathfrak h$.
Niech $W_j=\operatorname{span}(v,Tv,\ldots,T^jv)$
dla $j\ge0$, $W_{-1}=0$.
Dla wszystkich $y\in\mathfrak h$ wykażemy indukcyjnie
\[
 \rho(y)T^jv-\mu(y)T^jv\in W_{j-1}.
\]
Dla $j=0$ wynika to z wyboru $v$. Jeśli teza zachodzi
dla wykładnika $j-1$ i wszystkich $y\in\mathfrak h$,
to
\[
 \rho(y)T^jv=T\rho(y)T^{j-1}v+\rho([y,x])T^{j-1}v.
\]
Pierwszy składnik różni się od $\mu(y)T^jv$
o element $TW_{j-2}\subset W_{j-1}$.
Drugi należy do $W_{j-1}$, bo $[y,x]\in\mathfrak h$
i stosujemy to samo założenie indukcyjne dla $[y,x]$.
Daje to żądaną inkluzję.

Niech $r\ge1$ będzie najmniejszą liczbą, dla której
$T^rv$ należy do $\operatorname{span}(v,\ldots,T^{r-1}v)$;
istnieje dzięki skończonemu wymiarowi $V$.
Przestrzeń $W=W_{r-1}$ ma bazę $v,\ldots,T^{r-1}v$
i jest zachowana przez $T$ oraz wszystkie $\rho(y)$
dla $y\in\mathfrak h$ na mocy wykazanej inkluzji.
W tej bazie $\rho(y)|_W$ ma postać górnotrójkątną
z przekątną stale równą $\mu(y)$.
W szczególności dla $[x,y]\in\mathfrak h$
\[
 r\mu([x,y])=\operatorname{tr}(\rho([x,y])|_W)
 =\operatorname{tr}([T|_W,\rho(y)|_W])=0.
\]
Ostatnia równość wynika z cykliczności śladu,
wykazanej w Przykładzie~\ref{aa:prz:classical-lie}.
Charakterystyka zero gwarantuje $r\ne0$ w $K$,
więc $\mu([x,y])=0$.

Rozważmy teraz całą niezerową przestrzeń
$U=\{u\in V:\rho(y)u=\mu(y)u\ \forall y\in\mathfrak h\}$,
zawierającą $v$. Dla $u\in U,y\in\mathfrak h$
mamy
\[
 \rho(y)Tu=T\rho(y)u+\rho([y,x])u
 =\mu(y)Tu+\mu([y,x])u=\mu(y)Tu.
\]
Zatem $T$ zachowuje $U$. Ponieważ $K$ jest
algebraicznie domknięte, wielomian charakterystyczny
$T|_U$ dodatniego stopnia ma pierwiastek $\nu\in K$;
operator $T|_U-\nu I_U$ ma zerowy wyznacznik,
więc istnieje $w\in U\setminus\{0\}$ z $Tw=\nu w$.
Dla $y+ax\in\mathfrak h+Kx=\mathfrak g$ mamy
$\rho(y+ax)w=(\mu(y)+a\nu)w$.
Jednoznaczność zapisu $y+ax$ daje dobrze określony
liniowy funkcjonał $\lambda(y+ax)=\mu(y)+a\nu$.
\end{proof}

\begin{wniosek}[Wspólna postać trójkątna]
\label{aa:cor:lie-flag}
Przy założeniach Twierdzenia~\ref{aa:tw:lie-eigenvector}
istnieje baza $V$, w której wszystkie operatory
$\rho(x)$ są górnotrójkątne. Ponadto
$\rho([\mathfrak g,\mathfrak g])$ składa się
z macierzy ściśle górnotrójkątnych.
\end{wniosek}
\begin{proof}
Indukcja po wymiarze $V$, przy wymiarze zero
nie ma nic do wykazania. Twierdzenie Liego daje
niezmienniczą prostą $Kv$. Na ilorazie $V/Kv$
mamy reprezentację tej samej rozwiązalnej algebry.
Stosując indukcję, otrzymujemy pełną flagę
niezmienniczą na ilorazie. Jej przeciwobrazy
wraz z $0$ tworzą pełną flagę niezmienniczą na $V$.
Baza zgodna z flagą daje postać górnotrójkątną.
Komutatory macierzy górnotrójkątnych mają zerową
przekątną, jak wykazano w
Przykładzie~\ref{aa:prz:triangular-lie}.
Ponieważ $\rho$ zachowuje nawias i sumy liniowe,
teza o całym obrazie algebry pochodnej wynika
z zamkniętości przestrzeni macierzy ściśle
górnotrójkątnych na kombinacje liniowe.
\end{proof}

\begin{wniosek}[Nilpotentność algebry pochodnej]
\label{aa:cor:lie-derived-nilpotent}
Niech $\mathfrak g$ będzie skończenie wymiarową
rozwiązalną algebrą Liego nad algebraicznie domkniętym
ciałem charakterystyki zero. Wówczas ideał
$\mathfrak a=[\mathfrak g,\mathfrak g]$ jest nilpotentny.
\end{wniosek}
\begin{proof}
Stosujemy Wniosek~\ref{aa:cor:lie-flag} do reprezentacji
dołączonej $\operatorname{ad}$ na $\mathfrak g$.
Jeśli $\mathfrak g=0$, teza jest natychmiastowa.
Dla każdego $x\in\mathfrak a$ operator $\operatorname{ad}_x$
jest ściśle górnotrójkątny w pewnej wspólnej bazie,
a zatem nilpotentny: jego $d$-ta potęga jest zerowa,
gdy $d=\dim\mathfrak g$, na mocy własności filtracji
$F_r$ z Przykładu~\ref{aa:prz:triangular-lie}.
Ograniczenie do $\mathfrak a$ jest dobrze określone
(ideał) i nilpotentne. Jest to właśnie operator
dołączony elementu $x$ w algebrze $\mathfrak a$.
Twierdzenie~\ref{aa:tw:engel} daje nilpotentność
$\mathfrak a$.
\end{proof}

\begin{przyklad}[Znaczenie założeń twierdzenia Liego]
\label{aa:prz:lie-hypotheses}
Nad $\mathbb R$ jednowymiarowa abelowa algebra
operatorów generowana przez
$\begin{pmatrix}0&-1\\1&0\end{pmatrix}$
nie ma rzeczywistego niezerowego wektora własnego:
wielomian charakterystyczny $t^2+1$ nie ma pierwiastka
w $\mathbb R$.

W dodatniej charakterystyce $p$ weźmy algebraicznie
domknięte ciało $K$ i $V=K[t]/(t^p)$.
Operator $X$ mnożenia przez klasę $t$ i operator
$Y$ formalnego różniczkowania są dobrze określone:
pochodna $t^pq$ to $t^pq'$ w charakterystyce $p$.
Na klasie dowolnego wielomianu mamy
$[Y,X]v=Y(tv)-tY(v)=v$, więc $[Y,X]=I$.
Przestrzeń $\mathfrak l=\operatorname{span}(X,Y,I)$
jest algebrą Liego o algebrze pochodnej $KI$
i kolejnym wyrazie zerowym, a więc jest rozwiązalna.
Gdyby $v\ne0$ był wspólnym wektorem własnym $X,Y$,
to $[Y,X]v=0$ (iloczyny skalarów komutują),
co przeczyłoby $[Y,X]v=v$.
Brak zatem wspólnego wektora własnego, mimo
algebraicznej domkniętości ciała.
\end{przyklad}

Twierdzenie Engela dotyczy nilpotentności operatorów,
twierdzenie Liego dotyczy rozwiązalności algebry.
Nie należy utożsamiać tych założeń: algebra abelowa
generowana przez $I$ jest nilpotentna jako algebra
Liego, ale $I$ na niezerowej przestrzeni nie jest
operatorem nilpotentnym. Wniosek o ściśle trójkątnej
postaci z twierdzenia Engela wymaga nilpotentności
każdego operatora w rozważanej reprezentacji.
Kryteria półprostoty i związek z układami korzeni
będą przedmiotem dalszego rozdziału.

% AA-58
\chapter{Półproste algebry Liego. Forma Killinga i kryteria Cartana}
\label{aa:ch:semisimple-lie}

W rozdziale~\ref{aa:ch:lie-algebras} rozróżniliśmy
algebry nilpotentne i rozwiązalne. Teraz wyodrębnimy
największy ideał rozwiązalny i zbadamy algebry,
w których ideał ten jest zerowy. Kluczowym narzędziem
będzie forma dwuliniowa zbudowana ze śladów operatorów
reprezentacji dołączonej. W całym rozdziale algebry
Liego są skończenie wymiarowe.

\section{Radykał rozwiązalny}

\begin{lemat}[Rozwiązalność rozszerzenia]
\label{aa:lem:solvable-extension}
Niech $\mathfrak g$ będzie algebrą Liego nad ciałem $K$,
a $\mathfrak a$ jej ideałem. Jeśli $\mathfrak a$ oraz
$\mathfrak g/\mathfrak a$ są rozwiązalne, to
$\mathfrak g$ jest rozwiązalna.
\end{lemat}
\begin{proof}
Niech $q:\mathfrak g\to\mathfrak g/\mathfrak a$
będzie rzutem ilorazowym. Z rozwiązalności ilorazu
istnieje $r\ge0$, dla którego jego $r$-ty wyraz
szeregu pochodnego jest zerowy. Na mocy
Stwierdzenia~\ref{aa:prop:lie-series-images}
$q(\mathfrak g^{(r)})=(\mathfrak g/\mathfrak a)^{(r)}=0$,
czyli $\mathfrak g^{(r)}\subset\mathfrak a$.
Indukcja po $s\ge0$ daje
$\mathfrak g^{(r+s)}\subset\mathfrak a^{(s)}$:
dla $s=0$ już to wykazano, a przejście do $s+1$
wynika z wzięcia nawiasu obu stron z samymi sobą.
Ponieważ pewne $\mathfrak a^{(s)}=0$, otrzymujemy
$\mathfrak g^{(r+s)}=0$.
\end{proof}

\begin{twierdzenie}[Istnienie radykału rozwiązalnego]
\label{aa:tw:lie-radical}
Niech $\mathfrak g$ będzie skończenie wymiarową algebrą
Liego nad ciałem $K$. Istnieje dokładnie jeden
największy (względem inkluzji) ideał rozwiązalny
$\operatorname{rad}\mathfrak g$.
Każdy automorfizm $\mathfrak g$ zachowuje ten ideał.
Iloraz $\mathfrak g/\operatorname{rad}\mathfrak g$
nie ma niezerowych ideałów rozwiązalnych.
\end{twierdzenie}
\begin{proof}
Niech $\mathfrak a,\mathfrak b$ będą ideałami
rozwiązalnymi. Ich suma algebraiczna jest ideałem,
bo $[x,a+b]=[x,a]+[x,b]\in\mathfrak a+\mathfrak b$.
Odwzorowanie
$\mathfrak b\to(\mathfrak a+\mathfrak b)/\mathfrak a$,
$b\mapsto b+\mathfrak a$, jest surjektywnym
homomorfizmem: klasa $a+b$ jest klasą $b$.
Iloraz jest rozwiązalny na mocy
Stwierdzenia~\ref{aa:prop:lie-series-images}.
Lemat~\ref{aa:lem:solvable-extension} zastosowany
do $\mathfrak a+\mathfrak b$ i jej ideału
$\mathfrak a$ dowodzi rozwiązalności sumy.

Wybieramy ideał rozwiązalny $\mathfrak r$ o największym
wymiarze; istnieje, ponieważ ideał zerowy jest
rozwiązalny, a zbiór możliwych wymiarów jest skończony.
Dla dowolnego ideału rozwiązalnego $\mathfrak b$
suma $\mathfrak r+\mathfrak b$ jest rozwiązalna,
więc jej wymiar nie przekracza $\dim\mathfrak r$.
Z inkluzji $\mathfrak r\subset\mathfrak r+\mathfrak b$
wynika równość, a stąd $\mathfrak b\subset\mathfrak r$.
To dowodzi istnienia i jednoznaczności.
Jeśli $f$ jest automorfizmem, $f(\mathfrak r)$
jest ideałem (dla $x=f(y)$ i $a\in\mathfrak r$
mamy $[x,f(a)]=f([y,a])$) oraz jest rozwiązalny
jako obraz $\mathfrak r$. Z maksymalności
$f(\mathfrak r)\subset\mathfrak r$ i z równości
wymiarów wynika równość.

Niech $\mathfrak c$ będzie ideałem rozwiązalnym
w $\mathfrak g/\mathfrak r$. Jego przeciwobraz
$\mathfrak d$ przez rzut $q$ jest ideałem, bo
$q([x,d])=[q(x),q(d)]\in\mathfrak c$. Zawiera
$\mathfrak r$, a $\mathfrak d/\mathfrak r\simeq
\mathfrak c$ na mocy twierdzenia o izomorfizmie
(Twierdzenie~\ref{aa:tw:lie-quotient}).
Lemat~\ref{aa:lem:solvable-extension} daje rozwiązalność
$\mathfrak d$, stąd $\mathfrak d\subset\mathfrak r$,
czyli $\mathfrak c=0$.
\end{proof}

\begin{definicja}[Prostota i półprostota]
\label{aa:def:semisimple-lie}
Algebra Liego $\mathfrak g$ nad ciałem $K$ jest prosta,
jeśli jest nieabelowa i jej jedynymi ideałami są
$0$ oraz $\mathfrak g$. Skończenie wymiarowa algebra
jest półprosta, jeśli $\operatorname{rad}\mathfrak g=0$.
\end{definicja}

Algebra prosta jest półprosta: gdyby była rozwiązalna,
jej niezerowa algebra pochodna, będąca ideałem,
musiałaby być równa całej algebrze, co przeczy
rozwiązalności. Algebra zerowa jest półprosta,
lecz nie prosta. Każda półprosta algebra ma zerowe
centrum, ponieważ centrum jest abelowym ideałem
(na mocy Stwierdzenia~\ref{aa:prop:adjoint-lie} jest
jądrem reprezentacji dołączonej).

\begin{przyklad}[Radykał w przykładach z AA-57]
\label{aa:prz:radical-examples}
Nad ciałem charakterystyki zero algebra Heisenberga
i algebra $\mathfrak b_n$ macierzy górnotrójkątnych
są rozwiązalne, więc ich radykały są całymi algebrami
(Przykłady~\ref{aa:prz:triangular-lie}
i~\ref{aa:prz:heisenberg-sl2}).
Algebra $\mathfrak{sl}_2$ jest prosta. Istotnie,
w bazie $h,e,f$ z Przykładu~\ref{aa:prz:heisenberg-sl2}
operator $\operatorname{ad}_h$ ma wartości własne
$0,2,-2$. Dla dowolnego niezerowego ideału
$\mathfrak a$ i niezerowego $u\in\mathfrak a$
każdy wielomian w $\operatorname{ad}_h$ zachowuje
$\mathfrak a$. Projekcje na trzy proste własne
są takimi wielomianami: dla wartości własnej $\lambda$
wystarczy użyć $\prod_{\mu\ne\lambda}
(t-\mu)/(\lambda-\mu)$, gdzie $\mu$ przebiega
pozostałe dwie wartości. Mianowniki są niezerowe
w charakterystyce zero. Co najmniej jedna projekcja
$u$ jest niezerowa, więc $\mathfrak a$ zawiera
co najmniej jeden z $h,e,f$. Jeśli zawiera $e$ lub $f$,
nawias z drugim daje $h$ lub $-h$. Jeśli zawiera $h$,
nawiasy z $e,f$ dają $2e,-2f$. W każdym przypadku
$\mathfrak a$ zawiera bazę, zatem $\mathfrak a=
\mathfrak{sl}_2$. Radykał $\mathfrak{sl}_2$ jest zerowy.
\end{przyklad}

\begin{center}
\begin{tikzpicture}[every node/.style={font=\small}]
 \node[draw,rounded corners,minimum width=32mm,minimum height=10mm] (g) at (0,0) {$\mathfrak g$};
 \node[draw,rounded corners,minimum width=32mm,minimum height=10mm] (r) at (-4,0) {$\operatorname{rad}\mathfrak g$};
 \node[draw,rounded corners,minimum width=32mm,minimum height=10mm] (q) at (4,0) {$\mathfrak g/\operatorname{rad}\mathfrak g$};
 \draw[-{Stealth}] (r)--node[above]{inkluzja} (g);
 \draw[-{Stealth}] (g)--node[above]{rzut} (q);
 \node at (-4,-.8) {rozwiązalna};
 \node at (4,-.8) {półprosta};
\end{tikzpicture}

\small Radykał rozwiązalny i iloraz półprosty.
Rysunek nie zakłada rozkładu $\mathfrak g$ na sumę ideałów.
\end{center}

\section{Formy niezmiennicze i forma Killinga}

\begin{definicja}[Forma niezmiennicza]
\label{aa:def:invariant-lie-form}
Niech $\mathfrak g$ będzie algebrą Liego nad $K$.
Symetryczna forma dwuliniowa $B:\mathfrak g\times
\mathfrak g\to K$ jest niezmiennicza, jeśli
$B([x,y],z)=B(x,[y,z])$ dla wszystkich $x,y,z$.
Jej jądro to $\{x:B(x,y)=0\ \forall y\in\mathfrak g\}$.
Forma jest niezdegenerowana, jeśli jądro jest zerowe.
Dla podprzestrzeni $U$ oznaczamy
$U^\perp=\{x:B(x,u)=0\ \forall u\in U\}$.
\end{definicja}

\begin{lemat}
\label{aa:lem:lie-orthogonal-ideals}
Niech $B$ będzie symetryczną niezmienniczą formą
dwuliniową na skończenie wymiarowej algebrze Liego
$\mathfrak g$. Jeśli $\mathfrak a$ jest ideałem,
to $\mathfrak a^\perp$ jest ideałem.
W szczególności jądro $B=\mathfrak g^\perp$ jest ideałem.
Jeśli $B$ jest niezdegenerowana, to
$[\mathfrak a,\mathfrak a^\perp]=0$.
\end{lemat}
\begin{proof}
Dla $x\in\mathfrak g$, $u\in\mathfrak a^\perp$,
$a\in\mathfrak a$ mamy
$B([x,u],a)=-B(u,[x,a])=0$; pierwsza równość
wynika z antysymetrii i niezmienniczości,
druga z $[x,a]\in\mathfrak a$.
Zatem $[x,u]\in\mathfrak a^\perp$.
Dla dowolnego $y\in\mathfrak g$ zachodzi
$B([u,a],y)=B(u,[a,y])=0$.
Jeśli forma jest niezdegenerowana, to $[u,a]=0$.
\end{proof}

\begin{definicja}[Forma śladu i forma Killinga]
\label{aa:def:killing}
Dla reprezentacji $\rho:\mathfrak g\to\mathfrak{gl}(V)$
w skończenie wymiarowej przestrzeni $V$ określamy
$B_\rho(x,y)=\operatorname{tr}_V(\rho(x)\rho(y))$.
Formą Killinga algebry $\mathfrak g$ nazywamy
formę śladu reprezentacji dołączonej:
\[
 \kappa_{\mathfrak g}(x,y)=
 \operatorname{tr}_{\mathfrak g}(\operatorname{ad}_x
 \operatorname{ad}_y).
\]
\end{definicja}

\begin{stwierdzenie}
\label{aa:prop:killing-properties}
Każda forma $B_\rho$ z Definicji~\ref{aa:def:killing}
jest symetryczna i niezmiennicza. Jeśli
$\mathfrak a\triangleleft\mathfrak g$ jest ideałem,
to $\kappa_{\mathfrak g}|_{\mathfrak a\times\mathfrak a}
=\kappa_{\mathfrak a}$.
\end{stwierdzenie}
\begin{proof}
Dwuliniowość wynika z liniowości $\rho$ i śladu.
Symetria wynika z $\operatorname{tr}(AB)=
\operatorname{tr}(BA)$. Dla $A=\rho(x),B=\rho(y),
C=\rho(z)$, korzystając z reprezentacji i cykliczności
śladu, otrzymujemy
\[
 \operatorname{tr}([A,B]C)=\operatorname{tr}(ABC-BAC)
 =\operatorname{tr}(A(BC-CB)).
\]
Cykliczność dla trzech czynników wynika z tej dla
dwóch po potraktowaniu iloczynu dwóch macierzy
jako jednego czynnika. To daje niezmienniczość.

Dla $x\in\mathfrak a$ operator $\operatorname{ad}_x$
na $\mathfrak g$ ma obraz w $\mathfrak a$, ponieważ
$\mathfrak a$ jest ideałem. W bazie powstałej przez
uzupełnienie bazy $\mathfrak a$ do bazy $\mathfrak g$
macierz ma zatem postać
$\begin{pmatrix}\operatorname{ad}^{\mathfrak a}_x&*\\0&0\end{pmatrix}$.
Analogiczny zapis dla $y\in\mathfrak a$ pokazuje,
że iloczyn ma na przekątnej bloki
$\operatorname{ad}^{\mathfrak a}_x
\operatorname{ad}^{\mathfrak a}_y$ oraz $0$.
Równość śladów daje tezę.
\end{proof}

\section{Kryterium rozwiązalności Cartana}

Najpierw dowiedziemy wersji macierzowej.
Istotną trudnością jest przejście od zanikania
śladów do nilpotentności operatorów. Nie wystarcza
samo $\operatorname{tr}(X^2)=0$ nad ciałem zespolonym.

\begin{twierdzenie}[Kryterium macierzowe Cartana]
\label{aa:tw:cartan-linear}
Niech $K$ będzie algebraicznie domkniętym ciałem
charakterystyki zero, $V$ skończenie wymiarową
przestrzenią wektorową nad $K$, a $\mathfrak l$
podalgebrą Liego $\mathfrak{gl}(V)$. Wówczas $\mathfrak l$ jest
rozwiązalna wtedy i tylko wtedy, gdy
\[
 \operatorname{tr}_V(XY)=0\qquad
 (X\in[\mathfrak l,\mathfrak l],\ Y\in\mathfrak l).
\]
\end{twierdzenie}
\begin{proof}
Jeśli $V=0$, algebra jest zerowa i obie własności
zachodzą. Załóżmy $V\ne0$. Jeśli $\mathfrak l$
jest rozwiązalna, Wniosek~\ref{aa:cor:lie-flag}
stosuje się do reprezentacji naturalnej: wszystkie
operatory z $\mathfrak l$ są w pewnej bazie
górnotrójkątne, a z jej algebry pochodnej ściśle
górnotrójkątne. Iloczyn macierzy ściśle
górnotrójkątnej z górnotrójkątną ma zerową
przekątną: w $(i,i)$-tym wpisie jedyny możliwy
indeks niezerowego składnika to $i$, lecz pierwszy
czynnik ma wpis $(i,i)$ równy zero. Ślad znika.

Odwrotnie załóżmy warunek śladowy i ustalmy
$X\in[\mathfrak l,\mathfrak l]$. Niech
$\lambda_1,\ldots,\lambda_s$ będą parami różnymi
wartościami własnymi $X$, a $V_i$ odpowiednimi
uogólnionymi przestrzeniami własnymi. Nad $K$
mamy $V=\bigoplus_{i=1}^s V_i$, a
$X|_{V_i}=\lambda_i I+N_i$, gdzie $N_i$ jest
nilpotentny (rozkład Jordana operatora).
Dokładniej, dla $d=\dim V$ możemy przyjąć
$V_i=\ker(X-\lambda_i I_V)^d$. Wielomian
charakterystyczny rozkłada się nad $K$ na czynniki
liniowe, więc rozkład na uogólnione przestrzenie
własne obejmuje całe $V$. Każde $V_i$ jest
niezerowe, gdyż zawiera wektor własny odpowiadający
$\lambda_i$, i $N_i^d=0$. Oznaczamy $m_i=\dim V_i>0$.
Niech $E=\operatorname{span}_{\mathbb Q}
\{\lambda_1,\ldots,\lambda_s\}\subset K$;
ciało $K$ zawiera kopię $\mathbb Q$ dzięki
charakterystyce zero. Dla dowolnego liniowego
nad $\mathbb Q$ funkcjonału $f:E\to\mathbb Q$
określamy $Y_f\in\operatorname{End}_K(V)$
przez $Y_f|_{V_i}=f(\lambda_i)I_{V_i}$.
Suma prosta zapewnia jednoznaczność takiego operatora
i jego liniowość nad $K$, mimo że sam funkcjonał
$f$ jest liniowy tylko nad $\mathbb Q$.
Nie zakładamy, że $Y_f\in\mathfrak l$.

Wykażemy jednak
$[Y_f,\mathfrak l]\subset[\mathfrak l,\mathfrak l]$.
Przestrzeń $\operatorname{End}_K(V)$ rozkłada się
na sumę prostą bloków $\operatorname{Hom}_K(V_j,V_i)$
(zapisanych jako operatory zerowe na pozostałych
składnikach). Na takim bloku
\[
 \operatorname{ad}_X(A)=(\lambda_i-\lambda_j)A
                         +N_iA-AN_j.
\]
Operatory $P(A)=N_i A$ i $Q(A)=AN_j$ na tym bloku
komutują, gdyż $P(Q(A))=N_iAN_j=Q(P(A))$.
Jeśli $N_i^{r_i}=0$ i $N_j^{r_j}=0$, to dla
$k=r_i+r_j-1$ mamy
\[
 (P-Q)^k(A)=\sum_{a=0}^k(-1)^a
      \binom{k}{a}N_i^{k-a}AN_j^a=0.
\]
Istotnie, gdy $a<r_j$, zachodzi $k-a\ge r_i$,
a gdy $a\ge r_j$, znika prawy czynnik. W szczególności dla
$d=\dim V$ potęga $2d-1$ zeruje część nilpotentną
na każdym bloku. Na tym samym bloku
$\operatorname{ad}_{Y_f}$ działa przez mnożenie
przez $f(\lambda_i)-f(\lambda_j)=f(\lambda_i-\lambda_j)$.

Niech $\Delta$ będzie zbiorem wszystkich różnic
$\lambda_i-\lambda_j$; jest podzbiorem $E$
i zawiera $0$. Jeśli ta sama różnica pojawia się
dla kilku par indeksów, przypisujemy jej tę samą
wartość $f(\delta)$, więc warunki interpolacyjne
są zgodne.
W pierścieniu $K[t]$ wielomiany $(t-\delta)^{2d}$
dla różnych $\delta\in\Delta$ są parami względnie
pierwsze: ich jedyne nierozkładalne czynniki są
różne. Chińskie twierdzenie o resztach daje
wielomian $p(t)$ spełniający
\[
 p(t)\equiv f(\delta)\pmod{(t-\delta)^{2d}}
 \qquad(\delta\in\Delta).
\]
Dla $\delta=0$ mamy $f(0)=0$, więc $p(0)=0$.
Dla każdej $\delta$ istnieje $q_\delta\in K[t]$
taki, że $p(t)-f(\delta)=(t-\delta)^{2d}q_\delta(t)$.
Podstawmy $T=\operatorname{ad}_X$ ograniczony do
bloku o różnicy $\delta$. Na tym bloku
$(T-\delta I)^{2d}=0$, a podstawienie do równości
wielomianów daje
$p(T)-f(\delta)I=(T-\delta I)^{2d}q_\delta(T)=0$.
Zatem $p(T)$ jest dokładnie mnożeniem przez
$f(\delta)$, także gdy część nilpotentna $T$ jest
niezerowa. Na wszystkich blokach otrzymujemy
$p(\operatorname{ad}_X)=\operatorname{ad}_{Y_f}$.
Ponieważ $X\in\mathfrak l$, operator
$\operatorname{ad}_X$ przeprowadza $\mathfrak l$
w $[\mathfrak l,\mathfrak l]$, a ten ideał w siebie.
Dla każdego $k\ge1$ indukcja daje
$(\operatorname{ad}_X)^k(\mathfrak l)
\subset[\mathfrak l,\mathfrak l]$:
pierwszy krok wynika z $X\in\mathfrak l$, następne
z własności ideału. Ponieważ $p(0)=0$, operator
$p(\operatorname{ad}_X)$ jest kombinacją liniową
tych dodatnich potęg. Stąd
$[Y_f,\mathfrak l]=p(\operatorname{ad}_X)(\mathfrak l)
\subset[\mathfrak l,\mathfrak l]$.

Z definicji algebry pochodnej istnieją
$A_k,B_k\in\mathfrak l$ takie, że
$X=\sum_{k=1}^r[A_k,B_k]$ (współczynniki liniowe
można włączyć do pierwszych czynników).
Cykliczność śladu daje dla każdego $k$
\[
\begin{aligned}
 \operatorname{tr}([A_k,B_k]Y_f)
 &=\operatorname{tr}(A_kB_kY_f-B_kA_kY_f)\\
 &=\operatorname{tr}(A_k(B_kY_f-Y_fB_k)).
\end{aligned}
\]
Mamy $[B_k,Y_f]=-[Y_f,B_k]\in[\mathfrak l,\mathfrak l]$
oraz $A_k\in\mathfrak l$. Zamieniając kolejność
dwóch czynników pod śladem, możemy zastosować
warunek śladowy. Każdy składnik jest zerowy, więc
$\operatorname{tr}(XY_f)=0$.
Z drugiej strony na $V_i$ operator $XY_f$
ma ślad $m_i\lambda_i f(\lambda_i)$, ponieważ
$N_i$ ma ślad zero (jego wartości własne są zerowe).
Dla każdego $f$ otrzymujemy zatem równość w $K$:
\[
 \sum_{i=1}^s m_i\lambda_i f(\lambda_i)=0.
\]
Lewa strona należy do $E$, bo $f(\lambda_i)\in\mathbb Q$.
Stosując do niej $f$ i korzystając z liniowości
nad $\mathbb Q$, otrzymujemy w $\mathbb Q$
\[
 0=f\Bigl(\sum_i m_i f(\lambda_i)\lambda_i\Bigr)
   =\sum_i m_i f(\lambda_i)^2.
\]
Nie używamy żadnego porządku na $K$, lecz zwykłego
porządku na $\mathbb Q$. Wszystkie składniki
są nieujemnymi liczbami wymiernymi, a $m_i>0$,
więc $f(\lambda_i)=0$ dla każdego $i$.
Gdyby pewne $\lambda_i\ne0$, można by uzupełnić
je do bazy $E$ nad $\mathbb Q$ i wybrać funkcjonał
przyjmujący na nim wartość $1$, co jest sprzeczne
z ostatnią równością. Zatem wszystkie $\lambda_i=0$,
czyli $X$ na każdej uogólnionej przestrzeni własnej
jest równy nilpotentnemu $N_i$. Ponieważ $N_i^d=0$
i przestrzenie te rozpinają $V$, mamy $X^d=0$.

Wykazaliśmy nilpotentność każdego operatora z
$[\mathfrak l,\mathfrak l]$. Wniosek~\ref{aa:cor:engel-flag}
stosuje się do tej podalgebry na $V$ i daje wspólną
postać ściśle górnotrójkątną. Jest ona podalgebrą
algebry nilpotentnej $\mathfrak n_d$
(Przykład~\ref{aa:prz:triangular-lie}), zatem sama
jest nilpotentna, a więc rozwiązalna na mocy
Lematu~\ref{aa:lem:lie-series}. Ponieważ iloraz
$\mathfrak l/[\mathfrak l,\mathfrak l]$ jest abelowy,
Lemat~\ref{aa:lem:solvable-extension} dowodzi
rozwiązalności $\mathfrak l$.
\end{proof}

\begin{twierdzenie}[Kryterium rozwiązalności Cartana]
\label{aa:tw:cartan-solvability}
Niech $\mathfrak g$ będzie skończenie wymiarową algebrą
Liego nad ciałem $K$ charakterystyki zero.
Algebra $\mathfrak g$ jest rozwiązalna wtedy i tylko
wtedy, gdy $\kappa_{\mathfrak g}(x,y)=0$
dla wszystkich $x\in[\mathfrak g,\mathfrak g]$,
$y\in\mathfrak g$.
\end{twierdzenie}
\begin{proof}
Najpierw niech $K$ będzie algebraicznie domknięte.
Stosujemy Twierdzenie~\ref{aa:tw:cartan-linear}
do $\mathfrak l=\operatorname{ad}(\mathfrak g)$
na $V=\mathfrak g$. Ponieważ $\operatorname{ad}$
jest homomorfizmem,
$[\mathfrak l,\mathfrak l]=
\operatorname{ad}([\mathfrak g,\mathfrak g])$.
Warunek śladowy w twierdzeniu jest dokładnie
warunkiem w tezie. Algebra $\mathfrak l$ jest
rozwiązalna wtedy i tylko wtedy, gdy $\mathfrak g$
jest rozwiązalna: w jedną stronę stosujemy
Stwierdzenie~\ref{aa:prop:lie-series-images} do obrazu,
w drugą korzystamy z
$\mathfrak g/Z(\mathfrak g)\simeq\mathfrak l$
i Lematu~\ref{aa:lem:solvable-extension}, ponieważ
centrum jest abelowe.

Dla dowolnego $K$ charakterystyki zero bierzemy
jego domknięcie algebraiczne $L$ i rozszerzamy
skalary: $\mathfrak g_L=\mathfrak g\otimes_K L$,
z nawiasem określonym na tensorach prostych przez
$[x\otimes a,y\otimes b]=[x,y]\otimes ab$.
W bazie $\mathfrak g$ nad $K$ ta sama rodzina
wektorów z tensorem $1$ jest bazą nad $L$.
Nawias jest dobrze określony przez rozszerzenie
dwuliniowe stałych struktury; alternowanie
i tożsamość Jacobiego sprawdzamy na wektorach
bazowych, gdzie są tymi samymi równościami w $K$
odczytanymi w $L$ (alternowanie również wynika
z antysymetrii w charakterystyce zero).
Jeżeli $U,W\subset\mathfrak g$ są podprzestrzeniami,
to ich rozszerzenia skalarów utożsamiamy z
podprzestrzeniami $\mathfrak g_L$: bazę każdej
z nich można uzupełnić do bazy $\mathfrak g$.
Zachodzi równość
$[U\otimes_K L,W\otimes_K L]=[U,W]\otimes_K L$:
nawiasy tensorów prostych dają inkluzję w prawo,
a generatory $[u,w]\otimes a=[u\otimes1,w\otimes a]$
dają odwrotną. Indukcja pokazuje
$(\mathfrak g_L)^{(k)}=\mathfrak g^{(k)}\otimes L$.
Podprzestrzeń po rozszerzeniu skalarów znika
dokładnie wtedy, gdy przed rozszerzeniem jest zerowa,
co wynika z zachowania liczby elementów bazy.
Rozwiązalność jest więc równoważna przed i po
rozszerzeniu.

Macierze $\operatorname{ad}_{x\otimes1}$ w bazie
nad $L$ mają te same wpisy co macierze
$\operatorname{ad}_x$ nad $K$. Forma Killinga
na $\mathfrak g_L$ jest zatem rozszerzeniem
dwuliniowym $\kappa_{\mathfrak g}$.
Dokładniej, dla $x,y\in\mathfrak g$ i $a,b\in L$
mamy $\kappa_{\mathfrak g_L}(x\otimes a,y\otimes b)
=ab\,\kappa_{\mathfrak g}(x,y)$, z prawą stroną
odczytaną w $L$. Jeśli forma znika na
$[\mathfrak g,\mathfrak g]\times\mathfrak g$,
to znika po rozszerzeniu skalarów na
$[\mathfrak g_L,\mathfrak g_L]\times\mathfrak g_L$
przez dwuliniowość. Odwrotnie podstawiamy
$a=b=1$ i korzystamy z iniektywności $K\hookrightarrow L$.
Kryterium nad $L$ oraz wykazana równoważność
rozwiązalności przed i po rozszerzeniu dają tezę nad $K$.
\end{proof}

\section{Kryterium półprostoty Cartana}

\begin{twierdzenie}[Kryterium półprostoty Cartana]
\label{aa:tw:cartan-semisimple}
Niech $\mathfrak g$ będzie skończenie wymiarową algebrą
Liego nad ciałem charakterystyki zero.
Algebra $\mathfrak g$ jest półprosta wtedy i tylko
wtedy, gdy jej forma Killinga jest niezdegenerowana.
\end{twierdzenie}
\begin{proof}
Załóżmy najpierw półprostotę. Niech
$\mathfrak k=\ker\kappa_{\mathfrak g}$.
Z Lematu~\ref{aa:lem:lie-orthogonal-ideals} jest to
ideał. Na mocy Stwierdzenia~\ref{aa:prop:killing-properties}
forma Killinga $\mathfrak k$ jest ograniczeniem
formy Killinga $\mathfrak g$ do $\mathfrak k$,
więc jest identycznie zerowa. Kryterium
rozwiązalności Cartana (Twierdzenie~\ref{aa:tw:cartan-solvability})
stosuje się do $\mathfrak k$ i daje rozwiązalność.
Z półprostoty wynika $\mathfrak k=0$.

Odwrotnie, niech $\kappa_{\mathfrak g}$ będzie
niezdegenerowana. Gdyby radykał rozwiązalny
$\mathfrak r$ był niezerowy, jego ostatni niezerowy
wyraz szeregu pochodnego $\mathfrak a=
\mathfrak r^{(j)}$ byłby niezerowym abelowym
ideałem $\mathfrak g$. Aby uzasadnić własność
ideału w całej $\mathfrak g$, zaczynamy od
$\mathfrak r\triangleleft\mathfrak g$ i stosujemy
indukcyjnie fakt z dowodu Lematu~\ref{aa:lem:lie-series},
że nawias dwóch ideałów jest ideałem.
Dla $x\in\mathfrak a$ mamy
$\operatorname{ad}_x(\mathfrak g)\subset\mathfrak a$
oraz $\operatorname{ad}_x(\mathfrak a)=0$.
Dla dowolnego $y\in\mathfrak g$ operator
$\operatorname{ad}_y$ zachowuje $\mathfrak a$.
Zatem $\operatorname{ad}_x\operatorname{ad}_y$
ma obraz w $\mathfrak a$ i znika na $\mathfrak a$.
W bazie uzupełniającej bazę $\mathfrak a$
jego oba bloki diagonalne są zerowe, a więc
$\kappa_{\mathfrak g}(x,y)=0$.
Stąd $\mathfrak a\subset\ker\kappa_{\mathfrak g}=0$,
sprzeczność. Radykał jest zerowy.
\end{proof}

\section{Rozkład na ideały proste}

\begin{twierdzenie}[Rozkład algebry półprostej]
\label{aa:tw:semisimple-decomposition}
Niech $\mathfrak g$ będzie skończenie wymiarową
półprostą algebrą Liego nad ciałem $K$ charakterystyki
zero. Dla każdego ideału $\mathfrak a$ mamy rozkład
$\mathfrak g=\mathfrak a\oplus\mathfrak a^\perp$
na komutujące ideały, gdzie dopełnienie jest
brane względem formy Killinga.
W szczególności $\mathfrak g$ jest sumą prostą
skończenie wielu ideałów prostych. I odwrotnie,
suma prosta skończenie wielu algebr prostych nad
$K$ jest półprosta.
\end{twierdzenie}
\begin{proof}
Forma Killinga jest niezdegenerowana na mocy
Twierdzenia~\ref{aa:tw:cartan-semisimple}.
Z Lematu~\ref{aa:lem:lie-orthogonal-ideals}
$\mathfrak a^\perp$ jest ideałem oraz
$[\mathfrak a,\mathfrak a^\perp]=0$.
Przecięcie $\mathfrak a\cap\mathfrak a^\perp$
jest ideałem i jest abelowe, bo nawias dwóch
jego elementów należy do
$[\mathfrak a,\mathfrak a^\perp]=0$.
Półprostota wymusza zerowość przecięcia.
Niezdegenerowana forma indukuje izomorfizm
$\mathfrak g\to\mathfrak g^*$, $x\mapsto\kappa(x,\cdot)$.
Ograniczenie funkcjonałów do $\mathfrak a$
jest surjektywne (każdy funkcjonał przedłużamy
przez wybór bazy uzupełniającej).
Złożenie obu odwzorowań ma jądro
$\mathfrak a^\perp$ oraz obraz $\mathfrak a^*$.
Twierdzenie o wymiarze obrazu i jądra daje
$\dim\mathfrak a^\perp=\dim\mathfrak g-\dim\mathfrak a$.
Wraz z zerowym przecięciem dowodzi to rozkładu.
Ograniczenie formy do każdego składnika jest
niezdegenerowane, ponieważ element składnika
ortogonalny do własnego składnika jest również
ortogonalny do drugiego i tym samym do całej
$\mathfrak g$. Formy te są formami Killinga
składników na mocy Stwierdzenia~\ref{aa:prop:killing-properties}.
Kryterium półprostoty daje półprostotę obu składników.

Indukcja po wymiarze: dla wymiaru zero bierzemy
pustą sumę prostą. Niezerowa półprosta algebra
nie może być abelowa. Jeśli nie ma właściwego
niezerowego ideału, jest prosta. W przeciwnym
razie wybieramy taki ideał $\mathfrak a$,
a oba składniki rozkładu mają mniejszy wymiar
i są półproste. Indukcja daje rozkłady na ideały
proste w składnikach. Każdy ideał jednego
składnika jest także ideałem w całej algebrze:
nawias z drugim składnikiem jest zerowy,
a z pierwszym pozostaje w tym ideale.
Złożenie rozkładów daje tezę.

Odwrotnie niech $\mathfrak g=\bigoplus_i\mathfrak g_i$,
gdzie $\mathfrak g_i$ są proste. Dla ideału
rozwiązalnego $\mathfrak r$ rzut na każdy składnik
jest ideałem w tym składniku: dla $x\in\mathfrak g_i$
wybieramy jego przeciwobraz przez rzut, a następnie
stosujemy własność ideału $\mathfrak r$ i
homomorficzność rzutu. Obraz jest rozwiązalny.
Ponieważ każda algebra prosta jest półprosta,
każdy taki obraz jest zerowy. Wszystkie współrzędne
elementu $\mathfrak r$ znikają, więc $\mathfrak r=0$.
\end{proof}

\begin{wniosek}
\label{aa:cor:semisimple-perfect}
Niech $\mathfrak g$ będzie skończenie wymiarową
półprostą algebrą Liego nad ciałem charakterystyki
zero. Wtedy $[\mathfrak g,\mathfrak g]=\mathfrak g$.
\end{wniosek}
\begin{proof}
Z Twierdzenia~\ref{aa:tw:semisimple-decomposition}
mamy sumę prostą komutujących ideałów prostych.
Każdy składnik ma niezerową algebrę pochodną,
bo jest nieabelowy, a ta jest jego ideałem,
więc jest równa składnikowi. Nawias dwóch różnych
składników jest zerowy, a nawiasy wewnętrzne
rozpinają wszystkie składniki. Teza zachodzi
również dla pustej sumy, czyli algebry zerowej.
\end{proof}

\section{Derywacje algebry półprostej}

\begin{twierdzenie}[Derywacje wewnętrzne]
\label{aa:tw:semisimple-derivations}
Niech $\mathfrak g$ będzie skończenie wymiarową
półprostą algebrą Liego nad ciałem charakterystyki
zero. Każda derywacja $D$ algebry $\mathfrak g$
ma postać $\operatorname{ad}_z$ dla dokładnie
jednego $z\in\mathfrak g$.
\end{twierdzenie}
\begin{proof}
Forma Killinga jest niezdegenerowana na mocy
Twierdzenia~\ref{aa:tw:cartan-semisimple}.
Funkcjonał $x\mapsto\operatorname{tr}
(D\operatorname{ad}_x)$ jest liniowy,
zatem istnieje dokładnie jeden $z\in\mathfrak g$
taki, że dla wszystkich $x$
$\operatorname{tr}(D\operatorname{ad}_x)=\kappa(z,x)$:
jest to surjektywność i iniektywność izomorfizmu
z $\mathfrak g$ do $\mathfrak g^*$ wyznaczonego
przez formę niezdegenerowaną w skończonym wymiarze.
Połóżmy $D_0=D-\operatorname{ad}_z$.
Jest derywacją na mocy
Stwierdzenia~\ref{aa:prop:lie-derivations}
i spełnia $\operatorname{tr}(D_0\operatorname{ad}_x)=0$
dla każdego $x$. Dla dowolnych $x,y$ z tej samej
tożsamości o derywacjach mamy
$[D_0,\operatorname{ad}_x]=\operatorname{ad}_{D_0x}$.
Stąd
\[
\begin{aligned}
 \kappa(D_0x,y)
 &=\operatorname{tr}([D_0,\operatorname{ad}_x]
                   \operatorname{ad}_y)\\
 &=\operatorname{tr}(D_0[\operatorname{ad}_x,
                                  \operatorname{ad}_y])\\
 &=\operatorname{tr}(D_0\operatorname{ad}_{[x,y]})=0.
\end{aligned}
\]
Druga równość to cykliczność śladu, trzecia
to własność reprezentacji dołączonej.
Niezdegenerowanie formy daje $D_0x=0$ dla
każdego $x$, więc $D=\operatorname{ad}_z$.
Jednoznaczność wynika także z $Z(\mathfrak g)=0$.
\end{proof}

\section{Formy Killinga klasycznych algebr}

\begin{twierdzenie}
\label{aa:tw:classical-killing}
Niech $K$ będzie ciałem charakterystyki zero,
a $n\ge1$. W oznaczeniach
Przykładu~\ref{aa:prz:classical-lie} zachodzą wzory
\[
\begin{aligned}
 \kappa_{\mathfrak{gl}_n}(X,Y)&=
  2n\operatorname{tr}(XY)-2\operatorname{tr}(X)\operatorname{tr}(Y),\\
 \kappa_{\mathfrak{sl}_n}(X,Y)&=2n\operatorname{tr}(XY),\\
 \kappa_{\mathfrak{so}_n}(X,Y)&=(n-2)\operatorname{tr}(XY),\\
 \kappa_{\mathfrak{sp}_{2n}}(X,Y)&=(2n+2)\operatorname{tr}(XY),
\end{aligned}
\]
gdzie w każdym wierszu $X,Y$ należą do wskazanej
algebry, a ślady po prawej stronie są śladami
w reprezentacjach naturalnych.
\end{twierdzenie}
\begin{proof}
Na $M_n(K)$ piszemy $\operatorname{ad}_X=L_X-R_X$,
gdzie $L_X(A)=XA$, $R_X(A)=AX$.
W bazie jednostek macierzowych mamy
\[
 \operatorname{tr}_{M_n}(L_A)=n\operatorname{tr}(A),
 \quad \operatorname{tr}_{M_n}(R_A)=n\operatorname{tr}(A),
 \quad \operatorname{tr}_{M_n}(L_A R_B)
       =\operatorname{tr}(A)\operatorname{tr}(B).
\]
Istotnie współczynnik przy $E_{ij}$ w $AE_{ij}$
to $a_{ii}$, w $E_{ij}A$ to $a_{jj}$, a w
$AE_{ij}B$ to $a_{ii}b_{jj}$; sumujemy po $i,j$.
Ponieważ $L_XL_Y=L_{XY}$, $R_XR_Y=R_{YX}$
i lewe mnożenia komutują z prawymi,
rozwinięcie $(L_X-R_X)(L_Y-R_Y)$ daje pierwszy wzór.
Algebra $\mathfrak{sl}_n$ jest ideałem w $\mathfrak{gl}_n$,
więc Stwierdzenie~\ref{aa:prop:killing-properties}
daje drugi wzór przez ograniczenie.

Dla pozostałych dwóch przypadków niech $V=K^m$
i niech $b$ będzie niezdegenerowaną formą
dwuliniową symetryczną (gdy $m=n$, macierz $I_n$)
lub alternującą (gdy $m=2n$, macierz $J$).
Niech $\mathfrak l=\{X:b(Xu,v)+b(u,Xv)=0\
\forall u,v\in V\}$. Są to dokładnie rozważane
algebry ortogonalna i symplektyczna.
Odwzorowanie
\[
 T:V\otimes V\to\operatorname{End}_K(V),
 \qquad T(u\otimes v)(w)=b(v,w)u
\]
jest izomorfizmem: niezdegenerowanie $b$ daje
izomorfizm $V\to V^*$, $v\mapsto b(v,\cdot)$,
a w dowolnej bazie $e_i$ i bazie dualnej $e^j$
operatory $w\mapsto e^j(w)e_i$ tworzą bazę
jednostek macierzowych.
Dla $X\in\mathfrak l$ oraz $u,v,w\in V$ mamy
\[
 [X,T(u\otimes v)](w)
 = b(v,w)Xu-b(v,Xw)u
 = b(v,w)Xu+b(Xv,w)u.
\]
Zatem $T$ przeplata działanie dołączone z działaniem
$S_X=X\otimes I+I\otimes X$.

Niech $\tau(u\otimes v)=v\otimes u$.
Podprzestrzenie $V_+$ i $V_-$ odpowiadają wartościom
własnym $+1$ i $-1$ operatora $\tau$.
Ponieważ $\tau^2=I$ i $2$ jest odwracalne,
$P_\pm=(I\pm\tau)/2$ są rzutami na te przestrzenie,
a ich suma to $I$ i iloczyn zero.
Jeśli $b$ jest symetryczna, to
$T(V_-)=\mathfrak l$: dla $A=T(u\otimes v-v\otimes u)$
mamy
\[
 b(Aw,z)+b(w,Az)=
 b(v,w)b(u,z)-b(u,w)b(v,z)
 +b(v,z)b(w,u)-b(u,z)b(w,v)=0.
\]
Takie tensory rozpinają $V_-$, ponieważ współczynniki
w bazie $e_i\otimes e_j$ spełniają $a_{ij}=-a_{ji}$
i $a_{ii}=0$. Zatem obraz jest zawarty w $\mathfrak l$,
a oba wymiary są równe $m(m-1)/2$, co daje równość.
Jeśli $b$ jest alternująca, to dla
$A=T(u\otimes v+v\otimes u)$ analogiczna suma ma postać
\[
 b(v,w)b(u,z)+b(u,w)b(v,z)
 +b(v,z)b(w,u)+b(u,z)b(w,v)=0
\]
na mocy $b(w,u)=-b(u,w)$. Tensory te rozpinają $V_+$,
włącznie z $e_i\otimes e_i$, ponieważ $2$ jest
odwracalne. Wymiar $V_+$ wynosi $m(m+1)/2$;
dla $m=2n$ jest równy $n(2n+1)$, czyli wymiarowi
$\mathfrak l$ z Przykładu~\ref{aa:prz:classical-lie}.
Stąd $T(V_+)=\mathfrak l$.

Operator $S_X$ komutuje z $\tau$ (zamiana czynników
zamienia dwa składniki $S_X$ miejscami), więc
zachowuje obie przestrzenie własne. W konsekwencji
\[
 \operatorname{tr}_{V_\pm}(S_X S_Y)
 =\tfrac12\left(\operatorname{tr}_{V\otimes V}(S_X S_Y)
          \pm\operatorname{tr}_{V\otimes V}(\tau S_X S_Y)\right).
\]
W istocie lewa strona jest śladem $P_\pm S_X S_Y$
na całej przestrzeni, bo na drugim składniku
ten operator znika.

Dla dowolnych operatorów $A,B$ na $V$ mamy
$\operatorname{tr}(A\otimes B)=\operatorname{tr}(A)
\operatorname{tr}(B)$ oraz
$\operatorname{tr}(\tau(A\otimes B))=\operatorname{tr}(AB)$.
Pierwszą równość dostajemy sumując współczynniki
$a_{ii}b_{jj}$ przy $e_i\otimes e_j$, a drugą
sumując współczynniki $a_{ji}b_{ij}$.
Rozwijamy
\[
 S_XS_Y=XY\otimes I+X\otimes Y+Y\otimes X+I\otimes XY.
\]
Zatem jego ślad na całej przestrzeni wynosi
$2m\operatorname{tr}(XY)+2\operatorname{tr}(X)
\operatorname{tr}(Y)$, a po przemnożeniu przez
$\tau$ wynosi $4\operatorname{tr}(XY)$.
W algebrach ortogonalnej i symplektycznej
$\operatorname{tr}(X)=0$: w pierwszej
$X^{\mathsf T}=-X$, a w drugiej bloki diagonalne
mają postać $P,-P^{\mathsf T}$.
Stąd ślad ograniczenia na $V_-$ to
$(m-2)\operatorname{tr}(XY)$, a na $V_+$ to
$(m+2)\operatorname{tr}(XY)$.
Izomorfizm $T$ przeplata te ograniczenia
z reprezentacją dołączoną na $\mathfrak l$,
więc ślady są formami Killinga. Podstawiając
$m=n$ dla $V_-$ i $m=2n$ dla $V_+$,
otrzymujemy ostatnie dwa wzory.
\end{proof}

\begin{przyklad}[Forma Killinga i półprostota modeli klasycznych]
\label{aa:prz:killing-classical-consequences}
Nad ciałem charakterystyki zero forma
$(X,Y)\mapsto\operatorname{tr}(XY)$ na $M_n(K)$
jest niezdegenerowana: $\operatorname{tr}(XE_{ji})=x_{ij}$,
więc znikanie wszystkich takich śladów wymusza $X=0$.
Dla $X\in\mathfrak{sl}_n$ ortogonalnego do całej
$\mathfrak{sl}_n$ testy $Y=E_{ji}$, $i\ne j$,
dają zerowe wpisy poza przekątną, a testy
$Y=E_{ii}-E_{jj}$ dają równe wpisy diagonalne.
Zatem $X=cI_n$ i $nc=0$, czyli $c=0$.
Forma Killinga $\mathfrak{sl}_n$ jest więc
niezdegenerowana dla każdego $n\ge1$.
Dla $\mathfrak{gl}_n$ wzór z
Twierdzenia~\ref{aa:tw:classical-killing}
pokazuje, że $KI_n$ jest w jądrze, a każdy
element jądra jest skalarny na mocy identycznych
testów z macierzami śladu zero. Ponadto
$\mathfrak{gl}_n=KI_n\oplus\mathfrak{sl}_n$
(jawny rozkład $X=(\operatorname{tr}(X)/n)I_n+
(X-(\operatorname{tr}(X)/n)I_n)$). Pierwszy
składnik jest centralnym ideałem, drugi ideałem
półprostym na mocy kryterium Cartana.
Każdy ideał rozwiązalny ma zerowy obraz przez
rzut na $\mathfrak{sl}_n$, więc
$\operatorname{rad}\mathfrak{gl}_n=KI_n$.

Dla $\mathfrak{so}_n$ baza składa się z
$F_{ij}=E_{ij}-E_{ji}$, $i<j$. Ze wzoru
na mnożenie jednostek macierzowych
$\operatorname{tr}(F_{ij}F_{kl})=-2\delta_{ik}\delta_{jl}$
dla uporządkowanych par $i<j,k<l$. Forma śladu
ma diagonalną macierz z niezerowymi wpisami,
a forma Killinga jest niezdegenerowana dla $n\ge3$.
Dla $n=2$ algebra jest jednowymiarowa abelowa,
więc forma Killinga jest zerowa i algebra nie
jest półprosta; dla $n=1$ algebra jest zerowa.

Dla $\mathfrak{sp}_{2n}$ piszemy
$X=\begin{pmatrix}P&Q\\R&-P^{\mathsf T}\end{pmatrix}$,
$Y=\begin{pmatrix}A&B\\C&-A^{\mathsf T}\end{pmatrix}$,
gdzie $Q,R,B,C$ są symetryczne. Mnożenie bloków
i cykliczność śladu dają
\[
 \operatorname{tr}(XY)=2\operatorname{tr}(PA)
                    +\operatorname{tr}(QC)
                    +\operatorname{tr}(RB).
\]
Wybierając kolejno $B=C=0$ i dowolne $A$,
$A=B=0$ i dowolne symetryczne $C$ oraz $A=C=0$
i dowolne symetryczne $B$, wykazujemy, że
ortogonalność do całej algebry wymusza $P=Q=R=0$:
dla $P$ stosujemy test jednostkami macierzowymi,
dla macierzy symetrycznych testy $E_{ii}$
oraz $E_{ij}+E_{ji}$ dają wpisy diagonalne
i dwukrotności wpisów pozadiagonalnych.
Zatem forma śladu jest niezdegenerowana.
Współczynnik $2n+2$ jest niezerowy w charakterystyce
zero, więc forma Killinga także jest niezdegenerowana.
Kryterium~\ref{aa:tw:cartan-semisimple} daje
półprostotę $\mathfrak{sl}_n$, $\mathfrak{so}_n$
dla $n\ge3$ i $\mathfrak{sp}_{2n}$.
Nie dowodzi to jeszcze prostoty wszystkich
niezerowych modeli.
\end{przyklad}

\begin{przyklad}[Macierz formy Killinga na $\mathfrak{sl}_2$]
\label{aa:prz:killing-sl2}
W bazie $(h,e,f)$ z Przykładu~\ref{aa:prz:heisenberg-sl2}
korzystamy ze wzoru $4\operatorname{tr}(XY)$.
Mamy $h^2=I_2$, $ef=E_{11}$, $fe=E_{22}$,
$e^2=f^2=0$, a $he,hf$ mają ślad zero.
Macierz formy i jej wyznacznik wynoszą
\[
 \begin{pmatrix}8&0&0\\0&0&4\\0&4&0\end{pmatrix},
 \qquad -128.
\]
Forma jest niezdegenerowana, choć $\kappa(e,e)=0$
dla $e\ne0$. Niezdegenerowanie nie jest warunkiem
dodatniej określoności.
\end{przyklad}

Uzyskaliśmy rozkład półprostych algebr Liego i
charakterystykę poprzez formę Killinga. Związek
tych algebr z układami korzeni z rozdziału~\ref{aa:ch:root-systems}
wymaga jeszcze rozkładu na przestrzenie korzeniowe,
którym zajmiemy się w dalszej części tomu.

% AA-59
\chapter[Podalgebry Cartana]{Podalgebry Cartana i rozkład na przestrzenie korzeniowe}
\label{aa:ch:cartan-root-decomposition}

W rozdziale~\ref{aa:ch:semisimple-lie} wykazaliśmy niezdegenerowanie formy Killinga algebry półprostej oraz wewnętrzność jej derywacji. Teraz wyprowadzimy z tych wyników rozkład algebry na wspólne przestrzenie własne operatorów reprezentacji dołączonej. W całym rozdziale $K$ jest algebraicznie domkniętym ciałem charakterystyki zero, a $\mathfrak g$ jest skończenie wymiarową algebrą Liego nad $K$.

\section{Rozkład Jordana derywacji}

\begin{lemat}[Uogólnione przestrzenie własne derywacji]
\label{aa:lem:derivation-generalized-weights}
Niech $\mathfrak g$ będzie skończenie wymiarową algebrą Liego nad algebraicznie domkniętym ciałem $K$ charakterystyki zero, a $D$ jej derywacją. Dla każdego $\lambda\in K$ oznaczmy przez $V_\lambda$ uogólnioną przestrzeń własną $D$, czyli $\bigcup_{m\ge1}\ker(D-\lambda I)^m$ (równą zeru, jeśli $\lambda$ nie jest wartością własną). Wówczas dla wszystkich $\lambda,\mu\in K$ zachodzi $[V_\lambda,V_\mu]\subset V_{\lambda+\mu}$.
\end{lemat}
\begin{proof}
Dla $x\in V_\lambda$, $y\in V_\mu$ z tożsamości derywacji otrzymujemy
\[
 (D-(\lambda+\mu)I)[x,y]=
 [(D-\lambda I)x,y]+[x,(D-\mu I)y].
\]
Indukcja względem $m$ daje
\[
 (D-(\lambda+\mu)I)^m[x,y]=
 \sum_{j=0}^m\binom mj[(D-\lambda I)^jx,(D-\mu I)^{m-j}y].
\]
W kroku indukcyjnym stosujemy poprzednią tożsamość do każdego nawiasu; po zmianie indeksu współczynniki sumują się według wzoru Pascala. Wybieramy $p,q\ge1$ takie, że $(D-\lambda I)^px=0$ i $(D-\mu I)^qy=0$. Dla $m=p+q-1$ każdy składnik znika: albo $j\ge p$, albo $m-j\ge q$. Stąd teza, także gdy operator $D-(\lambda+\mu)I$ jest odwracalny.
\end{proof}

\begin{twierdzenie}[Rozkład Jordana w algebrze półprostej]
\label{aa:tw:inner-jordan}
Niech $\mathfrak g$ będzie skończenie wymiarową półprostą algebrą Liego nad algebraicznie domkniętym ciałem $K$ charakterystyki zero. Dla każdego $x\in\mathfrak g$ istnieją jednoznaczne $s,n\in\mathfrak g$ takie, że $x=s+n$, $[s,n]=0$, $\operatorname{ad}_s$ jest diagonalizowalny, a $\operatorname{ad}_n$ nilpotentny. Dla każdego $y\in\mathfrak g$, jeśli $[x,y]=0$, to $[s,y]=[n,y]=0$.
\end{twierdzenie}
\begin{proof}
Niech $D=\operatorname{ad}_x$. Rozkład na uogólnione przestrzenie własne daje $\mathfrak g=\bigoplus_\lambda V_\lambda$. Definiujemy $S|_{V_\lambda}=\lambda I$ oraz $N=D-S$. Operator $S$ jest diagonalizowalny, $N$ nilpotentny, a oba komutują, gdyż zachowują wszystkie $V_\lambda$ i $S$ jest na nich skalarny. Na mocy Lematu~\ref{aa:lem:derivation-generalized-weights}, dla $u\in V_\lambda$, $v\in V_\mu$ mamy
$S[u,v]=(\lambda+\mu)[u,v]=[Su,v]+[u,Sv]$.
Dwuliniowość i rozkład na sumę prostą rozszerzają tę równość na dowolne $u,v$. Zatem $S$ i $N=D-S$ są derywacjami.

Twierdzenie~\ref{aa:tw:semisimple-derivations} daje $s,n$ z $S=\operatorname{ad}_s$, $N=\operatorname{ad}_n$. Centrum $\mathfrak g$ jest zerowe, więc z $D=S+N$ i $[S,N]=0$ wynikają $x=s+n$ i $[s,n]=0$.

Dla jednoznaczności niech $D=S'+N'$ będzie innym takim rozkładem. Na przestrzeni własnej $W_\nu$ operatora $S'$ operator $N'$ zachowuje $W_\nu$ dzięki komutowaniu, a $D-\nu I=N'$ jest tam nilpotentny. Stąd $W_\nu\subset V_\nu$. Dla $v\in V_\lambda$ zapisujemy $v=\sum_\nu w_\nu$, gdzie $w_\nu\in W_\nu\subset V_\nu$ (przyjmujemy $W_\nu=0$, jeśli $\nu$ nie jest wartością własną $S'$). Jednoznaczność składowych w rozkładzie na $V_\nu$ daje $w_\nu=0$ dla $\nu\ne\lambda$ oraz $v=w_\lambda$. Zatem $V_\lambda\subset W_\lambda$, co wraz z wcześniejszą inkluzją daje równość. Zatem $S'=S$, $N'=N$, a zerowe centrum daje jednoznaczność $s,n$.

Jeśli $[x,y]=0$, operator $\operatorname{ad}_y$ komutuje z $D$, a zatem zachowuje każde $V_\lambda$ (komutuje z każdą potęgą $D-\lambda I$). Komutuje więc z $S$ i $N=D-S$. Ponownie zerowe centrum daje żądane równości.
\end{proof}

\section{Podalgebry toralne i Cartana}

\begin{definicja}
\label{aa:def:toral-cartan}
Niech $\mathfrak g$ będzie skończenie wymiarową algebrą Liego nad algebraicznie domkniętym ciałem $K$ charakterystyki zero. Podalgebra $\mathfrak t\subset\mathfrak g$ jest toralna, jeśli jest abelowa i dla każdego $t\in\mathfrak t$ operator $\operatorname{ad}_t$ na $\mathfrak g$ jest diagonalizowalny. Podalgebra $\mathfrak h$ jest podalgebrą Cartana, jeśli jest nilpotentna i jej normalizator
$N_{\mathfrak g}(\mathfrak h)=\{x\in\mathfrak g:[x,\mathfrak h]\subset\mathfrak h\}$ jest równy $\mathfrak h$. Centralizatorem podalgebry $\mathfrak h$ nazywamy $Z_{\mathfrak g}(\mathfrak h)=\{x\in\mathfrak g:[x,t]=0\ \forall t\in\mathfrak h\}$.
\end{definicja}

\begin{lemat}[Wspólny rozkład własny]
\label{aa:lem:toral-weight}
Niech $\mathfrak g$ będzie skończenie wymiarową algebrą Liego nad algebraicznie domkniętym ciałem $K$ charakterystyki zero, a $\mathfrak t$ jej podalgebrą toralną. Istnieje skończony zbiór $\Lambda\subset\mathfrak t^*$ takich funkcjonałów, że
\[
 \mathfrak g=\bigoplus_{\lambda\in\Lambda}\mathfrak g_\lambda,\qquad
 \mathfrak g_\lambda=\{x:[t,x]=\lambda(t)x\ \forall t\in\mathfrak t\},
\]
gdzie $\Lambda$ zawiera dokładnie te $\lambda$, dla których $\mathfrak g_\lambda\ne0$. Dla wszystkich $\lambda,\mu\in\mathfrak t^*$ (z zerowymi przestrzeniami dopuszczonymi) zachodzą
$[\mathfrak g_\lambda,\mathfrak g_\mu]\subset\mathfrak g_{\lambda+\mu}$ oraz
$\kappa(\mathfrak g_\lambda,\mathfrak g_\mu)=0$, jeśli $\lambda+\mu\ne0$, gdzie $\kappa$ jest formą Killinga.
\end{lemat}
\begin{proof}
Wybieramy bazę $t_1,\ldots,t_r$ przestrzeni $\mathfrak t$. Operatory $\operatorname{ad}_{t_i}$ komutują. Rozkład na przestrzenie własne pierwszego jest zachowany przez każdy pozostały operator. Ograniczenie operatora diagonalizowalnego do niezmienniczej podprzestrzeni pozostaje diagonalizowalne: jego wielomian minimalny dzieli wielomian o parami różnych pierwiastkach. Powtarzając rozkład na każdej składowej, otrzymujemy wspólną bazę własną. Wspólna lista wartości własnych na bazie $\mathfrak t$ wyznacza funkcjonał liniowy $\lambda$; liniowość reprezentacji dołączonej daje równość dla wszystkich $t$. Gdy $\mathfrak t=0$, jedynym możliwym funkcjonałem jest zero.

Ustalmy $\lambda,\mu\in\mathfrak t^*$, $x\in\mathfrak g_\lambda$, $y\in\mathfrak g_\mu$ oraz $t\in\mathfrak t$. Z tożsamości Jacobiego mamy
$[t,[x,y]]=[[t,x],y]+[x,[t,y]]=(\lambda+\mu)(t)[x,y]$.
Z niezmienniczości i symetrii formy Killinga wynika
$0=\kappa([t,x],y)+\kappa(x,[t,y])=(\lambda+\mu)(t)\kappa(x,y)$.
Jeżeli $\lambda+\mu\ne0$, wybieramy $t$, dla którego ten funkcjonał nie znika.
\end{proof}

\begin{twierdzenie}[Maksymalna podalgebra toralna]
\label{aa:tw:maximal-toral-cartan}
Niech $\mathfrak g$ będzie skończenie wymiarową półprostą algebrą Liego nad algebraicznie domkniętym ciałem $K$ charakterystyki zero. Istnieje maksymalna względem inkluzji podalgebra toralna $\mathfrak h$. Dla każdej takiej podalgebry zachodzą
$Z_{\mathfrak g}(\mathfrak h)=N_{\mathfrak g}(\mathfrak h)=\mathfrak h$, a ograniczenie formy Killinga do $\mathfrak h$ jest niezdegenerowane. W szczególności $\mathfrak h$ jest podalgebrą Cartana.
\end{twierdzenie}
\begin{proof}
Podalgebra zerowa jest toralna. Wybierając toralną podalgebrę największego wymiaru, otrzymujemy podalgebrę maksymalną. Niech teraz $\mathfrak h$ będzie dowolną maksymalną toralną podalgebrą. Jej centralizator $\mathfrak c=\mathfrak g_0$ jest podalgebrą (tożsamość Jacobiego). Na mocy kryterium półprostoty Cartana (Twierdzenie~\ref{aa:tw:cartan-semisimple}) forma Killinga $\kappa=\kappa_{\mathfrak g}$ jest niezdegenerowana. Wykażemy, że jej ograniczenie do $\mathfrak c$ również jest niezdegenerowane. Jeśli $u\in\mathfrak c$ i $\kappa(u,v)=0$ dla każdego $v\in\mathfrak c$, to na mocy Lematu~\ref{aa:lem:toral-weight} mamy także $\kappa(u,v)=0$ dla $v\in\mathfrak g_\lambda$, $\lambda\ne0$. Każdy element $\mathfrak g$ jest sumą składowej z $\mathfrak c$ i składowych z niezerowych przestrzeni wagowych; dwuliniowość daje więc $\kappa(u,\mathfrak g)=0$. Niezdegenerowanie $\kappa$ oznacza $u=0$. Chodzi tu o ograniczenie formy Killinga całej $\mathfrak g$, a nie o formę Killinga samej $\mathfrak c$.

Dla $x\in\mathfrak c$ rozkładamy $x=s+n$ według Twierdzenia~\ref{aa:tw:inner-jordan}. Obie składowe komutują z $\mathfrak h$. Przestrzeń $\mathfrak h+Ks$ jest abelową podalgebrą; każdy jej element działa diagonalizowalnie, gdyż komutujące operatory diagonalizowalne $\operatorname{ad}_s$ i $\operatorname{ad}_t$ dla $t\in\mathfrak h$ mają wspólną bazę własną. Maksymalność daje $s\in\mathfrak h$. Wobec tego $\operatorname{ad}_x|_{\mathfrak c}=\operatorname{ad}_n|_{\mathfrak c}$ jest nilpotentny. Twierdzenie Engela (Twierdzenie~\ref{aa:tw:engel}) zastosowane do $\mathfrak c$ daje jej nilpotentność, a więc rozwiązalność na mocy Lematu~\ref{aa:lem:lie-series}.

Na mocy Wniosku~\ref{aa:cor:lie-flag} reprezentacja $\mathfrak c\to\mathfrak{gl}(\mathfrak g)$, $u\mapsto\operatorname{ad}_u$ (ograniczenie reprezentacji dołączonej $\mathfrak g$) ma wspólną bazę górnotrójkątną. Operator $\operatorname{ad}_n$ jest nilpotentny, więc jego przekątna w tej bazie jest zerowa. Dla każdego $y\in\mathfrak c$ iloczyn $\operatorname{ad}_n\operatorname{ad}_y$ jest ściśle górnotrójkątny, zatem $\kappa(n,y)=0$. Niezdegenerowanie ograniczenia do $\mathfrak c$ i $n\in\mathfrak c$ dają $n=0$. Stąd $x=s\in\mathfrak h$ i $\mathfrak c=\mathfrak h$.

Niech $x\in N_{\mathfrak g}(\mathfrak h)$. Zapiszmy $x=x_0+\sum_{\lambda\ne0}x_\lambda$ według Lematu~\ref{aa:lem:toral-weight}. Dla każdego $t\in\mathfrak h$ mamy $[t,x]=\sum_{\lambda\ne0}\lambda(t)x_\lambda\in\mathfrak h=\mathfrak g_0$. Jednoznaczność składowych i wybór $t$ z $\lambda(t)\ne0$ dają $x_\lambda=0$ dla każdej $\lambda\ne0$. Zatem $x\in\mathfrak h$. Odwrotna inkluzja wynika z abelowości, która daje też nilpotentność $\mathfrak h$.
\end{proof}

\begin{uwaga}
W tym rozdziale pracujemy z podalgebrami Cartana otrzymanymi jako maksymalne podalgebry toralne. Odwrotność Twierdzenia~\ref{aa:tw:maximal-toral-cartan} oraz sprzężoność podalgebr Cartana nie są tu potrzebne i pozostają do osobnego dowodu.
\end{uwaga}

\section{Przestrzenie korzeniowe i forma Killinga}

\begin{definicja}
\label{aa:def:lie-root}
Niech $\mathfrak g$ będzie skończenie wymiarową półprostą algebrą Liego nad algebraicznie domkniętym ciałem $K$ charakterystyki zero, a $\mathfrak h$ jej maksymalną podalgebrą toralną. Korzeniem nazywamy niezerowy funkcjonał $\alpha\in\mathfrak h^*$ taki, że $\mathfrak g_\alpha\ne0$, gdzie $\mathfrak g_\alpha$ określono w Lemacie~\ref{aa:lem:toral-weight}. Zbiór korzeni oznaczamy $\Phi$. Rozkład korzeniowy ma postać
\[
 \mathfrak g=\mathfrak h\oplus\bigoplus_{\alpha\in\Phi}\mathfrak g_\alpha.
\]
Wynika on z Twierdzenia~\ref{aa:tw:maximal-toral-cartan}.
\end{definicja}

\begin{twierdzenie}
\label{aa:tw:root-killing-pairing}
Niech $\mathfrak g$ będzie skończenie wymiarową półprostą algebrą Liego nad algebraicznie domkniętym ciałem $K$ charakterystyki zero, $\mathfrak h$ jej maksymalną podalgebrą toralną, a $\Phi$ zbiorem korzeni względem $\mathfrak h$ z Definicji~\ref{aa:def:lie-root}. Korzenie rozpinają $\mathfrak h^*$. Dla każdego $\alpha\in\Phi$ mamy $-\alpha\in\Phi$, a forma Killinga $\kappa$ algebry $\mathfrak g$ określa niezdegenerowane parowanie $\mathfrak g_\alpha\times\mathfrak g_{-\alpha}\to K$. W szczególności $\dim\mathfrak g_\alpha=\dim\mathfrak g_{-\alpha}$.
\end{twierdzenie}
\begin{proof}
Jeśli $t\in\mathfrak h$ znika pod każdym korzeniem, to $\operatorname{ad}_t$ znika na wszystkich przestrzeniach korzeniowych oraz na $\mathfrak h$. Jest więc zerowy i $t=0$ dzięki zerowemu centrum. Gdyby rozpiętość korzeni była właściwą podprzestrzenią $\mathfrak h^*$, istniałby niezerowy $t$ w jej anihilatorze (uzupełniamy bazę tej podprzestrzeni do bazy $\mathfrak h^*$). Sprzeczność.

Dla $0\ne x\in\mathfrak g_\alpha$ niezdegenerowanie $\kappa$ daje $y\in\mathfrak g$ z $\kappa(x,y)\ne0$. Rozkład $y$ i Lemat~\ref{aa:lem:toral-weight} pokazują, że jedyny składnik mogący dać niezerową wartość należy do $\mathfrak g_{-\alpha}$. Dowodzi to istnienia przeciwnego korzenia i zerowości lewego jądra parowania. Stosując argument do $-\alpha$, otrzymujemy zerowość prawego jądra. Iniekcje do odpowiednich przestrzeni dualnych dają obie nierówności wymiarów.
\end{proof}

\begin{twierdzenie}[Trójka typu $\mathfrak{sl}_2$]
\label{aa:tw:root-sl2-triple}
Niech $\mathfrak g$ będzie skończenie wymiarową półprostą algebrą Liego nad algebraicznie domkniętym ciałem $K$ charakterystyki zero, $\mathfrak h$ jej maksymalną podalgebrą toralną, a $\Phi$ zbiorem korzeni względem $\mathfrak h$ z Definicji~\ref{aa:def:lie-root}. Niech $\kappa$ oznacza formę Killinga $\mathfrak g$. Dla $\alpha\in\Phi$ niech $t_\alpha\in\mathfrak h$ będzie jedynym elementem takim, że $\kappa(t_\alpha,t)=\alpha(t)$ dla wszystkich $t\in\mathfrak h$. Wtedy $\alpha(t_\alpha)\ne0$. Dla każdego $0\ne e\in\mathfrak g_\alpha$ istnieje $f\in\mathfrak g_{-\alpha}$ takie, że dla $h_\alpha=2t_\alpha/\alpha(t_\alpha)$ zachodzą
\[
 [h_\alpha,e]=2e,\quad [h_\alpha,f]=-2f,\quad [e,f]=h_\alpha.
\]
Elementy $e,f,h_\alpha$ rozpinają podalgebrę izomorficzną z $\mathfrak{sl}_2(K)$.
\end{twierdzenie}
\begin{proof}
Istnienie i jednoznaczność $t_\alpha$ wynikają z niezdegenerowania $\kappa|_{\mathfrak h\times\mathfrak h}$. Dla $x\in\mathfrak g_\alpha$, $y\in\mathfrak g_{-\alpha}$ mamy $[x,y]\in\mathfrak g_0=\mathfrak h$. Dla $t\in\mathfrak h$ niezmienniczość daje
$\kappa([x,y],t)=\kappa(x,[y,t])=\alpha(t)\kappa(x,y)$.
Niezdegenerowanie na $\mathfrak h$ daje tożsamość $[x,y]=\kappa(x,y)t_\alpha$.

Ustalmy $0\ne x\in\mathfrak g_\alpha$ i wybierzmy $y\in\mathfrak g_{-\alpha}$ z $\kappa(x,y)=1$, co jest możliwe dzięki Twierdzeniu~\ref{aa:tw:root-killing-pairing} i przeskalowaniu. Wtedy $[x,y]=t_\alpha$. Załóżmy, że $\alpha(t_\alpha)=0$. Wówczas $t_\alpha$ komutuje z $x$ i $y$, więc operatory $\operatorname{ad}_x,\operatorname{ad}_y$ zachowują wszystkie przestrzenie własne diagonalizowalnego operatora $T=\operatorname{ad}_{t_\alpha}$. Na każdej niezerowej przestrzeni własnej $W$ o wartości własnej $\nu$ zachodzi
\[
 [\operatorname{ad}_x|_W,\operatorname{ad}_y|_W]=\nu I_W.
\]
Ślad komutatora jest zerowy, więc $\nu\dim W=0$, a charakterystyka zero daje $\nu=0$. Z diagonalizowalności wynika $T=0$ i z zerowego centrum $t_\alpha=0$, co przeczy $\alpha\ne0$.

Dla ustalonego $e$ wybieramy $y$ z $\kappa(e,y)=1$ i kładziemy $f=2y/\alpha(t_\alpha)$. Tożsamość na nawias daje $[e,f]=h_\alpha$, a definicje przestrzeni korzeniowych dają dwie pozostałe równości. Trzy elementy są niezerowe i należą do różnych składowych sumy prostej (charakterystyka zero daje $\alpha\ne-\alpha$), więc są liniowo niezależne. Równości w tezie wraz z antysymetrią i $[v,v]=0$ dla każdego $v$ pokazują, że ich rozpiętość jest zamknięta na nawias. Ich nawiasy są dokładnie nawiasami standardowej bazy $\mathfrak{sl}_2(K)$ z Przykładu~\ref{aa:prz:heisenberg-sl2}. Odwzorowanie tej bazy na $h_\alpha,e,f$ jest zatem izomorfizmem.
\end{proof}

\section{Przykład macierzowy}

\begin{przyklad}[$\mathfrak{sl}_3(K)$]
\label{aa:prz:sl3-root-decomposition}
Niech $\mathfrak g=\mathfrak{sl}_3(K)$, a $\mathfrak h$ będzie przestrzenią macierzy diagonalnych o śladzie zero. Algebra $\mathfrak g$ jest półprosta: z Twierdzenia~\ref{aa:tw:classical-killing} forma Killinga wynosi $6\operatorname{tr}(XY)$. Jeżeli $X$ jest ortogonalna do wszystkich $E_{ij}$ dla $i\ne j$, to wszystkie jej wpisy pozadiagonalne znikają. Ortogonalność do $E_{ii}-E_{jj}$ wymusza równość wpisów diagonalnych; zerowy ślad i charakterystyka zero dają $X=0$. Stosujemy Twierdzenie~\ref{aa:tw:cartan-semisimple}.

Dla $H=\operatorname{diag}(a_1,a_2,a_3)$, gdzie $\sum_i a_i=0$, mamy
$[H,E_{ij}]=(a_i-a_j)E_{ij}$. Macierze diagonalne komutują, a sześć jednostek pozadiagonalnych wraz z bazą $\mathfrak h$ tworzy wspólną bazę własną. Stąd $\mathfrak h$ jest toralna. Jeżeli macierz $X$ komutuje ze wszystkimi $H$, dla każdego $i\ne j$ wybieramy $H$ z $a_i\ne a_j$; równość $[H,X]_{ij}=0$ daje $X_{ij}=0$. Centralizator jest zatem równy $\mathfrak h$, co wyklucza większą abelową podalgebrę.

Korzeniami są $\alpha_{ij}(H)=a_i-a_j$ dla $i\ne j$, a $\mathfrak g_{\alpha_{ij}}=KE_{ij}$. Funkcjonały są parami różne: ich wektory współczynników mają sumę zero i dwa takie wektory określają ten sam funkcjonał na płaszczyźnie $\sum_i a_i=0$ tylko wtedy, gdy ich różnica jest wielokrotnością $(1,1,1)$, a więc jest zerowa. Dla dowolnego $H\in\mathfrak h$ mamy $\kappa((E_{ii}-E_{jj})/6,H)=\operatorname{tr}((E_{ii}-E_{jj})H)=a_i-a_j$. Z jednoznaczności elementu reprezentującego funkcjonał względem formy Killinga dostajemy
$t_{\alpha_{ij}}=(E_{ii}-E_{jj})/6$ i $\alpha_{ij}(t_{\alpha_{ij}})=1/3$, zatem $h_{\alpha_{ij}}=E_{ii}-E_{jj}$. Standardowa trójka to $E_{ij},E_{ji},h_{\alpha_{ij}}$.

W modelu rzeczywistym płaszczyzny współczynników o sumie zero wektory odpowiadające korzeniom tworzą sześciokąt foremny. Rysunek przedstawia układ typu $A_2$ z rozdziału~\ref{aa:ch:root-systems}; nie utożsamia dowolnego ciała $K$ z podciałem $\mathbb R$.
\begin{center}
\begin{tikzpicture}[scale=1.45, every node/.style={font=\small}]
\foreach \ang/\lab in {0/{\alpha_{12}},60/{\alpha_{13}},120/{\alpha_{23}},180/{\alpha_{21}},240/{\alpha_{31}},300/{\alpha_{32}}}{
 \draw[-{Stealth}] (0,0)--(\ang:1);
 \node at (\ang:1.32) {$\lab$};
}
\fill (0,0) circle (0.025);
\end{tikzpicture}
\end{center}
Na przykład $[E_{12},E_{23}]=E_{13}$, gdyż $E_{12}E_{23}=E_{13}$ i $E_{23}E_{12}=0$. Odpowiada to równości $\alpha_{12}+\alpha_{23}=\alpha_{13}$. Z kolei $[E_{12},E_{13}]=0$, ponieważ obydwa iloczyny są zerowe.
\end{przyklad}

Wykazaliśmy istnienie rozkładu korzeniowego, parowanie przeciwległych przestrzeni oraz istnienie trójek typu $\mathfrak{sl}_2$. Jednowymiarowość przestrzeni korzeniowych, całkowitość liczb Cartana, ciągi korzeniowe i realizacja aksjomatów układu korzeni z rozdziału~\ref{aa:ch:root-systems} wymagają dalszych argumentów z teorii reprezentacji $\mathfrak{sl}_2$; nie zostały tutaj założone ani dowiedzione.

% AA-60
\chapter[Reprezentacje sl2 i układy korzeni]{Reprezentacje \texorpdfstring{$\mathfrak{sl}_2$}{sl2} i struktura układów korzeni}
\label{aa:ch:sl2-root-systems}

\providecommand{\atlasref}[2]{#2}


W rozdziale~\ref{aa:ch:cartan-root-decomposition} skonstruowaliśmy podalgebrę izomorficzną z $\mathfrak{sl}_2$ dla każdego korzenia. Teraz zbadamy jej reprezentacje i zastosujemy wyniki do reprezentacji dołączonej. W całym rozdziale $K$ jest algebraicznie domkniętym ciałem charakterystyki zero. Liczby całkowite utożsamiamy z ich obrazami w $K$.

\section{Reprezentacje nieprzywiedlne}
Niech $h,e,f$ będzie standardową bazą $\mathfrak{sl}_2(K)$, z relacjami $[h,e]=2e$, $[h,f]=-2f$, $[e,f]=h$. Operatory w reprezentacji oznaczamy przez $H,E,F$.

\begin{lemat}
\label{aa:lem:sl2-commutation}
Niech $K$ będzie algebraicznie domkniętym ciałem charakterystyki zero, a $\rho:\mathfrak{sl}_2(K)\to\mathfrak{gl}(V)$ reprezentacją. Oznaczmy $H=\rho(h)$, $E=\rho(e)$, $F=\rho(f)$, gdzie $h,e,f$ są standardowymi generatorami o relacjach podanych wyżej. Dla każdego całkowitego $j\ge1$ zachodzą
\[
 [H,F^j]=-2jF^j,\qquad [E,F^j]=jF^{j-1}(H-(j-1)I).
\]
\end{lemat}
\begin{proof}
Dla $j=1$ są to relacje definiujące. Z tożsamości $[A,BC]=[A,B]C+B[A,C]$ pierwszy wzór przechodzi z $j$ na $j+1$ przez dodanie $F^j[H,F]$. Dla drugiego wzoru otrzymujemy
\[
\begin{aligned}
 [E,F^{j+1}]&=jF^{j-1}(H-(j-1)I)F+F^jH\\
 &=jF^j(H-(j+1)I)+F^jH\\
 &=(j+1)F^j(H-jI),
\end{aligned}
\]
gdzie użyto $HF=FH-2F$.
\end{proof}

\begin{twierdzenie}
\label{aa:tw:sl2-simple-modules}
Niech $K$ będzie algebraicznie domkniętym ciałem charakterystyki zero. Każda niezerowa skończenie wymiarowa nieprzywiedlna reprezentacja $\mathfrak{sl}_2(K)$ jest izomorficzna z dokładnie jednym modułem $L(m)$, gdzie $m$ jest nieujemną liczbą całkowitą, o bazie $v_0,\ldots,v_m$ i działaniach
\[
 Hv_j=(m-2j)v_j,\quad Fv_j=v_{j+1},\quad
 Ev_j=j(m-j+1)v_{j-1},
\]
gdzie $v_{-1}=v_{m+1}=0$. Dla każdego $m\ge0$ wzory te definiują reprezentację nieprzywiedlną.
\end{twierdzenie}
\begin{proof}
Niech $V$ będzie taką reprezentacją. Skończone niepuste widmo $H$ zawiera $\lambda$ takie, że $\lambda+2$ nie należy do widma: w przeciwnym razie iterowane dodawanie $2$ dawałoby nieskończenie wiele wartości. Wybieramy $0\ne v$ z $Hv=\lambda v$. Relacja $HE=EH+2E$ daje $H(Ev)=(\lambda+2)Ev$, stąd $Ev=0$. Lemat~\ref{aa:lem:sl2-commutation} daje
\[
 HF^jv=(\lambda-2j)F^jv,\qquad
 EF^jv=j(\lambda-j+1)F^{j-1}v.
\]
Dla $j=0$ mamy $Hv=\lambda v$ i $Ev=0$; drugi wzór stosujemy dla $j\ge1$. Niezerowe $F^jv$ są liniowo niezależne jako wektory własne dla różnych wartości. Istnieje więc pierwszy indeks $m+1$ z $F^{m+1}v=0$. Stosując $E$, otrzymujemy $(m+1)(\lambda-m)F^mv=0$, czyli $\lambda=m$. Rozpiętość $v,Fv,\ldots,F^mv$ jest niezmiennicza względem wszystkich trzech operatorów, a więc równa $V$. Jednoznaczność $m$ wynika z $\dim V=m+1$.

Odwrotnie, w podanej bazie $[H,E]v_j=2Ev_j$, $[H,F]v_j=-2Fv_j$ oraz
$[E,F]v_j=((j+1)(m-j)-j(m-j+1))v_j=(m-2j)v_j$; wzory obejmują końce dzięki zerowym współczynnikom. Niezerowa podprzestrzeń niezmiennicza zawiera któryś $v_j$, ponieważ wielomianowe projekcje na przestrzenie własne $H$ ją zachowują. Kolejne zastosowania $E$ prowadzą do niezerowej wielokrotności $v_0$, gdyż współczynniki dla $1\le j\le m$ są niezerowe. Zastosowania $F$ dają całą bazę.
\end{proof}

\begin{przyklad}[Wielomiany jednorodne]
Na wielomianach jednorodnych stopnia $m$ w zmiennych $x,y$ kładziemy $H=x\partial_x-y\partial_y$, $E=x\partial_y$, $F=y\partial_x$. Dla $v_j=\frac{m!}{(m-j)!}x^{m-j}y^j$ otrzymujemy $Hv_j=(m-2j)v_j$, $Fv_j=v_{j+1}$ dla $j<m$, a $Fv_m=0$. Różniczkowanie po $y$ daje dla $j\ge1$ współczynnik $j(m-j+1)$ przy $v_{j-1}$, a $Ev_0=0$. Jest to model $L(m)$.
\end{przyklad}

\section{Pełna redukowalność}
\begin{lemat}[Operator Casimira]
\label{aa:lem:sl2-casimir}
Niech $K$ będzie algebraicznie domkniętym ciałem charakterystyki zero, a $H,E,F$ operatorami dowolnej reprezentacji $\mathfrak{sl}_2(K)$ odpowiadającymi $h,e,f$. Operator $C=H^2+2H+4FE$ komutuje z $H,E,F$. Na $L(m)$ działa przez $m(m+2)I$.
\end{lemat}
\begin{proof}
Mamy $[H,FE]=-2FE+2FE=0$ oraz
\[
\begin{aligned}
 [E,C]&=-2EH-2HE-4E+4HE=2[H,E]-4E=0,\\
 [F,C]&=2FH+2HF+4F-4FH=2[H,F]+4F=0.
\end{aligned}
\]
Na $v_0$ operator $C$ działa przez $m^2+2m$, bo $Ev_0=0$. Komutowanie z $F$ daje tę samą wartość na każdym $v_j=F^jv_0$.
\end{proof}

\begin{twierdzenie}
\label{aa:tw:sl2-complete-reducibility}
Niech $K$ będzie algebraicznie domkniętym ciałem charakterystyki zero. Każda skończenie wymiarowa reprezentacja $\mathfrak{sl}_2(K)$ w przestrzeni $V$ jest sumą prostą modułów $L(m)$, gdzie $m$ przebiega nieujemne liczby całkowite. Ich krotności są jednoznaczne.
\end{twierdzenie}
\begin{proof}
Indukujemy po wymiarze $V$, zaczynając od pustej sumy dla $V=0$. Wybieramy właściwy podmoduł $U$ największego wymiaru. Iloraz $V/U$ jest nieprzywiedlny, ponieważ przeciwobraz właściwego niezerowego podmodułu przeczyłby maksymalności. Zatem $V/U\simeq L(m)$, a z indukcji $U$ jest sumą prostą modułów $L(k)$.

Niech $q:V\to V/U$ będzie rzutem, $c=m(m+2)$, a $W$ uogólnioną przestrzenią własną operatora $C=H^2+2H+4FE$ dla $c$ (przyjmujemy przestrzeń zerową dla wartości spoza widma). Wszystkie uogólnione przestrzenie własne $C$ są podmodułami dzięki komutowaniu z $H,E,F$. Jeśli $(C-dI)^rx=0$ dla $d\ne c$, to $(c-d)^rq(x)=0$, więc $q(x)=0$. Z rozkładu na uogólnione przestrzenie własne wynika $q(W)=V/U$. Ponieważ $k(k+2)=m(m+2)$ dla nieujemnych liczb całkowitych zachodzi tylko dla $k=m$, podmoduł $U\cap W$ jest sumą kopii $L(m)$.

Widmo $H|_W$ zawiera się w $\{m,m-2,\ldots,-m\}$: w bazie dostosowanej do $U\cap W$ macierz jest blokowo górnotrójkątna, z blokami odpowiadającymi podmodułowi i ilorazowi. Wybieramy przeciwobraz $v$ wektora najwyższej wagi $v_0$ ilorazu w uogólnionej przestrzeni własnej $H|_W$ dla $m$. Taki przeciwobraz istnieje, ponieważ rozkład dowolnego przeciwobrazu na uogólnione przestrzenie własne ma obrazy zerowe poza składową dla $m$.

Relacje $(H-(\lambda+2)I)E=E(H-\lambda I)$ oraz $(H-(\lambda-2)I)F=F(H-\lambda I)$ pokazują po iteracji, że $E,F$ przesuwają uogólnione wagi odpowiednio o $2,-2$. Zatem $Ev=0$ oraz $F^{m+1}v=0$, bo wagi $m+2,-m-2$ nie występują. Wektor $u=(H-mI)v$ należy do $U\cap W$ i do uogólnionej przestrzeni własnej dla $m$. Na $U\cap W$ operator $H$ jest diagonalizowalny, więc $u$ jest zwykłym wektorem własnym lub zerem. Z Lematu~\ref{aa:lem:sl2-commutation} wynika
\[
 0=EF^{m+1}v=(m+1)F^mu.
\]
Na przestrzeni wagi $m$ w sumie kopii $L(m)$ operator $F^m$ jest iniektywny, gdyż na każdym składniku przekształca $v_0$ w $v_m$. Stąd $u=0$, czyli $Hv=mv$.

Wektory $v,Fv,\ldots,F^mv$ mają obrazy tworzące bazę ilorazu, więc są liniowo niezależne. Ich rozpiętość $L$ jest podmodułem dzięki wzorom komutacyjnym i $F^{m+1}v=0$. Rzut $q|_L$ jest izomorfizmem, zatem $V=U\oplus L$, gdzie $L\simeq L(m)$. To zamyka indukcję.

Jeśli $a_j$ jest wymiarem przestrzeni własnej $H$ dla $j\ge0$ (zero dla nieobecnej wagi), to każdy $L(k)$ wnosi jeden do $a_j$ dokładnie dla $k\ge j$ i $k-j$ parzystego. Zatem $a_j-a_{j+2}$ jest krotnością $L(j)$, niezależną od rozkładu.
\end{proof}

\section{Jednowymiarowość przestrzeni korzeniowych}
Niech odtąd $\mathfrak g$ będzie skończenie wymiarową półprostą algebrą Liego nad $K$, $\mathfrak h$ jej maksymalną podalgebrą toralną, a $\Phi$ zbiorem korzeni. Stosujemy oznaczenia $t_\alpha,h_\alpha,e_\alpha,f_\alpha$ z Twierdzenia~\ref{aa:tw:root-sl2-triple} oraz $\mathfrak g_0=\mathfrak h$, $\mathfrak g_\gamma=0$ dla niezerowych funkcjonałów niebędących korzeniami.

\begin{twierdzenie}
\label{aa:tw:root-dim-reduced}
Niech $\mathfrak g$ będzie skończenie wymiarową półprostą algebrą Liego nad algebraicznie domkniętym ciałem $K$ charakterystyki zero, $\mathfrak h$ jej maksymalną podalgebrą toralną, a $\Phi$ zbiorem korzeni względem $\mathfrak h$. Dla każdego $\alpha\in\Phi$ mamy $\dim\mathfrak g_\alpha=1$ oraz $K\alpha\cap\Phi=\{\alpha,-\alpha\}$. Dla wszystkich $\alpha,\beta\in\Phi$ liczba $\beta(h_\alpha)$ jest całkowita.
\end{twierdzenie}
\begin{proof}
Ograniczamy reprezentację dołączoną $\mathfrak g$ do podalgebry $\mathfrak{sl}_2$ związanej z $\alpha$. Z pełnej redukowalności i klasyfikacji jej modułów wszystkie wartości własne $\operatorname{ad}_{h_\alpha}$ są całkowite. Na $\mathfrak g_\beta$ wartością własną jest $\beta(h_\alpha)$.

Podmoduł $M=\mathfrak h\oplus\bigoplus_{k\in\mathbb Z\setminus\{0\}}\mathfrak g_{k\alpha}$ jest niezmienniczy na mocy relacji wagowych z AA-59. Jego wagi względem $h_\alpha$ są parzyste, równe $2k$, a przestrzenią wagi zero jest $\mathfrak h$. Rozkładamy $M$ na sumę prostą modułów nieprzywiedlnych. Ich najwyższe wagi są parzyste, bo występują w widmie $H|_M$. Każdy nietrywialny składnik nieprzywiedlny $L(2r)$ ma jednowymiarową przestrzeń wagi zero, na której $E$ działa niezerowo ze współczynnikiem $r(r+1)$. Składniki $L(0)$ są zerowane. Zatem liczba nietrywialnych składników równa się randze $\operatorname{ad}_{e_\alpha}|_{\mathfrak h}$, która wynosi jeden, ponieważ $[e_\alpha,t]=-\alpha(t)e_\alpha$.

Podalgebra związana z trójką jest podmodułem izomorficznym z $L(2)$: jej reprezentacja dołączona jest nieprzywiedlna dzięki prostocie $\mathfrak{sl}_2$, wykazanej w Przykładzie~\ref{aa:prz:lie-radical-examples}, i ma wymiar trzy. Rzuty tego podmodułu na składniki $L(0)$ są zerowe (niezerowy homomorfizm między nieprzywiedlnymi modułami jest izomorfizmem, gdyż jego jądro i obraz są podmodułami). Zatem jego rzut na jedyny nietrywialny składnik jest niezerowy i daje izomorfizm. Składnik jest więc $L(2)$. Wynika stąd $\dim\mathfrak g_\alpha=1$ i brak korzeni $k\alpha$ dla całkowitych $|k|\ge2$.

Jeśli $\beta=c\alpha\in\Phi$ dla $c\in K^\times$, to $t_\beta=ct_\alpha$ z jednoznaczności reprezentanta względem formy Killinga, a $h_\beta=c^{-1}h_\alpha$. Liczby $2c=\beta(h_\alpha)$ i $2/c=\alpha(h_\beta)$ są całkowite, a ich iloczyn to $4$. Stąd $c\in\{\pm1,\pm2,\pm1/2\}$. Przypadki $\pm2$ wykluczono, a $\pm1/2$ wykluczamy przez zastosowanie tego samego wyniku do $\beta$, gdyż wtedy $\alpha=\pm2\beta$. Korzenie $\alpha,-\alpha$ rzeczywiście należą do $\Phi$ na mocy Twierdzenia~\ref{aa:tw:root-killing-pairing}.
\end{proof}

\section{Ciągi korzeniowe}
\begin{twierdzenie}
\label{aa:tw:lie-root-string}
Niech $\mathfrak g$ będzie skończenie wymiarową półprostą algebrą Liego nad algebraicznie domkniętym ciałem $K$ charakterystyki zero, $\mathfrak h$ jej maksymalną podalgebrą toralną, a $\Phi$ zbiorem korzeni względem $\mathfrak h$. Dla $\alpha,\beta\in\Phi$ z $\beta\notin K\alpha$ istnieją jednoznaczne całkowite $p,q\ge0$ takie, że
\[
 \{k\in\mathbb Z:\beta+k\alpha\in\Phi\}=\{-p,\ldots,q\},
 \qquad \beta(h_\alpha)=p-q.
\]
Dla każdego $k\in\mathbb Z$, jeśli $\beta+k\alpha$ i $\beta+(k+1)\alpha$ są korzeniami, to $[e_\alpha,\mathfrak g_{\beta+k\alpha}]=\mathfrak g_{\beta+(k+1)\alpha}$.
\end{twierdzenie}
\begin{proof}
Suma $W=\bigoplus_{k\in\mathbb Z}\mathfrak g_{\beta+k\alpha}$ jest niezerowym podmodułem względem trójki związanej z $\alpha$, na mocy Lematu~\ref{aa:lem:toral-weight}. Żaden funkcjonał $\beta+k\alpha$ nie jest zerowy. Wagi operatora $\operatorname{ad}_{h_\alpha}|_W$ mają postać $b+2k$, gdzie $b=\beta(h_\alpha)$ jest całkowite; są parami różne, a każda występująca przestrzeń wagowa jest jednowymiarowa. Rozkładamy $W$ na moduły $L(m)$. Wszystkie ich najwyższe wagi mają tę samą parzystość co $b$. Dwa składniki miałyby wspólną wagę $0$ przy parzystym $b$ lub $1$ przy nieparzystym $b$, co przeczy jednowymiarowości. Stąd $W$ jest pojedynczym $L(m)$. Wagi tworzą ciąg od $-m$ do $m$ co dwa. Ponieważ $k=0$ występuje, indeksy tworzą przedział całkowity od $-p$ do $q$. Równania $b-2p=-m$, $b+2q=m$ dają $b=p-q$. Wzór na działanie $E$ mówi, że między sąsiednimi wagami działa niezerowo, poza najwyższą wagą, zatem jego obraz jest całą docelową prostą.
\end{proof}

\begin{wniosek}
\label{aa:cor:lie-root-reflection}
Dla $\alpha\in\Phi$ odwzorowanie $s_\alpha(\beta)=\beta-\beta(h_\alpha)\alpha$ na $\mathfrak h^*$ jest inwolucją zachowującą $\Phi$. Jeśli $\alpha,\beta,\alpha+\beta\in\Phi$, to $[\mathfrak g_\alpha,\mathfrak g_\beta]=\mathfrak g_{\alpha+\beta}$.
\end{wniosek}
\begin{proof}
Ponieważ $\alpha(h_\alpha)=2$, mamy $s_\alpha(\beta)(h_\alpha)=-\beta(h_\alpha)$, a więc $s_\alpha^2=I$. Dla korzeni proporcjonalnych do $\alpha$ odbicie zamienia $\alpha$ i $-\alpha$. W pozostałych przypadkach $-\beta(h_\alpha)=q-p$ leży w przedziale $[-p,q]$, więc $s_\alpha(\beta)\in\Phi$. Inwolucja daje równość $s_\alpha(\Phi)=\Phi$. Jeśli $\alpha+\beta\in\Phi$, to $\beta$ nie jest proporcjonalna do $\alpha$, ponieważ suma z $\pm\alpha$ daje zero lub $2\alpha$. Twierdzenie o ciągu korzeniowym dla $k=0$ i $\mathfrak g_\alpha=Ke_\alpha$ dają tezę o nawiasie.
\end{proof}

\section{Realizacja euklidesowa}

\noindent\emph{Powiązania tematyczne w Bibliotece matematycznej:} \atlasref{LA-12::ch:lin12-euklidesowe-unitarne}{LA-12: przestrzenie euklidesowe}.

\begin{twierdzenie}
\label{aa:tw:lie-root-realization}
Niech $\mathfrak g$ będzie niezerową skończenie wymiarową półprostą algebrą Liego nad algebraicznie domkniętym ciałem $K$ charakterystyki zero, $\mathfrak h$ jej maksymalną podalgebrą toralną, a $\Phi$ zbiorem korzeni względem $\mathfrak h$. Istnieje różnowartościowe odwzorowanie $\alpha\mapsto\bar\alpha$ zbioru $\Phi$ do rzeczywistej przestrzeni euklidesowej, którego obraz jest zredukowanym krystalograficznym układem korzeni w sensie Definicji~\ref{aa:def:root-system}, z relacją
\[
 \frac{2(\bar\beta,\bar\alpha)}{(\bar\alpha,\bar\alpha)}
 =\beta(h_\alpha)\qquad(\alpha,\beta\in\Phi).
\]
\end{twierdzenie}
\begin{proof}
Korzenie rozpinają $\mathfrak h^*$, a $\mathfrak h=0$ implikowałoby $\mathfrak g=0$ przez rozkład korzeniowy. Wybieramy korzenie $\alpha_1,\ldots,\alpha_r$ tworzące bazę $\mathfrak h^*$. Z niezdegenerowania formy Killinga na $\mathfrak h$ elementy $t_{\alpha_i}$ tworzą bazę $\mathfrak h$; przez przeskalowanie otrzymujemy bazę $h_i=h_{\alpha_i}$.

Niech $\kappa$ będzie formą Killinga $\mathfrak g$, a $B=(\kappa(h_i,h_j))_{i,j=1}^r$. Z rozkładu korzeniowego i jednowymiarowości przestrzeni korzeniowych mamy
\[
 B_{ij}=\sum_{\gamma\in\Phi}\gamma(h_i)\gamma(h_j).
\]
Wpisy są całkowite. Traktujemy $B$ jako macierz rzeczywistą. Dla $u\in\mathbb R^r$ zachodzi $u^{\mathsf T}Bu=\sum_\gamma(\sum_i u_i\gamma(h_i))^2$. Macierz $(\alpha_j(h_i))_{j,i}$ jest odwracalna nad $K$, bo opisuje parowanie dwóch baz. Ma całkowite wpisy, więc jej wyznacznik jest niezerową liczbą całkowitą i jest odwracalna nad $\mathbb R$. Dla $u\ne0$ co najmniej jeden składnik sumy kwadratów jest dodatni. Zatem $B$ jest dodatnio określona.

Na $V=\mathbb R^r$ kładziemy $(u,v)=u^{\mathsf T}B^{-1}v$. Macierz $B^{-1}$ jest symetryczna; dla $u\ne0$ wektor $w=B^{-1}u$ jest niezerowy i $u^{\mathsf T}B^{-1}u=w^{\mathsf T}Bw>0$. Jest to więc iloczyn skalarny. Definiujemy $\bar\gamma=(\gamma(h_i))_{i=1}^r$ jako rzeczywisty wektor kolumnowy. Różnowartościowość i niezerowość wynikają z tego, że $h_i$ jest bazą; wektory $\bar\alpha_i$ rozpinają $V$ dzięki powyższemu wyznacznikowi. Współrzędne $t_\alpha$ w bazie $h_i$ to $B^{-1}\bar\alpha$, ponieważ $\kappa(t_\alpha,h_j)=\alpha(h_j)$. Są wymierne, a rachunek na nich w $K$ zgadza się z rachunkiem rzeczywistym. Stąd $\beta(t_\alpha)=\bar\beta^{\mathsf T}B^{-1}\bar\alpha$. Wstawiając $h_\alpha=2t_\alpha/\alpha(t_\alpha)$, otrzymujemy żądany iloraz.

Odbicie euklidesowe względem hiperpłaszczyzny prostopadłej do $\bar\alpha$ odpowiada zatem $s_\alpha$, więc zachowuje obraz $\Phi$. Liczby Cartana są całkowite. Jeśli dwa wektory korzeni są proporcjonalne nad $\mathbb R$, współczynnik proporcjonalności jest wymierny jako iloraz dwóch niezerowych współrzędnych całkowitych. Korzenie są więc proporcjonalne również nad $K$, a współczynnik musi być równy $\pm1$. Wykazano wszystkie aksjomaty.
\end{proof}

\begin{przyklad}[Ciąg korzeniowy w $\mathfrak{sl}_3$]
W oznaczeniach Przykładu~\ref{aa:prz:sl3-root-decomposition} bierzemy $\alpha=\alpha_{12}$, $\beta=\alpha_{23}$. Mamy $\beta(h_\alpha)=-1$, $\beta+\alpha=\alpha_{13}$, natomiast $\beta-\alpha$ i $\beta+2\alpha$ nie są korzeniami. Zatem $p=0,q=1$. Nawiasy $[E_{12},E_{23}]=E_{13}$ i $[E_{21},E_{13}]=E_{23}$ pokazują, że te dwie przestrzenie tworzą $L(1)$.
\begin{center}
\begin{tikzpicture}[every node/.style={font=\small}]
\node (a) at (-2,0) {$\mathfrak g_\beta$};
\node (b) at (2,0) {$\mathfrak g_{\beta+\alpha}$};
\draw[-{Stealth},bend left=20] (a) to node[above] {$\operatorname{ad}_{e_\alpha}$} (b);
\draw[-{Stealth},bend left=20] (b) to node[below] {$\operatorname{ad}_{f_\alpha}$} (a);
\node at (-2,-.9) {waga $-1$};
\node at (2,-.9) {waga $1$};
\end{tikzpicture}
\end{center}
Odbicie $s_\alpha$ zamienia $\beta$ z $\beta+\alpha$.
\end{przyklad}

Teoria z rozdziału~\ref{aa:ch:root-systems} stosuje się do otrzymanego układu korzeni. Klasyfikacja układów korzeni nie jest jeszcze klasyfikacją algebr Liego: potrzebne są dowody odtworzenia algebry z danych korzeniowych i istnienia odpowiednich algebr.

% AA-61
\chapter{Generatory półprostych algebr Liego i relacje Serre'a}
\label{aa:ch:lie-serre-relations}

W AA-59 i AA-60 wyprowadziliśmy z półprostoty rozkład korzeniowy, jednowymiarowość przestrzeni korzeniowych i aksjomaty układu korzeni. Teraz wykorzystamy bazę układu z AA-56, aby opisać generatory algebry i rozkład na ideały proste. Nie zakładamy twierdzenia o prezentacji Serre'a: na razie wyprowadzamy relacje spełniane przez generatory już istniejącej algebry.

W całym rozdziale $K$ jest algebraicznie domkniętym ciałem charakterystyki zero, $\mathfrak g$ niezerową skończenie wymiarową półprostą algebrą Liego nad $K$, $\mathfrak h$ jej maksymalną podalgebrą toralną, a $\Phi$ zbiorem korzeni względem $\mathfrak h$. Korzystamy z realizacji rzeczywistej $\alpha\mapsto\bar\alpha$ z Twierdzenia~\ref{aa:tw:lie-root-realization}. Ponieważ realizacja wykorzystuje współrzędne względem bazy $\mathfrak h$, zachowuje ona wszystkie zależności liniowe o współczynnikach wymiernych. Wybieramy bazę układu rzeczywistego i oznaczamy odpowiadające jej korzenie przez $\Delta=(\alpha_1,\ldots,\alpha_r)$.

\section{Korzenie proste i rozkład trójkątny}

\begin{lemat}
\label{aa:lem:lie-simple-root-basis}
Niech $\mathfrak g$ będzie niezerową skończenie wymiarową półprostą algebrą Liego nad algebraicznie domkniętym ciałem $K$ charakterystyki zero, $\mathfrak h$ jej maksymalną podalgebrą toralną, $\Phi$ zbiorem korzeni, a $\Delta=(\alpha_1,\ldots,\alpha_r)$ bazą układu korzeni w realizacji z Twierdzenia~\ref{aa:tw:lie-root-realization}. Elementy $\alpha_1,\ldots,\alpha_r$ są bazą $\mathfrak h^*$ nad $K$. Elementy $h_i=h_{\alpha_i}$ z Twierdzenia~\ref{aa:tw:root-sl2-triple} tworzą bazę $\mathfrak h$. Każdy korzeń $\beta$ ma jednoznaczny zapis $\beta=\sum_i n_i\alpha_i$, w którym $n_i\in\mathbb Z$ są albo wszystkie nieujemne, albo wszystkie niedodatnie.
\end{lemat}
\begin{proof}
W konstrukcji realizacji z AA-60 wszystkie korzenie mają całkowite współrzędne względem bazy $\mathfrak h$ złożonej z elementów $h_\gamma$. Macierz kolumn odpowiadających $\bar\alpha_i$ jest kwadratowa i ma niezerowy wyznacznik nad $\mathbb R$, ponieważ tworzą one bazę przestrzeni rzeczywistej. Jej wyznacznik jest liczbą całkowitą, więc nie znika również nad $K$. Dowodzi to pierwszej tezy. Izomorfizm $\mathfrak h\to\mathfrak h^*$ zadany przez niezdegenerowaną formę Killinga przekształca $t_{\alpha_i}$ w $\alpha_i$. Zatem $t_{\alpha_i}$ tworzą bazę, podobnie jak ich niezerowe wielokrotności $h_i$. Całkowite współrzędne korzeni o odpowiednich znakach wynikają z Definicji~\ref{aa:def:root-base}; równości między wektorami całkowitymi przenoszą się do $K$. Jednoznaczność wynika z niezależności liniowej.
\end{proof}

\begin{definicja}
\label{aa:def:lie-positive-height}
Dla ustalonej bazy $\Delta$ z Lematu~\ref{aa:lem:lie-simple-root-basis} korzeń $\beta=\sum_i n_i\alpha_i$ nazywamy dodatnim, jeśli wszystkie $n_i\ge0$, i ujemnym, jeśli wszystkie $n_i\le0$. Zbiory takich korzeni oznaczamy przez $\Phi^+,\Phi^-$. Wysokością korzenia $\beta$ nazywamy liczbę całkowitą $\operatorname{ht}(\beta)=\sum_i n_i$. Dla każdego $\beta\in\Phi^+$ wysokość jest dodatnia, gdyż $\beta\ne0$.
\end{definicja}

\begin{lemat}[Odejmowanie korzenia prostego]
\label{aa:lem:lie-root-descent}
Niech $\mathfrak g$ będzie niezerową skończenie wymiarową półprostą algebrą Liego nad algebraicznie domkniętym ciałem $K$ charakterystyki zero, $\mathfrak h$ jej maksymalną podalgebrą toralną, $\Phi$ zbiorem korzeni, a $\Delta=(\alpha_1,\ldots,\alpha_r)$ bazą układu korzeni w realizacji z Twierdzenia~\ref{aa:tw:lie-root-realization}. Jeśli $\beta\in\Phi^+\setminus\Delta$, istnieje $\alpha_i\in\Delta$, dla którego $\beta-\alpha_i\in\Phi^+$ i $\operatorname{ht}(\beta-\alpha_i)=\operatorname{ht}(\beta)-1$.
\end{lemat}
\begin{proof}
Zapiszmy $\beta=\sum_i n_i\alpha_i$. W rzeczywistym modelu
\[
 0<(\bar\beta,\bar\beta)=\sum_i n_i(\bar\beta,\bar\alpha_i).
\]
Istnieje więc $i$ z $n_i>0$ i $(\bar\beta,\bar\alpha_i)>0$. Korzenie $\beta,\alpha_i$ nie są proporcjonalne, bo układ jest zredukowany, a $\beta$ nie jest prosty. Z Twierdzenia~\ref{aa:tw:lie-root-realization} liczba $b=\beta(h_i)$ jest obrazem dodatniej liczby całkowitej w $K$; nierówność dotyczy tej liczby całkowitej, a nie porządku w $K$. W Twierdzeniu~\ref{aa:tw:lie-root-string} mamy $b=p-q>0$, skąd $p\ge1$. Zatem $\beta-\alpha_i$ jest korzeniem; jego współczynniki w bazie są nieujemne, gdyż $n_i>0$. Suma współczynników zmniejsza się o jeden.
\end{proof}

\begin{twierdzenie}[Rozkład trójkątny]
\label{aa:tw:lie-triangular-decomposition}
Niech $\mathfrak g$ będzie niezerową skończenie wymiarową półprostą algebrą Liego nad algebraicznie domkniętym ciałem $K$ charakterystyki zero, $\mathfrak h$ jej maksymalną podalgebrą toralną, $\Phi$ zbiorem korzeni, a $\Delta=(\alpha_1,\ldots,\alpha_r)$ bazą układu korzeni w realizacji z Twierdzenia~\ref{aa:tw:lie-root-realization}. Niech $\mathfrak n^+=\bigoplus_{\alpha\in\Phi^+}\mathfrak g_\alpha$ oraz $\mathfrak n^-=\bigoplus_{\alpha\in\Phi^-}\mathfrak g_\alpha$. Są to nilpotentne podalgebry i
\[
 \mathfrak g=\mathfrak n^-\oplus\mathfrak h\oplus\mathfrak n^+
\]
jako suma prosta przestrzeni liniowych. Podalgebry $\mathfrak b^\pm=\mathfrak h\oplus\mathfrak n^\pm$ są rozwiązalne, a $\mathfrak n^\pm$ są w nich ideałami.
\end{twierdzenie}
\begin{proof}
Nawias przestrzeni o korzeniach $\alpha,\beta$ należy do przestrzeni o korzeniu $\alpha+\beta$ lub znika, jeśli suma nie jest korzeniem i nie jest zerem. Dla dwóch dodatnich korzeni suma, jeśli jest korzeniem, jest dodatnia; analogicznie dla ujemnych. Otrzymujemy podalgebry. Suma prosta wynika z rozkładu korzeniowego.

Z niezerowości $\mathfrak g$ i rozkładu korzeniowego oraz faktu, że korzenie rozpinają $\mathfrak h^*$, wynika $\Phi\ne\varnothing$. Zatem $\Phi^+\ne\varnothing$. Niech $M=\max_{\alpha\in\Phi^+}\operatorname{ht}(\alpha)$. Dla $k\ge1$ połóżmy $N_k=\bigoplus_{\alpha\in\Phi^+,\ \operatorname{ht}(\alpha)\ge k}\mathfrak g_\alpha$. Mamy $[N_i,N_j]\subset N_{i+j}$, ponieważ wysokość jest addytywna. Jeżeli $C^1=\mathfrak n^+$ i $C^{k+1}=[\mathfrak n^+,C^k]$, to indukcyjnie $C^k\subset N_k$. Ponieważ $N_{M+1}=0$, szereg centralny zstępujący kończy się zerem. Dla $\mathfrak n^-$ stosujemy tę samą argumentację do wysokości $-\operatorname{ht}$.

Ponadto $[\mathfrak h,\mathfrak n^\pm]\subset\mathfrak n^\pm$ i $[\mathfrak h,\mathfrak h]=0$. Stąd $\mathfrak b^\pm$ są podalgebrami z ideałami $\mathfrak n^\pm$, a $[\mathfrak b^\pm,\mathfrak b^\pm]\subset\mathfrak n^\pm$. Nilpotentność implikuje rozwiązalność na mocy Lematu~\ref{aa:lem:lie-series}; jeżeli $D^0(\mathfrak l)=\mathfrak l$ oraz $D^{k+1}(\mathfrak l)=[D^k(\mathfrak l),D^k(\mathfrak l)]$, to $D^1(\mathfrak b^\pm)\subset D^0(\mathfrak n^\pm)$, a indukcja daje $D^{k+1}(\mathfrak b^\pm)\subset D^k(\mathfrak n^\pm)$ dla każdego $k\ge0$. Zatem $\mathfrak b^\pm$ są rozwiązalne.
\end{proof}

\begin{center}
\begin{tikzpicture}[every node/.style={font=\small}]
\node[draw,rounded corners] (m) at (-3.2,0) {$\mathfrak n^-$};
\node[draw,rounded corners] (h) at (0,0) {$\mathfrak h$};
\node[draw,rounded corners] (p) at (3.2,0) {$\mathfrak n^+$};
\node at (-3.2,-.8) {ujemne wysokości};
\node at (0,-.8) {waga zero};
\node at (3.2,-.8) {dodatnie wysokości};
\node at (-1.6,0) {$\oplus$};
\node at (1.6,0) {$\oplus$};
\end{tikzpicture}

\small Rozkład przestrzeni liniowej; składniki nie muszą być ideałami $\mathfrak g$.
\end{center}

\section{Generatory związane z korzeniami prostymi}

Dla każdego $i$ wybieramy $0\ne e_i\in\mathfrak g_{\alpha_i}$ i element $f_i\in\mathfrak g_{-\alpha_i}$ taki, że $[e_i,f_i]=h_i$. Istnienie zapewnia Twierdzenie~\ref{aa:tw:root-sl2-triple}; jednoznaczność $f_i$ wynika z jednowymiarowości przestrzeni i niezerowości nawiasu.

\begin{twierdzenie}
\label{aa:tw:lie-simple-generators}
Niech $\mathfrak g$ będzie niezerową skończenie wymiarową półprostą algebrą Liego nad algebraicznie domkniętym ciałem $K$ charakterystyki zero, $\mathfrak h$ jej maksymalną podalgebrą toralną, $\Phi$ zbiorem korzeni, a $\Delta=(\alpha_1,\ldots,\alpha_r)$ bazą układu korzeni w realizacji z Twierdzenia~\ref{aa:tw:lie-root-realization}. Dla każdego $i$ wybierzmy $0\ne e_i\in\mathfrak g_{\alpha_i}$ i $f_i\in\mathfrak g_{-\alpha_i}$ spełniający $[e_i,f_i]=h_{\alpha_i}$. Podalgebra $\mathfrak n^+$ z Twierdzenia~\ref{aa:tw:lie-triangular-decomposition} jest generowana przez $e_1,\ldots,e_r$, podalgebra $\mathfrak n^-$ przez $f_1,\ldots,f_r$, a $\mathfrak g$ przez wszystkie $e_i,f_i$.
\end{twierdzenie}
\begin{proof}
Niech $\mathfrak u$ będzie podalgebrą generowaną przez $e_i$. Zawiera się w $\mathfrak n^+$. Indukujemy po wysokości dodatniego korzenia $\beta$. Dla wysokości jeden jest on prosty, więc $\mathfrak g_\beta=Ke_i\subset\mathfrak u$. Przy większej wysokości Lemat~\ref{aa:lem:lie-root-descent} daje $\gamma=\beta-\alpha_i\in\Phi^+$ o mniejszej wysokości. Z indukcji $\mathfrak g_\gamma\subset\mathfrak u$. Wniosek~\ref{aa:cor:lie-root-reflection} daje
$[\mathfrak g_{\alpha_i},\mathfrak g_\gamma]=\mathfrak g_\beta$, stąd $\mathfrak g_\beta\subset\mathfrak u$. Otrzymujemy $\mathfrak u=\mathfrak n^+$.
Dla ujemnych korzeni stosujemy ten sam dowód do bazy $-\Delta$ i generatorów $f_i$. Podalgebra generowana przez wszystkie $e_i,f_i$ zawiera także $h_i=[e_i,f_i]$, a więc $\mathfrak h$ na mocy Lematu~\ref{aa:lem:lie-simple-root-basis}. Rozkład trójkątny daje całą $\mathfrak g$.
\end{proof}

\section{Relacje Serre'a}

\begin{twierdzenie}
\label{aa:tw:lie-serre-necessary}
Niech $\mathfrak g$ będzie niezerową skończenie wymiarową półprostą algebrą Liego nad algebraicznie domkniętym ciałem $K$ charakterystyki zero, $\mathfrak h$ jej maksymalną podalgebrą toralną, $\Phi$ zbiorem korzeni, a $\Delta=(\alpha_1,\ldots,\alpha_r)$ bazą układu korzeni w realizacji z Twierdzenia~\ref{aa:tw:lie-root-realization}. Niech $h_i=h_{\alpha_i}$, $0\ne e_i\in\mathfrak g_{\alpha_i}$, $f_i\in\mathfrak g_{-\alpha_i}$ i $[e_i,f_i]=h_i$. Niech $a_{ij}=\alpha_j(h_i)$ dla $1\le i,j\le r$. Generatory $e_i,f_i,h_i$ spełniają relacje
\[
\begin{aligned}
 [h_i,h_j]&=0, &
 [h_i,e_j]&=a_{ij}e_j,\\
 [h_i,f_j]&=-a_{ij}f_j, &
 [e_i,f_j]&=\delta_{ij}h_i.
\end{aligned}
\]
Ponadto dla $i\ne j$ zachodzą relacje Serre'a
\[
 (\operatorname{ad}_{e_i})^{1-a_{ij}}e_j=0,\qquad
 (\operatorname{ad}_{f_i})^{1-a_{ij}}f_j=0.
\]
W obu relacjach potęga $-a_{ij}$ zamiast $1-a_{ij}$ daje wektor niezerowy.
\end{twierdzenie}
\begin{proof}
Macierz $(a_{ij})$ jest macierzą Cartana w konwencji AA-56, ponieważ $a_{ij}=2(\bar\alpha_i,\bar\alpha_j)/(\bar\alpha_i,\bar\alpha_i)$. W szczególności $a_{ii}=2$, a dla $i\ne j$ mamy $a_{ij}\le0$. Abelowość $\mathfrak h$, definicje przestrzeni korzeniowych i normalizacja trójek dają pierwsze trzy relacje i $[e_i,f_i]=h_i$. Jeśli $i\ne j$, różnica $\alpha_i-\alpha_j$ nie jest korzeniem, gdyż ma współczynniki o różnych znakach w bazie. Nie jest również zerem. Reguła nawiasów daje więc $[e_i,f_j]=0$.

W ciągu korzeniowym przez $\alpha_j$ w kierunku $\alpha_i$ nie ma korzenia $\alpha_j-\alpha_i$, więc $p=0$. Twierdzenie~\ref{aa:tw:lie-root-string} daje $a_{ij}=p-q=-q$. Występują dokładnie korzenie $\alpha_j+k\alpha_i$ dla $0\le k\le -a_{ij}$. Działanie $\operatorname{ad}_{e_i}$ na kolejnych przestrzeniach jest niezerowe i surjektywne poza ostatnią, a na ostatniej znika, bo następna waga nie występuje. Indukcja pokazuje, że $(\operatorname{ad}_{e_i})^k e_j\ne0$ dla $0\le k\le -a_{ij}$, zaś dla $k=1-a_{ij}$ dostajemy zero. Stosując ten argument do bazy $-\Delta$, której macierz Cartana jest taka sama, oraz trójek $(f_i,e_i,-h_i)$, otrzymujemy drugą relację i niezerowość poprzedniej potęgi.
\end{proof}

\begin{uwaga}
Relacje te są na razie warunkami koniecznymi. Jeśli skonstruujemy abstrakcyjną algebrę Liego z generatorami i dokładnie tymi relacjami, Twierdzenie~\ref{aa:tw:lie-simple-generators} zapewni surjektywny homomorfizm na $\mathfrak g$. Wykazanie jego iniektywności i istnienia odpowiedniej algebry dla każdej macierzy Cartana typu skończonego wymaga dalszych argumentów.
\end{uwaga}

\section{Ideały i składowe układu korzeni}

\begin{lemat}
\label{aa:lem:lie-ideal-weight-components}
Niech $\mathfrak g$ będzie skończenie wymiarową półprostą algebrą Liego nad algebraicznie domkniętym ciałem $K$ charakterystyki zero, $\mathfrak h$ jej maksymalną podalgebrą toralną, a $\Phi$ zbiorem korzeni. Jeżeli $\mathfrak a$ jest ideałem $\mathfrak g$, to
\[
 \mathfrak a=(\mathfrak a\cap\mathfrak h)\oplus
 \bigoplus_{\alpha\in\Phi}(\mathfrak a\cap\mathfrak g_\alpha).
\]
\end{lemat}
\begin{proof}
Wybieramy $t\in\mathfrak h$, dla którego wartości $\lambda(t)$ dla $\lambda\in\Phi\cup\{0\}$ są parami różne. Jest to możliwe: zbiór niedozwolony jest sumą mnogościową skończenie wielu jąder niezerowych funkcjonałów $\lambda-\mu$. Iloczyn ich równań liniowych jest niezerowym wielomianem, który nad nieskończonym ciałem nie może znikać wszędzie. Operator $\operatorname{ad}_t$ ma więc przestrzenie własne dokładnie równe składowym rozkładu korzeniowego. Ideał $\mathfrak a$ jest niezmienniczy względem tego operatora, a zatem względem każdej projekcji wielomianowej
\[
 P_\lambda=\prod_{\mu\in\Phi\cup\{0\},\ \mu\ne\lambda}
 \frac{\operatorname{ad}_t-\mu(t)I}{\lambda(t)-\mu(t)}.
\]
Na przestrzeni wagi $\lambda$ projekcja jest identycznością, a na pozostałych jest zerem. Składowe każdego elementu $\mathfrak a$ należą więc do $\mathfrak a$; odwrotna inkluzja wynika z liniowości ideału.
\end{proof}

\begin{twierdzenie}[Składowe proste]
\label{aa:tw:lie-dynkin-components}
Niech $\mathfrak g$ będzie niezerową skończenie wymiarową półprostą algebrą Liego nad algebraicznie domkniętym ciałem $K$ charakterystyki zero, $\mathfrak h$ jej maksymalną podalgebrą toralną, $\Phi$ zbiorem korzeni, a $\Delta=(\alpha_1,\ldots,\alpha_r)$ bazą układu korzeni w realizacji z Twierdzenia~\ref{aa:tw:lie-root-realization}. Niech $h_i=h_{\alpha_i}$, a $I_1,\ldots,I_s$ będą zbiorami indeksów składowych spójnych diagramu Dynkina bazy $\Delta$. Niech $\Phi_\ell$ będzie zbiorem korzeni, których współczynniki poza $I_\ell$ są zerowe, a $\mathfrak h_\ell=\operatorname{span}_K\{h_i:i\in I_\ell\}$. Wówczas
\[
 \mathfrak g_\ell=\mathfrak h_\ell\oplus
 \bigoplus_{\alpha\in\Phi_\ell}\mathfrak g_\alpha
\]
są niezerowymi prostymi ideałami, komutującymi parami, oraz $\mathfrak g=\bigoplus_{\ell=1}^s\mathfrak g_\ell$. W szczególności $\mathfrak g$ jest prosta wtedy i tylko wtedy, gdy jej układ korzeni jest nierozkładalny.
\end{twierdzenie}
\begin{proof}
Z dowodu Twierdzenia~\ref{aa:tw:root-cartan} korzenie należą do dokładnie jednej z przestrzeni rozpiętych przez proste korzenie poszczególnych składowych. Przestrzenie te są wzajemnie ortogonalne w modelu rzeczywistym. Zatem $\Phi$ jest rozłączną sumą $\Phi_\ell$, a Lemat~\ref{aa:lem:lie-simple-root-basis} daje $\mathfrak h=\bigoplus_\ell\mathfrak h_\ell$. Stąd rozkład $\mathfrak g$ na podprzestrzenie $\mathfrak g_\ell$ jest prosty.

Dla $\alpha\in\Phi_\ell$, $\beta\in\Phi_k$, $k\ne\ell$, suma $\alpha+\beta$ nie może być korzeniem: miałaby niezerowe składowe w dwóch różnych częściach bazy. Nie może też być zerem. Zatem $[\mathfrak g_\alpha,\mathfrak g_\beta]=0$. Ponadto $\beta(h_i)=0$ dla $i\in I_\ell$ dzięki ortogonalności, więc $[\mathfrak h_\ell,\mathfrak g_\beta]=0$. Abelowość $\mathfrak h$ daje komutowanie różnych $\mathfrak g_\ell$.

Sprawdzimy zamknięcie każdej $\mathfrak g_\ell$ na nawias. Dla korzeni z tej samej części ich suma, jeśli jest niezerowym korzeniem, należy do tej części. Jeśli suma jest zerem, wzór z dowodu Twierdzenia~\ref{aa:tw:root-sl2-triple} daje $[x,y]=\kappa(x,y)t_\alpha$ dla $x\in\mathfrak g_\alpha$, $y\in\mathfrak g_{-\alpha}$. Ponieważ $\alpha=\sum_{i\in I_\ell}n_i\alpha_i$, liniowość izomorfizmu wyznaczonego przez formę Killinga daje $t_\alpha=\sum_{i\in I_\ell}n_i t_{\alpha_i}\in\mathfrak h_\ell$. Nawiasy z $\mathfrak h_\ell$ zachowują przestrzenie korzeniowe, a sama $\mathfrak h_\ell$ jest abelowa. To dowodzi zamknięcia. Wraz z komutowaniem innych składników oznacza to, że $\mathfrak g_\ell$ jest ideałem. Nie jest abelowy, gdyż dla $i\in I_\ell$ zawiera trójkę z $[e_i,f_i]=h_i\ne0$.

Niech $\mathfrak a\ne0$ będzie ideałem $\mathfrak g_\ell$. Dzięki komutowaniu pozostałych składników jest ideałem $\mathfrak g$. Z Lematu~\ref{aa:lem:lie-ideal-weight-components} albo zawiera niezerowy element pewnej przestrzeni korzeniowej, albo jest zawarty w $\mathfrak h_\ell$. W drugim przypadku wybierzmy $0\ne t\in\mathfrak a$. Korzenie rozpinają $\mathfrak h^*$, więc istnieje $\beta$ z $\beta(t)\ne0$; ortogonalność części pokazuje, że $\beta\in\Phi_\ell$. Dla $0\ne x\in\mathfrak g_\beta$ mamy $[t,x]=\beta(t)x\in\mathfrak a$, zatem w obu przypadkach $\mathfrak a$ zawiera całą pewną przestrzeń korzeniową, bo jest jednowymiarowa.

Jeżeli zawarta przestrzeń odpowiada korzeniowi ujemnemu, wybieramy $0\ne x\in\mathfrak g_\beta\subset\mathfrak a$ i $y\in\mathfrak g_{-\beta}$ z $\kappa(x,y)\ne0$, korzystając z niezdegenerowanego parowania. Wtedy $[x,y]=\kappa(x,y)t_\beta\in\mathfrak a$, więc $h_\beta\in\mathfrak a$. Ponieważ $[h_\beta,y]=-2y$, cała prosta $\mathfrak g_{-\beta}$ należy do $\mathfrak a$. Możemy więc zacząć od dodatniego korzenia $\beta$. Jeśli $\beta$ nie jest prosty, Lemat~\ref{aa:lem:lie-root-descent} daje $\beta-\alpha_i\in\Phi^+$; z Wniosku~\ref{aa:cor:lie-root-reflection} nawias $[\mathfrak g_\beta,\mathfrak g_{-\alpha_i}]$ jest całą przestrzenią $\mathfrak g_{\beta-\alpha_i}$, zawartą w $\mathfrak a$. Po skończonej liczbie kroków otrzymujemy $e_i\in\mathfrak a$ dla pewnego $i\in I_\ell$. Wtedy $h_i=[e_i,f_i]\in\mathfrak a$ i $[h_i,f_i]=-2f_i$ daje $f_i\in\mathfrak a$.

Jeśli $i,j$ są sąsiednimi wierzchołkami diagramu, to $a_{ij}\ne0$. Nawiasy $[h_i,e_j]=a_{ij}e_j$ i $[h_i,f_j]=-a_{ij}f_j$ dają $e_j,f_j\in\mathfrak a$, a ich nawias daje $h_j$. Spójność pozwala przejść do wszystkich indeksów $I_\ell$. Dowód Twierdzenia~\ref{aa:tw:lie-simple-generators} zastosowany do korzeni tej części (odejmowanie prostego korzenia nie wychodzi poza część) daje wszystkie jej przestrzenie korzeniowe. Stąd $\mathfrak a=\mathfrak g_\ell$. Ostatnia teza wynika z równoważności spójności diagramu i nierozkładalności układu z AA-56.
\end{proof}

\section{Przykład typu A}
\begin{przyklad}
\label{aa:prz:lie-serre-sl3}
Dla $\mathfrak{sl}_3(K)$ wybieramy $\alpha_1=\alpha_{12}$, $\alpha_2=\alpha_{23}$ z Przykładu~\ref{aa:prz:sl3-root-decomposition}. Generatory to
$e_1=E_{12}$, $e_2=E_{23}$, $f_1=E_{21}$, $f_2=E_{32}$ i $h_1=E_{11}-E_{22}$, $h_2=E_{22}-E_{33}$. Macierz Cartana jest równa
$\left(\begin{smallmatrix}2&-1\\-1&2\end{smallmatrix}\right)$.
Dodatnia podalgebra składa się z macierzy ściśle górnotrójkątnych, ujemna z macierzy ściśle dolnotrójkątnych, a $\mathfrak h$ z macierzy diagonalnych o śladzie zero. Z wzoru $E_{ij}E_{kl}=\delta_{jk}E_{il}$ mamy
\[
 [e_1,e_2]=E_{13},\quad [f_1,f_2]=-E_{31},\quad
 [e_1,E_{13}]=[e_2,E_{13}]=0.
\]
Analogicznie $[f_1,E_{31}]=[f_2,E_{31}]=0$. Relacje Serre'a mają zatem wykładnik dwa i nie można go zmniejszyć do jednego. Dodatnia i ujemna podalgebra mają klasę nilpotentności dwa, gdyż ich nawiasy rozpinają odpowiednio $KE_{13}$ i $KE_{31}$, a te przestrzenie komutują ze swoimi podalgebrami. Diagram $A_2$ jest spójny, więc Twierdzenie~\ref{aa:tw:lie-dynkin-components} daje prostotę $\mathfrak{sl}_3(K)$.
\end{przyklad}

Rozdział przygotowuje dowód twierdzenia o prezentacji Serre'a: znamy generatory i wszystkie wskazane relacje w istniejącej algebrze, lecz nie wykazaliśmy jeszcze, że relacje te wyczerpują zależności między generatorami.

% AA-62
\chapter[Prezentacja Serre'a]{Twierdzenie o prezentacji Serre'a}
\label{aa:ch:serre-presentation}

W poprzednim rozdziale wyprowadziliśmy relacje spełniane przez generatory istniejącej półprostej algebry Liego. Teraz zmieniamy kierunek rozumowania: punktem wyjścia jest skończony zredukowany krystalograficzny układ korzeni. Zbudujemy algebrę za pomocą samych generatorów i relacji, wykażemy jej skończony wymiar oraz półprostotę. W szczególności dowód istnienia nie będzie zakładał wcześniejszego istnienia algebry danego typu.

\section{Algebra zadana generatorami i relacjami}

Niech $K$ będzie ciałem algebraicznie domkniętym charakterystyki zero, $\Phi$ niepustym skończonym zredukowanym krystalograficznym układem korzeni w rzeczywistej przestrzeni euklidesowej $V$ oraz $\Delta=(\alpha_1,\ldots,\alpha_r)$ jego bazą. Przyjmujemy konwencję
\[
 a_{ij}=\frac{2(\alpha_i,\alpha_j)}{(\alpha_i,\alpha_i)}.
\]
Macierz $A=(a_{ij})$ jest odwracalna, gdyż po pomnożeniu z lewej przez diagonalną macierz o wyrazach $(\alpha_i,\alpha_i)/2$ otrzymujemy dodatnio określoną macierz Grama. Ponieważ $A$ ma współczynniki całkowite, jest odwracalna także nad $K$.

\begin{definicja}
\label{aa:def:serre-algebra}
Algebrą $\mathfrak s(A)$ nazywamy iloraz wolnej algebry Liego nad $K$ o generatorach $e_i,f_i,h_i$ ($1\le i\le r$) przez ideał generowany przez relacje
\[
\begin{aligned}
 [h_i,h_j]&=0,&[h_i,e_j]&=a_{ij}e_j,\\
 [h_i,f_j]&=-a_{ij}f_j,&[e_i,f_j]&=\delta_{ij}h_i
\end{aligned}
\]
oraz, dla $i\ne j$, przez
\[
 (\operatorname{ad}e_i)^{1-a_{ij}}e_j=0,\qquad
 (\operatorname{ad}f_i)^{1-a_{ij}}f_j=0.
\]
Przez te same symbole oznaczamy obrazy generatorów w ilorazie.
\end{definicja}

Wolną algebrę Liego można zbudować z przestrzeni liniowej formalnych słów nawiasowych przez iloraz wyznaczony przez antysymetrię i tożsamość Jacobiego. Każde przypisanie generatorom elementów innej algebry Liego przedłuża się do homomorfizmu dokładnie wtedy, gdy spełnia wszystkie zadane relacje. Na tym etapie nie wiemy nawet, czy obrazy $h_i$ są liniowo niezależne.

Oznaczmy przez $\mathfrak u^+$ i $\mathfrak u^-$ podalgebry generowane odpowiednio przez $e_i$ i $f_i$, a przez $H$ przestrzeń rozpiętą przez $h_i$. Niech $Q=\bigoplus_i\mathbb Z\alpha_i$ będzie kratą korzeniową. Nadajemy generatorom stopnie $\alpha_i,-\alpha_i,0$ odpowiednio dla $e_i,f_i,h_i$. Wszystkie relacje są jednorodne, więc iloraz ma gradację przez $Q$.

\begin{lemat}[Rozkład trójkątny przed dowodem istnienia]
\label{aa:lem:serre-triangular}
Dla algebry $\mathfrak s(A)$ z Definicji~\ref{aa:def:serre-algebra} zachodzi
\[
 \mathfrak s(A)=\mathfrak u^-\oplus H\oplus\mathfrak u^+.
\]
Przestrzeń stopnia zero jest równa $H$, a każdy niezerowy stopień ma współczynniki w bazie $\Delta$ albo wszystkie nieujemne, albo wszystkie niedodatnie.
\end{lemat}
\begin{proof}
Każda algebra Liego generowana przez zbiór $X$ jest rozpięta przez elementy $X$ oraz słowa postaci $[x_1,[x_2,\ldots,[x_{m-1},x_m]\ldots]]$, gdzie $x_i\in X$. Wynika to przez indukcję po długości pierwszego argumentu nawiasu, z tożsamości $[[x,y],z]=[x,[y,z]]-[y,[x,z]]$.

Dla każdego $i$ mamy $[e_i,\mathfrak u^-]\subset H+\mathfrak u^-$. Istotnie, dla słowa długości jeden jest to relacja $[e_i,f_j]=\delta_{ij}h_i$. Dla słowa $[f_j,w]$ stosujemy
\[
 [e_i,[f_j,w]]=[[e_i,f_j],w]+[f_j,[e_i,w]].
\]
Pierwszy składnik należy do $\mathfrak u^-$, ponieważ $[H,\mathfrak u^-]\subset\mathfrak u^-$ (co wynika indukcyjnie z relacji dla $h_i$). Drugi też należy do $\mathfrak u^-$ na mocy założenia indukcyjnego. Zamieniając $e_i$ z $f_i$ oraz $h_i$ z $-h_i$, otrzymujemy $[f_i,\mathfrak u^+]\subset H+\mathfrak u^+$.

Przestrzeń $T=\mathfrak u^-+H+\mathfrak u^+$ jest więc niezmiennicza względem nawiasu z każdym generatorem i zawiera wszystkie generatory. Z opisu słów nawiasowych wynika $T=\mathfrak s(A)$. Stopnie występujące w trzech składnikach są parami rozłączne: ujemne, zerowy oraz dodatnie. Suma jest zatem prosta.
\end{proof}

\section{Odbicia i ograniczenie wymiaru}

\begin{lemat}[Automorfizmy odpowiadające odbiciom]
\label{aa:lem:serre-reflections}
Niech $\mathfrak s(A)$ będzie algebrą z Definicji~\ref{aa:def:serre-algebra}. Operatory $\operatorname{ad}e_i$ i $\operatorname{ad}f_i$ są lokalnie nilpotentne: każdy element algebry jest zerowany przez pewną potęgę danego operatora. Automorfizm
\[
 R_i=\exp(\operatorname{ad}e_i)\,
      \exp(-\operatorname{ad}f_i)\,
      \exp(\operatorname{ad}e_i)
\]
przenosi składową stopnia $\beta\in Q$ na składową stopnia $s_i\beta$, gdzie $s_i$ jest odbiciem względem hiperpłaszczyzny prostopadłej do $\alpha_i$.
\end{lemat}
\begin{proof}
Lokalna nilpotentność na generatorach wynika z relacji: działanie na $h_j$ kończy się po dwóch krokach, na $f_i$ po trzech, na $f_j$ dla $j\ne i$ po jednym, a na $e_j$ dla $j\ne i$ z relacji Serre'a. Dla derywacji $D$ zachodzi wzór
\[
 D^m[x,y]=\sum_{k=0}^m\binom{m}{k}[D^kx,D^{m-k}y].
\]
Jeśli $D^p x=D^q y=0$, to $D^{p+q-1}[x,y]=0$. Indukcja po słowach nawiasowych i skończoność sum dowodzą tezy na całej algebrze; dla $f_i$ argument jest analogiczny. Szereg $\exp(D)x=\sum_{m\ge0}D^mx/m!$ jest zatem dla każdego $x$ skończony. Ten sam wzór dwumianowy daje $\exp(D)[x,y]=[\exp(D)x,\exp(D)y]$, zaś mnożenie skończonych szeregów na każdym elemencie daje odwrotność $\exp(-D)$.

Dla ustalonego $i$ bezpośrednio z relacji otrzymujemy
$\exp(\operatorname{ad}e_i)h_i=h_i-2e_i$,
$\exp(-\operatorname{ad}f_i)h_i=h_i-2f_i$ oraz
$\exp(-\operatorname{ad}f_i)e_i=e_i+h_i-f_i$.
Po złożeniu trzech odwzorowań wynika $R_i h_i=-h_i$. Element $h_j-(a_{ji}/2)h_i$ komutuje z $e_i,f_i$ i jest ustalony przez $R_i$. Zatem
\[
 R_i h_j=h_j-a_{ji}h_i,\qquad R_i^{-1}h_j=R_i h_j.
\]
Jeżeli $x$ ma stopień $\beta=\sum_k n_k\alpha_k$, to
$[h_j,x]=b_j x$, gdzie $b_j=\sum_k a_{jk}n_k$; wynika to indukcyjnie z relacji i tożsamości Jacobiego. W konsekwencji
\[
 [h_j,R_ix]=R_i[R_i^{-1}h_j,x]
 =(b_j-a_{ji}b_i)R_ix.
\]
Współrzędne $s_i\beta=\beta-b_i\alpha_i$ dają właśnie te wartości. Różne stopnie mają różne rodziny wartości własnych, ponieważ $A$ jest odwracalna nad $K$, a charakterystyka jest zero. Rozkładając $R_ix$ na skończenie wiele jednorodnych składowych, widzimy, że wszystkie poza stopniem $s_i\beta$ muszą znikać. Odwrotność $R_i^{-1}$ daje przeciwne zawieranie.
\end{proof}

\begin{twierdzenie}[Górne ograniczenie wymiaru]
\label{aa:tw:serre-upper-dimension}
Dla algebry z Definicji~\ref{aa:def:serre-algebra} niezerowe składowe o stopniu różnym od zera mogą wystąpić jedynie w stopniach należących do $\Phi$, a każda ma wymiar co najwyżej jeden. W szczególności
\[
 \dim_K\mathfrak s(A)\le r+|\Phi|.
\]
\end{twierdzenie}
\begin{proof}
Niech $\beta=\sum_i n_i\alpha_i$ będzie dodatnim stopniem o niezerowej składowej. Jeżeli jego nośnik składa się tylko z indeksu $i$, to wszystkie słowa tego stopnia są słowami w jednym generatorze $e_i$. Każde takie słowo długości większej niż jeden jest zerem, ponieważ $[e_i,e_i]=0$. Zatem jedyny możliwy stopień o jednoelementowym nośniku to $\alpha_i$, a jego składowa ma wymiar co najwyżej jeden.

Gdy nośnik ma co najmniej dwa elementy, z
$(\beta,\beta)=\sum_i n_i(\beta,\alpha_i)>0$
wynika istnienie $i$ z $n_i>0$ i $(\beta,\alpha_i)>0$. Liczba
$b_i=2(\beta,\alpha_i)/(\alpha_i,\alpha_i)$ jest dodatnią liczbą całkowitą. Stopień $s_i\beta=\beta-b_i\alpha_i$ ma niezerową składową na mocy Lematu~\ref{aa:lem:serre-reflections}. Współczynnik przy pewnym innym korzeniu prostym pozostaje dodatni. Z Lematu~\ref{aa:lem:serre-triangular} wszystkie współczynniki nowego stopnia są więc nieujemne. Jego wysokość zmniejszyła się o $b_i\ge1$.

Po skończonej liczbie takich kroków dochodzimy do stopnia prostego. Odbicia zachowują $\Phi$, więc wyjściowy stopień należy do $\Phi$, a jego składowa ma ten sam wymiar co składowa stopnia prostego. Ujemne stopnie badamy przez zamianę generatorów $e_i\leftrightarrow f_i$, $h_i\mapsto-h_i$, która zachowuje relacje i odwraca stopnie. Składowa zerowa ma wymiar co najwyżej $r$.
\end{proof}

\section{Konstrukcja modułów i niezależność generatorów}

Do wykazania dolnego ograniczenia wymiaru nie możemy użyć istnienia półprostej algebry Liego o zadanym układzie korzeni. Skonstruujemy moduły bezpośrednio z relacji.

\begin{twierdzenie}[Moduły z wektorem najwyższej wagi]
\label{aa:tw:serre-highest-module}
Niech $\mathfrak s(A)$ będzie algebrą z Definicji~\ref{aa:def:serre-algebra}. Dla każdej rodziny $\lambda=(\lambda_1,\ldots,\lambda_r)\in K^r$ istnieje jej moduł $M_\lambda$ z niezerowym wektorem $v_\lambda$ spełniającym
\[
 e_i v_\lambda=0,\quad h_i v_\lambda=\lambda_i v_\lambda,\quad
 f_i v_\lambda\ne0\qquad(1\le i\le r).
\]
\end{twierdzenie}
\begin{proof}
Niech $T=K\langle x_1,\ldots,x_r\rangle$ będzie wolną łączną algebrą z jedynką. Oznaczmy przez $F_i$ mnożenie z lewej przez $x_i$. Dla jednorodnego słowa $w$ z $n_j$ wystąpieniami $x_j$ definiujemy
\[
 H_iw=\left(\lambda_i-\sum_j a_{ij}n_j\right)w.
\]
Definiujemy liniowo operator $E_i$ rekurencją $E_i(1)=0$ oraz
\[
 E_i(x_jw)=x_jE_i(w)+\delta_{ij}H_iw.
\]
Każde niepuste słowo ma jednoznaczną pierwszą literę, więc rekurencja jest dobrze określona. Operator $E_i$ usuwa jedno wystąpienie $x_i$ (lub daje zero), co sprawdza się indukcyjnie. Wobec tego na $T$ zachodzą
\[
 [H_i,H_j]=0,\quad [H_i,F_j]=-a_{ij}F_j,\quad
 [H_i,E_j]=a_{ij}E_j,\quad [E_i,F_j]=\delta_{ij}H_i.
\]

Niech $J$ będzie dwustronnym ideałem $T$ generowanym przez
$S_{jk}=(\operatorname{ad}x_j)^{1-a_{jk}}x_k$ dla $j\ne k$, gdzie nawias w $T$ jest komutatorem. Jest to ideał jednorodny względem liczby wystąpień każdej litery, generowany w stopniach długości co najmniej dwa. Iloraz $T/J$ ma niezerową jedynkę oraz niezerowe obrazy wszystkich $x_i$.

Sprawdzimy, że $J$ jest niezmienniczy względem $E_i$. Oznaczmy $Y_m=(\operatorname{ad}F_j)^m F_k$ dla $j\ne k$. Wtedy $Y_m$ jest mnożeniem z lewej przez $(\operatorname{ad}x_j)^m x_k$. Mamy
$[H_j,Y_m]=(-a_{jk}-2m)Y_m$.
Dla $i=j$ indukcja daje
\[
 [E_j,Y_m]=m(-a_{jk}-m+1)Y_{m-1}\quad(m\ge1).
\]
Rzeczywiście dla $m=1$ otrzymujemy $[H_j,F_k]=-a_{jk}F_k$. W kroku indukcyjnym, z $Y_{m+1}=[F_j,Y_m]$, współczynnik wynosi
$-a_{jk}-2m+m(-a_{jk}-m+1)=(m+1)(-a_{jk}-m)$.
Dla $m=1-a_{jk}$ komutator znika. Jeśli $i$ jest różny od $j,k$, znika na mocy komutowania z $F_j,F_k$. Jeśli $i=k$, to $[E_k,Y_m]=(\operatorname{ad}F_j)^m H_k$. Dla $m\ge2$ wyrażenie jest zerem, ponieważ $[F_j,H_k]=a_{kj}F_j$. Gdy $m=1-a_{jk}=1$, mamy $a_{jk}=0$, a więc także $a_{kj}=0$. Ostatecznie $E_i$ komutuje z mnożeniem z lewej przez każdy $S_{jk}$.

Ideał $J$ jest rozpięty przez elementy $u S_{jk} w$ dla słów $u,w$. Jeżeli $u=1$, to
$E_i(S_{jk}w)=S_{jk}E_i(w)\in J$. Jeśli $u=x_lu'$, to rekurencja i niezmienniczość $J$ względem $H_i$ dają
\[
 E_i(x_lu'S_{jk}w)=x_l E_i(u'S_{jk}w)
                +\delta_{il}H_i(u'S_{jk}w)\in J.
\]
Zatem $E_i$, podobnie jak $H_i,F_i$, indukuje operator na $M_\lambda=T/J$.

Na ilorazie relacje Serre'a dla $F_i$ zachodzą z definicji $J$. Pozostają relacje dla $E_i$. Niech $Z=(\operatorname{ad}E_j)^{1-a_{jk}}E_k$. Ten sam rachunek komutatorów co wyżej, po zamianie $E_i$ z $F_i$ i $H_i$ z $-H_i$, daje $[F_i,Z]=0$ dla każdego $i$. Rachunek korzystał wyłącznie z czterech rodzin relacji liniowych, nie z relacji Serre'a. Ponadto $Z$ zeruje $1+J$, bo każdy iloczyn dodatniej liczby operatorów $E_i$ zeruje ten wektor. Wektory $F_{i_1}\cdots F_{i_l}(1+J)$ rozpinają $M_\lambda$. Komutowanie z $F_i$ pokazuje więc $Z=0$ na całym module. Wszystkie relacje z Definicji~\ref{aa:def:serre-algebra} są spełnione. Wystarczy przyjąć $v_\lambda=1+J$.
\end{proof}

\begin{wniosek}
\label{aa:cor:serre-dimension}
W algebrze $\mathfrak s(A)$ z Definicji~\ref{aa:def:serre-algebra} elementy $h_i$ są liniowo niezależne, wszystkie $e_i,f_i$ są niezerowe oraz
\[
 \dim_K\mathfrak s(A)=r+|\Phi|,
 \qquad \dim_K\mathfrak s(A)_\beta=1\quad(\beta\in\Phi).
\]
\end{wniosek}
\begin{proof}
Jeśli $\sum_i c_i h_i=0$, działanie na $v_\lambda$ daje $\sum_i c_i\lambda_i=0$ dla każdej $\lambda\in K^r$. Wybierając wektory jednostkowe, otrzymujemy $c_i=0$ dla każdego $i$. Elementy $f_i$ działają niezerowo na $v_\lambda$, a $e_i\ne0$ wynika z $[e_i,f_i]=h_i\ne0$. Każdy korzeń należy do orbity pewnego korzenia prostego względem grupy generowanej przez odbicia proste, na mocy Twierdzenia~\ref{aa:tw:weyl-chambers}. Stosując automorfizmy z Lematu~\ref{aa:lem:serre-reflections}, otrzymujemy niezerowe składowe dla wszystkich $\beta\in\Phi$. Twierdzenie~\ref{aa:tw:serre-upper-dimension} daje ich jednowymiarowość i wymaganą równość wymiarów.
\end{proof}

\section{Półprostota i twierdzenie o prezentacji}

\begin{twierdzenie}[Istnienie algebry o zadanym układzie korzeni]
\label{aa:tw:serre-existence}
Niech $K$ będzie algebraicznie domkniętym ciałem charakterystyki zero, $\Phi$ niepustym skończonym zredukowanym krystalograficznym układem korzeni, a $A$ macierzą Cartana jego bazy. Algebra $\mathfrak s(A)$ jest skończenie wymiarowa i półprosta. Podalgebra $H=\operatorname{span}_K\{h_1,\ldots,h_r\}$ jest maksymalną podalgebrą toralną, a układ korzeni względem $H$ odpowiada $\Phi$ przez przypisanie $\alpha_j(h_i)=a_{ij}$.
\end{twierdzenie}
\begin{proof}
Skończony wymiar wykazano w Twierdzeniu~\ref{aa:tw:serre-upper-dimension}. Z Wniosku~\ref{aa:cor:serre-dimension} elementy $h_i$ są liniowo niezależne. Działanie $\operatorname{ad}H$ jest diagonalne w rozkładzie na stopnie, a $H$ jest abelowa. Jeśli element komutuje z $H$, to wszystkie jego składowe o niezerowym stopniu muszą znikać, ponieważ macierz $A$ jest odwracalna. Centralizator $H$ jest więc równy $H$; żadna większa abelowa podalgebra nie może go zawierać. Korzenie w sensie rozkładu wagowego są dokładnie funkcjonałami odpowiadającymi stopniom $\beta\in\Phi$.

Niech $\mathfrak a$ będzie niezerowym rozwiązalnym ideałem $\mathfrak s(A)$. Wybierzmy $t\in H$ rozdzielający wszystkie wagi, łącznie z zerem. Istnieje on, ponieważ w przestrzeni $H$ nad nieskończonym ciałem skończona suma mnogościowa właściwych hiperpłaszczyzn nie pokrywa całej przestrzeni. Projekcje na składowe stopni są wielomianami w $\operatorname{ad}t$, jak w Lemacie~\ref{aa:lem:lie-ideal-weight-components}. Zatem ideał zawiera składowe każdego swojego elementu.

Jeśli zawiera niezerowy element składowej o stopniu $\beta\ne0$, to zawiera tę całą składową. Ideał jest niezmienniczy względem $\operatorname{ad}e_i,\operatorname{ad}f_i$, więc także względem ich skończonych wykładników i automorfizmów $R_i$ oraz odwrotnych. Możemy przenieść składową do stopnia prostego, otrzymując $e_i\in\mathfrak a$. Wówczas $[e_i,f_i]=h_i\in\mathfrak a$ i $[h_i,f_i]=-2f_i\in\mathfrak a$. Elementy $e_i,f_i,h_i$ rozpinają podalgebrę izomorficzną z $\mathfrak{sl}_2(K)$: relacje są standardowe, a niezależność wynika z różnych stopni i niezerowości. Jej algebra pochodna jest nią samą, co przeczy rozwiązalności $\mathfrak a$.

W przeciwnym przypadku $\mathfrak a\subset H$. Dla $0\ne t\in\mathfrak a$ odwracalność $A$ i niezależność $h_i$ dają indeks $j$ z $\alpha_j(t)\ne0$. Wtedy $[t,e_j]=\alpha_j(t)e_j\in\mathfrak a$, sprzeczność. Nie ma więc niezerowego rozwiązalnego ideału, czyli $\mathfrak s(A)$ jest półprosta.

Sprawdźmy jeszcze zgodność iloczynów skalarnych w realizacji korzeni. Dla pierwotnego układu $\Phi$ wartości $\beta(h_i)$ są liczbami całkowitymi wyznaczonymi przez $A$ i współrzędne $\beta$. Forma Killinga na $H$ jest zatem określona macierzą
$B_{ij}=\sum_{\beta\in\Phi}\beta(h_i)\beta(h_j)$.
W modelu rzeczywistym forma $B$ jest dodatnio określona, gdyż korzenie rozpinają przestrzeń. Jest niezmiennicza względem odbić, ponieważ odbicia permutują korzenie. W przestrzeni dualnej forma o macierzy $B^{-1}$ jest również niezmiennicza względem odbić. Odbicie $s_i$ ma kierunek zmiany $\alpha_i$ i hiperpłaszczyznę punktów stałych $\{\beta:\beta(h_i)=0\}$. Z ortogonalności względem tej formy otrzymujemy dla każdego $\beta$ stosunek Cartana równy $\beta(h_i)$. Realizacja euklidesowa z Twierdzenia~\ref{aa:tw:lie-root-realization} daje więc tę samą macierz Cartana $A$. Odwzorowanie liniowe przypisujące pierwotnym korzeniom prostym odpowiadające im wagi jest odwracalne, ponieważ $A$ jest odwracalna. Przenosi ono zbiór korzeni na zbiór wag niezerowych, jak już wykazaliśmy, i przeplata odbicia proste. Ponieważ każdy korzeń należy do orbity korzenia prostego, przeplata również odbicia odpowiadające wszystkim korzeniom, a tym samym zachowuje wszystkie liczby Cartana. Na każdej składowej nierozkładalnej można też zastosować Twierdzenie~\ref{aa:tw:root-uniqueness}: iloczyny skalarne różnią się jedynie dodatnim mnożnikiem, który może być różny dla różnych składowych.
\end{proof}

\begin{twierdzenie}[Prezentacja Serre'a]
\label{aa:tw:serre-presentation}
Niech $\mathfrak g$ będzie niezerową skończenie wymiarową półprostą algebrą Liego nad algebraicznie domkniętym ciałem $K$ charakterystyki zero, $\mathfrak h$ jej maksymalną podalgebrą toralną, $\Delta$ bazą układu korzeni, a $A$ odpowiadającą macierzą Cartana. Niech $e_i,f_i,h_i$ będą generatorami z Twierdzenia~\ref{aa:tw:lie-serre-necessary}. Istnieje dokładnie jeden izomorfizm $\mathfrak s(A)\to\mathfrak g$ przypisujący formalnym generatorom odpowiadające im elementy $\mathfrak g$.
\end{twierdzenie}
\begin{proof}
Relacje z AA-61 dają homomorfizm na mocy własności ilorazu wolnej algebry Liego. Jest on surjektywny, ponieważ obrazy $e_i,f_i$ generują $\mathfrak g$ (Twierdzenie~\ref{aa:tw:lie-simple-generators}). Rozkład korzeniowy w $\mathfrak g$ ma $r$ wymiarów w stopniu zero i po jednym wymiarze dla każdego korzenia. Stąd
$\dim\mathfrak g=r+|\Phi|=\dim\mathfrak s(A)$
na mocy Wniosku~\ref{aa:cor:serre-dimension}. Surjektywny homomorfizm jest zatem izomorfizmem. Jednoznaczność wynika z generowania.
\end{proof}

\section{Konsekwencje dla klasyfikacji}

\begin{wniosek}
\label{aa:cor:serre-classification-step}
Dla każdego niepustego skończonego zredukowanego krystalograficznego układu korzeni istnieje półprosta algebra Liego nad $K$ realizująca ten układ. Jeżeli dwie półproste algebry Liego nad $K$ mają, względem wybranych maksymalnych podalgebr toralnych i baz układów, tę samą macierz Cartana z dokładnością do jednoczesnej permutacji wierszy i kolumn, to są izomorficzne.
\end{wniosek}
\begin{proof}
Pierwsza teza wynika z Twierdzenia~\ref{aa:tw:serre-existence}. W drugiej po zmianie numeracji obie algebry są izomorficzne z tą samą $\mathfrak s(A)$ na mocy Twierdzenia~\ref{aa:tw:serre-presentation}.
\end{proof}

Dla nierozkładalnych układów z AA-56 uzyskujemy algebry proste, na mocy Twierdzenia~\ref{aa:tw:lie-dynkin-components}. Aby otrzymać pełną klasyfikację niezależną od wyboru maksymalnej podalgebry toralnej, pozostaje wykazać sprzężoność takich podalgebr. Nie używamy jej w niniejszym dowodzie i nie zakładamy, że została już udowodniona.

\begin{przyklad}
\label{aa:prz:serre-rank-two}
Dla macierzy $\left(\begin{smallmatrix}2&-3\\-1&2\end{smallmatrix}\right)$ typu $G_2$ niech $\alpha$ będzie krótszym, a $\beta$ dłuższym korzeniem prostym. Dodatnie korzenie to $\alpha,\beta,\alpha+\beta,2\alpha+\beta,3\alpha+\beta,3\alpha+2\beta$. Wniosek~\ref{aa:cor:serre-dimension} daje wymiar $2+12=14$. Relacje dodatnie to
\[
 (\operatorname{ad}e_1)^4e_2=0,\qquad
 (\operatorname{ad}e_2)^2e_1=0.
\]
W algebrze istniejącej na mocy Twierdzenia~\ref{aa:tw:serre-existence} poprzednie potęgi nie znikają, zgodnie z AA-61. Rozkład trójkątny ma wymiary $6+2+6$.
\end{przyklad}

\begin{center}
\begin{tikzpicture}[every node/.style={font=\small}]
\node[draw,rounded corners,align=center] (a) at (0,0) {układ korzeni $\Phi$\\macierz $A$};
\node[draw,rounded corners,align=center] (b) at (0,-1.5) {generatory i relacje\\$\mathfrak s(A)$};
\node[draw,rounded corners,align=center] (c) at (0,-3) {półprosta algebra\\$\dim=r+|\Phi|$};
\draw[->] (a)--(b);
\draw[->] (b)--(c);
\node[anchor=west,align=left] at (2.25,-.75) {iloraz wolnej\\algebry Liego};
\node[anchor=west,align=left] at (2.25,-2.25) {odbicia, moduły\\i argument o ideałach};
\end{tikzpicture}

\small Konstrukcja od układu korzeni do półprostej algebry Liego.
\end{center}

\end{document}
